% Options for packages loaded elsewhere
% Options for packages loaded elsewhere
\PassOptionsToPackage{unicode}{hyperref}
\PassOptionsToPackage{hyphens}{url}
\PassOptionsToPackage{dvipsnames,svgnames,x11names}{xcolor}
%
\documentclass[
  9.5pt,
  letterpaper,
  open=right,
  numbers=noenddot,
  captions=tableheading]{scrbook}
\usepackage{xcolor}
\usepackage[paperwidth=8in,paperheight=10in,top=0.5in,bottom=0.625in,inner=0.875in,outer=1.75in,marginparwidth=1.25in,marginparsep=0.25in,headheight=14pt,headsep=18pt,includehead,footskip=24pt,heightrounded]{geometry}
\usepackage{amsmath,amssymb}
\setcounter{secnumdepth}{3}
\usepackage{iftex}
\ifPDFTeX
  \usepackage[T1]{fontenc}
  \usepackage[utf8]{inputenc}
  \usepackage{textcomp} % provide euro and other symbols
\else % if luatex or xetex
  \usepackage{unicode-math} % this also loads fontspec
  \defaultfontfeatures{Scale=MatchLowercase}
  \defaultfontfeatures[\rmfamily]{Ligatures=TeX,Scale=1}
\fi
\usepackage{lmodern}
\ifPDFTeX\else
  % xetex/luatex font selection
  \setmainfont[BoldFont=STIX Two Text Bold,ItalicFont=STIX Two Text
Italic,BoldItalicFont=STIX Two Text Bold Italic]{STIX Two Text}
  \setsansfont[Scale=0.90]{Avenir Next}
  \setmonofont[Scale=0.78]{Menlo}
\fi
% Use upquote if available, for straight quotes in verbatim environments
\IfFileExists{upquote.sty}{\usepackage{upquote}}{}
\IfFileExists{microtype.sty}{% use microtype if available
  \usepackage[]{microtype}
  \UseMicrotypeSet[protrusion]{basicmath} % disable protrusion for tt fonts
}{}
\usepackage{setspace}
\makeatletter
\@ifundefined{KOMAClassName}{% if non-KOMA class
  \IfFileExists{parskip.sty}{%
    \usepackage{parskip}
  }{% else
    \setlength{\parindent}{0pt}
    \setlength{\parskip}{6pt plus 2pt minus 1pt}}
}{% if KOMA class
  \KOMAoptions{parskip=half}}
\makeatother
% Make \paragraph and \subparagraph free-standing
\makeatletter
\ifx\paragraph\undefined\else
  \let\oldparagraph\paragraph
  \renewcommand{\paragraph}{
    \@ifstar
      \xxxParagraphStar
      \xxxParagraphNoStar
  }
  \newcommand{\xxxParagraphStar}[1]{\oldparagraph*{#1}\mbox{}}
  \newcommand{\xxxParagraphNoStar}[1]{\oldparagraph{#1}\mbox{}}
\fi
\ifx\subparagraph\undefined\else
  \let\oldsubparagraph\subparagraph
  \renewcommand{\subparagraph}{
    \@ifstar
      \xxxSubParagraphStar
      \xxxSubParagraphNoStar
  }
  \newcommand{\xxxSubParagraphStar}[1]{\oldsubparagraph*{#1}\mbox{}}
  \newcommand{\xxxSubParagraphNoStar}[1]{\oldsubparagraph{#1}\mbox{}}
\fi
\makeatother

\usepackage{color}
\usepackage{fancyvrb}
\newcommand{\VerbBar}{|}
\newcommand{\VERB}{\Verb[commandchars=\\\{\}]}
\DefineVerbatimEnvironment{Highlighting}{Verbatim}{commandchars=\\\{\}}
% Add ',fontsize=\small' for more characters per line
\usepackage{framed}
\definecolor{shadecolor}{RGB}{241,243,245}
\newenvironment{Shaded}{\begin{snugshade}}{\end{snugshade}}
\newcommand{\AlertTok}[1]{\textcolor[rgb]{0.68,0.00,0.00}{#1}}
\newcommand{\AnnotationTok}[1]{\textcolor[rgb]{0.37,0.37,0.37}{#1}}
\newcommand{\AttributeTok}[1]{\textcolor[rgb]{0.40,0.45,0.13}{#1}}
\newcommand{\BaseNTok}[1]{\textcolor[rgb]{0.68,0.00,0.00}{#1}}
\newcommand{\BuiltInTok}[1]{\textcolor[rgb]{0.00,0.23,0.31}{#1}}
\newcommand{\CharTok}[1]{\textcolor[rgb]{0.13,0.47,0.30}{#1}}
\newcommand{\CommentTok}[1]{\textcolor[rgb]{0.37,0.37,0.37}{#1}}
\newcommand{\CommentVarTok}[1]{\textcolor[rgb]{0.37,0.37,0.37}{\textit{#1}}}
\newcommand{\ConstantTok}[1]{\textcolor[rgb]{0.56,0.35,0.01}{#1}}
\newcommand{\ControlFlowTok}[1]{\textcolor[rgb]{0.00,0.23,0.31}{\textbf{#1}}}
\newcommand{\DataTypeTok}[1]{\textcolor[rgb]{0.68,0.00,0.00}{#1}}
\newcommand{\DecValTok}[1]{\textcolor[rgb]{0.68,0.00,0.00}{#1}}
\newcommand{\DocumentationTok}[1]{\textcolor[rgb]{0.37,0.37,0.37}{\textit{#1}}}
\newcommand{\ErrorTok}[1]{\textcolor[rgb]{0.68,0.00,0.00}{#1}}
\newcommand{\ExtensionTok}[1]{\textcolor[rgb]{0.00,0.23,0.31}{#1}}
\newcommand{\FloatTok}[1]{\textcolor[rgb]{0.68,0.00,0.00}{#1}}
\newcommand{\FunctionTok}[1]{\textcolor[rgb]{0.28,0.35,0.67}{#1}}
\newcommand{\ImportTok}[1]{\textcolor[rgb]{0.00,0.46,0.62}{#1}}
\newcommand{\InformationTok}[1]{\textcolor[rgb]{0.37,0.37,0.37}{#1}}
\newcommand{\KeywordTok}[1]{\textcolor[rgb]{0.00,0.23,0.31}{\textbf{#1}}}
\newcommand{\NormalTok}[1]{\textcolor[rgb]{0.00,0.23,0.31}{#1}}
\newcommand{\OperatorTok}[1]{\textcolor[rgb]{0.37,0.37,0.37}{#1}}
\newcommand{\OtherTok}[1]{\textcolor[rgb]{0.00,0.23,0.31}{#1}}
\newcommand{\PreprocessorTok}[1]{\textcolor[rgb]{0.68,0.00,0.00}{#1}}
\newcommand{\RegionMarkerTok}[1]{\textcolor[rgb]{0.00,0.23,0.31}{#1}}
\newcommand{\SpecialCharTok}[1]{\textcolor[rgb]{0.37,0.37,0.37}{#1}}
\newcommand{\SpecialStringTok}[1]{\textcolor[rgb]{0.13,0.47,0.30}{#1}}
\newcommand{\StringTok}[1]{\textcolor[rgb]{0.13,0.47,0.30}{#1}}
\newcommand{\VariableTok}[1]{\textcolor[rgb]{0.07,0.07,0.07}{#1}}
\newcommand{\VerbatimStringTok}[1]{\textcolor[rgb]{0.13,0.47,0.30}{#1}}
\newcommand{\WarningTok}[1]{\textcolor[rgb]{0.37,0.37,0.37}{\textit{#1}}}

\usepackage{longtable,booktabs,array}
\newcounter{none} % for unnumbered tables
\usepackage{calc} % for calculating minipage widths
% Correct order of tables after \paragraph or \subparagraph
\usepackage{etoolbox}
\makeatletter
\patchcmd\longtable{\par}{\if@noskipsec\mbox{}\fi\par}{}{}
\makeatother
% Allow footnotes in longtable head/foot
\IfFileExists{footnotehyper.sty}{\usepackage{footnotehyper}}{\usepackage{footnote}}
\makesavenoteenv{longtable}
\usepackage{graphicx}
\makeatletter
\newsavebox\pandoc@box
\newcommand*\pandocbounded[1]{% scales image to fit in text height/width
  \sbox\pandoc@box{#1}%
  \Gscale@div\@tempa{\textheight}{\dimexpr\ht\pandoc@box+\dp\pandoc@box\relax}%
  \Gscale@div\@tempb{\linewidth}{\wd\pandoc@box}%
  \ifdim\@tempb\p@<\@tempa\p@\let\@tempa\@tempb\fi% select the smaller of both
  \ifdim\@tempa\p@<\p@\scalebox{\@tempa}{\usebox\pandoc@box}%
  \else\usebox{\pandoc@box}%
  \fi%
}
% Set default figure placement to htbp
\def\fps@figure{htbp}
\makeatother


% definitions for citeproc citations
\NewDocumentCommand\citeproctext{}{}
\NewDocumentCommand\citeproc{mm}{%
  \begingroup\def\citeproctext{#2}\cite{#1}\endgroup}
\makeatletter
 % allow citations to break across lines
 \let\@cite@ofmt\@firstofone
 % avoid brackets around text for \cite:
 \def\@biblabel#1{}
 \def\@cite#1#2{{#1\if@tempswa , #2\fi}}
\makeatother
\newlength{\cslhangindent}
\setlength{\cslhangindent}{1.5em}
\newlength{\csllabelwidth}
\setlength{\csllabelwidth}{3em}
\newenvironment{CSLReferences}[2] % #1 hanging-indent, #2 entry-spacing
 {\begin{list}{}{%
  \setlength{\itemindent}{0pt}
  \setlength{\leftmargin}{0pt}
  \setlength{\parsep}{0pt}
  % turn on hanging indent if param 1 is 1
  \ifodd #1
   \setlength{\leftmargin}{\cslhangindent}
   \setlength{\itemindent}{-1\cslhangindent}
  \fi
  % set entry spacing
  \setlength{\itemsep}{#2\baselineskip}}}
 {\end{list}}
\usepackage{calc}
\newcommand{\CSLBlock}[1]{\hfill\break\parbox[t]{\linewidth}{\strut\ignorespaces#1\strut}}
\newcommand{\CSLLeftMargin}[1]{\parbox[t]{\csllabelwidth}{\strut#1\strut}}
\newcommand{\CSLRightInline}[1]{\parbox[t]{\linewidth - \csllabelwidth}{\strut#1\strut}}
\newcommand{\CSLIndent}[1]{\hspace{\cslhangindent}#1}



\setlength{\emergencystretch}{3em} % prevent overfull lines

\providecommand{\tightlist}{%
  \setlength{\itemsep}{0pt}\setlength{\parskip}{0pt}}



 


% =============================================================================
% PHYSICAL AI THEME COLORS & PALETTE
% =============================================================================
% Sets the primary accent color and semantic layer for Physical AI
% Include this file BEFORE header-includes.tex
% =============================================================================

% Primary branding colors
\definecolor{harvardcrimson}{HTML}{A51C30}
\definecolor{ethdarkblue}{HTML}{1F407A}
\definecolor{ethblue}{HTML}{215CAF}
\definecolor{ethpetrol}{HTML}{007A87}
\definecolor{ethbronze}{HTML}{B87333}
\definecolor{ethpurple}{HTML}{5B4B8A}

% Canonical MLSysBook variables
\definecolor{accentcolor}{HTML}{3E2A68}     % Primary accent for headings and crossrefs (ETH Petrol)
\definecolor{toccolor}{HTML}{000000}        % TOC entries: black (matches Vol1/Vol2)

% Semantic layer for Physical AI diagrams, links, and callouts
\colorlet{structural}{ethdarkblue}
\colorlet{measured}{ethpetrol}
\colorlet{proposal}{ethbronze}
\colorlet{consequence}{harvardcrimson}
\colorlet{authority}{ethpurple}

% Ink levels for body, secondary, and hairline material
\definecolor{ink}{HTML}{1A202C}
\definecolor{softink}{HTML}{4A5568}
\definecolor{paleink}{HTML}{A0AEC0}

% Surfaces
\definecolor{cardbg}{HTML}{F8FAFC}
\definecolor{shadecolor}{HTML}{F4F6F9}

% Quarto standard callouts harmonized to the palette
\definecolor{quarto-callout-note-color}{HTML}{1F407A}
\definecolor{quarto-callout-note-color-frame}{HTML}{1F407A}
\definecolor{quarto-callout-tip-color}{HTML}{007A87}
\definecolor{quarto-callout-tip-color-frame}{HTML}{007A87}
\definecolor{quarto-callout-warning-color}{HTML}{A51C30}
\definecolor{quarto-callout-warning-color-frame}{HTML}{A51C30}
\definecolor{quarto-callout-important-color}{HTML}{5B4B8A}
\definecolor{quarto-callout-important-color-frame}{HTML}{5B4B8A}
\definecolor{quarto-callout-caution-color}{HTML}{A51C30}
\definecolor{quarto-callout-caution-color-frame}{HTML}{A51C30}
% ===========================================================================
% Physical AI Systems, LaTeX preamble
% ===========================================================================
%
% Engine:  LuaLaTeX
% Class:   KOMA-Script scrbook (set in _quarto.yml)
% Loaded:  first of the include-in-header files. tex/cover.tex follows and
%          depends on the palette defined here.
%
% RULES FOR EDITING THIS FILE
%
%   1. Every setting is declared exactly once. To change a value, change it
%      where it lives. Do not append an override at the bottom. This file
%      previously carried two \setlist declarations, two \arraystretch values,
%      two \emergencystretch values, and fvextra loaded twice, all because
%      changes were bolted on rather than integrated.
%   2. Sections run in dependency order: packages, palette, page, type,
%      floats, code, notes, links, utilities. Put new settings in the section
%      they belong to.
%   3. Anything that draws something belongs in cover.tex or in a figure
%      source, not here.
%   4. Comment the intent, not the syntax. A future reader needs to know why a
%      value is what it is, not that \setlength sets a length.
%
% Page geometry, font size, and leading live in _quarto.yml, because Quarto
% owns the geometry package options. This file assumes them, never sets them.
%
% ===========================================================================


% ---------------------------------------------------------------------------
% 1. Packages
% ---------------------------------------------------------------------------

\usepackage{microtype}          % Optical margins and glyph expansion.
\usepackage{float}              % Strict float placement [H].
\usepackage{morefloats}         % Expand LaTeX float registers.
\usepackage{placeins}           % Prevent floats from overshooting section/chapter boundaries.
\usepackage{needspace}          % Dynamic space reservation to prevent awkward table breaks.
\usepackage{scrlayer-scrpage}   % KOMA page styles, for running heads.
\usepackage{etoolbox}           % Patching KOMA and Quarto internals.
\usepackage{array}              % Column specifiers used by pandoc tables.
\usepackage{booktabs}           % \toprule and friends in pandoc tables.
\usepackage{enumitem}           % List spacing control.
\usepackage{framed}             % Backing for the Shaded code environment.
\usepackage{fancyvrb}           % Verbatim engine that pandoc emits into.
\usepackage{fvextra}            % Line breaking inside listings.
\usepackage{ragged2e}            % Balanced ragged-right typesetting with hyphenation.
\usepackage{marginnote}         % Margin column content.
\usepackage{environ}            % \NewEnviron, used by the fullwidth utility.
\usepackage{lettrine}           % Drop caps for chapter openers.
\usepackage{colortbl}           % Alternating row tints and table styling.
\usepackage{tikz}               % Cover artwork.
\usetikzlibrary{arrows.meta,positioning,shapes.geometric,fit,backgrounds,calc}


% ---------------------------------------------------------------------------
% 2. Palette
%
% Harvard crimson and ETH Zurich blue, plus a semantic layer. Prose, callouts,
% and figures reference the semantic names, so a palette change happens in one
% place. The raw names stay because the TikZ figure sources use them directly.
% ---------------------------------------------------------------------------

\definecolor{harvardcrimson}{HTML}{A51C30}
\definecolor{ethdarkblue}{HTML}{1F407A}
\definecolor{ethblue}{HTML}{215CAF}
\definecolor{ethpetrol}{HTML}{007A87}
\definecolor{ethbronze}{HTML}{B87333}
\definecolor{ethpurple}{HTML}{5B4B8A}

% Semantic layer. These encode the book's argument: what is structural, what
% is measured, what is only proposed, what has consequences, who has authority.
\colorlet{structural}{ethdarkblue}
\colorlet{measured}{ethpetrol}
\colorlet{proposal}{ethbronze}
\colorlet{consequence}{harvardcrimson}
\colorlet{authority}{ethpurple}

% Ink levels for body, secondary, and hairline material.
\definecolor{ink}{HTML}{1A202C}
\definecolor{softink}{HTML}{4A5568}
\definecolor{paleink}{HTML}{A0AEC0}

% Surfaces. shadecolor is named by the framed/Shaded contract, not by us.
\definecolor{cardbg}{HTML}{F8FAFC}
\definecolor{shadecolor}{HTML}{F4F6F9}

% Quarto's standard callouts, harmonized to the semantic layer.
\definecolor{quarto-callout-note-color}{HTML}{1F407A}
\definecolor{quarto-callout-note-color-frame}{HTML}{1F407A}
\definecolor{quarto-callout-tip-color}{HTML}{007A87}
\definecolor{quarto-callout-tip-color-frame}{HTML}{007A87}
\definecolor{quarto-callout-warning-color}{HTML}{A51C30}
\definecolor{quarto-callout-warning-color-frame}{HTML}{A51C30}
\definecolor{quarto-callout-important-color}{HTML}{5B4B8A}
\definecolor{quarto-callout-important-color-frame}{HTML}{5B4B8A}
\definecolor{quarto-callout-caution-color}{HTML}{A51C30}
\definecolor{quarto-callout-caution-color-frame}{HTML}{A51C30}


% ---------------------------------------------------------------------------
% 3. Page furniture (Running headers & footers)
%
% Standard two-sided academic textbook layout (Cambridge / MIT Press):
%   Even (verso / left) pages:  [Page #] ........................ [Chapter Title]
%   Odd (recto / right) pages:  [Section Title] ........................ [Page #]
% Subtitles and linebreaks are strictly stripped from headers.
% ---------------------------------------------------------------------------

\clearpairofpagestyles

\KOMAoptions{headwidth=textwithmarginpar,footwidth=textwithmarginpar}
\KOMAoptions{headsepline=0.4pt}
\KOMAoptions{markcase=noupper}
\addtokomafont{headsepline}{\color{accentcolor}}

\renewcommand*{\chaptermarkformat}{}
\renewcommand*{\sectionmarkformat}{}

\setkomafont{pagehead}{\normalfont\rmfamily\fontsize{8.5pt}{10.5pt}\selectfont\color{accentcolor}}
\setkomafont{pagenumber}{\normalfont\rmfamily\fontsize{8.5pt}{10.5pt}\selectfont\color{accentcolor}}

% Verso (even / left) header: Left is Folio, Right is Section mark (matching MLSysBook)
\lehead{\pagemark}
\rehead{\small\rightmark}

% Recto (odd / right) header: Left is Chapter mark, Right is Folio (matching MLSysBook)
\lohead{\small\leftmark}
\rohead{\pagemark}

% Chapter opening pages (plain style): centered folio at bottom in accent color
\cfoot[\pagemark]{}
\setkomafont{pagefoot}{\normalfont\rmfamily\fontsize{8.5pt}{10.5pt}\selectfont\color{accentcolor}}

\automark[section]{chapter}



% ---------------------------------------------------------------------------
% 4. Typography
%
% Serif body at a 4.85in measure, with matching classical serif headings
% matching MLSysBook Volumes 1 and 2.
%
% Paragraphs are indented rather than separated by vertical space. Vertical
% space marks structure here, not every paragraph break.
% ---------------------------------------------------------------------------

\KOMAoptions{fontsize=9.5pt}

\frenchspacing
\microtypesetup{protrusion=true, expansion=true, final}

\setlength{\parskip}{0pt plus 0.6pt}
\setlength{\parindent}{1.2em}

% Justification. Long technical compounds at this measure produce ladders of
% hyphens. Loosen the badness tolerance rather than the word spacing, so the
% grey of the page stays even.
\hyphenpenalty=180
\tolerance=1400
\setlength{\emergencystretch}{2.2em}

% Widow and orphan control (Cambridge / MIT Press standard)
\clubpenalty=10000
\widowpenalty=10000
\displaywidowpenalty=10000
\raggedbottom

% Pandoc emits \linebreak inside margin material where it is invalid.
\let\origlinebreak\linebreak
\renewcommand{\linebreak}[1][4]{\ifhmode\origlinebreak[#1]\fi}

% --- Headings (Matching MLSysBook Volumes 1 and 2) ------------------------

\setkomafont{disposition}{\normalfont\rmfamily\color{accentcolor}}
\setkomafont{chapter}{\normalfont\rmfamily\bfseries\fontsize{20pt}{24pt}\selectfont\color{accentcolor}}
\setkomafont{section}{\normalfont\rmfamily\bfseries\fontsize{13pt}{16pt}\selectfont\color{accentcolor}}
\setkomafont{subsection}{\normalfont\rmfamily\bfseries\fontsize{11pt}{14pt}\selectfont\color{accentcolor}}
\setkomafont{subsubsection}{\normalfont\rmfamily\bfseries\fontsize{10pt}{13pt}\selectfont\color{accentcolor}}
\setkomafont{paragraph}{\normalfont\rmfamily\bfseries\fontsize{9.5pt}{12pt}\selectfont\color{accentcolor}}

% Drop caps configuration for chapter openings (Matching MLSysBook)
\renewcommand{\LettrineFontHook}{\rmfamily\bfseries\color{accentcolor}}
\renewcommand{\LettrineTextFont}{\normalfont\scshape}
\setlength{\DefaultFindent}{2pt}
\setlength{\DefaultNindent}{0pt}

% Display math spacing
\setlength{\abovedisplayskip}{6pt plus 2pt minus 2pt}
\setlength{\belowdisplayskip}{6pt plus 2pt minus 2pt}
\setlength{\abovedisplayshortskip}{0pt plus 2pt}
\setlength{\belowdisplayshortskip}{4pt plus 2pt minus 2pt}

% Chapter subtitle command for chapter openers (appears before the horizontal rule)
\DeclareRobustCommand{\chaptersubtitle}[1]{%
  \texorpdfstring{%
    \newline{\normalfont\fontsize{11pt}{14pt}\selectfont\color{ethdarkblue}\textit{#1}}%
  }{}%
}

% --- Chapter & Section Formatting (Matching Volume 1 & Volume 2) ---------
\usepackage[explicit]{titlesec}

% Giant chapter-number font in margin
\newcommand{\mlsysChapNumFont}{\normalfont\rmfamily\fontsize{70pt}{70pt}\selectfont\bfseries}

% Numbered chapters (Large serif title with giant margin number)
\titleformat{\chapter}[block]
  {\normalfont\rmfamily\bfseries\color{accentcolor}}
  {}
  {0pt}
  {%
    \begin{tikzpicture}[remember picture, overlay]
      \node[anchor=north east,
            inner sep=0pt,
            font={\mlsysChapNumFont},
            text=accentcolor]
        at ([xshift=-40pt]current page.north east |- current page text area.north)
        {\thechapter};
    \end{tikzpicture}%
    {\fontsize{34pt}{38pt}\selectfont #1}%
  }
  []

% Unnumbered chapters (Preface, Notation, Syllabus): clean large serif title without horizontal bars
\titleformat{name=\chapter,numberless}
  {\normalfont\rmfamily\huge\bfseries\color{accentcolor}}
  {}
  {0pt}
  {\Huge #1}
  []

\titlespacing*{\chapter}{0pt}{0pt}{24pt}

% Section (Large, bold serif, accentcolor)
\titleformat{\section}
  {\normalfont\rmfamily\Large\bfseries\color{accentcolor}\raggedright}
  {\thesection}
  {0.5em}
  {#1}
\titlespacing*{\section}{0pc}{12pt plus 3pt minus 3pt}{4pt plus 1pt minus 1pt}[0pc]

% Subsection (large, bold serif, accentcolor)
\titleformat{\subsection}
  {\normalfont\rmfamily\large\bfseries\color{accentcolor}\raggedright}
  {\thesubsection}
  {0.5em}
  {#1}
\titlespacing*{\subsection}{0pc}{10pt plus 3pt minus 3pt}{4pt plus 1pt minus 1pt}[0pc]

% Subsubsection (normal size, bold serif, accentcolor)
\titleformat{\subsubsection}
  {\normalfont\rmfamily\normalsize\bfseries\color{accentcolor}\raggedright}
  {\thesubsubsection}
  {0.5em}
  {#1}
\titlespacing*{\subsubsection}{0pc}{9pt plus 3pt minus 3pt}{3pt plus 1pt minus 1pt}[0pc]

% --- Lists ----------------------------------------------------------------

\setlist[itemize]{leftmargin=1.1em,topsep=2pt,itemsep=1pt,parsep=0pt}
\setlist[enumerate]{leftmargin=1.3em,topsep=2pt,itemsep=1pt,parsep=0pt}

% --- Table of Contents & Structure Dividers (MLSysBook Modern Layout) ------

\makeatletter
% State flags for TOC divisions
\newif\ifmlsysinbackmatter
\mlsysinbackmatterfalse
\newif\ifsynthesiswritten
\synthesiswrittenfalse
\newif\if@firstnumbered
\@firstnumberedtrue
\newcounter{mlsysappendixcount}

% Pure TeX helper to distinguish Chapter numbers (1..17) from Appendix letters (A..G)
\newcommand{\format@toc@number}[1]{%
  \if\relax\detokenize\expandafter{\romannumeral-0#1}\relax
    Chapter~#1%
  \else
    Appendix~#1%
  \fi
}

% Start frontmatter with Roman numerals and suppress subsection noise in frontmatter
\AtBeginDocument{%
  \pagenumbering{roman}%
  \renewcommand{\thepage}{\roman{page}}%
  \renewcommand*\contentsname{{\normalfont\bfseries\fontsize{18pt}{22pt}\selectfont\color{consequence}Contents}}%
  \addtocontents{toc}{\protect\setcounter{tocdepth}{0}}%
  % Hook addcontentsline AFTER hyperref initialization
  \let\orig@addcontentsline\addcontentsline
  \renewcommand{\addcontentsline}[3]{%
    \Ifstr{#1}{toc}{%
      \Ifstr{#2}{part}{%
        \ifmlsysinbackmatter
          % Suppress Appendices part divider; Back Matter division handles it
        \else
          \orig@addcontentsline{#1}{#2}{#3}%
        \fi
      }{%
        \Ifstr{#2}{chapter}{%
          \ifmlsysinbackmatter
            \Ifstr{#3}{References}{%
              \orig@addcontentsline{#1}{#2}{#3}%
            }{%
              \Ifstr{#3}{Index}{%
                \orig@addcontentsline{#1}{#2}{#3}%
              }{%
                \stepcounter{mlsysappendixcount}%
                \orig@addcontentsline{toc}{chapter}{\protect\numberline{\Alph{mlsysappendixcount}}#3}%
              }%
            }%
          \else
            \orig@addcontentsline{#1}{#2}{#3}%
          \fi
        }{%
          \orig@addcontentsline{#1}{#2}{#3}%
        }%
      }%
    }{%
      \orig@addcontentsline{#1}{#2}{#3}%
    }%
  }%
}

\let\orig@appendix\appendix
\renewcommand{\appendix}{%
  \orig@appendix
  \global\mlsysinbackmattertrue
  \addtocontents{toc}{\protect\tocdivisionentry{Back Matter}}%
  \addtocontents{toc}{\protect\setcounter{tocdepth}{0}}%
}

% Patch \mainmatter so Quarto does not prematurely switch to Arabic before Chapter 1
\let\orig@mainmatter\mainmatter
\renewcommand{\mainmatter}{%
  % Defer pagenumbering{arabic} until Chapter 1 is reached
}

\newif\ifiscontentstoc
\iscontentstocfalse

% Patch \tableofcontents to guarantee Crimson/Accent "Contents" heading
\let\orig@tableofcontents\tableofcontents
\renewcommand{\tableofcontents}{%
  \begingroup
    \iscontentstoctrue
    \DeclareRobustCommand{\chaptersubtitle}[1]{}%
    \renewcommand*\contentsname{{\normalfont\huge\bfseries\color{accentcolor}Contents}}%
    \orig@tableofcontents
  \endgroup
}

% Major division banners in TOC (e.g. "Main Text", "Synthesis", "Back Matter")
\DeclareRobustCommand{\tocdivisionentry}[1]{%
  \par\addvspace{18pt}%
  \noindent\hfil{\normalfont\sffamily\bfseries\fontsize{10.5pt}{13.5pt}\selectfont\color{ink}#1}\hfil\par%
  \vspace{3pt}%
  \noindent\hfil{\color{paleink!80!softink}\rule{0.45\textwidth}{0.45pt}}\hfil\par%
  \addvspace{8pt}%
}

% Part banners in TOC (e.g. "Part I: The Machine")
\DeclareRobustCommand{\tocpartentry}[1]{%
  \par\addvspace{16pt}%
  \noindent\hfil{\normalfont\sffamily\bfseries\fontsize{10pt}{13pt}\selectfont\color{consequence}#1}\hfil\par%
  \addvspace{7pt}%
}

% Hook into \addparttocentry to emit clean centered Crimson Part entries and division lines
\renewcommand*{\addparttocentry}[2]{%
  \Ifstr{#1}{}{%
    \Ifstr{#2}{Appendices}{}{%
      \addtocontents{toc}{\protect\tocdivisionentry{#2}}%
    }%
  }{%
    \Ifstr{#2}{}{%
      \addtocontents{toc}{\protect\tocpartentry{Part~#1}}%
    }{%
      \addtocontents{toc}{\protect\tocpartentry{#2}}%
    }%
  }%
}

% Suppress raw Appendices part entry and format part lines cleanly in TOC
\renewcommand*{\l@part}[2]{%
  \in@{Part}{#1}%
  \ifin@
    \par\addvspace{16pt}%
    \noindent\hfil{\normalfont\sffamily\bfseries\fontsize{10pt}{13pt}\selectfont\color{consequence}%
      \begingroup
        \renewcommand{\numberline}[1]{}%
        #1%
      \endgroup
    }\hfil\par%
    \addvspace{7pt}%
  \fi
}

% Intercept \chapter command to cleanly handle frontmatter vs mainmatter transitions
\let\old@chapter\chapter
\renewcommand{\chapter}{%
  \@ifstar{\unnumbered@chapter}{\numbered@chapter}%
}

\newcommand{\cleanmark}[1]{%
  \begingroup
    \def\chaptersubtitle##1{}%
    \def\newline{}%
    \def\texorpdfstring##1##2{##2}%
    \protected@xdef\@clean@title{#1}%
  \endgroup
  \markboth{\@clean@title}{}%
}

\newcommand{\unnumbered@chapter}[1]{%
  \cleardoublepage
  \old@chapter*{#1}%
  \cleanmark{#1}%
}

\newcommand{\numbered@chapter}[2][]{%
  \cleardoublepage
  \if@firstnumbered
    \pagenumbering{arabic}%
    \renewcommand{\thepage}{\arabic{page}}%
    \setcounter{page}{1}%
    \addtocontents{toc}{\protect\tocdivisionentry{Main Text}}%
    \addtocontents{toc}{\protect\setcounter{tocdepth}{1}}%
    \global\@firstnumberedfalse
  \fi
  \refstepcounter{chapter}%
  \setcounter{section}{0}%
  \if\relax\detokenize{#1}\relax
    \old@chapter*{{\chapterformat}#2}%
    \cleanmark{#2}%
    \addcontentsline{toc}{chapter}{\protect\numberline{\thechapter}#2}%
  \else
    \old@chapter*[#1]{{\chapterformat}#2}%
    \cleanmark{#1}%
    \addcontentsline{toc}{chapter}{\protect\numberline{\thechapter}#1}%
  \fi
}

% Definitive Chapter Entry layout in TOC (MLSysBook layout)
\renewcommand*{\l@chapter}[2]{%
  \ifnum \c@tocdepth >\m@ne
    \addpenalty{-\@highpenalty}%
    \vskip 0.60em\relax
    {%
      \leftskip 0pt\relax
      \rightskip \@tocrmarg
      \parfillskip -\rightskip
      \parindent 0pt\relax
      \@afterindenttrue
      \interlinepenalty\@M
      \leavevmode
      \begingroup
        \renewcommand{\numberline}[1]{{\normalfont\rmfamily\bfseries\fontsize{9.5pt}{12pt}\selectfont\color{ink}\format@toc@number{##1}\hspace{0.55em}}}%
        {\normalfont\rmfamily\bfseries\fontsize{9.5pt}{12pt}\selectfont\color{ink}#1}\nobreak
      \endgroup
      \TOCLineLeaderFill
      \nobreak
      {\normalfont\rmfamily\bfseries\fontsize{9.5pt}{12pt}\selectfont\color{ink}#2}%
      \par
    }%
  \fi
}

% Definitive Section Entry layout in TOC (MLSysBook layout)
\renewcommand*{\l@section}[2]{%
  \ifnum \c@tocdepth >\z@
    \vskip 0.10em\relax
    {%
      \leftskip 1.4em\relax
      \rightskip \@tocrmarg
      \parfillskip -\rightskip
      \parindent 0pt\relax
      \@afterindenttrue
      \interlinepenalty\@M
      \leavevmode
      \begingroup
        \renewcommand{\numberline}[1]{\makebox[2.9em][l]{##1\hspace{0.4em}}}%
        {\normalfont\rmfamily\fontsize{8.8pt}{11pt}\selectfont\color{ink}#1}\nobreak
      \endgroup
      \TOCLineLeaderFill
      \nobreak
      {\normalfont\rmfamily\fontsize{8.8pt}{11pt}\selectfont\color{ink}#2}%
      \par
    }%
  \fi
}
\makeatother


% ---------------------------------------------------------------------------
% 5. Floats: captions and tables
%
% Captions lead with a bold takeaway, so the label sits in structural blue and
% the caption text drops to secondary ink.
% ---------------------------------------------------------------------------

\setkomafont{captionlabel}{\normalfont\sffamily\bfseries\fontsize{8.2pt}{10.5pt}\selectfont\color{structural}}
\setkomafont{caption}{\normalfont\sffamily\fontsize{8.0pt}{10.2pt}\selectfont\color{softink}}

% Tables: Clean, academic, compact styling (Cambridge/MIT Press monograph style)
\renewcommand{\arraystretch}{1.18}
\setlength{\tabcolsep}{3.8pt}

% Table rules styling
\setlength{\heavyrulewidth}{1.0pt}
\setlength{\lightrulewidth}{0.5pt}

% Prevent awkward orphan headers: require at least 6 baselines of vertical space
% before starting a longtable, allowing tables to render directly under their introducing prose
\BeforeBeginEnvironment{longtable}{%
  \par\penalty0\needspace{6\baselineskip}%
}
\AtBeginEnvironment{longtable}{%
  \normalfont\sffamily\fontsize{7.6pt}{9.6pt}\selectfont%
}
\AfterEndEnvironment{longtable}{%
  \normalfont\selectfont%
}
\AtBeginEnvironment{tabular}{%
  \normalfont\sffamily\fontsize{7.6pt}{9.6pt}\selectfont%
  \setlength{\tabcolsep}{3.8pt}%
  \renewcommand{\arraystretch}{1.18}%
}
\AfterEndEnvironment{tabular}{%
  \normalfont\selectfont%
}
\AtBeginEnvironment{tabular*}{%
  \normalfont\sffamily\fontsize{7.6pt}{9.6pt}\selectfont%
  \setlength{\tabcolsep}{3.8pt}%
  \renewcommand{\arraystretch}{1.18}%
}
\AfterEndEnvironment{tabular*}{%
  \normalfont\selectfont%
}

% Floats: Academic textbook standard float placement with placeins barrier
\makeatletter
\def\fps@figure{tbp}
\def\fps@table{htbp}
\makeatother
\setcounter{topnumber}{3}
\setcounter{bottomnumber}{2}
\setcounter{totalnumber}{5}
\setcounter{dbltopnumber}{3}
\renewcommand{\topfraction}{0.90}       % max fraction of page for floats at top
\renewcommand{\bottomfraction}{0.60}    % allow bottom floats
\renewcommand{\textfraction}{0.10}      % min fraction of page for text
\renewcommand{\floatpagefraction}{0.65} % min fraction of float page occupied
\renewcommand{\dbltopfraction}{0.90}
\renewcommand{\dblfloatpagefraction}{0.65}
\setlength{\floatsep}{10pt plus 2pt minus 2pt}
\setlength{\textfloatsep}{12pt plus 2pt minus 4pt}
\setlength{\intextsep}{10pt plus 2pt minus 2pt}
\setlength{\dblfloatsep}{10pt plus 2pt minus 2pt}
\setlength{\dbltextfloatsep}{12pt plus 2pt minus 4pt}






% ---------------------------------------------------------------------------
% 6. Code listings
%
% commandchars is load bearing. Pandoc emits \CommentTok{...} and
% \NormalTok{...} inside Highlighting, and without commandchars those macros
% are typeset as literal text instead of applied. Redefining Highlighting
% without it silently breaks every syntax-highlighted block in the book,
% including the dossier YAML, which is exactly what happened before.
% ---------------------------------------------------------------------------

% Shaded environment: tighten padding and frame for listings
\ifcsname Shaded\endcsname\else
  \newenvironment{Shaded}{}{}
\fi
\renewenvironment{Shaded}{%
  \begin{tcolorbox}[
    enhanced,
    breakable,
    colback=paleink!10!white,
    colframe=paleink!30!white,
    boxrule=0.4pt,
    arc=1.0pt,
    left=3pt,
    right=3pt,
    top=2pt,
    bottom=2pt,
    before skip=3pt,
    after skip=4pt
  ]%
}{%
  \end{tcolorbox}%
}

\DefineVerbatimEnvironment{Highlighting}{Verbatim}{
  commandchars=\\\{\},
  fontsize=\fontsize{6.6pt}{8.0pt}\ttfamily\selectfont,
  breaklines=true,
  breakanywhere=true,
  xleftmargin=1pt,
  xrightmargin=1pt
}

% Plain verbatim carries no tokens, so it takes no commandchars.
\RecustomVerbatimEnvironment{verbatim}{Verbatim}{
  fontsize=\fontsize{6.6pt}{8.0pt}\ttfamily\selectfont,
  breaklines=true,
  breakanywhere=true,
  xleftmargin=1pt,
  xrightmargin=1pt
}


% ---------------------------------------------------------------------------
% 7. Margin notes and footnotes
%
% Footnotes are set in the margin (reference-location: margin in _quarto.yml),
% so both read as one register: small, sans, secondary ink, sized compactly
% for the narrow margin column.
% ---------------------------------------------------------------------------

\renewcommand*{\marginfont}{%
  \RaggedRight\normalfont\sffamily\fontsize{7.5pt}{9.8pt}\selectfont\color{softink}}
\setkomafont{footnote}{%
  \normalfont\sffamily\fontsize{7.8pt}{10.2pt}\selectfont\color{softink}}
\setkomafont{footnotemark}{%
  \normalfont\sffamily\bfseries\fontsize{7.8pt}{10.2pt}\selectfont\color{consequence}}

% Sidenotes package hook: Quarto emits \sidenote{\footnotesize ...} when
% reference-location: margin is used. We define custom handlers to neutralize
% pandoc's \footnotesize and enforce compact sans-serif typography in margin notes.
\makeatletter
\NewDocumentCommand \customsidenotetext { o o +m }
{%
  \IfNoValueOrEmptyTF{#1}%
    {%
      \marginnote{%
        \RaggedRight\normalfont\sffamily\fontsize{6.8pt}{8.8pt}\selectfont\color{softink}%
        {\bfseries\color{consequence}\textsuperscript{\thesidenote}}\nobreakspace
        {\let\footnotesize\relax #3}\par
      }%
      \refstepcounter{sidenote}%
    }%
    {%
      \marginnote{%
        \RaggedRight\normalfont\sffamily\fontsize{6.8pt}{8.8pt}\selectfont\color{softink}%
        {\bfseries\color{consequence}\textsuperscript{#1}}\nobreakspace
        {\let\footnotesize\relax #3}\par
      }%
    }%
}

\NewDocumentCommand \customsidenotemark { m }
{%
  \leavevmode
  \ifhmode
      \edef \@x@sf {\the \spacefactor }%
      \nobreak
  \fi
  \hbox {\textsuperscript{\normalfont\sffamily\bfseries\fontsize{6.8pt}{8.8pt}\selectfont\color{consequence}#1}}%
  \ifhmode
      \spacefactor \@x@sf
  \fi
  \relax
}

\AtBeginDocument{%
  \renewcommand*{\@sidenotes@placemarginal}[2]{%
    \IfNoValueOrEmptyTF{#1}{\marginnote{#2}}{\marginnote[#1]{#2}}%
  }%
  \let\sidenotetext\customsidenotetext
  \let\@sidenotes@thesidenotemark\customsidenotemark
}
\makeatother


% ---------------------------------------------------------------------------
% 8. Links
%
% Deferred to \AtBeginDocument so hyperref has loaded and Quarto has finished
% its own \hypersetup.
% ---------------------------------------------------------------------------

\AtBeginDocument{%
  \hypersetup{
    colorlinks=true,
    linkcolor=ethblue,
    urlcolor=ethpetrol,
    citecolor=ethpetrol,
    filecolor=ethpetrol
  }%
}


% ---------------------------------------------------------------------------
% 9. Utilities & Margin Definition Badges
% ---------------------------------------------------------------------------

% Margin Definition Badge: Sleek, compact card for margin glossary terms
\NewDocumentCommand{\defmargin}{ m +m }{%
  \marginnote{%
    \begin{tcolorbox}[
      enhanced,
      breakable=false,
      colback=structural!4!white,
      colframe=structural,
      boxrule=0pt,
      leftrule=2pt,
      arc=1pt,
      left=3.5pt,
      right=3pt,
      top=2pt,
      bottom=2.5pt,
      before skip=2pt,
      after skip=3pt
    ]%
      {\fontsize{6.5pt}{8pt}\sffamily\bfseries\scshape\color{structural}Definition}\\[1pt]
      {\fontsize{7.2pt}{9.2pt}\sffamily\bfseries\color{ink}#1}\\[1.5pt]
      {\fontsize{6.8pt}{8.8pt}\sffamily\color{softink}#2}%
    \end{tcolorbox}%
  }%
}

% fullwidth spans the text block plus the margin column, for an artifact the
% reader scans across rather than reads down.
%
%   \begin{fullwidth} ... \end{fullwidth}
%
% Use it sparingly, and only for genuinely wide material. A table of sentences
% does not become legible by getting wider; it needs shorter cells.
\newlength{\fullwidthlen}
\setlength{\fullwidthlen}{\dimexpr\textwidth+\marginparsep+\marginparwidth\relax}
\NewEnviron{fullwidth}{%
  \par\noindent
  \makebox[\textwidth][l]{%
    \begin{minipage}{\fullwidthlen}\BODY\end{minipage}%
  }\par
}


% ---------------------------------------------------------------------------
% 10. Margin Alignment Guides (development aid — toggle)
%
% Draws light frames around the text block, header/footer, and the margin-note
% column, so text, figures, tables, and notes can be aligned during authoring.
% Set \MarginDebugfalse for production builds; \MarginDebugtrue to show guides.
% ---------------------------------------------------------------------------

\newif\ifMarginDebug
\MarginDebugfalse% <<< set to \MarginDebugtrue to show margin alignment guides
\ifMarginDebug
  \usepackage{showframe}
  \renewcommand*\ShowFrameLinethickness{0.2pt}
  \renewcommand*\ShowFrameColor{\color{red!45}}
\fi



% ===========================================================================
% Physical AI Systems, cover and title pages
% ===========================================================================
%
% Loaded after tex/preamble.tex, which defines the palette this file uses.
% This file draws; the preamble configures. Keep it that way.
%
% \maketitle produces two pages:
%
%   1. The cover. Full-bleed ETH blue, crimson rules top and bottom, and the
%      proposal-permission loop drawn as the book's signature diagram. The
%      loop is on the cover because it is the argument: an untrusted proposer,
%      a trusted permitter, an irreversible world, and the feedback that makes
%      the whole thing endogenous.
%   2. The title page. Quiet, typographic, no artwork.
%
% Text on the cover is set here rather than pulled from \@title so the cover
% can be composed as a piece of design. If the book is retitled, change it in
% both places.
%
% ===========================================================================

\makeatletter
\renewcommand*{\maketitle}{%

% --- Page 1: cover ---------------------------------------------------------
\begin{titlepage}\thispagestyle{empty}%
  \newgeometry{margin=0pt}%
  \begin{tikzpicture}[remember picture,overlay]

    % Pure white background
    \fill[white] (current page.south west) rectangle (current page.north east);

    % Floating 3D sculpture artwork centered in upper 58%
    \node[anchor=center, inner sep=0pt] at ([yshift=1.75in]current page.center) {%
      
\includegraphics[width=0.86\paperwidth]{/Users/VJ/GitHub/MLSysBook-integrate-physical/physical/book/assets/covers/cover-physical-loop-recommended-art.png}%
    };

    % Left-aligned Subtitle in Roman Serif (Matching Vol 1 & Vol 2)
    \node[anchor=south west, inner sep=0pt] at ([xshift=0.875in, yshift=3.35in]current page.south west) {%
      {\fontsize{26pt}{30pt}\rmfamily\selectfont\color{softink}Physical AI}%
    };

    % Large Two-Line Title in Book's Native Serif (Matching Vol 1 & Vol 2)
    \node[anchor=north west, inner sep=0pt] at ([xshift=0.875in, yshift=3.15in]current page.south west) {%
      \linespread{1.05}\selectfont%
      {\fontsize{50pt}{56pt}\rmfamily\selectfont\color{ink}%
       \begin{tabular}{@{}l@{}}%
         \fontsize{50pt}{56pt}\rmfamily\selectfont Machine Learning\\[0.10em]%
         \fontsize{50pt}{56pt}\rmfamily\selectfont Systems\\[0.40em]{\fontsize{13.5pt}{17pt}\rmfamily\itshape\color{accentcolor}That Sense and Act}%
       \end{tabular}%
      }%
    };

    % Right-aligned Author Block at bottom right (Matching Vol 1 & Vol 2)
    \node[anchor=south east, inner sep=0pt] at ([xshift=-0.875in, yshift=0.95in]current page.south east) {%
      \linespread{1.1}\selectfont%
      {\fontsize{24pt}{30pt}\rmfamily\selectfont\color{ink}%
       \begin{tabular}{@{}r@{}}%
         \fontsize{24pt}{30pt}\rmfamily\selectfont Vijay\\[0.15em]%
         \fontsize{24pt}{30pt}\rmfamily\selectfont Janapa Reddi%
       \end{tabular}%
      }%
    };

  \end{tikzpicture}%
  \restoregeometry
\end{titlepage}%

% --- Page 2: title page ----------------------------------------------------
\begin{titlepage}\thispagestyle{empty}%
  \raggedright
  \vspace*{0.6in}
  {\textcolor{harvardcrimson}{\rule{2.2in}{1.2pt}}\par}
  \vspace{0.25in}
  {\sffamily\bfseries\fontsize{22}{26}\selectfont\color{structural}\@title\par}
  \vspace{0.08in}
  {\sffamily\fontsize{11}{15}\selectfont\color{consequence}\@subtitle\par}
  \vspace{0.55in}
  {\sffamily\large\bfseries\color{structural}\@author\par}
  \vspace{0.04in}
  {\sffamily\small\color{softink}Harvard University \ $\cdot$\  ETH Zurich\par}
  \vspace{0.18in}
  {\sffamily\footnotesize\color{paleink}\@date\par}
  \vfill
  {\sffamily\small\bfseries\color{consequence}$W_t \longrightarrow W_{t+1}$\par}
  \vspace{0.08in}
  {\sffamily\scriptsize\color{softink}%
   Learned proposals, deterministic permission, and defensible physical release.\par}
  \vspace{0.35in}
\end{titlepage}%

}
\makeatother
\usepackage{needspace}
\let\Needspace\needspace
\makeatletter
\@ifpackageloaded{bookmark}{}{\usepackage{bookmark}}
\makeatother
\makeatletter
\@ifpackageloaded{caption}{}{\usepackage{caption}}
\AtBeginDocument{%
\ifdefined\contentsname
  \renewcommand*\contentsname{Table of contents}
\else
  \newcommand\contentsname{Table of contents}
\fi
\ifdefined\listfigurename
  \renewcommand*\listfigurename{List of Figures}
\else
  \newcommand\listfigurename{List of Figures}
\fi
\ifdefined\listtablename
  \renewcommand*\listtablename{List of Tables}
\else
  \newcommand\listtablename{List of Tables}
\fi
\ifdefined\figurename
  \renewcommand*\figurename{Figure}
\else
  \newcommand\figurename{Figure}
\fi
\ifdefined\tablename
  \renewcommand*\tablename{Table}
\else
  \newcommand\tablename{Table}
\fi
}
\@ifpackageloaded{float}{}{\usepackage{float}}
\floatstyle{ruled}
\@ifundefined{c@chapter}{\newfloat{codelisting}{h}{lop}}{\newfloat{codelisting}{h}{lop}[chapter]}
\floatname{codelisting}{Listing}
\newcommand*\listoflistings{\listof{codelisting}{List of Listings}}
\makeatother
\makeatletter
\makeatother
\makeatletter
\@ifpackageloaded{caption}{}{\usepackage{caption}}
\@ifpackageloaded{subcaption}{}{\usepackage{subcaption}}
\makeatother
\makeatletter
\@ifpackageloaded{sidenotes}{}{\usepackage{sidenotes}}
\@ifpackageloaded{marginnote}{}{\usepackage{marginnote}}
\makeatother
\newcommand{\fbxIconPath}{assets/images/icons/callouts}
\newcommand{\fbxIconFormat}{svg}
\makeatletter
\@ifpackageloaded{tcolorbox}{}{\usepackage[many]{tcolorbox}}
\makeatother
%%%% --- foldbox preamble for Physical AI Systems --- %%%%%
\RequirePackage{xstring}
\RequirePackage{etoolbox}

% Icons are disabled: clean tinted title bar + colored left rule.
\newcommand{\fbxIncludeIcon}[2]{}

% -----------------------------------------------------------------------------
% Harvard Crimson & ETH Zurich Semantic Palette
% -----------------------------------------------------------------------------
\definecolor{harvardcrimson}{HTML}{A51C30}
\definecolor{ethdarkblue}{HTML}{1F407A}
\definecolor{ethpetrol}{HTML}{007A87}
\definecolor{ethbronze}{HTML}{B87333}
\definecolor{ethpurple}{HTML}{5B4B8A}

% Fallback colors
\definecolor{fbx-default-color1}{HTML}{F0F4FA}
\definecolor{fbx-default-color2}{HTML}{1F407A}

% Custom Callout Group Colors
\definecolor{callout-objective-color1}{HTML}{F0F4FA}
\definecolor{callout-objective-color2}{HTML}{1F407A}

\definecolor{callout-definition-color1}{HTML}{EDF7F7}
\definecolor{callout-definition-color2}{HTML}{007A87}

\definecolor{callout-incident-color1}{HTML}{FBF2F3}
\definecolor{callout-incident-color2}{HTML}{A51C30}

\definecolor{callout-math-color1}{HTML}{EEF4FB}
\definecolor{callout-math-color2}{HTML}{215CAF}

\definecolor{callout-rule-color1}{HTML}{FDF8F0}
\definecolor{callout-rule-color2}{HTML}{B87333}

\definecolor{callout-algorithm-color1}{HTML}{F4F6F9}
\definecolor{callout-algorithm-color2}{HTML}{1A365D}

\definecolor{callout-contract-color1}{HTML}{FDF8F0}
\definecolor{callout-contract-color2}{HTML}{B87333}

\definecolor{callout-law-color1}{HTML}{F0F4FA}
\definecolor{callout-law-color2}{HTML}{1F407A}

\definecolor{callout-decision-color1}{HTML}{EDF7F7}
\definecolor{callout-decision-color2}{HTML}{007A87}

\definecolor{callout-lab-color1}{HTML}{F0F4FA}
\definecolor{callout-lab-color2}{HTML}{1F407A}

\definecolor{callout-teaching-color1}{HTML}{F3F0F8}
\definecolor{callout-teaching-color2}{HTML}{5B4B8A}

\definecolor{callout-archetype-color1}{HTML}{F0F7F4}
\definecolor{callout-archetype-color2}{HTML}{1B6B48}

\definecolor{callout-takeaways-color1}{HTML}{FBF2F3}
\definecolor{callout-takeaways-color2}{HTML}{A51C30}

\definecolor{callout-autopsy-color1}{HTML}{FBF2F3}
\definecolor{callout-autopsy-color2}{HTML}{A51C30}

\definecolor{callout-case-color1}{HTML}{FDF8F0}
\definecolor{callout-case-color2}{HTML}{B87333}

\definecolor{callout-fallacy-color1}{HTML}{FBF2F3}
\definecolor{callout-fallacy-color2}{HTML}{A51C30}

% Legacy alias definitions for safety
\definecolor{measured-color1}{HTML}{EDF7F7}
\definecolor{measured-color2}{HTML}{007A87}
\definecolor{proposal-color1}{HTML}{FDF8F0}
\definecolor{proposal-color2}{HTML}{B87333}
\definecolor{consequence-color1}{HTML}{FBF2F3}
\definecolor{consequence-color2}{HTML}{A51C30}
\definecolor{structural-color1}{HTML}{F0F4FA}
\definecolor{structural-color2}{HTML}{1F407A}
\definecolor{authority-color1}{HTML}{F3F0F8}
\definecolor{authority-color2}{HTML}{5B4B8A}

% Clean, publication-grade academic tcolorbox callout styling
% Designed for 7x10in format: compact, elegant left rule, subtle tint, no heavy visual clutter
\newenvironment{fbx}[3]{%
  \begin{tcolorbox}[
    enhanced,
    breakable,
    before skip=5pt,
    after skip=6pt,
    colframe=#1-color2,
    colback=#1-color1!15!white,
    colbacktitle=#1-color1!35!white,
    coltitle=#1-color2,
    fonttitle=\normalfont\sffamily\bfseries\fontsize{8.2pt}{10.5pt}\selectfont,
    fontupper=\normalfont\fontsize{8.5pt}{11pt}\selectfont,
    fontlower=\normalfont\fontsize{8.5pt}{11pt}\selectfont,
    boxrule=0pt,
    leftrule=1.5pt,
    titlerule=0pt,
    arc=1.0pt,
    auto outer arc,
    left=5pt,
    right=5pt,
    top=3.5pt,
    bottom=3.5pt,
    toptitle=2.0pt,
    bottomtitle=2.0pt,
    title={%
      \normalfont\sffamily\bfseries\fontsize{8.2pt}{10.5pt}\selectfont%
      \ifblank{#3}{#2}{#2\hspace{0.4em}#3}%
    }
  ]%
}{%
  \end{tcolorbox}%
}

\newenvironment{fbxSimple}[3]{%
  \begin{tcolorbox}[
    enhanced,
    breakable,
    before skip=5pt,
    after skip=6pt,
    colframe=#1-color2,
    colback=#1-color1!15!white,
    colbacktitle=#1-color1!35!white,
    coltitle=#1-color2,
    fonttitle=\normalfont\sffamily\bfseries\fontsize{8.2pt}{10.5pt}\selectfont,
    fontupper=\normalfont\fontsize{8.5pt}{11pt}\selectfont,
    fontlower=\normalfont\fontsize{8.5pt}{11pt}\selectfont,
    boxrule=0pt,
    leftrule=1.5pt,
    titlerule=0pt,
    arc=1.0pt,
    auto outer arc,
    left=5pt,
    right=5pt,
    top=3.5pt,
    bottom=3.5pt,
    toptitle=2.0pt,
    bottomtitle=2.0pt,
    title={%
      \normalfont\sffamily\bfseries\fontsize{8.2pt}{10.5pt}\selectfont%
      \ifblank{#3}{#2}{#2\hspace{0.4em}#3}%
    }
  ]%
}{%
  \end{tcolorbox}%
}
%%%% --- end foldbox preamble --- %%%%%
%%==== colors from yaml ===%
\definecolor{callout-case-color1}{HTML}{FDF8F0}
\definecolor{callout-case-color2}{HTML}{B87333}
\definecolor{callout-takeaways-color1}{HTML}{FBF2F3}
\definecolor{callout-takeaways-color2}{HTML}{A51C30}
\definecolor{callout-autopsy-color1}{HTML}{FBF2F3}
\definecolor{callout-autopsy-color2}{HTML}{A51C30}
\definecolor{callout-algorithm-color1}{HTML}{F4F6F9}
\definecolor{callout-algorithm-color2}{HTML}{1A365D}
\definecolor{callout-fallacy-color1}{HTML}{FBF2F3}
\definecolor{callout-fallacy-color2}{HTML}{A51C30}
\definecolor{callout-definition-color1}{HTML}{EDF7F7}
\definecolor{callout-definition-color2}{HTML}{007A87}
\definecolor{callout-contract-color1}{HTML}{FDF8F0}
\definecolor{callout-contract-color2}{HTML}{B87333}
\definecolor{callout-incident-color1}{HTML}{FBF2F3}
\definecolor{callout-incident-color2}{HTML}{A51C30}
\definecolor{callout-decision-color1}{HTML}{EDF7F7}
\definecolor{callout-decision-color2}{HTML}{007A87}
\definecolor{callout-math-color1}{HTML}{EEF4FB}
\definecolor{callout-math-color2}{HTML}{215CAF}
\definecolor{callout-objective-color1}{HTML}{F0F4FA}
\definecolor{callout-objective-color2}{HTML}{1F407A}
\definecolor{callout-lab-color1}{HTML}{F0F4FA}
\definecolor{callout-lab-color2}{HTML}{1F407A}
\definecolor{callout-archetype-color1}{HTML}{F0F7F4}
\definecolor{callout-archetype-color2}{HTML}{1B6B48}
\definecolor{callout-law-color1}{HTML}{F0F4FA}
\definecolor{callout-law-color2}{HTML}{1F407A}
\definecolor{callout-teaching-color1}{HTML}{F3F0F8}
\definecolor{callout-teaching-color2}{HTML}{5B4B8A}
\definecolor{callout-rule-color1}{HTML}{FDF8F0}
\definecolor{callout-rule-color2}{HTML}{B87333}
%=============%
\usepackage{bookmark}
\IfFileExists{xurl.sty}{\usepackage{xurl}}{} % add URL line breaks if available
\urlstyle{same}
\hypersetup{
  pdftitle={Physical AI},
  pdfauthor={Vijay Janapa Reddi},
  colorlinks=true,
  linkcolor={structural},
  filecolor={Maroon},
  citecolor={measured},
  urlcolor={measured},
  pdfcreator={LaTeX via pandoc}}


\title{Physical AI}
\usepackage{etoolbox}
\makeatletter
\providecommand{\subtitle}[1]{% add subtitle to \maketitle
  \apptocmd{\@title}{\par {\large #1 \par}}{}{}
}
\makeatother
\subtitle{Machine Learning Systems That Sense and Act}
\author{Vijay Janapa Reddi}
\date{September 2, 2026}
\begin{document}
\frontmatter
\maketitle

\renewcommand*\contentsname{Table of contents}
{
\setcounter{tocdepth}{0}
\tableofcontents
}

\setstretch{1.14}
\mainmatter
\bookmarksetup{startatroot}

\chapter*{Preface}\label{preface}
\addcontentsline{toc}{chapter}{Preface}

\markboth{Preface}{Preface}

A physical AI system is a machine with three parts.

It has a \textbf{body}, which moves mass through space, applies force
through contact, or regulates a flow of energy. The body obeys mechanics
and thermodynamics rather than instructions. It has momentum that must
be arrested somewhere, windings that accumulate heat faster than they
shed it, and transmissions that absorb whatever discontinuity a command
hands them.

It has a \textbf{brain}, which perceives, remembers, reasons, and plans.
Today that brain is built from learned models, and their value is
precisely that nobody wrote down what they do. A learned component
recognizes an object nobody enumerated and follows an instruction
phrased in a way nobody anticipated. It also arrives at its answer after
a delay that varies by an order of magnitude between one invocation and
the next, and it cannot reliably say when it is wrong.

And it has a \textbf{nervous system}, which carries signals between the
two. It runs on deadlines the body sets rather than the brain, it
decides what reaches an actuator and what does not, and it is the only
part of the machine that can refuse.

The engineering problem this book takes up is that these three parts
have incompatible properties, and they have to work as one machine
anyway. The brain is capable and unaccountable. The body is unforgiving
and cannot be argued with. Between them sits a nervous system that must
be fast enough to matter, predictable enough to trust, and independent
enough to keep working when the brain does not.

\section*{What This Book Is}\label{what-this-book-is}
\addcontentsline{toc}{section}{What This Book Is}

\markright{What This Book Is}

\lettrine{T}{his} is a systems engineering book. It draws on machine
learning, on real-time and embedded systems, and on control, but it is
not a book about any of the three and it does not attempt to teach them.
Its subject is composition: what each part of the machine must guarantee
to the others, how authority over physical action is allocated, what
happens to those guarantees when everything shares one piece of silicon,
and what evidence would justify letting the result operate.

The question underneath every chapter is what the surrounding system
must know, measure, enforce, and be able to show before an unverified
proposal is permitted to move matter.

That question is new. Machine learning systems engineering has spent a
decade on training, serving, and deployment, and it has become a mature
discipline, but almost all of it stops at a digital output. This book
takes up what happens when the output is a current through a gate
driver, and the error is not a wrong label but a sheared gear.

\subsection*{The Three Paradigms}\label{the-three-paradigms}
\addcontentsline{toc}{subsection}{The Three Paradigms}

{\def\LTcaptype{none} % do not increment counter
\begin{longtable}[]{@{}
  >{\raggedright\arraybackslash}p{(\linewidth - 6\tabcolsep) * \real{0.2500}}
  >{\raggedright\arraybackslash}p{(\linewidth - 6\tabcolsep) * \real{0.2500}}
  >{\raggedright\arraybackslash}p{(\linewidth - 6\tabcolsep) * \real{0.2500}}
  >{\raggedright\arraybackslash}p{(\linewidth - 6\tabcolsep) * \real{0.2500}}@{}}
\toprule\noalign{}
\begin{minipage}[b]{\linewidth}\raggedright
Property
\end{minipage} & \begin{minipage}[b]{\linewidth}\raggedright
Traditional Digital ML
\end{minipage} & \begin{minipage}[b]{\linewidth}\raggedright
Classical Robotics \& Control
\end{minipage} & \begin{minipage}[b]{\linewidth}\raggedright
Physical AI
\end{minipage} \\
\midrule\noalign{}
\endhead
\bottomrule\noalign{}
\endlastfoot
\textbf{Observation Stream} & Exogenous, static
(\(x \sim P_{\text{data}}\)). Model errors never alter future test data.
& Dynamic feedback (\(y = Cx\)). Fixed sensor geometry with known noise
models. & \textbf{Endogenous \& Non-Stationary}
(\(s_{t+1} \sim P(\cdot \mid s_t, a_t)\)). Actions shift camera
viewpoints and contact states. \\
\textbf{Error Cost \& Boundary} & Low/Advisory: incorrect tokens or
labels cost compute time or screen pixels. & Bounded by analytical
transfer functions, gain margins, and Lyapunov proofs. &
\textbf{Thermodynamically Irreversible:} Bad actions cost kinetic
collisions, stator thermal damage, and mechanical yield. \\
\textbf{Policy Authority} & High-capacity statistical approximation;
unverified out-of-distribution. & Explicit mathematical formulas (PID,
LQR, H-infinity, analytic MPC). & \textbf{High-capacity learned
proposer} filtered by \textbf{deterministic hardware enforcers}. \\
\textbf{Temporal Deadlines} & Elastic / best-effort
(throughput-optimized batch processing). & Hard real-time cycle
determinism (\(100\text{ Hz}\text{--}1000\text{ Hz}\)). &
\textbf{Multi-Rate Tiered Hierarchy} (\(1000\text{ Hz}\) inner enforcer,
\(10\text{--}100\text{ Hz}\) perception, \(1\text{--}10\text{ Hz}\)
neural policy). \\
\textbf{Safety Certification} & Validation loss, accuracy, perplexity on
i.i.d. splits. & Formal stability proofs and gain/phase margins. &
\textbf{Control Barrier Functions (\(h(x) \ge 0\))}, hardware isolation,
and Claim-Argument-Evidence trees. \\
\end{longtable}
}

\section*{What This Book Is Not}\label{what-this-book-is-not}
\addcontentsline{toc}{section}{What This Book Is Not}

\markright{What This Book Is Not}

It is not a robotics text. Kinematics, rigid-body dynamics, and
manipulator geometry appear only where a systems decision rests on them,
and readers who need those subjects properly should go to Craig, to
Lynch and Park, or to Siciliano.

It is not a controls text. Stability analysis and frequency-domain
methods appear only where they decide an engineering question, and
Slotine and Åström remain the places to learn them.

It is not a machine learning text. No model is discussed apart from what
its use costs in time, memory, energy, and physical consequence, and
Goodfellow, Bishop, and Murphy cover the rest.

It is also not a book about physical AI as the phrase is often used,
which tends to mean generative models applied to robots and treats
physical execution as a downstream detail. This book is about the
engineering that execution demands.

\section*{Who This Book Is For}\label{who-this-book-is-for}
\addcontentsline{toc}{section}{Who This Book Is For}

\markright{Who This Book Is For}

Physical AI is the convergence of three distinct engineering traditions,
and this book is designed so that no reader is left behind and no reader
is talked down to:

\begin{itemize}
\tightlist
\item
  \textbf{The Machine Learning \& AI Practitioner:} You arrive with deep
  mastery of gradient descent, transformer architectures, diffusion
  decoders, and cloud training pipelines. This book builds the physical
  and systems bridge: showing why open-loop statistical accuracy fails
  in dynamic environments, how memory bandwidth limits streaming
  inference, and how to interface unprivileged learned policies with
  hard real-time hardware enforcers.
\item
  \textbf{The Robotics \& Control Systems Engineer:} You arrive with
  rigorous training in rigid-body dynamics, differential kinematics,
  Kalman filtering, and Lyapunov stability. This book shows how to
  augment analytical state pipelines with high-capacity learned
  foundation models, replacing fragile handcrafted feature detectors
  with open-vocabulary spatial perception while preserving formal
  Control Barrier Function guarantees.
\item
  \textbf{The Embedded Systems \& Silicon Architect:} You arrive with
  command over bare-metal C, real-time microcontrollers, interrupt
  service routines, bus protocols, and zero-allocation determinism. This
  book shows how to architect heterogeneous System-on-Chip (SoC)
  platforms where Linux-based neural accelerators and hard real-time
  safety cores share silicon, arbitrating memory crossbars without
  starving physical safety reflexes.
\end{itemize}

Every unfamiliar mechanism in this book is introduced first by the
\textbf{physical phenomenon it governs}, second by its \textbf{systems
interface}, and only third by its \textbf{disciplinary name}---ensuring
that engineers from any background build shared physical and
architectural intuition from first principles.

The book is written for advanced undergraduates, graduate students, and
practicing systems engineers, assuming basic familiarity with calculus,
linear algebra, probability, and standard programming, but requiring no
prior robotics or embedded systems background.

\section*{How the Book Is Organized}\label{how-the-book-is-organized}
\addcontentsline{toc}{section}{How the Book Is Organized}

\markright{How the Book Is Organized}

The book is organized into four sequential parts:

\begin{itemize}
\tightlist
\item
  \textbf{Part I: Physical AI Foundations (Chapters 1--4):} Establishes
  the physical container. Chapter 1 introduces the causal boundary and
  the three canonical machine archetypes. Chapter 2 derives the
  non-negotiable physical limits of the mechanical body. Chapter 3
  formalizes the neural brain as a high-capacity, unprivileged compute
  component under memory bandwidth walls. Chapter 4 derives the nervous
  system's multi-rate execution clocks, lock-free memory transport, and
  the master 5x5 synthesis grid.
\item
  \textbf{Part II: Data and Learning Systems (Chapters 5--7):} Opens the
  policy lifecycle. Chapter 5 accounts for the collection,
  teleoperation, timestamp synchronization, and compounding error of
  physical demonstration data. Chapter 6 examines the training systems
  trade-offs between imitation learning (ACT, Diffusion Policy),
  reinforcement learning, and sim-to-real physics. Chapter 7 establishes
  the systems limits of offline evaluation and constructs closed-loop
  physical benchmarking protocols.
\item
  \textbf{Part III: The Runtime Embodied Stack (Chapters 8--13):} Walks
  the runtime pipeline stage by stage: Perception Systems (Chapter 8),
  Spatial Memory \& State (Chapter 9), Intent \& Task Reasoning (Chapter
  10), Trajectory Planning (Chapter 11), and Real-Time Reflexes (Chapter
  12). Chapter 13 co-locates stochastic neural accelerators and hard
  real-time safety enforcers on shared silicon, proving when temporal
  and memory isolation holds.
\item
  \textbf{Part IV: Integration, Verification, and Deployment (Chapters
  14--17):} Assembles the systems engineering discipline for shipping
  and operating physical machines. Chapter 14 formalizes shared
  autonomy, human intervention authority, and bumpless transfer. Chapter
  15 ascends the qualification ladder through hardware-in-the-loop (HIL)
  testing and synthetic fault injection. Chapter 16 structures
  production release engineering, telemetry monitoring, and over-the-air
  updates. Chapter 17 explores open frontier challenges across
  foundation models and physical systems.
\end{itemize}

\section*{A Note on What Is Not Yet
Known}\label{a-note-on-what-is-not-yet-known}
\addcontentsline{toc}{section}{A Note on What Is Not Yet Known}

\markright{A Note on What Is Not Yet Known}

A book that teaches engineers to demand evidence cannot itself
overclaim. Where a number in these pages was measured, the measurement
is described. Where it was derived, the derivation is shown. Where it is
illustrative, it says so. And where the honest answer is that the field
does not yet know, the chapter says that instead of reaching for a
citation that does not settle it.

The separation of proposal from permission, which runs through the whole
book, is a teaching architecture and an engineering discipline. It is
not a claim of certified functional safety, and nothing here should be
read as a certification method. What the book teaches is the reasoning
that has to happen before any such process begins.

\bookmarksetup{startatroot}

\chapter*{Course Syllabus: Physical AI Systems}\label{sec-syllabus}
\addcontentsline{toc}{chapter}{Course Syllabus: Physical AI Systems}

\markboth{Course Syllabus: Physical AI Systems}{Course Syllabus:
Physical AI Systems}

\textbf{Project Seminar \& Hardware Studio} · Harvard University · ETH
Zurich

{\def\LTcaptype{none} % do not increment counter
\begin{longtable}[]{@{}
  >{\raggedright\arraybackslash}p{(\linewidth - 2\tabcolsep) * \real{0.5000}}
  >{\raggedright\arraybackslash}p{(\linewidth - 2\tabcolsep) * \real{0.5000}}@{}}
\toprule\noalign{}
\begin{minipage}[b]{\linewidth}\raggedright
Parameter
\end{minipage} & \begin{minipage}[b]{\linewidth}\raggedright
Specification
\end{minipage} \\
\midrule\noalign{}
\endhead
\bottomrule\noalign{}
\endlastfoot
\textbf{Course} & Physical AI Systems (Harvard CS/EE 288 · ETH Zurich
151-0851) \\
\textbf{Credits} & 6 ECTS / 4 Harvard Credits
\emph{(\textasciitilde150--180 total student hours)} \\
\textbf{Format} & Weekly Interactive Seminar + Hands-On Hardware Studio
· \textbf{No Written Exam} \\
\textbf{Language} & English \\
\textbf{Level} & Advanced Undergraduate \& Graduate (MSc / PhD) \\
\textbf{Hardware Kit} & \textbf{Physical AI Kit} (Arduino UNO Q
Dual-Brain: Linux MPU + Real-Time MCU) \\
\textbf{Textbook} & \emph{Physical AI: Machine Learning Systems That
Sense and Act} \\
\textbf{Course Portal} & \url{https://physical.mlsysbook.ai} \\
\end{longtable}
}

\begin{center}\rule{0.5\linewidth}{0.5pt}\end{center}

\section*{Pedagogical Vision}\label{pedagogical-vision}
\addcontentsline{toc}{section}{Pedagogical Vision}

\markright{Pedagogical Vision}

\lettrine{C}{lassical} machine learning curricula operate entirely
\textbf{behind glass}: a misplaced class label or high prediction error
results in a software retry or log entry. \textbf{Physical AI begins
when software commands physical actuators}---mass moves, friction
generates heat, kinetic energy is spent, and the physical world
permanently transitions (\(s_{t+1} \sim P(s \mid s_t, u_t)\)). You
cannot \texttt{ctrl+z} kinetic energy.

This course teaches the systems engineering disciplines required for
that transition:
\[\text{Observe} \longrightarrow \text{Estimate} \longrightarrow \text{Intend} \longrightarrow \text{Plan} \longrightarrow \text{Enforce}\]

The central invariant governing the course is the
\textbf{Proposal--Permission Boundary}: 1. \textbf{The Cognitive Cortex
(Linux MPU):} Runs deep multimodal foundation models (VLMs), spatial
tokenizers (DINOv2), and generative action chunking decoders (Diffusion
Policy / ACT). It generates timestamped, expiring intent proposals. 2.
\textbf{The Safety Enforcer (Real-Time MCU):} Runs a bare-metal,
zero-allocation \(1000\text{ Hz}\) control loop. It verifies Control
Barrier Functions (\(h(x) \ge 0\)), monitors dynamic stopping distances
(\(d_{\text{stop}}\)), and commands motor gate inverters.

\subsection*{Multi-Disciplinary Curriculum
Pathways}\label{multi-disciplinary-curriculum-pathways}
\addcontentsline{toc}{subsection}{Multi-Disciplinary Curriculum
Pathways}

Because students enter this curriculum from diverse engineering
backgrounds, the course provides three tailored conceptual pathways:

\begin{itemize}
\tightlist
\item
  \textbf{Pathway A · The Machine Learning Practitioner:} Emphasizes the
  thermodynamic and kinematic consequences of model decisions. Focuses
  on why open-loop statistical accuracy fails in dynamic environments,
  how memory bandwidth bounds inference frequency
  (\(t_{\text{stream}} = M_{\text{weights}} / B_{\text{mem}}\)), why
  compounding covariate drift requires closed-loop exposure bounds, and
  how to interface unprivileged neural models with deterministic
  real-time supervisors (Key Chapters: Chapter~\ref{sec-boundary},
  Chapter~\ref{sec-body}, Chapter~\ref{sec-brain},
  Chapter~\ref{sec-data}, Chapter~\ref{sec-training},
  Chapter~\ref{sec-intent}, Chapter~\ref{sec-enforcement},
  Chapter~\ref{sec-release}).
\item
  \textbf{Pathway B · The Robotics \& Control Engineer:} Emphasizes
  augmenting classical analytical dynamics with high-capacity learned
  foundation models. Focuses on open-vocabulary 3D semantic grounding,
  learning multimodal action distributions via trajectory diffusion
  without mode collapse, and enforcing Control Barrier Function (CBF-QP)
  safety shields over learned trajectory proposals (Key Chapters:
  Chapter~\ref{sec-boundary}, Chapter~\ref{sec-brain},
  Chapter~\ref{sec-nervous}, Chapter~\ref{sec-perception},
  Chapter~\ref{sec-memory}, Chapter~\ref{sec-planning},
  Chapter~\ref{sec-enforcement}, Chapter~\ref{sec-verification}).
\item
  \textbf{Pathway C · The Embedded \& Silicon Architect:} Emphasizes
  heterogeneous system partitioning, hardware arbitration, and
  deterministic safety enforcement. Focuses on multi-rate clock
  cadences, lock-free seqlocks with memory barriers, AXI QoS bus
  arbitration under neural DMA bursts, inductive voltage droop
  (\(L \, dI/dt\)), and hardware fault injection (Key Chapters:
  Chapter~\ref{sec-boundary}, Chapter~\ref{sec-body},
  Chapter~\ref{sec-nervous}, Chapter~\ref{sec-memory},
  Chapter~\ref{sec-enforcement}, Chapter~\ref{sec-placement},
  Chapter~\ref{sec-intervention}, Chapter~\ref{sec-verification}).
\end{itemize}

\begin{center}\rule{0.5\linewidth}{0.5pt}\end{center}

\section*{The 14-Week Curriculum \& Laboratory
Spine}\label{the-14-week-curriculum-laboratory-spine}
\addcontentsline{toc}{section}{The 14-Week Curriculum \& Laboratory
Spine}

\markright{The 14-Week Curriculum \& Laboratory Spine}

The course is structured across 17 textbook chapters and 14 laboratory
studio sessions:

\subsection*{Part I: The Laws of Physical Action (Weeks
1--4)}\label{part-i-the-laws-of-physical-action-weeks-14}
\addcontentsline{toc}{subsection}{Part I: The Laws of Physical Action
(Weeks 1--4)}

\begin{itemize}
\tightlist
\item
  \textbf{Week 1 · Kit Bring-Up \& The Closed Loop
  (Chapter~\ref{sec-boundary}):}

  \begin{itemize}
  \tightlist
  \item
    \emph{Lecture:} The causal loop, state mutation, and the
    first-exhausted physical budget.
  \item
    \emph{Studio Lab (\texttt{00-kit-bringup},
    \texttt{01-close-the-loop}):} Dual-brain IPC setup, safe idle
    transitions, open-loop vs closed-loop telemetry.
  \item
    \emph{Deliverable:} \textbf{Loop Charter}.
  \end{itemize}
\item
  \textbf{Week 2 · Reflected Inertia \& Actuator Dynamics
  (Chapter~\ref{sec-body}):}

  \begin{itemize}
  \tightlist
  \item
    \emph{Lecture:} Reflected rotor inertia (\(N^2 J_{\text{rotor}}\)),
    QDD vs.~Harmonic drives, and thermal impedance networks
    (\(\Delta T = P \cdot \theta_{\text{JA}}\)).
  \item
    \emph{Studio Lab (\texttt{02-freshness-wall}):} Inverter current
    shunt measurement, thermal step response, stopping distance bounds.
  \end{itemize}
\item
  \textbf{Week 3 · The Two-Speed Brain \& Multi-Rate Pipelines
  (Chapter~\ref{sec-brain}, Chapter~\ref{sec-nervous}):}

  \begin{itemize}
  \tightlist
  \item
    \emph{Lecture:} Slow deliberative foundation models
    (\(1\text{--}5\text{ Hz}\)) vs.~fast reflex enforcers
    (\(1000\text{ Hz}\)); lock-free seqlocks and zero-allocation SRAM.
  \item
    \emph{Studio Lab (\texttt{03-measure-both-brains},
    \texttt{04-runtime-fault-containment}):} \(P_{99}\) latency
    profiling, inter-processor seqlock buffers, MPU crash survival.
  \item
    \emph{Deliverable:} \textbf{Requirements Ledger}.
  \end{itemize}
\item
  \textbf{Week 4 · Data Foundations \& Teleoperation
  (Chapter~\ref{sec-data}):}

  \begin{itemize}
  \tightlist
  \item
    \emph{Lecture:} Compounding covariate shift, distribution drift,
    bilateral teleoperation latency, and Open X-Embodiment datasets.
  \item
    \emph{Studio Lab (\texttt{04-dma-tokens}):} Human demonstration
    capture, teleoperation jitter compensation, action space
    normalization.
  \end{itemize}
\end{itemize}

\subsection*{Part II: The Architectural Spine (Weeks
5--8)}\label{part-ii-the-architectural-spine-weeks-58}
\addcontentsline{toc}{subsection}{Part II: The Architectural Spine
(Weeks 5--8)}

\begin{itemize}
\tightlist
\item
  \textbf{Week 5 · Perception \& 3D Spatial Encoders
  (Chapter~\ref{sec-perception}):}

  \begin{itemize}
  \tightlist
  \item
    \emph{Lecture:} Sensor latency waterfalls, rolling shutter
    distortion, on-manifold IMU preintegration, and open-vocabulary 3D
    scene graphs.
  \item
    \emph{Studio Lab (\texttt{05-perception-frontier}):} MIPI CSI-2 DMA
    ring buffers, camera-to-IMU temporal calibration, spatial feature
    tokenization.
  \item
    \emph{Deliverable:} \textbf{Observation Contract}.
  \end{itemize}
\item
  \textbf{Week 6 · Spatial Memory \& Epistemic Drift
  (Chapter~\ref{sec-memory}):}

  \begin{itemize}
  \tightlist
  \item
    \emph{Lecture:} \(SE(3)\) coordinate frame trees, volumetric
    raycasting (3D-DDA), TSDFs, and covariance expansion.
  \item
    \emph{Studio Lab (\texttt{06-belief-drift}):} Timestamp skew
    mitigation, occupancy voxel updates, belief lease expiration.
  \item
    \emph{Deliverable:} \textbf{State \& Timing Model}.
  \end{itemize}
\item
  \textbf{Week 7 · Intent \& Open-Vocabulary Grounding
  (Chapter~\ref{sec-intent}):}

  \begin{itemize}
  \tightlist
  \item
    \emph{Lecture:} Vision-Language Models (VLMs) as unprivileged
    proposal engines, spatial affordances, and expiring intent leases.
  \item
    \emph{Studio Lab (\texttt{07-two-speed-intent}):} Edge VLM prompt
    parsing, 3D target bounding, intent lease TTL enforcement.
  \item
    \emph{Deliverable:} \textbf{Intent Schema}.
  \end{itemize}
\item
  \textbf{Week 8 · The Signature Enforcer \& Midterm Milestone
  (Chapter~\ref{sec-enforcement}):}

  \begin{itemize}
  \tightlist
  \item
    \emph{Lecture:} Control Barrier Functions (CBFs), active-set barrier
    quadratic programming (QP), and minimal intervention safety shields.
  \item
    \emph{Studio Lab (\texttt{08-mcu-enforcer}):} \textbf{Signature
    Lab:} Bare-metal \(1000\text{ Hz}\) MCU CBF veto filter overriding
    invalid MPU trajectory proposals.
  \item
    \emph{Deliverable:} \textbf{Enforcement Design}.
  \end{itemize}
\end{itemize}

\subsection*{Part III: Realization, Governance \& Release (Weeks
9--14)}\label{part-iii-realization-governance-release-weeks-914}
\addcontentsline{toc}{subsection}{Part III: Realization, Governance \&
Release (Weeks 9--14)}

\begin{itemize}
\tightlist
\item
  \textbf{Week 9 · Heterogeneous Silicon Placement
  (Chapter~\ref{sec-placement}):}

  \begin{itemize}
  \tightlist
  \item
    \emph{Lecture:} AXI crossbar contention, memory bandwidth walls,
    DRAM bank conflicts, and PDN inductive voltage droop
    (\(\Delta V = L \cdot dI/dt\)).
  \item
    \emph{Studio Lab (\texttt{09-placement-ripple}):} Memory bandwidth
    throttling, bus stress injection, real-time core isolation.
  \item
    \emph{Deliverable:} \textbf{Placement Ledger}.
  \end{itemize}
\item
  \textbf{Week 10 · Human Intervention \& Shared Autonomy
  (Chapter~\ref{sec-intervention}):}

  \begin{itemize}
  \tightlist
  \item
    \emph{Lecture:} Shared control policy blending, bumpless control
    transfer, and human takeover reaction times.
  \item
    \emph{Studio Lab (\texttt{10-shadow-and-faults},
    \texttt{11-authority-paths}):} Bumpless joystick override,
    state-conditioned authority arbitration.
  \item
    \emph{Deliverable:} \textbf{Authority Design}.
  \end{itemize}
\item
  \textbf{Week 11 · Verification \& Synthetic Falsification
  (Chapter~\ref{sec-verification}):}

  \begin{itemize}
  \tightlist
  \item
    \emph{Lecture:} The 4-stage qualification ladder (Sim \(\to\) PIL
    \(\to\) HIL \(\to\) In-Situ), temporal logic falsification, and
    hardware-in-the-loop stress testing.
  \item
    \emph{Studio Lab (\texttt{12-learning-turn},
    \texttt{13-ship-gate}):} Cross-layer fault injection (dropped
    packets, voltage sags, optical occlusion).
  \end{itemize}
\item
  \textbf{Week 12 · Release Assurance \& Safety Cases
  (Chapter~\ref{sec-release}):}

  \begin{itemize}
  \tightlist
  \item
    \emph{Lecture:} Claim-Argument-Evidence (CAE) safety cases, Goal
    Structuring Notation (GSN), and UL 4600 / ISO standard alignment.
  \item
    \emph{Studio Lab (\texttt{99-design-review} prep):} Cryptographic
    artifact signing, safety case construction, release gate audit.
  \item
    \emph{Deliverable:} \textbf{Release Case}.
  \end{itemize}
\item
  \textbf{Weeks 13--14 · Capstone Jury \& Unannounced Fault Defense
  (Chapter~\ref{sec-frontier}):}

  \begin{itemize}
  \tightlist
  \item
    \emph{Studio Defense:} Live capstone demonstration. The teaching
    team injects unannounced real-time hardware/sensor faults. Teams
    must prove that their physical agent either safely achieves the task
    or transitions into a certified safe holding state without
    mechanical impact.
  \item
    \emph{Deliverable:} \textbf{Final Signed Release Verdict}.
  \end{itemize}
\end{itemize}

\begin{center}\rule{0.5\linewidth}{0.5pt}\end{center}

\section*{Grading \& Assessment Policy}\label{grading-assessment-policy}
\addcontentsline{toc}{section}{Grading \& Assessment Policy}

\markright{Grading \& Assessment Policy}

This course employs a \textbf{Project-First, Continuous-Evidence
Evaluation} model without classical written exams:

{\def\LTcaptype{none} % do not increment counter
\begin{longtable}[]{@{}
  >{\raggedright\arraybackslash}p{(\linewidth - 4\tabcolsep) * \real{0.3077}}
  >{\centering\arraybackslash}p{(\linewidth - 4\tabcolsep) * \real{0.3846}}
  >{\raggedright\arraybackslash}p{(\linewidth - 4\tabcolsep) * \real{0.3077}}@{}}
\toprule\noalign{}
\begin{minipage}[b]{\linewidth}\raggedright
Component
\end{minipage} & \begin{minipage}[b]{\linewidth}\centering
Weight
\end{minipage} & \begin{minipage}[b]{\linewidth}\raggedright
Assessment Criteria
\end{minipage} \\
\midrule\noalign{}
\endhead
\bottomrule\noalign{}
\endlastfoot
\textbf{Studio Labs \& Checkpoints} & 25\% & Timely execution of weekly
bench milestones, hardware wiring integrity, and working firmware
builds. \\
\textbf{Engineering Notebook} & 35\% & Rigorous, version-controlled
records of latency distributions (\(P_{50}, P_{90}, P_{99}\)), energy
budgets, and system logs across all chapter milestones. \\
\textbf{Midterm Enforcement Gate} & 15\% & Formal verification and live
bench demonstration of the Lab 08 1 kHz MCU Control Barrier Function
enforcer. \\
\textbf{Capstone System Defense} & 25\% & Live final demonstration
defending the Physical Agent against unannounced seeded faults before
the faculty jury. \\
\end{longtable}
}

\bookmarksetup{startatroot}

\chapter*{Mathematical and Systems Notation}\label{sec-notation}
\addcontentsline{toc}{chapter}{Mathematical and Systems Notation}

\markboth{Mathematical and Systems Notation}{Mathematical and Systems
Notation}

To ensure complete clarity and consistency across all seventeen
chapters, this textbook adopts a unified mathematical, computational,
and physical systems notation.

\begin{center}\rule{0.5\linewidth}{0.5pt}\end{center}

\section*{1. Physical State and
Dynamics}\label{physical-state-and-dynamics}
\addcontentsline{toc}{section}{1. Physical State and Dynamics}

\markright{1. Physical State and Dynamics}

{\def\LTcaptype{none} % do not increment counter
\begin{longtable}[]{@{}
  >{\raggedright\arraybackslash}p{(\linewidth - 4\tabcolsep) * \real{0.3333}}
  >{\raggedright\arraybackslash}p{(\linewidth - 4\tabcolsep) * \real{0.3333}}
  >{\raggedright\arraybackslash}p{(\linewidth - 4\tabcolsep) * \real{0.3333}}@{}}
\toprule\noalign{}
\begin{minipage}[b]{\linewidth}\raggedright
Symbol
\end{minipage} & \begin{minipage}[b]{\linewidth}\raggedright
Definition
\end{minipage} & \begin{minipage}[b]{\linewidth}\raggedright
Canonical Units
\end{minipage} \\
\midrule\noalign{}
\endhead
\bottomrule\noalign{}
\endlastfoot
\(s_t \in \mathcal{S}\) & True physical state of the machine and
environment at time \(t\) & Mixed / Continuous \\
\(o_t \in \mathcal{O}\) & Sensor observation ingested by the computing
pipeline & Pixels / Volts / Rad \\
\(a_t \in \mathcal{A}\) & Actuator action command (e.g., target torque
or joint setpoint) & \(\text{N}\cdot\text{m}\) / \(\text{rad}\) \\
\(\mathbf{q} \in \mathbb{R}^n\) & Generalized joint coordinate
configuration vector & \(\text{rad}\) or \(\text{m}\) \\
\(\dot{\mathbf{q}} \in \mathbb{R}^n\) & Generalized joint velocity
vector & \(\text{rad/s}\) or \(\text{m/s}\) \\
\(\ddot{\mathbf{q}} \in \mathbb{R}^n\) & Generalized joint acceleration
vector & \(\text{rad/s}^2\) or \(\text{m/s}^2\) \\
\(\boldsymbol{\tau} \in \mathbb{R}^n\) & Generalized joint actuator
torque/force vector & \(\text{N}\cdot\text{m}\) or \(\text{N}\) \\
\(\mathbf{M}(\mathbf{q})\) & Generalized inertia / mass matrix &
\(\text{kg}\cdot\text{m}^2\) \\
\(\mathbf{C}(\mathbf{q}, \dot{\mathbf{q}})\) & Coriolis and centrifugal
force matrix & \(\text{N}\cdot\text{m}\cdot\text{s/rad}\) \\
\(\mathbf{g}(\mathbf{q})\) & Gravitational generalized force vector &
\(\text{N}\cdot\text{m}\) \\
\(\mathbf{J}(\mathbf{q})\) & Kinematic manipulator Jacobian matrix &
\(\text{m/rad}\) \\
\(N\) & Transmission / gearbox reduction ratio (\(N \gg 1\)) &
Dimensionless \\
\(N^2 J\) & Reflected rotor inertia seen at the gearbox output shaft &
\(\text{kg}\cdot\text{m}^2\) \\
\end{longtable}
}

\begin{center}\rule{0.5\linewidth}{0.5pt}\end{center}

\section*{2. Geometry and Spatial
Transformations}\label{geometry-and-spatial-transformations}
\addcontentsline{toc}{section}{2. Geometry and Spatial Transformations}

\markright{2. Geometry and Spatial Transformations}

{\def\LTcaptype{none} % do not increment counter
\begin{longtable}[]{@{}
  >{\raggedright\arraybackslash}p{(\linewidth - 4\tabcolsep) * \real{0.3333}}
  >{\raggedright\arraybackslash}p{(\linewidth - 4\tabcolsep) * \real{0.3333}}
  >{\raggedright\arraybackslash}p{(\linewidth - 4\tabcolsep) * \real{0.3333}}@{}}
\toprule\noalign{}
\begin{minipage}[b]{\linewidth}\raggedright
Symbol
\end{minipage} & \begin{minipage}[b]{\linewidth}\raggedright
Definition
\end{minipage} & \begin{minipage}[b]{\linewidth}\raggedright
Canonical Dimension
\end{minipage} \\
\midrule\noalign{}
\endhead
\bottomrule\noalign{}
\endlastfoot
\(SE(3)\) & Special Euclidean Group of 3D rigid-body poses
(\(\mathbb{R}^3 \times SO(3)\)) & Lie Group \\
\(SO(3)\) & Special Orthogonal Group of 3D rotational matrices &
\(3 \times 3\) Orthogonal \\
\(\mathbf{p} \in \mathbb{R}^3\) & Translation position vector
\([x, y, z]^\top\) & \(\text{m}\) \\
\(\mathbf{R} \in SO(3)\) & 3D rotation matrix
(\(\mathbf{R}^\top \mathbf{R} = \mathbf{I}, \det(\mathbf{R}) = +1\)) &
\(3 \times 3\) Matrix \\
\(\mathbf{q}_{\text{quat}} \in \mathbb{H}\) & Unit quaternion rotation
representation \([w, x, y, z]^\top\) &
\(\|\mathbf{q}_{\text{quat}}\| = 1\) \\
\(\mathbf{T}_{A \to B}\) & Homogeneous rigid transform mapping
coordinates from Frame \(A\) to Frame \(B\) & \(4 \times 4\) Matrix \\
\(\boldsymbol{\Sigma}_{\text{depth}}\) & 3D spatial covariance matrix of
perception depth estimate & \(\text{m}^2\) \\
\end{longtable}
}

\begin{center}\rule{0.5\linewidth}{0.5pt}\end{center}

\section*{3. Cognitive Deliberation and Policy
Models}\label{cognitive-deliberation-and-policy-models}
\addcontentsline{toc}{section}{3. Cognitive Deliberation and Policy
Models}

\markright{3. Cognitive Deliberation and Policy Models}

{\def\LTcaptype{none} % do not increment counter
\begin{longtable}[]{@{}
  >{\raggedright\arraybackslash}p{(\linewidth - 4\tabcolsep) * \real{0.3333}}
  >{\raggedright\arraybackslash}p{(\linewidth - 4\tabcolsep) * \real{0.3333}}
  >{\raggedright\arraybackslash}p{(\linewidth - 4\tabcolsep) * \real{0.3333}}@{}}
\toprule\noalign{}
\begin{minipage}[b]{\linewidth}\raggedright
Symbol
\end{minipage} & \begin{minipage}[b]{\linewidth}\raggedright
Definition
\end{minipage} & \begin{minipage}[b]{\linewidth}\raggedright
Canonical Units
\end{minipage} \\
\midrule\noalign{}
\endhead
\bottomrule\noalign{}
\endlastfoot
\(\pi_\theta(a \mid o)\) & Parametric neural policy with weights
\(\theta\) & Probability density \\
\(\mathbf{z} \in \mathbb{R}^d\) & Latent spatial-semantic representation
vector & Dimensionless \\
\(H\) & Action chunk temporal prediction horizon (e.g., \(H=16\) or
\(H=32\)) & Discrete time steps \\
\(\mathbf{A}_{t:t+H}\) & Predicted action chunk trajectory vector
\([a_t, a_{t+1}, \dots, a_{t+H-1}]\) & Sequence of actions \\
\(\mathcal{D}_{\text{demo}}\) & Dataset of nominal human/expert
demonstrations & Sequence of tuples \\
\(\mathcal{D}_{\text{rec}}\) & Dataset of human recovery and
intervention trajectories & Sequence of tuples \\
\(\epsilon_{\text{BC}}\) & Single-step imitation learning policy error
rate & Dimensionless (\(\le 1\)) \\
\(O(T^2 \epsilon)\) & Compounding quadratic covariate drift error under
open-loop policy rollouts & Cumulative error \\
\(\mathcal{L}_{\text{intent}}\) & Expiring Intent Lease with bounded
Time-To-Live (TTL) & Spatial-temporal contract \\
\end{longtable}
}

\begin{center}\rule{0.5\linewidth}{0.5pt}\end{center}

\section*{4. Real-Time Control, Safety, and
Enforcement}\label{real-time-control-safety-and-enforcement}
\addcontentsline{toc}{section}{4. Real-Time Control, Safety, and
Enforcement}

\markright{4. Real-Time Control, Safety, and Enforcement}

{\def\LTcaptype{none} % do not increment counter
\begin{longtable}[]{@{}
  >{\raggedright\arraybackslash}p{(\linewidth - 4\tabcolsep) * \real{0.3333}}
  >{\raggedright\arraybackslash}p{(\linewidth - 4\tabcolsep) * \real{0.3333}}
  >{\raggedright\arraybackslash}p{(\linewidth - 4\tabcolsep) * \real{0.3333}}@{}}
\toprule\noalign{}
\begin{minipage}[b]{\linewidth}\raggedright
Symbol
\end{minipage} & \begin{minipage}[b]{\linewidth}\raggedright
Definition
\end{minipage} & \begin{minipage}[b]{\linewidth}\raggedright
Canonical Units
\end{minipage} \\
\midrule\noalign{}
\endhead
\bottomrule\noalign{}
\endlastfoot
\(\mathcal{C} \subset \mathcal{S}\) & Forward-invariant safe operating
set & State space subset \\
\(h(\mathbf{x}) \ge 0\) & Continuously differentiable Control Barrier
Function (CBF) & Real scalar \\
\(L_f h(\mathbf{x}), L_g h(\mathbf{x})\) & Lie derivatives of the
barrier function along drift \(f\) and control \(g\) & Real scalars \\
\(\gamma\) & Extended class-\(\mathcal{K}\) barrier decay rate parameter
& \(\text{s}^{-1}\) \\
\(u_{\text{nom}}\) & Unprivileged nominal control action proposed by
neural policy & \(\text{N}\cdot\text{m}\) \\
\(u^*\) & Verified, minimally-intervened control action emitted by
CBF-QP filter & \(\text{N}\cdot\text{m}\) \\
\(d_{\text{stop}}\) & Dynamic stopping distance envelope &
\(\text{m}\) \\
\(v_{\max}, a_{\max}, j_{\max}\) & Maximum physical velocity,
acceleration, and jerk limits &
\(\text{m/s}, \text{m/s}^2, \text{m/s}^3\) \\
\end{longtable}
}

\begin{center}\rule{0.5\linewidth}{0.5pt}\end{center}

\section*{5. Hardware Silicon and Electrical
Budgets}\label{hardware-silicon-and-electrical-budgets}
\addcontentsline{toc}{section}{5. Hardware Silicon and Electrical
Budgets}

\markright{5. Hardware Silicon and Electrical Budgets}

{\def\LTcaptype{none} % do not increment counter
\begin{longtable}[]{@{}
  >{\raggedright\arraybackslash}p{(\linewidth - 4\tabcolsep) * \real{0.3333}}
  >{\raggedright\arraybackslash}p{(\linewidth - 4\tabcolsep) * \real{0.3333}}
  >{\raggedright\arraybackslash}p{(\linewidth - 4\tabcolsep) * \real{0.3333}}@{}}
\toprule\noalign{}
\begin{minipage}[b]{\linewidth}\raggedright
Symbol
\end{minipage} & \begin{minipage}[b]{\linewidth}\raggedright
Definition
\end{minipage} & \begin{minipage}[b]{\linewidth}\raggedright
Canonical Units
\end{minipage} \\
\midrule\noalign{}
\endhead
\bottomrule\noalign{}
\endlastfoot
\(B_{\text{mem}}\) & Main system memory bus streaming bandwidth &
\(\text{GB/s}\) \\
\(M_{\text{weights}}\) & Neural policy total parameter model memory
footprint & \(\text{GB}\) \\
\(t_{\text{stream}}\) & Minimum memory bandwidth parameter streaming
time (\(M_{\text{weights}}/B_{\text{mem}}\)) & \(\text{s}\) or
\(\text{ms}\) \\
\(\Delta t_{\text{fresh}}\) & Physical freshness deadline for sensor
ingestion & \(\text{ms}\) \\
\(t_{\text{lag}}\) & End-to-end sense-to-torque physical execution
latency & \(\text{ms}\) \\
\(f_{\text{ctrl}}\) & Inner reflex loop execution frequency (canonical:
1000 Hz) & \(\text{Hz}\) \\
\(I\) & Motor phase electrical current & \(\text{Amperes (A)}\) \\
\(R\) & Electrical phase winding or power bus resistance &
\(\text{Ohms } (\Omega)\) \\
\(C_{\text{th}}\) & Stator thermal capacitance & \(\text{J/K}\) \\
\(\Delta T\) & Joulean stator thermal rise
(\(\frac{I^2 R \Delta t}{C_{\text{th}}}\)) & \(\text{Kelvin (K)}\) \\
\(\Delta V\) & Inductive power rail dynamic droop
(\(L \frac{dI}{dt} + IR\)) & \(\text{Volts (V)}\) \\
\end{longtable}
}

\bookmarksetup{startatroot}

\chapter{\texorpdfstring{The Causal
Boundary\chaptersubtitle{Introduction to Physical AI Systems}}{The Causal Boundary}}\label{sec-boundary}

\begin{figure}[H]

\centering{

\includegraphics[width=1\linewidth,height=\textheight,keepaspectratio]{index_files/mediabag//Users/VJ/GitHub/MLSysBook-integrate-physical/physical/book/chapters/01-boundary/figures/fig_locator.pdf}

}

\caption{\label{fig-locator}\textbf{Physical AI Systems Architecture
Stack.} The master systems hierarchy spanning System Governance (Layer
4), Cognitive Deliberation (Layer 3), Hard Real-Time Reflex \& Silicon
(Layer 2), and the Physical Body (Layer 1). In this introductory
chapter, our active focus centers on the causal boundary at the
interface between digital software and the physical plant.}

\end{figure}%

Across humanoid robotics, autonomous transport fleets, and surgical
workcells, high-capacity neural foundation models are stepping out of
cloud data centers to command physical machines. Yet inside engineering
design reviews at companies from Tesla and Figure to Boston Dynamics and
Waymo, teams find themselves locked in an architectural crisis: machine
learning researchers celebrate 99\% virtual benchmark scores, while
robotics engineers watch hardware shake, overheat, and strip precision
gearboxes on the physical floor. An error that costs nothing on a screen
destroys hardware in the physical world.

This book is the founding curriculum for \textbf{Physical AI Systems
Architecture}---the discipline of embedding high-capacity learned
statistical models into closed-loop machines governed by classical
mechanics, thermodynamics, and hard real-time silicon determinism. It
bridges open-world cognitive intelligence and dynamical physics,
teaching engineers how to build dependable, certifiable machines that
sense, reason, and act in the physical world.

\newpage{}

\protect\phantomsection\label{callout-objectiveux2a-1.1}
\begin{fbxSimple}{callout-objective}{Learning Objectives}{}
\phantomsection\label{callout-objective*-1.1}

After studying this chapter, you will be able to:

\begin{enumerate}
\def\labelenumi{\arabic{enumi}.}
\tightlist
\item
  Locate the causal boundary in a system architecture and state what
  changes on each side of it.
\item
  Diagnose the ``Tower of Babel'' conflict between ML, Systems,
  Controls, and Embedded disciplines.
\item
  Classify an embodied machine into one of three canonical physical
  archetypes and identify which physical budget binds it first.
\item
  Apply the three-pillar scope test to determine whether a given system
  requires the engineering discipline of physical AI.
\item
  Identify the Four Pillars of an embodied machine: Dynamical Body,
  Cognitive Brain, Real-Time Nervous System, and Governance Envelope.
\item
  Trace how latency, classification instability, and missing authority
  boundaries caused the 2018 Tempe autonomous vehicle collision.
\item
  State the Four Bedrock Laws that govern physical machines across the
  causal boundary.
\item
  Explain how the heterogeneous Dual-Brain architecture (MPU proposer +
  MCU referee across shared SRAM) enforces physical safety invariants.
\item
  Trace the four-part architectural arc of this textbook and map each
  chapter to its corresponding tier in the physical AI systems stack.
\end{enumerate}

\emph{Core chapter takeaway.} The causal boundary is a physical
location, not an abstraction. Below it, the machine shapes its own
future sensory observations, computational delay is paid in physical
displacement, and energy released into the world cannot be recalled.

\end{fbxSimple}

\section{The Industrial Frontier: Artificial Intelligence Beyond
Glass}\label{sec-boundary-1-1}

\lettrine{E}{very}\lettrine{m}{ajor} epoch in computing has been defined
by the boundary across which bits travel. In personal computing, bits
crossed from mainframe backplanes to human eyes through glowing
cathode-ray displays. In cloud computing, bits crossed global
fiber-optic networks to synchronize distributed relational databases. In
modern generative artificial intelligence, bits traverse
trillion-parameter neural networks to generate text, synthesize images,
or recommend financial actions. Yet across all these transformations,
one fundamental property of computer software remained inviolate: the
computation lived safely behind glass.

Today, a massive global industrial race is underway to bring artificial
intelligence beyond the glass. Companies developing humanoid robots
(such as Tesla Optimus, Figure, and Boston Dynamics), autonomous vehicle
fleets (such as Waymo and Zoox), intelligent aerial drones (such as
Skydio), and robotic surgery platforms are deploying deep neural
policies directly aboard mobile, articulated hardware. Yet as these
systems transition from virtual benchmarks to real-world deployment,
engineering teams encounter a fundamental architectural barrier: the
\textbf{Tower of Babel} crisis.

In a typical embodied AI design review, four distinct engineering
cultures clash:

\begin{itemize}
\tightlist
\item
  \textbf{The Machine Learning Perspective:} Focuses on training deep
  transformer backbones and diffusion policies on millions of simulated
  steps, optimizing statistical loss functions and celebrating high
  validation scores, often assuming that real-world deployment is
  primarily a challenge of scaling compute and collecting broader
  teleoperation datasets.
\item
  \textbf{The Systems and Computer Architecture Perspective:} Focuses on
  memory bus bandwidth, DRAM bank conflicts, cache contention, and tail
  latency percentiles, recognizing that streaming multi-gigabyte
  foundation model weights at \(50\text{ Hz}\) threatens to saturate the
  memory bus and starve real-time camera DMA pipelines.
\item
  \textbf{The Robotics and Classical Control Perspective:} Focuses on
  Newton-Euler equations, coordinate transforms, actuator torque
  saturation, and stability margins, recognizing that unconstrained
  neural action chunks risk injecting high-frequency jerk that degrades
  mechanical gearboxes and excites unmodeled resonance.
\item
  \textbf{The Embedded Silicon and Functional Safety Perspective:}
  Focuses on bare-metal interrupt service routines, watchdog timers,
  hardware PWM trip zones, and safety standards (such as ISO 26262 and
  UL 4600), ensuring that non-deterministic neural weights never hold
  unverified direct write access to motor inverter registers.
\end{itemize}

None of these four perspectives can succeed in isolation. You cannot
solve a broken gearbox or an inductive rail voltage collapse with
gradient descent, and you cannot achieve open-vocabulary semantic task
generalization with a hand-tuned PID controller.

\protect\phantomsection\label{callout-definitionux2a-1.2}
\begin{fbxSimple}{callout-definition}{Definition:}{Physical Artificial Intelligence (Physical AI)}
\phantomsection\label{callout-definition*-1.2}

\textbf{Physical Artificial Intelligence (Physical AI)} is the systems
engineering science of embedding high-capacity learned statistical
models into closed-loop physical machines governed by classical
mechanics, thermodynamics, and hard real-time silicon determinism.

\end{fbxSimple}

Physical AI is the systems engineering synthesis that unifies these four
cultures into a single coherent architecture. It trains the
\textbf{Physical AI Systems Architect}---the bilingual systems leader
who designs the complete end-to-end container, ensuring that
high-capacity neural deliberation (Layer 3) safely and deterministically
coordinates with real-time silicon (Layer 2) and physical dynamics
(Layer 1) under formal governance (Layer 4).

To understand why this closed causal loop requires a fundamentally new
systems architecture, we must examine what physically occurs at the
hardware register level when software commands an actuator.

\section{Software That Acts: The Causal
Boundary}\label{sec-boundary-1-2}

Advisory software operates inside an informational container where
errors carry negligible physical momentum. When a recommender model
serves an irrelevant product, a classifier mislabels an image, or a
language model emits an inaccurate token, the computational pipeline
terminates in a display buffer, a network socket, or a sandboxed
database table. The failure remains bounded by the digital medium. A
user dismisses the recommendation, a human analyst closes the alert, or
a database transaction coordinator issues an abort that restores memory
pages from an undo log. \emph{Advisory software is immune to its own
errors because of where it stops, not because of how it was built.} The
safety of the pipeline does not stem from statistical guarantees in the
training objective or rigorous validation on test splits, but from the
physical detachment of the interface. Between the output of the neural
network and any change in the physical world, an air gap or a human
operator intercepts the signal.

Physical artificial intelligence eliminates that detachment. In a
physical machine, the final instruction of a control loop does not write
to a screen or populate a web response. It executes a memory-mapped bus
transaction that updates a hardware register on an actuator controller.
Depending on the machine architecture, that hardware register is a
digital-to-analog converter (DAC) latch holding a commanded reference
voltage, a pulse-width modulation (PWM) compare register (such as
\texttt{TIMx\_CCR} on an ARM Cortex-M or \texttt{ePWM\ CMPA} on a TI
C2000) setting the duty cycle of a power inverter stage, or a serial
control word configuring an H-bridge gate driver.\sidenote{\footnotesize In canonical
  safety microcontrollers, hardware trip zones (\texttt{EPWM\_TZFRC},
  \texttt{TIMx\_BKIN}) operate independently of the CPU. Hardware break
  circuits force complementary PWM outputs into high-impedance or
  low-side clamp states within \(<50\text{ ns}\), preventing
  shoot-through before software ISRs execute.} At that exact memory
address, informational abstraction ends. The bit pattern written to the
register immediately modulates gate voltages across power MOSFETs,
drawing current from a DC bus into the copper windings of a motor.
Current develops magnetic flux, which generates mechanical Lorentz
torque across the rotor. That torque accelerates the rotor and
transmission against reflected load inertia, imparting physical momentum
to the mechanism.

\protect\phantomsection\label{callout-definitionux2a-1.3}
\begin{fbxSimple}{callout-definition}{Definition:}{The Causal Boundary}
\phantomsection\label{callout-definition*-1.3}

The \textbf{causal boundary} is the physical and electrical interface
where digital computational bits convert into physical forces, currents,
or flows. Above the boundary, software state is digital, sandboxed, and
computationally idempotent; below the boundary, actions permanently
alter the physical environment, dissipate irreversible energy, and
endogenously reshape all subsequent sensory observations.

\end{fbxSimple}

When a digital neural network misclassifies an image in a cloud
benchmark, the consequence is confined to an incorrect integer stored in
a memory register. When an autoregressive language model generates an
ungrounded hallucination, the failure manifests as superfluous text on a
browser screen. Software engineers can catch runtime errors, discard
invalid outputs, roll back transactions to a clean checkpoint, or rely
on human users to ignore an answer. In computer systems engineering,
this property is known as \emph{idempotency}: the state of the
surrounding physical universe remains completely decoupled from the
intermediate outputs of the software pipeline.\sidenote{\footnotesize \textbf{For the
  Physical Engineer:} \emph{Idempotency} in digital software means an
  operation can be executed multiple times without changing the final
  state (e.g., refreshing a webpage or retrying a database write). In
  physical systems, no operation is idempotent---every motor command
  consumes irreversible energy and alters the physical environment.}
Memory can be cleared, registers zeroed, and processes restarted without
permanent side effects on the environment.

Physical artificial intelligence begins at the exact instant this
isolation breaks. The moment an algorithmic prediction leaves the
digital domain to modulate a pulse-width modulation motor drive, fire a
valve solenoid, or adjust an electrical grid inverter, computation
ceases to be advice and becomes physical force. At this threshold, an
unhandled exception is no longer a benign runtime error; it is a burnt
motor winding, a sheared gearbox, or a collision with an obstacle in the
physical world.

This transition from digital state to physical force destroys the
statistical foundation of classical machine learning. In offline vision,
natural language processing, and tabular prediction, data collection is
exogenous: the world generates a distribution of inputs, and the model
processes them without altering the mechanism that created them. An
incorrect classification on one image has zero effect on the
distribution governing the next image. In a physical system, the
environment evolves in a closed causal loop: the next physical state and
sensor observation depend directly on the action produced by the prior
model inference.

The sensory input stream is endogenous, generated by the machine's own
past actions. When a policy produces an erroneous torque command, the
joint rotates, the platform shifts, and mounted cameras or lidar units
swing into new viewpoints. The machine drives itself into an unfamiliar
region of the state space where training data is sparse, increasing the
probability of subsequent inference errors and compounding the
divergence from nominal operation. This fundamental architectural
divergence is illustrated in Figure~\ref{fig-01-causal-loop},
contrasting the open-loop, idempotent execution of advisory digital
machine learning with the closed causal loop of Physical AI crossing the
actuator register boundary.

\begin{figure}[t]

\centering{

\includegraphics[width=0.95\linewidth,height=\textheight,keepaspectratio]{index_files/mediabag//Users/VJ/GitHub/MLSysBook-integrate-physical/physical/book/chapters/01-boundary/figures/fig01_causal_loop.pdf}

}

\caption{\label{fig-01-causal-loop}\textbf{The Closed Causal Loop of
Physical AI versus Open-Loop Digital Computation.} (a) Open-loop
advisory machine learning operates behind glass on exogenous benchmark
distributions, where computational errors are idempotent and bounded by
software exception handling or transactional rollbacks. (b) In Physical
AI, neural deliberation proposals pass through a deterministic nervous
system safety gate before committing to hardware actuator registers.
Once memory-mapped pulse-width modulation compare registers latch,
electrical current converts into Lorentz force, mechanical torque,
kinetic momentum, and Joule heat, permanently mutating environment state
and endogenously shifting all subsequent sensory observations.}

\end{figure}%

Because physical mass possesses inertia and computation takes finite
time, latency in physical AI translates directly into physical
displacement. As shown on the autonomous testbed in
Figure~\ref{fig-01-sensor-rack}, real-world robotic systems rely on
dense multi-modal sensor suites---including spinning LiDAR scanners,
camera arrays, and radar units---that continuously stream high-bandwidth
telemetry to onboard processors. Every stage in this pipeline---from
photodiode exposure and bus serialization (MIPI CSI-2,
GMSL2),\sidenote{\footnotesize Automotive SerDes links (TI FPD-Link III, ADI GMSL2/3)
  serialize multi-gigabit MIPI CSI-2 sensor streams over coaxial cables
  with sub-\(15\text{ }\mu\text{s}\) transfer latency, hardware CRC-32
  verification, and bidirectional backchannel \texttt{FSYNC} triggers
  that align multi-camera exposure start times within sub-microsecond
  jitter.} to neural inference on accelerator silicon and packet
transmission over fieldbuses (CAN-FD, EtherCAT)---consumes precious
milliseconds. During this aggregate delay window, the machine cannot
incorporate new sensory information and continues moving along the
trajectory dictated by its existing momentum.

\protect\phantomsection\label{callout-ruleux2a-1.4}
\begin{fbxSimple}{callout-rule}{Architect’s Rule of Thumb:}{Latency is Paid in Millimeters}
\phantomsection\label{callout-rule*-1.4}

In advisory machine learning, computational delay only frustrates a
user; in physical AI, latency is paid directly in physical displacement.
At velocity \(v\), an unbudgeted delay \(\Delta t\) consumes unguided
motion that no downstream neural capacity can ever recover.

\end{fbxSimple}

As a concrete example, consider a mobile platform traveling at a modest
walking speed of \(2.0\text{ m/s}\). If the perception pipeline, neural
policy inference, and motor controller communication consume a combined
latency of \(150\text{ ms}\), the platform traverses
\((2.0\text{ m/s}) \times (0.15\text{ s}) = 0.30\text{ m}\)
(\(30\text{ cm}\)) before the first motor response can occur. During
those \(30\text{ cm}\), the vehicle is entirely blind to any
environmental change that occurred after the initial sensor exposure. If
an obstacle enters the path \(20\text{ cm}\) ahead during that inference
window, no amount of model capacity or predictive accuracy can alter the
outcome, because the physical displacement has already been accumulated
by the body.

Once digital commands transition into mechanical work, the state change
becomes thermodynamically irreversible. In standard software systems,
erroneous mutations are contained through exception handling, memory
isolation, and atomic rollback protocols. When an unhandled exception
occurs, the runtime unwinds the call stack, reclaims heap allocations,
and restores the machine state to a known clean checkpoint. The physical
world provides no rollback mechanism. When an actuator controller
applies current to a motor, electrical energy transfers from the battery
into kinetic energy in mechanical linkages or thermal energy dissipated
through resistive Joule heating in the stator windings. If an incorrect
joint torque drives a manipulator into a rigid steel fixture, the
kinetic energy converts into plastic deformation of the tool, mechanical
strain on the gearbox teeth, and heat. No software interrupt can pull
that kinetic energy out of the crushed structure or restore sheared gear
teeth to their original geometry.

\begin{figure}[t]

\centering{

\includegraphics[width=0.85\linewidth,height=\textheight,keepaspectratio]{/Users/VJ/GitHub/MLSysBook-integrate-physical/physical/book/chapters/01-boundary/figures/fig01_real_lidar_sensor_rack.jpg}

}

\caption{\label{fig-01-sensor-rack}\textbf{Spinning Multi-Beam LiDAR
Sensor Hardware for Real-Time Perception.} Close-up photograph of the
64-beam spinning optical LiDAR sensor assembly (Velodyne HDL-64E)
mounted aboard an autonomous research testbed. In Physical AI,
high-bandwidth sensor streams (1.33 million points/second) must be
serialized over real-time interfaces (e.g., FPD-Link III, GMSL2,
Automotive Ethernet), synchronized via hardware timestamps, and
processed through neural perception backbones within strict end-to-end
latency budgets, because every millisecond of transport and compute
delay converts directly into unguided physical displacement. (Source:
Steve Jurvetson, CC BY 2.0).}

\end{figure}%

Understanding why Physical AI requires a distinct systems engineering
discipline requires contrasting its core properties against both
advisory digital machine learning and classical control theory. As
summarized in Table~\ref{tbl-01-paradigm-comparison}, the three
paradigms are sorted from pure software isolation on the left, to pure
classical dynamics in the middle, to the unified physical AI synthesis
on the right.

\begin{longtable}[]{@{}
  >{\raggedright\arraybackslash}p{(\linewidth - 6\tabcolsep) * \real{0.2500}}
  >{\raggedright\arraybackslash}p{(\linewidth - 6\tabcolsep) * \real{0.2500}}
  >{\raggedright\arraybackslash}p{(\linewidth - 6\tabcolsep) * \real{0.2500}}
  >{\raggedright\arraybackslash}p{(\linewidth - 6\tabcolsep) * \real{0.2500}}@{}}
\caption{\textbf{Comparison Matrix of Computing Paradigms across the
Causal Boundary.} Systematic comparison of Advisory Digital Machine
Learning (left), Classical Robotics \& Control Theory (middle), and
Physical Artificial Intelligence (right) across operational substrate,
model formulation, feedback dynamics, execution cadence, actuation
authority, failure semantics, and safety certification
primitives.}\label{tbl-01-paradigm-comparison}\tabularnewline
\toprule\noalign{}
\begin{minipage}[b]{\linewidth}\raggedright
Architectural Dimension
\end{minipage} & \begin{minipage}[b]{\linewidth}\raggedright
Digital Machine Learning (Advisory)
\end{minipage} & \begin{minipage}[b]{\linewidth}\raggedright
Classical Robotics \& Control
\end{minipage} & \begin{minipage}[b]{\linewidth}\raggedright
Physical Artificial Intelligence (Closed-Loop)
\end{minipage} \\
\midrule\noalign{}
\endfirsthead
\toprule\noalign{}
\begin{minipage}[b]{\linewidth}\raggedright
Architectural Dimension
\end{minipage} & \begin{minipage}[b]{\linewidth}\raggedright
Digital Machine Learning (Advisory)
\end{minipage} & \begin{minipage}[b]{\linewidth}\raggedright
Classical Robotics \& Control
\end{minipage} & \begin{minipage}[b]{\linewidth}\raggedright
Physical Artificial Intelligence (Closed-Loop)
\end{minipage} \\
\midrule\noalign{}
\endhead
\bottomrule\noalign{}
\endlastfoot
\textbf{Operational Substrate} & Decoupled virtual memory buffers (GPU
VRAM, DRAM, web endpoints) & Continuous analog state space governed by
differential equations & Hybrid: high-dimensional neural representations
mapped to analog physical actuators \\
\textbf{Model Formulation} & High-capacity statistical approximators
(LLMs, ViTs, diffusion policies) & Analytically specified transfer
functions and deterministic state-space matrices & Foundation policies
paired with deterministic runtime safety filters \\
\textbf{Feedback Dynamics} & Exogenous and static; independent and
identically distributed samples & Continuous deterministic state
feedback with Gaussian sensor noise & Endogenous, non-stationary closed
loop; actions shape future observations \\
\textbf{Execution Cadence} & Throughput-optimized batching; soft or
unbounded latency budgets & Hard real-time single-rate loops on
microcontrollers (100--1000 Hz) & Asynchronous multi-rate cadences:
System 2 (1 Hz), System 1.5 (50 Hz), System 1 (1000 Hz) \\
\textbf{Actuation Authority} & Advisory; sandboxed APIs, display
buffers, or human-in-the-loop validation & Hardcoded feedback directly
driving power inverter timers and PWM registers & Decoupled
Proposal--Permission architecture; untrusted proposals gated by MCU
safety referee \\
\textbf{Failure Semantics} & Idempotent; call-stack unwind, database
rollback, or human alert dismissal & Actuator saturation or tracking
error within analytically bounded stability envelope & Irreversible
thermodynamic work, plastic deformation, kinetic collision, or stator
overheating \\
\textbf{Safety Assurance} & Offline statistical loss minimization on
validation splits & Formal mathematical stability proofs (Lyapunov
functions, gain/phase margins) & Layered runtime invariant enforcement,
forward-invariant barrier filters, and fault injection \\
\end{longtable}

\section{Three Machine Archetypes}\label{sec-boundary-1-3}

Once digital predictions cross the causal boundary into physical force,
they interact with physical machines of every conceivable shape and
scale. An autonomous delivery drone navigating turbulent wind gusts
seems to share nothing with a six-axis robotic arm assembling an
automotive transmission, and neither appears related to a containerized
battery management plant regulating liquid coolant. A practitioner or
student might wonder: how can a single systems engineering curriculum
possibly encompass such wildly diverse physical systems without
dissolving into a collection of bespoke, ad-hoc recipes?

The key is to look past cosmetic form factors and examine the underlying
physical conservation law that binds the machine. Beneath the surface
diversity of shapes, sizes, and applications, every embodied AI system
can be systematically categorized into one of three canonical physical
archetypes based on the physical quantity its actuators manipulate
(Figure~\ref{fig-01-three-archetypes}): moving mass across space,
applying force through physical contact, or regulating continuous fluid
and thermal flows. By viewing physical AI through this three-archetype
lens, a systems architect can immediately determine which resource
budget is exhausted first during a software stall or model anomaly, and
design the real-time nervous system accordingly.

\subsection{Class 1: Mass-Dominated Mobility Systems (Moving Through
Space)}\label{class-1-mass-dominated-mobility-systems-moving-through-space}

Class 1 machines move mass across space, where execution is bound first
by momentum conservation and kinetic energy dissipation
(Figure~\ref{fig-01-three-archetypes} a). In autonomous transport
platforms, quadruped robots navigating rugged terrain, mobile
manipulators, and aerial drones, the actuator bridge applies force to
translate or rotate a body with mass \(m\) at velocity \(v\). Because
kinetic energy \(E_k = \frac{1}{2}m v^2\) cannot be recalled once
delivered to the chassis, safety is measured as remaining spatial
clearance. When braking begins at maximum allowable deceleration
\(a_{\max}\), the physical floor for stopping distance follows the
fundamental kinematic relation:

\[d_{\text{stop}} = v \cdot t_{\text{lag}} + \frac{v^2}{2a_{\max}}\]

where \(t_{\text{lag}}\) is the total sense-to-actuation transport and
inference delay. In a mass-moving system, computational latency converts
directly into unguided spatial displacement
(\(v \cdot t_{\text{lag}}\)), while braking distance scales
quadratically with velocity (\(v^2 / 2a_{\max}\)). The first boundary
breached during a software stall or model misclassification is the
physical distance to an external obstacle.

\subsection{Class 2: Contact-Dominated Manipulation Systems (Applying
Contact
Force)}\label{class-2-contact-dominated-manipulation-systems-applying-contact-force}

Class 2 machines apply mechanical force through contact, where the
primary binding limit is actuator torque saturation, current loop
bandwidth, and structural stiffness
(Figure~\ref{fig-01-three-archetypes} b). In articulated manipulators,
precision assembly fixtures, and robot learning workcells executing
contact-rich insertion and grasping, the end-effector interacts directly
with a rigid or compliant workpiece. Unlike a mobile platform operating
in open space, an arm in contact with a surface cannot resolve an
unexpected displacement command by rolling forward. When an actuator
commands position into a rigid boundary, displacement is constrained and
the interaction is governed by mechanical impedance. By Hooke's law, the
contact force generated by a spatial penetration error \(\Delta x\)
against an environment of stiffness \(k\) is
\(F = k \Delta x\).\sidenote{\footnotesize Rigid contact obeys Signorini-Coulomb
  complementarity (\(0 \le \delta \perp F_N \ge 0\)). Real-time stacks
  approximate compliance via Kelvin-Voigt models to prevent force loop
  divergence.}

Because contact stiffness in metal and tooling fixtures is extremely
high (\(k \ge 10^5\text{ N/m}\)), small geometric trajectory deviations
convert into massive contact forces within milliseconds. If the nervous
system fails to clamp motor current within sub-millisecond response
windows, the motor drives the transmission past its mechanical yield
stress. For example, a tooling tip moving at just \(0.20\text{ m/s}\)
that encounters a stiff workcell fixture will accumulate
\(30\text{ mm}\) of uncorrected displacement during a \(150\text{ ms}\)
software stall---generating thousands of Newtons of force that instantly
shears precision gear teeth. The critical floor in contact systems is
not stopping distance, but the time rate of force rise relative to
actuator current loop bandwidth and structural yield stress.

\subsection{Class 3: Flow-Dominated Process and Energy Systems
(Regulating Energy and Fluid
Flows)}\label{class-3-flow-dominated-process-and-energy-systems-regulating-energy-and-fluid-flows}

Beyond moving linkages and rigid contacts, Class 3 machines regulate
continuous fluid, thermal, or electrical power flows---a domain that
serves both as dedicated process plants and as the foundational
electrical and thermal substrate underlying all high-performance mobile
and manipulation systems (Figure~\ref{fig-01-three-archetypes} c). In
grid-scale battery storage facilities, datacenter liquid cooling loops,
and grid-tied power inverters managing renewable generation, actuators
manipulate coolant pumps, flow valves, and high-frequency power
semiconductors rather than mechanical joints. The state of the system is
governed by mass transport, thermal capacitance, and energy balance
rather than kinematic positions. When a coolant valve or pump stalls,
accumulated heat drives rapid temperature excursions toward chemical
decomposition thresholds. Similarly, a grid-tied inverter operating on a
\(60\text{ Hz}\) utility network faces an electrical phase cycle of
\(16.6\text{ ms}\) (half-cycle zero crossing of \(8.33\text{ ms}\)),
where an unhandled scheduling delay creates alternating-current phase
angle errors that drive destructive surge currents through substation
transformers. The deadline in flow architectures is set by thermal
capacitance, transport delay, and grid phase timing, defining physical
budgets where accumulated energy cannot be quickly reversed once
started.

\subsection{The Integrative Capstone: Multi-Class Embodiments and
Humanoid
Robotics}\label{the-integrative-capstone-multi-class-embodiments-and-humanoid-robotics}

A central question for physical AI systems architects is: \emph{Where do
complex general-purpose embodiments, such as humanoid robots, fit into
this tri-archetype framework?}

A humanoid robot is not a disjoint fourth archetype. Rather, it
represents the \textbf{Integrative Capstone Embodiment}---the ultimate
cyber-physical convergence where all three canonical physical classes
operate and contend simultaneously on a single, tightly constrained
mobile chassis:

\begin{enumerate}
\def\labelenumi{\arabic{enumi}.}
\tightlist
\item
  \textbf{Class 1 Dynamics in Locomotion:} Underactuated bipedal
  balancing, dynamic footstep placement, momentum conservation
  (\(m\mathbf{v}\)), and ground contact slip during walking and stair
  climbing require continuous spatial clearance and kinetic energy
  management.
\item
  \textbf{Class 2 Dynamics in Manipulation:} High-DOF bimanual arms and
  dexterous tactile hands interacting with rigid workpieces, opening
  doors, and grasping tools generate instantaneous contact wrenches
  (\(\mathbf{F} = k \Delta \mathbf{x}\)) bounded by gearbox yield stress
  and joint impedance control bandwidth.
\item
  \textbf{Class 3 Dynamics in Actuation \& Power Delivery:} Driving 20
  to 50 high-torque actuators simultaneously from a single onboard DC
  battery pack pushes inverter thermal dissipation (\(I^2 R\)) to its
  limits and creates severe inductive power-rail voltage droop
  (\(L \frac{dI}{dt}\)) during simultaneous leg shock absorption and
  heavy arm lifting.
\end{enumerate}

While a single-class machine (such as a wheeled AMR or a fixed
industrial arm) isolates its physical budget to a single dominant
conservation law, a humanoid robot forces the systems architect to
manage \textbf{cross-class coupled contention}. When a bipedal robot
catches a falling heavy object, the contact wrench (Class 2) shifts its
center of mass, requiring an instantaneous balancing recovery step
(Class 1) that draws peak phase currents threatening to collapse the
electrical bus and overheat the knee actuators (Class 3). Throughout
this book, we will use individual archetypes to isolate specific
architectural principles in depth, while using humanoid robotics as the
capstone testbed where all physical constraints collide.

\begin{figure}[t]

\centering{

\includegraphics[width=1\linewidth,height=\textheight,keepaspectratio]{/Users/VJ/GitHub/MLSysBook-integrate-physical/physical/book/chapters/01-boundary/figures/fig01_three_archetypes.jpg}

}

\caption{\label{fig-01-three-archetypes}\textbf{The Three Canonical
Machine Archetypes of Physical AI.} (a) Mass-dominated mobility systems
(such as the Boston Dynamics Spot quadruped robot navigating rugged
outdoor terrain) bound first by momentum, kinetic energy dissipation,
and spatial clearance. (b) Contact-dominated manipulation systems (such
as the Franka Emika Panda articulated arm) bound first by contact
impedance, joint torque saturation, and structural yield stress. (c)
Flow-dominated process and thermal plants (such as containerized
grid-scale battery storage systems) bound first by lumped thermal
capacity, transport lag, and electrical phase timing. Complex
embodiments like humanoid robots serve as the integrative capstone
uniting all three classes on a single chassis. (Sources: Wikimedia
Commons, CC BY-SA 4.0).}

\end{figure}%

For any physical AI system, engineering the nervous system begins by
determining which physical budget is exhausted first during a failure.
In software systems, resource constraints center on processor cycles,
memory allocation, and network bandwidth. In physical systems, the
systems architect must identify whether spatial clearance, structural
stress, or heat dissipation creates the critical timing deadline. In a
high-speed vehicle or quadruped, spatial clearance fails within
fractions of a second while motor winding temperatures remain safely
below thermal limits. In an assembly press, mechanical stress exceeds
the yield point of structural members within tens of milliseconds long
before the chassis moves perceptibly or heats up. In a chemical reactor
or battery pack, mechanical stress is negligible, but accumulated
thermal energy breaches safe operational envelopes over several seconds.
The shortest time to failure across these three domains defines the
deterministic execution deadline that the nervous system must meet.

Yet understanding that a machine belongs to one of these three physical
archetypes does not automatically mean it requires the full systems
discipline of Physical AI. A classical hydraulic press or a
deterministic PID drone autopilot manipulates mass and contact force,
but operates entirely within classical control theory. To determine when
a machine crosses into Physical AI, we must establish a rigorous
diagnostic scope test.

\section{The Physical AI Scope Test}\label{sec-boundary-1-4}

Before designing or deploying an embodied architecture, a system
architect or practitioner must establish whether a machine actually
requires the full systems engineering discipline of Physical AI. Without
a rigorous scope boundary, engineering teams make two symmetric
mistakes: they over-engineer classical deterministic machines with
bloated neural models where closed-form PID control is provably optimal,
or they dangerously under-engineer stochastic neural models by granting
them unverified motor authority in safety-critical physical loops.

The scope of Physical AI is defined by the simultaneous intersection of
three structural criteria (Figure~\ref{fig-01-scope-venn}):

\begin{enumerate}
\def\labelenumi{\arabic{enumi}.}
\tightlist
\item
  \textbf{Learned Statistical Model (Unspecifiable Policy):} The system
  relies on a policy, perception encoder, or world model acquired from
  data (e.g., neural networks, foundation models, diffusion policies)
  whose transfer function cannot be written in closed-form analytical
  equations and whose worst-case outputs cannot be mathematically
  bounded offline.
\item
  \textbf{Consequential Physical Feedback:} Actuator actions alter the
  material state of the physical environment, endogenously determining
  what the sensors observe next under the governing laws of classical
  mechanics, thermodynamics, or electromagnetism.
\item
  \textbf{Delegated Physical Authority (Direct Actuator Write):}
  Automated software routines write directly to hardware actuator
  registers without synchronous human approval in the control loop,
  operating within the temporal window where delay or error converts
  directly into irreversible physical momentum.
\end{enumerate}

\begin{figure}[t]

\centering{

\includegraphics[width=0.9\linewidth,height=\textheight,keepaspectratio]{index_files/mediabag//Users/VJ/GitHub/MLSysBook-integrate-physical/physical/book/chapters/01-boundary/figures/fig01_scope_venn.pdf}

}

\caption{\label{fig-01-scope-venn}\textbf{The Physical AI Scope Test.}
Physical AI exists strictly at the three-way intersection of learned
models (unspecifiable policies), consequential physical feedback
(actions govern future sensory observations under mechanics and
thermodynamics), and delegated physical authority (automated register
writes without synchronous human intervention). Systems satisfying
exactly two criteria define adjacent engineering fields: Advisory
Decision Support (no physical authority), Classical Control Theory
(analytical transfer functions with Lyapunov certificates), and Digital
Autonomous Systems (transactional databases with idempotency and
rollback). The shared autonomy reflex gap marks the operational boundary
where human neuromuscular reflex latency (\(150\text{--}250\text{ ms}\))
delegates transient authority to the learned model.}

\end{figure}%

A closed feedback loop is necessary for physical interaction, but
feedback alone does not place a machine in this domain. Classical
proportional-integral-derivative controllers run at \(1000\text{ Hz}\)
across aerospace and industrial automation, regulating valve positions
and motor torques through tight sensory feedback. Their feedback loops
are continuous and physical, but their control laws are specified
entirely by fixed mathematical formulas whose stability bounds can be
proven analytically from transfer functions. They hold physical
authority and experience physical feedback, but they lack a learned
model. Conversely, high-frequency algorithmic trading engines execute
automated market orders under learned predictive models within
microseconds. They combine learned policies with delegated transactional
authority, but the feedback they receive is purely digital and
informational. When an order executes, it alters an electronic order
book rather than momentum, friction, or thermal energy. An unhandled
stall in algorithmic trading costs financial capital, but it does not
cause an inertia-driven collision with a structural fixture.

A learned model is equally insufficient on its own to define the field.
Modern engineering workflows frequently deploy deep neural networks to
optimize robot trajectories offline or inspect manufactured parts on
assembly lines. An offline neural trajectory optimizer may compute
time-optimal paths for a six-axis robotic manipulator over several
seconds of cluster compute, but if that trajectory is verified by an
engineer or executed through a classical deterministic controller with
hard-coded envelope stops, the learned model holds no delegated physical
authority over the plant. Similarly, a convolutional neural network
monitoring an assembly line may classify surface defects at
\(60\text{ frames/s}\). If the classification output merely logs an
entry in a database or flags a display for a human operator, the model
operates outside the causal loop. A false negative or an inference
timeout produces a corrupted record, but it cannot directly command an
actuator to over-torque a mechanical joint.

Systems that satisfy exactly two criteria demonstrate where the
structural boundary lies. A real-time driver drowsiness monitor uses a
learned computer vision model to track eyelid movement and vehicle lane
position, receiving continuous physical feedback from the vehicle state.
Yet it lacks delegated physical authority because its only output is an
auditory warning chime. If the vision model fails completely during an
inference stall, the speaker emits no sound, but the steering rack and
brake calipers remain under human control. An analog bimetallic
thermostat possesses delegated physical authority over a gas furnace
valve and operates in continuous feedback with the thermal mass of a
building, but its transfer function is a fixed thermodynamic property of
two bonded metals rather than a learned representation. An automated
trading bot possesses both a learned policy and delegated execution
authority, but its actions propagate through network packets and
databases without interacting with the physical laws of mass and motion.

The most demanding edge cases occur in shared autonomy, where physical
authority is divided dynamically between a human operator and a learned
model. In robotic surgery or advanced driver assistance, a human
operator provides continuous steering or tool guidance while a learned
policy injects corrections to filter hand tremors, avoid anatomical
boundaries, or compensate for vehicle slip. These architectures satisfy
both the learned component and physical feedback criteria. Whether they
fall within the scope of Physical AI depends strictly on the temporal
authority granted to the model. If the learned policy can command even
\(2.0\text{ N}\) of force or adjust steering angle without requiring
human confirmation within the \(200\text{ ms}\) threshold of human
reflex latency, the model possesses delegated physical authority. During
that \(200\text{ ms}\) window, an erroneous model output or a delayed
inference packet acts on the physical world before the operator can
intercede.

Evaluating an ambiguous machine requires a three-step decision rubric
that tests each criterion independently
(Table~\ref{tbl-01-scope-matrix}). The first test determines whether the
system contains a policy component whose mapping is learned from data
rather than specified analytically in closed form. If it does not, the
machine belongs to classical control theory. The second test determines
whether actuator commands alter a physical plant such that subsequent
sensor observations are governed by mechanics, thermodynamics, or fluid
dynamics. If they do not, the machine is a digital machine learning
system. The third test determines whether the learned component writes
commands to physical actuators without a human supervisor or an
independent deterministic barrier capable of vetoing the action before
physical momentum accumulates. If it does not, the machine is an offline
decision-support tool. When all three tests pass, the system operates as
Physical AI, and its engineering must be governed by the laws of the
causal boundary.

\begin{longtable}[]{@{}
  >{\raggedright\arraybackslash}p{(\linewidth - 10\tabcolsep) * \real{0.1667}}
  >{\raggedright\arraybackslash}p{(\linewidth - 10\tabcolsep) * \real{0.1667}}
  >{\raggedright\arraybackslash}p{(\linewidth - 10\tabcolsep) * \real{0.1667}}
  >{\raggedright\arraybackslash}p{(\linewidth - 10\tabcolsep) * \real{0.1667}}
  >{\raggedright\arraybackslash}p{(\linewidth - 10\tabcolsep) * \real{0.1667}}
  >{\raggedright\arraybackslash}p{(\linewidth - 10\tabcolsep) * \real{0.1667}}@{}}
\caption{\textbf{The Physical AI Scope Test Decision Matrix.} Systematic
evaluation of representative engineering systems across the three joint
criteria (Learned Model, Consequential Physical Feedback, and Delegated
Physical Authority), establishing the governing engineering discipline,
limiting physical factors, and primary failure modes for each
candidate.}\label{tbl-01-scope-matrix}\tabularnewline
\toprule\noalign{}
\begin{minipage}[b]{\linewidth}\raggedright
System Candidate / Case
\end{minipage} & \begin{minipage}[b]{\linewidth}\raggedright
Criterion 1: Learned Model (Unspecifiable Policy)
\end{minipage} & \begin{minipage}[b]{\linewidth}\raggedright
Criterion 2: Consequential Physical Feedback
\end{minipage} & \begin{minipage}[b]{\linewidth}\raggedright
Criterion 3: Delegated Physical Authority (Direct Actuator Write)
\end{minipage} & \begin{minipage}[b]{\linewidth}\raggedright
Scope Classification \& Discipline
\end{minipage} & \begin{minipage}[b]{\linewidth}\raggedright
Limiting Physics \& Critical Failure Mode
\end{minipage} \\
\midrule\noalign{}
\endfirsthead
\toprule\noalign{}
\begin{minipage}[b]{\linewidth}\raggedright
System Candidate / Case
\end{minipage} & \begin{minipage}[b]{\linewidth}\raggedright
Criterion 1: Learned Model (Unspecifiable Policy)
\end{minipage} & \begin{minipage}[b]{\linewidth}\raggedright
Criterion 2: Consequential Physical Feedback
\end{minipage} & \begin{minipage}[b]{\linewidth}\raggedright
Criterion 3: Delegated Physical Authority (Direct Actuator Write)
\end{minipage} & \begin{minipage}[b]{\linewidth}\raggedright
Scope Classification \& Discipline
\end{minipage} & \begin{minipage}[b]{\linewidth}\raggedright
Limiting Physics \& Critical Failure Mode
\end{minipage} \\
\midrule\noalign{}
\endhead
\bottomrule\noalign{}
\endlastfoot
\textbf{Autonomous Warehouse AMR (Class 1)} & \textbf{Yes} (Learned
navigation policy / costmap) & \textbf{Yes} (Kinetic momentum, chassis
inertia, and wheel traction) & \textbf{Yes} (Direct motor drive packets
without human in loop) & \textbf{Physical AI} & Kinetic impact; traction
slip and odometry loss \\
\textbf{Articulated Contact Manipulator (Class 2)} & \textbf{Yes}
(Diffusion action chunking for assembly) & \textbf{Yes} (Mechanical
contact impedance and reaction forces) & \textbf{Yes} (Direct joint
torque and velocity register writes) & \textbf{Physical AI} & Structural
yield stress exceedance; gearbox tooth shear \\
\textbf{High-Speed Drone Classical Autopilot} & \textbf{No}
(Analytically tuned 1000 Hz PID / cascaded loop) & \textbf{Yes}
(Aerodynamic lift, drag, and gyroscopic dynamics) & \textbf{Yes} (Direct
PWM motor speed control) & \textbf{Classical Control Theory} & Control
margin erosion; rotor thrust saturation \\
\textbf{High-Frequency Market Maker} & \textbf{Yes} (Deep transformer
sequence predictor) & \textbf{No} (Digital order-book state; memory and
network queues) & \textbf{Yes} (Automated limit order execution over 10
GbE) & \textbf{Digital Machine Learning / FinTech} & Financial capital
loss; market flash crash (idempotent rollback) \\
\textbf{Advisory Driver Drowsiness Monitor} & \textbf{Yes} (Facial
landmark vision network) & \textbf{Yes} (Vehicle kinematics and driver
physiology) & \textbf{No} (Auditory warning buzzer and dashboard LED
only) & \textbf{Advisory Decision Support} & Human perception-reaction
latency; zero physical authority \\
\textbf{Offline Trajectory Optimizer} & \textbf{Yes} (Deep neural
network for time-optimal path search) & \textbf{No} (Static CAD workcell
model; offline batch execution) & \textbf{No} (Motion plan verified and
loaded via human engineer) & \textbf{Offline Optimization Tool} &
Numerical non-convergence; zero direct runtime plant coupling \\
\textbf{Shared-Autonomy Surgical Assist} & \textbf{Yes} (Learned haptic
tremor filter / virtual fixtures) & \textbf{Yes} (Anatomical tissue
compliance and reaction forces) & \textbf{Conditional} (Model injects
forces within human reflex window) & \textbf{Borderline / Shared
Physical AI} & Human neuromuscular reflex latency vs.~tool force rise
rate \\
\end{longtable}

\section{The Four Pillars of an Embodied Agent}\label{sec-boundary-1-5}

Before exploring how physical systems fail or examining the silicon
circuits that control them, we must map the complete anatomical
landscape of an embodied machine. Every physical AI system---whether a
humanoid biped, an autonomous tractor, or an automated surgical
console---is composed of four fundamental architectural pillars that
operate across distinct physical timescales:

\begin{enumerate}
\def\labelenumi{\arabic{enumi}.}
\tightlist
\item
  \textbf{Pillar 1: The Dynamical Body (Layer 1):} The physical plant
  itself. Comprising mechanical linkages, structural mass, gear
  transmissions, motor windings, power electronic inverters, and battery
  packs. The body does not obey software abstractions, but rather the
  governing conservation laws of classical mechanics, thermodynamics,
  and electromagnetism. It accumulates fatigue, dissipates resistive
  Joule heat (\(I^2 R\)), and resists instantaneous acceleration through
  reflected rotor inertia (\(N^2 J\)).
\item
  \textbf{Pillar 2: The Cognitive Brain (Layer 3):} The high-capacity
  learned intelligence. Comprising multi-modal vision-language-action
  (VLA) foundation models, diffusion policies, and semantic world models
  running on high-performance compute accelerators (GPUs, NPUs). The
  brain provides open-vocabulary generalization, spatial scene
  understanding, and long-horizon affordance reasoning. Its competence
  arises from high-dimensional statistical approximation across massive
  datasets, meaning it is fundamentally stochastic, non-deterministic,
  and prone to out-of-distribution hallucinations.
\item
  \textbf{Pillar 3: The Real-Time Nervous System (Layer 2):} The
  deterministic cyber-physical bridge. Comprising hard real-time
  microcontrollers (bare-metal MCUs), deterministic fieldbus networks
  (CAN-FD, EtherCAT, TSN Ethernet), and lock-free shared memory fabrics.
  The nervous system enforces sub-millisecond execution deadlines
  (\(1000\text{ Hz}\)), guarantees microsecond clock synchronization,
  evaluates candidate commands, and retains exclusive physical authority
  over actuator registers.
\item
  \textbf{Pillar 4: The Governance Envelope (Layer 4):} The formal
  safety, verification, and certification harness. Comprising
  forward-invariant safety monitors (Control Barrier Functions),
  decaying cryptographic operating leases, cross-layer hardware fault
  injection, and auditable Claim-Argument-Evidence (CAE) safety cases
  that qualify the complete machine for deployment.
\end{enumerate}

These four pillars form an integrated physical container. The Cognitive
Brain (Pillar 2) generates high-level intentions and proposed action
trajectories; the Real-Time Nervous System (Pillar 3) evaluates those
proposals against the physical constraints of the Dynamical Body (Pillar
1); and the Governance Envelope (Pillar 4) bounds the operational
envelope and provides formal assurance.

When any one of these four pillars is missing or compromised, the causal
boundary breaks. To see this catastrophic failure mode in forensic
detail, we examine one of the most famous real-world failures in
autonomous systems history.

\section{Anatomy of a Crash: When the Boundary
Fails}\label{sec-boundary-1-6}

To understand why the boundary conditions of Physical AI are not mere
theoretical classifications, but life-critical systems requirements, we
must examine what happens when an embodied machine violates them. When a
high-capacity learned model is granted unverified motor authority over
physical mass, software failure does not remain behind glass---it
converts into kinetic destruction. A physical crash in an autonomous
machine is rarely caused by an inaccurate prediction alone, but by the
absence of an architectural mechanism with the authority to refuse one.

\subsection{The 2018 Tempe Autonomous
Collision}\label{the-2018-tempe-autonomous-collision}

On a clear night in March 2018, an automated test vehicle---a modified
Volvo XC90 SUV operating Uber Advanced Technologies Group's
developmental Automated Driving System (ADS)---was traveling at
forty-three miles per hour (\(v_0 = 19.2\text{ m/s}\), mass
\(m \approx 2100\text{ kg}\)) along a public multi-lane roadway in
Tempe, Arizona (\citeproc{ref-ntsb2019uber}{National Transportation
Safety Board 2019})\sidenote{\footnotesize NTSB Report HAR-19/03
  (\citeproc{ref-ntsb2019uber}{National Transportation Safety Board
  2019}) revealed that disabling the factory autonomous emergency
  braking system eliminated the deterministic physical safety enforcer,
  relying on unmonitored human supervision.}. With a total kinetic
energy of:

\[E_k = \frac{1}{2} m v_0^2 = \frac{1}{2}(2100\text{ kg})(19.2\text{ m/s})^2 \approx 387\text{ kJ}\]

the moving vehicle struck and fatally injured a pedestrian pushing a
bicycle across the asphalt.

The vehicle was equipped with an expansive, state-of-the-art physical
sensor suite (Figure~\ref{fig-01-tempe-vehicle}): a roof-mounted 64-beam
Velodyne LiDAR (360° azimuth, \(> 100\text{ m}\) range), ten optical
cameras, and eight medium- and long-range radar units. At timestamp
\(T - 5.6\text{ s}\), while the vehicle was still \(82.0\text{ m}\)
away, the primary LiDAR and radar sensors established positive contact
with the obstacle, initializing a dynamic track in host computer memory.

\subsection{Perception Toggling and State
Erasure}\label{perception-toggling-and-state-erasure}

Despite continuous sensor contact over more than five seconds, the
vehicle failed to brake. Between \(T - 4.4\text{ s}\) and
\(T - 1.2\text{ s}\), the stochastic neural perception classifier
repeatedly toggled its semantic classification label between an
\emph{unknown vehicle}, a \emph{bicycle}, and an \emph{unclassified
static obstacle}.

Because the software tracking architecture associated Kalman filter
state estimators with discrete semantic categories, each classification
label switch caused the runtime to discard the existing kinematic track
and re-initialize a new filter instance with an assumed velocity vector
of zero. As a direct consequence of this classification instability, the
downstream trajectory planner never accumulated persistent evidence of
an imminent collision course.

\subsection{Latency Suppression and Kinematic
Depletion}\label{latency-suppression-and-kinematic-depletion}

At \(T - 1.2\text{ s}\), the software finally consolidated the track but
triggered an internal \(1.2\text{ s}\) ``suppression delay''---a
software counter designed by developers to prevent false-positive
emergency braking alarms from disrupting passenger comfort.

During this suppression window, the vehicle traversed an unguided
displacement of:

\[d_{\text{lag}} = v_0 \cdot t_{\text{suppress}} = (19.2\text{ m/s}) \times (1.2\text{ s}) \approx 23.0\text{ m}\]

When the planner finally completed the suppression window and requested
emergency braking at \(T - 0.2\text{ s}\), the remaining physical
clearance had collapsed to just \(3.8\text{ m}\). Bringing a
\(2100\text{ kg}\) vehicle to a complete halt on dry asphalt (assuming a
maximum deceleration \(a_{\max} \approx 8.0\text{ m/s}^2\)) requires a
physical braking distance of:

\[d_{\text{brake}} = \frac{v_0^2}{2 a_{\max}} = \frac{(19.2\text{ m/s})^2}{2 \times 8.0\text{ m/s}^2} \approx 23.0\text{ m}\]

Total stopping distance required at \(T - 1.2\text{ s}\) exceeded
\(d_{\text{stop}} = d_{\text{lag}} + d_{\text{brake}} = 23.0\text{ m} + 23.0\text{ m} = 46.0\text{ m}\).
Because available clearance had eroded to \(3.8\text{ m}\), kinetic
collision at full cruising speed was mathematically sealed two full
seconds before impact.

\begin{figure}[t]

\centering{

\includegraphics[width=0.95\linewidth,height=\textheight,keepaspectratio]{/Users/VJ/GitHub/MLSysBook-integrate-physical/physical/book/chapters/01-boundary/figures/fig01_real_ntsb_tempe_vehicle.jpg}

}

\caption{\label{fig-01-tempe-vehicle}\textbf{Autonomous Test Vehicle
Sensor Configuration from NTSB Collision Investigation.} The modified
2017 Volvo XC90 sport utility vehicle operating Uber ATG's developmental
Automated Driving System (ADS) during the March 2018 fatal collision in
Tempe, Arizona. The physical sensor suite comprised a roof-mounted
64-beam Velodyne LiDAR (360° azimuth, \textgreater100 m range), a
10-camera optical sensing suite including forward-facing stereo and
wide-angle cameras, and eight medium- and long-range radar units
providing 360° perimeter coverage. Despite continuous sensor contact
starting at 82 meters (\(T-5.6\text{ s}\)), classification instability
in the perception pipeline and an unverified 1.2-second suppression
delay prevented the planner from requesting braking until 0.2 seconds
before impact. (Source: NTSB Highway Accident Report NTSB/HAR-19/03
(\citeproc{ref-ntsb2019uber}{National Transportation Safety Board
2019}), Public Domain).}

\end{figure}%

\protect\phantomsection\label{callout-autopsyux2a-1.5}
\begin{fbxSimple}{callout-autopsy}{Incident Autopsy:}{The Missing Safety Referee}
\phantomsection\label{callout-autopsy*-1.5}

\begin{itemize}
\tightlist
\item
  \textbf{System:} Developmental Automated Driving System (ADS) on a
  2017 Volvo XC90 SUV (\(2100\text{ kg}\), \(19.2\text{ m/s}\)) equipped
  with roof-mounted 64-beam LiDAR, radar, and optical cameras.
\item
  \textbf{Incident:} Struck and fatally injured a pedestrian pushing a
  bicycle across an urban roadway without executing emergency braking
  before impact.
\item
  \textbf{Telemetry Snapshot:} Sensor contact established at
  \(T - 5.6\text{ s}\) (\(82\text{ m}\)); classifier toggled across
  three semantic labels for \(3.2\text{ s}\); \(1.2\text{ s}\)
  false-alarm suppression delay expired at \(T - 0.2\text{ s}\)
  (\(3.8\text{ m}\) remaining clearance).
\item
  \textbf{Root Cause:} Semantic classification toggling caused the
  tracker to repeatedly discard kinematic state histories and
  re-initialize velocity estimators to zero, preventing the trajectory
  planner from recognizing an imminent collision.
\item
  \textbf{Architectural Flaw:} Complete absence of an isolated
  deterministic safety referee running on dedicated real-time silicon
  with the authority to evaluate geometric clearance invariants directly
  against raw sensor points.
\item
  \textbf{Provenance:} National Transportation Safety Board Accident
  Report NTSB/HAR-19/03 (2019).
\end{itemize}

\end{fbxSimple}

This investigation autopsy demonstrates why refining perception models
or expanding offline training datasets cannot compensate for an
architectural absence of authority boundaries. An open-world statistical
model will inevitably encounter out-of-distribution scenes, sensor
degradations, and transient classification oscillations. The Tempe
vehicle crashed not because its neural network was imperfect, but
because the autonomy architecture lacked an isolated nervous system
running on dedicated real-time silicon with the independent authority to
evaluate spatial clearance and execute an emergency stop. If a
deterministic safety enforcer had continuously checked the stopping
clearance invariant at the hardware floor, the machine would have
executed emergency braking at \(T - 2.5\text{ s}\), bringing the vehicle
to a controlled halt long before available clearance
collapsed.\sidenote{\footnotesize ISO 26262-9 Clause 5 defines Freedom From
  Interference (FFI) across ASIL boundaries. Proving FFI requires
  verified spatial partitioning (MPU/IOMMU memory protection), temporal
  partitioning (hardware bandwidth limiters and watchdog timers), and
  authenticated communication contracts.}

The forensic breakdown of the Tempe crash reveals four universal bedrock
laws that govern all physical machines across the causal boundary.

\section{The Four Bedrock Laws}\label{sec-boundary-1-7}

Four bedrock laws govern every machine whose nervous system translates
compute cycles into physical action. Notice how these four laws map
symmetrically to the \textbf{Four Pillars of an Embodied Agent} and the
\textbf{Four Parts of this Textbook}:

\protect\phantomsection\label{callout-lawux2a-1.6}
\begin{fbxSimple}{callout-law}{Bedrock Law:}{1. Physical Causality and Irreversibility (The Body Law — Part I)}
\phantomsection\label{callout-law*-1.6}

Software state can be caught with an exception block, isolated in a
virtual container, or restored from a checkpoint. Physical actions
governed by mass, inertia, friction, and gravity are permanent and
irreversible. Every actuation permanently alters the physical
environment, dissipates energy, and alters all future sensory
observations endogenously.

\end{fbxSimple}

The first bedrock law is the thermodynamic irreversibility of physical
action. In a pure software system, an erroneous operation is reversed by
rolling back a database transaction, restoring a file from a snapshot,
or restarting a process with clean memory. In a physical machine, an
actuator converts electrical charge into kinetic energy, structural
strain, or fluid pressure. Once that energy crosses the causal boundary
into the environment, it cannot be recalled. When an erroneous command
accelerates a warehouse robot toward an obstacle, the kinetic energy
deposited into the chassis can only be counteracted by applying opposing
forces through friction brakes, which dissipates the energy as heat. If
an over-current command drives a six-axis manipulator into an assembly
jig, the mechanical stress exceeds the yield point of the metal within
milliseconds, permanently deforming the tooling. Software transactions
can be rolled back, but thermodynamic work cannot be undone.

\protect\phantomsection\label{callout-lawux2a-1.7}
\begin{fbxSimple}{callout-law}{Bedrock Law:}{2. Endogenous Data and Compounding Drift (The Policy Law — Part II)}
\phantomsection\label{callout-law*-1.7}

A physical agent shapes its own future data distribution. Naive
behavioral cloning on open-loop benchmarks suffers compounding covariate
drift (\(\mathcal{O}(T^2 \epsilon)\)) in closed-loop execution. An error
on step \(t\) creates an out-of-distribution observation at step
\(t+1\), causing policy confidence to collapse into mechanical
divergence.

\end{fbxSimple}

The second bedrock law governs how embodied policies learn and fail. In
traditional computer vision or natural language processing, data is
exogenous: whether a model correctly identifies an image has zero effect
on the distribution of subsequent images. In Physical AI, every torque
command rotates joints, shifts camera angles, and alters contact forces.
The data stream is endogenous---created entirely by the machine's own
past actions. A single erroneous inference drives the hardware into an
unvisited state where neural predictions exhibit higher variance,
compounding errors until the machine exceeds physical stability limits.

\protect\phantomsection\label{callout-lawux2a-1.8}
\begin{fbxSimple}{callout-law}{Bedrock Law:}{3. Proposal–Permission Privilege and Freshness (The Nervous System Law — Part III)}
\phantomsection\label{callout-law*-1.8}

No stochastic neural network may hold direct motor authority. Learned
models running on application processors are unprivileged proposal
services. Real-time physical permission belongs exclusively to dedicated
microcontrollers executing deterministic safety filters at 1000 Hz.
Latency is physical displacement; stale proposals older than their
freshness deadline must be discarded.

\end{fbxSimple}

The third bedrock law requires strict separation between proposal and
permission. The subsystem of the nervous system that generates complex,
goal-directed behavior operates by optimization over high-dimensional
learned representations. Because deep neural networks are
uninterpretable and produce erratic outputs under out-of-distribution
inputs, a learned policy must never hold direct write access to physical
actuator registers. The role of the learned model is strictly to propose
candidate trajectories. The authority to grant permission, clip commands
to safe physical envelopes, or reject the proposed action entirely must
belong to a separate, deterministic supervisory layer running on
isolated execution hardware.

\protect\phantomsection\label{callout-lawux2a-1.9}
\begin{fbxSimple}{callout-law}{Bedrock Law:}{4. Governed Authority and Defensible Release (The Governance Law — Part IV)}
\phantomsection\label{callout-law*-1.9}

Physical AI systems are never `generally safe'; operating authority is a
decaying lease strictly conditioned on bounded evidence. A cherry-picked
video demonstration is not evidence of safety; defensible release
requires cross-layer hardware fault injection, formal barrier
guarantees, and an auditable safety case.

\end{fbxSimple}

The fourth bedrock law dictates that deployment readiness cannot be
justified by empirical validation loss. In cloud computing, software
reliability is evaluated statistically: an engineer measures mean
accuracy, accepts an uptime target of \(99.9\%\), and resolves rare
corner-case crashes through automated restarts and rolling updates. In
physical AI, the tail of the error distribution manifests as mechanical
destruction. Safe operation cannot be certified by average-case
performance metrics gathered over finite training sets. It requires
formal, worst-case envelope bounds verified against the physical laws of
the plant, ensuring that even if the learned proposal violates safety
limits, the machine remains within guaranteed physical boundaries.

The four bedrock laws establish the physical constraints that govern
every machine. But of these four, the Law of Proposal--Permission
Privilege dictates how we must architect and physically build the
computer system. How do we take an untrusted, stochastic neural policy
and safely translate its high-level intelligence into physical motion on
real hardware?

\section{The Dual-Brain Architecture}\label{sec-boundary-1-8}

To safely bridge high-capacity learned policies and real-time physical
machines, the systems architect implements the
\textbf{Proposal--Permission Architecture} across heterogeneous silicon.
The system partitions compute into two physically separated execution
domains: an unprivileged application processor (Linux MPU or Edge NPU)
generating candidate action proposals at deliberative cadences
(\(10\text{--}50\text{ Hz}\)), paired with an isolated, zero-allocation
real-time microcontroller (bare-metal MCU) executing deterministic
safety filters and physical fallbacks at \(1000\text{ Hz}\) with
exclusive electrical authority over actuator drive stages.

Authority over actuation must belong to a mechanism that can refuse
execution even when the proposer is broken. This separation between
proposal and permission is a foundational structural principle: in a
pure software system, an unprivileged process can be granted direct
write access to application memory because an erroneous write costs only
compute cycles and memory bandwidth to correct. In a physical machine,
state changes move mass, dissipate heat, and expend momentum that cannot
be recalled. Because a learned policy is an inductive statistical model
whose worst-case outputs under distribution shift cannot be
mathematically bounded, the neural network must function strictly as an
unprivileged proposer. Routing learned model outputs directly to
actuator registers couples the mechanical integrity of the machine to
the unverified tail of a probability distribution.

Training-time optimization cannot substitute for runtime invariant
enforcement. In reinforcement learning and imitation learning pipelines,
developers frequently attempt to enforce safety by augmenting the loss
function with penalty terms, barrier losses, or negative reward shaping.
These numerical penalties influence gradient descent during training,
reducing the empirical frequency of undesirable states across the
training distribution. But they provide zero hard guarantees at runtime.
When the deployed policy encounters novel sensor noise, lens occlusion,
or out-of-distribution state combinations, offline optimization cannot
prevent the network from generating an unsafe control vector. Soft
optimization incentives shape probability distributions, but physical
safety requires deterministic boundary enforcement. Safety constraints
must therefore live downstream of the model as an explicit, hard
architectural gate positioned between the policy output and the actuator
registers.

In computer systems engineering, developers often ask: \emph{why can't
we run this proposal--permission gate inside the same high-performance
Linux processor that runs the neural network?} The answer lies in the
fundamental conflict between high-capacity statistical computing and
hard real-time determinism. A modern vision-language-action model or
diffusion policy requires billions of floating-point operations and
gigabytes of memory bandwidth, necessitating multi-core application
processors (MPUs) or neural accelerators (NPUs) running complex
operating systems like Linux. However, general-purpose operating systems
introduce non-deterministic scheduling preemption, dynamic memory
allocation locks, and interrupt latencies ranging from \(10\text{ ms}\)
to over \(100\text{ ms}\). In a physical system, a \(50\text{ ms}\) OS
scheduling hiccup is not a dropped animation frame; it is
\(50\text{ ms}\) of unguided kinetic displacement that drives mechanical
linkages past structural limits. Conversely, one cannot run
high-capacity neural models inside a microcontroller (MCU), because
microcontrollers lack the memory bandwidth and tensor execution units
required for high-dimensional inference.

Therefore, every robust physical AI machine physically partitions its
compute across heterogeneous silicon (Figure~\ref{fig-01-dual-brain}):

\begin{enumerate}
\def\labelenumi{\arabic{enumi}.}
\tightlist
\item
  \textbf{The Cognitive Proposer (Linux MPU / Edge NPU):} Operates in
  unprivileged user space at a deliberative cadence
  (\(10\text{--}50\text{ Hz}\)). It executes high-capacity neural
  policies, multi-modal perception backbones, and trajectory generators.
  It possesses \textbf{zero direct electrical wiring to physical
  actuator registers}. It can only write candidate trajectory proposals
  into a shared SRAM mailbox or lock-free circular ring buffer across an
  inter-processor communication (IPC) channel.
\item
  \textbf{The Real-Time Permission Referee (Bare-Metal / RTOS MCU):}
  Operates on isolated, dedicated silicon at a hard real-time cadence
  (\(1000\text{ Hz}\), \(1\text{ ms}\) deadline,
  \(< 50\text{ }\mu\text{s}\) jitter) with static, zero-allocation
  memory (\texttt{malloc} is strictly forbidden at runtime). It holds
  \textbf{exclusive physical authority} over the PWM compare registers,
  digital-to-analog converters, and gate drivers. The MCU continuously
  checks physical safety invariants and stopping clearance. If the
  proposed command satisfies all physical constraints, the MCU latches
  the values to hardware; if the proposal is unsafe, stalls, or panics,
  the MCU vetoes the command and executes a deterministic physical
  fallback.
\end{enumerate}

To bridge the asynchronous Linux OS and deterministic bare-metal
firmware without priority inversion or lock contention, the MPU and MCU
communicate across a shared SRAM mailbox using lock-free atomic sequence
locks (\emph{seqlocks}) or double-buffered ping-pong ring
buffers.\sidenote{\footnotesize In heterogeneous SoCs (e.g., TI Sitara AM64x, NXP
  i.MX8M), Linux MPUs and real-time Cortex-M/R MCUs exchange data across
  non-cached dual-port SRAM via OpenAMP/RPMsg mailboxes. Out-of-order
  MPU cores require explicit Data Memory Barriers (\texttt{DMB}) and
  cache flushes before asserting Inter-Processor Interrupts
  (\texttt{IPI}) to prevent torn memory reads.} The MPU publishes
candidate action chunks with a monotonically incrementing sequence
version; the MCU reads the trajectory atomically without ever blocking
on a mutex. If an MPU core stalls or experiences an OS scheduling delay,
the sequence version fails to update, immediately signaling the MCU that
proposed commands have aged beyond their freshness deadline.

\begin{figure}[t]

\centering{

\includegraphics[width=1\linewidth,height=\textheight,keepaspectratio]{index_files/mediabag//Users/VJ/GitHub/MLSysBook-integrate-physical/physical/book/chapters/01-boundary/figures/fig01_dual_brain_soc.pdf}

}

\caption{\label{fig-01-dual-brain}\textbf{The Physical AI Dual-Brain
Privilege Separation Architecture.} Heterogeneous silicon partitioning
between the untrusted proposal engine (Linux MPU / Edge NPU) executing
high-capacity foundation models and the trusted permission authority
(bare-metal real-time MCU) executing zero-allocation 1 kHz safety
filters over shared SRAM mailboxes. The permission authority exercises
sole veto over actuator phase current setpoints, guaranteeing physical
invariants even during total model stalls or numerical divergence.}

\end{figure}%

\subsection{The Three Real-Time Referee
Gates}\label{the-three-real-time-referee-gates}

Because a machine learning model running on an application processor is
strictly a proposal generator, exclusive authority to commit physical
register writes belongs to the real-time referee. The referee operates
on dedicated microcontroller silicon at a fixed \(1\text{ kHz}\)
interrupt cadence (\(1.0\text{ ms}\) deadline) and enforces three
mandatory gates before latching any command to motor drivers:

\begin{enumerate}
\def\labelenumi{\arabic{enumi}.}
\tightlist
\item
  \textbf{Temporal Freshness Gate:} If the proposal timestamp is older
  than the strict freshness deadline
  (\(\Delta t > \Delta t_{\text{fresh}}\)), the referee drops the stale
  proposal.
\item
  \textbf{Physical Invariant Gate:} If the commanded torque, velocity,
  or trajectory breaches dynamic stopping clearance
  (\(d_{\text{stop}}\)) or current limits, the referee mathematically
  projects the candidate command back to the safe admissible boundary.
  Because microcontrollers have strict cycle budgets, this projection is
  computed using lightweight, closed-form Control Barrier Functions
  (CBFs) or box-constrained quadratic programs (QP) solved in
  microseconds via static, zero-allocation active-set routines without
  non-deterministic matrix inversions.
\item
  \textbf{Deterministic Fallback Gate:} If the neural proposer crashes,
  hangs, or outputs numerical anomalies (\(\text{NaNs} / \pm\infty\)),
  the referee autonomously engages verified regenerative braking or
  passive dynamic centering.
\end{enumerate}

When a candidate action is refused, the permission authority must
execute a deterministic intervention strategy that maintains the
physical machine within its safe operating envelope. Refusal cannot
consist of simply dropping the command and allowing actuator registers
to float. If a mobile base moving at \(2.0\text{ m/s}\) suddenly
receives zero motor current without controlled deceleration, momentum
drives the unconstrained mass into obstacles or induces wheel slip that
invalidates odometry. The permission authority applies a graded
hierarchy of interventions matched to the severity of the violation:

\begin{itemize}
\tightlist
\item
  \emph{Rate Clamping \& Admissible Projection:} For minor envelope
  exceedances, the gate mathematically projects the candidate vector to
  the nearest admissible boundary point, enforcing current limits and
  voltage ceilings.
\item
  \emph{Kinematic Slicing:} For dynamic trajectory proposals, the gate
  clips acceleration and jerk to maximum allowable mechanical
  transmission bounds.
\item
  \emph{Deterministic Physical Fallback:} When the proposer suffers a
  numerical exception (NaNs/Infinities), latency timeout, or heartbeat
  loss, the MCU engages an autonomous physical fallback---commanding
  controlled regenerative braking or executing a verified deceleration
  trajectory to bring the mass to a stationary, passive state.
\end{itemize}

Understanding this dual-brain realization of proposal and permission is
essential because it fundamentally restructures how we develop, verify,
and govern the entire machine. In traditional software services,
engineering teams treat functional safety, cyber-security, and audit
logging as three completely isolated disciplines handled by separate
organizations across disjoint software layers. In Physical AI, the
dual-brain proposal--permission boundary collapses all three concerns
onto a single physical primitive.

\section{Safety Convergence}\label{sec-boundary-1-9}

Safety asks whether the machine caused harm. Security asks whether an
external agent forced it to. Accountability asks whether an auditor can
determine afterwards what occurred and why. In software services, these
three concerns belong to separate engineering organizations and distinct
runtimes, dividing along application logic, network firewalls, and audit
logging databases. In a machine that moves mass and commands current,
all three questions collapse into a single physical primitive at the
causal boundary, determining whether an independent gate possessed the
state and the authority to refuse a register write before torque reached
the plant.

Safety failures arise without an adversary. A camera lens accumulates
dust over \(500\text{ h}\) of operation, shifting the input distribution
away from the training domain until an attention head generates an
out-of-distribution steering torque. A joint bearing develops mechanical
play, introducing unmodeled compliance that converts a stable learned
position trajectory into a resonant limit cycle. Or a sudden change in
surface friction drops tire traction below the friction coefficient
assumed by the learned policy. In each case, the proposer generates an
action that is statistically plausible to the network but physically
hazardous to the body. If the permission gate lacks the independent
kinematic model or thermal budget to detect the exceedance, the
unrefused command writes to the motor drive registers, and stored
kinetic energy is released into the workspace.

Security failures introduce an adversary who deliberately exploits the
gap between statistical inference and physical reality. An attacker does
not need to compromise cryptographic keys or gain root access on the
host operating system if they can manipulate sensory inputs. Projected
structured light can spoof a depth sensor, acoustic injection can
resonate MEMS gyroscopes at their resonant frequency of
\(27\text{ kHz}\) to corrupt inertial odometry,\sidenote{\footnotesize Resonant
  acoustic frequencies (\(18\text{--}32\text{ kHz}\)) can saturate MEMS
  gyroscopes, requiring physical damping and MCU kinematic cross-checks.}
or adversarial visual perturbations on a factory floor can induce a
vision-language-action model to output maximum actuator velocity. On a
traditional server, an injection attack compromises data confidentiality
or integrity within an abstract memory space. In physical AI, an
injection attack commands physical work, turning an unvalidated
inference into mechanical force. The security perimeter cannot terminate
at the network interface or the operating system system-call boundary.
It terminates at the permission gate, which must evaluate every
candidate command against physical invariants regardless of whether the
proposer acted in good faith or under adversarial control.

Traditional software authorization models fail in physical machines
because capability tokens and access control lists assume discrete,
memoryless transactions. In a web service, possessing a valid credential
grants permission to write a record to a database, and the cost of the
operation is bounded by compute cycles and storage bytes. In physical
space, permission is not a static binary entitlement. In an analytical
robotic joint, a command to apply \(50\text{ N}\cdot\text{m}\) of torque
is valid when the link is stationary at mid-stroke, but the identical
command causes structural yield if the joint is \(2\text{ mm}\) from its
hard stop moving at \(1.5\text{ rad/s}\). Physical authority is
state-dependent, continuous, and dynamic, governed by velocity,
momentum, thermal accumulation, and spatial occupancy. A permission gate
cannot evaluate an action in isolation as a stateless remote procedure
call. It must integrate the action over time against the current
physical state of the body, validating that the applied forces remain
within the safe invariant reachable set of the plant.

Accountability requires proving the exact causal chain that led to
actuation when a machine damages hardware or breaches an envelope. In
standard computing systems, logging is asynchronous, best-effort, and
coarse-grained, dropping records during buffer pressure or recording
timestamps with millisecond jitter. At the causal boundary,
accountability demands deterministic, high-rate telemetry captured
synchronously with the control loop. If a joint actuator commands
\(1000\text{ Hz}\) updates with a \(1\text{ ms}\) deadline, the
permission gate must record the raw candidate vector proposed by the
model, the state estimate used by the gate, the mathematical criterion
applied by the safety filter, and the exact vector written to the
hardware registers. Sealing these records in tamper-evident, append-only
circular buffers in non-volatile SRAM/FRAM provides post-incident
reconstruction. When an anomaly occurs, an engineer can verify
mathematically whether the failure originated in a corrupted sensor
reading, a degraded policy proposal, an erroneous permission rule, or a
mechanical component that failed to produce the commanded
counter-torque.

\section{Architecture of This Book: The Journey
Ahead}\label{sec-boundary-1-10}

Engineering dependable physical AI systems requires mastering how
statistical learning, physical dynamics, and hard real-time computer
systems compose across every tier of the machine. To build this
discipline systematically from first principles, this textbook is
organized as an unbroken intellectual journey across four foundational
movements (Table~\ref{tbl-01-book-roadmap}):

\subsection{The Four Movements of Physical
AI}\label{the-four-movements-of-physical-ai}

\textbf{Part I: The Physical Container (Chapter~\ref{sec-boundary}
through Chapter~\ref{sec-nervous}).} Before writing a single line of
policy code or launching a neural network training cluster, an engineer
must first step onto the hardware floor. We begin by characterizing what
the physical world will not negotiate---the unyielding laws of momentum,
reflected gear inertia, thermal dissipation, and microsecond clock
jitter that bound the body, the brain, and the real-time nervous system.
By treating the physical plant as an immutable container, we establish
the non-negotiable temporal and dynamic budgets that every upstream
algorithm must satisfy.

\textbf{Part II: The Policy Lifecycle (Chapter~\ref{sec-data} through
Chapter~\ref{sec-evaluation}).} Once the physical container is defined,
we confront where an embodied policy's intelligence comes from. We trace
the messy, endogenous reality of physical data collection, evaluate why
naive behavioral cloning suffers compounding drift in closed-loop
deployment, and establish why statistical benchmark accuracy is
powerless to guarantee physical survival without rigorous,
non-asymptotic testing and an understanding of the Butler \& Finelli
exposure wall.

\textbf{Part III: The Runtime Architecture (Chapter~\ref{sec-perception}
through Chapter~\ref{sec-placement}).} With a trained policy in hand, we
construct the live, multi-rate software engine that runs aboard the
machine. We follow the continuous flow of signals from photons striking
image sensors, through 3D spatial memory representations and expiring
intent leases, down through receding-horizon trajectory planning and the
bare-metal safety shields that execute on heterogeneous silicon without
starving under bus contention or voltage droop.

\textbf{Part IV: Governance and Defensible Release
(Chapter~\ref{sec-intervention} through Chapter~\ref{sec-frontier}).} In
the final movement, we address the human and societal reality of
physical embodiment. We establish how machines safely hand authority
back to human supervisors without mode confusion, how we break the
assembled system through systematic hardware fault injection across the
qualification ladder, how engineering teams construct auditable
Claim-Argument-Evidence safety cases under emerging international
standards, and how we confront the fundamental epistemic limits of what
AI can and cannot know in the physical world.

\begin{longtable}[]{@{}
  >{\raggedright\arraybackslash}p{(\linewidth - 8\tabcolsep) * \real{0.2000}}
  >{\raggedright\arraybackslash}p{(\linewidth - 8\tabcolsep) * \real{0.2000}}
  >{\raggedright\arraybackslash}p{(\linewidth - 8\tabcolsep) * \real{0.2000}}
  >{\raggedright\arraybackslash}p{(\linewidth - 8\tabcolsep) * \real{0.2000}}
  >{\raggedright\arraybackslash}p{(\linewidth - 8\tabcolsep) * \real{0.2000}}@{}}
\caption{\textbf{The Architecture of Physical AI Systems:
Chapter-by-Chapter Systems Matrix.} Complete structural roadmap of the
textbook across the four foundational parts, mapping each chapter to its
operational focus, primary physical/silicon constraint, and core
engineering mechanism.}\label{tbl-01-book-roadmap}\tabularnewline
\toprule\noalign{}
\begin{minipage}[b]{\linewidth}\raggedright
Part
\end{minipage} & \begin{minipage}[b]{\linewidth}\raggedright
Chapter
\end{minipage} & \begin{minipage}[b]{\linewidth}\raggedright
Subsystem / Focus
\end{minipage} & \begin{minipage}[b]{\linewidth}\raggedright
Core Constraint or Limit
\end{minipage} & \begin{minipage}[b]{\linewidth}\raggedright
Key Systems Mechanism
\end{minipage} \\
\midrule\noalign{}
\endfirsthead
\toprule\noalign{}
\begin{minipage}[b]{\linewidth}\raggedright
Part
\end{minipage} & \begin{minipage}[b]{\linewidth}\raggedright
Chapter
\end{minipage} & \begin{minipage}[b]{\linewidth}\raggedright
Subsystem / Focus
\end{minipage} & \begin{minipage}[b]{\linewidth}\raggedright
Core Constraint or Limit
\end{minipage} & \begin{minipage}[b]{\linewidth}\raggedright
Key Systems Mechanism
\end{minipage} \\
\midrule\noalign{}
\endhead
\bottomrule\noalign{}
\endlastfoot
\textbf{Part I: The Constraints} & Chapter~\ref{sec-boundary} & The
Causal Boundary \& Scope & Idempotency loss; actuation registers &
Proposal--Permission dual-brain contract \\
& Chapter~\ref{sec-body} & Mechanical \& Thermal Dynamics & Reflected
inertia (\(N^2 J\)); Joulean heat & Dynamic stopping envelope
(\(d_{\text{stop}}\)) \\
& Chapter~\ref{sec-brain} & Cognitive Deliberation & Memory wall
(\(M_{\text{weights}}/B_{\text{mem}}\)) & VLA models \& action
chunking \\
& Chapter~\ref{sec-nervous} & Real-Time Computing Bridge & Multi-rate
deadlines \& jitter & Lock-free \texttt{seqlock} \& 5x5 synthesis
grid \\
\textbf{Part II: Teaching the System} & Chapter~\ref{sec-data} &
Endogenous Telemetry & Teleoperation lag; covariate shift & 4-stream
synchronous logging ring buffer \\
& Chapter~\ref{sec-training} & Policy Synthesis Regimes & Compounding
drift (\(\mathcal{O}(T^2 \epsilon)\)) & Diffusion policy \& sim-to-real
domain rand. \\
& Chapter~\ref{sec-evaluation} & Closed-Loop Survival Testing & Butler
\& Finelli exposure wall & Clopper-Pearson exact confidence bounds \\
\textbf{Part III: Running the System} & Chapter~\ref{sec-perception} &
Spatial Perception Waterfall & Sensor transduction lag \& shear &
Geometric clearance insets \& depth cones \\
& Chapter~\ref{sec-memory} & Persistent Spatial Belief & Dynamic
occlusion; transform staleness & TSDF volumetric fusion \& temporal
decay \\
& Chapter~\ref{sec-intent} & Grounded Task Affordances & Semantic
ambiguity \& open horizons & Expiring Intent Leases
(\(\mathcal{L}_{\text{intent}}\)) \\
& Chapter~\ref{sec-planning} & Trajectory Synthesis \& Seams & Seam
torque spikes; execution lag & \(\mathcal{C}^2\) jerk spline blending \&
stop suffixes \\
& Chapter~\ref{sec-enforcement} & Deterministic Safety Shield &
Microcontroller cycle budget (\(1\text{ ms}\)) & Bare-metal 1 kHz CBF-QP
safety filter \\
& Chapter~\ref{sec-placement} & Heterogeneous Silicon QoS & Crossbar
starvation; \(L dI/dt\) droop & AXI QoS bandwidth arbitration \&
fences \\
\textbf{Part IV: Governance \& Assurance} &
Chapter~\ref{sec-intervention} & Shared Human-Machine Autonomy &
Operator reaction delay (\(1\text{ s}\)) & 4-phase handshake \&
\(\mathcal{C}^1\) bumpless transfer \\
& Chapter~\ref{sec-verification} & Hardware Qualification Ladder &
Untested environmental states & SIL \(\to\) PIL \(\to\) HIL fault
injection harness \\
& Chapter~\ref{sec-release} & Safety Cases \& Deployment & Open-world
ODD boundary drift & Claim-Argument-Evidence trees \& canary OTA \\
& Chapter~\ref{sec-frontier} & Epistemic Limits \& Research &
Detectorless physical failure gaps & Multi-timescale governor \&
epistemic bounds \\
\end{longtable}

\section{Systems Stress-Tests}\label{sec-boundary-1-11}

With the core architectural principles, physical archetypes, bedrock
laws, and dual-brain silicon partitioning established, an engineer must
be able to diagnose complex, borderline systems in the field. The
following four systems stress-tests serve as capstone analytical
exercises, challenging you to apply the scope test, physical limiting
budgets, and proposal--permission contracts to real-world edge cases:

\subsection*{Stress-Test 1: The Advisory Observer (Perception Without
Actuation
Authority)}\label{stress-test-1-the-advisory-observer-perception-without-actuation-authority}
\addcontentsline{toc}{subsection}{Stress-Test 1: The Advisory Observer
(Perception Without Actuation Authority)}

\emph{The scenario.} An automotive driver assistance module executes a
50-layer deep convolutional network to detect highway lane markings,
sounding an auditory buzzer whenever the vehicle drifts across a
boundary.

\emph{The systems dilemma.} The system satisfies Criterion 1 (learned
perception model) and Criterion 2 (vehicle motion continuously shifts
camera viewpoints). However, it completely lacks Criterion 3 (delegated
physical authority). Because its only output is an auditory chime, an
algorithmic failure or inference stall cannot directly alter steering
torque or vehicle momentum.

\emph{The systems verdict: Advisory Decision Support.} The system
operates entirely behind glass, relying on the human driver to provide
physical closed-loop actuation.

\subsection*{Stress-Test 2: The Deterministic Autopilot (Actuation
Without Learned
Models)}\label{stress-test-2-the-deterministic-autopilot-actuation-without-learned-models}
\addcontentsline{toc}{subsection}{Stress-Test 2: The Deterministic
Autopilot (Actuation Without Learned Models)}

\emph{The scenario.} An acrobatic quadrotor executes extreme trajectory
tracking through turbulent outdoor wind gusts, calculating motor PWM
commands at \(1000\text{ Hz}\) on a bare-metal microcontroller
(\citeproc{ref-maxwell1868governors}{Maxwell 1868}).

\emph{The systems dilemma.} The system satisfies Criterion 2 (rotor
thrust directly accelerates physical mass) and Criterion 3
(microcontroller latches write directly to motor inverter registers
without human intervention). However, it lacks Criterion 1 (learned
model). The autopilot control law is derived from analytical
Newton-Euler equations and linear state-space feedback whose stability
margins are proven mathematically in closed form.

\emph{The systems verdict: Classical Control Theory.} It requires hard
real-time execution, but does not involve stochastic neural
representations.

\subsection*{Stress-Test 3: The Virtual Agent (Autonomy Without Physical
Feedback)}\label{stress-test-3-the-virtual-agent-autonomy-without-physical-feedback}
\addcontentsline{toc}{subsection}{Stress-Test 3: The Virtual Agent
(Autonomy Without Physical Feedback)}

\emph{The scenario.} An automated high-frequency trading engine executes
thousands of market transactions per second using deep transformer
sequence models to predict price fluctuations.

\emph{The systems dilemma.} The system satisfies Criterion 1 (deep
learned predictive model) and Criterion 3 (automated execution of
transactions without human approval). However, it lacks Criterion 2
(consequential physical feedback). Its outputs alter electronic memory
buffers and order-book records, but do not interact with momentum,
contact friction, or thermodynamic heat.

\emph{The systems verdict: Digital Autonomous Systems.} If an algorithm
diverges, market circuit breakers can freeze accounts, unwind
transactions, and restore clean database states---privileges unavailable
in the physical universe.

\subsection*{Stress-Test 4: The Shared-Autonomy Reflex Gate (The
Neuromuscular Latency
Threshold)}\label{stress-test-4-the-shared-autonomy-reflex-gate-the-neuromuscular-latency-threshold}
\addcontentsline{toc}{subsection}{Stress-Test 4: The Shared-Autonomy
Reflex Gate (The Neuromuscular Latency Threshold)}

\emph{The scenario.} A robotic surgical scalpel applies learned haptic
guidance forces to keep a surgeon's tool tip within a safe anatomical
resection zone.

\emph{The systems dilemma.} The system satisfies Criterion 1 (learned
anatomical spatial models) and Criterion 2 (tool forces directly deform
living tissue). Whether it satisfies Criterion 3 depends strictly on the
actuator force rise time relative to human neuromuscular reflex latency
(\(150\text{--}250\text{ ms}\)). If the haptic actuator can apply forces
within that \(200\text{ ms}\) window, the learned policy exercises
transient, uncontested physical authority over tissue before the human
surgeon can react.

\emph{The systems verdict: Physical AI (Shared Autonomy).} The system
requires a deterministic bare-metal safety filter to clip tool forces
and prevent tissue damage during neural inference anomalies.

\section{Summary}\label{sec-boundary-1-12}

The causal boundary is the physical and electrical interface where
digital information irreversibly converts into thermodynamic action.
Informational abstraction ends at the memory-mapped registers of motor
inverters, valve solenoids, and power stages. Below this boundary,
computational delay is paid in unguided physical displacement, and
released kinetic energy cannot be undone with a software exception
handler.

Surviving this boundary requires an integrated physical container
governed by four bedrock laws:

\begin{enumerate}
\def\labelenumi{\arabic{enumi}.}
\tightlist
\item
  \textbf{The Physical Body} moves mass, exerts contact force, and
  dissipates heat under strict conservation laws.
\item
  \textbf{The Cognitive Brain} executes high-capacity neural models on
  application processors, functioning as an unprivileged proposal
  engine.
\item
  \textbf{The Real-Time Nervous System} executes deterministic safety
  filters on isolated microcontrollers, holding exclusive electrical
  authority over motor registers to verify physical invariants before
  every commit.
\item
  \textbf{The Governance Envelope} enforces decaying operating leases,
  dynamic safety boundaries, and auditable verification cases.
\end{enumerate}

\protect\phantomsection\label{callout-takeawaysux2a-1.10}
\begin{fbxSimple}{callout-takeaways}{Key Chapter Takeaways}{}
\phantomsection\label{callout-takeaways*-1.10}

\begin{itemize}
\tightlist
\item
  \textbf{The Causal Boundary is a Physical Address:} Informational
  abstraction ends at the memory-mapped registers of motor gate drivers.
  Digital bits convert irreversibly into magnetic flux, kinetic
  momentum, and Joule heat (\(I^2 R\)), where software exceptions cannot
  unwind thermodynamic work.
\item
  \textbf{Latency is Paid Directly in Millimeters:} Information decays
  with time (\(t_{\text{age}}\)). Every millisecond of compute and
  transport latency at velocity \(v\) accumulates unguided physical
  displacement (\(\Delta x = \int v \, dt\)), requiring
  sensing-to-actuation budgets to be bounded by strict worst-case tail
  latencies (\(P_{99.9}\)).
\item
  \textbf{The Tri-Archetype Physical Resource Floors:} Physical machines
  divide into Mass-Dominated (Class 1, bound by kinetic stopping
  distance \(v^2/2a_{\max}\)), Contact-Dominated (Class 2, bound by
  contact impedance \(k \Delta x\) and shear stress), and Flow-Dominated
  (Class 3, bound by lumped thermal capacitance and grid phase timing).
\item
  \textbf{The Three-Pillar Scope Test:} An architecture requires the
  full systems discipline of Physical AI if and only if it
  simultaneously combines an unspecifiable learned model, closed-loop
  physical feedback, and delegated actuator authority.
\item
  \textbf{The Four Anatomy Pillars:} Every embodied machine integrates a
  Dynamical Body (Layer 1), a Cognitive Brain (Layer 3), a Real-Time
  Nervous System (Layer 2), and a Governance Envelope (Layer 4).
\item
  \textbf{The Dual-Brain Proposal--Permission Contract:} High-capacity
  neural policies on application processors (Linux MPUs) operate
  strictly as unprivileged proposal engines (\(p_t\)), while isolated
  real-time microcontrollers (bare-metal MCUs) retain exclusive hardware
  authority to evaluate physical barrier invariants before committing
  PWM gate pulses.
\end{itemize}

\end{fbxSimple}

\protect\phantomsection\label{callout-labux2a-1.11}
\begin{fbxSimple}{callout-lab}{Hands-On Lab:}{The Educational Desktop Twin}
\phantomsection\label{callout-lab*-1.11}

To explore these proposal--permission dynamics on an engineer's desk
without requiring industrial factory floors, this course utilizes a
canonical heterogeneous dual-brain System-on-Chip (SoC) as an accessible
desktop reference twin. The architecture pairs an unprivileged Linux
microprocessor (MPU) executing neural vision and planning models with an
isolated, deterministic bare-metal microcontroller (MCU) running
\(1\text{ kHz}\) safety filters over shared SRAM mailboxes. This
desk-bench platform mirrors the exact proposal--permission boundary used
in commercial autonomous vehicles and industrial workcells.

\end{fbxSimple}

The next chapter (Chapter~\ref{sec-body}) transitions from conceptual
boundaries to physical metrology on the hardware floor. By attaching
measuring instruments directly to the compute substrate, the power
electronics, and the mechanical linkages, we derive the governing
equations and measurement disciplines for time, inertia, actuation,
energy, and silicon. That quantitative accounting establishes the firm
numerical budgets that every subsequent chapter must respect.

\part{The Machine Anatomy}

\lettrine{M}{achine} learning was born behind glass. For six decades,
computer science evolved under the protective abstraction of
idempotency: algorithms processed digital symbols in sandboxed memories
where time was elastic, memory was cheap in the cloud, errors printed
stack traces, and computation existed apart from physical consequence.
If an inference pass took an extra hundred milliseconds, a user simply
waited; if a database transaction aborted, the runtime rolled back to a
clean checkpoint; if a model produced an incorrect token, the physical
universe remained untouched.

Physical AI begins at the exact boundary where computation ceases to be
symbolic and becomes physical force. The moment an algorithm writes to a
memory-mapped register on a motor timer, an H-bridge switches, current
surges into copper windings, and mass accelerates through physical space
under momentum and gravity.

At that boundary, all the forgiving conventions of digital software
engineering vanish. Matter will not pause while an operating system
schedules a thread. Inertia does not negotiate with the tail latency of
a neural network. And physical action cannot be undone: kinetic energy
delivered into an environment permanently alters the physical state of
the world, heat accumulates irreversibly in stator windings, and a
mechanical collision cannot be caught by a software exception handler.

To engineer systems that can operate across this divide, Part I
establishes the \textbf{Constraint-Driven Architecture} of physical AI.
Rather than treating an autonomous machine as an abstract software
pipeline, we ground the entire discipline in the fundamental Triad that
comprises every physical machine:

\begin{itemize}
\tightlist
\item
  \textbf{The Body (Chapter 2):} The physical, mechanical, and thermal
  plant. The body obeys the conservation laws of classical mechanics,
  thermodynamics, and electromagnetism rather than software
  instructions. It is governed by five non-negotiable physical budgets:
  sensory measurement freshness, dynamic stopping distance, reflected
  rotor inertia (\(N^2 J\)), stator thermal dissipation (\(I^2 R\)), and
  power delivery voltage droop.
\item
  \textbf{The Brain (Chapter 3):} The learned cognitive deliberative
  engine. We define what the brain is \emph{for}---high-capacity
  open-world semantic generalization, spatial affordance extraction, and
  continuous trajectory synthesis that classical hand-crafted state
  machines cannot achieve. Crucially, we also define what \emph{bounds}
  the brain: the memory wall, Roofline arithmetic intensity limits,
  stochastic fat-tailed latency, and uncalibrated confidence. Because of
  these constraints, the brain must operate as an \textbf{unprivileged
  proposal engine}, emitting advisory action chunks rather than direct
  actuator commands.
\item
  \textbf{The Nervous System (Chapter 4):} The deterministic
  architectural bridge. Sitting between the slow, stochastic brain
  (\(1\text{--}5\text{ Hz}\)) and the fast, unforgiving body
  (\(1000\text{ Hz}\)), the nervous system arbitrates time and
  authority. It translates disparate operational cadences into a
  coherent multi-rate hierarchy, enforces physical invariants at
  microsecond cadences via Control Barrier Functions, and ensures that a
  delay, panic, or hallucination in high-level software never leads to
  unmanaged kinetic damage.
\end{itemize}

Part I provides the mental models, mathematical derivations, and
architectural principles required to understand this Triad. Before we
collect demonstrations, train neural policies, or build runtime
perception stacks, we must first master the physical constraints that
dictate how a machine survives in the material world.

\chapter{\texorpdfstring{The Physical
Body\chaptersubtitle{What the Physical World Will Not Allow}}{The Physical Body}}\label{sec-body}

\begin{figure}[H]

\centering{

\includegraphics[width=1\linewidth,height=\textheight,keepaspectratio]{index_files/mediabag//Users/VJ/GitHub/MLSysBook-integrate-physical/physical/book/chapters/02-body/figures/fig_locator.pdf}

}

\caption{\label{fig-locator-ch02}\textbf{Systems Organ Focus: The Body.}
The four-tier physical AI systems architecture highlights the physical
world at its foundation, bounded by mechanics, thermodynamics, and
electromagnetism. Above the causal boundary, the nervous system and
learned brain must budget all sensing latency, stopping distances, and
actuation requests against these non-negotiable physical constraints.}

\end{figure}%

In digital software systems, execution speed is treated as an elastic
resource: computation pauses while waiting for data, buffers absorb
bursts, and downstream consumers process results whenever the CPU yields
a time slice. In the physical world, time is an unyielding dynamical
state variable governed by the external universe. Every microsecond
spent deliberating in silicon allows moving mass to travel forward
unchecked, degrades sensory freshness as uncertainty envelopes expand,
accumulates irreversible thermal energy in copper windings, and injects
phase lag that systematically erodes closed-loop stability. When learned
models command physical actuators, the survival of the machine is
governed not by average computational throughput or nominal simulation
benchmarks, but by the strict, measured physical constraints of the
body.

\newpage{}

\protect\phantomsection\label{callout-objectiveux2a-2.1}
\begin{fbxSimple}{callout-objective}{Learning Objectives}{}
\phantomsection\label{callout-objective*-2.1}

After studying this chapter, you will be able to:

\begin{enumerate}
\def\labelenumi{\arabic{enumi}.}
\tightlist
\item
  Measure the 8-stage sense-to-actuation latency waterfall and enforce
  the physical freshness deadline
  \(\Delta t_{\text{fresh}} \le e_{\max}/v_{\max}\) using hardware
  timestamps.
\item
  Solve the dynamic clearance quadratic to bound machine velocity
  against linear reaction lag and quadratic braking distance.
\item
  Calculate reflected rotor inertia
  (\(J_{\text{ref}} = N^2 J_{\text{rotor}}\)) and apply \(C^2\) spline
  rate limits to protect transmissions from destructive torque shock.
\item
  Derive sustainable motor duty cycles (\(D\)) from thermodynamic
  Joulean heating and real-time \(I^2 t\) thermal observers.
\item
  Budget electrical rail transients (\(\Delta V_{\text{droop}}\)) and
  kinetic regenerative surges (\(V_{\text{final}}\)) across shared power
  delivery networks.
\item
  Apply the Cross-Archetype Physical Constraint Matrix to diagnose
  multi-variable failure avalanches across mobile, manipulation, and
  thermal systems.
\end{enumerate}

\emph{Core chapter takeaway.} The limits that matter are measured on the
machine in front of you, not read from a datasheet, and it is the tail
of the empirical distribution (\(P_{99.9}\)) that decides whether the
machine survives in the physical world.

\end{fbxSimple}

\section{The Five Physical Budgets}\label{sec-body-2-1}

\lettrine{A}{ software bug}\lettrine{i}{n} a cloud service produces an
exception, drops a request, or rolls back a database transaction to an
uncorrupted checkpoint. In a physical machine, the body obeys mechanics,
thermodynamics, and electromagnetism rather than software instructions.
Once current passes through an inverter into a motor winding, magnetic
flux generates torque, accelerating a mechanical link and imparting
kinetic energy to mass. If that acceleration drives an end-effector
toward a rigid barrier, software cannot issue an abort command that
retroactively absorbs momentum. The energy must dissipate through the
structure, the workpiece, or the mechanical joint. In pure computation,
state transitions are reversible abstractions. In physical AI, every
actuation is an irreversible physical event governed by conservation
laws that software can neither pause nor negotiate.

As established in the systems architecture locator
(Figure~\ref{fig-locator-ch02}), the physical body occupies the
foundational layer beneath the real-time nervous system and the learned
brain. While higher-level planning policies propose trajectories in
unconstrained software coordinates, the physical body enforces the
immutable boundary conditions of matter, momentum, and thermal
dissipation. Every subsequent layer in the physical AI stack---from the
millisecond reflex enforcers to the deliberative neural
planners---functions as a resource allocation against these physical
realities. The body sets terms nobody negotiates, and every later
decision in this book is an allocation against them.

To evaluate these constraints systematically, any learned policy output
draws against five non-negotiable physical budgets that define what the
hardware can survive:

\begin{enumerate}
\def\labelenumi{\arabic{enumi}.}
\tightlist
\item
  \textbf{Measurement Freshness}: How rapidly the physical state of the
  world changes relative to the age of the sensor data reaching the
  controller.
\item
  \textbf{Kinetic Momentum}: The minimum stopping distance and braking
  force required to arrest moving mass before colliding with obstacles
  or workcell boundaries.
\item
  \textbf{Actuator Bandwidth and Reflected Inertia}: How quickly
  delivered mechanical torque and velocity can change before winding
  inductance, inverter current clamps, or transmission inertia choke
  mechanical acceleration.
\item
  \textbf{Thermal Capacity}: How much Joulean heat can accumulate in
  stator coils and power silicon before insulation degrades or thermal
  throttling intervenes.
\item
  \textbf{Electrical Power Integrity}: The transient voltage droop, peak
  current limits, and regenerative energy surges that occur across the
  shared direct-current distribution rail during dynamic maneuvers.
\end{enumerate}

\protect\phantomsection\label{callout-lawux2a-2.2}
\begin{fbxSimple}{callout-law}{Bedrock Law:}{The Physical Body and the Five Conservation Budgets}
\phantomsection\label{callout-law*-2.2}

The \textbf{physical body} is the foundational electromechanical,
thermodynamic, and structural plant---comprising actuators,
transmissions, structural linkages, power distribution buses, and
transducers---that interacts directly with the physical universe. In
physical AI, every actuation is an irreversible physical event that must
be budgeted against five non-negotiable physical constraints:
measurement freshness, kinetic momentum, actuator bandwidth and
reflected inertia, thermal capacity, and electrical power integrity.

\end{fbxSimple}

\protect\phantomsection\label{callout-labux2a-2.3}
\begin{fbxSimple}{callout-lab}{Hands-On Lab:}{The Educational Desktop Twin}
\phantomsection\label{callout-lab*-2.3}

To explore body dynamics without heavy industrial machinery, this course
utilizes a canonical desktop reference twin. An unprivileged Linux
microprocessor (MPU) running high-level vision and trajectory generators
pairs with an isolated, deterministic bare-metal microcontroller (MCU)
that samples encoder phase lines and current shunts at \(1\text{ kHz}\)
over shared SRAM mailboxes. This modular twin allows measuring
electrical rise time, reflected inertia, and bus voltage sag directly on
an engineer's workbench.

\end{fbxSimple}

Engineers routinely fail physical machines because they read component
datasheets as if they were operational promises. A motor vendor tests an
isolated actuator on an open dynamometer bench at an ambient temperature
of \(25^\circ\text{C}\), backed by an infinite heatsink and powered by a
regulated laboratory bench supply through short, heavy-gauge cables. In
that isolated test setup, the motor delivers its published numbers.
Inside a functional machine, that same actuator operates inside a sealed
chassis alongside power converters and compute modules, surrounded by
stagnant air, mechanically coupled to gearboxes with non-zero friction,
and powered through a flexible cable assembly whose parasitic resistance
creates voltage drop under peak load.

Consider a brushless frameless motor rated on its datasheet for a
continuous torque of \(4.2\text{ N}\cdot\text{m}\) and a peak torque of
\(12.0\text{ N}\cdot\text{m}\) at \(25^\circ\text{C}\) with forced
cooling. When six of these units are enclosed within an unventilated
robotic limb alongside an onboard inference computer dissipating
\(60\text{ W}\), the internal cavity temperature rises to
\(55^\circ\text{C}\). Because the thermal resistance from the stator to
the ambient air is higher than the vendor test fixture by a factor of
three, the continuous torque the motor can sustain without exceeding its
insulation limit drops from \(4.2\text{ N}\cdot\text{m}\) down to
\(1.8\text{ N}\cdot\text{m}\). If an unconstrained high-level policy
commands the motor at its datasheet continuous rating of
\(4.2\text{ N}\cdot\text{m}\), copper losses generate heat at nearly
four times the dissipation rate of the chassis, forcing the joint into
thermal shutdown in less than ninety seconds. The true operating
envelope of a component is never the boundary drawn on the datasheet; it
is the restricted envelope measured on the fully assembled machine under
worst-case environmental conditions.

Systems engineering for physical AI requires establishing hard empirical
ceilings that upper software layers must budget against rather than
treating hardware limits as optimization penalties. If a safety filter
or an execution planner treats maximum torque or thermal dissipation as
a soft cost term in an objective function, an out-of-distribution neural
network output will eventually exploit that elasticity, demanding
current the battery cannot supply or commanding deceleration rates that
induce bus overvoltage. When software exceeds a physical limit, the
hardware enforces its own response without consulting the operating
system. Overcurrent trips protection circuits, unmanaged regenerative
current blows inverter MOSFETs, and thermal accumulation triggers hard
shutoffs that drop loads under gravity. Safe operation requires knowing
the exact empirical boundary of the physical machine at the \(P_{99.9}\)
tail of operating conditions, and reserving an explicit safety margin
below it.

Every physical allocation begins with an accurate accounting of state.
Before an actuator can safely exert torque to decelerate a mass or
balance a link, the machine must establish where that mass is and how
fast it is moving. Yet the physical state of the world continues to
evolve continuously while electrons traverse cables, digital converters
sample voltages, and bus controllers arbitrate packets. A measurement
does not represent the present state of the body; it records a
historical state whose physical validity begins decaying at the precise
microsecond of transduction.

\section{Measurement Age and Freshness}\label{sec-body-2-2}

On a fast-moving logistics sorting line, a vision-guided robotic
manipulator tracks parcels rushing past on a conveyor belt at
\(1.2\text{ m/s}\). To pick a moving target without missing or colliding
with the guide rails, the robot's control policy must know precisely
where the object is in physical space. Yet in the time it takes for
photons to expose the image sensor, for pixel data to stream across the
PCIe bus, for a neural network to infer a grasp pose, and for motor
windings to generate electromagnetic torque, the physical world has
moved forward.

In traditional software benchmarking, latency is often reported as the
pure execution time of a neural network running on an accelerator. In an
embodied physical machine, that compute slice is only an intermediate
portion of a much longer causal chain. \textbf{A measurement begins
decaying at transduction, and what matters is its age at the moment
something acts on it.} The physical age of an observation spans the full
interval from when photons strike a photodiode array or magnetic flux
rotates past a Hall sensor, to when motor coils generate mechanical
torque that alters the acceleration of a link. If a perception model
computes an obstacle location in \(15\text{ ms}\) but camera
integration, sensor readout, transport, queue dispatch, and motor
current rise time add another \(45\text{ ms}\), the physical age of the
measurement when the joint moves is \(60\text{ ms}\). During that entire
window, the physical world has continued to evolve.

Tracking measurement age requires decomposing the complete physical path
into eight distinct stages (Figure~\ref{fig-02-latency-waterfall}).
First, physical integration and exposure capture the environment over an
exposure window, setting the timestamp to the midpoint of integration.
Second, sensor readout serializes analog charge into digital frames over
a physical interface such as MIPI CSI-2 or SPI. Third, hardware
transport moves the packet into host memory through direct memory access
controllers and PCIe/AXI buses. Fourth, kernel buffering and driver
ingestion place the frame in user-space queues. Fifth, model computation
executes neural network inference, filtering, or kinematic projection on
an accelerator. Sixth, operating system scheduling and inter-process
communication dispatch the resulting control setpoint across the
processor boundary. Seventh, command transmission serializes the
setpoint across a real-time fieldbus or shared-memory mailbox to the
motor driver. Eighth, actuator response converts the command into
pulse-width modulation signals, building magnetic field flux in the
stator windings to develop mechanical force. On a canonical
heterogeneous dual-brain platform, this eight-stage pipeline spans the
CMOS image sensor \(\to\) MIPI DMA into host memory \(\to\) neural
policy inference running on the Linux application processor (MPU)
\(\to\) inter-processor communication (IPC) over shared SRAM mailboxes
via lock-free ring buffers \(\to\) deterministic \(1\text{ kHz}\) safety
monitor evaluation on the real-time microcontroller (MCU) \(\to\) PWM
timer register updates \(\to\) gate driver switching and stator coil
current rise (\(\tau_e = L/R\)). Every stage consumes a slice of the
total time budget.

\begin{figure}[t]

\centering{

\includegraphics[width=1\linewidth,height=\textheight,keepaspectratio]{index_files/mediabag//Users/VJ/GitHub/MLSysBook-integrate-physical/physical/book/chapters/02-body/figures/fig02_latency_waterfall.pdf}

}

\caption{\label{fig-02-latency-waterfall}\textbf{End-to-End
Sense-to-Actuation Latency Waterfall and Measurement Freshness Budget.}
(Top) The eight-stage sensory-motor causal chain spanning photonic CMOS
exposure, MIPI/PCIe DMA ingestion, vision patch tokenization, action
chunk policy inference, shared SRAM IPC mailbox transfer, 1 kHz
microcontroller barrier filter evaluation, EtherCAT fieldbus
serialization, and stator coil current rise (\(\tau_e = L/R\)). (Bottom
Left) Calibrated cumulative delay breakdown comparing nominal execution
(\(P_{50} = 35.0\text{ ms}\)) against a heavy-tail latency spike
(\(P_{99.9} = 68.4\text{ ms}\)) relative to the statutory freshness
deadline (\(\Delta t_{\text{fresh}} = 40.0\text{ ms}\)). (Bottom Right)
Corresponding physical consequence: unguided blind coasting travel
(\(d_{\text{react}} = v_0 \Delta t\)) expands by \(4.0\text{ cm}\) at
\(1.2\text{ m/s}\), eroding obstacle clearance margins into unsafe
collision risk.}

\end{figure}%

The rate of world-state change dictates how long a measurement remains
valid before acting on it introduces intolerable error. Let \(e_{\max}\)
denote the maximum acceptable error in a tracked state, such as joint
position or distance to an obstacle, and let \(v_{\max}\) denote the
fastest credible velocity at which that state can evolve. The freshness
deadline \(\Delta t_{\text{fresh}}\) is the upper bound on measurement
age before a control action based on that measurement becomes invalid,
defined by the inequality
\[\Delta t_{\text{fresh}} \le \frac{e_{\max}}{v_{\max}}\]

For our conveyor-tracking manipulator operating with a maximum relative
velocity \(v_{\max} = 1.2\text{ m/s}\) and a geometric clearance
tolerance \(e_{\max} = 24\text{ mm}\), the freshness deadline is
\(24\text{ mm} / (1200\text{ mm/s}) = 20\text{ ms}\). If an actuation
command reaches the motor based on an image captured \(35\text{ ms}\)
earlier, the workpiece has displaced by \(42\text{ mm}\), exceeding the
position allowance by \(18\text{ mm}\) and missing the grasp. Deriving
the freshness deadline from physical dynamics establishes a strict
threshold for the nervous system, defining the exact age beyond which
sensor data turns from an asset into a hazard.

Engineers often attempt to compute total latency by summing the
\(P_{99}\) latency budgets of individual stages, but the percentile of a
sum is not the sum of the percentiles.\sidenote{\footnotesize Quantiles are
  non-subadditive (\(Q_p(X+Y) \ne Q_p(X) + Q_p(Y)\)). Bus contention
  creates positive tail correlations where empirical \(P_{99.9}\)
  latency exceeds the sum of medians.} When stages are statistically
independent, the probability that all eight stages simultaneously
experience their \(P_{99}\) tail in the same cycle is roughly
\(10^{-16}\), meaning that a simple sum of \(P_{99}\) values produces an
overly pessimistic bound that wastes hardware capacity. Conversely, when
stages share underlying hardware resources, such as direct memory access
contention during model weight paging or kernel lock contention during
high-throughput packet reception, delays correlate across stages. Under
correlated load, the empirical \(P_{99.9}\) end-to-end latency can
exceed the sum of individual medians by a wide margin. Because
statistical independence cannot be assumed across shared buses and
operating system queues, end-to-end age must be measured directly on the
integrated machine from physical stimulus to current rise.

Consider the empirical characterization of this vision-guided
manipulator performing visual servoing over an EtherCAT bus under
continuous inference load, evaluated across \(N = 100{,}000\) recorded
cycles. At \(60\text{ Hz}\) camera operation with rolling shutter
readout, the sensor exposure and readout stage takes a median
\(P_{50} = 11.2\text{ ms}\), with a \(P_{99} = 16.6\text{ ms}\) and a
maximum observed value of \(16.7\text{ ms}\). Transport over PCIe to
host memory consumes \(P_{50} = 1.4\text{ ms}\),
\(P_{99} = 3.2\text{ ms}\), and a maximum of \(6.8\text{ ms}\). The
vision backbone and policy evaluation on the onboard accelerator require
\(P_{50} = 18.5\text{ ms}\), \(P_{99} = 24.1\text{ ms}\), and a maximum
of \(38.4\text{ ms}\) during memory throttling. Operating system
scheduling and inter-process communication account for
\(P_{50} = 0.6\text{ ms}\), \(P_{99} = 2.1\text{ ms}\), and a maximum of
\(12.3\text{ ms}\) during kernel preemption. Fieldbus command
transmission over EtherCAT operates at \(P_{50} = 1.0\text{ ms}\),
\(P_{99} = 1.0\text{ ms}\), and a maximum of \(2.0\text{ ms}\) upon
packet retransmission. Actuator current rise and torque generation
contribute \(P_{50} = 2.3\text{ ms}\), \(P_{99} = 3.0\text{ ms}\), and a
maximum of \(4.5\text{ ms}\).

Measuring the complete trajectory directly from photodiode illumination
to stator current produces an empirical end-to-end age distribution of
\(P_{50} = 35.0\text{ ms}\), \(P_{99} = 46.2\text{ ms}\),
\(P_{99.9} = 68.4\text{ ms}\), and a maximum observed age of
\(80.7\text{ ms}\). Summing the individual \(P_{99}\) stages gives
\(50.0\text{ ms}\), which overestimates the true \(P_{99}\) latency of
\(46.2\text{ ms}\) due to partial independence among intermediate
queues. If the physical task imposes a freshness deadline of
\(\Delta t_{\text{fresh}} = 40.0\text{ ms}\), the system operates within
budget at median latency (\(35.0\text{ ms}\)), but fails the deadline on
more than \(10\%\) of cycles where end-to-end age exceeds
\(40.0\text{ ms}\).

For the physical AI systems engineer, these dynamics dictate how the
nervous system must be architected. Executing an actuation command
computed from expired sensor data injects severe phase lag into the
feedback loop, degrading phase margin and exciting resonant mechanical
chatter.\sidenote{\footnotesize Delay \(\tau\) injects phase lag \(-\omega \tau\). By
  Bode's integral, delayed feedback inflates peak sensitivity
  \(|S(j\omega)|_\infty\), exciting mechanical chatter.} Rather than
assuming that high-frequency model inference guarantees real-time
responsiveness, the engineer must tag every incoming observation with a
hardware timestamp at the exact instant of physical transduction (using
IEEE 1588 gPTP or hardware strobe capture at the Ethernet PHY / MIPI
CSI-2 interface, rather than uncalibrated operating system clocks like
\texttt{time.now()}) and carry that timestamp through the entire
inference pipeline. Immediately before committing an action setpoint to
the motor registers, the execution runtime must compute the elapsed age
\(t_{\text{now}} - t_{\text{transduction}}\) and compare it against the
physical freshness deadline \(\Delta t_{\text{fresh}}\). When latency
jitter pushes measurement age past this threshold, the nervous system
must refuse the stale policy command and transition to a deterministic
holding reflex, protecting the mechanical plant from the destructive
phase lag of outdated beliefs.

A bounded measurement age guarantees that the machine acts on an
accurate representation of its current state, but it does not eliminate
the physical delay between sensing an obstacle and changing trajectory.
While signals propagate through the nervous system, the body continues
moving along its existing vector, carrying momentum that no software
interrupt can cancel.

\section{Momentum and Stopping Distance}\label{sec-body-2-3}

A warehouse transport robot carrying \(120\text{ kg}\) of palletized
freight navigates a distribution floor at \(2.0\text{ m/s}\). When a
worker steps unexpectedly into the aisle \(1.5\text{ m}\) ahead,
software cannot simply pause execution or instantaneously zero its
velocity vector. The vehicle's steel chassis and cargo carry unyielding
mechanical momentum, and bringing that mass to rest requires dissipating
kinetic energy across physical distance. \textbf{Stopping distance grows
linearly with delay and quadratically with speed, and that difference
sets the speed limit.}

The total distance traveled between sensing a hazard and reaching zero
velocity divides into two physical phases: reaction travel and braking
travel. During the reaction interval, the nervous system transduces the
scene, transfers data over internal buses, executes model inference, and
dispatches the brake setpoint. Throughout this full end-to-end
sense-to-actuation delay \(\Delta t_{\text{delay}}\), the actuators
apply no retarding force, and the vehicle continues moving along its
existing trajectory at its initial velocity \(v_0\). Assuming velocity
remains constant across this window, the reaction travel evaluates to
\[d_{\text{react}} = v_0 \Delta t_{\text{delay}}\] Because reaction
travel is proportional to delay, any latency tail in compute or
communication extends reaction distance in direct proportion to
velocity.

Once the brake command arrives and the actuators engage, the body enters
the braking phase. The mechanical work done by the braking force
\(F_{\text{brake}}\) over the braking distance \(d_{\text{brake}}\) must
dissipate the initial kinetic energy of the machine, \(W = \Delta K\).
For a machine of mass \(m\) decelerating to rest under a sustained
retarding force \(F_{\text{brake}} = m a_{\text{brake}}\), the
work-energy balance yields
\[m a_{\text{brake}} d_{\text{brake}} = \frac{1}{2} m v_0^2\] Dividing
by mass and solving for distance gives the kinematic braking travel
\[d_{\text{brake}} = \frac{v_0^2}{2 a_{\text{brake}}}\] The parameter
\(a_{\text{brake}}\) is the credible deceleration the body can exert
against the environment, not a nominal rating printed on an actuator
datasheet. It represents the sustained negative acceleration the
mechanical structure, contact interface, and transmission can produce
without inducing wheel slip, tipping, or structural yield.

This closed-form expression for \(d_{\text{brake}}\) relies on four
simplifying assumptions: constant deceleration throughout the maneuver,
instantaneous torque buildup at the contact interface, zero grade, and
rigid contact mechanics. In physical hardware, these assumptions
frequently fail. Actuator magnetic flux requires tens of milliseconds to
reach peak force, tire-ground friction drops nonlinearly if slip exceeds
the peak of the friction curve, and structural compliance in linkages
deflects under high inertial load. When operating on uncertain terrain,
grades, or compliant surfaces, kinematic equations underestimate true
stopping distance. Safety budgeting requires an empirical stopping trial
on the physical machine across its operating envelope, measuring actual
stopping trajectories with optical tracking or high-frequency wheel
encoders to establish the empirical lower bound of \(a_{\text{brake}}\)
at \(P_{99}\) rather than relying on ungrounded dry-friction models.

The differing exponents of the two terms create an asymmetric scaling
penalty when adjusting system parameters
(Figure~\ref{fig-02-stopping-distance}). In an analytical evaluation of
our \(120\text{ kg}\) transport robot operating with a measured
\(P_{99}\) sense-to-brake delay of
\(\Delta t_{\text{delay}} = 50.0\text{ ms}\) and a measured deceleration
limit of \(a_{\text{brake}} = 2.0\text{ m/s}^2\), moving at
\(v_0 = 1.0\text{ m/s}\) produces a reaction travel of
\(d_{\text{react}} = (1.0\text{ m/s})(0.050\text{ s}) = 0.050\text{ m}\)
and a braking travel of
\(d_{\text{brake}} = (1.0\text{ m/s})^2 / (2 \cdot 2.0\text{ m/s}^2) = 0.250\text{ m}\),
yielding a total stopping distance of \(0.300\text{ m}\). If software
scheduling jitter doubles the delay to
\(\Delta t_{\text{delay}} = 100.0\text{ ms}\) while keeping velocity at
\(1.0\text{ m/s}\), reaction travel doubles to \(0.100\text{ m}\),
increasing total stopping distance by \(17\%\) to \(0.350\text{ m}\). If
instead the system doubles velocity to \(v_0 = 2.0\text{ m/s}\) while
holding delay at \(50.0\text{ ms}\), reaction travel doubles to
\(0.100\text{ m}\) while braking travel quadruples to
\(1.000\text{ m}\), pushing total stopping distance to
\(1.100\text{ m}\), an increase of \(267\%\). System latency sets the
baseline reaction overhead, but velocity governs the growth of the
stopping envelope.

In an unstructured environment, a machine cannot allocate its entire
clearance to nominal stopping dynamics. Obstacle tracking carries state
estimation uncertainty, sensor calibration drifts over time, and
localization algorithms introduce position error. The system must defend
a stopping envelope \(d_{\text{stop}}\) that includes a spatial
localization uncertainty bound \(\delta_{\text{loc}}\) and a minimum
physical clearance margin \(\delta_{\text{margin}}\)\sidenote{\footnotesize ISO/TS
  15066 (Clause 5.5.4) and ISO 13849-1 mandate this protective
  separation distance for Speed and Separation Monitoring (SSM):
  \(S = (v_r \cdot T_r) + (v_h \cdot T_r) + S_r + S_s + C\), where
  \(T_r\) is worst-case processing delay, \(S_r\) is braking travel, and
  \(C\) is clearance margin (\(100\text{--}200\text{ mm}\)).}
\[d_{\text{stop}} = v_0 \Delta t_{\text{delay}} + \frac{v_0^2}{2 a_{\text{brake}}} + \delta_{\text{loc}} + \delta_{\text{margin}}\]
The sum
\(\delta_{\text{overhead}} = \delta_{\text{loc}} + \delta_{\text{margin}}\)
represents fixed spatial overhead that must remain unoccupied when the
body comes to a complete rest.

\begin{figure}[t]

\centering{

\includegraphics[width=1\linewidth,height=\textheight,keepaspectratio]{index_files/mediabag//Users/VJ/GitHub/MLSysBook-integrate-physical/physical/book/chapters/02-body/figures/fig02_stopping_distance.pdf}

}

\caption{\label{fig-02-stopping-distance}\textbf{Defended Dynamic
Stopping Envelope: Linear Reaction Lag vs.~Quadratic Kinetic Braking
Distance.} (Top) Analytical decomposition of stopping distance
\(d_{\text{stop}}(v_0, \Delta t_{\text{delay}}) = v_0 \Delta t_{\text{delay}} + v_0^2 / (2 a_{\text{brake}}) + \delta_{\text{overhead}}\)
defending against clearance \(D_{\text{clear}}\). (Left) Mathematical
parameter sensitivity curves illustrating why doubling delay
\(\Delta t_{\text{delay}}\) produces a linear \(17\%\) increase in
stopping distance, whereas doubling velocity \(v_0\) quadruples braking
travel, causing a \(267\%\) expansion in the stopping envelope. (Right)
Physical trajectory comparisons on an instrumented \(60.0\text{ cm}\)
obstacle track: Case 1 operates with certified safe margin
(\(+20.0\text{ cm}\)); Case 2 illustrates clearance margin compression
under a \(200\text{ ms}\) tail latency spike; and Case 3 demonstrates
catastrophic obstacle barrier breach (\(+26.3\text{ cm}\) collision)
when higher speed couples with tail compute lag.}

\end{figure}%

Given a detected obstacle at distance \(D_{\text{clear}}\), physical
safety requires that the total defended stopping envelope not exceed the
available clear path, \(d_{\text{stop}} \le D_{\text{clear}}\). Rather
than treating maximum velocity as a static configuration parameter, the
nervous system must continuously solve this clearance inequality for the
maximum permitted instantaneous speed \(v_{\text{max}}\). Grouping the
constant spatial terms into available physical distance
\(D_{\text{avail}} = D_{\text{clear}} - \delta_{\text{loc}} - \delta_{\text{margin}}\)
yields the quadratic inequality
\[\frac{1}{2 a_{\text{brake}}} v_0^2 + \Delta t_{\text{delay}} v_0 - D_{\text{avail}} \le 0\]
Multiplying by \(2 a_{\text{brake}}\) produces
\(v_0^2 + 2 a_{\text{brake}} \Delta t_{\text{delay}} v_0 - 2 a_{\text{brake}} D_{\text{avail}} \le 0\).
Applying the quadratic formula and taking the positive physical root
defines the velocity limit
\[v_{\text{max}} = -a_{\text{brake}} \Delta t_{\text{delay}} + \sqrt{(a_{\text{brake}} \Delta t_{\text{delay}})^2 + 2 a_{\text{brake}} (D_{\text{clear}} - \delta_{\text{loc}} - \delta_{\text{margin}})}\]
When sensor range degrades or an obstacle enters the trajectory,
\(D_{\text{clear}}\) contracts, forcing \(v_{\text{max}}\) down along a
square-root trajectory. If
\(D_{\text{clear}} \le \delta_{\text{loc}} + \delta_{\text{margin}}\),
available distance reaches zero, the square-root term collapses to
\(a_{\text{brake}} \Delta t_{\text{delay}}\), and \(v_{\text{max}}\)
evaluates to zero, commanding an immediate stop reflex.

\protect\phantomsection\label{callout-mathux2a-2.4}
\begin{fbxSimple}{callout-math}{Fermi Math:}{Stopping Distance Budgeting under High-Tail Compute Latency}
\phantomsection\label{callout-math*-2.4}

\emph{Problem:} A \(150\text{ kg}\) autonomous mobile robot (AMR)
operates at forward speed \(v_0 = 1.8\text{ m/s}\) in a logistics
facility. The real-time perception pipeline runs at \(20\text{ Hz}\)
(\(50\text{ ms}\) nominal compute), but experiencing an inference tail
spike, direct memory access crossbar contention, and fieldbus
serialization pushes total sense-to-actuation lag to
\(\Delta t_{\text{delay}} = 65\text{ ms}\) (\(P_{99.9}\)). Maximum
emergency friction deceleration is measured at
\(a_{\text{brake}} = 3.0\text{ m/s}^2\), localization uncertainty is
bounded by \(\delta_{\text{loc}} = 5.0\text{ cm}\), and the statutory
obstacle clearance margin is
\(\delta_{\text{margin}} = 10.0\text{ cm}\). What is the minimum
physical defended stopping envelope \(d_{\text{stop}}\), and how does it
scale if the operational speed increases by \(33\%\) to
\(2.4\text{ m/s}\)?

\textbf{1. Reaction Travel (\(d_{\text{react}}\)):} During the
\(65\text{ ms}\) sensing and compute delay, the chassis coasts unguided:
\[d_{\text{react}} = v_0 \Delta t_{\text{delay}} = (1.8\text{ m/s}) \times (0.065\text{ s}) = 0.117\text{ m} = 11.7\text{ cm}\]

\textbf{2. Braking Travel (\(d_{\text{brake}}\)):} Dissipating kinetic
energy (\(E_k = \frac{1}{2} m v_0^2 = 243\text{ J}\)) under sustained
mechanical deceleration:
\[d_{\text{brake}} = \frac{v_0^2}{2 a_{\text{brake}}} = \frac{(1.8\text{ m/s})^2}{2 \times (3.0\text{ m/s}^2)} = \frac{3.24}{6.0} = 0.540\text{ m} = 54.0\text{ cm}\]

\textbf{3. Total Defended Stopping Envelope (\(d_{\text{stop}}\)):}
\[d_{\text{stop}} = d_{\text{react}} + d_{\text{brake}} + \delta_{\text{loc}} + \delta_{\text{margin}} = 11.7\text{ cm} + 54.0\text{ cm} + 5.0\text{ cm} + 10.0\text{ cm} = 80.7\text{ cm}\]

\textbf{4. Scaling Impact of \(+33\%\) Speed Increase
(\(v_0 = 2.4\text{ m/s}\)):}
\[d_{\text{react}} = (2.4\text{ m/s}) \times (0.065\text{ s}) = 0.156\text{ m} = 15.6\text{ cm} \quad (+33.3\%)\]
\[d_{\text{brake}} = \frac{(2.4\text{ m/s})^2}{2 \times (3.0\text{ m/s}^2)} = \frac{5.76}{6.0} = 0.960\text{ m} = 96.0\text{ cm} \quad (+77.8\%)\]
\[d_{\text{stop}} = 15.6\text{ cm} + 96.0\text{ cm} + 15.0\text{ cm} = 126.6\text{ cm} \quad (+56.9\%)\]
A \(33\%\) increase in travel speed expands total stopping distance by
nearly \(57\%\), eroding \(45.9\text{ cm}\) of clearance buffer.

\textbf{Architect's Rule of Thumb:} \emph{Doubling travel speed
quadruples kinetic braking distance while doubling reaction travel;
never increase velocity setpoints without dynamically solving the
clearance quadratic.}

\end{fbxSimple}

\protect\phantomsection\label{callout-algorithmux2a-2.5}
\begin{fbxSimple}{callout-algorithm}{Algorithm:}{Dynamic Velocity Ceiling Governor (1 kHz Microcontroller Reflex)}
\phantomsection\label{callout-algorithm*-2.5}

\textbf{Inputs:}

\begin{itemize}
\tightlist
\item
  \(D_{\text{clear}} \in \mathbb{R}_{\ge 0}\): Detected obstacle
  distance along travel vector \([\text{m}]\).
\item
  \(\Delta t_{\text{delay}} \in \mathbb{R}_{> 0}\): Measured
  \(P_{99.9}\) sense-to-actuation latency \([\text{s}]\).
\item
  \(a_{\text{brake}} \in \mathbb{R}_{> 0}\): Measured minimum credible
  brake deceleration capacity \([\text{m/s}^2]\).
\item
  \(\delta_{\text{loc}}, \delta_{\text{margin}} \in \mathbb{R}_{\ge 0}\):
  Localization uncertainty bound and physical clearance margin
  \([\text{m}]\).
\item
  \(v_{\text{cand}} \in \mathbb{R}_{\ge 0}\): Candidate forward velocity
  setpoint proposed by neural policy \([\text{m/s}]\).
\end{itemize}

\textbf{Output:}

\begin{itemize}
\tightlist
\item
  \(v_{\text{safe}} \in \mathbb{R}_{\ge 0}\): Clamped safe velocity
  setpoint committed to motor registers \([\text{m/s}]\).
\end{itemize}

\textbf{Structural Invariants:}

\begin{enumerate}
\def\labelenumi{\arabic{enumi}.}
\tightlist
\item
  Clearance Invariant:
  \(d_{\text{stop}}(v_{\text{safe}}) = v_{\text{safe}} \Delta t_{\text{delay}} + \frac{v_{\text{safe}}^2}{2 a_{\text{brake}}} + \delta_{\text{loc}} + \delta_{\text{margin}} \le D_{\text{clear}}\)
  for all time \(t\).
\item
  Non-Negativity: \(v_{\text{safe}} \ge 0.0\).
\end{enumerate}

\textbf{Algorithm Steps:}

\begin{enumerate}
\def\labelenumi{\arabic{enumi}.}
\item
  Compute available physical braking distance:
  \[D_{\text{avail}} \leftarrow D_{\text{clear}} - \delta_{\text{loc}} - \delta_{\text{margin}}\]
\item
  If \(D_{\text{avail}} \le 0.0\), return
  \(v_{\text{safe}} \leftarrow 0.0\) (trigger immediate emergency
  holding reflex).
\item
  Compute quadratic discriminant:
  \[\Delta \leftarrow (a_{\text{brake}} \Delta t_{\text{delay}})^2 + 2 a_{\text{brake}} D_{\text{avail}}\]
\item
  Solve positive physical root for dynamic velocity ceiling:
  \[v_{\max} \leftarrow -a_{\text{brake}} \Delta t_{\text{delay}} + \sqrt{\Delta}\]
\item
  Enforce dynamic clamping against policy proposal:
  \[v_{\text{safe}} \leftarrow \min(v_{\text{cand}}, \max(0.0, v_{\max}))\]
\item
  Return \(v_{\text{safe}}\).
\end{enumerate}

\end{fbxSimple}

\protect\phantomsection\label{callout-archetypeux2a-2.6}
\begin{fbxSimple}{callout-archetype}{Archetype Focus:}{Class 1 Mass-Dominated Mobility Systems}
\phantomsection\label{callout-archetype*-2.6}

In Class 1 machines (autonomous mobile robots, warehouse transport
carts, and dynamic quadrupeds), the primary physical failure mode is
\textbf{spatial boundary breach under unmanaged momentum}. Because
kinetic energy \(E_k = \frac{1}{2}m v^2\) cannot be recalled once
delivered to the chassis, any tail latency spike in the neural policy or
communication bus extends reaction travel linearly
(\(d_{\text{react}} = v_0 \Delta t_{\text{delay}}\)), eroding the
physical clearance buffer before mechanical brakes can even engage. The
deterministic real-time MCU must dynamically clamp maximum forward
velocity to
\(v_{\text{max}}(D_{\text{clear}}, \Delta t_{\text{delay}})\) to
guarantee that the machine can always come to a complete halt within the
detected obstacle clearance.

\end{fbxSimple}

For the physical AI systems engineer, this quadratic formulation
transforms speed limits from static configuration constants into dynamic
runtime safety invariants. Rather than assuming a fixed operational
velocity in high-level software, the engineer must program the
deterministic nervous system to continuously solve the clearance
quadratic, dynamically clamping candidate policy proposals to
\(v_{\text{max}}(D_{\text{clear}}, \Delta t_{\text{delay}})\) on every
control cycle. When sensory range degrades, when compute latency spikes,
or when environmental clearance shrinks, the nervous system
automatically contracts the robot's velocity ceiling, guaranteeing that
moving mass never enters a dynamical regime where its physical stopping
distance outstrips available space.

Enforcing this speed limit assumes the nervous system can demand the
deceleration \(a_{\text{brake}}\) whenever safety requires it. Yet
commanding rapid deceleration requires sudden, large mechanical torque
from the drive motors and gearboxes. If the nervous system requests a
torque step beyond what the stator windings can magnetically sustain or
the gear teeth can structurally support, the mechanical body cannot
deliver the required deceleration, turning an analytical stopping
guarantee into a mechanical failure.

\section{Actuator and Transmission Limits}\label{sec-body-2-4}

During a high-speed precision assembly cycle, a 6-axis articulated
industrial robot arm attempts to seat a machined steel pin into a
tight-tolerance bore. In digital software, if an instruction or API call
receives an invalid parameter, it throws an exception, drops a packet,
or returns an error status. An electromechanical actuator possesses no
such software abstraction. When a learned neural policy commands a step
change in joint torque beyond what the magnetic field can generate, or
requests a discontinuous acceleration that strains gear teeth, the
physical actuator does not report a software error. It pulls maximum
phase current from the power stage, twists transmission shafts, and
strains gear teeth while digital command registers in memory continue to
report nominal execution. \textbf{An actuator will refuse a command in
ways that damage the mechanism rather than reporting an error.}

An actuator contract cannot be specified in terms of digital setpoints.
It must be written at the physical interface in terms of delivered
force, torque, or mass flow, response latency, slew rate limits, and
output error. A digital setpoint is merely a software request; physical
action is bounded by winding inductance, bus voltage ceilings, magnetic
saturation, and transmission elasticity.

The delay between commanding a torque and delivering it at the motor
shaft begins in the stator windings. Consider a motor phase winding
modeled as a series resistance \(R\) and inductance \(L\) under an
applied voltage step \(V_{\text{step}}\). Assuming the rotor is
initially at rest so that back-electromotive force is zero, the phase
current \(i(t)\) obeys the first-order differential relation
\[L \frac{di(t)}{dt} + R i(t) = V_{\text{step}}\] Integrating this
differential equation from zero initial current yields
\[i(t) = \frac{V_{\text{step}}}{R} \left( 1 - e^{-t / \tau_e} \right)\]
where \(\tau_e = L / R\) is the electrical time constant.\sidenote{\footnotesize In
  digital motor drives, FOC current loops run at
  \(10\text{--}25\text{ kHz}\), where current rise
  \(di/dt = (V_{\text{bus}} - e_{\text{bemf}})/L\) is bounded by phase
  inductance \(L\) and back-EMF.} Electromagnetic torque is proportional
to current through the torque constant \(K_t\), such that
\(\tau_{\text{m}}(t) = K_t i(t)\). Reaching a fraction
\(\alpha \in (0, 1)\) of the steady-state current
\(I_{\text{ss}} = V_{\text{step}} / R\) requires a budgetable electrical
rise time
\[t_{\text{rise}}(\alpha) = -\tau_e \ln(1 - \alpha) = -\frac{L}{R} \ln(1 - \alpha)\]
To reach \(95\%\) of the commanded torque (\(\alpha = 0.95\)), the rise
time evaluates to
\(t_{\text{rise}}(0.95) = -\tau_e \ln(0.05) \approx 3.0 \tau_e\). In an
analytical evaluation of a joint actuator in our 6-axis arm with winding
resistance \(R = 0.50\,\Omega\) and inductance \(L = 1.5\text{ mH}\),
the electrical time constant is
\(\tau_e = 1.5\times 10^{-3}\text{ H} / 0.50\,\Omega = 3.0\text{ ms}\).
Delivering \(95\%\) of target torque requires
\(3.0 \times 3.0\text{ ms} = 9.0\text{ ms}\). If an upstream policy
assumes torque appears instantaneously, this \(9.0\text{ ms}\) interval
passes as unguided drift and phase erosion before the arm can initiate
corrective acceleration or braking.

Even after the electrical rise time elapses, delivered torque is bounded
by hard current limits in the drive electronics. Power inverters enforce
a maximum current ceiling \(I_{\text{max}}\) to protect switching
transistors from thermal destruction and to prevent ferromagnetic
saturation in the stator core. When the motor drives a mechanical
transmission with gear reduction ratio \(N\) and transmission efficiency
\(\eta\), the ceiling on delivered output torque is:
\[\tau_{\text{out, max}} = \eta N K_t I_{\text{max}}\]

When a high-level policy demands \(120\text{ N}\cdot\text{m}\) of
braking torque to avoid an obstacle, it assumes the motor drive will
deliver it. But on real hardware, inverter current limits
(\(I_{\max} = 25\text{ A}\)) and gear efficiency cap peak output torque
at \(85\text{ N}\cdot\text{m}\). The missing
\(35\text{ N}\cdot\text{m}\) is silently clipped at the power stage.
Because the software planner remains unaware of the current clamp, the
machine's actual stopping distance expands beyond its calculated safety
boundary.

Transmissions introduce a severe mechanical asymmetry that software
planners routinely ignore. To deliver high torque from a compact motor,
robot joints rely on gear reductions with ratios \(N\) ranging from
\(50:1\) to \(160:1\). While gear reduction scales delivered torque
linearly by \(N\), it forces the motor rotor to spin \(N\) times faster
than the joint.

Because kinetic energy scales with the square of speed
(\(\frac{1}{2} J \dot{\theta}^2\)), the effective resistance of the
motor's rotor to joint acceleration scales with the \textbf{square of
the gear ratio}:

\[J_{\text{ref}} = N^2 J_{\text{rotor}}\]

This \textbf{reflected rotor inertia} means that in a \(100:1\)
transmission, a tiny \(50\text{ gram}\) motor rotor acts at the joint as
if it has the inertia of a \(500\text{ kg}\) flywheel
(\citeproc{ref-lynch2017modern}{Lynch and Park 2017}).

This quadratic scaling creates a hard physical boundary for machine
learning policies. Neural policies trained in physics simulators (such
as Isaac Gym or MuJoCo) frequently output piecewise-constant action
setpoints at discrete control intervals (\(20\text{ Hz}\) to
\(50\text{ Hz}\)). Across timestep boundaries, commanded velocity jumps
discontinuously, producing near-infinite requested acceleration
(\(\ddot{q} \to \infty\)). In an actuator with a high gear reduction
ratio (such as \(N = 50\) to \(100\), typical in strain wave gears and
multi-stage planetary gearboxes), \(N^2\) amplifies the rotor's inertia
by a factor of \(2500\times\) to \(10{,}000\times\). Purely to
accelerate the spinning rotor, the motor demands a massive current
surge. This demand instantly drives the inverter into its current
ceiling \(I_{\max}\) and saturates phase voltage against
back-electromotive force.

Because the actuator cannot deliver the requested torque step, the
reflected inertia acts as a dynamic mechanical low-pass filter.
Delivered output motion severely lags behind the policy's candidate
setpoint, injecting unmodeled phase lag that erodes feedback stability,
excites structural resonance, and spikes mechanical stress across the
gear teeth. Motion-profile filtering (such as quintic polynomial splines
or S-curve rate limiters) must be enforced by the real-time
microcontroller to guarantee that requested joint acceleration
\(\ddot{q}\) remains within the current and voltage envelopes of the
hardware (\citeproc{ref-lynch2017modern}{Lynch and Park 2017}).

Furthermore, transmissions introduce backlash and compliance that
decouple motor and load motion. When an actuator reverses direction, the
motor shaft must rotate across mechanical clearance before the gear
teeth engage the opposite flank. Coupled with torsional compliance in
gear teeth and harmonic flexsplines
(Figure~\ref{fig-02-real-strain-wave-gear}), reversal manifests at the
output as lost motion, hysteresis, and phase lag. Measuring position
solely at the motor encoder conceals this deadband, because the motor
turns smoothly while the physical output remains stationary until
contact is made. Under reversing loads, the actual joint position
\(q_{\text{out}}\) departs from the ideal kinematic projection
\(q_{\text{motor}} / N\) by a load-dependent error
\(\delta_{\text{lost}}\). The nervous system perceives continuous
trajectory tracking, yet the physical end-effector hesitates and then
experiences an impulsive shock load as tooth engagement snaps closed.
Under severe repetitive shock loads or high reflected inertia impacts,
this localized stress concentration leads to fatigue fracture and
catastrophic tooth shear (Figure~\ref{fig-02-real-broken-gear-tooth}).

\protect\phantomsection\label{callout-ruleux2a-2.7}
\begin{fbxSimple}{callout-rule}{Architect’s Rule of Thumb:}{Impedance Control: Virtual Spring-Dampers for Contact Safety}
\phantomsection\label{callout-rule*-2.7}

When a learned policy commands joint positions in rigid environments
(\(k \ge 10^5\text{ N/m}\)), small trajectory errors
(\(\Delta x = 1\text{ mm}\)) generate massive contact forces
(\(F = k \Delta x = 100\text{ N}\)) within milliseconds.

To prevent tooth shear and gearbox overload, the bare-metal
microcontroller converts candidate position setpoints into an active
\textbf{impedance control law}:
\[\tau = J^T \left( K_p (x_{\text{des}} - x) + K_d (\dot{x}_{\text{des}} - \dot{x}) \right) + g(q)\]
By emulating a programmable virtual spring-damper system, the joint
yields compliantly upon contact, capping interaction forces and
absorbing impact shocks before destructive stress accumulates in the
mechanical linkages.

\end{fbxSimple}

\protect\phantomsection\label{callout-incidentux2a-2.8}
\begin{fbxSimple}{callout-incident}{Landmark Engineering Incident:}{Mars Climate Orbiter and Physical Transmission Overload}
\phantomsection\label{callout-incident*-2.8}

\emph{Incident overview.} On September 23, 1999, NASA lost the
\(\$327.6\text{M}\) Mars Climate Orbiter during its orbital insertion
burn (\citeproc{ref-nasa1999mco}{NASA Mishap Investigation Board 1999}).
A software unit discrepancy in ground navigation models propagated
through thruster impulses into catastrophic aerodynamic shear and total
structural destruction.

\emph{Software force calculation \& dimensional invariant breach.}
During the nine-month cruise phase, the ground station trajectory
software (\texttt{SM\_FORCES}) computed small-force thruster
firings---intended for Angular Momentum Desaturation (AMD) of the
spacecraft's reaction wheels---in English engineering units (pound-force
seconds, \(\text{lbf}\cdot\text{s}\)). The onboard flight navigation
software expected these momentum adjustments in metric SI units
(newton-seconds, \(\text{N}\cdot\text{s}\)). Because
\(1\text{ lbf}\cdot\text{s} = 4.44822\text{ N}\cdot\text{s}\), the
ground software underestimated the thruster momentum change by a factor
of \(4.45\).

\emph{Trajectory dispersion to atmospheric kinetic destruction.} Over
hundreds of desaturation maneuvers, this uncorrected unit mismatch
accumulated a systematic trajectory error. When the spacecraft executed
its Mars Orbital Insertion burn, its periapsis altitude dropped from the
planned safe trajectory of \(226\text{ km}\) down to approximately
\(57\text{ km}\), plunging the probe deep into the dense Martian
atmosphere. At \(57\text{ km}\), hypersonic dynamic pressure
(\(q = \frac{1}{2} \rho v^2\)) and thermal friction exerted aerodynamic
shear loads that far exceeded the structural yield strength of the solar
array linkages and high-gain antenna, tearing the spacecraft apart.

\emph{Systems engineering root cause.} In physical AI, mathematical
tokens generated in software (forces, torques, setpoints) are directly
converted into physical stresses (\(F = m a\),
\(\tau = N^2 J \ddot{q}\), \(\sigma = \tau / W\)) by actuators and
linkages. The Mars Climate Orbiter failure
(\citeproc{ref-nasa1999mco}{NASA Mishap Investigation Board 1999})
illustrates the fundamental physical AI invariant: \textbf{software
models cannot assume abstract unit agnosticism}. When software issues
torque or force setpoints without hardware-floor dimensional
verification, structural strain limiters, and physical plausibility
monitors, numerical calculation bugs translate directly into mechanical
shear fracture and physical kinetic destruction.

\end{fbxSimple}

\begin{figure}[t]

\centering{

\includegraphics[width=0.45\linewidth,height=\textheight,keepaspectratio]{/Users/VJ/GitHub/MLSysBook-integrate-physical/physical/book/chapters/02-body/figures/fig02_real_strain_wave_gear.jpg}

}

\caption{\label{fig-02-real-strain-wave-gear}\textbf{Strain Wave
(Harmonic Drive) Gear Component Set and Mechanical Coupling.} Authentic
component set of a robotic strain wave gear transmission (Harmonic Drive
AG). The assembly comprises the elliptical wave generator with thin-race
ball bearings (top left), the thin-walled flexible cup flexspline with
external gear teeth (top right), and the rigid circular spline with
internal gear teeth (bottom). Elastic deformation of the flexspline
allows high single-stage gear reduction (typically \(50:1\) to
\(160:1\)) with near-zero initial backlash, but introduces significant
torsional compliance, mechanical hysteresis, and structural fatigue
limits under dynamic shock loads.}

\end{figure}%

\begin{figure}[t]

\centering{

\includegraphics[width=0.55\linewidth,height=\textheight,keepaspectratio]{/Users/VJ/GitHub/MLSysBook-integrate-physical/physical/book/chapters/02-body/figures/fig02_real_broken_gear_tooth.jpg}

}

\caption{\label{fig-02-real-broken-gear-tooth}\textbf{Catastrophic Gear
Tooth Root Fatigue Fracture from Dynamic Shock Loads.} Macro-fracture of
a high-torque transmission gear tooth undergoing cyclic stress testing
at NASA Glenn Research Center. When a high-ratio geared transmission
(\(N \ge 50\)) experiences discontinuous policy torque steps or sudden
external impact, reflected rotor inertia
(\(J_{\text{ref}} = N^2 J_{\text{rotor}}\)) prevents the rotor from
backdriving. The kinetic energy concentrates entirely as mechanical
shear stress at the tooth root fillet, causing catastrophic fatigue
fracture rather than reporting a software error.}

\end{figure}%

Selecting an appropriate mechanical transmission architecture
establishes the fundamental systems trade-offs between continuous torque
density and backdrivability, as summarized in
Table~\ref{tbl-02-mechanical-drives}.

\begin{longtable}[]{@{}
  >{\raggedright\arraybackslash}p{(\linewidth - 6\tabcolsep) * \real{0.2500}}
  >{\raggedright\arraybackslash}p{(\linewidth - 6\tabcolsep) * \real{0.2500}}
  >{\raggedright\arraybackslash}p{(\linewidth - 6\tabcolsep) * \real{0.2500}}
  >{\raggedright\arraybackslash}p{(\linewidth - 6\tabcolsep) * \real{0.2500}}@{}}
\caption{\textbf{Comparative Systems Matrix of Actuator Transmission
Mechanisms.} Quantitative trade-offs across robotic drive mechanisms,
mapping reduction ratio (\(N\)), reflected rotor inertia scaling
(\(N^2\)), and operational implications for learned control
policies.}\label{tbl-02-mechanical-drives}\tabularnewline
\toprule\noalign{}
\begin{minipage}[b]{\linewidth}\raggedright
Transmission Archetype
\end{minipage} & \begin{minipage}[b]{\linewidth}\raggedright
Gear Ratio (\(N\))
\end{minipage} & \begin{minipage}[b]{\linewidth}\raggedright
Reflected Inertia (\(N^2\))
\end{minipage} & \begin{minipage}[b]{\linewidth}\raggedright
Implication for Physical AI Policies
\end{minipage} \\
\midrule\noalign{}
\endfirsthead
\toprule\noalign{}
\begin{minipage}[b]{\linewidth}\raggedright
Transmission Archetype
\end{minipage} & \begin{minipage}[b]{\linewidth}\raggedright
Gear Ratio (\(N\))
\end{minipage} & \begin{minipage}[b]{\linewidth}\raggedright
Reflected Inertia (\(N^2\))
\end{minipage} & \begin{minipage}[b]{\linewidth}\raggedright
Implication for Physical AI Policies
\end{minipage} \\
\midrule\noalign{}
\endhead
\bottomrule\noalign{}
\endlastfoot
\textbf{Direct Drive (No Gearing)} & \(1:1\) (\(N = 1\)) & \(N^2 = 1\)
(Minimal) & Transparently backdrivable; absorbs impact shocks; limited
continuous torque density
(\(\sim 2\text{--}5\text{ N}\cdot\text{m/kg}\)). \\
\textbf{Quasi-Direct Drive (QDD)} & \(3:1\text{--}10:1\) (\(N \le 10\))
& \(N^2 \le 100\) (Low) & High backdrivability; enables proprioceptive
torque sensing via motor current; high impact tolerance for dynamic
quadrupeds and arms. \\
\textbf{High-Ratio Geared (Planetary / Harmonic)} &
\(30:1\text{--}160:1\) (\(N \ge 30\)) & \(N^2 \ge 900\text{--}25{,}000\)
(Extreme) & High torque density for heavy payloads; non-backdrivable;
high-frequency policy chatter causes inverter current saturation and
tooth shear. \\
\textbf{Open-Loop Stepper Motor} & \(1:1\) (High pole count) &
\(N^2 = 1\) & High static holding torque at zero speed; silent state
loss (``lost steps'') when commanded acceleration exceeds magnetic
pull-out torque. \\
\end{longtable}

\protect\phantomsection\label{callout-archetypeux2a-2.9}
\begin{fbxSimple}{callout-archetype}{Archetype Focus:}{Class 2 Contact-Dominated Manipulation Systems}
\phantomsection\label{callout-archetype*-2.9}

In Class 2 machines (articulated robot arms, multi-fingered hands, and
contact-rich assembly workcells), the primary physical failure mode is
\textbf{destructive contact force spikes and drive-train tooth shear}.
When a robot arm makes unexpected contact with a rigid environment of
stiffness \(k\), contact force spikes at \(F = k \Delta x\). If the
joint employs a high gear reduction (\(N \ge 50\)), the reflected rotor
inertia
\(J_{\text{ref}} = N^2 J_{\text{rotor}} \ge 2500 J_{\text{rotor}}\)
prevents external impact forces from backdriving the motor. The kinetic
impact cannot be converted into electrical regeneration; instead, the
shock load concentrates entirely as mechanical shear stress across the
gear teeth. For contact-dominated manipulation, neural policies must
either be paired with backdrivable Quasi-Direct Drive (QDD) actuators
(which enable proprioceptive current sensing without external load
cells) or filtered through real-time torque limiters and impedance
controllers.

\end{fbxSimple}

For the physical AI systems engineer, these electromechanical realities
dictate two foundational design rules. First, raw neural network outputs
must never be mapped directly to joint motor inverters without smooth
continuous motion profiling; the real-time microcontroller must
interpolate discrete action chunks through \(C^2\)-continuous splines
(such as quintic polynomials or S-curve limiters) to bound commanded
jerk and prevent current saturation against reflected rotor inertia.
Second, closed-loop safety verification cannot rely solely on motor-side
encoders; the engineer must place absolute encoders and torque sensors
directly on the output load side of the transmission, ensuring that
silent mechanical failures---such as flexspline tooth fracture, shaft
shear, or transmission backlash---are immediately detected and
intercepted before damaging the workcell.

Operating continuously near current and torque limits extracts an
irreversible thermodynamic cost. Every ampere driven through the stator
winding to overcome rotor inertia or sustain peak output torque
dissipates resistive heat (\(I^2 R\)) directly into the copper coils.
While magnetic flux builds within milliseconds, the heat it generates
accumulates in the physical thermal mass of the actuator over seconds
and minutes, creating a thermal clock that bounds how long peak torque
can be delivered before insulation breaks down.

\section{Thermal Limits and Duty Cycle}\label{sec-body-2-5}

Holding a \(25\text{ kg}\) cast iron workpiece motionless in mid-air
places a deceptive demand on an articulated robot. In simulation, an
artificial intelligence agent can sustain static holding torque
indefinitely without penalty. On physical hardware, holding a static
load against gravity requires continuous phase current through the
stator coils. \textbf{Heat arrives in seconds and leaves in minutes,
making duty cycle a primary design variable rather than a thermal
footnote.}

While mechanical momentum can be redirected and magnetic flux decays
within milliseconds of opening a switch, thermal energy has no fast
physical exit. Thermal capacity is the physical integral of power loss
over time. What that capacity protects is a specific, inflexible
material threshold, namely the dielectric rating of the stator winding
insulation and the Curie temperature of permanent rotor magnets
(Figure~\ref{fig-02-real-burnt-motor-stator}). Modern motor windings use
polymer enamel coatings classified under standards such as IEC 60085,
where Class F insulation withstands a maximum continuous hotspot
temperature of \(155^\circ\text{C}\) and Class H withstands
\(180^\circ\text{C}\). Exceeding this rating softens the polymer matrix,
causing inter-turn shorts that destroy the phase winding.
Neodymium-iron-boron (\(\text{NdFeB}\)) rotor magnets face a companion
limit, where temperatures between \(80^\circ\text{C}\) and
\(150^\circ\text{C}\) cause irreversible demagnetization under strong
stator demagnetizing fields.

\begin{figure}[t]

\centering{

\includegraphics[width=0.65\linewidth,height=\textheight,keepaspectratio]{/Users/VJ/GitHub/MLSysBook-integrate-physical/physical/book/chapters/02-body/figures/fig02_real_burnt_motor_stator.jpg}

}

\caption{\label{fig-02-real-burnt-motor-stator}\textbf{Irreversible
Stator Winding Insulation Breakdown from Unmanaged Overheating.}
Internal view of an electric motor stator suffering severe thermal
degradation. Continuous high-current operation (\(I^2 R\)) without
sufficient thermal recovery intervals (\(t_{\text{off}}\)) heats stator
copper beyond the dielectric temperature rating of its Class F/H polymer
enamel coating (\(155^\circ\text{C}\)--\(180^\circ\text{C}\)). The
charred slot liners and darkened copper coils indicate inter-turn
insulation failure, leading to catastrophic phase-to-phase short
circuits. (Source: Industrial Electromechanical Diagnostics Archive, CC
BY 4.0).}

\end{figure}%

The rate of heat generation follows directly from the current commanded
through the motor phases. For an instantaneous phase current \(I(t)\)
and winding resistance \(R(T)\), resistive Joulean dissipation dumps
thermal power into the stator copper according to the relation
\(P_{\text{loss}}(t) = I^2(t) R(T)\).\sidenote{\footnotesize Stator copper resistance
  increases by \(+47.1\%\) between \(20^\circ\text{C}\) and
  \(140^\circ\text{C}\) (\(R(T) = R_0[1 + \alpha_{\text{Cu}}(T-T_0)]\)),
  inflating \(I^2 R\) heat generation during heavy duty cycles. Winding
  temperatures are monitored via RTDs or NTC thermistors linearized via
  Steinhart-Hart equations on microcontroller ADCs.} This thermal power
divides between heating the physical thermal capacitance of the winding
\(C_{\text{th}}\) and conducting heat across the stator-to-ambient
thermal resistance \(\theta_{JA}\) (also denoted \(R_{\text{th}}\)) to
the motor casing and surrounding air. The governing first-order
thermodynamic differential equation is
\[C_{\text{th}} \frac{dT(t)}{dt} = P_{\text{loss}}(t) - \frac{T(t) - T_{\text{amb}}}{\theta_{JA}} = I^2(t) R(T) - \frac{T(t) - T_{\text{amb}}}{\theta_{JA}}\]
Defining the thermal time constant \(\tau = \theta_{JA} C_{\text{th}}\),
the temperature trajectory for a constant power loss \(P_{\text{loss}}\)
starting from an initial temperature \(T(0)\) integrates to
\[T(t) = T_{\text{amb}} + (T(0) - T_{\text{amb}}) e^{-t/\tau} + P_{\text{loss}} \theta_{JA} \left(1 - e^{-t/\tau}\right)\]
In the steady-state limit as \(t \to \infty\), the equilibrium
temperature rise above ambient is
\[\Delta T_{\text{ss}} = T_{\text{ss}} - T_{\text{amb}} = P_{\text{loss}} \theta_{JA} = I^2 R(T_{\text{ss}}) \theta_{JA}\]

This formulation exposes why stator thermal limits create a hard,
non-negotiable physical boundary for machine learning policies. Neural
policies trained in simulation environments (such as Isaac Gym, MuJoCo,
or PyBullet) optimize reward functions where torque is treated as an
instantaneous, massless mathematical variable. In simulation, a policy
can hold a heavy payload indefinitely, stiffen joint impedance through
continuous high-torque antagonist co-contraction, or chatter rapidly
across contact boundaries without incurring any physical consequence. On
real physical hardware, sustained current \(I\) pumps thermal energy
into stator copper at an adiabatic initial rate
\(dT/dt \approx I^2 R / C_{\text{th}}\).

This thermodynamic cost is compounded by the positive temperature
coefficient of copper
(\(\alpha_{\text{Cu}} \approx 0.00393\text{ K}^{-1}\)). As the stator
temperature climbs from \(20^\circ\text{C}\) toward its
\(140^\circ\text{C}\) engineering ceiling, phase resistance increases by
over \(47\%\) according to the linear relation
\(R(T) = R_0 [1 + \alpha_{\text{Cu}}(T - T_0)]\). To sustain identical
commanded torque (\(\tau = K_t I\)), the inverter must deliver higher
phase voltage, dissipating \(47\%\) more resistive power (\(I^2 R(T)\))
at the exact moment thermal cooling is least effective. If an
unconstrained neural policy commands continuous peak torque, this
positive thermal feedback accelerates the winding toward dielectric
breakdown in tens of seconds.

\protect\phantomsection\label{callout-archetypeux2a-2.10}
\begin{fbxSimple}{callout-archetype}{Archetype Focus:}{Class 3 Flow-Dominated Process and Thermal Plants}
\phantomsection\label{callout-archetype*-2.10}

In Class 3 systems (grid-scale battery energy storage, liquid cooling
distribution units, and continuous thermal regulation plants), the
primary physical failure mode is \textbf{irreversible thermal
accumulation and runaway}. Unlike mechanical momentum which can be
redirected or arrested by mechanical brakes, thermal energy has no fast
physical exit: heat arrives in seconds
(\(dT/dt \approx P_{\text{loss}} / C_{\text{th}}\)) but leaves across
minutes through thermal resistance (\(\theta_{JA}\)). For physical AI
controllers regulating fluid and power plants, learned policies cannot
optimize short-term throughput at the expense of steady-state heat
balance without inducing unrecoverable thermal trips.

\end{fbxSimple}

Datasheet thermal specifications provide an optimistic baseline because
manufacturers measure them under idealized laboratory conditions,
typically mounting an unconstrained motor to a large aluminum heat sink
in free air at \(20^\circ\text{C}\) or \(25^\circ\text{C}\). On an
operational machine, the actuator operates inside a sealed chassis
alongside power electronics, battery packs, and adjacent joint motors.
Enclosing the motor eliminates ambient convective airflow, increasing
the effective thermal resistance \(\theta_{JA}\). Furthermore,
neighboring computation and power conversion stages dump heat into the
shared structural chassis, raising the local internal ambient
temperature \(T_{\text{amb}}\) well above external room temperature. A
joint actuator operating near a graphics processor or battery pack can
experience a local ambient of \(45^\circ\text{C}\) to
\(55^\circ\text{C}\) before its own coils carry current. Accurately
budgeting thermal margins requires measuring the installed heating and
cooling time constants with the enclosure fully sealed and all adjacent
electrical loads operating at their \(P_{99}\) power draw.

Because thermal capacity integrates power over time, a peak torque or
peak current rating on a motor datasheet is physically meaningless
without four accompanying parameters, specifically the initial winding
temperature \(T(0)\), the maximum allowable duration \(t_{\text{on}}\),
the operating ambient temperature \(T_{\text{amb}}\), and the required
recovery duration \(t_{\text{off}}\). An actuator capable of producing
three times its continuous rated torque can sustain that load only while
the integral of \(I^2 R\) remains within the thermal margin below the
insulation limit. If the winding begins at room temperature, that burst
can last tens of seconds. If the winding is already soaked at its
continuous operating temperature from sustained motion, the identical
current pulse reaches the insulation threshold in a fraction of that
time.

Consider the thermal budgeting for the shoulder joint of our
heavy-payload arm, with winding resistance \(R = 1.2\ \Omega\), lumped
thermal capacitance \(C_{\text{th}} = 120\text{ J/K}\), and an installed
thermal resistance \(\theta_{JA} = 2.5\text{ K/W}\), yielding a thermal
time constant \(\tau = 300\text{ s}\) (\(5.0\text{ min}\)). The motor
uses Class F insulation rated to \(155^\circ\text{C}\) and operates in
an enclosed internal ambient \(T_{\text{amb}} = 40^\circ\text{C}\). At
its continuous rated current \(I_{\text{cont}} = 5.0\text{ A}\), the
steady-state resistive loss is
\(P_{\text{cont}} = (5.0\text{ A})^2 \times 1.2\ \Omega = 30.0\text{ W}\),
producing an equilibrium temperature rise of
\(\Delta T_{\text{cont}} = 30.0\text{ W} \times 2.5\text{ K/W} = 75.0\text{ K}\)
and an equilibrium operating temperature of \(115^\circ\text{C}\),
safely below the material rating. When a learned policy commands peak
torque requiring \(I_{\text{peak}} = 15.0\text{ A}\) (\(3.0\times\)
continuous rating) to accelerate the workpiece, resistive power
dissipation jumps by a factor of nine to
\(P_{\text{peak}} = (15.0\text{ A})^2 \times 1.2\ \Omega = 270.0\text{ W}\).
If sustained indefinitely, this loss would drive asymptotic temperature
to
\(T_{\text{amb}} + P_{\text{peak}} \theta_{JA} = 40^\circ\text{C} + 270.0\text{ W} \times 2.5\text{ K/W} = 715^\circ\text{C}\),
destroying the winding. Starting cold at \(T(0) = 40^\circ\text{C}\),
the time \(t\) to reach the \(155^\circ\text{C}\) insulation limit
satisfies \(155 = 40 + 675 (1 - e^{-t / 300})\), which evaluates to
\(t = -300 \ln(1 - 115 / 675) \approx 56.0\text{ s}\). If the motor was
already operating at its continuous thermal equilibrium of
\(115^\circ\text{C}\), the same \(15.0\text{ A}\) command drives the
winding to \(155^\circ\text{C}\) according to
\(155 = 715 - (715 - 115) e^{-t / 300}\), reaching the failure boundary
in just \(t = -300 \ln(560 / 600) \approx 20.7\text{ s}\).

To prevent winding degradation under repeated motion cycles with period
\(t_{\text{cycle}} = t_{\text{on}} + t_{\text{off}}\), the machine must
operate within a safe cyclic steady-state temperature profile. Enforcing
an engineering ceiling \(T_{\text{max}} = 140^\circ\text{C}\) to
preserve a \(15^\circ\text{C}\) safety buffer below the insulation limit
restricts the allowable temperature rise above ambient to
\(\Delta T_{\text{allow}} = 140^\circ\text{C} - 40^\circ\text{C} = 100\text{ K}\).
The maximum permissible average power dissipation over the cycle is
\(\bar{P}_{\text{max}} = \Delta T_{\text{allow}} / \theta_{JA} = 100\text{ K} / 2.5\text{ K/W} = 40.0\text{ W}\).
For a profile that alternates between peak torque dissipation
\(P_{\text{peak}} = 270.0\text{ W}\) during \(t_{\text{on}}\) and zero
current during \(t_{\text{off}}\), the time-averaged power is
\(\bar{P} = D P_{\text{peak}}\), where
\(D = t_{\text{on}} / t_{\text{cycle}}\) is the duty cycle. Bounding
\(\bar{P} \le 40.0\text{ W}\) constrains the allowable duty cycle to:
\[D \le \frac{T_{\text{max}} - T_{\text{amb}}}{P_{\text{peak}} \theta_{JA}} = \frac{\Delta T_{\text{allow}}}{I_{\text{peak}}^2 R(T_{\text{max}}) \theta_{JA}} = \frac{40.0\text{ W}}{270.0\text{ W}} \approx 0.148 \quad (14.8\%)\]

Just as a computer architect cannot run a processor at peak boost clock
indefinitely without violating its Thermal Design Power (TDP), a
physical AI architect cannot command peak actuator torque without
depleting the thermal capacity of the motor windings. Because Joulean
heat arrives in seconds (\(P = I^2 R\)) while dissipation through
conduction and natural convection takes minutes
(\(\tau = \theta_{JA} C_{\text{th}} \approx 300\text{ s}\)), peak bursts
require a mandatory thermal recovery interval:
\[t_{\text{off}} = t_{\text{on}} \left( \frac{1 - D}{D} \right)\] If a
learned policy commands a \(3.0\text{ s}\) peak torque burst to brace
against an unpredicted load disturbance, the nervous system must enforce
a recovery interval of at least \(t_{\text{off}} \approx 17.3\text{ s}\)
before peak torque can be safely granted again.

\protect\phantomsection\label{callout-mathux2a-2.11}
\begin{fbxSimple}{callout-math}{Fermi Math:}{Adiabatic Stator Thermal Spike under Peak Torque Stall}
\phantomsection\label{callout-math*-2.11}

\emph{Problem:} A robotic manipulator shoulder actuator features stator
copper windings with total mass \(m_{\text{Cu}} = 0.35\text{ kg}\)
(copper specific heat capacity
\(c_{\text{Cu}} = 385\text{ J}/(\text{kg}\cdot\text{K})\), yielding
lumped thermal capacitance
\(C_{\text{th}} = m_{\text{Cu}} c_{\text{Cu}} = 134.75\text{ J/K}\)).
Nominal room-temperature phase resistance is
\(R(20^\circ\text{C}) = 0.60\ \Omega\), and installed thermal resistance
to ambient (\(T_{\text{amb}} = 40^\circ\text{C}\)) is
\(\theta_{JA} = 2.2\text{ K/W}\). The motor operates with Class F
insulation (\(T_{\text{ins}} = 155^\circ\text{C}\), enforced limit
\(T_{\max} = 140^\circ\text{C}\)). When a neural policy commands peak
holding torque drawing \(I_{\text{peak}} = 22.0\text{ A}\) from a
pre-warmed operating temperature \(T(0) = 65^\circ\text{C}\), what is
the initial heating rate, and how long until the winding reaches thermal
breakdown?

\textbf{1. Initial Joulean Dissipation (\(P_{\text{loss}}\)) with Copper
Temperature Effect:} At \(T(0) = 65^\circ\text{C}\), winding resistance
has increased under the copper temperature coefficient
(\(\alpha_{\text{Cu}} = 0.00393\text{ K}^{-1}\)):
\[R(65^\circ\text{C}) = R_0 [1 + \alpha_{\text{Cu}}(65 - 20)] = 0.60 \times [1 + 0.00393 \times 45] = 0.60 \times 1.177 = 0.706\ \Omega\]
Instantaneous Joulean thermal dissipation is:
\[P_{\text{loss}} = I_{\text{peak}}^2 R(65^\circ\text{C}) = (22.0\text{ A})^2 \times 0.706\ \Omega = 484 \times 0.706 = 341.7\text{ W}\]

\textbf{2. Initial Adiabatic Heating Rate (\(dT/dt\)):} Before thermal
conduction to ambient takes effect
(\(t \ll \tau = \theta_{JA} C_{\text{th}} \approx 296\text{ s}\)), heat
accumulation is near-adiabatic:
\[\left.\frac{dT}{dt}\right|_{t=0} \approx \frac{P_{\text{loss}}}{C_{\text{th}}} = \frac{341.7\text{ W}}{134.75\text{ J/K}} \approx 2.54^\circ\text{C/s}\]

\textbf{3. Time to Insulation Breakdown Threshold
(\(T_{\max} = 140^\circ\text{C}\)):} Asymptotic steady-state temperature
if sustained would reach
\(T_{\text{ss}} = T_{\text{amb}} + P_{\text{loss}} \theta_{JA} = 40 + 341.7 \times 2.2 = 791.7^\circ\text{C}\).
Solving the transient thermal trajectory
\(140 = 791.7 - (791.7 - 65) e^{-t / 296}\):
\[e^{-t / 296} = \frac{791.7 - 140}{791.7 - 65} = \frac{651.7}{726.7} \approx 0.8968\]
\[t = -296 \times \ln(0.8968) \approx -296 \times (-0.1089) \approx 32.2\text{ seconds}\]
In just \(32.2\text{ seconds}\) of peak stall, the motor destroys its
insulation unless derated.

\end{fbxSimple}

Because thermal accumulation is an integral state with time constants
spanning hundreds of seconds, an agent trained in simulation has no
intrinsic incentive to preserve winding health. The systems engineer
must implement a real-time discrete \(I^2 t\) thermal observer directly
on the deterministic microcontroller, sampling phase currents across
low-side shunt resistors (\(R_{\text{shunt}}\)) at \(1000\text{ Hz}\) to
track estimated coil hotspot temperatures:
\[T_{k+1} = T_k + \frac{\Delta t}{C_{\text{th}}} \left[ I_{\text{meas}, k}^2 R(T_k) - \frac{T_k - T_{\text{amb}}}{\theta_{JA}}\right]\]
where \(R(T_k) = R_0 [1 + \alpha_{\text{Cu}}(T_k - T_0)]\). If estimated
temperature \(T_{k+1}\) crosses the derating ceiling
\(T_{\max} = 140^\circ\text{C}\), the safety enforcer scales candidate
policy currents by a dynamic derating factor
\(\gamma = \max\left(0.0, \frac{T_{\text{trip}} - T_{k+1}}{T_{\text{trip}} - T_{\max}}\right)\),
clamping commands to
\(I_{\text{cmd}} = \min(I_{\text{cand}}, \gamma I_{\text{cont}})\). If
\(T_{k+1} \ge T_{\text{trip}} = 155^\circ\text{C}\), the microcontroller
trips an emergency thermal shutdown (\(I_{\text{cmd}} = 0\)), vetoing
the learned policy before dielectric breakdown occurs.

Delivering the large current pulses that generate these thermal loads
simultaneously stresses the upstream electrical supply, transferring the
physical burden from thermal dissipation in copper to transient current
delivery across the power bus.

\section{Power and Electrical Faults}\label{sec-body-2-6}

When an agile quadruped or mobile manipulator executes an explosive
evasive maneuver---a rapid multi-joint lunge to avoid an oncoming
collision---every drive motor, joint inverter, and neural accelerator
demands maximum current simultaneously from a single direct-current
distribution bus. In simulation, current flows frictionlessly from an
idealized numerical source. On physical hardware, a machine that browns
out mid-motion fails physically, and the electrical supply is a shared
resource governed by finite capacity and transient limits. \textbf{A
machine that browns out mid-motion fails physically, and the electrical
supply is a shared resource like any other.}

Rail sag, regenerative overvoltage, and thermal accumulation are not
separate electrical phenomena, but different temporal manifestations of
the same bounded energy network. While winding heating reflects energy
accumulated over seconds as thermal dissipation in copper, voltage
fluctuations occur across microseconds and milliseconds whenever current
demand exceeds instantaneous source capacity. Onboard compute units,
neural network accelerators, sensing payloads, and motor inverters all
draw from a common direct-current distribution bus. When a learned
policy commands aggressive, coordinated accelerations across multiple
joints, the resulting sum of simultaneous current draws alters the
operating voltage delivered to every node on the machine.

The total current demand on the power rail is the sum of instantaneous
consumer currents,
\(I_{\text{bus}}(t) = I_{\text{compute}}(t) + I_{\text{sensors}}(t) + \sum_{k=1}^N I_{\text{actuator},k}(t)\).
The physical supply cannot deliver infinite current instantaneously.
Every electrical distribution bus exhibits a non-zero source impedance
\(R_{\text{bus}}\),\sidenote{\footnotesize Total bus impedance \(R_{\text{bus}}\)
  combines battery internal resistance
  (\(12\text{--}20\text{ m}\Omega/\text{cell}\)), wiring harness
  resistance, and power connectors. When high-acceleration current
  spikes pull bus potential below microcontroller Brownout Reset (BOR)
  thresholds (e.g., canonical microcontroller BOR registers), hardware
  asserts an asynchronous reset on \texttt{NRST}.} combining the
internal resistance of the battery or power supply and the resistance of
wiring harnesses, connectors, and protection switches. Bulk decoupling
capacitors with total capacitance \(C_{\text{bus}}\) provide local
charge storage to smooth fast load transients. When the policy demands a
rapid step change in total current \(\Delta I_{\text{step}}\) with
risetime \(\Delta t\), the current ramps at a rate
\(dI/dt \approx \Delta I_{\text{step}} / \Delta t\). Before the primary
power supply or battery chemistry can adjust its chemical reaction rate
or control loop over its response window \(\tau_{\text{reg}}\), the bus
voltage drops by an amount determined by the resistive drop and the
capacitive charge depletion. The total transient droop is given by:
\[\Delta V_{\text{droop}} = \Delta I_{\text{step}} R_{\text{bus}} + \frac{1}{C_{\text{bus}}} \int_0^{\Delta t} I_{\text{step}}(t) \, dt \approx \Delta I_{\text{step}} R_{\text{bus}} + \frac{\Delta I_{\text{step}} \Delta t}{2 C_{\text{bus}}}.\]

The opposite boundary occurs during deceleration, where kinetic energy
stored in moving mass returns to the electrical bus. \textbf{Braking
does not discard energy, but converts mechanical kinetic energy back
into electrical energy that the supply network must either absorb or
dissipate.} If the machine is powered from an AC-to-DC supply with
output diodes that prevent reverse current, or from a battery pack whose
management system disconnects charging under cold or high-charge states,
this energy accumulates entirely in the bus capacitance
\(C_{\text{bus}}\), causing final bus voltage \(V_{\text{final}}\) to
surge according to:
\[V_{\text{final}} = \sqrt{\frac{2(E_{C,\text{init}} + E_{\text{regen}})}{C_{\text{bus}}}}.\]

\protect\phantomsection\label{callout-mathux2a-2.12}
\begin{fbxSimple}{callout-math}{Fermi Math:}{DC Distribution Rail Droop and Regenerative Voltage Surge}
\phantomsection\label{callout-math*-2.12}

\textbf{1. Transient Rail Droop under Coordinated Acceleration Step:}
For a \(48.0\text{ V}\) bus with loop resistance
\(R_{\text{bus}} = 65.0\text{ m}\Omega\), bulk capacitance
\(C_{\text{bus}} = 2200\ \mu\text{F}\), and load step
\(\Delta I_{\text{step}} = 80.0\text{ A}\) in
\(\Delta t = 2.0\text{ ms}\):
\[\Delta V_R = \Delta I_{\text{step}} R_{\text{bus}} = 80.0\text{ A} \times 0.065\ \Omega = 5.20\text{ V}\]
\[\Delta V_C = \frac{\Delta I_{\text{step}} \Delta t}{2 C_{\text{bus}}} = \frac{80.0\text{ A} \times 2.0\times 10^{-3}\text{ s}}{2 \times 2.2 \times 10^{-3}\text{ F}} \approx 36.36\text{ V}\]
\[\Delta V_{\text{droop}} = \Delta V_R + \Delta V_C = 41.56\text{ V} \implies V_{\text{bus}} = 48.0\text{ V} - 41.56\text{ V} = 6.44\text{ V}\]
The bus collapses below the Brownout Reset (BOR) threshold before the
supply can compensate.

\textbf{2. Kinetic Regenerative Surge during Emergency Braking:}
Decelerating moving mass (\(E_{\text{kin}} = 35.0\text{ J}\)) in
\(t_{\text{brake}} = 120\text{ ms}\) with efficiency
\(\eta_{\text{regen}} = 0.82\) injects
\(E_{\text{regen}} = 28.7\text{ J}\). In unbuffered bus capacitance
(\(C_{\text{bus}} = 2200\ \mu\text{F}\),
\(V_{\text{init}} = 48.0\text{ V}\),
\(E_{C,\text{init}} = 2.53\text{ J}\)):
\[V_{\text{final}} = \sqrt{\frac{2(2.53\text{ J} + 28.7\text{ J})}{2.2 \times 10^{-3}\text{ F}}} \approx 168.5\text{ V}\]
Average regenerative dump power:
\(\bar{P}_{\text{regen}} = 28.7\text{ J} / 0.120\text{ s} \approx 239.2\text{ W}\),
which destroys power MOSFETs unless diverted into a brake chopper.

\end{fbxSimple}

\protect\phantomsection\label{callout-ruleux2a-2.13}
\begin{fbxSimple}{callout-rule}{Architect’s Rule of Thumb:}{Power Bus Integrity \& Transient Decoupling}
\phantomsection\label{callout-rule*-2.13}

\emph{Power rails are shared dynamical systems; simultaneous multi-axis
torque steps will collapse bus voltage across source impedance
(\(R_{\text{bus}}\)) and trip brownout resets unless clamped by
coordinated current slew-rate limiters.}

\end{fbxSimple}

This transient surge exceeds the breakdown voltage of standard
\(60.0\text{ V}\) or \(80.0\text{ V}\) power MOSFET switches, destroying
inverter stages or triggering overvoltage protection unless a dynamic
brake chopper diverts regenerative power into a resistive
dump.\sidenote{\footnotesize During high-deceleration braking (Category 1 SS1),
  kinetic energy charges DC bus capacitors
  (\(\Delta E = \frac{1}{2} m v^2\)). Hardware comparators activate
  dynamic braking chopper MOSFETs to dissipate regenerative current
  through ceramic resistors, preventing bus overvoltage faults.}

When bus voltage crosses an undervoltage threshold under peak
acceleration, the electrical fault cascades into a mechanical failure
through hardware supervisor resets. A supply rail dropping below the
undervoltage lockout threshold causes voltage supervisors to pull the
hardware reset line on microcontrollers and motor gate drivers.
Microcontroller clocks halt immediately, volatile register
configurations vanish, and digital communication buses drop into
high-impedance states. However, an acceleration command issued by the
policy just before the reset remains physically active in the hardware
state until the gate driver stage collapses. If the gate driver enters a
high-impedance unpowered state, phase current ceases and the motor
produces zero electromagnetic torque. The physical body, carrying
velocity and momentum from the preceding command, coasts ballistically
under gravity and inertia, colliding with workpieces or joint hardstops.
If the driver instead defaults to shorting low-side MOSFETs for dynamic
braking, back-EMF creates immediate counter-torque that jerks the
mechanics at the joint limits.

Power removal is an unmanaged state transition rather than an automatic
return to safety. Engineers frequently assume that dropping bus power or
tripping an undervoltage circuit brings a machine to a safe halt. In
physical reality, cutting electrical power leaves an actuator in one of
three unpowered states, namely free coasting, phase-short dynamic
braking, or spring-engaged mechanical braking. Each state enforces a
distinct mechanical behavior. An electromechanical friction brake
requires \(30\text{ ms}\) to \(80\text{ ms}\) for its mechanical spring
to overcome magnetic coil decay and engage its friction pads, during
which gravity accelerates an overhead arm downward. Phase-short dynamic
braking provides resistance proportional to velocity, generating zero
holding torque at zero speed. Free coasting allows gravity to accelerate
unconstrained joints without resistance. The actual trajectory a machine
follows after an electrical cutoff must be measured under full
mechanical payload rather than deduced from circuit diagrams. The choice
of which safe fallback state the machine should enter belongs to system
safety architecture in Chapter~\ref{sec-enforcement}, but the physical
dynamics of an unpowered joint are unalterable mechanical constraints of
the body.

For the physical AI systems engineer, power integrity is a foundational
architectural constraint bridging software planning and electrical
physics. Rather than treating power rails as infinite voltage sources,
the engineer must enforce coordinated power budgets and current
slew-rate limits across compute accelerators and motor drives. The
deterministic real-time safety layer must arbitrate simultaneous
multi-axis acceleration commands to prevent compound \(dI/dt\)
transients, while dynamic brake choppers protect inverter electronics
from regenerative feedback during emergency stops. Crucially, because
power loss leaves actuators in unpowered physical states governed by
gravity and momentum, the engineer must empirically characterize every
unpowered fallback trajectory under full mechanical payload, ensuring
that electrical trips resolve into controlled, predictable deceleration
rather than unguided kinetic collisions. Characterizing these dynamic
electrical and mechanical failure boundaries cannot rely on nominal
catalog specs or idealized CAD models; it demands disciplined bench
measurement.

\section{Measuring a Machine's Own Limits}\label{sec-body-2-7}

Before a high-torque actuator is integrated into an autonomous robot
chassis, it must undergo rigorous characterization on a dedicated
dynamometer test stand (Figure~\ref{fig-02-real-motor-dyno-bench}). In
software design, parameters are often treated as clean nominal constants
imported from a configuration file. In physical systems engineering,
\textbf{a measurement that cannot state its instrument, its percentile,
and its sample count cannot support an argument later.}

When an engineer characterizes a machine's body, the measurement must
begin from the specific decision it will govern in the nervous system or
safety layer. That decision determines the physical quantity, the
physical observation point, the required sensor bandwidth, the
measurement resolution, and the reference clock. If the nervous system
must clamp motor velocity to prevent dynamic braking from exceeding an
optical boundary, measuring steady-state motor current with an averaging
bench meter yields no usable number. The measurement must instead
capture the instantaneous peak phase current and deceleration trajectory
at the exact point of electrical cutoff, using an isolated current probe
sampled at \(100\text{ kHz}\) and an optical encoder sampled
synchronously on a common hardware timebase. As demonstrated on the NASA
Armstrong electric propulsion dynamometer test bench
(Figure~\ref{fig-02-real-motor-dyno-bench}), selecting an instrument
whose analog bandwidth falls below the electrical rise time attenuates
the true transient peak, leaving software engineers to build safety
margins on top of measurement artifacts rather than physical reality.

\begin{figure}[t]

\centering{

\includegraphics[width=0.85\linewidth,height=\textheight,keepaspectratio]{/Users/VJ/GitHub/MLSysBook-integrate-physical/physical/book/chapters/02-body/figures/fig02_real_motor_dyno_bench.jpg}

}

\caption{\label{fig-02-real-motor-dyno-bench}\textbf{Instrumented
Electric Motor Dynamometer and Propulsion Test Bench.} An electric motor
undergoing characterization on a dedicated high-bandwidth test stand at
NASA Armstrong Flight Research Center. Direct empirical measurement of
peak torque saturation, electrical time constants (\(\tau_e = L/R\)),
continuous thermal rise (\(\Delta T = P_{\text{loss}} \theta_{JA}\)),
and regenerative power feedback requires synchronizing multi-channel
load cells, optical encoders, isolated current probes, and thermocouple
arrays on a common hardware timebase rather than relying on manufacturer
catalog ratings. (Source: NASA Armstrong Flight Research Center,
Propulsion Systems Laboratory, Public Domain).}

\end{figure}%

Every probe intrudes on the physical process it observes. In software,
instrumentation often introduces context switches, cache eviction, bus
contention, and core heating, altering the execution timing and
scheduling jitter under test. In hardware, attaching a passive voltage
probe adds capacitive loading that distorts high-frequency digital
communication lines, inserting a current shunt in series with a motor
drive drops the rail voltage enough to trigger an early brownout under
load, and placing a thermocouple inside a motor housing alters heat
dissipation paths. On a desktop development platform, motor driver
current sensing uses precision low-side current sense shunts
(\(R_{\text{shunt}} = 10\text{ m}\Omega\text{ to }50\text{ m}\Omega\))
routed to dedicated analog-to-digital converter (ADC) channels on the
real-time microcontroller. Instrument uncertainty is therefore a
compound budget rather than a single tolerance number from a calibration
sheet. It combines sensor linearity, analog-to-digital quantization
error, frequency-dependent phase lag, probe loading effects, and clock
drift between distributed acquisition channels. When an optical encoder
with \(0.05^\circ\) quantization is sampled at \(1000\text{ Hz}\) to
compute joint velocity through discrete backward differences, numerical
differentiation noise introduces velocity variance of several degrees
per second that exists entirely in the measurement pipeline rather than
in the mechanical joint.

A physical limit is meaningless without the complete environmental and
computational context of the machine during the test. Every recorded
limit must state the mechanical payload, ambient temperature, supply
rail voltage, enclosure state, firmware version, and concurrent
computational workload present during the run. An actuator thermal time
constant measured on an open bench at \(20^\circ\text{C}\) with free
convection differs by an order of magnitude from the same actuator
operating inside a sealed composite fairing under direct sunlight at
\(45^\circ\text{C}\). If an electrical characterization omits whether
vision-language model inference was running on the adjacent neural
accelerator during motor stalls, the measured power rail stability
reflects an idle board rather than the combined electrical draw of full
compute and peak actuator torque. Documenting these conditions prevents
engineers from treating a laboratory measurement taken under ideal
cooling and fresh batteries as an absolute physical boundary of the
machine in field deployment.

Systems engineering demands strict discipline in distinguishing measured
physical limits from derived models, manufacturer ratings, and
illustrative figures. Four distinct categories of numbers govern a
machine, comprising direct measurements, derived values, datasheet
ratings, and illustrative examples. Direct measurements are raw,
calibrated time series with specified instruments, sample counts, and
empirical percentiles such as \(P_{50}\), \(P_{99}\), and \(P_{99.9}\).
Derived quantities are computed from direct measurements through
physical models, such as estimating rotor winding temperature from
terminal voltage, phase current across current sense shunts, and an
ambient thermistor trace, which requires propagating the uncertainties
of all three underlying sensors through the thermal model. Datasheet
ratings are nominal manufacturer claims reflecting idealized laboratory
conditions or continuous thermal equilibrium, which rarely hold during
transient impulse loads. Illustrative numbers are analytical examples
used to demonstrate arithmetic or scaling principles in design
calculations, and they must be explicitly labeled as analytical examples
so that no safety margin or architectural timeout relies on them as
measured ground truth.

The absence of a failure during testing establishes a statistical bound
on the failure rate, not an assertion of absolute safety.\sidenote{\footnotesize In
  non-parametric reliability estimation, zero failures across \(n\)
  i.i.d. Bernoulli trials yields the rule-of-three upper bound
  \(p_{\max} \approx 3/n\) at \(95\%\) confidence (exact
  Clopper-Pearson: \(p = 1 - \alpha^{1/n}\)). This bound collapses if
  thermal drift or wear introduces non-stationary confounding.} When an
empirical test of \(n\) independent trials yields zero limit
exceedances, the underlying physical system is not proven to be
defect-free. Let \(p\) represent the true probability of a failure or
limit violation on any single trial. The probability of observing zero
failures across \(n\) independent, identically distributed trials from a
stationary distribution is \((1 - p)^n\). To assert with statistical
confidence \((1 - \alpha)\), such as a \(95\%\) confidence level where
\(\alpha = 0.05\), that the true per-trial failure rate is bounded above
by an upper limit \(p_{\text{max}}\), the observation must satisfy
\((1 - p_{\text{max}})^n \le \alpha\). Taking the natural logarithm of
both sides gives \(n \ln(1 - p_{\text{max}}) \le \ln \alpha\). Applying
the standard linear approximation \(\ln(1 - p) \approx -p\) for small
failure probabilities yields \(-n p_{\text{max}} \le \ln \alpha\), which
simplifies to \(p_{\text{max}} \ge \frac{-\ln \alpha}{n}\). For a
\(95\%\) confidence level where
\(-\ln(0.05) \approx 2.996 \approx 3.0\), the upper bound on the failure
rate is \(p_{\text{max}} \approx \frac{3}{n}\). If an engineer runs
\(n = 3000\) consecutive pick-and-place trajectories on a robotic
manipulator without a brownout or joint collision, this empirical data
supports only the claim that the true per-cycle failure rate is below
\(1.0 \times 10^{-3}\) at \(95\%\) confidence, or roughly one failure
per thousand operations. This derivation holds only when trials are
independent and drawn from a stationary distribution. If motor windings
accumulate heat across successive trials, or if mechanical wear
increases gear backlash over time, consecutive cycles are correlated and
non-stationary. Testing a machine for \(3000\) cycles in a cool morning
laboratory provides zero mathematical guarantee about its failure rate
in the afternoon, when thermal equilibrium increases winding resistance,
lowers peak torque, and shifts the physical system into a brownout
regime.

\protect\phantomsection\label{callout-ruleux2a-2.14}
\begin{fbxSimple}{callout-rule}{Architect’s Rule of Thumb:}{Zero-Failure Statistical Reliability Bound ($p_{\max} \approx 3/n$)}
\phantomsection\label{callout-rule*-2.14}

When an empirical reliability qualification trial of \(n\) independent,
identically distributed operations yields zero limit violations, the
upper bound on the true per-cycle failure probability \(p_{\max}\) at
confidence level \((1 - \alpha)\) is:
\[p_{\max} \approx \frac{-\ln \alpha}{n} \quad \left(p_{\max} \approx \frac{3}{n} \text{ at } 95\%\text{ confidence}\right)\]
This mathematical bound remains valid only while operating conditions
(temperature, battery state of charge, and mechanical wear) remain
strictly stationary across trials.

\end{fbxSimple}

For the physical AI systems engineer, empirical characterization is not
a one-time validation checkbox, but the creation of an immutable
physical contract. Furthermore, this contract must account for
\textbf{lifecycle degradation}: in deployed fleets, physical limits are
non-stationary over operational time. Brake pad friction glazing and
tire tread wear degrade the sustainable deceleration limit
\(a_{\text{brake}}\) by \(15\text{--}30\%\); chemical aging in
lithium-ion battery packs increases internal DC bus impedance
\(R_{\text{bus}}\), exacerbating voltage droop
(\(\Delta V_{\text{droop}}\)) during dynamic maneuvers; and mechanical
gear tooth wear widens backlash clearance \(\delta_{\text{lost}}\).
Every parameter stored in a machine specification---whether a current
ceiling, thermal time constant, or braking deceleration---must record
its physical units, observation points, instrument bandwidth, sample
count, tail percentiles (\(P_{99}\) and \(P_{99.9}\)), and provenance,
with periodic recalibration across the hardware lifecycle. Without this
rigorous physical grounding, safety envelopes remain analytical fictions
that dissolve under the unyielding dynamics of the real world. Yet even
perfectly characterized limits do not remain isolated in operation; on
an integrated machine, physical constraints couple together, creating
compounding stress cascades.

\section{The Cross-Archetype Matrix}\label{sec-body-2-8}

The five physical constraints analyzed across this chapter do not act as
isolated boundaries; they couple together into a unified system
constraint matrix (Table~\ref{tbl-02-cross-archetype-matrix}). In
digital computing, software operates under an abstraction of
virtualization: if an inference pass runs late, the user experiences
minor interface lag; if memory is tight, the operating system pages data
to disk. But in an embodied machine, the physical body provides zero
elasticity. When an unconstrained neural policy issues an action
command, that command draws immediate electrical current, generates
non-negotiable kinetic momentum, accelerates reflected rotor inertia,
and dissipates non-reversible Joulean heat into motor stators.

How these physical constraints bind depends directly on the machine's
physical archetype:

\begin{itemize}
\tightlist
\item
  \textbf{Class 1 (Mass-Dominated Mobility):} For autonomous mobile
  robots, automated guided vehicles, and legged platforms,
  \textbf{Momentum and Stopping Distance} is the primary binding limit.
  Kinetic energy grows quadratically with travel speed
  (\(E_{\text{kin}} = \frac{1}{2} m v^2\)), making velocity the master
  multiplier of spatial risk. Sensor age directly expands unguided
  reaction displacement (\(v_0 \Delta t_{\text{delay}}\)), while
  simultaneous multi-axis drive acceleration threatens battery bus
  collapse.
\item
  \textbf{Class 2 (Contact-Dominated Manipulation):} For articulated
  robotic arms, dexterous hands, and automated assembly fixtures,
  \textbf{Actuator Bandwidth, Reflected Rotor Inertia
  (\(N^2 J_{\text{rotor}}\)), and Contact Stiffness} form the primary
  binding envelope. Discontinuous policy setpoints impact rigid fixtures
  with steep force rise rates (\(F = k \Delta x\)), risking gear tooth
  fatigue fracture. Holding heavy payloads against gravity draws
  continuous static stall current, rapidly accumulating thermal stress.
\item
  \textbf{Class 3 (Flow-Dominated Process \& Thermal Plants):} For
  battery energy storage systems, server cooling distribution units, and
  power electronics converters, \textbf{Thermal Dissipation and
  Electrical Power Integrity} govern survival. Long thermal time
  constants (\(\tau = \theta_{JA} C_{\text{th}}\)) hide continuous heat
  accumulation, where a small continuous duty-cycle overload slowly
  drives winding insulation or battery cells into thermal runaway.
\end{itemize}

\begin{longtable}[]{@{}
  >{\raggedright\arraybackslash}p{(\linewidth - 8\tabcolsep) * \real{0.2000}}
  >{\raggedright\arraybackslash}p{(\linewidth - 8\tabcolsep) * \real{0.2000}}
  >{\raggedright\arraybackslash}p{(\linewidth - 8\tabcolsep) * \real{0.2000}}
  >{\raggedright\arraybackslash}p{(\linewidth - 8\tabcolsep) * \real{0.2000}}
  >{\raggedright\arraybackslash}p{(\linewidth - 8\tabcolsep) * \real{0.2000}}@{}}
\caption{\textbf{Cross-Archetype Physical Constraint Matrix.} Systematic
mapping of the five non-negotiable physical constraints to their
governing laws, software illusions, archetype failure modes, and
deterministic systems engineering
takeaways.}\label{tbl-02-cross-archetype-matrix}\tabularnewline
\toprule\noalign{}
\begin{minipage}[b]{\linewidth}\raggedright
Physical Budget / Factor
\end{minipage} & \begin{minipage}[b]{\linewidth}\raggedright
Governing First Principle \& Scaling Law
\end{minipage} & \begin{minipage}[b]{\linewidth}\raggedright
Software Illusion vs.~Hardware Ground Truth
\end{minipage} & \begin{minipage}[b]{\linewidth}\raggedright
Primary Binding Archetype
\end{minipage} & \begin{minipage}[b]{\linewidth}\raggedright
Systems Architect Takeaway \& Runtime Invariant
\end{minipage} \\
\midrule\noalign{}
\endfirsthead
\toprule\noalign{}
\begin{minipage}[b]{\linewidth}\raggedright
Physical Budget / Factor
\end{minipage} & \begin{minipage}[b]{\linewidth}\raggedright
Governing First Principle \& Scaling Law
\end{minipage} & \begin{minipage}[b]{\linewidth}\raggedright
Software Illusion vs.~Hardware Ground Truth
\end{minipage} & \begin{minipage}[b]{\linewidth}\raggedright
Primary Binding Archetype
\end{minipage} & \begin{minipage}[b]{\linewidth}\raggedright
Systems Architect Takeaway \& Runtime Invariant
\end{minipage} \\
\midrule\noalign{}
\endhead
\bottomrule\noalign{}
\endlastfoot
\textbf{1. Measurement Freshness \& Age (§2.2)} &
\(\Delta t_{\text{fresh}} \le \frac{e_{\max}}{v_{\max}}\)\(\Delta t_{\text{e2e}} = \sum_{k=1}^8 \Delta t_k\)
& \emph{Illusion:} High inference rate (\(50\text{ Hz}\)) guarantees
real-time response.\emph{Truth:} Age begins at optical exposure;
pipeline delay adds phase lag (\(-\omega \tau\)). & \textbf{Class 1 \&
2:} Erodes clearance buffer and destabilizes high-gain contact loops. &
Tag frames with hardware timestamps at transduction; check
\(t_{\text{now}} - t_{\text{transduction}} \le \Delta t_{\text{fresh}}\)
at motor registers; reject stale commands. \\
\textbf{2. Momentum \& Stopping Distance (§2.3)} &
\(d_{\text{stop}} = v_0 \Delta t_{\text{delay}} + \frac{v_0^2}{2 a_{\text{brake}}} \le D_{\text{clear}}\)
& \emph{Illusion:} Deceleration can be commanded
instantaneously.\emph{Truth:} Braking distance scales quadratically
(\(v_0^2\)); kinetic energy requires friction work (\(W = F d\)). &
\textbf{Class 1 (Mobility):} Primary constraint. Doubling speed
quadruples braking distance (\(251\%\) total increase). & Solve dynamic
clearance quadratic
\(v_{\max}(D_{\text{clear}}, \Delta t_{\text{delay}})\) at
\(1\text{ kHz}\) in deterministic MCU loop; clamp policy velocity. \\
\textbf{3. Actuator \& Reflected Inertia (§2.4)} &
\(J_{\text{ref}} = N^2 J_{\text{rotor}}\)\(\tau_e = \frac{L}{R}\) ;
\(I(t) \le I_{\max}\) & \emph{Illusion:} Motor torque tracks policy
setpoints without mechanical resistance.\emph{Truth:} Gearbox squares
rotor inertia (\(N^2\)); inductance delays current rise; shock steps
shear teeth. & \textbf{Class 2 (Manipulation):} Primary constraint.
Rigid contacts create extreme \(dF/dt\) force spikes and gear fatigue. &
Enforce \(C^2\) spline (S-curve) trajectory rate limits; sample
load-side optical encoders; deploy quasi-direct drive (QDD)
actuators. \\
\textbf{4. Thermal Limits \& Duty Cycle (§2.5)} &
\(C_{\text{th}} \frac{dT}{dt} = I^2 R(T) - \frac{\Delta T}{\theta_{JA}}\)\(D \le \frac{T_{\max} - T_{\text{amb}}}{P_{\text{peak}} \theta_{JA}}\)
& \emph{Illusion:} Motors deliver peak catalog torque
continuously.\emph{Truth:} Joulean heat (\(I^2 R\)) has positive thermal
feedback (\(\alpha_{\text{Cu}}\)); heat accumulates in sealed housings.
& \textbf{Class 3 (Thermal) \& Class 2:} Static holding draws stall
current, overheating coils in seconds. & Run discrete \(I^2 t\) thermal
observers at \(1\text{ kHz}\) on MCU; dynamically derate speed and
torque when hotspot models approach limits. \\
\textbf{5. Electrical Power Integrity (§2.6)} &
\(\Delta V_{\text{droop}} \approx \Delta I R_{\text{bus}} + \frac{\Delta I \Delta t}{2 C_{\text{bus}}}\)\(V_{\text{final}} = \sqrt{\frac{2(E_{C} + E_{\text{regen}})}{C_{\text{bus}}}}\)
& \emph{Illusion:} Power rails act as infinite, ideal voltage
sources.\emph{Truth:} Bus impedance collapses rail voltage under
combined compute+motor bursts; braking surges overvoltage. &
\textbf{Class 1 \& 3:} Peak multi-axis acceleration trips
microcontroller BOD resets; regen surges destroy inverter FETs. &
Enforce coordinated compute-actuation power budgets and \(dI/dt\)
limits; install dynamic brake choppers; characterize unpowered
fallbacks. \\
\end{longtable}

When an engineer integrates a machine, these five budgets do not fail in
isolation. They form a coupled constraint hypersurface where exceeding a
margin in one domain triggers an immediate failure cascade across
adjacent domains.

\section{Systems Stress-Tests}\label{sec-body-2-9}

To transition from individual physical equations to whole-system
engineering judgment, consider four foundational dilemmas where physical
limits interact across subsystem boundaries to challenge an embodied
machine.

\subsection*{Stress-Test 1: The Thermal-Latency
Avalanche}\label{stress-test-1-the-thermal-latency-avalanche}
\addcontentsline{toc}{subsection}{Stress-Test 1: The Thermal-Latency
Avalanche}

\emph{The scenario.} An autonomous mobile delivery robot operates in a
high-temperature warehouse (\(T_{\text{amb}} = 45^\circ\text{C}\)).
Continuous high-torque maneuvering heats motor windings to
\(110^\circ\text{C}\), elevating the internal chassis air to
\(55^\circ\text{C}\). The onboard neural network accelerator, executing
vision-language navigation policies under heavy compute load, approaches
its \(95^\circ\text{C}\) silicon junction ceiling. This thermal stress
triggers Dynamic Voltage and Frequency Scaling (DVFS), dropping
accelerator clock frequency from \(2.0\text{ GHz}\) to
\(1.3\text{ GHz}\) and inflating inference latency from \(20\text{ ms}\)
to \(31\text{ ms}\).

\emph{The systems dilemma.} Software power managers treat DVFS as an
invisible, transparent background optimization. But on an embodied
machine operating at \(v_0 = 1.5\text{ m/s}\), that \(+11\text{ ms}\)
compute delay injects an unguided reaction displacement of
\(\Delta d_{\text{react}} = 1.65\text{ cm}\). Tracking controllers
expecting \(50\text{ Hz}\) state updates miss deadlines, phase margin
erodes, and the machine exceeds its defended clearance before mechanical
braking can even begin.

\emph{The systems verdict: Coupled Thermal-Timing Invariant.} Actuator
heat dissipation directly inflates computing latency across shared
chassis enclosures. A robust physical AI system must integrate thermal
telemetry with real-time velocity derating, dynamically lowering travel
speed when compute silicon enters thermal throttling.

\subsection*{Stress-Test 2: The Velocity-Clearance-Thermal
Cascade}\label{stress-test-2-the-velocity-clearance-thermal-cascade}
\addcontentsline{toc}{subsection}{Stress-Test 2: The
Velocity-Clearance-Thermal Cascade}

\emph{The scenario.} An automated logistics cart increases its travel
speed from \(v_0 = 1.0\text{ m/s}\) to \(v_0 = 2.0\text{ m/s}\) to meet
warehouse throughput targets, defending a fixed obstacle clearance
\(D_{\text{clear}} = 1.50\text{ m}\) with deceleration limit
\(a_{\text{brake}} = 2.0\text{ m/s}^2\), motor resistance
\(R = 0.8\ \Omega\), and thermal resistance
\(\theta_{JA} = 2.0\text{ K/W}\).

\emph{The systems dilemma.} Doubling travel speed doubles reaction
travel linearly (\(d_{\text{react}} \propto v\)), but quadruples braking
travel quadratically (\(d_{\text{brake}} \propto v^2\)), expanding total
stopping distance from \(0.33\text{ m}\) to \(1.16\text{ m}\) (\(251\%\)
increase) and compressing available clearance buffer from
\(1.17\text{ m}\) down to just \(0.34\text{ m}\). Simultaneously,
high-torque emergency braking demands \(20\text{ A}\) of current,
dumping \(320\text{ W}\) into stator coils (\(4\times\) heating) and
restricting allowable cyclic duty cycle to \(D \le 12.5\%\).

\emph{The systems verdict: Velocity Scaling Law.} Velocity is the master
amplifier of physical risk: it simultaneously expands stopping distance
quadratically, demands peak current against reflected rotor inertia
(\(N^2 J_{\text{rotor}}\)), and accelerates thermal breakdown along
exponential curves.

\subsection*{Stress-Test 3: The Peak-Inference Rail
Collapse}\label{stress-test-3-the-peak-inference-rail-collapse}
\addcontentsline{toc}{subsection}{Stress-Test 3: The Peak-Inference Rail
Collapse}

\emph{The scenario.} A mobile robot powered by a nominal
\(24.0\text{ V}\) lithium-ion battery with internal loop resistance
\(R_{\text{bus}} = 60.0\text{ m}\Omega\) navigates a raised doorway
threshold. A vision transformer inference burst draws \(15.0\text{ A}\)
from the compute rail at the exact millisecond all drive and arm motors
accelerate over the lip, demanding an additional \(55.0\text{ A}\) of
motor current.

\emph{The systems dilemma.} Software assumes compute accelerators and
motor inverters draw from independent, ideal power rails. In reality,
the combined \(70.0\text{ A}\) current transient across
\(R_{\text{bus}} = 60.0\text{ m}\Omega\) induces an instantaneous
voltage droop of \(4.2\text{ V}\). If the battery has partially
discharged to \(21.5\text{ V}\), the bus collapses to \(17.3\text{ V}\),
breaching the \(18.0\text{ V}\) dropout threshold of onboard buck
regulators and triggering a hardware Brownout Detector (BOD) reset on
the real-time microcontroller.

\emph{The systems verdict: Shared Rail Invariant.} Electrical bus
capacity is a shared, instantaneous physical resource. Peak neural
compute bursts coinciding with aggressive mechanical maneuvers can
collapse rail voltages, resetting safety microcontrollers at the exact
instant of maximum kinematic momentum.

\subsection*{Stress-Test 4: The Compound Stress
Protocol}\label{stress-test-4-the-compound-stress-protocol}
\addcontentsline{toc}{subsection}{Stress-Test 4: The Compound Stress
Protocol}

\emph{The scenario.} An engineering team qualifies an articulated robot
arm via two independent bench trials: testing maximum rated payload
(\(15\text{ kg}\)) in a \(20^\circ\text{C}\) laboratory, and testing
empty-gripper trajectory tracking in a \(45^\circ\text{C}\)
environmental chamber. Both isolated tests pass with nominal margins.
The robot is then deployed in a foundry handling maximum payload at
\(45^\circ\text{C}\) with a low battery pack.

\emph{The systems dilemma.} Isolated single-variable testing creates a
false sense of safety. Near battery depletion cutoff, inverters must
draw higher phase current to sustain identical mechanical power
(\(P = V I\)). This elevated current dumps \(I^2 R(T)\) resistive heat
into motor windings whose resistance has already increased by \(47\%\)
due to the \(45^\circ\text{C}\) ambient, driving the arm into
catastrophic thermal shutdown and bus undervoltage within minutes.

\emph{The systems verdict: Compound Qualification Law.} Physical limits
form an interconnected multidimensional hypersurface rather than
independent bounding boxes. Validating an embodied machine requires
subjecting the complete system to simultaneous worst-case thermal,
electrical, and kinematic stress.

\section{Summary}\label{sec-body-2-10}

These five physical limits bound everything built across the system
architecture in Part II. Everything that follows must operate inside the
envelope where mechanics, thermodynamics, and electromagnetism govern
physical hardware. As summarized in the Cross-Archetype Physical
Constraint Matrix (Table~\ref{tbl-02-cross-archetype-matrix}), an
engineer cannot negotiate with momentum, dissipate heat faster than
thermal conductance allows, or transport information across a wire
before voltage propagates. The physical body presents five boundaries
that define what the machine can execute without structural damage or
state loss:

\begin{enumerate}
\def\labelenumi{\arabic{enumi}.}
\tightlist
\item
  \textbf{Observation Age:} Sensor signals arrive with non-zero age
  (\(t_{\text{age}}\)), bounded by the physical freshness deadline
  \(\Delta t_{\text{fresh}} \le e_{\max} / v_{\max}\).
\item
  \textbf{Stopping Envelope:} Moving masses sweep through space along an
  unavoidable stopping envelope
  \(d_{\text{stop}} = v_0 \Delta t_{\text{delay}} + \frac{v_0^2}{2 a_{\text{brake}}} \le D_{\text{clear}}\).
\item
  \textbf{Delivered Actuation:} Actuators produce bounded delivered
  torque governed by reflected rotor inertia
  (\(J_{\text{ref}} = N^2 J_{\text{rotor}}\)), winding rise time
  (\(\tau_e = L/R\)), and gear tooth shear limits.
\item
  \textbf{Energy \& Thermal State:} Joulean heat accumulation
  \(P_{\text{loss}} = I^2 R(T)\) under positive thermal feedback
  (\(\alpha_{\text{Cu}}\)) restricts allowable cyclic duty cycles
  (\(D\)).
\item
  \textbf{Electrical Power Integrity:} Shared DC bus impedance droop
  (\(\Delta V_{\text{droop}}\)) and regenerative overvoltage surges
  (\(V_{\text{final}}\)) dictate power delivery stability.
\end{enumerate}

To design a system around these constraints, an engineer must replace
qualitative notions of hardware capacity with exact budgetable
quantities. Observation age is measured as the duration
\(t_{\text{age}}\) from the physical capture of a state change at a
transducer to the moment that measurement updates an internal state
estimate. Stopping envelope is budgeted as the swept physical distance
\(d_{\text{stop}}\) or volume required to bring all kinematic momentum
to zero under maximum braking effort, including sensing delay, brake
engagement latency, and deceleration dynamics. Delivered actuation is
quantified as the usable force, torque, or mechanical bandwidth
delivered at the load interface across joint velocities, bounded by peak
current limits and transmission stiffness. Energy state is budgeted
through two coupled physical mechanisms, specifically the voltage sag
\(\Delta V_{\text{rail}}\) on electrical supply lines under peak motor
transients and the temperature rise \(\Delta T\) across winding
insulation and semiconductor junctions under continuous duty cycles.
Worst-case response time is measured as the tail latency bound,
evaluated at \(P_{99.9}\) or bounded by a deterministic hardware timer,
spanning signal acquisition, transport, compute execution, and actuator
drive update.

Every physical limit is conditional on a declared operating
configuration and loses authority when that configuration changes. A
measured stopping distance or thermal time constant is not a static
property of a component, and it cannot be transferred from a component
datasheet into a safety case without qualification. In an analytical
example of a robotic mobile base carrying a \(50.0\text{ kg}\) payload
at \(2.0\text{ m/s}\), the stopping distance doubles if the payload mass
increases to \(100.0\text{ kg}\) on the same friction surface, or if
surface contaminants reduce the tire traction coefficient by half.
Similarly, a motor winding that operates at a steady-state temperature
of \(75.0^{\circ}\text{C}\) in a \(20.0^{\circ}\text{C}\) laboratory
reaches its thermal destruction threshold if deployed in a
\(45.0^{\circ}\text{C}\) industrial facility under the same duty cycle.
When ambient temperature rises, payload mass increases, supply voltage
drops near battery depletion cutoff, or structural wear degrades
mechanical stiffness, the physical boundaries contract. Operating
outside declared environmental and mechanical baselines invalidates
every timing, force, and energy budget derived from those baselines.

Maintaining integrity across a physical system requires preserving a
complete engineering record for every limit before integrating software
components. That record specifies six mandatory elements for each
boundary. It documents the numerical value or empirical distribution,
such as a \(P_{99.9}\) latency of \(4.2\text{ ms}\) or a maximum
deceleration torque of \(120.0\text{ N}\cdot\text{m}\). It specifies the
exact operating conditions under which that number was established,
including payload, ambient temperature, supply voltage, and surface
properties. It records the measurement uncertainty, identifying the
error bounds of the bench instrumentation used to capture the data. It
allocates a deliberate engineering margin that separates normal
operating targets from absolute physical breakdown thresholds. It
designates a hardware or low-level software detector, such as a
phase-current sensor, optical interrupter, or supply-rail voltage
monitor, responsible for catching boundary excursions in real time. It
records provenance, distinguishing whether a bound is derived from
analytical mechanics or measured directly on physical hardware under
stress.

\protect\phantomsection\label{callout-takeawaysux2a-2.15}
\begin{fbxSimple}{callout-takeaways}{Key Chapter Takeaways}{}
\phantomsection\label{callout-takeaways*-2.15}

\begin{itemize}
\tightlist
\item
  \textbf{Freshness is Measured at Transduction:} Observation age begins
  at photon capture on sensor silicon, not at model inference. Carry
  hardware timestamps to the motor register boundary and reject stale
  commands
  (\(t_{\text{now}} - t_{\text{transduction}} \le \Delta t_{\text{fresh}}\)).
\item
  \textbf{Velocity is the Master Risk Multiplier:} Travel speed scales
  reaction travel linearly (\(d_{\text{react}} \propto v\)), braking
  distance quadratically (\(d_{\text{brake}} \propto v^2\)), and stator
  heat dissipation quadratically (\(I^2 R\)).
\item
  \textbf{Gearboxes Reflect Rotor Inertia Quadratically:} Gear ratio
  \(N\) multiplies rotor inertia by \(N^2\), turning rapid motor
  acceleration into violent mechanical shock unless filtered with
  \(C^2\) continuous splines.
\item
  \textbf{Thermal Heat Accumulates Silently:} Joulean dissipation
  (\(I^2 R\)) creates positive thermal feedback
  (\(\alpha_{\text{Cu}}\)). Enforce cyclic duty cycles (\(D\)) and
  real-time \(I^2 t\) thermal observers.
\item
  \textbf{Electrical Power is a Shared Resource:} Simultaneous compute
  and motor acceleration transients collapse DC bus voltages, triggering
  brownout resets at the moment of peak kinematic momentum.
\item
  \textbf{Tail Distributions Rule Physical Safety:} Physical safety is
  governed by tail percentiles (\(P_{99.9}\)) and zero-failure
  statistical bounds (\(p_{\max} \approx 3/n\)), not by median bench
  performance.
\end{itemize}

\end{fbxSimple}

Chapter~\ref{sec-brain} takes custody of this constrained physical body
and addresses the machine from the inside out. With physical envelopes
established, the engineering objective is not to modify the physical
laws that govern the machine, but to build a nervous system and a
learned brain that respect them under all operating conditions. A neural
network policy possesses high representational capacity, but it has no
innate awareness of winding temperatures, bus saturation, or stopping
distances. The nervous system must enforce these five limits
deterministically, providing the safety and timing guarantees that the
learned model cannot provide for itself. The chapters ahead examine what
a learned component contributes to physical autonomy once its authority
is bounded within these measured physical limits.

\chapter{\texorpdfstring{The Cognitive
Brain\chaptersubtitle{The Cognitive Mental Model and Information Flow}}{The Cognitive Brain}}\label{sec-brain}

\begin{figure}[H]

\centering{

\includegraphics[width=1\linewidth,height=\textheight,keepaspectratio]{index_files/mediabag//Users/VJ/GitHub/MLSysBook-integrate-physical/physical/book/chapters/03-brain/figures/fig_locator.pdf}

}

\caption{\label{fig-locator}\emph{Systems organ focus: The brain.} The
four-tier physical AI systems architecture highlights the cognitive
deliberative layer. Operating above the real-time nervous system and
physical body, the learned brain processes high-dimensional multimodal
streams to generate candidate proposals, with execution authority
reserved for downstream deterministic safety governors.}

\end{figure}%

Every embodied agent requires a cognitive core---a brain that perceives
physical reality, maintains spatial memory, and synthesizes candidate
motion trajectories. In classical robotics, this layer was authored
through explicit differential equations and state machines, which break
down in unstructured environments. Physical AI introduces high-capacity
foundation models to generalize across open worlds. Yet this statistical
power introduces critical vulnerabilities: neural models forfeit formal
verifiability, output uncalibrated confidence, and collide with the
physical Memory Wall during real-time weight streaming. To engineer
dependable machines, we must treat the learned brain as an unprivileged
proposal engine---harnessing its generalizability while delegating all
execution authority to deterministic safety governors.

\newpage{}

\protect\phantomsection\label{callout-objectiveux2a-3.1}
\begin{fbxSimple}{callout-objective}{Learning Objectives}{}
\phantomsection\label{callout-objective*-3.1}

After studying this chapter, you will be able to:

\begin{enumerate}
\def\labelenumi{\arabic{enumi}.}
\tightlist
\item
  Identify the five ideal properties required of an embodied cognitive
  brain and trace its role in the four-tier physical AI stack.
\item
  Explain the historical and practical motivations behind the shift from
  classical analytical models to learned representations in physical
  systems.
\item
  Contrast classical analytical methods with learned foundation models
  across generalizability, sensory robustness, and formal mathematical
  verifiability.
\item
  Trace the five foundational stages of the learned cognitive pipeline:
  Ingestion, Perception, Memory, Intent, and Planning.
\item
  Ground each cognitive stage in physical robot archetypes (mobile
  systems, manipulators, process plants) and evaluate their systems
  engineering implications.
\item
  Diagnose the four core cognitive failure modes: semantic
  hallucinations, out-of-distribution overconfidence, perceptual
  chattering, and temporal expiration.
\item
  Apply the Proposal--Permission Principle and Two-Speed Brain
  architecture to decouple statistical deliberation from real-time
  safety enforcement.
\end{enumerate}

\emph{Core chapter takeaway.} Learned models provide open-world
generalizability and sensory robustness that classical equations cannot
achieve, but they lack formal safety guarantees. In physical AI, the
brain earns the capability to propose what the machine might attempt,
but never the authority to grant permission for physical action.

\end{fbxSimple}

\section{The Embodied Brain}\label{sec-brain-3-1}

\lettrine{A}{s}\lettrine{h}{ighlighted} in the systems organ locator at
the opening of this chapter (Figure~\ref{fig-locator}), the brain
occupies the topmost deliberative tier in the physical AI architecture.
In Chapter~\ref{sec-boundary}, we established that physical AI begins
the instant digital computation crosses the causal boundary to command
physical mass, while Chapter~\ref{sec-body} quantified the unyielding
budgets of freshness, stopping distance, reflected inertia, and thermal
dissipation that govern the physical body. Operating above the real-time
nervous system and the dynamical body, the brain serves as the cognitive
command center: it is the organ responsible for making sense of incoming
sensory evidence and deciding what the physical body should attempt to
do next.

To understand the engineering requirements of this cognitive layer, we
must first contrast an embodied brain with a disembodied datacenter
model. When a large language model in the cloud generates text, a
hallucinated token carries no physical inertia; a user simply re-prompts
the model, and the only cost incurred is digital compute. When an
embodied machine operates in the physical world, its cognitive decisions
govern moving mass that carries kinetic energy
(\(E_k = \frac{1}{2} m v^2\)). In the physical domain, a cognitive error
is not a software exception that can be caught in a \texttt{try-catch}
block; it is a bent manipulator arm, a broken conveyor fixture, or a
catastrophic collision with a human coworker.

To govern an embodied machine effectively in the open world, an embodied
brain must satisfy five ideal properties:

First, it must possess open-world semantic generalizability. Rather than
being confined to a narrow, pre-programmed lookup table of known
objects, the brain must interpret open-vocabulary human language and
generalize zero-shot across novel shapes, textures, and spatial
arrangements.

Second, it must exhibit continuous sensory robustness. Real-world
physical sensors operate under shifting sunlight, specular reflections,
dynamic shadows, and camera lens dust. The cognitive architecture must
smoothly absorb high-dimensional sensory noise without suffering
discrete, brittle failures.

Third, it must maintain metric spatial-temporal grounding. Disembodied
language models reason over abstract semantic concepts, but an embodied
machine moves through three-dimensional Euclidean space. The brain must
anchor its semantic concepts in metric coordinates, track objects
through temporary visual occlusions, and model the continuous velocities
of dynamic bodies. The metric pose of an object or robot link is
described by the Special Euclidean Lie group
\(SE(3) = \mathbb{R}^3 \rtimes SO(3)\), parameterized by a translation
vector \(\mathbf{p} \in \mathbb{R}^3\) and an orthonormal rotation
matrix \(\mathbf{R} \in SO(3)\).

Fourth, it must achieve temporal multi-rate cadence decoupling.
High-level reasoning and task planning are computationally demanding,
requiring multi-billion-parameter models that deliberate at slow, human
cadences (\(1\text{--}5\text{ Hz}\)). Yet controlling physical mass
requires continuous, reactive trajectory setpoints at
\(20\text{--}50\text{ Hz}\) to prevent dynamic drift. The cognitive
architecture must cleanly bridge these disparate operational cadences
without stalling the actuators.

Fifth, it must operate under bounded execution authority. Because
statistical models are inherently uncalibrated under out-of-distribution
shifts and prone to semantic hallucinations, the brain cannot hold
direct authority over physical actuators. It must emit candidate
proposals that remain strictly containable by downstream deterministic
safety monitors.

The fundamental engineering question is: \emph{what architectural
paradigm can deliver these five ideal properties?}

\section{The Classical Wall}\label{sec-brain-3-2}

For nearly a century, automated control and robotics approached
cognition as an explicit geometric and analytical physics problem
(\citeproc{ref-wiener1948cybernetics}{Wiener 1948};
\citeproc{ref-kalman1960new}{Kalman 1960};
\citeproc{ref-craig2005introduction}{Craig 2005}). Under this classical
paradigm, engineers constructed mathematical representations from first
principles.

Analytical kinematics and differential dynamics mapped desired
end-effector motions into actuator torques using closed-form equations
(\(\mathbf{M}(\mathbf{q})\ddot{\mathbf{q}} + \mathbf{C}(\mathbf{q}, \dot{\mathbf{q}})\dot{\mathbf{q}} + \mathbf{g}(\mathbf{q}) = \boldsymbol{\tau}\)).
Perception relied on matching hand-crafted 3D Computer-Aided Design
(CAD) models against sensor point clouds using iterative closest point
(ICP) or sample consensus (RANSAC) algorithms. High-level task logic was
authored as explicit boolean state machines that dictated deterministic
transitions between operational phases, such as approaching, aligning,
grasping, and retracting.

\begin{figure}[t]

\centering{

\includegraphics[width=0.85\linewidth,height=\textheight,keepaspectratio]{/Users/VJ/GitHub/MLSysBook-integrate-physical/physical/book/chapters/03-brain/figures/fig03_real_industrial_welding.jpg}

}

\caption{\label{fig-industrial-welding}\emph{Structured classical
robotics regime (automotive body shop).} Multiple 6-DoF articulated
industrial robots execute pre-programmed spot welding paths on stamped
automotive chassis at the BMW Leipzig plant. Operating in a strictly
calibrated environment with sub-millimeter repeatable fixtures and
physical interlock safety cages, classical analytical controllers
provide deterministic cycle timing and formal stability proofs. (Credit:
BMW Group / Wikimedia Commons, CC BY-SA 2.0 de).}

\end{figure}%

Where the operational environment is strictly structured and bounded
(Figure~\ref{fig-industrial-welding}), classical methods represent the
pinnacle of dependable engineering. On an automated automotive assembly
line, an industrial robot welds stamped steel panels where every
workpiece arrives at an exact coordinate within sub-millimeter
tolerances. Ambient lighting is tightly regulated by industrial
fixtures, and human workers are excluded by physical interlock safety
fences. In this controlled regime, classical algorithms provide complete
mathematical transparency, microsecond execution timing, and provable
stability certificates. An engineer can formally prove Lyapunov
stability, guarantee that trajectories stay within joint torque bounds,
and achieve structural code coverage under ISO 26262 without collecting
a single demonstration.

Yet the moment an embodied machine leaves the structured factory cage to
operate in homes, warehouses, construction sites, or agricultural
fields, this classical paradigm hits what roboticists describe as the
Classical Wall (\citeproc{ref-brooks1991intelligence}{Brooks 1991};
\citeproc{ref-levine2018learning}{Levine et al. 2018}). The physical
world is not a sterile CAD drawing; it is continuous, deformable,
uncataloged, and endlessly variable.

The first breakdown occurs in the combinatorial explosion of geometry.
In a commercial logistics warehouse handling hundreds of thousands of
retail goods, authoring CAD models, mesh templates, and grasp points for
every bottle, wrapped package, and flexible garment is impossible.
Discretizing orientations, frictional properties, and contact
permutations across overlapping, deformable items creates a
combinatorial space exceeding billions of distinct geometric
configurations that cannot be captured in hand-crafted templates.

The second breakdown arises from sensory brittleness. Classical computer
vision relies on hand-crafted feature extractors, such as edge filters,
Harris corners, or surface normal thresholds. When ambient sunlight
shifts, water puddles create specular reflections, or an object is
partially obscured by dust, discrete feature detectors fracture. A
single missing edge causes a planar segmentation algorithm to bridge
across two distinct physical objects, commanding a gripper to pinch an
empty void.

\begin{figure}[t]

\centering{

\includegraphics[width=0.85\linewidth,height=\textheight,keepaspectratio]{/Users/VJ/GitHub/MLSysBook-integrate-physical/physical/book/chapters/03-brain/figures/fig03_real_unstructured_rescue.jpg}

}

\caption{\label{fig-unstructured-rescue}\emph{Unstructured open-world
robotics regime (urban search and rescue).} A mobile rescue robot
navigates an obstacle course with irregular inclines, loose wooden
debris, and unpredictable visual occlusions. In such open-ended
environments, pre-authored 3D CAD models and deterministic state
machines shatter, requiring physical AI systems to infer spatial
affordances and terrain traversability directly from multimodal sensory
streams. (Credit: TU Darmstadt Rescue Robotics / Wikimedia Commons, CC
BY-SA 3.0).}

\end{figure}%

The third breakdown is the engineering bottleneck of hand-authored
logic. In classical systems, the machine's operational capability is
bounded by the human engineer's ability to anticipate every physical
edge case in advance. When an unexpected variation occurs---such as a
bent bracket, a scuffed label, or an unfamiliar obstacle---the
hand-authored rulebook contains no matching branch, causing the machine
to halt, drop its payload, or fail catastrophically.

\section{Learned Representations}\label{sec-brain-3-3}

Physical AI arose to dismantle the classical wall by replacing rigid
human-authored heuristics with deep neural representations fitted
directly to multimodal sensory data. In the embodied AI and robotics
community, this transition was driven by four foundational pillars:

The first pillar is open-world semantic generalizability. Classical
controllers are inherently narrow: an inverse kinematics solver can move
an end-effector to coordinate \((x, y, z)\), but it possesses zero
semantic understanding of what resides at that location or how
surrounding objects relate to human tasks. Modern vision-language-action
foundation models project open-vocabulary natural language instructions
and high-resolution camera streams into shared semantic embeddings
(\citeproc{ref-driess2023palme}{{Driess et al.} 2023};
\citeproc{ref-brohan2023rt1}{{Brohan, Brown, Carbajal, Chebotar, Dabis,
et al.} 2023}; \citeproc{ref-brohan2023rt2}{{Brohan, Brown, Carbajal,
Chebotar, Chen, et al.} 2023}; \citeproc{ref-kim2024openvla}{{Kim et
al.} 2024}). When commanded to ``wipe the coffee spill from the counter
and stow the mug in the drying rack,'' a learned model parses the
unstructured scene into functional affordances, identifying the
graspable handle of an unfamiliar mug and the planar boundaries of the
liquid stain without requiring an engineer to provide a 3D CAD model or
write a single line of procedural code.

The second pillar is continuous sensory robustness. While classical
algorithms suffer from boolean brittleness---where a minor shift in
lighting or surface texture flips a discrete classification
branch---deep neural networks project high-dimensional sensory inputs
onto smooth, continuous latent manifolds. Variations in ambient
illumination, background clutter, shadow transitions, and camera
perspective produce bounded displacements in latent space rather than
catastrophic categorical jumps.

The third pillar is the ability to handle complex multi-contact and
deformable dynamics. Many physical manipulation tasks---such as folding
laundry, routing flexible wire harnesses through car doors, or scooping
granular media---involve non-smooth contact mechanics, variable
friction, and infinite-dimensional deformations that cannot be modeled
in closed-form differential equations. By learning directly from
physical teleoperation demonstrations, neural policies acquire intuitive
physical behaviors that naturally exploit compliance and surface
contacts.

\begin{figure}[t]

\centering{

\includegraphics[width=0.85\linewidth,height=\textheight,keepaspectratio]{/Users/VJ/GitHub/MLSysBook-integrate-physical/physical/book/chapters/03-brain/figures/fig03_real_aloha_manipulator.png}

}

\caption{\label{fig-aloha-manipulator}\emph{Authentic bimanual robotic
policy testbed.} A canonical dual-arm bimanual teleoperation workcell
featuring two 6-DoF follower arms, two leader arms for kinesthetic
teleoperation, multi-view RGB cameras (overhead and wrist-mounted
RealSense modules), and mechanical parallel grippers mounted to an
extruded aluminum structural frame. The setup demonstrates how physical
AI policies ingest synchronized high-dimensional multi-camera streams to
generate continuous end-effector trajectories, requiring deterministic
safety governors and reachability filters to prevent self-collisions and
motor over-torquing during closed-loop execution. (Credit: Stanford
University \& Google DeepMind, CC-BY / Apache 2.0).}

\end{figure}%

The fourth pillar is empirical scaling with data and compute. As
articulated by Richard Sutton in \emph{The Bitter Lesson}
(\citeproc{ref-sutton2019bitter}{Sutton 2019}), methods that leverage
general computation and large-scale learning consistently surpass
hand-engineered domain heuristics over time. By training high-capacity
transformer architectures across millions of diverse robotic
teleoperation demonstrations (Figure~\ref{fig-aloha-manipulator}) and
web-scale multimodal datasets (\citeproc{ref-zhao2023learning}{Zhao et
al. 2023}; \citeproc{ref-chi2023diffusionpolicy}{Chi et al. 2025}),
physical AI systems steadily improve in dexterity and resilience without
requiring engineers to manually re-author control logic for each new
task.

\begin{figure}[t]

\centering{

\includegraphics[width=0.95\linewidth,height=\textheight,keepaspectratio]{index_files/mediabag//Users/VJ/GitHub/MLSysBook-integrate-physical/physical/book/chapters/03-brain/figures/fig03_vla_architecture.pdf}

}

\caption{\label{fig-vla-architecture}\emph{Vision-Language-Action (VLA)
architecture and unprivileged proposal pipeline.} (1) Multimodal
ingestion tokenizes multi-view camera streams
(\(224 \times 224 \times 3\) RGB), natural language task instructions,
and proprioceptive joint states (\(\mathbf{q}, \dot{\mathbf{q}}\)). (2)
A high-capacity transformer backbone (\(7\text{B--}55\text{B}\)
parameters) fuses spatial and semantic tokens across self-attention
layers, operating on embedded edge silicon (such as a canonical
heterogeneous application MPU with dedicated tensor coprocessors). (3)
The policy emits unprivileged candidate action proposals via discrete
token binning (single-step waypoints) or diffusion action chunking
(continuous \(H\)-step trajectory splines). (4) Downstream real-time
safety monitors on isolated microcontrollers evaluate physical
invariants at \(1000\text{ Hz}\) before granting actuator execution
permission.}

\end{figure}%

As structured in the modern Vision-Language-Action (VLA) paradigm
(Figure~\ref{fig-vla-architecture}), an embodied foundation model is not
an opaque monolith, but a structured four-stage neural pipeline that
transforms sensory fluxes into motor candidates:

\begin{enumerate}
\def\labelenumi{\arabic{enumi}.}
\tightlist
\item
  \textbf{Spatial \& Semantic Tokenization (Information from the Left):}
  Raw camera frames (\(224 \times 224 \times 3\)) enter a spatial vision
  backbone (such as DINOv2 or SigLIP), which partitions the image into a
  \(14 \times 14\) grid of \(16 \times 16\) pixel patches. Each patch is
  flattened and linearly projected into a \(D\)-dimensional visual
  embedding \(\mathbf{z}_{\text{vision}}\). Concurrently, natural
  language task instructions are parsed into subword tokens
  \([\mathbf{t}_1, \dots, \mathbf{t}_L]\), while continuous
  proprioceptive joint encoders (\(\mathbf{q}, \dot{\mathbf{q}}\)) are
  projected into matching state tokens \(\mathbf{z}_{\text{state}}\).
\item
  \textbf{Multimodal Alignment \& Fusion (The Backbone):} The visual,
  linguistic, and physical state tokens are concatenated into a unified
  sequence
  \([\mathbf{z}_{\text{vision}} \,|\, \mathbf{z}_{\text{text}} \,|\, \mathbf{z}_{\text{state}}]\)
  and processed through 32--80 transformer self-attention layers.
  Cross-attention mechanisms bind linguistic intent (``grasp the
  screwdriver'') directly to the spatial coordinates of visual patches.
\item
  \textbf{Action Decoding (Information to the Right):} Instead of
  generating autoregressive text words, the final hidden activation
  \(\mathbf{h}_t\) feeds an \textbf{Unprivileged Action Decoder Head}.
  This head projects the latent state into continuous motor action
  candidates---either single-step delta poses or multi-step continuous
  trajectory chunks
  \(\hat{\mathbf{a}}_{t:t+H} \in \mathbb{R}^{H \times d_a}\) (comprising
  3D position deltas \(\Delta \mathbf{p}\), rotation increments
  \(\Delta \boldsymbol{\theta}\), and gripper actuation flags).
\item
  \textbf{The Privilege Boundary (Isolation from Motors):} Crucially,
  the model's output is an unverified candidate proposal (\(p_t\)). It
  carries zero direct hardware permission to energize motor coils until
  downstream real-time enforcers on isolated silicon verify kinematic
  and kinetic invariants at \(1000\text{ Hz}\).\sidenote{\footnotesize Unprivileged
    proposal isolates actuators behind deterministic Control Barrier
    Functions.}
\end{enumerate}

\section{The Cognitive Pipeline}\label{sec-brain-3-4}

To architect a learned cognitive pipeline, we must trace how information
flows from raw physical sensations to actionable trajectory setpoints.
Rather than treating the brain as an unstructured end-to-end black box,
a dependable physical AI architecture organizes cognition into five
sequential stages (Figure~\ref{fig-cognitive-pipeline} and
\textbf{?@fig-five-handoffs}): Ingestion, Perception, Memory, Intent,
and Planning.

\begin{figure}[t]

\centering{

\includegraphics[width=0.95\linewidth,height=\textheight,keepaspectratio]{index_files/mediabag//Users/VJ/GitHub/MLSysBook-integrate-physical/physical/book/chapters/03-brain/figures/fig03_cognitive_pipeline_horizontal.pdf}

}

\caption{\label{fig-cognitive-pipeline}\emph{The cognitive flow of
physical AI: End-to-end left-to-right pipeline.} Information flows
sequentially from physical sensors through Ingestion (tokenization),
Perception (spatial geometry \& affordances), Memory (persistent world
model belief), Intent (deliberative goal leases), and Planning
(multi-step action chunking), crossing the Proposal--Permission boundary
to downstream real-time reflexes.}

\end{figure}%

\subsection*{Ingestion and
Representation}\label{ingestion-and-representation}
\addcontentsline{toc}{subsection}{Ingestion and Representation}

Cognition begins at the physical boundary where environmental energy is
transduced into digital bitstreams. Photons strike active-pixel CMOS
sensors across 2D camera grids, infrared pulses measure photon
time-of-flight in depth cameras, and micro-electromechanical systems
(MEMS) in inertial measurement units (IMUs) sample accelerations and
angular velocities.

The ingestion stage translates these heterogeneous, high-bandwidth
bitstreams into normalized numerical embeddings. Modern embodied
foundation models utilize Vision Transformers (ViTs) to process spatial
visual inputs (\citeproc{ref-dosovitskiy2020image}{{Dosovitskiy et al.}
2020}). A raw image frame
\(\mathbf{I} \in \mathbb{R}^{H \times W \times C}\) is partitioned into
a sequence of non-overlapping 2D patches
\(\mathbf{x}_p \in \mathbb{R}^{N \times (P^2 \cdot C)}\), where
\(P \times P\) is the patch resolution (typically \(14 \times 14\) or
\(16 \times 16\) pixels) and \(N = (HW)/P^2\) is the resulting token
sequence length. Each flattened patch is linearly projected into a
continuous latent space of dimension \(D\), concatenated with learned
positional embeddings to preserve spatial geometry, and interleaved with
tokenized natural language instructions. Continuous proprioceptive state
vectors (such as joint positions \(\mathbf{q}\), velocities
\(\dot{\mathbf{q}}\), and gripper opening widths) are similarly
normalized and projected into matching token spaces, producing a unified
multimodal token sequence that serves as the common input for downstream
reasoning.\sidenote{\footnotesize \textbf{Curriculum Architecture.} While this chapter
  establishes the overarching mental model of embodied cognition, deep
  algorithmic formulations for each stage are detailed in Part III
  (Chapter~\ref{sec-perception} through Chapter~\ref{sec-enforcement}).}

Consider how this operates across the three machine archetypes:

\begin{itemize}
\tightlist
\item
  \emph{Mobile systems:} A quadruped traversing rubble synchronizes four
  RGB-D cameras streaming at \(30\text{ fps}\) (\(>740\text{ MB/s}\) raw
  data) alongside 12 joint encoders sampled at \(1000\text{ Hz}\).
  Ingestion tokenizes each visual frame into hundreds of patch
  embeddings while downsampling and embedding the high-frequency joint
  state stream into a unified temporal sequence.
\item
  \emph{Manipulators:} A dual-arm workstation ingests synchronized
  stereo wrist camera feeds (\(60\text{ fps}\)) and 14-DoF joint
  positions, tokenizing visual patches and motor states into aligned
  input tokens for end-effector coordination.
\item
  \emph{Process / energy systems:} A continuous chemical reactor ingests
  thousands of analog thermocouple, pressure transducer, and mass flow
  meter channels at \(100\text{ Hz}\), normalizing heterogeneous sensor
  calibrations into a rolling multivariate state tensor.
\end{itemize}

This ingestion pipeline places substantial demands on data bandwidth and
memory throughput. Streaming multiple high-definition camera feeds
generates gigabytes of raw pixel data every second. Systems engineers
must implement efficient frame buffering, sensor exposure
synchronization, and spatial patch pooling to keep data moving smoothly
before cognitive reasoning begins.

\subsection*{Perception and Spatial
Grounding}\label{perception-and-spatial-grounding}
\addcontentsline{toc}{subsection}{Perception and Spatial Grounding}

The perception stage extracts metric structure, semantic identity, and
spatial affordances from the ingested tokens. Rather than treating
sensory data as an unstructured bag of visual features, perception
grounds representations in the 3D Euclidean space (\(SE(3)\)) of the
physical workspace.

Using dense self-supervised representations such as DINOv2
(\citeproc{ref-oquab2023dinov2}{{Oquab et al.} 2023}) alongside learned
spatial decoders, the perception layer isolates distinct physical
entities, segments object boundaries, estimates metric 3D bounding
boxes, and reconstructs surface normal vectors. Crucially, perception
extracts spatial affordances: functional geometric fields that describe
where and how an object can be physically engaged. For a drill,
perception identifies not merely a generic class label, but the 3D
position and orientation of the handle grasp axis, the trigger contact
surface, and the drill bit tip. This transforms ungrounded pixel arrays
into actionable spatial entities referenced to the robot's base
coordinate frame.

When a mobile manipulator is commanded to retrieve a torque wrench from
a cluttered workbench, perception isolates the tool from surrounding
background clutter, predicts its 6-DoF pose (\(SE(3)\) transformation
matrix \(\mathbf{T}_{\text{base}}^{\text{tool}}\)), and outputs a
continuous heatmap over the handle indicating valid gripper pinch
contacts.

In practice, perception is highly vulnerable to sensor calibration
drift. If camera mounting screws loosen by half a millimeter or
structural arm links deflect under payload weight, the transformation
matrix between the camera optical center and the robot base frame
drifts. This introduces metric scale errors that cause precision
grippers to miss targets or strike work surfaces. Systems engineers must
implement online hand-eye recalibration routines and maintain spatial
uncertainty covariances alongside raw point estimates.

\subsection*{Memory and Spatial-Temporal
Belief}\label{memory-and-spatial-temporal-belief}
\addcontentsline{toc}{subsection}{Memory and Spatial-Temporal Belief}

A physical agent cannot act reliably from a single instantaneous
snapshot of the world. Objects pass behind visual occlusions, cameras
pan away from manipulated targets, and physical bodies possess dynamic
momentum that cannot be inferred from a single frame.

The memory stage maintains a persistent, rolling spatial-temporal belief
of the environment. It tracks object identities across time, correlates
visual features with geometric coordinate frames, and updates 3D
volumetric representations such as Truncated Signed Distance Fields
(TSDFs), Language Embedded Radiance Fields (LERF)
(\citeproc{ref-kerr2023lerf}{Kerr et al. 2023}), or 3D Gaussian
Splatting maps. By fusing historical observations across a sliding
temporal window, memory maintains accurate state estimates even during
complete visual occlusion, computes linear and angular object
velocities, and tracks operational context across multi-step tasks.

For a quadruped traversing a rocky ledge, the forward-facing camera
points toward the horizon, losing sight of the immediate ground beneath
its front feet. The memory stage preserves the reconstructed 3D surface
geometry of the terrain, enabling the robot to securely place its hind
feet onto footholds it observed several seconds earlier.

Maintaining dense volumetric belief states introduces significant memory
footprints. Storing high-resolution 3D voxel grids or radiance field
parameters requires substantial storage capacity, competing directly
with the memory needed to run large neural network models. Furthermore,
belief representations degrade in freshness over time: if a dynamic
obstacle moves while occluded, an un-decayed memory belief causes the
robot to plan around a ``phantom'' obstacle or collide with an
unobserved intruder. Systems engineers must enforce explicit temporal
decay rates on spatial belief states to reflect physical entropy.

\subsection*{Intent and Deliberative
Reasoning}\label{intent-and-deliberative-reasoning}
\addcontentsline{toc}{subsection}{Intent and Deliberative Reasoning}

Given a persistent belief of the world and a high-level mission prompt,
the intent stage performs semantic reasoning and task decomposition.
Because open-vocabulary foundation models operate over vast parameter
spaces (typically 7B to 70B parameters), deliberative reasoning occurs
at slow, human-scale cadences (\(1\text{--}5\text{ Hz}\)).

To understand how intent operates, imagine a human operator issues a
natural language mission prompt:

\[\text{"Clear the table by placing the ceramic mug in the sink and throwing the paper napkin into the trash bin."}\]

The intent layer cannot execute this directive in a single leap. Using
embodied reasoning architectures such as SayCan
(\citeproc{ref-ahn2022saycan}{{Ahn et al.} 2023}), VoxPoser
(\citeproc{ref-huang2023voxposer}{Huang et al. 2023}), or Code as
Policies (\citeproc{ref-liang2023code}{{Liang et al.} 2023}), the
deliberative model queries the active world memory and decomposes the
mission into an ordered sequence of discrete subgoals: first approach
the mug, grasp the handle, transport it to the sink coordinate frame,
and release contact, before turning toward the napkin.

Rather than commanding actuator voltages directly, the intent stage
packages each subgoal into a formal, cryptographically bounded
\textbf{Intent Lease} (\(\mathcal{L}_{\text{intent}}\)). An intent lease
specifies the active target entity, defines a bounded spatial corridor
outside which motion is unauthorized, and establishes an explicit
expiration timestamp (\(t_{\text{expire}}\)).

\begin{longtable}[]{@{}
  >{\raggedright\arraybackslash}p{(\linewidth - 6\tabcolsep) * \real{0.2500}}
  >{\raggedright\arraybackslash}p{(\linewidth - 6\tabcolsep) * \real{0.2500}}
  >{\raggedright\arraybackslash}p{(\linewidth - 6\tabcolsep) * \real{0.2500}}
  >{\raggedright\arraybackslash}p{(\linewidth - 6\tabcolsep) * \real{0.2500}}@{}}
\caption{Intent Lease Contract Schema
(\(\mathcal{L}_{\text{intent}}\)).}\label{tbl-03-intent-lease-schema}\tabularnewline
\toprule\noalign{}
\begin{minipage}[b]{\linewidth}\raggedright
Field Name
\end{minipage} & \begin{minipage}[b]{\linewidth}\raggedright
Type
\end{minipage} & \begin{minipage}[b]{\linewidth}\raggedright
Unit
\end{minipage} & \begin{minipage}[b]{\linewidth}\raggedright
Description \& Structural Invariants
\end{minipage} \\
\midrule\noalign{}
\endfirsthead
\toprule\noalign{}
\begin{minipage}[b]{\linewidth}\raggedright
Field Name
\end{minipage} & \begin{minipage}[b]{\linewidth}\raggedright
Type
\end{minipage} & \begin{minipage}[b]{\linewidth}\raggedright
Unit
\end{minipage} & \begin{minipage}[b]{\linewidth}\raggedright
Description \& Structural Invariants
\end{minipage} \\
\midrule\noalign{}
\endhead
\bottomrule\noalign{}
\endlastfoot
\texttt{lease\_id} & \texttt{uint64\_t} & --- & Monotonically increasing
64-bit sequence identifier \\
\texttt{target\_entity\_id} & \texttt{uint32\_t} & --- & Hash of target
entity in persistent spatial memory \\
\texttt{pose\_target\_se3} & \texttt{float32{[}4{]}{[}4{]}} &
\(\text{m}, \text{rad}\) & Valid \(\text{SE}(3)\) pose matrix:
\(\mathbf{R} \in \text{SO}(3), \det(\mathbf{R})=1, \mathbf{p} \in \mathbb{R}^3\) \\
\texttt{spatial\_corridor\_box} & \texttt{float32{[}6{]}} & \(\text{m}\)
& Authorized 3D bounding box
\([\mathbf{p}_{\min}, \mathbf{p}_{\max}] \in \mathbb{R}^6\) \\
\texttt{t\_generated\_ns} & \texttt{uint64\_t} & \(\text{ns}\) & PTP
monotonic hardware timestamp at deliberation completion \\
\texttt{t\_expire\_ns} & \texttt{uint64\_t} & \(\text{ns}\) & Hard
expiration timestamp; enforces
\(t_{\text{expire}} - t_{\text{gen}} \le \tau_{\text{valid}}\) \\
\texttt{max\_permissible\_jerk} & \texttt{float32} & \(\text{m/s}^3\) &
Maximum allowable end-effector / joint jerk magnitude \\
\texttt{hmac\_signature} & \texttt{uint8\_t{[}32{]}} & bytes & SHA-256
deliberation cryptographic contract signature \\
\end{longtable}

This deliberative cadence creates an inherent freshness gap. If a large
vision-language foundation model requires \(400\text{ ms}\) to generate
an intent lease, the physical world has moved during that inference
window. If the target mug was bumped or a coworker walked between the
robot and the sink, the intent lease is partially stale before it
reaches the motion planner. Systems engineers must bound intent lease
lifetimes (\(\tau_{\text{valid}}\)) and enforce strict timeout rejection
gates at execution time.

\subsection*{Planning and Action
Chunking}\label{planning-and-action-chunking}
\addcontentsline{toc}{subsection}{Planning and Action Chunking}

An intent lease specifies what operational goal to pursue, but it
provides no continuous motor trajectories and no guarantee that
intervening motion satisfies physical kinematic limits. The planning
stage converts the active intent lease into candidate joint trajectories
or end-effector motion setpoints via an \textbf{Unprivileged Action
Decoder Head}.

To understand how planning is executed on physical silicon, we must
treat deep neural policies not as abstract mathematical functions, but
as physical compute workloads running on embedded accelerators---such as
a canonical high-performance heterogeneous edge System-on-Chip (SoC)
with integrated GPU tensor cores and dedicated deep learning
accelerators---bounded by memory bus bandwidth, silicon thermals, and
electrical power.

Consider what happens if a robot attempts to generate continuous motor
actions one discrete step at a time through an autoregressive foundation
model (such as a 7-billion parameter Vision-Language-Action network). In
autoregressive decoding, generating each individual action token
requires the processor to stream every parameter of the network from
external Dynamic Random-Access Memory (DRAM) across silicon bus traces
into on-chip cache and vector registers.

For a 7B parameter model stored in 16-bit floating-point precision
(FP16), the weight matrix occupies:

\[M_{\text{weights}} = 7 \times 10^9 \text{ parameters} \times 2\text{ bytes/parameter} = 14.0\text{ GB}\]

On an embedded edge compute module equipped with a 256-bit wide LPDDR5
memory interface delivering a theoretical peak bandwidth of
\(B_{\text{mem}} = 204.8\text{ GB/s}\), the minimum physical time
required simply to sweep those weights across the bus once
is:\sidenote{\footnotesize \textbf{LPDDR5 Memory Subsystem Throughput.} The
  \(204.8\text{ GB/s}\) peak bandwidth on a canonical heterogeneous SoC
  (256-bit LPDDR5 at \(6400\text{ MT/s}\)) drops to \(70\text{--}75\%\)
  sustained efficiency (\(145\text{ GB/s}\)) due to DRAM bank turnaround
  and auto-refresh (\(t_{\text{RFC}}\)).}

\[t_{\text{stream}} = \frac{M_{\text{weights}}}{B_{\text{mem}}} = \frac{14.0\text{ GB}}{204.8\text{ GB/s}} \approx 68.4\text{ ms}\]

This physical reality imposes an unyielding frequency ceiling: even if
arithmetic computation consumed zero nanoseconds, the memory bus alone
caps single-step token generation at an absolute maximum of
\(f_{\text{max}} = 1 / 0.0684\text{ s} \approx 14.6\text{ Hz}\). When we
account for sensory attention ingestion, KV-cache lookup, and operating
system scheduling, live single-token inference drops to
\(1\text{--}5\text{ Hz}\). If an embodied manipulator attempted to
control moving mass at \(5\text{ Hz}\), the actuators would freeze and
stutter between consecutive ticks like a buffering video stream, causing
dynamic drift, severe contact instability, and dropped payloads.

We can analyze this bottleneck formally through the \textbf{Roofline
Model} (\citeproc{ref-williams2009roofline}{Williams et al.
2009}),\sidenote{\footnotesize \textbf{For AI \& Theory Readers:} The Patterson
  Roofline proves why scaling raw FLOP capacity on edge SoCs cannot
  increase inference frequency if the memory bus bandwidth
  \(B_{\text{mem}}\) remains the operational bottleneck.} which bounds
attainable compute performance \(P\) (in FLOP/s) by the processor's peak
arithmetic throughput \(P_{\text{peak}}\) and its memory bandwidth
\(B_{\text{mem}}\):

\[P = \min\left(P_{\text{peak}}, \; I \cdot B_{\text{mem}}\right)\]

where \(I\) is the operational arithmetic intensity defined as:

\[I = \frac{\text{Floating-Point Operations (FLOPs)}}{\text{Memory Data Transferred (Bytes)}}\]

In single-step autoregressive generation, reading a 2-byte weight to
perform a single multiply-accumulate operation (2 FLOPs) yields an
arithmetic intensity of:

\[I_{\text{single}} \approx \frac{2\text{ FLOPs}}{2\text{ bytes}} = 1.0\text{ FLOP/byte}\]

Because this intensity is orders of magnitude below the accelerator's
saturation knee
(\(I_{\text{knee}} = P_{\text{peak}} / B_{\text{mem}} \approx 1340\text{ FLOP/byte}\)
for a canonical \(275\text{ TFLOP/s}\) heterogeneous SoC), single-step
inference is trapped deep in the \textbf{memory-bandwidth-bound
regime}.\sidenote{\footnotesize \textbf{Roofline Hardware Ridge Point.} With
  \(P_{\text{peak}} = 275\text{ TFLOP/s}\) and
  \(B_{\text{mem}} = 204.8\text{ GB/s}\), the ridge point is
  \(I_{\text{knee}} = 1342.8\text{ FLOP/byte}\). Single-token generation
  (\(I \approx 1\)) achieves \(<0.1\%\) peak compute efficiency
  (\(99.9\%\) idle time stalled on DRAM).} The tensor accelerator
achieves less than \(0.1\%\) of its peak compute throughput, spending
over \(99\%\) of its time stalled waiting for DRAM memory transfers.

\protect\phantomsection\label{callout-definitionux2a-3.2}
\begin{fbxSimple}{callout-definition}{Definition:}{The Memory Wall Ceiling}
\phantomsection\label{callout-definition*-3.2}

The \textbf{Memory Wall Ceiling} is the absolute upper bound on model
evaluation frequency
(\(f_{\max} = B_{\text{mem}} / M_{\text{weights}}\)) imposed by external
DRAM bus bandwidth when streaming non-resident neural network weights
into on-chip vector registers. In single-step autoregressive inference,
operational arithmetic intensity collapses to
\(I \approx 1.0\text{ FLOP/byte}\), stranding execution units in
memory-stalled idle states and forcing the control rate below the
mechanical resonance of the physical body.

\end{fbxSimple}

\begin{figure}[t]

\centering{

\includegraphics[width=0.95\linewidth,height=\textheight,keepaspectratio]{index_files/mediabag//Users/VJ/GitHub/MLSysBook-integrate-physical/physical/book/chapters/03-brain/figures/fig03_patterson_roofline.pdf}

}

\caption{\label{fig-patterson-roofline}\emph{The Patterson roofline
model and embodied memory wall.} Operational arithmetic intensity
(\(I = \text{FLOPs}/\text{Byte}\)) plotted against attainable execution
throughput on embedded edge silicon (a canonical heterogeneous SoC,
\(204.8\text{ GB/s}\) LPDDR5 DRAM bandwidth, \(275\text{ TFLOP/s}\)
compute ceiling). Single-step autoregressive foundation model inference
operates at \(I \approx 1.0\text{ FLOP/Byte}\), trapped deep in the
memory-bound regime where execution units stall \(>99\%\) of the time
waiting for DRAM weight sweeps
(\(t_{\text{stream}} = 68.4\text{ ms} \implies f_{\max} \le 14.6\text{ Hz}\)).
Action chunking (\(H = 32\text{--}64\)) increases operational arithmetic
intensity \(50\times\) to \(I \approx 50\text{--}150\text{ FLOP/Byte}\),
amortizing parameter transfers across time to bridge slow deliberative
foundation models with fast \(1000\text{ Hz}\) real-time reflexes.
INT8/FP8 quantization further doubles arithmetic intensity while halving
DRAM bus power dissipation.}

\end{figure}%

To break through this memory wall and deliver smooth, high-frequency
continuous control, modern physical AI policies employ \textbf{Action
Chunking} (\citeproc{ref-zhao2023learning}{Zhao et al. 2023};
\citeproc{ref-chi2023diffusionpolicy}{Chi et al. 2025}):\sidenote{\footnotesize \textbf{For
  MLSys Readers:} A VLA model replaces discrete text tokens with
  continuous trajectory chunks \(\mathbf{a}_{t:t+H}\) commanding joint
  angles or twists.} predicting a sequence of \(H\) future continuous
actions
\(\mathbf{A}_{t:t+H} = [\mathbf{a}_t, \mathbf{a}_{t+1}, \dots, \mathbf{a}_{t+H}]\)
in a single neural forward pass or fixed diffusion denoising schedule,
amortizing parameter memory transfers across multiple real-time
execution cycles.

Instead of fetching billions of weights from DRAM to predict a single
scalar action for the immediate millisecond, the planning policy (such
as Action Chunking with Transformers {[}ACT{]}, Diffusion Policy, or
Flow Matching) generates an entire temporal horizon of continuous
actions
\(\mathbf{A}_{t:t+H} = [\mathbf{a}_t, \mathbf{a}_{t+1}, \dots, \mathbf{a}_{t+H}] \in \mathbb{R}^{H \times d_a}\)
(typically \(H = 16\text{--}64\) steps spanning
\(320\text{--}1280\text{ ms}\)) in a single forward pass or fixed
denoising schedule (\(K = 10\text{--}16\) diffusion steps).\sidenote{\footnotesize \textbf{Action
  Chunk Horizon Bounds.} Horizon \(H\) balances latency amortization
  against open-loop state drift
  \(\|\mathbf{e}_{t+H}\| \le \frac{\epsilon_{\text{tot}}}{L_f}(e^{L_f H \Delta t} - 1)\).
  Setting \(H \in [16, 64]\) yields stable
  \(320\text{--}1280\text{ ms}\) windows.}

\begin{longtable}[]{@{}
  >{\raggedright\arraybackslash}p{(\linewidth - 6\tabcolsep) * \real{0.2500}}
  >{\raggedright\arraybackslash}p{(\linewidth - 6\tabcolsep) * \real{0.2500}}
  >{\raggedright\arraybackslash}p{(\linewidth - 6\tabcolsep) * \real{0.2500}}
  >{\raggedright\arraybackslash}p{(\linewidth - 6\tabcolsep) * \real{0.2500}}@{}}
\caption{Action Chunk Trajectory Packet Schema
(\(\mathbf{A}_{\text{chunk}}\)).}\label{tbl-03-action-chunk-schema}\tabularnewline
\toprule\noalign{}
\begin{minipage}[b]{\linewidth}\raggedright
Field Name
\end{minipage} & \begin{minipage}[b]{\linewidth}\raggedright
Type
\end{minipage} & \begin{minipage}[b]{\linewidth}\raggedright
Unit
\end{minipage} & \begin{minipage}[b]{\linewidth}\raggedright
Description \& Structural Invariants
\end{minipage} \\
\midrule\noalign{}
\endfirsthead
\toprule\noalign{}
\begin{minipage}[b]{\linewidth}\raggedright
Field Name
\end{minipage} & \begin{minipage}[b]{\linewidth}\raggedright
Type
\end{minipage} & \begin{minipage}[b]{\linewidth}\raggedright
Unit
\end{minipage} & \begin{minipage}[b]{\linewidth}\raggedright
Description \& Structural Invariants
\end{minipage} \\
\midrule\noalign{}
\endhead
\bottomrule\noalign{}
\endlastfoot
\texttt{chunk\_id} & \texttt{uint64\_t} & --- & Monotonically increasing
trajectory sequence ID \\
\texttt{parent\_lease\_id} & \texttt{uint64\_t} & --- & Must match
active unexpired Intent Lease ID \\
\texttt{horizon\_steps} (\(H\)) & \texttt{uint32\_t} & count & Temporal
horizon length of candidate waypoints (e.g., \(32\)) \\
\texttt{step\_duration\_dt} & \texttt{float32} & \(\text{s}\) & Fixed
inter-waypoint control sampling period \(\Delta t = 0.02\text{ s}\) \\
\texttt{trajectory\_q} & \texttt{float32{[}H{]}{[}N{]}} & \(\text{rad}\)
or \(\text{m}\) & Planned joint positions for all \(N\) degrees of
freedom \\
\texttt{trajectory\_qdot} & \texttt{float32{[}H{]}{[}N{]}} &
\(\text{rad/s}\) or \(\text{m/s}\) & Planned velocities satisfying
\(\|\dot{q}_j\| \le \dot{q}_{\max,j}\) \\
\texttt{trajectory\_qddot} & \texttt{float32{[}H{]}{[}N{]}} &
\(\text{rad/s}^2\) & Planned accelerations satisfying
\(\|\ddot{q}_j\| \le \ddot{q}_{\max,j}\) \\
\texttt{t\_generation\_ns} & \texttt{uint64\_t} & \(\text{ns}\) & PTP
synchronized exposure timestamp of source observation \\
\end{longtable}

As structured in the trajectory packet schema
(Table~\ref{tbl-03-action-chunk-schema}), by processing the entire
trajectory horizon simultaneously across parallel matrix
multiplications, action chunking increases the operational arithmetic
intensity by orders of magnitude
(\(I \ge 50\text{--}100\text{ FLOP/byte}\)). This shifts execution out
of the memory-bound bottleneck and into the \textbf{compute-bound
roofline plateau}, saturating tensor cores and allowing generative
policies to synthesize candidate trajectory chunks at
\(20\text{--}50\text{ Hz}\) on compact embedded power envelopes
(\(15\text{--}60\text{ W}\)).

For example, when a canonical bimanual manipulator threads a flexible
USB cable (Figure~\ref{fig-aloha-manipulator}), it executes a 32-step
action chunk. The policy emits coordinated 14-DoF joint setpoints that
smoothly align both grippers, flex the cable, and insert the connector
without pausing between individual control cycles.

However, action chunking introduces the critical systems challenge of
\textbf{chunk seam continuity}. Because each action chunk is generated
independently from rolling sensory observations, the tail of chunk \(k\)
and the head of chunk \(k+1\) may exhibit slight positional and velocity
discrepancies. If a controller concatenated these chunks directly, the
resulting step-discontinuities in acceleration would inject massive jerk
(\(\dddot{\mathbf{q}}\)) into the physical linkages, causing acoustic
vibration, cycloidal pin wear, and motor overcurrent faults.

To guarantee smooth physical execution, systems implement
\textbf{Temporal Ensembling} (exponentially weighted averaging across
overlapping chunk predictions):

\[\mathbf{a}_t = \frac{\sum_{i=0}^{\min(t, H-1)} w_i \, \mathbf{a}_{t \mid t-i}}{\sum_{i=0}^{\min(t, H-1)} w_i}, \quad w_i = \exp(-m \cdot i)\]

where \(\mathbf{a}_{t \mid t-i}\) represents the action predicted for
time step \(t\) by the action chunk initiated at time \(t-i\), and
\(m > 0\) is a tunable temporal discount factor. The resulting smoothed
waypoints are interpolated via quintic (\(C^2\)) splines by the
downstream real-time nervous system at \(1000\text{ Hz}\), ensuring
continuous positions, velocities, and accelerations across all actuator
joints while preserving valid \(SE(3)\) transformations (using Lie
algebra \(\mathfrak{se}(3)\) exponential mapping or Spherical Linear
Interpolation to keep rotation matrices strictly orthonormal).

The complete operational loop spanning intent reasoning, action chunk
synthesis, and deterministic safety arbitration is formalized in
Algorithm 3.1.

\protect\phantomsection\label{callout-algorithmux2a-3.3}
\begin{fbxSimple}{callout-algorithm}{Algorithm:}{Algorithm 3.1: Two-Speed Brain Deliberation and Safety Enforcement}
\phantomsection\label{callout-algorithm*-3.3}

\textbf{Inputs:} Multi-view visual frames
\(\mathbf{I}_t \in \mathbb{R}^{H \times W \times C}\), proprioceptive
state \((\mathbf{q}_t, \dot{\mathbf{q}}_t)\), mission prompt \(M\),
spatial corridor \(\text{Box}_{\text{safe}}\), stopping distance model
\(d_{\text{stop}}(v)\).

\textbf{Outputs:} Real-time motor command \(\mathbf{u}_t\) written to
hardware PWM registers.

\textbf{Structural Invariants:}

\begin{enumerate}
\def\labelenumi{\arabic{enumi}.}
\tightlist
\item
  \emph{Authority Invariant:} No learned proposal \(\mathbf{p}_t\)
  directly modifies actuator registers without passing deterministic
  safety validation.
\item
  \emph{Temporal Freshness Invariant:} If
  \((t_{\text{now}} - t_{\text{transduction}} > \Delta t_{\text{fresh}})\),
  candidate proposal is rejected and deterministic holding reflex is
  engaged.
\item
  \emph{Dynamic Clearance Invariant:} Motor torque is energized if and
  only if \(d_{\text{stop}}(v) \le D_{\text{clear}}\).
\end{enumerate}

\textbf{Execution Pipeline:}

\begin{enumerate}
\def\labelenumi{\arabic{enumi}.}
\tightlist
\item
  \textbf{System 2: Intent Deliberation (\(1\text{--}5\text{ Hz}\) on
  unprivileged Host MPU):}

  \begin{itemize}
  \item
    If (\(\text{active\_lease} == \text{NULL}\) or
    \(t_{\text{now}} \ge \text{active\_lease}.t_{\text{expire}}\)):

    \begin{itemize}
    \tightlist
    \item
      Tokenize visual observation \(\mathbf{I}_t\) and query 3D world
      memory (TSDF / 3DGS).
    \item
      Evaluate Vision-Language-Action backbone to decompose task prompt
      \(M\) into next subgoal.
    \item
      Construct Intent Lease \(\mathcal{L}_{\text{intent}}\) with
      spatial corridor \(\text{Box}_{\text{safe}}\) and expiration
      \(t_{\text{expire}} = t_{\text{now}} + \tau_{\text{valid}}\).
    \item
      Cryptographically sign \(\mathcal{L}_{\text{intent}}\) and
      dispatch to System 1.5 trajectory planner.
    \end{itemize}
  \end{itemize}
\item
  \textbf{System 1.5: Action Chunk Trajectory Synthesis
  (\(20\text{--}50\text{ Hz}\) on Neural Accelerator):}

  \begin{itemize}
  \item
    Ingest joint states \((\mathbf{q}_t, \dot{\mathbf{q}}_t)\) and
    active unexpired \(\mathcal{L}_{\text{intent}}\).
  \item
    Evaluate Diffusion Policy / ACT backbone to synthesize \(H\)-step
    chunk \(\mathbf{A}_{t:t+H}\).
  \item
    Apply temporal ensembling across overlapping chunk horizons:
    \[\mathbf{a}_t = \frac{\sum_{i=0}^{\min(t, H-1)} w_i \, \mathbf{a}_{t \mid t-i}}{\sum w_i}, \quad w_i = \exp(-m \cdot i)\]
  \item
    Fit \(\mathcal{C}^2\) quintic splines through waypoints and package
    into \(\mathbf{A}_{\text{chunk}}\).
  \end{itemize}
\item
  \textbf{System 1: Hard Real-Time Deterministic Reflex Gate
  (\(1000\text{ Hz}\) on Bare-Metal MCU):}

  \begin{itemize}
  \item
    On every \(1\text{ ms}\) timer interrupt:

    \begin{itemize}
    \tightlist
    \item
      Sample calibrated joint encoders and hardware distance sensors.
    \item
      Fetch candidate setpoint \(\mathbf{a}_t\) from active
      \(\mathbf{A}_{\text{chunk}}\).
    \item
      \emph{Evaluate Rejection Gates:}

      \begin{itemize}
      \tightlist
      \item
        If
        \((t_{\text{now}} > \mathcal{L}_{\text{intent}}.t_{\text{expire}}) \implies \text{REJECT}(\text{LEASE\_EXPIRED})\)
      \item
        If
        \((\mathbf{a}_t \notin \text{Box}_{\text{safe}} \lor |\dot{\mathbf{q}}| > \dot{\mathbf{q}}_{\max}) \implies \text{REJECT}(\text{KINEMATIC\_LIMIT})\)
      \item
        If
        \((d_{\text{stop}}(v) > D_{\text{clear}}) \implies \text{REJECT}(\text{COLLISION\_IMMINENT})\)
      \end{itemize}
    \item
      \emph{Arbitration Verdict:}

      \begin{itemize}
      \tightlist
      \item
        If all checks pass: Permission \textbf{GRANTED}
        \(\implies \mathbf{u}_t = \mathbf{a}_t \to \text{PWM Registers}\).
      \item
        Else: Permission \textbf{DENIED}
        \(\implies \mathbf{u}_t = \text{SafeStopReflex}() \to \text{Controlled Braking}\).
      \end{itemize}
    \end{itemize}
  \end{itemize}
\end{enumerate}

\end{fbxSimple}

\protect\phantomsection\label{callout-mathux2a-3.4}
\begin{fbxSimple}{callout-math}{Fermi Math:}{The LPDDR5 Memory Wall and Action Chunking Speedup}
\phantomsection\label{callout-math*-3.4}

How does Action Chunking overcome the memory bus bottleneck on an
embedded physical AI compute module?

\emph{1. The single-token memory bottleneck.} Consider a 7-billion
parameter Vision-Language-Action (VLA) policy stored in 16-bit precision
(FP16) on an embedded edge System-on-Chip (e.g., a canonical
heterogeneous SoC with a 256-bit LPDDR5 interface):

\begin{itemize}
\tightlist
\item
  Weight footprint:
  \(M_{\text{weights}} = 7 \times 10^9 \text{ params} \times 2\text{ bytes} = 14.0\text{ GB}\).
\item
  Theoretical memory bandwidth: \(B_{\text{peak}} = 204.8\text{ GB/s}\).
\item
  Realistic sustained DRAM bandwidth (accounting for \(\sim 70\%\) bus
  scheduling efficiency and row-turnaround delays):
  \(B_{\text{sustained}} \approx 143.4\text{ GB/s}\).
\end{itemize}

Minimum streaming latency to fetch the weights from DRAM across the
silicon bus once:
\[t_{\text{stream}} = \frac{14.0\text{ GB}}{143.4\text{ GB/s}} \approx 97.6\text{ ms}\]

Maximum attainable control frequency for single-step autoregressive
token generation:
\[f_{\text{single}} = \frac{1}{t_{\text{stream}}} \approx \frac{1}{0.0976\text{ s}} \approx 10.2\text{ Hz}\]
When accounting for ViT image tokenization (\(15\text{ ms}\)) and
KV-cache lookups, live single-step frequency drops to
\(1\text{--}5\text{ Hz}\). This is far too slow to stabilize a robotic
limb or mobile chassis without severe mechanical stuttering.

\emph{2. The action chunking amortization.} Instead of predicting one
action step per DRAM weight sweep, the policy generates an action chunk
of \(H = 32\) future time steps (\(\Delta t = 20\text{ ms}\), spanning a
\(640\text{ ms}\) trajectory horizon) in a single neural forward pass.

\begin{itemize}
\tightlist
\item
  Number of weight memory sweeps: 1 sweep (\(97.6\text{ ms}\)).
\item
  Parallel matrix multiplications across \(H\) steps reuse weights
  resident in on-chip SRAM/caches, elevating arithmetic intensity from
  \(I \approx 1.0\text{ FLOP/byte}\) to \(I \ge 65\text{ FLOP/byte}\).
\item
  Amortized compute latency per control step:
  \[\tau_{\text{step}} = \frac{t_{\text{stream}}}{H} = \frac{97.6\text{ ms}}{32\text{ steps}} \approx 3.05\text{ ms/step}\]
\end{itemize}

\emph{3. Physical AI systems impact.} By amortizing weight streaming
across the temporal chunk horizon, the effective control generation rate
climbs from \(10.2\text{ Hz}\) to:
\[f_{\text{effective}} = \frac{1}{\tau_{\text{step}}} \approx 328\text{ Hz (amortized)}\]
The embedded compute module synthesizes a new 32-step chunk every
\(100\text{ ms}\) (\(10\text{ Hz}\) background cadence), providing a
continuous stream of overlapping waypoints to the \(1000\text{ Hz}\)
spline interpolator with over \(80\%\) spare processor duty cycle for
real-time safety monitoring.

\end{fbxSimple}

\protect\phantomsection\label{callout-ruleux2a-3.5}
\begin{fbxSimple}{callout-rule}{Architect’s Rule of Thumb:}{Embodied Cognition Guidelines}
\phantomsection\label{callout-rule*-3.5}

\begin{enumerate}
\def\labelenumi{\arabic{enumi}.}
\tightlist
\item
  \textbf{The Autoregressive Actuator Stall:} Never drive physical motor
  loops directly from single-token autoregressive generation; DRAM
  streaming latency
  (\(t_{\text{stream}} = M_{\text{weights}}/B_{\text{mem}}\)) will
  starve the actuators below the body's resonant frequency.
\item
  \textbf{The Arithmetic Intensity Threshold:} If your policy forward
  pass achieves an arithmetic intensity below \(50\text{ FLOP/byte}\) on
  embedded silicon, you are paying for high-performance tensor cores
  that spend over \(90\%\) of their clock cycles idle in DRAM read
  queues.
\item
  \textbf{The Proposal--Permission Boundary:} The learned brain proposes
  candidate intent and trajectories in unconstrained software
  coordinates; the deterministic reflex engine grants physical execution
  permission in verified SI units.
\item
  \textbf{Softmax Is Not Physical Probability:} Exponentiated logits sum
  to one by mathematical definition, not physical truth; never gate
  safety-critical actuators on raw neural confidence without empirical
  calibration (ECE) and deterministic geometric fallback monitors.
\end{enumerate}

\end{fbxSimple}

\includegraphics[width=0.9\linewidth,height=\textheight,keepaspectratio,alt={The canonical five handoffs architecture and boundary rejection gates vs.~collapsed end-to-end latent pipeline. (Panel A) The canonical modular chain enforces five distinct functional contracts: Observe (timestamped raw evidence), Estimate (metric state belief and covariance), Intend (semantic goal lease with expiration), Plan (candidate action chunk with C\^{}2 spline bounds), and Enforce (active-set barrier QP solver). Each interface features a dedicated rejection gate (RG-1 through RG-5) that intercepts dropped packets, tracking divergence, expired leases, kinematic infeasibilities, or barrier violations before unprivileged proposals (dashed paths) can reach physical actuators. (Panel B) A collapsed end-to-end policy maps raw pixels directly to uninspected motor torques, forfeiting four critical intermediate inspection points and transferring the entire burden of sensory validation, state estimation, and lease verification onto low-level reflex filters.}]{index_files/mediabag//Users/VJ/GitHub/MLSysBook-integrate-physical/physical/book/chapters/03-brain/figures/fig03_five_handoffs.pdf}

\section{Cognitive Vulnerabilities}\label{sec-brain-3-5}

While learned foundation models dismantle the classical wall of
generalizability, they introduce what control theorists and safety
certifiers consider a profound liability: the complete loss of formal
mathematical verifiability.\sidenote{\footnotesize \textbf{Connectionist vs.~Symbolic
  Duality.} Physical AI unites the statistical learning of connectionist
  neural networks (used for high-capacity proposal synthesis) with
  deterministic symbolic logic (used for real-time safety enforcement).}

A deep neural network is a high-dimensional, non-linear statistical
function. Unlike classical linear or convex controllers, a deep network
is generally non-Lipschitz, meaning that infinitesimal perturbations in
sensor inputs can induce massive, unpredictable excursions in output
actions. An engineer cannot prove Lyapunov stability, cannot guarantee
worst-case execution bounds, and cannot achieve formal structural code
coverage under ISO 26262. A model that achieves \(99.9\%\) empirical
success on benchmark datasets will still produce an ungrounded
hallucination on the \(0.1\%\) distributional tail that can destroy
physical hardware or injure human operators.\sidenote{\footnotesize \textbf{Non-Lipschitz
  Neural Mappings.} Deep ReLU/GELU networks exhibit global Lipschitz
  constants \(L\) that scale exponentially with depth
  (\(\prod \|W_l\|\)). Infinitesimal sensor perturbations can induce
  massive trajectory jumps without pointwise safety guarantees.}

Because a learned brain operates through statistical approximation, its
failure modes are fundamentally distinct from traditional software
defects. A classical program fails by throwing an exception, hanging in
an infinite loop, or crashing. A neural network fails silently by
emitting pristine, mathematically bounded tensors that command
catastrophic physical motions. Systems architects must design around
four primary cognitive vulnerabilities:

\subsection*{Semantic Hallucinations}\label{semantic-hallucinations}
\addcontentsline{toc}{subsection}{Semantic Hallucinations}

A semantic hallucination occurs when a model's latent representation
generates a plausible prediction that has no physical counterpart in
reality. When an autonomous mobile manipulator approaches a workstation
in an industrial cleanroom, sunlight reflecting off a polished
transparent acrylic safety shield may be processed by the visual
transformer as open traversable space. The policy outputs a
high-confidence intent lease and trajectory chunk commanding the arm to
reach straight through the barrier. The output tensor passes all
numerical range checks, but its underlying physical premise is entirely
false.

\subsection*{The Softmax Illusion and Out-of-Distribution
Overconfidence}\label{the-softmax-illusion-and-out-of-distribution-overconfidence}
\addcontentsline{toc}{subsection}{The Softmax Illusion and
Out-of-Distribution Overconfidence}

Deployment interfaces frequently mistake mathematical normalization for
physical probability. In deep neural networks, output logits are passed
through a softmax function:

\[p_k = \frac{\exp(z_k)}{\sum_{j=1}^K \exp(z_j)}\]

The resulting values sum to \(1.0\) by mathematical definition, not
because the model is physically certain. When an embodied machine
encounters out-of-distribution (OOD) observations---such as novel
optical textures, deep shadows, or lens water droplets---deep networks
routinely assign confidence scores exceeding \(0.95\) to completely
erroneous classifications.

Evaluating the true relationship between reported model confidence and
empirical physical success requires measuring \textbf{Expected
Calibration Error (ECE)}:{[}\^{}fn-math-ece-formulation{]}

\[\text{ECE} = \sum_{m=1}^M \frac{|B_m|}{N} |\text{acc}(B_m) - \text{conf}(B_m)|\]

As detailed in Table~\ref{tbl-03-calibration-degradation}, under severe
distribution shift, calibration collapses: reported confidence remains
high (\(\text{conf} > 0.85\)) while empirical physical task success
drops toward zero (\(\text{acc} < 0.10\)). Without independent
supervisory verification, trusting raw softmax confidence leads directly
to physical failure.

\begin{longtable}[]{@{}
  >{\raggedright\arraybackslash}p{(\linewidth - 8\tabcolsep) * \real{0.2000}}
  >{\raggedright\arraybackslash}p{(\linewidth - 8\tabcolsep) * \real{0.2000}}
  >{\raggedright\arraybackslash}p{(\linewidth - 8\tabcolsep) * \real{0.2000}}
  >{\raggedright\arraybackslash}p{(\linewidth - 8\tabcolsep) * \real{0.2000}}
  >{\raggedright\arraybackslash}p{(\linewidth - 8\tabcolsep) * \real{0.2000}}@{}}
\caption{\textbf{Expected Calibration Error (ECE) and Failure Profiles
across Distribution Shift Regimes.} Breakdown of nominal model
confidence versus empirical physical task success rates, calibration
error degradation, and deterministic supervisory actions under
progressive environmental and sensor
anomalies.}\label{tbl-03-calibration-degradation}\tabularnewline
\toprule\noalign{}
\begin{minipage}[b]{\linewidth}\raggedright
Distribution Shift Regime
\end{minipage} & \begin{minipage}[b]{\linewidth}\raggedright
Physical \& Environmental Trigger
\end{minipage} & \begin{minipage}[b]{\linewidth}\raggedright
Reported Conf. vs Measured Acc. (\(\text{conf} \text{ vs } \text{acc}\))
\end{minipage} & \begin{minipage}[b]{\linewidth}\raggedright
Calibration Gap (\(\Delta \text{ECE}\))
\end{minipage} & \begin{minipage}[b]{\linewidth}\raggedright
Downstream Supervisory Action
\end{minipage} \\
\midrule\noalign{}
\endfirsthead
\toprule\noalign{}
\begin{minipage}[b]{\linewidth}\raggedright
Distribution Shift Regime
\end{minipage} & \begin{minipage}[b]{\linewidth}\raggedright
Physical \& Environmental Trigger
\end{minipage} & \begin{minipage}[b]{\linewidth}\raggedright
Reported Conf. vs Measured Acc. (\(\text{conf} \text{ vs } \text{acc}\))
\end{minipage} & \begin{minipage}[b]{\linewidth}\raggedright
Calibration Gap (\(\Delta \text{ECE}\))
\end{minipage} & \begin{minipage}[b]{\linewidth}\raggedright
Downstream Supervisory Action
\end{minipage} \\
\midrule\noalign{}
\endhead
\bottomrule\noalign{}
\endlastfoot
\textbf{In-Distribution} (Nominal Baseline) & Clean optics, structured
laboratory lighting, nominal workpiece textures &
\(\text{conf}=0.91, \, \text{acc}=0.89\) & \(\text{ECE} \le 0.035\)
(Calibrated) & Full autonomous execution permitted; nominal velocity
profiles \\
\textbf{Near-OOD Shift} (Covariate Shift) & Specular reflections, lens
water droplets, dynamic shadow transitions &
\(\text{conf}=0.89, \, \text{acc}=0.52\) &
\(\Delta \text{ECE} = +0.370\) (Overconfident) & Selective abstention
triggered: clamp velocity, initiate exploratory tactile probe \\
\textbf{Far-OOD Shift} (Semantic Shift) & Uncataloged deformable
obstacles, human limb intrusion into workspace &
\(\text{conf}=0.78, \, \text{acc}=0.08\) &
\(\Delta \text{ECE} = +0.700\) (Pathological) & Immediate proposal
refusal: hold position, engage controlled braking \\
\textbf{Sensor Degradation} (Hardware Drift) & Photodetector thermal
noise (\(>75^\circ\text{C}\)), focal blur from vibration &
\(\text{conf}=0.84, \, \text{acc}=0.31\) &
\(\Delta \text{ECE} = +0.530\) (Silent failure) & Trigger operational
downgrade; restrict kinematics to certified safety envelope \\
\end{longtable}

\subsection*{Perceptual Chattering and Covariate
Shift}\label{perceptual-chattering-and-covariate-shift}
\addcontentsline{toc}{subsection}{Perceptual Chattering and Covariate
Shift}

When an agent operates in ambiguous sensory environments, discrete
classification heuristics can oscillate rapidly across consecutive time
steps. In the 2018 Tempe autonomous vehicle collision
(\citeproc{ref-ntsb2019uber}{National Transportation Safety Board
2019}), the perception system reclassified a pedestrian pushing a
bicycle multiple times over \(5.6\text{ seconds}\) between
\emph{vehicle}, \emph{bicycle}, and \emph{unknown}, resetting trajectory
path predictions to zero on every discrete state transition.\sidenote{\footnotesize NTSB
  Report HAR-19/03 (\citeproc{ref-ntsb2019uber}{National Transportation
  Safety Board 2019}) revealed that disabling the factory autonomous
  emergency braking system eliminated the deterministic physical safety
  enforcer, relying on unmonitored human supervision.} In physical
manipulation, perceptual chattering injects high-frequency direction
reversals into trajectory planners, commanding destructive torque
chatter that exceeds the mechanical slew rate of actuators.

\subsection*{Temporal Expiration of
Intent}\label{temporal-expiration-of-intent}
\addcontentsline{toc}{subsection}{Temporal Expiration of Intent}

Because physical reality is continuously dynamic, a cognitive thought is
valid only for a finite duration (\(\tau_{\text{valid}}\)). If a
high-level vision-language model requires \(300\text{ ms}\) to
deliberate, the physical state of moving obstacles, robotic arms, or
chemical reactor valves has already evolved by the time the plan is
ready. If downstream communication or execution is delayed, acting upon
an expired intent lease means commanding the physical body based on an
obsolete mental representation of the universe.

\protect\phantomsection\label{callout-incidentux2a-3.6}
\begin{fbxSimple}{callout-incident}{Landmark Engineering Incident:}{The 2016 Williston Tesla Autopilot Collision Under Specular Glare}
\phantomsection\label{callout-incident*-3.6}

On May 7, 2016, near Williston, Florida, a Tesla Model S operating under
Autopilot collided with the side of a refrigerated tractor-semitrailer
crossing an uncontrolled highway intersection at \(74\text{ mph}\)
(\(33\text{ m/s}\)), resulting in a fatal collision
(\citeproc{ref-ntsb2017tesla}{National Transportation Safety Board
2017}).\sidenote{\footnotesize National Transportation Safety Board (NTSB) statutory
  highway accident investigation report NTSB/HAR-17/02
  (\citeproc{ref-ntsb2017tesla}{National Transportation Safety Board
  2017}), \emph{Collision Between a Car Operating With Automated Vehicle
  Control Systems and a Tractor-Semitrailer Truck Near Williston,
  Florida, May 7, 2016}.}

The National Transportation Safety Board (NTSB) investigation traced how
sensory and architectural failures propagated directly into kinetic
destruction:

\begin{enumerate}
\def\labelenumi{\arabic{enumi}.}
\tightlist
\item
  \textbf{Specular Glare and Vision Segmentation Failure:} The vehicle
  approached the intersection against a brightly lit midday sky. The
  high-reflectance white paint on the broad side panel of the crossing
  semitrailer produced intense specular glare that washed out optical
  contrast against the sky. The monocular camera's neural perception
  classifier failed to detect or segment the lateral trailer as a
  physical obstacle, misinterpreting the brightly illuminated panel as
  open overhead sky.
\item
  \textbf{Sensor Fusion and Radar False-Positive Suppression:} The
  forward millimeter-wave radar detected the metallic trailer structure.
  However, to prevent ``phantom braking'' chattering from stationary
  overhead highway signs and overpasses, the sensor fusion logic
  suppressed radar returns that lacked confirming bounding boxes from
  the vision pipeline.
\item
  \textbf{Unchecked Proposal Authority:} Because the learned perception
  system held uninhibited authority over the vehicle's speed controller
  without independent geometric occupancy verification, the driving
  policy operated under the false conviction that the path was clear. It
  maintained full cruise velocity (\(74\text{ mph}\)), commanding
  continuous torque to the drive inverters without initiating
  pre-collision braking or alerting the driver.
\item
  \textbf{Kinetic Energy Propagation:} The vehicle struck the
  semitrailer at full speed. Because the trailer bed clearance exceeded
  the vehicle's hood height, the trailer penetrated directly into the
  passenger cabin, shearing the A-pillars and roof structure. The
  vehicle traveled \(886\text{ ft}\) (\(270\text{ m}\)) beyond the
  impact point before coming to rest.
\end{enumerate}

\emph{Systems Engineering Lesson:} A learned neural network cannot
certify its own perceptual truth. Under optical degradation (specular
glare, high dynamic range, backlighting), statistical vision models fail
through silent semantic omission. Unvalidated perception proposals must
never exercise uninhibited execution authority over vehicle kinetic
energy (\(E_k = \frac{1}{2} m v^2 \approx 1.14\text{ MJ}\)) without
independent, deterministic geometric boundary validation.

\end{fbxSimple}

\subsection*{The Architectural Resolution: The Proposal--Permission
Principle}\label{the-architectural-resolution-the-proposalpermission-principle}
\addcontentsline{toc}{subsection}{The Architectural Resolution: The
Proposal--Permission Principle}

Physical AI reconciles this fundamental tension by establishing the
\textbf{Proposal--Permission Principle}: a learned model possesses the
capability to propose candidate goals, spatial estimations, and motion
trajectories based on statistical pattern matching, but it possesses
\textbf{zero authority} to grant physical execution permission. All
learned outputs are treated as \textbf{unprivileged proposals} (\(p_t\))
that must pass downstream deterministic safety validation before
energizing actuators.

This principle organizes the machine into a \textbf{Two-Speed Brain
Architecture} (\textbf{?@fig-two-speed-brain}). At the top tier,
high-capacity Vision-Language-Action foundation models perform slow
deliberative cognition (System 2) at \(1\text{--}5\text{ Hz}\), parsing
open-world multimodal streams to interpret semantic commands and emit
structured intent leases. At the intermediate tier, generative action
chunking diffusion policies and trajectory transformers (System 1.5)
operate at \(20\text{--}50\text{ Hz}\), unrolling high-level intent into
continuous candidate motion chunks. At the base reflex tier, hard
real-time safety supervisors (System 1) execute deterministically at
\(1000\text{ Hz}\) on isolated safety silicon,\sidenote{\footnotesize Invariant
  violations force PWMs to Hi-Z in \(<100\text{ ns}\) (Safe Torque Off).}
evaluating candidate proposals against physical kinematic invariants and
dynamic stopping clearances (\(d_{\text{stop}} \le D_{\text{clear}}\)).
If a proposal is safe, the supervisor grants execution permission
(\(u_t\)); if it threatens collision or mechanical instability, the
supervisor rejects the proposal (\(\bot\)) and executes a deterministic
safe stop.

Treating the learned brain as an unprivileged proposal generator
guarantees that even when a model hallucinates, suffers
out-of-distribution drift, or generates invalid trajectory setpoints,
the physical machine remains protected by deterministic physical
invariants. When the System 1 supervisor rejects a candidate proposal or
executes a safe holding reflex, the architecture engages a closed-loop
\textbf{Operational Recovery Loop}: the real-time enforcer immediately
asserts an asynchronous notification across the shared IPC mailbox,
causing the System 2 deliberative model to flush its pending trajectory
buffer, refresh its spatial belief map with the current stationary plant
state, and re-prompt the policy with the specific refusal reason (e.g.,
clearance obstacle detected or kinematic boundary reached), preventing
deadlock and continuous chattering.

\includegraphics[width=0.9\linewidth,height=\textheight,keepaspectratio,alt={The two-speed brain architecture and proposal--permission privilege boundary. The systems architecture partitions into an untrusted deliberative cognitive realm (System 2 vision-language model intent proposer at 1\textbackslash text\{-\/-\}5\textbackslash text\{ Hz\} and System 1.5 generative trajectory planner at 20\textbackslash text\{-\/-\}50\textbackslash text\{ Hz\}) and a trusted hard real-time safety realm (System 1 deterministic barrier enforcer at 1000\textbackslash text\{ Hz\}). All unprivileged candidate proposals (p\_t, dashed path) must cross the non-negotiable proposal--permission privilege boundary, where deterministic safety monitors validate dynamic stopping clearances and state constraints before granting permission (u\_t, solid path) to energize physical actuators, falling back to deterministic emergency-stop refusals (\textbackslash bot) upon any contract violation.}]{index_files/mediabag//Users/VJ/GitHub/MLSysBook-integrate-physical/physical/book/chapters/03-brain/figures/fig03_two_speed_brain.pdf}

\section{Systems Stress-Tests}\label{sec-brain-3-6}

To evaluate whole-system engineering judgment, consider five
foundational dilemmas where the cognitive brain confronts the unyielding
realities of the physical world across Mobile Systems, Manipulators, and
Process/Energy Systems.

\subsection*{Stress-Test 1: The Verifiability Dilemma on a Factory
Floor}\label{stress-test-1-the-verifiability-dilemma-on-a-factory-floor}
\addcontentsline{toc}{subsection}{Stress-Test 1: The Verifiability
Dilemma on a Factory Floor}

\emph{The scenario.} An automotive assembly plant evaluates two
competing approaches for installing wiring harnesses into car doors.
Approach A uses classical trajectory optimization and explicit CAD
models, achieving \(100\%\) formal MC/DC safety certification under ISO
26262, but failing on \(42\%\) of flexible harness insertions due to
unmodeled cable deformations. Approach B uses an end-to-end diffusion
policy trained on 5,000 human teleoperation demonstrations, achieving a
\(96\%\) empirical insertion success rate, but providing zero
mathematical proof of worst-case joint velocities or contact forces.

\emph{The systems dilemma.} Plant management cannot deploy Approach A
because it stops the production line on nearly half the assembly cycles.
However, safety compliance officers refuse to sign off on Approach B
because deep neural policies cannot provide formal ISO safety
certificates.

\emph{The systems verdict: Architectural partitioning.} Neither pure
extreme is viable. The physical AI solution wraps the high-capacity
diffusion policy (Approach B) inside a deterministic Control Barrier
Function safety shield (Approach A). The learned brain provides the
open-world deformable compliance to seat the harness, while the
real-time supervisor deterministically clamps actuator current and joint
limits to guarantee physical safety.

\subsection*{Stress-Test 2: Semantic Hallucination under Specular
Reflections}\label{stress-test-2-semantic-hallucination-under-specular-reflections}
\addcontentsline{toc}{subsection}{Stress-Test 2: Semantic Hallucination
under Specular Reflections}

\emph{The scenario.} An autonomous mobile manipulator in a cleanroom
uses a Vision-Language-Action policy to transport chemical reagent
vials. While approaching a transfer station, sunlight reflects off a
polished transparent acrylic safety barrier. The visual transformer
processes the specular glare as empty traversable space, outputting a
high-confidence intent lease and trajectory chunk that commands the arm
to reach directly through the physical barrier plane.

\emph{The systems dilemma.} The output tensor satisfies all structural
checks: coordinates are within mathematical bounds, values contain no
NaNs, and the packet arrives within its \(40.0\text{ ms}\) deadline. The
learned model operates under the statistical conviction that the path is
unobstructed. Because semantic truth cannot be determined from neural
activations alone, executing the proposal would shatter the acrylic
enclosure and destroy the manipulator payload.

\emph{The systems verdict: Semantic opacity invariant.} A neural network
cannot certify its own perceptual truth. In an embodied architecture,
all learned trajectory proposals must be checked against independent,
deterministic geometric boundary models before actuation permission is
granted.

\subsection*{Stress-Test 3: Uncalibrated Confidence in Precision
Contact}\label{stress-test-3-uncalibrated-confidence-in-precision-contact}
\addcontentsline{toc}{subsection}{Stress-Test 3: Uncalibrated Confidence
in Precision Contact}

\emph{The scenario.} An articulated six-axis robotic manipulator
performs precision bin-picking and placement of machined metal castings.
The learned grasping policy outputs candidate end-effector poses
accompanied by a predicted success probability \(\hat{p} \in [0, 1]\).
In an analytical evaluation across 100 grasp attempts where the model
reports confidence in the narrow range of \(\hat{p} \approx 0.90\),
physical execution on the fixture yields 64 successful seating
operations and 36 dropped or misaligned parts.

\emph{The systems dilemma.} The system supervisor uses predicted
confidence to trigger an abstention policy, delegating low-confidence
grasps to an operator or an exploratory tactile scan. The \(0.26\)
calibration gap between reported confidence (\(0.90\)) and empirical
success (\(0.64\)) destroys the monotonic probability ranking required
for selective execution, passing hazardous proposals downstream to
mechanical actuators while providing an illusion of safety.

\emph{The systems verdict: Calibration gating law.} Raw softmax outputs
and logit scores represent ungrounded network activation states rather
than true physical probabilities. A learned model cannot hold selective
authority over mechanical actuators without empirical calibration curves
validated across the target operational domain.

\subsection*{Stress-Test 4: Covariate Shift and Perceptual
Chattering}\label{stress-test-4-covariate-shift-and-perceptual-chattering}
\addcontentsline{toc}{subsection}{Stress-Test 4: Covariate Shift and
Perceptual Chattering}

\emph{The scenario.} A visual navigation policy on an automated delivery
robot encounters a dynamic construction zone where traffic cones are
draped with wet reflective tarps. The perception pipeline reclassifies
the scene on every \(50.0\text{ ms}\) cycle, alternating rapidly between
``static obstacle,'' ``movable drape,'' and ``free pathway.''

\emph{The systems dilemma.} Discretized categorization heuristics reset
motion planning on every classification change, creating severe
chattering setpoints that demand instantaneous direction reversals
beyond the torque slew limits of the wheel drive motors. The learned
brain oscillates in its semantic hypothesis, injecting high-frequency
jerk into the physical chassis.

\emph{The systems verdict: Latent continuity invariant.} Physical bodies
cannot tolerate discrete cognitive chattering. Robust sensorimotor
architectures must use continuous latent representations and multi-step
action chunking to maintain kinodynamic trajectory smoothness across
ambiguous perceptual states.

\subsection*{Stress-Test 5: The Stale Intent Dilemma in Continuous
Process
Operations}\label{stress-test-5-the-stale-intent-dilemma-in-continuous-process-operations}
\addcontentsline{toc}{subsection}{Stress-Test 5: The Stale Intent
Dilemma in Continuous Process Operations}

\emph{The scenario.} A learned optimization policy manages reagent feed
rates and cooling flows on a continuous exothermic chemical reactor
column. The deliberative model generates an intent lease to increase
monomer feed by \(12\%\) over a validity window of
\(\tau_{\text{lease}} = 500\text{ ms}\). Due to thread preemption on the
application host, the trajectory chunk reaches the deterministic
fieldbus gateway \(580\text{ ms}\) after generation (\(80\text{ ms}\)
past expiration). In the intervening interval, an exothermic thermal
pulse elevated core temperature near the thermal runaway threshold
(\(T_{\text{core}} \to T_{\text{crit}}\)).

\emph{The systems dilemma.} The proposed feed rate curve is smooth,
numerically valid, and would optimize yield under nominal conditions.
However, the intent lease is stale. Applying additional monomer feed
under an unmonitored thermal excursion would accelerate exothermic
heating beyond the cooling jacket's heat dissipation capacity.

\emph{The systems verdict: Temporal validity invariant.} The real-time
safety enforcer evaluates authority and temporal validity, never
semantic plausibility. An expired proposal is an invalid proposal; the
nervous system must reject late packets deterministically and engage
emergency cooling reflux.

\section{Summary}\label{sec-brain-3-7}

Physical AI represents neither the blind abandonment of classical
control theory nor the uncritical worship of machine learning black
boxes. It is the deliberate, principled synthesis of both traditions.

Where an engineering problem is bounded, static, and analytically
tractable---such as enforcing joint torque clamps, verifying stopping
distances, or solving quadratic programs at \(1000\text{ Hz}\)---we rely
unconditionally on deterministic classical algorithms. Where an
engineering problem requires parsing the infinite visual and geometric
complexity of the open world---such as open-vocabulary language
grounding, deformable object manipulation, and high-dimensional
trajectory synthesis---we embrace learned foundation models.

\protect\phantomsection\label{callout-takeawaysux2a-3.7}
\begin{fbxSimple}{callout-takeaways}{Key Chapter Takeaways}{}
\phantomsection\label{callout-takeaways*-3.7}

\begin{itemize}
\tightlist
\item
  \textbf{Five Ideal Properties:} Embodied cognition requires semantic
  generalizability, continuous sensory robustness, metric
  spatial-temporal grounding, multi-rate cadence decoupling, and bounded
  execution authority.
\item
  \textbf{Why We Learn (The Four Pillars):} Physical AI adopts learned
  models for open-world generalizability, continuous sensory robustness,
  complex multi-contact/deformable dynamics, and empirical scaling with
  compute and demonstrations.
\item
  \textbf{The Five Stages of Cognition:} Information flows sequentially
  through
  \(\text{Ingestion} \to \text{Perception} \to \text{Memory} \to \text{Intent} \to \text{Planning}\).
  Detailed mathematical formulations are explored in Part III
  (Chapter~\ref{sec-perception} through Chapter~\ref{sec-enforcement}).
\item
  \textbf{Two Cognitive Cadences:} Multi-billion-parameter foundation
  models operate at slow cadences (\(1\text{--}5\text{ Hz}\)) for intent
  reasoning, while \textbf{Action Chunking} (ACT / Diffusion Policy)
  amortizes inference to emit trajectory chunks at
  \(20\text{--}50\text{ Hz}\).
\item
  \textbf{The Proposal--Permission Principle:} The learned brain
  operates strictly as an \textbf{unprivileged proposal engine}
  (\(p_t\)), possessing capability to suggest actions, but zero
  authority to command actuators without deterministic validation.
\item
  \textbf{Softmax is Not Probability:} Exponentiated logits sum to one
  by mathematical definition, not physical truth. Models exhibit severe
  overconfidence under out-of-distribution shift, requiring empirical
  calibration (ECE) and selective abstention.
\end{itemize}

\end{fbxSimple}

The foundational trade-offs between these paradigms are synthesized in
Table~\ref{tbl-03-classical-vs-learned}.

\begin{longtable}[]{@{}
  >{\raggedright\arraybackslash}p{(\linewidth - 6\tabcolsep) * \real{0.2500}}
  >{\raggedright\arraybackslash}p{(\linewidth - 6\tabcolsep) * \real{0.2500}}
  >{\raggedright\arraybackslash}p{(\linewidth - 6\tabcolsep) * \real{0.2500}}
  >{\raggedright\arraybackslash}p{(\linewidth - 6\tabcolsep) * \real{0.2500}}@{}}
\caption{\textbf{Master Comparison: Classical Analytical Methods
vs.~Learned Foundation Models in Physical AI.} Architectural trade-offs
across generalizability, sensory robustness, verifiability, pacing, and
fallback modes, highlighting the physical AI synthesis that partitions
proposal capability from execution
authority.}\label{tbl-03-classical-vs-learned}\tabularnewline
\toprule\noalign{}
\begin{minipage}[b]{\linewidth}\raggedright
Systems Dimension
\end{minipage} & \begin{minipage}[b]{\linewidth}\raggedright
Classical Analytical Methods
\end{minipage} & \begin{minipage}[b]{\linewidth}\raggedright
Learned Foundation Models
\end{minipage} & \begin{minipage}[b]{\linewidth}\raggedright
Physical AI Architectural Synthesis
\end{minipage} \\
\midrule\noalign{}
\endfirsthead
\toprule\noalign{}
\begin{minipage}[b]{\linewidth}\raggedright
Systems Dimension
\end{minipage} & \begin{minipage}[b]{\linewidth}\raggedright
Classical Analytical Methods
\end{minipage} & \begin{minipage}[b]{\linewidth}\raggedright
Learned Foundation Models
\end{minipage} & \begin{minipage}[b]{\linewidth}\raggedright
Physical AI Architectural Synthesis
\end{minipage} \\
\midrule\noalign{}
\endhead
\bottomrule\noalign{}
\endlastfoot
\textbf{Foundational Principle} & Analytical differential equations,
geometry, CAD models, Lyapunov stability & High-capacity neural
representations fitted to empirical multimodal data & \textbf{Hybrid
Hierarchy:} Learned intent proposals gated by analytical safety
shields \\
\textbf{Generalizability \& Perception} & \textbf{Brittle \& Low:}
Confined to pre-modeled geometries; fails under novel textures or
optical flare & \textbf{Unprecedented \& Continuous:} Open-vocabulary
language grounding; latent manifolds absorb noise & \textbf{Unprivileged
Deliberation:} Foundation models synthesize continuous semantic goals \\
\textbf{Mathematical Verifiability} & \textbf{Complete:} Formal proofs
of stability, MC/DC code coverage, ISO compliance & \textbf{Zero:}
Non-Lipschitz black boxes; no formal worst-case execution proofs &
\textbf{Proposal--Permission Boundary:} Safety is guaranteed at the
deterministic gate \\
\textbf{Temporal Pacing} & \textbf{Hard Real-Time:} Deterministic
microsecond execution loops (\(1\text{ kHz}\)) & \textbf{Multi-Rate:}
Deliberative intent (\(1\text{--}5\text{ Hz}\)) vs Action Chunking
(\(20\text{--}50\text{ Hz}\)) & \textbf{Two-Speed Brain:} Action chunks
amortize slow deliberative foundation models \\
\textbf{Failure Modes \& Fallback} & Boolean branch crashes; failure on
uncataloged objects; rule exhaustion & Semantic hallucinations;
out-of-distribution overconfidence; cognitive drift &
\textbf{Deterministic Fallback:} Rejection gates catch invalid proposals
and safe-stop \\
\end{longtable}

These cognitive stages define how an embodied agent perceives,
remembers, reasons, and plans. Yet a mental model alone cannot move
physical mass. In Chapter~\ref{sec-nervous}, we construct the
\textbf{Real-Time Nervous System}---the heterogeneous computing
platform, shared memory buses, lock-free communication protocols, and
deterministic hardware watchdogs that physically bridge the deliberative
Brain to the dynamical Body.

\chapter{\texorpdfstring{The Real-Time Nervous
System\chaptersubtitle{The Fabric of Timing, Signals, and Authority}}{The Real-Time Nervous System}}\label{sec-nervous}

\begin{figure}[H]

\centering{

\includegraphics[width=1\linewidth,height=\textheight,keepaspectratio]{index_files/mediabag//Users/VJ/GitHub/MLSysBook-integrate-physical/physical/book/chapters/04-nervous/figures/fig_locator.pdf}

}

\caption{\label{fig-locator}\textbf{Systems Organ Focus: The Nervous
System.} The four-tier physical AI systems architecture highlights the
real-time nervous system (Layer 2). Operating between the deliberative
brain (Layer 3) and dynamical body (Layer 1), the deterministic nervous
system enforces multi-rate synchronization, transports sensorimotor
signals under hard real-time deadlines, and maintains the exclusive
privilege boundary over actuator energy delivery.}

\end{figure}%

A physical AI system is an uneasy alliance of fundamentally incompatible
physical and computational regimes. The physical body obeys continuous
Newtonian dynamics where moving mass will not pause for computation. The
cognitive brain operates through high-capacity statistical models that
generalize across unstructured environments, but whose execution
latencies are heavy-tailed and whose outputs are unverified. The nervous
system is the deterministic silicon and communication fabric that
bridges these two worlds: it synchronizes multi-rate cadences across
four orders of temporal magnitude, manages lock-free data exchange
across asynchronous clock domains, and enforces the physical invariants
that guarantee digital decisions never violate the mechanical,
electrical, and thermal limits of the machine.

\newpage{}

\protect\phantomsection\label{callout-objectiveux2a-4.1}
\begin{fbxSimple}{callout-objective}{Learning Objectives}{}
\phantomsection\label{callout-objective*-4.1}

After studying this chapter, you will be able to:

\begin{enumerate}
\def\labelenumi{\arabic{enumi}.}
\tightlist
\item
  Explain the role of the nervous system in reconciling stochastic
  high-level deliberation with deterministic physical dynamics.
\item
  Derive the execution rate of each pipeline tier from physical time
  constants rather than arbitrary software benchmarks.
\item
  Architect lock-free single-producer single-consumer (SPSC) seqlock
  buffers to exchange trajectory proposals across rate boundaries in
  bounded time.
\item
  Distinguish empirical latency percentiles (\(P_{99}\)) from structural
  worst-case execution time (WCET) guarantees (\(R = C + B + I \le D\)).
\item
  Establish hardware privilege boundaries that prevent unverified
  learned models from directly writing motor drive registers.
\item
  Derive physical watchdog lease horizons (\(\tau_{\text{lease}}\)) from
  stopping distance mechanics and kinetic energy dissipation limits.
\item
  Apply the Part I 5x5 Handoffs-by-Limits Synthesis Grid to diagnose
  timing jitter, bandwidth bottlenecks, and fault propagation.
\end{enumerate}

\emph{Core chapter takeaway.} Authority over physical actuation is
settled structurally in hardware, and the execution rate of each
subsystem follows from mechanics rather than computational preference.
In physical AI, the nervous system serves as the immutable referee: it
translates the slow, unverified thoughts of the brain into bounded, safe
currents delivered to the body.

\end{fbxSimple}

\section{The Multi-Rate Bridge}\label{sec-nervous-4-1}

\lettrine{A}{s}\lettrine{h}{ighlighted} in the systems organ locator at
the opening of this chapter (Figure~\ref{fig-locator}), the nervous
system occupies the intermediate architectural tier (Layer 2) that binds
the deliberative cognitive brain (Layer 3) to the dynamical physical
body (Layer 1). In Chapter~\ref{sec-boundary}, we established that
physical AI begins the instant digital computation crosses the causal
boundary to command moving mass. In Chapter~\ref{sec-body}, we
quantified the unyielding physical budgets that govern that mass:
reflected rotor inertia (\(N^2 J\)), stopping distance dynamics
(\(d_{\text{stop}}\)), feedback phase lag (\(\Delta \phi\)), and thermal
dissipation limits (\(I^2 R\)). Mass will not wait for computation. In
Chapter~\ref{sec-brain}, we demonstrated that the cognitive brain
requires high-capacity learned foundation models to generalize across
unstructured environments, but whose outputs are stochastic,
non-Lipschitz, and subject to variable inference latencies.

This creates the fundamental temporal conflict of embodied intelligence:
\emph{the cognitive brain needs time to deliberate, but the physical
body cannot wait to survive.}

A high-capacity foundation model or diffusion policy requires tens to
hundreds of milliseconds to ingest multi-camera visual streams, evaluate
cross-attention over spatial tokens, and synthesize future trajectories.
In stark contrast, an electromagnetic motor coil possesses an electrical
time constant (\(\tau_e = L/R\)) measured in fractions of a millisecond,
and moving mechanical linkages carry kinetic energy
(\(E_k = \frac{1}{2} m v^2\)) that will shear gearboxes or collide with
obstacles if current vectors are delayed or miscalculated. If a
vision-language model hallucinates an extreme joint command, stalls
during memory management, or experiences an accelerator latency spike,
the physical body does not encounter a software exception. Mass
continues traveling under its own momentum across all three physical
machine archetypes:

\begin{itemize}
\tightlist
\item
  In \emph{Mobile Systems} (Class 1), an autonomous chassis or quadruped
  traveling at \(2.0\text{ m/s}\) covers \(0.30\text{ m}\) of unguided
  displacement during a \(150\text{ ms}\) inference stall before wheel
  braking can engage.
\item
  In \emph{Manipulators} (Class 2), an articulated robot arm swinging an
  end-effector carries reflected rotor inertia (\(N^2 J\)) that will
  crush tooling or breach fixturing envelopes if high-rate impedance
  loops are interrupted.
\item
  In \emph{Process and Energy Systems} (Class 3), high-power battery
  inverters and thermal cooling circuits experience thermal runaway or
  power stage shoot-through if gate switching vectors drift outside
  microsecond commutation deadlines.
\end{itemize}

Physical safety cannot be guaranteed through software discipline,
operating system thread priorities, or runtime exception handlers alone.
It requires an explicit, deterministic computing, memory, and
communication fabric: the \textbf{real-time nervous system}.

Every signal that crosses the boundary between the brain and the body is
governed by five non-negotiable physical limits:

\begin{enumerate}
\def\labelenumi{\arabic{enumi}.}
\tightlist
\item
  \textbf{Latency (\(\tau\)):} The total transport and computation age
  of a signal. Sensor exposure windows, neural inference runtimes, and
  fieldbus serialization latencies accumulate along the pipeline; if
  cumulative latency exceeds physical margins, control loops
  destabilize.
\item
  \textbf{Bandwidth (\(\beta\)):} The data volume throughput across
  interconnects. Uncompressed high-resolution multi-camera feeds demand
  gigabytes per second, whereas motor actuation setpoints require only
  kilobits per second of compact parameterization.
\item
  \textbf{Thermal and Energy (\(\mathcal{E}\)):} The physical
  dissipation of heat. Massively parallel tensor accelerators consume
  tens of watts within tightly enclosed chassis, while motor windings
  dissipate non-reversible resistive heat (\(I^2 R\)) during continuous
  torque exertion.
\item
  \textbf{Temporal Determinism (\(\sigma_t\)):} The cycle-to-cycle
  execution predictability. A timing jitter of even a few hundred
  microseconds erodes feedback phase margin
  (\(\Delta \phi = -\frac{\omega T_s}{2}\)), transforming stabilizing
  feedback into acoustic chatter and mechanical resonance.
\item
  \textbf{Fault Domain Isolation (\(\mathcal{F}\)):} The structural
  containment of hardware and software failures. A memory leak in an
  operating system daemon, an out-of-memory purge in a neural runtime,
  or an electrical current surge on an accelerator power rail must never
  propagate into the motor control path.
\end{enumerate}

To reconcile the brain's need for deliberative compute time with the
body's need for microsecond physical reflexes across these five limits,
the nervous system must satisfy five ideal properties:

First, it must deliver microsecond cycle determinism. Feedback loops
regulating motor currents and joint torques must execute between
\(1000\text{ Hz}\) and \(20\text{ kHz}\) with sub-microsecond jitter,
ensuring invariant phase margins and ripple-free torque.

Second, it must achieve asynchronous multi-rate decoupling. High-level
cognitive models deliberate at slow, variable cadences
(\(1\text{--}50\text{ Hz}\)), while actuator bridges require continuous,
cycle-precise current vectors at \(1000\text{ Hz}\). The nervous system
must transport trajectory waypoints across these disparate clock domains
in bounded time without allowing compute stalls to block low-level
control loops.

Third, it must enforce isolated fault containment. The failure of
high-level software services must leave the real-time reflex core
running undisturbed on physically segregated hardware.

Fourth, it must maintain uncompromised privilege authority.
High-capacity learned models act strictly as untrusted proposal engines.
Direct physical write access to pulse-width modulation (PWM) registers,
gate-driver timers, and inverter power stages is restricted exclusively
to deterministic hardware enforcers that evaluate physical safety
invariants before permitting energy delivery.

Fifth, it must guarantee autonomous energy dissipation and safe halting.
When communication fails or proposals expire, moving mass continues
traveling under its own momentum. The nervous system must autonomously
detect timeout events via physical watchdog leases
(\(\tau_{\text{lease}}\)), evaluate dynamic stopping distances
(\(d_{\text{stop}} \le D_{\text{clear}}\)), and safely convert kinetic
energy into thermal dissipation through dynamic braking or hardwired
Safe Torque Off (STO) interlocks.

The nervous system is the architectural referee that settles authority
over physical actuation. The central engineering question of this
chapter is: \emph{what silicon partitioning, memory transport contracts,
and deterministic fieldbus fabrics allow us to build this bridge in
physical hardware?}

\section{Silicon Partitioning}\label{sec-nervous-4-2}

The computational workload of a physical AI machine divides cleanly into
two regimes with opposing requirements. The first regime is
deliberation, where learned models evaluate sensory context, infer
environmental state, and propose future action plans. The second regime
is reactive enforcement, where deterministic control logic samples
joints, validates kinematic limits, interpolates motion commands, and
regulates electrical currents to the motors. These two regimes differ in
their fundamental mathematical profiles, hardware architectures, and
tolerance for temporal error.

Deliberate work is defined by deep computational graphs, large memory
footprints, and variable execution times. A vision-language-action model
processing multi-camera video streams, tactile histories, and natural
language prompts requires billions of multiply-accumulate operations
across hundreds of neural network layers. Storing the parameters and
intermediate activation tensors demands gigabytes of high-bandwidth
memory. To achieve acceptable arithmetic throughput, deliberate
computation executes on massively parallel accelerators such as graphics
processing units (GPUs) or dedicated neural processing units (NPUs).
These processors maximize aggregate throughput across wide vector
registers and batched memory pipelines at the direct expense of
deterministic latency. Memory access stalls, dynamic tensor kernel
selection, bus transfers, and runtime memory management cause the
execution time of consecutive inferences to vary across a wide
distribution. If a deliberate model experiences a temporary delay, the
failure mode is soft, resulting in a delayed proposal or a stale
trajectory that can be held or discarded without immediate hardware
damage.

Reactive work operates under the inverse constraints. An inner control
loop computing inverse dynamics, joint limit monitoring, or
field-oriented motor control requires a few thousand arithmetic
operations per update. Its memory footprint is static, measured in
kilobytes rather than gigabytes, with every variable allocated at
compile time in tightly-coupled on-chip static memory to eliminate heap
fragmentation, dynamic paging, and runtime garbage collection. Reactive
routines execute on cycle-counted microcontrollers or field-programmable
gate arrays (FPGAs) designed for bounded execution rather than maximum
parallel throughput. The deadlines governing this work are hard
real-time constraints dictated by the mechanical inertia and electrical
inductance of the physical system, typically running at rates between
\(1000\text{ Hz}\) and \(10000\text{ Hz}\) on update budgets of
\(100\text{ }\mu\text{s}\) to \(1000\text{ }\mu\text{s}\). If a reactive
loop misses its execution deadline, the failure mode is hard, resulting
in lost current control, destabilizing torque ripple, or a physical
collision before the next command can execute.

Because learned inference runtimes depend on memory access patterns,
cache residency, dynamic tensor shapes, and hardware prefetching, their
execution times are stochastic. Consider an analytical example where an
inference pipeline runs on a dedicated edge accelerator with an
execution time \(T_{\text{inf}}\) modeled by a log-normal distribution,
where the natural logarithm of latency follows
\(\mathcal{N}(\mu_{\ln}, \sigma_{\ln}^2)\). Suppose the nominal median
runtime is \(m = 50\text{ ms}\), corresponding to a log-scale location
parameter \(\mu_{\ln} = \ln(50) \approx 3.912\), with a log-scale scale
parameter \(\sigma_{\ln} = 0.35\), and the deliberate loop is assigned a
fixed execution deadline of \(T_{\text{deadline}} = 100\text{ ms}\) at
\(10\text{ Hz}\). The probability of an inference exceeding this
deadline is given by the complementary cumulative distribution function:

\[P(T_{\text{inf}} > T_{\text{deadline}}) = 1 - \Phi\left(\frac{\ln T_{\text{deadline}} - \mu_{\ln}}{\sigma_{\ln}}\right)\]

where \(\Phi(\cdot)\) denotes the standard normal cumulative
distribution function. Evaluating the standard normal variable \(z\)
yields:

\[z = \frac{\ln(100) - 3.912}{0.35} = \frac{4.605 - 3.912}{0.35} \approx 1.98\]

This standard score gives an overrun probability of approximately
\(1 - \Phi(1.98) = 0.0238\), or \(2.38\%\). At an operating rate of
\(10\text{ Hz}\), the system attempts \(36000\) planning cycles per
hour, producing an expected \(857\) deadline violations every hour. Even
if accelerator hardware delivers an average latency well below the
deadline, heavy-tail variance guarantees frequent misses. While an
average latency of \(53\text{ ms}\) and a \(P_{95}\) latency of
\(89\text{ ms}\) appear acceptable under standard software benchmarks,
the \(P_{99}\) tail extends to \(113\text{ ms}\) and the \(P_{99.9}\)
tail reaches \(148\text{ ms}\), creating intervals where the motor
actuators receive no updated commands unless an independent,
deterministic process bridges the gap.

General-purpose operating systems and real-time kernel patches
(\texttt{PREEMPT\_RT}) cannot reconcile these two execution profiles on
a shared processor. Although real-time extensions allow static task
priorities and preemptible kernels, they cannot eliminate shared
microarchitectural resources. When a thread executing a learned model
streams gigabytes of weight matrices across the memory bus, it evicts
cache lines belonging to high-priority control tasks. If the deliberate
runtime triggers a translation lookaside buffer (TLB) shootdown, enters
a driver-level spinlock, or saturates the memory bus crossbar with
direct memory access (DMA) requests, the hardware pipeline
stalls.\sidenote{\footnotesize On non-coherent SoCs, software must execute cache
  invalidation barriers around peripheral DMA transfers to prevent stale
  cache lines from corrupting motor buffers.} A high-priority real-time
interrupt can preempt software execution, but it cannot preempt a
hardware memory transaction already occupying the memory bus.
Furthermore, Dynamic Voltage and Frequency Scaling (DVFS) governors
adjust core clock frequencies based on thermal dissipation and processor
utilization, altering instruction cycle times dynamically. The resulting
jitter ranges from tens of microseconds to several milliseconds, easily
exceeding the entire time budget of a \(1000\text{ Hz}\) control loop.

The hardware floor exposes this physical separation at the interface
boundary. A host processor running an accelerator communicates across
high-latency bus topologies such as PCIe or system fabric crossbars,
where DMA transfers require ring-buffer synchronization, host interrupt
service routines, and cache coherency handshakes that consume
\(10\text{ }\mu\text{s}\) to \(100\text{ }\mu\text{s}\) per transaction.
In contrast, an embedded microcontroller or dedicated real-time motor
controller IC (Figure~\ref{fig-real-motor-controller}) interfaces
directly with sensor transceivers and power stages through memory-mapped
registers, single-cycle input and output pins, and tightly-coupled
static memory with zero wait states. Reading an absolute joint encoder
over a serial peripheral interface (SPI) or updating a field-oriented
pulse-width modulation counter takes less than \(1\text{ }\mu\text{s}\)
of bounded execution time,\sidenote{\footnotesize Dedicated motor control timers (TI
  \texttt{ePWM}, STM32 \texttt{TIM1}) inject programmable dead-time
  delays (\(100\text{--}500\text{ ns}\)) between complementary high-side
  and low-side gate signals (\texttt{TIMx\_BDTR}, \texttt{DBRED}),
  preventing catastrophic direct-current bus shoot-through across power
  inverter phase legs.} completely isolated from system bus contention.

\begin{figure}[t]

\centering{

\includegraphics[width=0.85\linewidth,height=\textheight,keepaspectratio]{/Users/VJ/GitHub/MLSysBook-integrate-physical/physical/book/chapters/04-nervous/figures/fig04_real_motor_drive_controller.jpg}

}

\caption{\label{fig-real-motor-controller}\textbf{Dedicated Real-Time
Motor Drive Controller Silicon.} Macro view of an embedded motor drive
IC with precision current-sense shunt resistors (\texttt{R110}),
multi-layer power routing, and direct gate-drive outputs. Unlike
general-purpose application processors, dedicated real-time control
silicon executes cycle-deterministic current commutation loops without
operating system schedulers, cache hierarchies, or bus crossbar
contention. (Credit: Wikimedia Commons, CC BY-SA 4.0).}

\end{figure}%

In heterogeneous Dual-Brain System-on-Chip (SoC) architectures, this
physical reality enforces a strict dual-processor design requirement.
Deliberate planning and reactive enforcement cannot safely share a
processor core, a cache hierarchy, or a common memory bus. The slow,
stochastic path requires parallel hardware (Linux host application
processors and tensor accelerators) optimized for raw throughput across
wide data streams, while the fast, deterministic path requires isolated
hardware (bare-metal microcontrollers or FPGA cores) optimized for
worst-case execution time and cycle-counted predictability.

\protect\phantomsection\label{callout-ruleux2a-4.2}
\begin{fbxSimple}{callout-rule}{Architect’s Rule of Thumb:}{Silicon Privilege Partitioning}
\phantomsection\label{callout-rule*-4.2}

Never allow a high-capacity learned model to execute in the same fault
domain, address space, or power rail as the actuator commutation loop.
Deliberation is an unprivileged proposal generator; physical register
privilege belongs exclusively to cycle-deterministic real-time silicon.

\end{fbxSimple}

Safety in physical AI is an architectural partition in silicon, not an
operating system scheduling priority. The high-throughput deliberative
engine acts strictly as an untrusted client, while dedicated real-time
silicon acts as the immutable hardware gatekeeper.

\section{Multi-Rate Cadences}\label{sec-nervous-4-3}

The physical separation of compute domains reflects a fundamental
constraint in which no single execution frequency satisfies both the
physical dynamics of the body and the computational complexity of
perception and learned inference. A physical AI nervous system operates
across four distinct rate tiers
(Figure~\ref{fig-04-multirate-cadence-buffer}), each determined by a
different physical process:

The inner reflex and current regulation loop runs between
\(1000\text{ Hz}\) and \(20\text{ kHz}\). The motion interpolation and
tracking loop runs between \(100\text{ Hz}\) and \(500\text{ Hz}\). The
perception and policy planning pipeline processes incoming sensor
streams between \(10\text{ Hz}\) and \(50\text{ Hz}\). The high-level
deliberate reasoning tier updates semantic task targets between
\(0.5\text{ Hz}\) and \(2\text{ Hz}\). Among these tiers, only the
fastest rate has an origin that is non-negotiable.

The inner loop rate is governed by the physical time constants of the
actuator and the structural resonance of the machine. In an
electromagnetic actuator, winding inductance \(L\) and phase resistance
\(R\) define the electrical time constant \(\tau_e = L/R\). As an
analytical example, consider a brushless permanent magnet motor with
phase resistance \(R = 0.60\text{ }\Omega\) and winding inductance
\(L = 120\text{ }\mu\text{H}\). The electrical time constant is
\(\tau_e = 120\text{ }\mu\text{H} / 0.60\text{ }\Omega = 0.20\text{ ms}\).
Regulating current requires sampling phase currents and updating
pulse-width modulation duty cycles before winding current drifts beyond
acceptable control tolerances, setting an execution interval of
\(T_s \le 1.0\text{ ms}\) (\(1000\text{ Hz}\)), with inner
field-oriented current loops executing up to \(20\text{ kHz}\)
(\(50\ \mu\text{s}\)). At the joint level, the mechanical time constant
\(\tau_m\) and structural stiffness dictate the required closed-loop
bandwidth. Sampled-data control theory
(\citeproc{ref-astrom1997computer}{Åström and Wittenmark 1997})
establishes that a discrete sample-and-hold loop introduces an effective
phase delay of \(\Delta \phi \approx \omega (T_s / 2 + T_d)\), where
\(T_s\) is the sampling interval and \(T_d\) is the computation
latency.\sidenote{\footnotesize The zero-order hold introduces an unavoidable
  half-sample phase lag \(\Delta \phi = -\omega T_s / 2\), degrading
  phase margin independently of computational delay \(T_d\).} If a joint
with a target closed-loop crossover frequency of
\(\omega_c = 100\text{ rad/s}\) experiences an effective feedback delay
of \(T_s / 2 + T_d = 5.0\text{ ms}\) (for example, with
\(T_s = 2.0\text{ ms}\) and \(T_d = 4.0\text{ ms}\)), the discrete
sampling process injects:

\[\Delta \phi = 100\text{ rad/s} \times 0.0050\text{ s} = 0.50\text{ rad} \approx 28.6^\circ\]

This phase lag directly erodes the phase margin of the feedback
controller, driving the mechanical assembly toward resonance or
uncontrolled oscillations when loop execution slows below
\(500\text{ Hz}\).

The perception tier executes at an intermediate rate dictated by photon
collection physics, exposure windows, and target velocity relative to
the sensor field of view. An optical image sensor operating under
typical indoor illumination (\(300\text{ lux}\)) requires an integration
window of \(10\text{ ms}\) to \(30\text{ ms}\) to accumulate sufficient
charge on each pixel without amplifying thermal read noise. If a
manipulator end-effector moves at \(1.5\text{ m/s}\), an exposure time
of \(15\text{ ms}\) causes an edge feature to blur across
\(\Delta x = 1.5\text{ m/s} \times 0.015\text{ s} = 22.5\text{ mm}\) on
the image plane. Increasing the frame rate beyond \(60\text{ Hz}\) or
\(100\text{ Hz}\) provides diminishing returns because the physical
shutter cannot shorten without inducing underexposure, while reducing
the frame rate below \(10\text{ Hz}\) allows moving objects to traverse
tens of centimeters between frames, breaking visual tracking
associations. The rate of the perception pipeline is bounded from above
by photon flux and from below by spatial displacement across consecutive
frames.

The deliberate planning and policy tier runs between \(10\text{ Hz}\)
and \(50\text{ Hz}\) based on the horizon of environmental change and
the thermal limits of onboard compute. A vision-language-action policy
evaluates hundreds of millions of parameters, consuming tens of joules
of electrical energy per inference step. The external environment,
including moving obstacles, human collaborators, and changing contact
surfaces, evolves over intervals of \(20\text{ ms}\) to
\(200\text{ ms}\). Evaluating a learned policy at \(20\text{ Hz}\)
(\(50\text{ ms}\) period) to generate multi-step action chunks
(\(K = 16\) to \(50\) future waypoints) matches this environmental
evolution rate while maintaining accelerator dissipation within an
embedded power envelope of \(15\text{ W}\) to \(35\text{ W}\).

Because these tiers execute on isolated hardware domains at different
frequencies, data exchange across rate boundaries must be asynchronous,
non-blocking, and zero-copy
(Figure~\ref{fig-04-multirate-cadence-buffer}). If a \(1000\text{ Hz}\)
control loop waits on a mutex lock held by a \(20\text{ Hz}\) neural
network runtime, a single scheduling delay in the operating system
immediately causes the real-time controller to miss its deadline.

To prevent blocking, the nervous system implements \textbf{lock-free
single-producer single-consumer (SPSC) seqlock buffers} in shared static
memory.

\protect\phantomsection\label{callout-definitionux2a-4.3}
\begin{fbxSimple}{callout-definition}{Definition:}{Lock-Free Sequence Lock (`seqlock`)}
\phantomsection\label{callout-definition*-4.3}

A synchronization primitive for single-writer multi-reader shared memory
that eliminates reader blocking and priority inversion by coupling an
atomic monotonic sequence counter with memory barriers; odd counter
values signal active writer mutation, while identical even counter
values before and after reading guarantee that payload words were
latched without tear or inconsistency.

\end{fbxSimple}

The shared memory region maintains an atomic 64-bit sequence counter
alongside double-buffered payload slots.

When the untrusted application processor prepares a new action chunk
proposal, it begins by reading the current sequence counter and
atomically incrementing it to an odd value, signalling to the rest of
the system that a write operation is actively in progress. It then
copies the proposal payload---including a monotonic sequence identifier,
timestamp, validity lease duration, multi-step trajectory waypoints, and
a cyclic redundancy checksum---into the inactive memory buffer. After
enforcing an explicit memory barrier (\texttt{DMB\ ISHST} /
\texttt{DMB\ ISH}) to ensure all payload bytes are globally visible
across the bus,\sidenote{\footnotesize In SPSC seqlock buffers on weakly-ordered ARM
  cores, producers must issue Data Memory Barriers (\texttt{DMB\ ISHST}
  / \texttt{DMB\ ISH}) and store-release (\texttt{STLR}) instructions
  when updating odd/even sequence counters. Omitted barriers allow
  out-of-order execution to expose torn trajectory payloads before the
  counter update commits.} the processor increments the sequence counter
once more to an even value with store-release semantics (\texttt{STLR}),
atomically publishing the completed frame.

On the receiving side, the deterministic real-time microcontroller
latches the newest proposal without ever acquiring a lock or stalling
execution. At the start of its control cycle, the microcontroller reads
the sequence counter with load-acquire semantics (\texttt{LDAR}). If the
counter is odd, indicating that the producer is midway through updating
the buffer, the consumer immediately aborts the read, retaining its
existing validated trajectory spline. If the counter is even, the
microcontroller copies the payload into its private, tightly-coupled
static memory and re-reads the sequence counter. If the counter changed
during the copy, a torn read occurred due to a concurrent write, and the
frame is safely discarded. If the sequence counter remained identical
and even throughout the read, the latch is verified: the microcontroller
validates the cyclic redundancy checksum (CRC-32), checks timestamp
freshness against its hardware timer, and accepts the proposal. Under
this protocol, read latencies on the real-time microcontroller remain
strictly bounded below \(1.5\ \mu\text{s}\), completely decoupled from
the scheduling state or operating system load of the host processor.

\protect\phantomsection\label{callout-algorithmux2a-4.4}
\begin{fbxSimple}{callout-algorithm}{Algorithm:}{Algorithm 4.1: Lock-Free Trajectory Seqlock Exchange Protocol}
\phantomsection\label{callout-algorithm*-4.4}

\textbf{Input (Writer):} Proposal payload
\(P = \{\text{seq\_id}, \text{timestamp\_ns}, \text{lease\_us}, \mathbf{w}_{1..K}, \text{crc32}\}\).
\textbf{Input (Reader):} Target local buffer \(B_{\text{local}}\) in
tightly-coupled static SRAM. \textbf{State:} Shared SRAM mailbox: atomic
\texttt{uint64} \(c_{\text{seq}}\) (init 0), payload slot
\(\text{Slot}_{\text{shared}}\). \textbf{Output (Reader):} Verified
action chunk \(P_{\text{latched}}\) or status rejection code.

\textbf{Procedure \texttt{PUBLISH\_ACTION\_CHUNK}(\(P\)):}

\begin{enumerate}
\def\labelenumi{\arabic{enumi}.}
\tightlist
\item
  \(c_{\text{curr}} \leftarrow \text{ATOMIC\_LOAD\_EXPLICIT}(c_{\text{seq}}, \text{memory\_order\_relaxed})\)
\item
  \(\text{ATOMIC\_STORE\_EXPLICIT}(c_{\text{seq}}, c_{\text{curr}} + 1, \text{memory\_order\_relaxed})\)
  \emph{(Mark odd: write in progress)}
\item
  \(\text{HARDWARE\_MEMORY\_BARRIER}(\text{DMB\_ISHST})\) \emph{(Store
  barrier pushes counter)}
\item
  \(\text{COPY\_MEMORY}(\text{Slot}_{\text{shared}}.\text{payload}, P, \text{sizeof}(P))\)
  \emph{(Copy multi-word trajectory)}
\item
  \(\text{HARDWARE\_MEMORY\_BARRIER}(\text{DMB\_ISH})\) \emph{(Full
  barrier: payload visible)}
\item
  \(\text{ATOMIC\_STORE\_EXPLICIT}(c_{\text{seq}}, c_{\text{curr}} + 2, \text{memory\_order\_release})\)
  \emph{(Mark even: publish committed)}
\end{enumerate}

\textbf{Procedure
\texttt{LATCH\_LATEST\_ACTION\_CHUNK}(\(B_{\text{local}}, t_{\text{now}}\)):}

\begin{enumerate}
\def\labelenumi{\arabic{enumi}.}
\tightlist
\item
  \(c_{\text{start}} \leftarrow \text{ATOMIC\_LOAD\_EXPLICIT}(c_{\text{seq}}, \text{memory\_order\_acquire})\)
  \emph{(Load initial counter)}
\item
  \textbf{if} \(c_{\text{start}}\) is \textbf{ODD then return}
  \(\text{STATUS\_WRITER\_ACTIVE}\) \emph{(Writer active; retain prior
  spline)}
\item
  \(\text{HARDWARE\_MEMORY\_BARRIER}(\text{DMB\_ISHLD})\) \emph{(Load
  barrier prevents reordering)}
\item
  \(\text{COPY\_MEMORY}(B_{\text{local}}, \text{Slot}_{\text{shared}}.\text{payload}, \text{sizeof}(\text{Payload}))\)
  \emph{(Latch to local SRAM)}
\item
  \(\text{HARDWARE\_MEMORY\_BARRIER}(\text{DMB\_ISHLD})\) \emph{(Ensure
  payload reads finish)}
\item
  \(c_{\text{end}} \leftarrow \text{ATOMIC\_LOAD\_EXPLICIT}(c_{\text{seq}}, \text{memory\_order\_acquire})\)
  \emph{(Re-read sequence counter)}
\item
  \textbf{if} \(c_{\text{start}} \neq c_{\text{end}}\) \textbf{then
  return} \(\text{STATUS\_TORN\_READ}\) \emph{(Concurrent write
  detected)}
\item
  \textbf{if}
  \(\text{COMPUTE\_CRC32}(B_{\text{local}}) \neq B_{\text{local}}.\text{crc32}\)
  \textbf{then return} \(\text{STATUS\_CHECKSUM\_INVALID}\)
\item
  \textbf{if}
  \((t_{\text{now}} - B_{\text{local}}.\text{timestamp\_ns}) > B_{\text{local}}.\text{lease\_us} \times 1000\)
  \textbf{then return} \(\text{STATUS\_LEASE\_EXPIRED}\)
\item
  \textbf{return} \(\text{STATUS\_VALID\_LATCHED}\) \emph{(Proposal
  verified for spline synthesis)}
\end{enumerate}

\end{fbxSimple}

The trajectory payload exchanged through this lock-free mailbox follows
the strict schema defined in Table~\ref{tbl-04-schema-action-chunk},
establishing explicit data types, SI units, and structural invariants.

\begin{longtable}[]{@{}
  >{\raggedright\arraybackslash}p{(\linewidth - 6\tabcolsep) * \real{0.2500}}
  >{\raggedright\arraybackslash}p{(\linewidth - 6\tabcolsep) * \real{0.2500}}
  >{\raggedright\arraybackslash}p{(\linewidth - 6\tabcolsep) * \real{0.2500}}
  >{\raggedright\arraybackslash}p{(\linewidth - 6\tabcolsep) * \real{0.2500}}@{}}
\caption{\textbf{Abstract Data Schema for Real-Time Action Chunk
Exchange (\protect\texttt{TrajectoryProposalChunk}).} Field layout,
scalar types, physical units, and structural invariants for trajectory
proposals exchanged across the proposal--permission boundary via
lock-free shared
memory.}\label{tbl-04-schema-action-chunk}\tabularnewline
\toprule\noalign{}
\begin{minipage}[b]{\linewidth}\raggedright
Field Name
\end{minipage} & \begin{minipage}[b]{\linewidth}\raggedright
Scalar Type
\end{minipage} & \begin{minipage}[b]{\linewidth}\raggedright
Physical SI Units
\end{minipage} & \begin{minipage}[b]{\linewidth}\raggedright
Structural Invariant \& Description
\end{minipage} \\
\midrule\noalign{}
\endfirsthead
\toprule\noalign{}
\begin{minipage}[b]{\linewidth}\raggedright
Field Name
\end{minipage} & \begin{minipage}[b]{\linewidth}\raggedright
Scalar Type
\end{minipage} & \begin{minipage}[b]{\linewidth}\raggedright
Physical SI Units
\end{minipage} & \begin{minipage}[b]{\linewidth}\raggedright
Structural Invariant \& Description
\end{minipage} \\
\midrule\noalign{}
\endhead
\bottomrule\noalign{}
\endlastfoot
\texttt{sequence\_id} & \texttt{uint64\_t} & dimensionless &
Monotonically strictly increasing integer (\(s_k > s_{k-1}\)). \\
\texttt{timestamp\_ns} & \texttt{uint64\_t} & nanoseconds
(\(\text{ns}\)) & Hardware capture timestamp on master distributed
clock. \\
\texttt{validity\_lease\_us} & \texttt{uint32\_t} & microseconds
(\(\mu\text{s}\)) & Physical validity lease window
(\(T_{\text{lease}} \le 87.6\text{ ms}\)). \\
\texttt{num\_waypoints} & \texttt{uint16\_t} & dimensionless & Chunk
horizon count (\(K \in [16, 50]\)). \\
\texttt{joint\_positions} & \texttt{float32{[}K{]}{[}N\_j{]}} & radians
(\(\text{rad}\)) & Commanded joint angles within reach envelope
(\(q_{\min} \le q \le q_{\max}\)). \\
\texttt{joint\_velocities} & \texttt{float32{[}K{]}{[}N\_j{]}} &
radians/sec (\(\text{rad/s}\)) & Commanded joint velocities within
certified ceilings (\(\|\dot{q}\| \le \dot{q}_{\max}\)). \\
\texttt{feedforward\_torques} & \texttt{float32{[}K{]}{[}N\_j{]}} &
Newton-meters (\(\text{N}\cdot\text{m}\)) & Bounded feedforward motor
torque vectors (\(\|\tau\| \le \tau_{\max}\)). \\
\texttt{crc32} & \texttt{uint32\_t} & bitmask & IEEE 802.3 CRC-32
polynomial over all preceding payload bytes. \\
\end{longtable}

\begin{figure}[t]

\centering{

\includegraphics[width=0.95\linewidth,height=\textheight,keepaspectratio]{index_files/mediabag//Users/VJ/GitHub/MLSysBook-integrate-physical/physical/book/chapters/04-nervous/figures/fig04_multirate_cadence_buffer.pdf}

}

\caption{\label{fig-04-multirate-cadence-buffer}\textbf{Multi-Rate
Execution Cadence and Lock-Free Triple-Buffer Memory Exchange.} (Panel
A) Execution timeline across three heterogeneous rate tiers: stochastic
policy deliberation on the application processor (\(50\text{ Hz}\),
\(T=20\text{ ms}\), variable inference
\(T_{\text{inf}} \sim \text{LogNormal}\) with occasional tail overruns),
optical perception and state tokenization (\(50\text{--}100\text{ Hz}\),
\(10\text{ ms}\) sensor exposure), and deterministic real-time reflex
control on the bare-metal microcontroller (\(1000\text{ Hz}\),
\(T_s=1.0\text{ ms}\)). Inset breakdown shows microsecond-scale slot
budgeting (\(C \le 240\ \mu\text{s} \ll D = 1000\ \mu\text{s}\)) with a
\(760\ \mu\text{s}\) safety jitter margin. (Panel B) Single-producer
single-consumer (SPSC) seqlock triple-buffering memory architecture. The
producer writes new action chunks into an inactive buffer with an odd
sequence counter (\(c_{\text{seq}} = 2k+1\)) and commits via an atomic
store of an even counter (\(c_{\text{seq}} = 2k+2\)). The real-time
microcontroller latches the newest committed buffer in
\(<1.5\ \mu\text{s}\) without holding mutexes or risking priority
inversion, safely falling back to local spline extrapolation or
watchdog-driven emergency braking (\(T_{\text{wd}} \le 87.6\text{ ms}\))
upon buffer starvation.}

\end{figure}%

Rate mismatches introduce three failure modes: phase lag accumulation,
aliasing, and buffer starvation. In a discrete-to-continuous conversion,
a Zero-Order Hold (ZOH) holds each sample constant for duration
\(T_{\text{plan}}\), introducing an intrinsic half-sample phase lag
\(\Delta \phi = -\frac{\omega T_{\text{plan}}}{2}\). Directly applying a
planner setpoint across multiple inner loop cycles creates a zero-order
hold delay of \(T_{\text{plan}}/2\). At a \(10\text{ Hz}\) planning
rate, this delay amounts to \(50\text{ ms}\) of uncompensated lag,
degrading control margins if applied without local interpolation. The
inner loop avoids this by evaluating continuous cubic or quintic
polynomial splines (\(q(t) = a_0 + a_1 t + a_2 t^2 + a_3 t^3\)) across
the latched waypoints at \(1000\text{ Hz}\), generating smooth,
differentiable position, velocity, and acceleration setpoints on every
\(1.0\text{ ms}\) tick.

Aliasing occurs when high-frequency sensor noise folds into low-rate
estimators, which the system prevents through analog anti-aliasing
filters and digital decimation within the high-rate microcontroller
domain. Following classical nonlinear control results
(\citeproc{ref-slotine1991applied}{Slotine and Li 1991}), stability
guarantees hold under the assumption that inner-loop execution jitter
remains bounded within a narrow fraction of the sampling period and that
setpoint updates arrive before the local trajectory buffer drains. The
nervous system establishes these assumptions by enforcing
cycle-deterministic interrupt timing on the microcontroller and tracking
buffer depth. If the planner stalls and the buffer drains to its final
valid waypoint, the operational assumption no longer holds. At that
moment, the nervous system rejects further planner inputs and executes
an autonomous, deterministic stopping trajectory computed locally at
\(1000\text{ Hz}\), bringing all joints safely to rest.

\protect\phantomsection\label{callout-ruleux2a-4.5}
\begin{fbxSimple}{callout-rule}{Architect’s Rule of Thumb:}{Lock-Free Cadence Decoupling}
\phantomsection\label{callout-rule*-4.5}

Never allow a high-frequency real-time control loop to acquire a mutex
or spin on a lock shared with a stochastic process; communicate across
clock domains exclusively via lock-free atomic sequence mailboxes or
triple-buffers.

\end{fbxSimple}

Communication between cognitive models and motor loops cannot rely on
blocking queues, mutexes, or network sockets. Every interface boundary
must be bridged by asynchronous, lock-free shared memory or
deterministic fieldbuses that guarantee the high-rate motor loop never
waits for the slow, stochastic brain.

The operational characteristics, physical time constants, hardware
allocation, and failure protocols governing each tier across this
multi-rate hierarchy are formalized in
Table~\ref{tbl-04-cadence-hierarchy}.

\begin{longtable}[]{@{}
  >{\raggedright\arraybackslash}p{(\linewidth - 12\tabcolsep) * \real{0.1429}}
  >{\raggedright\arraybackslash}p{(\linewidth - 12\tabcolsep) * \real{0.1429}}
  >{\raggedright\arraybackslash}p{(\linewidth - 12\tabcolsep) * \real{0.1429}}
  >{\raggedright\arraybackslash}p{(\linewidth - 12\tabcolsep) * \real{0.1429}}
  >{\raggedright\arraybackslash}p{(\linewidth - 12\tabcolsep) * \real{0.1429}}
  >{\raggedright\arraybackslash}p{(\linewidth - 12\tabcolsep) * \real{0.1429}}
  >{\raggedright\arraybackslash}p{(\linewidth - 12\tabcolsep) * \real{0.1429}}@{}}
\caption{\textbf{Multi-Rate Cadence Hierarchy of Physical AI Execution.}
Formal specification of execution tiers, physical time constants, jitter
tolerances, hardware substrates, memory transport contracts, and
deterministic failure protocols. The hierarchy bridges four orders of
temporal magnitude, deriving update rates directly from actuator
electrical time constants (\(\tau_e = L/R\)), mechanical phase margins
(\(\Delta \phi \approx \omega(T_s/2 + T_d)\)), optical exposure windows,
and thermal compute dissipation limits. Hard real-time reflexes at
\(\ge 1000\text{ Hz}\) execute in bare-metal silicon with
sub-microsecond determinism to prevent hardware damage, while stochastic
high-level models at \(10\text{--}50\text{ Hz}\) generate candidate
action proposals across lock-free atomic buffers under explicit intent
leases.}\label{tbl-04-cadence-hierarchy}\tabularnewline
\toprule\noalign{}
\begin{minipage}[b]{\linewidth}\raggedright
Execution Tier \& Frequency
\end{minipage} & \begin{minipage}[b]{\linewidth}\raggedright
Nominal Period (\(T\))
\end{minipage} & \begin{minipage}[b]{\linewidth}\raggedright
Max Jitter Budget (\(\Delta T_{\max}\))
\end{minipage} & \begin{minipage}[b]{\linewidth}\raggedright
Hardware \& Compute Domain
\end{minipage} & \begin{minipage}[b]{\linewidth}\raggedright
Payload \& State Representation
\end{minipage} & \begin{minipage}[b]{\linewidth}\raggedright
Inter-Tier Transport Contract
\end{minipage} & \begin{minipage}[b]{\linewidth}\raggedright
Overrun \& Failure Protocol
\end{minipage} \\
\midrule\noalign{}
\endfirsthead
\toprule\noalign{}
\begin{minipage}[b]{\linewidth}\raggedright
Execution Tier \& Frequency
\end{minipage} & \begin{minipage}[b]{\linewidth}\raggedright
Nominal Period (\(T\))
\end{minipage} & \begin{minipage}[b]{\linewidth}\raggedright
Max Jitter Budget (\(\Delta T_{\max}\))
\end{minipage} & \begin{minipage}[b]{\linewidth}\raggedright
Hardware \& Compute Domain
\end{minipage} & \begin{minipage}[b]{\linewidth}\raggedright
Payload \& State Representation
\end{minipage} & \begin{minipage}[b]{\linewidth}\raggedright
Inter-Tier Transport Contract
\end{minipage} & \begin{minipage}[b]{\linewidth}\raggedright
Overrun \& Failure Protocol
\end{minipage} \\
\midrule\noalign{}
\endhead
\bottomrule\noalign{}
\endlastfoot
\textbf{Reflex \& Safety} (\(\ge 1000\text{ Hz}\)) &
\(T \le 1.0\text{ ms}\) (\(50\text{ }\mu\text{s}\) at \(20\text{ kHz}\))
& \(\Delta T_{\max} < 10\text{ }\mu\text{s}\) (hard real-time) &
Dedicated real-time MCU or FPGA fabric & Phase currents, encoder ticks,
PWM duty cycles, emergency clamp flags & Single-cycle MMIO registers,
tightly-coupled static memory & Immediate low-side MOSFET phase short
(dynamic braking) or mechanical brake drop \\
\textbf{Motion \& Tracking} (\(100\text{--}500\text{ Hz}\)) &
\(T = 2.0\text{--}10.0\text{ ms}\) &
\(\Delta T_{\max} \le 200\text{ }\mu\text{s}\) (firm real-time) &
Deterministic real-time RTOS / microcontroller core & Continuous
polynomial splines, \(SE(3)\) wrenches, feedforward torques & Lock-free
SPSC circular ring buffer with atomic sequence counters & Local spline
extrapolation (\(50\text{ ms}\) horizon); controlled deceleration \\
\textbf{Planning \& Policy} (\(10\text{--}50\text{ Hz}\)) &
\(T = 20.0\text{--}100.0\text{ ms}\) &
\(\Delta T_{\max} \le 15\text{ ms}\) (\(P_{99}\) tail bound) & Edge
tensor accelerator, host application processor & Action chunks
(\(K\)-step horizon waypoints), visual embeddings & Shared static memory
SPSC seqlock buffers & Expiration of intent lease
(\(T_{\text{lease}} \le 87.6\text{ ms}\)) revokes proposal authority \\
\textbf{Reasoning \& Deliberation} (\(0.5\text{--}2\text{ Hz}\)) &
\(T = 500\text{--}2000\text{ ms}\) & Asynchronous
(\(\Delta T > 500\text{ ms}\)) & Multi-core host CPU, workstation, cloud
VLA coprocessor & Semantic task goals, topological sub-graphs, natural
language tokens & Asynchronous inter-process messaging (gRPC / zero-copy
IPC) & Planner pauses; local motion layer retains stationary impedance
hold \\
\end{longtable}

\section{Deterministic Fieldbuses}\label{sec-nervous-4-4}

A processor promises far less about timing than most software assumes,
and the gap between an observed latency distribution and a structural
guarantee is where deadlines are missed. In statistical performance
analysis, an engineer measures ten million executions of a task, notes
that the ninety-ninth percentile (\(P_{99}\)) latency is
\(135\text{ }\mu\text{s}\) and the maximum observed latency is
\(150\text{ }\mu\text{s}\), and infers that the computation fits
comfortably inside a \(1.0\text{ ms}\) control deadline. That inference
is invalid. An empirical distribution is a record of what happened under
a finite sequence of test conditions, whereas a hard real-time guarantee
is a proof of what can happen under all legal states of the hardware.
When an execution platform includes out-of-order execution pipelines,
shared cache hierarchies, speculative prefetchers, and concurrent direct
memory access traffic, a finite latency sample cannot establish an upper
bound. The missing tail events do not vanish; they simply wait for the
rare hardware collision that triggers them in production.

Establishing a deterministic timing guarantee requires bounding every
term that contributes to task response time. For a periodic workload
with period \(T\) and relative deadline \(D \le T\), the
\textbf{worst-case response time} \(R\) is defined as the maximum
elapsed duration from task release to completion across all valid
hardware and software states:

\[R = C + B + I \le D\]

where \(C\) is the \textbf{worst-case execution time (WCET)} of the
computation running on the core in total isolation, with zero bus
contention and ideal cache hits;\sidenote{\footnotesize Liu and Layland (1973) proved
  that for \(n\) periodic tasks under Rate Monotonic scheduling,
  deadlines are guaranteed if CPU utilization
  \(\sum C_i/T_i \le n(2^{1/n}-1) \to \ln 2 \approx 69.3\%\). Exact
  Response Time Analysis (RTA) extends schedulability to \(100\%\) if
  \(R_i \le D_i\).} \(B\) is the maximum \textbf{blocking duration}
incurred while waiting for shared hardware peripherals, mutexes, or
lower-priority tasks holding critical resources (governed by priority
ceiling protocols); and \(I\) is the cumulative \textbf{interference}
from higher-priority interrupt service routines (ISRs), direct memory
access (DMA) bus arbitration delays, memory crossbar stalls, dynamic RAM
refresh cycles, and cache line evictions. If any one of these three
terms lacks a provable upper bound, the response time \(R\) is unbounded
and the system possesses no timing guarantee, regardless of how small
the mean latency appears in bench benchmarks.

To see how unanalyzed interference breaks an empirical guarantee,
consider an analytical example of a robotic torque computation running
at \(1000\text{ Hz}\) with a hard deadline
\(D = 1000\text{ }\mu\text{s}\). On an isolated \(2.0\text{ GHz}\)
application processor with empty bus queues and hot caches, the workload
exhibits an unloaded baseline execution time of
\(C = 120\text{ }\mu\text{s}\). Across millions of undisturbed bench
cycles, the measured tail latency never exceeds
\(P_{99.99} = 150\text{ }\mu\text{s}\), leaving an apparent safety
margin of \(850\text{ }\mu\text{s}\).

When the same binary runs on an integrated system-on-chip amidst normal
machine activity, the operating environment introduces contention across
shared silicon resources. High-resolution camera streams transfer image
frames via DMA at \(1.2\text{ GB/s}\), an onboard neural network
accelerator reads weight matrices in large bursts from shared memory,
and ambient thermal buildup causes the dynamic voltage and frequency
scaling (DVFS) governor to drop core clock frequency from
\(2.0\text{ GHz}\) to \(1.2\text{ GHz}\). The core slowdown expands the
base compute time by a factor of \(2.0 / 1.2 = 1.67\), raising execution
to \(C = 200\text{ }\mu\text{s}\). Concurrent bus traffic invalidates
cache lines and stalls load instructions at the memory crossbar,
introducing \(I_{\text{bus}} = 750\text{ }\mu\text{s}\) of memory
interference, while unmasked device driver interrupts and lock
contention add \(B = 180\text{ }\mu\text{s}\) of blocking. Summing these
terms yields a total response time of:

\[R = 200\text{ }\mu\text{s} + 180\text{ }\mu\text{s} + 750\text{ }\mu\text{s} = 1130\text{ }\mu\text{s}\]

The task crosses the \(1000\text{ }\mu\text{s}\) deadline by
\(130\text{ }\mu\text{s}\), not because the algorithmic complexity grew,
but because shared silicon resources lacked isolation.

Taming this latency spread requires classifying every source of
execution variance into three structural categories based on whether its
worst-case behavior can be formally bounded:

\emph{Bounded variance sources} encompass deterministic hardware
mechanisms whose execution cycles can be modeled with exact algebraic
upper bounds: fixed-pipeline single-cycle instructions, deterministic
branch decisions, round-robin bus arbitration, tightly-coupled static
memory access with zero wait states, and vectored interrupt dispatch on
dedicated bare-metal microcontrollers.

\emph{Experimentally characterized sources} encompass shared
architectural mechanisms whose maximum latencies cannot be proven from
first-principles schematics alone, but whose worst-case envelopes can be
bounded through exhaustive stress profiling under locked hardware
configurations: shared last-level cache misses under hardware line
locking, write-buffer drain stalls, dynamic RAM refresh intervals, and
thermal throttling ceilings enforced by active cooling.

\emph{Unbounded variance sources} encompass non-deterministic software
and hardware mechanisms whose maximum latency has no finite upper bound:
demand-paged virtual memory page faults, preemptive operating system
thread scheduling without priority inheritance,\sidenote{\footnotesize The 1997 Mars
  Pathfinder mission suffered unbounded priority inversion when a
  low-priority meteorological task blocked high-priority attitude
  control, tripping the system watchdog. The anomaly was resolved by
  enabling priority inheritance, as formalized by Sha, Rajkumar, and
  Lehoczky (1990).} dynamic heap allocations traversing fragmented free
lists, speculative branch mispredictions that pollute pipeline caches,
and unscheduled peripheral direct memory access traffic. A hard
real-time safety path must strictly eliminate every mechanism from this
third category.

\protect\phantomsection\label{callout-incidentux2a-4.6}
\begin{fbxSimple}{callout-incident}{Landmark Engineering Incident:}{The 1997 Mars Pathfinder Priority Inversion}
\phantomsection\label{callout-incident*-4.6}

In July 1997, shortly after the Mars Pathfinder spacecraft touched down
on Ares Vallis, the lander's single-board computer began experiencing
periodic total system resets, threatening mission survival.

The lander's avionics computer ran the VxWorks real-time operating
system with preemptive priority-based scheduling. A high-priority
attitude control thread (\texttt{bc\_dist}, the MIL-STD-1553 Bus
Controller task) ran periodically to distribute spacecraft attitude
telemetry and sensor packets. A low-priority meteorological data
acquisition thread (\texttt{ASI/MET}) periodically acquired a mutual
exclusion semaphore (\texttt{pipeMutex}) to write atmospheric
temperature and pressure readings to an inter-process pipe.

During surface operations, a medium-priority communications task
(performing long-running image compression and telemetry formatting
without accessing \texttt{pipeMutex}) preempted the low-priority
\texttt{ASI/MET} task while it held \texttt{pipeMutex}. When the
high-priority \texttt{bc\_dist} thread was subsequently released by its
periodic timer, it attempted to access the information pipe and blocked
on \texttt{pipeMutex}. Because the medium-priority communications thread
had higher priority than \texttt{ASI/MET}, the RTOS scheduler allowed
the communications thread to continue running, starving \texttt{ASI/MET}
of CPU cycles and preventing it from releasing \texttt{pipeMutex}.
Consequently, the high-priority attitude control thread remained blocked
indefinitely.

The spacecraft's real-time safety architecture included a
hardware/software watchdog timer task running at absolute highest
priority. When \texttt{bc\_dist} failed to execute within its prescribed
timing deadline, the watchdog detected the task starvation and forced a
total system reboot.

JPL engineers reproduced the failure on a terrestrial
hardware-in-the-loop (HIL) testbed using recorded execution trace
telemetry. They transmitted a software patch to Mars, remotely
reconfiguring the VxWorks RTOS mutex initialization parameter to enable
\emph{priority inheritance}
(\texttt{priority\ inversion\ inheritance\ =\ TRUE}). Under priority
inheritance, whenever the high-priority \texttt{bc\_dist} task blocked
on \texttt{pipeMutex}, the low-priority \texttt{ASI/MET} task
temporarily inherited the high priority, allowing it to complete its
critical section, release \texttt{pipeMutex}, and yield immediately to
the attitude controller before the watchdog tripped.

\end{fbxSimple}

A common misconception in real-time software design asserts that dynamic
memory allocation is categorically forbidden on a hard timing path. The
prohibition is imprecise. Unbounded dynamic allocation, exemplified by
general-purpose heap allocators (\texttt{malloc}/\texttt{free}) that
traverse fragmented free lists or execute variable-length page
coalescing, is incompatible with a hard path because its search latency
depends on allocation history and heap fragmentation. In contrast, a
bounded allocator, such as a fixed-size block pool or a Two-Level
Segregated Fit (TLSF) allocator with guaranteed constant-time
\(\mathcal{O}(1)\) allocation and deallocation bounds, can be integrated
safely into a real-time system. On the most safety-critical bare-metal
path, the gold standard remains \textbf{zero-allocation static
architecture}, where all buffers, ring queues, and state variables are
statically declared and mapped to physical memory addresses at compile
time, eliminating memory management execution entirely during runtime.

The choice of physical communication bus directly dictates the magnitude
of the interference term \(I\) and blocking term \(B\) in the worst-case
response time equation. As summarized in
Table~\ref{tbl-04-fieldbus-comparison}, robotics interconnects divide
across a trade-off surface spanning raw throughput, transport latency,
cycle jitter, and hardware determinism. In industrial robotic platforms,
this deterministic physical layer is implemented through dedicated
fieldbus hardware architectures
(Figure~\ref{fig-industrial-ethercat-controller}), which couple hard
real-time host processors to modular input/output and motor drive
terminals over synchronized Ethernet frames using hardware Distributed
Clocks (DC).\sidenote{\footnotesize Under IEC 61158-4-12, EtherCAT Distributed Clocks
  (DC) synchronize slave controllers (e.g., Beckhoff ET1100, TI
  PRU-ICSS) to a master reference clock. Hardware phase-locked loops
  trigger \texttt{SYNC0} physical interrupt pulses directly into motor
  timer pins, holding cyclic synchronization jitter strictly below
  \(100\text{ ns}\).}

\protect\phantomsection\label{callout-mathux2a-4.7}
\begin{fbxSimple}{callout-math}{Fermi Math:}{Fieldbus Serialization Latency and WCET Schedulability}
\phantomsection\label{callout-math*-4.7}

\emph{Practical Engineering Question:} Why can an EtherCAT network
synchronize 8 robot arm joint axes within a \(1.0\text{ ms}\) control
loop with \(< 100\text{ ns}\) jitter, whereas a CAN-FD bus experiences
severe deadline pressure under the identical workload?

\emph{1. CAN-FD Serialization and Arbitration Latency:} Consider an
8-axis manipulator where each joint node exchanges a 64-byte payload
(command torques, encoder positions, velocity feedback, error flags). In
CAN-FD, frames use a \(1\text{ Mbps}\) arbitration phase (header/CRC,
\(\approx 35\text{ bits}\)) and a \(5\text{ Mbps}\) data phase
(\(64\text{ bytes} \times 8\text{ bits} = 512\text{ bits}\) plus stuff
bits \(\approx 550\text{ bits}\)).

\begin{itemize}
\tightlist
\item
  Data transmission time per frame:
  \(t_{\text{data}} = \frac{550\text{ bits}}{5\times 10^6\text{ bps}} = 110\text{ }\mu\text{s}\).
\item
  Header/arbitration time:
  \(t_{\text{arb}} = \frac{35\text{ bits}}{1\times 10^6\text{ bps}} = 35\text{ }\mu\text{s}\).
\item
  Total frame time per axis:
  \(t_{\text{frame}} \approx 145\text{ }\mu\text{s}\).
\end{itemize}

For 8 sequential axis transactions over a single shared bus:
\[t_{\text{total, CAN-FD}} = 8 \times 145\text{ }\mu\text{s} = 1160\text{ }\mu\text{s} = 1.16\text{ ms}\]
Because \(1.16\text{ ms} > 1.0\text{ ms}\), sequential polling overflows
the \(1000\text{ Hz}\) control cycle. Furthermore, bitwise priority
arbitration causes lower-priority axes (e.g.~Axis 8) to suffer
cumulative transmission jitter \(\Delta T > 1.0\text{ ms}\) whenever
higher-priority chassis frames contend for the bus.

\emph{2. EtherCAT On-The-Fly Processing:} In EtherCAT (100BASE-TX,
\(100\text{ Mbps}\)), a single Ethernet telegram traverses all 8 nodes
in a physical daisy chain. Each node's EtherCAT Slave Controller (ESC)
reads its 32-byte Process Data Object (PDO) output and writes its
32-byte input on the fly with a hardware processing delay of
\(t_{\text{node}} \approx 1.0\text{ }\mu\text{s}\).

\begin{itemize}
\tightlist
\item
  Total frame payload:
  \(8 \times 32\text{ bytes} + 64\text{ bytes (Ethernet overhead)} = 320\text{ bytes} = 2560\text{ bits}\).
\item
  Frame wire time:
  \(t_{\text{wire}} = \frac{2560\text{ bits}}{100\times 10^6\text{ bps}} = 25.6\text{ }\mu\text{s}\).
\item
  Total 8-node daisy-chain propagation delay:
  \(t_{\text{prop}} \approx 8 \times 1.0\text{ }\mu\text{s} = 8.0\text{ }\mu\text{s}\).
\item
  Total bus cycle time:
  \(t_{\text{total, EtherCAT}} = 25.6\text{ }\mu\text{s} + 8.0\text{ }\mu\text{s} = 33.6\text{ }\mu\text{s}\).
\end{itemize}

\emph{Systems Verdict:} EtherCAT finishes the entire 8-axis cycle in
\(33.6\text{ }\mu\text{s}\) (consuming only \(3.4\%\) of the
\(1.0\text{ ms}\) servo budget) with hardware Distributed Clocks holding
synchronization jitter below \(100\text{ ns}\). CAN-FD is optimal for
asynchronous, distributed chassis management (BMS, thermal pumps,
auxiliary grippers), while EtherCAT or TSN is mandatory for multi-axis
coordinated kinematic joints. While TSN Ethernet offers the allure of
unifying gigabit camera feeds with microsecond servo frames over a
single physical link, configuring and formally verifying its IEEE
802.1Qbv Time-Aware Shaper (TAS) gate-control lists (GCL) across
asynchronous network switches requires exhaustive network-calculus
verification to prevent high-bandwidth video bursts from colliding with
time-critical control slots.

\end{fbxSimple}

Real-time dependability cannot be achieved by buying faster clock speeds
or measuring statistical latency percentiles on an idle bench.
Dependability requires structural analysis of the execution path:
eliminating all unbounded variance sources
(\(I_{\text{unbounded}} = 0\)), bounding shared hardware blocking
(\(B\)), enforcing zero-allocation static memory layouts, and selecting
fieldbuses with cycle-deterministic distributed clocks.

\begin{figure}[t]

\centering{

\includegraphics[width=0.85\linewidth,height=\textheight,keepaspectratio]{/Users/VJ/GitHub/MLSysBook-integrate-physical/physical/book/chapters/04-nervous/figures/fig04_real_industrial_ethercat_controller.jpg}

}

\caption{\label{fig-industrial-ethercat-controller}\textbf{Industrial
Fieldbus Controller for Real-Time Actuator Control.} An industrial
EtherCAT DIN-rail controller module featuring high-density I/O
terminals, synchronized distributed clock circuitry, and shielded RJ-45
fieldbus ports. The hardware provides sub-microsecond cycle determinism
over synchronized Ethernet frames, bridging host control software to
motor drive amplifiers on automated assembly lines. (Credit: Wikimedia
Commons, CC BY-SA 4.0).}

\end{figure}%

\begin{longtable}[]{@{}
  >{\raggedright\arraybackslash}p{(\linewidth - 12\tabcolsep) * \real{0.1429}}
  >{\raggedright\arraybackslash}p{(\linewidth - 12\tabcolsep) * \real{0.1429}}
  >{\raggedright\arraybackslash}p{(\linewidth - 12\tabcolsep) * \real{0.1429}}
  >{\raggedright\arraybackslash}p{(\linewidth - 12\tabcolsep) * \real{0.1429}}
  >{\raggedright\arraybackslash}p{(\linewidth - 12\tabcolsep) * \real{0.1429}}
  >{\raggedright\arraybackslash}p{(\linewidth - 12\tabcolsep) * \real{0.1429}}
  >{\raggedright\arraybackslash}p{(\linewidth - 12\tabcolsep) * \real{0.1429}}@{}}
\caption{\textbf{Fieldbus and Interconnect Protocol Trade-Off Matrix for
Physical AI Subsystems.} Comparative breakdown of physical signaling,
bandwidth, transfer latency, cycle jitter, determinism classification,
physical reach, and canonical subsystem placement across robotics
interconnects. SPI provides sub-microsecond register access within a
single board; CAN-FD provides robust priority-arbitrated messaging
across chassis wiring; EtherCAT delivers microsecond-level hardware
cycle synchronization across distributed joint motor drives; and TSN
Ethernet unifies high-bandwidth vision streams with scheduled real-time
control frames over a single physical
backbone.}\label{tbl-04-fieldbus-comparison}\tabularnewline
\toprule\noalign{}
\begin{minipage}[b]{\linewidth}\raggedright
Interconnect Protocol
\end{minipage} & \begin{minipage}[b]{\linewidth}\raggedright
Physical Topology \& Signaling
\end{minipage} & \begin{minipage}[b]{\linewidth}\raggedright
Payload Throughput (\(\beta\))
\end{minipage} & \begin{minipage}[b]{\linewidth}\raggedright
Transfer Latency \& Cycle Jitter
\end{minipage} & \begin{minipage}[b]{\linewidth}\raggedright
Determinism Classification
\end{minipage} & \begin{minipage}[b]{\linewidth}\raggedright
Physical Reach \& Wiring Harness
\end{minipage} & \begin{minipage}[b]{\linewidth}\raggedright
Physical AI Subsystem Placement
\end{minipage} \\
\midrule\noalign{}
\endfirsthead
\toprule\noalign{}
\begin{minipage}[b]{\linewidth}\raggedright
Interconnect Protocol
\end{minipage} & \begin{minipage}[b]{\linewidth}\raggedright
Physical Topology \& Signaling
\end{minipage} & \begin{minipage}[b]{\linewidth}\raggedright
Payload Throughput (\(\beta\))
\end{minipage} & \begin{minipage}[b]{\linewidth}\raggedright
Transfer Latency \& Cycle Jitter
\end{minipage} & \begin{minipage}[b]{\linewidth}\raggedright
Determinism Classification
\end{minipage} & \begin{minipage}[b]{\linewidth}\raggedright
Physical Reach \& Wiring Harness
\end{minipage} & \begin{minipage}[b]{\linewidth}\raggedright
Physical AI Subsystem Placement
\end{minipage} \\
\midrule\noalign{}
\endhead
\bottomrule\noalign{}
\endlastfoot
\textbf{SPI} (Serial Peripheral Interface) & Synchronous 4-wire
full-duplex (SCLK, MOSI, MISO, CS), point-to-point &
\(10\text{--}50\text{ Mbps}\) (up to \(80\text{ Mbps}\) in QSPI/OSPI) &
Latency: \(0.5\text{--}2.0\text{ }\mu\text{s}\) \newline Jitter:
\(< 50\text{ ns}\) & \textbf{Bounded} (cycle-deterministic register
reads) & \(< 0.3\text{ m}\) (intra-board PCB traces / flat flex) & Joint
magnetic angle encoders, 3-phase current-sense ADCs, gate drivers \\
\textbf{CAN-FD} (Controller Area Network FD) & 2-wire differential
multidrop with bitwise priority arbitration & \(1\text{ Mbps}\)
arbitration, \(5\text{--}8\text{ Mbps}\) data phase (64-byte payload) &
Latency: \(25\text{--}150\text{ }\mu\text{s}\) \newline Jitter:
\(10\text{--}50\text{ }\mu\text{s}\) & \textbf{Analytically Bounded}
(strictly bounded for top IDs) & \(\le 40\text{ m}\) at
\(1\text{ Mbps}\) / \(\le 15\text{ m}\) at \(5\text{ Mbps}\) (twisted
pair) & Chassis management, battery management systems (BMS), auxiliary
actuators \\
\textbf{EtherCAT} (Ethernet Control Automation) & 100BASE-TX full-duplex
daisy-chain; on-the-fly hardware frame processing & \(100\text{ Mbps}\)
(\(1\text{ Gbps}\) in EtherCAT G), payload efficiency \(>90\%\) &
Latency: \(1.0\text{--}5.0\text{ }\mu\text{s}\) per node
\newline Jitter: \(< 100\text{ ns}\) (distributed clocks) & \textbf{Hard
Real-Time} (hardware-scheduled cyclic PDO frames) & \(\le 100\text{ m}\)
between nodes (shielded Cat5e/Cat6) & Multi-axis robot arm joint servos,
high-rate force-torque sensors \\
\textbf{TSN Ethernet} (IEEE 802.1Qbv / Time-Sensitive Networking) &
Star/mesh switched Ethernet with time-aware traffic shapers &
\(1\text{--}10\text{ Gbps}\) full-duplex & Latency:
\(10\text{--}50\text{ }\mu\text{s}\) (scheduled slots) \newline Jitter:
\(< 1.0\text{ }\mu\text{s}\) (PTP / IEEE 802.1AS) & \textbf{Scheduled
Deterministic} (time-gated queue transmission) & \(\le 100\text{ m}\)
per link & Centralized backbone bridging multi-camera vision, lidar, and
safety reflexes \\
\end{longtable}

\section{Actuation Authority}\label{sec-nervous-4-5}

Authority in a physical AI system is not an abstract software concept;
it is the physical ability to energize motor coils and deliver
mechanical power. In a correctly structured machine, this authority is
concentrated exclusively within the deterministic real-time enforcer
(Figure~\ref{fig-04-privilege-enforcer-registers}).

\begin{figure}[t]

\centering{

\includegraphics[width=0.9\linewidth,height=\textheight,keepaspectratio]{index_files/mediabag//Users/VJ/GitHub/MLSysBook-integrate-physical/physical/book/chapters/04-nervous/figures/fig04_privilege_enforcer_registers.pdf}

}

\caption{\label{fig-04-privilege-enforcer-registers}\textbf{Hardware
Privilege Boundary and Exclusive Motor Drive Register Access.} (Panel A)
Strict physical isolation architecture. The application processor
executes untrusted vision-language-action policies in unprivileged user
space, submitting candidate trajectory proposals (\(p_t\), dashed path)
across an isolated shared-memory transport link. The deterministic
real-time microcontroller retains exclusive physical write privilege
over motor drive registers, evaluating Control Barrier Functions and
stopping invariants before dispatching verified setpoints (\(u_t\),
solid path) to gate drivers, falling back to dynamic braking (\(\bot\))
on violation. (Panel B) Unsafe collapsed architecture, where an
unverified policy writes actuator registers directly, allowing memory
corruption or model hallucinations to command dangerous physical
motions.}

\end{figure}%

A learned foundation model or trajectory planner running on an
application processor possesses \textbf{zero physical authority} to
write pulse-width modulation (PWM) timer registers, set gate-driver
current references, or toggle motor bridge enable pins. Instead,
high-level software generates unprivileged trajectory proposals
(\(p_t\)). These proposals pass across the memory transport interface
into the real-time microcontroller, which functions as a strict hardware
gatekeeper.

The microcontroller continuously evaluates physical safety invariants:
verifying that commanded joint positions lie within physical reach
envelopes, that joint velocities do not exceed certified limits, and
that projected decelerations remain safely within dynamic stopping
clearance (\(d_{\text{stop}} \le D_{\text{clear}}\)). If a proposal
satisfies all invariants, the microcontroller grants execution
permission, interpolates the waypoints into high-frequency current
references (\(u_t\)), and writes the physical gate-driver registers. If
a proposal threatens collision, contains numerical anomalies, or arrives
past its expiration lease, the enforcer revokes permission (\(\bot\))
and executes a deterministic safe stop by commanding hardware trip zone
registers (e.g.~TI \texttt{EPWM\_TZFRC} or STM32
\texttt{TIMx\_BDTR}).\sidenote{\footnotesize Microcontroller hardware trip zones (TI
  \texttt{EPWM\_TZFRC}, STM32 \texttt{TIMx\_BDTR} Break) respond
  autonomously to overcurrent comparators or software safety flags
  within tens of nanoseconds, forcing PWM outputs to high-impedance
  (Hi-Z) or low-side clamp states, bypassing CPU pipelines entirely.}

Consider the physical energy dynamics of this intervention. Suppose an
articulated joint with effective rotor inertia
\(J = 0.050\text{ kg}\cdot\text{m}^2\) rotates at
\(\omega_0 = 10.0\text{ rad/s}\), carrying kinetic energy
\(E_k = \frac{1}{2} J \omega_0^2 = 2.5\text{ J}\). If the enforcer
rejects an invalid proposal and commands emergency dynamic braking under
a maximum deceleration limit \(\alpha_{\max} = 50.0\text{ rad/s}^2\),
bringing the joint to rest requires a braking time:

\[t_{\text{brake}} = \frac{\omega_0}{\alpha_{\max}} = \frac{10.0\text{ rad/s}}{50.0\text{ rad/s}^2} = 0.20\text{ s} = 200\text{ ms}\]

The angular displacement traversed during this braking maneuver is:

\[\Delta \theta = \frac{\omega_0^2}{2 \alpha_{\max}} = \frac{(10.0\text{ rad/s})^2}{2 \times 50.0\text{ rad/s}^2} = 1.00\text{ rad}\]

Over this \(200\text{ ms}\) interval, the entire \(2.5\text{ J}\) of
kinetic energy converts into heat within the motor winding resistance
and switching transistors, yielding an average thermal dissipation rate:

\[P_{\text{avg}} = \frac{2.5\text{ J}}{0.20\text{ s}} = 12.5\text{ W}\]

The physical plant must absorb this energy without exceeding winding
insulation thermal limits or allowing the joint to exceed its geometric
clearance boundary.

At the enforcer boundary, the origin of an aberrant proposal is
irrelevant. A remote memory exploit in an operating system daemon, an
adversarial camera perturbation inducing model hallucination, and an
arithmetic overflow in a neural matrix multiplier produce the identical
boundary event: an invalid proposal packet. Because the enforcer
evaluates invariant bounds against physical sensor feedback from optical
encoders and current shunts rather than inspecting neural network
weights, physical safety bounds hold even under complete compromise of
the host operating system. The privilege boundary guarantees that
software compromise cannot manifest as a violation of physical limits.

To guard against internal hardware faults or microcontroller latch-ups
within the enforcer itself, the system incorporates out-of-band
supervisory channels and physical disconnects. These channels consist of
unprogrammable analog circuits, hardwired emergency stop mushroom push
buttons (Figure~\ref{fig-real-estop}), and electromechanical contactors
that sever power to the motor inverter bridges independently of all
digital logic, such as Safe Torque Off (STO, IEC 61800-5-2).\sidenote{\footnotesize Safe
  Torque Off (STO, IEC 61800-5-2 / SIL 3) uses dual-channel redundant
  \(24\text{ V}\) optocouplers. De-energizing either channel severs gate
  driver bias, allowing pull-down resistors to clamp MOSFET gates
  (\(V_{gs} = 0\text{ V}\)), physically disabling motor torque
  regardless of operating system state or CPU freeze.} While these
disconnects provide absolute power interruption, they impose a severe
cost on system availability: re-energizing a tripped high-voltage
direct-current bus requires pre-charging the inverter filter capacitors
through current-limiting resistors to prevent contactor welding,
consuming hundreds of milliseconds before motor control can resume.
Physical disconnects therefore trade system availability and continuous
tracking for unconditional energy removal. Ensuring that these
supervisory circuits and enforcer channels remain capable of acting when
the primary controller fails requires verifying that they share no
hidden dependencies in their clocks, power rails, or silicon substrates.

\begin{figure}[t]

\centering{

\includegraphics[width=0.85\linewidth,height=\textheight,keepaspectratio]{/Users/VJ/GitHub/MLSysBook-integrate-physical/physical/book/chapters/04-nervous/figures/fig04_real_estop_interlock.jpg}

}

\caption{\label{fig-real-estop}\textbf{Physical Out-of-Band Emergency
Stop Interlock.} A panel-mounted, hardwired emergency stop mushroom push
button with dual-channel positive break contacts. The interlock provides
unprogrammable power isolation directly to motor gate drivers or main
line contactors, operating completely outside digital software pipelines
and microcontroller firmware. (Credit: Wikimedia Commons, CC BY-SA
3.0).}

\end{figure}%

\protect\phantomsection\label{callout-ruleux2a-4.8}
\begin{fbxSimple}{callout-rule}{Architect’s Rule of Thumb:}{Hardware Authority Invariant}
\phantomsection\label{callout-rule*-4.8}

If an unverified learned policy possesses direct memory-mapped write
access to actuator gate-driver registers, your system has no safety
layer---it has only an accident waiting for an out-of-distribution
input.

\end{fbxSimple}

The Proposal--Permission boundary is the fundamental law of actuation
authority. The cognitive brain proposes behaviors across
high-dimensional semantic spaces, but possesses zero authority to
energize coils. Only the deterministic real-time enforcer retains
physical write privilege over the hardware registers that convert
digital words into moving mass.

\section{Hardware Isolation}\label{sec-nervous-4-6}

Independence is a claim about shared clocks, buses, rails, and memory,
and it is usually wrong the first time it is checked. System
architectures routinely label the brain and the nervous system as
isolated because software tasks execute in separate processes,
containers, or hypervisor partitions. True physical isolation requires
four distinct hardware dimensions: independent execution silicon with
private instruction pipelines, separate timing references derived from
physically segregated crystal oscillators, isolated power rails fed by
independent regulator stages, and segregated memory domains that share
neither physical dynamic random-access memory chips nor internal
interconnect arbitration buses.\sidenote{\footnotesize ISO 26262 Part 9 and IEC 61508
  formalize \emph{Freedom from Interference} (FFI). An unverified
  software element (e.g., a neural policy) may share compute or
  communication fabrics with safety-critical tasks only if spatial,
  temporal, and electrical hardware isolation prevents cascading faults.}
If any of these four physical boundaries collapses onto a shared
substrate, a transient failure in the unverified learned model
propagates directly into the deterministic nervous system through the
physics of the underlying electronics.

Shared electrical and physical infrastructure creates cross-domain
failure modes that bypass all software access controls. During a sudden
burst of dense matrix multiplication, a deep learning accelerator can
pull a transient current spike of \(50\text{ A}\) with a slew rate
exceeding \(100\text{ A}/\mu\text{s}\). If the accelerator and the
safety enforcer share a direct-current power rail without dedicated
isolation, the resulting voltage sag causes the enforcer rail to droop
below its microcontroller reset threshold, dropping below
\(3.0\text{ V}\) for \(5\ \mu\text{s}\) and resetting the safety
supervisor while the plant is moving under power. On a shared silicon
package or common thermal heat sink, sustained computational load in the
brain elevates junction temperatures beyond \(105^\circ\text{C}\),
triggering thermal throttling or silicon shutdown in adjacent real-time
cores. When communication relies on an arbitrated interface such as a
shared peripheral interconnect or memory bus, a high-bandwidth visual
pipeline or a misbehaving direct memory access controller can saturate
bus bandwidth, starving the nervous system of encoder feedback and
sensor telemetry. Even a common clock oscillator creates an unbudgeted
single point of failure, where thermal drift, mechanical vibration, or
crystal cracking desynchronizes both domains at the same instant.

A deterministic nervous system must distinguish between two
fundamentally distinct failure modes in the brain: unresponsiveness and
corruption.

Unresponsiveness occurs when an inference pipeline hangs, a scheduling
deadline is missed, or operating system memory fragmentation blocks
execution, producing no output. The nervous system detects
unresponsiveness through a monotonic heartbeat lease, where the brain
must deliver an authenticated, sequentially incrementing token across an
isolated channel before a hardware countdown timer reaches zero.

Corruption occurs when the brain remains fully active and passes its
heartbeat deadline, but generates kinematically impossible setpoints,
discontinuous velocity requests, or numerical infinities resulting from
out-of-distribution visual inputs or matrix overflow. The enforcer
treats these failures through separate mechanisms: unresponsiveness
triggers a time-based lease expiration, whereas corruption triggers an
immediate invariant rejection on the arrival of the malformed proposal
packet.

The duration of the watchdog heartbeat lease is not an arbitrary
software tuning parameter, but a bound dictated by the physical dynamics
of the uncontrolled plant. Consider an analytical mobile platform or
robotic link traveling with an initial velocity
\(v_0 = 1.00\text{ m/s}\) toward a geometric constraint, possessing a
remaining safe clearance margin \(d_{\text{margin}} = 0.25\text{ m}\).
In the event of a silent brain stall, assume the actuators continue
applying their last commanded torque, producing a worst-case uncommanded
acceleration \(a_{\text{fault}} = 2.00\text{ m/s}^2\) until the enforcer
detects the timeout and commands emergency braking with deceleration
\(a_{\text{brake}} = 5.00\text{ m/s}^2\). If the enforcer hardware
requires a physical actuation delay \(t_{\text{act}} = 10\text{ ms}\) to
engage the braking circuits, the plant accelerates unchecked for a total
duration \(\tau = T_{\text{wd}} + t_{\text{act}}\), where
\(T_{\text{wd}}\) is the watchdog timeout. During this unbraked
interval, the plant travels a displacement
\(d_1 = v_0 \tau + \frac{1}{2} a_{\text{fault}} \tau^2\) and reaches a
peak velocity \(v_1 = v_0 + a_{\text{fault}} \tau\). The subsequent
stopping distance under maximum deceleration is
\(d_{\text{brake}} = v_1^2 / (2 a_{\text{brake}})\). Enforcing the
geometric boundary condition
\(d_1 + d_{\text{brake}} \le d_{\text{margin}}\) yields the inequality:

\[v_0 \tau + \frac{1}{2} a_{\text{fault}} \tau^2 + \frac{(v_0 + a_{\text{fault}} \tau)^2}{2 a_{\text{brake}}} \le d_{\text{margin}}\]

Substituting the analytical parameters produces
\(1.00 \tau + 1.00 \tau^2 + 0.10(1.00 + 2.00 \tau)^2 \le 0.25\), which
simplifies to \(1.40 \tau^2 + 1.40 \tau - 0.15 \le 0\). Solving the
positive root gives a maximum allowable unbraked interval
\(\tau \le 97.6\text{ ms}\). Subtracting the \(10\text{ ms}\) actuation
delay yields a strict upper bound on the watchdog timeout of
\(T_{\text{wd}} \le 87.6\text{ ms}\). Setting the lease any higher
permits the physical plant to breach its dynamic safety envelope before
the nervous system is permitted to intervene.

When a lease expires or an enforcer rejects a corrupted proposal, the
nervous system executes an explicit degradation strategy rather than an
unconditioned power cut. In a moving mechanical system, abruptly
removing motor power turns active control into unconstrained ballistic
coasting, where stored kinetic and potential energy can back-drive
gearboxes into mechanical stops. Degradation follows a deterministic
state hierarchy calculated within the nervous system. The first tier is
envelope-bounded dead reckoning, in which the nervous system
extrapolates the trajectory for a fixed horizon of \(50\text{ ms}\)
using a low-order kinematic model constrained by certified velocity
limits. If nominal communication does not resume within this horizon,
the nervous system transitions to controlled deceleration, commanding
torque profiles that bring joint velocities to zero at the maximum
allowable rate that avoids mechanical tipping or thermal damage to the
motor windings. If the platform cannot maintain static equilibrium under
zero velocity, such as a multi-axis arm under gravitational loading, the
nervous system engages electromechanical holding brakes to lock the
joints mechanically.

Once the physical plant is brought to a complete, stationary halt, the
architecture executes a \textbf{Recovery Handshake} before restoring
autonomous proposal authority. The bare-metal microcontroller asserts a
persistent \texttt{FAULT\_HELD} latch in the shared SRAM mailbox and
reports the exact physical rest pose
\((\mathbf{q}_{\text{rest}}, \dot{\mathbf{q}} = \mathbf{0})\). The
application processor must flush its stale action chunks, clear its
autoregressive key-value (KV) attention cache, re-ingest fresh optical
observations tagged with new hardware timestamps, and publish a
re-synchronized intent proposal. The microcontroller verifies that the
initial waypoint of the new proposal matches the physical rest state
within a strict spatial epsilon
(\(||\mathbf{a}_0 - \mathbf{q}_{\text{rest}}|| \le \epsilon_{\text{sync}}\))
before releasing the fault latch and re-energizing PWM modulation.

Verifying that these isolation boundaries function as designed requires
physical falsification on the test bench. Four specific fault-injection
regimes must be executed against the physical prototype: power-rail
droop testing (injecting maximum-slew current steps to verify less than
\(50\text{ mV}\) of rail cross-coupling noise), clock glitching
(verifying independent timing domains maintain stable cycle counts under
frequency drift), communication bus flooding (saturating links to verify
heartbeat telemetry retains deterministic transmission
slots),\sidenote{\footnotesize CAN arbitration uses wired-AND differential signaling
  where dominant zeros overwrite recessive ones. A babbling-node defect
  transmitting continuous frame ID \texttt{0x000} starves all other
  nodes unless intercepted by an active hardware bus guardian or CAN-FD
  transceiver with dominant-timeout clamps.} and numerical fault
injection (streaming synthetic proposals containing NaNs and infinite
floats to verify single-cycle rejection without execution stalls).

Hardware independence must be proven across all four physical dimensions
(silicon cores, timing crystals, power rails, and memory fabrics), and
watchdog lease durations must be derived directly from the physical
stopping distance of the plant.

\section{Invariant Checking}\label{sec-nervous-4-7}

Locomotion policies trained with reinforcement learning and compact
visuomotor diffusion models can be compiled to execute on edge tensor
processors at rates between \(100\text{ Hz}\) and \(500\text{ Hz}\). At
\(500\text{ Hz}\), the model emits a new actuator command every
\(2.0\text{ ms}\). Because this latency is on the same order as the
mechanical time constant of small electric motors and matches the update
rate of conventional field-oriented current controllers, software
designs often fall into an architectural trap by collapsing the
distinction between the brain and the nervous system. The argument made
for this collapse is computational efficiency: if the policy directly
produces joint torques at the native rate of the plant, intermediate
setpoint generators, trajectory interpolators, and tracking controllers
appear to add latency without adding expressiveness.

Rate matching does not confer authority. A neural network emitting
torques at \(500\text{ Hz}\) is governed by the same statistical
uncertainty as a large model generating waypoints at \(5\text{ Hz}\). It
remains an unverified function approximator operating over
high-dimensional input spaces where safety guarantees cannot be proven
statically. Granting the policy direct access to motor-drive registers
eliminates the independent physical layer that enforces mechanical
boundaries. The requirement for separate proposal and permission stages
is not an artifact of mismatched clock rates, but a structural necessity
for maintaining invariant boundaries. Even when the policy and the
enforcer execute within the same \(2.0\text{ ms}\) scheduling period,
the output of the neural network is strictly a proposal that must obtain
permission before changing current in a physical winding.

Operating an unverified learned policy inside the fast control loop
exposes the physical plant to three distinct failure modes: compute
underrun, out-of-distribution torque spikes, and high-frequency
structural limit cycling.

\emph{Compute underrun} occurs when execution jitter in the accelerator
pipeline causes an inference cycle to exceed the actuation deadline. If
memory bus contention or transient thermal throttling extends inference
duration from a nominal \(1.4\text{ ms}\) to \(2.2\text{ ms}\), a model
directly driving the motor bridges leaves the hardware without a valid
command. In a decoupled nervous system, this deadline miss is detected
deterministically by a hardware watchdog timer mapped to the inference
completion interrupt: if the model fails to post a validated setpoint
before the \(2.0\text{ ms}\) boundary, the nervous system drops the
proposal and applies zero-torque freewheeling or shorts the motor phases
for dynamic braking.

\emph{Out-of-distribution torque spikes} occur when an unexpected
physical contact or obscured camera lens pushes the neural network into
an untrained region of its latent space, inducing an instantaneous
torque reversal. In a direct-drive or quasi-direct drive actuator with
rotor inertia \(J_{\text{rotor}} = 0.001\text{ kg}\cdot\text{m}^2\) and
gear reduction \(N = 10\)
(\(J_{\text{ref}} = N^2 J_{\text{rotor}} = 0.10\text{ kg}\cdot\text{m}^2\)),
stepping the commanded torque from \(+30\text{ N}\cdot\text{m}\) to
\(-30\text{ N}\cdot\text{m}\) across a single \(2.0\text{ ms}\) tick
represents an angular jerk of \(3\times 10^5\text{ rad/s}^3\). This
explosive rate of change can shear gear teeth or destroy inverter
switching transistors before any physical displacement is detected by
optical vision. The nervous system prevents this destruction by
evaluating synchronous rate-of-change (\(\frac{d\tau}{dt}\)) and torque
limits against the state of the previous cycle.

\emph{High-frequency structural limit cycling} occurs when unmodeled
link flexibility, gearbox compliance, or sensor phase lag creates
positive feedback with the learned policy, driving the mechanical
assembly into destructive acoustic resonance. Because this resonance
cannot be detected from the command tensor alone, the nervous system
evaluates a continuous energy tank model: monitoring phase-lagged
acceleration on load-side joint encoders and clamping any proposal that
causes net kinetic energy accumulation.

To prevent these failures without degrading loop frequency, the
architecture implements \textbf{synchronous invariant checking} directly
inline with the high-rate inference pipeline. Under a \(2.0\text{ ms}\)
cycle budget, the system allocates \(1.5\text{ ms}\) to sensor
acquisition and neural network execution on the host accelerator,
reserving the final \(500\text{ }\mu\text{s}\) for the nervous system to
evaluate physical invariants on an independent deterministic core. The
enforcer evaluates fixed-point arithmetic expressions that check
kinematic boundaries, maximum torque, maximum jerk, and power envelopes.
It also computes an energy tank model that measures the net mechanical
energy injected into the system over time, clamping any proposal that
would drive the plant into unbounded kinetic accumulation.

When a proposed torque violates an invariant, the enforcer does not
attempt numerical optimization, which would introduce unpredictable
execution latency. Instead, it clips the command to the boundary of the
safe set or substitutes a pre-computed safe control law, such as gravity
compensation combined with joint damping. The fast loop achieves high
throughput for task performance, but physical safety remains governed by
deterministic hardware with independent execution guarantees.

High inference frequency does not grant register write privilege to
unverified neural outputs. Microsecond-scale invariant enforcers must
gate every command before energy reaches the physical plant.

\section{Supervisory Teleoperation}\label{sec-nervous-4-8}

The nervous system places a human operator on the exact same structural
footing as an autonomous policy. A human operator is not an external
master switch with unmediated access to motor amplifiers, but a third
authority holder whose inputs are proposals subject to finite transport
delays and physical execution limits. On a common execution timing
diagram, the three authority paths occupy distinct bands. The nervous
system reflex enforcer executes deterministically on an internal core at
\(1000\text{ Hz}\), evaluating safety invariants every
\(1.0\text{ ms}\). The autonomous learned policy executes on an onboard
accelerator at \(20\text{ Hz}\) to \(50\text{ Hz}\), delivering a fresh
proposal every \(20\text{ ms}\) to \(50\text{ ms}\) with an internal
compute latency of \(15\text{ ms}\). A remote human operator operates at
visual-motor bandwidths between \(5\text{ Hz}\) and \(10\text{ Hz}\),
transmitting setpoints across an external communications link that adds
variable packet transport delays. Because biological nervous systems
exhibit an intrinsic neuromuscular reaction latency of \(150\text{ ms}\)
to \(250\text{ ms}\), an operator viewing a video stream is issuing
commands based on physical states that existed hundreds of milliseconds
in the past.

Treating human inputs as unconditional commands leads directly to
hardware damage. In an emergency or under high cognitive load, an
operator can pull a joystick to its mechanical stop, commanding an
angular acceleration that exceeds inverter current ratings or commanding
a rapid velocity reversal that shears gear teeth. Even during nominal
teleoperation, network jitter and latency can transform an operator's
stabilizing correction into phase-lagged positive feedback, driving the
mechanical structure into divergent oscillations. To prevent these
failures, the system routes all operator inputs through an arbitration
stage termed the \textbf{control multiplexer}. The control multiplexer
selects whether the active proposal originates from the autonomous
policy stream or the human supervisory channel (such as an industrial
teach pendant with deadman switches,
Figure~\ref{fig-real-teach-pendant}) based on operational mode and
heartbeat validity. The output of the multiplexer does not connect
directly to the motor gate drivers. Instead, it passes into the
identical pre-actuation invariant checker that inspects autonomous
proposals, where physical constraints on joint velocity, rotor jerk, and
energy accumulation are evaluated at \(1000\text{ Hz}\). If a human
operator commands a trajectory that violates an invariant boundary, the
enforcer clips the command to the safe set boundary or substitutes a
stabilizing damping law.

\begin{figure}[t]

\centering{

\includegraphics[width=0.85\linewidth,height=\textheight,keepaspectratio]{/Users/VJ/GitHub/MLSysBook-integrate-physical/physical/book/chapters/04-nervous/figures/fig04_real_teach_pendant.jpg}

}

\caption{\label{fig-real-teach-pendant}\textbf{Industrial Teach Pendant
for Human Operator Supervisory Input.} An industrial robotic teach
pendant equipped with a color touch display, 3-axis proportional
joystick, integrated hardwired emergency stop mushroom button, and a
3-position deadman enabling grip switch. Human supervisory jog commands
generated through such interfaces constitute asynchronous proposal
streams that must be arbitrated by the control multiplexer and verified
against physical joint invariants at \(1000\text{ Hz}\) before power is
applied to motor windings. (Credit: Wikimedia Commons, CC BY-SA 3.0).}

\end{figure}%

Because human inputs arrive across a communications channel subject to
latency and packet loss, the physical mechanics of the plant define an
upper bound on tolerable network delay before the nervous system must
intervene. Consider an analytical example of a mobile manipulator moving
along a linear track at velocity \(v = 1.2\text{ m/s}\) toward an
obstacle or workspace limit, with a maximum emergency braking
deceleration \(a_{\max} = 3.0\text{ m/s}^2\). The physical distance
required to bring the moving mass to rest from velocity \(v\) is given
by the kinetic stopping distance
\(d_{\text{stop}} = v^2 / (2 a_{\max})\), which equals
\((1.2\text{ m/s})^2 / (2 \times 3.0\text{ m/s}^2) = 0.24\text{ m}\). If
the local sensing horizon or reserved spatial margin is
\(d_{\text{margin}} = 0.60\text{ m}\), the vehicle travels an
unmonitored distance \(d_{\text{travel}} = v \cdot T_{\text{delay}}\)
during a communications drop or latency spike of duration
\(T_{\text{delay}}\) before corrective braking begins. Total forward
displacement before coming to a complete halt cannot exceed the spatial
margin, yielding the inequality:

\[v \cdot T_{\text{delay}} + \frac{v^2}{2 a_{\max}} \le d_{\text{margin}}\]

Solving for the maximum tolerable latency yields:

\[T_{\text{delay}} \le \frac{d_{\text{margin}}}{v} - \frac{v}{2 a_{\max}}\]

Substituting the analytical parameters produces
\(T_{\text{delay}} \le (0.60\text{ m} / 1.2\text{ m/s}) - (1.2\text{ m/s} / (2 \times 3.0\text{ m/s}^2)) = 0.50\text{ s} - 0.20\text{ s} = 0.30\text{ s}\),
or \(300\text{ ms}\). This total allowable delay budget must accommodate
the onboard video capture and encoding latency
\(T_{\text{cam}} = 30\text{ ms}\), the bidirectional network round-trip
time \(T_{\text{RTT}}\), and the human visual-motor reaction time
\(T_{\text{human}} = 200\text{ ms}\). The maximum allowable network
round-trip latency before teleoperation becomes physically
uncontrollable is therefore:

\[T_{\text{RTT}} \le 300\text{ ms} - 200\text{ ms} - 30\text{ ms} = 70\text{ ms}\]

If \(P_{99}\) network round-trip latency exceeds \(70\text{ ms}\), or if
a command packet fails to arrive within a watchdog interval of
\(300\text{ ms}\), the nervous system cannot wait for an operator
update. The onboard enforcer revokes teleoperation authority and
executes an autonomous braking profile to bring the machine to rest
within the reserved margin.

Placing the human inside the proposal hierarchy clarifies the
distinction between administrative priority and physical authority. An
operator holds higher administrative priority than an autonomous model,
meaning the control multiplexer grants preemption rights to human inputs
whenever an active override signal is received. However, administrative
priority does not grant immunity from physical invariants. Both the
human operator and the machine learning model sit strictly upstream of
the invariant enforcer, differing only in the temporal characteristics
and origin of their proposals. The broader systems problems of human
supervision, including formal state transitions during authority
transfer, cryptographic authentication over wireless uplinks, operator
workload limits, and supervisory user interfaces, constitute the safety
case examined in Chapter~\ref{sec-intervention}. Within the nervous
system, human intervention introduces only the arrival latency of the
proposal stream and the invariant envelope that restricts its physical
realization.

Human teleoperation is an asynchronous proposal stream, not an
unconstrained master override. While human operators hold administrative
priority over autonomous models, their commands must pass through the
identical invariant checks, and network round-trip latencies exceeding
\(T_{\text{RTT}} \le 70\text{ ms}\) trigger autonomous safe halting.

\section{The Handoffs Matrix}\label{sec-nervous-4-9}

The nervous system governs five canonical handoffs that bridge physical
mechanics and computational inference. Upstream, sensing transfers raw
transducer signals to perception pipelines. Perception yields
intermediate representations that fuse into an estimated world state.
The state feeds the learned policy to produce an action plan. The plan
submits candidate trajectories to the safety enforcer. Finally, the
approved command passes from the enforcer to the physical actuators.
Across each boundary, five physical limits govern feasibility: latency
(\(\tau\)), bandwidth (\(\beta\)), thermal energy (\(\mathcal{E}\)),
temporal determinism (\(\sigma_t\)), and fault domain isolation
(\(\mathcal{F}\)). These two axes form a five-by-five matrix that maps
every signal path to the physical budgets of the machine.

This matrix explicitly synthesizes the foundational pillars of Part I:

\emph{The Brain rows (\(\mathcal{H}_1, \mathcal{H}_2, \mathcal{H}_3\))}
inherit the computational and sensory properties established in
Chapter~\ref{sec-brain}. At the sensing-to-perception boundary
(\(\mathcal{H}_1\)), bandwidth and latency bind simultaneously as image
sensors stream raw 4K pixels over high-speed SerDes lanes, where
physical photon integration windows dictate sensor age. Between
perception and state estimation (\(\mathcal{H}_2\)), and state and
planning (\(\mathcal{H}_3\)), energy and latency form the primary
binding constraints. Neural network inference on edge tensor
accelerators demands tens of joules per forward pass, generating heat
within an enclosed chassis, while deep transformer backbones and
diffusion policies incur tens of milliseconds of tokenization and
denoising latency.

\emph{The Body row (\(\mathcal{H}_5\))} inherits the unyielding physical
mechanics established in Chapter~\ref{sec-body}. At the terminal handoff
from the enforcer to the actuators, temporal determinism and thermal
dissipation bind. Brushless motor drivers require cycle-deterministic
field-oriented current loops (\(50\ \mu\text{s}\) at \(20\text{ kHz}\))
over hard real-time fieldbuses to prevent torque ripple, while
continuous current demands dissipate resistive heat (\(I^2 R\)) directly
into stator windings and power MOSFETs under strict thermal clocks.

\emph{The Nervous System row (\(\mathcal{H}_4\))} is the architectural
bridge that ties the brain and body together. The boundary between the
learned planner and the enforcer carries an action proposal (\(p_t\)),
which consists of a compact spline parameterization or waypoint chunk.
Transmitting this payload consumes negligible bandwidth
(\(\beta < 100\text{ kB/s}\)), ensuring that communications never
bottleneck planner-to-enforcer handoffs. Instead, this boundary binds
strictly on temporal determinism (\(\sigma_t\)) and fault domain
isolation (\(\mathcal{F}\)). The nervous system isolates the stochastic,
unverified outputs of the brain, gating them through lock-free seqlock
mailboxes, time-bounded intent leases
(\(\tau_{\text{lease}} \le 87.6\text{ ms}\)), and microsecond-scale
Control Barrier Functions before granting permission to energize
physical mass.

When a physical AI system fails, the failure manifests as an unbudgeted
limit crossing at a specific handoff cell rather than an abstract
deficiency in learned intelligence. If an autonomous manipulator drops a
payload during high-speed tracking, tracing the signal pipeline reveals
that the \(P_{99}\) latency of the perception-to-state estimator drifted
from an analytical budget of \(30\text{ ms}\) to an unbudgeted
\(85\text{ ms}\) under thermal throttling, or that direct memory access
transfers on the main system bus caused jitter in the enforcer loop that
exceeded the \(1\text{ ms}\) deadline. Evaluating field failures against
the grid isolates the exact cell where a physical invariant was
violated, turning diffuse behavioral anomalies into concrete hardware
and timing diagnostics.

The complete five-by-five synthesis grid mapping every canonical signal
handoff against the five physical limits is formalized in
Table~\ref{tbl-04-synthesis-grid}.

\begin{longtable}[]{@{}
  >{\raggedright\arraybackslash}p{(\linewidth - 10\tabcolsep) * \real{0.1667}}
  >{\raggedright\arraybackslash}p{(\linewidth - 10\tabcolsep) * \real{0.1667}}
  >{\raggedright\arraybackslash}p{(\linewidth - 10\tabcolsep) * \real{0.1667}}
  >{\raggedright\arraybackslash}p{(\linewidth - 10\tabcolsep) * \real{0.1667}}
  >{\raggedright\arraybackslash}p{(\linewidth - 10\tabcolsep) * \real{0.1667}}
  >{\raggedright\arraybackslash}p{(\linewidth - 10\tabcolsep) * \real{0.1667}}@{}}
\caption{\textbf{Master 5x5 Handoffs-by-Limits Synthesis Grid for
Physical AI Systems.} The foundational architectural matrix mapping five
canonical signal handoffs (\(\mathcal{H}_1\text{--}\mathcal{H}_5\))
across the Brain (Chapter~\ref{sec-brain}), Body
(Chapter~\ref{sec-body}), and Nervous System (Chapter~\ref{sec-nervous})
against five physical limits
(\(\tau, \beta, \mathcal{E}, \sigma_t, \mathcal{F}\)). Cells marked
\textbf{BINDING} indicate non-negotiable physical constraints dictating
hardware selection, timing budgets, and safety mechanisms. Cells
designated \emph{Non-binding} are architecturally defended by structural
decoupling (such as compact spline parameterization reducing bandwidth
demands or dedicated SerDes hardware eliminating software fault
propagation). This matrix serves as the unified diagnostic rubric across
Part I, formalizing why safety and determinism are properties of
isolated hardware boundaries rather than software
conventions.}\label{tbl-04-synthesis-grid}\tabularnewline
\toprule\noalign{}
\begin{minipage}[b]{\linewidth}\raggedright
Canonical Handoff Boundary
\end{minipage} & \begin{minipage}[b]{\linewidth}\raggedright
Latency Limit (\(\tau\))
\end{minipage} & \begin{minipage}[b]{\linewidth}\raggedright
Bandwidth Limit (\(\beta\))
\end{minipage} & \begin{minipage}[b]{\linewidth}\raggedright
Energy \& Thermal Limit (\(\mathcal{E}\))
\end{minipage} & \begin{minipage}[b]{\linewidth}\raggedright
Determinism Limit (\(\sigma_t\))
\end{minipage} & \begin{minipage}[b]{\linewidth}\raggedright
Fault Domain Isolation (\(\mathcal{F}\))
\end{minipage} \\
\midrule\noalign{}
\endfirsthead
\toprule\noalign{}
\begin{minipage}[b]{\linewidth}\raggedright
Canonical Handoff Boundary
\end{minipage} & \begin{minipage}[b]{\linewidth}\raggedright
Latency Limit (\(\tau\))
\end{minipage} & \begin{minipage}[b]{\linewidth}\raggedright
Bandwidth Limit (\(\beta\))
\end{minipage} & \begin{minipage}[b]{\linewidth}\raggedright
Energy \& Thermal Limit (\(\mathcal{E}\))
\end{minipage} & \begin{minipage}[b]{\linewidth}\raggedright
Determinism Limit (\(\sigma_t\))
\end{minipage} & \begin{minipage}[b]{\linewidth}\raggedright
Fault Domain Isolation (\(\mathcal{F}\))
\end{minipage} \\
\midrule\noalign{}
\endhead
\bottomrule\noalign{}
\endlastfoot
\textbf{\(\mathcal{H}_1\): Sensing \(\to\) Perception} \emph{(Brain)} &
\textbf{BINDING} \newline \(\tau \le 16.7\text{ ms}\) (exposure \&
readout sets sensor age) & \textbf{BINDING}
\newline \(\beta \ge 1.2\text{ GB/s}\) (multi-camera 4K uncompressed
streams) & \emph{Non-binding} \newline \(P < 3\text{ W}\) (sensor head
dissipation isolated) & \textbf{BINDING}
\newline \(\sigma_t \le 1.0\text{ ms}\) (frame-sync triggers prevent
skew) & \emph{Non-binding} \newline Defended by SerDes hardware \& DMA
ring buffers \\
\textbf{\(\mathcal{H}_2\): Perception \(\to\) State} \emph{(Brain)} &
\textbf{BINDING} \newline \(\tau \le 30.0\text{ ms}\) (visual-inertial
odometry \& feature window) & \emph{Non-binding}
\newline \(\beta \approx 10\text{--}50\text{ MB/s}\) (compressed
keypoints \& voxels) & \textbf{BINDING}
\newline \(P \approx 15\text{--}35\text{ W}\) (edge vision backbone
accelerator envelope) & \emph{Non-binding} \newline Defended by
sliding-window state estimator buffers & \textbf{BINDING}
\newline Spatial memory isolation: MPUs prevent memory corruption \\
\textbf{\(\mathcal{H}_3\): State \(\to\) Plan} \emph{(Brain)} &
\textbf{BINDING} \newline \(\tau \le 100.0\text{ ms}\) (autoregressive
tokenization \& diffusion horizon) & \emph{Non-binding}
\newline \(\beta < 5\text{ MB/s}\) (low-dimensional state \& prompt
embeddings) & \textbf{BINDING}
\newline \(P \approx 30\text{--}75\text{ W}\) (transformer/diffusion
matrix accelerators) & \emph{Non-binding} \newline Defended by
receding-horizon chunking \& spline buffers & \textbf{BINDING}
\newline Sandboxed policy container prevents host panics halting
telemetry \\
\textbf{\(\mathcal{H}_4\): Plan \(\to\) Enforcer} \emph{(Nervous
System)} & \textbf{BINDING} \newline \(\tau \le 5.0\text{ ms}\) (intent
lease validation \& CBF projection deadline) & \emph{Non-binding}
\newline Defended by compact parameterized splines
(\(\beta < 100\text{ kB/s}\)) & \emph{Non-binding}
\newline \(P < 0.5\text{ W}\) (fixed-point constraint evaluation on MCU)
& \textbf{BINDING} \newline \(\sigma_t \le 100\text{ }\mu\text{s}\)
(cycle-deterministic packet reception on MCU) & \textbf{BINDING}
\newline Hardware boundary: private bus, separate oscillator, isolated
rail \\
\textbf{\(\mathcal{H}_5\): Enforcer \(\to\) Actuator} \emph{(Body)} &
\textbf{BINDING} \newline \(\tau \le 1.0\text{ ms}\) (servo interval
\(1.0\text{ ms}\); PWM loop \(50\text{ }\mu\text{s}\)) &
\emph{Non-binding} \newline \(\beta \approx 1\text{--}5\text{ Mbps}\)
(cyclic PDOs over deterministic fieldbus) & \textbf{BINDING}
\newline \(I^2 R\) copper winding dissipation, MOSFET switching losses &
\textbf{BINDING} \newline \(\sigma_t < 10\text{ }\mu\text{s}\) (hard
real-time PWM phase alignment \& FOC) & \textbf{BINDING}
\newline Galvanic isolation, out-of-band analog E-stops, STO circuits \\
\end{longtable}

To ensure these budgets survive the lifecycle of the machine, every
boundary condition in the grid is recorded in an engineering notebook.
The notebook serves as the single verifiable contract between hardware
designers, systems engineers, and model developers. Each entry in the
notebook follows the strict schema detailed in
Table~\ref{tbl-04-schema-notebook}, establishing mandatory fields for
physical dimension, SI units, bound category, provenance, and
verification timestamps.

\begin{longtable}[]{@{}
  >{\raggedright\arraybackslash}p{(\linewidth - 6\tabcolsep) * \real{0.2500}}
  >{\raggedright\arraybackslash}p{(\linewidth - 6\tabcolsep) * \real{0.2500}}
  >{\raggedright\arraybackslash}p{(\linewidth - 6\tabcolsep) * \real{0.2500}}
  >{\raggedright\arraybackslash}p{(\linewidth - 6\tabcolsep) * \real{0.2500}}@{}}
\caption{\textbf{Abstract Data Schema for Engineering Notebook Records
(\protect\texttt{EngineeringNotebookRecord}).} Standard schema for
documenting handoff boundary limits, physical units, provenance, and
verification state across all machine
subsystems.}\label{tbl-04-schema-notebook}\tabularnewline
\toprule\noalign{}
\begin{minipage}[b]{\linewidth}\raggedright
Field Name
\end{minipage} & \begin{minipage}[b]{\linewidth}\raggedright
Scalar Type
\end{minipage} & \begin{minipage}[b]{\linewidth}\raggedright
Physical SI Units
\end{minipage} & \begin{minipage}[b]{\linewidth}\raggedright
Structural Invariant \& Description
\end{minipage} \\
\midrule\noalign{}
\endfirsthead
\toprule\noalign{}
\begin{minipage}[b]{\linewidth}\raggedright
Field Name
\end{minipage} & \begin{minipage}[b]{\linewidth}\raggedright
Scalar Type
\end{minipage} & \begin{minipage}[b]{\linewidth}\raggedright
Physical SI Units
\end{minipage} & \begin{minipage}[b]{\linewidth}\raggedright
Structural Invariant \& Description
\end{minipage} \\
\midrule\noalign{}
\endhead
\bottomrule\noalign{}
\endlastfoot
\texttt{signal\_id} & \texttt{string} & dimensionless & Unique
alphanumeric identifier (e.g.~\texttt{SIG\_VIO\_01},
\texttt{SIG\_TRQ\_AX1}). \\
\texttt{source\_subsystem} & \texttt{string} & dimensionless &
Originating hardware or software module (\texttt{Camera\_SerDes},
\texttt{Planner\_VLA}). \\
\texttt{destination\_subsystem} & \texttt{string} & dimensionless &
Receiving hardware or software module (\texttt{State\_Estimator},
\texttt{Safety\_MCU}). \\
\texttt{physical\_dimension} & \texttt{enum} & dimensionless & Target
limit category (\(\tau\), \(\beta\), \(\mathcal{E}\), \(\sigma_t\),
\(\mathcal{F}\)). \\
\texttt{nominal\_value} & \texttt{float64} & SI units (\(\text{ms}\),
\(\text{MB/s}\), \(\text{W}\), \(\mu\text{s}\)) & Measured nominal
operating value under bench stress conditions. \\
\texttt{bound\_category} & \texttt{enum} & dimensionless &
Classification: \texttt{HARD\_MECHANICAL\_INVARIANT},
\texttt{BOUNDED\_WCET}, or \texttt{STATISTICAL\_P99}. \\
\texttt{provenance} & \texttt{string} & dimensionless & Derivation
origin (\texttt{ANALYTICAL\_PROOF}, \texttt{PHYSICS\_SIM},
\texttt{HIL\_TESTBENCH}). \\
\texttt{verification\_timestamp} & \texttt{uint64\_t} & UTC timestamp &
Epoch time of physical bench falsification and certification
sign-off. \\
\end{longtable}

For example, a concrete handoff entry for a visual-inertial state
estimation pipeline is recorded as:

\[\begin{aligned}
\big[\, &\text{Signal ID: } \text{SIG\_VIO\_01} \;\mid\; \text{Path: } \text{Camera\_SerDes} \to \text{State\_Estimator} \\
&\mid\; \text{Limit: } \text{Latency} = 24.5\text{ ms} \;\mid\; \text{Bound: } \text{Statistical } P_{99} \le 30.0\text{ ms} \\
&\mid\; \text{Provenance: } \text{HIL Bench Profiling} \;\mid\; \text{Status: } \text{VERIFIED} \,\big]
\end{aligned}\]

Provenance explicitly records whether a number originates from an
analytical first-principles calculation, a physics simulation, or direct
hardware measurement on a physical test bench. When a subsystem is
revised and a value is remeasured, the notebook updates the entry by
creating a versioned record that preserves historical measurements and
links the revision to the modified hardware or firmware configuration.
Later chapters in this text build upon this exact schema, adding
domain-specific fields for policy distributions, safety integrity
levels, and formal safety cases without restating or redefining the core
data structure.

The 5x5 Handoffs-by-Limits Grid is the master architectural contract of
the physical AI machine. Every signal crossing must have its latency,
bandwidth, thermal dissipation, determinism, and fault domain formally
bounded and recorded in the engineering notebook before physical
actuators are energized.

\section{Systems Stress-Tests}\label{sec-nervous-4-10}

A separation that is correct on an architectural block diagram is often
absent in the physical machine. In design schematics, functional modules
are partitioned by clean margins and labeled arrows, encouraging the
assumption that a deliberative compute engine and a real-time safety
enforcer occupy independent fault domains. On physical hardware,
however, these subsystems share copper traces, power distribution
networks, ground planes, interrupt controllers, and diagnostic buses.
When electrical and firmware architectures fail to enforce physical
isolation down to the circuit board floor, subtle cross-couplings create
bypass pathways that disable or circumvent the safety enforcer.

The following five analytical stress-tests evaluate where block-diagram
independence breaks down in embodied machines and define the diagnostic
protocols required to expose hidden couplings before field deployment:

\subsection*{Stress-Test 1: The Shared Rail Trap (Coupled Power
Distribution and Brownout
Resets)}\label{stress-test-1-the-shared-rail-trap-coupled-power-distribution-and-brownout-resets}
\addcontentsline{toc}{subsection}{Stress-Test 1: The Shared Rail Trap
(Coupled Power Distribution and Brownout Resets)}

\emph{The scenario.} An autonomous quadruped robot operates from a
nominal \(24.0\text{ V}\) lithium-ion battery pack with an internal
discharge state yielding an open-circuit bus voltage of
\(21.6\text{ V}\). The power distribution network routes current across
a shared bus bar with trace resistance
\(R_{\text{bus}} = 35.0\text{ m}\Omega\) and parasitic wiring inductance
\(L_{\text{bus}} = 120.0\text{ nH}\). The main board supplies both a
high-throughput neural accelerator and a hard real-time microcontroller
running the \(1000\text{ Hz}\) invariant enforcer. The microcontroller
is powered through a step-down buck converter with an undervoltage
lockout threshold of \(V_{\text{uvlo}} = 18.0\text{ V}\). During a
dynamic leaping maneuver, the neural accelerator initiates an
autoregressive vision burst, drawing a current step of
\(\Delta I_{\text{gpu}} = 25.0\text{ A}\) with a rise time of
\(\Delta t = 2.0\text{ }\mu\text{s}\), while twelve joint actuators
simultaneously accelerate, drawing a steady
\(I_{\text{motor}} = 45.0\text{ A}\).

\emph{The systems dilemma.} The total instantaneous voltage drop across
the shared power distribution network follows
\(\Delta V = I_{\text{total}} R_{\text{bus}} + L_{\text{bus}} \frac{di}{dt}\).
For \(I_{\text{total}} = 70.0\text{ A}\) and
\(di/dt = 25.0\text{ A} / (2.0\times 10^{-6}\text{ s}) = 1.25\times 10^7\text{ A/s}\),
the resistive drop is \(2.45\text{ V}\) and the inductive spike is
\(1.50\text{ V}\), yielding a total transient droop of
\(\Delta V = 3.95\text{ V}\). This sag pulls bus voltage to
\(21.6\text{ V} - 3.95\text{ V} = 17.65\text{ V}\), breaching the
\(18.0\text{ V}\) regulator threshold and causing an uncommanded
brownout reset of the safety microcontroller during peak torque
delivery.

\emph{The systems verdict: Power Isolation Invariant.} Digital
processors must never share unbuffered power rails with high-current
switching loads. Real-time safety silicon requires dedicated sub-rail
regulation, isolated power domains, or galvanic isolation to prevent
compute current spikes from resetting safety enforcers during critical
maneuvers.

\subsection*{Stress-Test 2: The Self-Feeding Watchdog (Interrupt
Servicing and Task
Starvation)}\label{stress-test-2-the-self-feeding-watchdog-interrupt-servicing-and-task-starvation}
\addcontentsline{toc}{subsection}{Stress-Test 2: The Self-Feeding
Watchdog (Interrupt Servicing and Task Starvation)}

\emph{The scenario.} An industrial mobile manipulator incorporates an
external windowed hardware watchdog timer that de-energizes the main
actuator contactor if it does not receive a service pulse every
\(10.0\text{ ms}\). The engineering team places the watchdog refresh
call inside a \(1000\text{ Hz}\) hardware timer interrupt service
routine (ISR) running at the highest hardware priority on the real-time
microcontroller. The safety-critical invariant enforcement task runs as
a lower-priority RTOS thread. While evaluating a workspace keep-out
boundary, the enforcement task deadlocks on a shared mutex.

\emph{The systems dilemma.} With the enforcement thread blocked, the
controller ceases updating joint torque limits, validating obstacle
clearance, or checking intent lease timestamps. However, because the
hardware timer ISR executes independently of thread scheduling, it
continues firing every \(1.0\text{ ms}\), consistently petting the
watchdog and preventing the supervisor from tripping the emergency power
relay while the manipulator moves uncontrolled toward a hardstop.

\emph{The systems verdict: Execution-Verified Watchdog Law.} A watchdog
must monitor the physical completion of safety invariants, not the
electrical ticking of a hardware timer. Servicing a watchdog from a
timer ISR certifies only that the clock generator is alive while masking
application deadlock.

\subsection*{Stress-Test 3: The Bypass Backdoor (Diagnostic Interfaces
and Unmediated
Authority)}\label{stress-test-3-the-bypass-backdoor-diagnostic-interfaces-and-unmediated-authority}
\addcontentsline{toc}{subsection}{Stress-Test 3: The Bypass Backdoor
(Diagnostic Interfaces and Unmediated Authority)}

\emph{The scenario.} An autonomous forklift enforces workspace safety
through an isolated microcontroller that validates all velocity
proposals before dispatching current setpoints over a CAN bus. For
factory testing and motor calibration, the wheel drive units include an
auxiliary serial diagnostic interface accessible over an internal
service socket. During an edge runtime failure on the host computer, a
diagnostic daemon writes direct register overrides to the motor drives
to read rotor thermal resistance.

\emph{The systems dilemma.} The diagnostic write command overwrites the
inverter current target register with an unvalidated test value of
\(35.0\text{ A}\), commanding forward torque that accelerates the
vehicle to \(1.8\text{ m/s}\). The primary safety enforcer, reading zero
velocity on its command mailbox, detects no safety violation in the
proposal stream and leaves its emergency brake relays open, moving the
machine without kinematic gating or obstacle avoidance.

\emph{The systems verdict: Exclusive Physical Authority Invariant.}
Actuation authority is the sum of every physical wire and register that
can command power stages. Auxiliary diagnostic ports must be physically
interlocked with hardware jumpers, blown fuses, or power disconnect
relays during operational mode.

\subsection*{Stress-Test 4: The Phantom Feedback Loop (Command Echo
versus Measured
State)}\label{stress-test-4-the-phantom-feedback-loop-command-echo-versus-measured-state}
\addcontentsline{toc}{subsection}{Stress-Test 4: The Phantom Feedback
Loop (Command Echo versus Measured State)}

\emph{The scenario.} A six-axis articulated manipulator
(\(m = 90.0\text{ kg}\)) executes a discrete Control Barrier Function at
\(1000\text{ Hz}\) to guarantee that the end-effector remains outside a
forbidden zone (\(h(\mathbf{x}) \ge 0\)). To minimize serial bus traffic
and filtering latency, the enforcer evaluates the barrier invariant
using the commanded joint position vector \(\mathbf{q}_{\text{cmd}}\)
from the trajectory generator rather than the measured optical encoder
vector \(\mathbf{q}_{\text{meas}}\) sampled at the load side.

\emph{The systems dilemma.} While carrying a \(20.0\text{ kg}\) payload
at \(\dot{\mathbf{q}} = 1.5\text{ rad/s}\), a cycloidal drive gear tooth
fractures, mechanically decoupling the motor rotor from the link.
Gravity accelerates the link downward toward an operator workstation at
\(9.81\text{ m/s}^2\). Because the trajectory interpolator continues
generating valid setpoints within bounds, the enforcer evaluates
\(h(\mathbf{q}_{\text{cmd}}) > 0\) and reports nominal operation,
holding emergency brakes disengaged while the physical link falls
\(0.35\text{ m}\) onto the workstation.

\emph{The systems verdict: Physical State Invariant.} An invariant
enforcer that evaluates software commands rather than physical sensor
telemetry validates its own intentions rather than reality. Safety
enforcement requires an unbroken, isolated sensory channel to load-side
physical transducers.

\subsection*{Stress-Test 5: The Diagnostic Bench Protocol
(Hardware-in-the-Loop Fault
Injection)}\label{stress-test-5-the-diagnostic-bench-protocol-hardware-in-the-loop-fault-injection}
\addcontentsline{toc}{subsection}{Stress-Test 5: The Diagnostic Bench
Protocol (Hardware-in-the-Loop Fault Injection)}

\emph{The scenario.} An engineering team prepares an autonomous mobile
manipulator for factory qualification, executing a four-stage diagnostic
fault-injection suite prior to release: power rail transient injection
(\(di/dt = 150.0\text{ A}/\mu\text{s}\)), communication bus packet
flooding, memory bit-flips in high-level policy runtimes, and mechanical
shaft locking on a dynamometer.

\emph{The systems dilemma.} Nominal operating logs cannot prove physical
independence. Pushing the machine across electrical, computational, and
mechanical failure boundaries verifies that voltage sag maintains a
\(20\%\) margin above brownout thresholds, communication flooding
induces zero deadline jitter on the safety loop, and locked-rotor
divergence exceeding \(\Delta \theta = 1.5^\circ\) clamps motor current
within \(2.0\text{ ms}\).

\emph{The systems verdict: Catastrophic Isolation Law.} A safety case is
valid only if the real-time invariant enforcer survives the catastrophic
collapse of every surrounding software and electrical component.

\newpage{}

\section{Summary}\label{sec-nervous-4-11}

The architecture constructed across Part I establishes the mechanical
and computational boundary conditions for physical autonomy. A physical
AI system begins with the body, whose mass, inertia, friction, and
thermal dissipation define the non-negotiable laws of the physical
plant. Above this physical plant sits the brain, composed of
high-capacity learned models that ingest high-dimensional sensory
observations and generate stochastic candidate plans across extended
horizons. Between these two tiers sits the nervous system, operating as
a deterministic arbiter on isolated hardware. Where the brain is
unverified and statistical, the nervous system is bounded and
predictable. It measures plant telemetry directly from isolated sensors,
evaluates candidate actions against a hard invariant envelope, and
guarantees that actuator commands remain strictly within the thermal,
electrical, and kinematic limits of the machine.

The central rule governing this three-part division is that proposal is
never permission. A learned policy never drives motor windings or opens
hydraulic valves directly. Instead, the policy submits candidate
trajectories through time-bounded intent leases, treating high-level
compute as an untrusted client requesting access to physical actuators.
Execution rates across the machine follow directly from physical time
constants rather than computational convenience. A brushless motor drive
must update its field-oriented current loop at \(20\text{ kHz}\) to
prevent winding saturation, and a joint position controller must close
its loop at \(1000\text{ Hz}\) to track commanded paths against
unexpected load disturbances. By contrast, a transformer-based
vision-language-action model running on an accelerator may take
\(50\text{ ms}\) to process an image pair and output an action chunk.
The nervous system bridges this three-order-of-magnitude frequency
disparity by interpolating valid chunks into smooth reference
trajectories, continually monitoring tracking error, and asserting local
control whenever the lease expires or the state diverges.

With the physical boundaries, timing contracts, and arbitration hardware
established, the machine no longer risks hardware destruction when
executing experimental software. Part I treated the learned policy as a
black box, assuming only that it generates stochastic proposals at low
frequency and occasionally fails to meet its deadline or produces an
invalid trajectory. Part II opens this black box to examine how learned
policies are constructed, trained, and validated. The systems question
shifts from who is permitted to move the motors to how the machine
acquires the capacity to propose sensible physical actions in the first
place.

Part II populates the architecture grid derived in
Chapter~\ref{sec-nervous}, addressing the deliberate layers of the
system row by row:

In Chapter~\ref{sec-data}, we examine the data systems that feed the
policy, tracing the end-to-end pipeline required to collect,
synchronize, and represent physical demonstrations. It treats human
teleoperation, multisensory logging, and trajectory curation not as
abstract datasets, but as high-throughput storage and network workloads
where sensor timestamp skew, dropped frames, and proprioceptive drift
introduce systematic errors into downstream training.

In Chapter~\ref{sec-training}, we address the computational and physical
physics of training and simulation. It analyzes what it costs to train
policies over millions of environment steps, contrasting the
computational throughput of hardware-accelerated physics simulators
against the reality gap of unmodeled contacts, friction hysteresis, and
actuator dynamics.

In Chapter~\ref{sec-evaluation}, we complete the policy foundation by
establishing how to evaluate a learned model before it receives
authority over physical hardware. Because statistical average-case
metrics such as validation loss or mean reward hide catastrophic failure
modes, the chapter defines systems-oriented evaluation protocols,
covering out-of-distribution detection, state-space coverage analysis,
and stress-testing methodologies designed to measure tail-risk
performance under synthetic perturbations.

\protect\phantomsection\label{callout-takeawaysux2a-4.9}
\begin{fbxSimple}{callout-takeaways}{Key Chapter Takeaways}{}
\phantomsection\label{callout-takeaways*-4.9}

\begin{itemize}
\tightlist
\item
  \textbf{Proposal Is Never Permission:} High-level learned models act
  strictly as unprivileged proposal engines (\(p_t\)). Physical write
  privilege over motor drive registers is reserved exclusively for the
  deterministic real-time enforcer on isolated safety silicon.
\item
  \textbf{Cadence is Dictated by Physics, Not Software:} Execution
  frequencies across the four tiers (\(20\text{ kHz}\) current loops,
  \(1000\text{ Hz}\) barrier enforcement, \(20\text{--}50\text{ Hz}\)
  action chunking, \(1\text{--}5\text{ Hz}\) deliberation) derive
  directly from winding time constants (\(\tau_e = L/R\)), dynamic phase
  margins, and sensory exposure windows.
\item
  \textbf{Lock-Free Sequence Locks Guarantee Zero-Wait Memory
  Transport:} Asynchronous data exchange across rate boundaries is
  mediated by single-producer single-consumer (SPSC) seqlocks in shared
  static memory, ensuring read operations on the \(1000\text{ Hz}\) core
  execute in \(<150\text{ ns}\) without mutex contention, priority
  inversion, or dynamic heap allocation.
\item
  \textbf{Structural Timing Guarantees Outrank Empirical Distributions:}
  A statistical latency percentile (\(P_{99}\)) cannot prove real-time
  safety. Deterministic certification requires analytical Response Time
  Analysis (\(R_i = C_i + B_i + I_i \le D_i\)), eliminating unbounded
  bus contention and memory allocations from the safety path.
\item
  \textbf{Four Dimensions of Hardware Isolation:} Dependable safety
  enforcement requires independent execution silicon, segregated timing
  crystals, isolated power rails, and dedicated memory buses to prevent
  catastrophic faults from crossing domain boundaries.
\item
  \textbf{Stopping Dynamics Dimension Watchdog Horizons:} Watchdog lease
  durations (\(\tau_{\text{lease}} \le 87.6\text{ ms}\)) are derived
  from kinetic stopping envelopes
  (\(d_{\text{stop}} \le D_{\text{clear}}\)), bounding the maximum
  unguided displacement traversed during an uncommanded compute stall.
\item
  \textbf{The 5x5 Synthesis Matrix:} The 5x5 Handoffs-by-Limits Grid
  establishes the diagnostic rubric mapping the five signal handoffs
  (\(\mathcal{H}_1\text{--}\mathcal{H}_5\)) against five physical limits
  (\(\tau, \beta, \mathcal{E}, \sigma_t, \mathcal{F}\)).
\end{itemize}

\end{fbxSimple}

\part{Teaching the Machine}

\lettrine{I}{n} conventional machine learning, data is passive and
computational cost is paid in cloud kilowatt-hours. Datasets are scraped
from the internet, curated, partitioned into training and validation
splits, and treated as an exogenous ground truth
(\(x \sim P_{\text{data}}\)). A vision classifier or language model
inspects this data from behind glass; its predictions have no physical
momentum and no causal power to alter the distribution of the documents
it will encounter tomorrow. If a classifier mispredicts a label, the
error is isolated---the physical world remains completely unaffected.

In Physical AI, this assumption collapses entirely. The data
distribution of a physical machine is \textbf{endogenous}: the machine
participates in a closed causal loop where every motor command changes
the configuration of the physical plant, directly dictating what its
sensors will observe next (\(s_{t+1} \sim P(s \mid s_t, a_t)\)).

Because the policy creates its own future sensory inputs, training an
embodied model is not an exercise in curve-fitting; it is the deliberate
management of \textbf{latent failure debts}:

\begin{itemize}
\tightlist
\item
  \textbf{Data as Physical Experience (Chapter 5):} Every demonstration
  collected in the physical world consumes mechanical life---wearing
  cycloidal gearboxes, draining battery energy, heating motor stators,
  and accumulating sensor calibration drift. We examine the systems
  physics of bilateral teleoperation, multi-camera SerDes timestamp
  synchronization, and the Four-Stream Action Tap Tuple. We prove why
  logging human recovery interventions is mathematically required to
  prevent compounding covariate drift (\(O(T^2 \epsilon)\)) and keep the
  machine on its demonstration manifold.
\item
  \textbf{Policy Training and Latent Debts (Chapter 6):} Formulating
  policy architectures that can generate high-dimensional continuous
  trajectories without suffering mode collapse. We analyze the
  fundamental policy synthesis trade-offs across Behavioral Cloning,
  Offline RL, and Sim-to-Real RL. We explain why naive Mean Squared
  Error (MSE) fails on multimodal human actions (averaging two valid
  paths into an obstacle) and why Action Chunking with Transformers
  (ACT) and Diffusion Policies succeed by predicting continuous
  trajectory splines. Crucially, we decompose the sim-to-real reality
  gap into friction hysteresis, structural compliance, and motor phase
  lag.
\item
  \textbf{Evaluation as Closed-Loop Reality (Chapter 7):} Confronting
  why offline test-set accuracy is fundamentally decoupled from physical
  survival. A policy achieving \(99\%\) open-loop accuracy on a static
  validation set can fail within seconds when closed in a real feedback
  loop due to non-Lipschitz sensitivity and compounding error. We
  establish the rigorous systems discipline of closed-loop
  evaluation---deriving non-asymptotic Clopper-Pearson coverage bounds,
  Butler \& Finelli exposure limits, and Extreme Value Theory to
  evaluate tail-risk safety.
\end{itemize}

Part II teaches students how to curate physical data, train generative
action policies, and measure closed-loop survival without falling into
the trap of superficial benchmark optimism.

\chapter{\texorpdfstring{Physical Data
Collection\chaptersubtitle{Where a Physical Policy's Training Data Comes From}}{Physical Data Collection}}\label{sec-data}

\begin{figure}[H]

\centering{

\includegraphics[width=1\linewidth,height=\textheight,keepaspectratio]{index_files/mediabag//Users/VJ/GitHub/MLSysBook-integrate-physical/physical/book/chapters/05-data/figures/fig_locator.pdf}

}

\caption{\label{fig-locator}\textbf{Demonstration Acquisition Cost
Scales with Physical Embodiment.} In the physical AI architecture, data
generation feeds the learned Brain organ through physical movement of
the Body. Because every recorded transition consumes wall-clock time,
manual operator attention, and hardware fatigue life, dataset volume is
bounded by physical wear and reset overhead rather than storage
capacity.}

\end{figure}%

\emph{What does the physical machine that collected a dataset leave
permanently stamped into the policy that trains on it?}

In digital machine learning, data is a static asset scraped in bulk and
curated offline. In physical AI, data is an endogenous trace of physical
embodiment coupled to the exact kinematics, sensor latencies, and
mechanical backlash of the recording platform. Teleoperated
demonstrations capture narrow nominal corridors while omitting recovery
dynamics, inducing compounding covariate drift under physical
disturbances. Furthermore, data collected on one machine carries an
idiosyncratic dynamical signature that cannot transfer unconditionally
to another. Building dependable physical AI requires treating
demonstration data as an expensive, perishable asset governed by
explicit kinematic contracts and synchronous multi-stream telemetry.

\newpage{}

\protect\phantomsection\label{callout-objectiveux2a-5.1}
\begin{fbxSimple}{callout-objective}{Learning Objectives}{}
\phantomsection\label{callout-objective*-5.1}

After studying this chapter, you will be able to:

\begin{enumerate}
\def\labelenumi{\arabic{enumi}.}
\tightlist
\item
  Account for what one hour of demonstration data actually costs, in
  wall-clock time, usable trajectory yield, and irreversible mechanical
  wear.
\item
  Explain why the collecting policy determines dataset coverage, derive
  the quadratic compounding error of behavioral cloning
  (\(O(T^2 \epsilon)\)) against the linear bound of corrective
  aggregation (\(O(T \epsilon)\)), and explain why an episode ending in
  a rescue must be cut where the policy diverged (\(t_{\text{div}}\))
  rather than where the human intervened (\(t_{\text{takeover}}\)).
\item
  Identify which properties of a dataset are bound to the embodiment
  that produced it, and evaluate the architectural challenges of
  cross-embodiment datasets such as Open X-Embodiment.
\item
  Formalize demonstration provenance across the Dual-Brain MPU/MCU
  boundary, separating requested, mapped, enforced, and measured actions
  while establishing explicit operational design domain (ODD) support
  boundaries.
\end{enumerate}

\emph{Core chapter takeaway.} A physical dataset is made by moving a
machine, so every sample costs wall-clock time and wear, and the policy
that collected it decides what is in it. There is no scraping your way
out of this.

\end{fbxSimple}

\section{Endogenous Experience}\label{sec-data-5-1}

\lettrine{W}{ithin}\lettrine{t}{he} physical AI architecture
(Figure~\ref{fig-locator}), data acquisition forms the physical
substrate feeding the Brain organ, bridging the unyielding mechanics of
the Body with the statistical representations of learned policies. In
digital machine learning, datasets exist independently of the model:
text corpora, image archives, and audio waveforms are scraped,
tokenized, and shuffled with zero impact on the physical world. In
Physical AI, dataset generation is an endogenous physical process
governed by the closed causal loop \(s_{t+1} \sim P(s \mid s_t, a_t)\).
Every recorded frame, joint angle, and force-torque vector requires
commanding electrical currents into motor coils, overcoming frictional
stiction, and pushing physical mass through space.

Because every training sample requires physical movement, demonstration
data cannot be scaled by scraping or passive logging. A single hour of
scheduled facility time on a robotic workcell yields only a fraction of
active, usable trajectory records once startup overhead, sensor
calibration routines, human fatigue, and physical scene resets are
accounted for. Furthermore, because a robot cannot observe an
environment without executing actions, the recorded data distribution is
strictly closed under the control policy of the collecting agent. If an
operator or scripted routine navigates only narrow, collision-free
corridors, the dataset remains entirely silent on how to recover from
boundary excursions, baking severe compounding drift into downstream
imitation policies.

Training dependable physical AI systems requires treating demonstration
data as an expensive, perishable physical asset with explicit kinematic
contracts and documented operational support boundaries. Characterizing
a physical dataset demands accounting for four fundamental systems
dimensions:

\begin{enumerate}
\def\labelenumi{\arabic{enumi}.}
\tightlist
\item
  \textbf{Physical Acquisition Cost and Fatigue Life:} The true
  accounting of scheduled facility time, active motion duty cycles,
  usable episode yield, and mechanical wear on transmissions and
  actuators.
\item
  \textbf{Collector-Induced Support Manifolds:} The mathematical reality
  that the collecting policy dictates state-space coverage, creating the
  quadratic error growth of naive behavioral cloning
  (\(O(T^2 \epsilon)\)) and dictating where intervention episodes must
  be truncated.
\item
  \textbf{Embodiment Idiosyncrasy:} The physical coupling between
  kinematic geometry, sensor latencies, actuator dynamics, and mass
  distributions that prevents uncalibrated transfer across different
  robot platforms.
\item
  \textbf{Boundary Support and Provenance Contracts:} The multi-tier
  action tap tuple
  (\((a_{\text{req}}, a_{\text{cmd}}, a_{\text{enf}}, a_{\text{meas}})\))
  across the MPU/MCU boundary and the explicit operational design domain
  (ODD) ledger defining where data provides physical support versus
  where execution must be refused.
\end{enumerate}

The mathematical formulation of training algorithms, simulation
environments, and domain randomization belongs to
Chapter~\ref{sec-training}, while the statistical methods required to
verify whether a policy behaves safely on physical hardware belong to
Chapter~\ref{sec-evaluation}. Before a policy can be trained or
evaluated, physical records must exist, and every sample in the archive
begins with mass moving under power in the real world across mobile
systems (burning battery capacity and scrubbing tire treads),
articulated manipulators (cycling gearbox teeth and flexing internal
harnesses), or process plants (subjecting piping to thermal stresses and
valve packing wear).

\textbf{Endogenous physical demonstration data} is the empirical
sequence of sensory observations and control actions generated by
physical movement of an embodied machine
(\(s_{t+1} \sim P(s \mid s_t, a_t)\)). Because every transition requires
physical displacement, data volume is strictly bounded by wall-clock
facility time, manual reset overhead, and component fatigue life, while
state-space coverage is closed under the dynamical authority of the
collecting policy.

\section{Where Physical Data Comes From}\label{sec-data-5-2}

A physical dataset is not found, logged from passive internet traffic,
or assembled from an existing public repository. Every transition tuple
in an archive required burning electrical energy, pushing against
mechanical friction, and moving physical mass through a coordinate
frame. Physical data sources separate according to who chooses the
action:

\begin{enumerate}
\def\labelenumi{\arabic{enumi}.}
\tightlist
\item
  In \textbf{kinesthetic teaching}, a colocated demonstrator physically
  grasps and moves the unpowered or gravity-compensated joints of the
  mechanism.
\item
  In \textbf{teleoperation}, a remote operator guides an input device
  such as a master arm, a multi-axis space mouse, or a tracked hand
  controller, transmitting commands across a software interface to the
  follower robot.\sidenote{\footnotesize Raymond Goertz developed the first bilateral
    mechanical teleoperation manipulators at Argonne National Laboratory
    between 1948 and 1954 to handle radioactive materials, establishing
    the kinematic master-follower architecture still used in robot
    demonstration capture today.}
\item
  In \textbf{scripted collection}, an engineer executes a deterministic
  control routine, such as a state machine tracking interpolated splines
  or force-thresholded primitives.
\item
  In \textbf{exploratory collection}, an active policy selects actions
  to optimize an information-theoretic objective, sample state space, or
  maximize a reward signal.
\item
  In \textbf{autonomous logging during deployment}, the production
  policy commands the actuators as it performs work in a factory or
  warehouse, storing incoming sensor streams in situ.
\end{enumerate}

For each collection source, the engineer must decide which physical or
digital variable is recorded as the action label, and that choice alters
what the learned model attempts to reproduce. In teleoperation and
autonomous systems, data logging taps into multiple layers across the
control pipeline. On heterogeneous dual-brain platforms, such as a
canonical heterogeneous dual-brain System-on-Chip (SoC) bench twin or a
bimanual leader-follower teleoperation workstation
(\citeproc{ref-zhao2023learning}{Zhao et al. 2023};
\citeproc{ref-wu2023gello}{Wu et al. 2024}),\sidenote{\footnotesize Canonical dual-arm
  bimanual teleoperation workcells (\citeproc{ref-wu2023gello}{Wu et al.
  2024}; \citeproc{ref-zhao2023learning}{Zhao et al. 2023}) utilize
  3D-printed, kinematically isomorphic leader arms with digital joint
  encoders to provide direct joint-space retargeting without inverse
  kinematics singularities.} this pipeline spans two distinct
computational domains:

\begin{itemize}
\tightlist
\item
  The unprivileged \textbf{Linux Application Processor (MPU)} captures
  high-bandwidth camera streams (e.g., active stereo RGB-D at
  \(30\text{ Hz}\)), ingests master teleoperation inputs or neural
  policy forward passes, and maps coordinates.
\item
  The real-time \textbf{Microcontroller Unit (MCU)} executes at
  \(1000\text{ Hz}\) on bare metal, running low-level PID/impedance
  control, forward-invariant safety barrier functions, velocity
  clamping, and motor current limiters over CAN-FD or PWM
  buses.\sidenote{\footnotesize On real-time MCU architectures, \(a_{\text{meas}}\) is
    captured via dedicated quadrature encoder peripheral registers
    (e.g., \texttt{TIMx\_CNT} / \texttt{eQEP}) and phase-current shunt
    ADCs triggered synchronously by PWM counter reload events
    (\texttt{ePWM\ CMPA}), with telemetry packed into DMA ring buffers
    over CAN-FD (\texttt{FDCAN}) to eliminate OS interrupt latency
    jitter.}
\end{itemize}

\begin{figure}[t]

\centering{

\includegraphics[width=1\linewidth,height=\textheight,keepaspectratio]{index_files/mediabag//Users/VJ/GitHub/MLSysBook-integrate-physical/physical/book/chapters/05-data/figures/fig05_four_stream_action_taps.pdf}

}

\caption{\label{fig-05-action-taps}\textbf{The Four-Stream Action Tap
Tuple across the Dual-Brain MPU/MCU Boundary.} On a heterogeneous edge
architecture, data logging captures four distinct physical semantics:
(1) \(a_{\text{req}}\) (raw operator or policy intent on the Linux MPU
prior to software checks), (2) \(a_{\text{cmd}}\) (kinematically
retargeted setpoint), (3) \(a_{\text{enf}}\) (safety-governed command
post-CBF barrier enforcement and motor torque clamping on the real-time
MCU), and (4) \(a_{\text{meas}}\) (measured plant response via joint
optical encoders and current shunts). Logging only \(a_{\text{req}}\)
bakes unexecutable actuator drift into training loss, while logging only
\(a_{\text{meas}}\) bakes mechanical compliance, link deflection, and
backlash as demonstrated human skill. Source: Architectural design
adapted from Zhao et al. (\citeproc{ref-zhao2023learning}{2023}) and Wu
et al. (\citeproc{ref-wu2023gello}{2024}).}

\end{figure}%

The four taps along this control path capture fundamentally different
physical semantics:

\begin{itemize}
\tightlist
\item
  The \textbf{requested action} (\(a_{\text{req}}\)): the raw input from
  the operator's tracking device or policy proposal before coordinate
  transformations or software boundary checks.
\item
  The \textbf{mapped command} (\(a_{\text{cmd}}\)): the Cartesian
  setpoint or joint target computed after kinematic retargeting and
  workspace boundary checks.
\item
  The \textbf{enforced command} (\(a_{\text{enf}}\)): the setpoint after
  real-time safety filters, velocity limits, and barrier functions have
  clamped the candidate trajectory.
\item
  The \textbf{measured response} (\(a_{\text{meas}}\)): the actual
  kinematic and kinetic response reported by joint optical encoders and
  current shunts after mechanical contact and structural deflection have
  occurred.
\end{itemize}

If the logging pipeline records \(a_{\text{req}}\) prior to software
limit enforcement, the dataset contains dynamically infeasible
trajectories that command accelerations beyond motor current limits. If
it logs \(a_{\text{meas}}\) after contact, the recorded label reflects
structural compliance and gear backlash rather than commanded intent,
turning a feedforward control signal into an observation of material
deformation. Logging the full tuple
\((a_{\text{req}}, a_{\text{cmd}}, a_{\text{enf}}, a_{\text{meas}})\)
across the MPU/MCU boundary is necessary to distinguish operator intent
from governor intervention.

\protect\phantomsection\label{callout-definitionux2a-5.2}
\begin{fbxSimple}{callout-definition}{Definition:}{The Four-Stream Action Tap Tuple}
\phantomsection\label{callout-definition*-5.2}

The \textbf{four-stream action tap tuple}
\((a_{\text{req}}, a_{\text{cmd}}, a_{\text{enf}}, a_{\text{meas}})\) is
the synchronized multi-tier record logged across the Dual-Brain MPU/MCU
boundary that distinguishes raw operator or policy intent
(\(a_{\text{req}}\)), kinematically retargeted setpoints
(\(a_{\text{cmd}}\)), safety-governor clamped commands
(\(a_{\text{enf}}\)), and delivered physical actuator response
(\(a_{\text{meas}}\)), preventing governor interventions or structural
compliance from being falsely attributed to demonstrated physical skill.

\end{fbxSimple}

\subsection{Synchronous 4-Stream Action Tap and Telemetry Ring-Buffer
Protocol}\label{synchronous-4-stream-action-tap-and-telemetry-ring-buffer-protocol}

To operationalize the four-stream action tap in hardware without
introducing non-deterministic logging jitter or buffer contention, the
system executes a synchronized inter-core capture loop:

\textbf{Input Interfaces:}

\begin{itemize}
\tightlist
\item
  Raw observation stream \(\mathbf{o}_t \in \mathcal{O}\) from the host
  perception pipeline (active stereo RGB-D at \(30\text{ Hz}\), hardware
  timestamp \(t_{\text{exp}}\)).
\item
  Raw operator input or policy action candidate
  \(a_{\text{req}}(t) \in \mathcal{A}_{\text{raw}}\) from MPU shared
  memory (\(50\text{ Hz}\)).
\item
  Kinematic retargeting and workspace boundary mapping
  \(\mathcal{K}: \mathcal{A}_{\text{raw}} \to \mathcal{A}_{\text{cartesian}}\).
\item
  Forward-invariant Control Barrier Function (CBF) and velocity limiter
  \(\mathcal{B}: \mathcal{X} \times \mathcal{A}_{\text{cartesian}} \to \mathcal{A}_{\text{safe}}\)
  on the real-time MCU.
\item
  Bare-metal hardware registers for quadrature encoder counts
  (\texttt{TIMx\_CNT} / \texttt{eQEP}) and phase current shunts
  (\texttt{ADC\_DR}).
\end{itemize}

\textbf{Execution Protocol (Periodic \(1.0\text{ ms}\) MCU Cycle):}

\begin{enumerate}
\def\labelenumi{\arabic{enumi}.}
\item
  \textbf{Perception and Intent Ingestion (Host MPU Domain):} Ingest
  synchronized camera frame \(\mathbf{o}_t\) with hardware exposure
  timestamp \(t_{\text{exp}}\), sample raw master input
  \(a_{\text{req}}(t)\), evaluate forward kinematics
  \(a_{\text{cmd}}(t) \leftarrow \mathcal{K}(a_{\text{req}}(t))\), and
  commit
  \(\langle t_{\text{exp}}, \mathbf{o}_t, a_{\text{req}}, a_{\text{cmd}} \rangle\)
  to shared memory DMA TX descriptors.
\item
  \textbf{Safety Barrier Evaluation and Actuation (MCU Domain,
  \(1000\text{ Hz}\) ISR):} Latch hardware monotonic clock
  \(t_{\text{mono}} \leftarrow \text{ReadTimerRegister}(\text{TIMx\_CNT})\),
  ingest \(a_{\text{cmd}}(t)\) across inter-core SPSC buffers with
  memory fence (\texttt{DMB\ ISH}), and evaluate dynamic barrier
  constraints:
  \[a_{\text{enf}}(t) \leftarrow \mathcal{B}(\mathbf{x}_{\text{current}}, a_{\text{cmd}}(t)).\]
  If \(a_{\text{enf}}(t) \neq a_{\text{cmd}}(t)\), assert governor
  intervention flag \(b_{\text{clamped}} \leftarrow \text{true}\) with
  clamp code \(\kappa\), and dispatch \(a_{\text{enf}}(t)\) to motor PWM
  registers (\texttt{ePWM\ CMPA}) over the fieldbus.
\item
  \textbf{Physical Response Capture and Commit:} Sample optical joint
  encoders and phase current shunts at PWM reload
  (\(a_{\text{meas}}(t) \leftarrow \text{ReadPlantSensors}()\)), compute
  tracking discrepancy
  \(\delta_{\text{track}} \leftarrow \|a_{\text{enf}}(t) - a_{\text{meas}}(t)\|_2\),
  assemble synchronous telemetry packet
  \(\mathcal{P}_t = \langle t_{\text{mono}}, \mathbf{o}_t, a_{\text{req}}, a_{\text{cmd}}, a_{\text{enf}}, a_{\text{meas}}, \mathbf{h}_{\text{plant}}, [b_{\text{clamped}}, \kappa, \delta_{\text{track}}] \rangle\),
  and commit \(\mathcal{P}_t\) into a lock-free DMA circular ring buffer
  for non-blocking asynchronous persistence.
\end{enumerate}

Because the collecting entity dictates how the machine approaches the
environment, kinesthetic teaching, teleoperation, unguided play, and
autonomous logging occupy disjoint state-action manifolds even when
attempting the exact same physical contact task. Consider inserting a
cylindrical steel pin into a chamfered hole with a clearance of
\(50\text{ }\mu\text{m}\). A colocated human guiding the tool directly
senses contact reaction forces through their hands and makes
instantaneous, fine-grained compliance adjustments at their own wrist
joints. A remote teleoperator watching a two-dimensional video feed
lacks direct haptic feedback, creating visual latency delays, command
overshoot, and characteristic hunting motions as the operator pecks at
the contact boundary. A scripted controller follows a programmed helical
search pattern at constant linear velocity, generating a dense geometric
spiral with predictable force spikes against the chamfer face.
Autonomous logging from an exploratory or partially trained policy
wanders across off-target regions, repeatedly binding the pin at oblique
angles until motor current thresholds trigger a software fault. These
logs do not constitute noisy samples of a single underlying trajectory.
They represent outputs from distinct dynamical systems operating under
different closed-loop bandwidths, latency profiles, and mechanical
compliances.

A common accounting error in physical machine learning is treating every
recorded control step as an independent unit of experience. When a joint
controller logs state-action pairs at \(50\text{ Hz}\), one hour of
continuous recording produces \(180{,}000\text{ steps}\). Multiplying
sample frequency by recording time yields large step counts, but
adjacent steps exhibit extreme temporal autocorrelation. Two camera
frames logged \(20\text{ ms}\) apart differ by sub-millimeter
end-effector displacements and minor sensor noise. Distributional
coverage scales with the count of independent task episodes and distinct
physical configurations, not with the sampling clock. An archive
containing \(100{,}000\text{ steps}\) distributed across five hundred
distinct workpiece positions and lighting conditions provides vastly
broader support than \(1{,}000{,}000\text{ steps}\) logged during a
single continuous demonstration where the end-effector engaged a
workpiece at a single bench location.

Evaluating the throughput of a physical collection campaign requires
calculating the usable episode yield from the mechanical and operational
parameters of the workcell. Consider an analytical example of an
articulated manipulator collecting contact demonstrations for a
precision assembly task. Let an attempt require an average duration
\(t_{\text{att}} = 20\text{ s}\). After each attempt, the cell must
return to a valid initial state, requiring the manipulator to transit
home and an automated fixture to reset the workpiece, adding a reset
duration \(t_{\text{rst}} = 15\text{ s}\). The total cycle time for a
single attempt is
\(t_{\text{cyc}} = t_{\text{att}} + t_{\text{rst}} = 35\text{ s}\). In
one wall-clock hour (\(3600\text{ s}\)), a single rig completes at most
\(102\text{ attempts}\). Physical collection suffers from trial failures
where workpieces slip, parts jam, or safety watchdogs trip on excessive
contact force. If the trial failure rate is \(25\%\), and
post-collection validation rejects an additional \(10\%\) of logs due to
camera occlusion or dropped network packets, the retention rate is:
\[R = (1 - 0.25) \times (1 - 0.10) = 0.675.\] The usable episode yield
for a single rig is:
\[Y_{\text{usable}} = \left\lfloor\frac{3600\text{ s}}{t_{\text{cyc}}}\right\rfloor \times R = 102 \times 0.675 \approx 68.8\text{ usable episodes per hour}.\]

If training a policy requires \(5000\text{ validated episodes}\) to
achieve target success rates across geometric variations, collecting
that data on one workstation requires \(72.7\text{ hours}\) of active
execution time. Across those \(72.7\text{ hours}\), the machine executes
\(7415\text{ physical cycles}\), subjecting gearbox teeth to repeated
directional reversals, stressing flex cables inside articulated joints
through thousands of bends, and accumulating thermal load in motor
windings under sustained stall currents. Operating three parallel
collection rigs increases the throughput to approximately
\(206.4\text{ usable episodes per hour}\) and compresses calendar time
to \(24.2\text{ hours}\), but it scales hardware wear and maintenance
linearly across three sets of mechanical transmissions and sensors.

Simulated trajectories generated in rigid-body physics engines avoid
direct hardware wear and can generate millions of transitions per hour
across distributed compute clusters. The mechanics, approximations, and
domain transfer boundaries of simulated data are treated in
Chapter~\ref{sec-training}. When synthetic trajectories are combined
with physical logs, the physical provenance of every record must be
tracked explicitly. A model trained on an undifferentiated mixture
inherits the contact unfaithfulness of the physics engine wherever
synthetic data dominates, while inheriting the unmodeled sensor noise
and mechanical backlash of the hardware wherever physical data
dominates. Without tracking the collecting agent and the recording tap
point for each sample, an engineer cannot determine whether an on-robot
failure stems from simulation error or from physical collection bias.
Human teleoperation remains the standard technique for gathering
high-quality physical trajectories, yet the interface between the human
and the machine shapes the recorded dynamics long before the data
reaches a model.

\section{Teleoperation and What It Costs}\label{sec-data-5-3}

Human teleoperation closes a feedback loop that spans the operator, the
input device, the communication bus, the low-level motor controller, and
the physical environment (Figure~\ref{fig-05-teleop-rig}). In a typical
teleoperation workstation, an image sensor captures the workspace,
encodes the video stream, and transmits the frames across a network to a
display console. The human operator perceives the visual feedback,
decides on a corrective trajectory, and applies physical displacement to
an input device. The workstation samples the device state, maps the
spatial displacement into joint or end-effector coordinates, and
transmits the resulting setpoint across an industrial bus to the robot's
motor drives. The actuators generate torque, accelerating the mechanical
linkages against inertia, joint friction, and contact constraints. The
resulting displacement alters the physical scene, producing new optical
signals that the image sensor records on the next cycle.

\begin{figure}[t]

\centering{

\includegraphics[width=0.95\linewidth,height=\textheight,keepaspectratio]{/Users/VJ/GitHub/MLSysBook-integrate-physical/physical/book/chapters/05-data/figures/fig05_real_gello_teleop.png}

}

\caption{\label{fig-05-teleop-rig}\textbf{Kinematically Isomorphic
Leader-Follower Teleoperation Architecture across Heterogeneous
Manipulators.} Demonstration acquisition on physical robot platforms
using low-cost, kinematically isomorphic leader controllers (canonical
architecture (\citeproc{ref-wu2023gello}{Wu et al. 2024})). (A) Bimanual
teleoperation setup with twin 6-DoF leader arms mapped directly to dual
heavy-payload follower manipulators during a compliant liquid-pouring
task. (B) Single-arm leader-follower configuration controlling a
torque-controlled manipulator with an overhead task camera. (C)
Precision pick-and-place manipulation with passive gravity compensation
and joint-angle direct drive. Because the leader device mirrors follower
kinematics without requiring numerical inverse kinematics (IK), the
interface eliminates IK singularities while preserving the operator's
direct joint-level coordination. However, human reaction latency
(\(200\text{--}300\text{ ms}\)) and communication delays still introduce
phase lag into the recorded \((s_t, a_t)\) demonstration tuples. Source:
Adapted from Wu et al. (\citeproc{ref-wu2023gello}{Wu et al. 2024}), CC
BY 4.0.}

\end{figure}%

Every stage in this loop introduces latency that alters the timing of
the demonstrated correction. In an analytical reference setup with a
high-speed camera and a wired local network, display processing requires
\(t_{\text{disp}} = 50\text{ ms}\) across exposure, H.264 decoding, and
monitor refresh. Human visual perception and neuromuscular transmission
introduce an operator response latency of
\(t_{\text{hum}} = 200\text{ ms}\) for continuous path corrections.
Command quantization, packetization, and bus transport to the robot
controller add \(t_{\text{cmd}} = 10\text{ ms}\). Actuator torque
ramp-up and structural inertia impose a mechanical response latency of
\(t_{\text{mech}} = 80\text{ ms}\) before the robot achieves the
commanded acceleration. The total round-trip latency is:
\[t_{\text{loop}} = t_{\text{disp}} + t_{\text{hum}} + t_{\text{cmd}} + t_{\text{mech}} = 50\text{ ms} + 200\text{ ms} + 10\text{ ms} + 80\text{ ms} = 340\text{ ms}.\]

When a contact misalignment occurs at time \(t\), the operator sees the
error at \(t + 50\text{ ms}\), issues a motor command at
\(t + 250\text{ ms}\), and the arm exerts corrective force against the
workpiece at \(t + 340\text{ ms}\). During those \(340\text{ ms}\), the
physical system continues moving under momentum and external contact
forces. If the data logger records the current sensor observation
\(s_t\) alongside the human command \(a_t\) issued at that same clock
instant, the recorded demonstration pairs a physical state with an
action chosen in response to a state that existed \(250\text{ ms}\) in
the past. This temporal misalignment bakes an unmodeled phase lag
directly into the training targets, producing policies that over-correct
or oscillate when deployed autonomously at \(50\text{--}100\text{ Hz}\).

\textbf{Teleoperation round-trip latency}
(\(t_{\text{loop}} = t_{\text{disp}} + t_{\text{hum}} + t_{\text{cmd}} + t_{\text{mech}}\))
is the total elapsed time between a physical contact event and the
delivered corrective actuator response, embedding biological reaction
delays (\(200\text{--}300\text{ ms}\)) as artificial phase lag in
demonstration timestamps. \textbf{Demonstration yield}
(\(Y_{\text{usable}} = \lfloor T_{\text{active}} / t_{\text{cyc}} \rfloor \cdot R\))
is the fraction of scheduled facility time that converts into validated,
defect-free trajectory records after deducting startup, calibration,
operator rest, physical resets, and quality filtering losses.

\protect\phantomsection\label{callout-mathux2a-5.3}
\begin{fbxSimple}{callout-math}{Fermi Math:}{Ingestion Bandwidth and Storage Budgets for Multi-Camera Teleoperation}
\phantomsection\label{callout-math*-5.3}

\textbf{Question:} What is the sustained data ingestion rate, NVMe
storage bandwidth, and daily storage consumption for a bimanual
leader-follower teleoperation rig capturing multi-modal demonstration
data for physical imitation learning?

\textbf{Parameters:}

\begin{itemize}
\item
  \textbf{Perception:} 4 active stereo RGB-D cameras (2 static
  eye-to-hand environment cameras + 2 wrist eye-in-hand cameras):

  \begin{itemize}
  \tightlist
  \item
    RGB streams: \(1920 \times 1080\text{ px}\) at \(30\text{ Hz}\),
    \(24\text{ bits/px}\) (\(3\text{ bytes/px}\)).
  \item
    Depth streams: \(1280 \times 720\text{ px}\) at \(30\text{ Hz}\),
    \(16\text{ bits/px}\) (raw \texttt{uint16} depth in millimeters,
    \(2\text{ bytes/px}\)).
  \end{itemize}
\item
  \textbf{Kinematics \& Dynamics:} 2 leader arms + 2 follower arms
  (14-DoF total):

  \begin{itemize}
  \tightlist
  \item
    Joint positions \(\mathbf{q}\), velocities \(\dot{\mathbf{q}}\),
    commanded torques \(\boldsymbol{\tau}_{\text{cmd}}\), measured
    torques \(\boldsymbol{\tau}_{\text{meas}}\), and 4-stream action
    tuples logged at \(1000\text{ Hz}\) (\(128\text{ bytes/step}\)).
  \end{itemize}
\item
  \textbf{Tactile \& Haptics:} 2 dual-finger fingertip tactile sensor
  arrays (\(16 \times 16\) pressure taxels
  \(\times 2\text{ bytes} \times 4\text{ pads} = 2048\text{ bytes/step}\))
  logged at \(100\text{ Hz}\).
\end{itemize}

\textbf{Calculation:}

\begin{enumerate}
\def\labelenumi{\arabic{enumi}.}
\item
  \textbf{Raw Video Bandwidth:}
  \[\text{BW}_{\text{RGB}} = 4 \times (1920 \times 1080 \times 3\text{ bytes}) \times 30\text{ s}^{-1} = 4 \times 6.2208\text{ MB} \times 30 = 746.50\text{ MB/s}.\]
  \[\text{BW}_{\text{Depth}} = 4 \times (1280 \times 720 \times 2\text{ bytes}) \times 30\text{ s}^{-1} = 4 \times 1.8432\text{ MB} \times 30 = 221.18\text{ MB/s}.\]
  \[\text{BW}_{\text{Vision, total}} = 746.50\text{ MB/s} + 221.18\text{ MB/s} = 967.68\text{ MB/s} \approx 0.968\text{ GB/s}.\]
\item
  \textbf{Kinematic and Tactile Bandwidth:}
  \[\text{BW}_{\text{Kinematics}} = 128\text{ bytes} \times 1000\text{ s}^{-1} = 0.128\text{ MB/s}.\]
  \[\text{BW}_{\text{Tactile}} = 2048\text{ bytes} \times 100\text{ s}^{-1} = 0.205\text{ MB/s}.\]
\item
  \textbf{Total Sustained Ingestion Rate:}
  \[\text{BW}_{\text{total}} = 967.68 + 0.13 + 0.21 = 968.02\text{ MB/s} \approx 0.968\text{ GB/s}.\]
\item
  \textbf{Hourly and Shift Storage Consumption:}
  \[\text{Storage}_{\text{1 hour}} = 968.02\text{ MB/s} \times 3600\text{ s} = 3{,}484{,}872\text{ MB} \approx 3.48\text{ TB/hr} \text{ (raw uncompressed)}.\]
  When recording at a standard \(1280 \times 720\) RGB target
  (\(331.8\text{ MB/s}\)) alongside depth (\(221.2\text{ MB/s}\)) with
  synchronous 4-stream metadata, sustained ingestion requires
  \textbf{\(740\text{ MB/s}\)}, accumulating exactly
  \textbf{\(2.66\text{ TB/hr}\)} of uncompressed telemetry. Across an
  8-hour collection shift with 5 parallel workstations, the raw
  telemetry ingestion generates:
  \[\text{Storage}_{\text{shift}} = 5 \times (2.66\text{ TB/hr} \times 8\text{ hr}) = 106.4\text{ TB per shift}.\]
\end{enumerate}

\textbf{Systems Significance:} A high-capacity data acquisition
workstation cannot log uncompressed multi-view streams over standard USB
3.0 buses (capped at \(400\text{ MB/s}\) real-world throughput) or write
to SATA SSDs (capped at \(550\text{ MB/s}\)). Ingesting
\(740\text{ MB/s}\) to \(968\text{ MB/s}\) requires dedicated PCIe Gen4
NVMe storage (capable of \(>5000\text{ MB/s}\) sequential writes) or
hardware-accelerated NVENC video compression (H.265 intra-frame at CRF
18, compressing streams \(8\times\) to \(\approx 92.5\text{ MB/s}\) /
\(333\text{ GB/hr}\)) to prevent OS buffer overruns and silent frame
drops.

\end{fbxSimple}

In \textbf{bilateral teleoperation}, where contact forces measured at
the follower arm are reflected back as haptic torques to the leader arm,
transport delay creates a severe stability hazard. Consider the net
energy exchanged across the communication channel over time horizon
\(t\):
\[\Delta E_{\text{channel}}(t) = \int_0^t \left( f_{\text{master}}(\sigma) \dot{x}_{\text{master}}(\sigma) - f_{\text{follower}}(\sigma) \dot{x}_{\text{follower}}(\sigma) \right) d\sigma.\]
In an ideal, zero-latency system, \(\Delta E_{\text{channel}} = 0\),
satisfying the passivity condition \(\Delta E \le 0\) (meaning the
communication channel acts as a passive transmission line that cannot
generate energy). When transport delay \(\tau(t) > 0\) varies
dynamically across the network, the delayed force reflection acts
out-of-phase with operator hand velocity. The interface behaves like a
delayed spring: when the operator pushes into contact and begins to
retract, the delayed reaction force arrives late, pushing the operator's
hand in the direction of motion. The product
\(f_{\text{master}} \cdot \dot{x}_{\text{master}}\) does positive
mechanical work, injecting virtual energy (\(\Delta E > 0\)) into the
feedback loop.\sidenote{\footnotesize In bilateral teleoperation with force feedback,
  variable network transport latency injects virtual energy into the
  control loop (\(\Delta E > 0\)). Without time-domain passivity
  controllers (TDPC) or wave-variable transformations, this non-passive
  energy triggers destructive high-frequency mechanical limit cycles and
  contact chatter.} This non-passive energy drives high-frequency
mechanical limit cycles, causing the master and follower arms to chatter
violently against rigid contact fixtures. Stabilizing the rig requires
Time-Domain Passivity Controllers (TDPC) or wave-variable
transformations, which inject virtual damping into the teleoperation
link. While passivity controllers prevent hardware damage, the added
damping acts as a low-pass filter on operator hand motions, attenuating
high-frequency contact cues and suppressing subtle tactile feedback.

\protect\phantomsection\label{callout-ruleux2a-5.4}
\begin{fbxSimple}{callout-rule}{Architect’s Rule of Thumb:}{Bilateral Teleoperation Passivity Condition}
\phantomsection\label{callout-rule*-5.4}

In bilateral teleoperation with force reflection, variable network
latency \(\tau(t)\) injects virtual energy (\(\Delta E > 0\)) into the
operator interface. To prevent contact instability and mechanical limit
cycles, the real-time nervous system must enforce the passivity
condition
\(\int_0^t (f_{\text{master}}\dot{x}_{\text{master}} - f_{\text{follower}}\dot{x}_{\text{follower}}) \, d\tau \le 0\)
via Time-Domain Passivity Observers (TDPO).

\end{fbxSimple}

The physical interface through which the operator delivers commands acts
as a hardware filter on the demonstration data. Consider an assembly
task requiring the insertion of a metal peg into a tight clearance
fixture. When collected using a direct kinesthetic leader arm with six
degrees of freedom and bilateral force feedback, the operator feels the
contact forces transmitted from the follower arm and adjusts insertion
angles smoothly at velocities exceeding \(0.15\text{ m/s}\). When the
same task is teleoperated using a rate-controlled six-axis mouse or
planar joystick with visual-only feedback, the operator has no physical
sensation of contact force or mechanical binding. To avoid generating
excessive torques that trigger hardware safety cutoffs, the operator
switches to an open-loop step-and-wait strategy, advancing the tool in
small increments of \(1.0\text{ mm}\), pausing to inspect the video feed
for part deflection, and then commanding another step. The two
interfaces produce fundamentally different datasets for the exact same
physical task. The kinesthetic leader arm populates the state-action
space with continuous velocity profiles and high-frequency force
adaptations during contact. The rate-controlled interface records a
dataset characterized by low average velocities, prolonged periods of
zero motion while the operator inspects the visual feed, and an absence
of corrective torques during the insertion phase. A policy trained by
imitation on the rate-controlled dataset learns this stop-and-go
behavior and stalls whenever visual latency prevents it from confirming
contact. The demonstration dataset captures the narrow operational
envelope permitted by the human interface rather than the physical
capability of the mechanical body. As categorized in
Table~\ref{tbl-05-teleop-taxonomy}, the teleoperation interface class
fundamentally governs control bandwidth, round-trip transport latency,
operator fatigue limits, and the resulting dynamical artifacts embedded
in the recorded demonstration corpus.

\begin{longtable}[]{@{}
  >{\raggedright\arraybackslash}p{(\linewidth - 12\tabcolsep) * \real{0.1429}}
  >{\raggedright\arraybackslash}p{(\linewidth - 12\tabcolsep) * \real{0.1429}}
  >{\raggedright\arraybackslash}p{(\linewidth - 12\tabcolsep) * \real{0.1429}}
  >{\raggedright\arraybackslash}p{(\linewidth - 12\tabcolsep) * \real{0.1429}}
  >{\raggedright\arraybackslash}p{(\linewidth - 12\tabcolsep) * \real{0.1429}}
  >{\raggedright\arraybackslash}p{(\linewidth - 12\tabcolsep) * \real{0.1429}}
  >{\raggedright\arraybackslash}p{(\linewidth - 12\tabcolsep) * \real{0.1429}}@{}}
\caption{\textbf{Teleoperation Master-Follower Interface Taxonomy for
Physical Demonstration Collection.} Comparison of input interface
paradigms across kinematic mapping, transport latency, haptic fidelity,
control bandwidth, operator fatigue limits, and characteristic
demonstration artifacts. The choice of teleoperation hardware acts as a
mechanical and dynamical filter on the recorded state-action
distribution.}\label{tbl-05-teleop-taxonomy}\tabularnewline
\toprule\noalign{}
\begin{minipage}[b]{\linewidth}\raggedright
Interface Paradigm \& Hardware Class
\end{minipage} & \begin{minipage}[b]{\linewidth}\raggedright
Kinematic \& Control Mapping
\end{minipage} & \begin{minipage}[b]{\linewidth}\raggedright
Round-Trip Latency (\(\tau_{\text{loop}}\))
\end{minipage} & \begin{minipage}[b]{\linewidth}\raggedright
Haptic / Force Feedback Fidelity
\end{minipage} & \begin{minipage}[b]{\linewidth}\raggedright
Effective Control Bandwidth (\(f_{\text{ctrl}}\))
\end{minipage} & \begin{minipage}[b]{\linewidth}\raggedright
Operator Fatigue Limit (\(T_{\text{fatigue}}\))
\end{minipage} & \begin{minipage}[b]{\linewidth}\raggedright
Characteristic Demonstration Artifacts \& Policy Bias
\end{minipage} \\
\midrule\noalign{}
\endfirsthead
\toprule\noalign{}
\begin{minipage}[b]{\linewidth}\raggedright
Interface Paradigm \& Hardware Class
\end{minipage} & \begin{minipage}[b]{\linewidth}\raggedright
Kinematic \& Control Mapping
\end{minipage} & \begin{minipage}[b]{\linewidth}\raggedright
Round-Trip Latency (\(\tau_{\text{loop}}\))
\end{minipage} & \begin{minipage}[b]{\linewidth}\raggedright
Haptic / Force Feedback Fidelity
\end{minipage} & \begin{minipage}[b]{\linewidth}\raggedright
Effective Control Bandwidth (\(f_{\text{ctrl}}\))
\end{minipage} & \begin{minipage}[b]{\linewidth}\raggedright
Operator Fatigue Limit (\(T_{\text{fatigue}}\))
\end{minipage} & \begin{minipage}[b]{\linewidth}\raggedright
Characteristic Demonstration Artifacts \& Policy Bias
\end{minipage} \\
\midrule\noalign{}
\endhead
\bottomrule\noalign{}
\endlastfoot
\textbf{Bilateral Kinesthetic Leader Arm} (1:1 Active Replica Puppet,
e.g., canonical dual-arm bimanual teleoperation workcell
(\citeproc{ref-zhao2023learning}{Zhao et al. 2023})) & Direct
joint-to-joint position copying
(\(q_{\text{master}} \mapsto q_{\text{follower}}\)) with active gravity
compensation & \(15\text{ to }35\text{ ms}\) (Direct bus / EtherCAT /
CANopen) & High bilateral force reflection; senses contact stiffness and
binding directly & \(500\text{ to }1000\text{ Hz}\) telemetry
(\(5\text{ to }10\text{ Hz}\) human neuromuscular) &
\(45\text{ to }60\text{ min}\) (Arm tremor \(>2.0\text{ mm}\) after
prolonged holding) & Continuous velocity profiles
(\(>0.15\text{ m/s}\)); fluid compliance; operator motor tremors
recorded in joint velocity derivatives \\
\textbf{6-DoF 3D Spacemouse} (Isometric Elastic Puck) & Elastic
force/torque to Cartesian velocity setpoint
(\(\dot{\mathbf{x}}_{\text{cmd}} = \mathbf{K}_v \mathbf{f}_{\text{mouse}}\));
resolved by IK & \(60\text{ to }120\text{ ms}\) (USB HID polling + host
IK solver + motor ramp) & Zero force reflection; visual-only estimation
of workpiece contact loads & \(50\text{ to }100\text{ Hz}\) Cartesian
rate (\(<2\text{ Hz}\) human correction rate) &
\(90\text{ to }120\text{ min}\) (Low physical effort, desk-supported
wrist rest) & Step-and-wait hunting motions; low average velocity
(\(<0.03\text{ m/s}\)); prolonged zero-motion pauses during visual
verification; absent reactive contact forces \\
\textbf{Optical / VR 6-DoF Spatial Controller} (Tracked Wand / Spatial
Poses) & Absolute Cartesian pose tracking
(\(\mathbf{T}_{\text{hand}} \in SE(3)\)) with clutching and workspace
scale & \(80\text{ to }180\text{ ms}\) (Optical exposure + SLAM pose
solve + wireless link) & Low (Vibrotactile ERM/LRA pulses only); no
directional force reflection & \(60\text{ to }120\text{ Hz}\) spatial
pose (\(\sigma_{\text{jitter}} \approx 0.8\text{ mm}\)) &
\(30\text{ to }45\text{ min}\) (``Gorilla arm'' shoulder fatigue from
ungrounded holding) & Over-correction oscillations from visual-motor
phase lag; clutching discontinuities; out-of-plane drift; exaggerated
retraction motions upon visual contact \\
\textbf{Direct Instrumented Exoskeleton} (Wearable Master Glove / Arm
Linkage) & High-DoF joint-angle mapping with kinematic retargeting to
multi-finger hand & \(25\text{ to }60\text{ ms}\) (Wearable bus +
retargeting solver + low-level drive) & Moderate-to-high active tendon
force feedback on individual finger linkages &
\(200\text{ to }500\text{ Hz}\) joint state
(\(7\text{ to }12\text{ Hz}\) finger dexterity) &
\(20\text{ to }35\text{ min}\) (Mechanical weight burden, pressure
points, harness sweating) & Biomechanical calibration mismatch;
non-linear retargeting artifacts; rich contact force profiles during
multi-point precision grasps \\
\end{longtable}

Treating the recorded stream as a single sequence of states and actions
obscures the signal transformations that occur between human intent and
mechanical motion. A physical logging pipeline must record four distinct
quantities at each time step:

\begin{enumerate}
\def\labelenumi{\arabic{enumi}.}
\tightlist
\item
  \(a_{\text{req}}\): the operator requested action from the master
  input device.
\item
  \(a_{\text{cmd}}\): the mapped setpoint produced after inverse
  kinematics and workspace checks on the Linux MPU.
\item
  \(a_{\text{enf}}\): the enforced command emitted by the real-time MCU
  safety layer after velocity limits, current clamps, and barrier
  functions have evaluated.
\item
  \(a_{\text{meas}}\): the measured physical response from joint
  encoders and current sensors.
\end{enumerate}

Logging only one of these quantities corrupts the training signal. If an
operator commands an excessive velocity toward a joint limit and the
low-level controller clamps the setpoint, recording \(a_{\text{req}}\)
teaches the policy to demand unsafe, unexecutable motions that violate
hardware boundaries. If the logger records \(a_{\text{enf}}\) or
\(a_{\text{meas}}\) while discarding \(a_{\text{req}}\), the recorded
state-action pairs attribute the safety filter's intervention to the
human demonstration, teaching the model that a gentle stop was the
operator's deliberate choice rather than the result of an automated
emergency clamp. Preserving the full chain from \(a_{\text{req}}\)
through \(a_{\text{meas}}\) enables an engineer to isolate whether an
execution failure during deployment stems from policy prediction error,
kinematic retargeting artifacts, safety interventions, or mechanical
tracking error.

A nominal hour of scheduled teleoperation yields only a fraction of that
time in usable training data. In an analytical baseline for a
\(T_{\text{sched}} = 60.0\text{ min}\) (\(3600\text{ s}\)) collection
shift, initial cell initialization, bus synchronization, and safety
watchdog checks consume \(t_{\text{setup}} = 8.0\text{ min}\)
(\(480\text{ s}\)). Optical sensor calibration, multi-camera extrinsic
sweeps, and 6-axis force-torque sensor gravity nulling require
\(t_{\text{calib}} = 7.0\text{ min}\) (\(420\text{ s}\)). Ergonomic
fatigue limits human operator concentration, requiring mandatory rest
breaks totaling \(t_{\text{rest}} = 10.0\text{ min}\) (\(600\text{ s}\))
per hour to prevent motor tremor drift. The remaining active motion
window is
\(T_{\text{motion}} = 60.0 - 8.0 - 7.0 - 10.0 = 35.0\text{ min}\)
(\(2100\text{ s}\)). Within this active window, each manipulation
attempt requires an average execution time of
\(t_{\text{att}} = 25.0\text{ s}\), followed by a reset phase of
\(t_{\text{rst}} = 15.0\text{ s}\) to return the robot to its home
configuration and reset the physical workpiece fixture. The total cycle
time per attempt is
\(t_{\text{cyc}} = t_{\text{att}} + t_{\text{rst}} = 40.0\text{ s}\),
yielding at most
\(N_{\text{att}} = \lfloor 2100\text{ s} / 40.0\text{ s} \rfloor = 52\text{ attempts}\)
per session. This represents \(52 \times 25.0\text{ s} = 1300\text{ s}\)
(\(21.67\text{ min}\), or \(36.1\%\) of the hour) of gross attempt
motion and \(52 \times 15.0\text{ s} = 780\text{ s}\)
(\(13.0\text{ min}\), or \(21.7\%\)) of physical reset time. Human
teleoperation is subject to operational errors, including dropped
objects, jamming, and accidental contact limit violations. Assuming an
in-trial failure rate of \(20\%\), \(41.6\text{ attempts}\) complete
successfully. Post-collection quality review removes an additional
\(15\%\) of completed trials due to operator hesitation, visual
occlusions, or borderline task success metrics. The net retention rate
across all collection stages is
\(R = (1 - 0.20) \times (1 - 0.15) = 0.68\), resulting in
\(N_{\text{usable}} = 52 \times 0.68 \approx 35.36\text{ usable episodes}\).
At \(25.0\text{ s}\) of demonstration trajectory per episode,
\(35.36\text{ episodes}\) produce \(884\text{ s}\), or approximately
\(14.73\text{ min}\) of usable demonstration time per scheduled hour
(\(24.6\%\) net yield). A dataset target of
\(5000\text{ validated episodes}\) requires
\(5000 / 35.36 \approx 141.4\text{ scheduled operator hours}\), scaling
labor requirements and collection schedules by a factor of \(4.07\)
relative to the raw \(34.7\text{ h}\) of retained trajectory duration.

\protect\phantomsection\label{callout-ruleux2a-5.5}
\begin{fbxSimple}{callout-rule}{Architect’s Rule of Thumb:}{The 4-to-1 Facility-to-Yield Ratio}
\phantomsection\label{callout-rule*-5.5}

Never budget physical demonstration campaigns by scheduled facility
hours. Due to sensor calibration, workpiece re-fixturing, operator
ergonomics, and post-collection QA filtering, one hour of scheduled
robot cell time yields at most 15 minutes of validated, defect-free
trajectory data (a \(4:1\) overhead factor).

\end{fbxSimple}

Hardware degradation during collection depends on the accumulation of
physical cycles and directional stress reversals rather than elapsed
wall-clock time. Consider the wrist joint of an articulated arm equipped
with a strain-wave harmonic drive and internal routing flex cables. In
an analytical wear model, the harmonic gear transmission carries a rated
fatigue life of \(L_{\text{gear}} = 5{,}000{,}000\text{ stress cycles}\)
under rated output torque, and the internal cable bundle is rated for
\(L_{\text{cable}} = 200{,}000\text{ bending cycles}\). During the
\(25.0\text{ s}\) execution of a peg-insertion attempt, contact
alignment and search motions introduce an average of
\(n_{\text{gear}} = 14\text{ torque reversals}\) and
\(n_{\text{cable}} = 6\text{ joint flex cycles}\). Resetting the rig
adds an additional \(2\text{ full flex cycles}\) during the return
stroke. To collect \(5000\text{ validated episodes}\) at a net retention
rate of \(R = 0.68\), the hardware must execute
\(N_{\text{total}} = 5000 / 0.68 \approx 7353\text{ physical attempts}\).
Across these \(7353\text{ attempts}\), the wrist transmission sustains
\(7353 \times 14 = 102{,}942\text{ stress cycles}\), consuming \(2.1\%\)
of its rated fatigue life. Concurrently, the internal cable bundle
undergoes \(7353 \times (6 + 2) = 58{,}824\text{ bending cycles}\),
consuming \(58{,}824 / 200{,}000 = 29.4\%\) of its total service life on
a single dataset collection campaign. If three collection campaigns run
sequentially on the same rig without component replacement, the flex
cable reaches \(88.2\%\) of its fatigue limit, where conductor strand
breakage and intermittent signal loss begin to corrupt sensor streams.
The mechanical wear budget must be computed directly from cycle counts
and load reversals to schedule component replacement before hardware
failure invalidates collection sessions. The operational time breakdown
and mechanical wear budget across a declared \(60\text{ min}\)
teleoperation session are synthesized in
Table~\ref{tbl-05-collection-budget}.

\begin{longtable}[]{@{}
  >{\raggedright\arraybackslash}p{(\linewidth - 10\tabcolsep) * \real{0.1667}}
  >{\raggedright\arraybackslash}p{(\linewidth - 10\tabcolsep) * \real{0.1667}}
  >{\raggedright\arraybackslash}p{(\linewidth - 10\tabcolsep) * \real{0.1667}}
  >{\raggedright\arraybackslash}p{(\linewidth - 10\tabcolsep) * \real{0.1667}}
  >{\raggedright\arraybackslash}p{(\linewidth - 10\tabcolsep) * \real{0.1667}}
  >{\raggedright\arraybackslash}p{(\linewidth - 10\tabcolsep) * \real{0.1667}}@{}}
\caption{\textbf{Physical Demonstration Collection Budget Breakdown
across a Scheduled 60-Minute Session.} Partitioning of facility clock
time, active motion yield, quality filtering losses, and mechanical
component fatigue across collection stages for a precision peg-in-hole
assembly workcell (\(T_{\text{sched}} = 60.0\text{ min}\),
\(N_{\text{att}} = 52\text{ attempts}\), nominal attempt
\(t_{\text{att}} = 25.0\text{ s}\), reset
\(t_{\text{rst}} = 15.0\text{ s}\)). The nominal hour yields only
\(14.73\text{ min}\) of retained demonstration data (\(24.6\%\) net
yield), while imposing significant cumulative fatigue on mechanical
transmissions and wiring flex
lines.}\label{tbl-05-collection-budget}\tabularnewline
\toprule\noalign{}
\begin{minipage}[b]{\linewidth}\raggedright
Session Stage / Activity
\end{minipage} & \begin{minipage}[b]{\linewidth}\raggedright
Allocated Duration (\(\Delta t\))
\end{minipage} & \begin{minipage}[b]{\linewidth}\raggedright
Wall-Clock Share (\%)
\end{minipage} & \begin{minipage}[b]{\linewidth}\raggedright
Physical Operation \& Control State
\end{minipage} & \begin{minipage}[b]{\linewidth}\raggedright
Hardware Stress \& Wear Mechanism
\end{minipage} & \begin{minipage}[b]{\linewidth}\raggedright
Net Trajectory Output (\(N_{\text{episodes}}\) / \(T_{\text{yield}}\))
\end{minipage} \\
\midrule\noalign{}
\endfirsthead
\toprule\noalign{}
\begin{minipage}[b]{\linewidth}\raggedright
Session Stage / Activity
\end{minipage} & \begin{minipage}[b]{\linewidth}\raggedright
Allocated Duration (\(\Delta t\))
\end{minipage} & \begin{minipage}[b]{\linewidth}\raggedright
Wall-Clock Share (\%)
\end{minipage} & \begin{minipage}[b]{\linewidth}\raggedright
Physical Operation \& Control State
\end{minipage} & \begin{minipage}[b]{\linewidth}\raggedright
Hardware Stress \& Wear Mechanism
\end{minipage} & \begin{minipage}[b]{\linewidth}\raggedright
Net Trajectory Output (\(N_{\text{episodes}}\) / \(T_{\text{yield}}\))
\end{minipage} \\
\midrule\noalign{}
\endhead
\bottomrule\noalign{}
\endlastfoot
\textbf{Cell Startup \& Bus Sync} & \(8.0\text{ min}\)
(\(480\text{ s}\)) & \(13.3\%\) & Bus synchronization
(EtherCAT/CANopen), safety watchdog verification, optical homing &
Power-on electrical inrush; thermal warm-up in motor windings and sensor
bridges & \(N = 0\text{ ep}\), \(0.0\text{ min}\) yield \\
\textbf{Sensor Calibration \& Nulling} & \(7.0\text{ min}\)
(\(420\text{ s}\)) & \(11.7\%\) & Camera intrinsic/extrinsic sweep,
6-axis F/T gravity bias nulling, tactile matrix zeroing & Static holding
torque; sensor bridge thermal stabilization; zero dynamic gear wear &
\(N = 0\text{ ep}\), \(0.0\text{ min}\) yield \\
\textbf{Operator Rest \& Ergonomics} & \(10.0\text{ min}\)
(\(600\text{ s}\)) & \(16.7\%\) & Mandatory operator breaks to mitigate
visual fatigue, cognitive saturation, and tremor drift & Passive
standby; electromagnetic brakes engaged or zero-velocity holding current
& \(N = 0\text{ ep}\), \(0.0\text{ min}\) yield \\
\textbf{Physical Scene Resets} & \(13.0\text{ min}\) (\(780\text{ s}\))
& \(21.7\%\) & Part unclamping, receptacle clearing, peg re-fixturing,
manipulator return stroke (\(15.0\text{ s/attempt}\)) &
\(52\text{ return strokes} \times 2\text{ bends} = 104\text{ cable flex cycles}\);
zero contact payload & \(N = 0\text{ ep}\), \(0.0\text{ min}\) yield \\
\textbf{Active Teleoperated Attempts} & \(22.0\text{ min}\)
(\(1320\text{ s}\)) & \(36.6\%\) & Real-time human teleoperation guiding
peg approach, chamfer search, and insertion (\(25.0\text{ s/attempt}\))
& \(52 \times 14 = 728\text{ torque reversals}\) on harmonic drive;
\(52 \times 6 = 312\text{ cable flex cycles}\); contact loads
\(\le 15\text{ N}\) & \(N = 52\text{ raw attempts}\),
\(21.67\text{ min}\) gross motion \\
\textbf{In-Trial Failures (Rejected)} & \(-4.33\text{ min}\)
(\(-260\text{ s}\)) & \(-7.2\%\) & Part slippage, jamming, watchdog
force trip (\(F_N > 15.0\text{ N}\)), dropped pegs (\(20\%\) loss) &
High-force contact transients; abrasive wear on gripper silicone pads &
\(-10.4\text{ ep}\) discarded \\
\textbf{Post-Collection QA Filtering} & \(-2.60\text{ min}\)
(\(-156\text{ s}\)) & \(-4.3\%\) & Offline filtering: link latency
\(>100\text{ ms}\), visual occlusion, operator hesitation (\(15\%\)
loss) & Zero hardware wear (offline computational validation pass) &
\(-6.24\text{ ep}\) discarded \\
\textbf{Net Validated Dataset Yield} & \textbf{\(14.73\text{ min}\)}
(\(884\text{ s}\)) & \textbf{\(24.6\%\)} & High-quality synchronized
multi-modal trajectories (\(s_t, a_t\)) admitted into training buffer &
Consumes \(2.1\%\) of wrist gear life and \(29.4\%\) of internal flex
cable life per \(5000\text{ episode}\) campaign &
\textbf{\(N_{\text{usable}} = 35.36\)} episodes \\
\end{longtable}

Teleoperation records trajectories that are bounded by what a human
operator can perceive through a display, coordinate through an input
interface, and stabilize against system latencies. The resulting dataset
reflects the operational policy of the human-interface composite rather
than an unconstrained sample of physical dynamics. Every demonstration
collected in this manner carries the structural limitations, response
delays, and conservative motion biases of the collecting apparatus into
the training distribution.

\section{Coverage and the Collecting Policy}\label{sec-data-5-4}

A physical dataset is generated by executing a sequence of control
actions on hardware and logging the resulting state transitions over
time. When an agent or human operator executes trajectories
\(\tau = \{(s_0, a_0), (s_1, a_1), \dots, (s_T, a_T)\}\) at a fixed
logging rate such as \(100\text{ Hz}\), the collected samples define an
\textbf{empirical state-action occupancy distribution
(\citeproc{ref-spencer2021feedback}{Spencer et al. 2021})}
\(d^{\pi_{\text{col}}}(s, a)\). The training loss in imitation learning
minimizes prediction error weighted by this distribution, optimizing the
learned policy \(\pi_\theta\) exclusively on states where
\(d^{\pi_{\text{col}}}(s) > 0\). The machine that gathers the data
determines what is in the dataset, which is the operational endogeneity
of physical systems entering the learning pipeline at training time. The
dataset does not represent an unconstrained sampling of the physical
state space. It represents the trajectory manifold traced by the
collecting policy \(\pi_{\text{col}}\) under its specific control
biases, interface latencies, and risk constraints.

Consider an analytical example of a robotic gripper aligning with a
workpiece edge using tactile sensing and force feedback. The physical
workspace requires managing contact angles over
\(\theta \in [-15.0^\circ, +15.0^\circ]\) and normal contact forces up
to \(F_N = 20.0\text{ N}\). Two different collecting policies gather
data on this exact task using identical six-axis force-torque sensors
and tactile arrays logging at \(f_s = 100\text{ Hz}\). Each policy
records \(N = 1000\text{ episodes}\) lasting \(10.0\text{ s}\) per
episode, producing identical dataset sizes of
\(1{,}000{,}000\text{ transition records}\)
(\(1.00 \times 10^6\text{ samples}\)).

The first collector is a cautious human teleoperator who approaches the
workpiece along a single nominal axis at \(v = 0.01\text{ m/s}\),
maintaining contact strictly perpendicular to the surface
(\(\theta \in [-2.0^\circ, +2.0^\circ]\)) under light normal force
(\(F_N \in [4.0\text{ N}, 6.0\text{ N}]\)). The second collector is an
exploratory compliance policy that approaches across a broader spread of
approach angles (\(\theta \in [-15.0^\circ, +15.0^\circ]\)),
encountering transient slips and varied contact loads
(\(F_N \in [1.0\text{ N}, 18.0\text{ N}]\)). Discretizing the
two-dimensional state space of contact angle and normal force into bins
of \(1.0^\circ\) and \(1.0\text{ N}\) produces a deployment grid of
\(30 \times 20 = 600\text{ discrete cells}\).

Reconstructing the visitation counts reveals the structural difference
between these datasets. The cautious collector deposits all
\(1{,}000{,}000\text{ samples}\) into \(4 \times 2 = 8\text{ cells}\),
concentrating an average of \(125{,}000\text{ samples}\) in each visited
bin while leaving \(592\text{ cells}\) (\(98.7\%\) of the required
operational grid) with a visitation count of zero. The exploratory
collector distributes its \(1{,}000{,}000\text{ samples}\) across
\(510\text{ cells}\) with a mean density of
\(1961\text{ samples per cell}\), leaving \(90\text{ cells}\)
(\(15.0\%\)) unvisited. Despite identical byte counts, logging
frequencies, and sensor streams, these datasets answer different
physical questions. The cautious dataset provides dense supervision on
maintaining an already aligned contact, whereas the exploratory dataset
provides information on how tactile signatures and reaction forces
evolve during off-axis engagement.

An action label in a physical dataset exists only because the collecting
policy chose to visit that state and issued a command. When a state has
zero visitation count, the dataset contains neither a preferred action
nor physical evidence that an action cannot be found. Supervised
learning algorithms cannot distinguish between a physical configuration
that is inherently unstable, a configuration that is safe but unvisited,
and a configuration that causes mechanical damage. To the loss function,
all unvisited regions appear identically as a total absence of
constraints.

This absence creates the compounding failure mechanism of behavioral
cloning. Suppose a policy trained on the cautious dataset executes on
the physical hardware. At \(t = 0.00\text{ s}\), an unmodeled friction
variation or mechanical vibration perturbs the gripper to an angle of
\(\theta = +4.5^\circ\), placing the state outside the 8 demonstrated
cells. The policy evaluates its neural network on an unvisited input.
Because empirical risk minimization provides zero constraint outside the
support of the training distribution, the policy outputs an action
\(\hat{a}\) with a small bounded angular error. Applying \(\hat{a}\)
causes the next state at \(t = 0.01\text{ s}\) to reach
\(\theta = +6.0^\circ\) with a contact force of \(11.0\text{ N}\).
Because the dataset contains zero examples showing how to reduce an
off-axis angle under high contact force back to the nominal corridor,
the policy cannot generate a corrective command. The error shifts the
system further from the demonstrated manifold on every control step, and
the gripper jams against the fixture.

The \textbf{compounding covariate shift flywheel} is the closed-loop
destabilization dynamic in behavioral cloning where unmodeled
disturbances push the physical state outside the empirical demonstration
support \(d^{\pi_{\text{col}}}(s)\), prompting out-of-distribution
policy commands that amplify tracking errors, scale state drift
quadratically (\(O(T^2 \epsilon)\)) over execution horizon \(T\), and
culminate in boundary collisions, wheel lockup, or joint overcurrent
trips.

\protect\phantomsection\label{callout-ruleux2a-5.6}
\begin{fbxSimple}{callout-rule}{Architect’s Rule of Thumb:}{The Boundary Attractor Law}
\phantomsection\label{callout-rule*-5.6}

An imitation policy cannot interpolate across an unvisited physical
state. Every boundary state omitted from the collection distribution
\(d^{\pi_{\text{col}}}(s)\) becomes an unconstrained dynamical attractor
toward physical collision or motor overcurrent tripping when evaluated
under closed-loop feedback.

\end{fbxSimple}

Evaluating dataset quality by computing extrema over the collected
values creates a circular evaluation metric. If a collector only records
contact forces between \(4.0\text{ N}\) and \(6.0\text{ N}\), measuring
coverage relative to the minimum and maximum logged values reports
\(100\%\) completeness across that interval. Coverage must instead be
evaluated against an external \textbf{scenario ledger}, which defines
the required deployment envelope across object poses, surface friction
ranges, contact velocities, payload masses, and external disturbance
bounds.

Unsupported deployment regions are detected by performing a relational
join between the required scenario bins and the empirical dataset
visitation counts. For each cell in the ledger, the analysis reports the
exact transition count \(N_{\text{cell}}\). Cells with
\(N_{\text{cell}} = 0\) represent unmonitored blind spots where the
policy will execute without training supervision, while cells with low
counts, such as \(N_{\text{cell}} < 50\text{ samples}\), represent
high-variance training regimes.

This coverage deficit cannot be detected by evaluating validation loss
on a held-out test split of the collected data. Because the held-out
split is drawn from the same trajectories generated by
\(\pi_{\text{col}}\), it shares the identical empirical occupancy
\(d^{\pi_{\text{col}}}\). The validation loss can reach
\(0.001\text{ rad}\), indicating statistical convergence, while
\(98.7\%\) of the required operational envelope contains zero training
examples. Verifying coverage requires computing the ratio of supported
scenario cells to total required cells before a policy is deployed on
physical actuators (Figure~\ref{fig-05-collector-coverage-ledger}).

\begin{figure}[t]

\centering{

\includegraphics[width=1\linewidth,height=\textheight,keepaspectratio]{index_files/mediabag//Users/VJ/GitHub/MLSysBook-integrate-physical/physical/book/chapters/05-data/figures/fig05_collector_coverage_ledger.pdf}

}

\caption{\label{fig-05-collector-coverage-ledger}\textbf{Collector
Empirical Occupancy and Scenario Ledger Relational Join.} (Left)
Comparison of empirical state-action visitation distributions between a
cautious teleoperator (concentrating \(1{,}000{,}000\) transitions into
8 nominal cells, leaving 98.7\% of the 600-cell operational envelope
unvisited despite near-zero validation loss
\(\mathcal{L}_{\text{val}} = 0.001\text{ rad}\)) and an exploratory
policy (visiting 510 cells with 85.0\% coverage across the full contact
angle and force envelope). (Right) Performing a relational join between
the external scenario ledger \(\mathcal{S}_{\text{req}}\) and empirical
visitation \(d^{\pi_{\text{col}}}\) partitions the state space into
supported nominal regimes (\(N_{\text{cell}} \ge 500\)), high-variance
transition zones (\(0 < N_{\text{cell}} < 50\)), and zero-support blind
spots (\(N_{\text{cell}} = 0\)) where autonomous control must be
refused. Systems significance: Demonstrates why statistical convergence
on a training split cannot verify operational coverage on physical
hardware.}

\end{figure}%

\section{Intervention, Censoring, and Recovery Data}\label{sec-data-5-5}

When a hazardous or boundary state does not appear in a training
dataset, an engineer cannot conclude that the physical system avoids
that state naturally. An absent state in the empirical log corresponds
to one of four distinct physical mechanisms during data collection:

\begin{enumerate}
\def\labelenumi{\arabic{enumi}.}
\tightlist
\item
  \textbf{Unvisited nominal path}: The collecting policy never visited
  the state because nominal rollouts remained entirely within the
  primary task corridor. This absence represents unmonitored nominal
  operation, detected by cross-referencing trajectory logs against the
  scenario ledger.
\item
  \textbf{Safety governor interception}: An active safety barrier or
  deterministic reflex in the nervous system detected an impending
  boundary violation and truncated the motion before the state was
  reached. This absence is detected directly in the safety monitor logs,
  which record the barrier trip.
\item
  \textbf{Human supervisory takeover}: A human safety operator perceived
  a developing hazard and disengaged the autonomous policy before
  contact occurred. This absence is detected through the takeover
  disengagement flag on the vehicle or manipulator bus.
\item
  \textbf{Post-collection censoring}: The machine actually entered the
  hazardous state and collided with an obstacle or jammed a joint, but
  an automated post-processing script discarded the episode as corrupt
  or incomplete. This absence is not detected by dataset inspection,
  because the failure transitions were purged before the training arrays
  were written to disk.
\end{enumerate}

Each mechanism produces an identical zero count in a scenario cell, yet
treating them identically leads to incorrect conclusions about policy
stability.

Human takeovers and automated disengagements introduce a selection trap
into fleet data. An intervention occurs precisely when the autonomous
policy has entered a region of state space that it handles poorly,
causing tracking error or model uncertainty to exceed an operational
threshold. Because disengagements are triggered by failing rollouts, the
collected intervention data is not drawn from the nominal state
distribution of the task. If a safety driver intervenes when a mobile
robot drifts \(0.15\text{ m}\) outside its lane corridor at
\(2.0\text{ m/s}\), the trajectory recorded after the takeover captures
an expert recovering from an edge case disturbance. The distribution of
states visited during these rescues, denoted \(d^{\text{takeover}}\), is
conditioned on policy failure rather than standard task execution.
Treating intervention logs as if they were drawn from the nominal
distribution \(d^{\pi^*}\) distorts the training distribution toward
states that only exist because the prior policy failed.

Appending a rescued episode directly to a behavioral cloning dataset
teaches the policy that moving toward a hazard is the correct precursor
to completing the task. When an episode is recorded continuously across
a takeover, the pre-intervention segment contains erroneous control
outputs generated by the learned policy, whereas the post-intervention
segment contains corrective actions generated by the human operator that
guide the system back to the goal. When a neural network trains on this
full sequence under a mean squared error or cross-entropy imitation
objective, it minimizes loss across all transitions simultaneously. The
optimization associates the uncorrected drift (such as an increasing
heading error or excessive contact force) with subsequent task success.
When evaluated on the physical machine, the resulting policy actively
steers toward the dangerous boundary because the training demonstrations
showed that entering that boundary preceded successful goal attainment.

Eliminating this spurious correlation requires finding the
\textbf{divergence point} where the autonomous policy parted from the
nominal manifold, rather than the takeover point where the physical
override occurred. On a machine operating at a control frequency of
\(100\text{ Hz}\), human perceptual and neuromuscular latency introduces
a reaction delay between \(300\text{ ms}\) and \(800\text{ ms}\). During
those \(30\) to \(80\) execution cycles, the policy outputs divergent
actions that drive the physical state away from the target trajectory
while the logger continues recording them as autonomous operations. If
the training pipeline truncates the episode at the takeover timestamp
\(t_{\text{takeover}}\), those \(30\) to \(80\) corrupted transitions
remain in the dataset. To construct a clean training sample, the
pipeline must cut the autonomous rollout at the earlier divergence
timestamp \(t_{\text{div}}\) and discard the transitional interval from
\(t_{\text{div}}\) to \(t_{\text{takeover}}\). The recovery
demonstration recorded after \(t_{\text{takeover}}\) is retained as a
separate, corrective sequence.

Detecting \(t_{\text{div}}\) in offline ingestion pipelines requires
searching backwards in time from the physical override timestamp across
a synchronized pre-trigger telemetry buffer: \[\begin{aligned}
\hat{t}_{\text{div}} = \min \Bigl\{ t \in [t_{\text{takeover}} - \Delta t_{\text{lookback}}, t_{\text{takeover}}] \;\Big|\; &\|\mathbf{x}_{\text{plant}}(t) - \mathbf{x}_{\text{ref}}(t)\|_2 > \delta_{\text{tol}} \\
&\lor\; \text{Tr}(\boldsymbol{\Sigma}_{\pi}(t)) > \sigma_{\text{thresh}}^2 \Bigr\}.
\end{aligned}\] The pipeline scans backwards across a lookback window
(\(\Delta t_{\text{lookback}} = 1.0\text{--}2.0\text{ s}\)) until both
the kinematic tracking residual
\(\|\mathbf{x}_{\text{plant}} - \mathbf{x}_{\text{ref}}\|_2\) and the
policy ensemble action covariance \(\text{Tr}(\boldsymbol{\Sigma}_\pi)\)
return within nominal \(3\sigma\) confidence bounds. All transitions
recorded between \(\hat{t}_{\text{div}}\) and \(t_{\text{takeover}}\)
are pruned from the training corpus.

\textbf{Intervention curation} is the causal data-cleaning protocol that
truncates an autonomous rollout at its divergence timestamp
\(t_{\text{div}}\) rather than the physical override timestamp
\(t_{\text{takeover}}\), excising human perceptual reaction delay
(\(300\text{--}800\text{ ms}\)) and unmodeled drift transitions to
prevent learned policies from acquiring spurious attractor basins toward
safety boundaries.

The necessity of isolating divergence points and training on recovery
trajectories is governed by how compounding errors accumulate over an
operational horizon of \(T\) control steps
(\citeproc{ref-ross2011reduction}{Ross et al. 2011}).\sidenote{\footnotesize Dean
  Pomerleau's ALVINN (1989) first encountered behavioral cloning drift.
  Stéphane Ross and J. Andrew Bagnell formalized this in 2011 with
  DAgger (\citeproc{ref-ross2011reduction}{Ross et al. 2011}), proving
  that querying the expert over the learner's induced distribution
  collapses quadratic compounding error (\(O(T^2 \epsilon)\)) to linear
  bounds (\(O(T \epsilon)\)).} Let the autonomous policy \(\hat{\pi}\)
have an expected single-step decision error:
\[\epsilon = \mathbb{E}_{s \sim d^{\pi^*}} [\|\hat{\pi}(s) - \pi^*(s)\|]\]
relative to an expert policy \(\pi^*\). The mathematical contrast
between naive behavioral cloning and corrective aggregation follows a
clear three-step derivation:

\begin{enumerate}
\def\labelenumi{\arabic{enumi}.}
\tightlist
\item
  \textbf{Probability of First Excursion}: At time step \(t\), the
  probability that the policy made at least one mistake in the preceding
  \(t\) steps is \(1 - (1 - \epsilon)^t \le t \epsilon\).
\item
  \textbf{Unconstrained Loss Outside Support}: Under standard behavioral
  cloning trained only on nominal demonstrations, the expert provides
  zero supervision outside \(d^{\pi^*}\). Once the system exits the
  demonstrated tube, the expected downstream cost per step is
  unconstrained, bounded only by the maximum task penalty
  \(C_{\max} = 1\).
\item
  \textbf{Quadratic Cumulative Bound}: Integrating these failure
  probabilities across the full operational horizon of \(T\) steps
  yields a cumulative trajectory error that grows quadratically:
  \[\mathbb{E}[\text{Error}_{\text{BC}}] \le \sum_{t=1}^{T} t \epsilon = \epsilon \frac{T(T+1)}{2} = O(T^2 \epsilon).\]
\end{enumerate}

When interventions are logged, cut at \(t_{\text{div}}\), and ingested
through \textbf{corrective aggregation} (such as DAgger
(\citeproc{ref-ross2011reduction}{Ross et al. 2011})), the expert
provides demonstrations initialized directly from the empirical state
distribution \(d^{\hat{\pi}}\) induced by the running policy. Because
the dataset now contains corrective demonstrations for off-manifold
states, an error at step \(t\) does not cause unrecoverable divergence
on step \(t+1\). The policy learns to return to the nominal tube,
collapsing the per-step error to a constant \(\epsilon\) across the
entire horizon:
\[\mathbb{E}[\text{Error}_{\text{CA}}] \le \sum_{t=1}^T \epsilon = T \epsilon = O(T \epsilon).\]

For an assembly task running for \(T = 500\text{ steps}\)
(\(5.0\text{ s}\) at \(100\text{ Hz}\)) with a per-step error
probability of \(\epsilon = 0.01\), quadratic compounding yields a
theoretical error factor of \(T^2 \epsilon = 2500\), whereas corrective
aggregation holds the error factor to \(T \epsilon = 5\), representing a
\(500\times\) reduction in cumulative drift
(Figure~\ref{fig-05-compounding-error-flywheel}).

\begin{figure}[t]

\centering{

\includegraphics[width=1\linewidth,height=\textheight,keepaspectratio]{index_files/mediabag//Users/VJ/GitHub/MLSysBook-integrate-physical/physical/book/chapters/05-data/figures/fig05_compounding_error_flywheel.pdf}

}

\caption{\label{fig-05-compounding-error-flywheel}\textbf{Covariate
Shift Compounding Error Flywheel and Intervention Episode Curation.}
(Left) Cyclic feedback dynamics of naive Behavioral Cloning (BC) showing
\(O(T^2 \epsilon)\) quadratic error explosion as unobserved
perturbations trigger out-of-distribution commands, contrasted against
DAgger Corrective Aggregation (CA) which expands dataset support to
perturbed states and bounds trajectory drift to \(O(T \epsilon)\).
(Right) Trajectory phase space showing an autonomous policy
\(\hat{\pi}\) diverging from nominal path \(\pi^*\) at timestamp
\(t_{\text{div}}\), the subsequent human reaction delay window from
\(t_{\text{div}}\) to \(t_{\text{takeover}}\)
(\(300\text{--}800\text{ ms}\)) where divergent commands are falsely
logged as autonomous execution, the takeover point
\(t_{\text{takeover}}\), and the expert recovery trajectory returning
the state to the nominal tube. Preserving data integrity requires
truncating rollouts at \(t_{\text{div}}\), discarding the corrupted
drift window from \(t_{\text{div}}\) to \(t_{\text{takeover}}\), and
ingesting the recovery segment from \(t_{\text{takeover}}\) to
\(t_{\text{rec}}\) into the corrective training buffer. Source:
Theoretical formulation based on Ross et al.
(\citeproc{ref-ross2011reduction}{Ross et al. 2011}) and Spencer et al.
(\citeproc{ref-spencer2021feedback}{Spencer et al. 2021}).}

\end{figure}%

Deploying successive policy generations on physical hardware without
cutting divergence intervals triggers a retraining cascade:

\begin{enumerate}
\def\labelenumi{\arabic{enumi}.}
\tightlist
\item
  In \textbf{Generation 1}, a policy operates with a slight calibration
  error, causing a robotic arm to drift toward a joint torque limit
  until a human operator grabs the override pendant.
\end{enumerate}

\begin{itemize}
\tightlist
\item
  In \textbf{Generation 2}, a policy trained on these uncurated episodes
  treats the drifting states as if they were deliberate precursors to
  success. The resulting neural network develops phantom attractor
  basins near the intervention boundary, driving the arm aggressively
  toward the torque limit before attempting an abrupt, uncoordinated
  correction.
\item
  Operators respond by intervening earlier in subsequent deployments. By
  \textbf{Generation 3}, the intervention rate per operating hour
  increases while validation loss on the collected dataset continues to
  drop, creating a false metric of improvement while physical autonomy
  degrades.
\end{itemize}

Preventing the retraining cascade requires the nervous system to emit a
formal \textbf{intervention record} whenever a takeover or barrier trip
occurs.\sidenote{\footnotesize Under ISO 13849 PLe and ISO 26262 ASIL D,
  disengagements require deterministic hardware fault logging into
  non-volatile ring buffers before software task termination.} As
specified in the safety architecture of Chapter~\ref{sec-intervention},
this record cannot be a simple binary flag attached to a video stream.
It must preserve a circular pre-trigger buffer containing raw sensor
frames, actuator effort commands, and policy latent states spanning at
least \(2.0\text{ s}\) prior to disengagement. It must include a
monotonic hardware timestamp marking the exact microsecond the override
asserted, an identifier indicating whether the override came from a
human input device or a deterministic safety monitor, and the full
recovery trajectory until the system returns to the nominal scenario
envelope. Without this record, a training pipeline cannot compute the
divergence timestamp \(t_{\text{div}}\), cannot prune the transitional
error window, and cannot transform physical rescues into valid training
data.

\section{Embodiment and Transfer}\label{sec-data-5-6}

Data collected from one physical machine is not data about the task.
Every recorded trajectory binds the physical characteristics of the
collecting hardware into its numerical values, including link lengths,
joint velocity envelopes, sensor optical paths, communication delays,
and contact surface compliances. When an engineer saves an
observation-action trace \((o_t, a_t)\), the record does not capture an
abstract manipulation skill. It captures the \textbf{embodiment
contract} established between that specific machine and its environment
(governed by the ISO/IEC 5259 data quality standard series). This
contract encompasses the sensor geometry and optical calibration, the
sampling period and transport latency of the perception bus, the
coordinate frames and kinematic limits of the actuators, the
torque-speed saturation curves of the motors (including reflected rotor
inertia \(N^2 J_{\text{rotor}}\) in geared transmissions), and the
friction and compliance of the contact interface. Changing any parameter
in this contract alters the mapping between commanded inputs and
physical state transitions.

The \textbf{cross-embodiment compatibility contract} is the formal
geometric, dynamical, and temporal specification (such as those
underpinning the Open X-Embodiment
(\citeproc{ref-open_x_embodiment_collaboration2024open}{{Open
X-Embodiment Collaboration et al.} 2024}) and DROID
(\citeproc{ref-khazatsky2024droid}{{Khazatsky et al.} 2024}) datasets)
that defines the forward kinematics mapping, control rate normalization
(\(\Delta t\)), and actuator effort bounds required before multi-robot
trajectories can be transferred to a target physical plant.

Engineering a data reuse pipeline requires separating task-level
invariants from embodiment-specific realizations. A task-level invariant
is a geometric or dynamical condition required by the environment, such
as bringing a compliant probe into contact with a terminal pad at an
angle within \(5^\circ\) of the surface normal while maintaining a
normal force between \(2.0\text{ N}\) and \(5.0\text{ N}\). An
embodiment-specific realization is the sequence of joint encoder counts,
pulse-width modulation duty cycles, bus transport delays, and camera
pixel rasters that achieved that condition on a specific test rig. While
the contact target exists in world coordinates, the recorded dataset
contains only the realization. A downstream model trained directly on
those raw records learns the idiosyncratic dynamics of the source rig,
including its structural flex, drive backlash, and camera mount
vibration.

In this analytical example, because no bench measurement exists yet,
consider transferring a dataset recorded on Embodiment A to Embodiment B
for an identical touch-and-grasp task:

\begin{itemize}
\tightlist
\item
  \textbf{Embodiment A} uses a parallel-jaw gripper with a
  \(0\text{ to }100\text{ mm}\) stroke range, a \(50\text{ Hz}\) command
  rate (\(20\text{ ms}\) period), a continuous force limit of
  \(40\text{ N}\), and a wrist camera offset \(45\text{ mm}\) behind the
  tool center point with a \(15\text{ ms}\) transport delay.
\item
  \textbf{Embodiment B} uses a compact gripper with a
  \(0\text{ to }60\text{ mm}\) stroke range, a \(20\text{ Hz}\) command
  rate (\(50\text{ ms}\) period), a continuous force limit of
  \(15\text{ N}\), and a wrist camera offset \(20\text{ mm}\) behind the
  tool center point with a \(35\text{ ms}\) transport delay.
\end{itemize}

Mapping the dataset from Embodiment A through the limits of Embodiment B
reveals the points where transfer fails mechanically. For an object
width of \(75\text{ mm}\), Machine A opens to \(80\text{ mm}\), a
command that exceeds the mechanical travel of Machine B by
\(20\text{ mm}\) and drives its leadscrew against the physical endstop.
During the hold phase, Machine A applies \(30\text{ N}\) of clamping
force. That same command exceeds Machine B's continuous rating of
\(15\text{ N}\) by a factor of two, driving the motor amplifier into
current saturation. In the time domain, subsampling the \(50\text{ Hz}\)
trajectory onto Machine B's \(20\text{ Hz}\) bus aliases high-frequency
contact transitions and introduces a zero-order hold phase lag of
\(15\text{ ms}\). Combined with the \(20\text{ ms}\) increase in camera
transport delay (\(35\text{ ms} - 15\text{ ms}\)), Machine B perceives
contact events \(35\text{ ms}\) after they occur. At an approach
velocity of \(0.10\text{ m/s}\), this sensor-to-actuation latency
carries the gripper \(3.5\text{ mm}\) past the intended contact plane
before the first corrective command can be dispatched, while the
\(25\text{ mm}\) camera mounting disparity introduces visual parallax
that shifts the apparent position of the workpiece boundary in the image
frame.

Standard machine learning pipelines attempt to reconcile these physical
discrepancies by normalizing observation and action channels to the
interval \([0, 1]\) or \([-1, 1]\). This normalization produces a
cosmetic alignment of numerical ranges while leaving physical
incompatibility untouched. If an opening command is normalized by
dividing by the maximum stroke, a normalized command \(u = 0.75\)
corresponds to \(75\text{ mm}\) on Machine A, which clears a
\(70\text{ mm}\) workpiece. On Machine B, that same normalized command
\(u = 0.75\) commands an opening of \(45\text{ mm}\)
(\(0.75 \times 60\text{ mm}\)), causing the fingers to collide with the
workpiece during the approach trajectory. Normalization rescales
numbers, but it cannot expand actuator stroke, cannot lower electrical
resistance, and cannot eliminate transport latency. When the data
distributions of both embodiments are plotted in normalized space, the
demonstration trajectories appear identical, yet only the subset lying
within the physical intersection (openings under \(60\text{ mm}\),
forces under \(15\text{ N}\), and command bandwidths under
\(20\text{ Hz}\)) can execute on the target hardware without saturation
or impact damage.

\protect\phantomsection\label{callout-ruleux2a-5.7}
\begin{fbxSimple}{callout-rule}{Architect’s Rule of Thumb:}{The Non-Invariance of Unitless Scaling}
\phantomsection\label{callout-rule*-5.7}

Unitless range normalization (\([-1, 1]\)) aligns tensor dimensions on
disk but destroys physical invariants. Normalizing an actuator command
masks torque saturation, stroke limits, and transport delays,
guaranteeing dynamic failure when transferred to hardware with differing
reflected inertia or bus latency.

\end{fbxSimple}

\subsection{Cross-Embodiment Datasets and the Open X-Embodiment
Frontier}\label{cross-embodiment-datasets-and-the-open-x-embodiment-frontier}

The emergence of large-scale multi-robot datasets, spearheaded by the
\textbf{Open X-Embodiment (OXE) initiative}
(\citeproc{ref-open_x_embodiment_collaboration2024open}{{Open
X-Embodiment Collaboration et al.} 2024}) and the \textbf{DROID dataset}
(\citeproc{ref-khazatsky2024droid}{{Khazatsky et al.} 2024}), highlights
both the potential and the systems bottlenecks of cross-embodiment
learning. The Open X-Embodiment repository aggregates over
\(1{,}000{,}000\text{ demonstration trajectories}\) collected across 22
distinct robot embodiments from more than 20 academic and industrial
institutions, powering generalist robot policies such as RT-X, Octo, and
OpenVLA. However, ingesting heterogeneous physical telemetry into a
single training corpus exposes four fundamental systems mismatches:

\begin{enumerate}
\def\labelenumi{\arabic{enumi}.}
\item
  \textbf{Action Space Incompatibility}: A canonical 7-DoF articulated
  research arm records joint positions \(\mathbf{q} \in \mathbb{R}^7\),
  a bimanual leader-follower workstation records dual 6-DoF arm vectors
  \(\mathbf{q} \in \mathbb{R}^{14}\), and a mobile manipulator logs base
  planar velocities \((\dot{x}, \dot{y}, \dot{\theta})\) paired with a
  5-DoF arm. Because joint-space representations are strictly
  non-transferable across different kinematic topologies,
  multi-embodiment architectures are forced to project actions into a
  universal Cartesian end-effector delta representation:
  \[\Delta \mathbf{x}_t = (\Delta x, \Delta y, \Delta z, \Delta \phi_{\text{roll}}, \Delta \phi_{\text{pitch}}, \Delta \phi_{\text{yaw}}) \in SE(3), \quad a_{\text{gripper}} \in [0, 1].\]
  While Cartesian deltas provide geometric standardization, they discard
  actuator dynamics. A Cartesian displacement
  \(\Delta \mathbf{x} = [0.03\text{ m}, 0.0, 0.0]\) that is trivial for
  an industrial arm near its nominal configuration will drive a
  lightweight manipulator through a kinematic singularity, commanding
  infinite joint velocities that trigger immediate overcurrent
  protection.
\item
  \textbf{Control Loop Frequency Discrepancies}: Across the Open
  X-Embodiment datasets, recording frequencies vary widely, spanning
  \(3\text{ Hz}\) (e.g., Bridge-v2), \(10\text{ Hz}\) (e.g., Google
  Fractal), \(20\text{ Hz}\) (e.g., DROID), and \(50\text{ Hz}\) (e.g.,
  canonical bimanual high-frequency datasets). A step delta
  \(\Delta x = 20\text{ mm}\) recorded at \(3\text{ Hz}\) corresponds to
  a linear velocity of \(0.06\text{ m/s}\), whereas the same numerical
  delta at \(50\text{ Hz}\) corresponds to a violent \(1.0\text{ m/s}\)
  stroke. If an ingestion pipeline trains a policy on unscaled delta
  tokens without explicit timebase conditioning (\(\Delta t\)), the
  learned model emits chaotic velocity commands that depend on which
  dataset dominated the training mini-batch.
\item
  \textbf{Perception and Viewpoint Disparity}: Heterogeneous datasets
  combine disparate camera viewpoints, including high-mounted
  third-person room cameras, angled workcell cameras, and wrist-mounted
  eye-in-hand cameras, each with distinct optical focal lengths, sensor
  aspect ratios, and dynamic exposure ranges. Training an invariant
  visual representation requires multi-view camera extrinsic calibration
  and coordinate frame reprojection
  (\citeproc{ref-lynch2017modern}{Lynch and Park 2017}).
\item
  \textbf{Gripper Mechanics and Contact Compliance}: Gripper actions
  range from binary open/close flags (pneumatic chucks) to continuous
  position setpoints (\(0\text{--}100\text{ mm}\)) and force-controlled
  tendon tensioning. Furthermore, differences in contact pad compliance
  (e.g., soft silicone vs rigid 3D-printed PLA) alter the friction
  coefficient \(\mu\) by a factor of four, rendering identical gripping
  force setpoints either secure or prone to catastrophic slip.
\item
  \textbf{Null-Space and Kinematic Redundancy Projection}: When mapping
  a 7-DoF articulated arm (e.g., a 7-DoF torque-controlled manipulator)
  or a 14-DoF bimanual platform into a 6-DoF Cartesian delta
  (\(\Delta \mathbf{x} \in SE(3)\)), the 1-DoF null-space redundancy
  (the arm angle or elbow swivel parameter \(\psi\)) is projected out
  and lost. When the learned Cartesian policy executes on target
  hardware, the local inverse kinematics solver must resolve the
  unconstrained null-space independently (e.g., via joint-limit
  avoidance or manipulability maximization). Without the original
  demonstrator's elbow posture bias, the target arm can drift into
  mechanical joint limits or sweep its elbow into surrounding fixtures
  during maneuvers that were completely collision-free in the source
  demonstrations.
\end{enumerate}

Reusing physical data across embodiments requires a formal compatibility
check at the ingestion boundary before any sample enters the training
buffer. The ingestion pipeline must transform the recorded trace through
forward kinematics to evaluate coordinate consistency
(\citeproc{ref-lynch2017modern}{Lynch and Park 2017}), verify that all
commanded joint velocities and accelerations lie within the target
machine's dynamic limits (\citeproc{ref-tedrake2023manipulation}{Tedrake
2023}), and verify that perceptual features match the target sensor's
field of view and resolution. Actions that exceed kinematic reach or
actuator effort limits cannot simply be clamped at execution time;
clamping distorts the velocity profile and destroys the causal
relationship between observed errors and corrective motions. Samples
that pass the compatibility check are admitted; samples that can be
geometrically re-targeted through inverse kinematics are relabeled; and
samples that demand unavailable physical authority or unobservable
states are rejected. When an ingestion pipeline admits unverified
records across an embodiment boundary, it injects invalid dynamics into
the model's training distribution. The resulting policy does not fail
because the neural network underfitted the training loss. It fails
because the training dataset contained actions that were physically
sound only on the machine that collected them.

\section{How Datasets Go Wrong}\label{sec-data-5-7}

A physical dataset contains no internal markers to indicate what it
fails to measure. When an engineer inspects a static corpus of logged
robot interactions, standard distribution metrics like feature variance,
sample density, and covariance evaluate only the states the machine
actually reached. They cannot detect states the machine never
encountered. Consider an empirical dataset
\(\mathcal{D} = \{(\mathbf{s}_k, \mathbf{a}_k)\}\) collected during
peg-in-hole assembly, where all recorded contact forces fall within
\(f_z \in [0.5, 2.0]\text{ N}\) and lateral offsets within
\(\Delta x \in [-1.5, 1.5]\text{ mm}\). If Deployment Requirement A
defines a nominal operational envelope \(\mathcal{S}_{\text{req}, A}\)
bounded by \(f_z \le 2.0\text{ N}\) and
\(|\Delta x| \le 1.5\text{ mm}\), the dataset appears fully populated.
If Deployment Requirement B instead specifies
\(\mathcal{S}_{\text{req}, B}\) with \(f_z \le 8.0\text{ N}\) and
\(|\Delta x| \le 5.0\text{ mm}\) to account for pallet tolerances,
\(\mathcal{D}\) leaves large regions of the required state space
completely unvisited. Yet \(\mathcal{D}\) itself is identical in both
cases. Because the recorded data distribution does not change when the
deployment specification expands, internal analysis of \(\mathcal{D}\)
cannot establish whether coverage is sufficient. Coverage failure is an
external mismatch between the required operational envelope
\(\mathcal{S}_{\text{req}}\) partitioned into discrete evaluation cells
and the empirical visitation distribution \(\mathcal{S}_{\text{obs}}\).
Without an independent, external specification of
\(\mathcal{S}_{\text{req}}\), coverage failure cannot be detected.

When training data originates from human teleoperation, the data
collection policy drifts over time as human operators adapt to the
interface. Early in a collection campaign, an operator moves
deliberately, introducing corrective sub-movements and pauses while
searching for visual cues. After hundreds of repetitions, the same
operator develops muscle memory, increasing speed, rounding corners, and
executing aggressive feedforward maneuvers. In an analytical example of
\(500\) teleoperated assembly demonstrations gathered over two weeks,
median completion time drops from \(12.4\text{ s}\) in the first \(50\)
episodes to \(5.8\text{ s}\) in the final \(50\), while the \(P_{99}\)
commanded angular velocity increases from \(0.42\text{ rad/s}\) to
\(1.85\text{ rad/s}\). The resulting late-stage demonstrations rely on
dynamic momentum and precise timing that the physical robot cannot
replicate when running under closed-loop latency at deployment.
Detecting operator drift requires partitioning the dataset
chronologically into sliding temporal windows and comparing distribution
shifts across episode duration, corrective jerk count, mean command
magnitude, and teleoperation intervention rate. If these metrics drift
across windows, treating the dataset as an identically distributed
sample pool injects multi-modal, non-stationary behavior into the
policy.

File containers such as HDF5, Zarr, or Apache Parquet enforce tabular
schema consistency, but they conceal hardware state transitions that
occur between recording runs. A robot arm may be shut down overnight,
recalibrated, or bumped by a technician, altering the zero-position
offsets of its joint encoders or the extrinsics of an overhead camera.
If a camera mount shifts by \(3\text{ mm}\), the visual position of the
target workpiece changes relative to the robot base, yet the resulting
image array matches the recorded tensor shape and type definitions
perfectly. Similarly, clock synchronization drifts across distributed
microcontrollers over long sessions, causing the actual temporal offset
between a force-torque measurement and a camera frame to vary by tens of
milliseconds.\sidenote{\footnotesize While high-end perception systems use IEEE 1588
  Precision Time Protocol (PTP) for sub-microsecond hardware
  timestamping across Ethernet subnets, uncoordinated USB and CAN buses
  introduce variable OS scheduling jitter that silently corrupts
  cross-sensor temporal alignment.} As illustrated in
Figure~\ref{fig-05-camera-calibration}, multi-camera demonstration rigs
deployed across heterogeneous environments must continuously audit the
extrinsic spatial transformation
\(\mathbf{T}_{\text{cam} \rightarrow \text{base}}\) by projecting the
robot's kinematic URDF (Unified Robot Description Format) wireframe
directly over raw camera streams. Detecting calibration and timebase
discontinuities requires joining every recorded episode against an
immutable hardware audit trail containing sensor calibration
identifiers, optical reference checks, and clock synchronization health
metrics. When a dataset is loaded without verifying that every
trajectory shares a common, valid calibration identifier, the learning
algorithm is forced to fit conflicting physical mappings within
identical numerical coordinates.

\begin{figure}[t]

\centering{

\includegraphics[width=0.95\linewidth,height=\textheight,keepaspectratio]{/Users/VJ/GitHub/MLSysBook-integrate-physical/physical/book/chapters/05-data/figures/fig05_real_camera_calibration.jpg}

}

\caption{\label{fig-05-camera-calibration}\textbf{Multi-Camera-to-Robot
Extrinsic Calibration and Visual Verification in Diverse Demonstration
Environments.} Data collection across heterogeneous physical
environments requires continuous validation of the extrinsic spatial
transformation \(\mathbf{T}_{\text{cam} \rightarrow \text{base}}\)
between camera optical frames and the robot base frame. Six
representative real-world collection scenes from the DROID platform
depict autonomous extrinsic calibration verification using 3D kinematic
reprojection (CtRNet-X / PyTorch3D)
(\citeproc{ref-khazatsky2024droid}{{Khazatsky et al.} 2024}). Forward
kinematics computed from optical joint encoders render the robot's URDF
mesh directly over the camera video stream. If mechanical vibration or
physical disturbances displace camera brackets by even a few
millimeters, the visual discrepancy between the kinematic wireframe and
the physical arm exposes extrinsic drift before corrupted trajectories
enter the training pipeline. Documenting immutable calibration hashes in
the dataset record prevents learning algorithms from attempting to fit
contradictory spatial mappings within identical numerical coordinates.
Source: Adapted from Khazatsky et al.
(\citeproc{ref-khazatsky2024droid}{{Khazatsky et al.} 2024}), CC BY
4.0.}

\end{figure}%

Physical data collected across multiple nominally identical workstations
frequently harbors silent embodiment mismatch. Two robotic arms from the
same manufacturer can have different gearbox wear, thermal dissipation
profiles, or firmware velocity filters. In an analytical comparison, Rig
1 with new strain-wave gears exhibits \(0.02^\circ\) backlash, whereas
Rig 2 after \(2000\text{ hours}\) of operation exhibits \(0.35^\circ\)
backlash and a \(22\text{ ms}\) phase lag in joint torque response.
Actions recorded on Rig 2 include compensatory lead times and higher
gain commands that produce severe oscillations when executed on Rig 1.
Detecting embodiment mismatch requires verifying a machine fingerprint
before any trajectory enters the training pipeline. This fingerprint
binds the kinematics, actuator serial numbers, firmware versions, and
dynamic transfer function models to the recorded data file. If a
trajectory fails the interface contract for the target deployment
hardware, the ingestion pipeline must flag and exclude it before
training begins.

A physical collection pipeline generates two distinct classes of failed
data that must not be conflated. A failed task attempt occurs when the
machine executes a complete physical sequence but fails to achieve the
objective, such as dropping a connector during insertion. A failed
recording occurs when the logging infrastructure drops network packets,
corrupts an image frame, or experiences a buffer overrun on the logging
storage bus. Discarding failed task attempts creates severe survivorship
bias, stripping the dataset of the recovery behaviors and error
distributions needed to train a stabilizing policy. Discarding or
retaining failed recordings without accounting corrupts the training
timebase, presenting non-uniform sample periods \(\Delta t\) as constant
control intervals. To prevent both failure modes, the data ingestion
pipeline must enforce strict reconciliation accounting where the total
attempt count equals the sum of retained successful trajectories,
retained failed task trajectories, and excluded corrupted recordings,
with every exclusion tied to a documented hardware or network error
code.

The canonical failure mode of silent coordinate and unit contract
mismatch is documented in the loss of the Mars Climate Orbiter
(\citeproc{ref-nasa1999mco}{NASA Mishap Investigation Board 1999}).
Ground software computed angular momentum desaturation thruster impulses
in imperial pound-force-seconds (\(\text{lbf}\cdot\text{s}\)), while the
onboard autonomous navigation software ingested the velocity change
values assuming metric Newton-seconds (\(\text{N}\cdot\text{s}\)).
Because \(1\text{ lbf}\cdot\text{s} = 4.45\text{ N}\cdot\text{s}\), the
onboard trajectory integration underestimated the cumulative velocity
adjustment by a factor of \(4.45\) over nine months of cruise. The
numerical files satisfied all schema and format checks, yet the silent
unit contract mismatch accumulated an unobserved \(170\text{ km}\)
orbital trajectory error, causing the spacecraft to enter the Martian
atmosphere at an altitude of \(57\text{ km}\) rather than its intended
\(226\text{ km}\) periapsis, resulting in atmospheric breakup.

When a silent coordinate or unit failure passes through an ingestion
pipeline undetected, the resulting policy failure is difficult to
diagnose from model loss curves alone. In multi-rig robotic collection,
a minor physical shift in sensor mounting or an uncalibrated
tool-center-point transform introduces a systematic coordinate
discrepancy into the training corpus. The neural network optimizes its
objective by averaging the conflicting kinematic demonstrations,
producing a policy that converges mathematically on training loss but
generates indecisive, oscillatory motions at the physical workpiece
interface. Because the model was trained on data containing conflicting
physical ground truths, the failure manifests on the real machine as an
unexpected limit-cycle oscillation rather than a detectable numerical
error.

\protect\phantomsection\label{callout-incidentux2a-5.8}
\begin{fbxSimple}{callout-incident}{Landmark Engineering Incident:}{Cruise Robotaxi Pedestrian Drag (2023) and Missing Recovery Covariate Shift}
\phantomsection\label{callout-incident*-5.8}

On October 2, 2023, in San Francisco, California, an autonomous robotaxi
operated by Cruise LLC was involved in a severe multi-agent collision
that resulted in catastrophic physical injury and the subsequent
grounding of the entire commercial fleet
(\citeproc{ref-nhtsa2023cruise}{National Highway Traffic Safety
Administration (NHTSA) 2023}).\sidenote{\footnotesize National Highway Traffic Safety
  Administration (NHTSA), \emph{Office of Defects Investigation (ODI)
  Preliminary Evaluation PE23-018: Automated Driving System Post-Crash
  Maneuver Collision and Pedestrian Injury}, San Francisco, CA, October
  2023; California Department of Motor Vehicles (DMV), \emph{Order of
  Suspension for Cruise LLC Autonomous Vehicle Deployment and Testing
  Permits}, October 24, 2023.}

\emph{The physical sequence.} A human-driven vehicle traveling in an
adjacent lane struck a pedestrian crossing against the traffic signal,
violently catapulting her across the lane divider directly into the path
of the Cruise autonomous vehicle (AV). The AV executed an emergency
stop, applying \(1.4g\) maximum deceleration and stopping in
\(0.46\text{ s}\), coming to rest directly above the pedestrian who was
now pinned underneath the vehicle's floorpan and rear axle.

\emph{The catastrophic covariate drift.} Following the primary emergency
stop, the onboard contextual planner evaluated its state: the vehicle
was stationary in an active urban traffic lane. The autonomous driving
policy executed its nominal post-collision ``pull-over'' maneuver to
clear the travel lane and pull to the curb. In doing so, the AV
accelerated forward, dragging the trapped pedestrian underneath the
chassis for \(20\text{ feet}\) (\(6.1\text{ meters}\)) at speeds up to
\(7\text{ mph}\) (\(11\text{ km/h}\)) before coming to a final stop.

\emph{Root-cause systems and dataset failure analysis.}

\begin{enumerate}
\def\labelenumi{\arabic{enumi}.}
\item
  \textbf{The Absent Recovery State in Training Distribution
  (\(d^{\pi_{\text{col}}}(s) = 0\)):} The autonomous driving policy had
  been trained on millions of miles of nominal driving and thousands of
  simulated and logged roadside pull-over maneuvers. However, the
  training corpus contained \emph{zero} demonstration trajectories or
  recovery maneuvers representing states where an injured human was
  lodged underneath the vehicle chassis.
\item
  \textbf{Occlusion Tracker Blind Spot and Covariate Shift:} When the
  pedestrian was launched underneath the chassis, she passed entirely
  out of the line-of-sight of all vehicle LiDARs, RADARs, and perimeter
  cameras. The perception state estimator, lacking an under-chassis
  tactile or ultrasonic ground-plane sensor, cleared the pedestrian
  bounding-box track after several frames of total occlusion. The
  high-level planner evaluated the scene assuming the road ahead was
  clear.
\item
  \textbf{Execution Outside the Empirical Support Envelope without
  Refusal:} Because empirical risk minimization produces arbitrary,
  ungrounded actions outside the training support
  \(d^{\pi_{\text{col}}}\), the autonomous policy evaluated its nominal
  pull-over routine. The vehicle lacked a deterministic safety barrier
  or hardware governor prohibiting vehicle re-engagement after an airbag
  deployment or severe collision event without explicit remote human
  operator clearance.
\end{enumerate}

\emph{The physical AI lesson.} The Cruise incident is a tragic
demonstration of compounding covariate shift in physical systems. When a
machine enters an extreme edge state that was omitted during dataset
collection, naive policy execution does not default to safe
passivity---it executes nominal primitives that can convert a survivable
incident into a fatal mechanical interaction. Autonomous machines must
maintain an immutable scenario support ledger and enforce deterministic
control refusal when sensor occlusions and state uncertainty cross the
empirical boundary.

\end{fbxSimple}

No statistical operation performed on a tensor array can verify that the
underlying machine was calibrated, that the human demonstrator did not
drift, or that unobserved regions of the workspace remain safe. A
physical dataset is reliable only when its collection history and
physical unit contracts are recorded as rigorously as the numerical
features themselves.

A \textbf{support envelope relational join} is the set-theoretic
cross-referencing between an empirical state-action occupancy
distribution \(d^{\pi_{\text{col}}}(s,a)\) and an external scenario
ledger \(\mathcal{S}_{\text{req}}\) that partitions the operational
design domain into supported nominal regimes
(\(N_{\text{cell}} \ge N_{\text{min}}\)), high-variance transition
zones, and zero-support blind spots (\(N_{\text{cell}} = 0\)) where
autonomous control authority must be refused.

\section{The Dataset Schema}\label{sec-data-5-8}

The general notebook schema, cryptographic hashing, and immutable
logging contracts established in Chapter~\ref{sec-nervous} govern how a
machine records its execution events. This chapter adds the fields that
bind physical movement to a training dataset. A raw array of joint
trajectories and camera frames cannot speak for its own physical
validity. Without an explicit data record, an ingestion pipeline treats
numbers recorded on a vibrating, misaligned rig as equivalent to numbers
collected on a rigid, calibrated testbed. The data record documents what
was collected, by whom, on which physical instance, and under what
operational limits, establishing the boundary between valid training
evidence and unmodeled mechanical error.

The \textbf{demonstration provenance record} is the immutable metadata
tuple
\((a_{\text{req}}, a_{\text{cmd}}, a_{\text{enf}}, a_{\text{meas}}, \mathbf{h}_{\text{plant}}, \tau_{\text{bus}}, \mu_{\text{contact}})\)
logged at every control step that binds behavioral trajectories to
specific actuator kinematics, sensor calibration hashes, bus latencies,
and contact wear states across the Dual-Brain MPU/MCU boundary.

\protect\phantomsection\label{callout-ruleux2a-5.9}
\begin{fbxSimple}{callout-rule}{Architect’s Rule of Thumb:}{The Absence Ledger Principle}
\phantomsection\label{callout-rule*-5.9}

A physical dataset is defined as much by what it omits as by what it
contains. If a dataset does not explicitly catalog its unvisited angles,
unattempted payloads, and excluded failure trips in an immutable support
ledger, downstream verification must treat the missing physics as
actively unsafe rather than successfully mastered.

\end{fbxSimple}

Every trajectory in the training corpus must identify its collection
source, collector policy, and execution authority. When data originates
from human teleoperation, the record logs the operator identifier, the
interface hardware, and the transport timing, such as a 6-DOF master arm
sampled at \(100\text{ Hz}\) with \(0.8\text{ ms}\) of input jitter.
When data originates from an autonomous or mixed controller, the record
stores the exact policy checkpoint, the planner version, and the active
enforcement configuration, including Cartesian bounding boxes, joint
velocity limits, and contact-force clamping thresholds. Trajectories
must also record the active retention rule, documenting whether an
episode was committed automatically upon task completion, retained after
an explicit operator review, or flagged for inclusion as an instructive
failure recovery.

A physical machine cannot be identified by a generic model name alone.
The data record specifies the individual plant through its extrinsic and
intrinsic sensor calibrations, action contract, motor drive firmware
revisions, and component wear state. Replaceable contact surfaces alter
the governing dynamics. A parallel-jaw gripper fitted with compliant
\(30\text{ A}\) durometer silicone pads generates contact compliance and
friction characteristics that differ entirely from the same gripper
fitted with worn \(60\text{ A}\) polyurethane pads under identical joint
commands. Maintenance milestones, such as cable retensioning, drive belt
replacements, gearbox swaps, and camera remounts, are logged directly
into the machine state record. If an arm underwent joint recalibration
between morning and afternoon collection shifts, the dataset must expose
that boundary so downstream training does not attempt to resolve two
conflicting kinematic mappings into a single policy.

Accounting for dataset volume requires explicit denominators across both
time and trial counts. Reporting only the count of retained frames or
successful episodes obscures the true yield of the collection process.
The data record tracks scheduled cell time, active motion time,
attempted episodes, retained successes, retained task failures, excluded
recordings, and the count of independent physical task configurations.
In an analytical example of a precision peg-insertion campaign scheduled
for \(40.0\text{ h}\) of cell time, \(28.5\text{ h}\) of net motion time
produces 1000 total attempts across 12 distinct workpiece fixtures. The
resulting tally records 820 retained successes, 110 retained task
failures containing recovery attempts, and 70 excluded recordings
rejected due to logging packet drops. Dividing motion time by attempted
episodes yields an average cycle budget of \(102.6\text{ s}\) per
attempt, exposing the actual operational pacing and the proportion of
time spent resetting the physical environment.

A single action column in a dataset conceals the causal filtering
performed by the nervous system. The data record preserves four distinct
action streams at every control step: the requested action from the
operator or learned policy, the mapped action produced after coordinate
transforms and kinematics, the enforced action output by the safety
filter after clamping and barrier evaluations, and the measured action
realized by the physical actuators. Logging these streams alongside
boolean override flags, truncation markers, and numeric exclusion codes
reveals where the collector commanded kinematically impossible motions,
where the safety enforcer intervened to prevent collisions, and where
tracking lag separated command from execution. Discarding the
intermediate representations leaves the training pipeline unable to
distinguish between an intentional human demonstration and an action
heavily modified by an underlying governor.

Every physical dataset operates within a bounded envelope, and the data
record must explicitly catalog both supported conditions and known
absences. Supported conditions define the verified operating domain,
encompassing payload masses between \(0.10\text{ kg}\) and
\(1.50\text{ kg}\), ambient illumination from \(300\text{ lx}\) to
\(1200\text{ lx}\), workbench surface friction coefficients, and
workpiece dimensional tolerances. Known absences define what was omitted
or unattempted: unvisited workspace volumes, approach angles outside a
\(30^\circ\) cone, ambient temperatures outside the range of
\(18\text{ }^\circ\text{C}\) to \(26\text{ }^\circ\text{C}\), excluded
failure modes such as unrecoverable motor stalls, missing operators who
did not participate in the shift, and hardware embodiments that lack
kinematic parity. Stating what is absent prevents an engineer from
treating the absence of training failures as evidence of policy
competence. The complete metadata schema and physical contract fields
required to preserve this dataset provenance are formalized in
Table~\ref{tbl-05-provenance-schema}, defining the metadata structures,
calibration contracts, multi-stream action tap points, and operational
envelope boundaries required before any trajectory is admitted into
training pipelines.

\begin{longtable}[]{@{}
  >{\raggedright\arraybackslash}p{(\linewidth - 8\tabcolsep) * \real{0.2000}}
  >{\raggedright\arraybackslash}p{(\linewidth - 8\tabcolsep) * \real{0.2000}}
  >{\raggedright\arraybackslash}p{(\linewidth - 8\tabcolsep) * \real{0.2000}}
  >{\raggedright\arraybackslash}p{(\linewidth - 8\tabcolsep) * \real{0.2000}}
  >{\raggedright\arraybackslash}p{(\linewidth - 8\tabcolsep) * \real{0.2000}}@{}}
\caption{\textbf{Demonstration Dataset Provenance Schema for Physical AI
Systems.} Formal specification of metadata blocks, data types, physical
invariants, and enforcement obligations required to bind physical
movement logs to training pipelines (extending ISO/IEC 5259 and the
Chapter~\ref{sec-nervous} notebook schema). Trajectories omitting these
fields fail ingestion compatibility
checks.}\label{tbl-05-provenance-schema}\tabularnewline
\toprule\noalign{}
\begin{minipage}[b]{\linewidth}\raggedright
Provenance Domain
\end{minipage} & \begin{minipage}[b]{\linewidth}\raggedright
Schema Field / Data Structure
\end{minipage} & \begin{minipage}[b]{\linewidth}\raggedright
Representation \& Units
\end{minipage} & \begin{minipage}[b]{\linewidth}\raggedright
Physical \& Invariant Purpose
\end{minipage} & \begin{minipage}[b]{\linewidth}\raggedright
Ingestion Enforcement \& Downstream Consumer
\end{minipage} \\
\midrule\noalign{}
\endfirsthead
\toprule\noalign{}
\begin{minipage}[b]{\linewidth}\raggedright
Provenance Domain
\end{minipage} & \begin{minipage}[b]{\linewidth}\raggedright
Schema Field / Data Structure
\end{minipage} & \begin{minipage}[b]{\linewidth}\raggedright
Representation \& Units
\end{minipage} & \begin{minipage}[b]{\linewidth}\raggedright
Physical \& Invariant Purpose
\end{minipage} & \begin{minipage}[b]{\linewidth}\raggedright
Ingestion Enforcement \& Downstream Consumer
\end{minipage} \\
\midrule\noalign{}
\endhead
\bottomrule\noalign{}
\endlastfoot
\textbf{Session \& Authority} & \texttt{session\_authority} &
\texttt{\{collector\_type:\ enum,\ operator\_id:\ str,\ policy\_hash:\ sha256,\ sampling\_rate\_hz:\ uint16,\ p99\_jitter\_ms:\ float32\}}
& Records commanding authority (teleop/scripted/autonomous),
demonstrator identity, and bus timing jitter &
Chapter~\ref{sec-intervention} Operator Audit; rejects sessions with
\(P_{99}\text{ jitter} > 2.0\text{ ms}\) or uncertified operators \\
\textbf{Kinematic Plant \& Embodiment} & \texttt{embodiment\_contract} &
\texttt{\{robot\_serial:\ str,\ dh\_parameters:\ float32{[}N,4{]},\ backlash\_deg:\ float32{[}N{]},\ gear\_ratios:\ uint16{[}N{]},\ firmware\_build:\ str\}}
& Binds trajectory to exact link geometry, harmonic drive backlash
(\(\Delta \theta_{\text{backlash}}\)), and motor firmware filters &
Ingestion Forward Kinematics Validator; verifies dynamic feasibility and
joint angle limits against target hardware \\
\textbf{Perception \& Calibration State} & \texttt{sensor\_calibration}
&
\texttt{\{camera\_serials:\ str{[}M{]},\ intrinsic\_K:\ float32{[}M,3,3{]},\ extrinsic\_T:\ SE(3){[}M{]},\ timebase\_sync:\ enum(PTP,\ NTP),\ ft\_bias\_null:\ float32{[}6{]}\}}
& Documents optical intrinsics, eye-in-hand extrinsics, clock
synchronization (IEEE 1588 PTP), and static gravity wrench offsets &
Spatial Perception Pipeline; rejects data with expired extrinsic
calibration (\(>24\text{ h}\)) or clock drift \(>5\text{ ms}\) \\
\textbf{Contact Interface \& Wear State} & \texttt{contact\_wear\_state}
&
\texttt{\{end\_effector\_sn:\ str,\ pad\_material:\ str,\ shore\_hardness:\ uint8,\ friction\_coeff:\ float32,\ flex\_cycle\_count:\ uint32,\ maintenance\_epoch:\ uint16\}}
& Logs fingertip elastomer durometer (\(30\text{ A}\) silicone
vs.~\(60\text{ A}\) polyurethane), friction (\(\mu\)), and accumulated
flex fatigue & Dynamics Transfer Engine; flags contact demonstrations as
non-transferable if elastomer compliance shifted \(>15\%\) \\
\textbf{Four-Stream Action Tap Tuple} & \texttt{action\_stream\_tuple} &
\texttt{\{a\_req:\ float32{[}D{]},\ a\_cmd:\ float32{[}D{]},\ a\_enf:\ float32{[}D{]},\ a\_meas:\ float32{[}D{]},\ override\_flag:\ bool,\ clamp\_code:\ uint8\}}
& Preserves full causal chain: raw operator intent, mapped setpoint,
safety-barrier clamped command, and measured actuator response &
Behavioral Cloning Loss Aggregator; isolates safety governor
interventions and excludes clamped transitions from nominal imitation
loss \\
\textbf{Scenario Ledger \& Support Bounds} & \texttt{support\_envelope}
&
\texttt{\{payload\_bounds\_kg:\ {[}float32,\ 2{]},\ reach\_cone\_deg:\ float32,\ illum\_lux:\ {[}uint16,\ 2{]},\ temp\_c:\ {[}float32,\ 2{]},\ known\_absences:\ str{[}{]}\}}
& Defines explicit operational design domain (ODD) support bounds and
catalogs unattempted failure modes, workspace blind spots, and missing
angles & Chapter~\ref{sec-release} Safety Case Release; ensures
deployment state queries remain within empirically supported convex
hull \\
\end{longtable}

A short record audit demonstrates why two files with identical tensor
schemas and sample counts cannot support the same deployment claims.
Consider an analytical audit of two dataset files, File A and File B,
each containing 500,000 state-action pairs recorded at \(50\text{ Hz}\)
(\(2.78\text{ h}\) of motion) for a dual-arm assembly task:

\begin{itemize}
\tightlist
\item
  \textbf{File A} records four certified operators across eight fixture
  layouts, reports zero safety enforcer overrides, and verifies
  consistent fingertip durometer across all trajectories.
\item
  \textbf{File B} records a single operator at a single fixed jig, logs
  620 enforcer interventions where joint velocity limits clamped the
  requested trajectories, and notes a sensor re-zeroing event midway
  through collection without recalibrating the extrinsics.
\end{itemize}

While an automated validation script checking tensor shapes and value
ranges accepts both files equally, File A supports a deployment claim
across variable workstations, whereas File B documents an
overconstrained demonstration corrupted by unmodeled kinematic shifts.

Downstream lifecycle stages consume the data record for distinct
engineering guarantees:

\begin{itemize}
\tightlist
\item
  The human and override provenance feeds directly into the supervisory
  auditing of Chapter~\ref{sec-intervention}, where safety engineers
  evaluate demonstrator consistency, teleoperation latency tails, and
  the frequency of governor interventions to assess human reliability.
\item
  The coverage boundaries, exclusion ledgers, and embodiment calibration
  histories pass forward to Chapter~\ref{sec-release}, where they form
  the empirical baseline for the safety case, verifying that the
  training distribution matches the target operational design domain
  before the learned policy is granted execution authority on a physical
  machine.
\end{itemize}

\section{Systems Stress-Tests}\label{sec-data-5-9}

A dataset that satisfies every static schema validator and passes every
tensor shape assertion can remain completely invalid for physical
deployment. In physical systems, data quality is not a property of
tensor shape consistency, missing-value percentages, or raw sample
volume. Because physical data is gathered by moving a mechanical body
through a workspace under the authority of a specific collecting policy,
the recorded trajectories reflect only the narrow manifold of states
that the collector visited, the physical parameters of the collection
rig, and the filtering rules applied during logging. Evaluating a
candidate training corpus requires testing it against the external
physical requirements of the operational design domain rather than
inspecting its internal statistics. The following four analytical
stress-tests evaluate datasets that appear large and clean while
concealing systematic operational deficits:

\subsection*{Stress-Test 1: The Volume Illusion in Repetitive
Demonstrations (Nominal Attractor
Overfitting)}\label{stress-test-1-the-volume-illusion-in-repetitive-demonstrations-nominal-attractor-overfitting}
\addcontentsline{toc}{subsection}{Stress-Test 1: The Volume Illusion in
Repetitive Demonstrations (Nominal Attractor Overfitting)}

\emph{The scenario.} An articulated six-axis manipulator executes a
tactile surface-tracking task. The operational scenario ledger requires
verifying performance across 500 discrete cells spanning contact normal
forces (\(F_N = 2.0\text{--}10.0\text{ N}\)), surface friction
(\(\mu = 0.10\text{--}0.40\)), and contact angles
(\(\theta = -10.0^\circ\text{ to }+10.0^\circ\)). The baseline archive
contains \(200{,}000\) control steps across 45 cells. An operator
completes a \(1.0\text{ h}\) (\(3600\text{ s}\)) session repeating the
nominal motion
(\(v = 0.05\text{ m/s}, F_N = 5.0\text{ N}, \theta = 0.0^\circ\)),
logging \(72{,}000\) new steps without a single dropped packet.

\emph{The systems dilemma.} The ingestion pipeline reports a \(36.0\%\)
sample volume increase with zero schema errors. However, all
\(72{,}000\) new transitions fall entirely inside the 3 nominal cells
already populated by the baseline archive
(\(\Delta N_{\text{cells}} = 0\)). Empirical loss minimization on this
corpus overfits the nominal attractor while leaving \(91.0\%\) of the
required operational envelope completely unsupported.

\emph{The systems verdict: Coverage Delta Law.} High logging frequencies
generate correlated temporal samples rather than new physical
experience. Raw sample count must never qualify a dataset; engineering
acceptance requires evaluating coverage deltas
(\(\Delta N_{\text{cells}}\)) against the operational scenario ledger.

\subsection*{Stress-Test 2: Embodiment Drift after Hardware Maintenance
(Decoupled Optical vs Contact
Mechanics)}\label{stress-test-2-embodiment-drift-after-hardware-maintenance-decoupled-optical-vs-contact-mechanics}
\addcontentsline{toc}{subsection}{Stress-Test 2: Embodiment Drift after
Hardware Maintenance (Decoupled Optical vs Contact Mechanics)}

\emph{The scenario.} An automated assembly workcell uses a parallel-jaw
gripper with compliant fingertip pads and an eye-in-hand camera to
perform peg-in-hole insertions (\(2000\) historical episodes,
\(500{,}000\) transitions). During maintenance, a technician replaces
worn \(30\text{ A}\) silicone pads with stiffer \(60\text{ A}\)
polyurethane pads (shifting friction from \(\mu = 0.85\) to
\(\mu = 0.35\) and tripling contact stiffness), while tightening the
camera bracket (shifting extrinsic pose by \(\Delta x = +4.2\text{ mm}\)
and \(\Delta \theta = +1.8^\circ\)).

\emph{The systems dilemma.} The legacy archive passes array schema
checks, but the physical plant has mutated across both optical and
contact mechanics. Camera extrinsic offsets can be resolved by
relabeling and reprojecting visual coordinate frames through the known
kinematic delta. However, altered fingertip compliance and friction
fundamentally change the physical transfer function between commanded
force and resulting slip, rendering legacy force-torque and tactile
labels physically incorrect.

\emph{The systems verdict: Embodiment Invariant.} A dataset is bound to
the physical mechanics of the body that generated it. Optical
coordinates can be reprojected across geometric calibrations, but
historical contact demonstration labels cannot be reused when physical
contact compliance changes.

\subsection*{Stress-Test 3: Survivorship Bias from Safety Filter
Censoring (The Missing Boundary
Dynamics)}\label{stress-test-3-survivorship-bias-from-safety-filter-censoring-the-missing-boundary-dynamics}
\addcontentsline{toc}{subsection}{Stress-Test 3: Survivorship Bias from
Safety Filter Censoring (The Missing Boundary Dynamics)}

\emph{The scenario.} A robotic arm executes a tactile insertion task
requiring seating an aluminum connector into a socket with
\(100\text{ }\mu\text{m}\) clearance under a real-time
\(1000\text{ Hz}\) safety force barrier (\(F_{\max} = 15.0\text{ N}\)).
Across 1000 attempted collection episodes, contact forces reach
\(15.0\text{ N}\) in 280 trials, triggering the enforcer to clamp
velocity and abort. An automated curation script discards all 280
aborted attempts, retaining only the 720 pristine trials where forces
remained below \(8.0\text{ N}\).

\emph{The systems dilemma.} The retained 720-episode corpus exhibits
zero force violations and near-zero training loss. However, when
deployed, minor vibrations push contact force to \(9.5\text{ N}\).
Because the model has never observed a state transition above
\(8.0\text{ N}\), it outputs unconditioned forward commands that
escalate contact force until the hardware emergency stop trips.
Discarding safety aborts eliminates the boundary transitions and
recovery dynamics required for closed-loop stabilization.

\emph{The systems verdict: Boundary Retention Law.} Deleting safety
interventions creates severe survivorship bias. A policy trained only on
unperturbed successes learns that high contact forces never occur,
depriving the system of the negative feedback and recovery behaviors
required to survive real-world disturbances.

\subsection*{Stress-Test 4: Aggregation Holes across Asymmetric
Interfaces (Master Device Kinematic
Blindspots)}\label{stress-test-4-aggregation-holes-across-asymmetric-interfaces-master-device-kinematic-blindspots}
\addcontentsline{toc}{subsection}{Stress-Test 4: Aggregation Holes
across Asymmetric Interfaces (Master Device Kinematic Blindspots)}

\emph{The scenario.} A teleoperated manipulation rig collects
demonstrations for a snap-fit assembly task requiring roll, pitch, and
yaw alignment. Operator A collects 2000 episodes using a 3-DOF
translation master with a 2-position wrist toggle
(\(\theta_{\text{roll}} \in \{0.0^\circ, 90.0^\circ\}\)). Operator B
collects 2000 episodes using a 6-DOF master arm whose mechanical gimbal
stop restricts roll to \([-15.0^\circ, +15.0^\circ]\). The team merges
both shifts into a single dataset of 4000 episodes (\(1{,}200{,}000\)
transitions).

\emph{The systems dilemma.} Mapping empirical state occupancy onto the
required orientation ledger reveals that roll angles between
\(20.0^\circ\) and \(70.0^\circ\) contain exactly zero samples. When
encountering intermediate roll perturbations, behavioral cloning
policies emit erratic out-of-distribution commands.

\emph{The systems verdict: Interface Aggregation Law.} Aggregating data
across kinematically constrained master devices increases statistical
volume while preserving hardware blind spots. An unvisited degree of
freedom cannot be filled by adding more demonstrations from an interface
that cannot express it.

\section{Summary}\label{sec-data-5-10}

Physical data collection converts mechanical work, operator hours, and
hardware fatigue into numerical records. A dataset of
\(1{,}000{,}000\text{ state-action pairs}\) logged at \(50\text{ Hz}\)
represents exactly \(20{,}000\text{ s}\), or \(5.56\text{ hours}\), of
active motion. Translating scheduled time into that quantity requires
budgeting across three operational stages and one hardware consumable:

\begin{itemize}
\tightlist
\item
  \textbf{Scheduled facility time} spans the full operational shift,
  including rig initialization, sensor calibration, and network staging.
\item
  \textbf{Usable trajectory time} is the narrow window during which the
  machine actively executes task transitions under valid logging
  conditions.
\item
  \textbf{Reset burden} represents the time spent repositioning
  manipulators, fixtures, and workpieces between episodes, which
  frequently demands three to four times the duration of the
  demonstration run itself.
\item
  \textbf{Consumed service life} tracks the irreversible mechanical
  fatigue inflicted on the physical system. Every trajectory consumes a
  finite portion of the rated operating hours of harmonic gearheads,
  cycles the thermal limits of motor windings, wears down mechanical
  cable linkages, and erodes compliant fingertip elastomers.
\end{itemize}

Coverage within a physical dataset is the empirical visitation record of
a specific collector operating under specific physical, kinematic, and
cognitive constraints. It is not an inherent attribute produced by large
sample volumes. Amassing millions of transitions from a teleoperator
manipulating a 3-DOF translation master cannot yield data on rotational
dynamics or pitch recovery. When the collector is a human operator, the
recorded trajectory distribution reflects biological reaction latencies
(\(P_{99} \approx 250\text{ ms}\)), hand-eye coordination limits, and
visual blind spots imposed by camera placement. When the collector is an
algorithmic baseline or an existing policy, the distribution reflects
the exact tracking gains, conservative safety margins, and kinematic
limits of that controller. Adding volume within familiar operational
envelopes merely increases statistical density over previously visited
paths. Beyond the physical and cognitive boundaries of the collection
setup, the dataset contains zero information.

Because every physical sample reflects the exact conditions under which
it was recorded, provenance cannot be treated as optional packaging. It
is a constitutive component of the dataset. A complete dataset ledger
preserves four distinct provenance dimensions alongside the raw
observation and action arrays:

\begin{enumerate}
\def\labelenumi{\arabic{enumi}.}
\tightlist
\item
  \textbf{Embodiment provenance} logs the physical instance, link
  lengths, joint backlash, sensor calibration parameters, and firmware
  version of the recording testbed.
\item
  \textbf{Interface provenance} captures the master controller type,
  mapping function, transport delay, and input filtering applied to
  operator commands.
\item
  \textbf{Safety enforcement provenance} records every instance where an
  underlying reflex, velocity governor, or Cartesian barrier clipped,
  filtered, or overrode commanded actions before current reached the
  motor windings across the MPU/MCU boundary.
\item
  \textbf{Exclusion provenance} records every trial terminated by safety
  trips, operator errors, or task aborts. Discarding failed runs without
  preserving their initial states, disturbance profiles, and divergence
  points artificially sanitizes the data, hiding the boundaries where
  the physical system failed.
\end{enumerate}

To prevent downstream training algorithms from interpolating across
unvisited regions of the state space, the dataset must carry an explicit
support boundary. This boundary is a formal geometric or statistical
ledger of the subspace where empirical observations actually exist. In
Cartesian coordinates, it bounds the envelope of reached positions and
approach vectors. In velocity and acceleration space, it establishes the
maximum dynamic rates verified on hardware. In wrench space, it records
the envelope of contact forces sustained by the end effector. When a
learned model is queried outside this support boundary during
deployment, the ensuing error is not a minor statistical variance to be
resolved by function approximation, but an ungrounded state transition
where the physical system has never operated. Propagating the support
boundary directly into the downstream pipeline provides the nervous
system with the envelope required to trigger control refusal or safety
interventions.

\protect\phantomsection\label{callout-takeawaysux2a-5.10}
\begin{fbxSimple}{callout-takeaways}{Key Chapter Takeaways}{}
\phantomsection\label{callout-takeaways*-5.10}

\begin{itemize}
\tightlist
\item
  \textbf{Coverage is an Empirical Visitation Record, Not a Sample
  Count:} Repetitive collection inside nominal attractors inflates
  tensor row counts without expanding state-space coverage
  (\(\Delta N_{\text{cells}} = 0\)). Training corpora must be certified
  against multi-dimensional scenario ledgers spanning the entire
  operational design domain.
\item
  \textbf{Physical Datasets Are Bound to Embodiment Mechanics:} Optical
  calibrations can be reprojected mathematically across geometric
  transforms, but contact mechanics (elastomer durometer, contact
  compliance, joint backlash) permanently bake physical transfer
  functions into force-torque labels that cannot transfer across mutated
  hardware.
\item
  \textbf{Safety Interventions Must Be Preserved (Boundary Retention
  Law):} Censoring aborted collection runs or governor-clamped
  trajectories deprives the policy of boundary recovery dynamics.
  Datasets must record the complete 4-stream action tuple
  (\(\mathbf{a}_{\text{req}}, \mathbf{a}_{\text{cmd}}, \mathbf{a}_{\text{enf}}, \mathbf{a}_{\text{meas}}\))
  to expose governor interventions.
\item
  \textbf{Master Interface Blind Spots Create Irrecoverable Voids:}
  Aggregating demonstrations across kinematically restricted master
  devices (e.g., 3-DOF translation pucks vs.~6-DOF arms) leaves
  structural voids across unvisited rotational degrees of freedom that
  no downstream scaling can recover.
\item
  \textbf{Training Cannot Synthesize Missing Physical Counterfactuals:}
  No loss function, diffusion schedule, or geometric data augmentation
  can generate physics that the collection regime omitted. Explicit
  support boundaries must accompany datasets to trigger selective
  runtime abstention outside the empirical convex hull.
\end{itemize}

\end{fbxSimple}

Chapter~\ref{sec-training} receives both the logged dataset and this
complete mechanical record, shifting the systems problem from
acquisition to training. The governing constraint for every optimization
method covered next is that no training choice can recover missing
physics. No neural network architecture, loss formulation, diffusion
schedule, or data augmentation heuristic can synthesize state
transitions, reactive counterfactuals, or control authorities that the
collection regime omitted. If the collection process never visited a
high-slip friction regime, never demonstrated recovery from a
\(40\text{ mm}\) position error, or operated entirely behind a force
governor that capped contact loads to \(10.0\text{ N}\), the resulting
policy cannot learn to manage those physical interactions.
Chapter~\ref{sec-training} examines the computational, architectural,
and safety costs incurred when converting this finite, constrained
record into an executable policy.

\chapter{\texorpdfstring{Policy Synthesis \&
Training\chaptersubtitle{What Each Way of Making a Policy Charges You}}{Policy Synthesis \& Training}}\label{sec-training}

\begin{figure}[H]

\centering{

\includegraphics[width=1\linewidth,height=\textheight,keepaspectratio]{index_files/mediabag//Users/VJ/GitHub/MLSysBook-integrate-physical/physical/book/chapters/06-training/figures/fig_locator.pdf}

}

\caption{\label{fig-locator}\textbf{Systems Organ Focus: Brain.}
Architectural trade-off landscape across policy synthesis regimes,
comparing the operational capabilities purchased during training against
the latent failure modes and physical costs charged during deployment.}

\end{figure}%

\emph{Why does a policy that achieves 100\% task reward in simulation
destroy real gearboxes within minutes of deployment?}

The central dilemma of physical policy synthesis is that learning
requires exploratory interaction, but physical hardware punishes
unconstrained exploration with mechanical fracture, thermal overheating,
and catastrophic yield. While physics simulators offer fast, scalable
training environments, the sim-to-real reality gap is not a single
scalar error that can be optimized away; it is a coupled vector of
unmodeled joint friction, sensor latency jitter, contact compliance
mismatches, and visual discrepancies. A policy trained in simulation
frequently discovers unphysical shortcuts that violently shake physical
hardware. Training dependable physical AI policies demands deliberate
architectural co-design between conservative behavioral cloning,
domain-randomized simulation, and expressive action decoders that
explicitly account for the physical cost of exploration.

\newpage{}

\protect\phantomsection\label{callout-objectiveux2a-6.1}
\begin{fbxSimple}{callout-objective}{Learning Objectives}{}
\phantomsection\label{callout-objective*-6.1}

After studying this chapter, you will be able to:

\begin{enumerate}
\def\labelenumi{\arabic{enumi}.}
\tightlist
\item
  State what each policy synthesis regime buys in capability and what
  failure mode it charges at deployment.
\item
  Explain why compounding error in behavioral cloning is structural
  (\(\mathcal{O}(T^2 \epsilon)\)) rather than a hyperparameter tuning
  problem.
\item
  Compare the three primary action decoder architectures---Discrete
  Token Binning, Continuous Trajectory Diffusion, and Latent Variable
  Action Chunking---and explain how each resolves multimodal
  demonstration averaging.
\item
  Treat a physics simulator as a versioned system component, specifying
  its contract, numerical relaxations, and physical omissions.
\item
  Decompose the sim-to-real reality gap vector across dynamics, sensing,
  appearance, and timing channels using paired-trace metrics.
\item
  Predict a physical deployment failure mode from a training choice
  before the machine is assembled.
\item
  Construct an immutable Policy Lineage Manifest documenting training
  data support, simulator provenance, and hardware wear logs.
\end{enumerate}

\emph{Core chapter takeaway.} Every way of making a policy buys a
capability and charges a deployment failure mode. Cloning fails silently
off-distribution. Trial-based learning imports a simulator, and
everything wrong with the simulator becomes wrong with the policy.

\end{fbxSimple}

\section{Policy Synthesis}\label{sec-training-6-1}

\lettrine{C}{onsider}\lettrine{t}{wo} learned controllers designed to
regulate liquid flow through a proportional valve before a mechanical
placement head seats a component onto a seal. Both models achieve an
identical validation score of \(0.002\) mean squared error on their
respective offline benchmarks. The first policy was trained by
behavioral cloning on recorded operator demonstrations, while the second
was trained through trial and error inside a physics simulator. When
deployed to identical physical test benches, both policies regulate
steady flow at \(1.50\text{ L/min}\) under nominal supply pressure. When
a supply pump transient drops line pressure by \(12\text{ kPa}\), their
failure modes diverge immediately. The demonstrated policy issues a
correction that undershoots slightly, landing in a pressure-flow state
that never appeared in the human demonstrations. Because the operator
never visited this offset, the model generates an erratic valve command,
compounds its tracking error over consecutive \(20\text{ ms}\) control
steps, and overflows the reservoir. The simulated policy responds to the
same pressure drop by commanding high-frequency valve reversals at
\(50\text{ Hz}\), reflecting the simulator's assumption that the valve
actuator accelerates instantaneously. On physical hardware, where the
solenoid spool requires \(35\text{ ms}\) to transit between positions,
the valve cannot track the commanded reversals, causing pressure
oscillations that blow out an upstream pipe seal. Both policies
succeeded on their training metrics, but each inherited a distinct
physical destruction mode determined entirely by how it was trained.

From a systems engineering perspective, a \textbf{policy synthesis
regime} is defined as the specific boundary conditions of its data
pipeline, optimization objective, and execution environment. These
boundary conditions comprise three specific elements: the manifold of
physical states the policy is permitted to explore during optimization,
the sensory observations and disturbance profiles it records across
those states, and the mechanism that provides the target signal. Whether
that target signal is an operator's joint torque log, a scalar reward
evaluating liquid throughput, an analytical cost function, or a
self-supervised reconstruction loss on next-frame camera telemetry, the
learned model remains completely ignorant of physical mechanics outside
the envelope of that evidence.

Every training regime represents an economic trade between capability
and vulnerability. What a training regime buys is a deployment
capability, such as the ability to coordinate high-dimensional sensory
inputs without an explicit state estimator, or the ability to discover
non-obvious control strategies across millions of simulated iterations.
What it charges is a runtime dependency on every simplifying assumption
embedded in the data-generating process, reward formulation, numerical
solver, or adaptation environment. Success on a training objective
measures only how thoroughly an optimization algorithm minimized error
across its designated evidence set. It provides zero guarantees
regarding how the policy behaves under the state distribution it
generates once connected to physical actuators across all three physical
machine archetypes:

\begin{itemize}
\tightlist
\item
  In \textbf{Mobile Systems} (Class 1), an autonomous ground vehicle or
  quadruped trained purely on centered highway logs or flat-ground
  simulation fails when encountering shoulder gravel, wheel slip on ice,
  or camera sensor blur during high-speed yaw.
\item
  In \textbf{Manipulators} (Class 2), an articulated assembly arm
  trained via supervised cloning commands non-functional intermediate
  grasp poses when presented with off-nominal workpiece orientations,
  dropping parts or jamming tooling into rigid fixtures.
\item
  In \textbf{Process and Energy Systems} (Class 3), a battery thermal
  management controller trained on nominal heat cycles commands
  high-frequency pump oscillations during sudden thermal spikes,
  cavitating the coolant loop and tripping overtemperature protection.
\end{itemize}

This discrepancy arises because physical execution is closed-loop.
During training, states are either passively retrieved from static disks
or computed by numerical integrators that enforce convenient
simplifications. Once deployed, the policy's own control outputs
determine the subsequent physical state of the machine. If a policy
trained on demonstrations makes a single actuation error that shifts the
machine into an unrecorded state, the training distribution no longer
governs the inputs. If a policy trained in simulation relies on
unmodeled physical phenomena, such as thermal resistance changes in a
motor winding or friction hysteresis in a gearbox, the physical plant
diverges from the simulated plant at every integration step. The debt
incurred during training is deferred until deployment, where physics
collects the full balance at line rate.

The primary training regimes used in physical systems are best
understood as distinct engineering trades rather than competing
algorithmic paradigms. As mapped in Figure~\ref{fig-locator}, within the
Brain organ, policy synthesis is structured around four primary
regimes---learning from demonstration, trial-based reinforcement
learning, simulation-based exploration, and physical fine-tuning.
Learning from human demonstration buys immediate task competency without
requiring mathematical reward engineering, but charges an acute
vulnerability to error compounding when disturbances push the system off
the demonstrated trajectory. Simulation-based reinforcement learning
buys exhaustive exploration across billions of synthetic state
transitions without endangering physical hardware, but charges a
structural dependency on simulator fidelity that breaks when actuator
dynamics or contact mechanics diverge from the model. Real-machine
adaptation buys authentic contact dynamics and unmodeled physical
interactions, but charges physical wear, long collection times bounded
by wall-clock physics, and the operational risk of running unverified
policies on live equipment.

Designing a physical AI system requires treating the training process as
a component specification rather than an isolated machine learning task.
An engineering team deploying a learned policy must produce three
deliverables alongside the model weights, specifically a failure mode
prediction derived from the regime's structural assumptions, a runtime
detector configured to catch that failure before mechanical limits are
exceeded, and a complete provenance log detailing the data sources,
simulation bounds, and target criteria that shaped the policy. Because
learning from demonstration represents the most direct and widely
adopted method for transferring human skill to an automated actuator,
understanding its exact mathematical mechanics and error-propagation
characteristics establishes the baseline against which all other
training trades are measured.

\section{Policy Synthesis Regimes and Demonstration
Learning}\label{sec-training-6-2}

\begin{longtable}[]{@{}
  >{\raggedright\arraybackslash}p{(\linewidth - 10\tabcolsep) * \real{0.1667}}
  >{\raggedright\arraybackslash}p{(\linewidth - 10\tabcolsep) * \real{0.1667}}
  >{\raggedright\arraybackslash}p{(\linewidth - 10\tabcolsep) * \real{0.1667}}
  >{\raggedright\arraybackslash}p{(\linewidth - 10\tabcolsep) * \real{0.1667}}
  >{\raggedright\arraybackslash}p{(\linewidth - 10\tabcolsep) * \real{0.1667}}
  >{\raggedright\arraybackslash}p{(\linewidth - 10\tabcolsep) * \real{0.1667}}@{}}
\caption{\textbf{Comparative Trade-Off Matrix Across Policy Synthesis
Regimes.} Systematic evaluation of behavioral cloning, offline
reinforcement learning, sim-to-real reinforcement learning, and residual
RL/fine-tuning. Columns detail the primary operational capability
purchased during training, representative sample and compute budgets,
sample efficiency bounds, safety risks incurred during data acquisition,
and dominant failure signatures with corresponding nervous system
interceptors.}\label{tbl-06-regimes}\tabularnewline
\toprule\noalign{}
\begin{minipage}[b]{\linewidth}\raggedright
Policy Synthesis Regime
\end{minipage} & \begin{minipage}[b]{\linewidth}\raggedright
Core Purchased Capability
\end{minipage} & \begin{minipage}[b]{\linewidth}\raggedright
Data \& Compute Budget
\end{minipage} & \begin{minipage}[b]{\linewidth}\raggedright
Sample Efficiency \& Scale
\end{minipage} & \begin{minipage}[b]{\linewidth}\raggedright
Hardware Safety \& Wear Risk
\end{minipage} & \begin{minipage}[b]{\linewidth}\raggedright
Dominant Failure Mode \& Interceptor
\end{minipage} \\
\midrule\noalign{}
\endfirsthead
\toprule\noalign{}
\begin{minipage}[b]{\linewidth}\raggedright
Policy Synthesis Regime
\end{minipage} & \begin{minipage}[b]{\linewidth}\raggedright
Core Purchased Capability
\end{minipage} & \begin{minipage}[b]{\linewidth}\raggedright
Data \& Compute Budget
\end{minipage} & \begin{minipage}[b]{\linewidth}\raggedright
Sample Efficiency \& Scale
\end{minipage} & \begin{minipage}[b]{\linewidth}\raggedright
Hardware Safety \& Wear Risk
\end{minipage} & \begin{minipage}[b]{\linewidth}\raggedright
Dominant Failure Mode \& Interceptor
\end{minipage} \\
\midrule\noalign{}
\endhead
\bottomrule\noalign{}
\endlastfoot
\textbf{Behavioral Cloning (BC)} & Direct imitation of demonstrated
trajectories without reward engineering or system identification &
\(10^2\)--\(10^3\) teleop trajectories (\(\sim 10\)--\(50\text{ h}\)
data); \(<50\text{ GPU-h}\) compute & Supervised \(\mathcal{O}(N)\) over
expert demonstrations; fixed \(1\times\) real-time data throughput &
Minimal exploration risk during logging; zero hardware exploration
during optimization & Quadratic compounding error
(\(\mathcal{O}(T^2 \epsilon)\)) via covariate shift. \emph{Interceptor}:
Kernel/Mahalanobis density monitors, convex hull boundary clamps \\
\textbf{Offline / Batch RL} & Value-guided policy optimization over
fixed, heterogeneous sub-optimal logs; stitches multi-trajectory
transitions & \(10^5\)--\(10^7\) logged transitions with state-action
coverage; \(100\)--\(500\text{ GPU-h}\) compute & Offline
\(\mathcal{O}(N)\) over static datasets; sample-efficient but demands
high state-action overlap & Zero active exploration hazards during
offline training; logging requires caged rigs or safe scripted
controllers & Out-of-distribution action extrapolation and Q-value
overestimation. \emph{Interceptor}: Conservative Q-learning (CQL)
penalty bounds, action support filters \\
\textbf{Sim-to-Real RL} & Discovery of dynamic, non-prehensile
manipulation and agile locomotion via billions of unconstrained rollouts
& Zero physical data; \(10^7\)--\(10^9\) simulated transitions across
\(10^3\)--\(10^4\) GPU environments; \(500\)--\(5000\text{ GPU-h}\) &
Low algorithmic sample efficiency (\(10^7\)--\(10^9\) steps), offset by
\(10^3\times\)--\(10^5\times\) simulation acceleration & Zero hardware
wear or collision damage during training; severe deployment risk if
plant dynamics diverge & Exploitation of numerical solver artifacts
(contact interpenetration, zero latency). \emph{Interceptor}:
Paired-trace discrepancy budgets, force-derivative (\(dF/dt\)) clamps \\
\textbf{Residual RL \& Real Fine-Tuning} & Direct policy correction
against physical contact compliance, gearbox hysteresis, and true sensor
noise & \(10^3\)--\(10^5\) physical steps (\(0.5\)--\(10\text{ h}\)
wall-clock time); \(<20\text{ GPU-h}\) compute & High sample efficiency
required (SAC/TD3 or residual delta-action formulation); strictly
\(1\times\) real-time & High physical exposure; repetitive impact shocks
degrade gearboxes; exploratory current dither elevates winding
temperature & Catastrophic forgetting of nominal transit stability,
local policy collapse. \emph{Interceptor}: Parameter drift bounds
(\(\Vert\Delta \theta\Vert_2 \le \delta\)), frozen regression suites,
thermal limiters \\
\end{longtable}

As compared across synthesis paradigms (Table~\ref{tbl-06-regimes}),
\textbf{behavioral cloning} is the cheapest route to a working policy
and the one whose failure mode is best understood. It frames policy
synthesis as supervised regression or classification over a collected
dataset of expert state-action pairs \(\mathcal{D} = \{(o_t, a_t)\}\).
In its standard formulation, the optimization objective minimizes a
one-step prediction loss, fitting an action conditional on the current
observation \(o_t\) or a finite observation history
\((o_{t-H}, \dots, o_t)\) drawn from demonstration
trajectories:{[}\^{}fn-inc-alvinn-cloning{]}
\[\mathcal{L}_{\text{MSE}}(\theta) = \mathbb{E}_{(o, a) \sim \mathcal{D}} \left[ \| \pi_\theta(o) - a \|_2^2 \right]\]
Because this loss penalizes only the discrepancy between the commanded
action and the recorded demonstration at observed states, the training
objective does not evaluate task completion, terminal cost, or the
stability of the closed-loop dynamical system. It fits actions to
observations, and it fits nothing else.

\subsubsection*{The Multimodality Breakdown and Mode
Averaging}\label{the-multimodality-breakdown-and-mode-averaging}
\addcontentsline{toc}{subsubsection}{The Multimodality Breakdown and
Mode Averaging}

The primary mathematical vulnerability of standard behavioral cloning
under Mean Squared Error (MSE) loss is its inability to represent
multimodal action distributions. When an optimization objective
minimizes expected squared \(L_2\) error, the theoretical minimizer for
a deterministic neural network \(\pi_\theta(o)\) is the conditional
expectation of the action distribution:\sidenote{\footnotesize Under an \(L_2\)
  regression loss, the expected risk
  \(\mathcal{R}(\hat{a}) = \int \|\hat{a} - a\|_2^2 p(a \mid o) da\) is
  strictly convex with respect to \(\hat{a}\). Differentiating with
  respect to \(\hat{a}\) and setting the gradient to zero yields
  \(2 \int (\hat{a} - a) p(a \mid o) da = 0 \implies \hat{a}^* = \int a p(a \mid o) da = \mathbb{E}[a \mid o]\).
  By Jensen's inequality, for any non-linear downstream forward physical
  plant \(f_{\text{plant}}(a)\),
  \(\mathbb{E}[f_{\text{plant}}(a)] \ne f_{\text{plant}}(\mathbb{E}[a])\),
  meaning that commanding the average action leads to physical states
  that were never intended or observed during data logging.}
\[\pi^*(o) = \arg\min_{\hat{a}} \int \| \hat{a} - a \|_2^2 \, p(a \mid o) \, da = \mathbb{E}_{a \sim p(a \mid o)} [a] = \int a \, p(a \mid o) \, da\]
When human demonstrators perform a physical task, their strategy is
frequently multimodal. Consider an obstacle avoidance scenario where a
mobile base or robotic arm must navigate around a central pillar. A
human demonstrator may choose to pass around the pillar to the left
(\(a_1 = -v_{\text{lat}}\)) in \(50\%\) of the demonstrations, or to the
right (\(a_2 = +v_{\text{lat}}\)) in the remaining \(50\%\). The true
demonstration distribution \(p(a \mid o)\) is bimodal with distinct
probability peaks at \(-v_{\text{lat}}\) and \(+v_{\text{lat}}\), and
zero probability density in the center (\(p(0 \mid o) \approx 0\)).

When trained on this dataset with an MSE regression loss, the
deterministic network computes the conditional expectation:
\[\mathbb{E}[a \mid o] = 0.5(-v_{\text{lat}}) + 0.5(+v_{\text{lat}}) = 0\]
The policy outputs zero lateral velocity, commanding the machine to
drive directly into the obstacle. The conditional mean represents an
unviable, catastrophic action that never occurred in the demonstration
dataset. Across physical machine archetypes:

\begin{itemize}
\tightlist
\item
  In \emph{Mobile Systems}, the vehicle drives into pillars, bridge
  abutments, or dividing medians.
\item
  In \emph{Manipulators}, mode averaging causes end-effector grippers to
  command non-functional intermediate grasp poses that straddle an
  object's outer boundary, dropping workpieces or jamming tooling into
  rigid assembly fixtures.
\item
  In \emph{Process Systems}, averaging between two discrete valve
  operating regimes (\(10\%\) bypass vs.~\(90\%\) main feed) commands an
  uncalibrated \(50\%\) intermediate flow that cavitates the supply pump
  and violates cooling deadlines.
\end{itemize}

Furthermore, in high-frequency single-step control loops
(\(50\text{--}100\text{ Hz}\)), Markovian policies mapping single
observations \(o_t \to a_t\) suffer from severe temporal inconsistency.
Minor sensor noise or slight state perturbations cause the predicted
output to alternate between distinct modes across consecutive timesteps.
The policy commands rapid alternating positive and negative actuator
torques
(\(+10\text{ N}\cdot\text{m} \leftrightarrow -10\text{ N}\cdot\text{m}\)),
exciting mechanical drive resonances, accelerating gearbox tooth wear,
and driving high Joule heating in motor windings.

\subsubsection*{Overcoming Multimodality: The Three Action Decoder
Paradigms}\label{overcoming-multimodality-the-three-action-decoder-paradigms}
\addcontentsline{toc}{subsubsection}{Overcoming Multimodality: The Three
Action Decoder Paradigms}

Modern imitation learning resolves the multimodality breakdown by
replacing deterministic single-step regression with expressive action
decoder architectures. Three dominant paradigms define the state of the
art:

\begin{enumerate}
\def\labelenumi{\arabic{enumi}.}
\item
  \textbf{Discrete Token Binning (RT-1, RT-2, OpenVLA)}: Discrete token
  binning discretizes each dimension of a continuous action vector
  \(a_t \in [-1, 1]^{D_a}\) into \(B\) uniform categorical bins (e.g.,
  \(B = 256\) tokens)
  \cite{brohan2023rt1, brohan2023rt2, kim2024openvla}. The policy treats
  action prediction as autoregressive sequence generation over discrete
  action tokens \((a_{t,1}, \dots, a_{t,D_a})\), optimizing standard
  cross-entropy loss:
  \[\mathcal{L}_{\text{CE}}(\theta) = -\sum_{d=1}^{D_a} \log p_\theta(a_{t,d} \mid o_t, a_{t,<d})\]
  \emph{What it buys:} Natively represents arbitrary multimodal
  probability distributions over discrete bins without mode collapse,
  and integrates directly into pretrained Vision-Language-Action (VLA)
  foundation models. \emph{What it charges:} Discretization quantization
  error bounded by
  \(\Delta a = \frac{a_{\text{max}} - a_{\text{min}}}{B}\) (e.g.,
  \(\Delta a = 2/256 \approx 0.0078\) normalized units), generating
  stair-stepped commands; high inference latency due to autoregressive
  sequential token generation (\(D_a\) forward passes per control step);
  and lack of native inter-step kinematic smoothness, requiring
  downstream low-pass filtering on the Real-Time MCU.
\item
  \textbf{Continuous Trajectory Diffusion (Diffusion Policy)}: Diffusion
  Policy treats policy synthesis as conditional denoising diffusion over
  continuous multi-step action trajectory chunks
  \(A_t = (a_t, a_{t+1}, \dots, a_{t+K-1}) \in \mathbb{R}^{K \times D_a}\)
  conditioned on an observation history
  \(O_t = (o_{t-H+1}, \dots, o_t)\) \cite{chi2023diffusionpolicy}. An
  apparent paradox often confuses practitioners: \emph{if Mean Squared
  Error causes mode averaging in direct policy regression, why does
  Diffusion Policy optimize an MSE loss?} The crucial distinction lies
  in the target of the loss function. In standard BC, MSE is applied
  directly across the multimodal action space \(p(a \mid o)\), forcing
  the network to predict the geometric centroid. In Diffusion Policy,
  the neural network \(\epsilon_\theta(A^k, O_t, k)\) is trained via MSE
  to predict the standard Gaussian noise vector
  \(\epsilon \sim \mathcal{N}(0, \mathbf{I})\) added at diffusion step
  \(k \in \{1, \dots, N_{\text{diff}}\}\):{[}\^{}fn-hw-diffusion-latency{]}
  \[\mathcal{L}_{\text{diff}}(\theta) = \mathbb{E}_{k, A^0, \epsilon} \left[ \| \epsilon - \epsilon_\theta(\sqrt{\bar{\alpha}_k} A^0 + \sqrt{1 - \bar{\alpha}_k}\epsilon, O_t, k) \|_2^2 \right]\]
  Because the noise prediction target \(\epsilon\) is strictly unimodal
  Gaussian, MSE is the exact maximum-likelihood objective. At inference
  time, sampling begins by drawing pure Gaussian noise
  \(A^{N_{\text{diff}}} \sim \mathcal{N}(0, \mathbf{I})\) and
  iteratively denoises it through the learned score field
  \(\nabla_A \log p(A \mid O_t)\) over \(N_{\text{diff}}\) reverse
  steps. The non-linear stochastic integration breaks symmetry, causing
  the trajectory to settle into one distinct demonstrated mode without
  mode collapse or centroid averaging. Furthermore, trajectory-level
  diffusion inherently enforces kinematic smoothness across physical
  inertia, amortizing the effect of high-frequency observation noise.
\item
  \textbf{Latent Variable Action Chunking (ACT)}: ACT reformulates
  policy synthesis by predicting an entire sequence of future actions
  \(A_t = (a_t, a_{t+1}, \dots, a_{t+K-1}) \in \mathbb{R}^{K \times D_a}\)
  across a receding horizon \(K\) given the current observation \(o_t\)
  \cite{zhao2023learning}. To capture multimodality, ACT employs a
  Conditional Variational Autoencoder (CVAE) parameterized by a
  Transformer encoder-decoder. During training, a Transformer encoder
  compresses demonstrated action chunks and current observations into a
  low-dimensional latent style variable
  \(z \sim q_\phi(z \mid o_t, A_t)\) regularized by KL divergence
  against a unit Gaussian prior \(\mathcal{N}(0, \mathbf{I})\). At
  inference time, the latent code is sampled from
  \(z \sim \mathcal{N}(0, \mathbf{I})\) and fed into the Transformer
  decoder \(\pi_\theta(o_t, z)\). Sampling a specific latent code \(z\)
  selects a single coherent demonstrated intention (e.g., ``navigate
  left''), generating an internally consistent, continuous multi-step
  trajectory chunk without averaging across distinct modes. To prevent
  discontinuous jerks across chunk boundaries, ACT applies
  \emph{temporal ensembling} across overlapping receding horizons. At
  global execution time \(t\), multiple prior chunks \(A_{t-i}\) (for
  \(i \in \{0, \dots, \min(t, K-1)\}\)) provide overlapping predictions
  for step \(t\). ACT computes the executed action as an exponentially
  decaying weighted blend:
  \[a_t = \frac{\sum_{i=0}^{\min(t, K-1)} w_i A_{t-i}[i]}{\sum_{i=0}^{\min(t, K-1)} w_i}, \quad w_i = \exp(-m \cdot i)\]
  where \(i\) represents the lookback age of the prediction chunk and
  \(m \ge 0\) parameterizes the decay weight, giving highest priority to
  fresher observations while smoothing out single-step actuator jitter.
\end{enumerate}

\subsubsection*{Algorithmic Formulation: Trajectory Diffusion and
Asynchronous
Decoupling}\label{algorithmic-formulation-trajectory-diffusion-and-asynchronous-decoupling}
\addcontentsline{toc}{subsubsection}{Algorithmic Formulation: Trajectory
Diffusion and Asynchronous Decoupling}

Deploying generative imitation architectures introduces significant
computational latency (\(20\text{--}80\text{ ms}\) for multi-step
diffusion denoising or multi-head Transformer self-attention on embedded
GPUs/NPUs). In a canonical heterogeneous dual-brain edge System-on-Chip
(SoC) architecture, this latency requires strict physical decoupling:
the unprivileged Linux MPU runs the generative policy asynchronously at
a lower rate (\(10\text{--}20\text{ Hz}\)), populating an action chunk
ring buffer. Meanwhile, the deterministic Real-Time MCU samples the
chunk buffer at \(500\text{--}1000\text{ Hz}\), executing smooth cubic
spline interpolation, velocity envelope governor checks, and hard
torque-slew limiters (\(d\tau/dt \le \dot{\tau}_{\text{max}}\)). This
complete pipeline is formalized in Algorithm 6.1.

\protect\phantomsection\label{callout-algorithmux2a-6.2}
\begin{fbxSimple}{callout-algorithm}{Algorithm:}{Algorithm 6.1: Diffusion Policy Trajectory Chunk Denoising and Dual-Rate Actuation}
\phantomsection\label{callout-algorithm*-6.2}

\textbf{Inputs:} Observation history buffer
\(\mathbf{O}_t = (\mathbf{o}_{t-H+1}, \dots, \mathbf{o}_t) \in \mathbb{R}^{H \times D_{\text{obs}}}\),
noise prediction network
\(\boldsymbol{\epsilon}_\theta(\mathbf{A}^k, \mathbf{O}_t, k)\), noise
schedule
\(\{\alpha_k, \bar{\alpha}_k, \sigma_k\}_{k=1}^{N_{\text{diff}}}\),
chunk horizon \(K\) (e.g., \(16\)), action dimension \(D_a\), MCU loop
period \(\Delta t_{\text{mcu}} = 1.0\text{ ms}\).

\textbf{Outputs:} Line-rate joint torque setpoint / action vector
\(\mathbf{a}_t \in \mathbb{R}^{D_a}\).

\textbf{Structural Invariants:}

\begin{enumerate}
\def\labelenumi{\arabic{enumi}.}
\tightlist
\item
  \emph{Slew-Rate Bound:}
  \(\|\mathbf{a}_t - \mathbf{a}_{t-1}\|_2 / \Delta t_{\text{mcu}} \le \dot{\mathbf{a}}_{\max}\).
\item
  \emph{Non-Blocking Ring Buffer:} SPSC ring buffer never underflows
  during nominal steady-state operation.
\end{enumerate}

\textbf{Execution Pipeline:}

\begin{enumerate}
\def\labelenumi{\arabic{enumi}.}
\tightlist
\item
  \textbf{Linux MPU --- Asynchronous Trajectory Chunk Generation
  (\(10\text{--}20\text{ Hz}\)):}

  \begin{itemize}
  \item
    Ingest observation history \(\mathbf{O}_t\) and compute visual/state
    embedding:
    \(\mathbf{E}_t \leftarrow \text{VisualEncoder}(\mathbf{O}_t)\).
  \item
    Initialize reverse diffusion trajectory from Gaussian noise:
    \(\mathbf{A}^{N_{\text{diff}}} \sim \mathcal{N}(\mathbf{0}, \mathbf{I}_{K \times D_a})\).
  \item
    For \(k = N_{\text{diff}}\) down to \(1\):

    \begin{itemize}
    \tightlist
    \item
      Predict noise component:
      \(\hat{\boldsymbol{\epsilon}} \leftarrow \boldsymbol{\epsilon}_\theta(\mathbf{A}^k, \mathbf{E}_t, k)\).
    \item
      Compute mean update:
      \(\boldsymbol{\mu}_k \leftarrow \frac{1}{\sqrt{\alpha_k}} \left( \mathbf{A}^k - \frac{1 - \alpha_k}{\sqrt{1 - \bar{\alpha}_k}} \hat{\boldsymbol{\epsilon}} \right)\).
    \item
      Sample \(\mathbf{z} \sim \mathcal{N}(\mathbf{0}, \mathbf{I})\) for
      \(k > 1\), else \(\mathbf{z} \leftarrow \mathbf{0}\).
    \item
      Step reverse diffusion:
      \(\mathbf{A}^{k-1} \leftarrow \boldsymbol{\mu}_k + \sigma_k \mathbf{z}\).
    \end{itemize}
  \item
    Assign synthesized trajectory chunk:
    \(\mathbf{A}_{\text{chunk}} \leftarrow \mathbf{A}^0 \in \mathbb{R}^{K \times D_a}\).
  \item
    Lock-free deposit \(\mathbf{A}_{\text{chunk}}\) into Single-Producer
    Single-Consumer (SPSC) ring buffer.
  \end{itemize}
\item
  \textbf{Real-Time MCU --- Synchronous Deterministic Actuation
  (\(1000\text{ Hz}\)):}

  \begin{itemize}
  \item
    Poll SPSC ring buffer; if new chunk available, update local spline
    knots.
  \item
    Compute interpolated nominal action:
    \(\mathbf{a}_{\text{raw}} \leftarrow \text{CubicSplineInterpolate}(\mathbf{A}_{\text{chunk}}, t_{\text{current}})\).
  \item
    Enforce kinematic slew-rate governor:
    \[\Delta \mathbf{a} \leftarrow \text{clamp}\left(\mathbf{a}_{\text{raw}} - \mathbf{a}_{\text{prev}}, -\dot{\mathbf{a}}_{\max} \Delta t_{\text{mcu}}, +\dot{\mathbf{a}}_{\max} \Delta t_{\text{mcu}}\right), \quad \mathbf{a}_{\text{clamped}} \leftarrow \mathbf{a}_{\text{prev}} + \Delta \mathbf{a}\]
  \item
    Enforce nervous-system safety envelope:
    \(\mathbf{a}_t \leftarrow \text{SafetyGovernorClamp}(\mathbf{a}_{\text{clamped}})\).
  \item
    Transmit \(\mathbf{a}_t\) to motor bridge PWM / inverter registers
    via EtherCAT / SPI.
  \item
    Update state memory:
    \(\mathbf{a}_{\text{prev}} \leftarrow \mathbf{a}_t\).
  \end{itemize}
\end{enumerate}

\end{fbxSimple}

\subsubsection*{Offline Reinforcement Learning and Action Distributional
Shift}\label{offline-reinforcement-learning-and-action-distributional-shift}
\addcontentsline{toc}{subsubsection}{Offline Reinforcement Learning and
Action Distributional Shift}

When physical demonstration datasets
\(\mathcal{D} = \{(s_t, a_t, r_t, s_{t+1})\}\) contain heterogeneous or
sub-optimal trajectories (e.g., mixtures of teleoperation, scripted
routines, and exploratory human play), behavioral cloning struggles
because it treats all demonstrations with equal supervised weight.
\textbf{Offline Reinforcement Learning (Batch RL)} aims to optimize a
policy directly from logged transition archives without active hardware
exploration, using reward signals to stitch together high-performing
composite trajectories from disjoint, imperfect demonstrations
\cite{kumar2020conservative}.

However, applying standard dynamic programming or Q-learning to static
offline datasets triggers a catastrophic failure mode known as
\textbf{Out-of-Distribution (OOD) Action Overestimation}. In standard
temporal-difference learning, the Bellman optimality operator evaluates:
\[Q(s_t, a_t) \leftarrow r(s_t, a_t) + \gamma \max_{a'} Q(s_{t+1}, a')\]
To evaluate the target value, the \(\max_{a'}\) operator queries the
parameterized critic network \(Q_\phi(s', a')\) across the continuous
action space. For actions \(a'\) that were never executed in the logging
dataset \(\mathcal{D}\), the function approximator has received zero
training signal. Due to neural network generalization variance, the
unconstrained critic outputs erroneously large Q-values for some OOD
actions. The \(\max\) operator greedily selects these spurious peaks,
propagating over-optimistic value errors backwards through the Bellman
graph. When deployed on the physical machine, the actor greedily selects
these ungrounded OOD actions, commanding physically unviable motor
voltages that cause immediate joint stalls or boundary collisions.

To prevent value explosion without querying the physical environment,
\textbf{Conservative Q-Learning (CQL)} adds a regularization penalty
directly to the critic loss.{[}\^{}fn-math-cql-bound{]} CQL penalizes
high Q-values on out-of-distribution actions while maximizing the
expected Q-value on transitions within the empirical dataset:
\[\begin{aligned}
\min_Q &\;\alpha \left( \mathbb{E}_{s \sim \mathcal{D}, a \sim \mu(a \mid s)} [Q(s, a)] - \mathbb{E}_{(s, a) \sim \mathcal{D}} [Q(s, a)] \right) \\
&+ \frac{1}{2} \mathbb{E}_{(s, a, s') \sim \mathcal{D}} \left[ \left( Q(s, a) - \hat{\mathcal{B}}^{\pi} \hat{Q}(s, a) \right)^2 \right]
\end{aligned}\] where \(\mu(a \mid s)\) is a wide sampling distribution
(such as a uniform distribution over the actuator limits or the current
policy). By explicitly pushing down Q-values outside the data manifold,
CQL mathematically guarantees that the learned critic is a lower bound
on the true policy performance (\(\hat{Q}^\pi(s, a) \le Q^\pi(s, a)\)),
preventing dangerous extrapolation beyond the physical support of the
dataset.

\subsubsection*{Compounding Covariate Error and State Support
Drift}\label{compounding-covariate-error-and-state-support-drift}
\addcontentsline{toc}{subsubsection}{Compounding Covariate Error and
State Support Drift}

Even when multimodality is handled by generative architectures,
behavioral cloning remains structurally vulnerable to closed-loop error
compounding. Consider a positioning system that regulates the height of
a dispensing nozzle over a moving substrate at \(100\text{ Hz}\) before
touching the surface to deposit adhesive. In a clean demonstration
dataset, an expert operator or an engineered classical controller
maintains the nozzle within a narrow operating band of
\(\pm 0.2\text{ mm}\) around the reference trajectory. A learned policy
evaluated on a held-out test split of this dataset can achieve a mean
squared error near zero, predicting the nominal motor current with high
fidelity. That low held-out validation loss evaluates accuracy only on
the states the demonstrator visited. It contains zero samples of the
machine recovering from actuator current saturation, sensor noise,
transport delay, or an incorrect command issued three time steps
earlier. The policy learns how to remain near the nominal setpoint when
already there, but learns nothing about how to steer back toward it once
displaced.

Let the policy operate in discrete time over an execution horizon of
\(T\) control steps, issuing an action \(a_t\) at each interval
\(\Delta t\). Suppose that under the demonstrated distribution of
observations, the policy has an independent probability \(\epsilon\) of
making an action error at any single time step, where an error is
defined as an action that deviates from the expert envelope by enough to
perturb the physical state beyond nominal bounds. The probability that
the policy executes a single step without an error is \(1 - \epsilon\).
If each decision were independent, the probability that the policy
completes an entire rollout of \(T\) consecutive steps without
committing a single error would be \((1 - \epsilon)^T\). The probability
of committing at least one error across the trajectory is therefore
\[P(\text{at least one error}) = 1 - (1 - \epsilon)^T\] To derive the
standard cumulative cost bound analytically, let
\(c(s_t, a_t) \in [0, 1]\) represent a bounded instantaneous cost
incurred at step \(t\), where \(c = 0\) corresponds to matching expert
behavior and \(c = 1\) denotes a task failure. At step \(t=1\), the
policy starts on the demonstrated distribution, so the probability of
making an error is bounded by \(\epsilon\). At step \(t\), if the policy
has made an error at any previous step \(\tau < t\), the machine may
have drifted off-distribution, potentially incurring a cost of up to
\(1\) at every subsequent step. Under the assumption that making an
error introduces a transition to an off-distribution state and that the
policy cannot recover on its own, the probability that the system has
left the demonstrated support by step \(t\) is bounded by the union
bound of errors over preceding steps, which grows as \(t \epsilon\).
Summing the expected cost across all \(T\) steps yields the cumulative
trajectory cost \(J\):
\[J = \sum_{t=1}^T \mathbb{E}[c(s_t, a_t)] \le \sum_{t=1}^T t \epsilon = \epsilon \frac{T(T+1)}{2} = \mathcal{O}(T^2 \epsilon)\]
The quadratic dependency on time \(T^2\) represents the structural price
of open-loop supervised training applied to a closed-loop dynamical
system. The first error incurs a direct penalty, and every subsequent
step spends the remainder of the rollout accumulating cost because the
demonstration dataset never provided examples of how to return to the
nominal path.

\protect\phantomsection\label{callout-mathux2a-6.3}
\begin{fbxSimple}{callout-math}{Fermi Math:}{Compounding Covariate Drift vs. Closed-Loop Recovery}
\phantomsection\label{callout-math*-6.3}

\emph{Problem statement:} An autonomous mobile delivery robot operates
at a control frequency of \(50\text{ Hz}\) (\(\Delta t = 20\text{ ms}\))
across a \(10.0\text{ s}\) transit corridor (\(T = 500\) steps). A
learned behavioral cloning policy achieves a \(99.5\%\) one-step
accuracy on validation data (\(\epsilon = 0.005\)). If an uncorrected
single-step error introduces an unmodeled lateral acceleration drift
bias of \(a_{\text{bias}} = 0.05\text{ m/s}^2\), calculate:

\begin{enumerate}
\def\labelenumi{\arabic{enumi}.}
\tightlist
\item
  The probability \(P(\text{drift})\) that the robot encounters at least
  one off-distribution perturbation during transit.
\item
  The expected cumulative trajectory lateral displacement drift under
  open-loop cloning versus closed-loop on-policy recovery (where
  recovery time constant
  \(\tau_{\text{rec}} = 5\Delta t = 0.10\text{ s}\)).
\end{enumerate}

\emph{Step 1: Probability of at least one error}
\[P(\text{at least one error}) = 1 - (1 - \epsilon)^T = 1 - (1 - 0.005)^{500} = 1 - 0.0816 = 91.84\%\]
Even with \(99.5\%\) one-step precision, an off-distribution departure
is virtually guaranteed over a \(10\text{ s}\) horizon.

\emph{Step 2: Cumulative trajectory drift under open-loop cloning} Under
open-loop cloning, once perturbed at step \(\tau\), lateral displacement
error at step \(t > \tau\) grows uncorrected as
\(d(t) \approx \frac{1}{2} a_{\text{bias}} ((t - \tau)\Delta t)^2\).
Summing the expected drift across all steps under compounding covariate
departure yields:
\[D_{\text{cloning}} \approx \frac{1}{2} a_{\text{bias}} \Delta t^2 \sum_{t=1}^T t^2 \epsilon \approx \frac{1}{6} a_{\text{bias}} \Delta t^2 T^3 \epsilon\]
Substituting parameters: \[\begin{aligned}
D_{\text{cloning}} &\approx \frac{1}{6} \times (0.05\text{ m/s}^2) \times (0.02\text{ s})^2 \times (500)^3 \times 0.005 \\
&= \frac{1}{6} \times 0.05 \times (4.0\times 10^{-4}) \times (1.25\times 10^8) \times 0.005 \approx 2.08\text{ m}
\end{aligned}\]

\emph{Step 3: Cumulative trajectory drift under on-policy recovery}
Under on-policy recovery training, the policy actively steers back to
the corridor within \(\tau_{\text{rec}} = 5\Delta t\), bounding local
displacement per error to
\(d_{\text{local}} = \frac{1}{2} a_{\text{bias}} (5\Delta t)^2 = 0.5 \times 0.05 \times (0.10)^2 = 2.5\times 10^{-4}\text{ m} = 0.25\text{ mm}\).
Cumulative drift scales linearly as \(\mathcal{O}(T \epsilon)\):
\[D_{\text{recovery}} = T \epsilon \cdot d_{\text{local}} = 500 \times 0.005 \times 0.25\text{ mm} = 0.625\text{ mm}\]

\emph{Systems verdict:} The ratio between open-loop cloning drift
(\(2.08\text{ m}\)) and closed-loop recovery drift (\(0.625\text{ mm}\))
exceeds \(3300\times\). A \(99.5\%\) accurate supervised policy is
completely unusable on physical mobile or robotic hardware without
active closed-loop recovery training or runtime support monitors.

\end{fbxSimple}

\protect\phantomsection\label{callout-ruleux2a-6.4}
\begin{fbxSimple}{callout-rule}{Architect’s Rule of Thumb:}{The Imitation Horizon Limit}
\phantomsection\label{callout-rule*-6.4}

\emph{A behavioral cloning policy cannot safely execute beyond
\(T \approx 1/\epsilon\) control cycles without active closed-loop
support monitoring or recovery demonstrations.}

\end{fbxSimple}

As illustrated in Figure~\ref{fig-06-compounding-error}, this dynamic
feedback creates an endogenous cascade where a single perturbation at
step \(\tau\) shifts the physical state outside demonstrated support
\(\mathcal{S}_{\text{demo}}\), transferring execution into an
unconstrained out-of-distribution regime. While on-policy recovery
training bounds cumulative cost linearly as \(\mathcal{O}(T\epsilon)\),
behavioral cloning compounds cost quadratically as
\(\mathcal{O}(T^2\epsilon)\) until intercepted by a runtime support
monitor in the nervous system.

\begin{figure}[t]

\centering{

\includegraphics[width=1\linewidth,height=\textheight,keepaspectratio]{index_files/mediabag//Users/VJ/GitHub/MLSysBook-integrate-physical/physical/book/chapters/06-training/figures/fig06_compounding_error.pdf}

}

\caption{\label{fig-06-compounding-error}\textbf{Compounding Covariate
Shift and State Support Drift in Behavioral Cloning.} Open-loop
supervised training on demonstration dataset \(\mathcal{D}\) creates an
acute structural vulnerability under closed-loop physical execution. (a)
State-space trajectory and distribution drift: Expert demonstrations
remain tightly bound within demonstrated support
\(\mathcal{S}_{\text{demo}}\) (\(\pm 0.2\text{ mm}\) around nominal). At
step \(t = \tau\), an action error (\(P = \epsilon\)) perturbs the
physical plant, pushing subsequent state \(s_{\tau+1}\)
out-of-distribution (OOD) where the network loss is unconstrained. The
trajectory drifts quadratically until tripping a nervous-system
Mahalanobis support monitor (\(D_M \ge \gamma_{\text{trip}}\)) or
striking a mechanical stop, whereas an on-policy/recovery policy returns
to the support tube. (b) Cumulative error scaling: The expected
cumulative cost scales quadratically as
\(J = \mathcal{O}(T^2 \epsilon)\) under cloning versus linearly
\(\mathcal{O}(T \epsilon)\) under on-policy recovery. For horizon
\(T = 500\) steps with single-step error \(\epsilon = 0.005\), the
probability of at least one error reaches
\(P = 1 - (1-\epsilon)^T \approx 91.8\%\). (c) Causal feedback loop: The
learned policy emits syntactically valid floating-point actions that
pass standard range checks, but physical actuation shifts the sensor
observation \(o_{t+1} \notin \text{supp}(p_{\text{demo}})\), querying
the network off-distribution and amplifying subsequent tracking error.}

\end{figure}%

Increasing dataset size, scaling model parameter capacity, or tuning
architectures can drive the one-step error \(\epsilon\) down by orders
of magnitude, but this only postpones the divergence. If a policy
operating at \(100\text{ Hz}\) has its error rate reduced from
\(\epsilon = 10^{-3}\) to \(\epsilon = 10^{-5}\), the expected time
until the first distribution shift increases from \(1.0\text{ s}\) to
\(100.0\text{ s}\). For a short-duration task, this delay may allow the
machine to reach the terminal setpoint before the divergence manifests.
Yet the underlying vulnerability remains unaddressed. Any unmodeled
external disturbance, such as an uncalibrated payload mass, a voltage
sag at the motor driver, or a mechanical jolt, shifts the state into the
unvisited region immediately, triggering the same
\(\mathcal{O}(T^2 \epsilon)\) cascade regardless of how low the static
training loss was driven.

The operational hazard of this failure mode is that it occurs silently.
Unlike a software crash or an unhandled exception, a neural network
evaluated on an out-of-distribution observation does not halt execution
or raise an alarm. It continues forward inference, multiplying matrices
and outputting numerical tensor values that are syntactically valid
floating-point numbers within normal actuator voltage limits. The
nervous system sees valid commands, while the physical state drifts
steadily away from the demonstration envelope until the mechanism
encounters a hard mechanical stop, stalls against a surface, or triggers
an overcurrent fault. Detecting this condition before physical contact
occurs cannot be accomplished by monitoring the offline training loss.
It requires runtime telemetry that tracks state space support directly,
measuring Mahalanobis distance in feature space, estimating model
epistemic uncertainty, or enforcing an explicit envelope monitor in the
nervous system that trips when state trajectories leave the convex hull
of the demonstration dataset.

\section{Learning by Trial}\label{sec-training-6-3}

Learning by trial addresses the blind spots of imitation by generating
state-action trajectories under the policy's own closed-loop execution.
Instead of mimicking a curated sequence of nominal operator
demonstrations, the optimization algorithm exposes the learning system
to the consequences of its own imperfect outputs. When an exploratory
action perturbs an end-effector five millimeters off its nominal path,
the subsequent observation lands outside the demonstration envelope,
forcing the optimization loop to evaluate recovery actions from that
exact perturbed state. The policy accumulates gradient signals across
near-misses, degraded alignments, and unstable contacts that an expert
human or calibrated baseline controller never exhibits. By sampling
transitions directly from the state distribution induced by the current
weights, trial-based learning eliminates the covariate shift that causes
open-loop supervised imitation to compound errors quadratically over
long execution horizons.

The capability bought by closed-loop exploration comes with a steep
physical billing rate. Computing the hardware feasibility boundary
begins with the wall-clock time required to collect a training dataset.
If a policy requires \(N\) complete rollouts to converge on an effective
control law, each lasting an average duration of \(t_{\text{rollout}}\)
seconds and followed by an automated or manual reset taking
\(t_{\text{reset}}\) seconds, the total wall-clock duration
\(T_{\text{wall}}\) across \(P\) safe parallel physical test benches is
governed by
\[T_{\text{wall}} = \frac{N (t_{\text{rollout}} + t_{\text{reset}})}{P}\]
As an analytical example, consider a contact-rich assembly task where a
manipulator regulates contact force while seating a connector over an
active horizon of \(t_{\text{rollout}} = 8.0\text{ s}\). Reinforcement
learning algorithms operating on high-dimensional sensory inputs
routinely require \(N = 10^6\) trials to discover coordinated
force-velocity profiles. On a single physical machine with \(P = 1\),
the raw execution time alone accounts for \(8.0 \times 10^6\text{ s}\),
which represents roughly \(92.6\text{ days}\) of uninterrupted
round-the-clock operation before accounting for any reset overhead.

Factoring in the reset phase reveals the true temporal bottleneck of
physical experimentation. If an automated fixture returns the workpiece
and robot arm to a valid initial state in
\(t_{\text{reset}} = 4.0\text{ s}\), the total collection time for one
million trials expands to \(1.2 \times 10^7\text{ s}\), or
\(138.9\text{ days}\). When exploratory trajectories drop parts, jam
pins, or trigger joint overcurrent trips, automated resetting becomes
impossible, requiring a human operator to enter the work cell, clear the
obstruction, and re-index the mechanism. If manual interventions take an
average of \(t_{\text{reset}} = 30.0\text{ s}\) across failed attempts,
the wall-clock requirement exceeds \(440.0\text{ days}\). Scaling
physical parallelism to clusters of \(P = 14\) synchronized rigs
(Figure~\ref{fig-06-robot-arm-farm}) compresses calendar duration to
several weeks (\citeproc{ref-levine2018learning}{Levine et al. 2018}),
but it multiplies the physical footprint, capital investment, electrical
power delivery, and human supervisory labor proportionally.

\begin{figure}[t]

\centering{

\includegraphics[width=0.85\linewidth,height=\textheight,keepaspectratio]{/Users/VJ/GitHub/MLSysBook-integrate-physical/physical/book/chapters/06-training/figures/fig06_real_robot_arm_farm.png}

}

\caption{\label{fig-06-robot-arm-farm}\textbf{Physical Policy Training
Farm and Hardware Parallelism.} Large-scale robotic data collection
cluster comprising 14 custom articulated manipulators operating
concurrently over months to collect over 800,000 physical grasp attempts
for vision-based hand-eye coordination (Levine et al., 2018). Each
robotic station consists of a 7-DoF arm, a two-finger compliant gripper,
an over-the-shoulder RGB camera, and a dedicated bin containing diverse
household objects. Scaling hardware parallelism to \(P=14\) compresses
the wall-clock calendar time
\(T_{\text{wall}} = N(t_{\text{rollout}} + t_{\text{reset}})/P\), but
incurs heavy capital expenditure, continuous mechanical wear across
gearboxes and linkage bearings, thermal buildup in motor windings, and
substantial supervisory overhead to clear dropped objects and mechanical
jams. Attribution: Levine et al.~(2018), IJRR.}

\end{figure}%

Converting that same sample count into physical fatigue metrics
demonstrates why time is often not the primary limiting constraint. In
our analytical example, an 8-second insertion sequence operating under a
\(100\text{ Hz}\) inner control loop executes 800 torque commands per
trial, subjecting each joint to approximately \(k = 4\) reciprocating
load cycles. Across \(N = 10^6\) trials, the drivetrain accumulates
\(4.0 \times 10^6\) bidirectional load reversals. If the joint relies on
a standard strain-wave gear with a manufacturer-specified fatigue life
of \(L_{10} = 1.5 \times 10^6\text{ cycles}\) at rated output torque
(ISO 281), the training campaign consumes the certified operational life
of nearly three full gearheads per articulated joint before the policy
learns to seat the connector. Simultaneously, one million physical
contact events crush compliant silicone gripper pads, generate
micro-notches in aluminum alignment jigs, and cycle electrical wiring
looms toward fatigue fracture. At an average electrical power demand of
\(350\text{ W}\) per test bench, executing \(1.2 \times 10^7\text{ s}\)
of automated trials draws \(1167\text{ kWh}\) of electrical energy,
transforming policy training into an industrial-scale manufacturing
operation.

These wear mechanics highlight a fundamental difference between software
tensors and mechanical hardware. In a purely computational environment,
evaluating an exploratory branch that produces numerical instability
costs nothing beyond discarding a GPU memory buffer. In a physical
machine, exploration is thermodynamic, irreversible, and potentially
destructive across all machine archetypes:

\begin{itemize}
\tightlist
\item
  In \emph{Mobile Systems}, exploratory steering torque spikes skid
  tires across asphalt, rapidly abrading tread rubber and overheating
  motor inverters during stall recoveries.
\item
  In \emph{Manipulators}, unconstrained exploratory torque commands
  drive end-effectors into rigid aluminum tables, bending load-cell
  flexures, stripping harmonic drive teeth, and heating brushless DC
  motor stators past their \(155\,^\circ\text{C}\) Class F insulation
  threshold (where copper resistance rises \(+39.3\%\) at
  \(125^\circ\text{C}\), depressing peak torque and accelerating
  \(I^2 R\) Joule heating).
\item
  In \emph{Process Systems}, exploratory pressure pulsing degrades valve
  seat elastomeric seals, accelerates cavitation pitting on pump
  impellers, and cycles high-voltage DC contactors toward premature arc
  erosion.
\end{itemize}

Furthermore, physical degradation alters the plant during the training
process itself. As contact surfaces burnish, gear backlash expands, and
structural fasteners loosen under repetitive vibration, the physical
transition distribution \(p(s_{t+1} \mid s_t, a_t)\) at trial
\(n = 900{,}000\) differs measurably from the pristine dynamics at trial
\(n = 1\). The machine being controlled mutates under the physical cost
of collecting its own training data.

Compounding the physical burden of data collection is the structural
fragility of the reward function. In trial-based learning, the scalar
reward serves as the formal interface specification between the
engineer's operational intent and the policy's numerical optimization.
Any behavior, physical constraint, or protective boundary not explicitly
measured and penalized in that scalar signal represents an unconstrained
degree of freedom that the optimization algorithm will exploit. If a
reward function incentivizes rapid setpoint tracking without
incorporating high-order penalty terms for instantaneous jerk and
winding temperature, the policy converges toward aggressive bang-bang
switching that vibrates mounting brackets and overheats drive
electronics. If contact force is unmeasured, the policy may discover
that driving an actuator to its maximum current against a rigid
mechanical stop creates a kinematic brace that artificially eliminates
tracking jitter. The learned network does not satisfy the designer's
unstated common sense. It satisfies the exact mathematical formula
presented to its objective function.

\protect\phantomsection\label{callout-ruleux2a-6.5}
\begin{fbxSimple}{callout-rule}{Architect’s Rule of Thumb:}{The Contact Fidelity Invariant}
\phantomsection\label{callout-rule*-6.5}

\emph{Never trust a simulated contact policy whose success depends on
instantaneous momentum absorption or zero-delay friction engagement; if
it penetrates the digital surface, it will shatter the physical
workpiece.}

\end{fbxSimple}

Trial-based learning on physical hardware is not categorically
impossible. For low-dimensional systems with trivial reset dynamics,
negligible wear rates, and wide safety margins, direct on-hardware
reinforcement learning remains technically viable. Simulation becomes
mandatory at the threshold where the required sample count exceeds the
allowable time, budget, structural life, or risk tolerance of the
physical apparatus. To collect the \(10^6\) to \(10^9\) transitions
needed for complex motor policies without destroying physical hardware,
practitioners shift the rollout loop into a digital physics engine
executed in parallel across computing clusters.

This transition substitutes computational hardware for physical
mechanisms, but it alters the contract between training and reality. A
policy trained inside a simulator avoids bending linkages or burning out
stators during exploration, but it inherits both the omissions of the
reward function and the simulator's specific mathematical model of
mechanics. Every unmodeled friction discontinuity, idealized rigid-body
contact impulse, uncharacterized actuator latency, and simplified cable
compliance in the physics engine becomes an assumption the policy learns
to rely upon. When the converged network is transferred to the physical
robot, those simulation inaccuracies re-emerge as unmodeled disturbances
and tracking failures at the hardware interface.

\section{The Simulator as a System Component}\label{sec-training-6-4}

Treating the simulator as a development tool external to the machine is
an architectural error. A compiler produces machine instructions from
source code, but the resulting binary does not retain a hidden
dependency on the compiler's internal memory allocator once executing on
standard hardware. A simulator behaves differently. It is an upstream
policy-production component whose physical assumptions, numerical
approximations, and boundary simplifications become permanently embedded
in the policy parameters. When the trained neural network is copied onto
the target machine, the simulator is absent from the execution path, yet
the policy acts as though every unmodeled dynamic in that simulator
remains true. The simulator is an active component of the system
architecture, with its own specification, tolerances, and failure modes.

\protect\phantomsection\label{callout-definitionux2a-6.6}
\begin{fbxSimple}{callout-definition}{Definition:}{Sim-to-Real Reality Gap Vector}
\phantomsection\label{callout-definition*-6.6}

The \textbf{sim-to-real reality gap vector} is the coupled,
multi-channel decomposition of physical discrepancies between simulated
and hardware environments spanning dynamics and contact mechanics
(\(\Delta_{\text{dyn}}\)), sensor transduction and noise
(\(\Delta_{\text{sens}}\)), perceptual feature representations
(\(\Delta_{\text{vis}}\)), and asynchronous transport latency and jitter
(\(\Delta_{\text{lat}}\)), each carrying distinct physical units,
transfer limits, and failure modes.

\end{fbxSimple}

Like any hardware subsystem, the simulator requires a formal interface
contract defined at the first budgetable quantities. This contract
specifies observation dimensions, action bounds, simulation time step
\(\Delta t\), transport latency between command issuance and physical
actuation, actuator torque and current limits, sensor noise
distributions, reset distributions, and termination conditions. If the
simulator specifies an observation loop executing at \(100\text{ Hz}\)
with an assumed zero-delay sensor channel and Gaussian noise with
variance \(\sigma^2 = 1.0 \times 10^{-4}\text{ (rad/s)}^2\), the policy
learns feedback gains tuned to those precise signal statistics. If the
physical machine delivers those state observations at \(100\text{ Hz}\)
but with an unmodeled \(15\text{ ms}\) bus transport delay, the phase
margin built during training is depleted before the policy executes its
first physical control cycle. As demonstrated in the robotic
manipulation testbed in Figure~\ref{fig-06-sim2real-dactyl},
transferring a policy synthesized inside a numerical engine to physical
hardware requires enforcing an exact match between the real
sensor-actuator cage and its randomized simulated twin
(\citeproc{ref-andrychowicz2020learning}{{Andrychowicz et al.} 2020}).

\begin{figure}[t]

\centering{

\includegraphics[width=0.95\linewidth,height=\textheight,keepaspectratio]{/Users/VJ/GitHub/MLSysBook-integrate-physical/physical/book/chapters/06-training/figures/fig06_real_sim2real_dactyl.jpg}

}

\caption{\label{fig-06-sim2real-dactyl}\textbf{Sim-to-Real Robotic
Transfer Setup with Digital Twin Simulation.} Authentic experimental
apparatus for vision-based dexterous in-hand manipulation from OpenAI's
Dactyl project (Andrychowicz et al., 2020). \emph{(Left)} The physical
experimental cage housing a 24-DoF anthropomorphic Shadow Dexterous
Hand, surrounded by 16 PhaseSpace infrared motion-tracking cameras and 3
Basler RGB cameras for ground-truth state validation and multi-view
vision. \emph{(Right)} The high-throughput MuJoCo simulated digital twin
environment used to train reinforcement learning policies across
billions of synthetic transitions. Robust sim-to-real transfer is
achieved by applying extensive domain randomization over surface
friction, restitution, joint damping, actuator time constants, visual
lighting, and camera calibration parameters, forcing the neural policy
to treat unmodeled physical variations as noise while operating within
the supported deployment envelope. Attribution: Andrychowicz et
al.~(2020), IEEE T-RO.}

\end{figure}%

\protect\phantomsection\label{callout-incidentux2a-6.7}
\begin{fbxSimple}{callout-incident}{Landmark Engineering Incident:}{OpenAI Dactyl In-Hand Dexterity Friction Collapse}
\phantomsection\label{callout-incident*-6.7}

\emph{System:} Anthropomorphic 24-DoF Shadow Dexterous Hand controlled
by an end-to-end recurrent neural network policy (LSTM) trained in
MuJoCo to rotate a physical cube to arbitrary target orientations using
three RGB camera feeds and joint encoder telemetry.

\emph{Incident:} During early hardware deployment trials, policies that
achieved \(>98\%\) task success in simulation suffered catastrophic
physical grasp failure within \(3\text{--}8\text{ s}\) on hardware. The
robotic fingers either dropped the object immediately or squeezed the
cube with extreme, rigid force that stalled the tendon drive actuators
and tripped thermal overcurrent protection within \(45\text{ s}\).

\emph{Root cause and physics of failure:} The MuJoCo physics engine
modeled fingertip contact using an idealized isotropic Coulomb friction
cone with instantaneous, memoryless velocity transitions
(\(F_f = \mu F_N \operatorname{sgn}(v)\)) and soft constraint
regularization. In physical reality, the silicone rubber fingertip pads
exhibit:

\begin{enumerate}
\def\labelenumi{\arabic{enumi}.}
\tightlist
\item
  Non-linear elastomeric compliance and contact area growth under normal
  load.
\item
  Velocity-dependent Stribeck friction curves where static friction
  (\(\mu_s \approx 0.85\)) drops precipitously to dynamic friction
  (\(\mu_k \approx 0.45\)) upon microscopic slip.
\item
  Severe friction hysteresis, where the effective friction coefficient
  during unloading lags behind that during loading due to polymer chain
  relaxation.
\end{enumerate}

The simulated policy discovered high-frequency finger gaiting maneuvers
(\(>35\text{ Hz}\) micro-repositioning), assuming that sticking friction
would instantly re-engage the moment fingertip sliding velocity crossed
zero. On physical hardware, the delayed contact re-engagement and
elastomeric shear relaxation caused the cube to slip free from the
grasp.\sidenote{\footnotesize OpenAI's Dactyl project bridged physical elastomeric
  fingertip friction gaps by combining domain randomization with an
  LSTM. The recurrent hidden state accumulated sensory residuals across
  timesteps, performing implicit online system identification of
  friction and mass.}

\emph{Systems resolution:} The engineering team had to introduce
Automatic Domain Randomization (ADR) across fingertip friction
coefficients (\(\mu \in [0.1, 2.0]\)), contact damping, joint damping,
tendon elasticity, and actuator latencies (\(0\text{--}40\text{ ms}\)),
forcing the LSTM hidden states to act as an implicit real-time system
identification estimator capable of adapting to unmodeled contact
compliance and friction hysteresis online.

\emph{Systems takeaway:} Optimization algorithms aggressively exploit
the non-dissipative shortcuts of idealized contact models. Never trust a
contact policy whose success depends on instantaneous friction
transitions without modeling contact compliance, stiction hysteresis,
and velocity-dependent Stribeck curves.

\end{fbxSimple}

A clear distinction exists between an inaccurate parameter value and an
omitted physical mechanism. When a simulator uses an inaccurate rotor
inertia, assigning \(J = 2.0 \times 10^{-3}\text{ kg}\cdot\text{m}^2\)
instead of a true \(2.5 \times 10^{-3}\text{ kg}\cdot\text{m}^2\), the
equations of motion still compute an acceleration-dependent inertial
torque. The parameter error alters the magnitude of the response, but
the causal structure remains intact. When a simulator omits an entire
physical mechanism, such as drive train backlash, structural compliance,
or inter-channel communication jitter, the simulator does not set the
corresponding physical effect to zero in a physically supported sense.
It provides no mechanism at all. An omitted delay or compliance is
unsupported. A policy trained in an environment with omitted joint
compliance assumes infinite structural stiffness, discovering trajectory
strategies that command instantaneous torque reversals that shatter gear
teeth or excite high-frequency structural resonances on physical
hardware.

\subsubsection*{Non-Smooth Contact Physics and Solver
Relaxations}\label{non-smooth-contact-physics-and-solver-relaxations}
\addcontentsline{toc}{subsubsection}{Non-Smooth Contact Physics and
Solver Relaxations}

The most severe structural omission in physical simulators arises from
the mathematical formulation of contact mechanics. In physical reality,
contact between rigid or semi-rigid bodies is fundamentally non-smooth.
Let \(\phi(q) \ge 0\) denote the signed distance (gap function) between
two surfaces in generalized coordinates \(q\), and let
\(\lambda_N \ge 0\) represent the normal contact force. Continuous
contact mechanics is governed by the unilateral \textbf{Signorini
complementarity condition}:
\[0 \le \phi(q) \perp \lambda_N \ge 0 \quad \iff \quad \phi(q) \ge 0, \quad \lambda_N \ge 0, \quad \phi(q) \cdot \lambda_N = 0\]
This complementarity condition enforces two physical truths: when bodies
are separated (\(\phi(q) > 0\)), the normal contact force must be
identically zero (\(\lambda_N = 0\)). When contact is established
(\(\phi(q) = 0\)), the normal contact force can take any non-negative
compressive value (\(\lambda_N \ge 0\)) required to prevent
interpenetration. Simultaneously, tangential friction is governed by the
Coulomb friction cone:
\[\|\lambda_T\|_2 \le \mu \lambda_N, \quad v_T \ne 0 \implies \lambda_T = -\mu \lambda_N \frac{v_T}{\|v_T\|_2}\]
where \(v_T\) is the relative tangential velocity at the contact point
and \(\mu\) is the friction coefficient. When sliding occurs
(\(v_T \ne 0\)), friction acts at the boundary of the cone opposing
motion; when sticking (\(v_T = 0\)), friction acts as an inequality
constraint residing strictly within the cone.

This formulation introduces two profound computational and physical
difficulties:

\begin{enumerate}
\def\labelenumi{\arabic{enumi}.}
\tightlist
\item
  \textbf{Mathematical Non-Smoothness and Painlevé Paradoxes}: At the
  instant of collision, velocity changes discontinuously
  (\(v^+ \ne v^-\)), acceleration mappings become set-valued and
  non-differentiable, and rigid Coulomb friction models can yield
  non-unique or non-existent acceleration solutions (the classical
  Painlevé paradox).\sidenote{\footnotesize Rigid contact with Coulomb friction is
    non-smooth and exhibits Painlevé paradoxes. Physics engines resolve
    ill-posedness via discrete Linear Complementarity Problems (LCP) or
    convex penalty relaxations, introducing numerical softness and
    penetration artifacts that policies exploit during training.}
\item
  \textbf{Computational Solvers vs.~Numerical Relaxations}: Rigid-body
  physics engines resolve these difficulties through one of two
  numerical paths:

  \begin{itemize}
  \tightlist
  \item
    \emph{Discrete Linear Complementarity Problem (LCP) / Mixed
    Complementarity Problem (MCP) Solvers} (e.g., Bullet, ODE, Drake):
    Discretize the contact dynamics into inequality-constrained
    quadratic programs at each timestep \(\Delta t\). While maintaining
    strict non-penetration, LCP solvers are computationally expensive
    (\(\mathcal{O}(m^3)\) in contact points), non-differentiable across
    contact mode switches, and susceptible to solver stalling when
    geometries pinch.
  \item
    \emph{Convex Relaxations and Soft Penalty Formulations} (e.g.,
    MuJoCo's elliptic contact cones and soft constraint regularizations,
    PhysX penalty models): Smooth the non-smooth boundary by replacing
    strict complementarity with compliant potential energies or
    regularized dual dissipation functions:
    \[\lambda_N = k_p \max(0, -\phi(q)) - k_d \dot{\phi}(q)\] where
    \(k_p\) and \(k_d\) represent numerical stiffness and damping
    parameters, allowing slight interpenetration (\(\phi(q) < 0\)) to
    guarantee rapid, continuous solver convergence at high frame rates
    (\(>1000\text{ Hz}\)).
  \end{itemize}
\end{enumerate}

\subsubsection*{The Sim-to-Real Contact
Exploit}\label{the-sim-to-real-contact-exploit}
\addcontentsline{toc}{subsubsection}{The Sim-to-Real Contact Exploit}

Reinforcement learning algorithms optimize ruthlessly against the exact
physics presented by the solver. In simulators that implement soft
contact relaxations or penalty springs, the optimizer discovers that
``sinking'' into the surface by \(1.0\text{--}3.0\text{ mm}\)
(\(\phi(q) < 0\)) provides two unphysical affordances:

\begin{enumerate}
\def\labelenumi{\arabic{enumi}.}
\tightlist
\item
  The artificial spring-damper reaction acts as an instantaneous energy
  sink, arresting downward momentum in a single integration step without
  bouncing or plastic deformation.
\item
  Interpenetration increases the effective contact surface area in the
  solver, generating unphysical lateral traction forces that prevent
  slipping during aggressive maneuvers.
\end{enumerate}

When transferred to physical hardware---where a machined steel or
aluminum fixture has an elastic stiffness \(k > 10^7\text{ N/m}\) and
zero allowable penetration---the physical plant provides no such
compliance. A policy that executes an unbraked plunge at
\(v_z = 0.35\text{ m/s}\) does not experience a gentle \(2.0\text{ ms}\)
penalty deceleration; instead, the physical collision decelerates the
moving mass in less than \(0.5\text{ ms}\). By the impulse-momentum
theorem
(\(F_{\text{peak}} \approx m \Delta v / \Delta t_{\text{impact}}\)), the
contact force spikes past \(450\text{ N}\) at a rate
\(dF/dt > 3.8\times 10^4\text{ N/s}\), shattering load-cell flexures,
shearing alignment pins, and tripping motor driver overcurrent
protection.

For a machine that regulates motion and then touches a surface, the
fidelity of the simulator cannot be evaluated through the internal
mathematics of its contact solver. It must be budgeted through
observable mechanical quantities, specifically contact onset timing,
peak normal force, total transferred impulse, tangential slip distance,
and contact release delay. Consider an analytical example of a robotic
end-effector approaching a rigid workpiece at \(0.20\text{ m/s}\). In a
standard penalty-based physics engine using a discrete time step of
\(2.0\text{ ms}\), the collision resolves across a single time step,
producing a simulated peak force of \(450\text{ N}\) and an impulse of
\(0.90\text{ N}\cdot\text{s}\). On physical hardware, contact pad
deformation and structural compliance distribute that same momentum
transfer over \(18\text{ ms}\), yielding a true peak force of
\(85\text{ N}\) with an impulse of \(0.88\text{ N}\cdot\text{s}\). The
integrated impulse matches within \(2.2\%\), but the peak force differs
by more than a factor of five. A policy that learned to detect contact
onset from the sharp simulated force spike will fail to trigger on
physical hardware, continuing to push until the drive motors hit current
limits.

Candidate omissions must be identified and cataloged based on the
dynamics required by the task. Dry Coulomb friction discontinuities at
zero velocity, nonlinear structural compliance, gearbox backlash,
voltage saturation under high back-EMF, bus transport delays, and sensor
packet loss represent standard omissions in rigid-body physics engines.
A policy that operates purely at high continuous velocity may tolerate
coarse friction models because the system never lingers near the
zero-velocity stiction threshold. Conversely, a policy that regulates
force while holding a static contact will enter uncontrolled limit-cycle
oscillations if static friction, stick-slip transitions, and actuator
deadbands are omitted from the simulator.

Establishing the validity of a simulator component requires paired
probing. The engineer applies identical open-loop command sequences to
both the simulator and the physical hardware, measuring the difference
between the resulting state trajectories against explicit error budgets.
For a regulating actuator, injecting a step torque command of
\(1.5\text{ N}\cdot\text{m}\) into both systems yields concrete
comparative metrics. If the simulator predicts a position settling time
of \(40\text{ ms}\) with zero overshoot, while the physical drive
exhibits a \(110\text{ ms}\) settling time and a \(2.2\text{ mm}\)
position overshoot due to unmodeled stator heating and structural
flexure, the discrepancy is budgeted directly in milliseconds and
millimeters. When a measured difference exceeds the allowable tolerance
for the operational envelope, the simulator is formally defective for
that training task.

Because the structural assumptions of the simulator are frozen into the
neural network weights during optimization, the simulator specification
forms an immutable element of the policy's configuration and identity.
Every trained policy artifact must document its complete simulator
provenance, including the physics engine build and version, integrator
method and step size, contact solver parameters, geometry mesh and
collision hull identifiers, domain randomization distributions, and
documented regimes where the simulation is known to be physically
invalid. Deploying a learned policy without its simulator provenance is
equivalent to shipping compiled firmware without recording the target
architecture, clock speed, or register map for which it was built.

\section{What the Gap Is Made Of}\label{sec-training-6-5}

The sim-to-real gap is routinely treated as a single scalar distance, as
though one could measure the divergence between a simulated world and a
physical machine with an aggregate loss or a probability metric. That
abstraction fails because sim-to-real is not one gap. It is several
distinct mismatches, and they are not equally hard to close. Errors in
different physical channels carry different units, propagate through
different paths in the closed loop, and produce fundamentally different
failure modes. Sim-to-real mismatch is therefore a vector of coupled
components spanning dynamics, sensing, appearance, and timing. Sizing
this vector requires paired command-response traces recorded under
identical open-loop and closed-loop excitations. Sizing is not an
abstract qualification; it is an arithmetic accounting of physical
errors against the tolerances of the task. As structured in
Table~\ref{tbl-06-reality-gap}, decomposing the reality gap across
constitutive physical channels reveals that each dimension introduces
distinct numerical artifacts, exploits specific control shortcuts, and
demands tailored hardware and nervous system interceptors. Consider an
analytical example of a linear drive tasked with approaching a rigid
surface at \(0.20\text{ m/s}\) and regulating normal contact force to
\(10.0\text{ N}\), where no bench measurements yet exist and every
discrepancy must be budgeted from first principles.

Dynamics mismatch measures differences between commanded actuator effort
and the resulting physical motion or contact force across the
operational envelope. In the analytical example, applying a step torque
command of \(2.00\text{ N}\cdot\text{m}\) reveals how the physics engine
diverges from hardware. The simulator models a rigid link with zero
backlash and instantaneous contact stiffness, predicting a force rise
rate of \(500\text{ N/s}\), reaching the target \(10.0\text{ N}\) in
\(20\text{ ms}\) with zero steady-state deflection. On physical
hardware, a gearbox compliance of
\(1.5\times 10^{-3}\text{ rad/N}\cdot\text{m}\) and an unmodeled stator
drive delay of \(8\text{ ms}\) introduce phase lag, while soft surface
compliance limits the force rise rate to \(250\text{ N/s}\). Reaching
the target force requires \(40\text{ ms}\) on hardware, producing a
transient contact overshoot impulse of \(0.15\text{ N}\cdot\text{s}\).
If a policy learned to associate the simulated \(20\text{ ms}\)
stiffness boundary with stable contact, it interprets the delayed force
rise as a failure to touch, commands additional motor torque, and drives
the physical mechanism \(3.0\text{ mm}\) past the nominal boundary into
a current-limit shutdown.

Sensing mismatch measures discrepancies between the true physical state
and the observation vector delivered to the policy network. In the
analytical example, the simulated depth sensor provides clean range
measurements with a uniform quantization step of \(1.0\text{ mm}\). The
physical depth camera, measured under matched ambient illumination of
\(450\text{ lux}\), exhibits a static calibration bias of
\(+3.2\text{ mm}\), zero-mean Gaussian noise with standard deviation
\(\sigma = 1.8\text{ mm}\), and periodic optical dropout bursts where
four consecutive frames return null values during specular reflections.
Simultaneously, the physical force sensor exhibits an uncalibrated
thermal drift of \(0.40\text{ N}\) per \(10^\circ\text{C}\) temperature
rise and an analog-to-digital quantization step of \(0.05\text{ N}\).
When the policy executes its descent, the \(+3.2\text{ mm}\) range bias
causes the controller to switch from rapid approach to fine force
regulation \(3.2\text{ mm}\) above the physical surface. The four-frame
dropout burst, lasting \(133\text{ ms}\) at a \(30\text{ Hz}\) sampling
rate, leaves the policy blind to the true collision boundary during
descent, converting a controlled deceleration into an unbraked impact.

Appearance mismatch is not a mean squared error computed over raw camera
pixels. Image differences between photorealistic rendering and physical
camera frames are dominated by high-frequency background textures and
minor sensor noise patterns that may have no effect on policy action.
Appearance mismatch is sized by measuring the deviation in
policy-relevant perceptual outputs, such as estimated surface normals,
object keypoints, or spatial feature activations, under matched physical
geometry, illumination, materials, and camera exposure settings. In the
analytical example, rendering the contact surface under simulated
lighting produces an image whose raw pixels differ from the physical
camera frame by a mean intensity of 14 on an 8-bit scale (\(5.5\%\)).
However, passing both images through the policy spatial encoder yields
an estimated contact patch centroid offset of \(12\text{ pixels}\),
corresponding to a \(6.0\text{ mm}\) spatial localization error on the
physical fixture. Sizing the gap by the \(5.5\%\) pixel difference
suggests an acceptable visual match, but the \(6.0\text{ mm}\) spatial
feature error exceeds the \(2.0\text{ mm}\) chamfer tolerance of the
contact surface, causing the end-effector to strike the fixture
shoulder.

Timing mismatch governs the temporal alignment of observation capture,
policy evaluation, and actuation. Identical state transitions occurring
at different moments in physical time do not produce equivalent
closed-loop dynamics. The simulator assumes a synchronous discrete-time
step of \(\Delta t = 20\text{ ms}\) (\(50\text{ Hz}\)) in which state
observation, neural network inference, and motor torque application
execute instantaneously at the beginning of each cycle. On physical
hardware, camera exposure requires \(16.7\text{ ms}\), bus serialization
adds \(4.2\text{ ms}\), inference requires \(12.0\text{ ms}\) at
\(P_{50}\) and stretches to \(26.5\text{ ms}\) at \(P_{99}\) due to
shared memory contention, and motor bus transport adds
\(2.1\text{ ms}\). Total observation-to-actuation latency on hardware is
\(35.0\text{ ms}\) at \(P_{50}\) and \(49.5\text{ ms}\) at \(P_{99}\),
with a cycle-to-cycle jitter of \(\pm 6.2\text{ ms}\). Because physical
actuation arrives up to \(29.5\text{ ms}\) later than the simulator
assumed, the drive travels an unobserved distance of
\(0.20\text{ m/s} \times 0.0295\text{ s} = 5.9\text{ mm}\) before
corrective action begins, altering the damping and phase margin of the
closed-loop system.

As illustrated in Figure~\ref{fig-06-sim2real-gap}, sizing the
sim-to-real gap decomposes the divergence across these four distinct
physical axes, exposing why scalar loss metrics fail to capture
localized hardware failure modes.

\begin{figure}[t]

\centering{


\includegraphics[width=1\linewidth,height=\textheight,keepaspectratio]{index_files/mediabag//Users/VJ/GitHub/MLSysBook-integrate-physical/physical/book/chapters/06-training/figures/fig06_sim2real_gap.pdf}

}

\caption{\label{fig-06-sim2real-gap}\textbf{The Multidimensional
Sim-to-Real Reality Gap Vector.} Sizing sim-to-real transfer requires
paired-trace telemetry across four heterogeneous physical channels
rather than an aggregate scalar loss. (a) Dynamics and contact mismatch
(\(\Delta_{\text{dyn}}\) in
\(\text{N}, \text{rad}/\text{N}\cdot\text{m}\)): An \(8\text{ ms}\)
stator drive delay and soft compliance cause the physical force to rise
at \(250\text{ N/s}\) instead of the simulated \(500\text{ N/s}\),
resulting in a \(0.15\text{ N}\cdot\text{s}\) overshoot impulse and
\(3.0\text{ mm}\) over-travel into an overcurrent trip. (b) Actuator and
pipeline latency (\(\Delta_{\text{lat}}\) in \(\text{ms}, \text{mm}\)):
Segmented hardware latency (camera exposure \(16.7\text{ ms}\), bus
serialization \(4.2\text{ ms}\), neural inference \(12.0\text{ ms}\)
with a \(P_{99}\) tail of \(26.5\text{ ms}\), and motor bus transport
\(2.1\text{ ms}\)) creates a \(35.0\text{ ms}\) (\(P_{50}\)) to
\(49.5\text{ ms}\) (\(P_{99}\)) delay, producing \(5.9\text{ mm}\) of
unobserved blind travel that depletes control phase margin. (c) Sensing
and transduction (\(\Delta_{\text{sens}}\) in \(\text{mm}, \text{N}\)):
A \(+3.2\text{ mm}\) calibration bias triggers premature regulation,
while a 4-frame (\(133\text{ ms}\)) specular dropout burst blinds the
policy during descent. (d) Perceptual feature shift
(\(\Delta_{\text{vis}}\) in feature space): A seemingly trivial
\(5.5\%\) raw pixel discrepancy masks a \(12\text{ px}\)
(\(6.0\text{ mm}\)) spatial token centroid shift that exceeds the
\(2.0\text{ mm}\) chamfer tolerance, driving the tool into the fixture
shoulder.}

\end{figure}%

\begin{longtable}[]{@{}
  >{\raggedright\arraybackslash}p{(\linewidth - 10\tabcolsep) * \real{0.1667}}
  >{\raggedright\arraybackslash}p{(\linewidth - 10\tabcolsep) * \real{0.1667}}
  >{\raggedright\arraybackslash}p{(\linewidth - 10\tabcolsep) * \real{0.1667}}
  >{\raggedright\arraybackslash}p{(\linewidth - 10\tabcolsep) * \real{0.1667}}
  >{\raggedright\arraybackslash}p{(\linewidth - 10\tabcolsep) * \real{0.1667}}
  >{\raggedright\arraybackslash}p{(\linewidth - 10\tabcolsep) * \real{0.1667}}@{}}
\caption{\textbf{Multidimensional Sim-to-Real Reality Gap Taxonomy.}
Decomposition of the reality gap into constitutive physical channels
across dynamics, contact mechanics, sensing, and temporal
synchronization. Columns specify the unmodeled physical phenomenon,
representative quantitative discrepancy, the latent control shortcut
exploited by the optimizer, paired-trace verification metrics, and
architectural hardware/nervous system
mitigations.}\label{tbl-06-reality-gap}\tabularnewline
\toprule\noalign{}
\begin{minipage}[b]{\linewidth}\raggedright
Reality Gap Dimension
\end{minipage} & \begin{minipage}[b]{\linewidth}\raggedright
Unmodeled Physical Mechanism
\end{minipage} & \begin{minipage}[b]{\linewidth}\raggedright
Nominal Simulation Spec vs.~Real Hardware
\end{minipage} & \begin{minipage}[b]{\linewidth}\raggedright
Latent Policy Exploitation / Behavioral Failure
\end{minipage} & \begin{minipage}[b]{\linewidth}\raggedright
Paired-Trace Diagnostic Metric \& Tolerance
\end{minipage} & \begin{minipage}[b]{\linewidth}\raggedright
Hardware \& Nervous System Interceptor
\end{minipage} \\
\midrule\noalign{}
\endfirsthead
\toprule\noalign{}
\begin{minipage}[b]{\linewidth}\raggedright
Reality Gap Dimension
\end{minipage} & \begin{minipage}[b]{\linewidth}\raggedright
Unmodeled Physical Mechanism
\end{minipage} & \begin{minipage}[b]{\linewidth}\raggedright
Nominal Simulation Spec vs.~Real Hardware
\end{minipage} & \begin{minipage}[b]{\linewidth}\raggedright
Latent Policy Exploitation / Behavioral Failure
\end{minipage} & \begin{minipage}[b]{\linewidth}\raggedright
Paired-Trace Diagnostic Metric \& Tolerance
\end{minipage} & \begin{minipage}[b]{\linewidth}\raggedright
Hardware \& Nervous System Interceptor
\end{minipage} \\
\midrule\noalign{}
\endhead
\bottomrule\noalign{}
\endlastfoot
\textbf{Actuator \& Drivetrain Dynamics} & Gearbox harmonic compliance,
Coulomb stiction, velocity-dependent Stribeck friction, and back-EMF
voltage saturation & Idealized direct torque source
(\(J_{\text{rotor}}=0\), \(\tau_{\text{latency}}=0\text{ ms}\))
vs.~physical harmonic drive
(\(k_c=1.5\times 10^3\text{ N}\cdot\text{m/rad}\), backlash
\(\theta_b=0.08^\circ\), bus delay \(\tau=8\text{ ms}\)) & Policy
commands high-frequency bang-bang torque reversals (\(>50\text{ Hz}\)),
inducing mechanical chatter, gear fatigue, and stator winding thermal
runaway (\(>90\,^\circ\text{C}\)) & Step-torque response discrepancy:
phase lag \(\Delta \phi \le 5^\circ\) at \(10\text{ Hz}\); settling time
discrepancy \(\Delta t_{\text{settle}} \le 15\text{ ms}\); overshoot
\(\le 5\%\) & Low-pass action filtering
(\(f_{\text{cutoff}} \le 20\text{ Hz}\)), motor current RMS thermal
estimators, and actuator slew-rate limiters
(\(d\tau/dt \le \tau_{\text{max}}/\Delta t\)) \\
\textbf{Contact Mechanics \& Compliance} & Viscoelastic surface
deformation, contact hysteresis, non-smooth stick-slip transitions, and
finite restitution phase & Point-mass rigid contact via LCP/penalty
spring (\(k_p=10^6\text{ N/m}\), restitution
\(\Delta t_{\text{res}}=2\text{ ms}\)) vs.~physical contact pad
(\(k=8.0\times 10^4\text{ N/m}\), deformation span
\(\Delta t_{\text{def}}=18\text{ ms}\)) & Policy executes unbraked
high-velocity plunge (\(v_z > 0.4\text{ m/s}\)), relying on simulated
instantaneous energy absorption; shatters physical load-cell flexures
(\(F_{\text{peak}} > 450\text{ N}\)) & Contact impulse error
\(\Delta I \le 0.05\text{ N}\cdot\text{s}\); peak contact force
discrepancy \(\Delta F_{\text{peak}} \le 15\text{ N}\); contact onset
rise-time error \(\Delta t_{\text{rise}} \le 5\text{ ms}\) & Approach
velocity envelope governor
(\(v_{\text{safe}}(z) \propto \sqrt{2 a_{\text{max}} z}\)),
force-derivative clamp (\(dF/dt \le 1.0\times 10^4\text{ N/s}\)),
compliant mechanical flexures \\
\textbf{Sensor Transduction \& Noise Profiles} & Optical specular
dropout, CMOS rolling-shutter distortion, ambient lux-dependent noise,
and force-torque gauge thermal drift & Synthetic RGB-D with zero-bias
Gaussian noise (\(\sigma = 0.5\text{ mm}\), zero dropout) vs.~physical
structured-light depth (\(+3.2\text{ mm}\) bias, 4-frame dropout bursts,
drift \(0.4\text{ N}/10\,^\circ\text{C}\)) & Policy triggers premature
deceleration above surface or remains blind during multi-frame dropouts,
driving end-effector past hard fixture boundaries into stall overcurrent
& Sensor distribution Kolmogorov-Smirnov distance
\(D_{\text{KS}} \le 0.05\); spatial feature activation centroid error
\(\le 2.0\text{ mm}\); zero sustained dropout frames
(\(N_{\text{drop}} = 0\)) & Multi-sensor fusion (IMU/encoders + vision),
optical dropout holdover filters, thermal drift autocalibration,
independent proximity break-beams \\
\textbf{Temporal Latency \& Bus Asynchrony} & Variable camera exposure
times, PCIe/EtherCAT bus serialization, neural inference tail latency
(\(P_{99}\)), and non-deterministic OS scheduling & Synchronous
discrete-time lockstep (\(\Delta t = 20.0\text{ ms}\), zero transport
latency) vs.~asynchronous physical pipeline
(\(\tau_{\text{total}} = 35.0\text{ ms}\) at \(P_{50}\),
\(49.5\text{ ms}\) at \(P_{99}\), jitter \(\pm 6.2\text{ ms}\)) & Severe
phase margin erosion (\(>30^\circ\) loss); policy experiences
underdamped oscillations, hunting limit-cycles, and unobserved travel
during latency spikes (\(>5.9\text{ mm}\) blind drift) & End-to-end
command-response latency jitter
\(\Delta \tau_{\text{jitter}} \le 2.0\text{ ms}\); worst-case
observation age \(\tau_{\text{obs}} \le 25\text{ ms}\); zero deadline
overruns at \(50\text{ Hz}\) & Real-time PREEMPT\_RT / Xenomai
scheduling, action buffer delay queue injection in simulation, kinematic
extrapolation (Smith predictor) \\
\end{longtable}

Domain randomization is often treated as a method for erasing the
sim-to-real gap, but parameter randomization operates within strict
structural limits.\sidenote{\footnotesize Domain randomization optimizes expected
  return over \(\xi \sim \mathcal{P}(\Xi)\). Generalization holds only
  within the support of \(\mathcal{P}(\Xi)\), giving zero protection
  against omitted physical mechanisms.} Sampling friction coefficients
across a uniform distribution from \(\mu = 0.2\) to \(\mu = 0.8\), or
link masses across \(\pm 20\%\), broadens the coverage of training over
mechanisms that the simulator already represents. It cannot synthesize
an omitted physical mechanism, a missing state correlation, or an
unmodeled rare event. If the simulator differential equations omit
gearbox backlash, thermal resistance in motor windings, or thread
preemption in the operating system, no distribution over link mass will
force the policy to learn compensation for those dynamics. Adding
unstructured Gaussian noise to simulated observations to mask unmodeled
effects degrades performance, training the policy to lower its control
gains and narrow its operational bandwidth to avoid actions that excite
the unmodeled modes. As cataloged in
Table~\ref{tbl-06-domain-randomization}, parameter randomization defines
a narrow operating corridor bounded on one side by over-randomization
(which induces hyper-conservative, low-gain control laws) and on the
other by under-randomization (which leaves the policy brittle against
physical hardware dynamics, risking structural overload, motor thermal
saturation, and statutory contact force violations).

The \textbf{domain randomization corridor} is defined as the narrow
viable parameter distribution envelope in physics engine training
bounded on one side by under-randomization (which leaves the policy
brittle against unmodeled dynamics, risking mechanical shock or thermal
tripping) and on the other by over-randomization (which induces
hyper-conservative, low-gain control laws that degrade operational
throughput). As cataloged in Table~\ref{tbl-06-domain-randomization},
parameter randomization defines a narrow operating corridor bounded by
these physical constraints.

\subsubsection*{The Reproducible Domain Randomization
Workflow}\label{the-reproducible-domain-randomization-workflow}
\addcontentsline{toc}{subsubsection}{The Reproducible Domain
Randomization Workflow}

To prevent domain randomization from degenerating into ad-hoc parameter
tuning, robust physical AI engineering requires a structured, five-stage
reproducible workflow:

\begin{enumerate}
\def\labelenumi{\arabic{enumi}.}
\tightlist
\item
  \textbf{Stage 1: Hardware System Identification (SysID) Baseline}:
  Before initializing simulation, execute empirical calibration suites
  on the physical test bench. Measure nominal parameters
  \(\xi_0 = [m_0, \mathbf{I}_0, K_{t,0}, \tau_0, \mu_0]\) using
  frequency-response sweeps, static torque calibrations, and optical
  latency profilers.
\item
  \textbf{Stage 2: Physics-Grounded Parameter Bounds
  \(\mathcal{P}(\Xi)\)}: Ground randomization bounds directly in
  verified physical mechanisms rather than arbitrary heuristic
  intervals. Set motor torque constant variation
  \(K_t \in [0.105, 0.135]\text{ N}\cdot\text{m/A}\) to match the
  \(+39.3\%\) electrical resistance growth of copper windings
  (\(R_{\text{Cu}}(T) = R_0[1 + \alpha_{\text{Cu}}(T-T_0)]\)) and
  permanent-magnet thermal demagnetization across an operating
  temperature range of \(25^\circ\text{C}\) to \(125^\circ\text{C}\).
  Set friction bounds \(\mu \in [0.30, 0.85]\) to span dry clean
  aluminum down to lubricated or oxidized contact pads.
\item
  \textbf{Stage 3: Curriculum and Adaptive Domain Randomization (ADR)}:
  Training directly on a maximally randomized parameter distribution
  frequently causes optimization collapse, forcing the neural network
  into an ultra-low-gain regime that degrades operational velocity by
  \(>60\%\). Under ADR, simulation begins in a narrow corridor around
  nominal physics \(\xi_0\). As policy evaluation return exceeds
  declared performance thresholds (e.g., \(>95\%\) success across
  evaluation seeds), the environment dynamically expands the sampling
  boundaries
  \(\Xi_{\text{min}} \leftarrow \Xi_{\text{min}} - \Delta\Xi\) and
  \(\Xi_{\text{max}} \leftarrow \Xi_{\text{max}} + \Delta\Xi\),
  progressively exposing the learner to broader dynamical variation.
\item
  \textbf{Stage 4: Paired-Trace Discrepancy Auditing}: Apply identical
  open-loop and closed-loop command trajectories to both the simulated
  digital twin and the physical hardware bench. Evaluate discrepancy
  metrics across all four channels against declared tolerance budgets:
  phase lag \(\Delta \phi \le 5^\circ\) at \(10\text{ Hz}\), contact
  impulse error \(\Delta I \le 0.05\text{ N}\cdot\text{s}\), feature
  token offset \(\le 2.0\text{ mm}\), and latency jitter
  \(\Delta \tau \le 2.0\text{ ms}\). If discrepancies exceed budget, the
  simulation dynamic model must be updated before policy deployment.
\item
  \textbf{Stage 5: Dual-Brain Runtime Safety Enclosure}: Because domain
  randomization cannot protect against omitted physical mechanisms
  outside the simulation graph, the unprivileged neural policy executing
  on the Linux MPU must be bounded by deterministic monitors on the
  Real-Time MCU. The real-time core enforces hardware force-derivative
  clamps (\(dF/dt \le 1.0\times 10^4\text{ N/s}\)), torque slew-rate
  limits (\(d\tau/dt \le \dot{\tau}_{\text{max}}\)), and velocity
  governors at \(500\text{--}1000\text{ Hz}\), providing guaranteed
  physical protection regardless of out-of-distribution neural commands.
\end{enumerate}

\protect\phantomsection\label{callout-ruleux2a-6.8}
\begin{fbxSimple}{callout-rule}{Architect’s Rule of Thumb:}{The Domain Randomization Boundary}
\phantomsection\label{callout-rule*-6.8}

\emph{Randomizing parameters covers uncertainty within the model family,
but no distribution over friction or mass will teach a policy to
compensate for an omitted physical mechanism.}

\end{fbxSimple}

\begin{longtable}[]{@{}
  >{\raggedright\arraybackslash}p{(\linewidth - 10\tabcolsep) * \real{0.1667}}
  >{\raggedright\arraybackslash}p{(\linewidth - 10\tabcolsep) * \real{0.1667}}
  >{\raggedright\arraybackslash}p{(\linewidth - 10\tabcolsep) * \real{0.1667}}
  >{\raggedright\arraybackslash}p{(\linewidth - 10\tabcolsep) * \real{0.1667}}
  >{\raggedright\arraybackslash}p{(\linewidth - 10\tabcolsep) * \real{0.1667}}
  >{\raggedright\arraybackslash}p{(\linewidth - 10\tabcolsep) * \real{0.1667}}@{}}
\caption{\textbf{Domain Randomization Parameter Bounds and Hardware
Stress Envelopes.} Quantitative mapping of physics engine parameter
randomization distributions against operational performance and physical
damage risk. The table establishes the narrow viable corridor between
over-randomization (which induces over-conservative, low-gain
sluggishness) and under-randomization (which causes aggressive, brittle
control policies that induce dynamic shock, thermal saturation, or
statutory torque threshold
breaches).}\label{tbl-06-domain-randomization}\tabularnewline
\toprule\noalign{}
\begin{minipage}[b]{\linewidth}\raggedright
Parameter / Physical Modality
\end{minipage} & \begin{minipage}[b]{\linewidth}\raggedright
Nominal Physics Value
\end{minipage} & \begin{minipage}[b]{\linewidth}\raggedright
Recommended Randomization Bounds (\(\mathcal{P}(\Xi)\))
\end{minipage} & \begin{minipage}[b]{\linewidth}\raggedright
Over-Randomization Consequence (Performance Collapse)
\end{minipage} & \begin{minipage}[b]{\linewidth}\raggedright
Under-Randomization Consequence (Physical Damage Risk)
\end{minipage} & \begin{minipage}[b]{\linewidth}\raggedright
Hardware Interceptor \& Safety Budget
\end{minipage} \\
\midrule\noalign{}
\endfirsthead
\toprule\noalign{}
\begin{minipage}[b]{\linewidth}\raggedright
Parameter / Physical Modality
\end{minipage} & \begin{minipage}[b]{\linewidth}\raggedright
Nominal Physics Value
\end{minipage} & \begin{minipage}[b]{\linewidth}\raggedright
Recommended Randomization Bounds (\(\mathcal{P}(\Xi)\))
\end{minipage} & \begin{minipage}[b]{\linewidth}\raggedright
Over-Randomization Consequence (Performance Collapse)
\end{minipage} & \begin{minipage}[b]{\linewidth}\raggedright
Under-Randomization Consequence (Physical Damage Risk)
\end{minipage} & \begin{minipage}[b]{\linewidth}\raggedright
Hardware Interceptor \& Safety Budget
\end{minipage} \\
\midrule\noalign{}
\endhead
\bottomrule\noalign{}
\endlastfoot
\textbf{Surface Friction Coefficient (\(\mu\))} &
\(\mu_{\text{dry}} = 0.50\) (steel on aluminum) &
\(\mathcal{U}(0.30, 0.85)\) or
\(\text{LogNormal}(\mu{=}0.50, \sigma{=}0.25)\) & Policy converges to
hyper-conservative grasping with excessive clamping forces; velocity
throttles by \(>60\%\) to prevent slip & Instantaneous payload slip and
dropped workpieces upon slight surface contamination (oil film, dust,
oxidation); micro-impact damage to tool & Tactile slip-detection arrays,
dynamic grasp-force governors, compliant mechanical gripper suction
backup \\
\textbf{Payload Mass \& Link Inertia (\(m, \mathbf{I}\))} &
\(m_{\text{payload}} = 2.00\text{ kg}\);
\(I_{zz} = 0.015\text{ kg}\cdot\text{m}^2\) &
\(\mathcal{U}(1.60\text{ kg}, 2.40\text{ kg})\) (\(\pm 20\%\));
\(I_{zz} \in [0.80, 1.20] \times I_{\text{nom}}\) & Low feedback gains,
massive trajectory lag (\(>50\text{ mm}\) tracking offset), and failure
to track nominal velocity setpoints & Severe overshoot
(\(>25\text{ mm}\)) when handling uncalibrated payloads; excites
structural flexure, tripping joint torque limits
(\(\tau > \tau_{\text{peak}}\)) & In-flight payload mass estimator
(force-sensor telemetry), Cartesian tracking error tripwires
(\(\Delta x_{\text{max}} \le 10\text{ mm}\)), overcurrent breakers \\
\textbf{Actuator \& Transport Latency (\(\tau_{\text{lat}}\))} &
\(\tau_{\text{nom}} = 20.0\text{ ms}\) (\(1\) cycle at \(50\text{ Hz}\))
& Random queue delay
\(\tau \in \{20\text{ ms}, 40\text{ ms}, 60\text{ ms}\}\) with jitter
\(\delta \sim \mathcal{N}(0, 2\text{ ms}^2)\) & Destructive phase margin
collapse; policy learns ultra-slow deadbeat motions to avoid
instability, degrading throughput by \(>75\%\) & Closed-loop limit-cycle
oscillations on physical hardware; gantry instability, position
overshoot, and collision with workcell boundaries & Hardware watchdog
timer, cycle timestamp discrepancy monitor
(\(\Delta t > 30\text{ ms} \implies \text{safe hold}\)), phase-lead
compensator \\
\textbf{Motor Torque Constant (\(K_t, \eta\))} &
\(K_t = 0.120\text{ N}\cdot\text{m/A}\); efficiency \(\eta = 0.85\) &
\(K_t \in [0.105, 0.135]\text{ N}\cdot\text{m/A}\) (\(\pm 12\%\)
thermal/remanence); \(\eta \in [0.75, 0.95]\) & Policy commands
excessive current saturation at nominal states, causing motor coil
heating and power supply voltage sag & Inability to overcome static
friction when cold, or steady-state torque droop when hot
(\(R_{\text{Cu}}\) rise); stall at mechanical stops & Real-time motor
winding thermal estimator (\(I^2 R\)), driver bus voltage sag monitor
(\(V_{\text{bus}} \ge 42.0\text{ V}\)), hardware current limiter
(\(I \le 12\text{ A}\)) \\
\textbf{Sensor Noise \& Bias (\(b_{\text{sens}}, \sigma_n\))} & Bias
\(b = 0.0\text{ mm}\); noise std dev \(\sigma_n = 1.0\text{ mm}\) &
\(b \sim \mathcal{U}(-3.5\text{ mm}, +3.5\text{ mm})\);
\(\sigma_n \sim \text{Rayleigh}(1.5\text{ mm})\); dropout burst
\(P(\text{drop}) = 0.02\) & Policy ignores sensor feedback entirely,
falling back to blind, open-loop kinematic trajectories; fails fine
alignment & Sensor noise amplification in high-gain feedback;
high-frequency actuator jitter (\(>80\text{ Hz}\)), acoustic whine, and
rapid bearing wear & Outlier-rejection Kalman filter / median filter,
state covariance threshold monitor
(\(\text{Tr}(\mathbf{P}) \le P_{\text{thresh}}\)), multi-sensor
agreement check \\
\end{longtable}

Decomposing the gap into its constitutive physical channels yields two
concrete engineering artifacts, namely a supported deployment envelope
and an explicit residual mismatch list. The supported deployment
envelope defines the bounded domain within which physical hardware
matches simulation within validated tolerances. For the analytical
drive, this envelope specifies approach velocities below
\(0.25\text{ m/s}\), ambient illumination between \(300\text{ lux}\) and
\(600\text{ lux}\), and total loop latency below \(40\text{ ms}\). The
residual mismatch list itemizes every unclosed difference that remains
outside the training distribution, including the
\(1.5\times 10^{-3}\text{ rad/N}\cdot\text{m}\) drivetrain compliance,
the \(49.5\text{ ms}\) \(P_{99}\) latency tail, and the
\(0.40\text{ N}\) thermal sensor drift. Each entry on the residual list
defines an explicit requirement for runtime monitors and safety filters,
which must detect when physical execution drifts past the boundary where
the simulated policy remains valid.

\section{Fine-Tuning on the Real Machine}\label{sec-training-6-6}

When residual errors resist simulation tuning, fine-tuning directly on
physical hardware appears to be the direct path to closure. Instead of
constructing higher-fidelity friction models or identifying unmodeled
drivetrain dynamics through separate system identification routines,
gradient updates computed from physical rollouts adjust the policy
against the actual mechanics of the deployed machine. Adapting a policy
on the deployed machine closes an observed gap and charges physical
exposure for every sample it takes. The policy adjusts its parameters
against the true compliance of the chassis, the true contact dynamics of
the workpiece, and the actual latency profile of the bus. What makes
this trade exacting is that physical interaction is not free. In a
simulator, sampling states near the boundary of stability costs only
compute cycles. On hardware, every sample expends component life,
stresses mechanical interfaces, and risks driving the machine into
unrecoverable contact states.

\textbf{Real-machine fine-tuning debt} is the compounded operational
cost incurred when updating policy parameters on physical hardware,
where gradient steps solve localized contact errors but consume finite
mechanical fatigue life, induce thermal heating in motor windings, and
risk catastrophic forgetting of pre-verified nominal transit stability.

Consider an analytical example of a direct-drive assembly station
adapting an insertion policy to overcome an unmodeled \(1.5\text{ mm}\)
mounting offset. The policy executes at \(50\text{ Hz}\), corresponding
to a sample period of \(20\text{ ms}\). To adapt its final layer weights
to the contact dynamics of a hardened steel fixture, an on-robot
fine-tuning routine collects \(20{,}000\) interaction steps,
representing \(400\text{ s}\) of active contact and near-contact
operation. Because the policy must explore to improve, stochastic action
perturbations generate unpredictable force transients before low-level
feedback can intervene. If \(1.5\%\) of exploratory steps exceed the
nominal \(15.0\text{ N}\) seating force and strike the fixture at a
contact threshold of \(60.0\text{ N}\) before the torque limiters in the
nervous system clamp the command, the fine-tuning run injects \(300\)
impact shocks into the drivetrain.

In this analytical example, the consequences of those \(300\) impact
events can be calculated directly. At an approach velocity of
\(0.08\text{ m/s}\) with an effective moving mass of \(3.5\text{ kg}\),
each impact dissipates
\(\frac{1}{2} m v^2 = \frac{1}{2} \times 3.5\text{ kg} \times (0.08\text{ m/s})^2 = 1.12\times 10^{-2}\text{ J}\)
directly into the gear teeth and bearings. If the cycloidal drive on the
primary joint is rated for a fatigue threshold of \(5.0\times 10^4\)
shock cycles at that peak load, those seven minutes of policy
exploration consume \(300 / (5.0\times 10^4) = 0.60\%\) of the
mechanical life of the reducer. Thermal accumulation introduces an equal
constraint. High-frequency exploration dither elevates the
root-mean-square current through the motor coils from a nominal
\(1.8\text{ A}\) to \(3.6\text{ A}\). In a winding with an electrical
resistance of \(2.2\ \Omega\), Joule heating increases from
\(I^2 R = (1.8\text{ A})^2 \times 2.2\ \Omega = 7.13\text{ W}\) to
\((3.6\text{ A})^2 \times 2.2\ \Omega = 28.51\text{ W}\), quadrupling
thermal dissipation and bringing the motor core to its
\(80^\circ\text{C}\) thermal derating threshold within \(210\text{ s}\)
unless cooling intervals are budgeted into the training schedule.

Beyond mechanical wear, real-world adaptation undermines the statistical
verification established during pre-deployment testing. In simulation, a
policy can be subjected to \(1.0\times 10^7\) validation episodes across
randomized perturbations, establishing bounded tracking errors and tail
behavior out to \(P_{99.9}\). Once weight updates \(\Delta W\) are
applied on hardware, every statistical claim conditioned on the original
parameter vector \(W_0\) is void. The updated policy may resolve the
steady-state force error against the target fixture, yet the gradient
steps can simultaneously degrade free-space tracking accuracy, lower
phase margins during high-speed transit, or distort response profiles
during emergency stops. Because hardware throughput is bounded by
physical time, the machine cannot execute millions of physical trials to
re-establish those confidence bounds. The adapted policy enters
production with a narrower body of empirical evidence than the simulated
baseline it replaced.

Fine-tuning on the deployed machine therefore converts the policy from a
qualified, frozen software component into an evolving state. In an
engineered system, qualification records depend on strict configuration
control, where the exact binary or parameter hash corresponds to a
documented safety case. When online updates alter policy weights during
operation, that qualification record must be re-opened. The burden of
preventing failure shifts entirely from the learned policy to the
deterministic monitors and safety enforcers residing in the nervous
system. Those monitors must be parameterized to intercept the errors of
the baseline policy as well as the unpredictable intermediate behaviors
generated while the optimizer updates weights against physical feedback.

\section{How Training Choices Go Wrong}\label{sec-training-6-7}

A physical failure in a learned policy is rarely an arbitrary numerical
accident. When an articulated arm drops a workpiece or drives a tool
into a fixture, the command was generated by optimization over a
specific dataset, loss function, and observation space. Because the
mathematical objective selects for behaviors that minimize empirical
loss rather than behaviors that satisfy unstated physical requirements,
each training regime leaves an indelible structural footprint in the
resulting policy. An engineer can predict the failure signature of a
policy before the physical machine is assembled by tracing a
five-element inference chain. That chain begins with the training
exposure, which establishes a learned dependency within the network
weights. When a deployment condition violates that dependency, the
policy produces an observable failure signature in its output commands,
dictating the required detector in the nervous system.

The \textbf{policy vulnerability inference chain} is defined as the
causal five-stage framework
(\(\text{Exposure} \to \text{Dependency} \to \text{Violated Assumption} \to \text{Observable Signature} \to \text{Nervous Interceptor}\))
used to predict hardware destruction modes and synthesize protective
safety monitors before physical machine deployment.

In behavioral cloning, the training exposure consists of state
trajectories recorded under competent human demonstration or a classical
controller. Because competent demonstrators avoid unrecoverable errors,
the training dataset contains zero transitions originating in degraded,
stalled, or off-nominal states. The learned dependency is strict, as the
policy assumes that every input observation resides within the narrow
support of the demonstration distribution. When a mechanical jolt, a
thermal drift in a sensor, or a momentary torque limit shifts the
machine by a tracking error of \(15\text{ mm}\), that deployment
condition violates the distribution boundary. The resulting failure
signature is compounding covariate shift. The policy does not output
mathematical singularities or out-of-bounds numbers. Instead, it emits
continuous, kinematically plausible commands that steer the actuators
into regions of state space where the policy has never received training
signal. Within the brain, the command appears normal. If the nervous
system lacks an explicit distribution monitor or novelty filter at its
input interface, this failure mode cannot be detected until the machine
physically collides with an obstacle or exceeds an envelope limit.

In trial-based learning within simulation, the training exposure
consists of millions of synthetic interactions against a numerical
physics engine. The learned dependency is the set of mathematical
simplifications chosen to keep the simulation fast. When the simulator
models dry friction without velocity-dependent Stribeck curves, assumes
an unvarying communication latency of \(0\text{ ms}\), or calculates
contact forces through instantaneous penalty springs that allow
penetration without restitution, the policy exploits those numerical
artifacts to maximize its reward. During deployment, the machine
encounters real physical dynamics. In an analytical example with an
approach velocity of \(0.5\text{ m/s}\), a transport delay of
\(16\text{ ms}\) across the digital bus means the body travels
\((0.5\text{ m/s}) \times 0.016\text{ s} = 8.0\text{ mm}\) into the
physical boundary before the brain receives the first contact frame and
computes a deceleration command. In the simulator, where latency was set
to \(0\text{ ms}\) and contact was modeled as a soft penalty spring, the
policy learned to decelerate within that \(8.0\text{ mm}\) window
without bouncing. On physical hardware, the unyielding metal contact
stops the tool in less than \(1.0\text{ ms}\), driving transient contact
forces past \(450\text{ N}\) and tripping the actuator overcurrent
protection.\sidenote{\footnotesize Under ISO 10218 and ISO/TS 15066 (Power and Force
  Limiting), transient dynamic impact forces are capped between
  \(140\text{ N}\) and \(210\text{ N}\). Unmodeled simulation contact
  spikes exceeding \(450\text{ N}\) violate these statutory safety
  envelopes, requiring independent physical safety interlocks.} This
failure mode is detectable before deployment by evaluating the policy in
a simulation test environment with injected latency jitter, and at
runtime by tracking force derivatives and high-frequency command chatter
in the nervous system.

\protect\phantomsection\label{callout-incidentux2a-6.9}
\begin{fbxSimple}{callout-incident}{Landmark Engineering Incident:}{Tooling Collision from Simulation Contact and Latency Mismatch}
\phantomsection\label{callout-incident*-6.9}

\emph{System:} High-speed robotic pick-and-place assembly cell with
direct-drive Cartesian gantry, six-axis wrist force-torque sensor, and
end-to-end neural motion policy running at \(50\text{ Hz}\) over an
EtherCAT fieldbus.

\emph{Incident:} End-effector tooling collision during automated seating
of a machined aluminum connector, fracturing the precision ceramic
alignment guide, shearing the wrist load-cell mounting pins, and causing
a 6-hour assembly line shutdown.

\emph{Telemetry snapshot:} At \(t = 0.000\text{ s}\), the policy
initiates vertical descent from \(150\text{ mm}\) clearance at
\(v_z = 0.50\text{ m/s}\); at \(t = 0.280\text{ s}\), physical contact
occurs at \(z = 0.0\text{ mm}\); due to a \(16.7\text{ ms}\) camera
exposure time and an \(8.2\text{ ms}\) bus transmission delay, the first
contact frame reaches the policy at \(t = 0.305\text{ s}\)
(\(25\text{ ms}\) late); during this unobserved interval, the gantry
moves \(12.5\text{ mm}\) past the contact plane; at
\(t = 0.308\text{ s}\), peak normal force spikes to \(465\text{ N}\) at
a rate \(dF/dt > 3.8\times 10^4\text{ N/s}\), shearing mechanical load
pins before the drive controller trips on overcurrent at
\(t = 0.312\text{ s}\).

\emph{Root cause:} The rigid-body simulation engine modeled contact
using an idealized compliant penalty spring with zero transport delay,
allowing the policy to learn an unbraked plunge strategy that relied on
immediate synthetic force feedback to arrest momentum.

\emph{Architectural flaw:} Absence of an independent runtime
velocity-envelope governor and hardware force-derivative (\(dF/dt\))
clamp in the nervous system to restrict kinetic energy prior to
touch-down regardless of policy command outputs.

\emph{Provenance:} Precision manufacturing cell automated commissioning
incident log and post-mortem failure analysis report.

\end{fbxSimple}

Training through trial directly on physical hardware eliminates the
simulation gap, but it introduces an opposite structural dependency. The
training exposure consists of unconstrained physical interactions guided
by an engineered scalar reward. The learned dependency is the exact
mathematical formulation of that reward function. Engineers frequently
treat real dynamics as a natural safeguard against unrealistic
behaviors, assuming that physical execution prevents the policy from
finding shortcuts. The machine, however, discovers physically valid
strategies that satisfy the mathematical reward while violating unstated
operational constraints. If the reward function rewards rapid velocity
regulation toward a workpiece and penalizes final position error without
penalizing torque derivatives, the optimizer converges on aggressive
bang-bang control. The observable failure signature is unconstrained
mechanical stress and thermal overload. In an analytical example where a
joint motor with a winding resistance of \(1.5\ \Omega\) switches
between its current limits of \(+12\text{ A}\) and \(-12\text{ A}\) at
\(80\text{ Hz}\), the sustained root-mean-square current produces Joule
heating of
\(I^2 R = (12\text{ A})^2 \times 1.5\ \Omega = 216\text{ W}\), elevating
the stator core temperature past its \(90^\circ\text{C}\) thermal
derating limit within \(180\text{ s}\). The task-level reward monitor
records full success and detects no failure. Detection requires the
nervous system to monitor physical observables that were omitted from
the training reward, including motor coil temperatures, current spectra,
and high-frequency accelerometer energy.

When a base policy is fine-tuned on physical hardware to correct
localized errors, the training exposure is restricted to a narrow task
slice collected over a brief operating window. The learned dependency is
an over-representation of recent local transitions coupled with the
deliberate suppression of dangerous trajectories. When human operators
intervene or safety tripwires halt execution during training, those
intervention states are censored from the update buffer to avoid
reinforcing abort sequences. The observable failure signature is
regression and boundary instability. While the policy improves
steady-state force regulation against a specific production fixture,
gradient updates overwrite the parameter representations that maintained
stability during free-space transit or off-axis approaches. The policy
loses competence on previously verified nominal states, failing when
presented with lighting angles or approach vectors that were common in
the pretraining dataset but absent during the fine-tuning session.
Because the neural network assigns valid output magnitudes across all
inputs, this internal regression is undetectable by the policy itself.
Detecting it requires a pre-deployment regression test suite on frozen
hardware benchmarks to evaluate non-adapted operational modes before the
updated weights are granted authority over physical actuators.

Hybrid pipelines that combine simulation pretraining, demonstration
filtering, and real-world fine-tuning do not eliminate these structural
failure modes. Instead, they shift the dominant failure signature across
the lifecycle. A policy initialized from demonstrations and refined in
simulation shifts its primary risk from covariate drift to simulator
exploitation. When that policy is subsequently fine-tuned on hardware
for contact regulation, the updates suppress simulator exploitation at
the point of contact but induce forgetting in the upstream velocity
approach. The earlier failure mechanisms remain latent within the
network weights, ready to trigger whenever operating conditions push the
machine outside the fine-tuned envelope.

Predicting these signature failures establishes three concrete
engineering obligations for the system designer. First, every identified
training dependency requires a predeployment test obligation, consisting
of a verification test protocol designed to stress the exact boundary
where the training environment deviated from physical reality. Second,
where the physical failure manifests in measurable mechanical
quantities, the nervous system must incorporate a dedicated runtime
detector operating at the loop rate, whether monitoring tracking
deviation, force rate of change, or current spectra. Third, where
detection is physically or mathematically impossible, such as ungrounded
commands caused by covariate shift in pure cloning, the system
specification must declare an explicit unsupported operating envelope
and rely on deterministic safety interlocks to refuse policy authority.
Fulfilling these obligations requires that the engineering team maintain
an exact, immutable record of how the policy was constructed, detailing
the source data, the simulator dynamics, the reward terms, and every
modification applied across its training history.

\section{The Policy Manifest}\label{sec-training-6-8}

The training record establishes the causal link between how a policy was
constructed and how it behaves when granted authority over physical
actuators. Where Chapter~\ref{sec-nervous} defines the base engineering
notebook, including its immutable hashing, schema inheritance, and
versioning rules, this section specifies only the domain fields that
Chapter~\ref{sec-training} adds to that schema. A learned policy does
not fail generically. It fails along the exact seam lines where its
training regime, source demonstrations, or simulation environment
departed from physical reality. Capturing those seam lines requires
logging the sequence of optimization stages, the bounded support of the
training data, the simulator solver parameters, and every hardware
exposure incurred during real-world adaptation.

The \textbf{training provenance record} is defined as the immutable
notebook schema extension that registers optimization stage checkpoints,
empirical demonstration support manifolds, physics engine version
hashes, domain randomization intervals, solver omissions, and hardware
fine-tuning wear logs to predict physical deployment failure modes.

\begin{longtable}[]{@{}
  >{\raggedright\arraybackslash}p{(\linewidth - 6\tabcolsep) * \real{0.2500}}
  >{\raggedright\arraybackslash}p{(\linewidth - 6\tabcolsep) * \real{0.2500}}
  >{\raggedright\arraybackslash}p{(\linewidth - 6\tabcolsep) * \real{0.2500}}
  >{\raggedright\arraybackslash}p{(\linewidth - 6\tabcolsep) * \real{0.2500}}@{}}
\caption{\textbf{Physical Policy Training Provenance Record Schema.}
Immutable schema fields added to the base engineering notebook to record
policy synthesis stages, empirical demonstration support bounds, physics
engine numerical specifications, domain randomization intervals,
hardware fine-tuning wear logs, and declared operational
envelopes.}\label{tbl-06-training-schema}\tabularnewline
\toprule\noalign{}
\begin{minipage}[b]{\linewidth}\raggedright
Record Schema Field
\end{minipage} & \begin{minipage}[b]{\linewidth}\raggedright
Physical / Computational Data Modality
\end{minipage} & \begin{minipage}[b]{\linewidth}\raggedright
Provenance \& Immutability Guarantee
\end{minipage} & \begin{minipage}[b]{\linewidth}\raggedright
Downstream Verifier \& Usage
\end{minipage} \\
\midrule\noalign{}
\endfirsthead
\toprule\noalign{}
\begin{minipage}[b]{\linewidth}\raggedright
Record Schema Field
\end{minipage} & \begin{minipage}[b]{\linewidth}\raggedright
Physical / Computational Data Modality
\end{minipage} & \begin{minipage}[b]{\linewidth}\raggedright
Provenance \& Immutability Guarantee
\end{minipage} & \begin{minipage}[b]{\linewidth}\raggedright
Downstream Verifier \& Usage
\end{minipage} \\
\midrule\noalign{}
\endhead
\bottomrule\noalign{}
\endlastfoot
\textbf{\texttt{regime\_sequence}} & Array of optimization stages
(\(\text{BC} \to \text{SimRL} \to \text{Residual}\)), loss formulas,
reward weights, and stopping criteria hashes & Git commit SHA,
PyTorch/JAX checkpoint UUIDs, parameter hash \(H(\theta_k)\) & Tracing
which regime introduced latent failure modes or high-gain chatter \\
\textbf{\texttt{data\_support\_ledger}} & Ingested demonstration
datasets (\(\mathcal{D}_k\)), collecting policies \(\pi_{\text{col}}\),
sensor calibration transforms, and convex hull bounds & SHA-256 dataset
hash, ISO/IEC 5259 data sheet, empirical support bounding volume &
Parameterizing OOD novelty detectors and Mahalanobis support boundaries
in nervous system \\
\textbf{\texttt{negative\_event\_registry}} & Circular pre-trigger
buffers (\(2.0\text{ s}\)), human takeover flags, E-stop trips, and
divergence timestamps \(t_{\text{div}}\) & Monotonic microsecond
hardware timestamps, uncompressed sensor raw dumps & Isolating physical
failure boundaries and preventing corrupted drift rollouts from
training \\
\textbf{\texttt{simulation\_provenance}} & Physics engine build,
integrator method, timestep \(\Delta t\), contact solver penalty
parameters, and domain randomization ranges \(\mathcal{P}(\Xi)\) &
Physics engine binary hash, URDF/MJCF asset tree hash, declared
unmodeled mechanisms & Determining which dynamics were omitted
(backlash, compliance) and auditing paired-trace budgets \\
\textbf{\texttt{adaptation\_wear\_ledger}} & Serialized hardware IDs,
contact shock counts, motor current RMS thermal logs, modified layer
indices, and regression suite hash & Calibration certificate hashes,
hardware lifetime fatigue logs & Re-opening formal qualification cases
and preventing silent regression of nominal capabilities \\
\textbf{\texttt{declared\_odd\_boundary}} & Operational Design Domain
bounds (payload \(m \in [1, 3]\text{ kg}\),
\(v_{\text{approach}} \le 0.25\text{ m/s}\), temp
\(T \le 75^\circ\text{C}\)) & Formal statutory release gate bounds, UL
4600 / ISO 10218 safety envelope & Hard gating of autonomous policy
execution; trigger for deterministic fail-safe handoff \\
\end{longtable}

As detailed in the provenance record schema
(Table~\ref{tbl-06-training-schema}), every physical AI training
pipeline executes a sequence of distinct regimes where objectives, data
distributions, and stopping criteria change. The training record
registers each stage in this sequence, beginning with the source
checkpoint hash and detailing the mathematical formulation of the loss
function or reward vector. For each stage, the record documents the
stopping rule, whether governed by an epoch count, a validation loss
threshold, or a plateau in simulated return, alongside the engineering
rationale for advancing to the next regime. In an analytical example
where a manipulator is trained to approach a surface at
\(1.2\text{ m/s}\) and regulate contact force to \(8.0\text{ N}\),
pretraining on kinesthetic demonstrations might be terminated after
\(500\) epochs to prevent overfitting to a single human demonstrator. If
training then transitions to reinforcement learning in simulation to
expand angular coverage, the record logs that motivation and references
the exact initial weights ingested by the simulator. Without this stage
provenance, an engineer diagnosing an unstable velocity profile cannot
determine whether oscillations originated in the demonstration weights
or were introduced during subsequent simulation updates.

For every dataset ingested during imitation or offline optimization, the
record specifies the dataset version identifier, the data collection
policy, the environmental operating conditions, and the source of the
control commands. It records whether trajectories were generated by an
expert scripted controller, a human teleoperator, or rollouts from an
earlier policy checkpoint, alongside sensor calibration matrices and
physical workspace parameters present during logging. The record
explicitly quantifies the state-action support over which demonstrations
actually exist. If human demonstrations of a surface-wiping task contain
end-effector velocities between \(0.2\text{ m/s}\) and
\(0.8\text{ m/s}\) and contact angles within \(10^\circ\) of the surface
normal, the record marks the boundary of that support manifold. When the
deployed policy is presented with an approach angle of \(25^\circ\), the
record reveals that the network is operating on ungrounded extrapolation
rather than demonstrated competence.

Standard machine learning pipelines frequently discard aborted episodes,
operator vetoes, and safety interlock interventions to prevent
sub-optimal actions from entering the training loss. In physical
systems, filtering out these events discards the exact evidence needed
to identify where the policy required assistance. The training record
preserves every intervention, emergency stop assertion, and truncated
trajectory as structured negative metadata. If a human supervisor
intervened eighteen times across two hundred trial runs to prevent a
tool from striking an alignment fixture during touchdown, recording the
joint positions, velocities, and motor current spikes at the moment of
each override exposes the state space region where the policy lost
tracking stability. Retaining these aborted trajectories ensures the
ledger reflects the true failure boundaries encountered during training
rather than a curated subset of successful rollouts.

When policy weights are optimized in simulation, the training record
captures the physics engine build hash, the interface contract between
the policy and the simulator, the reward formulation, the numerical
integration method, and the solver timestep, such as a fixed
\(2.0\text{ ms}\) step at \(500\text{ Hz}\). It documents the domain
randomization distributions applied to link masses, surface friction
coefficients, actuator response latencies, and sensor noise profiles.
The record also logs known physical omissions, identifying the unmodeled
dynamics that were deliberately excluded from the simulation graph to
maintain training throughput. If a rigid-body simulator omits harmonic
drive compliance, structural deflection under a \(5.0\text{ kg}\)
payload, or thermal resistance changes in motor windings, the record
lists those omissions as unverified dynamic assumptions. When the
resulting policy exhibits limit-cycle oscillations during high-load
contact on the physical bench, the omission log directs the diagnosis
directly to joint compliance rather than ungrounded parameter drift.

Real-machine adaptation introduces mechanical wear and localized
parameter shifts that must be tracked across the hardware lifecycle. The
training record registers the hardware serial numbers, physical
configurations, and calibration offsets of the adaptation test bench,
along with the cumulative physical exposure measured in runtime hours,
contact cycles, or total work performed. It lists the exact parameter
subset modified during adaptation, distinguishing between full weight
updates, adapter layer tuning, and output-head fine-tuning, while
logging every anomaly or actuator limit reached during the session.
After adaptation updates are computed, the record captures the execution
results of a frozen regression test suite. If adapting a policy to
soften touchdown impact on a rigid surface increases tracking error
during the preceding free-space transit by \(14\text{ mm}\), the
regression ledger records that loss of nominal capability before the
modified weights are granted authority over physical actuators.

The training record separates empirical optimization results from the
declared operating envelope of the machine. High training reward or low
loss on an offline split does not demonstrate that the policy is safe to
execute across all environmental states. The record defines the claimed
operating envelope across payload masses, approach velocities, thermal
bounds, and surface geometries, while maintaining a separate ledger of
unobserved, excluded, or weakly supported conditions. If training data
includes surface approaches only under ambient illumination between
\(200\text{ lx}\) and \(500\text{ lx}\) with rigid flat workpieces, the
record declares curved surfaces, compliant substrates, and illumination
levels below \(100\text{ lx}\) as unsupported operational modes. This
separation prevents a convergent training metric from being mistaken for
an operational clearance across uncharacterized physical regimes.

Downstream engineering teams consume the training record throughout the
design, evaluation, and release lifecycle. Chapter~\ref{sec-evaluation}
uses the recorded support boundaries and simulation omissions to
construct targeted empirical evaluation suites that probe the specific
failure modes predicted by each training regime. Runtime safety monitors
and deterministic enforcement layers in Chapter~\ref{sec-enforcement}
use the recorded data support bounds to parameterize out-of-distribution
detectors and determine when policy authority must be curtailed.
Verification frameworks in Chapter~\ref{sec-verification} and the
multi-layer safety case in Chapter~\ref{sec-release} draw on the record
to establish the provenance of every parameter update and verify that
known simulation omissions are bounded by independent physical
interlocks. Maintaining this complete record ensures that when a policy
achieves high scores across its training benchmarks, engineers do not
mistake optimization convergence for physical reliability.

\section{Systems Stress-Tests}\label{sec-training-6-9}

A policy that is better by every training metric can be worse on the
physical machine. In physical systems, minimizing empirical loss across
an offline demonstration split or maximizing return inside a numerical
simulator evaluates compliance with the training environment rather than
competence on physical hardware. When an optimization objective rewards
behaviors that exploit data curation filters or simulator
approximations, the ranking of candidate policies established during
training inverts upon deployment. Evaluating a trained policy before
granting it authority over actuators requires stress-testing the causal
assumptions that connect training metrics to physical outcomes.

\begin{figure}[t]

\centering{

\includegraphics[width=1\linewidth,height=\textheight,keepaspectratio]{index_files/mediabag//Users/VJ/GitHub/MLSysBook-integrate-physical/physical/book/chapters/06-training/figures/fig06_policy_ranking_inversion.pdf}

}

\caption{\label{fig-06-policy-ranking-inversion}\textbf{Policy
Performance Ranking Inversion across Training and Physical Hardware
Outcomes.} Three analytical case studies demonstrating how optimization
gains in training invert upon closed-loop physical deployment. (1) Flow
Regulation: Sanitizing 50 operator recovery runs lowers validation loss
\(4\times\) (\(0.042\text{ V}^2 \to 0.011\text{ V}^2\)), but collapses
physical closed-loop recovery (\(100\% \to 0\%\)) when a supply pressure
drop occurs. (2) Surface Touchdown: An unbraked plunge achieves \(28\%\)
higher simulated return by exploiting soft penalty interpenetration, but
shatters physical load-cell flexures (\(68.0\text{ N}\) peak force
vs.~\(11.4\text{ N}\) for the braked baseline). (3) Regime Transition:
Replacing behavioral cloning with simulation RL achieves \(100\%\)
simulated task success, but commands \(65\text{ Hz}\) high-frequency
torque chatter that excites mechanical resonance, trips motor winding
thermal protection (\(95^\circ\text{C}\)), and reduces hardware
completion to \(12\%\).}

\end{figure}%

The following three analytical stress-tests evaluate machines that
regulate motion and touch surfaces, demonstrating how nominal training
gains convert into physical failure modes:

\subsection*{Stress-Test 1: Recovery Curation and Policy Reversal in
Flow Regulation (State Support
Truncation)}\label{stress-test-1-recovery-curation-and-policy-reversal-in-flow-regulation-state-support-truncation}
\addcontentsline{toc}{subsection}{Stress-Test 1: Recovery Curation and
Policy Reversal in Flow Regulation (State Support Truncation)}

\emph{The scenario.} Two behavioral-cloning policies are trained to
regulate fluid flow via a proportional valve at \(50\text{ Hz}\).
Dataset A contains 500 demonstrated runs including 50 recovery maneuvers
where an operator corrected for transient supply pressure drops to
\(120\text{ kPa}\). Dataset B filters out the 50 recovery runs to
maximize smoothness, retaining 450 nominal runs near \(200\text{ kPa}\).
Policy A achieves an offline validation loss of
\(\text{MSE} = 0.042\text{ (V)}^2\) across the full distribution, while
Policy B achieves \(\text{MSE} = 0.011\text{ (V)}^2\) on its held-out
split, appearing nearly four times more accurate.

\emph{The systems dilemma.} Under a physical supply pressure transient
dropping to \(130\text{ kPa}\), Policy A restores nominal flow within
\(140\text{ ms}\) (\(\Delta P = 15\text{ kPa}\)). Policy B enters an
unrecoverable diverging oscillation that breaches the reservoir overflow
limit within \(800\text{ ms}\). Sanitizing recovery episodes
artificially lowers offline validation error by eliminating
high-variance, multi-modal demonstration segments, blinding the policy
to the off-nominal states its own compounding errors will generate.

\emph{The systems verdict: State Support Invariant.} Offline validation
error measures agreement with recorded history rather than closed-loop
stability. Pruning recovery trajectories improves the training score
while removing the negative feedback manifold required to survive
physical disturbances.

\subsection*{Stress-Test 2: Contact Model Exploitation in Surface
Touchdown (The Simulation Energy
Sink)}\label{stress-test-2-contact-model-exploitation-in-surface-touchdown-the-simulation-energy-sink}
\addcontentsline{toc}{subsection}{Stress-Test 2: Contact Model
Exploitation in Surface Touchdown (The Simulation Energy Sink)}

\emph{The scenario.} A direct-drive vertical actuator
(\(m = 2.5\text{ kg}\)) approaches a rigid fixture at
\(0.40\text{ m/s}\) to establish a steady normal force
\(F_N = 10.0\text{ N}\). In simulation, Policy 2 achieves \(28\%\)
higher return than Policy 1, reaching target force in \(12\text{ ms}\)
(versus \(45\text{ ms}\) for Policy 1) by plunging into the surface at
\(0.35\text{ m/s}\), exploiting an idealized linear spring-damper
contact solver that arrests momentum in a single integration step
(\(2.0\text{ ms}\)).

\emph{The systems dilemma.} On physical hardware, structural compliance
(\(k = 8.0\times 10^4\text{ N/m}\)) and \(12\text{ ms}\) sensory lag
prevent instantaneous force generation. When Policy 2 strikes the
surface at \(0.35\text{ m/s}\), physical contact force spikes to
\(68.0\text{ N}\) before the first sensor measurement arrives,
shattering the load-cell flexure. Policy 1 decelerates to
\(0.06\text{ m/s}\) prior to touchdown, seating safely in
\(52\text{ ms}\) with a peak force of \(11.4\text{ N}\).

\emph{The systems verdict: Simulator Exploitation Law.} Optimization
algorithms treat simulator simplifications as physical affordances. When
a numerical contact solver permits instantaneous momentum transfer,
learned policies develop high-velocity impact strategies that destroy
physical hardware.

\subsection*{Stress-Test 3: Failure Migration across Regime Transitions
(Data Debt vs Modeling
Debt)}\label{stress-test-3-failure-migration-across-regime-transitions-data-debt-vs-modeling-debt}
\addcontentsline{toc}{subsection}{Stress-Test 3: Failure Migration
across Regime Transitions (Data Debt vs Modeling Debt)}

\emph{The scenario.} An assembly arm executes a \(200\text{ mm}\)
free-space transit at \(v = 0.50\text{ m/s}\) before preloading a gasket
to \(12.0\text{ N}\). Behavioral cloning (Policy A) exhibits \(34\%\)
task failure due to compounding lateral drift. Replacing cloning with
rigid-body reinforcement learning in simulation (Policy B) achieves
\(100\%\) simulated success and eliminates free-space drift across 100
physical trials.

\emph{The systems dilemma.} To seat the seal rapidly, Policy B learned
in simulation to apply high-frequency torque chatter at
\(65\text{ Hz}\), exploiting the simulator's massless rotor assumption.
On hardware, \(65\text{ Hz}\) chatter excites drivetrain resonance,
heating motor windings from \(28^{\circ}\text{C}\) to the
\(95^{\circ}\text{C}\) thermal limit in \(45\text{ s}\) and collapsing
physical completion rate to \(12\%\).

\emph{The systems verdict: Failure Migration Law.} Changing a training
regime does not eliminate engineering debt; it converts one failure
currency into another. Shifting from human demonstrations to simulation
replaces data-support limits with unmodeled physical dynamics limits.

\section{Summary}\label{sec-training-6-10}

Behavioral cloning purchases operational capability at minimal compute
cost by mapping recorded sensory observations directly to demonstrated
actions. It bypasses the combinatorial complexity of reward design and
avoids running dangerous trial trajectories on physical mechanisms. The
cost of this simplicity is an absolute dependence on the state-space
support of the demonstration dataset. Because the training objective
penalizes deviations only along demonstrated paths, the resulting policy
learns nothing about states outside that narrow manifold. When an
external disturbance, a worn bearing, or a dropped camera frame pushes
the machine into an unrecorded configuration, the policy encounters
covariate shift. In these unrepresented states, the network emits
ungrounded actuator commands with high numerical confidence, compounding
the tracking error with each control cycle until the physical system
collides with an obstacle or exceeds a joint limit.

Trial-based reinforcement learning purchases closed-loop recovery
experience. By exposing the agent to errors during training, it produces
policies that actively correct perturbations and steer the mechanism
back toward the task manifold when disturbances occur. This resilience
is paid for through mechanical wear, clock time, exploration risk, and
reward specification debt. On physical hardware, accumulating the
millions of transitions required for policy convergence subjects
gearboxes, ball screws, and motor windings to repetitive mechanical and
thermal stress. Every exploration step risks driving an end-effector
into a hard fixture or drawing stall current long enough to trip a
thermal breaker. When developers construct a scalar reward function to
guide this search, any misalignment between the mathematical objective
and the physical design intent invites policy exploitation, producing
aggressive, chatter-inducing torque profiles that satisfy the numerical
metric while destroying the physical transmission.

Simulation decouples experience generation from the constraints of
real-time physics and mechanical wear, allowing billions of synthetic
transitions to accumulate across parallel compute clusters at zero
hardware risk. What simulation delivers in throughput, it charges in
modeling debt. Every physics engine represents a versioned mathematical
compromise containing an explicit set of modeled kinematic links, an
explicit set of omitted physical phenomena, and a collection of
numerical approximations. Effects such as motor back-electromotive
force, rotor inertia, contact friction hysteresis, structural
compliance, and sensor thermal noise are routinely simplified or
neglected to sustain high simulation frame rates. The learned policy
optimizes against the exact numerical artifact presented by the solver.
If the simulator treats a joint as a zero-inertia actuator capable of
instantaneous torque reversal, the policy will exploit that unphysical
affordance, leaving the machine vulnerable to resonance, motor
overheating, and mechanical failure when transferred to real hardware.

Physical fine-tuning attempts to reconcile these discrepancies by taking
gradient steps on small batches of rollouts collected directly on the
machine. This step can resolve local tracking errors, adjust for
unmodeled contact compliance at an assembly jig, or tune gain responses
to match physical surface friction. Yet localized fine-tuning carries
its own systemic tax. Updating network weights against a narrow
distribution of hardware rollouts risks eroding the broad recovery
behaviors acquired during large-scale pretraining or simulation.
Overwriting parameters to solve a specific contact mismatch can compress
the effective support of the policy, inducing catastrophic forgetting
across unvisited regions of the state space and degrading the machine's
ability to handle nominal kinematic variations or payload changes.

The durable output of the training process is therefore not simply a
binary weights file, but the complete training record. This artifact
defines the physical and operational boundaries of the learned policy by
recording its provenance. For cloned models, the training record
documents the demonstrated state-space envelope, operator skill
variance, and sensor calibration baselines. For trial-based and
simulated models, it documents the reward formulations, exploration
boundaries, simulator version identifiers, integrator time-steps,
contact solver tolerances, and explicitly unmodeled physical effects.
Without this provenance, a systems engineer cannot predict the failure
modes of the policy, budget appropriate thermal and current margins in
the nervous system, or construct meaningful stress tests to probe the
boundaries of safe authority.

\protect\phantomsection\label{callout-takeawaysux2a-6.10}
\begin{fbxSimple}{callout-takeaways}{Key Chapter Takeaways}{}
\phantomsection\label{callout-takeaways*-6.10}

\begin{itemize}
\tightlist
\item
  \textbf{Optimization Inversion (Good in Training, Bad on Hardware):}
  Offline loss minimization and simulated return evaluate numerical
  compliance rather than physical stability. Policies that score highest
  on sanitized datasets often fail catastrophic first-contact trials due
  to unmodeled contact compliance or omitted rotor inertia.
\item
  \textbf{Behavioral Cloning Compounds Off-Support Drift:} Supervised
  imitation penalizes deviations only along demonstrated trajectories.
  Small tracking errors push the body into unvisited states where the
  policy emits ungrounded high-confidence commands, requiring continuous
  \(\mathcal{C}^2\) action chunking or recovery demonstrations.
\item
  \textbf{Simulators Charge Modeling Debt for Scaled Throughput:}
  Physics engines accelerate throughput by simplifying contact friction
  hysteresis, motor back-EMF, and structural compliance. Learned
  policies aggressively exploit these numerical simplifications,
  commanding unphysical torque chatter that overheats real actuators.
\item
  \textbf{Hardware Fine-Tuning Risks Catastrophic Forgetting:} Taking
  gradient steps on narrow physical hardware batches can fix local
  contact errors while silently destroying global recovery manifolds and
  inflating tracking error across nominal transit regimes.
\item
  \textbf{The Training Record is the True Engineering Output:} A
  checkpoint file without provenance is an uncharacterizable liability.
  Dependable physical AI requires an immutable training record detailing
  dataset support boundaries, simulator solver assumptions, and declared
  operational envelopes.
\end{itemize}

\end{fbxSimple}

Chapter~\ref{sec-evaluation} receives this policy alongside its
inherited modeling assumptions, unrepresented state volumes, and
actuator demands. A training record documenting zero regression loss, a
simulated success rate of \(100\%\), or a high cumulative reward in a
software sandbox establishes only that the optimization algorithm
converged under its training conditions. It does not establish that the
policy will move real mass without exceeding motor thermal limits,
violating timing deadlines, or causing mechanical damage. Proving
whether the learned policy can be trusted with physical authority
requires empirical evaluation on hardware under strict physical and
temporal constraints, testing the machine against the very failures that
its training regime chose to defer.

\chapter{\texorpdfstring{Closed-Loop
Evaluation\chaptersubtitle{What Can Be Established About a Policy, and What Cannot}}{Closed-Loop Evaluation}}\label{sec-evaluation}

\begin{figure}[H]

\centering{

\includegraphics[width=1\linewidth,height=\textheight,keepaspectratio]{index_files/mediabag//Users/VJ/GitHub/MLSysBook-integrate-physical/physical/book/chapters/07-evaluation/figures/fig_locator.pdf}

}

\caption{\label{fig-locator}\textbf{Systems Organ Focus: Brain.} Exact
one-sided \(95\%\) lower confidence bounds on true success probability
as a function of zero-failure physical trial count \(n\), illustrating
how statistical confidence scales from \(86.1\%\) at \(n=20\) to
\(99.9\%\) at \(n=2{,}995\) and exposing the wide tail risk that remains
under small testing budgets.}

\end{figure}%

\emph{Why does passing twenty consecutive physical trials without
failure tell you surprisingly little about whether the machine is safe
to deploy?}

Standard machine learning workflows evaluate models by measuring average
accuracy or loss across static, held-out test splits. In physical
systems, open-loop prediction metrics are almost completely uncorrelated
with closed-loop survival. A policy with negligible offline tracking
error can diverge catastrophically in physical execution because its own
actions shift the observation stream into unfamiliar territory. Physical
AI demands closed-loop evaluation methodologies grounded in
non-asymptotic statistics and extreme-value tail analysis. Verifying an
embodied agent requires measuring performance across worst-case
operating envelope boundaries, calculating exact Clopper-Pearson
binomial coverage bounds (\(86.1\%\) at \(n=20\)), and enforcing
pre-registered stopping oracles to expose heavy-tailed physical risks.

\newpage{}

\protect\phantomsection\label{callout-objectiveux2a-7.1}
\begin{fbxSimple}{callout-objective}{Learning Objectives}{}
\phantomsection\label{callout-objective*-7.1}

After studying this chapter, you will be able to:

\begin{enumerate}
\def\labelenumi{\arabic{enumi}.}
\tightlist
\item
  State what a success rate on \(N\) physical trials actually supports,
  with its exact non-asymptotic confidence interval.
\item
  Explain why an open-loop dataset metric misleads in a specific,
  predictable direction in closed-loop physical execution.
\item
  Compute the exact physical trial count and wall-clock testing budget
  required to certify a target reliability claim under zero failures.
\item
  Apply Extreme Value Theory (EVT) and the Butler \& Finelli
  infeasibility bound to identify physical tail hazards that empirical
  hours cannot cover.
\item
  Construct an immutable physical evaluation record and formalize a
  pre-registered stopping oracle to prevent benchmark gaming and
  censored trial fallacies.
\end{enumerate}

\emph{Core chapter takeaway.} No pre-deployment evaluation can establish
deployment behavior, because the policy determines the distribution it
will meet. And twenty clean trials support a far weaker claim than the
number feels like it does.

\end{fbxSimple}

\section{Closed-Loop Survival}\label{sec-evaluation-7-1}

\lettrine{A}{s mapped}\lettrine{i}{n} the systems architecture locator
(Figure~\ref{fig-locator}), the evaluation framework sits within the
Brain organ, establishing what can and cannot be mathematically and
empirically established about a learned policy before granting it
authority over physical actuators. Consider an autonomous mobile robot
that completes twenty consecutive navigation runs across a test
facility, traveling at \(1.5\text{ m/s}\) along a cluttered corridor
without striking a wall or requiring human intervention. In a robotics
laboratory, that result is frequently celebrated as a perfect
demonstration: twenty out of twenty runs passed, yielding an observed
empirical success rate of \(\hat{p} = 100\%\). In classical machine
learning, where held-out test splits contain thousands of static images,
an empirical score of \(100\%\) suggests near-perfection. In a physical
machine, that conclusion is an illusion created by small-sample
statistics.

To determine what twenty clean trials actually prove, we model the
sequence of physical trials as independent and identically distributed
Bernoulli experiments governed by an unknown true operational success
probability \(p\). As derived in the margin, setting the binomial
likelihood of twenty consecutive successes equal to our statistical
significance threshold (\(\alpha = 0.05\)) reveals that the minimum
success probability that would produce twenty clean runs at least
\(5\%\) of the time is only \(86.1\%\). This calculation---the exact
Clopper-Pearson lower bound---reveals the stark evidentiary deficit of
small testing budgets.\sidenote{\footnotesize For zero-failure testing, standard Wald
  confidence intervals degenerate to zero width. The exact
  Clopper-Pearson interval yields
  \(p_{\text{upper}} = 1 - \alpha^{1/n}\), preventing asymptotic
  Gaussian approximations from understating catastrophic risk in
  rare-event tail distributions.} If the true operational reliability of
the vehicle were \(86.1\%\), the probability of observing twenty
straight successes would be exactly \(5\%\). If the true reliability
were below \(86.1\%\), the probability of observing twenty clean runs
would fall below \(5\%\), allowing us to reject the hypothesis
\(p < 0.861\) at the \(95\%\) confidence level.

\protect\phantomsection\label{callout-definitionux2a-7.2}
\begin{fbxSimple}{callout-definition}{Definition:}{Clopper-Pearson Non-Asymptotic Bound}
\phantomsection\label{callout-definition*-7.2}

The \textbf{Clopper-Pearson non-asymptotic bound} is an exact
statistical lower confidence limit on true operational success
probability derived from the inverted cumulative binomial distribution
(evaluating to \(p_{\text{lower}} = \alpha^{1/n}\) for zero failures
across \(n\) independent trials at confidence level \(1-\alpha\)),
establishing the mathematical ceiling of empirical testing without
assuming asymptotic normality.

\end{fbxSimple}

Crucially, twenty successes in twenty trials does not prove that the
robot is \(100\%\) reliable, nor even \(95\%\) reliable. It proves only
that the true success probability is at least \(86.1\%\) at \(95\%\)
confidence. A physical system that fails once in every eight attempts
(\(p = 87.5\%\)) will successfully complete twenty consecutive test runs
more than \(6.9\%\) of the time (\(0.875^{20} \approx 0.0692\)). A
flawless score in twenty runs is entirely consistent with a physical
machine that suffers a catastrophic boundary violation nearly once every
seven deployments.

\protect\phantomsection\label{callout-ruleux2a-7.3}
\begin{fbxSimple}{callout-rule}{Architect’s Rule of Thumb:}{Twenty Clean Runs}
\phantomsection\label{callout-rule*-7.3}

Twenty consecutive successes in a laboratory test cell establish an
\(86.1\%\) lower confidence bound, not a \(100\%\) reliability
certificate. A flawless score in twenty runs is mathematically
consistent with a physical system that fails once in every seven
deployments.

\end{fbxSimple}

\protect\phantomsection\label{callout-mathux2a-7.4}
\begin{fbxSimple}{callout-math}{Fermi Math:}{Derivation of the 20-Run Lower Bound}
\phantomsection\label{callout-math*-7.4}

\emph{Problem:} An autonomous robot completes \(n = 20\) physical
navigation runs with zero collisions or human interventions
(\(k = 20\)). What is the exact one-sided \(95\%\) lower confidence
bound on its true operational success probability \(p\)?

\emph{Mathematical Derivation:}

\begin{enumerate}
\def\labelenumi{\arabic{enumi}.}
\item
  Under independent Bernoulli trials, the probability of \(n\)
  consecutive successes is \(P(k = n) = p^n\).
\item
  To find the lower bound \(p_{\text{lower}}\) at confidence level
  \(1 - \alpha = 0.95\) (\(\alpha = 0.05\)), solve:
  \[p_{\text{lower}}^n = \alpha \implies p_{\text{lower}} = \alpha^{1/n}\]
\item
  Substitute \(n = 20\) and \(\alpha = 0.05\):
  \[\ln(p_{\text{lower}}) = \frac{\ln(0.05)}{20} = \frac{-2.99573}{20} = -0.149787\]
  \[p_{\text{lower}} = e^{-0.149787} \approx 0.86089 \implies 86.09\%\]
\item
  \emph{Implied Maximum Failure Risk:}
  \[\text{Failure Probability} \le 1 - p_{\text{lower}} = 1 - 0.8609 = 0.1391 \implies 13.91\%\]
  \[\text{Mean Failure Interval} = \frac{1}{0.1391} \approx 7.19 \text{ trials}\]
\end{enumerate}

\emph{Conclusion:} The observation ``20 out of 20 passed'' permits the
machine to fail once every \(7.2\) runs under sustained operation.

\end{fbxSimple}

This gap separates an empirical record from a statistical population
claim. The statement that twenty out of twenty trials passed is an
immutable historical observation. It describes what occurred on a
particular physical chassis on a specific surface, under a fixed set of
lighting conditions, with a freshly charged battery pack, on a single
afternoon. The statement that the policy succeeds with probability at
least \(p\) is a forward-looking claim about future performance over an
operational design domain. Translating the historical record into a
forward-looking guarantee requires assuming that the test environment
matches the operational domain, that physical dynamics do not drift, and
that twenty trials provide enough statistical mass to characterize rare
failure modes.

An evaluation measurement never characterizes a learned policy in
isolation. Every claim about policy performance depends jointly on five
concrete elements: the specific compiled model weights (\(\pi_\theta\)),
the individual physical machine and its mechanical wear (\(m\)), the
distribution of environmental conditions (\(e\)), the formal outcome
rule that defines what constitutes success (\(\phi\)), and the exact
sampling protocol used to initialize each run
(\(s_0 \sim \mathcal{D}\)). Altering the ambient room temperature,
substituting a tire with slightly lower friction, or shifting the
initial yaw angle by two degrees changes the distribution of state
transitions. Without explicitly conditioning on all five elements, an
observed success rate is an unmoored number.

This principle holds across all three canonical Physical AI archetypes:

\begin{itemize}
\tightlist
\item
  \emph{Mobile systems (e.g., wheeled rovers, quadrotors, autonomous
  vehicles).} A twenty-run obstacle avoidance test on dry concrete at
  \(20^\circ\text{C}\) provides zero statistical license for operations
  on wet epoxy flooring (\(\Delta \mu = -0.55\)) or under direct solar
  glare (\(10^5\text{ lx}\)) that saturates optical cameras.
\item
  \emph{Manipulators (e.g., 6-DoF/7-DoF robotic arms, insertion
  mechanisms).} Twenty flawless peg-in-hole insertions on clean steel
  components reveal nothing about performance when thermal expansion
  shifts link lengths by \(0.15\text{ mm}\) or when surface burrs
  increase insertion friction.
\item
  \emph{Process \& energy systems (e.g., thermal management loops,
  chemical reactors, power grid inverters).} Twenty successful
  temperature regulation cycles under nominal electrical loads provide
  no guarantee against thermal runaway when ambient cooling capacity
  degrades by \(20\%\) during peak summer ambient temperatures.
\end{itemize}

Pre-deployment evaluation encounters three structural limits when
applied to machines that interact with the physical world. First, static
or open-loop evaluations on recorded datasets fail to reflect feedback.
A policy executing in the physical loop alters its own state trajectory,
where minor tracking errors accumulate over time and push the machine
into unmodeled regions of state space where the learned policy behaves
unpredictably. Second, finite physical testing budgets leave wide
statistical bounds. Demonstrating that a machine satisfies a high
reliability target, such as a failure rate below one in ten thousand
runs (\(P_{99.99}\)), requires tens of thousands of physical trials that
are impractical to execute before release. Third, deployment changes the
encountered distribution. Because a learned policy actively chooses
which states the machine visits, the policy itself shapes the
operational distribution it encounters, creating closed-loop
distribution shift.

The \textbf{predeployment evaluation boundary} is the formal statistical
and physical limitation governing empirical verification, establishing
that an observed sample success rate \(\hat{p}\) on finite trials
(\(n\)) supports only a conditional claim over explicitly sampled
environmental distributions (\(c = (s_0, e, m, \pi_\theta, \kappa)\))
and non-asymptotic confidence bounds
(\(p_{\text{lower}} = \alpha^{1/n}\)), leaving unvisited dynamical
states and endogenous closed-loop feedback uncertified prior to
deployment.

These constraints do not diminish the necessity of pre-deployment
evaluation, but they bound what it can establish. Evaluation before
deployment provides valuable negative evidence and supports well-defined
conditional claims. It systematically surfaces kinematic limits,
identifies perception failures caused by sensor noise, reveals numerical
instabilities in the inference pipeline, and validates behavior within
the specific envelopes covered by the test protocol. Evaluation
demonstrates that the policy operates correctly across tested instances,
but it cannot prove that the policy will remain safe across unvisited
states.

Because testing before deployment cannot verify safety across the
unbounded variety of the physical world, the limits of evaluation define
an architectural requirement for the machine. What cannot be proved
before deployment must be observed and constrained during operation. The
nervous system cannot treat the learned policy as a guaranteed
controller based on an empirical pass rate. Instead, it must
continuously monitor physical states, observe system invariants, and
enforce safety boundaries in real time, retaining the authority to
reject actions whenever the machine approaches the limits of its
physical operating envelope.

\section{What You Are Actually Measuring}\label{sec-evaluation-7-2}

An evaluation benchmark routinely reports a single summary statistic,
such as an eighty-five percent success rate on an obstacle avoidance
task. In machine learning practice, that number is often treated as an
intrinsic score belonging to the learned model weights, analogous to
top-one accuracy on a static image classification dataset. In a physical
machine, that interpretation is invalid. A learned policy has no
independent physical performance. What the benchmark measured was the
closed-loop execution of an entire physical system across a specific
sequence of days and environmental states. Success rate is an empirical
property of the coupled system, comprising the policy, the mechanical
plant, the onboard computational hardware, the operating environment,
and the trial protocol.

To evaluate a policy, an engineering team executes a sequence of
discrete physical runs. A single trial is fully determined only when
every variable that can alter the physical state trajectory is fixed or
explicitly sampled. Formally, let a single trial condition be defined by
a tuple \(c = (s_0, e, m, \pi_\theta, \kappa)\), drawn from an
operational testing distribution \(\mathcal{D}\). The vector \(s_0\)
represents the initial state of the machine, including position, linear
velocity, angular orientation, and actuator positions. The environment
vector \(e\) captures external variables such as floor friction, ambient
illuminance, obstacle geometry, and ambient air temperature. The machine
state \(m\) captures physical and hardware properties, including battery
state of charge, payload mass, tire wear, mechanical backlash in the
drivetrain, and sensor calibration parameters. The computational
configuration \(\pi_\theta\) encompasses the frozen policy weights
\(\theta\), inference quantization, scheduling deadlines, and software
versions. Finally, the termination protocol
\(\kappa = (T_{\max}, g, \phi)\) defines the execution time limit
\(T_{\max}\), the spatial goal region \(g\), and the binary evaluation
function \(\phi\) that maps the resulting trajectory to a terminal
outcome \(Y \in \{0, 1\}\).

The true success probability of the system under this evaluation
protocol is the expectation of the binary outcome over the joint
distribution of trial conditions:
\[p(\mathcal{D}) = \mathbb{E}_{c \sim \mathcal{D}}\left[ \phi(\text{Traj}(c)) \right] = \int \phi(\text{Traj}(c)) \, p(c) \, dc\]
In this formulation, \(\text{Traj}(c)\) is the continuous physical
trajectory produced when the policy \(\pi_\theta\) executes within the
physical environment and hardware state defined by \(c\). This
expectation demonstrates why \(p(\mathcal{D})\) cannot be treated as a
measurement of \(\theta\) alone. The model weights \(\theta\) constitute
only one parameter inside the condition vector \(c\). If the
distribution over floor friction \(p(e)\) shifts, or if the distribution
over initial battery charge \(p(m)\) changes, the value of the integral
\(p(\mathcal{D})\) changes while \(\theta\) remains identical. The
sampling protocol is not an external apparatus measuring an isolated
model; it is part of the mathematical quantity being estimated.

The physical mechanisms that alter \(c\) between trials operate
independently of the policy:

\begin{itemize}
\tightlist
\item
  In \emph{Mobile systems.}, a lithium-ion battery pack at ninety
  percent state of charge provides higher terminal voltage to the motor
  drivers than a depleted pack at twenty percent state of charge. As
  terminal voltage drops under high load, the maximum achievable motor
  torque decreases according to the electromechanical torque-speed
  curve. If a trajectory demands a peak acceleration of
  \(1.2\text{ m/s}^2\) to clear a moving obstacle, the robot clears the
  obstacle when the battery is fully charged, but understeers and
  collides when the battery is depleted. Similarly, ambient sunlight
  shifting across an afternoon changes camera exposure times from
  \(2\text{ ms}\) to \(20\text{ ms}\). This tenfold increase introduces
  motion blur during turning maneuvers, corrupting feature extraction in
  the visual backbone and causing state estimation drift. Floor dust
  accumulation lowers the static friction coefficient from
  \(\mu = 0.80\) to \(\mu = 0.55\), extending the stopping distance from
  an initial speed of \(2.0\text{ m/s}\) from \(0.26\text{ m}\) to
  \(0.37\text{ m}\).
\item
  In \emph{Manipulator systems.}, continuous high-speed pick-and-place
  cycles generate Joule heating in harmonic drive gearboxes, heating
  lubricant and altering friction damping coefficients by over \(25\%\).
  As structural linkages warm, thermal expansion shifts
  tool-center-point (TCP) kinematic calibration by
  \(0.15\text{–}0.30\text{ mm}\). A policy trained to insert a peg into
  a chamfered hole with \(0.20\text{ mm}\) clearance succeeds during
  early morning cold starts but jams repeatedly in the afternoon as
  thermal growth exhausts kinematic tolerance margins.
\item
  In \emph{Process \& energy systems.}, ambient temperature swings alter
  the thermal gradient across cooling heat exchangers. A neural policy
  regulating chiller pump speeds that maintains coolant temperature
  below \(45^\circ\text{C}\) in a \(20^\circ\text{C}\) laboratory
  environment fails when ambient air hits \(38^\circ\text{C}\),
  triggering high-pressure compressor trips despite executing identical
  control laws.
\end{itemize}

In all these cases, the policy weights and input tensor shapes are
identical, yet the physical trajectory transitions from success to
failure.

Non-model variation also enters through the reset procedure and operator
handling. When human operators manually reposition a machine between
trials, they introduce subtle spatial biases. If an operator
instinctively resets the robot facing open space after a difficult run,
the empirical distribution of initial heading angles shifts toward
favorable states. Conversely, an automated reset procedure that
repeatedly drives the machine along a fixed path induces cyclic heating
in the actuator windings, increasing electrical resistance and reducing
actuator response speed across successive trials. If an operator
intervenes to catch a falling machine and discards the run as an
unrecorded reset, the hardest failure modes are erased from the sample
history, artificially inflating the observed success rate.

Statistical evaluation requires distinguishing ordinary sampling
variation from a change in the underlying estimand. When \(N\) trials
are drawn from a fixed distribution \(\mathcal{D}\), the sample
proportion \(\hat{p} = \frac{1}{N} \sum_{i=1}^N Y_i\) fluctuates around
the true parameter \(p(\mathcal{D})\) due to random sampling noise. This
uncertainty can be bounded using standard statistical estimators.
However, if an evaluation changes the payload mass from
\(0.0\text{ kg}\) to \(5.0\text{ kg}\), or shifts testing from dry
concrete to wet asphalt, the experiment is no longer drawing samples
from \(\mathcal{D}\). It is sampling from an entirely different
distribution \(\mathcal{D}'\), defined by a new parameter
\(p(\mathcal{D}')\). Treating data collected under \(\mathcal{D}'\) as
an extension of the original sample pool conflates sampling variance
with structural bias, combining observations from two different physical
machines.

Because success probability is defined only with respect to a specific
joint distribution \(\mathcal{D}\), the intended operating envelope must
be declared before trials are sampled. In systems engineering, this
envelope is the \textbf{operational design domain} (ODD). If an
engineering team executes fifty physical runs across ad hoc trajectories
and reports a ninety percent pass rate, the denominator corresponds to
an undefined population. A measured success rate is statistically
meaningful only when the support and density of \(\mathcal{D}\) are
specified in advance, establishing explicit bounds on illuminance,
payload ranges, surface friction coefficients, and allowable initial
state offsets. Without an explicit prior specification of
\(\mathcal{D}\), an evaluation can be made arbitrarily successful by
restricting test runs to benign regions of the operating envelope.

Extracting valid confidence bounds from empirical trials requires three
operational conditions:

\begin{enumerate}
\def\labelenumi{\arabic{enumi}.}
\tightlist
\item
  \emph{Independent and Identically Distributed Sampling:} Trials must
  be independent realizations drawn from \(\mathcal{D}\). If thermal
  accumulation in motor windings or progressive mechanical wear degrades
  hardware efficiency across consecutive runs, trial outcomes become
  negatively correlated over time, violating independence.
\item
  \emph{Fixed Outcome Predicates:} The binary evaluation function
  \(\phi\) and the termination criteria \(\kappa\) must remain fixed
  throughout the campaign. Modifying the time horizon \(T_{\max}\) or
  adjusting the spatial tolerance of the goal region \(g\) midway
  through testing alters the definition of success.
\item
  \emph{Full-Denominator Accounting:} Every initiated run that aborts
  due to hardware faults, communication timeouts, or emergency stops
  must be recorded as a failure (\(Y = 0\)). Dropping aborted trials
  substitutes a measure of conditional execution on surviving hardware
  for the true operational reliability of the machine.
\end{enumerate}

When these conditions are strictly enforced, physical trials yield an
unbiased sample mean for the closed-loop system under the tested
distribution \(\mathcal{D}\). Yet physical testing remains constrained
by hardware wear, operator labor, and wall-clock execution time. To
bypass these constraints, evaluation workflows often turn to open-loop
evaluations on recorded datasets, calculating how accurately a candidate
policy reproduces logged expert actions offline. That substitution
assumes that a model that predicts demonstrated actions with high
statistical accuracy will exhibit equivalent competence when granted
real-time authority over a physical machine.

\section{Open Loop and Closed Loop}\label{sec-evaluation-7-3}

In an \textbf{open-loop evaluation}, a candidate policy \(\pi_\theta\)
is scored across a static dataset of recorded demonstrations
\(\mathcal{D}_{\mathrm{demo}} = \{(s_t^*, a_t^*)\}_{t=1}^T\). At each
discrete timestep \(t\), the policy receives a recorded demonstrator
state \(s_t^*\) and outputs a predicted action
\(\hat{a}_t = \pi_\theta(s_t^*)\). An error metric, such as mean squared
error or cross-entropy, measures the distance between \(\hat{a}_t\) and
the recorded expert action \(a_t^*\). At step \(t+1\), the evaluation
pipeline supplies the next recorded state \(s_{t+1}^*\) from the
dataset. The action \(\hat{a}_t\) produced by the policy is discarded
without executing on hardware, keeping the sequence of input states
independent of the policy's past outputs. In a \textbf{closed-loop
evaluation}, the policy is coupled directly to the state-transition
recurrence \(s_{t+1} = f(s_t, a_t)\). The action
\(a_t = \pi_\theta(s_t)\) commanded by the policy executes on the
physical plant, altering the state trajectory and making the policy
responsible for the observations it receives at subsequent timesteps.

Because open-loop evaluation tests actions only along the demonstrator's
state trajectory, a policy that scores low prediction error offline can
fail immediately when granted authority over physical motion
(\citeproc{ref-ross2011reduction}{Ross et al. 2011};
\citeproc{ref-pomerleau1989alvinn}{Pomerleau 1989}). Consider an
analytical example of a planar moving platform tracking a straight
corridor of length \(L = 10\text{ m}\) at a constant forward speed of
\(v = 1.0\text{ m/s}\). The nominal expert demonstration follows the
centerline with heading \(\theta^* = 0\text{ rad}\) and lateral offset
\(y^* = 0\text{ m}\) for a duration of \(10\text{ s}\). The policy runs
at a control frequency of \(10\text{ Hz}\), executing updates every
\(\Delta t = 0.1\text{ s}\) across \(100\) discrete steps. Suppose the
candidate policy has a systematic heading bias of
\(\epsilon_\theta = 1.0^\circ\) (\(0.01745\text{ rad}\)) on every step.

In an open-loop evaluation, the policy receives the demonstrator state
\(s_t^* = (x_t^*, 0, 0)\) at each step and predicts a forward velocity
with a \(1.0^\circ\) steering offset. The offline scoring pipeline
records a mean absolute angular error of \(1.0^\circ\), or a mean
squared error of
\(\text{MSE} = (0.01745)^2 \approx 3.05 \times 10^{-4}\text{ rad}^2\).
By standard regression metrics, this policy appears to match the expert
with \(99.7\%\) directional accuracy. In closed-loop execution, however,
the state evolves under the kinematic relation
\(\dot{y} = v \sin(\theta)\). A constant heading deviation of
\(1.0^\circ\) produces a lateral velocity of:
\[v_y = 1.0\text{ m/s} \times \sin(1.0^\circ) \approx 0.01745\text{ m/s}\]
Over the \(10\text{ m}\) transit (\(10\text{ s}\) duration), this
velocity produces a cumulative lateral displacement:
\[y(10\text{ s}) = \int_0^{10} v \sin(\epsilon_\theta) \, dt = 10\text{ s} \times 0.01745\text{ m/s} \approx 0.1745\text{ m} = 17.5\text{ cm}\]
If the corridor provides \(15.0\text{ cm}\) of clearance to a side
boundary, the platform collides with the wall at \(t = 8.6\text{ s}\)
after traveling \(8.6\text{ m}\)
(\(y(8.6\text{ s}) = 8.6 \times 0.01745 = 0.150\text{ m}\)). The policy
achieves an open-loop error of \(0.0175\text{ rad}\) and a closed-loop
task success rate of zero percent.

\protect\phantomsection\label{callout-ruleux2a-7.5}
\begin{fbxSimple}{callout-rule}{Architect’s Rule of Thumb:}{Open-Loop Blindness}
\phantomsection\label{callout-rule*-7.5}

Open-loop loss measures dataset mimicry; closed-loop dynamics decide
physical survival. A policy with near-zero offline prediction error will
diverge catastrophically the instant an uncorrected action bias
integrates through the physical plant.

\end{fbxSimple}

The \textbf{offline-online evaluation disconnect} is the fundamental
decoupling of open-loop dataset loss metrics (such as action mean
squared error or imitation log-likelihood on static logs) from
closed-loop physical stability and task survival, caused by the
inability of static evaluation to reflect compounding kinematic
integration errors
(\(y(t) \approx \int_0^t v \sin \epsilon_\theta \, d\tau\)) and
self-induced distribution shift during embodied execution.

This divergence occurs because open-loop evaluation resets the machine
to an expert state at every clock cycle, masking the compounding nature
of integration errors. In closed loop, the state at index \(t\) is the
expanded composition of all previous actions:
\[s_t = f(f(\dots f(s_0, a_0)\dots, a_{t-2}), a_{t-1})\] The first minor
action error shifts the physical machine from the demonstrated state
\(s_1^*\) to an unvisited state \(s_1\). Once the platform drifts
\(2.0\text{ cm}\) off nominal, it enters a state that does not appear
anywhere in \(\mathcal{D}_{\mathrm{demo}}\). Because the policy was
evaluated only on demonstrator states where \(y = 0\text{ m}\), the
offline metric contains no evidence regarding what action
\(\pi_\theta(s)\) the policy will produce when \(y \neq 0\text{ m}\). If
the training distribution lacked demonstrations of recovery maneuvers
from off-center coordinates, the policy can output larger corrective
errors, accelerating the divergence instead of arresting it.

This disconnect manifests across all physical machine archetypes:

\begin{itemize}
\tightlist
\item
  \emph{Mobile Systems:} As derived above, a \(1.0^\circ\) angular error
  integrates across \(10\text{ m}\) into a \(17.5\text{ cm}\) lateral
  offset, breaching corridor clearance limits.
\item
  \emph{Manipulator Systems:} An end-effector position error of
  \(0.5\text{ mm}\) per step during a contact assembly task (e.g.,
  chamfered pin insertion) appears negligible on offline Cartesian MSE.
  In closed-loop contact, that \(0.5\text{ mm}\) displacement drives the
  pin against the chamfer rim, generating high lateral reaction forces
  that jam the joint and trip the robot's torque limits
  (\(F_z > F_{\max}\)).
\item
  \emph{Process \& Energy Systems:} A temperature prediction error of
  \(0.5^\circ\text{C}\) per minute in an industrial reactor cooling
  policy seems minor on offline logs. In closed loop, failing to
  increase coolant flow compounds thermal accumulation, driving reactor
  core temperature past the safe \(95^\circ\text{C}\) trip threshold
  within twenty minutes.
\end{itemize}

The relationship between offline and online metrics is directional.
Open-loop evaluation is optimistically blind to compounding trajectory
divergence and unvisited state instabilities, because it assumes an
external mechanism continuously restores the machine to the expert path.
Conversely, open-loop action metrics can be pessimistic when the
physical task permits multiple valid actions. If a human demonstrator
avoided an obstacle by steering left, an open-loop evaluation penalizes
a candidate policy that steers right, recording a large action error
against \(a_t^*\). In closed-loop execution, passing on the right clears
the obstacle and reaches the target state with equal validity. An
action-level prediction error measures divergence from a specific
demonstration recording, not the satisfaction of the physical task.

As illustrated in Figure~\ref{fig-07-open-vs-closed-loop}, open-loop
evaluation evaluates candidate actions only on demonstrator states where
errors are immediately discarded, obscuring the cumulative trajectory
divergence that drives closed-loop physical execution into catastrophic
wall collisions.

\begin{figure}[t]

\centering{

\includegraphics[width=1\linewidth,height=\textheight,keepaspectratio]{index_files/mediabag//Users/VJ/GitHub/MLSysBook-integrate-physical/physical/book/chapters/07-evaluation/figures/fig07_open_vs_closed_loop.pdf}

}

\caption{\label{fig-07-open-vs-closed-loop}\emph{Open-Loop Error Masking
versus Closed-Loop Compounding State Divergence.} (a) In open-loop
evaluation (dataset trace replay), a candidate policy is scored
exclusively against recorded demonstrator states
\(s_t^* \in \mathcal{D}_{\mathrm{demo}}\), where a constant
\(1.0^\circ\) (\(0.0175\text{ rad}\)) heading bias produces an apparent
\(99.7\%\) directional accuracy
(\(\text{MSE} = 3.05 \times 10^{-4}\text{ rad}^2\)) because the
evaluation harness resets the platform to the centerline at every step.
(b) In closed-loop physical execution, the unmodeled angular bias
continuously integrates through kinematic state transitions
(\(\dot{y} = v \sin \theta\)), causing cumulative lateral displacement
that breaches the \(15\text{ cm}\) clearance boundary at
\(t = 8.6\text{ s}\) (\(x = 8.6\text{ m}\)), reducing closed-loop task
success to \(0\%\).}

\end{figure}%

This fundamental disconnect between open-loop statistical agreement and
physical closed-loop survival is systematized in
Table~\ref{tbl-07-offline-closed-loop-disconnect}, which maps common
machine learning validation metrics to their underlying physical failure
modes and corresponding closed-loop evaluators.

\begin{longtable}[]{@{}
  >{\raggedright\arraybackslash}p{(\linewidth - 8\tabcolsep) * \real{0.2000}}
  >{\raggedright\arraybackslash}p{(\linewidth - 8\tabcolsep) * \real{0.2000}}
  >{\raggedright\arraybackslash}p{(\linewidth - 8\tabcolsep) * \real{0.2000}}
  >{\raggedright\arraybackslash}p{(\linewidth - 8\tabcolsep) * \real{0.2000}}
  >{\raggedright\arraybackslash}p{(\linewidth - 8\tabcolsep) * \real{0.2000}}@{}}
\caption{\emph{Offline Open-Loop Prediction Metrics versus Closed-Loop
Physical Survival Disconnect Matrix.} Comparing standard machine
learning benchmark metrics evaluated on static demonstration logs
against the physical dynamics of embodied execution. Open-loop metrics
implicitly assume that demonstrator state distributions remain invariant
to policy errors; in closed loop, unmodeled actuator biases integrate
into compounding spatial drift, multimodal action alternatives incur
spurious penalties, and distribution shifts trigger unrecoverable
boundary
violations.}\label{tbl-07-offline-closed-loop-disconnect}\tabularnewline
\toprule\noalign{}
\begin{minipage}[b]{\linewidth}\raggedright
Offline Evaluation Metric
\end{minipage} & \begin{minipage}[b]{\linewidth}\raggedright
Open-Loop Statistical Assumption
\end{minipage} & \begin{minipage}[b]{\linewidth}\raggedright
Physical Disconnect Mechanism
\end{minipage} & \begin{minipage}[b]{\linewidth}\raggedright
Closed-Loop Hardware Failure Mode
\end{minipage} & \begin{minipage}[b]{\linewidth}\raggedright
Mandatory Closed-Loop Evaluator \& Metric
\end{minipage} \\
\midrule\noalign{}
\endfirsthead
\toprule\noalign{}
\begin{minipage}[b]{\linewidth}\raggedright
Offline Evaluation Metric
\end{minipage} & \begin{minipage}[b]{\linewidth}\raggedright
Open-Loop Statistical Assumption
\end{minipage} & \begin{minipage}[b]{\linewidth}\raggedright
Physical Disconnect Mechanism
\end{minipage} & \begin{minipage}[b]{\linewidth}\raggedright
Closed-Loop Hardware Failure Mode
\end{minipage} & \begin{minipage}[b]{\linewidth}\raggedright
Mandatory Closed-Loop Evaluator \& Metric
\end{minipage} \\
\midrule\noalign{}
\endhead
\bottomrule\noalign{}
\endlastfoot
\textbf{Action Mean Squared
Error}\(\text{MSE} = \frac{1}{T}\sum_{t=1}^T \|\hat{\mathbf{a}}_t - \mathbf{a}_t^*\|_2^2\)
& Small \(L_2\) prediction residuals produce proportionally small
trajectory deviations. & Unmodeled kinematic bias
\(\boldsymbol{\epsilon}_\theta\) integrates over time
(\(y(t) \approx \int_0^t v \sin \boldsymbol{\epsilon}_\theta \, d\tau\)),
driving states into out-of-distribution regions. & Catastrophic boundary
collision, curb strike, or lane departure despite \(<0.01\text{ rad}^2\)
MSE. & \textbf{Closed-Loop Rollout Tracking
Error}\(\max_{t} \|\mathbf{p}_t - \mathbf{p}_{\text{ref}}\|\) and
Time-to-Intervention (\(T_{\text{TTI}}\)). \\
\textbf{Top-\(1\) / Discrete Action
Accuracy}\(\text{Acc} = \frac{1}{T}\sum_{t=1}^T \mathbb{I}(\hat{a}_t = a_t^*)\)
& High classification accuracy on logged states ensures identical
trajectory execution. & Equal penalty for benign alternative maneuvers
vs catastrophic boundary excursions; blind to temporal correlation of
errors. & High-frequency actuator chatter, rapid control sign-flipping,
or sudden structural joint over-torque. & \textbf{Smoothness \& Dynamic
Feasibility}Peak jerk \(\|\dddot{\mathbf{q}}\|_{\infty}\) and actuator
torque limits
\(\|\boldsymbol{\tau}\|_{\infty} \le \boldsymbol{\tau}_{\max}\). \\
\textbf{Cosine Similarity / Directional
Agreement}\(\cos \phi = \frac{\hat{\mathbf{a}}_t \cdot \mathbf{a}_t^*}{\|\hat{\mathbf{a}}_t\|_2 \|\mathbf{a}_t^*\|_2}\)
& Collinear action proposals produce dynamically stable physical paths.
& Metric ignores control magnitude errors; under-actuation fails to
overcome static friction or over-actuation triggers tire slip. &
Limit-cycle hunting, velocity stall on inclines, or spin-out during
dynamic cornering maneuvers. & \textbf{Phase Margin \& Slip
Ratio}Effective damping ratio \(\zeta \ge 0.707\) and wheel slip ratio
\(s \le 0.15\). \\
\textbf{Negative Log-Likelihood
(NLL)}\(\mathcal{L}_{\text{NLL}} = -\frac{1}{T}\sum_{t=1}^T \log \pi_\theta(\mathbf{a}_t^* \mid \mathbf{s}_t^*)\)
& High likelihood on expert demonstrations guarantees unimodal
trajectory convergence. & Penalizes valid multimodal bypass trajectories
(e.g., navigating left vs right around a symmetric obstacle) by
averaging modes. & Policy outputs the unfeasible arithmetic mean of
modes, steering straight into the obstacle center. & \textbf{Closed-Loop
Task Success Rate}Terminal goal completion
\(\phi(\text{Traj}) \in \{0, 1\}\) over randomized goal offsets. \\
\textbf{Sensory Reconstruction / Chamfer
Loss}\(\mathcal{L}_{\text{recon}} = \|\hat{\mathbf{o}}_{t+1} - \mathbf{o}_{t+1}^*\|_2^2\)
& Accurate world-model state prediction ensures safe receding-horizon
plan synthesis. & Visual/feature reconstruction loss weights background
pixels over sub-pixel spatial obstacles and contact manifolds. & Phantom
obstacle braking, clipping thin obstacles (cables, glass), or false
clearance estimation. & \textbf{Spatial Clearance \& CBF
Invariance}Minimum obstacle distance \(d_{\min} \ge 0.50\text{ m}\) and
barrier condition \(h(\mathbf{x}) \ge 0\). \\
\end{longtable}

Recorded replay remains widely used because it evaluates model
checkpoints across stored arrays in seconds, consuming compute cycles
without occupying physical test cells or incurring mechanical wear.
Closed-loop simulation restores the causal recurrence
\(s_{t+1} = f_{\text{sim}}(s_t, a_t)\), allowing compounding drift to
manifest and testing whether a policy can stabilize off-nominal states
without risking physical hardware. Yet simulation replaces the open-loop
distribution assumption with a dynamical modeling assumption. The
simulation transition operator \(f_{\text{sim}}\) approximates real
physics \(f_{\text{real}}\), but it abstracts away unmodeled friction
non-linearities, contact mechanics, and physical transport latencies.
Establishing that a policy stabilizes \(f_{\text{sim}}\) demonstrates
mathematical convergence inside the simulator, but it does not establish
how the policy behaves when executed against physical hardware.

On heterogeneous physical hardware---such as a canonical heterogeneous
dual-brain edge System-on-Chip (SoC) bench twin---this open-loop versus
closed-loop disconnect maps directly across the compute boundary. When
an unprivileged application processor (Linux MPU) runs offline trace
replay, it measures only whether neural inference executes within its
memory and compute budgets. In physical closed-loop execution, the Linux
MPU continuously dispatches candidate action proposals via shared-SRAM
mailboxes (RPMSG) to the deterministic microcontroller (bare-metal MCU).
If policy proposals drift or chatter, the real-time MCU enforcer detects
the spatial boundary breach or lease expiration, rejecting proposals and
asserting emergency deceleration. Evaluating a policy offline without
the physical plant and the real-time permission referee in the loop
evaluates tensor mathematics, not cyber-physical machine survival.

\section{The Evidence Regimes and Their
Boundaries}\label{sec-evaluation-7-4}

Evaluating a learned policy requires selecting an execution environment,
and systems engineering divides these environments into four distinct
evidence regimes: recorded replay, closed-loop simulation, shadow
operation, and constrained physical testing. Each regime establishes a
different relationship between the policy, its sensory inputs, and the
state of the machine. Consequently, each regime supports a strictly
bounded claim. Recorded replay evaluates numerical prediction error over
static data. Simulation evaluates closed-loop stability against a
synthetic transition operator. Shadow operation evaluates onboard
execution timing and behavioral divergence against an active baseline.
Constrained physical testing evaluates closed-loop dynamics within an
artificially restricted operational envelope. None of these four regimes
establishes how the unassisted policy behaves during unconstrained
physical deployment, because each regime alters either the causal
feedback of the control loop or the distribution of states the machine
encounters.

The \textbf{four evidence regimes} are the hierarchical testing
environments spanning recorded replay (evaluating numerical loss on
static data), closed-loop simulation (evaluating stability against
synthetic transition models), shadow operation (evaluating runtime
timing and divergence against an active baseline), and constrained
physical testing (evaluating real dynamics inside safety supervisor
bounds), each supporting bounded claims while remaining blind to
open-deployment tail risk.

In recorded replay, the policy executes as a feedforward computation on
stored sensor buffers. The environment does not respond to the policy
outputs, and the state sequence remains fixed to the trajectory recorded
during data collection. This regime supports claims regarding
computational throughput, memory allocation, and the statistical loss of
the model relative to a logged dataset. It establishes whether the
policy executes within a target latency budget, such as \(18\text{ ms}\)
on an onboard accelerator, and whether its numerical outputs match the
recorded actions. It cannot support any claim about closed-loop
stability, because the difference between the predicted action and the
logged action never alters the subsequent state observation.

Closed-loop simulation restores the causal feedback loop by substituting
an explicit mathematical model, \(f_{\text{sim}}\), for the physical
machine and its surroundings. When the policy outputs an action, the
simulation integrates the equations of motion, updates the scene, and
renders synthetic sensor inputs for the next cycle at index \(t+1\).
This regime supports claims about the policy's algorithmic behavior
under the modeled dynamics. An engineer can run thousands of parallel
simulator instances to verify that a mobile base recovers when
initialized \(0.5\text{ m}\) off its nominal path, or that an
articulated arm avoids self-collision across its workspace. However,
this evidence holds only with respect to \(f_{\text{sim}}\). Unmodeled
mechanical backlash, tire slip non-linearities, contact friction
transitions, and lighting variations are excluded by construction.

Shadow operation mounts the candidate policy onto an active physical
machine while a validated baseline controller, or a human operator,
retains authority over the actuators. The candidate policy receives live
sensor frames, runs onboard inference under real mechanical vibration
and bus traffic, and computes candidate actions
\(a_t^{\text{cand}}\).\sidenote{\footnotesize Hardware-in-the-loop buses introduce
  physical contention: CAN frame arbitration, PCIe DMA cache-snooping
  stalls during unaligned tensor copies, and TSN clock jitter introduce
  \(5\text{--}25\text{ ms}\) age variations, eroding phase margin even
  when neural inference execution is deterministic.} These candidate
actions are logged but never dispatched to the motor drives. Shadow
operation supports claims about physical execution timing, validating
whether the inference pipeline sustains a \(50\text{ Hz}\) update rate
under operational thermal conditions, and quantifying how often
\(a_t^{\text{cand}}\) diverges from the baseline action
\(a_t^{\text{base}}\) by more than a specified threshold, such as a
steering difference of \(0.1\text{ rad}\). Yet shadow operation shares
the fundamental limitation of recorded replay. Because the baseline
controller steers the platform, the platform remains on the baseline's
trajectory. If the candidate policy outputs an action that would have
driven the wheels off a curb, the physical platform does not leave the
road, and the sensor stream never reflects the candidate's impending
failure.

Constrained physical testing closes the causal loop on real hardware by
granting the candidate policy authority over the actuators inside a
bounded envelope. Boundaries take the form of physical tethers, velocity
clamps, padded test cells, or an automated safety supervisor running at
\(1000\text{ Hz}\) in the nervous system that vetoes unsafe commands.
This regime supports claims about physical contact dynamics, actuator
saturation, and sensor noise within the tested envelope. It establishes
whether the machine can track a trajectory at \(1.5\text{ m/s}\) on
concrete without breaking traction. What it cannot establish is the tail
behavior of the policy in open deployment. Every time a safety tether
tensions or a supervisory controller intervenes to arrest a
\(2.0\text{ m/s}\) skid, the evaluation truncates the state trajectory.
The observed data measures the performance of the combined
policy-plus-supervisor system within the constraint boundary, not the
unassisted policy in an unconstrained environment.

Distinguishing between these regimes requires tracking what properties
transfer across regime boundaries and what properties stay behind.
Computational properties transfer reliably. A tensor operation that
consumes \(1.2\text{ MB}\) of memory and requires \(18\text{ ms}\) of
compute on an edge device in replay will consume the same memory and
approximately the same compute in shadow mode or on constrained
hardware. Algorithmic invariants also transfer. A numerical overflow
that causes an output to saturate in simulation will manifest
identically on physical hardware given the same input tensor. What does
not transfer across regimes is the state distribution. The moment the
environment model changes from \(f_{\text{sim}}\) to
\(f_{\text{real}}\), or the authority switches from a safety supervisor
to the autonomous policy, the sequence of states
\(s_0, s_1, \dots, s_t\) diverges.

Treating these regimes as a simple ladder of fidelity creates an
engineering trap. Teams frequently assume that advancing from replay to
simulation, from simulation to shadow mode, and from shadow mode to
constrained testing progressively strengthens a single unified
confidence claim. In reality, each transition trades one experimental
distortion for another. Simulation trades real sensor noise and physical
contact for closed-loop causality. Shadow mode trades closed-loop
causality back for real sensor streams and real hardware execution.
Constrained testing restores causality and physical reality, but
introduces state censorship through safety boundaries. A policy that
achieves zero interventions across one hundred trials in a constrained
test cell has not established a low failure rate in the wild. It has
established that its trajectory deviations remained within the envelope
where the constraint mechanisms did not engage.

Rigorous evaluation requires matching each engineering question to the
specific regime capable of answering it, rather than defaulting to the
most expensive apparatus available. Questions about inference latency,
memory footprint, and numerical consistency belong in recorded replay or
on target hardware in shadow mode. Questions about kinematic envelope
coverage, recovery from synthetic disturbances, and multi-agent
coordination belong in simulation, where millions of boundary conditions
can be tested without hardware wear. Questions about electrical noise,
bus jitter, and mechanical vibration belong in shadow operation.
Questions about traction, contact compliance, and actuator thermal rise
under load belong in constrained physical trials. Each regime provides
valid evidence for its own domain, while remaining blind to the failure
modes of the others.

\section{What Physical Trials Cost, and What They Cannot
Cover}\label{sec-evaluation-7-5}

A physical trial does not execute in the abstract time of a simulation
step or the zero-cost replay of a log file. It consumes wall-clock time,
stresses mechanical linkages, heats motor windings, and risks physical
damage to the hardware. Consider an analytical example of a mobile robot
evaluating a trajectory policy over a \(50\text{ m}\) course at a
nominal velocity of \(1.0\text{ m/s}\). The active run occupies
\(50\text{ s}\). However, budgeting the trial by active run time ignores
the physical mechanics of the test cell:

\begin{itemize}
\tightlist
\item
  \emph{Active Navigation Transit:}
  \(50\text{ m} / 1.0\text{ m/s} = 50\text{ s}\).
\item
  \emph{Return \& Starting Pose Reset:} Autonomous or teleoperated
  return to initial dock waypoint consumes \(120\text{ s}\).
\item
  \emph{Telemetry Flush \& Sensor Re-Zeroing:} Committing high-rate ring
  buffers to NVMe storage and re-calibrating IMU biases consumes
  \(30\text{ s}\).
\item
  \emph{Actuator Thermal Rest Window:} Dissipating Joule heating from
  motor stator windings to keep temperatures below their
  \(80^\circ\text{C}\) insulation ceiling requires \(60\text{ s}\).
\end{itemize}

The true wall-clock cost of one \(50\text{ s}\) physical trial is:
\[T_{\text{trial}} = 50\text{ s} + 120\text{ s} + 30\text{ s} + 60\text{ s} = 260\text{ s} \approx 4.33\text{ minutes}\]
Completing one hundred runs requires over \(7.2\text{ hours}\) of
continuous test cell occupancy and \(14.4\text{ person-hours}\) from two
safety operators monitoring manual emergency stops.

The physical plant changes while these trials accumulate. Statistical
evaluation assumes that repeated trials are independent and identically
distributed draws from a stationary data-generating process. On physical
machinery, mechanical wear violates stationarity directly. As a mobile
chassis completes hundreds of trials, planetary gearheads shed surface
material, expanding backlash from \(0.05^\circ\) to \(0.20^\circ\).
Rubber drive tires wear down, altering the effective wheel radius and
shifting the peak friction coefficient by over \(15\%\) across different
floor segments. Structural joints loosen, damping characteristics
change, and motor commutator brushes abrade. By trial one thousand, the
policy is interacting with a mechanical system whose transfer function
differs measurably from the system present at trial one. Pooling results
across one thousand trials treats a nonstationary progression of
degrading hardware states as if it were a repeated test of a single
invariant machine. Early successes do not guarantee that the worn
machine can execute the same trajectory, and late failures often reflect
physical degradation rather than an algorithmic deficiency.

This physical degradation is compounded by the problem of state
censorship. The operational states that an engineering team most needs
to evaluate are those at the stability boundary of the vehicle,
including sudden loss of traction, high-speed collision avoidance, and
sensor blinding. Yet these are precisely the states that operators
cannot permit the physical machine to explore freely. Entering an
unrecoverable skid or colliding with an obstacle damages the chassis,
destroys external sensing rigs, and takes the test facility offline for
repairs. Consequently, physical testing relies on safety cages,
geofences, and human operators with physical emergency stops who
intervene whenever a trajectory drifts outside an arbitrary conservative
envelope. These protective measures keep the hardware intact, but they
censor the empirical evidence. The trial records an external
intervention, but the true downstream consequence of the policy output
remains unobserved. Physical trials cannot sample the severe tails of
the operational distribution because the financial and physical cost of
entering those tails is prohibitive.\sidenote{\footnotesize Extreme Value Theory
  (Fisher-Tippett-Gnedenko) proves asymptotic sample extrema follow
  Gumbel, Fréchet, or Weibull distributions. Evaluating Physical AI via
  mean tracking error or empirical standard deviation implicitly assumes
  light symmetric tails, drastically underestimating catastrophic
  boundary excursion probabilities.}

While the Central Limit Theorem (CLT) governs the average operational
energy of a machine---driving nominal behavior in the bulk
distribution---cyber-physical machine survival is a weakest-link problem
governed by the maximum single peak disturbance
(\(M_n = \max\{X_1, \dots, X_n\}\)). To characterize the severe
disturbance tails that physical trials cannot safely explore, systems
engineers apply extreme value theory (EVT) under the
Fisher-Tippett-Gnedenko theorem to model the asymptotic distribution of
operational extrema, as detailed in
Table~\ref{tbl-07-extreme-value-tails}.

\textbf{Extreme-value tail risk} is the operational hazard governed by
rare, non-Gaussian disturbance extrema (such as bus transport latency
spikes, shock impact forces, or localized friction collapse modeled
under the Fisher-Tippett-Gnedenko theorem via Gumbel, Fréchet, or
Weibull distributions) whose catastrophic consequences reside in severe
distribution tails that central-tendency metrics systematically conceal.

\begin{longtable}[]{@{}
  >{\raggedright\arraybackslash}p{(\linewidth - 10\tabcolsep) * \real{0.1667}}
  >{\raggedright\arraybackslash}p{(\linewidth - 10\tabcolsep) * \real{0.1667}}
  >{\raggedright\arraybackslash}p{(\linewidth - 10\tabcolsep) * \real{0.1667}}
  >{\raggedright\arraybackslash}p{(\linewidth - 10\tabcolsep) * \real{0.1667}}
  >{\raggedright\arraybackslash}p{(\linewidth - 10\tabcolsep) * \real{0.1667}}
  >{\raggedright\arraybackslash}p{(\linewidth - 10\tabcolsep) * \real{0.1667}}@{}}
\caption{\emph{Extreme Value Tail Distribution Taxonomy across Physical
AI Disturbance Channels.} Parametric classification of operational
disturbance tails governing physical AI systems under the
Fisher-Tippett-Gnedenko theorem. Standard validation protocols that
report mean tracking errors or Gaussian confidence intervals
underestimate catastrophic tail risks by orders of magnitude; mitigating
extreme excursions requires dedicated deterministic nervous-system
interceptors, mechanical compliances, and hardware safety
supervisors.}\label{tbl-07-extreme-value-tails}\tabularnewline
\toprule\noalign{}
\begin{minipage}[b]{\linewidth}\raggedright
Disturbance Channel \& Physical AI Domain
\end{minipage} & \begin{minipage}[b]{\linewidth}\raggedright
EVT Tail Family \& Shape Parameter
\end{minipage} & \begin{minipage}[b]{\linewidth}\raggedright
Physical Generating Mechanism \& Hardware Reality
\end{minipage} & \begin{minipage}[b]{\linewidth}\raggedright
Central Metric Pitfall (Gaussian Fallacy)
\end{minipage} & \begin{minipage}[b]{\linewidth}\raggedright
Asymptotic Tail Hazard
\end{minipage} & \begin{minipage}[b]{\linewidth}\raggedright
Nervous System Interceptor \& Mitigator
\end{minipage} \\
\midrule\noalign{}
\endfirsthead
\toprule\noalign{}
\begin{minipage}[b]{\linewidth}\raggedright
Disturbance Channel \& Physical AI Domain
\end{minipage} & \begin{minipage}[b]{\linewidth}\raggedright
EVT Tail Family \& Shape Parameter
\end{minipage} & \begin{minipage}[b]{\linewidth}\raggedright
Physical Generating Mechanism \& Hardware Reality
\end{minipage} & \begin{minipage}[b]{\linewidth}\raggedright
Central Metric Pitfall (Gaussian Fallacy)
\end{minipage} & \begin{minipage}[b]{\linewidth}\raggedright
Asymptotic Tail Hazard
\end{minipage} & \begin{minipage}[b]{\linewidth}\raggedright
Nervous System Interceptor \& Mitigator
\end{minipage} \\
\midrule\noalign{}
\endhead
\bottomrule\noalign{}
\endlastfoot
\textbf{Compute \& Bus Transport Latency}(Autonomous Mobility \& Mobile
Manipulation) & \textbf{Gumbel (Type I)}\(\xi = 0\) (Exponential
Tail)\(F(x) = \exp(-e^{-(x-\mu)/\beta})\) & Unaligned DMA transfers,
PCIe cache-snoop stalls, and CAN bus frame arbitration collisions under
\(>70\%\) load. & Reporting mean latency \(\bar{\tau} = 18\text{ ms}\);
ignores \(P_{99.9}\) queuing tails reaching \(65\text{ ms}\). &
Accumulated phase lag \(\Delta \phi = -\omega_c \tau\) consumes phase
margin, inducing violent steering oscillation. & Hard real-time RTOS
priority inheritance, IEEE 802.1AS TSN, and deterministic watchdog drop
counters. \\
\textbf{Contact Impact Force \& Shock Torques}(Legged Locomotion \&
Heavy Manipulation) & \textbf{Fréchet (Type II)}\(\xi > 0\) (Heavy /
Power-Law Tail)\(F(x) = \exp(-x^{-\alpha})\) & Non-smooth mechanical
contact against unyielding rigid obstacles; foot strikes on uneven
terrain steps. & Assuming normal force distribution
\(\mathcal{N}(\mu, \sigma^2)\); underestimates \(10\times\) peak impact
transients. & Dynamic shear stress exceeds yield limit, stripping
harmonic drive gear teeth or fracturing linkages. & Series Elastic
Actuators (SEA), mechanical shear pins, and high-rate (\(1\text{ kHz}\))
analog current limits. \\
\textbf{Traction Coefficient \& Surface Slip}(Wheeled Mobile Robots \&
Delivery Rovers) & \textbf{Weibull (Type III)}\(\xi < 0\) (Bounded
Support)\(F(x) = \exp(-(-(x-\omega)/\beta)^\alpha)\) & Abrupt transition
onto localized friction anomalies: spilled industrial lubricants, wet
epoxy, or ice patches (\(\Delta \mu = -0.55\)). & Evaluating average
braking distance on dry concrete; ignores sudden left-tail support
collapse. & Sudden loss of lateral tire grip causing unrecoverable
understeer, skid excursions, or rollover. & Model Predictive Control
with dynamic friction estimation, wheel-slip traction clamps
(\(s \le 0.15\)). \\
\textbf{Ambient Photometric Glare \& Contrast}(Vision-Based Aerial \&
Field Robotics) & \textbf{Gumbel / GPD}\(\xi \ge 0\)
(Peaks-Over-Threshold)\(F(x) = 1 - (1 + \xi x / \sigma)^{-1/\xi}\) &
Low-angle direct solar irradiance (\(>10^5\text{ lx}\)) entering camera
optics paired with specular floor reflections. & Evaluating policy on
diffuse, uniform lab lighting (\(500\text{ lx}\)); blind to localized
sensor saturation. & Image sensor blooming and dynamic range collapse,
zeroing visual feature gradients across \(80\%\) of pixels. &
Multi-exposure HDR imaging, redundant multi-modal sensing (active FMCW
radar / lidar), analog photodiode tripwire. \\
\textbf{Aerodynamic Gusts \& Wind Shear}(UAVs, Quadrotors \& Delivery
Drones) & \textbf{Fréchet (Type II)}\(\xi > 0\) (Heavy Power-Law
Tail)\(F(x) = \exp(-x^{-\alpha})\) & Atmospheric convective turbulence,
urban canyon vortex shedding, and ground-effect recirculation bursts. &
Budgeting motor thrust against average wind speed; ignores cubic power
relationship (\(P \propto v^3\)). & Aerodynamic torque exceeds rotor
torque authority, inducing rotor stall, attitude loss, and crash. &
Dynamic airspeed sensing (Pitot arrays), Control Barrier Function thrust
clamps, ballistic parachute. \\
\end{longtable}

Attempting to bypass these physical bottlenecks by distributing trials
across a fleet of identical machines introduces spatial and
environmental correlation. When ten mobile robots run concurrently
across different bays of a testing warehouse, their operational
histories are not independent. The machines share identical
manufacturing tolerances from the same production batch, run on
identical concrete floor finishes, share ambient temperature profiles
governed by a single building ventilation system, and interact with
human operators trained under identical protocols. If an environmental
anomaly such as floor glazing or ambient lighting glare is absent from
the facility, none of the ten machines will experience it. As
illustrated by the multi-bay robotic qualification test facility in
Figure~\ref{fig-07-nist-test-arena}, standardized physical test fixtures
allow teams to isolate specific locomotion or perception capabilities
under repeatable physical geometry, but they also bound the state space
to that precise fixture envelope. Accumulating ten thousand fleet hours
across ten machines in one facility yields ten parallel observations of
a narrow environmental distribution, not ten thousand hours of
independent operational diversity. Treating aggregated fleet hours as a
direct substitute for independent statistical exposure obscures the
common-mode biases that dominate physical operations.

\begin{figure}[t]

\centering{

\includegraphics[width=0.75\linewidth,height=\textheight,keepaspectratio]{/Users/VJ/GitHub/MLSysBook-integrate-physical/physical/book/chapters/07-evaluation/figures/fig07_real_nist_test_arena.jpg}

}

\caption{\label{fig-07-nist-test-arena}\emph{Standardized Robotic
Evaluation Facility and Obstacle Bays.} Standardized physical test
arenas---such as the National Institute of Standards and Technology
(NIST) Robot Test Facility configured under ASTM E54.09 test
methods---provide repeatable physical obstacle stations (including
calibrated step fields, pitch and roll ramps, drop-offs, and constrained
passages) to benchmark mobile robot locomotion, sensor payload
stability, and operator control across discrete, repeatable trials.
While such structured arenas isolate mechanical and algorithmic
variables under controlled geometry and lighting, they bound the state
space the robot encounters: zero-failure runs across standardized bays
do not establish out-of-distribution robustness or cover environmental
anomalies (such as specular glare, surface contaminants, or novel
obstacle dynamics) present in unconstrained operational domains.}

\end{figure}%

The resulting mismatch between what physical trials yield and what
statistical confidence requires exposes the fundamental inadequacy of
classical estimation tools when applied to physical autonomy. In
standard machine learning benchmarks, practitioners frequently compute
confidence intervals using the asymptotic Wald normal approximation:
\[\hat{p} \pm z_{1-\alpha/2}\sqrt{\frac{\hat{p}(1-\hat{p})}{n}}\] When
testing high-reliability physical systems where zero failures are
observed across \(n\) runs (\(\hat{p} = 1.0\)), the estimated variance
collapses to zero (\(\hat{p}(1-\hat{p}) = 0\)), yielding the absurd and
degenerate confidence interval \([1.0, 1.0]\). The normal approximation
falsely claims absolute certainty precisely where sample evidence is
most incomplete.

Valid statistical inference requires exact, non-asymptotic interval
estimation via the Clopper-Pearson method
(\citeproc{ref-clopper1934use}{Clopper and Pearson 1934}). Inverting the
binomial relation determines the minimum number of consecutive
zero-failure trials \(n\) required to certify a target success
probability \(p_{\text{target}}\) at confidence level \(1-\alpha\):
\[n = \frac{\ln \alpha}{\ln p_{\text{target}}}\]

\protect\phantomsection\label{callout-mathux2a-7.6}
\begin{fbxSimple}{callout-math}{Fermi Math:}{Statistical Scaling of High-Reliability Assurance}
\phantomsection\label{callout-math*-7.6}

\emph{Problem:} An industrial qualification standard mandates certifying
a robotic policy with \(P_{99.9}\) reliability
(\(p_{\text{target}} = 0.999\), allowing a failure probability at most
\(10^{-3}\)) at \(95\%\) confidence (\(\alpha = 0.05\)). If each
physical trial requires \(T_{\text{trial}} = 260\text{ s}\), how many
zero-failure runs and how much wall-clock testing time are required?

\emph{Calculations:}

\begin{enumerate}
\def\labelenumi{\arabic{enumi}.}
\item
  \emph{Required zero-failure sample size (\(n\)).}
  \[n = \frac{\ln(0.05)}{\ln(0.999)} = \frac{-2.995732}{-0.0010005} \approx 2{,}994.23 \implies n = 2{,}995 \text{ trials}\]
\item
  \emph{Total Wall-Clock Testing Budget:}
  \[T_{\text{total}} = 2{,}995 \times 260\text{ s} = 778{,}700\text{ s} \approx 216.31\text{ hours} \approx 9.01\text{ days}\]
\item
  \emph{Personnel \& shift budget (2-operator safety protocol).}
  \[\text{Operator Hours} = 216.31\text{ h} \times 2 = 432.62\text{ person-hours} \approx 54.1\text{ eight-hour shifts}\]
\end{enumerate}

\emph{Conclusion:} Establishing \(99.9\%\) reliability requires nearly
\(3{,}000\) consecutive zero-failure trials and over nine continuous
24-hour days of test cell occupancy. On a single physical chassis,
continuous operation for \(216\text{ hours}\) under aggressive dynamics
exceeds the mean time between maintenance for drive belts, gearhead
backlash limits, and tire wear, guaranteeing that non-stationary
hardware degradation will invalidate the experiment before the
statistical target is reached.

\end{fbxSimple}

When reliability requirements scale to continuous-time physical systems
governed by an operational failure rate \(\lambda\) (failures per
operating hour), the exposure deficit becomes insurmountable. The
required zero-failure physical exposure time is governed by the Butler
\& Finelli exposure bound. As shown in the detailed callout box, at
standard \(95\%\) statistical confidence, bounding an operational
failure rate requires testing for roughly three times the target
mean-time-to-failure continuous hours
(\(T pprox 3.0 / \lambda\)).\sidenote{\footnotesize Butler and Finelli (1993) proved
  mathematically that validating ultra-high software reliability
  (\(\lambda \le 10^{-9}\text{ h}^{-1}\)) via operational exposure alone
  is physically impossible (\(T = -\ln(1-C)/\lambda\)), requiring
  architectural fault tolerance, runtime invariants, and deterministic
  fault injection over black-box burn-in.}

The \textbf{Butler \& Finelli infeasibility bound} is the mathematical
theorem establishing that certifying ultra-high cyber-physical
reliability (e.g., catastrophic failure rates \(\lambda \le 10^{-6}\) to
\(10^{-9}\text{ h}^{-1}\)) through black-box operational exposure alone
is physically impossible due to astronomical duration requirements
(\(T = -\ln(1-C)/\lambda\)), mathematically necessitating deterministic
runtime invariant enforcement and structured fault injection.

\protect\phantomsection\label{callout-incidentux2a-7.7}
\begin{fbxSimple}{callout-incident}{Landmark Engineering Incident:}{Butler \& Finelli Flight Control Software Reliability Infeasibility Study}
\phantomsection\label{callout-incident*-7.7}

\emph{Study Overview:} In their landmark NASA Langley investigation,
Ricky W. Butler and George B. Finelli proved mathematically that
validating ultra-high reliability requirements for life-critical
cyber-physical systems (such as commercial fly-by-wire flight control
software requiring a catastrophic failure rate
\(\lambda \le 10^{-9}\text{ failures/hour}\)) through empirical
operational testing alone is physically impossible
(\citeproc{ref-butler1993infeasibility}{Butler and Finelli 1993})
\sidenote{\footnotesize Butler and Finelli (1993) proved mathematically that
  validating ultra-high software reliability
  (\(\lambda \le 10^{-9}\text{ h}^{-1}\)) via operational exposure alone
  is physically impossible (\(T = -\ln(1-C)/\lambda\)), requiring
  architectural fault tolerance, runtime invariants, and deterministic
  fault injection over black-box burn-in.}.

\emph{Mathematical Proof \& Exposure Deficit:} Under a continuous-time
Poisson failure process with failure rate \(\lambda\), the probability
of observing zero catastrophic failures over test duration \(T\) is
\(P(0 \text{ failures in } T) = e^{-\lambda T}\). To certify that the
true failure rate is at most \(\lambda\) at confidence level
\(C = 1 - \alpha\), the required zero-failure physical exposure time is:
\[T = \frac{-\ln(1-C)}{\lambda} = \frac{\ln(1/\alpha)}{\lambda}\] For
standard \(95\%\) statistical confidence (\(\alpha = 0.05\)),
\(\ln(1/0.05) = \ln(20) \approx 2.996 \approx 3.0\), yielding:
\[T \approx \frac{3.0}{\lambda}\] Evaluating this bound across safety
integrity standards reveals an insurmountable physical deficit:

\begin{itemize}
\tightlist
\item
  \emph{Industrial machinery (IEC 61508 SIL 2 / PL d,
  \(\lambda = 10^{-6}\text{ h}^{-1}\)).} Requires
  \(T = 3.0 \times 10^6\text{ hours}\) (\(342\text{ years}\)) of
  continuous zero-failure execution.
\item
  \emph{Automotive autonomous driving (ISO 26262 ASIL D,
  \(\lambda = 10^{-8}\text{ h}^{-1}\)).} Requires
  \(T = 3.0 \times 10^8\text{ hours}\)
  (\(34{,}200\text{ years}\)).\sidenote{\footnotesize ISO 26262 ASIL D targets random
    hardware failures \(\lambda \le 10^{-8}\text{ h}^{-1}\). IEC 61508
    SIL 3 requires \(\lambda \in [10^{-8}, 10^{-7}]\text{ h}^{-1}\).}
\item
  \emph{Commercial flight control (DO-178C DAL A,
  \(\lambda = 10^{-9}\text{ h}^{-1}\)).} Requires
  \(T = 3.0 \times 10^9\text{ hours}\) (\(342{,}000\text{ years}\)).
\end{itemize}

\emph{Fleet Testing Fallacy:} Distributing tests across a fleet of one
hundred physical testbeds operating 24 hours a day, 365 days a year
still requires over \(3,420\text{ years}\) of continuous operation to
certify a \(10^{-9}\text{ h}^{-1}\) claim. Furthermore, fleet testbeds
share common-mode software builds, environmental biases, and fixture
setups, violating the statistical independence assumption.

\emph{Systems Engineering Impact:} The Butler \& Finelli Infeasibility
Theorem proves that ultra-reliability cannot be tested into a machine
through operational burn-in, fleet mileage accumulation, or passive
sampling. High-assurance physical AI systems must instead enforce safety
through deterministic runtime architectural supervisors, formal
invariant filtering, and active unannounced fault injection.

\end{fbxSimple}

\protect\phantomsection\label{callout-ruleux2a-7.8}
\begin{fbxSimple}{callout-rule}{Architect’s Rule of Thumb:}{The Butler-Finelli Invariant}
\phantomsection\label{callout-rule*-7.8}

Ultra-reliability cannot be tested into a machine; it must be
architected into its deterministic safety supervisor. When certifying
failure rates below \(10^{-6}\text{ h}^{-1}\), empirical testing
establishes nothing; deterministic runtime invariant filtering
guarantees survival.

\end{fbxSimple}

The mathematical scaling of these exact confidence lower bounds across
trial counts and confidence levels is visualized in
Figure~\ref{fig-07-confidence-bounds}, highlighting the widening gap
between achievable testing budgets and high-assurance reliability
targets.

\begin{figure}[t]

\centering{

\includegraphics[width=1\linewidth,height=\textheight,keepaspectratio]{index_files/mediabag//Users/VJ/GitHub/MLSysBook-integrate-physical/physical/book/chapters/07-evaluation/figures/fig07_confidence_bounds.pdf}

}

\caption{\label{fig-07-confidence-bounds}\emph{Non-Asymptotic Confidence
Lower Bounds and Physical Testing Budget Scaling.} Exact Clopper-Pearson
one-sided lower confidence bounds on operational success probability
\(p\) as a function of zero-failure physical trial count \(n\) across
confidence levels \(1-\alpha \in \{0.90, 0.95, 0.99\}\). While twenty
clean trials (\(n=20\)) yield only an \(86.1\%\) lower bound at \(95\%\)
confidence (permitting a \(13.9\%\) failure rate), establishing
high-assurance industrial reliability (\(P_{99.9}\), \(p \ge 0.999\))
demands \(2{,}995\) consecutive zero-failure runs
(\(216.3\text{ hours}\) of test cell occupancy at
\(260\text{ s}\)/trial), exceeding the mechanical wear life of physical
actuators and demonstrating why statistical certification alone cannot
replace deterministic runtime safeguards.}

\end{figure}%

As summarized in Table~\ref{tbl-07-confidence-bounds}, the
non-asymptotic scaling of confidence bounds reveals the severe
evidentiary deficit of small testing budgets across common trial sample
sizes, establishing the mathematical ceiling of empirical certification.

\begin{longtable}[]{@{}
  >{\raggedright\arraybackslash}p{(\linewidth - 14\tabcolsep) * \real{0.1250}}
  >{\raggedright\arraybackslash}p{(\linewidth - 14\tabcolsep) * \real{0.1250}}
  >{\raggedright\arraybackslash}p{(\linewidth - 14\tabcolsep) * \real{0.1250}}
  >{\raggedright\arraybackslash}p{(\linewidth - 14\tabcolsep) * \real{0.1250}}
  >{\raggedright\arraybackslash}p{(\linewidth - 14\tabcolsep) * \real{0.1250}}
  >{\raggedright\arraybackslash}p{(\linewidth - 14\tabcolsep) * \real{0.1250}}
  >{\raggedright\arraybackslash}p{(\linewidth - 14\tabcolsep) * \real{0.1250}}
  >{\raggedright\arraybackslash}p{(\linewidth - 14\tabcolsep) * \real{0.1250}}@{}}
\caption{\emph{Non-Asymptotic Binomial Confidence Lower Bounds under
Zero Failures.} Exact one-sided Clopper-Pearson lower confidence limits
on true operational success probability \(p\) as a function of
zero-failure physical trial count \(n\) across confidence levels
\(1-\alpha \in \{0.90, 0.95, 0.99\}\), calculated via
\(p_{\text{lower}} = \alpha^{1/n}\). The table pairs statistical bounds
with empirical test budgets (evaluated at the baseline \(260\text{ s}\)
per physical run derived in Section~\ref{sec-evaluation-7-5}) and
hardware wear limits, demonstrating why empirical black-box testing
alone cannot license ultra-reliable autonomy claims without
deterministic runtime
containment.}\label{tbl-07-confidence-bounds}\tabularnewline
\toprule\noalign{}
\begin{minipage}[b]{\linewidth}\raggedright
Zero-Failure Trials (\(n\))
\end{minipage} & \begin{minipage}[b]{\linewidth}\raggedright
\(90\%\) Confidence Lower Bound (\(p_{0.10}\))
\end{minipage} & \begin{minipage}[b]{\linewidth}\raggedright
\(95\%\) Confidence Lower Bound (\(p_{0.05}\))
\end{minipage} & \begin{minipage}[b]{\linewidth}\raggedright
\(99\%\) Confidence Lower Bound (\(p_{0.01}\))
\end{minipage} & \begin{minipage}[b]{\linewidth}\raggedright
Max Implied Failure Risk (\(1 - p_{0.05}\))
\end{minipage} & \begin{minipage}[b]{\linewidth}\raggedright
Implied Failure Interval (1 in \(K\) Runs)
\end{minipage} & \begin{minipage}[b]{\linewidth}\raggedright
Wall-Clock Test Budget (\(260\text{ s}\)/run)
\end{minipage} & \begin{minipage}[b]{\linewidth}\raggedright
Systems Engineering Takeaway \& Statistical License
\end{minipage} \\
\midrule\noalign{}
\endfirsthead
\toprule\noalign{}
\begin{minipage}[b]{\linewidth}\raggedright
Zero-Failure Trials (\(n\))
\end{minipage} & \begin{minipage}[b]{\linewidth}\raggedright
\(90\%\) Confidence Lower Bound (\(p_{0.10}\))
\end{minipage} & \begin{minipage}[b]{\linewidth}\raggedright
\(95\%\) Confidence Lower Bound (\(p_{0.05}\))
\end{minipage} & \begin{minipage}[b]{\linewidth}\raggedright
\(99\%\) Confidence Lower Bound (\(p_{0.01}\))
\end{minipage} & \begin{minipage}[b]{\linewidth}\raggedright
Max Implied Failure Risk (\(1 - p_{0.05}\))
\end{minipage} & \begin{minipage}[b]{\linewidth}\raggedright
Implied Failure Interval (1 in \(K\) Runs)
\end{minipage} & \begin{minipage}[b]{\linewidth}\raggedright
Wall-Clock Test Budget (\(260\text{ s}\)/run)
\end{minipage} & \begin{minipage}[b]{\linewidth}\raggedright
Systems Engineering Takeaway \& Statistical License
\end{minipage} \\
\midrule\noalign{}
\endhead
\bottomrule\noalign{}
\endlastfoot
\textbf{\(n = 10\)} & \(79.43\%\) & \(74.11\%\) & \(63.10\%\) &
\(25.89\%\) & \(1\) in \(3.9\) & \(0.7\text{ h}\) (\(43.3\text{ min}\))
& \emph{Trivial Proof-of-Concept.} Consistent with a \(1\)-in-\(4\)
catastrophic failure rate; provides zero deployment license. \\
\textbf{\(n = 20\)} & \(89.13\%\) & \(86.09\%\) & \(79.43\%\) &
\(13.91\%\) & \(1\) in \(7.2\) & \(1.4\text{ h}\) (\(86.7\text{ min}\))
& \emph{Typical Lab Benchmark.} Consistent with a \(1\)-in-\(7\) failure
rate; licenses only supervised, tethered testing. \\
\textbf{\(n = 50\)} & \(95.50\%\) & \(94.18\%\) & \(91.20\%\) &
\(5.82\%\) & \(1\) in \(17.2\) & \(3.6\text{ h}\) (\(216.7\text{ min}\))
& \emph{Screening Baseline.} Permits detecting coarse behavioral bugs;
insufficient for unconstrained field trials. \\
\textbf{\(n = 100\)} & \(97.72\%\) & \(97.05\%\) & \(95.50\%\) &
\(2.95\%\) & \(1\) in \(33.9\) & \(7.2\text{ h}\) (\(433.3\text{ min}\))
& \emph{Full Single-Day Battery.} Licenses human-in-the-loop pilot
deployment with physical safety override. \\
\textbf{\(n = 250\)} & \(99.08\%\) & \(98.81\%\) & \(98.17\%\) &
\(1.19\%\) & \(1\) in \(84.0\) & \(18.1\text{ h}\)
(\(1{,}083\text{ min}\)) & \emph{Multi-Day Rig Testing.} Approaches
\(99\%\) operational reliability bound; wear non-stationarity begins to
manifest. \\
\textbf{\(n = 500\)} & \(99.54\%\) & \(99.40\%\) & \(99.08\%\) &
\(0.60\%\) & \(1\) in \(167.4\) & \(36.1\text{ h}\)
(\(2{,}167\text{ min}\)) & \textbf{Automated Hardware-in-the-Loop (HIL)
Qualification.} Licenses geofenced commercial deployment; exceeds
single-component wear lifecycle. \\
\textbf{\(n = 1{,}000\)} & \(99.77\%\) & \(99.70\%\) & \(99.54\%\) &
\(0.30\%\) & \(1\) in \(334.3\) & \(72.2\text{ h}\)
(\(4{,}333\text{ min}\)) & \emph{Extensive Fleet Aggregate.} Fails the
\(99.9\%\) reliability target (\(0.30\% > 0.10\%\) failure budget);
requires \(3\times\) more trials. \\
\textbf{\(n = 2{,}995\)} & \(99.92\%\) & \(99.90\%\) & \(99.85\%\) &
\(0.10\%\) & \(1\) in \(1{,}000.3\) & \(216.3\text{ h}\)
(\(12{,}978\text{ min}\)) & \textbf{\(P_{99.9}\) Industrial Standard.}
First threshold establishing \(<10^{-3}\) failure rate; physically
unattainable on a single chassis. \\
\end{longtable}

These physical and statistical ceilings force an unavoidable
architectural conclusion: physical AI safety cannot rest on statistical
confidence in unconstrained model weights. When the exposure required to
validate an autonomous policy cannot be purchased through physical
testing because of wall-clock time, mechanical wear, and safety hazards,
systems engineering mandates two foundational strategies:

\begin{enumerate}
\def\labelenumi{\arabic{enumi}.}
\tightlist
\item
  \textbf{Dual-Brain Runtime Containment}: Bounding the physical
  authority of the learned policy by delegating safety invariants,
  acceleration envelopes, and emergency stops to an independent,
  deterministic nervous system (such as the Linux MPU proposal engine
  and bare-metal MCU permission referee introduced in
  Chapter~\ref{sec-boundary}). Bounding policy authority ensures that
  the safety of the machine does not depend on empirical claims that
  physical trials are fundamentally incapable of establishing.
\item
  \textbf{Unannounced Fault Injection on Physical Benches}: Rather than
  waiting millions of operational hours for passive disturbances to
  strike, verification engineers actively attack system invariants on
  hardware test benches (such as a canonical heterogeneous dual-brain
  desk bench twin). By deterministically injecting unannounced bus
  latency spikes, memory transport contention, sensor frame corruption,
  and localized friction drops during active rollouts, engineers
  directly measure whether the real-time supervisor intercepts
  excursions within bounded physical margins.
\end{enumerate}

\textbf{Unannounced fault injection} is the active evaluation
methodology of deterministically introducing hardware, timing, and
sensory disturbances (such as bus transport latency spikes, packet
corruption, clock skew, and localized friction collapse) during live
test bench execution without prior model notification, directly
verifying that real-time safety supervisors intercept excursions within
bounded physical margins.

The complete operational protocol for executing physical trials,
enforcing full-denominator accounting, and evaluating the exact
non-asymptotic stopping rule without falling into the optional stopping
fallacy is formalized below:

\subsection{Closed-Loop Physical Trial Exposure and Stopping Oracle
Protocol}\label{closed-loop-physical-trial-exposure-and-stopping-oracle-protocol}

\emph{Execution Context \& Invariants:}

\begin{itemize}
\tightlist
\item
  Automated Physical Test Cell / Hardware-in-the-Loop (HIL)
  Qualification Rig.
\item
  \emph{Statistical Invariant:} Pre-registered required sample size
  \(N_{\text{req}} = \lceil \ln(\alpha) / \ln(p_{\text{target}}) \rceil\);
  zero informal looks to prevent optional stopping inflation
  (\(P(\text{false alarm}) \le \alpha\)).
\item
  \emph{Physical Invariant:} Complete-denominator ledger accounting;
  every trial initiated under policy authority increments
  \(n_{\text{attempted}}\) and cannot be discarded.
\end{itemize}

\emph{Execution Protocol:}

\begin{enumerate}
\def\labelenumi{\arabic{enumi}.}
\item
  \textbf{Pre-Registration \& Exposure Sizing (Stage 1):} Compute
  required zero-failure sample count
  \(N_{\text{req}} \leftarrow \lceil \ln(\alpha) / \ln(p_{\text{target}}) \rceil\),
  initialize ledger counters
  (\(n_{\text{attempted}} \leftarrow 0, n_{\text{success}} \leftarrow 0, n_{\text{supervisor}} \leftarrow 0, n_{\text{timeout}} \leftarrow 0, n_{\text{fault}} \leftarrow 0\)),
  and generate an immutable experiment pre-registration record with
  cryptographic digest
  \(\text{SHA256}(\pi_\theta \parallel \Theta_{\text{safe}} \parallel N_{\text{req}})\).
\item
  \textbf{Thermal \& Fixture Reset (Stage 2):} For each trial
  \(i = 1, 2, \dots, N_{\text{req}}\):

  \begin{itemize}
  \tightlist
  \item
    Sample condition tuple \(c_i \sim \mathcal{D}_{\text{ODD}}\) from
    the stratified parameter grid (payload, surface friction, lighting).
  \item
    Poll motor winding temperatures \(T_{\text{motor}}\); if
    \(T_{\text{motor}} > 40^\circ\text{C}\), enter thermal rest loop
    until cooled.
  \item
    Reposition physical chassis to initial pose \(s_0\), zero IMU
    biases, and flush telemetry ring buffers.
  \end{itemize}
\item
  \textbf{Closed-Loop Execution \& Real-Time Supervision (Stage 3):}

  \begin{itemize}
  \tightlist
  \item
    Increment attempted trial count:
    \(n_{\text{attempted}} \leftarrow n_{\text{attempted}} + 1\).
  \item
    Launch closed-loop trajectory execution under policy \(\pi_\theta\)
    at \(50\text{ Hz}\) control cadence.
  \item
    Concurrently, execute deterministic supervisor
    \(\mathcal{M}_{\text{referee}}\) at \(1000\text{ Hz}\) on the
    bare-metal MCU:

    \begin{itemize}
    \tightlist
    \item
      Check tracking error
      \(\|s(t) - s_{\text{ref}}(t)\| \le e_{\max}\).
    \item
      Check proposal freshness lease
      \(\tau_{\text{age}} \le \Delta t_{\max}\).
    \item
      Check spatial obstacle clearance
      \(d_{\text{clear}}(t) \ge d_{\min}\).
    \item
      Check actuator current/torque limits
      \(\|\boldsymbol{\tau}(t)\| \le \boldsymbol{\tau}_{\max}\).
    \end{itemize}
  \end{itemize}
\item
  \textbf{Trial Outcome Classification (Stage 4):}

  \begin{itemize}
  \tightlist
  \item
    If spatial goal \(g\) reached within \(T_{\max}\) with zero
    supervisor overrides: record success (\(Y_i \leftarrow 1\);
    \(n_{\text{success}} \leftarrow n_{\text{success}} + 1\)).
  \item
    Else if \(\mathcal{M}_{\text{referee}}\) asserted safety override:
    record censored failure (\(Y_i \leftarrow 0\);
    \(n_{\text{supervisor}} \leftarrow n_{\text{supervisor}} + 1\)).
  \item
    Else if execution elapsed time \(t > T_{\max}\): record timeout
    failure (\(Y_i \leftarrow 0\);
    \(n_{\text{timeout}} \leftarrow n_{\text{timeout}} + 1\)).
  \item
    Else if bus dropped or hardware faulted: record hardware failure
    (\(Y_i \leftarrow 0\);
    \(n_{\text{fault}} \leftarrow n_{\text{fault}} + 1\)).
  \item
    Append immutable trial metadata record to ledger \(\mathcal{L}\).
  \end{itemize}
\item
  \textbf{Stopping Oracle \& Certification (Stage 5):}

  \begin{itemize}
  \tightlist
  \item
    If \(Y_i = 0\) (any failure during zero-failure certification
    protocol): terminate evaluation early, compute Clopper-Pearson bound
    with failures
    (\(p_{\text{lower}} \leftarrow \text{BetaQuantile}(\alpha; n_{\text{success}}, n_{\text{attempted}} - n_{\text{success}} + 1)\)),
    and return
    \((\text{FAIL\_REJECTED}, p_{\text{lower}}, \mathcal{L})\).
  \item
    If \(i = N_{\text{req}}\) and
    \(n_{\text{success}} = N_{\text{req}}\): compute exact zero-failure
    lower bound
    \(p_{\text{lower}} \leftarrow \alpha^{1/N_{\text{req}}}\) and return
    \((\text{PASS\_CERTIFIED}, p_{\text{lower}}, \mathcal{L})\).
  \end{itemize}
\end{enumerate}

Chapter~\ref{sec-frontier} develops the complete safety case arithmetic
for claims that exceed any physical testing budget.

\section{What Evaluation Cannot Establish}\label{sec-evaluation-7-6}

In a conventional software benchmark, the test framework evaluates a
model against a fixed dataset whose distribution does not depend on the
predictions the model makes. When an image classifier labels a stop sign
as a speed limit sign, the test runner records the error, loads the next
image from disk, and continues along its predetermined sequence. In a
physical machine, the input distribution is endogenous because the
machine operates in a closed loop with the physical world.\sidenote{\footnotesize Under
  Signorini contact and friction cones, transitions are non-smooth. The
  state distribution \(d_t^\pi(s)\) is non-Lipschitz; small weight
  shifts cause discontinuous bifurcations in state support.} At time
step \(t\), the machine occupies a physical state
\(s_t \in \mathcal{S}\) and receives an observation \(o_t = h(s_t)\)
through its sensors. The learned policy selects an action
\(a_t \sim \pi(a \mid o_t)\), the actuators apply physical forces to the
body, and the environment transitions to a new state
\(s_{t+1} \sim \mathcal{T}(s_{t+1} \mid s_t, a_t)\). The observation at
the next step, \(o_{t+1} = h(s_{t+1})\), is directly caused by the
control action the policy selected at step \(t\). Every action alters
the physical state of the machine, and that state dictates what the
sensors observe next.

Unrolling this recursive interaction across a finite horizon of \(T\)
control steps yields the joint trajectory distribution induced by policy
\(\pi\):
\[p_\pi(s_0, a_0, s_1, a_1, \dots, s_T) = p(s_0) \prod_{t=0}^{T-1} \pi(a_t \mid h(s_t)) \, \mathcal{T}(s_{t+1} \mid s_t, a_t)\]
The marginal state occupancy distribution at step \(t\), denoted
\(d_t^\pi(s)\), is obtained by integrating over all preceding states and
actions:
\[d_t^\pi(s) = \int p_\pi(s_0, a_0, \dots, s_t = s) \, d s_0 \dots d s_{t-1} \, d a_0 \dots d a_{t-1}\]
The state occupancy distribution carries the policy \(\pi\) inside the
product at every preceding time step. Because the distribution of states
depends on the policy itself, two policies \(\pi_A\) and \(\pi_B\)
starting from the identical initial state \(s_0\) generate divergent
state distributions as soon as their actions differ. In an analytical
example where a mobile robot travels forward at \(2.0\text{ m/s}\) with
a control frequency of \(50\text{ Hz}\) (\(20\text{ ms}\) per step), a
steering discrepancy of \(0.05\text{ rad}\) at step \(t=0\) alters the
lateral position by \(2.0\text{ mm}\) at step \(t=1\). That displacement
changes the camera perspective, which alters the feature activations in
the neural network at step \(t=1\), producing further steering
divergences. After \(50\) control steps (\(1.0\text{ s}\) of physical
operation), the state distributions \(d_{50}^{\pi_A}\) and
\(d_{50}^{\pi_B}\) share little overlapping support. A static
pre-deployment evaluation benchmark, whether an offline log of expert
demonstrations or a fixed series of test injections, samples from an
arbitrary fixed distribution \(d_{\text{test}}(s)\) that matches neither
policy. Evaluating a policy against a fixed dataset measures prediction
error on states that an external logging policy visited, while providing
no evidence about the states that the evaluated policy will create when
placed in authority over the actuators.

Executing the policy in closed-loop simulation allows the model to
generate its own trajectory, but it replaces one distributional
assumption with another. In simulation, the next state is generated by
an approximate transition model
\(\widehat{\mathcal{T}}(s_{t+1} \mid s_t, a_t)\) and a synthetic sensor
model \(\widehat{h}(s)\). The simulator samples from the simulated
trajectory distribution \(p_{\pi, \widehat{\mathcal{T}}}\), not the
physical deployment distribution \(p_{\pi, \mathcal{T}}\). If the
simulator models surface traction with a constant Coulomb friction
coefficient of \(\mu = 0.80\) while the physical floor has a localized
oil patch with \(\mu = 0.25\), or if the simulated actuator responds
instantaneously while the physical motor exhibits a \(15\text{ ms}\)
electrical time constant in its windings, the physical trajectory
immediately diverges from \(\widehat{\mathcal{T}}\). Small mechanical
discrepancies compound across successive time steps, driving the policy
into physical states that the simulator never generated. Closed-loop
simulation proves that the policy stabilizes the simulator transition
model \(\widehat{\mathcal{T}}\). It cannot prove that the policy
stabilizes the physical machine operating under the mechanics of
\(\mathcal{T}\).

Because the deployed distribution is endogenous and tied to physical
mechanics, the strongest conclusion an empirical evaluation can support
is narrow and conditional. A pre-deployment evaluation establishes that
the specific software build and hardware revision tested achieved a
measured outcome rate under the exact sampled initial conditions, within
the specific environmental envelope present during the trial, subject to
the stated confidence interval. For an analytical example of \(500\)
completed evaluation runs yielding zero catastrophic incidents, the
arithmetic of one-sided interval estimation establishes with \(95\%\)
confidence that the true incident rate under those exact test conditions
is below \(1 - 0.05^{1/500} \approx 0.597\%\) per run (a failure
probability under \(6.0 \times 10^{-3}\)). Evaluation does not establish
that the failure rate remains below \(0.597\%\) when the ambient
temperature drops by \(10^\circ\text{C}\), when mechanical gear backlash
increases by \(0.2\text{ mm}\) through wear, or when optical lens
contamination degrades camera contrast by \(15\%\).

Accumulating larger sample counts reduces sampling uncertainty, which
tightens the statistical error bounds around the tested distribution,
but it cannot reduce structural uncertainty about unmodeled physics or
unvisited domains. If a mobile robot completes \(10{,}000\) consecutive
navigation trials inside a climate-controlled test facility with an
illumination of \(500\text{ lux}\), the resulting \(95\%\) lower
confidence bound confirms a failure probability below
\(3.0 \times 10^{-4}\) per trial in that facility. That statistical
confidence is mathematically valid for the tested distribution
\(d_{\text{test}}\), but it provides zero coverage for an unmodeled
condition such as a low-angle sun glare at \(45\text{ lux}\) entering
through an open doorway. More trials inside the test envelope merely
measure the known distribution with finer precision. They cannot repair
missing environmental variations, unmodeled actuator dynamics, or the
physical wear of mechanical bearings over the operating life of the
machine.

Every boundary of the evaluation claim must therefore be converted into
a runtime monitoring obligation inside the nervous system. Because
evaluation cannot establish how the policy behaves outside the sampled
distribution, the nervous system must continuously verify that the
physical machine and its environment remain within the envelope where
evaluation data exists. An autonomous machine requires five distinct
monitoring channels operating in real time:

\begin{enumerate}
\def\labelenumi{\arabic{enumi}.}
\tightlist
\item
  \textbf{Input Support / Out-of-Distribution (OOD) Monitor:} Checks
  whether incoming sensor observations fall within the support of the
  training and evaluation distribution, measuring sensory reconstruction
  error or feature embedding distance.
\item
  \emph{Execution Latency Monitor:} Measures the computation time of
  each inference pass, verifying that the \(P_{99}\) latency does not
  exceed the real-time control deadline (\(20\text{ ms}\) on a
  \(50\text{ Hz}\) control loop).
\item
  \emph{Constraint Margin Monitor:} Tracks the distance between the
  current physical state and hardware safety limits, including joint
  torque limits, motor winding temperature rated to
  \(85^\circ\text{C}\), and clearance to obstacles budgeted at a minimum
  of \(0.50\text{ m}\).
\item
  \emph{Intervention Rate Monitor:} Calculates the frequency of safety
  supervisor overrides, detecting when policy instability causes
  excessive corrections.
\item
  \emph{Environmental Assumption Monitor:} Measures external operating
  conditions, such as surface traction estimated from wheel slip and
  ambient illumination measured by photometric sensors.
\end{enumerate}

A runtime monitor that merely records an anomaly to a disk log provides
no protection to a moving body. When a monitor detects that a deployment
parameter has crossed outside the evaluated envelope, the nervous system
must act on the actuators by scaling down policy authority, switching to
a deterministic fallback controller, or executing a controlled stop.
Mathematically, this selective prediction mechanism is parameterized as
a runtime reject-option gate:
\[\hat{a}(o_t) = \begin{cases} \pi_\theta(o_t) & \text{if } u(o_t) \le \tau_{\text{reject}} \;\land\; \Delta t_{\text{infer}} \le \Delta t_{\max} \\ \pi_{\text{fallback}}(s_t) & \text{if } u(o_t) > \tau_{\text{reject}} \;\lor\; \Delta t_{\text{infer}} > \Delta t_{\max} \end{cases}\]
where \(u(o_t)\) is a calibrated epistemic uncertainty or conformal
non-conformity score, \(\Delta t_{\text{infer}}\) is the measured
inference latency, \(\Delta t_{\max}\) is the real-time deadline
(e.g.~\(20\text{ ms}\) at \(50\text{ Hz}\)), and
\(\pi_{\text{fallback}}\) is a deterministic, pre-certified safe stop or
low-speed compliance law. When execution latency exceeds the
\(20\text{ ms}\) deadline, the nervous system discards the overdue
inference output and commands a deterministic deceleration trajectory,
bringing the machine to rest within \(1.2\text{ m}\). When the input
support monitor detects that sensor observations have drifted into an
uncharacterized regime (\(u(o_t) > \tau_{\text{reject}}\)), the nervous
system restricts maximum speed to \(0.30\text{ m/s}\) and requests human
intervention.

Certain failure modes cannot be detected by runtime monitoring. When a
learned policy generates a control command that satisfies every latency
deadline, respects all joint torque limits, and operates on familiar
sensor inputs, but nevertheless represents an incorrect action, such as
steering into a glass partition that the vision backbone misidentified
as free space, the runtime monitors observe no violation. The input
appears supported, the computation finishes on schedule, and the state
margins appear clear until contact occurs. Because this semantic failure
mode produces no anomalous telemetry prior to impact, runtime monitoring
cannot detect it. Safety for undetectable failures cannot be established
through empirical evaluation or software monitors. It requires physical
zoning, mechanical current limiters, and independent safety barriers
whose guarantees do not depend on the correctness of the learned model.

\section{How Evaluation Goes Wrong}\label{sec-evaluation-7-7}

The systematic errors that undermine policy evaluation rarely stem from
deliberate fraud. They arise when an evaluation protocol measures a
well-chosen environment, terminates on a well-timed run, and tracks a
surrogate metric that answers an easier question than physical task
completion. In a physical trial, environment selection occurs whenever
the operating domain of the test bench is narrower than the claimed
deployment envelope. Consider a mobile robot whose navigation policy is
evaluated down a central \(3.0\text{ m}\) wide hallway with diffuse
overhead lighting of \(450\text{ lx}\), uniform concrete flooring, and
zero dynamic obstacles. When the evaluation report aggregates twenty
consecutive runs along this route into a single success rate, the claim
silently expands to cover the entire facility, including
\(1.2\text{ m}\) wide side corridors with direct afternoon glare and
dense pedestrian traffic. The policy succeeds on the bench because its
sensory representations are never tested outside a benign, high-contrast
slice of state space. Environment selection is detected by comparing a
structured condition ledger, which records measured illumination,
passage width, surface friction, and obstacle velocities for every
trial, directly against the boundaries of the claimed operational design
domain. If the ledger covers only a tight cluster within that envelope,
the evaluation establishes nothing about the remaining volume.

A second failure mode occurs when the evaluation metric drifts away from
physical task completion toward offline surrogate quantities. In
open-loop evaluation, a policy is scored on action agreement, measuring
the mean squared error or cosine distance between predicted control
commands and expert demonstrations recorded on a static dataset. A
policy can achieve an action agreement metric above \(98\%\) on recorded
sensor logs, yet fail immediately when given authority over a physical
machine. Action agreement does not account for the closed-loop
compounding of execution error. An unmodeled steering bias of
\(0.02\text{ rad}\) produces negligible loss on an offline validation
batch, but over a \(10\text{ m}\) traversal at \(1.5\text{ m/s}\), the
uncorrected lateral drift pushes the vehicle \(0.20\text{ m}\) off
course, colliding with a curb before completing the maneuver. Metric
drift is detected by requiring the primary evaluation metric to measure
closed-loop terminal outcomes, such as final pose error within
\(0.05\text{ m}\) and zero supervisor interventions, rather than
open-loop prediction error. When an evaluation relies on open-loop
metrics, it detects whether the network reproduced demonstration
vectors, not whether the closed-loop physical system can reach a goal
state.

Even when physical trials are conducted in closed loop, repeated tuning
against a fixed evaluation benchmark degrades the statistical
independence of the result. When an engineering team adjusts policy
hyperparameters, observation filtering, or neural network checkpoints
after every trial on a designated evaluation course, the nominal test
set ceases to be an independent holdout. Each configuration change
incorporates information about the specific obstacle positions, wall
textures, and lighting angles of that exact track. After forty design
iterations, the final model checkpoint reflects implicit training on the
test course itself, masking how the policy responds to novel geometry.
This test set contamination cannot be detected by inspecting the final
model weights or running a single verification trial on the same course.
It is detected exclusively through provenance tracking, analyzing the
chronological history of policy versions, tuning commits, and test dates
to confirm that the reported evaluation was executed on a fresh,
quarantined configuration that was never exposed to iterative tuning.

Evaluation protocols are equally vulnerable to unintended cherry-picking
through optional stopping. When an engineer monitors a continuous
sequence of physical runs and halts the experiment as soon as a target
performance threshold or statistical significance criterion is reached,
the reported confidence is distorted. If an evaluation team conducts an
experiment where each look carries a nominal false-positive threshold of
\(\alpha = 0.05\), stopping early upon observing an apparent success
across multiple evaluation checkpoints inflates the true error rate. If
the experiment is checked at five sequential intervals under the null
hypothesis of no genuine capability difference, the probability of
obtaining at least one spurious passing result is:
\[P(\text{Spurious Pass}) = 1 - (1 - 0.05)^5 = 1 - (0.95)^5 = 1 - 0.77378 \approx 0.2262 \implies 22.6\%\]
What appears to be an evaluation conducted at a \(5\%\) significance
level has degenerated into a protocol with more than a one-in-five
probability of false acceptance. Optional stopping cannot be detected
from summary tables or final pass rates. It is detected only by auditing
the pre-registered trial protocol, verifying that sample sizes, stopping
conditions, and evaluation milestones were fixed before the first
physical run began.

A related statistical distortion occurs when the denominator of an
evaluation is quietly filtered. In physical testing, runs frequently
terminate abnormally due to software timeouts, emergency stop
actuations, sensor communication drops, or safety supervisor
interventions. When an evaluation team discards these aborted runs as
invalid test artifacts, calculating success as the ratio of completed
goals to non-aborted trials, the reported metric misrepresents physical
reliability. If a mobile base completes nine runs out of ten attempts,
but the tenth attempt required an operator takeover \(4.0\text{ s}\)
into execution to prevent a collision with a wall, removing that attempt
raises the apparent success rate from \(90\%\) to \(100\%\). Every
initiated trial, operator takeover, emergency stop, and hardware timeout
must remain in the denominator. When physical constraints or safety
limits truncate a run before the nominal time limit of \(60\text{ s}\)
expires, the trial must be recorded explicitly as a censored failure
rather than omitted from the ledger. Denominator filtering is detected
by cross-referencing low-level actuator telemetry and safety supervisor
logs against the trial ledger, confirming that every power-on cycle and
autonomy engagement corresponds to a recorded outcome.

The \textbf{censored trial fallacy} is the systematic artificial
inflation of empirical reliability metrics caused by discarding aborted
runs, manual supervisor interventions, or emergency stop triggers from
the evaluation denominator, mistaking conditional survival on truncated
trajectories for genuine operational competence.

The temptation to substitute visual plausibility for statistical
sampling represents a pervasive vulnerability in physical systems
engineering. A single video recording of a robotic arm grasping an
irregular workpiece or a mobile robot navigating a cluttered room is a
demonstration, not a measurement. A visually convincing demonstration
proves that the system state trajectory intersects the success set for
at least one sequence of initial conditions and environmental
disturbances, but it provides no information about the volume of that
success set. Because human perception is tuned to smooth motion and
coordinated kinematics, an observer easily mistakes kinematic fluidity
for policy competence. An execution that appears deliberate and smooth
can operate within \(10\text{ mm}\) of a catastrophic kinematic
singularity or carry an unmeasured \(P_{95}\) tracking error that
produces a collision in nine out of ten repeats. A representative
demonstration has no statistical standing. A representative sample
requires randomized initial poses, stratified environmental conditions
spanning the full operational envelope, and a pre-calculated sample size
capable of bounding tail failure rates.

\protect\phantomsection\label{callout-incidentux2a-7.9}
\begin{fbxSimple}{callout-incident}{Landmark Engineering Incident:}{Tow-Tractor Sidewall Collision from Censored Acceptance Testing}
\phantomsection\label{callout-incident*-7.9}

\emph{System:} Autonomous \(450\text{ kg}\) industrial electric
tow-tractor with forward stereo vision, 2D planar lidar, and end-to-end
vision-action policy running at \(50\text{ Hz}\) in factory material
handling.

\emph{Incident:} High-speed sidewall collision against a steel
structural column at \(1.6\text{ m/s}\) in a production corridor,
shearing the sensor mast, bending the steering knuckle, and causing
\(\$14{,}000\) in hardware damage within 48 hours of commissioning.

\emph{Telemetry Snapshot:} At \(t = 0.000\text{ s}\), the vehicle enters
a \(2.1\text{ m}\) corridor (\(0.15\text{ m}\) side clearance);
low-angle solar glare (\(35\text{ lx}\) ambient with \(1800\text{ lx}\)
floor reflections) attenuates camera dynamic range; an unmodeled
\(+0.018\text{ rad}\) steering bias accumulates across \(50\text{ Hz}\)
cycles, drifting the chassis \(0.22\text{ m}\) laterally over
\(12\text{ m}\) (\(7.5\text{ s}\)); at \(t = 7.500\text{ s}\), the
vehicle strikes the steel pillar at \(1.6\text{ m/s}\) before obstacle
braking engages.

\emph{Root Cause:} Acceptance testing was conducted exclusively in a
\(4.5\text{ m}\) wide testing bay under uniform \(500\text{ lx}\)
lighting, achieving an artificial \(100\%\) pass rate (\(40/40\)) by
silently discarding two lidar communication timeouts and three manual
emergency stops as ``testbed anomalies.''

\emph{Architectural Flaw:} Absence of an immutable evaluation ledger
that mandates full-denominator accounting and an independent
hardware-in-the-loop lateral clearance envelope monitor in the nervous
system.

\emph{Provenance:} Factory automation logistics vehicle commissioning
post-mortem report and sensor log audit.

\end{fbxSimple}

\begin{figure}[t]

\centering{

\includegraphics[width=0.85\linewidth,height=\textheight,keepaspectratio]{/Users/VJ/GitHub/MLSysBook-integrate-physical/physical/book/chapters/07-evaluation/figures/fig07_real_crash_barrier_test.jpg}

}

\caption{\label{fig-07-crash-barrier-test}\emph{Rigid Barrier Impact
Testing and Destructive Physical Limits.} High-speed photometric capture
of a full-scale vehicle impact against a load-cell instrumented rigid
concrete crash barrier during dynamic qualification testing (NHTSA /
Calspan facility). High-contrast optical fiducials along the vehicle
flank track chassis deceleration profiles and structural crumple
dynamics at millisecond resolution. In Physical AI systems, physical
impact testing reveals the catastrophic consequence when learned
perception or control policies fail at dynamical boundaries: unmodeled
latency, perceptual blinding, or compounding steering drift rapidly
consumes spatial clearance margins, converting kinetic energy
(\(E_k = \frac{1}{2}mv^2\)) into structural deformation that software
runtime monitors cannot mitigate once physical contact occurs.}

\end{figure}%

As captured in full-scale impact barrier testing
(Figure~\ref{fig-07-crash-barrier-test}), physical kinetic energy
transfers into destructive structural deformation within milliseconds
when spatial clearance is lost. Pre-deployment evaluations that censor
near-misses or test only benign, high-contrast environments leave
physical machinery vulnerable to unrecoverable collisions the moment
real-world operating conditions diverge from the tested envelope.

Understanding the mechanisms by which evaluation fails clarifies the
purpose of physical testing. An empirical trial is not a demonstration
to convince an observer that a model works. It is an instrument to
measure where policy authority fails, what physical conditions trigger
the failure, and how much statistical confidence the evidence supports.
Documenting those boundaries requires an immutable record that binds
every trial, condition, and intervention to a verifiable physical
quantity.

\section{Evaluation Logs}\label{sec-evaluation-7-8}

Chapter~\ref{sec-nervous} established the baseline schema, cryptographic
hashing rules, and provenance invariants for physical run ledgers. An
\textbf{evaluation record} extends that ledger with the formal fields
required to audit a learned policy. It does not record developer
sentiment or subjective impressions of performance. Instead, it serves
as a self-contained certificate that states what was measured, under
what physical conditions, with what statistical power, and which
operational claims the resulting evidence explicitly does not
support.\sidenote{\footnotesize ANSI/UL 4600 requires structured safety cases pairing
  testing with explicit ODD boundaries and monitors. ISO 21448 (SOTIF)
  targets performance limitations in the absence of hardware faults.}

The \textbf{evaluation record} is the immutable engineering notebook
schema extension that registers target operational design domains,
hardware and model cryptographic hashes, complete-denominator trial
ledgers (including all aborts and overrides), and exact non-asymptotic
confidence bounds to certify policy verification boundaries.

\begin{longtable}[]{@{}
  >{\raggedright\arraybackslash}p{(\linewidth - 6\tabcolsep) * \real{0.2500}}
  >{\raggedright\arraybackslash}p{(\linewidth - 6\tabcolsep) * \real{0.2500}}
  >{\raggedright\arraybackslash}p{(\linewidth - 6\tabcolsep) * \real{0.2500}}
  >{\raggedright\arraybackslash}p{(\linewidth - 6\tabcolsep) * \real{0.2500}}@{}}
\caption{\emph{Physical Policy Pre-Deployment Evaluation Record Schema.}
Immutable schema fields added to the base engineering notebook to record
target operational boundaries, hardware/model cryptographic hashes,
complete-denominator trial ledgers, non-asymptotic statistical
confidence bounds, and unsupported coverage
gaps.}\label{tbl-07-evaluation-record-schema}\tabularnewline
\toprule\noalign{}
\begin{minipage}[b]{\linewidth}\raggedright
Record Schema Field
\end{minipage} & \begin{minipage}[b]{\linewidth}\raggedright
Physical / Statistical Modality
\end{minipage} & \begin{minipage}[b]{\linewidth}\raggedright
Cryptographic \& Provenance Guarantee
\end{minipage} & \begin{minipage}[b]{\linewidth}\raggedright
Downstream Verifier \& Usage
\end{minipage} \\
\midrule\noalign{}
\endfirsthead
\toprule\noalign{}
\begin{minipage}[b]{\linewidth}\raggedright
Record Schema Field
\end{minipage} & \begin{minipage}[b]{\linewidth}\raggedright
Physical / Statistical Modality
\end{minipage} & \begin{minipage}[b]{\linewidth}\raggedright
Cryptographic \& Provenance Guarantee
\end{minipage} & \begin{minipage}[b]{\linewidth}\raggedright
Downstream Verifier \& Usage
\end{minipage} \\
\midrule\noalign{}
\endhead
\bottomrule\noalign{}
\endlastfoot
\textbf{\texttt{target\_odd\_boundary}} & Operational Design Domain
bounds: passage width \(w \ge 1.5\text{ m}\), illuminance
\([200, 1000]\text{ lx}\), surface friction \(\mu \ge 0.60\), payload
\(m \in [0, 15]\text{ kg}\) & Formal system specification SHA-256 hash,
UL 4600 / ISO 26262 boundary declaration & Enforcing automated gating;
setting bounds on runtime OOD anomaly detectors in nervous system \\
\textbf{\texttt{component\_hashes}} & Frozen neural network weights
(\(H(\theta)\)), inference engine binary, chassis serial number, tire
batch, sensor calibration matrices, and firmware build & SHA-256 binary
digests, hardware UUID, ISO 9001 calibration certificate references &
Guaranteeing physical repeatability; preventing unrecorded hardware or
driver substitutions \\
\textbf{\texttt{sampling\_protocol}} & Stratification matrix across
payload/illumination, PRNG seed sequence for initial poses, reset
procedure, and trial independence audit & Automated test harness
configuration script hash, telemetry sequence log & Auditing against
optional stopping and testing against hidden autocorrelation across
trials \\
\textbf{\texttt{complete\_outcome\_ledger}} & Full-denominator raw
outcome counts: \(n_{\text{attempted}}\), \(n_{\text{success}}\),
\(n_{\text{estop}}\), \(n_{\text{timeout}}\), \(n_{\text{supervisor}}\),
and \(T_{\text{TTI}}\) time-series & Monotonic microsecond hardware
timestamps, uninterrupted ring-buffer telemetry & Preventing censored
trial fallacy; calculating uncorrupted empirical reliability \\
\textbf{\texttt{statistical\_certification}} & Exact Clopper-Pearson
lower bound (\(p_{\text{lower}} = \alpha^{1/n}\)), nominal confidence
\(1-\alpha\), and Poisson continuous exposure rate \(\lambda\) &
Cryptographic verification signature, reproducible Python evaluation
script & Establishing mathematical ceiling of empirical testing for
qualification cases in Chapter 15 \\
\textbf{\texttt{unsupported\_gap\_catalog}} & Explicit catalog of
unvisited state spaces, omitted dynamics (surface water, dynamic
crowds), and required runtime monitoring channels & Declared non-support
manifest, signed safety engineering sign-off & Parameterizing real-time
safety supervisors, CBF guardrails, and fallback policies in Part III \\
\end{longtable}

As specified in the pre-deployment evaluation record schema
(Table~\ref{tbl-07-evaluation-record-schema}), the evaluation record
begins with the operational claim and its formal predicates. It defines
the mathematical success criterion, such as completing a trajectory
between dock stations while maintaining a lateral lane error below
\(0.05\text{ m}\) and an end-point velocity below \(0.1\text{ m/s}\). It
bounds the operational design domain across physical parameters,
including ambient illuminance, surface friction coefficients, payload
mass distributions, and obstacle velocities. Finally, it names the
target population of physical environments to which the result is
intended to generalize, distinguishing between the specific warehouse
bays tested and the broader class of industrial facilities. Defining
these boundaries in the record before trials begin prevents post-hoc
adjustments to the pass criteria when anomalous behaviors appear.

To enable independent replication of the physical execution, the record
captures immutable cryptographic identifiers for every software and
hardware component. These identifiers include the content hashes of the
neural network weights, sensor preprocessing pipelines, policy
configuration files, simulation engine builds when synthetic runs are
incorporated, and the evaluation execution scripts. On the physical
system, the record logs the machine chassis serial number, wheel tire
compound batch, sensor calibration matrices, motor controller firmware
revisions, and communication bus configurations. Swapping a camera lens,
altering a driver firmware version, or replacing worn drive belts
changes the physical transfer function of the machine, invalidating the
historical data if the change is unrecorded.

The record preserves the trial-sampling protocol and its operational
execution. It documents the stratification strategy used to cover the
operating envelope across payload increments from \(0.0\text{ kg}\) to
\(15.0\text{ kg}\), the pseudorandom seed sequence governing initial
robot poses, and the reset sequence executed between consecutive runs.
When a mobile base returns to a starting waypoint under autonomous
control, battery thermal rise and odometry drift introduce statistical
dependence between adjacent attempts. The record specifies whether
trials were statistically independent or chained in series, and it
registers every manual reset, hardware abort, or deviation from the
pre-specified protocol.

Aggregate percentages obscure the physical failure modes that determine
system risk. The evaluation record reports raw outcome counts across all
categories, separating completed trials, kinematic boundary violations,
safety supervisor torque overrides, communication timeouts, and manual
emergency stops. In addition to summary tallies, the record preserves
per-trial time-series metadata, including peak motor winding
temperatures, maximum trajectory tracking errors, and task completion
durations. Logging this telemetry ensures that nominal successes
operating near thermal limits or within \(5\text{ mm}\) of a hard
boundary are visible during audit. Every trial initiated under policy
authority remains in the denominator.

The statistical specification defines the point estimator, the chosen
confidence method, the nominal confidence level, and the underlying
mathematical assumptions. Alongside the supported claim, the record
documents unsupported claims and known coverage gaps. If an evaluation
suite tested navigation on dry concrete at an ambient temperature of
\(20^\circ\text{C}\), the record explicitly notes that it provides zero
empirical support for performance on wet epoxy flooring, at temperatures
below \(5^\circ\text{C}\), or in the presence of unmapped obstacles
taller than \(0.15\text{ m}\).

Consider an analytical example using the twenty-trial evaluation
introduced earlier in this chapter. A policy controls a
\(25.0\text{ kg}\) differential-drive robot navigating a
\(30.0\text{ m}\) indoor corridor at \(1.5\text{ m/s}\). The test
fixture executes twenty trials across clean concrete with ambient
illumination between \(250\text{ lx}\) and \(350\text{ lx}\), recording
twenty completed runs with zero safety supervisor interventions. An
engineer examining this record recomputes the exact one-sided
Clopper-Pearson binomial lower bound at \(95\%\) confidence (see
Table~\ref{tbl-07-confidence-bounds}), obtaining
\(0.05^{1/20} \approx 0.861\), or a success rate of \(86.1\%\).

The completed record exposes the structural limits of the claim. The
lower bound of \(86.1\%\) means the evidence permits a failure rate as
high as \(13.9\%\), which corresponds to roughly one failure every seven
trials under sustained operation. When the engineer checks the
operational parameters in the record, the documentation confirms that
the payload mass was fixed at \(0.0\text{ kg}\) and ambient illuminance
was confined to a narrow \(100\text{ lx}\) interval. The evaluation
record proves that twenty successful runs cannot establish high
reliability, and it provides the quantitative basis to reject any
assertion that the policy is ready for unmonitored operation.

An evaluation record does not certify that an autonomous machine is
safe. It establishes the empirical boundaries and statistical
uncertainty of a single learned component operating under controlled
conditions. In Chapter~\ref{sec-verification}, the verification argument
ingests these evaluation records as formal evidence nodes, combining
policy bounds with physical guardrail guarantees and hardware supervisor
latencies to evaluate whether the integrated machine satisfies its
operational safety case.

\section{Systems Stress-Tests}\label{sec-evaluation-7-9}

An evaluation record does not certify that an autonomous machine is
safe. It establishes the empirical boundaries and statistical
uncertainty of a single learned component operating under controlled
conditions. In physical systems, an evaluation that is honest and
well-run can easily support nothing if the statistical and physical
assumptions connecting trial metrics to operational claims are broken.
The following four analytical stress-tests evaluate evaluation regimes
that appear rigorous while concealing systematic evidence deficits:

\subsection*{Stress-Test 1: The Homogeneous Route Illusion (Overfitting
Narrow Environmental
Slices)}\label{stress-test-1-the-homogeneous-route-illusion-overfitting-narrow-environmental-slices}
\addcontentsline{toc}{subsection}{Stress-Test 1: The Homogeneous Route
Illusion (Overfitting Narrow Environmental Slices)}

\emph{The scenario.} An autonomous delivery rover
(\(m = 35.0\text{ kg}\)) executes waypoint navigation at
\(1.8\text{ m/s}\) along a \(120.0\text{ m}\) asphalt path. An
evaluation team logs 20 consecutive zero-failure trials on a clear
Tuesday afternoon (\(T = 22.0^{\circ}\text{C}\), dry surface, fixed
\(5.0\text{ kg}\) payload). The team submits an engineering report
asserting campus-wide deployment readiness.

\emph{The systems dilemma.} Under identical dry pavement,
\(5.0\text{ kg}\) payload, and constant lighting, 20 zero-failure trials
establish a one-sided Clopper-Pearson lower bound of
\(p_{\text{lower}} = 0.05^{1/20} \approx 86.1\%\) at \(95\%\)
confidence. Asserting campus-wide readiness constitutes an unverified
extrapolation across unmeasured physical regimes: wet pavement
(\(\mu < 0.3\)), temperature extremes
(\(0^{\circ}\text{C}\text{ to }35^{\circ}\text{C}\)), payload variations
(\(0\text{--}15\text{ kg}\)), and dynamic obstacles.

\emph{The systems verdict: Operational Envelope Support Law.} A policy
evaluated across a narrow slice of physical conditions provides zero
empirical evidence for states where physical dynamics shift. Slicing
trials across environmental variations rather than accumulating
repetitive nominal runs is required to support deployment claims.

\subsection*{Stress-Test 2: The Statistical Tail Deficit in
High-Assurance Claims (The Finite-Sample
Trap)}\label{stress-test-2-the-statistical-tail-deficit-in-high-assurance-claims-the-finite-sample-trap}
\addcontentsline{toc}{subsection}{Stress-Test 2: The Statistical Tail
Deficit in High-Assurance Claims (The Finite-Sample Trap)}

\emph{The scenario.} A learned pick-and-place policy on a 6-axis
manipulator transfer packages at \(0.833\text{ Hz}\). The operational
requirement mandates a minimum success probability of \(99.9\%\)
(\(p \ge 0.999\)) at \(95\%\) confidence. An automated test rig runs
1,000 independent trials under randomized orientations, recording 1,000
successes and zero drops. The product team claims the \(99.9\%\)
specification is met.

\emph{The systems dilemma.} The exact one-sided Clopper-Pearson \(95\%\)
lower bound for zero failures in 1,000 trials is
\(p_{\text{lower}} = 0.05^{1/1000} \approx 0.9970\) (\(99.70\%\)). A
\(99.70\%\) lower bound permits a failure rate of \(0.30\%\) (3 failures
per thousand cycles)---triple the allowable failure budget. To
mathematically certify \(p \ge 0.999\) at \(95\%\) confidence without
failures, solving \(0.999^n = 0.05\) yields
\(n = \ln(0.05)/\ln(0.999) \approx 2{,}995\text{ trials}\). The test
fixture is 1,995 trials short of the required mathematical evidence.

\emph{The systems verdict: Finite-Sample Reliability Law.} Large trial
counts in software deceive systems engineers on physical hardware.
Substantial probability mass remains in the unobserved tail, requiring
formal exact confidence bounds rather than point-estimate assertions.

\subsection*{Stress-Test 3: The Open-Loop Deception in Closed-Loop
Navigation (Error
Compounding)}\label{stress-test-3-the-open-loop-deception-in-closed-loop-navigation-error-compounding}
\addcontentsline{toc}{subsection}{Stress-Test 3: The Open-Loop Deception
in Closed-Loop Navigation (Error Compounding)}

\emph{The scenario.} An autonomous mobile robot (\(m = 80.0\text{ kg}\),
track width \(0.60\text{ m}\)) navigates factory corridors at
\(2.0\text{ m/s}\). During offline validation against 50,000 expert
state-action pairs, the policy achieves \(98.8\%\) action agreement with
mean squared steering error of \(0.004\text{ rad}^2\). In closed-loop
physical deployment, the robot collides with structural pillars within
\(8.0\text{ m}\) of travel in 3 of 5 attempts.

\emph{The systems dilemma.} In offline evaluation, error is measured
against expert trajectories where errors do not compound. In closed
loop, the policy visits states generated by its own past actions. A tiny
systematic steering bias of \(+0.015\text{ rad}\) evaluated at
\(50\text{ Hz}\) (\(20\text{ ms}\) intervals) rapidly rotates the
vehicle heading, inducing \(>0.15\text{ m}\) lateral drift within
\(5.0\text{ m}\) of travel and driving the chassis into walls.

\emph{The systems verdict: Closed-Loop Decoupling Law.} Open-loop
validation error measures mimicry on static records where the model
cannot affect the next state. In physical AI, policy outputs govern the
plant state, and small systematic biases accumulate into collision.
Closed-loop physical stability cannot be inferred from open-loop offline
agreement.

\subsection*{Stress-Test 4: The Robustness Attribution Fallacy (Hardware
Enforcer vs Policy
Weights)}\label{stress-test-4-the-robustness-attribution-fallacy-hardware-enforcer-vs-policy-weights}
\addcontentsline{toc}{subsection}{Stress-Test 4: The Robustness
Attribution Fallacy (Hardware Enforcer vs Policy Weights)}

\emph{The scenario.} A planar robot powered by a heterogeneous
dual-brain platform (Linux MPU generating \(50\text{ Hz}\) trajectory
proposals and a deterministic real-time MCU evaluating safety barriers
at \(1000\text{ Hz}\)) achieves \(100\%\) task completion across 50
nominal runs. The team claims the learned policy is robust to transport
jitter. During qualification, an unannounced \(15.0\text{ ms}\) IPC
mailbox stall is injected while moving at \(0.8\text{ m/s}\) toward a
\(5.0\text{ mm}\) workspace boundary.

\emph{The systems dilemma.} The neural policy on the MPU possesses zero
internal awareness of transport stalls. When the \(15.0\text{ ms}\)
stall strikes, the proposal freshness lease
(\(\tau_{\text{lease}} = 10.0\text{ ms}\)) expires. The deterministic
MCU detects the stale lease and asserts local braking at
\(t = 1.210\text{ s}\), arresting motion with \(3.2\text{ mm}\) of
clearance. The survival of the machine is an architectural property of
the deterministic permission referee, not a capability of the neural
weights.

\emph{The systems verdict: Architectural Attribution Law.} Physical
fail-safety is guaranteed by the deterministic real-time supervisor,
never by the adaptability of learned neural weights. Unannounced
physical fault injection on testbed hardware separates genuine
supervisor guarantees from ungrounded policy assumptions.

\section{Summary}\label{sec-evaluation-7-10}

A measured success rate does not describe an isolated policy. When an
evaluation bench records that a robotic system completed ninety-eight
out of one hundred navigation runs, that metric measures the entire
physical assembly operating under a specific set of test conditions. The
success rate reflects sensor calibration, track friction, link flexure,
battery charge state, operating temperature, execution latency, and
policy weights operating as a coupled system. Modifying any single
mechanical or environmental parameter changes the measured metric. An
identical neural network deployed on a chassis with slightly degraded
wheel damping or evaluated under ambient lighting that increases image
sensor exposure times by \(10\text{ ms}\) can fail on trajectories it
previously traversed without error. A benchmark score is not an
intrinsic property of a weight file, but an empirical measurement of a
closed loop operating within defined operational boundaries.

Offline validation error cannot establish closed-loop physical
stability. In an offline dataset, a model evaluates predictions against
static state-action pairs recorded during prior runs. In that regime, an
action error is an isolated mathematical residual. In closed-loop
execution, the machine executes the policy action, moves through space,
and reaches a physical state that depends entirely on its own previous
commands. An offline mean angular error of \(0.015\text{ rad}\) appears
negligible on a static validation curve, but applied continuously at
\(50\text{ Hz}\) across \(20\text{ ms}\) control intervals, that
systematic bias rotates the heading vector and drives a mobile base
across a \(0.15\text{ m}\) lateral deviation in several meters of
travel. High accuracy on expert demonstrations confirms only that the
model mimics human data on human-selected paths. It does not establish
that the closed-loop system will recover once its own actions perturb it
away from the training distribution.

Finite trial counts severely restrict the strength of any reliability
claim, and correlated trials weaken that evidence further. Demonstrating
twenty consecutive successful experimental trials without a failure
yields an upper bound on failure probability of approximately \(14\%\)
at a \(95\%\) confidence level under independent and identically
distributed assumptions. If those twenty runs are conducted on the same
vehicle, in the same facility, over the course of an afternoon, they are
not independent. Shared battery discharge curves, static thermal
gradients in motor windings, and consistent lighting introduce
correlations that reduce the effective statistical sample size. A test
suite of fifty correlated runs may provide less evidentiary support than
five truly independent trials conducted across varied operational
envelopes. To verify that a system meets an industrial failure budget of
one incident per \(1000\text{ operating hours}\), an engineering team
cannot rely on unscaled physical sampling.

No evaluation regime drawing from a fixed distribution can determine how
a machine behaves under conditions its own policy generates over time.
In standard supervised machine learning, test instances are drawn from
an environment that remains unaffected by model inference. In physical
systems, policy decisions actively modify the environment and the
mechanical condition of the machine. Actuator commands generate heat,
induce gear backlash, wear down tire treads, and reposition objects
within the workspace. Over extended deployments, these physical changes
shift the sensor observation distribution into regions that were never
sampled during training or acceptance testing. Because a learned policy
continuously shapes the operational conditions it encounters,
pre-deployment testing establishes only a baseline on known
distributions, leaving long-term interaction dynamics unverified.

Part III receives a policy whose limits are known. The system
engineering task does not treat the learned model as an unconstrained
controller or assume that empirical testing has eliminated the
probability of severe failure. Instead, the policy arrives accompanied
by four required engineering deliverables:

\begin{enumerate}
\def\labelenumi{\arabic{enumi}.}
\tightlist
\item
  \emph{Bounded Operational Claim:} Specifying the precise operational
  design domain (ODD) boundaries where performance was empirically
  characterized.
\item
  \emph{Reproducible Evaluation Record:} Documenting the complete
  physical and computational configuration of every test, including
  cryptographic digests of weights, firmware, and calibration matrices.
\item
  \emph{Explicit Coverage Gap Catalog:} Detailing unverified corner
  cases, omitted environmental variations, and known failure modes.
\item
  \emph{Runtime Monitoring Specifications:} Defining the telemetry
  channels, latency bounds, OOD reject gates, and divergence thresholds
  required to detect anomalous behavior at runtime.
\end{enumerate}

\protect\phantomsection\label{callout-takeawaysux2a-7.10}
\begin{fbxSimple}{callout-takeaways}{Key Chapter Takeaways}{}
\phantomsection\label{callout-takeaways*-7.10}

\begin{itemize}
\tightlist
\item
  \emph{Success rates belong to closed-loop systems, not weights files.}
  A measured \(98\%\) completion rate reflects the coupled dynamics of
  sensor exposure latency, tire friction, battery charge state, and
  mechanical backlash. Modifying any single physical parameter
  invalidates the benchmark score.
\item
  \emph{Closed-loop feedback unmasks compounding error.} Open-loop
  validation error (\(0.015\text{ rad}\)) evaluates static predictions
  without state perturbation. In closed-loop execution at
  \(50\text{ Hz}\), uncorrected systematic steering errors compound into
  unrecoverable lateral drift and structural collisions within meters of
  travel.
\item
  \emph{Finite sample counts underestimate tail risk (\(3/N\) rule).}
  Achieving zero failures in \(N=1000\) trials proves only that failure
  probability \(p \le 3/N \approx 0.3\%\) at \(95\%\) confidence---a
  statistical gap that fails high-assurance industrial safety budgets
  (\(p \ge 99.9\%\), requiring \(2995\) zero-failure runs).
\item
  \emph{Physical fail-safety is an architectural property, not a model
  capability.} When injected communication stalls or out-of-distribution
  observations occur, safety is guaranteed exclusively by the
  deterministic real-time supervisor (\(1000\text{ Hz}\) MCU), never by
  the adaptability of learned neural weights.
\item
  \emph{The evaluation deliverable quadrant.} Every trained policy
  passing to runtime deployment must be accompanied by four formal
  engineering artifacts: Bounded Operational Claim (ODD), Reproducible
  Evaluation Record, Explicit Coverage Gap Catalog, and Runtime
  Monitoring Specifications.
\end{itemize}

\end{fbxSimple}

Engineering the surrounding machine requires designing directly for
these limitations. Because empirical evaluation cannot guarantee
stability across unvisited states or prevent compounding actuation
drift, the nervous system must provide the deterministic timing, safety
enforcement, and physical fallbacks that keep the machine safe. Part III
turns to the physical architecture, building the real-time compute
platforms, sensor pipelines, and hardware supervisors that govern
execution when a learned model operates mass in motion.

\part{Running the Machine}

\lettrine{B}{uilding} a physical AI machine is fundamentally an exercise
in \textbf{reconciling disparate timescales, representations, and
failure domains}.

No single algorithm, neural model, or silicon core can span the vast
continuum between high-level human reasoning and low-level physical
dynamics. At the deliberative tier, human intent and task semantics are
linguistic and symbolic, evolving over seconds
(\(0.1\text{--}1\text{ Hz}\)). In the middle tier, physical space is
continuous and metric, requiring 3D spatial representations and smooth
kinematic trajectories generated at tens of hertz
(\(20\text{--}50\text{ Hz}\)). At the execution tier, motor torque
control is electromagnetic and immediate, demanding cycle-deterministic
execution at kilohertz frequencies (\(1000\text{ Hz}\)) on bare-metal
silicon to prevent mechanical resonance and kinetic destruction.

Part III details \textbf{The Living Organism: The Millisecond Pipeline}
that transforms raw physical sensor photons into coordinated motor
currents:

\begin{itemize}
\tightlist
\item
  \textbf{From Photons to Spatial Claims (Chapter 8: Perception):} How
  raw, uncompressed sensor streams from SerDes lanes are ingested via
  DMA ring buffers, corrected for rolling-shutter distortion, calibrated
  against IMUs, and lifted into calibrated metric \(SE(3)\) spatial
  claims under strict information-age budgets.
\item
  \textbf{Preserving Belief Under Occlusion (Chapter 9: Memory):} How
  spatial representations (transform trees, TSDF voxel grids, 3D
  Gaussian Splats) persist when objects pass behind obstacles or cameras
  pan away, enforcing explicit temporal decay rates to prevent ``frozen
  transform'' collisions.
\item
  \textbf{From Semantic Goals to Reachable Affordances (Chapter 10:
  Intent):} How open-vocabulary foundation models decompose human
  language into spatial affordances, issuing cryptographically bounded,
  expiring \textbf{Intent Leases} with hard time-to-live (TTL) envelopes
  that fail-safe on silence.
\item
  \textbf{Synthesizing Continuous Motion (Chapter 11: Planning):} How
  continuous generative models predict \(H\)-step action chunks, smooth
  kinematic transitions across chunk boundaries using \(C^2\) jerk
  splines, and guarantee safe stopping suffixes when upstream inference
  experiences latency spikes.
\item
  \textbf{The Deterministic Gatekeeper (Chapter 12: Enforcement):} The
  bare-metal 1 kHz safety shield. We implement Control Barrier Functions
  (CBFs) and quadratic programming (QP) projection on dedicated
  microcontrollers, enforcing the Four-Tier Fallback Ladder and the
  Simplex Architecture Pattern to retain exclusive veto authority over
  power stages.
\item
  \textbf{The Silicon Substrate (Chapter 13: Placement):} How
  high-throughput neural deliberation and hard-real-time safety co-exist
  on shared heterogeneous System-on-Chip (SoC) platforms. We quantify
  on-chip memory crossbar QoS, DRAM bank contention, power distribution
  network (\(L \frac{dI}{dt}\)) voltage droop, and hysteretic thermal
  clock throttling.
\end{itemize}

Part III teaches students how to architect, synchronize, and isolate the
six computational stages that keep high-level machine intelligence
grounded in physical reality.

\chapter{\texorpdfstring{Sensor
Perception\chaptersubtitle{What an Observation Must Carry}}{Sensor Perception}}\label{sec-perception}

\begin{figure}[H]

\centering{

\includegraphics[width=1\linewidth,height=\textheight,keepaspectratio]{index_files/mediabag//Users/VJ/GitHub/MLSysBook-integrate-physical/physical/book/chapters/08-perception/figures/fig_locator.pdf}

}

\caption{\label{fig-locator}\textbf{Systems Organ Focus: Brain.}
Transduction to memory as a temporal progression, illustrating
cumulative latency along the perception pipeline from CMOS photodiodes
across MIPI CSI-2 serialization, DMA transfer, ISP debayering, and
neural backbone evaluation. Attribution: Textbook team (CC BY 4.0).}

\end{figure}%

\emph{When was this taken, where from, and how wrong might it be?}

A physical machine cannot act directly on raw photons or acoustic
reflections. It must translate asynchronous, noisy sensor streams into
calibrated metric claims in 3D space before physical clearance is
consumed. This transformation is bound by physical latency: while an
image frame is serialized over a MIPI CSI-2 bus and evaluated by a
vision backbone, the physical world continues to move. Mechanical
vibration, rolling-shutter scanline skew, and lens distortion degrade
geometric accuracy. Reliable perception requires bounding end-to-end
information age, attaching covariance ellipsoids to every spatial claim,
and packaging observations into immutable records with strict health
contracts.

\newpage{}

\protect\phantomsection\label{callout-objectiveux2a-8.1}
\begin{fbxSimple}{callout-objective}{Learning Objectives}{}
\phantomsection\label{callout-objective*-8.1}

After studying this chapter, you will be able to:

\begin{enumerate}
\def\labelenumi{\arabic{enumi}.}
\tightlist
\item
  State the three properties every observation must carry and explain
  what breaks when each is missing.
\item
  Account for the age accumulated between transduction and use across
  the latency waterfall, stage by stage.
\item
  Trace visual sensory processing from CMOS pixel arrays across MIPI
  CSI-2 DMA ring buffers on the Linux MPU to 3D foundation model
  backbones.
\item
  Model rolling shutter scanline skew and camera-to-IMU temporal and
  extrinsic calibration from first principles.
\item
  Compute what a capture configuration costs in memory bus bandwidth and
  name what it takes that from.
\item
  Name the failure modes of perception and identify which are detectable
  at runtime versus silent.
\end{enumerate}

\emph{Core chapter takeaway.} An observation without an age, a frame,
and an uncertainty is not usable, and the cost of producing one is taken
from the same bus everything else is competing for.

\end{fbxSimple}

\section{Photons to Spatial Claims}\label{sec-perception-8-1}

\lettrine{P}{art}\lettrine{I}{I} established what a learned policy can
deliver and what its evaluation record can prove. That record documents
empirical performance across a test distribution, bounds the
distribution shifts under which the policy was evaluated, and lists the
latency and compute budgets required to execute a forward pass. Yet a
policy residing in memory does not act on the physical world in
isolation. To govern physical mass, it must receive measurements that
reflect true mechanical states and issue commands that execute through
real actuators under strict temporal deadlines. The architectural
contracts developed across Part III take that evaluation record as their
baseline input. The task of the system architecture is to wrap the
learned model in the sensing pipelines, synchronization mechanisms,
state representations, and execution runtimes necessary to convert
statistical inference into safe physical movement.

A physical architecture can take many concrete shapes. An end-to-end
vision-to-action model maps raw camera bitstreams directly to joint
torque commands within a single neural network. A classical modular
pipeline partitions that same computation across independent processes
for sensor ingestion, state estimation, mapping, trajectory generation,
and motor control. While these implementations look entirely distinct on
a software architecture diagram, they represent the same machine
analyzed at different levels of granularity. Perception, state
estimation, spatial reasoning, planning, tracking, and actuation are not
mandatory software modules that must run in separate address spaces.
They are fundamental responsibilities that physics requires the machine
to discharge somewhere along the signal path between physical stimulus
and physical response.\sidenote{\footnotesize Marr (1982) decomposed vision into
  theory, algorithm, and hardware. In Physical AI, perception is a
  real-time pipeline where transport latency directly erodes stability
  margins.}

When an engineering team consolidates these responsibilities into an
end-to-end model, the stages do not vanish. They are absorbed into the
weight matrix, where they operate implicitly across hidden layers. What
changes is what the system engineer can inspect, bound, expire, or
refuse. In an explicit pipeline, an engineer can verify that an
intermediate coordinate transform satisfies rigid-body geometry, attach
an expiration deadline to a stale feature vector, or intercept an
estimation error before it reaches the control loop. When a stage is
implicit, intermediate representations cannot be isolated or
independently measured. If an implicit perception representation
confuses specular reflection on a wet floor for open clearance, the
downstream layers act on that representation directly, with no
intermediate interface capable of detecting the error or invalidating
the data before a command reaches an actuator.

This distinction identifies the single responsibility that can never be
absorbed into a learned policy. While perception, estimation, and
nominal trajectory generation can all be learned from data, enforcement
must remain external and deterministic. A learned model is inherently
statistical; its evaluation record guarantees performance in expectation
over a test set, but provides no absolute mathematical invariant for any
single inference cycle. When a policy encounters an out-of-distribution
input, suffers an internal numerical failure, or stalls on a compute
bottleneck, the authority to halt the machine, clamp actuator current,
or invoke a reflexive fallback must rest with a separate execution
monitor.

In modern physical AI systems, this division of responsibility is
physically embodied in a \textbf{dual-brain architecture}. The
high-bandwidth, unprivileged sensory perception stack---ingesting
multi-camera streams via Direct Memory Access (DMA) ring buffers,
running hardware image signal processors (ISPs), and evaluating
billion-parameter 3D vision foundation models---resides on a
high-performance Linux Microprocessor Unit (MPU) or System-on-Chip (SoC)
accelerator. Concurrently, a separate, real-time Microcontroller Unit
(MCU) runs deterministic safety monitors, inertial strapdown filters,
and control barrier functions on bare metal. The MPU emits candidate
perceptual interpretations and unprivileged motion proposals; the
real-time MCU enforces physical clearance invariants and holds the
unconditional hardware authority to refuse or override them. That
enforcement boundary is the structural invariant that Part III traces
from initial sensor ingress down to motor commutators.

The architectural pipeline begins at the boundary where physical energy
converts into digital data. As illustrated in the systems architectural
locator (Figure~\ref{fig-locator}), perception constitutes the first
structural handoff in the sensory-to-actuation chain, tracking the
timeline from transduction at the detector array to state estimation in
memory while accounting for the information age accumulated at each
step. Perception determines what the machine is permitted to believe
about its immediate environment and precisely when that physical state
existed. For downstream computation to consume a sensory measurement
without causing mechanical instability, every observation must carry
three distinct properties. It must carry an acquisition timestamp that
records when the physical phenomenon occurred, an identified coordinate
frame that fixes the measurement in physical space, and an uncertainty
estimate that bounds the reliability of the signal. Without these three
metadata fields, a tensor of detected objects, foundation model patch
embeddings, or surface normals is not a physical measurement; it is an
ungrounded array of numbers.

The need for these properties becomes concrete when tracing the latency
accumulated across an embedded perception pipeline. Consider an
autonomous mobile base operating in an industrial facility at
\(2.0\text{ m/s}\). Photons striking the photodiode array of a CMOS
camera do not instantly appear as feature embeddings in an inference
engine. In an analytical breakdown of a typical embedded vision system,
the sensor array requires \(33\text{ ms}\) for exposure and pixel
readout at \(30\text{ Hz}\). Direct memory access transfers across a
serial bus into system memory add \(8\text{ ms}\). Operating system
kernel scheduling and driver buffer handling introduce \(4\text{ ms}\)
of queuing delay. Image debayering, rectification, and color correction
require \(5\text{ ms}\) of preprocessing time. Finally, running a vision
backbone on an onboard accelerator requires \(25\text{ ms}\) at
\(P_{99}\) latency.

By the time the perception stage yields an observation tensor,
\(75\text{ ms}\) of total latency have accumulated between photon
capture and model consumption. Over that \(75\text{ ms}\) window, the
vehicle has traversed \(0.15\text{ m}\) of physical space. If the
perception stage hands downstream components an object detection
bounding box without a timestamp and a spatial transform, a downstream
planner must evaluate that obstacle relative to the current position of
the chassis. It will assume the obstacle sits \(0.15\text{ m}\) closer
or further than its true current coordinate, steering around a physical
ghost or driving directly into an obstacle that moved into its
trajectory during the processing interval. Downstream software cannot
safely consume a perception output that does not explicitly declare when
it was captured, what coordinate system it references, and how wide its
error distribution remains.

\protect\phantomsection\label{callout-definitionux2a-8.2}
\begin{fbxSimple}{callout-definition}{Definition:}{Spatial Information Age}
\phantomsection\label{callout-definition*-8.2}

The \textbf{spatial information age} (\(\tau_{\text{age}}\)) is the
total elapsed physical time between the energy integration midpoint
(\(t_{\text{exp}}/2\)) at the transducer and the downstream consumption
of the observation by a state estimator or policy, translating directly
into unmodeled spatial displacement:
\[\Delta x = \int_{t_0}^{t_0+\tau_{\text{age}}} v(t)\,dt\] as the
embodied platform traverses physical space.

\end{fbxSimple}

\section{What Perception Hands Over}\label{sec-perception-8-2}

Raw sensor hardware yields electrical measurements such as photodiode
voltages, time-of-flight counter values, or capacitive displacements.
These unstructured signals bear little resemblance to the state
variables required by downstream autonomy. An actuator controller cannot
track a desired force using a two-dimensional grid of raw pixel
intensities, and a trajectory optimizer cannot plan collision-free paths
using unorganized time-of-flight returns. Downstream estimation,
planning, and memory subsystems require structured representations of
physical state, including metric positions, velocities, bounding
volumes, surface normals, and semantic classifications. Translating raw
sensor streams into these representations is the operational
responsibility of the perception pipeline. However, downstream consumers
cannot safely interpret a perceived state without explicit metadata that
anchors the measurement in physical reality. To be consumable by
downstream software, an observation must carry three formal properties,
comprising a capture timestamp, an assigned reference frame, and an
uncertainty envelope.

The first mandatory property is the \textbf{capture timestamp} \(t_0\),
which specifies the temporal midpoint of the physical energy integration
window. Sensors do not take instantaneous snapshots; a camera integrates
photons across an exposure duration, and a mechanical lidar sweeps laser
pulses over an azimuth arc across several milliseconds. Downstream state
estimators require \(t_0\) to integrate the measurement into the
machine's dynamic equations of motion. A common architectural failure in
embedded systems is recording the arrival timestamp \(t_{\text{recv}}\),
the moment the operating system kernel completes the direct memory
access transfer, and passing \(t_{\text{recv}}\) downstream as the
observation epoch. Because pipeline execution, serial bus transfers, and
driver scheduling introduce a transport delay
\(\tau = t_{\text{recv}} - t_0\), substituting arrival time for capture
time corrupts the state estimator.

The magnitude of this corruption is a direct function of platform
kinematics. Consider a vehicle traversing a trajectory with
instantaneous velocity \(v(t)\). When an observation taken at \(t_0\) is
evaluated by a downstream filter at \(t_{\text{recv}}\), the physical
state of the machine has shifted. The resulting spatial tracking error
\(\Delta x\) is the integral of the platform velocity over the transport
delay, \(\Delta x = \int_{t_0}^{t_0+\tau} v(t)\,dt\). In an analytical
case of constant linear velocity \(v\), this relationship simplifies to
\(\Delta x = v \cdot \tau\). For a mobile robot traveling at
\(v = 3.0\text{ m/s}\) subject to an end-to-end perception delay
\(\tau = 80\text{ ms}\), treating the arrival time as the measurement
time shifts the perceived obstacle boundary by
\(\Delta x = 0.24\text{ m}\). In closed-loop tracking, feeding delayed
observations without timestamp compensation introduces a phase lag into
the feedback path. This phase lag erodes the phase margin of the
controller, driving the physical system toward oscillatory tracking
errors or complete mechanical divergence.

The second mandatory property is the \textbf{reference frame} identifier
\(\mathcal{F}_{\text{sensor}}\). Every spatial measurement is natively
resolved in the local coordinate system of the physical transducer. A
camera measures bearing angles relative to its optical center, while a
radar measures range and Doppler velocity relative to its antenna array.
Downstream trajectory planners and obstacle maps operate in a chassis
body frame \(\mathcal{F}_{\text{body}}\) or an inertial navigation frame
\(\mathcal{F}_{\text{world}}\). Transforming an observation from
\(\mathcal{F}_{\text{sensor}}\) to \(\mathcal{F}_{\text{body}}\)
requires applying a rigid-body spatial transformation consisting of a
rotation matrix and a translation vector that define the sensor mounting
extrinsics.\sidenote{\footnotesize REP 105 (\emph{Coordinate Frames for Mobile
  Platforms}) standardizes transform trees
  (\(\mathcal{F}_{\text{earth}} \to \mathcal{F}_{\text{map}} \to \mathcal{F}_{\text{odom}} \to \mathcal{F}_{\text{body}} \to \mathcal{F}_{\text{sensor}}\)).
  It isolates high-frequency continuous odometry from discrete map
  corrections, preventing step discontinuities in low-level motor
  controllers during global relocalization.} Across the three machine
archetypes, this spatial binding takes distinct physical forms: in
Mobile Systems, \(\mathcal{F}_{\text{body}}\) represents the mobile
chassis \texttt{base\_link} tracking through an odometric frame
\texttt{odom}; in Manipulator Systems, it maps across serial kinematic
chains from the mounting flange \texttt{tool0} through articulated joint
angles to the base origin \texttt{base\_link}; in Process / Energy
Systems, it maps local pyrometry or laser galvanometer scan coordinates
to the fixed workcell or reactor vessel datum
\(\mathcal{F}_{\text{vessel}}\). As structured in the coordinate frame
transformation tree (Figure~\ref{fig-08-coordinate-frame-tree}), the
standard spatial hierarchy formalizes this spatial progression by
nesting coordinate frames:
\(\mathcal{F}_{\text{earth}} \to \mathcal{F}_{\text{map}} \to \mathcal{F}_{\text{odom}} \to \mathcal{F}_{\text{body}} \to \mathcal{F}_{\text{sensor}}\).
When perception outputs omit the frame identifier or rely on nominal CAD
dimensions rather than calibrated physical extrinsics, small angular
offsets create large spatial errors that scale with distance. An
uncalibrated pitch or yaw error of \(\theta = 1.0^\circ\) on a sensor
detecting an object at a range \(d = 15.0\text{ m}\) induces a
transverse position error of
\(\Delta y \approx d \cdot \sin(1.0^\circ) \approx 0.26\text{ m}\). In
tight warehouse aisles or cluttered physical environments, a lateral
error of \(0.26\text{ m}\) causes the planner to steer into physical
racking or reject valid corridors as obstructed.

The formal mathematical specifications, continuity classes, and
covariance growth characteristics of this coordinate hierarchy are
codified in Table~\ref{tbl-08-rep105-tree}. By establishing explicit
parent-child relations and separating continuous dead reckoning from
discontinuous global localization, the transform tree ensures that
high-frequency control loops and downstream state estimators maintain
mathematically consistent reference origins without incurring
step-discontinuity shocks.

\begin{longtable}[]{@{}
  >{\raggedright\arraybackslash}p{(\linewidth - 12\tabcolsep) * \real{0.1429}}
  >{\raggedright\arraybackslash}p{(\linewidth - 12\tabcolsep) * \real{0.1429}}
  >{\raggedright\arraybackslash}p{(\linewidth - 12\tabcolsep) * \real{0.1429}}
  >{\raggedright\arraybackslash}p{(\linewidth - 12\tabcolsep) * \real{0.1429}}
  >{\raggedright\arraybackslash}p{(\linewidth - 12\tabcolsep) * \real{0.1429}}
  >{\raggedright\arraybackslash}p{(\linewidth - 12\tabcolsep) * \real{0.1429}}
  >{\raggedright\arraybackslash}p{(\linewidth - 12\tabcolsep) * \real{0.1429}}@{}}
\caption{\textbf{REP 105 Spatial Transform Tree Standard and Covariance
Propagation Parameters.} System-wide coordinate frame hierarchy defining
parent-child transform relationships, temporal continuity classes,
estimation sources, and covariance dynamics per ROS Enhancement Proposal
105 (REP 105) and Smith--Cheeseman spatial uncertainty principles
(\citeproc{ref-smith1986estimating}{Smith and Cheeseman 1986}). The
architecture enforces a strict decoupling between discontinuous global
map relocalization (\protect\texttt{map}) and smooth, continuous dead
reckoning (\protect\texttt{odom}), guaranteeing that inner-loop
trajectory controllers and safety barrier functions never encounter
mathematical step discontinuities during loop-closure
corrections.}\label{tbl-08-rep105-tree}\tabularnewline
\toprule\noalign{}
\begin{minipage}[b]{\linewidth}\raggedright
Coordinate Frame (\(\mathcal{F}\))
\end{minipage} & \begin{minipage}[b]{\linewidth}\raggedright
Parent Frame
\end{minipage} & \begin{minipage}[b]{\linewidth}\raggedright
Continuity Class
\end{minipage} & \begin{minipage}[b]{\linewidth}\raggedright
Estimator Source
\end{minipage} & \begin{minipage}[b]{\linewidth}\raggedright
Covariance Growth Dynamic (\(\mathbf{\Sigma}(t)\))
\end{minipage} & \begin{minipage}[b]{\linewidth}\raggedright
Permitted Downstream Consumers
\end{minipage} & \begin{minipage}[b]{\linewidth}\raggedright
Step-Discontinuity \& Safety Policy
\end{minipage} \\
\midrule\noalign{}
\endfirsthead
\toprule\noalign{}
\begin{minipage}[b]{\linewidth}\raggedright
Coordinate Frame (\(\mathcal{F}\))
\end{minipage} & \begin{minipage}[b]{\linewidth}\raggedright
Parent Frame
\end{minipage} & \begin{minipage}[b]{\linewidth}\raggedright
Continuity Class
\end{minipage} & \begin{minipage}[b]{\linewidth}\raggedright
Estimator Source
\end{minipage} & \begin{minipage}[b]{\linewidth}\raggedright
Covariance Growth Dynamic (\(\mathbf{\Sigma}(t)\))
\end{minipage} & \begin{minipage}[b]{\linewidth}\raggedright
Permitted Downstream Consumers
\end{minipage} & \begin{minipage}[b]{\linewidth}\raggedright
Step-Discontinuity \& Safety Policy
\end{minipage} \\
\midrule\noalign{}
\endhead
\bottomrule\noalign{}
\endlastfoot
\texttt{earth} (\(\mathcal{F}_{\text{earth}}\)) & \emph{None} (Root) &
Piecewise \(C^0\) & Multi-band GNSS / RTK / Geodetic survey & Bounded
absolute error: \(\sigma \approx 0.02\text{--}5.0\text{ m}\) depending
on RTK carrier lock & Global mission routing, statutory geofencing &
Step jumps on satellite re-acquisition; \textbf{forbidden} in
closed-loop control \\
\texttt{map} (\(\mathcal{F}_{\text{map}}\)) & \texttt{earth} & Piecewise
\(C^0\) & Global LiDAR/Visual SLAM, landmark factor graphs & Bounded
global variance; discrete contraction on loop closure:
\(\mathbf{\Sigma} \leftarrow (\mathbf{\Sigma}^{-1} + \mathbf{H}^T\mathbf{R}^{-1}\mathbf{H})^{-1}\)
& Global path planners, topological graph search & Loop closure induces
instantaneous coordinate shifts; \textbf{forbidden} in high-rate
tracking \\
\texttt{odom} (\(\mathcal{F}_{\text{odom}}\)) & \texttt{map} & Strictly
\(C^1\) & Wheel encoders, Visual-Inertial Odometry, IMU strapdown &
Unbounded monotonic drift:
\(\mathbf{\Sigma}_{\text{odom}}(t) \propto t\); short-term variance
minimal (\(\sigma_{\Delta p} < 1\text{ mm}\) over \(100\text{ ms}\)) &
High-rate trajectory tracking, MPC, Control Barrier Functions, velocity
estimation & Transform must remain strictly continuous; discrete jumps
are treated as catastrophic faults \\
\texttt{base\_link} (\(\mathcal{F}_{\text{body}}\)) & \texttt{odom} &
\(C^\infty\) (Rigid) & Platform kinematic definition / center of mass &
Identity variance relative to structural mass
(\(\mathbf{\Sigma}_{\text{base}} = \mathbf{0}_{6\times 6}\)) &
Forward/inverse kinematics, swept-volume collision checking, CoM balance
& Fixed origin; updates solely via kinematic joint state and dynamic
integration \\
\texttt{sensor\_link} (\(\mathcal{F}_{\text{sensor}}\)) &
\texttt{base\_link} & \(C^\infty\) (or \(C^1\) gimbal) & Optical/RF CAD
calibration + online extrinsic estimator & Extrinsic covariance
\(\mathbf{\Sigma}_{\text{mount}} \in \mathbb{R}^{6\times 6}\); expands
under thermal expansion (\(\alpha = 23\times 10^{-6}\text{ K}^{-1}\))
and vibration & Perception unprojection, raycasting, optical flow
extraction & Extrinsic drift must be continuously bounded; unprojected
points must transform to \texttt{base\_link} \\
\end{longtable}

The \textbf{REP 105 spatial transform tree} formalizes this
standardized, acyclic kinematic coordinate hierarchy
(\(\mathcal{F}_{\text{earth}} \to \mathcal{F}_{\text{map}} \to \mathcal{F}_{\text{odom}} \to \mathcal{F}_{\text{body}} \to \mathcal{F}_{\text{sensor}}\)),
strictly isolating continuous, high-rate local dead reckoning
(\(\mathcal{F}_{\text{odom}}\)) used in inner control loops from
discrete, discontinuous global map relocalization corrections
(\(\mathcal{F}_{\text{map}}\)) used by global planners, preventing
step-discontinuity shocks in physical tracking controllers.

\begin{figure}[t]

\centering{

\includegraphics[width=0.95\linewidth,height=\textheight,keepaspectratio]{index_files/mediabag//Users/VJ/GitHub/MLSysBook-integrate-physical/physical/book/chapters/08-perception/figures/fig08_coordinate_frame_tree.pdf}

}

\caption{\label{fig-08-coordinate-frame-tree}\textbf{REP 105 Spatial
Transform Tree and Uncertainty Composition.} The spatial coordinate
frame hierarchy formalizes the separation between continuous local dead
reckoning (\(\mathcal{F}_{\text{odom}} \to \mathcal{F}_{\text{body}}\))
consumed by high-frequency tracking loops and discontinuous global map
corrections (\(\mathcal{F}_{\text{map}} \to \mathcal{F}_{\text{odom}}\))
consumed by global planners. Spatial covariance compounds across the
kinematic chain
(\(\mathbf{\Sigma}_{\text{map}} \approx \mathbf{J}_{\text{odom}}\mathbf{\Sigma}_{\text{drift}}\mathbf{J}_{\text{odom}}^T + \mathbf{R}_{\text{body}}\mathbf{\Sigma}_{\text{ext}}\mathbf{R}_{\text{body}}^T + \mathbf{R}_{\text{tot}}\mathbf{\Sigma}_z\mathbf{R}_{\text{tot}}^T\)),
requiring downstream consumers to reject observations that lack an
explicit coordinate key \(\mathcal{F}_{\text{sensor}}\), acquisition
epoch \(t_0\), and covariance envelope \(\mathbf{\Sigma}_z\).
Attribution: Book kinematic systems diagram, CC BY 4.0.}

\end{figure}%

The third mandatory property is the \textbf{uncertainty envelope}
\(\mathbf{\Sigma}_z\). A physical measurement without a variance is an
assertion of absolute certainty. Environmental degradation such as low
illumination, atmospheric scattering, surface specularities, and model
epistemic uncertainty mean that perception outputs are stochastic
estimates. The uncertainty envelope \(\mathbf{\Sigma}_z\), parameterized
as a covariance matrix over the observation dimensions, bounds the
spread of the measurement error. State estimators such as Kalman filters
and nonlinear factor graph optimizers rely on \(\mathbf{\Sigma}_z\) to
weight sensory innovations against the internal dynamic model of the
body. If perception hands over an observation without
\(\mathbf{\Sigma}_z\), downstream algorithms must assume a static noise
parameter or treat the measurement as ground truth. When the physical
sensor degrades, treating a noisy or out-of-distribution perception
output as certain forces the optimizer to pull state estimates toward
spurious measurements. The resulting high-frequency corrections
propagate directly to motor drives, causing severe torque chatter, rapid
thermal rise in actuator windings, and mechanical wear.

The \textbf{perception handoff contract} is the formal interface
invariant requiring every observation packet crossing from the
perception pipeline into the nervous system to carry a capture timestamp
(\(t_0\)), a verified reference coordinate frame
(\(\mathcal{F}_{\text{sensor}}\)), and an explicit uncertainty
covariance matrix (\(\mathbf{\Sigma}_z\)). Downstream components must
reject any ungrounded packet that omits these fields.

\protect\phantomsection\label{callout-ruleux2a-8.3}
\begin{fbxSimple}{callout-rule}{Architect’s Rule of Thumb:}{The Three Invariants of Observation}
\phantomsection\label{callout-rule*-8.3}

Never pass a perception tensor downstream without an acquisition
timestamp latched at the exposure midpoint, an unambiguous coordinate
frame identifier, and an explicit spatial uncertainty covariance matrix.

\end{fbxSimple}

These three properties define the formal handoff contract between
perception and the remainder of the autonomy stack. Downstream state
estimators, planners, and safety watchdogs must reject any observation
packet that lacks an explicit capture timestamp \(t_0\), a verified
coordinate frame \(\mathcal{F}_{\text{sensor}}\), and an uncertainty
covariance \(\mathbf{\Sigma}_z\). Dropping an incomplete or unbounded
observation allows a downstream estimator to propagate its state forward
using physical dynamics, degrading uncertainty smoothly over time.
Accepting an ungrounded packet, by contrast, injects unmodeled delays,
spatial distortions, and false confidence directly into the control
loop. Enforcing this contract requires tracing how an observation
originates, beginning at the hardware interface where physical phenomena
are first sampled.

\section{Transduction, Hardware Sampling, and Temporal
Calibration}\label{sec-perception-8-3}

Lag accumulates before software receives anything at all, beginning at
the photosite where physical phenomena are converted into electrical
charge. A camera does not capture an instantaneous snapshot of the
physical world. Photons collect on a photodiode over a finite exposure
duration \(t_{\text{exp}}\), which means the recorded pixel intensity
represents the temporal integral of scene radiance \(I(t)\) over that
interval: \[I_{\text{pixel}} = \int_{0}^{t_{\text{exp}}} I(t)\,dt\] When
the sensor or the observed environment moves at velocity \(v\) during
exposure, a sharp physical boundary spreads across multiple photosites.
The resulting motion blur width in the direction of motion is given by
\(\Delta x_{\text{blur}} = v \cdot t_{\text{exp}}\). Because light
accumulates continuously across the exposure window, the effective point
in time that corresponds to the accumulated charge is the temporal
midpoint of the interval, introducing an unavoidable physical
measurement lag of \(t_{\text{exp}}/2\) relative to the start of
exposure.

\subsection{Rolling Shutter Distortion and Motion
Compensation}\label{rolling-shutter-distortion-and-motion-compensation}

How that exposure window is scheduled across the pixel array determines
the spatial geometry of the resulting image. In a \textbf{global
shutter} sensor, every pixel photodiode accumulates charge during the
identical time interval and transfers the accumulated electrons to
shielded storage nodes simultaneously. In a \textbf{rolling shutter}
sensor, silicon area and cost constraints dictate that rows are reset,
exposed, and read out sequentially from the top scanline \(y = 0\) to
the bottom scanline \(y = H - 1\). Each successive scanline begins
exposure after a row period \(t_{\text{row}}\), creating a progressive
temporal offset across the image plane:\sidenote{\footnotesize In rolling shutter
  sensors, sequential row readout
  (\(t_{\text{readout}} = (H-1)t_{\text{row}}\)) warps epipolar lines
  under angular velocity \(\boldsymbol{\omega}\), breaking rigid
  multi-view geometry unless continuous-time pose interpolation
  \(\mathbf{T}_{wb}(t(y))\) is modeled.}
\[t(y) = t_0 + y \cdot t_{\text{row}}\] where \(t_0\) is the start of
exposure for the top row. When the body translates at linear velocity
\(v\) perpendicular to an observed vertical feature, scanline \(y\)
captures that feature at a spatial displacement:
\[\Delta x(y) = v \cdot y \cdot t_{\text{row}}\]

Under rotational motion, rolling shutter distortion becomes non-linear.
If the camera undergoes angular rotation at angular velocity vector
\(\boldsymbol{\omega} = [\omega_x, \omega_y, \omega_z]^T\) (e.g., rapid
yaw or pitch during aggressive robot maneuvers), the effective camera
attitude changes continuously as scanlines are read. For a scanline at
vertical coordinate \(y\), the differential rotation relative to the
frame start is
\(\Delta \boldsymbol{\theta}(y) \approx \boldsymbol{\omega} \cdot y \cdot t_{\text{row}}\).
This scanline-dependent rotation warps straight physical lines into
curves, induces focal-length dilation (scaling distortion during pitch),
and shears vertical walls into slanted surfaces.

As an analytical example, consider a ground robot traversing an aisle at
\(v = 3.0\text{ m/s}\) equipped with a rolling shutter camera operating
at \(30\text{ Hz}\) with vertical resolution \(H = 1080\text{ rows}\).
If the active readout time across all rows is
\(t_{\text{readout}} = 33.3\text{ ms}\), the inter-row readout interval
is
\(t_{\text{row}} = 33.3\text{ ms} / 1080 \approx 30.8\text{ }\mu\text{s}\).
A vertical pallet edge spanning the full height of the frame does not
appear vertical in the captured array. The bottom scanline
(\(y = 1079\)) is sampled \(33.3\text{ ms}\) after the top scanline
(\(y = 0\)), producing a total physical spatial skew of:
\[\Delta x = 3.0\text{ m/s} \times 0.0333\text{ s} \approx 0.10\text{ m} = 100\text{ mm}\]
With an exposure duration of \(t_{\text{exp}} = 10.0\text{ ms}\), each
individual scanline also suffers a motion blur width of
\(\Delta x_{\text{blur}} = 3.0\text{ m/s} \times 0.010\text{ s} = 0.03\text{ m} = 30\text{ mm}\).
If downstream geometry reconstruction treats this image as an
instantaneous rigid projection, the perceived \(100\text{ mm}\) tilt
causes the motion planner to identify an artificial intrusion into the
free corridor. Correcting this distortion requires
\textbf{rolling-shutter motion compensation}, the continuous-time
geometric unwarping algorithm that utilizes high-rate
(\(1000\text{ Hz}\)) inertial angular velocities
(\(\boldsymbol{\omega}\)) and linear velocities (\(\mathbf{v}\)) to
interpolate the instantaneous 6-DoF camera pose
\(\mathbf{T}_{wb}(t(y))\) for each individual scanline \(y\), restoring
rigid epipolar geometry prior to 3D triangulation.

\protect\phantomsection\label{callout-incidentux2a-8.4}
\begin{fbxSimple}{callout-incident}{Landmark Engineering Incident:}{Mobile Robot Rolling-Shutter Image Shear Fracture}
\phantomsection\label{callout-incident*-8.4}

\textbf{System:} Autonomous mobile warehouse transport robot equipped
with a forward-facing rolling-shutter CMOS camera (\(1080\text{p}\) at
\(30\text{ Hz}\), row readout time
\(t_{\text{row}} = 30.8\ \mu\text{s}\), total frame readout
\(t_{\text{readout}} = 33.3\text{ ms}\)) and an onboard 6-DoF
Visual-Inertial Odometry (VIO) pipeline running a sliding-window factor
graph optimizer.

\textbf{Incident:} While executing a rapid \(90^\circ\) turn at an
angular velocity of \(\omega_z = 2.4\text{ rad/s}\)
(\(137.5^\circ/\text{s}\)) alongside an industrial storage rack, the
robot's state estimator suffered an instantaneous divergence fracture
(\(>0.52\text{ m}\) lateral position jump). The motion planner commanded
a corrective steer directly into the steel racking upright at
\(2.2\text{ m/s}\), shearing the camera mount and causing structural
chassis deformation.

\textbf{Root Cause \& Physics of Failure:} Rolling shutter line-by-line
electronic scanning exposed the top scanline (\(y = 0\))
\(33.3\text{ ms}\) earlier than the bottom scanline (\(y = 1079\)).
During the \(2.4\text{ rad/s}\) yaw maneuver, vertical pallet uprights
suffered a differential rotational skew of:
\[\Delta \theta = \omega_z \cdot t_{\text{readout}} = 2.4\text{ rad/s} \times 0.0333\text{ s} \approx 0.080\text{ rad} = 4.58^\circ\]
This angular skew bent straight vertical shelf beams into tilted
hyperbolas across the image plane. The visual frontend extracted 2D
feature tracks under the assumption of a rigid pinhole projection.
Because the factor graph optimizer lacked scanline-level continuous-time
pose interpolation (\(\mathbf{T}_{wb}(t(y))\)), the epipolar
reprojection residuals exploded. To minimize reprojection error across
warped features, the non-linear solver forced the scale factor to
collapse and estimated a false lateral translation of \(0.52\text{ m}\),
pulling the vehicle's belief away from the real physical obstacle
(\citeproc{ref-forster2017onmanifold}{Forster et al. 2017};
\citeproc{ref-li2014online}{Li and Mourikis 2014}).

\textbf{Systems Resolution:} Replaced rolling-shutter sensors with
global-shutter CMOS imagers equipped with hardware GPIO exposure
strobing synchronized to the IMU sampling register (\texttt{TIMx\_CCR}),
and integrated continuous-time B-spline trajectory compensation into the
VIO factor graph to interpolate exact scanline poses during transient
sensor exposure.

\textbf{Systems Takeaway:} A 2D camera image is not an instantaneous
spatial snapshot unless global shutter hardware or continuous-time
kinematic motion compensation explicitly guarantees that all pixels
share an identical temporal epoch.

\end{fbxSimple}

\subsection{MIPI CSI-2 DMA Ingress on the Linux
MPU}\label{mipi-csi-2-dma-ingress-on-the-linux-mpu}

Once analog-to-digital converters on the sensor die digitize the pixel
charges, the frame must travel across a physical interconnect to host
memory. In a canonical heterogeneous System-on-Chip (SoC) comprising an
unprivileged application processor (Linux MPU) and an isolated hard
real-time safety microcontroller (bare-metal MCU), camera streams enter
via the \textbf{Mobile Industry Processor Interface Camera Serial
Interface (MIPI CSI-2)} bus.\sidenote{\footnotesize MIPI CSI-2 transmits serial
  packets over D-PHY (up to \(4.5\text{ Gbps/lane}\)) or C-PHY
  (\(3.5\text{ Gsym/s/lane}\)), protected by ECC headers and CRC
  checksums.}

\protect\phantomsection\label{callout-decisionux2a-8.5}
\begin{fbxSimple}{callout-decision}{Design Decision:}{Silicon Ingestion Dataflow Architecture}
\phantomsection\label{callout-decision*-8.5}

\begin{enumerate}
\def\labelenumi{\arabic{enumi}.}
\tightlist
\item
  \textbf{D-PHY / C-PHY Physical Layer}: Differential data lanes
  transmit high-speed serialized pixel packets
  (\(1.5\text{--}2.5\text{ Gbps}\) per lane), protected by header error
  correction.
\item
  \textbf{SoC CSI-2 Host Controller}: Deserializes byte packets, strips
  headers, and validates CRC integrity.
\item
  \textbf{Scatter-Gather DMA Controller}: A dedicated DMA engine
  transfers pixel byte words directly into DRAM over the internal system
  interconnect (e.g., AXI crossbar).
\item
  \textbf{Kernel DMA Ring Buffers}: Under Linux (\texttt{v4l2}), the
  driver allocates pre-mapped non-cacheable DMA buffers
  (\texttt{V4L2\_MEMORY\_MMAP}). The DMA populates pages without
  user-space copies.\sidenote{\footnotesize Zero-copy V4L2 DMA
    (\texttt{V4L2\_MEMORY\_DMABUF}) maps frames to accelerator DRAM,
    bypassing CPU copies. Memory bus contention can still inject
    millisecond tail latencies under heavy load.}
\end{enumerate}

\end{fbxSimple}

Once DMA write completion occurs, the kernel driver issues an interrupt.
Pointers to these contiguous memory regions are shared as zero-copy file
descriptors (\texttt{dma\_buf}) directly with the hardware Image Signal
Processor (ISP) and Neural Processing Unit (NPU), bypassing CPU memory
bus hops. By contrast, interfaces relying on Universal Serial Bus (USB3)
or GigE Vision over Ethernet package pixel streams into discrete
packets, incurring frame assembly, kernel network stack traversal, and
driver memory copies that inflate transport latency from
\(t_{\text{transport}} < 0.5\text{ ms}\) (MIPI CSI-2) up to
\(5\text{--}20\text{ ms}\) (USB3/GigE), with tail latencies exceeding
\(P_{99} > 35\text{ ms}\) under CPU thread scheduling contention.

\subsection{Camera-to-IMU Temporal and Extrinsic
Calibration}\label{camera-to-imu-temporal-and-extrinsic-calibration}

Visual perception does not operate in isolation; it functions alongside
an Inertial Measurement Unit (IMU) running at high frequencies
(\(1000\text{ Hz}\)). Fusing high-rate angular velocity
\(\boldsymbol{\omega}\) and linear acceleration \(\mathbf{a}\) from the
IMU with low-rate visual frames (\(30\text{--}60\text{ Hz}\)) in
Visual-Inertial Odometry (VIO) requires rigorous spatial and temporal
calibration (\citeproc{ref-furgale2013unified}{Furgale et al. 2013};
\citeproc{ref-li2014online}{Li and Mourikis 2014};
\citeproc{ref-forster2017onmanifold}{Forster et al. 2017}).

The spatial relationship between the IMU reference frame
\(\mathcal{F}_{\text{imu}}\) and the camera reference frame
\(\mathcal{F}_{\text{cam}}\) is defined by the rigid extrinsic
transformation
\(\mathbf{T}_{\text{imu}}^{\text{cam}} = (\mathbf{R}_{\text{imu}}^{\text{cam}}, \mathbf{p}_{\text{imu}}^{\text{cam}})\).
Because the camera and IMU are physically separated by a lever arm
\(\mathbf{p}_{\text{imu}}^{\text{cam}}\), an angular acceleration
\(\dot{\boldsymbol{\omega}}\) or angular velocity
\(\boldsymbol{\omega}\) experienced by the body induces an additional
tangential and centripetal acceleration at the camera:
\[\mathbf{a}_{\text{cam}} = \mathbf{R}_{\text{imu}}^{\text{cam}} (\mathbf{a}_{\text{imu}} - \mathbf{g}) + \boldsymbol{\omega} \times (\boldsymbol{\omega} \times \mathbf{p}_{\text{imu}}^{\text{cam}}) + \dot{\boldsymbol{\omega}} \times \mathbf{p}_{\text{imu}}^{\text{cam}}\]

Equally critical is the \textbf{temporal calibration offset}
\(\tau_{\text{cam-imu}} = t_{\text{cam}} - t_{\text{imu}}\). Standard
commercial crystal oscillators on independent sensor breakout boards
drift by roughly \(\pm 50\text{ ppm}\), accumulating
\(50\text{ }\mu\text{s}\) of error every second (\(4.3\text{ s}\) per
day). Moreover, analog anti-aliasing filters in the IMU, electronic
exposure delays in the camera, and DMA driver scheduling inject fixed
and variable temporal offsets between the two data streams.

If the temporal offset \(\Delta t\) is uncalibrated, platform rotation
couples directly into catastrophic spatial estimation errors. Suppose a
robot rotates at an angular velocity of \(\omega = 2.0\text{ rad/s}\)
(\(115^\circ/\text{s}\)). An undetected temporal lag of
\(\Delta t = 10\text{ ms}\) between the camera frame and the IMU
integration causes a rotational orientation mismatch of:
\[\Delta \theta = \omega \cdot \Delta t = 2.0\text{ rad/s} \times 0.010\text{ s} = 0.020\text{ rad} \approx 1.15^\circ\]
At an operating distance of \(z = 10.0\text{ m}\), this angular
misalignment projects into a transverse spatial error of:
\[\Delta p \approx z \cdot \Delta \theta = 10.0\text{ m} \times 0.020\text{ rad} = 0.20\text{ m} = 200\text{ mm}\]
This artificial \(200\text{ mm}\) displacement corrupts visual feature
triangulation, causing visual-inertial estimators to compute false scale
factors and driving the filter toward mathematical divergence.
Eliminating these timing errors requires hardware-level synchronization:
a hardware GPIO trigger strobe latches exposure start times across all
cameras and IMU sampling registers simultaneously,\sidenote{\footnotesize In canonical
  safety microcontrollers, hardware trip zones (\texttt{EPWM\_TZFRC},
  \texttt{TIMx\_BKIN}) operate independently of the CPU. Hardware break
  circuits force complementary PWM outputs into high-impedance or
  low-side clamp states within \(<50\text{ ns}\), preventing
  shoot-through before software ISRs execute.} while online
continuous-time calibration algorithms estimate residual sub-millisecond
offsets \(\tau_{\text{cam-imu}}\) within the joint factor graph
(\citeproc{ref-furgale2013unified}{Furgale et al. 2013}). As shown in
the physical multi-modal perception pod in
Figure~\ref{fig-08-sensor-suite}, co-locating multiple high-resolution
cameras and LiDAR modules on a rigid vehicle bracket requires routing
these deterministic hardware sync strobes alongside high-speed
GMSL2/FPD-Link serializer links to bound clock drift and transport
latency across all sensing channels.

\textbf{Camera-to-IMU spatial-temporal calibration} is the joint
determination of the rigid 6-DoF extrinsic transformation
\(\mathbf{T}_{\text{imu}}^{\text{cam}} \in SE(3)\) and the temporal
latency offset
\(\tau_{\text{cam-imu}} = t_{\text{cam}} - t_{\text{imu}}\), bounding
lever-arm centripetal accelerations and rotational misalignment errors
(\(\Delta p \approx z \omega \Delta t\)) during dynamic motion.

\begin{figure}[t]

\centering{

\includegraphics[width=0.85\linewidth,height=\textheight,keepaspectratio]{/Users/VJ/GitHub/MLSysBook-integrate-physical/physical/book/chapters/08-perception/figures/fig08_real_sensor_suite.jpg}

}

\caption{\label{fig-08-sensor-suite}\textbf{Physical Multi-Modal
Perception Sensor Suite on an Autonomous Testbed.} Close-up view of an
integrated perception pod mounted on an autonomous vehicle chassis. The
assembly houses multiple high-dynamic-range CMOS cameras with distinct
focal lengths and fields of view alongside solid-state and spinning
time-of-flight LiDAR modules on a rigid machined aluminum mounting
bracket. Each physical transducer defines its own local reference frame
(\(\mathcal{F}_{\text{sensor}}\)), requiring precise extrinsic geometric
calibration (\(\mathbf{T}_{\text{sensor}\to\text{body}}\)) and thermal
compensation to prevent angular deviations (\(\theta\)) from projecting
into severe lateral errors (\(\Delta y \approx d \cdot \sin\theta\)) at
operating distance. Shielded coaxial serializers (GMSL2/FPD-Link) and
dedicated hardware sync lines route synchronized microsecond triggers to
prevent inter-sensor phase skew. Attribution: Open-source robotics
testbed platform, CC BY-SA 4.0.}

\end{figure}%

\subsection{The Latency Waterfall}\label{the-latency-waterfall}

Summing these physical stages yields the total age of the sensory data
as it cascades through the perception pipeline. That information age is
governed by the \textbf{latency waterfall}:
\[t_{\text{age}} = \frac{t_{\text{exp}}}{2} + t_{\text{readout}} + t_{\text{transport}} + t_{\text{DMA}} + t_{\text{ISP}} + t_{\text{backbone}} + t_{\text{IPC}}\]

\protect\phantomsection\label{callout-mathux2a-8.6}
\begin{fbxSimple}{callout-math}{Fermi Math:}{Sensor Transduction Latency and Spatial Information Age Budget}
\phantomsection\label{callout-math*-8.6}

\emph{The scenario.} An embodied platform navigates a dynamic workspace.
We evaluate the cumulative information age \(\tau_{\text{age}}\) across
the complete sensory-to-inference pipeline and compute the resulting
unmodeled physical displacement
\(\Delta x = v \cdot \tau_{\text{age}} + \frac{1}{2} a \cdot \tau_{\text{age}}^2\)
across three distinct physical machine archetypes.

\textbf{Stage-by-Stage Latency Waterfall:}

\begin{enumerate}
\def\labelenumi{\arabic{enumi}.}
\tightlist
\item
  \emph{Exposure Integration Midpoint (\(t_{\text{exp}}/2\)):} CMOS
  camera running at \(30\text{ Hz}\) with auto-exposure
  \(t_{\text{exp}} = 16.0\text{ ms} \implies \Delta t_1 = 8.0\text{ ms}\).
\item
  \emph{Rolling Shutter Line Readout (\(t_{\text{readout}}\)):} Column
  parallel ADC scanning across 1080 rows
  \(\implies \Delta t_2 = 16.6\text{ ms}\) (average scanline).
\item
  \emph{MIPI CSI-2 Serialization \& Transport
  (\(t_{\text{transport}}\)):} 4-lane D-PHY at
  \(2.5\text{ Gbps/lane} \implies \Delta t_3 = 0.8\text{ ms}\).
\item
  \emph{DMA Ingress \& Kernel V4L2 Buffer Queue (\(t_{\text{DMA}}\)):}
  Scatter-gather DRAM write and driver interrupt
  \(\implies \Delta t_4 = 1.2\text{ ms}\).
\item
  \emph{Hardware ISP Debayering \& Rectification (\(t_{\text{ISP}}\)):}
  Dedicated on-chip ISP hardware pass
  \(\implies \Delta t_5 = 2.5\text{ ms}\).
\item
  \emph{Vision Backbone Tensor Forward Pass (\(t_{\text{backbone}}\)):}
  INT8 quantized Vision Transformer on edge NPU
  \(\implies \Delta t_6 = 22.0\text{ ms}\).
\item
  \emph{Covariance Packaging \& Zero-Copy IPC (\(t_{\text{IPC}}\)):}
  Shared-memory circular ring buffer dispatch
  \(\implies \Delta t_7 = 0.5\text{ ms}\).
\end{enumerate}

\textbf{Total Spatial Information Age:}
\[\tau_{\text{age}} = 8.0 + 16.6 + 0.8 + 1.2 + 2.5 + 22.0 + 0.5 = 51.6\text{ ms}\]

\textbf{Archetype Displacement Analysis:}

\begin{itemize}
\item
  \textbf{Mobile Systems (Autonomous Mobile Robot):} At nominal cruising
  speed \(v = 2.5\text{ m/s}\) under acceleration
  \(a = 0\text{ m/s}^2\):
  \[\Delta x_{\text{AMR}} = 2.5\text{ m/s} \times 0.0516\text{ s} = 0.129\text{ m} = 129\text{ mm}\]
  \emph{Systems Impact:} Consumes \(>60\%\) of the nominal
  \(200\text{ mm}\) safety corridor clearance before the policy
  evaluates the frame.
\item
  \textbf{Manipulators (High-Speed Industrial Arm):} At end-effector
  Cartesian velocity \(v = 1.8\text{ m/s}\) during precision peg-in-hole
  assembly:
  \[\Delta x_{\text{arm}} = 1.8\text{ m/s} \times 0.0516\text{ s} = 0.0929\text{ m} = 92.9\text{ mm}\]
  \emph{Systems Impact:} Exceeds standard chamfer insertion tolerances
  (\(1\text{--}5\text{ mm}\)) by nearly two orders of magnitude,
  guaranteeing tool collision without historical pose interpolation.
\item
  \textbf{Process / Energy Systems (Laser Metal Additive Melt-Pool
  Scanner):} Galvanometer scan mirror traversing at
  \(v_{\text{scan}} = 0.8\text{ m/s}\) with thermal pyrometry:
  \[\Delta x_{\text{melt}} = 0.8\text{ m/s} \times 0.0516\text{ s} = 0.0413\text{ m} = 41.3\text{ mm}\]
  \emph{Systems Impact:} The thermal feedback loop controls laser power
  against a melt-pool position that has shifted by \(>80\) melt-pool
  diameters (\(d_{\text{pool}} \approx 0.5\text{ mm}\)), causing thermal
  runaway or void formation.
\end{itemize}

\end{fbxSimple}

Tracing this latency progression across the entire perception and
inference pipeline reveals how time accumulates from initial photon
arrival down to downstream token generation and control dispatch. As
budgeted stage-by-stage in Table~\ref{tbl-08-latency-budget}, every
hardware serialization step, memory bus arbitration stall, and neural
forward pass compounds the effective information age. Across typical
platform velocities, this accumulated latency translates directly into
centimeters or meters of unrecoverable physical displacement before a
control loop can react.

\begin{longtable}[]{@{}
  >{\raggedright\arraybackslash}p{(\linewidth - 14\tabcolsep) * \real{0.1250}}
  >{\raggedright\arraybackslash}p{(\linewidth - 14\tabcolsep) * \real{0.1250}}
  >{\raggedright\arraybackslash}p{(\linewidth - 14\tabcolsep) * \real{0.1250}}
  >{\raggedright\arraybackslash}p{(\linewidth - 14\tabcolsep) * \real{0.1250}}
  >{\raggedright\arraybackslash}p{(\linewidth - 14\tabcolsep) * \real{0.1250}}
  >{\raggedright\arraybackslash}p{(\linewidth - 14\tabcolsep) * \real{0.1250}}
  >{\raggedright\arraybackslash}p{(\linewidth - 14\tabcolsep) * \real{0.1250}}
  >{\raggedright\arraybackslash}p{(\linewidth - 14\tabcolsep) * \real{0.1250}}@{}}
\caption{\textbf{End-to-End Information Age Accumulation and Physical
Displacement Budget.} Stage-by-stage temporal accounting from initial
photon arrival at the sensor photosite through hardware deserialization,
direct memory access (DMA), neural backbone inference, and inter-process
communication (IPC) dispatch. Information age \(\Delta t\) reflects the
elapsed time from physical energy integration midpoint
(\(t_{\text{exp}}/2\)) to downstream availability. Nominal metrics
reflect optimized embedded hardware pipelines (e.g., \(60\text{ Hz}\)
global-shutter camera, hardware ISP, INT8 NPU acceleration); worst-case
bounds (\(P_{99.9}\)) reflect thermal clock throttling, memory bus
saturation, and OS scheduling jitter. Physical displacement
\(\Delta x = v \cdot t_{\text{age}}\) demonstrates how latency converts
directly into unrecoverable spatial drift across operational
speeds.}\label{tbl-08-latency-budget}\tabularnewline
\toprule\noalign{}
\begin{minipage}[b]{\linewidth}\raggedright
Pipeline Stage \& Physical Mechanism
\end{minipage} & \begin{minipage}[b]{\linewidth}\raggedright
Nominal Latency (\(\mu \pm \sigma\))
\end{minipage} & \begin{minipage}[b]{\linewidth}\raggedright
Worst-Case Bound (\(P_{99.9}\))
\end{minipage} & \begin{minipage}[b]{\linewidth}\raggedright
Added Age (\(\Delta t\))
\end{minipage} & \begin{minipage}[b]{\linewidth}\raggedright
Cumulative Age (\(t_{\text{age}}\))
\end{minipage} & \begin{minipage}[b]{\linewidth}\raggedright
Displacement at \(2.0\text{ m/s}\)
\end{minipage} & \begin{minipage}[b]{\linewidth}\raggedright
Displacement at \(10.0\text{ m/s}\)
\end{minipage} & \begin{minipage}[b]{\linewidth}\raggedright
Contention Source \& Hardware Mitigation
\end{minipage} \\
\midrule\noalign{}
\endfirsthead
\toprule\noalign{}
\begin{minipage}[b]{\linewidth}\raggedright
Pipeline Stage \& Physical Mechanism
\end{minipage} & \begin{minipage}[b]{\linewidth}\raggedright
Nominal Latency (\(\mu \pm \sigma\))
\end{minipage} & \begin{minipage}[b]{\linewidth}\raggedright
Worst-Case Bound (\(P_{99.9}\))
\end{minipage} & \begin{minipage}[b]{\linewidth}\raggedright
Added Age (\(\Delta t\))
\end{minipage} & \begin{minipage}[b]{\linewidth}\raggedright
Cumulative Age (\(t_{\text{age}}\))
\end{minipage} & \begin{minipage}[b]{\linewidth}\raggedright
Displacement at \(2.0\text{ m/s}\)
\end{minipage} & \begin{minipage}[b]{\linewidth}\raggedright
Displacement at \(10.0\text{ m/s}\)
\end{minipage} & \begin{minipage}[b]{\linewidth}\raggedright
Contention Source \& Hardware Mitigation
\end{minipage} \\
\midrule\noalign{}
\endhead
\bottomrule\noalign{}
\endlastfoot
\textbf{1. Photodiode Transduction \& Exposure}Photon integration on
CMOS pixel array & \(8.0 \pm 1.5\text{ ms}\) & \(16.5\text{ ms}\) &
\(4.0\text{--}16.5\text{ ms}\) & \(4.0\text{--}16.5\text{ ms}\) &
\(8\text{--}33\text{ mm}\) & \(40\text{--}165\text{ mm}\) & Low-light
integration stretch; clamp maximum exposure via driver register
(\texttt{SHUTTER\_MAX}) \\
\textbf{2. Sensor Readout \& Column ADC}Correlated double sampling \&
pixel digitization & \(11.0 \pm 0.5\text{ ms}\) & \(33.3\text{ ms}\) &
\(11.0\text{--}33.3\text{ ms}\) & \(15.0\text{--}49.8\text{ ms}\) &
\(30\text{--}100\text{ mm}\) & \(150\text{--}498\text{ mm}\) & Rolling
shutter line skew; deploy global shutter sensors with hardware GPIO
strobe triggers \\
\textbf{3. Physical Interconnect Transport}Differential serialization
over GMSL2 / CSI-2 & \(0.8 \pm 0.1\text{ ms}\) & \(2.2\text{ ms}\) &
\(0.8\text{--}2.2\text{ ms}\) & \(15.8\text{--}52.0\text{ ms}\) &
\(32\text{--}104\text{ mm}\) & \(158\text{--}520\text{ mm}\) &
Packetization \& protocol overhead; bypass USB/GigE in favor of
dedicated GMSL2/FPD-Link links \\
\textbf{4. DMA Ingress \& Kernel Ring Buffering}Scatter-gather DMA
transfer to host DRAM & \(1.2 \pm 0.3\text{ ms}\) & \(8.5\text{ ms}\) &
\(1.2\text{--}8.5\text{ ms}\) & \(17.0\text{--}60.5\text{ ms}\) &
\(34\text{--}121\text{ mm}\) & \(170\text{--}605\text{ mm}\) & Shared
DRAM bus contention; assign high AXI QoS priorities and lock
non-cacheable DMA buffers \\
\textbf{5. Hardware ISP \& Tensor Normalization}Debayering, geometric
rectify, color format & \(2.5 \pm 0.4\text{ ms}\) & \(7.0\text{ ms}\) &
\(2.5\text{--}7.0\text{ ms}\) & \(19.5\text{--}67.5\text{ ms}\) &
\(39\text{--}135\text{ mm}\) & \(195\text{--}675\text{ mm}\) & Software
fallback stalls; enforce dedicated on-chip fixed-function ISP hardware
pipelines \\
\textbf{6. Neural Backbone Forward Pass}INT8/FP16 tensor matrix multiply
on NPU & \(18.0 \pm 2.0\text{ ms}\) & \(38.0\text{ ms}\) &
\(18.0\text{--}38.0\text{ ms}\) & \(37.5\text{--}105.5\text{ ms}\) &
\(75\text{--}211\text{ mm}\) & \(375\text{--}1055\text{ mm}\) & Thread
preemption and thermal throttling (\(T_j > 85^\circ\text{C}\)); lock NPU
core clocks and pin weights in SRAM \\
\textbf{7. Covariance Packaging \& IPC Dispatch}Spatial bounding,
bitfield check, shmem ring & \(0.5 \pm 0.1\text{ ms}\) &
\(4.5\text{ ms}\) & \(0.5\text{--}4.5\text{ ms}\) &
\textbf{\(38.0\text{--}110.0\text{ ms}\)} &
\textbf{\(76\text{--}220\text{ mm}\)} &
\textbf{\(380\text{--}1100\text{ mm}\)} & Downstream lock contention;
deploy lock-free POSIX shared memory and enforce hard age-drop rules \\
\end{longtable}

\section{The Cost of Ingestion}\label{sec-perception-8-4}

Getting sensory data into main memory is not free, and it does not
happen in an isolated subsystem. When raw pixel streams arrive across
physical links such as MIPI CSI-2 or Ethernet, hardware engines write
those bytes directly into dynamic random-access memory (DRAM) through
direct memory access (DMA). That transfer shares the physical
interconnect, cache hierarchies, and memory controllers with every other
executing workload on the machine, including neural inference engines,
operating system drivers, and safety monitors. Every sensory stream
therefore imposes two numbers on the machine. The first number is the
age the transport pipeline adds to the observation before compute cores
can read it. The second number is the delay that incoming data stream
imposes on the safety enforcement path competing for the exact same
memory bus. The first number delays when the policy knows what happened.
The second number delays whether the safety reflex can veto an unsafe
actuation command before a physical deadline expires.

To see the scale of this competition across physical archetypes,
consider their respective ingress demands: an industrial manipulator
synchronizing dual high-speed eye-in-hand cameras with tactile arrays
(\(600\text{ MB/s}\) burst ingress); a laser powder-bed fusion process
scanner streaming coaxial NIR pyrometry at \(1000\text{ Hz}\)
(\(1.2\text{ GB/s}\) continuous ingress); or an autonomous mobile
machine equipped with four high-resolution surround cameras and dual
lidars. In the mobile configuration, suppose the four global shutter
cameras capture \(3840 \times 2160\) images at \(30\text{ Hz}\) using a
raw 12-bit format packed into \(1.5\text{ B}\) per pixel. Each
uncompressed frame occupies
\(3840 \times 2160 \times 1.5\text{ B} \approx 12.44\text{ MB}\) of
memory. Across four cameras at \(30\text{ Hz}\), the sustained write
traffic generated by frame ingestion alone is:
\[B_{\text{ingest}} = 4 \times 12.44\text{ MB} \times 30\text{ s}^{-1} \approx 1.493\text{ GB/s}\]
When the perception stack reads these frames to feed vision backbones,
operating system page caches duplicate intermediate buffers, and neural
network accelerators read weights and write activations, the memory
controller experiences traffic several times larger than the raw
ingestion rate. If the host platform uses a dual-channel LPDDR5 memory
subsystem with a theoretical peak bandwidth of \(51.2\text{ GB/s}\), the
raw camera ingress consumes roughly \(2.9\text{ percent}\) of
theoretical capacity. Evaluated purely as a steady-state average, this
ingress appears well within the limits of the hardware interface.

Steady-state averages, however, conceal the physical failure mode.
Physical sensors do not emit bytes as a continuous, uniformly
distributed trickle. They emit data in synchronized bursts when hardware
exposure gates finish readout. When four cameras triggered by a shared
hardware strobe complete their frame readout simultaneously over a
\(16.0\text{ ms}\) window, DMA controllers write large transaction
bursts across the system interconnect to maximize bus utilization. While
these high-priority DMA writes stream into DRAM, the memory controller
must switch bus directions from reads to writes, activate new DRAM
banks, and arbitrate competing transaction queues. If the safety
enforcement monitor needs to read joint state estimates, evaluate
control barrier functions, and submit a brake command during this
precise burst, its memory read requests stall behind the inbound sensory
flood. An average ingestion load that consumes only three percent of bus
capacity can still monopolize memory controller channels for several
continuous milliseconds, pushing the tail latency of the enforcement
path past its deadline.

The maximum admissible ingestion load must therefore be derived from the
enforcement deadline and the measured timing slack established in
Chapter~\ref{sec-body}, rather than from what the physical interface can
sustain. Let \(T_{\text{enforce}}\) denote the fixed period of the
safety enforcement loop, and let \(t_{\text{exec}}\) denote its
execution duration on an idle memory bus. The timing slack available to
absorb interference is \(S = T_{\text{enforce}} - t_{\text{exec}}\).
Under concurrent sensory ingestion, the execution duration expands by a
contention stall \(\Delta t_{\text{stall}}\), which depends on the burst
volume \(V_{\text{burst}}\) and the effective DRAM service rate
\(R_{\text{arb}}\) under read-write arbitration. Preventing deadline
violations requires that the worst-case stall stay strictly within the
timing slack:
\[\Delta t_{\text{stall}} = \frac{V_{\text{burst}}}{R_{\text{arb}}} \le S\]
Solving for the burst volume yields the maximum admissible ingestion
load: \[V_{\text{burst}} \le S \times R_{\text{arb}}\] The admissible
sensory payload is bounded by the smallest timing slack in the fastest
enforcement reflex. In an analytical example where an enforcement
monitor runs with a timing slack of \(S = 250\text{ }\mu\text{s}\) and
the memory controller delivers an effective service rate of
\(R_{\text{arb}} = 20.0\text{ GB/s}\) during write contention, the total
sensory ingestion burst cannot exceed:
\[V_{\text{burst}} \le 250\text{ }\mu\text{s} \times 20.0\text{ GB/s} = 5.0\text{ MB}\]
Any sensor configuration whose combined DMA write burst exceeds
\(5.0\text{ MB}\) will stall the memory controller long enough to
exhaust the slack, causing the safety monitor to miss its deadline
regardless of the average bandwidth.

\protect\phantomsection\label{callout-ruleux2a-8.7}
\begin{fbxSimple}{callout-rule}{Architect’s Rule of Thumb:}{Ingress Bandwidth as a Safety Parameter}
\phantomsection\label{callout-rule*-8.7}

Sensor resolution, frame rate, and channel count are safety parameters,
not perception parameters; every megabyte of sensory DMA burst consumes
memory bus service time directly from the safety refusal path.

\end{fbxSimple}

This relationship establishes that sensor resolution, frame rate, and
sensor count are safety parameters rather than purely perception
parameters. In machine learning pipelines, engineering teams routinely
increase camera resolution or add sensor modalities to improve open-loop
model accuracy. In a physical machine, every additional pixel, bit of
quantization depth, and exposure frame is an unavoidable withdrawal from
the memory bandwidth budget that pays for timely refusal. Upgrading four
cameras from \(1080\text{p}\) (\(1920 \times 1080\)) to \(4\text{K}\)
(\(3840 \times 2160\)) quadruples the burst volume \(V_{\text{burst}}\)
from \(3.11\text{ MB}\) to \(12.44\text{ MB}\) per frame. On a system
with the timing margins derived above, that single change pushes the
worst-case memory stall from \(155.5\text{ }\mu\text{s}\) to
\(622.0\text{ }\mu\text{s}\), exceeding the \(250\text{ }\mu\text{s}\)
slack and causing the safety enforcement loop to overrun its period. The
higher-resolution image may produce a more confident detection in the
learned model, but it removes the machine's guarantee of stopping in
time.\sidenote{\footnotesize To bound memory access latency under DMA bursts, SoCs
  implement ARM MPAM (ISO 26262 ASIL D) and AXI QoS priority levels to
  ensure safety monitor reads preempt camera streams.}

The hardware mechanisms governing how frames move into memory, including
PCIe scatter-gather descriptor rings, I/O memory management unit (IOMMU)
page tables, and DRAM bank conflict schedulers, are standard computer
systems topics treated in depth by Hennessy and Patterson
(\citeproc{ref-hennessy2019computer}{Hennessy and Patterson 2019}) as
well as Bryant and O'Hallaron (\citeproc{ref-bryant2015computer}{Bryant
and O'Hallaron 2015}). For the physical AI architect, the implementation
details of descriptor rings sit below the floor. The quantity to budget
is the latency penalty that sensory ingress inflicts on safety-critical
memory transactions. When perception traffic claims the bus, the delay
propagates directly into the physical domain. Every millisecond that
memory contention steals from the enforcement loop is a millisecond
during which the body continues along its existing trajectory without
permission, converting a digital stall in the memory controller into an
irrevocable loss of physical stopping distance.

\section{Spatial Error, 3D Foundation Models, and Physical
Clearance}\label{sec-perception-8-5}

A spatial error is not a number in a representation. When downstream
control logic uses an estimated coordinate to authorize motion, that
error becomes a physical displacement, and it is spent out of the same
clearance that keeps the machine off the obstacle. Transforming a sensor
observation into physical authority requires \textbf{metric grounding},
which maps raw detector measurements into geometric space. This
grounding relies on unprojecting two-dimensional sensor coordinates into
three dimensions, an inversion that assumes known optics, rigid
mounting, and predictable surface properties.

The mathematical principles of projection and triangulation, detailed by
Hartley and Zisserman (\citeproc{ref-hartley2003multiple}{Hartley and
Zisserman 2003}) as well as Szeliski
(\citeproc{ref-szeliski2022computer}{Szeliski 2022}), govern how
measurement noise converts into spatial displacement. In passive stereo
vision, depth \(z\) is computed from baseline \(b\), focal length \(f\),
and disparity \(d\) through the relationship \(z = \frac{f b}{d}\).
Differentiating this unprojection with respect to disparity yields the
depth uncertainty: \[\delta z \approx \frac{z^2}{f b} \delta d\] where
\(\delta d\) is the subpixel disparity matching error.\sidenote{\footnotesize While
  image disparity error is Gaussian
  (\(\delta d \sim \mathcal{N}(0, \sigma_d^2)\)), non-linear
  unprojection \(z = fb/d\) creates asymmetric heavy tails in 3D
  Euclidean space. At long ranges (\(d \to 0\)), depth variance
  diverges, requiring estimators to parameterize landmarks in inverse
  depth (\(\rho = 1/z\)).} Because depth error scales quadratically with
distance \(z\), doubling the operating range quadruples the physical
depth uncertainty. In contrast, time-of-flight sensors such as pulsed
LiDAR calculate distance from photon transit time
\(z = \frac{c \Delta t}{2}\), where \(c\) is the speed of light and
\(\Delta t\) is the measured round-trip time. The resulting range
uncertainty depends on clock jitter and detector rise time, remaining
approximately constant across the operational range of the sensor, as
demonstrated in the dense urban point cloud scan in
Figure~\ref{fig-08-lidar-pointcloud}.

\textbf{Inverse-depth metric grounding} is the geometric
parameterization of visual range observations in inverse-distance
coordinates (\(\rho = 1/z\)) to preserve linear Gaussian measurement
error properties and prevent Kalman filter covariance divergence under
the heavy-tailed, quadratic spatial error distributions
(\(\delta z \approx \frac{z^2}{fb}\delta d\)) inherent to passive stereo
triangulation and monocular unprojection.

Consider an analytical example of a stereo camera with baseline
\(b = 0.20\text{ m}\), focal length \(f = 800\text{ px}\), and a
disparity matching uncertainty of \(\delta d = 0.5\text{ px}\). At a
distance of \(z = 2.0\text{ m}\), the expected depth uncertainty is:
\[\delta z = \frac{(2.0\text{ m})^2}{800\text{ px} \times 0.20\text{ m}} \times 0.5\text{ px} = 0.0125\text{ m} = 12.5\text{ mm}\]
At a distance of \(z = 10.0\text{ m}\), the same disparity uncertainty
yields:
\[\delta z = \frac{(10.0\text{ m})^2}{800\text{ px} \times 0.20\text{ m}} \times 0.5\text{ px} = 0.3125\text{ m} = 312.5\text{ mm}\]
At ten meters, the depth uncertainty alone exceeds the total physical
envelope of many small mobile robots.

\begin{figure}[t]

\centering{

\includegraphics[width=0.85\linewidth,height=\textheight,keepaspectratio]{/Users/VJ/GitHub/MLSysBook-integrate-physical/physical/book/chapters/08-perception/figures/fig08_real_lidar_pointcloud.jpg}

}

\caption{\label{fig-08-lidar-pointcloud}\textbf{Registered 3D LiDAR
Point Cloud in an Urban Environment.} High-resolution orthographic
projection of 23 million time-of-flight points acquired by an Ouster
OS1-64 multi-beam LiDAR mounted on a moving vehicle navigating an urban
intersection (Folsom and Dore St, San Francisco). Points are registered
in real time into the global coordinate frame
(\(\mathcal{F}_{\text{world}}\)) via simultaneous localization and
mapping (SLAM) and colorized by calibrated return intensity scaled by
range. Because time-of-flight ranging uncertainty
(\(z = c \Delta t / 2\)) remains approximately constant across operating
range, LiDAR generates dense metric surfaces for physical clearance
estimation without the quadratic depth degradation
(\(\delta z \propto z^2\)) characteristic of passive stereo
triangulation. Continuous motion compensation across sweeping scanlines
is essential to prevent vehicle velocity from shearing the reconstructed
geometry. Attribution: Ouster Inc.~public dataset, CC BY 4.0.}

\end{figure}%

\subsection{3D Foundation Models in Physical AI: DINOv2, SAM, and
ConceptGraphs}\label{d-foundation-models-in-physical-ai-dinov2-sam-and-conceptgraphs}

Modern physical AI architectures do not rely solely on classical
hand-engineered geometric detectors or closed-set bounding box
classifiers. Instead, they increasingly leverage \textbf{3D vision
foundation models} that combine zero-shot open-vocabulary semantic
understanding with metric spatial grounding:

\begin{enumerate}
\def\labelenumi{\arabic{enumi}.}
\tightlist
\item
  \textbf{Self-Supervised Vision Backbones (DINOv2)}: DINOv2
  (\citeproc{ref-oquab2023dinov2}{{Oquab et al.} 2023}) uses vision
  transformer (ViT) architectures trained via self-distillation on
  massive visual datasets.\sidenote{\footnotesize Vision Transformers (DINOv2)
    partition images into \(P \times P\) patches
    (\(14 \times 14\text{ px}\)), projecting them into \(N = (HW)/P^2\)
    patch tokens. All-to-all self-attention preserves fine geometric
    correspondences and part-level semantics across views without
    requiring explicit task-specific supervised training.} Unlike
  supervised models that collapse fine spatial details into discrete
  class labels, DINOv2 produces dense patch token representations where
  feature vectors preserve powerful geometric and semantic properties.
  Patch tokens from DINOv2 naturally encode dense cross-view
  correspondences, part-level semantics, and relative depth cues without
  any task-specific fine-tuning.
\item
  \textbf{Promptable Zero-Shot Segmentation (SAM)}: The Segment Anything
  Model (SAM) (\citeproc{ref-kirillov2023segment}{{Kirillov et al.}
  2023}) delivers class-agnostic instance mask segmentation across
  arbitrary visual scenes. By running a heavyweight image encoder once
  per frame and evaluating fast, promptable mask decoders, SAM
  partitions complex physical scenes into coherent object surfaces.
\item
  \textbf{Open-Vocabulary 3D Scene Graphs (ConceptGraphs)}: To translate
  2D foundation model representations into actionable physical
  representations, systems such as ConceptGraphs
  (\citeproc{ref-kuwajerwala2024conceptgraphs}{{Kuwajerwala et al.}
  2024}; \citeproc{ref-jatavallabhula2023conceptfusion}{{Jatavallabhula
  et al.} 2023}) and Hydra (\citeproc{ref-hughes2022hydra}{Hughes et al.
  2022}) lift 2D instance masks and dense semantic features into 3D
  metric space:

  \begin{itemize}
  \tightlist
  \item
    For every video frame, 2D segmentation masks from SAM and dense
    semantic feature embeddings from DINOv2 / CLIP are extracted.
  \item
    Using calibrated depth maps (from active stereo or LiDAR) and camera
    poses \(\mathbf{T}_{wb}(t)\), the 2D masked pixels are unprojected
    into 3D point cloud clusters in the body or map frame.
  \item
    Redundant 3D object proposals across sequential camera viewpoints
    are merged by evaluating both geometric overlap (3D bounding box
    Intersection over Union, \(\text{IoU}_{3\text{D}}\)) and semantic
    cosine similarity in the foundation embedding space.
  \item
    The resulting output is an \textbf{object-centric 3D scene graph},
    where each node represents a distinct physical object carrying an
    oriented 3D bounding box, a centroid position
    \(\mathbf{p} \in \mathbb{R}^3\), a spatial covariance
    \(\mathbf{\Sigma}_p \in \mathbb{R}^{3\times 3}\), and an
    open-vocabulary text-aligned embedding vector.
  \end{itemize}
\end{enumerate}

\textbf{3D foundation metric lifting} is the multi-rate systems
architecture that projects dense 2D foundation model representations
(such as DINOv2 self-supervised patch tokens and SAM instance
segmentation masks) into 3D metric entities (such as ConceptGraphs
object-centric nodes) with calibrated spatial covariances
(\(\mathbf{\Sigma}_p\)) and temporal freshness bounds.

From an ML Systems perspective, deploying foundation models on physical
machines introduces a fundamental compute and latency trade-off. Running
DINOv2 and SAM requires between \(10^9\) and \(10^{11}\text{ FLOPs}\)
per forward pass, consuming hundreds of megabytes of memory bandwidth
and exhibiting execution latencies of \(100\text{--}300\text{ ms}\) on
embedded NPUs (\citeproc{ref-lin2018architectural}{{Lin et al.} 2018}).
An autonomous machine moving at \(2.0\text{ m/s}\) cannot close
high-rate feedback loops over a \(300\text{ ms}\) foundation model
cycle.

Physical AI architectures solve this by partitioning perception into a
\textbf{multi-rate hierarchy}:

\begin{itemize}
\tightlist
\item
  \textbf{High-Rate Metric Ingress (\(30\text{--}60\text{ Hz}\))}: The
  Linux MPU runs hardware ISP debayering and fast metric depth
  unprojection directly from MIPI CSI-2 DMA buffers, generating
  instantaneous 3D obstacle point sets and occupancy grids.
\item
  \textbf{Asynchronous Foundation Graph Updates
  (\(1\text{--}5\text{ Hz}\))}: In parallel worker threads, the NPU
  evaluates DINOv2 patch embeddings and SAM segmentation masks on
  keyframes, incrementally updating the semantic nodes of the 3D
  ConceptGraph.
\item
  \textbf{Deterministic Real-Time Enforcement (\(1000\text{ Hz}\))}: The
  real-time MCU queries the metric clearance envelope defended around
  the robot, verifying that no physical obstacle intrudes into the
  clearance inset, completely decoupled from the execution jitter of the
  MPU foundation model pipeline.
\end{itemize}

\subsection{Clearance Inset Margins and Transform Chain
Propagation}\label{clearance-inset-margins-and-transform-chain-propagation}

Spatial error does not remain confined to the sensor optical frame. To
govern physical motion, the estimated point must be propagated through a
chain of coordinate transformations to the robot base frame and out to
the end effector or vehicle perimeter. The mechanics of rigid body
transformations and kinematic frame composition are standard results
treated by Craig (\citeproc{ref-craig2005introduction}{Craig 2005}) as
well as Lynch and Park (\citeproc{ref-lynch2017modern}{Lynch and Park
2017}). Each link in this transform chain contributes scale errors,
mounting calibration drift, joint encoder quantization, and mechanical
compliance under load. Angular errors at the sensor mount amplify
linearly with lever-arm distance. In an analytical case where an
extrinsic calibration shifts by \(\theta = 0.005\text{ rad}\)
(\(0.29^\circ\)) between a camera mount and the vehicle chassis, the
resulting lateral displacement error is
\(\delta x \approx z \theta = 50.0\text{ mm}\) at a target distance of
\(z = 10.0\text{ m}\). When combined with quadratic depth uncertainty,
the total spatial error vector
\(\boldsymbol{\epsilon}_{\text{spatial}}\) at long range is dominated by
depth along the optical axis and angular lever-arm effects orthogonal to
it.

Every physical safety barrier operates by defending a geometric boundary
around the machine. If an enforcement monitor enforces a nominal
clearance margin \(C_{\text{nominal}}\) between the body envelope and a
detected obstacle, the true physical clearance \(C_{\text{true}}\) is
reduced by the upstream spatial error:
\[C_{\text{true}} = C_{\text{nominal}} - \|\boldsymbol{\epsilon}_{\text{spatial}}\|\]
To guarantee that the machine never breaches a physical safety limit
\(C_{\text{safe}}\), the safety enforcer must defend an \textbf{inset
clearance} \(C_{\text{inset}}\):
\[C_{\text{inset}} = C_{\text{safe}} + \|\boldsymbol{\epsilon}_{\text{spatial}}\|_{P_{99}}\]
where \(\|\boldsymbol{\epsilon}_{\text{spatial}}\|_{P_{99}}\) represents
the ninety-ninth percentile bound of the propagated spatial error.
Across physical archetypes, this inset defending \(C_{\text{safe}}\)
enforces distinct operational envelopes: in Mobile Systems, defending a
\(C_{\text{safe}} = 200\text{ mm}\) vehicle corridor buffer against
industrial racking; in Manipulators, defending a
\(C_{\text{safe}} = 5.0\text{ mm}\) tool clearance during high-speed
fixture approach where link compliance and joint backlash compound
optical errors; and in Process / Energy Systems, defending a
\(C_{\text{safe}} = 1.5\text{ mm}\) standoff distance for laser
deposition print heads over warped metallic substrate plates. As
illustrated in Figure~\ref{fig-08-depth-uncertainty-clearance}, depth
uncertainty divergence dictates how sensor noise converts into physical
encroachment. While time-of-flight LiDAR exhibits a constant range error
floor across distance, passive stereo triangulation uncertainty diverges
quadratically (\(\delta z \propto z^2\)). When compounded by angular
mounting compliance (\(\delta x \approx z \theta\)), distant spatial
uncertainty ellipsoids expand along the optical axis, directly
subtracting from physical clearance. A nominal safety margin that has
not been inset by the upstream perception error is a fiction. If an
enforcer commands a stop when nominal clearance reaches
\(100\text{ mm}\), but the perception pipeline carries an unmodeled
spatial uncertainty of \(150\text{ mm}\), the physical structure will
collide with the obstacle while the monitor registers a safe state.

The \textbf{clearance inset margin} is the mathematically enforced
buffer expansion
(\(C_{\text{inset}} = C_{\text{safe}} + \|\boldsymbol{\epsilon}_{\text{spatial}}\|_{P_{99}}\))
added to nominal obstacle clearance boundaries to account for upstream
spatial perception uncertainty, uncalibrated extrinsic lever arms
(\(\delta x \approx z\theta\)), and quadratic depth degradation
(\(\delta z \propto z^2\)), ensuring physical safety limits
(\(C_{\text{safe}}\)) are never breached.

\begin{figure}[t]

\centering{

\includegraphics[width=0.95\linewidth,height=\textheight,keepaspectratio]{index_files/mediabag//Users/VJ/GitHub/MLSysBook-integrate-physical/physical/book/chapters/08-perception/figures/fig08_depth_uncertainty_clearance.pdf}

}

\caption{\label{fig-08-depth-uncertainty-clearance}\textbf{Metric
Grounding, Depth Error Scaling, and Clearance Inset Geometry.} (A) Depth
uncertainty \(\delta z\) as a function of target distance \(z\),
contrasting quadratic stereo vision divergence
(\(\delta z \approx \frac{z^2}{fb}\delta d\)) with constant pulsed LiDAR
time-of-flight error (\(\delta z \approx \frac{c\Delta t}{2}\)). (B)
Overhead geometric clearance diagram of a mobile platform approaching an
obstacle. The propagated spatial uncertainty ellipsoid
\(\|\boldsymbol{\epsilon}_{\text{spatial}}\|_{P_{99}}\) (depth axis
error \(\delta z\) and transverse lever-arm error
\(\delta x \approx z\theta\)) directly subtracts from nominal clearance,
requiring the safety enforcer to defend an expanded inset margin
\(C_{\text{inset}} = C_{\text{safe}} + \|\boldsymbol{\epsilon}_{\text{spatial}}\|_{P_{99}}\)
to guarantee the physical collision threshold \(C_{\text{safe}}\) is
never breached. Attribution: Book geometric error analysis, CC BY 4.0.}

\end{figure}%

Because clearance depends on propagated error, downstream planners and
safety monitors cannot consume raw coordinate tuples. An observation
crossing the perception boundary must satisfy an explicit interface
contract defining its metric meaning. Every spatial estimate must
specify its target coordinate frame, its acquisition epoch timestamp,
and its bounded uncertainty representation, such as a spatial covariance
matrix \(\mathbf{\Sigma} \in \mathbb{R}^{3 \times 3}\) or a bounding
radius. An observation that delivers \((x, y, z) = (4.2, -1.1, 0.8)\)
without stating the reference frame, the time at which the photons
struck the detector, and the confidence ellipsoid cannot be integrated
into a safety case. Downstream control loops cannot deduce whether the
measurement warrants a \(10\text{ mm}\) inset or a \(500\text{ mm}\)
inset.

When spatial error is bounded and quantified, the nervous system can
compensate by expanding the clearance inset, trading vehicle speed or
workspace volume for physical separation. However, when spatial error is
unbounded or unknown, widening the margin is mathematically undefined.
If a stereo camera encounters a textureless wall, optical glare that
causes sensor saturation, or a lost calibration parameter, the disparity
variance diverges. Under unbounded or unquantified uncertainty, the
perception subsystem cannot furnish a valid spatial bound, and the
system must immediately revoke permission for autonomous motion.
Expanding a clearance margin assumes that the error distribution has a
known tail; when that assumption fails, continuing motion relies on
chance rather than engineering.

The complete operational protocol for metric unprojection, kinematic
covariance propagation, and dynamic clearance inset sizing is formalized
below:

\subsection{Spatial Clearance Inset and Depth Covariance Projection
Protocol}\label{spatial-clearance-inset-and-depth-covariance-projection-protocol}

\textbf{Cadence \& Invariants:}

\begin{itemize}
\tightlist
\item
  \emph{Execution Cadence:} Synchronous with perception frame ingress
  (\(30\text{--}60\text{ Hz}\) on Linux MPU / Hardware ISP).
\item
  \emph{Memory Invariant:} Zero dynamic heap allocation in inner
  unprojection loops; bounded matrix scratchpad buffers.
\item
  \emph{Interface Contract:} Emits bounded spatial claims
  \((\mathbf{p}_{\text{body}}, \mathbf{\Sigma}_{\text{body}}, C_{\text{inset}}, \tau_{\text{age}})\)
  to the shared-memory spatial scene graph.
\end{itemize}

\textbf{Execution Protocol:}

\begin{enumerate}
\def\labelenumi{\arabic{enumi}.}
\item
  \textbf{Information Age \& Temporal Freshness Verification (Stage 1):}
  Compute total information age
  \(\tau_{\text{age}} \leftarrow t_{\text{curr}} - t_0\). If
  \(\tau_{\text{age}} > \tau_{\text{max\_admissible}}\), assert
  \texttt{UNBOUNDED\_ERROR\_FAULT} and trigger stale-frame refusal.
  Otherwise, compute unmodeled kinematic displacement vector
  \(\Delta \mathbf{p}_{\text{age}} \leftarrow \mathbf{v} \cdot \tau_{\text{age}}\).
\item
  \textbf{Optical Metric Unprojection \& Sensor Covariance Synthesis
  (Stage 2):} For each measurement point
  \((u, v, z, \sigma_d^2) \in \mathcal{P}_{\text{sensor}}\):

  \begin{itemize}
  \item
    Compute 3D optical coordinates:
    \(x_s \leftarrow \frac{(u - c_x) z}{f_x}, \quad y_s \leftarrow \frac{(v - c_y) z}{f_y}, \quad z_s \leftarrow z\).
  \item
    For passive stereo triangulation (\(z = f b / d\)):
    \[\begin{aligned}
    \sigma_{z,s}^2 &\leftarrow \left( \frac{z^2}{f_x b} \right)^2 \sigma_d^2, \\
    \sigma_{x,s}^2 &\leftarrow \left(\frac{u - c_x}{f_x}\right)^2 \sigma_{z,s}^2 + \left(\frac{z}{f_x}\right)^2 \sigma_u^2, \\
    \sigma_{y,s}^2 &\leftarrow \left(\frac{v - c_y}{f_y}\right)^2 \sigma_{z,s}^2 + \left(\frac{z}{f_y}\right)^2 \sigma_v^2.
    \end{aligned}\] Construct diagonal sensor covariance:
    \(\mathbf{\Sigma}_{\text{sensor}} \leftarrow \operatorname{diag}(\sigma_{x,s}^2, \sigma_{y,s}^2, \sigma_{z,s}^2)\).
  \item
    For time-of-flight LiDAR (\(z = c \Delta t / 2\)): set range
    variance \(\sigma_{z,s}^2 \leftarrow \sigma_{\text{ToF}}^2\) and
    construct beam covariance \(\mathbf{\Sigma}_{\text{sensor}}\).
  \end{itemize}
\item
  \textbf{Rigid Kinematic Transform \& Spatial Covariance Projection
  (Stage 3):} Transform position to body coordinates
  \(\mathbf{p}_{\text{body}} \leftarrow \mathbf{R}_{\text{sensor}}^{\text{body}} \mathbf{p}_{\text{sensor}} + \mathbf{p}_{\text{sensor}}^{\text{body}} + \Delta \mathbf{p}_{\text{age}}\).
  Project spatial covariance via first-order error propagation across
  mounting extrinsics
  (\(\mathbf{J}_{\text{ext}} = [\mathbf{I}_{3\times 3}, -[\mathbf{R}_{\text{sensor}}^{\text{body}} \mathbf{p}_{\text{sensor}}]_\times]\)):
  \[\mathbf{\Sigma}_{\text{body}} \leftarrow \mathbf{J}_p \mathbf{\Sigma}_{\text{sensor}} \mathbf{J}_p^T + \mathbf{J}_{\text{ext}} \mathbf{\Sigma}_{\text{ext}} \mathbf{J}_{\text{ext}}^T + \|\mathbf{v}\|^2 \sigma_{\tau}^2 \mathbf{I}_{3\times 3}.\]
\item
  \textbf{Bounding Error Ellipsoid \& Dynamic Clearance Inset Sizing
  (Stage 4):} Compute maximum principal eigenvalue
  \(\lambda_{\max} \leftarrow \max \operatorname{eig}(\mathbf{\Sigma}_{\text{body}})\)
  and evaluate the \(99\text{th}\)-percentile error radius
  (\(\chi_3^2(0.99) \implies \kappa_{99} \approx 3.368\)):
  \[\|\boldsymbol{\epsilon}_{\text{spatial}}\|_{P_{99}} \leftarrow \kappa_{99} \sqrt{\lambda_{\max}}.\]
  If \(\lambda_{\max} > \lambda_{\text{divergence\_limit}}\) or
  \(\det(\mathbf{\Sigma}_{\text{body}}) \le 0\), assert
  \texttt{UNBOUNDED\_ERROR\_FAULT} and revoke motion authorization.
  Otherwise, size the dynamic clearance inset:
  \[C_{\text{inset}} \leftarrow C_{\text{safe}} + \|\boldsymbol{\epsilon}_{\text{spatial}}\|_{P_{99}}.\]
\item
  \textbf{Record Packaging \& Shared Memory Emission (Stage 5):}
  Populate claim record
  \(\mathcal{C}_{\text{body}} \leftarrow (\mathbf{p}_{\text{body}}, \mathbf{\Sigma}_{\text{body}}, C_{\text{inset}}, \tau_{\text{age}})\)
  and commit to the lock-free shared memory ring buffer with memory
  barrier (\texttt{DMB\ ISH}).
\end{enumerate}

A fundamental invariant emerges from this propagation chain: a nominal
clearance boundary that has not been expanded by the ninety-ninth
percentile propagated perception covariance is an optical fiction that
guarantees physical collision.

\section{How Observations Go Wrong}\label{sec-perception-8-6}

Perception failures divide into two distinct operational classes.
Detectable faults occur when the physical sensor or the communication
transport reports an explicit error condition. A serializer-deserializer
link drops a lock, a cyclical redundancy check detects a corrupted
payload, a frame buffer overflows, or a transport bus times out. These
failures are straightforward to handle because the hardware driver
asserts an error flag at the boundary of the nervous system. Downstream
planning algorithms immediately register the missing frame, invalidate
stale state estimates, and transition the machine into a degraded
operating mode or command a controlled stop. In contrast, \textbf{silent
failures} occur when the sensor hardware and transport pipeline function
without error, delivering a fully populated, structurally valid tensor
whose numerical values misrepresent the physical world.

Physical degradation of the optical channel represents the most common
source of silent error.\sidenote{\footnotesize Analogous silent degradation occurs in
  active time-of-flight LiDAR sensors using Single-Photon Avalanche
  Diodes (SPADs) or Avalanche Photodiodes (APDs). When a laser pulse
  strikes a highly retroreflective surface (such as a highway sign or
  safety garment), the intense photon return saturates the photodetector
  die, causing \emph{range walk} or detector blooming. The distorted
  pulse shape triggers the receiver timing discriminator prematurely,
  injecting range measurement biases of several centimeters without
  asserting any hardware failure flag.} When direct sunlight enters an
optical aperture, incident photons exceed the full-well capacity of the
CMOS photodiodes, clamping pixel intensities at their maximum
quantization limit across large patches of the array. This dynamic range
saturation destroys local spatial gradients. Feature extractors, optical
flow estimators, and disparity matching algorithms fail because
mathematical optimization on a uniform flat field yields zero
information, yet the driver emits a valid frame at nominal line rate.
Low-light environments introduce the opposite failure mode. As ambient
illumination drops, photon shot noise and sensor read noise dominate the
measurement. The signal-to-noise ratio collapses, causing downstream
neural networks to extract geometric features from random thermal
fluctuations rather than physical scene boundaries. Environmental
contaminants on the protective lens or dome window (such as water
droplets, airborne dust, or mud splatters) act as irregular refractive
elements and optical low-pass filters. They attenuate contrast, distort
spatial geometry, and obscure obstacles while the transport bus
continues to operate without asserting a single fault flag. Across
physical archetypes, these silent degradation modes take distinct
manifestations: in Mobile Systems, low-angle sun glare saturates forward
cameras on highway asphalt or rain droplets refract LiDAR beams; in
Manipulators, specular reflections off polished metallic workpieces
blind active stereo depth cameras and coolant mist coats eye-in-hand
lenses; and in Process / Energy Systems, metallic vapor plumes,
incandescent spatter, and soot deposits attenuate pyrometer optical
sight glasses in laser additive powder-bed fusion, causing temperature
estimators to report nominal states while melt pool temperatures surge
uncontrollably.

Timing and transport anomalies can also corrupt an observation silently.
If an unshielded physical interface experiences electromagnetic
interference, deserializer bit slips can alter payload semantics without
invalidating standard packet lengths. Far more destructive than a lost
packet is an observation delivered with a timestamp that misreports its
true acquisition epoch. If an image sensor captures photons at epoch
\(t_0\), but an unbuffered PCIe interrupt handler or an unsynchronized
system clock timestamps the resulting memory buffer at
\(t_0 + \delta t\), the downstream state estimator integrates spatial
features under an incorrect body pose. Consider a machine traveling with
instantaneous velocity \(v\) and linear acceleration \(a\) along a given
trajectory. When perception suffers an undetected timing
desynchronization \(\delta t\), the state estimator assigns the detected
position of an obstacle to the body coordinates at \(t_0 + \delta t\)
rather than \(t_0\). Over the unmodeled time interval \(\delta t\), the
physical state estimation error \(\Delta p\) between the asserted
position and the true position is governed by the kinematic relation:
\[\Delta p = v \cdot \delta t + \frac{1}{2} a \cdot (\delta t)^2\]

In an analytical example of an autonomous vehicle operating at a
velocity of \(v = 20.0\text{ m/s}\) under an active braking deceleration
of \(a = -4.0\text{ m/s}^2\), an undetected timing delay of
\(\delta t = 50\text{ ms}\) (the duration of a single delayed frame
across a congested internal bus) produces a positional estimation error
of:
\[\Delta p = (20.0\text{ m/s})(0.050\text{ s}) + \frac{1}{2}(-4.0\text{ m/s}^2)(0.050\text{ s})^2 = 1.0\text{ m} - 0.005\text{ m} = 0.995\text{ m}\]
An undetected positional error of \(0.995\text{ m}\) exceeds the nominal
safety clearance of an industrial aisle or a standard roadway lane.
Because the pipeline delivers a continuous stream of valid tensors, the
state estimator assumes a fresh observation, computing safe actuator
trajectories from stale physics. A frame that is entirely missing causes
a known loss of information that a Kalman filter can handle by expanding
its state covariance. A frame that arrives on time with a corrupt
timestamp causes the filter to contract its covariance around an
incorrect mean, driving actuators toward an obstacle that the model
believes has already been cleared.

Defending against silent failure requires dedicated runtime health
monitors operating at the sensor ingestion boundary. A pixel intensity
histogram monitor measures the distribution of raw digitized values
immediately after direct memory access transfer, identifying
over-exposure saturation when values concentrate at the maximum bit
depth or identifying lens occlusion when spatial variance drops below an
empirical noise floor. A hardware watchdog timer tracks inter-arrival
times across the direct memory access channel, flagging transport jitter
when inter-frame intervals exceed
\(T_{\text{nominal}} + 3\sigma_{\text{jitter}}\). Higher up the stack,
an optical flow plausibility monitor compares computed pixel
displacement vectors against the physical velocity envelope of the body,
rejecting observations that imply accelerations beyond mechanical
capability. In multi-modal configurations, a cross-sensor consistency
check compares high-rate (\(1000\text{ Hz}\)) inertial measurement unit
integrations against visual odometry. If the visual pipeline reports
that the machine is stationary while the accelerometer measures a
sustained \(2.0\text{ m/s}^2\) specific force, the monitor detects the
discrepancy and trips a perception fault.

Every runtime health monitor operates under a fundamental boundary. A
single sensing channel cannot validate its own semantic interpretations.
A camera assembly with clean optics, balanced exposure, jitter-free
transport, and synchronized microsecond clocking can still present an
observation that leads a learned model to hallucinate a clear path or
ignore an obstacle. Because deep vision models project high-dimensional
inputs through complex non-linear manifolds, in-distribution optical
anomalies (such as a high-contrast shadow across a road, a painted mural
on the side of a transport container, or specular reflection from wet
asphalt) generate plausible pixel statistics, valid optical flow fields,
and high model confidence. A single-sensor monitor evaluates whether the
data collection process operated within nominal limits, not whether the
inferred semantics match physical reality. Validating semantic truth
requires either cross-validation against an orthogonal physical modality
(such as pulsed lidar or radar that does not share optical failure
modes) or a safety architecture that verifies physical clearance through
independent mechanical guarantees.

A comprehensive architectural comparison of primary spatial sensor
modalities, their physical transduction principles, spatial-temporal
limits, and environmental failure modes is provided in
Table~\ref{tbl-08-sensor-modalities}. Understanding these fundamental
physical trade-offs is essential for constructing multi-modal perceptual
architectures capable of surviving orthogonal sensor degradation.

\begin{longtable}[]{@{}
  >{\raggedright\arraybackslash}p{(\linewidth - 12\tabcolsep) * \real{0.1429}}
  >{\raggedright\arraybackslash}p{(\linewidth - 12\tabcolsep) * \real{0.1429}}
  >{\raggedright\arraybackslash}p{(\linewidth - 12\tabcolsep) * \real{0.1429}}
  >{\raggedright\arraybackslash}p{(\linewidth - 12\tabcolsep) * \real{0.1429}}
  >{\raggedright\arraybackslash}p{(\linewidth - 12\tabcolsep) * \real{0.1429}}
  >{\raggedright\arraybackslash}p{(\linewidth - 12\tabcolsep) * \real{0.1429}}
  >{\raggedright\arraybackslash}p{(\linewidth - 12\tabcolsep) * \real{0.1429}}@{}}
\caption{\textbf{Spatial Sensor Modality Comparison Matrix.}
Quantitative comparison of spatial sensing modalities across physical
transduction mechanisms, spatial and angular resolutions, operational
range limits, depth uncertainty scaling, ingestion latencies, bandwidth
requirements, and physical degradation failure modes. Combining
complementary sensing channels with orthogonal physical failure
mechanisms (e.g., pairing high-resolution optical cameras vulnerable to
glare with all-weather 77 GHz radar or active pulsed LiDAR) is
mathematically necessary to prevent silent perception failures from
compromising physical safety
margins.}\label{tbl-08-sensor-modalities}\tabularnewline
\toprule\noalign{}
\begin{minipage}[b]{\linewidth}\raggedright
Sensor Modality \& Physical Principle
\end{minipage} & \begin{minipage}[b]{\linewidth}\raggedright
Spatial / Angular Resolution
\end{minipage} & \begin{minipage}[b]{\linewidth}\raggedright
Effective Range (\(z_{\text{max}}\)) \& Uncertainty
\end{minipage} & \begin{minipage}[b]{\linewidth}\raggedright
Ingestion Latency \& Refresh Rate
\end{minipage} & \begin{minipage}[b]{\linewidth}\raggedright
Ingress Bandwidth Budget
\end{minipage} & \begin{minipage}[b]{\linewidth}\raggedright
Environmental Degradation \& Silent Failure Modes
\end{minipage} & \begin{minipage}[b]{\linewidth}\raggedright
Primary Autonomy Role
\end{minipage} \\
\midrule\noalign{}
\endfirsthead
\toprule\noalign{}
\begin{minipage}[b]{\linewidth}\raggedright
Sensor Modality \& Physical Principle
\end{minipage} & \begin{minipage}[b]{\linewidth}\raggedright
Spatial / Angular Resolution
\end{minipage} & \begin{minipage}[b]{\linewidth}\raggedright
Effective Range (\(z_{\text{max}}\)) \& Uncertainty
\end{minipage} & \begin{minipage}[b]{\linewidth}\raggedright
Ingestion Latency \& Refresh Rate
\end{minipage} & \begin{minipage}[b]{\linewidth}\raggedright
Ingress Bandwidth Budget
\end{minipage} & \begin{minipage}[b]{\linewidth}\raggedright
Environmental Degradation \& Silent Failure Modes
\end{minipage} & \begin{minipage}[b]{\linewidth}\raggedright
Primary Autonomy Role
\end{minipage} \\
\midrule\noalign{}
\endhead
\bottomrule\noalign{}
\endlastfoot
\textbf{RGB Global Shutter Camera}CMOS photodiode array with
simultaneous pixel charge transfer (\(400\text{--}700\text{ nm}\)) &
\(1920\times 1080\) to
\(3840\times 2160\text{ px}\)(\(\sim 0.02^\circ\text{/px}\)) & Monocular
bearing only; scale ambiguous without geometric priors &
\(8\text{--}33\text{ ms}\) latency(\(30\text{--}120\text{ Hz}\)) &
\(1.5\text{--}6.0\text{ Gbps}\)(\(12\text{--}16\text{-bit}\) RAW via
MIPI / GMSL2) & Optical saturation under direct sun
(\(>100\text{ klux}\)); SNR collapse in low light (\(<1\text{ lux}\));
lens contamination (water, mud, dust) & 2D semantic classification, fine
feature tracking, traffic signal / sign recognition \\
\textbf{Active / Passive Stereo Camera}Dual calibrated CMOS imagers;
passive ambient or active IR speckle (\(850\text{--}940\text{ nm}\)) &
\(1280\times 720\) to \(1920\times 1080\text{ depth px}\) &
\(0.2\text{--}15.0\text{ m}\); quadratic depth uncertainty:
\(\delta z = \frac{z^2}{fb}\delta d\) (\(\sigma_z \approx 12\text{ mm}\)
at \(2\text{ m}\), \(312\text{ mm}\) at \(10\text{ m}\)) &
\(15\text{--}35\text{ ms}\) latency(\(30\text{--}90\text{ Hz}\)) &
\(0.5\text{--}2.5\text{ Gbps}\)(\(16\text{-bit}\) depth map via USB3 /
CSI-2) & Textureless surfaces; direct sun washing out active IR pattern;
thermal baseline expansion (\(\alpha = 23\times 10^{-6}\text{ K}^{-1}\))
& Close-range metric clearance, manipulation visual servoing, dense 3D
obstacle avoidance \\
\textbf{Pulsed Time-of-Flight LiDAR}Pulsed laser diode
(\(905\text{ / }1550\text{ nm}\)) with SPAD/APD array; transit time
\(z = c\Delta t / 2\) & \(32\text{--}128\text{ beams}\);
\(0.05^\circ\text{--}0.2^\circ\) angular resolution &
\(0.5\text{--}250.0\text{ m}\); range uncertainty constant:
\(\sigma_z \approx 10\text{--}30\text{ mm}\) independent of range \(z\)
& \(33\text{--}100\text{ ms}\) sweep(\(10\text{--}30\text{ Hz}\)) &
\(50\text{--}300\text{ Mbps}\)(\(10^5\text{--}2\times 10^6\text{ pts/s}\)
via UDP Ethernet) & Atmospheric backscatter in dense fog, rain
(\(>25\text{ mm/h}\)), and dust; detector blooming on retroreflectors
(range walk) & Metric 3D geometric reconstruction, global map
localization, direct envelope defense \\
\textbf{77 GHz FMCW Radar}Millimeter-wave radar
(\(76\text{--}81\text{ GHz}\)); beat frequency for range, Doppler shift
for \(\dot{z}\) & \(1^\circ\text{--}5^\circ\) azimuth,
\(3^\circ\text{--}10^\circ\) elevation; range resolution
\(\Delta R \approx 0.05\text{--}0.15\text{ m}\) &
\(0.5\text{--}300.0\text{ m}\); range error
\(\sigma_R \approx 0.1\text{ m}\); direct radial velocity
\(\sigma_v \approx 0.1\text{ m/s}\) & \(20\text{--}50\text{ ms}\) chirp
sequence(\(20\text{--}50\text{ Hz}\)) &
\(5\text{--}50\text{ Mbps}\)(Target clusters via CAN-FD / 100BASE-T1) &
Multipath ghost reflections; elevation angle ambiguity (overhead sign
clutter); specular bounce off angled sheet metal & All-weather
long-range obstacle tracking, direct radial velocity feedback in
fog/glare \\
\textbf{Neuromorphic Event Camera}Asynchronous pixel circuits detecting
log illumination changes \(\Delta \ln I \ge \pm C\) & \(640\times 480\)
to \(1280\times 720\text{ px}\) & Bearing and optical flow only;
requires ego-motion for depth triangulation &
\(1\text{--}10\text{ }\mu\text{s}\) latency(Continuous asynchronous
stream) & \(10\text{--}200\text{ Mbps}\)(Dynamic:
\(10^6\text{--}5\times 10^7\text{ events/s}\)) & Zero output on static
scenes without motion; vibration-induced background event noise;
high-frequency flicker saturation & Ultra-high-speed dynamic obstacle
avoidance, high-speed visual odometry under extreme dynamic range
(\(>120\text{ dB}\)) \\
\end{longtable}

The catastrophic potential of undetected, silent perception failures is
demonstrated in the investigation of the May 2016 collision near
Williston, Florida, involving an automated passenger vehicle and a
perpendicular tractor-semitrailer (\citeproc{ref-ntsb2017tesla}{National
Transportation Safety Board 2017}). The vehicle was cruising at
\(33.0\text{ m/s}\) (\(74\text{ mph}\)) under automated lateral and
longitudinal control on a divided highway under bright midday sun
(\(110{,}000\text{ lux}\)). As a 53-foot refrigerated semitrailer turned
left across the vehicle's path, the vision subsystem's CMOS camera
saturated against the high-contrast white side of the trailer set
against the brightly lit sky, failing to identify an obstacle contour.
Concurrently, the forward 77 GHz radar received strong electromagnetic
returns from the trailer body. However, because radar reflections from
elevated highway overpasses and overhead signs share similar geometric
signatures with high-riding trailer beds, the radar tracking pipeline's
heuristic filter discarded the returns as stationary overhead clutter to
prevent false-positive emergency braking events. Because neither sensor
pipeline reported a fault condition or asserted a degraded confidence
flag, the fusion module concluded the roadway was clear, commanding zero
braking torque until impact.

\protect\phantomsection\label{callout-incidentux2a-8.8}
\begin{fbxSimple}{callout-incident}{Landmark Engineering Incident:}{Optical Dynamic Range Saturation and Radar Heuristic Suppression Collision}
\phantomsection\label{callout-incident*-8.8}

\textbf{System:} Production automated passenger vehicle equipped with a
forward-facing CMOS camera (Mobileye EyeQ3), forward 77 GHz radar
(Continental ARS4-B), ultrasonic array, and electronic braking
controller.

\textbf{Incident:} Fatal perpendicular collision with a 53-foot
tractor-semitrailer at \(33.0\text{ m/s}\) (\(74\text{ mph}\)) under
bright midday sun on US Highway 27A near Williston, Florida, with zero
deceleration commanded prior to impact.

\textbf{Telemetry Snapshot:} At \(t = -9.0\text{ s}\), vehicle cruises
at \(33.0\text{ m/s}\) with lane-keeping and cruise active; at
\(t = -3.5\text{ s}\) (range \(\approx 115\text{ m}\)), a white
semitrailer enters perpendicular crossing path; camera optical array
saturates against the white trailer side under sky glare
(\(110{,}000\text{ lux}\)), yielding no obstacle bounding box; 77 GHz
radar detects trailer returns but tracking filter discards them as
stationary overhead sign clutter; at \(t = -0.5\text{ s}\), perception
continues asserting a clear roadway; at \(t = 0.0\text{ s}\), vehicle
strikes trailer without braking.

\textbf{Root Cause \& Physics of Failure:} Optical dynamic range
saturation eliminated edge contrast on the white trailer, while the
radar tracking filter's elevation ambiguity heuristic discarded real
mass as an overhead sign to suppress false braking.

\textbf{Architectural Flaw:} Absence of an explicit spatial uncertainty
contract and cross-modal disagreement escalation; the perception stack
allowed two independent silent failure modes (optical saturation and
radar elevation filtering) to produce an ungrounded consensus of an open
road without asserting a degraded confidence flag.\sidenote{\footnotesize National
  Transportation Safety Board (NTSB) statutory highway accident
  investigation report NTSB/HAR-17/02
  (\citeproc{ref-ntsb2017tesla}{National Transportation Safety Board
  2017}), \emph{Collision Between a Car Operating With Automated Vehicle
  Control Systems and a Tractor-Semitrailer Truck Near Williston,
  Florida, May 7, 2016}.}

\textbf{Provenance:} National Transportation Safety Board (NTSB)
statutory highway accident investigation report NTSB/HAR-17/02
(\citeproc{ref-ntsb2017tesla}{National Transportation Safety Board
2017}).

\end{fbxSimple}

\protect\phantomsection\label{callout-ruleux2a-8.9}
\begin{fbxSimple}{callout-rule}{Architect’s Rule of Thumb:}{The Timestamp Integrity Law}
\phantomsection\label{callout-rule*-8.9}

A sensor frame that arrives on time with a corrupted timestamp is far
more dangerous than a dropped frame; dropouts force estimators to expand
covariance, whereas false timestamps contract covariance around a
physical lie.

\end{fbxSimple}

When an observation crosses from the perception pipeline into the
nervous system, downstream tracking loops and safety enforcers cannot
deduce whether the underlying sensors were saturated, delayed, or
misaligned. The perception subsystem cannot simply output a tensor and
rely on downstream components to infer its integrity. To prevent silent
failures from corrupting physical control, every observation must arrive
packaged with an explicit, immutable record of its temporal origin,
reference frame, and validated health metrics.

\section{The Observation Contract}\label{sec-perception-8-7}

The boundary between the perception pipeline and the nervous system is a
memory interface governed by an explicit serialization contract. A
downstream tracking loop, state estimator, or safety monitor cannot
inspect camera registers, probe optical hardware, or reconstruct sensor
exposure history. If the perception pipeline provides only an isolated
array of floats or a raw pixel buffer, downstream consumers are forced
to make unverified assumptions about when the data was captured, where
the sensor was pointed, and how much noise corrupts the signal. To
eliminate these unstated assumptions, the perception pipeline packages
every observation into an immutable \textbf{observation record}, laid
out in contiguous, byte-aligned memory. The top-level message envelope,
transport header, and schema versioning adhere to the universal notebook
contract established in Chapter~\ref{sec-nervous}.
Chapter~\ref{sec-perception} extends that baseline envelope by defining
the perception payload, consisting of five mandatory physical channels,
three metadata blocks, and a fixed-width validity bitfield.

The \textbf{observation record} is the immutable engineering notebook
schema extension that encapsulates raw sensor measurements and learned
feature tokens with nanosecond-precise energy integration timestamps
(\(t_{\text{exp-mid}}\)), kinematic reference frame keys
(\(\mathcal{F}_{\text{sensor}}\)), spatial uncertainty covariance
matrices (\(\mathbf{\Sigma}_z\)), thermal calibration offsets, and a
\(64\text{-bit}\) hardware health bitmask.

\begin{longtable}[]{@{}
  >{\raggedright\arraybackslash}p{(\linewidth - 6\tabcolsep) * \real{0.2500}}
  >{\raggedright\arraybackslash}p{(\linewidth - 6\tabcolsep) * \real{0.2500}}
  >{\raggedright\arraybackslash}p{(\linewidth - 6\tabcolsep) * \real{0.2500}}
  >{\raggedright\arraybackslash}p{(\linewidth - 6\tabcolsep) * \real{0.2500}}@{}}
\caption{\textbf{Physical AI Observation Record Schema.} Immutable
schema fields added to the base engineering notebook to record physical
transduction epochs, coordinate frame bindings, spatial uncertainty
covariances, temperature compensation parameters, and hardware health
flags.}\label{tbl-08-observation-record-schema}\tabularnewline
\toprule\noalign{}
\begin{minipage}[b]{\linewidth}\raggedright
Record Schema Field
\end{minipage} & \begin{minipage}[b]{\linewidth}\raggedright
Physical / Semantic Modality
\end{minipage} & \begin{minipage}[b]{\linewidth}\raggedright
Cryptographic \& Provenance Guarantee
\end{minipage} & \begin{minipage}[b]{\linewidth}\raggedright
Downstream Verifier \& Usage
\end{minipage} \\
\midrule\noalign{}
\endfirsthead
\toprule\noalign{}
\begin{minipage}[b]{\linewidth}\raggedright
Record Schema Field
\end{minipage} & \begin{minipage}[b]{\linewidth}\raggedright
Physical / Semantic Modality
\end{minipage} & \begin{minipage}[b]{\linewidth}\raggedright
Cryptographic \& Provenance Guarantee
\end{minipage} & \begin{minipage}[b]{\linewidth}\raggedright
Downstream Verifier \& Usage
\end{minipage} \\
\midrule\noalign{}
\endhead
\bottomrule\noalign{}
\endlastfoot
\textbf{\texttt{capture\_timestamp\_epoch}} & \(64\text{-bit}\) integer
elapsed nanoseconds at physical photon integration midpoint
(\(t_{\text{exp-mid}}\)) & Hardware GPIO strobe latch / IEEE 1588 PTP
monotonic hardware counter & State estimator latency compensation;
computing spatial displacement \(\Delta x = v \tau_{\text{age}}\) \\
\textbf{\texttt{sequence\_index}} & Monotonically increasing
\(64\text{-bit}\) integer frame counter (\(N_{\text{seq}}\)) & Direct
hardware FPGA/CSI-2 host controller counter & Detecting dropped frames
and bus packet loss without clock dependence \\
\textbf{\texttt{sensor\_frame\_id}} & Kinematic transform key
(\(\mathcal{F}_{\text{sensor}}\)) binding observation to REP 105
hierarchy & Formal coordinate frame tree specification SHA-256 digest &
Unprojecting 2D/3D features into \texttt{base\_link} and \texttt{odom}
reference origins \\
\textbf{\texttt{tensor\_payload}} & Calibrated depth tensors, 3D point
clusters, or DINOv2 / SAM patch token embeddings & Fixed-width
contiguous memory buffer, zero-copy \texttt{dma\_buf} descriptor &
Feeding downstream 3D spatial scene graphs and neural policy
backbones \\
\textbf{\texttt{spatial\_covariance}} & Packed upper-triangular array
representing 3D spatial uncertainty
(\(\mathbf{\Sigma}_z \in \mathbb{R}^{3\times 3}\)) & Propagated optical
geometry: \(\sigma_z^2 = (z^2 / fb)^2 \sigma_d^2\) & Weighting Kalman
filter gain; expanding clearance insets
(\(C_{\text{inset}} = C_{\text{safe}} + \|\boldsymbol{\epsilon}\|_{P_{99}}\)) \\
\textbf{\texttt{calibration\_metadata}} & Focal lengths \((f_x, f_y)\),
principal point \((c_x, c_y)\), and active housing thermistor readings &
ISO 9001 factory calibration certificate hash, online temperature
coefficients & Dynamic thermal baseline drift compensation
(\(\Delta b = \alpha b \Delta T\)) \\
\textbf{\texttt{provenance\_digest}} & \(256\text{-bit}\) cryptographic
SHA-256 hash of sensor firmware build and model backbone weights &
Cryptographic binary signature, immutable weight tensor digest & Audit
verification against unrecorded model or firmware substitutions \\
\textbf{\texttt{health\_bitmask}} & \(64\text{-bit}\) unpadded bitmask:
lower \(32\text{ b}\) hardware health, upper \(32\text{ b}\) perceptual
validity & Real-time hardware watchdog monitors, single-instruction bit
check & Reflexive rejection of saturated, desynchronized, or occluded
frames \\
\end{longtable}

As formalized in the observation record schema
(Table~\ref{tbl-08-observation-record-schema}), an observation record is
not an ephemeral message passed between software threads. It is a
formal, immutable contract that specifies the physical validity of a
sensory measurement before any downstream process is permitted to trust
it. Where Chapter~\ref{sec-nervous} defines the base schema for all run
records, this section details the specific fields that the perception
subsystem adds to that foundation. The observation record combines raw
perceptual data with its spatial frame, uncertainty bounds, and hardware
health state in a single contiguous memory layout.

The formal schema of the observation record contains five mandatory
fields:

\begin{enumerate}
\def\labelenumi{\arabic{enumi}.}
\tightlist
\item
  \textbf{Exposure Midpoint Timestamp} (\(t_{\text{exp-mid}}\)): A
  \(64\text{-bit}\) integer recording nanoseconds since the UNIX epoch,
  latched at the midpoint of physical exposure by a hardware timer.
\item
  \textbf{Monotonic Sequence Index} (\(k\)): A \(64\text{-bit}\)
  unsigned integer incremented by the hardware capture engine on every
  frame, allowing downstream nodes to detect frame drops and
  out-of-order delivery.
\item
  \textbf{Reference Frame Identifier} (\(\mathcal{F}_{\text{sensor}}\)):
  A \(16\text{-bit}\) enumerated identifier that indexes into the
  spatial transform tree, resolving the sensor's physical mounting
  location relative to the body origin.
\item
  \textbf{Perceptual Payload}: The raw sensor tensor or spatial
  embedding, stored in a contiguous buffer with fixed byte alignment.
\item
  \textbf{Spatial Covariance Matrix} (\(\mathbf{\Sigma}_z\)): A
  \(6 \times 6\) positive semi-definite matrix encoding the spatial and
  angular uncertainty of the measurement in the sensor frame.
\end{enumerate}

Three metadata structures follow the mandatory physical fields to
preserve calibration provenance across varying operating environments:

\begin{enumerate}
\def\labelenumi{\arabic{enumi}.}
\tightlist
\item
  \textbf{Active Optical Intrinsics}: Records focal lengths
  \((f_x, f_y)\), optical center coordinates \((c_x, c_y)\), and radial
  distortion coefficients.
\item
  \textbf{Temperature Calibration Coefficients}: Stores thermal
  coefficients alongside real-time readings from thermistors mounted on
  the sensor housing. Mechanical structures expand with temperature
  changes. In an analytical calculation for an aluminum camera mounting
  bar with a thermal expansion coefficient of
  \(23 \times 10^{-6}\text{ K}^{-1}\), a temperature rise of
  \(25\text{ }^\circ\text{C}\) over an initial \(120\text{ mm}\)
  baseline causes a mechanical expansion of
  \(23 \times 10^{-6} \times 0.12\text{ m} \times 25\text{ K} \approx 6.9 \times 10^{-5}\text{ m} = 0.069\text{ mm}\).
  If left uncorrected, this thermal baseline drift induces systematic
  ranging errors in downstream triangulation.\sidenote{\footnotesize Thermal expansion
    of aluminum sensor brackets
    (\(\alpha_{\text{Al}} = 23\times 10^{-6}\text{ K}^{-1}\)) alters
    extrinsic baseline (\(b(T) = b_0(1+\alpha \Delta T)\)). In
    long-range stereo, an uncorrected \(0.1\text{ mm}\) shift at
    \(20\text{ m}\) induces a \(0.69\text{ m}\) range bias, requiring
    online extrinsic tracking.} The active temperature coefficients
  allow downstream estimators to adjust baseline and focal lengths
  dynamically during operation.
\item
  \textbf{Sensor Firmware and Model Revision Hash}: A \(256\text{-bit}\)
  cryptographic hash of the sensor firmware, register configurations,
  and foundation model backbone weights. This hash guarantees that
  downstream algorithms can verify the exact signal processing pipeline
  and model version used to acquire and encode the raw observations.
\end{enumerate}

Before reading the tensor payload or executing matrix operations,
downstream tracking loops and safety enforcers evaluate an unpadded
\(64\text{-bit}\) validity bitfield located at a fixed offset within the
record header. This bitfield operates as a formal execution contract
between the perception front-end and the nervous system. The lower
\(32\text{ bits}\) report hardware health status flags populated by the
runtime monitors established in Section~\ref{sec-perception-8-6},
including analog-to-digital saturation, rolling shutter synchronization
faults, optical interface transmission errors, and sensor temperature
limit trips. The upper \(32\text{ bits}\) convey perception validity
flags, such as partial lens occlusion markers, lens flare indicators,
and semantic confidence degradation flags. The contract dictates that if
any safety-critical health bit is set, downstream motion controllers
must treat the observation as invalid without evaluating the tensor
payload. Evaluating a \(64\text{-bit}\) bitmask requires one CPU
instruction (\(0.3\text{ ns}\) on a \(3.0\text{ GHz}\) core), allowing
reflex loops to reject corrupt sensory data with negligible latency
overhead before committing compute cycles to tensor operations.

The explicit byte structure of the observation record defines the fault
injection interface used for system verification. Rather than simulating
complex physical environments or modifying sensor hardware, fault
injectors in Chapter~\ref{sec-enforcement} manipulate observation
records directly on the shared memory bus before downstream delivery. An
injector can add a systematic offset of \(50\text{ ms}\) to
\(t_{\text{exp-mid}}\) to simulate clock drift, drop every fifth record
by introducing gaps in the sequence index, inflate the diagonal elements
of \(\mathbf{\Sigma}_z\) by a factor of ten, or set the partial
occlusion bit while leaving the raw tensor intact. Verification suites
in Chapter~\ref{sec-verification} execute these structured mutations
deterministically against candidate software builds. If a downstream
motion planner fails to transition into a safe fallback state when the
sensor health bitfield indicates a communication fault, or if an
estimator diverges because timestamp jitter reached \(100\text{ ms}\),
the verification suite pinpoints the failure directly in the software
interface. The record layout ensures that every failure mode observed on
physical test tracks can be reproduced, audited, and tested in
regression pipelines without physical hardware in the loop.

An observation record provides the complete spatial, physical, and
operational context of a measurement at a single instant in time. Yet
having a fully formed, mathematically validated observation record does
not guarantee that the measurement can be acted upon. Once the record is
committed to memory, its value decays with every microsecond of
transport delay, queuing stall, and processing latency across the system
bus.

\section{Systems Stress-Tests}\label{sec-perception-8-8}

An observation that arrives with zero sensor noise, perfect spatial
resolution, and high semantic classification confidence is useless if it
arrives late. In an embodied system moving through space, perceptual
validity decays with every millisecond of transport latency, queuing
stall, and processing jitter. When software queues or memory buses
decouple a measurement timestamp from the physical state of the machine,
acting on that measurement forces the controller to execute actions
against a world that no longer exists. The following three analytical
stress-tests examine the failure modes that emerge when timing,
synchronization, and interconnect bandwidth break down across the
perception pipeline:

\subsection*{Stress-Test 1: The High-Confidence Stale Observation (The
Expired Reality
Trap)}\label{stress-test-1-the-high-confidence-stale-observation-the-expired-reality-trap}
\addcontentsline{toc}{subsection}{Stress-Test 1: The High-Confidence
Stale Observation (The Expired Reality Trap)}

\emph{The scenario.} An autonomous mobile transport robot
(\(m = 450.0\text{ kg}\)) traverses an industrial corridor at
\(v = 3.0\text{ m/s}\). The front stereo perception pipeline feeds
bounding box proposals to a local collision avoidance planner at
\(20\text{ Hz}\) (\(50.0\text{ ms}\) period). A memory compaction pause
in the vision process delays tensor serialization by
\(\Delta t = 220.0\text{ ms}\), delivering an observation with a
\(0.98\) confidence score asserting that the corridor doorway remains
clear. During that \(220.0\text{ ms}\), an automated forklift moving
transversely at \(v_{\text{fork}} = 2.5\text{ m/s}\) advances
\(0.55\text{ m}\) into the robot's swept path.

\emph{The systems dilemma.} Sizing clearance against stale observations
converts pipeline delay into physical collision. Relying on the \(0.98\)
semantic confidence, the planner calculates a nominal clearance margin
of \(0.40\text{ m}\), yet true physical clearance has closed to zero.
For vehicle velocity \(v = 3.0\text{ m/s}\), obstacle speed
\(v_{\text{obs}} = 2.5\text{ m/s}\), and braking deceleration
\(a_{\text{max}} = 2.0\text{ m/s}^2\), the freshness ceiling is
\(t_{\text{max}} = 50.0\text{ ms}\). At \(\Delta t = 220.0\text{ ms}\),
an older observation is far more dangerous than an explicit sensor
dropout (\(o_t = \emptyset\)).

\emph{The systems verdict: Freshness Invariant.} Semantic certainty
cannot substitute for temporal freshness. In physical systems, the
nervous system must enforce a hard drop rule that discards any
observation exceeding its freshness ceiling (\(t_{\text{max}}\)),
falling back to dynamic model propagation under expanding covariance
envelopes.

\subsection*{Stress-Test 2: Inter-Sensor Clock Skew and Phantom
Obstacles (Temporal
Desynchronization)}\label{stress-test-2-inter-sensor-clock-skew-and-phantom-obstacles-temporal-desynchronization}
\addcontentsline{toc}{subsection}{Stress-Test 2: Inter-Sensor Clock Skew
and Phantom Obstacles (Temporal Desynchronization)}

\emph{The scenario.} A mobile platform tracks an oncoming transport
vehicle operating at \(v_{\text{target}} = 30.0\text{ m/s}\) using two
overlapping global-shutter cameras connected to separate interface
cards. Oscillator drift between the cards produces a steady clock skew
of \(\Delta t_{\text{skew}} = 15.0\text{ ms}\). Both cameras
independently detect the oncoming vehicle and tag records with local
timestamps.

\emph{The systems dilemma.} The spatial separation between camera
exposure instants is
\(\Delta d = v_{\text{target}} \Delta t_{\text{skew}} = 30.0\text{ m/s} \times 0.015\text{ s} = 0.45\text{ m}\).
When the central tracking filter fuses these observations assuming a
unified time base, the \(45.0\text{ cm}\) spatial discrepancy exceeds
the association gate for a single bounding box, splitting the target
into two distinct phantom tracks and forcing unnecessary emergency
evasive maneuvers.

\emph{The systems verdict: Clock Synchronization Law.} Spatial
calibration across multi-sensor suites is physically invalid without
microsecond-accurate temporal synchronization. To bound spatial
association error below \(\epsilon_{\text{pos}} = 0.05\text{ m}\), the
system must enforce hardware-level clock synchronization (IEEE 1588 PTP
or GPIO strobing) directly at sensor exposure circuits
(\(\Delta t_{\text{skew}} \le \epsilon_{\text{pos}} / v_{\text{max}}\)).

\subsection*{Stress-Test 3: Bus Saturation under Optical Transients
(Shared Interconnect
Starvation)}\label{stress-test-3-bus-saturation-under-optical-transients-shared-interconnect-starvation}
\addcontentsline{toc}{subsection}{Stress-Test 3: Bus Saturation under
Optical Transients (Shared Interconnect Starvation)}

\emph{The scenario.} An autonomous vehicle transitions from direct
sunlight (\(100{,}000\text{ lux}\)) into a dark loading bay
(\(20\text{ lux}\)) at \(v = 6.0\text{ m/s}\). The illumination change
triggers four cameras to adjust exposure gains simultaneously, bursting
raw pixel telemetry across the shared SoC interconnect at
\(4.8\text{ GB/s}\) against a memory controller limit of
\(5.2\text{ GB/s}\) (\(92\%\) bus utilization).

\emph{The systems dilemma.} Real-time safety monitors reading
high-frequency IMU and wheel encoder telemetry (\(128\text{ kB/s}\) at
\(1000\text{ Hz}\)) experience head-of-line blocking in the shared
memory controller, inflating access latency from \(15.0\ \mu\text{s}\)
to \(2.8\text{ ms}\) (\(P_{99.9}\)) and starving the closed-loop
traction controller.

\emph{The systems verdict: Interconnect Quality-of-Service Law.}
Perceptual data ingestion must never starve real-time safety reflexes.
The SoC architecture must enforce hardware AXI QoS priority registers,
isolating real-time safety channels on dedicated SRAM or high-priority
lanes while rate-limiting or dropping perception frames during transient
bus saturation.

\section{Summary}\label{sec-perception-8-9}

Perception in a physical AI system is the physical and computational
process of producing dated, localized, and bounded observation records.
An observation is not an abstract tensor of activations or a list of
uncalibrated bounding boxes generated by an isolated neural network. In
a machine that moves through physical space, sensory data begins as
physical energy interacting with a detector array, passes through
analog-to-digital converters and direct memory access controllers, and
undergoes neural processing before reaching downstream consumers. To be
usable by a learned policy or a state estimator, every observation
record must carry three physical invariants. First, it requires a
hardware-latched timestamp capturing the temporal midpoint of sensor
integration, allowing downstream loops to compute exactly how old the
observation is. Second, it requires an explicit geometric transform that
maps coordinates from the local sensor frame into the physical body
frame of the machine. Third, it requires an uncertainty envelope that
bounds measurement noise, lens distortion, and model dispersion. Without
these three metadata fields, perception delivers ungrounded numbers that
force downstream controllers to treat historical measurements as present
reality.

The resources required to ingest sensory data represent finite physical
budgets that must be accounted for during system design rather than
treated as software overhead. High-rate cameras, LiDAR units, and radar
arrays generate continuous streams of digital payloads that compete
directly for shared memory buses, direct memory access channels, and
cache capacity. Consider an analytical example of an autonomous vehicle
equipped with four high-resolution cameras operating at
\(30\text{ Hz}\). If each camera streams raw Bayer frames of
\(1920 \times 1080\) pixels at \(16\text{ bits}\) per pixel, the four
sensors continuously generate \(497.7\text{ MB/s}\) of raw ingress
traffic across the system interconnect. When image pipelines debayer
these frames, convert color spaces, and populate input tensors for
neural network backbones, the aggregate memory read and write volume
exceeds \(1.5\text{ GB/s}\). If this data stream shares a silicon
crossbar with the microcontroller bus without hardware-enforced
quality-of-service priorities, memory bus saturation can stall
safety-critical telemetry. If an inertial measurement unit interrupt
running at \(1000\text{ Hz}\) experiences a \(P_{99}\) scheduling delay
of \(5.0\text{ ms}\) because camera direct memory access transfers
congest the bus, a vehicle traveling at \(10.0\text{ m/s}\) travels
\(50\text{ mm}\) without active acceleration feedback. Ingestion latency
cannot be masked through asynchronous software queues, because software
buffering does not alter the physical time elapsed between photon
arrival and actuator response. Every millisecond consumed by
serialization, sensor driver dispatch, and neural inference ages the
observation record and widens the physical distance traveled by the
machine.

Downstream controllers, planners, and safety enforcers require
observation records expressed in metric coordinates with bounded
uncertainty to prevent physical control errors. A neural object detector
that outputs a scalar softmax classification probability of \(0.97\)
provides a metric of training distribution alignment, but it provides
zero information regarding the spatial variance of the object surface.
This creates a dangerous systems paradox: an AI model can be highly
confident in its semantic prediction while being entirely blind to
geometric drift. If an obstacle is reported at a nominal range of
\(5.0\text{ m}\) without a spatial covariance matrix
\(\mathbf{\Sigma}\), a trajectory optimization solver cannot evaluate
clearance bounds along a planned path. Consider an analytical example
where mechanical vibration or thermal expansion introduces an unmodeled
angular calibration drift of \(0.01\text{ rad}\) at a sensor bracket. At
a measurement distance of \(5.0\text{ m}\), this small angular deviation
produces a lateral translation error of \(50\text{ mm}\). If a mobile
machine plans a navigation trajectory through an opening with a nominal
clearance margin of \(80\text{ mm}\), the unmodeled calibration error
consumes over sixty percent of the available margin. Because the policy
relies on a high softmax confidence while ignoring geometric
uncertainty, the machine will drive smoothly into a catastrophic
structural collision, entirely unaware that its physical reality has
drifted from its mathematical belief. When perception algorithms fail to
output explicit covariance envelopes, downstream control systems cannot
distinguish between sensor noise and true physical displacement. This
missing bound leads either to phantom emergency braking when noise
spikes cross fixed thresholds, or to structural collisions when physical
drift remains undetected.

An observation record that knows its age, its spatial frame, and its
uncertainty fulfills the contract of the perception system, but an
individual observation remains an isolated, historical snapshot of a
dynamic world. Sensors sample the environment at discrete intervals,
introduce non-zero transport and inference delays, and experience
intermittent dropouts when physical objects pass behind obstacles or
enter sensor blind spots. A physical machine cannot achieve continuous,
stable operation by reacting directly to raw, asynchronous perception
packets as they arrive. If a high-speed vehicle updates its motor
torques solely in response to instantaneous, delayed vision outputs, the
feedback loop encounters phase lag that destabilizes tracking and
induces violent oscillations. To bridge the gap between discrete,
delayed sensory observations and continuous physical reality, the
architecture must maintain an internal representation of the physical
environment that persists across measurement intervals.

\protect\phantomsection\label{callout-takeawaysux2a-8.10}
\begin{fbxSimple}{callout-takeaways}{Key Chapter Takeaways}{}
\phantomsection\label{callout-takeaways*-8.10}

\begin{itemize}
\tightlist
\item
  \textbf{Observation Age Begins at Transduction Midpoint:} Sensory
  records must carry hardware-latched exposure timestamps
  (\(t_{\text{mid}}\)). Every millisecond of serialization, debayering,
  and neural inference inflates observation age (\(t_{\text{age}}\)),
  consuming unguided stopping clearance (\(\Delta d = \int v \, dt\)).
\item
  \textbf{Sensory Ingress Contends for Shared Silicon Bandwidth:}
  Uncompressed camera streams (\(>1.5\text{ GB/s}\)) saturate shared SoC
  interconnects. Ingestion pipelines must enforce hardware QoS channels
  and DMA prioritization to prevent bulk visual payloads from starving
  real-time IMU interrupts (\(1000\text{ Hz}\)).
\item
  \textbf{Softmax Probability is Not Geometric Uncertainty:} A
  classification score (\(0.97\)) gives zero spatial variance.
  Downstream kinodynamic planners require metric coordinate vectors with
  formal covariance ellipsoids
  (\(\mathbf{\Sigma} \in \mathbb{R}^{3 \times 3}\)) to evaluate spatial
  clearance margins against obstacle surfaces.
\item
  \textbf{Sensor Calibration is a Dynamic Physical Variable:} Mechanical
  shock, vibration, and thermal expansion continuously drift extrinsic
  transformation matrices (\(T_{\text{sensor}}^{\text{body}}\)). An
  unmodeled \(0.01\text{ rad}\) angular drift at \(5.0\text{ m}\) range
  induces \(50\text{ mm}\) of lateral error, requiring continuous online
  observability checks.
\item
  \textbf{Perception Provides Dated Snapshots, Not Continuous Reality:}
  Discrete sensory packets arrive with non-uniform latencies
  (\(20\text{--}100\text{ ms}\)). Raw perception outputs cannot directly
  drive motor inverters; they must feed persistent spatial memory
  representations that propagate physical state forward during sensor
  occlusion.
\end{itemize}

\end{fbxSimple}

This requirement defines the handoff to Chapter~\ref{sec-memory}. Where
perception produces isolated, dated observation records bounded in space
and uncertainty, the memory and state estimation subsystem fuses these
asynchronous records against forward kinematic and dynamic models. By
maintaining an explicit state distribution, the nervous system
propagates physical belief forward during sensory dead times, reconciles
delayed observation records through historical state buffers, and
preserves object permanence when obstacles are temporarily occluded.
Perception answers what physical quantities were measured, when they
occurred, and how imprecise they might be. The memory and state
architecture consumes those bounded records to determine what the world
is doing right now, providing the continuous, low-latency state vector
required to govern physical action safely.

\chapter{\texorpdfstring{Spatial Memory \&
Belief\chaptersubtitle{How Fast the Picture Goes Stale}}{Spatial Memory \& Belief}}\label{sec-memory}

\begin{figure}[H]

\centering{

\includegraphics[width=1\linewidth,height=\textheight,keepaspectratio]{index_files/mediabag//Users/VJ/GitHub/MLSysBook-integrate-physical/physical/book/chapters/09-memory/figures/fig_locator.pdf}

}

\caption{\label{fig-locator}\textbf{Systems Organ Focus: Brain.} Dynamic
spatial memory and temporal belief decay. As time elapses after the last
physical sensor measurement (\(t_0\)), unobserved kinematic drift and
process noise expand the spatial error covariance envelope \(E(t)\). The
memory system must enforce deterministic expiration
(\(t_{\mathrm{exp}}\)) when uncertainty crosses task clearance margins,
preventing stale or frozen transform representations from misleading
downstream planners.}

\end{figure}%

\emph{How fast does a spatial belief turn into dangerous fiction once a
sensor loses sight of an object?}

Physical interaction is fundamentally non-Markovian: an obstacle that
passes behind an occluding pillar, a tool resting in a closed bin, or a
surface whose compliance was probed seconds ago continues to exert
physical constraints on the workspace. An embodied agent must maintain a
coherent temporal belief state that tracks unseen entities without
confusing past observations with present physical reality. As occlusion
duration grows, spatial uncertainty expands, and unobserved objects
move. Dependable physical AI architectures separate geometric covariance
from existence probability, anchor perceptual tokens to structured
transform trees, and enforce explicit temporal decay horizons
(\(t_{\text{exp}}\)) so that phantom stale representations never induce
catastrophic planning collisions.

\newpage{}

\protect\phantomsection\label{callout-objectiveux2a-9.1}
\begin{fbxSimple}{callout-objective}{Learning Objectives}{}
\phantomsection\label{callout-objective*-9.1}

After studying this chapter, you will be able to:

\begin{enumerate}
\def\labelenumi{\arabic{enumi}.}
\tightlist
\item
  State what a belief must carry in addition to a value, and why each
  part is required.
\item
  Predict how long a belief remains usable after its last measurement,
  from the dynamics rather than from habit.
\item
  Set an expiry from a measured decay rate and defend the number against
  mechanical clearance budgets.
\item
  Identify which belief failures are visible downstream and which are
  silent.
\item
  Contrast volumetric representations (TSDF voxel grids, OctoMaps, 3D
  Gaussian Splats) across memory footprints and query latencies.
\item
  Design lock-free coordinate transform and raycasting pipelines that
  assert negative evidence to prevent the frozen transform trap.
\end{enumerate}

\emph{Core chapter takeaway.} Belief has a rate of decay, and it must
expire on that rate rather than on convenience. A frozen transform is
more dangerous than a missing one, because nothing downstream can tell.

\end{fbxSimple}

\section{Preserving Spatial Belief}\label{sec-memory-9-1}

\lettrine{I}{n the four-tier}\lettrine{a}{rchitectural} hierarchy of
physical AI (Chapter~\ref{sec-nervous}), the \textbf{Brain} organ sits
between incoming sensory observations and low-level nervous system
enforcement. Within this cognitive layer, \textbf{Memory} serves as the
temporal bridge that carries geometric, spatial, and physical beliefs
forward in time (Figure~\ref{fig-locator}). A physical observation
records what a sensor detected at a single discrete instant, fixing a
measurement to a past timestamp and storing raw values that do not alter
once written into memory. A \textbf{belief} is the estimate the machine
carries forward after that instant, representing what the policy assumes
to be true about the environment at the present moment. Between
consecutive sensor frames, no physical system possesses direct knowledge
of its surroundings. Instead, it operates entirely on belief, projecting
past data into the immediate present. Because the physical world does
not stop moving when an exposure ends or a packet drops, every belief
carries an intrinsic rate of decay, drifting away from physical reality
until it becomes fiction. As illustrated in the systems locator
(Figure~\ref{fig-locator}), the primary systems challenge of memory is
to bound this uncertainty growth against task-defined usability
thresholds before stale data can misguide downstream planning.

In a canonical heterogeneous dual-brain System-on-Chip (SoC)
architecture, this principle governs the division of responsibility
between compute tiers. The unprivileged application processor (Linux
MPU) maintains high-dimensional spatial memory structures, processes
multi-modal perception, and generates candidate trajectory proposals.
The real-time safety processor (deterministic bare-metal microcontroller
unit, MCU) acts as the permission referee, enforcing hard temporal
deadlines and physical validity limits. Because the Linux MPU
experiences scheduling jitter, memory bus contention, and inference
latency, it cannot be trusted to self-report temporal freshness. The
real-time MCU evaluates every incoming proposal against independent
hardware timers, revoking tracking authority the moment an unrefreshed
belief breaches task safety margins.

Consider an analytical example of a manipulator tasked with touching a
workpiece and regulating contact force against a moving assembly
fixture. The fixture drifts relative to the robot base at a bounded
velocity of at most \(20\text{ mm/s}\). The mechanical compliance of the
tooling permits an allowable positioning error of \(2\text{ mm}\) before
the tool either breaks contact entirely or generates excessive normal
forces that overload the wrist sensor. Dividing the allowable error by
the maximum drift velocity yields an operational lifetime of
\(100\text{ ms}\), calculated as
\(2\text{ mm} / (20\text{ mm/s}) = 0.100\text{ s}\). For the first
\(50\text{ ms}\) after the camera captures an image, the physical
position of the workpiece remains safely within the tolerance window. At
\(100\text{ ms}\), the unobserved displacement reaches the outer
boundary of the error budget. At \(101\text{ ms}\), without a new
measurement to constrain the state, the stored coordinate is no longer a
reliable guide for action.

In software systems, data age is frequently measured from the instant a
message was serialized, published over an inter-process bus, or updated
in a shared memory table. In a physical AI system, this convention is
dangerous. Propagating a state estimate through an extended Kalman
filter at \(1000\text{ Hz}\) or republishing a kinematic transform at
\(200\text{ Hz}\) creates a stream of fresh timestamps, but it does not
refresh the underlying information. The true age of a belief is
determined strictly by the timestamp of the last physical measurement
that constrained it. A downstream tracking loop that reads a position
estimate generated by dead reckoning or policy rollouts may observe a
publish timestamp only \(2\text{ ms}\) old, yet the physical error
envelope continues to expand at the velocity of the unobserved drift.
Timestamping an inference output or a cache lookup records when compute
occurred, but it tells the receiver nothing about when reality was last
checked.

In traditional computing architectures, an immutable variable in memory
retains its validity indefinitely until an explicit write modifies its
bits. A cached key-value pair remains accurate as long as the underlying
database record has not signaled an invalidation. In physical systems,
the failure mode is inverted. An unchanged value in memory becomes
progressively less true because physical state evolves under external
forces while the stored bits remain static. The bits do not spoil
through corruption; they spoil through divergence. When a camera is
occluded or an optical link drops frames, a register holding the last
observed object pose retains its floating-point value with perfect
numerical precision. If the machine continues executing trajectory
commands against that frozen value, the discrepancy between the stored
representation and physical geometry compounds with every millisecond of
unobserved motion, eventually driving the end-effector into collision.

To prevent a machine from acting on stale data, memory systems must
treat expiry as an explicit interface event. An expired belief is not a
missing variable or a null pointer in an operating system; it is a
formal boundary condition where the uncertainty envelope around a state
estimate exceeds the safety tolerance of the task. The \textbf{belief
state invalidation horizon}
(\(t_{\mathrm{exp}} = \inf \{ t \ge 0 : E(t) \ge E_{\max} \}\)) marks
the exact temporal deadline where unobserved kinematic drift and process
noise drive accumulated error past the mechanical clearance or safety
margin (\(E_{\max}\)), revoking tracking authority and forcing an
explicit transition to a safe fallback mode. Managing this boundary
requires treating every stored quantity not merely as a coordinate or a
vector, but as an active contract defined by when it was measured, the
reference frame in which it was established, and the bounded rate at
which its validity drains away.

\protect\phantomsection\label{callout-ruleux2a-9.2}
\begin{fbxSimple}{callout-rule}{Architect’s Rule of Thumb:}{The Wow Core Invariant}
\phantomsection\label{callout-rule*-9.2}

\emph{Extrapolation advances timestamps without refreshing evidence age.
Unmonitored belief drifts into physical failure. Never let an
intermediate predictor, filter, or transport bridge reset the evidence
acquisition epoch \(t_0\); compute freshness is a software artifact, but
physical freshness belongs strictly to the hardware transducer.}

\end{fbxSimple}

\section{What Memory Hands Over}\label{sec-memory-9-2}

A belief passed across a bus or retrieved from a buffer cannot simply be
a raw array of floating-point numbers. An unannotated vector such as
\([0.142, -0.051, 0.812]\) conveys nothing about what physical entity is
being described, what physical dimension is represented, or how that
quantity is scaled. A downstream controller receiving raw coordinates
without metadata must rely on implicit conventions established at
compile time. If an estimator produces positions in meters while an
actuator controller expects millimeters, or if an angular rate is
mistaken for an absolute orientation, the resulting command drives motor
current to saturation limits within milliseconds. To prevent
misinterpretation across software boundaries, a belief record must first
specify the entity identity, the physical quantity, and the measurement
units. The consumer must be able to verify that the value represents the
center of mass of a specific workpiece in meters, or the normal contact
force on a fingertip in newtons, before routing that number into a
control law.

Specifying the quantity and units remains incomplete without binding the
value to a concrete reference frame. Identical spatial coordinates
describe entirely different physical realities depending on the origin
and orientation of the coordinate system. A Cartesian position
\([0.142, -0.051, 0.812]\text{ m}\) designates a point suspended in free
space when evaluated in the robot base frame, but represents a point
several hundred millimeters inside a rigid fixture when interpreted in
the tool flange frame. Because kinematic links flex, joints move, and
coordinate transforms update at varying rates, a belief record must
carry an explicit frame identifier. Carrying the frame allows the
receiving node to look up the corresponding spatial transformation tree
and project the stored geometry into the local frame required by the
actuator, preventing the controller from acting on geometric
misalignments caused by moving kinematic chains.

Geometry and identity describe the physical state at a single moment,
but memory must also capture the temporal anchor of that state. The
timestamp attached to a belief must record the exact epoch \(t_0\) of
the physical interaction or optical exposure that last constrained the
estimate. In distributed architectures, intermediate nodes often
republish data, execute Kalman filter prediction steps, or run forward
kinematic propagations to maintain high-rate control streams. A
prediction step or cache refresh updates the transmission timestamp, but
it does not add physical evidence. If an intermediate process overwrites
\(t_0\) with the current system clock during a forward rollout,
downstream safety monitors observe a deceptive stream of fresh
timestamps while the underlying state continues to drift unchecked.
Preserving the original measurement epoch \(t_0\) across all
serialization and filtering stages ensures that consumers measure the
age of the information against the physical world rather than against
the scheduling frequency of the local computer.

Knowing the measurement epoch \(t_0\) establishes elapsed time
\(\Delta t = t - t_0\), but elapsed time alone cannot determine whether
a belief is valid for a given action. Two estimates of identical age
diverge at different rates depending on physical constraints. A
workpiece firmly held in a pneumatic vise maintains sub-millimeter
positional certainty after several seconds of sensor occlusion, whereas
the same workpiece sliding along a moving conveyor belt accumulates
substantial position error in fifty milliseconds. Therefore, a belief
record must include both the initial uncertainty bound \(e_0\)
established at the epoch of measurement and an explicit growth law that
quantifies how that uncertainty expands over time.

The minimum record fields necessary for a consumer to evaluate a stored
state can be derived directly from the requirement to bound current
error. Let \(e(t)\) represent the upper bound of the discrepancy between
the stored value and true physical state at time \(t\). If an estimator
establishes a state at time \(t_0\) with an initial measurement
uncertainty \(e_0\), and external disturbances or unmodeled velocities
cause the true state to diverge at a maximum rate \(r\), the upper bound
on error at current time \(t\) obeys the linear expansion
\(e(t) = e_0 + r(t - t_0)\). As an analytical example, consider a
tactile sensor array that establishes the initial slip displacement of a
grasped object with an uncertainty \(e_0 = 0.40\text{ mm}\) at
\(t_0 = 1.200\text{ s}\). If the object slides beneath the gripper at a
maximum unmeasured velocity \(r = 15.0\text{ mm/s}\), the error bound
after an unobserved interval of \(\Delta t = 80\text{ ms}\) expands to
\(e(t) = 0.40\text{ mm} + (15.0\text{ mm/s})(0.080\text{ s}) = 1.60\text{ mm}\).
Omitting either \(e_0\) or \(r\) leaves the consumer unable to compute
\(e(t)\), demonstrating that an initial bound and an expansion rate are
both structurally required.

Expressing uncertainty through physical bounds and expansion rates
distinguishes a physical AI belief from a probabilistic confidence
score. An object detection network might output a class confidence of
\(0.94\), but a downstream position controller cannot convert a unitless
probability into an actuator displacement margin. A confidence score of
\(0.94\) provides no information about whether an obstacle boundary is
uncertain by \(2.0\text{ mm}\) or by \(200\text{ mm}\), nor does it
indicate how rapidly that boundary moves when occlusion occurs. A
physical memory record must express its initial bound \(e_0\) and
expansion rate \(r\) in concrete physical dimensions such as meters,
radians, or newtons per second. Furthermore, these rates are valid only
within a defined operating envelope, such as a maximum assumed
acceleration or a verified contact friction coefficient. If physical
operating conditions exceed that envelope, the expansion rate is no
longer guaranteed, and the belief record must be treated as invalid.

When memory hands over a tuple containing value, reference frame,
measurement epoch \(t_0\), initial bound \(e_0\), and growth rate \(r\),
the consumer executes a deterministic verification step before applying
the data to physical control. The consumer possesses an application
tolerance \(\epsilon_{\text{tol}}\) dictated by the physical task. In a
precision insertion operation with a radial clearance of
\(0.80\text{ mm}\), the controller requires that positional uncertainty
remain below \(\epsilon_{\text{tol}} = 0.80\text{ mm}\). Upon reading
the belief record at current time \(t_{\text{now}}\), the consumer
calculates the elapsed time \(\Delta t = t_{\text{now}} - t_0\),
evaluates the present error bound
\(e(t_{\text{now}}) = e_0 + r \Delta t\), and compares the result
directly against \(\epsilon_{\text{tol}}\). If
\(e(t_{\text{now}}) \le \epsilon_{\text{tol}}\), the controller executes
the motion. If \(e(t_{\text{now}}) > \epsilon_{\text{tol}}\), the belief
has expired, and the controller transitions immediately to a safe
holding state or an active search routine without waiting for an
external invalidation signal.

\protect\phantomsection\label{callout-definitionux2a-9.3}
\begin{fbxSimple}{callout-definition}{Definition:}{The Spatial Belief Contract}
\phantomsection\label{callout-definition*-9.3}

A \textbf{spatial belief contract} is the explicit typed tuple
\((\mathbf{x}, \mathcal{F}, t_0, e_0, r, \mathcal{E})\) governing state
handoffs between memory and control, binding the estimated value
(\(\mathbf{x}\)) to an unambiguous coordinate frame (\(\mathcal{F}\)), a
physical acquisition epoch (\(t_0\)), an initial error bound (\(e_0\)),
a validated growth rate (\(r\)), and an operating envelope
(\(\mathcal{E}\)).

\end{fbxSimple}

\section{A State Has an Epoch}\label{sec-memory-9-3}

Every state reported by a sensor or estimated by a model corresponds to
a specific point in time, termed its \textbf{epoch}. In a physical
computing pipeline, that state passes through three distinct temporal
boundaries before a controller can act on it. The \textbf{acquisition
epoch} (\(t_0\)) is the nanosecond-precise timestamp recording the exact
physical instant when energy transduction occurred at the hardware
sensor boundary, such as photons striking a photodiode array or a
digital register latching an encoder count.\sidenote{\footnotesize Microcontroller
  timer input-capture units latch signal edges into shadow registers
  upon external interrupt. Reading capture registers without atomic
  synchronization against overflow interrupt flags induces 16-bit or
  32-bit counter aliasing, corrupting acquisition epochs by up to the
  full counter period (e.g., \(65.5\text{ ms}\)).} The
\textbf{estimation epoch} (\(t_{\text{est}}\)) marks the instant the
computational pipeline (whether a state estimator or neural network)
completed its mathematical inference over those measurements. The
\textbf{publication epoch} (\(t_{\text{pub}}\)) marks the instant the
serialized message was committed to a communication bus or network queue
for dispatch to downstream consumers.

Collapsing these three timestamps into a single receipt or transmission
timestamp destroys the physical validity of the state. If an inference
accelerator requires \(45\text{ ms}\) to execute an object detection
model and the operating system incurs another \(5\text{ ms}\) of
scheduling and serialization latency before publication, tagging the
resulting object pose with the publication time hides a \(50\text{ ms}\)
physical displacement. During that \(50\text{ ms}\) interval, gravity,
momentum, and contact forces continue to act on the physical system
without pause. If a conveyor belt carries a workpiece at
\(0.60\text{ m/s}\), a downstream gripper that treats the publication
timestamp as the measurement epoch assumes the workpiece is located
where it was \(50\text{ ms}\) in the past, introducing an unmodeled
position error of
\((0.60\text{ m/s})(0.050\text{ s}) = 30.0\text{ mm}\). A state that
records only when data became available tells the consumer nothing about
the physical instant the measurement represents, making temporal
alignment impossible.

Even when a system records true acquisition epochs, physical AI
architectures rely on distributed processing nodes whose local
oscillator clocks drift independently. Synchronizing these distributed
clocks over standard bus networks introduces a bounded synchronization
error, and that temporal uncertainty translates directly into spatial
uncertainty.\sidenote{\footnotesize Fieldbus protocols exhibit distinct clock
  synchronization jitter envelopes: IEEE 1588 PTP over Ethernet achieves
  \(\delta t < 1.0\ \mu\text{s}\), EtherCAT Distributed Clocks (DC)
  achieve \(<100\text{ ns}\) hardware SYNC0 jitter, and CAN-FD bit-phase
  synchronization yields \(\pm 10\text{--}50\ \mu\text{s}\).} Consider
an analytical example where an autonomous manipulator tracks a moving
assembly target. The vision processor and the joint motion controller
run on separate compute boards synchronized across an Ethernet bus using
the Precision Time Protocol (PTP, IEEE 1588), maintaining a bounded
clock offset of \(\delta t = \pm 1.5\text{ ms}\). If the target moves
relative to the camera at an operational velocity of
\(v = 1.20\text{ m/s}\), the clock offset alone creates an irrecoverable
spatial positioning uncertainty of
\(e_{\text{sync}} = v \cdot \delta t = (1.20\text{ m/s})(0.0015\text{ s}) = 1.80\text{ mm}\).
In an assembly operation with a mechanical clearance tolerance of
\(0.50\text{ mm}\), an uncalibrated clock synchronization bound of
\(1.5\text{ ms}\) causes physical interference and tool collision before
any controller tracking latency is added. The physical error budget must
therefore allocate margin for clock synchronization jitter alongside
sensor resolution and structural deflection.

Real physical systems never sample all relevant state variables
simultaneously. A tactile array might latch normal force readings at
\(t_1 = 12.0\text{ ms}\), an optical tracking camera might finish sensor
integration at \(t_2 = 18.5\text{ ms}\), and joint optical encoders
might sample angular positions at \(t_3 = 21.0\text{ ms}\). Because
these observations occur at different epochs, they cannot be combined
through direct algebraic operations. Fusing a camera observation at
\(t_2\) with an encoder observation at \(t_3\) requires the consumer to
project one state across the temporal gap
\(\Delta t = |t_3 - t_2| = 2.5\text{ ms}\) using explicit kinematic
models and bounded velocities. Without the true acquisition epoch for
each individual signal, the alignment calculation cannot determine the
sign or magnitude of \(\Delta t\). Extrapolating state without an
accurate epoch causes the fusion algorithm to compute spatial offsets
from fictitious temporal baselines, compounding errors across every
fusion cycle.

For a consumer to use a state safely, the memory record must provide
four explicit quantities: the physical state value, the precise
acquisition epoch referenced to a synchronized system clock, the upper
bound on clock synchronization uncertainty, and the maximum rate of
physical state change. Given these quantities, the consumer computes the
exact elapsed duration between the acquisition epoch and the execution
deadline, evaluates the worst-case expansion of the error bound, and
determines whether the state remains within task tolerances. Handing
over a state without its acquisition epoch leaves the consumer blind to
transport delays and scheduling jitter, turning deterministic control
into open-loop guesswork. Because physical geometry also changes as
bodies move through space, spatial transforms between reference frames
are subject to the same temporal decay, carrying their own distinct
measurement epochs.

\section{Reference Frames and Volumetric
Representations}\label{sec-memory-9-4}

Every geometric coordinate is meaningless without the reference frame in
which it was observed, and every reference frame in a moving machine
moves. Spatial transformation matrices express translations and
rotations between coordinate frames, allowing an onboard computer to map
a target detected by a camera into joint torques applied at an
end-effector (Figure~\ref{fig-09-robot-transform-tree}). In rigorous
spatial mechanics, rigid-body poses belong to the Special Euclidean Lie
group \(SE(3) = SO(3) \times \mathbb{R}^3\). A homogeneous
transformation matrix mapping coordinates from frame \(B\) to frame
\(A\) at time \(t\) is expressed as
\[\mathbf{T}_A^B(t) = \begin{bmatrix} \mathbf{R}_A^B(t) & \mathbf{p}_A^B(t) \\ \mathbf{0}^T & 1 \end{bmatrix} \in SE(3)\]
where \(\mathbf{R}_A^B(t) \in SO(3)\) is a \(3 \times 3\) orthonormal
rotation matrix satisfying \(\mathbf{R}^T \mathbf{R} = \mathbf{I}\) and
\(\det(\mathbf{R}) = +1\), and \(\mathbf{p}_A^B(t) \in \mathbb{R}^3\) is
the translation vector designating the origin of frame \(B\) relative to
frame \(A\).

A robot maintains its spatial relationships as a directed acyclic
coordinate transform tree \(\mathcal{T} = (\mathcal{V}, \mathcal{E})\),
where each vertex \(v \in \mathcal{V}\) denotes a physical or virtual
coordinate frame (e.g., base link, camera sensor, wrist flange, tool
center point, workpiece target), and each directed edge
\((i, j) \in \mathcal{E}\) stores a time-varying transform
\(\mathbf{T}_i^j(t)\) parameterized by joint encoder angles and static
link offsets (\citeproc{ref-lynch2017modern}{Lynch and Park 2017};
\citeproc{ref-craig2005introduction}{Craig 2005}).\sidenote{\footnotesize Traversing
  pointer-based spatial trees and voxel grids induces CPU cache misses.
  On unified-memory SoCs, concurrent DMA transfers from sensor ingest
  pipelines saturate shared DRAM arbiters, expanding worst-case spatial
  query latency tails from sub-microsecond bounds to tens of
  milliseconds.} What classical mechanics treats as a timeless geometric
relationship, however, is in a physical computing system a time-stamped
message subject to finite transport delays, processing latency, and
clock drift. A coordinate transform is valid only at the exact epoch
when its underlying physical geometry was measured. When an onboard
computer composes a multi-link kinematic chain across a directed path
\(0 \to 1 \to \dots \to K\), the end-to-end transform is computed by
sequential matrix multiplication:
\[\mathbf{T}_0^K(t) = \prod_{k=1}^K \mathbf{T}_{k-1}^k(t_k)\] When the
measurement epochs \(t_k\) across individual edges diverge (due to
asynchronous joint encoder sampling, bus serialization queues, or
dropped camera frames), the resulting matrix product does not describe
the physical machine at any single instant \(t\). It evaluates a
synthetic, hybrid pose that never existed in physical
reality.\sidenote{\footnotesize In heterogeneous dual-brain platforms, 6-DoF
  transformation matrices exceed atomic word sizes (\(64\text{ B}\)).
  Lock-free synchronization requires sequence locks (\texttt{seqlock})
  with inner-shareable memory barriers (\texttt{DMB\ ISH} on ARM,
  \texttt{fence} on RISC-V) to prevent torn reads between MPU updates
  and MCU reads.}

\begin{figure}[t]

\centering{

\includegraphics[width=0.85\linewidth,height=\textheight,keepaspectratio]{/Users/VJ/GitHub/MLSysBook-integrate-physical/physical/book/chapters/09-memory/figures/fig09_real_robot_transform_tree.png}

}

\caption{\label{fig-09-robot-transform-tree}\textbf{Kinematic Coordinate
Frame Tree on an Articulated Manipulator.} Physical and software
transformation frame hierarchy for a 7-DOF Franka Emika Panda
manipulator inside RViz. The coordinate tree anchors the base link
(\(\texttt{panda\_link0}\)) to the world origin and recursively chains
joint frames (\(\texttt{panda\_link1}\) through
\(\texttt{panda\_link7}\)) to the mechanical flange
(\(\texttt{panda\_link8}\)), nominal end-effector
(\(\texttt{panda\_NE}\)), tool center point (\(\texttt{panda\_EE}\)),
finger graspers (\(\texttt{panda\_hand}\),
\(\texttt{panda\_leftfinger}\), \(\texttt{panda\_rightfinger}\)), and
stiffness frame (\(\texttt{panda\_K}\)). Each directed edge represents a
time-varying homogeneous transformation matrix \(\mathbf{T}_{A}^{B}(t)\)
parameterized by encoder angles and kinematic offsets. Composing
transforms of disparate measurement epochs injects lever-arm positioning
error
\(e_{\mathrm{tool}} \approx e_{\mathrm{trans}} + \ell \Delta \theta\),
while a frozen transform edge in the tree silently guides motion against
phantom coordinates. (Credit: Franka Robotics GmbH / ROS Community, CC
BY 4.0).}

\end{figure}%

The cost of an expired transform is direct spatial displacement
(Figure~\ref{fig-frame-staleness-error}). Applying the delay-as-distance
principle from Chapter~\ref{sec-body} to kinematics, uncompensated
latency translates directly into geometric positioning error through the
lever arm of the linkages. For an articulated kinematic chain with \(N\)
moving joints, an uncompensated latency \(\Delta t_i\) at joint \(i\)
moving with angular velocity \(\boldsymbol{\omega}_i\) induces a spatial
velocity
\(\mathbf{v}_{\text{tool}} = \boldsymbol{\omega}_i \times (\mathbf{p}_{\text{tool}} - \mathbf{p}_i)\)
at the tool tip. Over an unobserved latency window, the cumulative
tool-point positioning error expands across the articulated chain
according to
\[\Delta \mathbf{p}_{\text{tool}} \approx \Delta \mathbf{p}_{\text{base}} + \sum_{i=1}^N \left( \boldsymbol{\omega}_i \Delta t_i \right) \times (\mathbf{p}_{\text{tool}} - \mathbf{p}_i)\]
where \(\ell_i = \|\mathbf{p}_{\text{tool}} - \mathbf{p}_i\|\)
represents the effective lever arm from joint \(i\) to the tool center
point.

Consider an analytical example of a manipulator holding a tool point at
a distance of \(\ell = 0.60\text{ m}\) from its wrist assembly. Suppose
the base and wrist rotate at an angular velocity of
\(\omega = 1.50\text{ rad/s}\) while the arm translates through space at
\(v = 0.80\text{ m/s}\). If a tracking pipeline publishes a
camera-to-tool transform with an uncompensated staleness of
\(\Delta t = 40\text{ ms} = 0.040\text{ s}\), the translation of the
frame drifts by
\(e_{\text{trans}} = v \cdot \Delta t = (0.80\text{ m/s})(0.040\text{ s}) = 0.032\text{ m} = 3.20\text{ cm}\).
Over the same interval, the angular orientation rotates through
\(\Delta \theta = \omega \cdot \Delta t = (1.50\text{ rad/s})(0.040\text{ s}) = 0.060\text{ rad}\).
Through the lever-arm relation
\(e_{\text{rot}} \approx \ell \Delta \theta\), this angular shift swings
the tool point across an arc of
\(e_{\text{rot}} \approx (0.60\text{ m})(0.060\text{ rad}) = 0.036\text{ m} = 3.60\text{ cm}\).
Summing the translational shift and the lever-arm deflection yields a
total tool-point positioning error of
\(e_{\text{tool}} \approx e_{\text{trans}} + \ell \Delta \theta = 3.20\text{ cm} + 3.60\text{ cm} = 6.80\text{ cm}\).
In a precision contact task where an end-effector regulates force
against a workpiece, a spatial discrepancy of \(6.80\text{ cm}\)
commands an actuator to drive through solid metal before tactile sensors
can register unexpected resistance.

\protect\phantomsection\label{callout-ruleux2a-9.4}
\begin{fbxSimple}{callout-rule}{Architect’s Rule of Thumb:}{The Leverage of Staleness}
\phantomsection\label{callout-rule*-9.4}

\emph{Transform staleness multiplies through kinematic link lengths; an
uncompensated millisecond of lag at the robot base swings the tool
center point through centimeters of unobserved deflection.}

\end{fbxSimple}

\begin{figure}[t]

\centering{

\includegraphics[width=1\linewidth,height=\textheight,keepaspectratio]{index_files/mediabag//Users/VJ/GitHub/MLSysBook-integrate-physical/physical/book/chapters/09-memory/figures/fig09_frame_staleness_error.pdf}

}

\caption{\label{fig-frame-staleness-error}\textbf{Coordinate Transform
Staleness and Lever-Arm Kinematic Error Propagation.} Rigid-body spatial
divergence induced by an uncompensated \(40\text{ ms}\) transform
latency on a manipulator arm translating at \(v = 0.80\text{ m/s}\) and
rotating at \(\omega = 1.50\text{ rad/s}\). Base translation shifts the
tool frame by \(e_{\text{trans}} = v \Delta t = 3.20\text{ cm}\), while
arm link rotation (\(\ell = 0.60\text{ m}\)) swings the tool point
through an arc
\(e_{\text{rot}} \approx \ell \omega \Delta t = 3.60\text{ cm}\),
yielding a cumulative tool-point positioning error of
\(e_{\text{tool}} \approx 6.80\text{ cm}\) (\(68.0\text{ mm}\)) that
breaches the \(\pm 0.50\text{ mm}\) task clearance by \(136\times\).
Whereas a missing transform triggers a loud exception and safe holding
reflex, a frozen transform republished with fresh system timestamps
evaluates with zero linear algebra exceptions, silently driving the tool
into collision.}

\end{figure}%

Beyond discrete kinematic frames, an embodied agent must also maintain
continuous volumetric spatial memory to track unobserved obstacles, plan
collision-free paths, and support real-time reactive grasping
(Figure~\ref{fig-09-octomap-volumetric-mapping}). By executing
\textbf{volumetric free-space raycasting} along sensor sightlines,
spatial memory actively updates 3D voxel hierarchies (such as OctoMap or
OpenVDB) and deterministically clears unoccupied voxels. The
foundational algorithm for volumetric raycasting is the 3D Differential
Digital Analyzer (3D-DDA) fast voxel traversal
(\citeproc{ref-amanatides1987fast}{Amanatides and Woo 1987}).\sidenote{\footnotesize On
  SIMT GPU architectures (e.g., NVIDIA Ampere/Ada Lovelace SMs),
  parallel 3D-DDA raycasting encounters thread divergence when rays
  within a 32-thread warp traverse octree voxels of differing depths.
  Modern spatial engines implement branchless DDA stepping, pack active
  voxel blocks into GPU Shared Memory (L1 cache), and utilize atomic
  integer bitmasks (\texttt{atomicOr}) to serialize voxel hash updates
  without global DRAM lock contention.} For a sensor sightline ray
\(\mathbf{r}(s) = \mathbf{o} + s\mathbf{d}\) (\(s \ge 0\)) originating
at optical origin \(\mathbf{o} \in \mathbb{R}^3\) with unit direction
\(\mathbf{d} \in \mathbb{R}^3\) through a discrete voxel grid of
isotropic pitch \(\Delta v\), the traversal steps through voxel
boundaries \((u, v, w)\) by maintaining next-plane crossing parameters:
\[t_{\text{next}, x} = \frac{(u + \operatorname{step}_x)\Delta v - o_x}{d_x}, \quad \Delta t_x = \frac{\Delta v}{|d_x|}\]
At each iteration, the algorithm advances along the dimension with
minimal \(t_{\text{next}}\), achieving \(O(1)\) arithmetic operations
per traversed voxel cell.

Crucially, as the ray traverses free space from \(s = 0\) to
\(s = z_{\text{depth}} - \delta\), the pipeline updates occupancy
log-odds \(L(\mathbf{m}) = \ln \frac{P(\mathbf{m})}{1 - P(\mathbf{m})}\)
by asserting \textbf{negative evidence}:
\[L(\mathbf{m}_i) \leftarrow L(\mathbf{m}_i) + l_{\text{free}}, \quad l_{\text{free}} = \ln \left( \frac{P(\text{occ} \mid \text{free})}{1 - P(\text{occ} \mid \text{free})} \right) < 0\]
At the terminal measurement interval
\([z_{\text{depth}} - \delta, z_{\text{depth}} + \delta]\), positive
evidence is integrated:
\[L(\mathbf{m}_{\text{end}}) \leftarrow L(\mathbf{m}_{\text{end}}) + l_{\text{occ}}, \quad l_{\text{occ}} = \ln \left( \frac{P(\text{occ} \mid \text{hit})}{1 - P(\text{occ} \mid \text{hit})} \right) > 0\]
This active assertion of negative evidence along sensor sightlines is
the exact computational mechanism that invalidates phantom obstacle
tracks when dynamic objects leave a region, pruning ghost tracks without
waiting for passive timeout decay.

For closed-loop trajectory optimization and reactive grasping, discrete
occupancy grids are insufficient because they lack continuous spatial
derivatives. Physical AI systems employ \textbf{Truncated Signed
Distance Fields (TSDFs)} (\citeproc{ref-curless1996volumetric}{Curless
and Levoy 1996}; \citeproc{ref-newcombe2011kinectfusion}{Newcombe et al.
2011}; \citeproc{ref-oleynikova2017voxblox}{Oleynikova et al.
2017}).\sidenote{\footnotesize In continuous TSDF fusion, the projective distance
  update assumes a pinhole camera geometry with calibrated focal lengths
  \((f_x, f_y)\) and optical center \((c_x, c_y)\). Spatial gradients
  \(\nabla d(\mathbf{x})\) extracted via central finite differences or
  analytical trilinear interpolation are guaranteed \(C^0\)-continuous
  across voxel boundaries, providing smooth Jacobians
  \(\mathbf{J}_{\text{dist}} = \nabla d(\mathbf{x})^T \mathbf{J}_{\text{kin}}(q)\)
  for gradient-based inverse kinematics and trajectory optimization.}
For any spatial query point \(\mathbf{x} \in \mathbb{R}^3\), let
\(\pi(\mathbf{x})\) denote the camera pinhole perspective projection and
\(\mathbf{p}_{\text{cam}}\) the optical center. The projective signed
distance is
\(D(\mathbf{x}) = (\mathbf{x} - \mathbf{p}_{\text{cam}})_z - z_{\text{depth}}(\pi(\mathbf{x}))\).
Truncating by distance metric margin \(\mu\) yields normalized signed
distance
\(d(\mathbf{x}) = \max\left(-1, \min\left(1, D(\mathbf{x})/\mu\right)\right)\).
Voxels integrate incoming frames through running weighted recursive
fusion:
\[D_k(\mathbf{x}) = \frac{W_{k-1}(\mathbf{x}) D_{k-1}(\mathbf{x}) + w_k d_k(\mathbf{x})}{W_{k-1}(\mathbf{x}) + w_k}, \quad W_k(\mathbf{x}) = \min(W_{k-1}(\mathbf{x}) + w_k, W_{\max})\]
Tri-linear interpolation across adjacent TSDF voxel vertices yields
exact continuous spatial distance metrics and analytical spatial
gradients
\(\nabla d(\mathbf{x}) = \frac{\mathbf{x} - \mathbf{x}_{\text{surf}}}{\|\mathbf{x} - \mathbf{x}_{\text{surf}}\|}\).
Gradient-based motion planners utilize this continuous derivative to
evaluate collision cost gradients
\(\nabla_{\mathbf{x}} \mathcal{C}_{\text{coll}}(\mathbf{x}) = -\nabla d(\mathbf{x})\),
pushing trajectory waypoints smoothly out of obstacle boundaries.

Moving beyond fixed voxel lattices, \textbf{3D Gaussian Splatting
(3DGS)} (\citeproc{ref-kerbl20233d}{Kerbl et al. 2023};
\citeproc{ref-matsuki2024gaussian}{Matsuki et al. 2024}) provides
continuous, spatially indexed, and differentiable spatial memory. A
scene is represented as a collection of \(M\) explicit anisotropic 3D
ellipsoidal Gaussians:
\[G_i(\mathbf{x}) = \exp\left( -\frac{1}{2} (\mathbf{x} - \boldsymbol{\mu}_i)^T \mathbf{\Sigma}_i^{-1} (\mathbf{x} - \boldsymbol{\mu}_i) \right)\]
where center \(\boldsymbol{\mu}_i \in \mathbb{R}^3\) and spatial
covariance
\(\mathbf{\Sigma}_i = \mathbf{R}_i \mathbf{S}_i \mathbf{S}_i^T \mathbf{R}_i^T\)
are parameterized by a unit quaternion \(\mathbf{q}_i \in SO(3)\) and
scale vector \(\mathbf{s}_i \in \mathbb{R}^3\), accompanied by opacity
\(\alpha_i \in [0, 1]\) and spherical harmonics color coefficients
\(\mathbf{c}_i\). During robot operation, online spatial memory updates
optimize Gaussian parameters in real time via tile-based differentiable
rasterization and gradient backpropagation under combined photometric
and geometric depth losses:
\[\mathcal{L} = (1 - \lambda) \|\mathbf{I} - \hat{\mathbf{I}}\|_1 + \lambda \mathcal{L}_{\text{D-SSIM}} + \lambda_d \|\mathbf{D} - \hat{\mathbf{D}}\|_1\]
Gaussians are adaptively densified in under-reconstructed areas and
pruned in empty regions, enabling real-time photorealistic novel-view
synthesis and continuous collision querying without the memory explosion
of uniform volumetric grids (Table~\ref{tbl-09-spatial-rep}).

\begin{figure}[t]

\centering{

\includegraphics[width=0.95\linewidth,height=\textheight,keepaspectratio]{/Users/VJ/GitHub/MLSysBook-integrate-physical/physical/book/chapters/09-memory/figures/fig09_real_octomap_volumetric_mapping.png}

}

\caption{\label{fig-09-octomap-volumetric-mapping}\textbf{Probabilistic
3D Volumetric Mapping and Hierarchical Octree Space Subdivision.}
Empirical 3D spatial memory representation reconstructed from a physical
mobile robot traversing an indoor corridor with a tilting LiDAR scanner
(\citeproc{ref-hornung2013octomap}{Hornung et al. 2013}). (a) Raw 3D
point cloud acquired during robot motion, subject to sensor noise and
finite field-of-view occlusions. (b) Volumetric OctoMap voxel
representation modeling 3D space as an 8-ary hierarchical tree, where
occupied volumes are probabilistically updated and stored across
variable resolutions. (c) Explicit free-space raycasting and lossless
octree pruning, differentiating verified traversable space (white
voxels) from occupied obstacles (dark voxels) and unobserved unknown
regions. By maintaining explicit free-space evidence along sensor
sightlines, spatial memory actively invalidates phantom obstacle tracks
without waiting for arbitrary software timeouts. (Credit: Hornung et
al.~/ Autonomous Intelligent Systems Lab, University of Freiburg, CC BY
3.0).}

\end{figure}%

\begin{longtable}[]{@{}
  >{\raggedright\arraybackslash}p{(\linewidth - 14\tabcolsep) * \real{0.1250}}
  >{\raggedright\arraybackslash}p{(\linewidth - 14\tabcolsep) * \real{0.1250}}
  >{\raggedright\arraybackslash}p{(\linewidth - 14\tabcolsep) * \real{0.1250}}
  >{\raggedright\arraybackslash}p{(\linewidth - 14\tabcolsep) * \real{0.1250}}
  >{\raggedright\arraybackslash}p{(\linewidth - 14\tabcolsep) * \real{0.1250}}
  >{\raggedright\arraybackslash}p{(\linewidth - 14\tabcolsep) * \real{0.1250}}
  >{\raggedright\arraybackslash}p{(\linewidth - 14\tabcolsep) * \real{0.1250}}
  >{\raggedright\arraybackslash}p{(\linewidth - 14\tabcolsep) * \real{0.1250}}@{}}
\caption{\textbf{Volumetric Spatial Memory Representation Trade-Offs in
Physical AI Architectures.} Quantitative comparison of continuous and
discrete 3D spatial memory structures across storage footprint
(evaluated for a nominal
\(10\text{ m} \times 10\text{ m} \times 3\text{ m}\) active manipulation
workcell at \(1\text{ cm}\) isotropic voxel resolution), insertion
latency, raycast collision query latency, continuous spatial gradient
availability \(\nabla d(\mathbf{x})\), differentiable rendering support,
and primary silicon execution
targets.}\label{tbl-09-spatial-rep}\tabularnewline
\toprule\noalign{}
\begin{minipage}[b]{\linewidth}\raggedright
Spatial Representation Paradigm
\end{minipage} & \begin{minipage}[b]{\linewidth}\raggedright
Underlying Geometric Primitive
\end{minipage} & \begin{minipage}[b]{\linewidth}\raggedright
Memory Footprint (\(300\,\text{m}^3 @ 1\,\text{cm}\))
\end{minipage} & \begin{minipage}[b]{\linewidth}\raggedright
Dynamic Update Latency (\(10^5\,\text{pts}\))
\end{minipage} & \begin{minipage}[b]{\linewidth}\raggedright
Collision / Raycast Query Latency
\end{minipage} & \begin{minipage}[b]{\linewidth}\raggedright
Continuous Gradient \(\nabla d(\mathbf{x})\)
\end{minipage} & \begin{minipage}[b]{\linewidth}\raggedright
Differentiable Fusion Support
\end{minipage} & \begin{minipage}[b]{\linewidth}\raggedright
Primary Execution Target \& Bottleneck
\end{minipage} \\
\midrule\noalign{}
\endfirsthead
\toprule\noalign{}
\begin{minipage}[b]{\linewidth}\raggedright
Spatial Representation Paradigm
\end{minipage} & \begin{minipage}[b]{\linewidth}\raggedright
Underlying Geometric Primitive
\end{minipage} & \begin{minipage}[b]{\linewidth}\raggedright
Memory Footprint (\(300\,\text{m}^3 @ 1\,\text{cm}\))
\end{minipage} & \begin{minipage}[b]{\linewidth}\raggedright
Dynamic Update Latency (\(10^5\,\text{pts}\))
\end{minipage} & \begin{minipage}[b]{\linewidth}\raggedright
Collision / Raycast Query Latency
\end{minipage} & \begin{minipage}[b]{\linewidth}\raggedright
Continuous Gradient \(\nabla d(\mathbf{x})\)
\end{minipage} & \begin{minipage}[b]{\linewidth}\raggedright
Differentiable Fusion Support
\end{minipage} & \begin{minipage}[b]{\linewidth}\raggedright
Primary Execution Target \& Bottleneck
\end{minipage} \\
\midrule\noalign{}
\endhead
\bottomrule\noalign{}
\endlastfoot
\textbf{Octree Occupancy Grid} (e.g., OctoMap) & Hierarchical 8-ary tree
with log-odds occupancy voxels & \(45 - 120\text{ MB}\) (pointer
overhead) & \(15 - 45\text{ ms}\) (CPU pointer recursion) &
\(2.5 - 8.0\,\mu\text{s}\) per ray & Discontinuous (finite difference
only) & None (discrete binary/probabilistic voting) & Host CPU (pointer
chasing \& L1/L2 cache misses) \\
\textbf{Hierarchical Sparse VDB} (e.g., OpenVDB / NanoVDB) & Dynamic
\(B\)-tree array (\(32^3 \to 16^3 \to 8^3\) leaf tiles) &
\(18 - 40\text{ MB}\) (linearized sparse leaf buffers) &
\(2.0 - 6.5\text{ ms}\) (SIMD/GPU bitmask traversal) &
\(0.15 - 0.45\,\mu\text{s}\) per ray (DDA hardware step) & \(C^0\)
tri-linear spatial gradient & Partial (voxel-grid differentiable
splatting) & Edge GPU / SIMD Host (DRAM bus bandwidth bound) \\
\textbf{Truncated SDF (TSDF)} (e.g., Voxblox / KinectFusion) &
Voxel-hashing spatial blocks with metric distance weights &
\(60 - 180\text{ MB}\) (dense allocated sub-blocks) &
\(1.2 - 4.0\text{ ms}\) (GPU projective ray integration) &
\(0.08 - 0.25\,\mu\text{s}\) (direct hash index lookup) & Exact metric
\(\nabla d(\mathbf{x}) = \frac{\mathbf{x} - \mathbf{x}_{\text{surf}}}{\|\mathbf{x} - \mathbf{x}_{\text{surf}}\|}\)
& Moderate (mesh zero-crossing extraction) & Embedded GPU / FPGA (Shared
SRAM/L2 capacity) \\
\textbf{Continuous Neural SDF} (e.g., Instant-NGP / NeuS) &
Multi-resolution spatial hash table + shallow MLP & \(8 - 25\text{ MB}\)
(quantized hash table weights) & \(30 - 120\text{ ms}\) (online gradient
backpropagation) & \(12 - 50\,\mu\text{s}\) (neural network batch
inference) & Analytic autograd
\(\nabla_{\mathbf{x}} f_\theta(\mathbf{x}) \in C^\infty\) & Full
(photometric / geometric end-to-end loss) & Dedicated NPU / GPU (Tensor
core compute bound) \\
\textbf{3D Gaussian Splatting (\citeproc{ref-kerbl20233d}{Kerbl et al.
2023}; \citeproc{ref-matsuki2024gaussian}{Matsuki et al. 2024}) (3DGS)}
(e.g., Spatially Indexed 3DGS) & Parametric anisotropic 3D ellipsoids
(\(\boldsymbol{\mu}, \mathbf{\Sigma}, \alpha, \mathbf{c}\)) &
\(35 - 90\text{ MB}\) (\(5 \times 10^5\) explicit Gaussians) &
\(8 - 22\text{ ms}\) (tile-based differentiable raster) &
\(0.05 - 0.18\,\mu\text{s}\) (radix-sorted tile rasterization) &
Explicit spatial derivative \(\nabla_{\mathbf{x}} G(\mathbf{x})\) & Full
(sub-millisecond photorealistic differentiable render) & Edge GPU /
High-End SoC (Memory bandwidth bound) \\
\end{longtable}

\protect\phantomsection\label{callout-mathux2a-9.5}
\begin{fbxSimple}{callout-math}{Fermi Math:}{Memory Footprint and DRAM Bandwidth of 3D TSDF Voxel Grids}
\phantomsection\label{callout-math*-9.5}

\emph{Practical Engineering Question:} What is the memory footprint and
DRAM bandwidth demand of 3D TSDF voxel grids across sparse vs naive
dense representations, and why does naive dense allocation starve
concurrent NPU policy inference on edge SoCs?

\emph{1. Naive Dense Volumetric Allocation:} Consider an active robotic
manipulation workcell measuring
\(10\text{ m} \times 10\text{ m} \times 3\text{ m} = 300\text{ m}^3\).

\begin{itemize}
\item
  At isotropic resolution \(\Delta v = 1.0\text{ cm} = 0.01\text{ m}\):
  \[\text{Voxel Count} = \left(\frac{10}{0.01}\right) \times \left(\frac{10}{0.01}\right) \times \left(\frac{3}{0.01}\right) = 1000 \times 1000 \times 300 = 3 \times 10^8\text{ voxels}\]
\item
  Memory footprint at \(4\text{ bytes}\) per voxel (\(16\text{-bit}\)
  TSDF distance \(+ 16\text{-bit}\) fusion weight):
  \[M_{\text{dense, 1cm}} = 3 \times 10^8 \times 4\text{ B} = 1.20 \times 10^9\text{ B} = 1.20\text{ GB}\]
\item
  If resolution is refined to
  \(\Delta v = 5.0\text{ mm} = 0.005\text{ m}\) for precision tool
  contact:
  \[\text{Voxel Count} = 2000 \times 2000 \times 600 = 2.4 \times 10^9\text{ voxels} \implies M_{\text{dense, 5mm}} = 9.60\text{ GB}\]
  This exceeds the entire physical DRAM capacity (\(8\text{ GB}\)) of
  the edge SoC!
\end{itemize}

\emph{2. Sparse Voxel Hashing (Occupancy Sparsity \(\approx 3\%\)):}
Physical obstacles and workcell surfaces occupy only a thin 2D manifold
within 3D volume. Grouping voxels into \(8^3 = 512\) voxel leaf blocks
and allocating only within the TSDF truncation band
(\(\mu = \pm 3\text{ cm}\)):

\begin{itemize}
\item
  Active occupied volume: \(\approx 3.0\%\) of
  \(300\text{ m}^3 = 9.0\text{ m}^3\).
\item
  Active voxel count at \(1\text{ cm}\):
  \(9.0 \times 10^6\text{ voxels} \implies \approx 17,578\text{ blocks}\).
\item
  Memory footprint:
  \[M_{\text{sparse, 1cm}} = 17,578\text{ blocks} \times (512\text{ voxels} \times 4\text{ B} + 32\text{ B header}) \approx 36.5\text{ MB}\]
\item
  At \(5.0\text{ mm}\) resolution:
  \(M_{\text{sparse, 5mm}} \approx 292\text{ MB}\) (easily fits in SoC
  LPDDR5).
\end{itemize}

\emph{3. DRAM Bus Bandwidth under \(60\text{ Hz}\) RGB-D Raycasting:} A
depth sensor running at \(640 \times 480\) resolution at
\(60\text{ fps}\) produces:
\[\text{Ray Rate} = 640 \times 480 \times 60 = 18,432,000\text{ rays/sec} \approx 18.43\text{ M rays/s}\]

\begin{itemize}
\item
  Traversing a dense grid requires reading and writing along the full
  \(3.0\text{ m}\) ray (\(300\) voxel lookups per ray):
  \[\text{Dense DRAM Bandwidth} = 18.43 \times 10^6 \times 300 \times 8\text{ B (R+W)} \approx 44.2\text{ GB/s}\]
  This consumes \(>85\%\) of the SoC's total shared LPDDR5 bandwidth
  (\(51.2\text{ GB/s}\)), causing severe memory bus throttling.
\item
  Sparse VDB with 3D-DDA only updates voxels within the narrow
  truncation margin (\(6\) voxels per ray) and caches active leaf blocks
  in GPU shared memory / L2 cache:
  \[\text{Sparse DRAM Bandwidth} = 18.43 \times 10^6 \times 6 \times 8\text{ B} \times (1 - 0.95\text{ cache hit rate}) \approx 44.2\text{ MB/s}\]
\end{itemize}

\emph{Systems Verdict:} Naive dense 3D grids are mathematically trivial
but silicon-infeasible. Dependable physical AI architectures must use
sparse hierarchical voxel hashing (e.g., OpenVDB / Voxblox) with GPU
thread-block localization to reduce memory footprint by \(33\times\) and
DRAM bus traffic by \(1000\times\), protecting the memory bus for
real-time neural policy execution.

\end{fbxSimple}

To integrate incoming sensory point clouds while actively evicting stale
obstacle geometry, the spatial memory engine executes the structured
raycast update detailed in Algorithm 9.1.

\protect\phantomsection\label{callout-algorithmux2a-9.6}
\begin{fbxSimple}{callout-algorithm}{Algorithm:}{Algorithm 9.1: Volumetric TSDF Raycast Update with Temporal Decay}
\phantomsection\label{callout-algorithm*-9.6}

\textbf{Inputs:} Point cloud
\(\mathcal{P} = \{\mathbf{p}_k\}_{k=1}^{N_{\text{pts}}}\) in sensor
frame \(\mathcal{F}_{\text{sensor}}\), sensor pose
\(\mathbf{T}_{\text{world}\leftarrow\text{sensor}}(t_0)\) at capture
epoch \(t_0\), truncation margin \(\mu\), max weight \(W_{\max}\), decay
half-life \(\tau_{\text{decay}}\).

\textbf{State:} Sparse TSDF voxel block hash table
\(\mathcal{V}_{\text{tsdf}} = \{\text{block\_idx} \to \text{Block}(8^3\text{ voxels})\}\).

\textbf{Outputs:} Updated continuous TSDF spatial belief and dynamic
negative evidence field.

\textbf{Execution Procedure (\texttt{UPDATE\_VOLUMETRIC\_TSDF}):}

\begin{enumerate}
\def\labelenumi{\arabic{enumi}.}
\tightlist
\item
  \(t_{\text{curr}} \leftarrow \text{GetPtpMonotonicTimeNs}()\);
  \(\Delta t_{\text{elapsed}} \leftarrow (t_{\text{curr}} - t_0) \times 10^{-9}\).
\item
  If \(\Delta t_{\text{elapsed}} > \tau_{\max,\text{latency}} \implies\)
  \textbf{return} \(\text{STATUS\_FRAME\_TOO\_STALE}\) (reject stale
  sensory frame).
\item
  \(\mathbf{p}_{\text{cam,world}} \leftarrow \mathbf{T}_{\text{world}\leftarrow\text{sensor}}.\text{translation}\)
  (optical center in world frame).
\item
  \textbf{For each point \(\mathbf{p}_k \in \mathcal{P}\) (parallel GPU
  threads):}

  \begin{itemize}
  \tightlist
  \item
    \(\mathbf{p}_{\text{world}} \leftarrow \mathbf{T}_{\text{world}\leftarrow\text{sensor}} \cdot \mathbf{p}_k\);
    \(\mathbf{r}_{\text{dir}} \leftarrow \frac{\mathbf{p}_{\text{world}} - \mathbf{p}_{\text{cam,world}}}{\|\mathbf{p}_{\text{world}} - \mathbf{p}_{\text{cam,world}}\|}\);
    \(r_{\text{len}} \leftarrow \|\mathbf{p}_{\text{world}} - \mathbf{p}_{\text{cam,world}}\|\).
  \item
    \emph{Phase 1 (Negative Free-Space Clearing via 3D-DDA):}

    \begin{itemize}
    \item
      For each voxel \(v\) along ray from \(0\) to
      \(r_{\text{len}} - \mu\):

      \begin{itemize}
      \tightlist
      \item
        \(\text{vox} \leftarrow \text{GetOrAllocateBlock}(\mathcal{V}_{\text{tsdf}}, v.\text{block\_index}).\text{voxels}[v.\text{local\_idx}]\)
      \item
        \(\text{vox}.\text{weight} \leftarrow \text{vox}.\text{weight} \cdot \exp\left(-\frac{t_0 - \text{vox}.t_{\text{last}}}{\tau_{\text{decay}}}\right)\)
      \item
        \(\text{vox}.t_{\text{last}} \leftarrow t_0\)
      \item
        If
        \(\text{vox}.d_{\text{tsdf}} < 0.0 \implies \text{vox}.d_{\text{tsdf}} \leftarrow \min(\text{vox}.d_{\text{tsdf}} + L_{\text{clear}}, 1.0)\)
      \end{itemize}
    \end{itemize}
  \item
    \emph{Phase 2 (Surface Boundary Truncated Signed Distance
    Integration):}

    \begin{itemize}
    \item
      For each voxel \(v\) along ray from \(r_{\text{len}} - \mu\) to
      \(r_{\text{len}} + \mu\):

      \begin{itemize}
      \item
        \(\text{dist} \leftarrow \|\text{vox}.\mathbf{x} - \mathbf{p}_{\text{cam,world}}\| - r_{\text{len}}\)
      \item
        If \(\text{dist} > -\mu\):

        \begin{itemize}
        \tightlist
        \item
          \(d_{\text{norm}} \leftarrow \text{clamp}(\text{dist} / \mu, -1.0, 1.0)\)
        \item
          \(w_{\text{obs}} \leftarrow \frac{1}{1 + \sigma_z^2 r_{\text{len}}^2}\)
          (distance-weighted variance)
        \item
          \(w_{\text{old}} \leftarrow \text{vox}.\text{weight} \cdot \exp\left(-\frac{t_0 - \text{vox}.t_{\text{last}}}{\tau_{\text{decay}}}\right)\)
        \item
          \(\text{vox}.d_{\text{tsdf}} \leftarrow \frac{w_{\text{old}} \cdot \text{vox}.d_{\text{tsdf}} + w_{\text{obs}} \cdot d_{\text{norm}}}{w_{\text{old}} + w_{\text{obs}}}\)
        \item
          \(\text{vox}.\text{weight} \leftarrow \min(w_{\text{old}} + w_{\text{obs}}, W_{\max})\);
          \(\text{vox}.t_{\text{last}} \leftarrow t_0\)
        \end{itemize}
      \end{itemize}
    \end{itemize}
  \end{itemize}
\item
  \textbf{return} \(\text{STATUS\_TSDF\_UPDATED}\)
\end{enumerate}

\end{fbxSimple}

The volumetric belief state produced by this pipeline adheres to the
abstract data schema specified in
Table~\ref{tbl-09-schema-volumetric-belief}, guaranteeing typed spatial
bounds and physical SI unit invariants for downstream motion planners.

\begin{longtable}[]{@{}
  >{\raggedright\arraybackslash}p{(\linewidth - 6\tabcolsep) * \real{0.2500}}
  >{\raggedright\arraybackslash}p{(\linewidth - 6\tabcolsep) * \real{0.2500}}
  >{\raggedright\arraybackslash}p{(\linewidth - 6\tabcolsep) * \real{0.2500}}
  >{\raggedright\arraybackslash}p{(\linewidth - 6\tabcolsep) * \real{0.2500}}@{}}
\caption{\textbf{Abstract Volumetric Spatial Belief Manifest Schema.}
Structural parameters, data types, physical SI units, and invariant
constraints governing 3D TSDF and occupancy state memory handoffs
between spatial estimation backbones and downstream motion planning
controllers.}\label{tbl-09-schema-volumetric-belief}\tabularnewline
\toprule\noalign{}
\begin{minipage}[b]{\linewidth}\raggedright
Field Name
\end{minipage} & \begin{minipage}[b]{\linewidth}\raggedright
Scalar Type
\end{minipage} & \begin{minipage}[b]{\linewidth}\raggedright
Physical SI Units
\end{minipage} & \begin{minipage}[b]{\linewidth}\raggedright
Structural Invariant \& Description
\end{minipage} \\
\midrule\noalign{}
\endfirsthead
\toprule\noalign{}
\begin{minipage}[b]{\linewidth}\raggedright
Field Name
\end{minipage} & \begin{minipage}[b]{\linewidth}\raggedright
Scalar Type
\end{minipage} & \begin{minipage}[b]{\linewidth}\raggedright
Physical SI Units
\end{minipage} & \begin{minipage}[b]{\linewidth}\raggedright
Structural Invariant \& Description
\end{minipage} \\
\midrule\noalign{}
\endhead
\bottomrule\noalign{}
\endlastfoot
\texttt{voxel\_grid\_origin} & \texttt{float32{[}3{]}} & meters
(\(\text{m}\)) & Bounded world coordinates of spatial memory anchor
(\(x, y, z \in [-50, 50]\)). \\
\texttt{voxel\_resolution\_m} & \texttt{float32} & meters (\(\text{m}\))
& Isotropic voxel pitch (\(\Delta v \in [0.005, 0.050]\)). \\
\texttt{evidence\_epoch\_t0} & \texttt{int64\_t} & nanoseconds
(\(\text{ns}\)) & Immutable hardware capture timestamp of newest fused
sensory frame. \\
\texttt{tsdf\_trunc\_margin\_m} & \texttt{float32} & meters
(\(\text{m}\)) & Metric distance margin for TSDF clamping
(\(\mu \in [0.02, 0.10]\)). \\
\texttt{active\_block\_count} & \texttt{uint32\_t} & dimensionless &
Number of dynamically allocated \(8^3\) sparse leaf blocks in GPU hash
map. \\
\texttt{surface\_gradient\_norm} & \texttt{float32{[}3{]}} &
dimensionless & Analytical spatial gradient vector
(\(\nabla d(\mathbf{x}) = \frac{\mathbf{x} - \mathbf{x}_{\text{surf}}}{\|\mathbf{x} - \mathbf{x}_{\text{surf}}\|}\),
\(\|\nabla d\| = 1.0\)). \\
\texttt{temporal\_decay\_tau\_s} & \texttt{float32} & seconds
(\(\text{s}\)) & Exponential decay half-life for unobserved volumetric
evidence (\(\tau_{\text{decay}} \le 2.0\text{ s}\)). \\
\texttt{free\_space\_sweep\_epoch} & \texttt{int64\_t} & nanoseconds
(\(\text{ns}\)) & Timestamp of most recent raycast sweep asserting
negative free-space evidence. \\
\end{longtable}

A corrupted system architecture fails more destructively when a
transform freezes than when it disappears. Known as the \textbf{frozen
transform trap}, this failure mode occurs when a sensor driver or
intermediate transport bridge caches its last computed pose and
repeatedly republishes that stale matrix under refreshed message
timestamps, allowing every mathematical operation downstream to succeed
without error. Matrix multiplications evaluate cleanly, linear algebra
libraries raise no exceptions, and tracking filters treat the static
values as proof of a stationary physical environment. While the physical
robot continues to articulate, the planner calculates goal vectors
relative to a phantom coordinate frame pinned to where the arm was
located hundreds of milliseconds earlier. Because numerical pipelines
cannot distinguish between a stationary object and a halted sensor
pipeline, the machine executes open-loop motions with complete
mathematical confidence.

The safety of a physical architecture depends on treating every spatial
transform as a decaying belief bounded by an explicit expiration budget.
A transform without an epoch is an unverified assertion, and a transform
whose age exceeds task tolerance must be discarded before it reaches
motor drivers. When direct geometric observations become unavailable,
the machine cannot simply halt every operation at the first missed
sensor frame. Extending action across brief sensing gaps requires
estimators that propagate physical belief forward until new measurements
arrive to confirm or correct the state.

\section{Belief Through Occlusion}\label{sec-memory-9-5}

When direct sensing fails, a physical AI system uses an internal
transition model to advance its estimate of the state. A Kalman filter
updates state estimates using plant dynamics,\sidenote{\footnotesize Kálmán
  (\citeproc{ref-kalman1960new}{Kalman 1960}) formulated state
  estimation via recursive projections. Moravec and Elfes
  (\citeproc{ref-moravec1985high}{Moravec and Elfes 1985}) extended
  Bayesian updates to spatial occupancy grids under explicit sensor
  noise models.} while a learned policy might predict the next relative
pose using an autoregressive network. In both implementations,
propagating a state across an unobserved interval does not generate new
information. Prediction is the mathematical extrapolation of old
evidence forward in time. It preserves the utility of an observation by
compensating for expected physical motion, but it does not reset the age
of the underlying measurement. If an optical sensor captured the
position of a workpiece at time \(t = 0\text{ ms}\), an estimator
projecting that position to \(t = 100\text{ ms}\) still relies entirely
on evidence that is \(100\text{ ms}\) old. Treating a predicted state as
an active observation confuses mathematical continuity with physical
truth.

This distinction becomes critical during physical contact. Consider an
articulated arm tasked with inserting a hardened steel pin into a
close-tolerance bearing sleeve. A wrist-mounted stereo camera tracks the
pin and the sleeve during the approach phase. When the parallel-jaw
gripper descends to seat the pin, the mechanical structure of the
gripper and the fingers completely occludes the camera line of sight
\(40\text{ mm}\) above the fixture. For the remaining distance, the
nervous system receives no direct optical measurements of the mating
interface. The tracking pipeline switches to forward propagation,
computing the relative pose between the tool center point and the sleeve
bore from the last verified camera frame combined with joint encoder
integration.

Because every physical model simplifies reality, the error between the
projected state and the physical world expands continuously throughout
the blind interval. In state estimation, this dispersion is governed by
\textbf{occlusion covariance expansion}---the continuous,
physics-governed mathematical widening of state estimation uncertainty
(governed by stochastic plant models such as the continuous Lyapunov
differential equation
\(\dot{\mathbf{P}} = \mathbf{A}\mathbf{P} + \mathbf{P}\mathbf{A}^T + \mathbf{Q}\))
during sensor line-of-sight dropouts, bounding the unobserved divergence
caused by plant velocity uncertainties and disturbance accelerations
before reacquisition. In this analytical example, suppose the vision
system provides an initial position uncertainty bound
\(\epsilon_0 = 0.30\text{ mm}\) at the moment of occlusion
(\(t = 0\text{ ms}\)). During the unobserved descent, two independent
noise sources degrade the prediction. First, unobserved mechanical
vibration and workpiece thermal expansion introduce an unmodeled drift
velocity bounded by \(v_{\text{drift}} = 4.0\text{ mm/s}\). Second,
minor compliance in the arm joints and unmodeled drive backlash create
an unobserved acceleration uncertainty bounded by
\(a_{\text{err}} = 30.0\text{ mm/s}^2\). The worst-case error growth
envelope \(\epsilon(t)\) over an occlusion duration \(t\) follows the
quadratic expansion

\[\epsilon(t) = \epsilon_0 + v_{\text{drift}} t + \frac{1}{2} a_{\text{err}} t^2\]

The insertion operation requires a maximum radial alignment error of
\(\epsilon_{\text{tol}} = 0.80\text{ mm}\) to clear the lead-in chamfer
of the sleeve. Substituting the parameters into the envelope yields the
quadratic relation \(15.0 t^2 + 4.0 t - 0.50 = 0\). Solving for the
positive root gives a maximum admissible rollout duration of
\(t_{\max} \approx 93\text{ ms}\). If the descent velocity is
\(150\text{ mm/s}\), the arm travels \(14.0\text{ mm}\) during this
\(93\text{ ms}\) window. Throughout this interval, the internal
simulation runs smoothly, floating-point coordinates update without
numerical instability, and trajectory generators report valid setpoints.
Yet at \(t = 94\text{ ms}\), the bound of the predicted error exceeds
the physical clearance of the chamfer. The prediction has stopped
functioning as evidence. Allowing the motor drivers to continue downward
motion on predicted coordinates past \(93\text{ ms}\) means accepting
that the pin may strike the rim of the sleeve, generating lateral forces
that can fracture the tooling. The nervous system must terminate the
open-loop motion precisely at this boundary, switching to a guarded
compliance routine or halting the descent before contact occurs.

When the sensor pipeline reacquires line of sight, the system cannot
simply overwrite the propagated state with the new measurement, nor can
it blindly pass the raw measurement into the downstream controller.
Sensor readings obtained immediately after occlusion often suffer from
transient lighting reflections, partial edge visibility, or
correspondence errors. Before blending or resetting the state, the
estimator must calculate the innovation, defined as the difference
between the returning measurement and the model's predicted pose. If the
new measurement falls inside the calculated error envelope
\(\epsilon(t)\), the measurement confirms the model's drift assumptions,
allowing the estimator to collapse the uncertainty and reset the
evidence clock. If the incoming measurement falls outside the envelope,
either the physical object moved unexpectedly or the sensor returned an
outlier. Snapping the coordinate frame directly to an erroneous outlier
introduces an instantaneous position step into the trajectory planner.
The resulting derivative demand requests infinite joint acceleration,
triggering overcurrent faults on the motor drives. The nervous system
must isolate the discrepancy, evaluate whether the measurement is valid,
and reject discontinuous jumps before updating actuation targets.

The identical failure mechanism governs continuous process machines that
regulate fluids or gases. In an automated chemical dispensing loop, a
motor-driven needle valve regulates fluid delivery into a reactor. A
Coriolis mass flow meter measures liquid throughput, but because the
sensor requires fluid mass to transit its measurement tubes, it sits
\(1.5\text{ m}\) downstream of the valve orifice, introducing a
\(120\text{ ms}\) transport delay at nominal flow rates. To maintain
stable closed-loop pressure without waiting for the delayed sensor
frame, an internal flow model predicts real-time throughput from valve
stem displacement and differential pressure. This short model-based
propagation allows the valve actuator to make rapid micro-adjustments
during normal operation. If downstream viscosity shifts or an inline
filter begins to clog, the physical plant diverges from the static
model. If the controller relies on its internal model beyond the
\(120\text{ ms}\) transport window without reconciling the prediction
against delayed meter readings, the calculated flow diverges from
physical reality. The controller continues to open the valve to satisfy
a phantom deficit, over-pressurizing the manifold and bursting the
upstream fitting.

In both discrete manipulation and continuous fluid control, model
prediction serves as a temporary bridge across missing or delayed
observations. The bridge holds only while the accumulated model error
remains smaller than the mechanical tolerance of the task. Forward
propagation does not manufacture certainty out of thin air, and its
authority over physical actuators must expire the moment the error
envelope outgrows the clearance of the physical world.

\section{The Rate of Belief Decay}\label{sec-memory-9-6}

Physical tasks impose hard clearance and error limits on the body. In a
precision peg-in-hole insertion, the chamfer geometry defines an
allowable radial alignment error of \(0.8\text{ mm}\) before the tool
collides with the shoulder of the bore. In a continuous dispensing
process, a fluid temperature error exceeding
\(2.5\text{ }^\circ\text{C}\) alters viscosity enough to double the bead
width and flood the channel. These physical bounds define the maximum
tolerable error budget \(E_{\max}\). The task of state estimation is not
to eliminate error, but to maintain an uncertainty envelope \(E(t)\)
that is guaranteed to remain strictly below \(E_{\max}\) throughout
execution.

The moment an observation frame ends, that envelope expands as
unmeasured dynamics accumulate. Between sensor updates, the nervous
system propagates the state forward by integrating commanded forces
through an internal model of the plant.

\protect\phantomsection\label{callout-mathux2a-9.7}
\begin{fbxSimple}{callout-math}{Fermi Math:}{Continuous-Time Covariance Propagation}
\phantomsection\label{callout-math*-9.7}

In continuous-time estimation theory, the physical system is modeled as
a linear stochastic differential equation:
\[\dot{\mathbf{x}}(t) = \mathbf{A}\mathbf{x}(t) + \mathbf{B}\mathbf{u}(t) + \mathbf{G}\mathbf{w}(t)\]
where \(\mathbf{x}(t) \in \mathbb{R}^n\) is the true state vector,
\(\mathbf{u}(t) \in \mathbb{R}^m\) is the control input, and
\(\mathbf{w}(t) \in \mathbb{R}^p\) is a zero-mean white Gaussian
disturbance with continuous process noise spectral density matrix
\(\mathbf{Q}_c \in \mathbb{R}^{p \times p}\), representing torque
ripple, unmodeled joint friction, and aerodynamic buffeting. The
propagation of the state error covariance matrix
\(\mathbf{P}(t) = \mathbb{E}[(\mathbf{x}(t) - \hat{\mathbf{x}}(t))(\mathbf{x}(t) - \hat{\mathbf{x}}(t))^T]\)
obeys the continuous-time Lyapunov differential equation:\sidenote{\footnotesize The
  term \(\frac{1}{3}q_c t^3\) represents cubic variance expansion from
  integrating acceleration noise through two state integrators, setting
  the fundamental limit on blind dead reckoning.}
\[\dot{\mathbf{P}}(t) = \mathbf{A}\mathbf{P}(t) + \mathbf{P}\mathbf{A}^T + \mathbf{G}\mathbf{Q}_c\mathbf{G}^T\]

For a canonical kinematic double-integrator plant tracking position
\(p\) and velocity \(v\) along a single axis
(\(\mathbf{x} = [p, v]^T\)), the system and disturbance coupling
matrices are:
\[\mathbf{A} = \begin{bmatrix} 0 & 1 \\ 0 & 0 \end{bmatrix}, \quad \mathbf{G} = \begin{bmatrix} 0 \\ 1 \end{bmatrix}, \quad \mathbf{Q}_c = [q_c]\]
where \(q_c\) is the continuous acceleration noise power spectral
density (\((\text{m/s}^2)^2/\text{Hz}\)). Integrating the Lyapunov
equation analytically from an initial covariance
\(\mathbf{P}(0) = \begin{bmatrix} \sigma_p^2(0) & \sigma_{pv}(0) \\ \sigma_{pv}(0) & \sigma_v^2(0) \end{bmatrix}\)
over an unobserved duration \(t\) yields the exact closed-form
covariance matrix:
\[\mathbf{P}(t) = \begin{bmatrix} \sigma_p^2(0) + 2\sigma_{pv}(0) t + \sigma_v^2(0) t^2 + \frac{1}{3} q_c t^3 & \sigma_{pv}(0) + \sigma_v^2(0) t + \frac{1}{2} q_c t^2 \\ \sigma_{pv}(0) + \sigma_v^2(0) t + \frac{1}{2} q_c t^2 & \sigma_v^2(0) + q_c t \end{bmatrix}\]

\end{fbxSimple}

This closed-form derivation exposes the three distinct physical
mechanisms governing spatial memory decay:

\begin{enumerate}
\def\labelenumi{\arabic{enumi}.}
\tightlist
\item
  \textbf{Initial Velocity Dispersion (\(\sigma_v^2(0) t^2\))}: An
  unmeasured velocity error at the moment of occlusion propagates
  linearly into position displacement (\(\sigma_v(0) t\)), driving
  quadratic variance expansion.
\item
  \textbf{Stochastic Acceleration Diffusion (\(\frac{1}{3} q_c t^3\))}:
  White process noise integrated through the double integrator induces
  cubic variance growth, causing rapid spatial dispersion over extended
  sensing dropouts.
\item
  \textbf{Deterministic Disturbance Acceleration
  (\(\frac{1}{2} a_{\mathrm{dist}} t^2\))}: Any unmodeled constant force
  or friction bias bounded by \(|a_{\mathrm{dist}}|\) deflects the
  physical body quadratically with time.\sidenote{\footnotesize Contact introduces
    non-smooth Coulomb friction and stick-slip transitions, causing step
    changes in error growth that require hybrid contact mode tracking.}
\end{enumerate}

As an analytical example (a bench realization will depend on specific
motor thermal curves and bearing preload), consider an articulated arm
holding a tool tip where the operation requires positioning error to
remain below \(E_{\max} = 1.0\text{ mm}\)
(Figure~\ref{fig-uncertainty-growth}). At the final visual update at
\(t = 0\), the estimator reports an initial 3-sigma position envelope of
\(E(0) = 0.1\text{ mm}\) and a velocity uncertainty
\(\sigma_v = 4.0\text{ mm/s}\). Unmodeled joint friction and
transmission backlash inject an unmeasured acceleration disturbance
bounded by \(a_{\mathrm{dist}} = 50\text{ mm/s}^2\). The validated
position error envelope expands over time according to the kinematic
growth bound
\[E(t) = E(0) + 3\sigma_v t + \frac{1}{2} a_{\mathrm{dist}} t^2\]
Evaluating this expression with the stated parameters yields
\(E(t) = 0.1 + 12.0 t + 25.0 t^2\), with \(t\) in seconds and \(E(t)\)
in millimeters. Setting \(E(t) = E_{\max} = 1.0\text{ mm}\) gives the
quadratic equation \(25.0 t^2 + 12.0 t - 0.9 = 0\). The single positive
real root gives \(t = 0.066\text{ s}\), establishing that the unobserved
propagation window cannot exceed \(66\text{ ms}\) before the tool
envelope breaches the task clearance.

An estimator's internal covariance matrix is a product of its own
mathematical assumptions, not a direct measurement of physical
divergence. Filter tuning parameters are frequently softened in
production software to smooth commanded trajectories or prevent filter
divergence, causing theoretical filter covariance to underestimate
actual tracking drift. A defensible decay rate must be established by
measuring empirical prediction error over staged sensor dropout
intervals during hardware characterization. By systematically
suppressing sensor frames for controlled durations from \(10\text{ ms}\)
to \(500\text{ ms}\) while recording ground truth with an external laser
tracker or rigid optical cage, the engineer collects an empirical
distribution of prediction residuals. Because system safety depends on
worst-case deviations rather than average tracking, the error envelope
\(E(t)\) must be fitted to the \(P_{99}\) or \(P_{99.9}\) empirical
error tail across thousands of dropout trials rather than the
theoretical mean covariance.

The \textbf{belief state invalidation horizon} (formalized as the expiry
time \(t_{\mathrm{exp}}\)) is defined as the earliest time at which the
validated error envelope reaches the physical task tolerance,
\[t_{\mathrm{exp}} = \inf \{ t \ge 0 : E(t) \ge E_{\max} \}\] Under this
definition, expiry is neither an arbitrary communication watchdog timer
nor a fixed \(100\text{ ms}\) heuristic chosen for scheduling
convenience. It is a physical deadline dictated by the intersection of
task geometry, actuator authority, and disturbance dynamics. When the
last valid observation occurs at timestamp \(t_0\), the belief state
carries an absolute expiration timestamp
\(t_{\mathrm{deadline}} = t_0 + t_{\mathrm{exp}}\). Beyond
\(t_{\mathrm{deadline}}\), downstream control loops cannot assume the
estimated pose lies within the required mechanical clearance.

\begin{figure}[t]

\centering{

\includegraphics[width=1\linewidth,height=\textheight,keepaspectratio]{index_files/mediabag//Users/VJ/GitHub/MLSysBook-integrate-physical/physical/book/chapters/09-memory/figures/fig09_uncertainty_growth.pdf}

}

\caption{\label{fig-uncertainty-growth}\textbf{Temporal Belief
Uncertainty Growth and Invalidation Horizon.} Covariance expansion
\(E(t) = E(0) + 3\sigma_v t + \frac{1}{2} a_{\mathrm{dist}} t^2\) under
unobserved forward propagation following measurement latch epoch
\(t_0\). The quadratic error envelope expands under unmodeled velocity
uncertainty (\(\sigma_v = 4.0\text{ mm/s}\)) and unmeasured disturbance
acceleration (\(a_{\mathrm{dist}} = 50.0\text{ mm/s}^2\)) until crossing
the task clearance limit \(E_{\max} = 1.00\text{ mm}\) at the
invalidation horizon \(t_{\mathrm{exp}} = 66\text{ ms}\). The lower
lifecycle bar transitions through \texttt{VALID} (full tracking
authority), \texttt{DEGRADED} (reduced speed and expanded safety
margins), and \texttt{EXPIRED} (action permissions revoked). Incoming
measurements within the innovation gate reset the evidence clock to
\(t_0'\), whereas republished cached matrices (the frozen transform
trap) create deceptive software freshness while true physical drift
compounds toward collision.}

\end{figure}%

Because disturbance magnitudes and plant dynamics vary with operating
conditions, \(t_{\mathrm{exp}}\) is not a single static scalar for the
entire machine. A robotic arm traversing free space at
\(0.2\text{ m/s}\) incurs small disturbance torques, yielding a
validated error growth that permits an unobserved window of
\(140\text{ ms}\). When the same arm moves at \(1.5\text{ m/s}\)
carrying an unmodeled \(5.0\text{ kg}\) payload, centrifugal forces and
joint compliance accelerate error accumulation, shrinking
\(t_{\mathrm{exp}}\) to \(22\text{ ms}\). Transitions between free
motion and surface contact introduce impulsive contact forces that
invalidate kinematic drift models within milliseconds. A production
state interface must therefore evaluate decay rates across distinct
operating regimes, selecting the active expiry limit based on payload
mass, velocity, and contact state.

The derived expiration deadline \(t_{\mathrm{exp}}\) spans multiple
orders of magnitude depending on the entity kinematics and task
clearances, as systematized in Table~\ref{tbl-09-belief-decay-horizons}.
While a rigid table fixture tolerates hundreds of milliseconds of
occlusion before thermal and base vibration drift breaches
micro-machining tolerances, a high-speed tool tip traversing near rigid
boundaries expires in less than \(7\text{ ms}\). Similarly, human
coworkers and dynamic mobile robots demand distinct temporal
invalidation horizons coupled to active fail-safe reflexes, proving that
memory freshness is fundamentally an operational physical contract
rather than a uniform software timer.

\begin{longtable}[]{@{}
  >{\raggedright\arraybackslash}p{(\linewidth - 12\tabcolsep) * \real{0.1429}}
  >{\raggedright\arraybackslash}p{(\linewidth - 12\tabcolsep) * \real{0.1429}}
  >{\raggedright\arraybackslash}p{(\linewidth - 12\tabcolsep) * \real{0.1429}}
  >{\raggedright\arraybackslash}p{(\linewidth - 12\tabcolsep) * \real{0.1429}}
  >{\raggedright\arraybackslash}p{(\linewidth - 12\tabcolsep) * \real{0.1429}}
  >{\raggedright\arraybackslash}p{(\linewidth - 12\tabcolsep) * \real{0.1429}}
  >{\raggedright\arraybackslash}p{(\linewidth - 12\tabcolsep) * \real{0.1429}}@{}}
\caption{\textbf{Temporal Belief Decay Horizons and Dynamic Expiry
Budgets across Physical Operational Classes.} Quantitative breakdown of
kinematic disturbance models, nominal dynamics, task clearance budgets
\(E_{\max}\), analytical error envelopes \(E(t)\), derived expiration
horizons \(t_{\text{exp}} = \inf\{t \ge 0 : E(t) \ge E_{\max}\}\), and
mandatory interface fallback responses across six embodied operational
regimes. Evaluated for unobserved blind intervals following direct
sensor dropout.}\label{tbl-09-belief-decay-horizons}\tabularnewline
\toprule\noalign{}
\begin{minipage}[b]{\linewidth}\raggedright
Physical Operational Entity Class
\end{minipage} & \begin{minipage}[b]{\linewidth}\raggedright
Kinematic Drift \& Disturbance Model
\end{minipage} & \begin{minipage}[b]{\linewidth}\raggedright
Nominal Velocity (\(v_0\)) \& Disturbance (\(a_{\text{dist}}\))
\end{minipage} & \begin{minipage}[b]{\linewidth}\raggedright
Task Clearance Margin (\(E_{\max}\))
\end{minipage} & \begin{minipage}[b]{\linewidth}\raggedright
Dynamic Error Growth Envelope \(E(t)\)
\end{minipage} & \begin{minipage}[b]{\linewidth}\raggedright
Validated Expiry Horizon (\(t_{\text{exp}}\))
\end{minipage} & \begin{minipage}[b]{\linewidth}\raggedright
Mandatory Interface Fallback at Expiry
\end{minipage} \\
\midrule\noalign{}
\endfirsthead
\toprule\noalign{}
\begin{minipage}[b]{\linewidth}\raggedright
Physical Operational Entity Class
\end{minipage} & \begin{minipage}[b]{\linewidth}\raggedright
Kinematic Drift \& Disturbance Model
\end{minipage} & \begin{minipage}[b]{\linewidth}\raggedright
Nominal Velocity (\(v_0\)) \& Disturbance (\(a_{\text{dist}}\))
\end{minipage} & \begin{minipage}[b]{\linewidth}\raggedright
Task Clearance Margin (\(E_{\max}\))
\end{minipage} & \begin{minipage}[b]{\linewidth}\raggedright
Dynamic Error Growth Envelope \(E(t)\)
\end{minipage} & \begin{minipage}[b]{\linewidth}\raggedright
Validated Expiry Horizon (\(t_{\text{exp}}\))
\end{minipage} & \begin{minipage}[b]{\linewidth}\raggedright
Mandatory Interface Fallback at Expiry
\end{minipage} \\
\midrule\noalign{}
\endhead
\bottomrule\noalign{}
\endlastfoot
\textbf{Rigid Workcell Fixture / Datum} & Thermal expansion \&
structural base vibration &
\(v_0 = 0.0\text{ mm/s}, \, a_{\text{dist}} = 2.0\text{ mm/s}^2\) &
\(0.20\text{ mm}\) (micrometer tooling clearance) &
\(E(t) = 0.02 + 1.0 t^2\text{ (mm)}\) & \(424\text{ ms}\) & Hold
trajectory; re-zero tool touch-probe against physical datum \\
\textbf{Workpiece on Transfer Conveyor} & Belt speed jitter \& contact
friction micro-slip &
\(v_0 = 600\text{ mm/s} \pm 25\text{ mm/s}, \, a_{\text{dist}} = 150\text{ mm/s}^2\)
& \(2.00\text{ mm}\) (suction cup compliance limit) &
\(E(t) = 0.30 + 25.0 t + 75.0 t^2\text{ (mm)}\) & \(58\text{ ms}\) &
Abort grasp descent; hold standoff clearance above belt \\
\textbf{Articulated Human Coworker} & Ballistic torso motion \&
unobserved limb reach &
\(v_0 = 1.20\text{ m/s}, \, a_{\text{dist}} = 6.00\text{ m/s}^2\) &
\(150\text{ mm}\) (ISO/TS 15066 collaborative margin) &
\(E(t) = 15.0 + 1200 t + 3000 t^2\text{ (mm)}\) & \(92\text{ ms}\) &
Revoke motion permission; engage Category 0/1 monitored stop \\
\textbf{Autonomous Mobile Robot (AMR)} & Unicycle planar kinematics \&
traction wheel slip &
\(v_0 = 1.80\text{ m/s}, \, a_{\text{dist}} = 2.50\text{ m/s}^2\) &
\(300\text{ mm}\) (aisle clearance corridor limit) &
\(E(t) = 30.0 + 1800 t + 1250 t^2\text{ (mm)}\) & \(137\text{ ms}\) &
Halt corridor advance; broadcast spatial yield token on inter-agent
bus \\
\textbf{High-Speed Tool Center Point} & Joint flexibility, centrifugal
torque \& gear backlash &
\(v_0 = 2.50\text{ m/s}, \, a_{\text{dist}} = 18.0\text{ m/s}^2\) &
\(0.80\text{ mm}\) (precision chamfer insertion limit) &
\(E(t) = 0.05 + 50.0 t + 9000 t^2\text{ (mm)}\) & \(6.8\text{ ms}\) &
Decouple feedforward trajectory; engage low-gain impedance compliance \\
\textbf{Free-Falling / Dropped Component} & Unconstrained gravitational
acceleration \& drag &
\(v_0 = 0.0\text{ mm/s}, \, a_{\text{dist}} = 9810\text{ mm/s}^2\) &
\(25.0\text{ mm}\) (dynamic catch intercept basket) &
\(E(t) = 2.0 + 4905 t^2\text{ (mm)}\) & \(68.5\text{ ms}\) & Cancel
dynamic intercept; park manipulator in protected safety home \\
\end{longtable}

When elapsed time reaches \(t_{\mathrm{exp}}\) without a confirming
sensor frame, the state estimation interface must publish an invalid or
degraded status code alongside the state estimate. Publishing stale
coordinates with an unflagged valid status violates the subsystem
contract and forces downstream planners into unguided execution. The
state estimator does not decide what recovery maneuver the machine
executes. The selection between engaging mechanical holding brakes,
falling back to a hard real-time impedance reflex, or executing a safe
stop belongs to the safety enforcement architecture detailed in
Chapter~\ref{sec-enforcement}. The state estimator fulfills its contract
entirely by revoking its claim of certainty the moment the validated
bound breaches \(E_{\max}\).

This expiration contract is an imported result from statistical
estimation theory and rigid-body mechanics, based on three explicit
assumptions: that the physical plant parameters remain within their
identified bounds, that external disturbances do not exceed the modeled
noise spectrum \(a_{\mathrm{dist}}\), and that every relevant error mode
is observable in the measurement stream. The nervous system monitors
these assumptions during runtime by tracking the innovation sequence
whenever sensor measurements arrive. If incoming measurement residuals
consistently exceed the expected distribution, the estimator marks its
own decay model as invalid, forcing immediate state expiration. This
mathematical guarantee holds only so long as divergence remains
observable. If an external disturbance or unmodeled physical interaction
enters a subspace that the sensors cannot detect, the physical error
accumulates without expanding the estimator's calculated uncertainty
envelope, leaving the machine unaware that its belief has decayed.

\section{What Can and Cannot Be Estimated}\label{sec-memory-9-7}

A learned model predicting a physical quantity is not evidence that the
quantity is physically observable. In state estimation, a state
trajectory is \textbf{observable} if and only if the history of
available sensor outputs over a finite time window uniquely determines
the initial state of the machine. When two distinct physical state
histories produce the exact same sequence of sensor readings, no
mathematical operation, filter, or neural network running on the nervous
system can determine which physical trajectory the body actually
followed. The machine receives identical bitstreams from its sensor
interfaces across both realities.

To see why predictability differs from observability, consider an
analytical example of a robotic end-effector regulating contact force
against a rigid surface. Let the true normal contact force be \(x(t)\),
and let the strain-gauge conditioning electronics introduce an unknown,
slowly changing sensor bias \(b(t)\). The raw measurement delivered to
the estimator at each sample interval is \(y(t) = x(t) + b(t)\). Suppose
the sensor outputs a constant stream \(y(t) = 5.0\text{ N}\). This
observation is satisfied by a true physical force
\(x(t) = 5.0\text{ N}\) with zero offset \(b(t) = 0.0\text{ N}\). It is
equally satisfied by a true force \(x(t) = 2.0\text{ N}\) with a sensor
bias \(b(t) = 3.0\text{ N}\), or by \(x(t) = 0.0\text{ N}\) with a bias
\(b(t) = 5.0\text{ N}\). For any arbitrary offset trajectory
\(\Delta(t)\), the state pair \((x(t) - \Delta(t), b(t) + \Delta(t))\)
produces the exact same measurement sequence \(y(t)\). The difference
between the true force and the sensor offset lies in an
\textbf{unobservable subspace} of the measurement system.

Because the measurement stream cannot distinguish between physical load
and sensor drift, an estimator must rely on prior assumptions about how
each state evolves. If a Kalman filter or a learned recurrent estimator
is configured with a belief that the sensor offset \(b(t)\) is static
with near-zero variance, it attributes every variation in \(y(t)\) to
changes in true contact force \(x(t)\). As new samples arrive at
\(1000\text{ Hz}\), the filter incorporates the incoming data,
repeatedly updating its internal state. Because the incoming
measurements match the forward projection of the state dynamics, the
calculated state covariance contracts, and the estimator reports an
increasingly narrow uncertainty bound, publishing an estimate of
\(5.0\text{ N} \pm 0.05\text{ N}\) at \(P_{99}\). Meanwhile, ambient
heat from the joint motors warms the sensor bridge, causing the physical
bias \(b(t)\) to drift at an analytical rate of \(0.2\text{ N/min}\).
Over \(15\text{ min}\), \(b(t)\) shifts by \(3.0\text{ N}\). The true
force on the surface has fallen from \(5.0\text{ N}\) to
\(2.0\text{ N}\), yet the estimator reports high confidence in a state
that has drifted by sixty percent of its nominal value.

The downstream control policy responds to this false certainty by
commanding motor torques to maintain what it believes is a
\(5.0\text{ N}\) preload. Because the true force is only
\(2.0\text{ N}\), the friction force at the contact patch drops below
the shear load of the workpiece, causing the part to slip from the
gripper and fall into the tooling fixture. This failure occurs without
triggering a single internal software diagnostic. The estimator monitors
the \textbf{innovation sequence}, which is the difference between
incoming measurements \(y(t)\) and predicted measurements
\(\hat{y}(t)\). In this drift scenario, the estimator predicted
\(\hat{y}(t) = \hat{x}(t) + \hat{b}(t) = 5.0\text{ N}\), and the
physical sensor returned \(y(t) = 5.0\text{ N}\). The innovation
residual is zero at every time step. An innovation check measures only
whether the physical sensors disagree with the forward projection of the
internal state. It cannot expose an unobservable error mode, because
along the unobservable manifold the sensor output is invariant to the
drift.

Recovering observability requires the system architecture to provide an
additional constraint that breaks the symmetry along the unobservable
subspace. An engineering design accomplishes this in one of three ways.
First, independent sensing can introduce a second measurement channel
with different failure physics, such as monitoring joint motor current
or an optical deflection sensor that does not share the strain gauge's
thermal sensitivity. Second, the machine can perform deliberate physical
excitation, commanding the gripper to briefly break contact with the
workpiece for \(40\text{ ms}\), creating a known boundary condition of
\(x(t) = 0.0\text{ N}\) where the measured sensor value directly exposes
\(b(t)\). Third, the machine can rely on periodic external anchoring,
bringing the tool into contact with a calibrated reference load cell
during scheduled tool-change cycles to reset the bias estimate.

When a physical state cannot be constrained by complementary sensing or
active excitation, the nervous system must treat that state as
unobservable and refuse to publish false precision. The estimation layer
must mark the corresponding state dimensions as unavailable, or expand
their uncertainty envelope according to the worst-case physical drift
rate of the underlying hardware rather than the contracted covariance of
a filtering algorithm. Formal observability analysis for linear and
nonlinear dynamic systems is an established foundation of control
theory, treated rigorously by Kalman (1960) for linear state spaces and
by Hermann and Krener (\citeproc{ref-hermann1977nonlinear}{Hermann and
Krener 1977}) for nonlinear systems via Lie derivatives. A systems
engineer does not rederive these rank conditions at runtime. The
engineering task is to identify which physical states lack full
observability rank under the current sensor suite, and to enforce that
the state interface never reports statistical convergence on an
unconstrained degree of freedom.

State estimation guarantees are imported results that hold only under
four explicit assumptions. First, the sensor forward model must capture
all physical couplings without unmodeled algebraic dependencies. Second,
physical parameters such as component stiffness and thermal coefficients
must remain within pre-calibrated envelopes. Third, sensor noise must
remain zero-mean and uncorrelated with actuator commands. Fourth, the
operating trajectory of the machine must avoid singular configurations
that cause the observability matrix to lose rank. The nervous system
validates these assumptions during operation by cross-checking redundant
measurement channels, checking environmental temperature sensors, and
enforcing active re-zeroing routines before uncertainty timers expire.
When an operating condition violates these assumptions, such as when an
optical sensor becomes occluded or a kinematic configuration reaches a
singular alignment, the nervous system revokes the valid status flag of
the estimate and forces downstream controllers to fall back to
conservative velocity and torque limits.

\section{How Belief Goes Wrong}\label{sec-memory-9-8}

A failure in a state estimation pipeline takes one of two forms. In a
loud failure, a component halts execution, drops a communication packet,
clears an output buffer, or asserts an explicit error bit. Downstream
consumers immediately observe the absence of data or the cleared
validity flag, enabling deterministic fallback routines to halt motion
or transition to open-loop compliance. In a silent failure, the memory
layer continues to publish a complete, numerically well-formed state
record with valid status bits, low variance indicators, and continuous
timestamps, while the numbers inside the record diverge from the
physical configuration of the machine. Downstream control loops consume
these corrupted records without raising an exception, computing feedback
corrections against a phantom geometry. Because physical actuators
execute commands derived from these valid-looking records, silent state
corruption produces high-velocity collisions and unconstrained force
delivery before software safety monitors register an anomaly.

The most deceptive silent failure is the \textbf{frozen transform},
which occurs when software replaces an aging spatial relationship with a
synthetic stream of current timestamps. Consider an optical tracking
node that computes a transformation matrix
\(T_{\text{camera}}^{\text{object}}\) at timestamp
\(t_0 = 0\text{ ms}\). If the optical sensor disconnects, drops frames,
or suffers thread starvation at \(t = 10\text{ ms}\), the underlying
tracking algorithm ceases to produce new spatial evidence. If the
message-passing middleware or an intermediate transport node retains the
last computed matrix in an unexpired buffer and republishes it at a loop
rate of \(1000\text{ Hz}\), generating a new transmission timestamp
\(t_k = t_0 + k \cdot (1\text{ ms})\), the temporal provenance of the
measurement is erased. A downstream tracking controller verifies
freshness by computing data age as
\(\Delta t = t_{\text{now}} - t_{\text{msg}}\). Because the publisher
stamped the message with the current system clock, \(\Delta t\)
evaluates to less than \(1\text{ ms}\), satisfying a strict freshness
threshold of \(\tau_{\text{max}} = 5\text{ ms}\). The controller accepts
the transform as current evidence. If the physical object moves at
\(0.5\text{ m/s}\) while the transform remains static, the controller
accumulates \(50\text{ mm}\) of spatial tracking error in
\(100\text{ ms}\), yet reports zero estimated tracking divergence.

\protect\phantomsection\label{callout-incidentux2a-9.8}
\begin{fbxSimple}{callout-incident}{Landmark Engineering Incident:}{Warehouse AMR Forklift Dynamic Occlusion Collision}
\phantomsection\label{callout-incident*-9.8}

\textbf{System:} Autonomous Mobile Robot (AMR) logistics transport fleet
(\(1.80\text{ m/s}\) differential-drive pallet transporter) operating in
an active high-density fulfillment facility, equipped with dual 2D
safety LiDAR scanners, forward depth cameras, a distributed coordinate
transform tree (\texttt{/tf}), and a 2D/3D dynamic costmap tracking
engine.

\textbf{Incident:} AMR traversing an aisle intersection collided with a
high-speed manual forklift (\(v_{\text{fork}} = 3.20\text{ m/s}\))
traveling along the orthogonal cross-aisle. The AMR suffered severe
frontal chassis buckling, crushing its primary safety LiDAR optical
housing (\(>\$45,000\) direct damage) and forcing a 6-hour facility-wide
logistics dispatch shutdown.

\textbf{Telemetry Snapshot:}

\begin{itemize}
\tightlist
\item
  At \(t = 0.000\text{ s}\), the AMR approaches the blind 4-way
  intersection at \(v_0 = 1.80\text{ m/s}\). High-bay metal pallet racks
  obstruct optical and LiDAR line of sight into the cross-aisle.
\item
  At \(t = 0.080\text{ s}\), the forward LiDAR briefly registers the
  forklift's steel mast through a narrow racking gap at a distance of
  \(d = 4.20\text{ m}\), closing at relative velocity
  \(v_{\mathrm{rel}} = 5.00\text{ m/s}\). The tracking pipeline
  instantiates track token \texttt{entity\_track\_942}.
\item
  At \(t = 0.120\text{ s}\), an unmodeled structural building column
  completely occludes the forklift from all AMR sensor sightlines.
\item
  From \(t = 0.120\text{ s}\) to \(t = 0.650\text{ s}\)
  (\(530\text{ ms}\) blind occlusion window), the dynamic obstacle
  tracker continues forward-propagating \texttt{entity\_track\_942}
  under a constant-velocity Kalman filter. However, the filter's process
  noise covariance \(\mathbf{Q}\) was artificially softened, preventing
  spatial covariance from expanding at the physical rate of unobserved
  acceleration (\(a_{\text{max}} = 2.50\text{ m/s}^2\)).
\item
  Simultaneously, the costmap layer lacked 3D-DDA free-space raycasting
  to assert negative evidence behind the column, while an auxiliary SLAM
  CPU core lockup froze an intermediate coordinate transform edge
  (\texttt{map} \(\to\) \texttt{odom}) in the shared \texttt{/tf} tree.
  The transport bridge continued republishing this frozen matrix under
  fresh local monotonic timestamps.
\item
  At \(t = 0.580\text{ s}\), the trajectory planner evaluated aisle
  clearance against the stale \texttt{/tf} coordinate frame and the
  unexpanded Kalman track, which placed the forklift \(1.85\text{ m}\)
  behind its true physical location. The planner commanded the AMR to
  accelerate through the intersection.
\item
  At \(t = 0.650\text{ s}\), the AMR emerged from behind the rack and
  struck the forklift's steel forks with a kinetic impact energy of
  \(2.92\text{ kJ}\).\sidenote{\footnotesize Industrial mobile robot safety standards
    (ANSI/ITSDF B56.5, ISO 3691-4) mandate uninterrupted spatial fields
    of view and dynamic braking distances under worst-case closing
    velocities. Decoupling existence decay from unobserved velocity
    expansion violates these standards by creating blind travel zones
    under dynamic occlusion.}
\end{itemize}

\textbf{Root Cause \& Physics of Failure:}

\begin{enumerate}
\def\labelenumi{\arabic{enumi}.}
\tightlist
\item
  \emph{Frozen Transform Trap:} A locked SLAM thread caused the
  \texttt{/tf} bridge to republish static coordinate frames with
  synthetic fresh timestamps
  (\(\Delta t = t_{\text{now}} - t_{\text{msg}} < 2\text{ ms}\)),
  completely blinding downstream collision checkers to \(530\text{ ms}\)
  of uncompensated spatial drift.
\item
  \emph{Decoupled Spatial State from Existence Probability:} The
  obstacle tracker held existence probability
  \(p_{\text{exist}} \ge 0.85\) over an arbitrary static timeout
  (\(2.0\text{ s}\)) rather than evaluating the physical invalidation
  horizon (\(t_{\mathrm{exp}} = 137\text{ ms}\) for
  \(E_{\max} = 0.30\text{ m}\) clearance at
  \(v_{\mathrm{rel}} = 5.00\text{ m/s}\)).
\item
  \emph{Absence of Negative Free-Space Evidence:} The volumetric memory
  failed to execute line-of-sight raycasting to clear or bound
  unobserved obstacle tracks as the AMR's viewing geometry evolved.
\end{enumerate}

\textbf{Systems Resolution:} Implemented hardware-timestamped spatial
belief records with immutable acquisition epochs (\(t_0\)), enforced
deterministic track invalidation deadlines
(\(t_{\mathrm{exp}} = \inf\{t : E(t) \ge E_{\max}\}\)), integrated
3D-DDA raycasting with explicit negative evidence clearing, and added a
cross-checking watchdog on \texttt{/tf} transform tree edges that
revokes motion authority whenever transform jitter exceeds
\(10\text{ ms}\).

\textbf{Systems Takeaway:} In dynamic environments, an unobserved
obstacle's spatial belief decays rapidly under closing velocities;
projecting a track forward without exponential covariance expansion and
strict physical expiry converts temporary sensor occlusion into
high-energy kinetic impact.

\end{fbxSimple}

Phantom-object tracks arise when the nervous system decouples the
lifetime of spatial coordinates from the lifetime of existence evidence.
A perception module maintains an object track comprising a kinematic
state vector \(\mathbf{x} = [p, v]^T\) and an existence probability
\(p_{\text{exist}} \in [0, 1]\). When the physical object is removed
from the workspace or an occlusion begins, raw detector hits drop to
zero. If the estimation layer continues to propagate \(\mathbf{x}\)
forward using process dynamics while existence probability decay is
governed by a long time constant (for instance, an analytical decay rate
where \(p_{\text{exist}}\) remains above an operating threshold for
\(1500\text{ ms}\) without reinforcement), the tracker maintains a
persistent phantom entity. The path planner routes the manipulator
around empty space, or an automated grasping routine closes its fingers
on an empty coordinate. Detecting a phantom track cannot rely on passive
measurement absence. It requires an active negative observation, where a
line-of-sight sensor sweeps the volume and confirms that the expected
occupancy is empty, or an unfulfilled contact expectation when a force
sensor registers zero load at the predicted contact boundary.\sidenote{\footnotesize Autonomous
  safety standards codify spatial belief validity. ANSI/UL 4600 Section
  8.3 requires explicit data age bounds, deterministic invalidation of
  unobserved dynamic tracks, and active negative-obstacle checks. ISO
  26262-6 Clause 7 mandates temporal freshness verification and ASIL
  fault isolation for shared transform states.}

Coordinate frame misattributions divide along a similar boundary between
syntax and semantics. A schema-detectable frame mismatch occurs when a
publisher labels a transform with an invalid or undefined frame
identifier in the transformation tree graph. A message bus with static
schema enforcement rejects the packet at the serialization boundary
because the named frame cannot be resolved to a known parent link. In
contrast, a semantically wrong frame assignment occurs when an engineer
or configuration file assigns a valid, type-correct frame identifier
that represents the wrong physical link. For example, commanding a tool
center point trajectory using the transform of the bare mounting flange
\(T_{\text{flange}}\) instead of the extended tool tip
\(T_{\text{tool}}\) satisfies every schema validation check, type
assertion, and tree traversal test. If the physical tool tip extends
\(150\text{ mm}\) beyond the mounting flange along the tool \(z\)-axis,
the kinematic solver computes joint angles that drive the physical metal
of the tool into the workspace surface before the commanded frame
reaches its target. Detecting a semantic frame error requires physical
consistency checks, such as verifying measured motor torques against an
analytical dynamic model of the arm, or cross-referencing multi-camera
optical tracker observations against joint encoder kinematics.

When sensor updates drop out, an estimator relies entirely on its
predictive forward model to advance state estimates across the blind
interval. A well-formulated filter updates the mean state estimate
\(\hat{\mathbf{x}}(t)\) while monotonically expanding the estimation
covariance \(\mathbf{\Sigma}(t)\) to reflect the growth of unconstrained
process noise. An overconfident predictor understates process variance
\(\mathbf{Q}\), holding \(\mathbf{\Sigma}(t)\) artificially compact over
an extended prediction duration exceeding \(200\text{ ms}\). When fresh
sensor measurements eventually arrive, the small covariance prevents the
filter from correcting its trajectory, treating valid measurements as
statistical outliers and rejecting them. Engineers identify
overconfidence during system qualification by executing staged synthetic
occlusions, systematically dropping sensor frames for intervals between
\(50\text{ ms}\) and \(500\text{ ms}\) and evaluating the reacquisition
innovation residual \(y(t) - \hat{y}(t)\). If the reacquisition residual
consistently exceeds the predicted \(3\sigma\) covariance bounds across
fifty analytical trials, the process noise model is invalid. However,
when drift occurs along unobservable degrees of freedom (such as
workpiece slippage inside compliant gripper pads during
high-acceleration transfers), no reacquisition residual can form because
no sensor directly measures the displacement. Unobservable physical
drift evades both residual checks and staged occlusion detectors,
persisting until the machine makes physical contact with an external
reference surface.

Preventing stale, misattributed, and unbounded estimates from
compromising the physical machine requires formal structure at the
boundary between estimation and control. An ad hoc dictionary of
floating-point values and raw transformation matrices cannot distinguish
between active evidence and cached momentum. Every element of spatial
memory must reside within an explicit state record that couples the
numerical value to its base coordinate frame, its sensor provenance, its
decay envelope, and a hard invalidation deadline that forces downstream
consumers to act on physical truth rather than numerical habit.

\section{The Spatial State Schema}\label{sec-memory-9-9}

The state record instantiates the notebook schema established in
Chapter~\ref{sec-nervous} by binding the physical state estimate to the
temporal and spatial constraints under which that estimate remains true.
Where Chapter~\ref{sec-nervous} defines generic message structures,
Chapter~\ref{sec-memory} introduces seven memory-specific metadata
fields required to govern decay. Every recorded state entry contains the
physical quantity, its SI units, the unique entity identity of the
physical object, the reference coordinate frame in which coordinates are
expressed, the evidence timestamp \(t_e\) recording when the originating
photons or forces touched a sensor, the estimate timestamp \(t_m\)
produced by the estimator forward model, and the explicit clock domain
identifier. Separating evidence time from estimate time ensures that a
predictor running at \(1000\text{ Hz}\) cannot disguise an old optical
measurement taken \(80\text{ ms}\) earlier as immediate ground truth.
Stating the clock domain prevents uncompensated bus delays or clock
drift across distributed microcontrollers from corrupting the delta
\(t_m - t_e\).

The \textbf{state record} is the immutable memory notebook extension
that encapsulates an estimated physical quantity \((\mathbf{x})\) with
its reference frame identifier \((\mathcal{F})\), physical evidence
epoch (\(t_e\)), computational estimate epoch (\(t_m\)), initial error
bound (\(\epsilon_0\)), validated growth law (\(g(\Delta t)\)),
operating envelope (\(\mathcal{E}\)), and discrete lifecycle status.

To govern how fast an estimate degrades across blind intervals, the
record stores the initial error bound \(\epsilon_0\), a validated error
growth law \(g(\Delta t)\), the operating envelope \(\mathcal{E}\), the
expiry threshold \(\tau\), and the statistical confidence level
supporting the bound. In an analytical example of a tactile peg
insertion, the state record tracks the lateral peg position
\(x_{\text{peg}}\) in the assembly frame with an initial uncertainty
bound \(\epsilon_0 = 0.05\text{ mm}\) derived from touch contact at
\(t_e\), supported at a \(P_{99.9}\) confidence level. When contact
breaks and optical tracking is occluded, the validated error growth law
models lateral displacement drift as
\(g(\Delta t) = \epsilon_0 + v_{\max} \Delta t + \frac{1}{2} a_{\text{slip}} (\Delta t)^2\),
where \(v_{\max} = 0.01\text{ m/s}\) is maximum steady gripper creep and
\(a_{\text{slip}} = 0.2\text{ m/s}^2\) is acceleration-induced pad slip.
The operating envelope specifies the kinematic limits under which this
growth law was verified, such as arm accelerations below
\(2.0\text{ m/s}^2\). If the arm exceeds those dynamics, the model
\(g(\Delta t)\) is unvalidated, and the bound collapses. The expiry
threshold \(\tau\) sets the maximum tolerable physical error before the
estimate becomes structurally unusable.

Overloading data absence or a single null pointer to represent every
estimation failure deprives downstream consumers of the physical reason
a belief cannot be trusted. The record enforces an explicit discrete
lifecycle status with four mandatory states comprising valid, degraded,
expired, and unobservable. An entry is valid when the elapsed duration
\(\Delta t = t_{\text{curr}} - t_e\) yields an accumulated error
\(g(\Delta t)\) within nominal operational tolerances and operating
envelope invariants hold. It transitions to degraded when \(\Delta t\)
causes \(g(\Delta t)\) to exceed high-precision operational thresholds
while remaining inside the hard safety limit \(\tau\), requiring
downstream planners to reduce velocity or widen clearance buffers. It
transitions to expired the instant \(\Delta t \ge \tau\) or when an
envelope invariant is violated, forbidding the consumer from executing
actions that depend on the value. Finally, the state is marked
unobservable when the estimator detects that the underlying physical
degree of freedom has decoupled from all active sensors, such as
workpiece rotation inside a compliant gripper where no internal sensor
measures angular slip, even if elapsed time has not yet reached
\(\tau\).

The numeric value of \(\tau\) is not an arbitrary timeout chosen for
computational convenience. It is strictly bounded by the physical
clearance of the tooling and the maximum closing velocity of the
actuators. In a chamferless insertion task with a radial clearance of
\(0.10\text{ mm}\), exceeding a lateral uncertainty of
\(0.10\text{ mm}\) guarantees physical collision between the workpiece
shoulder and the mating hole edge. Given the analytical slip model where
error reaches \(0.10\text{ mm}\) after \(45\text{ ms}\) of unobserved
motion, \(\tau\) must be set to \(45\text{ ms}\) or less. The state
record mandates the required consumer response upon entering degraded or
expired states without defining the lower-level execution machinery
reserved for Chapter~\ref{sec-enforcement}. For an expired target
position, the record specifies whether the consumer must hold current
joint positions, transition to open-loop impedance compliance, or
surrender authority to an upstream supervisor, ensuring that the
consumer cannot silently retain stale numbers as frozen targets.

Freshness cannot be manufactured by copying data between memory
locations or updating timestamps on unobserved states. To prevent
synthetic freshness loops, the record preserves the cryptographic hash
or identifier of the raw source observations, the exact software version
of the estimation model, the full transform tree provenance linking the
base frame to the sensor origin, and the designated update ownership
token. Only the originating estimator process possesses write authority
for a given entity record. If an intermediate planning node reads an
estimate and republishes it across a middleware bus, it cannot update
\(t_e\) or reset the status to valid without providing fresh sensor
observation hashes. Preserving transform provenance ensures that if a
parent frame in the kinematic chain becomes invalid, all child state
records derived through that transform branch invalidate simultaneously.

By isolating decay rates, provenance hashes, and discrete statuses in
standardized fields, the state record provides a direct target for
systematic verification. In Chapter~\ref{sec-verification},
fault-injection harnesses evaluate software resilience by manipulating
the record under test, such as artificially holding \(t_e\) constant
while advancing \(t_m\), corrupting frame identifiers, or injecting
sensor dropout to verify that consumers transition into their specified
fallback modes within one control cycle. In Chapter~\ref{sec-release},
the safety case utilizes the recorded bounds, operating envelopes, and
verified confidence levels as empirical evidence to defend the claim
that unobserved drift cannot exceed mechanical safety margins during
normal and degraded operation.

The complete memory interface contract is formalized in
Table~\ref{tbl-09-invalidation-epochs}, defining the binary encoding,
hardware timing authorities, provenance validation invariants,
deterministic invalidation triggers, and failure modes across all eight
state metadata fields. By strictly binding evidence epochs to hardware
input-capture timers, enforcing cryptographic entity tracking hashes,
and invalidating beliefs whenever the dynamic error bound \(E(t)\)
crosses task limits, this schema prevents intermediate software proxies
from manufacturing synthetic freshness.

\begin{longtable}[]{@{}
  >{\raggedright\arraybackslash}p{(\linewidth - 10\tabcolsep) * \real{0.1667}}
  >{\raggedright\arraybackslash}p{(\linewidth - 10\tabcolsep) * \real{0.1667}}
  >{\raggedright\arraybackslash}p{(\linewidth - 10\tabcolsep) * \real{0.1667}}
  >{\raggedright\arraybackslash}p{(\linewidth - 10\tabcolsep) * \real{0.1667}}
  >{\raggedright\arraybackslash}p{(\linewidth - 10\tabcolsep) * \real{0.1667}}
  >{\raggedright\arraybackslash}p{(\linewidth - 10\tabcolsep) * \real{0.1667}}@{}}
\caption{\textbf{Physical Memory Invalidation Epochs and Hardware
Timestamp Provenance Schema.} Formal specification of the eight
mandatory metadata fields governing belief lifecycle, hardware clock
domain bindings, error envelope parameterization, and deterministic
invalidation triggers required to prevent stale-state execution across
distributed physical AI
sub-brains.}\label{tbl-09-invalidation-epochs}\tabularnewline
\toprule\noalign{}
\begin{minipage}[b]{\linewidth}\raggedright
Schema Field Identifier
\end{minipage} & \begin{minipage}[b]{\linewidth}\raggedright
Binary Encoding \& Wire Type
\end{minipage} & \begin{minipage}[b]{\linewidth}\raggedright
Clock Domain \& Hardware Authority
\end{minipage} & \begin{minipage}[b]{\linewidth}\raggedright
Provenance Validation Invariant
\end{minipage} & \begin{minipage}[b]{\linewidth}\raggedright
Deterministic Invalidation Trigger
\end{minipage} & \begin{minipage}[b]{\linewidth}\raggedright
Physical Failure Mode on Staleness / Omission
\end{minipage} \\
\midrule\noalign{}
\endfirsthead
\toprule\noalign{}
\begin{minipage}[b]{\linewidth}\raggedright
Schema Field Identifier
\end{minipage} & \begin{minipage}[b]{\linewidth}\raggedright
Binary Encoding \& Wire Type
\end{minipage} & \begin{minipage}[b]{\linewidth}\raggedright
Clock Domain \& Hardware Authority
\end{minipage} & \begin{minipage}[b]{\linewidth}\raggedright
Provenance Validation Invariant
\end{minipage} & \begin{minipage}[b]{\linewidth}\raggedright
Deterministic Invalidation Trigger
\end{minipage} & \begin{minipage}[b]{\linewidth}\raggedright
Physical Failure Mode on Staleness / Omission
\end{minipage} \\
\midrule\noalign{}
\endhead
\bottomrule\noalign{}
\endlastfoot
\texttt{evidence\_epoch\_t0} & \texttt{int64} (64-bit ns since PTP
epoch) & Hardware Input-Capture Timer / Sensor PHY & Monotonically
strictly increasing (\(t_{0, k} > t_{0, k-1}\)); matches raw PHY frame
header & Stale when \(t_{\text{now}} - t_0 \ge \tau_{\text{exp}}\);
invalid if re-stamped by proxy & Silent frozen transform; blind motion
into displaced physical obstacles \\
\texttt{clock\_domain\_id} & \texttt{uint16} ENUM (\texttt{0x01}=PTP,
\texttt{0x02}=TCXO, \texttt{0x03}=CAN) & Grandmaster Clock / Hardware
Sync Bus & Must match recognized domain topology in active hardware
registry & Invalidate if sync lock lost or inter-domain offset
\(\delta t > \delta t_{\max}\) & Asynchronous time skew; uncalibrated
spatial shear \(e = v_{\text{rel}} \Delta t\) \\
\texttt{sync\_jitter\_bound} & \texttt{uint32} (32-bit ns, \(3\sigma\)
upper bound) & Hardware PLL / PTP Servo Filter & Monitored
\(\le 500\,\mu\text{s}\); verified by NIC hardware timestamp counter &
Invalidate when PLL unlocks or network packet delay variation (PDV)
spikes & Hidden positioning uncertainty
\(e_{\text{sync}} = v \cdot \delta t\) breaching task budget \\
\texttt{entity\_uuid\_hash} & \texttt{uint64} (Cryptographic SIPHash of
Track ID) & Central Spatial Registry / Track Server & Unique per
physical entity; bound to explicit coordinate frame tree node &
Invalidate upon negative line-of-sight sweep or track termination token
& Persistent phantom tracks; grasping empty space or avoiding cleared
zones \\
\texttt{coordinate\_frame\_id} & \texttt{uint16} (Unique TF graph node
token) & Kinematic Tree Server / Baseboard MPU & Must resolve to active
path in kinematic transform tree without cycles & Invalidate if any
parent transform in the kinematic branch expires or faults & Coordinate
misattribution; force applied along wrong physical axis \\
\texttt{growth\_envelope\_params} & \texttt{float32{[}4{]}}
(\([\epsilon_0, v_{\text{drift}}, a_{\text{max}}, P_{\text{conf}}]\)) &
Offline System ID / Characterization Lab & Verified against empirical
dropout distributions (\(P_{99.9}\) confidence) & Invalidate when
physical accelerations exceed verified envelope \(\mathcal{E}\) &
Overconfident state propagation; unbounded slip during blind motion \\
\texttt{lifecycle\_status} & \texttt{uint8} ENUM (\texttt{0x01}=VALID,
\texttt{0x02}=DEGRADED, \texttt{0x04}=EXPIRED,
\texttt{0x08}=UNOBSERVABLE) & Local State Record Invalidation Engine &
State transitions obey strict FSM; cannot revert to VALID without new
\(t_0\) & Automatic transition to EXPIRED when
\(E(t_{\text{now}}) \ge E_{\max}\) or fault detected & Unflagged
degraded execution; bypass of fail-safe impedance reflexes \\
\texttt{negative\_evidence\_epoch} & \texttt{int64} (64-bit ns of last
clear sweep) & Active Raycast Pipeline / LiDAR Grid & Ray sweep frustum
intersects entity volume with zero observed returns & Forces
\texttt{lifecycle\_status\ =\ EXPIRED} immediately upon volume clearance
& Ghost obstacle blocking path planning; unrecoverable workcell
starvation \\
\end{longtable}

When each individual state record faithfully accounts for its own decay
and strictly enforces its own expiry deadline, the nervous system avoids
the hazard of acting on single stale beliefs. However, physical
manipulation rarely depends on a solitary coordinate in isolation. A
machine holding an object against a fixture evaluates multiple spatial
relations simultaneously, drawing on separate state records populated by
distinct estimators operating across different sensor modalities and
clock domains. When these independent beliefs are evaluated together,
new failure modes emerge where every individual record is internally
valid, freshly observed, and mathematically sound, yet the combined
geometric picture describes a physical impossibility.

\section{Systems Stress-Tests}\label{sec-memory-9-10}

When individual state records are evaluated in isolation, each may
satisfy its freshness deadline, maintain bounded covariance, and present
clean cryptographic provenance. In a physical machine that touches and
regulates, failures rarely announce themselves through a single expired
record. Instead, hazardous behaviors emerge when multiple estimates,
each internally consistent and unexpired, are evaluated jointly by a
downstream policy or controller. The following five stress tests
evaluate how joint inconsistency across asynchronous sensing, spatial
frame fusion, track termination, common-cause sensor bias, and kinematic
epoch skew creates physical failure modes:

\subsection*{Stress-Test 1: Asynchronous Dual-Arm Assembly (Timestamp
Skew in Relative
Motion)}\label{stress-test-1-asynchronous-dual-arm-assembly-timestamp-skew-in-relative-motion}
\addcontentsline{toc}{subsection}{Stress-Test 1: Asynchronous Dual-Arm
Assembly (Timestamp Skew in Relative Motion)}

\emph{The scenario.} A dual-arm assembly cell seats a cylindrical shaft
into a moving collar. Arm 1 moves the shaft along the insertion axis at
\(v_1 = 0.35\text{ m/s}\), while Arm 2 tracks the collar translating at
\(v_2 = -0.15\text{ m/s}\) (relative closing velocity
\(v_{\mathrm{rel}} = 0.50\text{ m/s}\)). Both state estimators publish
records with a \(50.0\text{ ms}\) expiry. Camera 1 captures shaft
position at \(t_1 = 100.0\text{ ms}\) and Camera 2 captures collar
position at \(t_2 = 120.0\text{ ms}\). At evaluation time
\(t = 125.0\text{ ms}\), both records remain unexpired. The insertion
policy commands alignment based on a nominal radial clearance of
\(1.5\text{ mm}\).

\emph{The systems dilemma.} The controller evaluates spatial separation
without reconciling timestamps. Over the skew interval
\(\Delta t = |t_2 - t_1| = 20.0\text{ ms}\), relative displacement
shifts by
\(\Delta x = v_{\mathrm{rel}} \Delta t = 0.50\text{ m/s} \times 0.020\text{ s} = 10.0\text{ mm}\).
Because this \(10.0\text{ mm}\) skew error exceeds the \(1.5\text{ mm}\)
clearance by more than a factor of six, the commanded trajectory drives
the shaft directly into the collar lip.

\emph{The systems verdict: Joint Temporal Alignment Law.} Individual
freshness guarantees only that a measurement was true when acquired, but
multi-body physical interaction demands that multiple beliefs be
reconciled to the same microsecond instant. Multi-entity controllers
must forward-propagate asynchronous beliefs through their dynamic models
or enforce maximum allowable observation skew
\(\Delta t_{\mathrm{max}} = \delta_{\mathrm{clearance}} / v_{\mathrm{rel}}\).

\subsection*{Stress-Test 2: Conflicting Spatial Estimates in Robotic
Milling (The Fusion Midpoint
Fallacy)}\label{stress-test-2-conflicting-spatial-estimates-in-robotic-milling-the-fusion-midpoint-fallacy}
\addcontentsline{toc}{subsection}{Stress-Test 2: Conflicting Spatial
Estimates in Robotic Milling (The Fusion Midpoint Fallacy)}

\emph{The scenario.} An overhead vision system and an in-hand optical
profilometer estimate the origin of a titanium blank clamped to a
workbench. The overhead camera reports origin at
\(x_{\mathrm{base}} = 145.0\text{ mm} \pm 1.5\text{ mm}\) (\(3\sigma\)).
The in-hand profilometer reports the same feature at
\(x_{\mathrm{base}} = 149.0\text{ mm} \pm 1.0\text{ mm}\). The combined
error bound is \(2.5\text{ mm}\), but the spatial separation between
means is \(|145.0\text{ mm} - 149.0\text{ mm}| = 4.0\text{ mm}\).

\emph{The systems dilemma.} An unvalidated sensor fusion module computes
a variance-weighted average of \(147.4\text{ mm} \pm 0.83\text{ mm}\).
When the cutter approaches at rapid traverse, the tool plunges into
material \(2.4\text{ mm}\) deeper than intended, inducing severe chatter
and spindle overcurrent.

\emph{The systems verdict: Geometric Innovation Gating Law.} Averaging
conflicting spatial estimates whose separation exceeds the Minkowski sum
of their bounds does not resolve uncertainty; it fabricates a confident
midpoint where no physical surface exists. Spatial estimators must
reject fusion and halt motion when Chi-squared innovation gates fail.

\subsection*{Stress-Test 3: Track Termination and Ballistic Descent
(Covariance vs Physical
Existence)}\label{stress-test-3-track-termination-and-ballistic-descent-covariance-vs-physical-existence}
\addcontentsline{toc}{subsection}{Stress-Test 3: Track Termination and
Ballistic Descent (Covariance vs Physical Existence)}

\emph{The scenario.} A robotic arm picks steel components from a feed
belt moving at \(0.80\text{ m/s}\). An operator removes a component
before the pick zone. The tracking filter receives no subsequent visual
detections, but continues extrapolating component position using its
belt velocity model (\(\sigma_x(t) = \sigma_0 + \alpha t\)), maintaining
variance below the rejection threshold for \(300.0\text{ ms}\).

\emph{The systems dilemma.} At \(t = 200.0\text{ ms}\) after removal,
the gripper executes a ballistic descent to grasp the component at the
extrapolated coordinate, slamming fingers into an empty belt and
striking an alignment guide.

\emph{The systems verdict: Existence Decoupling Law.} Spatial covariance
represents geometric uncertainty assuming an entity exists. It contains
zero information regarding whether the entity remains present in the
physical world. Estimators must couple spatial propagation with an
existence probability field that decays upon negative visual evidence.

\subsection*{Stress-Test 4: Common-Cause Sensor Drift in Force
Regulation (The Illusion of
Redundancy)}\label{stress-test-4-common-cause-sensor-drift-in-force-regulation-the-illusion-of-redundancy}
\addcontentsline{toc}{subsection}{Stress-Test 4: Common-Cause Sensor
Drift in Force Regulation (The Illusion of Redundancy)}

\emph{The scenario.} An articulated press regulates \(200.0\text{ N}\)
normal force against an aircraft panel using two estimators: an
analytical joint-current model and a learned neural estimator ingesting
wrist strain-gauge voltages. A nearby radiant heater elevates the wrist
casing by \(18.0^{\circ}\text{C}\), introducing a \(+15.0\text{ N}\)
thermal zero-point bias into the shared strain-gauge bridge.

\emph{The systems dilemma.} Because both estimators draw from the same
uncompensated analog channel, both report an applied force of
\(200.0\text{ N}\) with zero residual, showing perfect mutual agreement.
In physical reality, the exerted force is only \(185.0\text{ N}\),
leaving the structural bond dangerously under-compacted.

\emph{The systems verdict: Sensor Provenance Invariant.} Redundant
software estimators provide zero safety margin if their noise models and
systematic biases share a common physical hardware transducer or analog
line. Genuine redundancy requires heterogeneous sensing modalities.

\subsection*{Stress-Test 5: Kinematic Epoch Skew in Turbine Blade
Finishing (Moving Coordinate
Frames)}\label{stress-test-5-kinematic-epoch-skew-in-turbine-blade-finishing-moving-coordinate-frames}
\addcontentsline{toc}{subsection}{Stress-Test 5: Kinematic Epoch Skew in
Turbine Blade Finishing (Moving Coordinate Frames)}

\emph{The scenario.} A mobile manipulator finishes a turbine blade while
its base translates at \(0.20\text{ m/s}\). An in-hand scanner acquires
surface profile points at \(t_{\mathrm{scan}} = 50.0\text{ ms}\). The
kinematics engine computes the base-to-tip coordinate transform at
\(t_{\mathrm{kin}} = 70.0\text{ ms}\) while the arm decelerates at
\(2.0\text{ m/s}^2\). Both records carry a \(100.0\text{ ms}\) expiry
window.

\emph{The systems dilemma.} The surface reconstruction node transforms
the \(50.0\text{ ms}\) point cloud using the \(70.0\text{ ms}\)
kinematic matrix because both records are unexpired. Over the
\(20.0\text{ ms}\) epoch gap, the base moves \(4.0\text{ mm}\) and the
arm flexes under deceleration, inducing a \(5.8\text{ mm}\) spatial
distortion in the computed surface normal that causes the grinder to
gouge the blade.

\emph{The systems verdict: Transform Epoch Invariant.} Kinematic
transforms are continuous functions of time. Applying an unexpired
transform from the wrong epoch distorts geometric truth as destructively
as corrupting raw sensor packets. State records must enforce strict
epoch matching (\(|t_{\mathrm{scan}} - t_{\mathrm{kin}}| \le \epsilon\))
across transform lookups.

\section{Summary}\label{sec-memory-9-11}

A state estimate handed across the interface to the planning and control
architecture is not a bare array of floating-point coordinates. It is a
bound data structure that carries its physical semantics, the exact
timestamp of the physical measurement that produced it, an explicit
physical-error bound, and the hardware provenance of its sensing path.
When a tactile probe or vision sensor measures the surface of a
workpiece, the value produced is meaningful only within the coordinate
frame where the physical interaction occurred, at the precise
microsecond of measurement. If a downstream consumer receives a position
vector without knowing that the underlying optical measurement occurred
\(45\text{ ms}\) in the past and relies on a kinematic chain calibrated
at a different epoch, the consumer acts on an illusion of geometry.
Downstream modules do not need raw voltage streams, but they must
receive a belief package whose provenance and temporal origins cannot be
stripped away. The policy and the nervous system must never consume a
coordinate without knowing which sensor captured it, through which bus
it traveled, and when that evidence ceased to be fresh.

State propagation algorithms, such as numerical integration through a
kinematic model or a forward dynamic filter, advance the timestamp of a
state vector to match the current execution cycle. This forward
projection bridges the latency gap between discrete sensor arrivals, but
it generates zero new physical evidence. Extrapolating a position from
an encoder reading taken at \(t_0 = 0\text{ ms}\) to
\(t_1 = 40\text{ ms}\) under an assumption of constant velocity yields a
computational state for \(t_1\), yet the evidence supporting that value
remains rooted entirely at \(t_0\). The age of the belief is
\(40\text{ ms}\), not \(0\text{ ms}\). Calling a kinematic predictor
twenty times in succession without a fresh measurement does not refresh
the memory of the machine. It merely inflates the spatial variance
envelope according to the physical limits of unmodeled acceleration.
Confusing the prediction horizon with evidence arrival is the root
mechanism of stale-state execution, where a controller acts with
mathematical certainty while tracking an extrapolation that has drifted
away from the physical workpiece.

The lifetime of a state estimate cannot be governed by convenient
software timeouts, such as an operating system watchdog or a
communications transport buffer set to a round \(100\text{ ms}\). Expiry
is a physical calculation dictated by the mechanical tolerances of the
workpiece and the validated growth rate of unobserved error. If a
milling cutter operates with an allowable geometric tolerance of
\(0.5\text{ mm}\), and the tool carriage can accelerate under unexpected
load disturbances at up to \(4.0\text{ m/s}^2\), then after an
observation dropout of \(\Delta t = 20\text{ ms}\), the unmeasured
displacement can reach \(0.8\text{ mm}\) via the second-order expansion
\(\frac{1}{2} a (\Delta t)^2\). At the threshold of \(15.8\text{ ms}\),
the unobserved error envelope crosses the \(0.5\text{ mm}\) tolerance
limit and the belief expires, regardless of whether the software message
queue remains healthy. The nervous system must invalidate that belief at
the exact physical budget boundary. A frozen transform or an expired
state held in a transport cache is far more hazardous than an explicit
missing frame, because a missing frame triggers an immediate reflex or a
controlled stop, whereas a stale transform allows the machine to drive
full power into a physical fixture.

A systems engineer must preserve the operational distinction between
modeled uncertainty and error that the available measurements cannot
observe. Modeled uncertainty expands predictably, widening the
covariance matrix as process noise accumulates during prediction steps.
Unobservable error, in contrast, occurs when the available sensors
cannot distinguish between physical states, such as when redundant force
sensors share an elastic mounting plate that deforms under thermal load,
or when an in-hand scanner experiences unmeasured structural vibration
along the tool arm. In these regimes, the estimator may report a tightly
bounded covariance and negligible residual between redundant channels,
falsely indicating high agreement while the physical tool is off target
by several millimeters. Guarding against unobservable drift requires
cross-modal validation across mechanically independent sensing channels
and explicit hardware lineage records, verifying that redundant
estimators do not share common-mode physical failure paths or
uncalibrated mechanical compliance.

\protect\phantomsection\label{callout-takeawaysux2a-9.9}
\begin{fbxSimple}{callout-takeaways}{Key Chapter Takeaways}{}
\phantomsection\label{callout-takeaways*-9.9}

\begin{itemize}
\tightlist
\item
  \textbf{Prediction Advances State but Generates Zero Evidence:}
  Forward-propagating a state vector through kinematic differential
  equations advances the timestamp, but evidence age
  (\(t_{\text{age}}\)) remains rooted at the raw measurement epoch.
  Calling a forward predictor without new measurements inflates spatial
  variance (\(\mathbf{\Sigma}(t)\)).
\item
  \textbf{Temporal Skew Across Moving Bodies Destroys Spatial
  Clearance:} In multi-entity relative motion
  (\(v_{\text{rel}} = 0.5\text{ m/s}\)), an unsynchronized measurement
  skew of \(\Delta t = 20\text{ ms}\) induces
  \(\Delta x = 10\text{ mm}\) of unmodeled geometric error, requiring
  all relative states to be strictly reconciled to the identical
  microsecond epoch.
\item
  \textbf{Innovation Gating Must Precede Sensor Fusion (The Midpoint
  Fallacy):} When two sensors disagree by more than their combined
  uncertainty margins (e.g., a camera sees a block at \(145\text{ mm}\)
  and a tactile probe feels it at \(149\text{ mm}\) with \(1\text{ mm}\)
  bounds), blindly averaging them computes a mathematically optimal
  midpoint (\(147\text{ mm}\)) where no physical object actually exists.
  Systems must reject fusion and halt motion when Chi-squared
  (\(\chi^2\)) innovation gates fail, rather than driving the machine
  into thin air.
\item
  \textbf{Spatial Covariance Does Not Represent Physical Existence:}
  Covariance matrices quantify geometric variance assuming an object
  exists. If a tracked object is removed or occluded, spatial
  extrapolation continues blindly, requiring an explicit existence
  probability field (\(P_{\text{exist}}(t)\)) that decays upon negative
  visual evidence.
\item
  \textbf{State Expiry is Dictated by Physics, Not Software Timeouts:}
  Belief validity windows are bounded by mechanical tolerances and
  disturbance acceleration
  (\(\tau_{\text{expire}} = \sqrt{2 \delta_{\text{tol}} / a_{\text{dist}}}\)).
  Stale state records held in transport buffers are far more hazardous
  than explicit missing frames.
\end{itemize}

\end{fbxSimple}

Chapter~\ref{sec-intent} takes up the synthesis of intent, where
planning architectures and learned policies construct trajectories and
select actions. Those planners must consume the state estimate under a
strict contract. An action is defensible only while the underlying
coordinate frame tree remains synchronized, the sensor epochs match
within an analytical tolerance \(\epsilon\), and the physical error
bound remains within the task budget. When structural components deflect
under rapid deceleration, applying a kinematic transform from one epoch
to a perceptual point cloud captured at another introduces geometric
shear that no learned policy can identify or compensate for at runtime.
By ensuring that every belief retains its evidence age, its kinematic
epoch token, and its validated physical bounds, the memory architecture
provides the policy with an honest boundary between what the machine
directly perceives and where physical ignorance begins.

\chapter{\texorpdfstring{Grounded Intent \&
Horizons\chaptersubtitle{Goals That Expire}}{Grounded Intent \& Horizons}}\label{sec-intent}

\begin{figure}[H]

\centering{

\includegraphics[width=1\linewidth,height=\textheight,keepaspectratio]{index_files/mediabag//Users/VJ/GitHub/MLSysBook-integrate-physical/physical/book/chapters/10-intent/figures/fig_locator.pdf}

}

\caption{\label{fig-locator}\textbf{Systems Organ Focus: Brain.}
High-level reasoning in the Brain formulates semantic objectives and
projects them across the architectural boundary as time-bounded,
expiring geometric intent leases (\(\mathcal{L}_{\text{intent}}\)) to
the deterministic Nervous System.}

\end{figure}%

\emph{What must an autonomous goal carry to prevent a slow reasoning
model from driving a machine into a situation that no longer exists?}

Natural human language commands and high-level mission goals are phrased
in abstract, ungrounded terms that contain zero information regarding
joint torque limits, obstacle clearances, or dynamic reachability. The
intent subsystem bridges this semantic gap, translating task-level
directives into metric spatial affordances and actionable geometric
constraints. However, an intent proposal is never an open-ended
authorization to act. Because external environments change dynamically
while high-level foundation models operate with high variance and
multi-second latency, every intent proposal must be formalized as an
expiring geometric lease (\(\mathcal{L}_{\text{intent}}\)) with explicit
spatial bounding boxes and strict time-to-live contracts that fail-safe
on silence.

\newpage{}

\protect\phantomsection\label{callout-objectiveux2a-10.1}
\begin{fbxSimple}{callout-objective}{Learning Objectives}{}
\phantomsection\label{callout-objective*-10.1}

After studying this chapter, you will be able to:

\begin{enumerate}
\def\labelenumi{\arabic{enumi}.}
\tightlist
\item
  State the three bounds every goal must carry and what each one
  prevents.
\item
  Set a validity window from the dynamics of the scene rather than from
  the model's cadence.
\item
  Filter a goal against reachability before anything plans for it.
\item
  Explain why an intent that expires on its own is safer than one that
  is cancelled.
\end{enumerate}

\emph{Core chapter takeaway.} A goal must carry its own expiry and its
own admissible envelope. An intent that has to be cancelled by something
else is already a hazard, because the thing that would cancel it may be
the thing that failed.

\end{fbxSimple}

\section{Grounded Affordances}\label{sec-intent-10-1}

\lettrine{A}{s illustrated in}\lettrine{t}{he} system architecture
locator (Figure~\ref{fig-locator}), intent generation originates in the
learned reasoning tier (the Brain) and must traverse the architectural
boundary to the real-time execution tier (the Nervous System). In
classical software, artificial intelligence operates entirely ``behind
glass'': computation is digital, virtual, and idempotent. If a language
model hallucinates an incorrect token, produces a malformed JSON
payload, or suffers an out-of-memory exception, the software simply
raises an error, logs a traceback, or retries the query with zero
physical consequence. In physical AI, computation directly commands
moving mass. An end-effector carrying an inertial load possesses kinetic
energy \(E_k = \frac{1}{2} m \|\mathbf{v}\|^2\) and irreversible
physical momentum. When an ungrounded or delayed model output commands
motor current, that energy dissipates into the physical environment
through mechanical deflection, acoustic radiation, friction, or
destructive plastic deformation.

\emph{The systems dilemma.} High-level reasoning pipelines---such as
vision-language models, spatial transformers, and neural task
planners---operate at human cadence, processing camera streams and
performing multimodal token projections at update rates between
\(1\text{ Hz}\) and \(5\text{ Hz}\). Because these models run on complex
operating systems (such as Linux) with deep memory hierarchies and GPU
coprocessors, their execution latencies exhibit high variance: an
inference engine with a nominal \(P_{50}\) latency of \(120\text{ ms}\)
can easily experience \(P_{99}\) latency spikes exceeding
\(500\text{ ms}\) under thermal throttling, PCIe memory bus contention,
or thread scheduling jitter. In contrast, the physical machine operates
in a continuous mechanical domain governed by rigid-body kinematics,
elasticity, and contact dynamics, requiring joint controllers and
trajectory generators that execute deterministically at \(50\text{ Hz}\)
to \(500\text{ Hz}\) (with inner torque and current loops running at
\(1000\text{ Hz}\)). Between these two disparate regimes lies the
boundary where semantic reasoning translates into physical action.

To reconcile these timing regimes without compromising physical safety,
modern physical AI systems adopt a \textbf{Dual-Brain} hardware
architecture (Chapter~\ref{sec-boundary}). As realized in a canonical
heterogeneous dual-brain System-on-Chip (SoC) architecture, computation
is partitioned across an asymmetric silicon boundary:

\begin{enumerate}
\def\labelenumi{\arabic{enumi}.}
\tightlist
\item
  \textbf{The Proposer Brain (Linux MPU)}: An unprivileged application
  processor (e.g., a multi-core MPU running Linux) executes the
  large-scale learned perception models, open-vocabulary grounding
  backends, and high-level reasoning policies. Operating in an
  asynchronous, non-deterministic domain, it formulates candidate
  semantic hypotheses and packages them into structured proposals.
\item
  \textbf{The Permission Referee (Real-Time MCU)}: A deterministic
  microcontroller (e.g., an on-chip ARM Cortex-M or RISC-V core running
  bare metal or a hard real-time RTOS) executes the safety shield,
  trajectory interpolation, and motor drive loops. Operating on
  millisecond hardware deadlines, it validates incoming proposals
  against invariant barrier certificates and hardware timer leases
  before granting actuation authority.
\end{enumerate}

The data structure that crosses this MPU-to-MCU inter-processor
communication (IPC) boundary is the \textbf{intent record}: an explicit,
strictly typed interface contract that converts high-level objectives
into bounded geometric targets that the real-time MCU can safely
evaluate, constrain, and track. On an articulated arm holding a
\(2.0\text{ kg}\) payload, a \(400\text{ ms}\) inference stall on the
Linux MPU leaves the machine without fresh guidance across twenty
consecutive control cycles. If the control architecture treats the
high-level goal as an unbounded, open-ended directive, momentum and
gravity continue to drive mass whether the target remains physically
valid or not.

Open-ended instructions derived from language models, such as ``tidy the
desk'' or ``pick up the mug,'' cannot be passed directly to a physical
controller because they omit the spatial, temporal, and mechanical
boundaries that define safe execution. A motor controller cannot
interpret abstract goals; it accepts setpoints, joint currents, and
trajectory polynomials. If an instruction provides only a semantic
target without geometric coordinates, velocity limits, and contact
thresholds, it functions as an open-ended command with an infinite time
horizon.\sidenote{\footnotesize Classical STRIPS planning formulated goals as static,
  infinite-duration logical predicates. In Physical AI, where operating
  environments evolve outside the perception horizon, goals cannot be
  persistent logical truths; they must be structured as decaying,
  transient geometric leases that expire automatically.} If a human hand
moves the mug while the arm is in flight, the low-level controller has
no perception of the change. Without an explicit bounding contract, the
arm sweeps through the empty workspace, or drives the gripper into the
tabletop at maximum current until an actuator hits its thermal trip
point or an optical encoder registers a stall.

When a machine design relies on an external cancellation message to
abort an invalid motion, it assumes that the subsystem responsible for
detecting failure will remain operational. In physical systems, this
assumption fails precisely when safety is most threatened. If an onboard
vision pipeline crashes due to an out-of-memory error, if an inference
kernel hangs on a direct memory access lock, or if a sensor cable
shears, the MPU node that issued the goal becomes unreachable. If the
real-time MCU requires an active cancel signal to halt execution, a
crash in the reasoning layer leaves the last dispatched goal active
indefinitely. The actuators will continue driving mass toward a target
that may no longer exist, through space that may no longer be clear. A
goal must carry its own expiry, because whatever would have canceled it
may be the exact component that failed.

To enforce physical safety across the Dual-Brain boundary, every intent
packet emitted by the learned MPU must carry three non-negotiable
bounds:

\begin{enumerate}
\def\labelenumi{\arabic{enumi}.}
\tightlist
\item
  \textbf{3D Spatial Tolerance Envelope}: An explicit target pose in
  \(SE(3)\) accompanied by a bounded metric tolerance volume (e.g., an
  oriented covariance ellipsoid), defining both the intended goal pose
  and the allowable positional dispersion within which the low-level
  controller is permitted to maneuver and declare goal attainment.
\item
  \textbf{Expiring Temporal Lease (TTL)}: An explicit time-to-live
  \(\tau\) measured in milliseconds. If an analytical model operating at
  \(4\text{ Hz}\) issues a lease of \(300\text{ ms}\), the real-time MCU
  tracks the target while the lease is fresh; if no renewal packet
  arrives before the deadline passes, a local hardware countdown timer
  automatically revokes motion authority and transitions the actuators
  to a controlled hold.
\item
  \textbf{Admissible Command and Force Envelope}: Explicit dynamic
  limits that cap the peak velocity (\(v_{\max}\)), joint acceleration
  (\(a_{\max}\)), and contact wrench (\(F_{\max}\)) the machine may
  exert while pursuing the target, ensuring that an end-effector cannot
  apply more than \(12\text{ N}\) of normal force to a surface
  regardless of tracking error.
\end{enumerate}

Establishing these protections requires a systematic interface pipeline
between the reasoning layer and the nervous system. The system must
resolve unstructured language and image features into unambiguous
spatial coordinates, attach strict self-expiring leases to every
command, filter candidate goals for kinematic and dynamic reachability
before admitting them to the trajectory queue, and enforce a formal
intent record schema on the inter-processor communication bus. When
every goal carries its own spatial boundary, its own temporal deadline,
and its own force ceiling, the physical machine remains stable and
bounded even when the learned brain stalls, crashes, or falls silent.

\protect\phantomsection\label{callout-definitionux2a-10.2}
\begin{fbxSimple}{callout-definition}{Definition:}{Expiring Intent Lease}
\phantomsection\label{callout-definition*-10.2}

An \textbf{expiring intent lease} (\(\mathcal{L}_{\text{intent}}\)) is a
time-bounded cryptographic and spatial contract
(\(t_{\text{exp}} = t_{\text{src}} + \tau\)) issued by an upstream
reasoning tier authorizing downstream trajectory tracking within a
finite duration \(\tau\), derived from maximum workspace drift velocity
and sensor noise; upon unrenewed expiration, motion authority is
automatically revoked at the local real-time execution tier without
requiring an external abort command.

\end{fbxSimple}

These bounding invariants apply directly across all Three Physical
Machine Archetypes:

\begin{itemize}
\tightlist
\item
  \textbf{Class 1 (Mass-Dominated Mobility Systems):} Autonomous mobile
  robots (AMRs) and quadruped platforms navigating crowded environments
  must bind intent with a spatial clearance corridor and a decaying
  temporal lease. If an obstacle enters the corridor while the
  high-level planner is stalled, the expiration of the lease forces
  immediate friction braking before the remaining spatial buffer
  collapses.
\item
  \textbf{Class 2 (Contact-Dominated Manipulation Systems):} Articulated
  industrial arms and multi-fingered grippers executing precision
  assembly must bind intent with an oriented covariance ellipsoid and a
  strict normal wrench cap (\(F_{\max}\)). If an insertion target
  shifts, the force ceiling prevents the gripper from driving tooling
  into rigid fixtures past mechanical yield limits.
\item
  \textbf{Class 3 (Flow-Dominated Process and Energy Systems):} Battery
  energy storage cooling loops and grid power inverters must bind intent
  with thermal and phase rate limits. If a learned optimizer hangs, the
  expiring lease automatically reverts coolant pumps and inverter gate
  drivers to deterministic baseline setpoints before thermal capacitance
  is exhausted.
\end{itemize}

\protect\phantomsection\label{callout-ruleux2a-10.3}
\begin{fbxSimple}{callout-rule}{Architect’s Rule of Thumb:}{Silence Is a State Transition, Not an Absence}
\phantomsection\label{callout-rule*-10.3}

In digital software, the absence of network packets leaves state
unchanged; in physical AI, silence is an active physical hazard. Every
intent lease must be designed such that when upstream communication
ceases, the downstream real-time controller autonomously executes a
deterministic safe state through local hardware timers without requiring
an external cancellation signal.

\end{fbxSimple}

\section{What Intent Hands Over}\label{sec-intent-10-2}

Every intent packet delivered across the architectural boundary must
establish a spatial contract before the nervous system accepts it for
execution. In physical manipulation, handing a downstream trajectory
generator an isolated coordinate is an incomplete instruction. The
intent record must specify both a target pose in \(SE(3)\), comprising a
position vector \(\mathbf{p} \in \mathbb{R}^3\) and an orientation
matrix \(\mathbf{R} \in SO(3)\), and an explicit tolerance volume that
defines acceptable terminal error. The \textbf{3D spatial tolerance
envelope} is a bounded metric volume in \(SE(3)\) (parameterized by
target pose \(\mathbf{p}_0 \in \mathbb{R}^3\) and oriented covariance
matrix \(\mathbf{\Sigma} \in \mathbb{R}^{3 \times 3}\)) defining the
allowable positional dispersion within which low-level controllers
maneuver and declare goal attainment. Without a defined tolerance
volume, a controller cannot determine when a motion has succeeded.
Consider an articulated arm commanded to align a tactile probe with a
contact pad on a circuit board. If the high-level policy specifies a
target pose but omits the tolerance volume, the low-level tracking loop
attempts to drive Euclidean tracking error to absolute zero. Because
sensor noise, link compliance, and kinematic calibration errors prevent
exact mathematical convergence, the integrator in the feedback loop
winds up, increasing motor currents until an overcurrent protection
circuit trips or the probe cracks the substrate. By supplying an
explicit tolerance geometry, such as an ellipsoid with semi-axes of
\(1.5\text{ mm}\) translationally and \(0.05\text{ rad}\) rotationally,
the intent packet provides the nervous system with a computable
termination condition. The local controller can declare the goal
reached, damp out residual momentum, and signal the next execution
stage.

Spatial validity in the physical world is transient. Because the
environment changes continuously, a physical goal must carry an absolute
expiration timestamp rather than relying on an asynchronous cancellation
message from the reasoning layer. When a vision-language model generates
a spatial target from a camera frame captured at \(t = 0\text{ ms}\),
that coordinate is valid only for the duration of the operational window
that justified it. If the inference pipeline operating at
\(5\text{ Hz}\) on the Linux MPU hangs on a tensor allocation or drops a
thread, the physical scene continues to evolve. A human operator reaches
into the workcell, or a secondary mechanism shifts the target object by
\(50\text{ mm}\). If the real-time MCU depends on an active abort
command to stop, the silence of a crashed vision process allows the
stale goal to persist indefinitely. The local tracking loop continues
driving mass toward a coordinate that is no longer unoccupied. Attaching
a temporal lease, such as a validity duration of \(250\text{ ms}\) for
an analytical model evaluated at \(4\text{ Hz}\), converts liveness into
a local invariant.\sidenote{\footnotesize Machinery safety standards (ISO 13849-1 PL
  d/e, IEC 61508 SIL 2/3) require independent hardware watchdogs that
  autonomously command safe states (Safe Torque Off, IEC 60204-1 Stop
  Category 1) upon loss of cyclic communication, without requiring
  active cooperation from non-deterministic host processors.} The MCU
evaluates the current hardware clock against the goal expiration
timestamp on every \(1000\text{ Hz}\) control cycle. If the deadline
passes without the receipt of an authenticated renewal packet over the
shared SRAM mailbox, the local controller automatically drops the active
goal and transitions the actuators into a controlled hold. This property
establishes the foundational architectural invariant of \textbf{failing
safe on silence}: the complete absence of upstream communication
passively transitions the physical plant into a deterministic safe state
via hardware countdown timers.

A target bounded in space and time remains hazardous if the controller
is permitted unlimited physical effort to achieve it. The intent packet
must therefore declare an admissible command envelope that caps the peak
contact forces, joint velocities, and actuator power that the planner
may utilize during execution. When a model computes a target pose that
lies slightly beneath a physical surface due to stereo disparity error,
an unconstrained position controller treats the discrepancy as tracking
error and commands maximum restoring torque. In a tactile search task
where a multi-jointed finger slides across an unknown object,
unconstrained execution generates excessive normal forces that can bend
linkages or shear tactile skins. By embedding explicit dynamic limits in
the goal record, such as a maximum normal contact force of
\(6.0\text{ N}\), a linear velocity cap of \(0.10\text{ m/s}\), and a
joint power limit of \(25\text{ W}\), the high-level policy sets the
physical budget for the action. The low-level impedance controller
projects these bounds directly into its torque saturation limits,
ensuring that mechanical interactions remain compliant regardless of
visual estimation errors.

When any of these three bounds is omitted, an intermittent communication
dropout or perception stall degrades the control system into an
open-loop runaway hazard. Consider an analytical example of a mobile
manipulator moving at \(1.2\text{ m/s}\) toward an unbounded target pose
located \(3.0\text{ m}\) ahead. At an operating speed of
\(1.2\text{ m/s}\), the base covers \(120\text{ mm}\) every
\(100\text{ ms}\). If an Ethernet link drops packets or an inference
accelerator stalls for \(400\text{ ms}\), an unbounded goal leaves the
motor drives executing the last received velocity profile, carrying the
machine through \(480\text{ mm}\) of unmonitored displacement. If an
obstacle enters the corridor during that interval, a collision occurs
before the host software can re-establish communication. With a bounded
intent packet carrying a \(150\text{ ms}\) expiration lease and a
commanded deceleration bound of \(3.0\text{ m/s}^2\), the failure
sequence proceeds differently. The communication loss at
\(t = 0\text{ ms}\) allows the machine to proceed for at most
\(150\text{ ms}\), covering \(180\text{ mm}\), at which point the MCU
detects lease expiration. The local reflex immediately commands maximum
deceleration, bringing the kinetic energy of the machine to zero in an
additional \(400\text{ ms}\) and \(240\text{ mm}\) of stopping distance.
Total unguided travel is strictly capped at \(420\text{ mm}\),
preventing kinetic runaway without requiring assistance from the failed
upstream node.

\protect\phantomsection\label{callout-mathux2a-10.4}
\begin{fbxSimple}{callout-math}{Fermi Math:}{Lease Duration vs. Collision Exposure Tradeoff Across the Three Archetypes}
\phantomsection\label{callout-math*-10.4}

To understand how an expiring intent lease
(\(\mathcal{L}_{\text{intent}}\)) bounds physical risk across diverse
machines, consider the physical consequences of an unserviced
\(200\text{ ms}\) upstream compute stall (\(\Delta t = 0.200\text{ s}\))
when comparing an unbounded goal (\(t_{\text{exp}} = \infty\)) against a
physics-derived lease
(\(\tau \le \frac{r_{\text{tol}} - \sigma_{\text{sensor}}}{v_{\text{drift}}}\)):

\begin{enumerate}
\def\labelenumi{\arabic{enumi}.}
\tightlist
\item
  \textbf{Class 1 (Mass-Dominated Mobility System):} A \(150\text{ kg}\)
  autonomous mobile robot (AMR) navigates a shared factory aisle at
  cruising velocity \(v = 1.5\text{ m/s}\) with maximum friction braking
  \(a_{\max} = 3.0\text{ m/s}^2\). The task tolerance radius is
  \(r_{\text{tol}} = 30\text{ mm}\), sensor noise is
  \(\sigma_{\text{sensor}} = 6\text{ mm}\), and dynamic obstacle drift
  velocity is \(v_{\text{drift}} = 0.40\text{ m/s}\).

  \begin{itemize}
  \tightlist
  \item
    \emph{Safe Lease Duration:}
    \(\tau = \frac{30\text{ mm} - 6\text{ mm}}{400\text{ mm/s}} = 0.060\text{ s} = 60\text{ ms}\).
  \item
    \emph{Bounded vs.~Unbounded Stopping Envelope:} Under a
    \(60\text{ ms}\) expiring lease, the chassis travels
    \(d_{\text{react}} = v \tau = (1.5\text{ m/s})(0.060\text{ s}) = 0.090\text{ m} = 90\text{ mm}\)
    before the hardware timer triggers maximum deceleration, bringing
    the robot to rest in
    \(d_{\text{brake}} = \frac{v^2}{2 a_{\max}} = \frac{(1.5)^2}{2(3.0)} = 0.375\text{ m} = 375\text{ mm}\)
    (total stopping distance \(d_{\text{stop}} = 465\text{ mm}\)). Under
    an unbounded \(200\text{ ms}\) stall, unguided travel expands to
    \(d_{\text{react}} = (1.5)(0.200) = 0.300\text{ m}\), increasing
    total stopping distance to \(675\text{ mm}\)---an unmonitored
    spatial overrun of \(210\text{ mm}\) that breaches pedestrian
    separation corridors.
  \end{itemize}
\item
  \textbf{Class 2 (Contact-Dominated Manipulation System):} A 7-DoF
  manipulator tooling tip approaches a rigid aluminum assembly fixture
  (\(k = 4.0 \times 10^5\text{ N/m}\)) at approach velocity
  \(v = 0.30\text{ m/s}\).

  \begin{itemize}
  \tightlist
  \item
    \emph{Rate of Force Rise:}
    \(\frac{dF}{dt} = k v = (4.0 \times 10^5\text{ N/m})(0.30\text{ m/s}) = 1.2 \times 10^5\text{ N/s} = 120\text{ N/ms}\).
  \item
    \emph{Bounded vs.~Unbounded Contact Force:} If the intent record
    specifies a \(40\text{ ms}\) lease with a wrench ceiling
    \(F_{\max} = 15\text{ N}\), the local current loop clamps torque
    within \(2\text{ ms}\) of contact, bounding peak force to
    \(F_{\text{peak}} \le 18\text{ N}\). Under an unbounded goal during
    a \(200\text{ ms}\) inference stall, unbraked motion generates an
    uncorrected penetration of
    \(\Delta x = v \Delta t = (0.30\text{ m/s})(0.200\text{ s}) = 0.060\text{ m} = 60\text{ mm}\),
    driving theoretical contact force to
    \(F = k \Delta x = 24{,}000\text{ N}\)---shattering the harmonic
    drive gear teeth and buckling the wrist link.
  \end{itemize}
\item
  \textbf{Class 3 (Flow-Dominated Process and Energy System):} A
  containerized battery energy storage system (BESS) liquid cooling loop
  with lumped thermal capacitance \(C_{\text{th}} = 900\text{ J/K}\)
  dissipates heat from cells generating
  \(P_{\text{in}} = 3600\text{ W}\) during rapid megawatt-level
  charging.

  \begin{itemize}
  \tightlist
  \item
    \emph{Thermal Runaway Rate:} With coolant valve closed
    (\(P_{\text{out}} = 0\text{ W}\)),
    \(\frac{dT}{dt} = \frac{P_{\text{in}}}{C_{\text{th}}} = \frac{3600\text{ W}}{900\text{ J/K}} = 4.0\text{ K/s}\).
  \item
    \emph{Safe Lease Duration:} If the cell chemistry permits a maximum
    temperature overshoot of \(\Delta T_{\max} = 2.0\text{ K}\) before
    cell degradation accelerates, the maximum allowable pump valve
    command lease is
    \(\tau = \frac{\Delta T_{\max}}{dT/dt} = \frac{2.0\text{ K}}{4.0\text{ K/s}} = 0.500\text{ s} = 500\text{ ms}\).
    If an ungrounded thermal optimization policy stalls for
    \(10.0\text{ s}\), the temperature rises by
    \(\Delta T = (4.0\text{ K/s})(10.0\text{ s}) = 40.0\text{ K}\),
    triggering irreversible thermal runaway and flammable electrolyte
    venting.
  \end{itemize}
\end{enumerate}

\end{fbxSimple}

These three bounds formalize the division of authority between
high-level reasoning and low-level reflexes. High-level reasoning on the
Linux MPU operates over unstructured data, long time horizons, and
semantic abstractions. It possesses the contextual breadth required to
define \textbf{what is permissible}, establishing which spatial volume
contains the target, what time window makes the action valid, and what
maximum forces are safe for the surrounding environment. Low-level
reflexes on the real-time MCU operate deterministically on millisecond
deadlines, measuring joint encoders and strain gauges. Reflexes possess
the high-frequency state feedback required to enforce \textbf{how
execution occurs}, synthesizing smooth trajectories, resolving kinematic
redundancies, managing actuator thermal budgets, and executing immediate
stops upon lease expiry. High-level models never issue raw torque or
voltage commands to the motors, because an unverified neural network
cannot guarantee real-time stability. Low-level controllers never
generate autonomous goals, because a feedback loop lacks semantic
context. The bounded intent record serves as the strictly typed contract
between them, ensuring that high-level intelligence governs the envelope
of operation while low-level determinism guarantees physical
containment. As illustrated on a physical 7-DoF manipulator executing
everyday manipulation tasks in
Figure~\ref{fig-10-real-grounded-affordance-franka}, high-level intent
translates language goals into grounded 3D affordance and constraint
value maps that govern the end-effector's spatial target envelopes
across sequential execution phases. Low-level tracking loops then
consume these spatial objective manifolds to synthesize collision-free
paths within granted time slices.

\begin{figure}[t]

\centering{

\includegraphics[width=0.95\linewidth,height=\textheight,keepaspectratio]{/Users/VJ/GitHub/MLSysBook-integrate-physical/physical/book/chapters/10-intent/figures/fig10_real_grounded_affordance_franka.jpg}

}

\caption{\label{fig-10-real-grounded-affordance-franka}\textbf{Real-Robot
3D Affordance Grounding and Value-Map Trajectory Execution.} A canonical
7-DoF torque-controlled manipulator executing diverse open-vocabulary
manipulation tasks guided by composed 3D spatial value maps in
real-world lab environments (\citeproc{ref-huang2023voxposer}{Huang et
al. 2023}). Across everyday language instructions---including sweeping
paper trash into a dustpan, closing cabinet drawers, turning on lamps,
opening capped vitamin bottles, and multi-stage pick-and-place
sequences---an upstream vision-language model translates semantic goals
into 3D voxelized affordance and constraint costmaps directly in the
robot's physical workspace. The low-level trajectory generator and
impedance controller evaluate these 3D spatial bounds to synthesize
collision-free setpoint paths across sequential time intervals
(\(t=1, 2, \dots\)), bounding the robot's kinetic authority within
task-specific geometric envelopes without requiring task-specific
reinforcement training. (Source: Huang et al., 2023
(\citeproc{ref-huang2023voxposer}{Huang et al. 2023}), CC BY 4.0).}

\end{figure}%

Grounding an intent packet into these three bounds requires a formal
transformation from ambiguous sensory inputs into structured, verifiable
proposals. When a system processes a high-level directive or a sequence
of visual tokens, the initial proposal is rarely a single, unambiguous
set of coordinates. It begins as an uncertain geometric hypothesis that
must carry its own spatial ambiguity and its own temporal limits before
it can be admitted to the execution pipeline.

\section{From Request to Expiring Geometric
Proposal}\label{sec-intent-10-3}

A natural language instruction becomes physically dangerous at the
moment an unverified semantic hypothesis acquires authority over motion.
When an operator or high-level supervisor issues a semantic directive
such as ``grasp the brass valve,'' an onboard vision-language model
parses the character stream, projects image features into a shared
embedding space, and selects a candidate visual bounding region. From
those pixels, the network extracts an estimated pose in \(SE(3)\) and
writes coordinate values into a memory buffer. At that exact boundary,
an ungrounded string of tokens becomes a physical destination that
commands mass, velocity, and actuator currents.\sidenote{\footnotesize The challenge
  of connecting syntactic tokens to physical entities is the
  \emph{symbol grounding problem} (Harnad, 1990). Open-vocabulary VLMs
  project tokens into metric coordinates; however, non-Lipschitz latent
  manifolds mean small sensory shifts can trigger discontinuous
  coordinate jumps.} If the visual backbone associates the language
token with an adjacent fixture or an unlabelled calibration target, the
resulting coordinate vector \(\mathbf{p} \in \mathbb{R}^3\) enters the
trajectory pipeline as an authoritative physical goal. Because a
low-level motor drive measures only encoder pulses, winding
temperatures, and phase currents, it cannot evaluate whether the
coordinate it tracks corresponds to the intended component or to an
empty tabletop.\sidenote{\footnotesize Vision-language backbones lack global Lipschitz
  continuity: small pixel variations \(\delta \mathbf{I}\) can induce
  topological embedding jumps
  (\(\|\Delta \phi\| \gg L\|\delta \mathbf{I}\|\)). Bounding ellipsoids
  (\(\mathbf{\Sigma}\)) and refusal gates
  (\texttt{ERR\_GROUNDING\_AMBIGUITY}) prevent step-discontinuous
  trajectory jumps to real-time motor drives.} The grounding process
must therefore produce an explicit geometric proposal that packages its
own spatial uncertainty, its selection evidence, and its temporal
expiration before any motion planner accepts it.

To understand how open-vocabulary language directives translate into
actionable geometric targets, modern physical AI systems rely on three
foundational grounding paradigms:

\begin{enumerate}
\def\labelenumi{\arabic{enumi}.}
\item
  \textbf{Affordance-Grounded Skill Selection (SayCan)}: Pioneered by
  Ahn et al. (\citeproc{ref-ahn2022saycan}{{Ahn et al.} 2023}), SayCan
  grounds large language model reasoning in physical feasibility by
  explicitly factoring task execution probabilities into semantic
  relevance and affordance probability:\sidenote{\footnotesize In SayCan, affordance
    probability
    \(P_{\text{affordance}}(\text{skill} \mid \mathbf{s}_t)\) is
    parameterized by policy value functions \(V(s) = \max_a Q(s, a)\).
    Weighting semantic relevance with physical feasibility prevents
    models from issuing commands that are semantically plausible but
    kinematically unreachable.}
  \[P(\text{skill}_k \mid \text{instruction}, \mathbf{s}_t) \propto P_{\text{LLM}}(\text{skill}_k \mid \text{instruction}) \times P_{\text{affordance}}(\text{skill}_k \mid \mathbf{s}_t)\]
  where \(P_{\text{LLM}}(\text{skill}_k \mid \text{instruction})\)
  represents the semantic prior that skill \(k\) makes logical progress
  toward the language goal, and
  \(P_{\text{affordance}}(\text{skill}_k \mid \mathbf{s}_t)\) is a
  learned value function (such as a Q-value
  \(Q(\mathbf{s}_t, \text{skill}_k)\) or visual success classifier)
  representing whether skill \(k\) is physically achievable from current
  state \(\mathbf{s}_t\). If an LLM proposes ``pick up the sponge'' when
  the sponge is trapped inside a closed drawer,
  \(P_{\text{affordance}}\) drops to zero, preventing the system from
  dispatching an impossible action and compelling the reasoner to select
  an intermediate pre-condition skill (``open the drawer'').
\item
  \textbf{3D Value-Map Affordance Fields (VoxPoser)}: Rather than
  selecting from discrete, pre-trained skill policies, VoxPoser
  (\citeproc{ref-huang2023voxposer}{Huang et al. 2023}) leverages
  vision-language models and LLM-generated Python code to synthesize
  continuous 3D spatial affordance and constraint value maps directly in
  the robot's physical workspace. By querying 2D visual grounding
  backends (e.g., OWL-ViT, SAM, Grounding DINO) across multi-view RGB-D
  streams, the system back-projects visual activations into a 3D voxel
  grid \(\mathcal{V} \subset \mathbb{R}^3\). The LLM writes code that
  computes 3D potential fields: an \textbf{attraction value map}
  \(\mathbf{V}_{\text{attr}}(\mathbf{x})\) that pulls the end-effector
  toward target contact affordances (e.g., grasping handles or rim
  geometries), and a \textbf{repulsion constraint map}
  \(\mathbf{V}_{\text{rep}}(\mathbf{x})\) that penalizes trajectories
  passing near obstacles. A downstream model-predictive controller (MPC)
  or trajectory optimizer then computes dynamic paths directly over this
  composed 3D cost manifold without requiring task-specific
  reinforcement training.
\item
  \textbf{Language Model Programs for Embodied Control (Code as
  Policies)}: Introduced by Liang et al.
  (\citeproc{ref-liang2023code}{{Liang et al.} 2023}), Code as Policies
  (CaP) prompts large language models to generate executable Python
  programs that compose perception APIs, geometric coordinate
  transformations, and parameterized motion primitives. Instead of
  outputting single static setpoints, the LLM writes algorithmic control
  structures containing spatial arithmetic, conditional branching
  (\texttt{if\ is\_grasped():\ ...}), and reactive loops that query
  real-time sensor feedback.
\end{enumerate}

Formally, \textbf{open-vocabulary affordance grounding} is the
computational process of projecting unconstrained natural language
instructions and multimodal sensor streams into verified 3D spatial
target manifolds and interaction affordance fields (such as SayCan
feasibility factors, VoxPoser value maps, or Code as Policies
primitives), parameterized with physical SI units and coordinate frame
attachments.

Across all three paradigms, transforming raw visual observations and
language tokens into executable intent records requires a multi-stage
perceptual filtering pipeline, as demonstrated on the physical
manipulation platform in Figure~\ref{fig-10-real-open-vocab-grounding}.
The system must transition from raw visual observations to dense
candidate grasps (via geometric grasp generators like AnyGrasp), isolate
the semantic referent through language-conditioned instance segmentation
masks (e.g., Lang-SAM combining Grounding DINO with Segment Anything),
prune ungrounded proposals outside the segmented geometric envelope, and
extract an oriented \(SE(3)\) target pose with defined spatial tolerance
bounds.

\begin{figure}[t]

\centering{

\includegraphics[width=0.95\linewidth,height=\textheight,keepaspectratio]{/Users/VJ/GitHub/MLSysBook-integrate-physical/physical/book/chapters/10-intent/figures/fig10_real_open_vocab_grounding.jpg}

}

\caption{\label{fig-10-real-open-vocab-grounding}\textbf{Real-World
Open-Vocabulary Language Grounding and Geometric Grasp Proposal
Pipeline.} On a physical robotic manipulation workstation,
open-vocabulary semantic instructions are grounded into executable
spatial targets through a multi-stage perceptual filtering pipeline
(\citeproc{ref-bahl2024okrobot}{Bahl et al. 2024}). (a) Robot
point-of-view RGB camera observation of a cluttered tabletop workspace
containing everyday tools and objects. (b) Dense geometric grasp
generation proposing unconditioned candidate 6D end-effector poses
across the scene using AnyGrasp. (c) Language-conditioned instance
segmentation using Lang-SAM (combining Grounding DINO with Segment
Anything) to isolate the target object mask from an open-vocabulary
query. (d) Bounding mask spatial filtering, pruning all grasp candidates
located outside the segmented geometric envelope. (e) Final selected
grasp pose with 3D orientation frame (\(SE(3)\) target) handed over to
the motion planner with explicit spatial tolerance bounds and an
expiration lease.}

\end{figure}%

To prevent an unverified neural network output from masquerading as a
deterministic setpoint, the grounded proposal must satisfy a formal
interface contract. The proposal cannot consist solely of an isolated
Cartesian pose. It must carry the referent identity, the perceptual
evidence that selected it, a spatial tolerance volume, and an expiration
lease. The referent identity specifies the semantic class and instance
identifier assigned by the upstream model. The selection evidence
records the bounding box coordinates, cross-attention weights, and
posterior class probabilities that justified the match. The spatial
tolerance volume, parameterized by an oriented covariance ellipsoid,
defines the physical region within which the referent is estimated to
lie. The expiration lease establishes an absolute deadline after which
the target coordinate is automatically discarded. In a physical
workcell, an object can shift, roll, or be occluded milliseconds after a
camera frame is exposed. A target coordinate detached from its
perceptual evidence and its expiration lease becomes an open-loop
hazard, directing mechanical power toward coordinates that no longer
contain the target.

When a vision-language model processes an ambiguous request in an
environment containing multiple candidate objects, its mathematical
objective encourages it to output an argmax prediction. In digital
applications, choosing the top-ranked candidate when two options exhibit
near-equal probabilities of \(0.51\) and \(0.49\) produces a harmless
user-interface recommendation. In a physical machine, breaking an
evidential tie silently produces an arbitrary trajectory. If two
identical mechanical components sit on an assembly bench separated by
\(180\text{ mm}\), an ungrounded model that guesses between them commits
the manipulator to a ballistic approach trajectory that passes through
the clearance volume of the adjacent part. Ambiguity must therefore be
treated as an explicit, first-class output state. When the posterior
entropy across \(K\) candidate referents,
\[H(\mathbf{z}) = -\sum_{i=1}^K p_i \log_2 p_i\] exceeds an
admissibility threshold \(H_{\text{thresh}}\), or when the top two
candidate softmax probabilities satisfy
\(|p_1 - p_2| \le \delta_{\text{margin}}\), or when the spatial
tolerance volumes of competing hypotheses overlap in \(SE(3)\), the
required system behavior is an explicit \textbf{grounding refusal}
(\texttt{ERR\_GROUNDING\_AMBIGUITY}) rather than an optimistic guess.
The grounding layer rejects the command, signals an unresolved reference
to the supervisor, and holds the actuators in a safe state until an
unambiguous directive or an active disambiguation scan resolves the
distribution.

Formally, \textbf{open-vocabulary grounding ambiguity} is the condition
where multimodal feature projections yield multi-peak posterior
distributions, overlapping geometric hypotheses, or unresolvable
evidential ties across candidate physical referents
(\(H(\mathbf{z}) > H_{\text{thresh}}\)), requiring an explicit grounding
refusal (\texttt{ERR\_GROUNDING\_AMBIGUITY}) rather than optimistic
argmax tie-breaking.

\protect\phantomsection\label{callout-ruleux2a-10.5}
\begin{fbxSimple}{callout-rule}{Architect’s Rule of Thumb:}{Ambiguity Demands Refusal, Never Argmax Guessing}
\phantomsection\label{callout-rule*-10.5}

In digital web applications, breaking an evidential tie between two
items via argmax is harmless; in physical AI, breaking a tie between two
spatial coordinates commands kinetic energy through the clearance volume
of an unselected object. When posterior entropy exceeds threshold, the
only safe physical action is an explicit grounding refusal
(\texttt{ERR\_GROUNDING\_AMBIGUITY}).

\end{fbxSimple}

Consider an analytical example where an articulated manipulator receives
an instruction to align its gripper with a target component located at a
nominal standoff distance of \(1.20\text{ m}\) along the camera optical
axis. The vision-language model extracts a nominal target pose
\(\mathbf{p}_0 = [0.450, -0.120, 0.220]^T\text{ m}\) relative to the
robot base frame, accompanied by an empirical localization covariance
matrix \(\mathbf{\Sigma} \in \mathbb{R}^{3 \times 3}\). Due to the
geometric disparity between sensor pixel resolution in the focal plane
and range uncertainty along the line of sight, the diagonal elements of
the covariance matrix in the principal sensor frame are
\(\sigma_x^2 = (4.0\text{ mm})^2\), \(\sigma_y^2 = (6.0\text{ mm})^2\),
and \(\sigma_z^2 = (15.0\text{ mm})^2\). For a three-dimensional
Gaussian hypothesis, the \(95\%\) confidence region is bounded by the
tolerance ellipsoid
\((\mathbf{p} - \mathbf{p}_0)^T \mathbf{\Sigma}^{-1} (\mathbf{p} - \mathbf{p}_0) \le \chi^2_{3, 0.95}\),
where the critical value from the chi-square distribution is
\(\chi^2_{3, 0.95} \approx 7.815\). As illustrated in
Figure~\ref{fig-10-intent-lease-envelope} (Panel A), evaluating the
principal semi-axes yields
\(a_x = \sqrt{7.815} \times 4.0\text{ mm} \approx 11.2\text{ mm}\),
\(a_y = \sqrt{7.815} \times 6.0\text{ mm} \approx 16.8\text{ mm}\), and
\(a_z = \sqrt{7.815} \times 15.0\text{ mm} \approx 41.9\text{ mm}\),
bounding a spatial tolerance volume of
\(V = \frac{4}{3} \pi a_x a_y a_z \approx 3.30 \times 10^{-5}\text{ m}^3\),
or \(33.0\text{ cm}^3\). If the physical gripper mechanism allows a
maximum mechanical positioning error of \(12.0\text{ mm}\) along the
approach axis, the unconstrained spatial uncertainty along the depth
axis (\(41.9\text{ mm}\)) exceeds the physical capture envelope. A
downstream planner receiving only the nominal point \(\mathbf{p}_0\)
attempts a precision grasp and collides with the component edges; a
planner receiving the oriented ellipsoid can either align its approach
vector with the narrow \(11.2\text{ mm}\) axis or reject the proposal as
ungrounded.

\begin{figure}[t]

\centering{

\includegraphics[width=1\linewidth,height=\textheight,keepaspectratio]{index_files/mediabag//Users/VJ/GitHub/MLSysBook-integrate-physical/physical/book/chapters/10-intent/figures/fig10_intent_lease_envelope.pdf}

}

\caption{\label{fig-10-intent-lease-envelope}\textbf{The Intent Lease
Contract: Spatial Tolerance Volume and Temporal Expiration Timelines.}
(Panel A) 3D spatial grounding geometry showing nominal target pose
\(\mathbf{p}_0 \in \mathbb{R}^3\), the oriented localization covariance
ellipsoid
\((\mathbf{p} - \mathbf{p}_0)^T \mathbf{\Sigma}^{-1} (\mathbf{p} - \mathbf{p}_0) \le \chi^2_{3, 0.95}\)
with principal semi-axes (\(a_x=11.2\text{ mm}\),
\(a_y=16.8\text{ mm}\), \(a_z=41.9\text{ mm}\)), the approach velocity
vector (\(v_{\text{approach}} \le 0.25\text{ m/s}\)), and the admissible
wrench envelope (\(F_{\max} \le 12\text{ N}\)). When physical scene
drift displaces the object past the tolerance boundary, a spatial fault
is registered. (Panel B) Multi-rate execution timelines contrasting:
(Case 1) nominal continuous lease renewal stream maintaining
\(C^2\)-smooth quintic spline tracking; (Case 2) silent host reasoner
hang at \(t=80\text{ ms}\), where the hardware TTL timer expires at
\(t_{\text{exp}}=150\text{ ms}\) and triggers autonomous local
controlled deceleration (\(t_{\text{decel}}=112\text{ ms}\),
\(d_{\text{stop}}=28\text{ mm}\)) to an active hold without requiring an
upstream abort message; and (Case 3) an unbounded infinite-horizon goal
anti-pattern that drives unmonitored mass into an open-loop impact at
\(1.4\text{ m/s}\) (\(F_{\text{peak}}=1420\text{ N}\)).}

\end{figure}%

A noisy detection with a large covariance volume is easily flagged and
rejected by spatial filters. The failure that defeats downstream safety
systems is a confident grounding to the wrong physical entity. If a
vision-language model evaluates an image containing an exposed glass
beaker adjacent to a machined steel housing and assigns a \(99.2\%\)
confidence score to the glass beaker when instructed to ``secure the
housing,'' the model emits a sharp, low-variance geometric proposal
centered on the glass. The trajectory planner receives the coordinate,
computes an inverse kinematics solution, verifies joint limits, and
confirms that the path clears all obstacles in the static environment
map. The low-level safety enforcer evaluates the trajectory against
dynamic velocity caps and finds every acceleration within admissible
boundaries. Because the motion is kinematically smooth and geometrically
valid, the machine executes the trajectory without a single software
exception, driving the end effector into the glass with the full
clamping force intended for steel.

This failure reveals the boundary between what a downstream safety
enforcer can verify and what it must accept as an unvalidated premise. A
deterministic safety enforcer checks coordinate bounds, kinematic
reachability, joint torque limits, and geometric collisions against
known occupancy grids. It can never determine whether a coordinate
corresponds to the intended semantic object. Semantic classification
depends on high-dimensional feature representations that cannot be
formally verified against real-time safety invariants. The mechanics of
multimodal feature extraction, cross-attention alignment, and
open-vocabulary phrase grounding belong to the computer vision and
language literature (\citeproc{ref-kamath2021mdetr}{Kamath et al. 2021};
\citeproc{ref-gu2022open}{Gu et al. 2022};
\citeproc{ref-radford2021learning}{{Radford et al.} 2021};
\citeproc{ref-liu2023grounding}{{Liu et al.} 2023}). A physical AI
architecture consumes these models without relying on their internal
activations for safety. The responsibility of the systems engineering
contract is to govern the boundary where these models meet the physical
machine, ensuring that every semantic hypothesis is bounded by an
explicit tolerance volume, gated by a grounding refusal on ambiguity,
and constrained by a strict temporal lease before it ever acquires the
authority to move.

To evaluate how different open-vocabulary perceptual frontends interface
with deterministic systems contracts,
Table~\ref{tbl-10-grounding-paradigms} compares the dominant
architectural categories of language grounding against their geometric
output representations, nominal localization uncertainties, compute
latency budgets, and deterministic verification readiness.

\begin{longtable}[]{@{}
  >{\raggedright\arraybackslash}p{(\linewidth - 10\tabcolsep) * \real{0.1667}}
  >{\raggedright\arraybackslash}p{(\linewidth - 10\tabcolsep) * \real{0.1667}}
  >{\raggedright\arraybackslash}p{(\linewidth - 10\tabcolsep) * \real{0.1667}}
  >{\raggedright\arraybackslash}p{(\linewidth - 10\tabcolsep) * \real{0.1667}}
  >{\raggedright\arraybackslash}p{(\linewidth - 10\tabcolsep) * \real{0.1667}}
  >{\raggedright\arraybackslash}p{(\linewidth - 10\tabcolsep) * \real{0.1667}}@{}}
\caption{\textbf{Open-Vocabulary Grounding Model Categories vs.~Spatial
Precision and Verification Readiness.} Comparison of
language-to-geometry perception paradigms across spatial output
representations, typical 3D localization uncertainty bounds
(\(\sigma_{xyz}\)), computational resource envelopes, deterministic
safety verifiability, and systems-level failure mitigations. Hardware
systems cannot verify internal model activations, requiring geometric
bounding contracts to govern the interface between learned
vision-language frontends and deterministic motion
planners.}\label{tbl-10-grounding-paradigms}\tabularnewline
\toprule\noalign{}
\begin{minipage}[b]{\linewidth}\raggedright
Grounding Architecture \& Paradigm
\end{minipage} & \begin{minipage}[b]{\linewidth}\raggedright
Geometric Output Representation
\end{minipage} & \begin{minipage}[b]{\linewidth}\raggedright
Nominal Localization Uncertainty (\(\sigma_{xyz}\))
\end{minipage} & \begin{minipage}[b]{\linewidth}\raggedright
Compute Latency \& Memory Envelope
\end{minipage} & \begin{minipage}[b]{\linewidth}\raggedright
Deterministic Verification Readiness
\end{minipage} & \begin{minipage}[b]{\linewidth}\raggedright
Primary Systems Failure Signature \& Hardware Fallback Reflex
\end{minipage} \\
\midrule\noalign{}
\endfirsthead
\toprule\noalign{}
\begin{minipage}[b]{\linewidth}\raggedright
Grounding Architecture \& Paradigm
\end{minipage} & \begin{minipage}[b]{\linewidth}\raggedright
Geometric Output Representation
\end{minipage} & \begin{minipage}[b]{\linewidth}\raggedright
Nominal Localization Uncertainty (\(\sigma_{xyz}\))
\end{minipage} & \begin{minipage}[b]{\linewidth}\raggedright
Compute Latency \& Memory Envelope
\end{minipage} & \begin{minipage}[b]{\linewidth}\raggedright
Deterministic Verification Readiness
\end{minipage} & \begin{minipage}[b]{\linewidth}\raggedright
Primary Systems Failure Signature \& Hardware Fallback Reflex
\end{minipage} \\
\midrule\noalign{}
\endhead
\bottomrule\noalign{}
\endlastfoot
\textbf{2D Bounding-Box Detectors}(e.g., GLIP, Grounding DINO, MDETR) &
Planar bounding box
\([u_{\min}, v_{\min}, u_{\max}, v_{\max}] \subset \mathbb{R}^2\) +
depth plane backprojection &
\(\sigma_{xy} \approx 5\text{--}15\text{ mm}\)\(\sigma_z \approx 25\text{--}60\text{ mm}\)
(depth noise) & Latency: \(80\text{--}250\text{ ms}\)VRAM:
\(1.8\text{--}4.2\text{ GB}\) (FP16) & \textbf{High:} Bounding box
frustum yields convex hull; depth discontinuity filters detect
background bleed & Over-inclusive bounding box encloses adjacent fixture
\(\rightarrow\) Triggers spatial tolerance covariance expansion refusal
(\(\lambda_{\max}(\mathbf{\Sigma}) > \sigma_{\text{admit}}^2\)) \\
\textbf{Dense Semantic Segmentation}(e.g., Grounded-SAM, CLIPSeg,
Lang-SAM) & Pixel-level binary mask \(\mathcal{M} \subset \mathbb{Z}^2\)
+ surface normal point cloud extraction &
\(\sigma_{xy} \approx 2\text{--}8\text{ mm}\)\(\sigma_z \approx 5\text{--}18\text{ mm}\)
(surface mesh) & Latency: \(200\text{--}650\text{ ms}\)VRAM:
\(3.5\text{--}7.5\text{ GB}\) (FP16) & \textbf{High:} Exact boundary
contour allows 3D contact mesh fitting and point-cloud occupancy
verification & Mask bleeds across thin contact boundaries
\(\rightarrow\) Triggers multi-hypothesis contour entropy abort
(\(H(\mathbf{z}) > H_{\text{thresh}}\)) \\
\textbf{3D Neural Fields \& VLM Point Clouds}(e.g., LERF, ConceptFusion,
3D-LLM, OpenShape) & Continuous 3D feature field
\(\phi(\mathbf{p}) \in \mathbb{R}^d\) or oriented 3D bounding box
\(\in SE(3)\) &
\(\sigma_{xy} \approx 8\text{--}20\text{ mm}\)\(\sigma_z \approx 10\text{--}25\text{ mm}\)
(voxel grid) & Latency: \(400\text{--}1800\text{ ms}\)VRAM:
\(6.0\text{--}16.0\text{ GB}\) (FP16) & \textbf{Moderate:} Ray-marching
density thresholding admits local geometric minima and spatial drift &
Feature field local minimum generates phantom free-space target
\(\rightarrow\) Cross-validation against raw 3D ToF point cloud fails \\
\textbf{Direct Spatial Token Heads \& VLAs}(e.g., RT-2, PaLM-E, Octo,
OpenVLA) & Discretized Cartesian coordinate bins
\(\mathbf{p}_{\text{bin}} \in \mathbb{Z}^3\) or direct action tokens &
\(\sigma_{xyz} \approx 20\text{--}50\text{ mm}\) (bin quantization
discretization error) & Latency: \(150\text{--}500\text{ ms}\)VRAM:
\(8.0\text{--}28.0\text{ GB}\) (FP8/FP16) & \textbf{Low:} Lacks
geometric bounding volume; token logits exhibit non-Lipschitz spatial
jumps near decision seams & Out-of-distribution visual artifact triggers
discontinuous coordinate jump \(\rightarrow\) Rejected at intent
ingestion filter by kinematic workspace check
(\(\mathbf{p} \notin \mathcal{W}\)) \\
\end{longtable}

\section{Goals That Expire}\label{sec-intent-10-4}

An intent must die on its own, because whatever would have cancelled it
may be the thing that failed. In distributed computing architectures, a
coordinator cancels obsolete jobs by transmitting explicit abort signals
across an interconnect. In a physical machine, relying on explicit
cancellation introduces an asymmetric safety vulnerability. An abort
command requires an intact communications bus, an uncorrupted scheduler,
and a sender process that remains sufficiently functional to recognize
its own failure. When a learned vision-language model crashes, deadlocks
on a memory allocation, or hangs inside a graphics driver, no
cancellation packet is ever assembled or dispatched. If the motion
planner treats an active goal as valid until explicitly revoked, the
actuator continues executing open-loop physical actions toward a stale
or invalid target while the supervising software is dead. Granting
intent through a self-expiring \textbf{temporal lease} inverts this
failure mode. By bounding every dispatched goal with a finite lifetime,
silence becomes safe. If an upstream model stops communicating, the
absence of renewal automatically revokes the machine's authority to
move, requiring zero cooperation from the failed subsystem.\sidenote{\footnotesize In
  safety silicon (ARM Cortex-R52, TI TMS570), lease enforcement is
  hardware-driven via real-time timers (RTT) and windowed watchdogs
  (WWDT). Counter underflows assert Non-Maskable Interrupts or de-assert
  PWM gate drivers within nanoseconds, bypassing the RTOS scheduler and
  stalled threads.}

The duration of a temporal lease cannot be arbitrary. It must be derived
mathematically from the rate at which the physical environment diverges
from the model's perception of it. Consider an object moving or settling
in the physical workspace with a maximum drift velocity
\(v_{\text{drift}}\) relative to the world frame. Let \(r_{\text{tol}}\)
denote the spatial tolerance radius of the task, defined as the maximum
positional error between the end effector and the target that still
permits successful execution without collision. At the moment of
perception, the sensor pipeline introduces measurement noise with an
upper error bound \(\sigma_{\text{sensor}}\). Over an elapsed duration
\(t\), the object can displace by at most \(v_{\text{drift}} t\),
causing the total worst-case spatial error to grow according to
\(\sigma_{\text{sensor}} + v_{\text{drift}} t\). For the goal to remain
geometrically valid throughout its execution window, this total error
must not exceed the allowable tolerance radius. Setting the worst-case
error equal to \(r_{\text{tol}}\) establishes the maximum safe lease
duration,
\[\tau \le \frac{r_{\text{tol}} - \sigma_{\text{sensor}}}{v_{\text{drift}}}\]
When sensor noise meets or exceeds the tolerance radius
(\(r_{\text{tol}} \le \sigma_{\text{sensor}}\)), the spatial uncertainty
of perception exhausts the entire error budget immediately, yielding
\(\tau \le 0\text{ s}\) and requiring the system to reject the goal
without moving.

In an analytical example of an articulated manipulator picking
components from a moving conveyor, let the belt move with a drift
velocity of \(v_{\text{drift}} = 0.20\text{ m/s}\)
(\(200\text{ mm/s}\)). The mechanical compliance of the gripper allows a
positioning tolerance radius of \(r_{\text{tol}} = 15\text{ mm}\), while
the overhead vision pipeline exhibits a calibrated perception noise
bound of \(\sigma_{\text{sensor}} = 3\text{ mm}\). Substituting these
values into the lease equation yields
\[\tau \le \frac{15\text{ mm} - 3\text{ mm}}{200\text{ mm/s}} = \frac{12\text{ mm}}{200\text{ mm/s}} = 0.060\text{ s} = 60\text{ ms}\]
As plotted along the dynamic operating curve in
Figure~\ref{fig-10-lease-dynamics-tradeoff} (Panel A), a lease duration
longer than \(60\text{ ms}\) permits the component to drift outside the
gripper's physical capture envelope while the planner is still driving
the arm toward the initial coordinate. If the conveyor slows to
\(v_{\text{drift}} = 0.02\text{ m/s}\) (\(20\text{ mm/s}\)), the
allowable validity window expands to \(600\text{ ms}\). Because physical
drift rates vary with workspace dynamics, the safe lease duration is a
property of the environment and the tooling, not an arbitrary parameter
tuned to software convenience.

\begin{figure}[t]

\centering{

\includegraphics[width=1\linewidth,height=\textheight,keepaspectratio]{index_files/mediabag//Users/VJ/GitHub/MLSysBook-integrate-physical/physical/book/chapters/10-intent/figures/fig10_lease_dynamics_tradeoff.pdf}

}

\caption{\label{fig-10-lease-dynamics-tradeoff}\textbf{Physics-Derived
Lease Validity Horizon and Ingestion Rejection Gateways.} (Panel A)
Maximum safe lease duration \(\tau\) as a function of scene drift
velocity \(v_{\text{drift}}\) across distinct task tolerance radii
(\(r_{\text{tol}} = 5\text{ mm}, 15\text{ mm}, 30\text{ mm}\)) derived
from
\(\tau \le (r_{\text{tol}} - \sigma_{\text{sensor}}) / v_{\text{drift}}\).
The green shaded region denotes the admissible safe execution manifold,
while the red shaded region marks the stale-goal collision hazard zone.
Coupling lease duration to reasoning model inference cadence
(\(P_{99} = 650\text{ ms}\)) violates the physical bound under dynamic
drift. (Panel B) The 4-stage early rejection pipeline: (1) monotonic
sequence and parent perception hash verification (\(<1\ \mu\text{s}\)),
(2) \(O(1)\) analytical kinematic workspace manifold \(\mathcal{W}\) and
static occupancy check (\(<5\ \mu\text{s}\)), (3) dynamic reachability
filter (\(t_{\min} = 2\sqrt{d/a_{\max}} \le \tau_{\text{rem}}\)), and
(4) semantic covariance ambiguity gate
(\(\lambda_{\max}(\mathbf{\Sigma}) \le \sigma^2_{\max}\)). Inadmissible
proposals are rejected in \(42\ \mu\text{s}\) (\(P_{99}\)) with
structured diagnostic feedback, eliminating \(15\text{–}80\text{ ms}\)
trajectory optimizer stalls and preventing planning queue starvation.}

\end{figure}%

Coupling the lease duration to the execution cadence of the reasoning
model introduces an open-loop drift hazard whenever inference latency
fluctuates. A common engineering error sets the lease duration equal to
the mean execution period of the high-level policy, granting a
\(250\text{ ms}\) validity window because the vision model operates at a
nominal \(4\text{ Hz}\). This conflates compute throughput with physical
state stability. Physical scene dynamics dictate the rate at which
reality diverges from the state estimate, whereas inference latency
reflects GPU memory bandwidth, batch size, and scheduler contention.
When an onboard accelerator encounters thermal throttling or kernel
preemption, inference latency stretches into long tails. If a policy
with a mean latency of \(250\text{ ms}\) exhibits a \(P_{99}\) latency
spike of \(650\text{ ms}\), a system that coupled lease validity to the
nominal cadence allows the machine to operate blindly for
\(400\text{ ms}\) past its safe envelope. The lease duration \(\tau\)
must remain anchored to physical dynamics, defining a hard deadline that
the inference pipeline must service, rather than an elastic budget that
expands to match computational delay.

To sustain continuous motion without coming to an unwanted stop, the
high-level reasoning system must refresh the active lease before it
expires. This renewal protocol requires the upstream model to transmit a
new geometric proposal and an updated duration timestamp at an arrival
period \(T_{\text{period}} \le \tau - t_{\text{latency}}\), where
\(t_{\text{latency}}\) accounts for the total time spent in image
capture, inference computation, serialization, and network transport.
When a renewal arrives while the current lease remains active, the
downstream trajectory planner must integrate the updated coordinates
without generating acceleration spikes. The planner evaluates the new
boundary conditions, computes a jerk-limited trajectory splice from the
current state vector \((\mathbf{p}, \mathbf{v}, \mathbf{a})\) to the
refreshed goal manifold, and extends the execution time horizon. If the
updated goal lies within the kinematic tolerance of the current path,
the controller smoothly updates the terminal waypoint. If the refreshed
goal shifts abruptly, indicating an object jump or semantic
reclassification, the planner rejects the trajectory splice and prepares
the actuator for an immediate controlled stop.

In a canonical heterogeneous dual-brain System-on-Chip (SoC)
architecture, this temporal lease is enforced at the silicon boundary.
The unprivileged Linux host processor (MPU) constructs the intent record
in user space, populating \(\tau\) based on environmental dynamics, and
transfers the payload into a shared SRAM ring buffer via the Remote
Processor Messaging (RPMSG) framework.\sidenote{\footnotesize Heterogeneous SoCs use
  RPMSG over shared-memory vrings with hardware mailboxes. When the
  Linux MPU writes an intent record to shared SRAM, it executes data
  synchronization barriers (\texttt{DSB\ ISH}). Real-time MCUs service
  mailbox interrupts in \(<500\text{ ns}\), avoiding bus lock
  contention.} The deterministic real-time microcontroller (MCU) reads
the record and immediately arms a dedicated hardware countdown timer
(such as a Windowed Watchdog Timer or high-resolution hardware counter).
If the Linux MPU freezes, suffers driver contention, or fails to deliver
an authenticated renewal packet before the hardware counter reaches
zero, the MCU timer underflows. This underflow triggers a non-maskable
hardware interrupt in \(< 1\ \mu\text{s}\), causing the MCU to de-assert
motor enable lines or switch the inverter gate drivers into a
deterministic fallback state, completely bypassing the stalled Linux
operating system.

When a lease expires without renewal, the machine executes the
deterministic terminal state specified within the original lease
contract. A physical system cannot simply hold its last commanded
actuator velocities, because unguided momentum carries mass across the
workspace without valid intent. The safety enforcer in the nervous
system triggers one of three deterministic actions the instant the timer
reaches zero. An active position hold engages closed-loop
proportional-derivative control to lock joint coordinates against
gravity, preserving the machine's position in free space. A controlled
deceleration applies maximum permissible joint deceleration along the
current path tangent vector, dissipating kinetic energy to bring the
body to rest before geometric boundaries are crossed. A zero-torque
compliance mode disables motor current or sets active stiffness to zero,
allowing external mechanical compliance or gravity to govern the body
when an abrupt freeze would risk pinching a human or damaging a brittle
workpiece. Because this fallback behavior is compiled into the
deterministic nervous system before motion begins, the transition occurs
without requiring communication with the failed upstream model.

Executing an expiring lease safely assumes that the commanded geometric
proposal is actually achievable by the physical body. If a proposed
coordinate lies beyond kinematic reach or demands joint velocities that
exceed actuator limits during the lease window, granting the lease
wastes control authority and risks triggering hardware fault trips.
Determining whether a goal is kinematically and dynamically admissible
belongs before the trajectory generator accepts the lease, establishing
a physical filter between high-level intent and motion execution.

\section{Reachability and Admissible Goals}\label{sec-intent-10-5}

Filtering a goal against the physical capability of the machine belongs
at the intent ingestion boundary, before the proposed target enters the
trajectory optimization queue. When an upstream vision-language model or
neural policy emits a coordinate, treating that proposal as an
unverified candidate protects the real-time motion pipeline from
executing or even evaluating impossible motions. Kinematic workspace
analysis, imported from classical formulations
(\citeproc{ref-craig2005introduction}{Craig 2005};
\citeproc{ref-lynch2017modern}{Lynch and Park 2017}), defines the closed
spatial manifold of reachable end-effector configurations given link
dimensions, joint range limits, and mounting geometry.\sidenote{\footnotesize The
  reachable workspace \(\mathcal{W} \subset SE(3)\) is the forward
  kinematics image over joint limits. Near kinematic singularities where
  Jacobian rank collapses (\(\det(J J^T) \to 0\)), small Cartesian steps
  demand unbounded joint velocities
  (\(\dot{\mathbf{q}} = J^\dagger \mathbf{v}\)), causing trajectory
  optimization failure.} An analytical reachability check evaluates
whether the target coordinate \(\mathbf{p}_{\text{target}}\) resides
inside this workspace manifold \(\mathcal{W}\), while an occupancy check
verifies that the bounding volume of the target does not intersect known
persistent obstacles. If a proposed touch target lies outside
\(\mathcal{W}\), no sequence of joint torques can reach it. Passing an
out-of-reach goal to a numerical trajectory optimizer forces the solver
to search an empty solution space, converting a geometric impossibility
into a computing failure.

The computational cost of rejecting an invalid goal at the ingestion
boundary is orders of magnitude smaller than rejecting it inside the
motion planner (Figure~\ref{fig-10-lease-dynamics-tradeoff}, Panel B).
An analytical workspace bounding check or hierarchical bounding-sphere
intersection executes in \(O(1)\) time, requiring less than
\(5\ \mu\text{s}\) of scalar processor time. In contrast, a trajectory
optimization loop solving a constrained non-linear program across \(N\)
discrete time steps consumes between \(15\text{ ms}\) and
\(60\text{ ms}\) of compute time, with a \(P_{99}\) tail latency
exceeding \(80\text{ ms}\) under complex collision geometry. When an
unachievable target enters the planner, numerical solvers run through
their maximum iteration budget before declaring infeasibility, driving
internal slack variables against their penalty boundaries. During that
calculation window, the planner cannot service valid incoming goals,
creating planning queue starvation. If the machine operates under a
\(50\text{ ms}\) control cycle deadline, a single unachievable goal
creates a pipeline bubble that starves downstream interpolation buffers
and forces the nervous system into a fallback stop. Early rejection
ensures that the trajectory planner only commits compute cycles to goals
that possess guaranteed kinematic solutions.

\protect\phantomsection\label{callout-ruleux2a-10.6}
\begin{fbxSimple}{callout-rule}{Architect’s Rule of Thumb:}{Reject at the Ingestion Gate, Not in the Trajectory Solver}
\phantomsection\label{callout-rule*-10.6}

Evaluating kinematic reachability and static occupancy in \(O(1)\)
scalar time (\(< 5\,\mu\text{s}\)) at the intent ingestion boundary
prevents invalid goals from entering \(O(N)\) trajectory optimization
loops (\(15\text{--}80\text{ ms}\)), eliminating solver starvation
bubbles and avoiding downstream pipeline deadline misses.

\end{fbxSimple}

A target that lies within the static kinematic workspace may still be
physically unachievable within the operational window. Static
reachability establishes only that a stationary pose exists where the
end-effector touches the target coordinate. In contrast, \textbf{dynamic
reachability} requires that the machine can transition from its current
kinematic state \((\mathbf{p}_0, \mathbf{v}_0)\) to the target state
\(\mathbf{p}_{\text{target}}\) within the remaining lease duration
\(t_{\text{rem}}\) without exceeding actuator velocity limits
\(v_{\max}\) and acceleration limits \(a_{\max}\). Consider an
analytical example of a single-axis translation over distance
\(d = \|\mathbf{p}_{\text{target}} - \mathbf{p}_0\|\) starting and
ending at rest. The minimum distance required to accelerate to maximum
velocity and decelerate back to zero is
\(d_{\text{crit}} = v_{\max}^2 / a_{\max}\). If \$d
\le \(d_{\text{crit}}\), the actuator executes a triangular velocity
profile without reaching \(v_{\max}\), yielding a minimum transit time
of \(t_{\min} = 2\sqrt{d / a_{\max}}\). If \(d > d_{\text{crit}}\), the
actuator reaches \(v_{\max}\) and cruises over distance
\(d - d_{\text{crit}}\), yielding
\(t_{\min} = (d / v_{\max}) + (v_{\max} / a_{\max})\). If the remaining
lease duration \(t_{\text{rem}}\) is less than \(t_{\min}\), the goal is
dynamically inadmissible. Admitting an inadmissible goal guarantees that
the lease will expire before the physical body reaches the target,
triggering a mid-trajectory deceleration that wastes mechanical work and
cycle time.

Formally, \textbf{dynamic reachability} is the physical condition
wherein an end-effector can transition from its instantaneous initial
state \((\mathbf{p}_0, \mathbf{v}_0)\) to a target pose
\(\mathbf{p}_{\text{target}}\) within the remaining lease validity
duration \(t_{\text{rem}}\)
(\(t_{\min}(\|\Delta \mathbf{p}\|, v_{\max}, a_{\max}) \le t_{\text{rem}}\))
while strictly respecting actuator velocity (\(v_{\max}\)), acceleration
(\(a_{\max}\)), and thermal torque limits.

These dynamic reachability bounds rest on imported rigid-body mechanics
under three assumptions, namely that actuator torque-speed curves remain
within nominal thermal envelopes, joint encoders accurately report the
initial state \((\mathbf{p}_0, \mathbf{v}_0)\), and the payload mass
matches the inertial model used in the acceleration bounds. The machine
verifies these assumptions at runtime through high-frequency telemetry
in the nervous system, sampling joint current, motor temperature, and
bus latency at \(1000\text{ Hz}\). When motor winding temperatures
exceed thermal thresholds, or when an external force measurement
indicates payload mass deviation, the nervous system scales down
\(a_{\max}\) and \(v_{\max}\) in the parameter registry. The intent
filter immediately evaluates subsequent incoming leases against these
reduced limits. At the instant thermal derating occurs, goals that were
previously admissible become dynamically unreachable, and the ingestion
filter drops them before they can command unattainable torques.

When the ingestion filter rejects an intent proposal, it must not simply
discard the packet. Discarding without feedback leaves the upstream
reasoning model unaware of why motion failed to occur, forcing it to
wait for an end-to-end task timeout. The intent filter generates an
immediate, structured diagnostic packet returned directly to the
reasoning layer over the inter-core mailbox. This diagnostic payload
contains the rejection classification, such as a kinematic workspace
violation (\texttt{ERR\_WORKSPACE\_OUT\_OF\_BOUNDS}), dynamic time
insufficiency (\texttt{ERR\_DYNAMIC\_TIME\_INSUFFICIENT}), or
swept-volume collision intersection (\texttt{ERR\_OCCUPANCY\_BLOCKED}),
along with the numerical margins of failure. For a dynamic reachability
failure, the feedback returns the shortfall
\(t_{\min} - t_{\text{rem}}\) and the maximum reachable distance along
the commanded vector within the allotted lease. For a static workspace
violation, the feedback provides the projection of the target onto the
boundary of \(\mathcal{W}\). The reasoning layer uses this structured
signal to adjust its proposal on the subsequent inference cycle, either
by extending the requested lease duration, choosing an alternate grasp
approach, or commanding base relocation.

Filtering goals for reachability guarantees that every intent handed to
the trajectory generator is geometrically and dynamically solvable by
the physical body. It prevents computational starvation in the motion
planner and eliminates dead time caused by attempting unachievable
trajectories. Yet reachability verifies only that the machine has the
physical capability to reach a commanded coordinate in space. It cannot
determine whether that coordinate corresponds to a meaningful physical
object or a phantom artifact of sensory noise.

\section{How Intent Goes Wrong}\label{sec-intent-10-6}

When a vision-language model outputs a three-dimensional target
coordinate for a physical manipulator, that coordinate can fail in
several distinct ways. It may be a hallucinated target coordinate
generated from visual attention patterns on reflective background
textures. It may be a stale target representing an object that a human
coworker bumped or displaced \(300\text{ ms}\) earlier while the camera
pipeline was buffering frames. It may be a reference ambiguity where the
model arbitrarily selects one of two identical machined parts on a tray
without resolving which one the task requires. Or it may stem from an
out-of-distribution scene shift, such as unexpected lighting or dust on
an optical lens, that degrades the spatial prediction accuracy of the
network. While their causal origins differ between learned
representation errors and physical timing mismatches, a hallucinated
target and a correct target that outlived its scene fail the same way
downstream.

To the trajectory generator and the nervous system, every intent arrives
as a coordinate, a bounding volume, or a target pose. If the brain
commands a grasp at a coordinate where an object previously sat, the
motion planner solves an inverse kinematics problem and plans an
obstacle-free trajectory to empty air. If the brain commands a grasp at
a coordinate generated by an attention hallucination inside a solid
steel fixture, the planner attempts to construct a path into occupied
space. In both cases, the downstream motion pipeline receives a goal
that is mathematically well-formed but physically unachievable or
hazardous. When the target lies in free space, the manipulator sweeps
through empty air, closing its gripper on nothing and stalling until a
task timeout expires. When the target intersects unexpected geometry,
the planned trajectory commands the arm into contact with solid mass,
converting a software perception error into an unintended impact at
\(1.5\text{ m/s}\) unless an independent layer halts motion.

The first defense against ungrounded coordinates is geometric
cross-validation against low-latency depth sensors. While large
vision-language models may require \(200\text{ ms}\) to
\(500\text{ ms}\) of inference time, a structured-light camera or a
time-of-flight sensor produces a raw point cloud at \(30\text{ Hz}\) or
\(60\text{ Hz}\), maintaining an occupancy map updated within
\(20\text{ ms}\). Before an intent proposal passes to the trajectory
generator, the ingestion filter cross-checks the commanded coordinate
against this occupancy map. In an analytical example where the brain
commands a grasp at a target coordinate, the filter queries the point
cloud within a bounding sphere of radius \(30\text{ mm}\) centered on
the target. If the sphere contains zero depth returns, the filter
rejects the target as a hallucination or a stale coordinate from an
object that was removed. If the depth returns indicate a solid surface
extending continuously across the approach vector, the filter rejects
the target as an unsegmented fixture. Geometric occupancy checks detect
gross spatial hallucinations and out-of-date free-space targets before
the planner spends compute generating a trajectory.

Geometric validation cannot detect errors where the coordinate lands
squarely on a real, solid object, but that object is the wrong one. When
a scene contains multiple identical workpieces and the prompt requests
the defective part, an overconfident model may emit the coordinate of
whichever object had a slightly higher logit in its spatial token head.
The nervous system can detect this semantic reference ambiguity by
evaluating the entropy across candidate object detections before issuing
a physical intent record. If the model outputs a multimodal probability
distribution over several candidate poses, where the difference in
softmax confidence between the top two candidates is within an
analytical margin of \(0.15\), the spatial reference is ambiguous.
Rather than allowing the physical body to execute an arbitrary choice
that may ruin a workpiece, the ingestion pipeline intercepts the
high-entropy goal and halts execution, returning an ambiguity diagnostic
to the reasoning layer or requesting disambiguation.

Intent failure modes organize along two axes, defined by whether the
error is detectable by geometric occupancy or requires semantic
verification, and how rapidly that failure converts into physical
hazard. Gross hallucinations and stale targets in occupied space
represent immediate kinematic collisions, detectable geometrically in
milliseconds by raw point clouds. Stale targets in free space represent
delayed starvation hazards, detectable geometrically by the absence of
mass. Reference ambiguities and misclassified affordances, such as
grasping a fragile ceramic vessel by its lip with excessive force or
picking up the wrong identical part, pass every geometric occupancy
filter because physical mass is present at the target. These errors have
a moderate time-to-hazard, governed by the duration of the transport
phase, and can only be caught prior to motion by evaluating semantic
uncertainty distributions. The complete taxonomy of intent failure
modes, their mathematical detection invariants, real-time sensor
channels, and corresponding deterministic hardware reflexes is
formalized in Table~\ref{tbl-10-grounding-failure-modes}.

\begin{longtable}[]{@{}
  >{\raggedright\arraybackslash}p{(\linewidth - 8\tabcolsep) * \real{0.2000}}
  >{\raggedright\arraybackslash}p{(\linewidth - 8\tabcolsep) * \real{0.2000}}
  >{\raggedright\arraybackslash}p{(\linewidth - 8\tabcolsep) * \real{0.2000}}
  >{\raggedright\arraybackslash}p{(\linewidth - 8\tabcolsep) * \real{0.2000}}
  >{\raggedright\arraybackslash}p{(\linewidth - 8\tabcolsep) * \real{0.2000}}@{}}
\caption{\textbf{Semantic Grounding Failure Modes, Detection Invariants,
and Real-Time Mitigation Actions.} Structured taxonomy of intent failure
modes categorized by physical root cause, mathematical detection
triggers, sensor tiers, deterministic hardware reflexes, and physical
stopping margins. A robust physical AI architecture ensures that every
failure mode has an independent, non-overridable deterministic
interceptor in the nervous
system.}\label{tbl-10-grounding-failure-modes}\tabularnewline
\toprule\noalign{}
\begin{minipage}[b]{\linewidth}\raggedright
Failure Taxonomy \& Root Cause
\end{minipage} & \begin{minipage}[b]{\linewidth}\raggedright
Mathematical Detection Invariant / Trigger Condition
\end{minipage} & \begin{minipage}[b]{\linewidth}\raggedright
Detection Tier \& Sensor Channel
\end{minipage} & \begin{minipage}[b]{\linewidth}\raggedright
Real-Time Mitigation Action \& Hardware Reflex
\end{minipage} & \begin{minipage}[b]{\linewidth}\raggedright
Worst-Case Stopping \& Recovery Margin
\end{minipage} \\
\midrule\noalign{}
\endfirsthead
\toprule\noalign{}
\begin{minipage}[b]{\linewidth}\raggedright
Failure Taxonomy \& Root Cause
\end{minipage} & \begin{minipage}[b]{\linewidth}\raggedright
Mathematical Detection Invariant / Trigger Condition
\end{minipage} & \begin{minipage}[b]{\linewidth}\raggedright
Detection Tier \& Sensor Channel
\end{minipage} & \begin{minipage}[b]{\linewidth}\raggedright
Real-Time Mitigation Action \& Hardware Reflex
\end{minipage} & \begin{minipage}[b]{\linewidth}\raggedright
Worst-Case Stopping \& Recovery Margin
\end{minipage} \\
\midrule\noalign{}
\endhead
\bottomrule\noalign{}
\endlastfoot
\textbf{Phantom Target Hallucination}(Visual attention artifact or
specular reflection generates target in empty space) &
\(\mathcal{N}_{\text{cloud}}(\mathbf{p}_{\text{target}}, r_{\text{sphere}}) = \emptyset\)(Zero
depth returns in local spatial bounding sphere) & Ingestion filter
cross-check against 3D ToF LiDAR / structured-light point cloud
(\(30\text{--}60\text{ Hz}\)) & \textbf{Pre-Planning Ingestion Drop:}
Rejects intent packet, emits structured diagnostic
\texttt{ERR\_EMPTY\_VOXEL}, keeps actuators stationary & \(0\text{ mm}\)
displacement (rejected before trajectory generation begins;
\(t_{\text{reject}} \le 42\,\mu\text{s}\)) \\
\textbf{Stale Target from Unobserved Motion}(Workpiece shifts or is
removed during upstream neural inference window) &
\((\mathbf{x}_{\text{meas}} - \mathbf{p}_{\text{target}})^T \mathbf{\Sigma}^{-1} (\mathbf{x}_{\text{meas}} - \mathbf{p}_{\text{target}}) > 1\)(Tracked
target state exceeds covariance ellipsoid) & Fast visual odometry /
local point cloud tracker in nervous system
(\(100\text{--}500\text{ Hz}\)) & \textbf{Dynamic Trajectory Abort:}
Invalidate active lease locally, command maximum deceleration
\(a_{\max}\) to active hold state &
\(t_{\text{stop}} \le 20\text{ ms}\); physical displacement bounded to
\(\Delta x \le \frac{v_0^2}{2 a_{\max}} + v_0 t_{\text{abort}} \le 5.2\text{ mm}\) \\
\textbf{Semantic Reference Ambiguity}(Multiple identical workpieces
produce multi-peak posterior logits) &
\(H(\mathbf{z}) = -\sum_{i=1}^K p_i \log_2 p_i > H_{\text{thresh}}\)or
\(\lvert p_1 - p_2 \rvert \le \delta_{\text{margin}}\) & Ingestion
filter softmax distribution and spatial covariance evaluator
(\(10\text{--}50\text{ Hz}\)) & \textbf{Grounding Refusal:} Intercept
command, hold actuators in safe position, transmit disambiguation
request to supervisor & \(0\text{ mm}\) displacement (refusal enforced
before trajectory synthesis; no arbitrary tie-breaking) \\
\textbf{Misclassified Surface Affordance}(Grasping rigid bolted fixture
or brittle structure with full clamping torque) &
\(\|\mathbf{F}_{\text{meas}}\| > F_{\text{max}}\) or
\(\|\boldsymbol{\tau}_{\text{meas}} - \mathbf{J}^T \mathbf{F}_{\text{cmd}}\| > \boldsymbol{\tau}_{\text{thresh}}\)(Excess
contact reaction force / torque mismatch) & Real-time 6-axis wrist
force-torque transducer and joint torque strain gauges
(\(1000\text{ Hz}\)) & \textbf{Zero-Torque Compliance / STO:} Hardware
analog/digital comparator switches inverter to compliance damping or
Safe Torque Off & \(t_{\text{reflex}} \le 2.0\text{ ms}\); peak contact
wrench bounded to \(F_{\text{peak}} \le k \Delta x \le 8.0\text{ N}\) \\
\textbf{Dynamic Reachability Violation}(Target kinematically valid but
unreachable within lease duration under \(a_{\max}, v_{\max}\)) &
\(t_{\min}(\|\Delta \mathbf{p}\|, v_{\max}, a_{\max}) > t_{\text{rem}}\)where
\(t_{\min} = \frac{d}{v_{\max}} + \frac{v_{\max}}{a_{\max}}\)
(trapezoidal profile) & Ingestion boundary kinematic and dynamic
feasibility checker (\(1000\text{ Hz}\)) & \textbf{Early Admission
Rejection:} Drop packet immediately, return shortfall margin
\(t_{\min} - t_{\text{rem}}\) to reasoning layer & \(0\text{ ms}\)
motion wasted; eliminates trajectory solver queue starvation (\(O(1)\)
scalar check \(< 5\,\mu\text{s}\)) \\
\end{longtable}

\protect\phantomsection\label{callout-incidentux2a-10.7}
\begin{fbxSimple}{callout-incident}{Landmark Engineering Incident:}{Industrial Robotic Sorting Voice Ambiguity Runaway}
\phantomsection\label{callout-incident*-10.7}

\textbf{System:} High-throughput automated e-commerce sorting cell
featuring a 6-axis articulated industrial manipulator
(\(16.0\text{ kg}\) payload capacity), multimodal vision-language model
(VLM) frontend for natural language operator tasking, and a high-speed
vacuum/clamp end-effector.

\textbf{Incident:} A verbal operator directive (``clear the blue
container near the bin'') possessed unresolvable semantic reference
ambiguity between two identical blue plastic totes placed
\(180\text{ mm}\) apart on a dual-lane sorting table. Rather than
refusing execution, the upstream reasoning pipeline broke the evidential
tie by computing the spatial centroid of both candidates, outputting an
ungrounded intent coordinate into free space adjacent to a rigid steel
frame. The unconstrained intent packet was dispatched without a spatial
covariance bound, expiring TTL lease, or wrench limit, causing the arm
to execute a high-speed ballistic sweep (\(1.8\text{ m/s}\)) directly
into the structural divider flange.\sidenote{\footnotesize OSHA and ISO 10218
  investigation archives document multiple incidents where multi-modal
  language or vision ambiguities triggered ungrounded trajectory
  generation, causing mechanical collisions with workcell fixtures.}

\textbf{Telemetry \& Physical Failure Sequence:}

\begin{itemize}
\tightlist
\item
  \(t = 0.0\text{ ms}\): Operator issues voice directive; acoustic
  reflection from nearby conveyor motors introduces noise, degrading
  phoneme extraction confidence.
\item
  \(t = 160.0\text{ ms}\): Multimodal VLM generates a bimodal posterior
  distribution (\(p_1 = 0.52, p_2 = 0.48\)); softmax argmax collapses
  the distribution into an intermediate spatial coordinate
  \(\mathbf{p}_0 = [0.650, 0.120, 0.180]^T\text{ m}\), located squarely
  in the physical clearance envelope of the central steel divider.
\item
  \(t = 210.0\text{ ms}\): Host MPU publishes an intent packet with an
  infinite expiration horizon (\(t_{\text{exp}} = \infty\)) and zero
  covariance bounds (\(\mathbf{\Sigma} = \mathbf{0}\)).
\item
  \(t = 280.0\text{ ms}\): \(500\text{ Hz}\) trajectory optimizer
  interpolates a minimum-jerk quintic spline toward the phantom
  coordinate at \(1.8\text{ m/s}\).
\item
  \(t = 465.0\text{ ms}\): End-effector impacts the rigid steel
  structural frame at \(1.65\text{ m/s}\), registering an instantaneous
  peak contact force \(F_{\text{peak}} = 2150\text{ N}\).
\item
  \(t = 485.0\text{ ms}\): Tracking error integrator in the position
  loop winds up attempting to overcome the mechanical obstruction,
  commanding maximum saturation current to Joint 3 and Joint 4 motor
  drives (\(I_q = 45.0\text{ A}\), \(320\%\) nominal rating).
\item
  \(t = 620.0\text{ ms}\): Joint 3 harmonic drive gear teeth shear under
  excessive reaction torque
  (\(\tau_{\text{react}} = 380\text{ N}\cdot\text{m}\) against a
  \(160\text{ N}\cdot\text{m}\) continuous rating); stator winding
  temperature rises at \(18.5^{\circ}\text{C/s}\) until hardware thermal
  shutdown trips at \(t = 1150.0\text{ ms}\).
\end{itemize}

\textbf{Root Cause:} The multimodal perception model broke evidential
ties through an optimistic argmax projection instead of issuing an
explicit grounding refusal (\texttt{ERR\_GROUNDING\_AMBIGUITY}) when
semantic entropy exceeded threshold
(\(H(\mathbf{z}) > H_{\text{thresh}}\)).

\textbf{Architectural Flaw:} The system architecture lacked three
foundational physical AI boundaries: (1) no covariance admission gate
(\(\lambda_{\max}(\mathbf{\Sigma}) \le \sigma_{\text{admit}}^2\)), (2)
no self-expiring temporal lease with local hardware watchdog countdown,
and (3) no trajectory-level wrench limit (\(F_{\max}\)) to saturate
motor currents upon unexpected rigid contact.

\textbf{Systems Invariant:} Ambiguity in physical AI is an actionable
refusal condition, never an argmax guess. Every intent lease must
encapsulate an oriented spatial tolerance volume, a finite TTL
countdown, and a physical wrench ceiling before motion authority is
committed.

\end{fbxSimple}

The most dangerous failure modes are those that pass both geometric
occupancy filters and semantic entropy checks, entering the motion
planner with valid status. A vision-language model may identify a clean,
reachable grasp pose on a tool, but misjudge the friction coefficient of
the handle or fail to recognize that the tool is mechanically bolted to
a table. Because mass exists at the predicted location and the target
classification exhibits high model confidence, no upstream software
check flags the intent. When the manipulator closes its fingers and
attempts to lift, the kinematic trajectory demands upward acceleration
that the mechanical constraint prevents. In these blind spots, software
intent validation fails completely. Safety falls entirely to low-level
contact reflexes in the nervous system, where joint torque sensors
detect an unexpected reaction force exceeding \(20\text{ N}\) within
\(2\text{ ms}\) and abort current to the motor windings, or to human
authority exercising a hard-wired emergency stop.

Because intent failures span this range of physical consequences and
detection mechanisms, the interface between the brain and the nervous
system cannot consist of naked Cartesian coordinates or raw joint
angles. An unadorned coordinate carries neither the physical assumptions
under which it was generated nor the temporal boundary beyond which it
becomes hazardous. To survive the inevitable failures of learned
perception, every goal handed to the physical machine must be
encapsulated within a formal structure that declares its spatial
envelope, its dynamic constraints, and its explicit lifetime.

\section{The Intent Lease}\label{sec-intent-10-7}

Building on the system schema, data types, and record versioning
established in Chapter~\ref{sec-nervous}, the \textbf{intent record}
defines the minimal payload required to bound a physical objective in
space, time, and force. An unadorned kinematic command fails because it
communicates only where a mechanism should move, omitting the physical
assumptions, force limits, and temporal boundaries that justify the
motion. The intent record resolves this vulnerability by wrapping the
target objective in an explicit mathematical envelope that the nervous
system can parse, validate, and enforce without consulting the model
that produced it. If an upstream planner stalls, crashes, or drops
offline, the record contains every parameter required to bring the
machine to a controlled physical stop before the lease boundary is
breached.

Formally, the \textbf{intent record} is the strictly typed IPC message
structure crossing the Dual-Brain boundary that bounds a high-level
cognitive goal within an explicit 3D target pose (\(T \in SE(3)\)),
spatial tolerance matrix (\(\mathbf{\Sigma}\)), temporal validity lease
(\(t_{\text{exp}}\)), maximum allowable wrench envelope
(\(F_{\text{max}}\)), cryptographic parent perception digest, and
deterministic terminal behavior flags.

In a serialized byte-level layout, an intent record occupies a fixed
footprint of \(168\text{ bytes}\), avoiding dynamic memory allocation in
real-time communication rings. The payload begins with an
\(8\text{-byte}\) monotonic sequence counter and an \(8\text{-byte}\)
perception timestamp \(t_{\text{src}}\) recorded in nanoseconds of UTC
epoch time. Kinematics are defined by a target pose \(T \in SE(3)\),
packed as three \(8\text{-byte}\) double-precision Cartesian coordinates
\((x, y, z)\) and a four-element unit quaternion
\((q_w, q_x, q_y, q_z)\) totaling \(56\text{ bytes}\). A diagonal
tolerance matrix \(\mathbf{\Sigma} \in \mathbb{R}^6\) specifies
allowable translational and rotational tracking deviations across six
single-precision floats (\(24\text{ bytes}\)). The record closes with an
\(8\text{-byte}\) lease expiration timestamp \(t_{\text{exp}}\), a
six-element wrench limit vector \(F_{\text{max}} \in \mathbb{R}^6\)
bounding allowable forces and torques using single-precision floats
(\(24\text{ bytes}\)), and an \(8\text{-byte}\) field containing
terminal behavior flags and memory alignment padding. The complete
byte-level serialization layout, physical units, mathematical validation
invariants, and hardware enforcement mechanisms of the intent record
schema are summarized in Table~\ref{tbl-10-intent-schema}.

\marginnote{\begin{footnotesize}

\textbf{Cryptographic Lineage \& Token Signing} A \(32\text{-byte}\)
cryptographic hash of the parent perception state binds the command
directly to the specific sensor frame from which the inference was
derived. This ensures that the real-time enforcer can verify
cryptographic lineage via \texttt{HMAC-SHA256} signatures against a
rolling buffer of authorized hardware sensor captures, rejecting
injected packets or stale inferences.

\end{footnotesize}}

\begin{longtable}[]{@{}
  >{\raggedright\arraybackslash}p{(\linewidth - 8\tabcolsep) * \real{0.2000}}
  >{\raggedright\arraybackslash}p{(\linewidth - 8\tabcolsep) * \real{0.2000}}
  >{\raggedright\arraybackslash}p{(\linewidth - 8\tabcolsep) * \real{0.2000}}
  >{\raggedright\arraybackslash}p{(\linewidth - 8\tabcolsep) * \real{0.2000}}
  >{\raggedright\arraybackslash}p{(\linewidth - 8\tabcolsep) * \real{0.2000}}@{}}
\caption{\textbf{Intent Lease Protocol Contract Schema and Serialization
Specification.} Byte-level memory layout (\(168\text{ bytes}\) total),
data types, physical coordinate frames, real-time validation invariants,
and low-level nervous system enforcement mechanisms for the intent
record. This strictly typed binary contract decouples high-level learned
perception from low-level deterministic motor actuation, ensuring that
unvalidated model hypotheses can never compromise physical hardware
safety.}\label{tbl-10-intent-schema}\tabularnewline
\toprule\noalign{}
\begin{minipage}[b]{\linewidth}\raggedright
Byte Offset \& Field Name
\end{minipage} & \begin{minipage}[b]{\linewidth}\raggedright
Data Type \& Size
\end{minipage} & \begin{minipage}[b]{\linewidth}\raggedright
Physical Units \& Frame
\end{minipage} & \begin{minipage}[b]{\linewidth}\raggedright
Real-Time Validation Invariant \& Pre-Condition
\end{minipage} & \begin{minipage}[b]{\linewidth}\raggedright
Hardware Interceptor \& Enforcement Mechanism
\end{minipage} \\
\midrule\noalign{}
\endfirsthead
\toprule\noalign{}
\begin{minipage}[b]{\linewidth}\raggedright
Byte Offset \& Field Name
\end{minipage} & \begin{minipage}[b]{\linewidth}\raggedright
Data Type \& Size
\end{minipage} & \begin{minipage}[b]{\linewidth}\raggedright
Physical Units \& Frame
\end{minipage} & \begin{minipage}[b]{\linewidth}\raggedright
Real-Time Validation Invariant \& Pre-Condition
\end{minipage} & \begin{minipage}[b]{\linewidth}\raggedright
Hardware Interceptor \& Enforcement Mechanism
\end{minipage} \\
\midrule\noalign{}
\endhead
\bottomrule\noalign{}
\endlastfoot
\texttt{0x00..0x07}\texttt{sequence\_id} & \texttt{uint64\_t}(8 bytes) &
Monotonic counter (dimensionless) &
\(\text{Seq} > \text{Seq}_{\text{last\_committed}}\) (strictly
monotonic; detects out-of-order or replayed packets) & Atomic register
comparison in NIC/DMA ring buffer (\(< 50\text{ ns}\)); drops replayed
packets \\
\texttt{0x08..0x0F}\texttt{t\_source} & \texttt{uint64\_t}(8 bytes) &
Nanoseconds (\(\text{ns}\)) UTC epoch time &
\(t_{\text{current}} - t_{\text{src}} \le \Delta t_{\text{max\_age}}\)
(verifies freshness of upstream camera exposure) & Real-time hardware
clock counter check; drops stale perception frames before decoding \\
\texttt{0x10..0x2F}\texttt{parent\_hash} & \texttt{uint8\_t{[}32{]}}(32
bytes) & SHA-256 Digest (dimensionless) &
\(\text{Match}(\text{parent\_hash}, \mathcal{B}_{\text{sensor\_ring}})\)
(verifies cryptographic perception lineage) & Fast SRAM hash-table
lookup against local sliding window of verified sensor captures \\
\texttt{0x30..0x67}\texttt{target\_pose\_SE3} &
\texttt{double{[}7{]}}(56 bytes) & \([x, y, z]\) in \(\text{m}\);
\([q_w, q_x, q_y, q_z]\) unit quaternion (Robot Base Frame) &
\(\left\lvert \|\mathbf{q}\|_2^2 - 1.0 \right\rvert \le 10^{-6}\)
(quaternion normality) and \(\mathbf{p} \in \mathcal{W}_{\text{reach}}\)
(kinematic workspace manifold) & Hardware FPU vector norm evaluator +
\(O(1)\) analytical boundary check against joint limit table \\
\texttt{0x68..0x7F}\texttt{tolerance\_sigma} &
\texttt{float32{[}6{]}}(24 bytes) & \([\sigma_x, \sigma_y, \sigma_z]\)
in \(\text{m}\); \([\sigma_\phi, \sigma_\theta, \sigma_\psi]\) in
\(\text{rad}\) & \(\mathbf{\Sigma} \succ 0\) (strictly
positive-definite) and
\(\lambda_{\max}(\mathbf{\Sigma}) \le \sigma_{\text{admit}}^2\)
(statistical compactness) & Diagonal Cholesky decomposition check;
rejects proposals with excessive spatial dispersion \\
\texttt{0x80..0x87}\texttt{t\_expire} & \texttt{uint64\_t}(8 bytes) &
Nanoseconds (\(\text{ns}\)) UTC epoch time &
\(t_{\text{exp}} \ge t_{\text{current}} + t_{\text{min\_exec}}\) and
\(t_{\text{exp}} - t_{\text{src}} \le \tau_{\text{safe}}\) (bounded
finite validity window) & Dedicated hardware countdown watchdog timer
(RTT / WWDT); triggers safe fallback upon underflow \\
\texttt{0x88..0x9F}\texttt{wrench\_limit\_Fmax} &
\texttt{float32{[}6{]}}(24 bytes) & \([F_x, F_y, F_z]\) in \(\text{N}\);
\([\tau_x, \tau_y, \tau_z]\) in \(\text{N}\cdot\text{m}\) &
\(\|F_{xyz}\| \le F_{\text{payload\_cap}}\) and
\(\|\tau_{xyz}\| \le \tau_{\text{motor\_limit}}\) (enforces physical
compliance envelope) & Inverter current saturation clamp + analog
threshold comparator on 6-axis wrist F/T sensor \\
\texttt{0xA0..0xA7}\texttt{terminal\_flags} & \texttt{uint32\_t} +
\texttt{uint8\_t{[}4{]}}(8 bytes) & Bitmask enum + 4-byte zero-padding &
\(\text{Flag} \in \{\text{ACTIVE\_HOLD}, \text{COMPLIANT\_DROP}, \text{DECEL\_STOP}\}\)
(deterministic fallback specification) & Hardware state machine
selector; directly configures motor PWM gate drivers upon lease
expiry \\
\end{longtable}

The nervous system evaluates monotonicity and provenance before passing
an incoming intent record to the trajectory generator. When an
asynchronous vision-language model requires \(250\text{ ms}\) of
inference time, a lighter geometric planner operating at
\(50\text{ Hz}\) may have already emitted a command derived from a newer
camera exposure. By comparing the incoming sequence identifier against
the highest sequence number committed to the actuation pipeline, the
receiver discards out-of-order or replayed packets in constant time. If
\(t_{\text{src}}\) indicates that the perception frame is older than the
current tracking epoch, or if the parent perception hash does not match
an entry in the local sliding ring buffer of verified sensor frames, the
packet is dropped immediately. This check executes in less than
\(1\text{ }\mu\text{s}\) of memory lookup and prevents the machine from
executing delayed commands generated from stale visual scenes.

An intent record governs execution through a five-state finite state
machine comprising the Pending, Active, Renewed, Expired, and Aborted
states. Upon arrival over the real-time bus, the record enters the
Pending state while the nervous system evaluates kinematic reachability
and verifies that the commanded motion remains within the wrench limit
\(F_{\text{max}}\). Once validated, the record transitions to Active,
granting the motor controllers authority to track toward the target pose
\(T\) within the tolerance envelope \(\mathbf{\Sigma}\). If an upstream
model generates a successor intent record with a higher sequence counter
and valid provenance before the local hardware clock reaches
\(t_{\text{exp}}\), the state machine transitions to Renewed, updating
the target pose and extending the expiration timestamp without
interrupting joint velocity profiles. If the mechanism satisfies the
target pose within the tolerance bounds \(\mathbf{\Sigma}\), the lease
completes; if joint torque sensors detect unexpected resistance
exceeding \(F_{\text{max}}\), the state machine transitions immediately
to Aborted, cutting motor authority.

When the local clock reaches \(t_{\text{exp}}\) without receiving a
valid renewal record, the intent transitions deterministically to the
Expired state. Because this transition is evaluated locally in the
nervous system at \(1000\text{ Hz}\), the machine does not depend on an
upstream cancellation message that might never arrive. The terminal
behavior flag in the intent record dictates the exact mechanical action
executed at the exact millisecond of expiration. Under an active hold
policy, the controller clamps the reference trajectory to the measured
position at \(t_{\text{exp}}\) and engages closed-loop position control,
drawing motor current to counteract gravity and hold a payload in place.
Under a compliant drop policy, the controller commands zero torque or
applies velocity damping, allowing kinetic energy to dissipate smoothly
as mechanical brakes engage. The selection is governed by the physical
load. A robotic arm holding a \(5\text{ kg}\) workpiece over an open bin
requires an active hold to prevent dropping the part, whereas a
high-speed spindle requires torque cutoff and damping to prevent
continuous friction heating during a software stall.

The complete deterministic execution loop that unifies intent record
ingestion, cryptographic lineage verification, kinematic workspace
gating, and hardware watchdog lease countdown on real-time
microcontroller silicon is formalized below:

\subsection{Intent Lease Validation and Time-To-Live Watchdog
Protocol}\label{intent-lease-validation-and-time-to-live-watchdog-protocol}

\textbf{Cadence \& Invariants:}

\begin{itemize}
\tightlist
\item
  \emph{Execution Cadence:} Hard real-time periodic interrupt
  (\(f = 1000\text{ Hz}\), period \(T_c = 1.0\text{ ms}\), maximum
  allowable jitter \(< 50\,\mu\text{s}\)).
\item
  \emph{Memory Invariant:} Zero dynamic heap allocation
  (\texttt{malloc}/\texttt{free} strictly forbidden; all state variables
  statically allocated in TCM SRAM).
\item
  \emph{Hardware Authority:} Direct write/evaluate access to hardware
  timer compare registers (\texttt{WWDT} / \texttt{RTT\_VR}), shared IPC
  SRAM mailbox (\texttt{RPMSG} virtio rings), and motor drive PWM enable
  lines (\texttt{TIMx\_CCR} / \texttt{ePWM\ CMPA}).
\end{itemize}

\textbf{Execution Protocol:}

\begin{enumerate}
\def\labelenumi{\arabic{enumi}.}
\item
  \textbf{Cryptographic Provenance \& Monotonic Sequence Check (Stage
  1):} If a new packet \(\mathcal{L}_{\text{cand}}\) is resident in the
  shared SRAM mailbox:

  \begin{itemize}
  \tightlist
  \item
    Sequence monotonicity: if
    \(\mathcal{L}_{\text{cand}}.\text{seq\_id} \le \mathcal{L}_{\text{active}}.\text{seq\_id}\),
    drop packet with \texttt{ERR\_SEQ\_NON\_MONOTONIC}.
  \item
    Perception freshness: if
    \((t_{\text{now}} - \mathcal{L}_{\text{cand}}.t_{\text{src}}) > \Delta t_{\text{max\_age}}\),
    drop packet with \texttt{ERR\_PERCEPTION\_STALE}.
  \item
    Perception lineage: if
    \(\text{Lookup}(\mathcal{L}_{\text{cand}}.\text{parent\_hash}, \mathcal{B}_{\text{sensor\_ring}}) = \text{FALSE}\),
    drop packet with \texttt{ERR\_PARENT\_HASH\_MISMATCH}.
  \item
    HMAC signature: if
    \(\text{HMAC\_SHA256}(\mathcal{L}_{\text{cand}}) \ne \mathcal{L}_{\text{cand}}.\text{hmac}\),
    drop packet with \texttt{ERR\_HMAC\_INVALID}.\sidenote{\footnotesize Hardware-isolated
      SHA-256 HMAC keys stored in secure enclave eFuses protect
      inter-core IPC rings from buffer injection, payload corruption,
      and replay attacks across compromised operating system driver
      domains.}
  \end{itemize}
\item
  \textbf{Static Workspace, Normality, and Occupancy Gating (Stage 2):}

  \begin{itemize}
  \tightlist
  \item
    Check unit quaternion normalization: if
    \(\left\lvert \|\mathcal{L}_{\text{cand}}.\mathbf{q}_{\text{quat}}\|_2^2 - 1.0 \right\rvert > 10^{-6}\),
    drop packet with \texttt{ERR\_QUATERNION\_NON\_NORMAL}.
  \item
    Check kinematic reachability: if
    \(\mathcal{L}_{\text{cand}}.\mathbf{p}_{\text{target}} \notin \mathcal{W}_{\text{reach}}\),
    drop packet with \texttt{ERR\_WORKSPACE\_OUT\_OF\_BOUNDS}.
  \item
    Check static occupancy: if
    \(\text{Intersects}(\mathcal{L}_{\text{cand}}.\mathbf{p}_{\text{target}}, \text{Octree}_{\text{static}})\),
    drop packet with \texttt{ERR\_OCCUPANCY\_BLOCKED}.
  \end{itemize}
\item
  \textbf{Dynamic Reachability \& Covariance Ambiguity Gate (Stage 3):}
  Evaluate spatial covariance dispersion by checking maximum eigenvalue
  \(\lambda_{\max}(\mathcal{L}_{\text{cand}}.\mathbf{\Sigma})\); if
  \(\lambda_{\max} > \sigma_{\text{admit}}^2\), drop packet with
  \texttt{ERR\_GROUNDING\_AMBIGUITY}. Compute minimum transit time under
  trapezoidal profile for displacement
  \(d = \|\mathcal{L}_{\text{cand}}.\mathbf{p}_{\text{target}} - \mathbf{p}_{\text{curr}}\|\):
  \[t_{\min} = \begin{cases} 2\sqrt{d / a_{\max}} & \text{if } d \le v_{\max}^2 / a_{\max} \\ (d / v_{\max}) + (v_{\max} / a_{\max}) & \text{if } d > v_{\max}^2 / a_{\max} \end{cases}\]
  If
  \(t_{\min} > (\mathcal{L}_{\text{cand}}.t_{\text{exp}} - t_{\text{now}})\),
  drop packet with \texttt{ERR\_DYNAMIC\_TIME\_INSUFFICIENT} and return
  margin shortfall \(t_{\min} - \tau_{\text{rem}}\).
\item
  \textbf{Hardware Watchdog Arming \& Trajectory Latch (Stage 4):} If
  Stages 1--3 pass: atomically update
  \(\mathcal{L}_{\text{active}} \leftarrow \mathcal{L}_{\text{cand}}\),
  set
  \(\mathcal{L}_{\text{active}}.\text{state} \leftarrow \text{ACTIVE}\)
  (or \(\text{RENEWED}\)), program countdown timer register
  \(\text{WRITE\_REG}(\text{RTT\_VR}, \mathcal{L}_{\text{active}}.t_{\text{exp}} - t_{\text{now}})\)\sidenote{\footnotesize Peripheral
    registers (ARM Cortex-R52 \texttt{RTT\_VR}, STM32
    \texttt{TIMx\_CCR}, TI ePWM \texttt{CMPA}) provide hardware
    down-counters driven by isolated crystal oscillators. Writing
    microsecond lease durations ensures hardware protection operates
    independently of CPU load or thread preemption.}, and dispatch
  \((\mathbf{p}_{\text{target}}, \dot{\mathbf{p}}_{\text{max}}, \mathcal{L}_{\text{active}}.\mathbf{F}_{\max})\)
  to the quintic trajectory spline generator.
\item
  \textbf{\(1\text{ kHz}\) Watchdog Evaluation \& Deterministic Fallback
  Trigger (Stage 5):} On every millisecond interrupt tick:

  \begin{itemize}
  \tightlist
  \item
    If measured contact wrench
    \(\|\mathbf{F}_{\text{meas}}\| > \mathcal{L}_{\text{active}}.\mathbf{F}_{\max}\)\sidenote{\footnotesize ISO
      10218-1/2 and ISO/TS 15066 define collaborative robot operation
      modes, including Power and Force Limiting (PFL) and Speed and
      Separation Monitoring (SSM), specifying transient and quasi-static
      contact force thresholds across anatomical body regions.}: set
    \(\mathcal{L}_{\text{active}}.\text{state} \leftarrow \text{ABORTED}\)
    and command immediate torque limit clamp or Safe Torque Off.
  \item
    Else if
    \(t_{\text{now}} \ge \mathcal{L}_{\text{active}}.t_{\text{exp}}\)
    (or watchdog underflow interrupt fires): set
    \(\mathcal{L}_{\text{active}}.\text{state} \leftarrow \text{EXPIRED}\),
    and execute the programmed fallback:

    \begin{itemize}
    \tightlist
    \item
      \texttt{ACTIVE\_HOLD}: clamp trajectory setpoint to
      \(\mathbf{p}_{\text{curr}}\) and engage closed-loop position hold.
    \item
      \texttt{DECEL\_STOP}: apply maximum deceleration along current
      path tangent (\(a = -a_{\max}\)).
    \item
      \texttt{COMPLIANT\_DROP}: set active stiffness to zero and apply
      velocity damping.
    \end{itemize}
  \end{itemize}
\end{enumerate}

Because the intent record defines an explicit, self-contained
mathematical contract, verification frameworks can evaluate system
safety without inspecting the weights of the learned model. In
Chapter~\ref{sec-verification}, the fields of the intent record schema
serve as the interface specification for automated runtime assertions
and hardware-in-the-loop fuzz testing. Invariant monitors verify that
the lease duration \(t_{\text{exp}} - t_{\text{src}}\) never exceeds the
lookahead horizon of the safety shield, and that \(\mathbf{\Sigma}\)
remains strictly positive-definite. Fuzzing harnesses inject corrupted
quaternions, expired timestamps, and wrench limits set to zero directly
into the communication ring at \(10\text{ kHz}\), verifying that the
nervous system rejects every non-compliant packet and enters a safe
terminal state within a single \(1\text{ ms}\) control tick.

The intent record guarantees that no goal can command physical actuators
without an active, verified lease, preventing runaway motions caused by
model stalls, network drops, or replayed packets. Yet enforcing strict
temporal and spatial bounds at the interface does not prevent the
physical environment from changing around a valid objective. When an
initially clear path is obstructed by a moving human or a shifting
workpiece, the intent record remains unexpired and mathematically
well-formed while the physical motion toward the target becomes
dangerous.

\section{Systems Stress-Tests}\label{sec-intent-10-8}

The failure mode that breaks a physical AI system is rarely a corrupted
memory buffer or a malformed floating-point payload, but a correct goal
that becomes dangerous without changing. When an upstream
vision-language model emits an intent record, the target coordinates,
spatial tolerances, and allowable contact wrenches accurately reflect
the state of the physical scene at the sampling timestamp
\(t_{\text{src}}\). If the physical world shifts after
\(t_{\text{src}}\), or if the reasoning compute node crashes
mid-trajectory, the intent record remains mathematically valid,
well-formed, and unmutated. The nervous system cannot inspect the
contents of the record to discover that the physical justification for
the motion has evaporated. If the downstream tracking loop treats an
intent as valid until an explicit cancellation arrives, a stalled model
leaves the physical body executing obsolete instructions at full motor
authority. The following four analytical stress-tests evaluate whether
the nervous system enforces temporal and spatial boundaries when
upstream intelligence fails silently:

\subsection*{Stress-Test 1: The Frozen Reasoner (Silent Kernel Deadlocks
and Dead-Man
Brakes)}\label{stress-test-1-the-frozen-reasoner-silent-kernel-deadlocks-and-dead-man-brakes}
\addcontentsline{toc}{subsection}{Stress-Test 1: The Frozen Reasoner
(Silent Kernel Deadlocks and Dead-Man Brakes)}

\emph{The scenario.} An articulated 6-axis manipulator
(\(12.0\text{ kg}\) payload) approaches a rigid aluminum assembly
fixture at \(0.40\text{ m/s}\). The vision-language model executes
asynchronously on a host accelerator, issuing intent records with lease
duration \(\Delta t_{\text{lease}} = 100.0\text{ ms}\). At
\(t = 340.0\text{ ms}\), the host accelerator encounters a kernel
deadlock during matrix multiplication, freezing the pipeline without
emitting error packets.

\emph{The systems dilemma.} If the \(1000\text{ Hz}\) nervous system
extrapolates the unrenewed trajectory past \(t_{\text{exp}}\), the
end-effector impacts the rigid fixture in \(65.0\text{ ms}\), shattering
the harmonic drives. If motor current is cut instantly, the payload
drops under gravity. The nervous system must detect
\(t \ge t_{\text{exp}}\) within \(1.0\text{ ms}\), executing a
controlled deceleration (\(a_{\max} = 2.5\text{ m/s}^2\)) that brings
velocity to zero in \(t_{\text{stop}} = 201.0\text{ ms}\)
(\(d_{\text{stop}} = 50.5\text{ mm}\)) and engages active gravity
holding.

\emph{The systems verdict: Passive Temporal Expiry Law.} Host software
crashes cannot be detected by waiting for network error packets. The
passive expiration of a local hardware timer is the only failure
detector guaranteed to survive the death of an upstream model.

\subsection*{Stress-Test 2: The Bumped Target (Spatial Contract
Invalidation)}\label{stress-test-2-the-bumped-target-spatial-contract-invalidation}
\addcontentsline{toc}{subsection}{Stress-Test 2: The Bumped Target
(Spatial Contract Invalidation)}

\emph{The scenario.} A 7-axis compliant arm lowers a gripper toward a
\(0.50\text{ kg}\) steel bracket. The active intent record specifies
target pose \(\mathbf{p}_{\text{target}}\) with a tolerance radius
\(r_{\text{tol}} = 12.0\text{ mm}\) and a lease duration of
\(120.0\text{ ms}\). At \(t = 45.0\text{ ms}\), conveyor vibration
shifts the bracket \(40.0\text{ mm}\) laterally.

\emph{The systems dilemma.} The intent record remains temporally valid
with \(75.0\text{ ms}\) remaining. The local perception loop evaluates
the spatial displacement condition
\((\mathbf{x}_{\text{measured}} - \mathbf{p}_{\text{target}})^T \mathbf{\Sigma}^{-1} (\mathbf{x}_{\text{measured}} - \mathbf{p}_{\text{target}}) > 1\)
at \(500\text{ Hz}\). Detecting that the physical target has escaped the
tolerance ball invalidates the lease locally within \(2.0\text{ ms}\),
clamping downward contact force below \(8.0\text{ N}\) before striking
the table.

\emph{The systems verdict: Spatial Contract Invariant.} A goal does not
need to be mathematically corrupt to be physically hazardous. Embedding
spatial bounds directly into the intent record allows the nervous system
to reject an obsolete goal the instant the environment moves, without
waiting for slow deliberative re-planning.

\subsection*{Stress-Test 3: Inference Latency Spike under Thermal
Throttling (Velocity
Scaling)}\label{stress-test-3-inference-latency-spike-under-thermal-throttling-velocity-scaling}
\addcontentsline{toc}{subsection}{Stress-Test 3: Inference Latency Spike
under Thermal Throttling (Velocity Scaling)}

\emph{The scenario.} A mobile manipulator carrying a \(4.0\text{ kg}\)
container encounters \(40.0^{\circ}\text{C}\) ambient heat, throttling
the coprocessor so reasoning latency spikes from
\(P_{50} = 75.0\text{ ms}\) to \(P_{99} = 260.0\text{ ms}\) against a
fixed lease of \(\Delta t_{\text{lease}} = 100.0\text{ ms}\). Renewal
packets arrive \(160.0\text{ ms}\) late.

\emph{The systems dilemma.} Tripping hard emergency dynamic brakes at
every expiration causes violent \(4\text{ Hz}\) shuddering that spills
payload. Automatically extending leases violates temporal safety. The
trajectory generator scales reference velocity along a
\(C^2\)-continuous curve as remaining lease \(\tau_{\text{remain}}\)
drops below stopping horizon \(t_{\text{brake}} = v / a_{\max}\),
settling to a stationary hold until fresh updates arrive.

\emph{The systems verdict: Latency Degradation Law.} Inference latency
spikes are physical realities driven by thermal limits and bus
contention. Trajectory generators must decelerate smoothly to a
controlled stationary hold rather than oscillating between emergency
trips and unmonitored execution.

\subsection*{Stress-Test 4: Semantic Distractor (Unimodal Fusion in
Bimodal
Scenes)}\label{stress-test-4-semantic-distractor-unimodal-fusion-in-bimodal-scenes}
\addcontentsline{toc}{subsection}{Stress-Test 4: Semantic Distractor
(Unimodal Fusion in Bimodal Scenes)}

\emph{The scenario.} A sorting robot descends at \(0.30\text{ m/s}\) to
pick a connector. A second identical connector is dropped
\(35.0\text{ mm}\) away. The vision model computes the spatial mean
between the two objects, emitting a target
\(\mathbf{p}_{\text{target}}\) in empty space with covariance expanding
to \(3\sigma = 55.0\text{ mm}\).

\emph{The systems dilemma.} Admitting this intent drives the gripper
into empty space and risks a side collision with both parts. The
real-time admission gate evaluates maximum eigenvalue
\(\lambda_{\max}(\mathbf{\Sigma}) \le \sigma_{\max}^2 = (10.0\text{ mm})^2\)
via Cholesky decomposition (\(42.0\ \mu\text{s}\) latency), rejecting
the ambiguous intent packet in under \(1.0\text{ ms}\) and maintaining a
safe holding state.

\emph{The systems verdict: Covariance Admission Law.} Semantic ambiguity
in upstream neural networks manifests physically as spatial dispersion.
Evaluating covariance at the admission boundary prevents actuators from
executing an averaged trajectory that corresponds to no real object in
the physical world.

\section{Summary}\label{sec-intent-10-9}

A physical AI system cannot permit a learned reasoning model to command
raw actuator trajectories directly. High-level vision-language models
and learned policies operate over nondeterministic inference graphs,
where context lengths, attention heads, and memory bandwidth produce
variable execution latencies. If an upstream model generates
point-by-point joint torques or Cartesian path splines, any scheduling
hiccup or inference pause leaves the physical body coasting on momentum
or executing stale motor commands against an unobserved obstacle. The
intent contract establishes a hard boundary between high-level reasoning
and physical execution by restricting the output of learned models to
bounded, expiring, and kinematically reachable goals.

High-level reasoning determines what target state the machine should
seek, but it delegates how the machine moves to deterministic downstream
control loops. By decoupling goal formulation from motion synthesis, the
system bounds the authority of the neural network. The model cannot
command an arbitrary path through the workspace, nor can it bypass
dynamic limits. Instead, the model issues an intent packet that
encapsulates the desired state alongside explicit admissibility
criteria, leaving the nervous system free to reject any packet that
violates the kinematic limits or safety envelopes of the physical body.

This interface relies on three core mathematical bounds that define the
validity envelope of any goal. First, target geometry is bounded by a
spatial tolerance ball
\(\mathcal{B}_\epsilon(x^*) = \{ x \in \mathbb{R}^3 \mid \|x - x^*\| \le \epsilon \}\),
where the radius \(\epsilon\) defines the admissible error volume around
the target pose \(x^*\). Second, temporal authority is bounded by the
lease duration \(\tau\), establishing an expiration timestamp
\(t_{\text{exp}} = t_{\text{issue}} + \tau\) after which the intent
record becomes null. Third, dynamic feasibility limits constrain the
spatial displacement \(\Delta x = \|x^* - x_0\|\) relative to the
current end-effector pose \(x_0\). For an analytical example of a
manipulator with a maximum allowable acceleration
\(a_{\max} = 2.0\text{ m/s}^2\) and maximum velocity
\(v_{\max} = 0.50\text{ m/s}\), traversing a displacement of
\(\Delta x = 0.20\text{ m}\) requires accelerating to peak velocity in
\(0.25\text{ s}\), cruising for \(0.15\text{ s}\), and decelerating to
rest in \(0.25\text{ s}\), yielding a minimum traversal time
\(t_{\min} = 650\text{ ms}\). If an incoming intent packet specifies a
lease duration \(\tau < 650\text{ ms}\) for that displacement, the
admission boundary rejects the intent before execution begins because
the goal cannot be reached without violating kinematic limits.

The trajectory planner in Chapter~\ref{sec-planning} receives this
intent record as a clean, self-contained specification. The planner
consumes the target coordinate \(x^*\), the tolerance radius
\(\epsilon\), the covariance matrix \(\mathbf{\Sigma}\), and the lease
duration \(\tau\) as numeric constraints. It operates without access to
the prompt tokens, image embeddings, or model weights that produced the
packet. The trajectory planner does not evaluate whether a target
represents a tool, a fastener, or an empty bin, because all semantic
requirements were already converted into spatial envelopes, velocity
limits, and contact force thresholds at the intent boundary. By
stripping semantic ambiguity from the real-time planning pipeline, the
system ensures that the trajectory generator solves a purely geometric
and dynamic optimization problem governed by deterministic execution
budgets.

\emph{The physical AI insight.} The underlying design rule for physical
safety in learned systems is that protection depends on structural
temporal leases rather than active software cancellations. In classical
distributed software, a supervising process detects an anomaly and
transmits an abort signal to interrupt a running thread. In physical AI,
relying on an active cancellation message introduces a dangerous failure
mode, because the very fault that causes the learned model to produce
incorrect actions (such as an out-of-distribution visual artifact, an
operating system thread lock, or an out-of-memory exception) often halts
the execution path that would otherwise issue the abort. A machine that
requires an explicit cancellation signal to halt will continue driving
its actuators until a mechanical stop or a physical collision arrests
the mechanism.

Structuring intent as a decaying temporal lease reverses this failure
vulnerability. A lease is an affirmative grant of execution authority
that automatically dissolves unless an uncorrupted, validated update
arrives before the deadline. When an upstream model hangs, stalls, or
fails validation, the expiration timestamp passes without renewal, and
the nervous system automatically triggers a controlled deceleration to a
zero-velocity state. Safety is preserved through the passive passage of
time rather than the active dispatch of a software message.

\protect\phantomsection\label{callout-takeawaysux2a-10.8}
\begin{fbxSimple}{callout-takeaways}{Key Chapter Takeaways}{}
\phantomsection\label{callout-takeaways*-10.8}

\begin{itemize}
\tightlist
\item
  \textbf{Intent Packages Emit Bounded Goals, Never Raw Actuator
  Trajectories:} Foundation vision-language models have variable
  inference latencies. Restricting high-level outputs to bounded target
  poses (\(x^*\)), tolerance balls (\(\mathcal{B}_\epsilon(x^*)\)), and
  lease durations (\(\tau_{\text{lease}}\)) prevents scheduling hiccups
  from translating into kinematic momentum runaway.
\item
  \textbf{Dynamic Feasibility Gating Precedes Kinematic Execution:} A
  target pose displacement \(\Delta x = \|x^* - x_0\|\) requires minimum
  kinematic traversal time
  \(t_{\min} = \Delta x / v_{\max} + v_{\max}/a_{\max}\). Intent leases
  specifying \(\tau_{\text{lease}} < t_{\min}\) are kinematically
  infeasible and must be rejected at the admission gate.
\item
  \textbf{Safety Requires Expiring Leases Rather Than Active Abort
  Messages:} In physical systems, software faults that cause incorrect
  commands (thread deadlocks, memory leaks, VLM stalls) kill the very
  execution path that would dispatch an abort message. Protection must
  rely on positive temporal leases that automatically expire upon
  silence.
\item
  \textbf{Semantic Ambiguity is Stripped at the Intent Boundary:}
  Trajectory planners and real-time safety referees evaluate pure
  geometry, covariance bounds, and dynamic constraints
  (\(x^*, \epsilon, \mathbf{\Sigma}, \tau\)), remaining entirely
  decoupled from natural language tokens, embeddings, or model weights.
\item
  \textbf{SayCan Affordance Gating Prevents Unreachable Commands:} High
  semantic probability from an LLM must be multiplicatively gated by
  physical affordance value functions
  (\(P = P_{\text{VLM}} \times V_{\text{affordance}}(s)\)) to prevent
  plausible-sounding goals from commanding kinematically singular
  configurations.
\end{itemize}

\end{fbxSimple}

Chapter~\ref{sec-planning} takes this bounded, expiring goal and
constructs the continuous time-parameterized trajectories that guide the
body to its target within the granted lease.

\chapter{\texorpdfstring{Kinodynamic
Planning\chaptersubtitle{Feasible Paths, and What Happens When the Next Is Late}}{Kinodynamic Planning}}\label{sec-planning}

\begin{figure}[H]

\centering{

\includegraphics[width=1\linewidth,height=\textheight,keepaspectratio]{index_files/mediabag//Users/VJ/GitHub/MLSysBook-integrate-physical/physical/book/chapters/11-planning/figures/fig_locator.pdf}

}

\caption{\label{fig-locator}\textbf{Systems Organ Focus: Brain.} In the
physical AI systems architecture, the deliberative brain proposes
time-parameterized action chunks across an unprivileged boundary to the
deterministic nervous system. Because neural inference latencies exhibit
heavy-tailed stochastic distributions, every plan must define its own
replacement timing, boundary derivatives, and precomputed stopping
suffix before execution begins.}

\end{figure}%

\emph{What does a physical machine do during the unbudgeted milliseconds
when the next neural trajectory chunk is late?}

The raw output of a neural policy is often a discrete sequence of
waypoints or a stochastic action chunk sampled from a diffusion model.
If these discrete references are streamed directly to joint motor
drives, the resulting step discontinuities command infinite jerk,
exciting structural resonances, shearing gearboxes, and tripping
overcurrent inverters. Trajectory planning in physical AI is the
discipline of transforming discrete neural proposals into continuous,
dynamically feasible trajectory splines that respect the physical limits
of the plant. A dependable planner must guarantee \(\mathcal{C}^2\) jerk
smoothness, enforce receding-horizon spline blending across chunk seams,
and execute precomputed emergency stopping suffixes whenever upstream
neural inference is delayed.

\newpage{}

\protect\phantomsection\label{callout-objectiveux2a-11.1}
\begin{fbxSimple}{callout-objective}{Learning Objectives}{}
\phantomsection\label{callout-objective*-11.1}

After studying this chapter, you will be able to:

\begin{enumerate}
\def\labelenumi{\arabic{enumi}.}
\tightlist
\item
  Set a planning horizon from an inference latency distribution rather
  than a round number.
\item
  Predict what a discontinuity at a chunk boundary does to a
  transmission.
\item
  Specify the behavior of a plan whose replacement never arrives, and
  defend it as part of the plan.
\item
  Identify which trajectory failures are visible before execution and
  which appear only during it.
\end{enumerate}

\emph{Core chapter takeaway.} A plan must define what happens when its
replacement is late, because the replacement will be late. The seam
between one plan and the next is where mechanisms break, not the middle
of either.

\end{fbxSimple}

\section{Continuous Trajectories}\label{sec-planning-11-1}

\lettrine{I}{n traditional machine learning,}\lettrine{c}{omputation}
operates behind glass. In the digital domain of large language models
and computer vision pipelines, compute is idempotent: if a model
requires an extra \(100\text{ ms}\) to evaluate an attention layer or
decode a token, the user experiences a minor perceptual lag, but the
underlying data structures remain intact. Time can be paused, frames can
be dropped, and operations can be retried without physical consequence.
In physical AI, that assumption collapses. At the wrist of a robotic
manipulator lowering a cutting tool onto a steel workpiece, the joint
tracking loop executes deterministically on a Real-Time Microcontroller
Unit (MCU) every \(10\text{ ms}\). At this \(100\text{ Hz}\) cadence
(with inner field-oriented current loops cycling at \(1000\text{ Hz}\)),
the nervous system queries joint optical encoders, evaluates closed-loop
PID or computed torque control laws, and writes pulse-width modulation
(PWM) current setpoints to the motor phase inverters. The physical body
possesses inertia \(J\) and kinetic energy
\(E_k = \frac{1}{2} J \dot{\mathbf{q}}^2\). Physical time cannot pause
for a slow memory bus or an operating system scheduling interrupt. If
the tracking loop receives a continuous, smooth stream of position and
velocity references, the tool makes contact with controlled force and
begins its traverse across the surface.

This temporal invariant governs machines across all Three Physical
Machine Archetypes:

\begin{enumerate}
\def\labelenumi{\arabic{enumi}.}
\tightlist
\item
  \textbf{Mobile Systems (AMRs / AGVs / Quadrupeds)}: An autonomous
  mobile robot navigating a high-density logistics corridor at
  \(v_0 = 1.5\text{ m/s}\) requires continuous steering angle and wheel
  torque setpoints every \(10\text{ ms}\). If an inference engine stalls
  for \(150\text{ ms}\), the unguided chassis drifts through
  \(d = v_0 \Delta t = 225\text{ mm}\) of dead reckoning, exhausting its
  spatial obstacle clearance and inducing severe tire slip if emergency
  braking is clamped abruptly.
\item
  \textbf{Manipulators (Articulated Arms / High-Speed Delta Pickers)}: A
  6-DOF industrial arm sweeping a \(4.0\text{ kg}\) payload at
  \(2.0\text{ m/s}\) relies on model-based feedforward torque
  \(\boldsymbol{\tau}_{\mathrm{ff}} = \mathbf{M}(\mathbf{q})\ddot{\mathbf{q}} + \mathbf{C}(\mathbf{q}, \dot{\mathbf{q}})\dot{\mathbf{q}} + \mathbf{g}(\mathbf{q})\).
  An unexpected missing reference forces low-level joint controllers to
  hold stale currents against gravity or trip mechanical spring brakes,
  producing shock torques that shear gearbox drive keys.
\item
  \textbf{Process/Energy Systems (Thermal Extrusion / Chemical Batch
  Reactors / Gas Turbines)}: In polymer fused-filament 3D printing or
  continuous flow reactors, mass feed rates and heater thermal powers
  must obey strict continuous derivative bounds
  (\(d\dot{m}/dt \le \dot{M}_{\max}\),
  \(d\dot{Q}/dt \le \dot{Q}_{\max}\)). A step discontinuity in commanded
  mass flow creates instantaneous nozzle melt-pressure surges
  (\(\Delta P = 4.0\text{--}8.0\text{ MPa}\)) that blow out barrel seals
  or fracture nozzle orifices.\sidenote{\footnotesize In continuous manufacturing,
    setpoints control mass feed valves (\(\dot{m}\)) or heaters (\(P\)).
    Unlike rigid joints, thermal nozzles exhibit transport lag
    (\(\tau = L/v\)). Slew limiters (\(d\dot{m}/dt \le \dot{M}_{\max}\))
    prevent hydraulic water hammer and thermal shock cracking.}
\end{enumerate}

Above this hard real-time execution loop sits the Brain tier of the
architectural hierarchy (Figure~\ref{fig-locator}), hosted on a
high-throughput Linux Microprocessor Unit (MPU) or embedded neural
accelerator (such as a canonical heterogeneous dual-brain System-on-Chip
(SoC) architecture). Here, learned policies---such as Action Chunking
with Transformers (ACT) (\citeproc{ref-zhao2023learning}{Zhao et al.
2023}) or Diffusion Policy (\citeproc{ref-chi2023diffusionpolicy}{Chi et
al. 2025})---process camera frames, point clouds, and tactile signals to
propose reference paths. Unlike the deterministic, microsecond-accurate
periodic cycle of the motor drives on the real-time MCU, policy
inference on the Linux MPU exhibits variable execution time. A forward
pass that completes in \(35\text{ ms}\) under nominal conditions
stretches to \(120\text{ ms}\) when image compression buffers clear, and
extends beyond \(200\text{ ms}\) when operating system context switches,
garbage collection pauses, or memory bus contention delay tensor
execution on the neural coprocessor. The real-time tracking loop never
waits for the upstream model to finish. It consumes whatever reference
is resident in its execution buffer when its hardware timer fires,
regardless of whether a fresh neural computation has arrived.

When an architecture generates only a single control step ahead
(\(H = 1\)), this timing mismatch immediately breaks physical control.
If a policy predicts a motor action for \(t + 10\text{ ms}\) but
requires \(25\text{ ms}\) of compute time to evaluate the network, the
output arrives \(15\text{ ms}\) after the target instant has already
passed into history. The real-time tracking loop arrives at
\(t + 10\text{ ms}\) with an empty or expired reference buffer. If the
low-level MCU controller holds the previous motor currents, the
manipulator accelerates unguided under gravitational and external
contact loads. If it clamps the mechanical brakes or zeros the drive
torques, the sudden step change in acceleration exerts massive shock
loads on the gearboxes and dislodges the workpiece. A single-step
prediction contract implicitly assumes that computation is
instantaneous, an assumption that the physical execution pipeline
violates on every single cycle.

To bridge this temporal chasm between slow, variable neural inference on
the MPU and high-frequency actuator control on the MCU, planning systems
output a trajectory spanning multiple future control intervals. This
sequence of future references is an \textbf{action chunk}---a finite,
time-parameterized sequence of future reference states or control
commands \([\mathbf{u}_t, \mathbf{u}_{t+1}, \dots, \mathbf{u}_{t+H-1}]\)
emitted as an atomic block by an upstream neural policy across horizon
\(H\), buffering the nervous system against inference latency variance.
In an analytical example, suppose a learned planner generates a
trajectory across a horizon of \(H = 200\text{ ms}\), providing \(20\)
discrete setpoints at \(10\text{ ms}\) intervals. At
\(t = 0\text{ ms}\), the tracking loop on the MCU starts executing the
first setpoint of this plan. The system cannot wait until the plan
finishes to calculate the next one, so it dispatches a replacement
request across the inter-processor communication (IPC) bus at
\(t_{\text{req}} = 50\text{ ms}\). If the nominal inference latency is
\(t_{\text{lat}} = 60\text{ ms}\) at \(P_{50}\), the replacement arrives
at \(t_{\text{arr}} = 110\text{ ms}\). The new plan arrives
\(90\text{ ms}\) before the original plan reaches its expiry time of
\(t_{\text{exp}} = 200\text{ ms}\), providing an overlap during which
the nervous system can splice the new trajectory into the current
execution path.

The nominal arrival time describes only the center of the distribution,
and physical systems fail in the tail. If memory contention, thermal
throttling, or bus contention extends the replacement latency to
\(P_{99} = 180\text{ ms}\), the replacement arrives at
\(t = 230\text{ ms}\). Because the existing plan expires at
\(t_{\text{exp}} = 200\text{ ms}\), an uncovered interval of
\(30\text{ ms}\) opens between the exhaustion of the valid trajectory
and the arrival of its successor. Lengthening the chunk horizon to
\(1000\text{ ms}\) increases the time buffer, but it forces the body to
execute actions planned on stale perceptual measurements. If the
workpiece shifts during that second, the stale trajectory drives the
tool into collision. Action chunking buys finite time, but it does not
eliminate the tail of the latency distribution or guarantee that a
replacement will arrive before the horizon expires.

As captured in Figure~\ref{fig-locator}, a plan is therefore not a
direct command issued to the joints, nor is it a guaranteed physical
path. It is an unprivileged proposed trajectory submitted across the
MPU/MCU boundary to the tracking controller and the real-time safety
enforcer, valid only over its stated temporal horizon. Because every
computing platform exhibits a nonzero tail of delayed inference, every
plan must define its own \textbf{seam behavior}. The seam is the
terminal transition that governs the actuators if the expiry time passes
without an accepted replacement. A well-formed plan specifies whether
the machine should decelerate to a controlled stop along the path
tangent, switch to compliant force regulation, or hold a static
position. The seam is where physical mechanisms break when computation
stalls, not the middle of a smoothly executing trajectory.

The planner is responsible for generating kinematic paths that the body
can physically follow, managing the timing windows of action chunks, and
specifying deterministic termination states for the moments when
replacements are late. Whether those proposed trajectories violate
environmental keep-out zones, exceed dynamic load limits, or remain safe
under unexpected obstacles is an independent determination. The planner
proposes what the body might do and defines how the trajectory ends when
computation fails. The safety enforcer independently evaluates whether
that proposal is permitted to reach the actuators
(Chapter~\ref{sec-enforcement}).

\section{What Planning Hands Over}\label{sec-planning-11-2}

A geometric path describes a sequence of configurations in space,
parameterized by arc length \(s \in [0, S]\) or a dimensionless
progression variable \(\sigma \in [0, 1]\), without reference to
time.\sidenote{\footnotesize Khatib (1985) formulated real-time obstacle avoidance via
  artificial potential fields, modeling obstacles as repulsive barriers
  and goals as attractive wells. While computationally efficient for
  reactive steering, potential fields lack temporal parameterization and
  cannot guarantee smooth derivative continuity across replanning seams.}
A path defines where the machine travels, but it contains no information
about velocity, acceleration, or actuation limits. To become an
\textbf{executable trajectory}, the geometric path must be mapped to
time,
\(t \mapsto \mathbf{x}(t) = [\mathbf{q}(t)^T, \dot{\mathbf{q}}(t)^T, \ddot{\mathbf{q}}(t)^T]^T\),
assigning concrete physical states, kinematic derivatives, or actuator
targets to explicit timestamps. Only when parameterized by time does the
path expose the rates, accelerations, and temporal deadlines that the
nervous system must evaluate. The data structure handed over from the
upstream planner to the execution system is therefore never an abstract
curve or an ungrounded tensor. It is a time-indexed \textbf{trajectory
packet} delivered across the MPU/MCU boundary to two distinct consumers:
the tracking controller and the safety enforcer. The tracking controller
consumes the trajectory to interpolate high-rate actuator setpoints
(\(1\text{ kHz}\)) and compute model-based feedforward torques. The
safety enforcer consumes the trajectory to independently verify that the
planned motion remains within physical, thermal, and environmental
boundaries before any command reaches the motor phase inverters
(Chapter~\ref{sec-enforcement}).

Every field in the trajectory handoff corresponds to a specific check
that the tracker or enforcer must execute. The first requirement is
spatial and temporal alignment. The trajectory packet must explicitly
declare its coordinate reference frame identifier \(\mathcal{F}\), its
generation timestamp \(t_0\), its planned initial state
\(\mathbf{x}_0\), and its validity interval
\([t_{\text{start}}, t_{\text{exp}}]\). Without an explicit coordinate
frame, the tracking controller cannot transform end-effector positions
and velocities into joint coordinates. Without a generation timestamp
and initial state, the enforcer cannot measure pipeline transport
latency or detect state drift. Consider an inference pipeline that
samples the robot state at \(t = 0\text{ ms}\), requires
\(60\text{ ms}\) to compute a trajectory chunk on the Linux MPU, and
delivers a trajectory starting at \(\mathbf{x}_0\). If the physical
mechanism was in motion during inference and drifted by
\(20\text{ mm}\), passing the trajectory directly to the tracker causes
an instantaneous step error in the control loop, commanding peak torque
and inducing severe mechanical jerk. By comparing the planned initial
state \(\mathbf{x}_0\) against the true state measured by MCU encoders
at \(t_{\text{start}}\), the enforcer can reject the chunk if the
tracking discrepancy exceeds an allowed tolerance, or the tracker can
synthesize a short transition spline to bridge the position offset. The
validity interval defines the temporal window over which the planner
guarantees internal consistency, setting the exact expiration time
\(t_{\text{exp}}\) when the chunk ceases to be authoritative.

The handoff represents a feasibility claim under declared body limits
rather than a proof that the machine can execute the trajectory under
every future disturbance. When a planner generates a trajectory, it
verifies that the requested joint velocities \(\dot{\mathbf{q}}(t)\),
accelerations \(\ddot{\mathbf{q}}(t)\), and torques
\(\boldsymbol{\tau}(t)\) fall within nominal operating bounds under an
assumed dynamic model. It cannot account for unmodeled friction,
external payloads, or sudden contact forces that arise during execution.
The safety enforcer must independently evaluate this feasibility claim
against conservative runtime limits before permitting the tracker to
execute the trajectory. If a learned planner generates a trajectory
assuming an empty gripper, demanding \(80\%\) of peak actuator torque
during high-speed acceleration, and the physical gripper is carrying a
\(4.0\text{ kg}\) payload, executing the trajectory would demand
\(130\%\) of rated motor torque. The resulting current draw would trip
thermal protection or cause severe tracking lag. The enforcer uses the
explicit state derivatives and required forces in the handoff to reject
plans that exceed safe margins under the actual physical configuration.

To allow continuous transitions between consecutive planning cycles, the
trajectory must specify boundary conditions at both ends of its temporal
horizon. A trajectory that guarantees only continuous position creates
physical damage at the replacement seam. If the planner produces a
trajectory whose initial velocity \(\dot{\mathbf{q}}(t_{\text{start}})\)
does not match the true velocity of the body at the moment of handoff,
the resulting step change in velocity requires infinite acceleration,
causing torque saturation and structural vibration. Even if velocity is
continuous, a step discontinuity in acceleration causes a step change in
motor torque, producing severe jerk that excites mechanical resonances
in gearboxes and linkages. The handoff must therefore enforce
second-order derivative continuity (\(C^2\) continuity) at the seam,
matching position, velocity, and acceleration to the terminal state of
the currently executing trajectory chunk. Specifying terminal boundary
conditions at \(t_{\text{exp}}\) similarly ensures that if a successor
plan arrives on time, the new trajectory can splice directly into the
existing motion without arresting the body.

Because feasibility depends on physical context, the handoff must
include the explicit assumptions under which the trajectory was
generated. These metadata fields include the assumed payload mass and
inertia, the expected contact state (such as free-space motion, sliding
contact, or rigid fixture engagement), the allocated actuation budget,
and the minimum spatial clearance to known obstacles. The safety
enforcer compares these declared assumptions against independent sensor
measurements in real time. If the handoff specifies a free-space motion
with an expected contact force of \(0\text{ N}\), but an onboard
six-axis force-torque sensor registers an unexpected load of
\(25\text{ N}\) along the tool axis, the enforcer immediately
invalidates the trajectory. The physical premise of the plan has been
violated, and continuing along the precomputed path risks driving the
tool into a rigid obstruction.

Every computing architecture exhibits latency tails where replacement
plans arrive late or fail to arrive entirely. An executable handoff must
therefore provide an explicit, finite \textbf{precomputed stopping
suffix} that governs the body if the clock reaches \(t_{\text{exp}}\)
without a valid successor. A precomputed stopping suffix is a
dynamically feasible deceleration trajectory resident in execution
memory upon chunk admission that safely brings the physical plant to a
controlled resting state along the path tangent if replanning stalls or
successor plans are exhausted. Leaving terminal behavior unspecified or
defaulting to an instantaneous position hold is physically dangerous. If
a machine traveling at \(1.2\text{ m/s}\) encounters an expired plan and
defaults to holding its final position sample, the tracking controller
attempts to bring the velocity to zero instantaneously. This step
command saturates actuator braking torque and can cause mechanical
transmission damage or tip-over. The terminal suffix is a precomputed
deceleration trajectory that begins at \(t_{\text{exp}}\) and brings the
machine to a safe resting state along the path tangent over a finite
stopping time \(T_{\text{stop}}\). Depending on the task and mechanical
structure, this terminal behavior may decelerate to a stop along the
path tangent, switch actuator control to passive damping, or descend to
a mechanically stable resting posture. The critical systems requirement
is that the terminal suffix is budgeted, checked for dynamic
feasibility, and resident in memory before execution of the chunk
begins.

On real-time embedded hardware, this handoff is realized inside an
\textbf{embedded execution buffer} resident in the MCU's static RAM
(SRAM).\sidenote{\footnotesize In dual-brain SoCs, exchanging trajectory packets
  across non-coherent shared SRAM requires memory barriers
  (\texttt{DMB\ ISH}, \texttt{DSB\ ISH}). The MPU cleans caches
  (\texttt{DCCMVAC}) before triggering mailbox interrupts
  (\texttt{HSEM}). Real-time MCUs evaluate packets in TCM and swap
  buffer pointers in \(<1\ \mu\text{s}\).} To decouple the asynchronous
Linux MPU from the microsecond-accurate MCU servo loop, the execution
engine employs a lockless double-buffering or ring-buffer architecture
with atomic status registers (\texttt{BUFFER\_EMPTY},
\texttt{BUFFER\_READY}, \texttt{BUFFER\_ACTIVE},
\texttt{BUFFER\_EXPIRED}). When the upstream policy on the MPU
synthesizes a new action chunk, it populates the inactive background
buffer over shared memory, direct memory access (DMA), or high-speed
SPI. At the admission boundary, the MCU safety enforcer performs an
\(O(1)\) algorithmic validation check on the candidate chunk. If the
packet passes all invariant tests, the MCU performs an atomic pointer
swap at the next servo tick, seamlessly committing the new trajectory
chunk to the active tracking loop without heap allocation, thread
synchronization locks, or timing jitter.

Structuring the plan handoff as a bounded, hardware-grounded contract
shifts the burden of safety and timing from the non-deterministic
planner to the deterministic nervous system. The tracking controller
receives the dense state targets, timestamps, and feedforward terms
required to drive the actuators along a smooth path. The safety enforcer
receives the generation metadata, dynamic assumptions, validity bounds,
and terminal suffix needed to guard the machine against latency tails,
state divergence, and physical model mismatches. When each field in the
handoff corresponds to an explicit runtime verification, the nervous
system can execute high-speed motions without relying on the assumption
that the next inference cycle will always succeed.

\section{From Goal to Trajectory}\label{sec-planning-11-3}

An intent names a desired physical outcome---such as touching a
workpiece with a specified contact force or placing an end-effector at a
target pose (Chapter~\ref{sec-intent})---but an actuator cannot execute
a goal. A goal defines where the machine should eventually arrive or
what constraint it should satisfy, without providing the joint
velocities, torque bounds, intermediate kinematic states, or continuous
time schedule required to move physical mass through space. The nervous
system requires a trajectory, which assigns concrete state targets,
feedforward actuation terms, and explicit timestamps across a bounded
temporal interval. If a learned policy could evaluate multimodal state
observations and output valid motor currents instantaneously at the
hardware control period (\(1000\text{ Hz}\)), planning as an
architectural stage would not exist; the machine would operate purely as
a high-frequency feedback controller. In physical systems, learned
neural policies require tens or hundreds of milliseconds to compute a
motion sequence, while the tracking controller and safety enforcer must
execute deterministically every few milliseconds. A plan exists
precisely because a goal cannot be executed directly by hardware and
neural inference operates far slower than the cadence of the control
loop.

To span this temporal gap between slow model inference and
high-frequency actuator control, modern robot learning architectures
formulate policy execution via \textbf{receding-horizon action
chunking}.\sidenote{\footnotesize Bellman (1957) formulated the Principle of
  Optimality: decisions from any state must be optimal regardless of
  prior states. In Physical AI, receding-horizon planning repeatedly
  reapplies this principle over rolling windows, balancing trajectory
  optimality against stochastic compute latency.} Two dominant paradigms
characterize this generation:

\begin{enumerate}
\def\labelenumi{\arabic{enumi}.}
\tightlist
\item
  \textbf{Action Chunking with Transformers (ACT)}: Formulated as a
  Conditional Variational Autoencoder (CVAE) paired with a Transformer
  encoder-decoder (\citeproc{ref-zhao2023learning}{Zhao et al. 2023}),
  ACT predicts a discrete sequence of future joint setpoints
  \(\mathbf{A} = [\mathbf{a}_1, \dots, \mathbf{a}_H] \in \mathbb{R}^{H \times D}\)
  simultaneously across a prediction horizon \(H\). To mitigate
  compounding trajectory drift across successive chunks, ACT applies
  \emph{temporal ensembling}, maintaining an exponential moving average
  over overlapping predicted chunks (\(w_i = \exp(-m \cdot i)\)). While
  effective in simulation and low-speed teleoperation, temporal
  ensembling averages discrete coordinate samples without enforcing
  derivative continuity, which can inject high-frequency derivative
  chatter into real-time motor drives.
\item
  \textbf{Diffusion Policy}: Formulating trajectory synthesis as an
  iterative score-based denoising process
  (\citeproc{ref-chi2023diffusionpolicy}{Chi et al. 2025}), Diffusion
  Policy rolls out future action trajectories by solving a score
  Stochastic Differential Equation (SDE) or Denoising Diffusion Implicit
  Model (DDIM) ODE conditioned on visual and proprioceptive histories. A
  diffusion policy predicts a chunk of horizon \(T_p\) (e.g., \(16\)
  steps), executes a committed subset of steps \(T_a\) (e.g., \(8\)
  steps), and replans on subsequent sensor ingress
  (\textbf{?@fig-11-action-chunking-trace}). While diffusion models
  capture rich multimodal distributions, their iterative denoising
  schedules (\(K = 8\text{--}16\) steps) introduce significant inference
  latency and stochastic latency variance.
\end{enumerate}

From an ML systems perspective, the parameter governing the handoff
across these models is the \textbf{end-to-end replacement latency},
denoted as random variable \(L \sim \mathcal{D}(L)\). Treating model
execution time on an isolated accelerator as the planning latency is an
engineering error that undercounts the true physical delay. End-to-end
replacement latency begins at the physical timestamp of the sensory
observation used to condition the planner and spans six pipeline stages:
\[L = t_{\mathrm{sense\_ingress}} + t_{\mathrm{ipc\_dma}} + t_{\mathrm{infer}} + t_{\mathrm{deser}} + t_{\mathrm{validate}} + t_{\mathrm{admit}}\]
where \(t_{\mathrm{sense\_ingress}}\) includes camera exposure and frame
grabber serialization (\(15\text{ ms}\)), \(t_{\mathrm{ipc\_dma}}\) is
direct memory access transfer to the accelerator (\(5\text{ ms}\)),
\(t_{\mathrm{infer}}\) is neural forward pass or diffusion denoising
(\(65\text{ ms}\)), \(t_{\mathrm{deser}}\) is trajectory deserialization
into engineering units (\(2\text{ ms}\)), \(t_{\mathrm{validate}}\) is
kinematic and constraint verification by the safety enforcer
(\(18\text{ ms}\)), and \(t_{\mathrm{admit}}\) is atomic buffer
admission into the MCU execution queue (\(1\text{ ms}\)). In this
representative deployment, while raw inference consumes
\(65\text{ ms}\), the true end-to-end replacement latency is
\(L = 106\text{ ms}\).

Because operating system scheduling, memory bus contention, and
accelerator thermal throttling introduce variability into every stage,
replacement latency exhibits an empirical distribution with a pronounced
tail. If the tracking controller consumes the final state sample of an
active chunk before a validated replacement is admitted into the
embedded buffer, the queue suffers starvation, forcing the nervous
system to abandon nominal tracking and execute its precomputed
deceleration suffix. Sizing the plan horizon against a conventional tail
percentile such as \(P_{99}\) without accounting for mission exposure
guarantees regular physical interruptions. Consider an articulated arm
performing an uninterrupted \(10\text{ minute}\) (\(600\text{ s}\))
manipulation task where new trajectory chunks are requested every
\(200\text{ ms}\), creating \(K = 3000\) consecutive handoffs across the
mission. If the trajectory horizon accommodates only the \(P_{99}\)
latency, the probability of completing the mission without exhausting
the buffer at any seam is
\((1 - 0.01)^{3000} \approx 7.2 \times 10^{-14}\). On average, the
machine encounters thirty late seams during a single ten-minute run,
repeatedly dropping into emergency deceleration. To achieve an
acceptable mission exhaustion failure probability
\(\epsilon_{\mathrm{mission}}\), the allowable per-handoff failure
probability \(\epsilon\) must satisfy
\((1 - \epsilon)^K = 1 - \epsilon_{\mathrm{mission}}\), which simplifies
to \(\epsilon \approx \epsilon_{\mathrm{mission}} / K\) when
\(\epsilon_{\mathrm{mission}} \ll 1\). For a mission failure budget of
\(\epsilon_{\mathrm{mission}} = 10^{-3}\) across \(3000\) handoffs, the
allowable per-plan exhaustion probability is
\(\epsilon \approx 3.33 \times 10^{-7}\), requiring the system to size
its horizon against the extreme tail quantile
\(Q_{1-\epsilon}(L) \approx P_{99.99997}\).

Determining the length of each trajectory chunk requires scheduling when
the replacement request is dispatched relative to the start of the
current chunk. Requesting the next chunk immediately upon entering the
active chunk generates motor plans from outdated sensor readings,
accumulating state drift. The system therefore delays the dispatch until
an execution lead time \(T_{\mathrm{request}}\) has elapsed within the
active chunk. From the instant of dispatch, the remaining execution time
available in the active buffer must accommodate the full replacement
latency at the target quantile \(Q_{1-\epsilon}(L)\), plus a
deterministic reserve \(T_{\mathrm{reserve}}\) dedicated to trajectory
splicing, safety validation, and timer jitter. If the control loop
operates at a period of \(T_c\), the minimum number of discrete
trajectory samples \(N\) that each generated chunk must contain is:
\[N = \left\lceil \frac{T_{\mathrm{request}} + Q_{1-\epsilon}(L) + T_{\mathrm{reserve}}}{T_c} \right\rceil\]
Every parameter in this equation must be measured under maximum
representative system load, including saturated direct memory access
channels, peak bus traffic from adjacent sensor streams, and thermal
throttling of the compute substrate.

In an analytical example for a robotic arm executing a sequence that
touches a surface and then slides a fixture, the tracking controller
runs at \(100\text{ Hz}\) (\(T_c = 10\text{ ms}\)). Empirical profiling
under full sensor and bus contention yields a median latency of
\(65\text{ ms}\), a \(P_{95}\) latency of \(110\text{ ms}\), and a tail
quantile \(Q_{1-\epsilon}(L) = 240\text{ ms}\) corresponding to the
mission reliability budget. The safety enforcer requires an admission
and validation reserve \(T_{\mathrm{reserve}} = 30\text{ ms}\) to
evaluate torque limits and geometric clearance. To incorporate updated
contact force readings gathered during the initial touch, the dispatch
delay is set to \(T_{\mathrm{request}} = 80\text{ ms}\). Sizing the
chunk according to the tail distribution yields:
\[N = \left\lceil \frac{80\text{ ms} + 240\text{ ms} + 30\text{ ms}}{10\text{ ms}} \right\rceil = \left\lceil \frac{350\text{ ms}}{10\text{ ms}} \right\rceil = 35\text{ samples}\]
If an engineer selects a round-number horizon of \(200\text{ ms}\)
(\(N = 20\text{ samples}\)) based on the median latency of
\(65\text{ ms}\) plus an informal margin, the budget fails under real
execution. With \(T_{\mathrm{request}} = 80\text{ ms}\), the time
remaining in the active chunk after dispatch is
\(200\text{ ms} - 80\text{ ms} = 120\text{ ms}\). Subtracting the
\(30\text{ ms}\) reserve leaves an allowable latency budget of only
\(90\text{ ms}\). In this analytical latency distribution, replacement
latency exceeds \(90\text{ ms}\) in approximately \(10\%\) of all cycles
(the \(P_{90}\) threshold). As a consequence, the robot runs out of
trajectory samples once every ten seams, causing the nervous system to
abort tracking, apply joint braking, and shudder to a halt mid-motion.

\protect\phantomsection\label{callout-ruleux2a-11.2}
\begin{fbxSimple}{callout-rule}{Architect’s Rule of Thumb:}{The Tail Latency Horizon Invariant}
\phantomsection\label{callout-rule*-11.2}

Size the action chunk horizon \(N\) against the mission tail quantile
\(Q_{1-\epsilon_{\mathrm{mission}}/K}(L)\) measured under maximum bus
and memory contention, never against the median \(P_{50}\) plus an
informal round-number margin.

\end{fbxSimple}

When a generative policy serves as the planner, iterative score
evaluation and denoising steps govern the shape of this latency
distribution. The internal mechanics of diffusion schedules, attention
caching, and token generation sit below the first budgetable quantity of
the physical AI system. The systems interface treats the learned model
strictly as a stochastic producer of motion chunks characterized by an
empirical latency profile and an output horizon \(N\). In physical
hardware deployments of receding-horizon action chunking, the robot
predicts a trajectory chunk of horizon \(T_p\), executes a committed
subset of steps \(T_a\), and replans on subsequent sensor ingress,
allowing the controller to absorb inference latency while reacting
dynamically to disturbances (\textbf{?@fig-11-action-chunking-trace}).
Sizing \(N\) against the latency tail guarantees that a replacement
chunk is resident in memory before the active plan runs out of samples.
Yet ensuring that a replacement trajectory arrives in time addresses
only the temporal continuity of the plan. The state samples contained
within the admitted chunk must also match the kinematic capabilities of
the joints and remain continuous with the physical velocity and
acceleration of the mechanism at the handoff seam.

\includegraphics[width=1\linewidth,height=\textheight,keepaspectratio,alt={Receding-Horizon Action Chunking and Closed-Loop Disturbance Recovery on Real Robot Hardware. Hardware deployment of a receding-horizon diffusion policy (T\_p = 16 predicted action steps, T\_a = 8 executed steps before replanning) on a canonical 6-DOF industrial manipulator performing the planar Push-T manipulation task under dynamic perturbations. (Left) Visual Occlusion: A human hand waving in front of the overhead tracking camera introduces perceptual noise and inference latency jitter, but the receding-horizon execution buffer prevents pipeline starvation and preserves continuous end-effector motion. (Middle) In-Flight State Divergence: A physical disturbance abruptly shifts the target T-block during the active pushing phase; the rolling replanning cycle detects state divergence from the recorded plan assumptions and immediately synthesizes a corrective trajectory chunk that re-acquires and aligns the block. (Right) Terminal Suffix Replanning: When the block is displaced after the manipulator has already initiated its terminal motion toward the designated end-zone, the policy aborts the terminal sequence, replans a return trajectory, and restores the block to the target alignment. Source: Chi et al. (2025).}]{/Users/VJ/GitHub/MLSysBook-integrate-physical/physical/book/chapters/11-planning/figures/fig11_real_action_chunking_trace.png}

\section{Feasibility and Continuity}\label{sec-planning-11-4}

A candidate trajectory produced by a learned policy or path planner is
merely a sequence of geometric propositions until the nervous system
verifies that the physical body can execute it. An admitted motion plan
must satisfy simultaneous constraints across multiple physical domains,
including joint position limits, angular velocity limits, actuator
acceleration envelopes, contact force thresholds, workspace obstacle
clearance, and the dynamic stopping distance required to prevent
collision if an emergency occurs. A systems engineer does not compute
full-body kinematics from first principles within the tracking loop
(\citeproc{ref-craig2005introduction}{Craig 2005};
\citeproc{ref-lynch2017modern}{Lynch and Park 2017}). Instead, the
safety enforcer evaluates a machinery-specific rigid-body model that
maps commanded joint positions \(\mathbf{q}(t)\), velocities
\(\dot{\mathbf{q}}(t)\), and accelerations \(\ddot{\mathbf{q}}(t)\)
directly to actuator torques
\(\boldsymbol{\tau}(t) = \mathbf{M}(\mathbf{q})\ddot{\mathbf{q}} + \mathbf{C}(\mathbf{q}, \dot{\mathbf{q}})\dot{\mathbf{q}} + \mathbf{g}(\mathbf{q})\)
and motor phase currents \(I_i(t) = \tau_i(t) / k_{\tau,i}\).
Within-plan \textbf{feasibility} is therefore established not by
geometric reachability, but by verifying that the predicted peak torque
remains within the continuous thermal limit of the motor windings and
the mechanical yield limits of the reduction gears across the entire
time horizon. In real hardware deployments, trajectory optimization
frameworks solve these constrained optimization problems over
parallelized GPU threads, generating time-parameterized minimum-jerk
splines that dynamically steer industrial manipulators around obstacles
while strictly enforcing kinematic and torque bounds
(Figure~\ref{fig-11-manipulator-trajectory}).

\begin{figure}[t]

\centering{

\includegraphics[width=0.95\linewidth,height=\textheight,keepaspectratio]{/Users/VJ/GitHub/MLSysBook-integrate-physical/physical/book/chapters/11-planning/figures/fig11_real_manipulator_trajectory.png}

}

\caption{\label{fig-11-manipulator-trajectory}\textbf{Real-Time
Trajectory Optimization and Dynamic Obstacle Avoidance on Industrial
Manipulators.} Experimental deployment of GPU-accelerated motion
generation (\emph{cuRobo}) executing time-parameterized, minimum-jerk
trajectory splines on physical hardware. (Top row) A canonical 6-DOF
industrial manipulator executes an online obstacle avoidance trajectory
around a dynamic workspace obstacle, continuously evaluating
continuous-time collision distances against a 3D Euclidean Signed
Distance Field (ESDF) voxel map generated via real-time depth
perception. (Middle row) A canonical 6-DOF collaborative arm executes a
collision-free pick-and-lift trajectory spline constrained by joint
velocity, acceleration, and torque envelopes. (Bottom row) Coordinated
multi-arm trajectory synthesis in simulation showing two canonical
industrial manipulators executing synchronized, collision-free obstacle
avoidance around shared workspace boundaries. Source: Sundaralingam et
al. (\citeproc{ref-sundaralingam2023curobo}{2023}).}

\end{figure}%

Validating feasibility solely at discrete planning waypoints creates a
false sense of compliance. When a planner outputs trajectory points
separated by a sampling interval \(T_s\), the physical mechanism moves
along an interpolated path between those samples. If the nervous system
checks constraints only at the discrete sample timestamps
\(t_k = k T_s\), high-order trajectory derivatives can exceed physical
bounds during the inter-sample interval \([t_k, t_{k+1}]\). Consider an
analytical example of a linear joint with a maximum acceleration limit
\(a_{\max} = 10.0\text{ m/s}^2\) and a velocity ceiling
\(v_{\max} = 2.00\text{ m/s}\). If the planner outputs consecutive
velocity samples \(\dot{q}_k = 1.95\text{ m/s}\) and
\(\dot{q}_{k+1} = 1.95\text{ m/s}\) separated by \(T_s = 20\text{ ms}\),
a cubic spline interpolator that applies an intermediate acceleration of
\(a_{\max}\) produces an intra-sample velocity overshoot bounded by
\(\Delta v \le \frac{1}{2} a_{\max} T_s\). Evaluating this bound yields:
\[\Delta v \le \frac{1}{2} (10.0\text{ m/s}^2)(0.020\text{ s}) = 0.10\text{ m/s}\]
The true physical velocity peaks at
\(1.95\text{ m/s} + 0.10\text{ m/s} = 2.05\text{ m/s}\) between the two
valid waypoints, exceeding the \(2.00\text{ m/s}\) limit and triggering
an over-speed clamp. To guarantee continuous-time feasibility without
evaluating every infinitesimal point, the admission filter must evaluate
an interior contracted bound,
\(|\dot{q}_k| \le v_{\max} - \frac{1}{2} a_{\max} T_s\), derived from
the maximum rate of change of the interpolator.

Concatenating two trajectories that are each individually feasible over
their respective time domains does not produce a feasible combined plan.
Consider an active trajectory chunk \(A\) defined over
\(t \in [0, T_A]\) and an incoming replacement chunk \(B\) defined over
\(t \in [0, T_B]\), joined at the handoff timestamp
\(t_{\mathrm{seam}}\). Even when chunk \(A\) and chunk \(B\) strictly
obey all kinematic and dynamic bounds in isolation, their intersection
at \(t = t_{\mathrm{seam}}\) can exhibit severe derivative
discontinuities (Figure~\ref{fig-11-action-chunk-seam-continuity}). The
level of continuity required across the seam is dictated by the
low-level actuator interface. A position-controlled actuator receiving a
continuous position command (\(C^0\) continuity) may still experience a
step change in velocity
\(\Delta \dot{q} = \dot{q}_B(0) - \dot{q}_A(t_{\mathrm{seam}})\). A
velocity-controlled or torque-controlled actuator receiving matching
velocity endpoints (\(C^1\) continuity) will still experience a step
change in commanded acceleration
\(\Delta \ddot{q} = \ddot{q}_B(0) - \ddot{q}_A(t_{\mathrm{seam}})\),
demanding an instantaneous jump in actuator torque
\(\Delta \boldsymbol{\tau} = \mathbf{M}(\mathbf{q})\Delta \ddot{\mathbf{q}}\).

\protect\phantomsection\label{callout-definitionux2a-11.3}
\begin{fbxSimple}{callout-definition}{Definition:}{The $\mathcal{C}^2$ Jerk Spline Seam}
\phantomsection\label{callout-definition*-11.3}

A \textbf{\(\mathcal{C}^2\) jerk spline seam} is the temporal,
kinematic, and dynamic boundary interface (\(t_{\mathrm{seam}}\))
between consecutive receding-horizon action chunks, where the nervous
system enforces simultaneous second-order continuity of position
(\(\mathbf{q}\)), velocity (\(\dot{\mathbf{q}}\)), and acceleration
(\(\ddot{\mathbf{q}}\)) via online polynomial spline synthesis, strictly
bounding dynamic jerk
(\(\dddot{\mathbf{q}} \le \dddot{\mathbf{q}}_{\mathrm{limit}}\)) to
eliminate step torque impulses and prevent exciting mechanical
transmission anti-resonances.

\end{fbxSimple}

When a discontinuous velocity command reaches the tracking controller,
the resulting shock load is governed by the reflected inertia of the
drive and the mechanical compliance of the transmission, not by the
nominal step size alone. In an analytical example of a rotary robotic
joint, a motor with rotor inertia
\(J_m = 1.2 \times 10^{-4}\text{ kg}\cdot\text{m}^2\) drives a link
inertia \(J_L = 1.5\text{ kg}\cdot\text{m}^2\) through a \(100:1\) speed
reducer with torsional stiffness
\(k_\theta = 2.0 \times 10^4\text{ N}\cdot\text{m/rad}\). The total
effective inertia reflected at the joint output is
\(J_{\mathrm{eff}} = J_L + N_g^2 J_m = 1.5\text{ kg}\cdot\text{m}^2 + (100)^2 (1.2 \times 10^{-4}\text{ kg}\cdot\text{m}^2) = 2.7\text{ kg}\cdot\text{m}^2\).
If replacement chunk \(B\) introduces a boundary velocity step
\(\Delta \dot{q} = 0.15\text{ rad/s}\) across a single digital control
cycle \(T_c = 2.0\text{ ms}\) (\(500\text{ Hz}\)), the tracking loop
attempts to close the velocity gap by commanding an apparent
acceleration
\(\ddot{q}_{\mathrm{cmd}} = \frac{\Delta \dot{q}}{T_c} = \frac{0.15\text{ rad/s}}{0.002\text{ s}} = 75.0\text{ rad/s}^2\).
The inertial torque required to accelerate the reflected mass across
this single period is:
\[\tau_{\mathrm{seam}} = J_{\mathrm{eff}} \ddot{q}_{\mathrm{cmd}} = (2.7\text{ kg}\cdot\text{m}^2)(75.0\text{ rad/s}^2) = 202.5\text{ N}\cdot\text{m}\]
If the joint transmission has a continuous torque rating of
\(45\text{ N}\cdot\text{m}\) and a transient peak ceiling of
\(120\text{ N}\cdot\text{m}\), this seam step demands \(169\%\) of the
maximum allowable transient rating. Across the elastic compliance of the
transmission, this dynamic torque produces an instantaneous angular
deflection
\(\theta_{\mathrm{deflect}} = \frac{\tau_{\mathrm{seam}}}{k_\theta} = \frac{202.5\text{ N}\cdot\text{m}}{20000\text{ N}\cdot\text{m/rad}} = 0.0101\text{ rad}\)
(\(0.58^\circ\)). The command step does not appear dangerous in an
unrendered coordinate plot, but the physical gearbox absorbs the entire
kinetic pulse as tooth bending stress and bearing brinelling.\sidenote{\footnotesize Modern
  learned policies stream action chunks to servo drives over fieldbuses
  (EtherCAT, CANopen). Unblended acceleration steps at chunk boundaries
  inject delta current impulses, causing inverter DC bus voltage
  pumping, desaturation trips, and mechanical bearing brinelling.}

\begin{figure}[t]

\centering{

\includegraphics[width=1\linewidth,height=\textheight,keepaspectratio]{index_files/mediabag//Users/VJ/GitHub/MLSysBook-integrate-physical/physical/book/chapters/11-planning/figures/fig11_action_chunk_seam_continuity.pdf}

}

\caption{\label{fig-11-action-chunk-seam-continuity}\textbf{Action Chunk
Seam Boundary Dynamics: Unchecked \(C^0\) Concatenation vs.~\(C^2\)
Jerk-Continuous Spline Synthesis.} Splicing an incoming action chunk
(\(B\)) onto an active chunk (\(A\)) without higher-order derivative
matching creates severe dynamic shocks at the handoff seam
(\(t_{\mathrm{seam}} = 50\text{ ms}\)). (Case A) In an unchecked \(C^0\)
concatenation, matching position alone allows a velocity step
\(\Delta \dot{q} = +0.60\text{ rad/s}\) across a single
\(T_c = 1.0\text{ ms}\) servo period, inducing an apparent joint
acceleration impulse \(\ddot{q}_{\mathrm{cmd}} = 600\text{ rad/s}^2\).
Through the joint's reflected inertia
(\(J_{\mathrm{eff}} = 2.7\text{ kg}\cdot\text{m}^2\)), this step
transmits an impulsive torque
\(\tau_{\mathrm{seam}} = 324\text{ N}\cdot\text{m}\) that exceeds the
continuous motor rating (\(45\text{ N}\cdot\text{m}\)), transient peak
ceiling (\(110\text{ N}\cdot\text{m}\)), and the
\(285\text{ N}\cdot\text{m}\) shear yield limit, resulting in
catastrophic drive pin fracture and bearing brinelling. (Case B)
Splicing via a \(5^{\text{th}}\)-order polynomial bridge over
\(T_{\mathrm{blend}} = 20\text{ ms}\) guarantees \(C^2\) continuity and
bounded minimum jerk (\(\dddot{q} \le 4500\text{ rad/s}^3\)), keeping
dynamic torque within the continuous operating zone
(\(\tau \le 38\text{ N}\cdot\text{m}\)) with zero resonant vibration.}

\end{figure}%

The physical consequences of an unmitigated seam discontinuity vary
drastically across actuator architectures and mechanical drive
technologies. As detailed in Table~\ref{tbl-11-transmission-dynamics},
the effective reflected inertia \(J_{\mathrm{eff}} = J_L + N_g^2 J_m\)
and torsional stiffness \(k_\theta\) determine both the natural resonant
frequency of the driveline
(\(f_{\mathrm{res}} \approx \frac{1}{2\pi}\sqrt{k_\theta / J_{\mathrm{eff}}}\))
and its vulnerability to impulsive torque spikes. In high-ratio strain
wave (harmonic) drives commonly employed in collaborative robots, the
rotor inertia is amplified by
\(N_g^2 \approx 10^4\text{--}2.5\times 10^4\), creating low resonant
frequencies (\(12\text{--}28\text{ Hz}\)) that are readily excited by
step acceleration commands, risking flexspline tooth fatigue.
Conversely, quasi-direct drive (QDD) actuators in dynamic legged robots
feature small gear ratios (\(6:1\text{--}10:1\)) and high mechanical
bandwidth (\(80\text{--}180\text{ Hz}\)); while resilient against
mechanical tooth shear, unblended command steps in QDD systems cause
immediate inverter DC bus voltage pumping and current loop desaturation
trips. Sizing trajectory derivative clamps therefore requires matching
the planner's boundary checker to the physical impedance and resonance
envelope of the specific actuator transmission.

\begin{longtable}[]{@{}
  >{\raggedright\arraybackslash}p{(\linewidth - 12\tabcolsep) * \real{0.1429}}
  >{\raggedright\arraybackslash}p{(\linewidth - 12\tabcolsep) * \real{0.1429}}
  >{\raggedright\arraybackslash}p{(\linewidth - 12\tabcolsep) * \real{0.1429}}
  >{\raggedright\arraybackslash}p{(\linewidth - 12\tabcolsep) * \real{0.1429}}
  >{\raggedright\arraybackslash}p{(\linewidth - 12\tabcolsep) * \real{0.1429}}
  >{\raggedright\arraybackslash}p{(\linewidth - 12\tabcolsep) * \real{0.1429}}
  >{\raggedright\arraybackslash}p{(\linewidth - 12\tabcolsep) * \real{0.1429}}@{}}
\caption{\textbf{Mechanical Transmission Parameters, Resonance
Frequencies, and Jerk Limits Across Physical AI Embodiments.} Drive
technologies exhibit divergent reflected inertia ratios
(\(N_g^2 J_m / J_L\)), torsional stiffness ratings (\(k_\theta\)), and
resonant frequencies
(\(f_{\mathrm{res}} \approx \frac{1}{2\pi}\sqrt{k_\theta / J_{\mathrm{eff}}}\)).
While high-ratio strain wave transmissions suffer mechanical flexspline
damage under low-frequency jerk (\(< 30\text{ Hz}\)), quasi-direct
drives transfer unmitigated seam acceleration steps into motor current
saturation and inverter bus voltage
surges.}\label{tbl-11-transmission-dynamics}\tabularnewline
\toprule\noalign{}
\begin{minipage}[b]{\linewidth}\raggedright
Embodiment \& Actuator Architecture
\end{minipage} & \begin{minipage}[b]{\linewidth}\raggedright
Gear Ratio (\(N_g\)) \& Torsional Stiffness (\(k_\theta\))
\end{minipage} & \begin{minipage}[b]{\linewidth}\raggedright
Reflected Inertia Ratio (\(N_g^2 J_m / J_L\))
\end{minipage} & \begin{minipage}[b]{\linewidth}\raggedright
Resonant Frequency (\(f_{\mathrm{res}}\))
\end{minipage} & \begin{minipage}[b]{\linewidth}\raggedright
Continuous Jerk Limit (\(\dddot{q}_{\mathrm{limit}}\))
\end{minipage} & \begin{minipage}[b]{\linewidth}\raggedright
Critical Seam Step Failure Mode
\end{minipage} & \begin{minipage}[b]{\linewidth}\raggedright
Hardware Interceptor \& Protective Filter
\end{minipage} \\
\midrule\noalign{}
\endfirsthead
\toprule\noalign{}
\begin{minipage}[b]{\linewidth}\raggedright
Embodiment \& Actuator Architecture
\end{minipage} & \begin{minipage}[b]{\linewidth}\raggedright
Gear Ratio (\(N_g\)) \& Torsional Stiffness (\(k_\theta\))
\end{minipage} & \begin{minipage}[b]{\linewidth}\raggedright
Reflected Inertia Ratio (\(N_g^2 J_m / J_L\))
\end{minipage} & \begin{minipage}[b]{\linewidth}\raggedright
Resonant Frequency (\(f_{\mathrm{res}}\))
\end{minipage} & \begin{minipage}[b]{\linewidth}\raggedright
Continuous Jerk Limit (\(\dddot{q}_{\mathrm{limit}}\))
\end{minipage} & \begin{minipage}[b]{\linewidth}\raggedright
Critical Seam Step Failure Mode
\end{minipage} & \begin{minipage}[b]{\linewidth}\raggedright
Hardware Interceptor \& Protective Filter
\end{minipage} \\
\midrule\noalign{}
\endhead
\bottomrule\noalign{}
\endlastfoot
\textbf{Precision Articulated Arm (Harmonic / Strain Wave)} &
\(N_g = 100\text{--}160\);
\(k_\theta = 1.5\text{--}4.0 \times 10^4\text{ N}\cdot\text{m/rad}\);
zero backlash & \(2.0\text{--}8.0 \times\) (rotor inertia dominates
driveline dynamics) & \(12\text{--}28\text{ Hz}\) (low structural
anti-resonance) & \(15\text{--}45\text{ rad/s}^3\) & Flexspline tooth
fatigue, wave generator brinelling, limit cycles &
\(5^{\text{th}}\)-order polynomial bridge + notch filter tuned to
\(f_{\mathrm{res}}\) \\
\textbf{Heavy Industrial Manipulator (Cycloidal Drive)} &
\(N_g = 40\text{--}120\);
\(k_\theta = 8.0\text{--}25.0 \times 10^4\text{ N}\cdot\text{m/rad}\);
backlash \(< 1.5'\) & \(1.0\text{--}4.0 \times\) (balanced motor-to-link
inertia) & \(35\text{--}75\text{ Hz}\) (stiff cycloidal pin assembly) &
\(50\text{--}120\text{ rad/s}^3\) & Input drive pin shear fracture,
output roller spalling, keyway shear & Model-based inverse dynamics
torque limiter + hardware rate clamp \\
\textbf{Dynamic Legged Robot / Humanoid (Quasi-Direct Drive)} &
\(N_g = 6\text{--}10\);
\(k_\theta = 0.8\text{--}2.0 \times 10^3\text{ N}\cdot\text{m/rad}\);
backdrivable & \(0.05\text{--}0.20 \times\) (link mass and impact
dynamics dominate) & \(80\text{--}180\text{ Hz}\) (high structural
bandwidth) & \(250\text{--}600\text{ rad/s}^3\) & Inverter DC bus
voltage pumping, current desaturation, stator heating & Microsecond
CBF-QP torque filter + regenerative bus clamp shunt \\
\textbf{High-Speed Linear Gantry (Direct-Drive Ironless)} & \(N_g = 1\)
(gearless); \(k_x = 5.0\text{--}15.0 \times 10^6\text{ N/m}\); zero
backlash & \(0.0 \times\) (rotor mass equals carriage payload mass) &
\(120\text{--}300\text{ Hz}\) (very high mechanical rigidity) &
\(40\text{--}100\text{ m/s}^3\) & Carriage linear bearing pitting,
magnetic pole slip, encoder loss & FPGA-based \(S\)-curve interpolator
with microsecond jerk limiter \\
\end{longtable}

Smoothing this boundary discontinuity requires online \textbf{\(C^2\)
trajectory spline blending}, which must be treated as the synthesis of
an entirely new motion plan rather than as cosmetic filtering. When the
nervous system inserts a polynomial bridge \(q_{\mathrm{blend}}(t)\)
over a transition interval \(T_{\mathrm{blend}}\) to connect chunk \(A\)
and incoming chunk \(B\), the resulting bridge segment must undergo the
exact same dynamic feasibility checks as the primary chunks.\sidenote{\footnotesize Flash
  and Hogan (1985) showed human reaching minimizes jerk. In robotic
  trajectory generation, bounding jerk (\(\dddot{\mathbf{q}}\)) is
  mandatory to prevent excitation of high-frequency torsional resonances
  in harmonic drives.} Standard trajectory generation theory proves that
a fifth-order (quintic) polynomial bridge can match position, velocity,
and acceleration boundary conditions simultaneously, guaranteeing
\(C^2\) continuity and continuous jerk across the transition
(\citeproc{ref-biagiotti2008trajectory}{Biagiotti and Melchiorri 2008};
\citeproc{ref-siciliano2009robotics}{Siciliano et al. 2009}).

Mathematically, the quintic blend bridge is parameterized on the
normalized transition interval \(t \in [0, T_{\mathrm{blend}}]\) as:
\[q(t) = a_0 + a_1 t + a_2 t^2 + a_3 t^3 + a_4 t^4 + a_5 t^5\] To
satisfy second-order derivative continuity, the polynomial coefficients
\(\mathbf{a} = [a_0, a_1, a_2, a_3, a_4, a_5]^T\) must satisfy six
simultaneous boundary conditions:
\[q(0) = q_0, \quad \dot{q}(0) = \dot{q}_0, \quad \ddot{q}(0) = \ddot{q}_0\]
\[q(T_{\mathrm{blend}}) = q_1, \quad \dot{q}(T_{\mathrm{blend}}) = \dot{q}_1, \quad \ddot{q}(T_{\mathrm{blend}}) = \ddot{q}_1\]
where \((q_0, \dot{q}_0, \ddot{q}_0)\) is the terminal physical state of
chunk \(A\) at the handoff instant, and \((q_1, \dot{q}_1, \ddot{q}_1)\)
is the target state on incoming chunk \(B\) at
\(t = T_{\mathrm{blend}}\). Solving this linear algebraic system yields
closed-form algebraic expressions for all six coefficients:
\[a_0 = q_0, \quad a_1 = \dot{q}_0, \quad a_2 = \frac{1}{2}\ddot{q}_0\]
\[a_3 = \frac{20(q_1 - q_0) - (8\dot{q}_1 + 12\dot{q}_0)T_{\mathrm{blend}} - (3\ddot{q}_0 - \ddot{q}_1)T_{\mathrm{blend}}^2}{2 T_{\mathrm{blend}}^3}\]
\[a_4 = \frac{30(q_0 - q_1) + (14\dot{q}_1 + 16\dot{q}_0)T_{\mathrm{blend}} + (3\ddot{q}_0 - 2\ddot{q}_1)T_{\mathrm{blend}}^2}{2 T_{\mathrm{blend}}^4}\]
\[a_5 = \frac{12(q_1 - q_0) - 6(\dot{q}_1 + \dot{q}_0)T_{\mathrm{blend}} - (\ddot{q}_1 - \ddot{q}_0)T_{\mathrm{blend}}^2}{2 T_{\mathrm{blend}}^5}\]
Because these equations require only basic arithmetic operations
(\(O(1)\) computational complexity), an embedded real-time MCU evaluates
them in less than \(5\,\mu\text{s}\), well within a \(1\text{ kHz}\)
servo timer interrupt.

The resulting dynamic \textbf{jerk profile} is linear and continuous:
\[\dddot{q}(t) = 6 a_3 + 24 a_4 t + 60 a_5 t^2\] Bounding maximum
dynamic jerk
(\(\max_{t \in [0, T_{\mathrm{blend}}]} |\dddot{q}(t)| \le \dddot{q}_{\mathrm{limit}}\))
is non-negotiable in physical AI: it ensures that the spectrum of
commanded torque does not contain high-frequency components near the
transmission anti-resonance frequency
\(f_{\mathrm{res}} = \frac{1}{2\pi}\sqrt{k_\theta / J_{\mathrm{eff}}}\),
preventing destructive limit cycles and structural vibration.

Selecting an appropriate mathematical representation for candidate
trajectories governs the trade-off between expressive policy generation,
real-time computational overhead, and hardware-level derivative
guarantees. As compared in
Table~\ref{tbl-11-trajectory-representations}, classical
parameterizations such as quintic polynomial splines and minimum-jerk
B-splines provide rigorous, closed-form continuity (\(C^2\) and \(C^4\))
and strict interior derivative bounds with sub-millisecond evaluation
times, making them directly consumable by motor drives. In contrast,
discrete action chunks produced by autoregressive transformers (such as
Action Chunking with Transformers) or iterative score-based diffusion
policies operate across multi-millisecond inference horizons and output
discrete waypoint sequences (\(C^{-1}\)). Because these raw generative
proposals lack inherent physical feasibility guarantees and exhibit
broad-spectrum derivative scatter, the nervous system must project them
through intermediate spline synthesizers or safety-barrier filters
before admission to the real-time execution queue.

\begin{longtable}[]{@{}
  >{\raggedright\arraybackslash}p{(\linewidth - 12\tabcolsep) * \real{0.1429}}
  >{\raggedright\arraybackslash}p{(\linewidth - 12\tabcolsep) * \real{0.1429}}
  >{\raggedright\arraybackslash}p{(\linewidth - 12\tabcolsep) * \real{0.1429}}
  >{\raggedright\arraybackslash}p{(\linewidth - 12\tabcolsep) * \real{0.1429}}
  >{\raggedright\arraybackslash}p{(\linewidth - 12\tabcolsep) * \real{0.1429}}
  >{\raggedright\arraybackslash}p{(\linewidth - 12\tabcolsep) * \real{0.1429}}
  >{\raggedright\arraybackslash}p{(\linewidth - 12\tabcolsep) * \real{0.1429}}@{}}
\caption{\textbf{Trajectory Representation Comparison Matrix for
Physical AI Systems.} Comparative analysis of classical continuous
spline parameterizations and modern learned discrete action
representations across continuity degree, synthesis latency, dynamic
feasibility guarantees, and mechanical jerk mitigation. While generative
models provide high-level semantic generalization, their discrete,
unconstrained outputs (\(C^{-1}\)) require downstream spline
interpolation (\(C^2\)/\(C^3\)) and constraint projection before
reaching motor drive current
loops.}\label{tbl-11-trajectory-representations}\tabularnewline
\toprule\noalign{}
\begin{minipage}[b]{\linewidth}\raggedright
Trajectory Paradigm
\end{minipage} & \begin{minipage}[b]{\linewidth}\raggedright
Mathematical Parameterization
\end{minipage} & \begin{minipage}[b]{\linewidth}\raggedright
Continuity Order (\(C^k\))
\end{minipage} & \begin{minipage}[b]{\linewidth}\raggedright
Synthesis Latency (\(T_{\mathrm{eval}}\))
\end{minipage} & \begin{minipage}[b]{\linewidth}\raggedright
Dynamic Feasibility Guarantee
\end{minipage} & \begin{minipage}[b]{\linewidth}\raggedright
Jerk Bounding \& Vibration Control
\end{minipage} & \begin{minipage}[b]{\linewidth}\raggedright
Seam Splicing Overhead
\end{minipage} \\
\midrule\noalign{}
\endfirsthead
\toprule\noalign{}
\begin{minipage}[b]{\linewidth}\raggedright
Trajectory Paradigm
\end{minipage} & \begin{minipage}[b]{\linewidth}\raggedright
Mathematical Parameterization
\end{minipage} & \begin{minipage}[b]{\linewidth}\raggedright
Continuity Order (\(C^k\))
\end{minipage} & \begin{minipage}[b]{\linewidth}\raggedright
Synthesis Latency (\(T_{\mathrm{eval}}\))
\end{minipage} & \begin{minipage}[b]{\linewidth}\raggedright
Dynamic Feasibility Guarantee
\end{minipage} & \begin{minipage}[b]{\linewidth}\raggedright
Jerk Bounding \& Vibration Control
\end{minipage} & \begin{minipage}[b]{\linewidth}\raggedright
Seam Splicing Overhead
\end{minipage} \\
\midrule\noalign{}
\endhead
\bottomrule\noalign{}
\endlastfoot
\textbf{Quintic Polynomial Spline (\(5^{\text{th}}\)-Order)} &
\(q(t) = \sum_{i=0}^5 a_i t^i\), algebraic BVP on \([0, T]\) & \(C^2\)
continuous (\(q, \dot{q}, \ddot{q}\) matched; piecewise-linear
\(\dddot{q}\)) & \(< 5\,\mu\text{s}\)
(\(\mathbf{A}_{6 \times 6}^{-1} \mathbf{b}\) on MCU/DSP) & Exact;
algebraic extrema via derivative polynomial roots & Closed-form linear
jerk profile; avoids acceleration step shocks & \(O(1)\) closed-form
boundary bridge matching \((q_0, \dot{q}_0, \ddot{q}_0)\) \\
\textbf{Minimum-Jerk B-Spline (Degree \(p \ge 5\))} &
\(q(t) = \sum_{j=0}^m \mathbf{P}_j N_{j,p}(t)\);
\(\min \int_0^T \Vert \dddot{q} \Vert^2 dt\) & \(C^{p-1}\) continuous
(e.g., \(C^4\) for quintic; smooth jerk \(\dddot{q}\) and snap
\(q^{(4)}\)) & \(50\text{--}500\,\mu\text{s}\) (banded linear system
solver) & Convex hull property guarantees bounds within control polygon
& Globally minimizes total jerk integral; provably suppresses resonance
& \(O(m)\) banded knot insertion and degree elevation \\
\textbf{Receding Horizon Action Chunk (ACT / Transformers)} & Discrete
sequence
\(\mathbf{A} = [\mathbf{a}_1, \dots, \mathbf{a}_H] \in \mathbb{R}^{H \times D}\)
& \(C^{-1}\) discrete / \(C^0\) under zero-order or linear hold &
\(15\text{--}60\text{ ms}\) (batched attention on accelerator) & None;
stochastic outputs may violate velocity and torque envelopes & Zero
native jerk bounding; exhibits inter-sample derivative jitter & Temporal
ensembling / exponential moving average over chunks \\
\textbf{Diffusion Policy / Flow-Matching Rollout} & Continuous score
ODE/SDE:
\(\mathbf{a}_{k-1} = \mathbf{a}_k - \gamma \mathbf{s}_\theta(\mathbf{a}_k) + \boldsymbol{\sigma}\boldsymbol{\xi}\)
& \(C^{-1}\) discrete waypoint sequence & \(30\text{--}180\text{ ms}\)
(\(K=8\text{--}16\) denoising steps on NPU) & Probabilistic;
out-of-distribution drift causes trajectory spikes & Unbounded output
chatter; score noise injects driveline harmonics & Receding-horizon
re-sampling with initial-state conditioning \\
\end{longtable}

\protect\phantomsection\label{callout-mathux2a-11.4}
\begin{fbxSimple}{callout-math}{Fermi Math:}{Seam Acceleration Jump and Reflected Inertia Torque Spike}
\phantomsection\label{callout-math*-11.4}

To quantify why \(C^0\) position matching destroys transmissions,
consider what happens when two trajectory chunks meet with a velocity
mismatch across three physical machine embodiments:

\textbf{1. Industrial Manipulator Joint (Cycloidal Reducer):}

\begin{itemize}
\item
  Parameters: Rotor inertia
  \(J_m = 1.2 \times 10^{-4}\text{ kg}\cdot\text{m}^2\), link inertia
  \(J_L = 1.5\text{ kg}\cdot\text{m}^2\), gear ratio \(N_g = 100\),
  servo cycle \(T_c = 1.0\text{ ms}\).
\item
  Effective reflected inertia:
  \[J_{\mathrm{eff}} = J_L + N_g^2 J_m = 1.5 + (100)^2 (1.2 \times 10^{-4}) = 2.7\text{ kg}\cdot\text{m}^2\]
\item
  Unblended seam jump: Velocity step
  \(\Delta \dot{q} = 0.60\text{ rad/s}\) across single tick
  \(T_c = 1.0\text{ ms}\).
  \[\ddot{q}_{\mathrm{cmd}} = \frac{0.60\text{ rad/s}}{0.001\text{ s}} = 600.0\text{ rad/s}^2 \implies \tau_{\mathrm{seam}} = J_{\mathrm{eff}} \ddot{q}_{\mathrm{cmd}} = (2.7)(600.0) = 1620\text{ N}\cdot\text{m}\]
  Even if the motor drive current loop saturates at its hardware limit
  (\(I_{\max} \implies \tau_{\mathrm{sat}} = 324\text{ N}\cdot\text{m}\)),
  this impulsive torque exceeds the \(285\text{ N}\cdot\text{m}\)
  cycloidal disk pin yield shear limit by \(14\%\), causing pin fracture
  within \(1.1\text{ ms}\).
\item
  \(C^2\) Quintic Spline Solution: Splicing over
  \(T_{\mathrm{blend}} = 30\text{ ms}\) limits peak acceleration to
  \(\ddot{q}_{\mathrm{peak}} \approx \frac{5.77 \Delta \dot{q}}{T_{\mathrm{blend}}} = \frac{5.77(0.60)}{0.030} = 115.4\text{ rad/s}^2\),
  requiring
  \(\tau_{\mathrm{peak}} = (2.7)(115.4) = 31.2\text{ N}\cdot\text{m}\),
  operating safely within the \(45\text{ N}\cdot\text{m}\) continuous
  envelope (\(31\%\) margin).
\end{itemize}

\textbf{2. Autonomous Mobile Robot (Differential-Drive AGV):}

\begin{itemize}
\tightlist
\item
  Parameters: Chassis mass \(m = 150\text{ kg}\), wheel radius
  \(r = 0.10\text{ m}\), tire-ground friction coefficient
  \(\mu = 0.70\), control cadence \(T_c = 10\text{ ms}\).
\item
  Unblended seam jump: Forward velocity step
  \(\Delta v = 0.35\text{ m/s}\) across \(10\text{ ms}\).
  \[a_{\mathrm{cmd}} = \frac{0.35}{0.010} = 35.0\text{ m/s}^2 \implies F_{\mathrm{traction}} = m a = (150\text{ kg})(35.0\text{ m/s}^2) = 5250\text{ N}\]
  The maximum available friction force is
  \(F_{\mathrm{fric},\max} = \mu m g = (0.70)(150)(9.81) = 1030\text{ N}\).
  The commanded force exceeds available traction by \(510\%\), causing
  total wheel lockup, tire skidding, and complete loss of odometric
  localization.
\end{itemize}

\textbf{3. Process Extrusion System (Polymer Flow Governor):}

\begin{itemize}
\tightlist
\item
  Parameters: Extruder screw barrel volume
  \(V_0 = 2.5 \times 10^{-4}\text{ m}^3\), polymer melt bulk modulus
  \(\beta = 1.2\text{ GPa}\), nominal flow
  \(\dot{V} = 1.0\times 10^{-5}\text{ m}^3/\text{s}\), servo tick
  \(T_c = 5\text{ ms}\).
\item
  Unblended seam jump: Volumetric feed rate step
  \(\Delta \dot{V} = 0.4\times 10^{-5}\text{ m}^3/\text{s}\) over
  \(5\text{ ms}\).
  \[\Delta P_{\mathrm{seam}} \approx \beta \frac{\Delta \dot{V} \cdot T_c}{V_0} = (1.2 \times 10^9) \frac{(0.4 \times 10^{-5})(0.005)}{2.5 \times 10^{-4}} = 9.6\times 10^7\text{ Pa} = 96.0\text{ MPa}\]
  This pressure pulse exceeds the \(45\text{ MPa}\) barrel burst rupture
  rating, blowing out pressure transducer seals.
\end{itemize}

\end{fbxSimple}

\protect\phantomsection\label{callout-ruleux2a-11.5}
\begin{fbxSimple}{callout-rule}{Architect’s Rule of Thumb:}{The $\mathcal{C}^2$ Continuity Invariant}
\phantomsection\label{callout-rule*-11.5}

A trajectory without continuous acceleration (\(\mathcal{C}^2\)) is not
a motion plan; unblended velocity steps across replanning seams are
converted by driveline inertia into destructive torque spikes.

\end{fbxSimple}

A related physical failure arises from output chatter in learned policy
proposals. High-frequency oscillations across successive planning chunks
are frequently produced by stochastic sampling, score-matching noise, or
autoregressive variance in generative policies. Even when every discrete
sample in a chattering trajectory satisfies absolute position and
velocity limits, the rapid sign reversals in commanded acceleration
force the actuators into continuous torque reversal. In an analytical
duty cycle where an actuator reverses a torque of
\(30\text{ N}\cdot\text{m}\) at \(25\text{ Hz}\), the rapid current
transitions generate sustained \(I^2 R\) resistive heating in the motor
stator and hammer the gear teeth across their backlash zone.\sidenote{\footnotesize Torque
  chatter accelerates Joule heating (\(P = 3 I_{\text{rms}}^2 R\)). With
  \(\alpha = 0.00393\text{ K}^{-1}\), stator resistance rises \(+51\%\)
  at \(155^\circ\text{C}\).} The systems engineer budgets this
high-frequency variation not as a perceptual roughness metric, but as an
actuator thermal rise and transmission fatigue cost. A trajectory that
excites structural resonances or exceeds the allowable thermal
dissipation of the motor drive is physically infeasible, regardless of
how accurately its geometric waypoints follow the nominal task
objective. Even when a candidate chunk satisfies every internal and
boundary feasibility check, the nervous system must still account for
the reality that the next replacement chunk may fail to arrive before
the active plan reaches its terminal sample.

\section{Behavior at the Seam}\label{sec-planning-11-5}

A replacement trajectory will eventually arrive late. Treating that
delay as an asynchronous communication error or a generic software
exception shifts responsibility away from the trajectory generator to
software layers that lack physical state context. In an execution
architecture governed by real-time deadlines, \textbf{seam lateness} is
not the event where inference time exceeds its median \(P_{50}\)
duration. It is the physical interval between the last admissible
replacement time and the actual timestamp at which a verified
replacement trajectory is accepted by the nervous system. When a learned
policy runs on a compute substrate subject to memory contention, kernel
launch jitter, and thermal throttling, tail latency at \(P_{99}\) or
\(P_{99.9}\) routinely exceeds the nominal replanning period. The
trajectory currently in execution must therefore contain its own
behavioral resolution for that latency. What the machine does when the
clock crosses the replanning deadline belongs to the plan itself,
budgeted in millimeters, millivolts, and thermal energy before the first
setpoint reaches the joint controllers.

Every active plan chunk operates across five distinct temporal
milestones, as mapped across the replanning regimes and clearance
budgets of Figure~\ref{fig-11-seam-timeline-fallback}. These are the
request point \(t_{\mathrm{req}}\), the blend deadline
\(t_{\mathrm{blend}}\), the fallback commitment point
\(t_{\mathrm{commit}}\), the nominal plan expiry \(t_{\mathrm{exp}}\),
and the terminal-state deadline \(t_{\mathrm{term}}\). Consider an
analytical example of a direct-drive manipulator carrying a
\(4.0\text{ kg}\) payload at an operating speed of
\(v_0 = 1.2\text{ m/s}\) toward an assembly fixture, with an available
downstream clearance of \(d_{\mathrm{clear}} = 0.30\text{ m}\). The
active chunk has a nominal execution duration of
\(T_{\mathrm{chunk}} = 500\text{ ms}\). At
\(t_{\mathrm{req}} = 0\text{ ms}\), the nervous system dispatches the
current state estimate to the policy to compute the next chunk. Under
nominal conditions, inference completes at \(P_{50} = 180\text{ ms}\),
leaving ample margin before the blend deadline at
\(t_{\mathrm{blend}} = 300\text{ ms}\). If the replacement chunk arrives
before \(t_{\mathrm{blend}}\), the controller blends the new waypoints
over a nominal \(100\text{ ms}\) window, achieving a continuous
transition before the active chunk reaches its nominal expiry at
\(t_{\mathrm{exp}} = 400\text{ ms}\). If inference encounters tail
delays and stretches to \(P_{99} = 340\text{ ms}\), the replacement
arrives after \(t_{\mathrm{blend}}\) but before
\(t_{\mathrm{commit}} = 360\text{ ms}\). In that late-but-blendable
interval, the blending window is compressed, forcing higher
instantaneous accelerations across the remaining overlap.

Once elapsed time reaches the fallback commitment point
\(t_{\mathrm{commit}}\), the window for evaluating and splicing an
incoming replacement closes completely. Any candidate trajectory
arriving after \(t_{\mathrm{commit}}\) is rejected, because executing
the validation arithmetic and interpolating an incoming boundary
condition would consume time that the physical mechanism cannot spare.
At \(t_{\mathrm{commit}}\), the controller branches unconditionally into
a precomputed \textbf{fallback suffix}. A fallback suffix is not an
open-ended idle routine, but a deterministic, finite trajectory segment
defined by an explicit entry condition, an execution duration, a bounded
resource allocation, and an exact terminal state. The plan cannot simply
instruct the nervous system to hold. In physical actuation, holding has
three conflicting interpretations that consume different energy and
clearance, and conflating them causes mechanical failure:

\begin{enumerate}
\def\labelenumi{\arabic{enumi}.}
\tightlist
\item
  \textbf{Holding the last commanded torque or current}: Maintains
  constant actuator effort, which allows the mechanism to accelerate
  uncontrollably if external contact loads shift or if gravitational
  moments vary as the arm drifts.
\item
  \textbf{Holding the last position setpoint}: Commands an instantaneous
  zero velocity at the final waypoint, which presents the low-level
  tracking loop with a step change in commanded velocity, drawing peak
  current, tripping motor drive overcurrent protection, and exciting
  structural resonances against stiff contact surfaces.
\item
  \textbf{Continuing at the last commanded velocity}: Avoids transient
  current spikes but expends remaining clearance at full operating
  speed, causing the mechanism to drive through its workspace boundary.
\end{enumerate}

A valid fallback suffix must specify a continuous deceleration profile
that brings joint velocities to zero while satisfying torque, voltage,
and thermal constraints throughout the braking maneuver.

The admissibility of the fallback deceleration profile is determined by
whether the braking distance and stopping time fit within the physical
clearance and available actuator torque. Continuing the analytical
example, suppose the motor drives can sustain a maximum linear
deceleration of \(a_{\mathrm{max}} = 6.0\text{ m/s}^2\) without
exceeding the thermal dissipation limits of the winding insulation.
Between the blend deadline \(t_{\mathrm{blend}} = 300\text{ ms}\) and
the fallback commitment point \(t_{\mathrm{commit}} = 360\text{ ms}\),
the machine accumulates a late latency interval of
\(\Delta t_{\mathrm{late}} = 60\text{ ms} = 0.060\text{ s}\) while
continuing along its nominal trajectory at \(v_0 = 1.2\text{ m/s}\).
During this unreplaced interval, the end-effector travels an initial
distance of
\(d_{\mathrm{late}} = v_0 \cdot \Delta t_{\mathrm{late}} = 1.2\text{ m/s} \times 0.060\text{ s} = 0.072\text{ m}\).
At \(t_{\mathrm{commit}}\), the fallback deceleration begins. Bringing
the mass from \(v_0 = 1.2\text{ m/s}\) to rest at
\(a_{\mathrm{max}} = 6.0\text{ m/s}^2\) requires a braking duration of
\(t_{\mathrm{brake}} = v_0 / a_{\mathrm{max}} = 1.2 / 6.0 = 0.20\text{ s} = 200\text{ ms}\),
covering a braking distance of
\(d_{\mathrm{brake}} = v_0^2 / (2 a_{\mathrm{max}}) = (1.44\text{ m}^2/\text{s}^2) / (12.0\text{ m/s}^2) = 0.120\text{ m}\).
The total physical distance consumed from the blend deadline to a
complete stop is
\(d_{\mathrm{total}} = d_{\mathrm{late}} + d_{\mathrm{brake}} = 0.072\text{ m} + 0.120\text{ m} = 0.192\text{ m}\).
Because \(d_{\mathrm{total}} = 0.192\text{ m}\) is less than the
available clearance \(d_{\mathrm{clear}} = 0.30\text{ m}\), the machine
halts safely at \(t_{\mathrm{term}} = 560\text{ ms}\).\sidenote{\footnotesize ISO
  10218-1/2 and ISO/TS 15066 mandate budgeting latency
  (\(T_{\text{sense}} + T_{\text{compute}} + T_{\text{brake}}\)) to
  ensure protective stoppage before violating safety zones.} If the
clearance had been restricted to \(0.150\text{ m}\), the plan would have
been physically inadmissible from the moment of synthesis, because its
fallback commitment point was scheduled too late to prevent contact.

\begin{figure}[t]

\centering{

\includegraphics[width=1\linewidth,height=\textheight,keepaspectratio]{index_files/mediabag//Users/VJ/GitHub/MLSysBook-integrate-physical/physical/book/chapters/11-planning/figures/fig11_seam_timeline_fallback.pdf}

}

\caption{\label{fig-11-seam-timeline-fallback}\textbf{Action Chunk Seam
Timeline, Replacement Ingestion Regimes, and Fallback Deceleration
Dynamics.} Execution timeline and physical clearance allocation across
the five critical temporal milestones of an active trajectory chunk
(\(T_{\mathrm{chunk}} = 500\text{ ms}\)). (Regime 1) On-time replacement
arrival (\(P_{50} = 180\text{ ms} \le t_{\mathrm{blend}}\)) allows a
full \(100\text{ ms}\) polynomial blend into Chunk \(B\) with zero
velocity penalty. (Regime 2) Late-but-blendable arrival
(\(P_{99} = 340\text{ ms} \in (t_{\mathrm{blend}}, t_{\mathrm{commit}}]\))
compresses the blend window to \(60\text{ ms}\), requiring runtime
algebraic peak-acceleration checks
(\(a_{\mathrm{peak}} \approx 6\Delta q / T_{\mathrm{blend}}^2\)) to
prevent over-torque conditions. (Regime 3) Once time crosses the
fallback commitment point \(t_{\mathrm{commit}} = 360\text{ ms}\)
without an admitted replacement, the controller branches unconditionally
into a precomputed \(C^2\) deceleration profile
(\(a_{\mathrm{max}} = 6.0\text{ m/s}^2\)). The total distance traversed
during late latency drift (\(d_{\mathrm{late}} = 72\text{ mm}\)) and
active braking (\(d_{\mathrm{brake}} = 120\text{ mm}\)) consumes
\(d_{\mathrm{total}} = 192\text{ mm}\), safely halting within the
available physical clearance (\(d_{\mathrm{clear}} = 300\text{ mm}\)) at
\(t_{\mathrm{term}} = 560\text{ ms}\) before engaging mechanical holding
brakes. The bottom cards contrast this structured fallback against the
failure modes of naive hold strategies (holding torque, holding
position, or continuing velocity).}

\end{figure}%

When a late replacement arrives before \(t_{\mathrm{commit}}\),
attempting to preserve continuous motion by blending the trajectories
consumes the remaining overlap duration and risks violating actuator
limits. If a replacement arrives at \(t = 340\text{ ms}\), the available
blending duration shrinks from the nominal \(100\text{ ms}\) to
\(T_{\mathrm{blend}} = 60\text{ ms}\). A bridge polynomial connecting
the active chunk state \((q_A, \dot{q}_A)\) to the replacement chunk
state \((q_B, \dot{q}_B)\) across a compressed duration
\(T_{\mathrm{blend}}\) requires peak accelerations that scale inversely
with \(T_{\mathrm{blend}}^2\). If the spatial separation between the
active trajectory and the new proposal at the blend boundary is
\(\Delta q = 0.015\text{ m}\), the additional acceleration induced
solely by the bridging interpolation scales as approximately
\(6 \Delta q / T_{\mathrm{blend}}^2 = 6(0.015\text{ m}) / (0.060\text{ s})^2 = 25.0\text{ m/s}^2\).
Added to the baseline chunk acceleration, this peak demand exceeds the
actuator limit of \(a_{\mathrm{max}} = 6.0\text{ m/s}^2\). Even though
Chunk A and Chunk B each satisfy kinematic and dynamic constraints when
evaluated in isolation, splicing them across a compressed temporal seam
causes an actuator overtorque condition. The trajectory generator must
verify the algebraic extrema of the spliced segment before committing to
the blend; if peak accelerations violate motor torque bounds, the
controller must reject the replacement and execute the fallback
deceleration.

In the terminal case where a replacement never arrives, the fallback
suffix must guide the machine into a bounded, statically stable
equilibrium before the terminal-state deadline \(t_{\mathrm{term}}\)
elapses. For a stationary manipulator, this equilibrium corresponds to
zero velocity, zero acceleration, and engagement of mechanical holding
brakes to dissipate motor current and eliminate winding heating. For
tasks where holding position in free space is impossible or hazardous,
such as balancing an inverted pendulum or transferring a molten material
that will solidify in a nozzle, reaching the terminal state without a
new plan constitutes a loss of operational authority. In such regimes,
the fallback suffix cannot resolve the failure on its own. It must
culminate in an explicit refusal request transmitted to the nervous
system's safety enforcer (Chapter~\ref{sec-enforcement}), triggering the
deterministic shutdown sequence or protective interlock that governs
unrecoverable state transitions.

To prevent catastrophic loss of physical authority when computation
delays breach timing margins, the nervous system enforces a
deterministic, multi-tier fallback hierarchy. As structured in
Table~\ref{tbl-11-seam-fallback-hierarchy}, each execution tier maps the
five temporal milestones
(\(t_{\mathrm{req}}, t_{\mathrm{blend}}, t_{\mathrm{commit}}, t_{\mathrm{exp}}, t_{\mathrm{term}}\))
to an explicit trajectory action, control mode, and enforcer invariant.
In Tier 0, on-time arrivals
(\(t_{\mathrm{arr}} \le t_{\mathrm{blend}}\)) execute nominal polynomial
blending with zero dynamic penalty. In Tier 1, replacement chunks
arriving during the late-but-blendable interval
(\(t_{\mathrm{blend}} < t_{\mathrm{arr}} \le t_{\mathrm{commit}}\))
undergo compressed spline bridging, subject to strict algebraic
peak-acceleration checks. Once time reaches the fallback commitment
point \(t_{\mathrm{commit}}\) without an admitted replacement, the
system branches unconditionally to Tier 2, engaging the precomputed
\(C^2\) path-tangent deceleration suffix. If the physical embodiment
encounters unexpected contact or workspace boundary violations during
deceleration, control escalates to Tier 3 (compliant passive damping) or
Tier 4 (Safe Torque Off and mechanical spring brake clamping). Under
this structured contract, late computation never results in undefined
hardware behavior or unguided momentum drift.

\begin{longtable}[]{@{}
  >{\raggedright\arraybackslash}p{(\linewidth - 12\tabcolsep) * \real{0.1429}}
  >{\raggedright\arraybackslash}p{(\linewidth - 12\tabcolsep) * \real{0.1429}}
  >{\raggedright\arraybackslash}p{(\linewidth - 12\tabcolsep) * \real{0.1429}}
  >{\raggedright\arraybackslash}p{(\linewidth - 12\tabcolsep) * \real{0.1429}}
  >{\raggedright\arraybackslash}p{(\linewidth - 12\tabcolsep) * \real{0.1429}}
  >{\raggedright\arraybackslash}p{(\linewidth - 12\tabcolsep) * \real{0.1429}}
  >{\raggedright\arraybackslash}p{(\linewidth - 12\tabcolsep) * \real{0.1429}}@{}}
\caption{\textbf{Action Chunk Seam Discontinuity and Replanning Latency
Fallback Hierarchy.} Five-tier deterministic safety and execution
escalation hierarchy mapping temporal arrival milestones
(\(t_{\mathrm{req}}, t_{\mathrm{blend}}, t_{\mathrm{commit}}, t_{\mathrm{exp}}, t_{\mathrm{term}}\))
to trajectory actions, control modes, and enforcer invariants. When
computational delays breach the fallback commitment point
\(t_{\mathrm{commit}}\), the nervous system unconditionally engages the
precomputed deceleration suffix rather than accepting tardy, unverified
neural proposals.}\label{tbl-11-seam-fallback-hierarchy}\tabularnewline
\toprule\noalign{}
\begin{minipage}[b]{\linewidth}\raggedright
Fallback Tier \& Execution Regime
\end{minipage} & \begin{minipage}[b]{\linewidth}\raggedright
Temporal Trigger Window
\end{minipage} & \begin{minipage}[b]{\linewidth}\raggedright
Active Trajectory Action
\end{minipage} & \begin{minipage}[b]{\linewidth}\raggedright
Kinematic \& Energy Budget Impact
\end{minipage} & \begin{minipage}[b]{\linewidth}\raggedright
Actuator Control Mode
\end{minipage} & \begin{minipage}[b]{\linewidth}\raggedright
Recovery \& Re-entry Path
\end{minipage} & \begin{minipage}[b]{\linewidth}\raggedright
Safety Enforcer Invariant \& Interlock
\end{minipage} \\
\midrule\noalign{}
\endfirsthead
\toprule\noalign{}
\begin{minipage}[b]{\linewidth}\raggedright
Fallback Tier \& Execution Regime
\end{minipage} & \begin{minipage}[b]{\linewidth}\raggedright
Temporal Trigger Window
\end{minipage} & \begin{minipage}[b]{\linewidth}\raggedright
Active Trajectory Action
\end{minipage} & \begin{minipage}[b]{\linewidth}\raggedright
Kinematic \& Energy Budget Impact
\end{minipage} & \begin{minipage}[b]{\linewidth}\raggedright
Actuator Control Mode
\end{minipage} & \begin{minipage}[b]{\linewidth}\raggedright
Recovery \& Re-entry Path
\end{minipage} & \begin{minipage}[b]{\linewidth}\raggedright
Safety Enforcer Invariant \& Interlock
\end{minipage} \\
\midrule\noalign{}
\endhead
\bottomrule\noalign{}
\endlastfoot
\textbf{Tier 0: Nominal Handover (Synchronous Splicing)} &
\(t_{\mathrm{arr}} \le t_{\mathrm{blend}}\) (On-time arrival) & Evaluate
\(5^{\text{th}}\)-order polynomial blend bridge over
\(T_{\mathrm{blend}}\) (\(100\text{ ms}\)) & Continuous \(C^2\)
transition; nominal torque profile with \(0\%\) margin erosion &
Feedforward torque + closed-loop trajectory tracking & Direct execution
of incoming Chunk \(B\); seamless continuation &
\(\Vert \Delta \dot{\mathbf{q}} \Vert \le \Delta \dot{\mathbf{q}}_{\mathrm{allow}}\),
\(\tau_{\mathrm{blend}} \le \tau_{\mathrm{cont}}\), zero tracking error
fault \\
\textbf{Tier 1: Compressed Overlap (Late-but-Blendable)} &
\(t_{\mathrm{blend}} < t_{\mathrm{arr}} \le t_{\mathrm{commit}}\)
(Delayed arrival) & Synthesize compressed spline bridge over
\(T_{\mathrm{blend}} = t_{\mathrm{commit}} - t_{\mathrm{arr}}\) &
Acceleration scales as \(O(1/T_{\mathrm{blend}}^2)\); consumes up to
\(85\%\) peak torque ceiling & High-gain tracking with dynamic
feedforward torque clamping & Execute Chunk \(B\) if algebraic peak
checks pass; else drop to Tier 2 &
\(\max \vert \ddot{q}(t) \vert \le a_{\mathrm{max}}\),
\(\Delta v \le v_{\mathrm{safe}} - v_{\mathrm{limit}}\), intra-sample
bound verified \\
\textbf{Tier 2: Path-Tangent Deceleration (Fallback Commitment)} &
\(t \ge t_{\mathrm{commit}}\) with no valid replacement chunk &
Unconditionally execute precomputed \(C^2\) stopping suffix over
\(T_{\mathrm{stop}}\) & Constant deceleration
\(a_{\mathrm{decel}} = -a_{\mathrm{max}}\); consumes clearance
\(d_{\mathrm{brake}} = \frac{v_0^2}{2 a_{\mathrm{max}}}\) & Closed-loop
deceleration profile to zero velocity (\(v \to 0\)) & Incoming chunks
discarded; replanning must resume from stationary state &
\(d_{\mathrm{total}} = d_{\mathrm{late}} + d_{\mathrm{brake}} < d_{\mathrm{clearance}}\);
ISO 10218-1 stopping envelope enforced \\
\textbf{Tier 3: Compliant Passive Damping (Contact / Stall Fallback)} &
\(t \ge t_{\mathrm{commit}}\) under contact load or boundary trip &
Terminate trajectory tracking; engage Cartesian damping:
\(\boldsymbol{\tau} = -\mathbf{D}\dot{\mathbf{q}} + \mathbf{g}(\mathbf{q})\)
& Kinetic energy dissipated into motor windings; bounds contact wrench &
Zero-stiffness joint torque control with virtual viscous damping &
Requires supervisor release clearance and workspace re-homing &
\(\Vert F_{\mathrm{ext}} \Vert \le F_{\mathrm{contact,max}}\)
(\(< 50\text{ N}\) per ISO/TS 15066); motor RMS thermal current within
envelope \\
\textbf{Tier 4: Mechanical Brake Clamping (Terminal Refusal / E-Stop)} &
\(t \ge t_{\mathrm{term}}\) (Plan exhausted) or watchdog trip & Planner
emits formal refusal; de-energize motor stages and drop spring brakes &
Actuator current zeroed; mechanical friction dissipates residual
momentum & Safe Torque Off (STO, IEC 61800-5-2 SIL 3) + spring brake
engagement & Hard fault latch; requires operator reset and safety
self-test & \(I_q = 0\text{ A}\); static holding torque
\(\tau_{\mathrm{brake}} \ge 1.5 \tau_{\mathrm{gravity}}\), zero
electrical power \\
\end{longtable}

The deterministic execution loop that unifies receding-horizon chunk
ingestion, \(\mathcal{C}^2\) polynomial blending, and stopping suffix
escalation on real-time microcontroller silicon is formalized below:

\subsection{\texorpdfstring{Receding-Horizon \(\mathcal{C}^2\) Action
Chunk Blending and Stopping Suffix
Protocol}{Receding-Horizon \textbackslash mathcal\{C\}\^{}2 Action Chunk Blending and Stopping Suffix Protocol}}\label{receding-horizon-mathcalc2-action-chunk-blending-and-stopping-suffix-protocol}

\textbf{Cadence \& Invariants:}

\begin{itemize}
\tightlist
\item
  \emph{Execution Cadence:} Hard real-time periodic interrupt
  (\(f = 1000\text{ Hz}\), period \(T_c = 1.0\text{ ms}\), maximum
  allowable timing jitter \(< 50\,\mu\text{s}\)).
\item
  \emph{Memory Invariant:} Zero dynamic heap allocation
  (\texttt{malloc}/\texttt{free} strictly forbidden; all state variables
  statically allocated in SRAM TCM).
\item
  \emph{Hardware Authority:} Direct evaluate/write access to motor phase
  PWM compare registers (\texttt{TIMx\_CCR} / \texttt{ePWM\ CMPA}) and
  fieldbus setpoint registers.
\end{itemize}

\textbf{Execution Protocol:}

\begin{enumerate}
\def\labelenumi{\arabic{enumi}.}
\item
  \textbf{Local State Sampling \& Temporal Milestone Evaluation (Stage
  1):} Sample current joint position \(\mathbf{q}_t\) and velocity
  \(\dot{\mathbf{q}}_t\) from optical encoder registers. Compute elapsed
  chunk time
  \(\tau_{\mathrm{chunk}} \leftarrow t_{\mathrm{now}} - \mathcal{C}_{\mathrm{active}}.t_{\mathrm{start}}\).
  If
  \(\tau_{\mathrm{chunk}} \ge \mathcal{C}_{\mathrm{active}}.t_{\mathrm{commit}}\)
  and \(\mathrm{BufferStatus} \ne \text{BUFFER\_SPLICED}\), transition
  immediately to Stage 4 (Stopping Suffix Escalation).
\item
  \textbf{Candidate Chunk Invariant Validation \& Boundary Check (Stage
  2):} If \(\mathcal{C}_{\mathrm{cand}}\) is resident with status
  \texttt{BUFFER\_READY} and
  \(t_{\mathrm{now}} \le \mathcal{C}_{\mathrm{active}}.t_{\mathrm{commit}}\):

  \begin{itemize}
  \item
    Evaluate boundary discontinuity at planned handover point
    \(t_{\mathrm{seam}}\):
    \[\Delta \mathbf{q} = \mathbf{Q}_B(0) - \mathbf{Q}_A(t_{\mathrm{seam}}), \quad \Delta \dot{\mathbf{q}} = \dot{\mathbf{Q}}_B(0) - \dot{\mathbf{Q}}_A(t_{\mathrm{seam}}), \quad \Delta \ddot{\mathbf{q}} = \ddot{\mathbf{Q}}_B(0) - \ddot{\mathbf{Q}}_A(t_{\mathrm{seam}}).\]
  \item
    Verify invariants:
    \(\|\Delta \dot{\mathbf{q}}\| \le \Delta \dot{\mathbf{q}}_{\mathrm{allow}}\)
    and \(\|\Delta \mathbf{q}\| \le \Delta \mathbf{q}_{\mathrm{tol}}\).
    If valid, proceed to Stage 3; else reject candidate and set
    \texttt{ERR\_SEAM\_DISCONTINUITY}.
  \end{itemize}
\item
  \textbf{Closed-Form \(\mathcal{C}^2\) Quintic Spline Synthesis (Stage
  3):} Determine available blend duration
  \(T_{\mathrm{blend}} \leftarrow \min(T_{\mathrm{nominal\_blend}}, \mathcal{C}_{\mathrm{active}}.t_{\mathrm{commit}} - t_{\mathrm{now}})\).
  Solve closed-form quintic coefficients
  \(\mathbf{a}_0 \dots \mathbf{a}_5\) for each active joint channel:
  \[\mathbf{a}_0 = \mathbf{q}_0, \quad \mathbf{a}_1 = \dot{\mathbf{q}}_0, \quad \mathbf{a}_2 = \frac{1}{2}\ddot{\mathbf{q}}_0,\]
  \[\mathbf{a}_3 = \frac{20(\mathbf{q}_1 - \mathbf{q}_0) - (8\dot{\mathbf{q}}_1 + 12\dot{\mathbf{q}}_0)T_{\mathrm{blend}} - (3\ddot{\mathbf{q}}_0 - \ddot{\mathbf{q}}_1)T_{\mathrm{blend}}^2}{2 T_{\mathrm{blend}}^3},\]
  \[\mathbf{a}_4 = \frac{30(\mathbf{q}_0 - \mathbf{q}_1) + (14\dot{\mathbf{q}}_1 + 16\dot{\mathbf{q}}_0)T_{\mathrm{blend}} + (3\ddot{\mathbf{q}}_0 - 2\ddot{\mathbf{q}}_1)T_{\mathrm{blend}}^2}{2 T_{\mathrm{blend}}^4},\]
  \[\mathbf{a}_5 = \frac{12(\mathbf{q}_1 - \mathbf{q}_0) - 6(\dot{\mathbf{q}}_1 + \dot{\mathbf{q}}_0)T_{\mathrm{blend}} - (\ddot{\mathbf{q}}_1 - \ddot{\mathbf{q}}_0)T_{\mathrm{blend}}^2}{2 T_{\mathrm{blend}}^5}.\]
  Verify interior peak acceleration
  \(\max_{t} |\ddot{\mathbf{q}}_{\mathrm{blend}}(t)| \le \mathbf{a}_{\max}\)
  and jerk
  \(\max_{t} |\dddot{\mathbf{q}}_{\mathrm{blend}}(t)| \le \dddot{\mathbf{q}}_{\max}\).
  If verified, atomically swap
  \(\mathcal{C}_{\mathrm{active}} \leftarrow \mathcal{C}_{\mathrm{blend\_then\_cand}}\),
  mark \texttt{BUFFER\_SPLICED}, and set
  \(\mathrm{Regime} \leftarrow \text{BLEND\_SPLINE}\). Else, drop
  candidate chunk and escalate to Stage 4.
\item
  \textbf{Stopping Suffix Fallback Engagement (Stage 4):} If
  \(\tau_{\mathrm{chunk}} \ge \mathcal{C}_{\mathrm{active}}.t_{\mathrm{commit}}\)
  without successful splice: set
  \(\mathrm{Regime} \leftarrow \text{STOPPING\_SUFFIX}\) and interpolate
  precomputed stopping polynomial \(\mathbf{q}_{\mathrm{stop}}(t)\)
  guaranteeing constant deceleration \(-\mathbf{a}_{\max}\) to zero
  velocity along path tangent. If remaining clearance
  \(d_{\mathrm{remain}} < d_{\mathrm{stop}}\) or external contact wrench
  \(\|F_{\mathrm{ext}}\| > F_{\max}\), assert Safe Torque Off
  (\texttt{STO}), engage mechanical spring holding brakes, and set
  \(\mathrm{Regime} \leftarrow \text{ESTOP\_CLAMP}\).
\item
  \textbf{Feedforward Torque Synthesis \& PWM Output Dispatch (Stage
  5):} Evaluate rigid-body inverse dynamics
  \(\boldsymbol{\tau}_{\mathrm{ff}} = \mathbf{M}(\mathbf{q}_t^*)\ddot{\mathbf{q}}_t^* + \mathbf{C}(\mathbf{q}_t^*, \dot{\mathbf{q}}_t^*)\dot{\mathbf{q}}_t^* + \mathbf{g}(\mathbf{q}_t^*)\),
  compute feedback voltage duty cycles
  \(\mathbf{u}_t^* = \boldsymbol{\tau}_{\mathrm{ff}} + \mathbf{K}_p(\mathbf{q}_t^* - \mathbf{q}_t) + \mathbf{K}_d(\dot{\mathbf{q}}_t^* - \dot{\mathbf{q}}_t)\),
  write \(\mathbf{u}_t^*\) to hardware PWM registers, and commit
  telemetry records to the lockless ring buffer.
\end{enumerate}

This real-time safety architecture enforces the precomputed suffix
invariant: every trajectory chunk must carry its own precomputed
stopping suffix resident in hardware memory upon admission, ensuring
that when the replacement is late, the nervous system has zero
microseconds of unbudgeted deliberation.

\section{How Plans Go Wrong}\label{sec-planning-11-6}

A plan can be feasible when issued and infeasible by the time it
executes. The latency between trajectory synthesis in the learned policy
and physical actuation by the motor drives creates a temporal window
during which the state of the machine drifts, the environment shifts, or
the geometric interface between consecutive trajectory chunks violates
dynamic constraints. Failures in planned motion distribute across four
distinct points of first visibility. A failure may become observable
before admission into the execution buffer, at concatenation during the
splicing of adjacent chunks, after execution begins as the physical body
interacts with its environment, or only across a sequence of successive
plans where no individual plan violates any local constraint.
Establishing where a failure first becomes visible dictates which
subsystem possesses the signal and the time budget required to detect
and reject it.

The earliest observable failure occurs when a generated chunk is
internally invalid before any command reaches the motor drivers. An
unconstrained neural policy or an imperfectly sampled trajectory can
propose waypoint sequences that exceed physical hardware limits. A
pre-execution bound checker evaluates candidate chunks against fixed
kinematic and dynamic limits during the admission phase. Rather than
evaluating arbitrary safety thresholds, the bound checker calculates the
exact margin \(M(t) = L_{\mathrm{limit}} - |x(t)|\) across every
discrete trajectory point for position \(q\), velocity \(\dot{q}\), and
acceleration \(\ddot{q}\). In an analytical example of a manipulator
joint governed by a velocity limit of
\(\dot{q}_{\mathrm{limit}} = 2.50\text{ rad/s}\), a candidate trajectory
containing a velocity sample of
\(|\dot{q}(120\text{ ms})| = 2.65\text{ rad/s}\) produces a negative
margin of \(M(120\text{ ms}) = -0.15\text{ rad/s}\). The pre-execution
checker intercepts the chunk before admission, emitting an explicit
fault record containing the violated joint channel, the numerical limit,
the timestamp of the violation, and the negative margin. Rejecting an
invalid plan at the admission boundary incurs zero dynamic disturbance
because the machine remains on its previously validated trajectory.

When two consecutive plans are each internally feasible, their
concatenation can still produce an instantaneous dynamic violation at
the temporal seam. If the terminal velocity or acceleration of the
active plan does not match the initial state of the candidate
replacement, splicing them across a discrete control cycle requires an
impulsive force. A boundary checker evaluates the kinematic mismatch
across the transition boundary \(t_{\mathrm{handover}}\) by comparing
the terminal derivatives of the active plan with the initial derivatives
of its candidate replacement. The checker evaluates the step differences
\(\Delta q = q_B(t_{\mathrm{handover}}) - q_A(t_{\mathrm{handover}})\),
\(\Delta \dot{q} = \dot{q}_B(t_{\mathrm{handover}}) - \dot{q}_A(t_{\mathrm{handover}})\),
and
\(\Delta \ddot{q} = \ddot{q}_B(t_{\mathrm{handover}}) - \ddot{q}_A(t_{\mathrm{handover}})\).
The admissible boundary velocity discontinuity
\(\Delta \dot{q}_{\mathrm{allow}}\) is derived from the peak torque
rating \(\tau_{\mathrm{max}}\) of the actuator and the total joint
inertia \(J\) over a single discrete control interval
\(\Delta t = 1.0\text{ ms}\). In an analytical example with link and
rotor inertia \(J = 0.080\text{ kg}\cdot\text{m}^2\) and an actuator
torque limit of \(\tau_{\mathrm{max}} = 40.0\text{ N}\cdot\text{m}\),
the maximum velocity step that the actuator can bridge across
\(\Delta t = 0.001\text{ s}\) without saturating torque is
\(\Delta \dot{q}_{\mathrm{allow}} = \tau_{\mathrm{max}} \Delta t / J = (40.0\text{ N}\cdot\text{m})(0.001\text{ s}) / (0.080\text{ kg}\cdot\text{m}^2) = 0.50\text{ rad/s}\).
If the boundary checker measures a velocity step exceeding
\(0.50\text{ rad/s}\), it rejects the candidate chunk before
concatenation, preventing torque saturation, gear tooth impact, and
mechanical fatigue in the drive train.

The physical severity of omitting boundary derivative checks is
demonstrated in the forensic autopsy of high-speed industrial mechanisms
during continuous receding-horizon replanning.

\protect\phantomsection\label{callout-incidentux2a-11.6}
\begin{fbxSimple}{callout-incident}{Landmark Engineering Incident:}{Trajectory Seam Step-Discontinuity Inverter Trip and Cycloidal Pin Shear}
\phantomsection\label{callout-incident*-11.6}

\textbf{Incident Overview:} In an automated high-speed food packaging
facility, a 4-axis delta pick-and-place manipulator running a learned
receding-horizon visuomotor policy suffered a catastrophic driveline
failure during a \(3.2\text{ m/s}\) transfer trajectory. The joint 1
input drive pin sheared completely, causing loss of kinematic
constraint, dropping a \(3.5\text{ kg}\) tool payload onto secondary
conveyor tooling, and destroying the motor drive's three-phase IGBT
inverter stage \sidenote{\footnotesize Step discontinuities at chunk handoffs excite
  resonance (\(15\text{--}40\text{ Hz}\)), tripping inverter protection
  (\(V_{CE(\mathrm{sat})} > 7\text{ V}\)) and shearing drive pins.}.

\textbf{Actuator Architecture \& Load Path:} The Joint 1 axis was driven
by a permanent magnet synchronous motor (PMSM, continuous rating
\(45\text{ N}\cdot\text{m}\), transient peak
\(110\text{ N}\cdot\text{m}\)) through a \(50:1\) cycloidal reduction
gearbox (\(J_{\mathrm{eff}} = 2.7\text{ kg}\cdot\text{m}^2\)). The
low-level field-oriented current controller executed at
\(1.0\text{ kHz}\) on a dedicated digital signal processor (DSP).

\textbf{Telemetry Snapshot \& Dynamic Overload:}

\begin{itemize}
\item
  At \(t = 0.0\text{ ms}\), Chunk \(A\) commanded Joint 1 to decelerate
  toward a target waypoint at
  \(\dot{q}_A(50\text{ ms}) = -0.45\text{ rad/s}\) with
  \(\ddot{q}_A(50\text{ ms}) = -8.5\text{ rad/s}^2\).
\item
  At \(t = 50.0\text{ ms}\), incoming Chunk \(B\) was admitted to the
  tracking loop. Chunk \(B\) began at identical position
  \(q_B(0) = q_A(50\text{ ms})\) (\(C^0\) continuous), but with an
  initial velocity of \(\dot{q}_B(0) = +0.15\text{ rad/s}\) (see Case A
  in Figure~\ref{fig-11-action-chunk-seam-continuity}).
\item
  Across a single \(T_c = 1.0\text{ ms}\) Fieldbus servo tick, the
  low-level position interpolator attempted to close the
  \(0.60\text{ rad/s}\) velocity discrepancy, commanding an apparent
  joint acceleration:
  \[\ddot{q}_{\mathrm{cmd}} = \frac{0.15 - (-0.45)}{0.001\text{ s}} = 600.0\text{ rad/s}^2\]
\item
  Through the joint's reflected inertia
  (\(J_{\mathrm{eff}} = 2.7\text{ kg}\cdot\text{m}^2\)), this commanded
  an instantaneous torque of
  \(\tau = (2.7)(600) = 1620\text{ N}\cdot\text{m}\).
\item
  The motor current loop slammed into its \(40\text{ A}\) hardware
  ceiling within \(150\,\mu\text{s}\), dumping peak saturation torque
  (\(\tau_{\mathrm{sat}} = 324\text{ N}\cdot\text{m}\)) into the
  transmission.
\item
  At \(t = 51.1\text{ ms}\), the cycloidal disk input pin sheared at its
  \(285\text{ N}\cdot\text{m}\) torsional yield limit. The sudden load
  release caused the motor rotor to accelerate uncontrollably
  (\(d\omega/dt > 15000\text{ rad/s}^2\)), generating massive back-EMF
  that pumped the DC bus voltage past \(480\text{ V}\), destroying the
  upper bridge IGBTs via desaturation overvoltage.
\end{itemize}

\textbf{Systems Root Cause:} The trajectory admission pipeline checked
only position continuity (\(C^0\) match) across the \(50\text{ ms}\)
replanning seam, ignoring the \(0.60\text{ rad/s}\) boundary velocity
step that demanded an instantaneous torque excursion \(295\%\) above the
continuous drive rating.

\textbf{Architectural Invariant:} Trajectory admission gates must
evaluate second-order derivative continuity (\(C^2\)) across all chunk
seams, enforcing closed-form quintic polynomial blending and hardware
acceleration rate clamps at the servo interpolator level.

\end{fbxSimple}

Even when a trajectory passes all pre-execution and boundary checks,
changes in the physical world during execution can render the remainder
of the plan infeasible. A grasped workpiece can slip, an unmodeled
obstacle can enter the workspace, or structural tracking error can
accumulate beyond allowable tolerances. Runtime monitors detect
changed-world infeasibility by observing four dedicated signals,
comprising joint state divergence, tracking error velocity, unexpected
contact wrenches, and the invalidation of geometric assumptions recorded
in the plan header. The tracking error allowance
\(\epsilon_{\mathrm{pos}}\) is derived directly from the geometric
clearance \(d_{\mathrm{clearance}}\) between the planned trajectory
envelope and the nearest obstacle boundary minus the structural
deflection allowance \(\delta_{\mathrm{deflect}}\). In an analytical
packaging cell with a budgeted clearance of
\(d_{\mathrm{clearance}} = 0.020\text{ m}\) and an unmodeled payload
deflection allowance of \(\delta_{\mathrm{deflect}} = 0.005\text{ m}\),
the position error trip threshold evaluates to
\(\epsilon_{\mathrm{pos}} = d_{\mathrm{clearance}} - \delta_{\mathrm{deflect}} = 0.015\text{ m}\).
If measured tracking error
\(|q_{\mathrm{actual}}(t) - q_{\mathrm{planned}}(t)|\) exceeds
\(0.015\text{ m}\), or if a six-axis force-torque sensor registers an
unmodeled contact wrench exceeding
\(F_{\mathrm{contact}} = 10.0\text{ N}\) during a free-space motion
phase, the runtime monitor invalidates the active plan and triggers
immediate reflex deceleration.

A plan also becomes infeasible when time elapses faster than replacement
plans arrive, leading to horizon exhaustion. Waiting for an execution
pointer to fall off the end of an array or relying on a zero-order hold
to repeat the final setpoint guarantees an abrupt loss of dynamic
authority, resulting in saturated position gains or violent impacts with
mechanical end stops. Horizon exhaustion must be detected by comparing
the remaining executable duration
\(T_{\mathrm{remain}}(t) = t_{\mathrm{end}} - t\) against the worst-case
stopping duration \(T_{\mathrm{stop}}\). The stopping budget is derived
from the maximum operational velocity and the constant emergency
deceleration limit \(a_{\mathrm{decel}}\) across all active joints,
defined by
\(T_{\mathrm{stop}}(t) = \max_i (|\dot{q}_i(t)| / a_{\mathrm{decel},i})\).
For an arm traveling at joint velocity \(\dot{q} = 1.20\text{ rad/s}\)
with an emergency deceleration limit of
\(a_{\mathrm{decel}} = 4.0\text{ rad/s}^2\), the minimum time required
to bring the joint to rest is
\(T_{\mathrm{stop}} = (1.20\text{ rad/s}) / (4.0\text{ rad/s}^2) = 0.300\text{ s}\)
(\(300\text{ ms}\)). The exhaustion detector fires precisely when
\(T_{\mathrm{remain}}(t) = 300\text{ ms}\), initiating controlled
deceleration toward a static equilibrium before the trajectory buffer
empties.

The most deceptive category of plan failure involves cross-plan
oscillation, which cannot be observed within any single trajectory
chunk. In this failure mode, every individual plan is smooth,
kinematically bounded, and collision-free, yet the sequence of plans
produces high-frequency reversals or spatial limit cycles. This behavior
arises when a policy alternates between competing local minima, such as
choosing between left and right avoidance paths around an obstacle,
flipping its intent every \(100\text{ ms}\) inference cycle. Single-plan
bound checkers and seam boundary checkers pass every individual proposal
because each chunk is locally feasible. Detecting cross-plan oscillation
requires a sequence monitor that maintains a bounded sliding history of
\(N\) previous plans over a rolling duration such as \(1.0\text{ s}\).
The monitor tracks the directional reversal count \(R_N\) and the total
acceleration variation
\(V_N = \sum_{k=1}^N \|\ddot{q}_k(0) - \ddot{q}_{k-1}(T)\|\). The
threshold for allowable directional reversals is derived from the
thermal dissipation limit of the motor windings under continuous torque
reversal. When \(R_N\) exceeds the thermal budget or when total
variation indicates thrashing without net spatial progress, the sequence
monitor rejects the oscillating candidate and commands the nervous
system to hold position or execute a deterministic fallback path.

Every detector in the planning pipeline relies on an explicit contract
between the trajectory generator and the execution runtime. If a plan is
delivered as a raw list of coordinates without metadata defining its
valid lifespan, its geometric clearance bounds, its boundary
derivatives, and its intended fallback behavior, the nervous system
cannot determine whether an ongoing motion remains safe or has drifted
into infeasibility. Making plan execution verifiable requires
encapsulating waypoints within a structured record that binds kinematic
trajectories to their temporal validity windows and monitoring
thresholds.

\section{The Trajectory Contract}\label{sec-planning-11-7}

The \textbf{trajectory record} is the immutable data structure that the
planning system hands to the execution runtime at every replanning
cycle. While the general schema, serialization format, and global
provenance rules for system logs are established in
Chapter~\ref{sec-nervous}, trajectory management introduces
domain-specific fields that enforce physical continuity and temporal
deadlines. Within its kinematic definition, the record stores the
geometric motion profile either as discrete kinematic setpoint arrays or
as continuous polynomial spline coefficients. It pairs this
representation with the sample cadence (such as
\(\Delta t = 10\text{ ms}\) for a \(100\text{ Hz}\) setpoint stream),
the total planned horizon \(T_{\mathrm{plan}}\) (such as
\(400\text{ ms}\)), the target coordinate frame identifier
\(\mathcal{F}\), the engineering SI units, and the exact initial-state
timestamp \(t_0\) matching the state measurement from which the path was
computed. Recording the coordinate frame and units directly in the
handoff prevents frame mismatches at the tracking layer, while the
explicit cadence and timestamp allow the low-level joint controller to
interpolate setpoints and evaluate kinematic derivatives without relying
on ambient system clocks.

Because physical feasibility depends on the state of the world when the
plan was formed, the record preserves the state estimate
\(\hat{\mathbf{x}}(t_0)\) and the environmental model against which the
path was validated. This snapshot captures the assumed obstacle
geometry, occupancy grid revisions, surface friction coefficients
(\(\mu = 0.80\) in an analytical dry-contact example), and estimated
payload properties (such as a mass of \(2.50\text{ kg}\)). If an
external object moves into the planned path during execution, an active
runtime monitor cannot determine whether the policy saw and bypassed the
obstacle or whether the path is in collision unless it can evaluate the
difference between the current sensor stream and the assumptions stored
in the record. Retaining this state snapshot allows the runtime
collision checker to compute relative displacement margins and trigger
path replanning only when the physical world diverges from the planner's
recorded assumptions.

Rather than storing a single boolean flag indicating that validation
passed, the trajectory record logs every checked kinematic and dynamic
limit alongside the validator software version and the computed safety
margins. These fields record the minimum obstacle clearance distance
(\(d_{\mathrm{clear}} = 42\text{ mm}\) in an analytical corridor
evaluation), the peak joint torque ratio
(\(\tau_{\mathrm{peak}}/\tau_{\mathrm{limit}} = 0.78\)), the maximum
anticipated Cartesian jerk, and any unresolved boundary conditions, such
as unverified contact states in downstream grasping phases. When
hardware operating conditions change, such as when an actuator heats up
and the thermal model decreases the maximum allowable torque
\(\tau_{\mathrm{limit}}\), the nervous system can inspect these recorded
margins directly. The runtime safety filter can immediately determine
whether an ongoing trajectory violates the newly constrained limit and
command a fallback maneuver without waiting for a complete replanning
cycle.

To manage the boundary between successive trajectory chunks, the record
defines the temporal contract and the terminal behavior of the motion.
This contract specifies the boundary state
\((\mathbf{q}(T_{\mathrm{plan}}), \dot{\mathbf{q}}(T_{\mathrm{plan}}), \ddot{\mathbf{q}}(T_{\mathrm{plan}}))\),
the required continuity order \(C^k\) (typically \(C^2\) continuity for
acceleration-bounded motor drives), the replacement deadline
\(t_{\mathrm{replace}}\), the fallback commitment point
\(t_{\mathrm{commit}}\), the terminal suffix trajectory, and the
terminal steady-state condition. Consider a trajectory chunk with an
active horizon of \(T_{\mathrm{plan}} = 400\text{ ms}\) and a stopping
duration of \(T_{\mathrm{stop}} = 150\text{ ms}\). The fallback
commitment point falls at
\(t_{\mathrm{commit}} = T_{\mathrm{plan}} - T_{\mathrm{stop}} = 250\text{ ms}\).
If a valid replacement plan is not admitted before
\(t = 250\text{ ms}\), the execution runtime abandons the expectation of
an incoming chunk and executes the precomputed terminal suffix, bringing
the arm to a controlled stop at \(t = 400\text{ ms}\) without exceeding
joint acceleration or jerk limits.

In embedded firmware on the MCU, this contract is materialized as a
contiguous, fixed-size memory structure
(\texttt{struct\ PlannedTrajectoryChunk}), detailed in
Table~\ref{tbl-11-trajectory-schema}. To guarantee deterministic
execution without heap fragmentation, the record allocates static arrays
for active trajectory setpoints and stopping suffix coefficients in
SRAM. The MCU execution engine maintains a lockless ping-pong buffer:
the Linux MPU populates the background record over shared memory or DMA,
while the real-time core consumes the active record. When validation
passes, an atomic pointer swap commits the new record to the tracking
loop at the next servo tick with zero copy overhead.

\begin{longtable}[]{@{}
  >{\raggedright\arraybackslash}p{(\linewidth - 8\tabcolsep) * \real{0.2000}}
  >{\raggedright\arraybackslash}p{(\linewidth - 8\tabcolsep) * \real{0.2000}}
  >{\raggedright\arraybackslash}p{(\linewidth - 8\tabcolsep) * \real{0.2000}}
  >{\raggedright\arraybackslash}p{(\linewidth - 8\tabcolsep) * \real{0.2000}}
  >{\raggedright\arraybackslash}p{(\linewidth - 8\tabcolsep) * \real{0.2000}}@{}}
\caption{\textbf{Abstract Data Schema for Real-Time Planned Trajectory
Chunk (\protect\texttt{PlannedTrajectoryChunk}).} Memory layout
(\(248\text{ bytes}\) total), data types, physical SI units, validation
invariants, and hardware enforcement mechanisms for trajectory proposals
exchanged across the MPU/MCU
boundary.}\label{tbl-11-trajectory-schema}\tabularnewline
\toprule\noalign{}
\begin{minipage}[b]{\linewidth}\raggedright
Byte Offset \& Field Name
\end{minipage} & \begin{minipage}[b]{\linewidth}\raggedright
Data Type \& Size
\end{minipage} & \begin{minipage}[b]{\linewidth}\raggedright
Physical Units \& Frame
\end{minipage} & \begin{minipage}[b]{\linewidth}\raggedright
Real-Time Validation Invariant \& Pre-Condition
\end{minipage} & \begin{minipage}[b]{\linewidth}\raggedright
Hardware Interceptor \& Enforcement Mechanism
\end{minipage} \\
\midrule\noalign{}
\endfirsthead
\toprule\noalign{}
\begin{minipage}[b]{\linewidth}\raggedright
Byte Offset \& Field Name
\end{minipage} & \begin{minipage}[b]{\linewidth}\raggedright
Data Type \& Size
\end{minipage} & \begin{minipage}[b]{\linewidth}\raggedright
Physical Units \& Frame
\end{minipage} & \begin{minipage}[b]{\linewidth}\raggedright
Real-Time Validation Invariant \& Pre-Condition
\end{minipage} & \begin{minipage}[b]{\linewidth}\raggedright
Hardware Interceptor \& Enforcement Mechanism
\end{minipage} \\
\midrule\noalign{}
\endhead
\bottomrule\noalign{}
\endlastfoot
\texttt{0x00..0x07}\texttt{sequence\_id} & \texttt{uint64\_t}(8 bytes) &
Monotonic Counter (dimensionless) &
\(\text{Seq} > \text{Seq}_{\text{last\_committed}}\) (strictly
monotonic; drops replayed packets) & Atomic register comparison in DMA
ring buffer (\(< 50\text{ ns}\)); drops stale chunks \\
\texttt{0x08..0x0F}\texttt{t\_generation\_ns} & \texttt{uint64\_t}(8
bytes) & Nanoseconds (\(\text{ns}\)) PTP Epoch &
\(t_{\mathrm{now}} - t_0 \le \Delta t_{\mathrm{transport\_max}}\)
(\(< 50\text{ ms}\); bounds pipeline transport lag) & Hardware timestamp
comparator; rejects proposals delayed by bus contention \\
\texttt{0x10..0x2F}\texttt{intent\_lease\_hash} &
\texttt{uint8\_t{[}32{]}}(32 bytes) & SHA-256 Digest (dimensionless) &
\(\text{Match}(\text{hash}, \mathcal{L}_{\mathrm{active\_intent}})\)
(validates intent lease lineage) & Fast SRAM hash lookup; prevents
trajectory execution for revoked intent leases \\
\texttt{0x30..0x37}\texttt{cadence\_dt\_ns} & \texttt{uint64\_t}(8
bytes) & Nanoseconds (\(\text{ns}\)) per sample &
\(\Delta t = 10\,000\,000\text{ ns}\) (\(10\text{ ms}\) nominal tracking
cadence, \(100\text{ Hz}\)) & Hardware timer interval match check;
prevents clock scale mismatch in interpolator \\
\texttt{0x38..0x3F}\texttt{horizon\_samples\_N} & \texttt{uint32\_t} +
pad(8 bytes) & Sample count (\(N \in [16, 64]\)) &
\(N \ge \lceil(T_{\mathrm{req}} + Q_{1-\epsilon}(L) + T_{\mathrm{res}})/T_c\rceil\)
(tail latency horizon invariant) & Pre-execution buffer size validator;
drops undersized action chunks before execution \\
\texttt{0x40..0x47}\texttt{t\_commit\_ns} & \texttt{uint64\_t}(8 bytes)
& Nanoseconds (\(\text{ns}\)) PTP Epoch &
\(t_{\mathrm{commit}} = t_{\mathrm{exp}} - T_{\mathrm{stop}}\) (hard
mechanical stopping deadline) & Hardware countdown timer; triggers
automatic deceleration suffix upon underflow \\
\texttt{0x48..0x5F}\texttt{boundary\_q\_vel\_acc} &
\texttt{float32{[}6{]}}(24 bytes) & \(\text{rad}\), \(\text{rad/s}\),
\(\text{rad/s}^2\) (Joint Frame) &
\(\|\dot{\mathbf{q}}(0) - \dot{\mathbf{q}}_{\mathrm{active}}(t_{\mathrm{seam}})\| \le \Delta \dot{\mathbf{q}}_{\mathrm{allow}}\)
(\(C^2\) continuity) & \(O(1)\) boundary derivative validator; prevents
cycloidal/harmonic pin torque shocks \\
\texttt{0x60..0xBF}\texttt{suffix\_poly\_coeffs} &
\texttt{float32{[}24{]}}(96 bytes) & Polynomial
\(\mathbf{a}_0 \dots \mathbf{a}_5\) per DOF &
\(T_{\mathrm{stop}} \ge v_{\mathrm{boundary}} / a_{\mathrm{decel\_limit}}\)
(stopping feasibility invariant) & MCU hardware interpolator registers;
directly drives deceleration on replan stall \\
\texttt{0xC0..0xEF}\texttt{assumed\_mass\_friction} &
\texttt{float32{[}12{]}}(48 bytes) & \(\text{kg}\), \(\mu\),
\(\text{N}\cdot\text{m}\), \(\text{mm}\) &
\(m_{\mathrm{payload}} \le m_{\mathrm{rated}}\),
\(\tau_{\mathrm{peak}} \le \tau_{\max}\),
\(d_{\mathrm{clear}} \ge d_{\mathrm{stop}}\) & Dynamic model parameter
validator; ensures model-based feedforward stability \\
\texttt{0xF0..0xF7}\texttt{prng\_seed\_hash} & \texttt{uint64\_t}(8
bytes) & Hash / Seed (dimensionless) & Link to exact diffusion denoising
/ ACT CVAE latent noise seed & Non-volatile provenance recorder;
guarantees deterministic simulation replay \\
\end{longtable}

Diagnosing cross-plan anomalies or mechanical failures requires offline
reproduction from logged data. The trajectory record links the generated
motion directly to the high-level intent identifier
(Chapter~\ref{sec-intent}), the planner binary version, the neural
network weights hash, the configuration parameter vector, and the
pseudorandom generator seed. Sampling-based motion planners and learned
diffusion policies exhibit stochastic behavior across runs. If an
execution failure occurs on the physical machine, attempting to
re-evaluate the planner in simulation with identical sensor inputs will
produce a different trajectory unless the original seed is supplied.
Logging the exact seed and configuration hash ensures that an engineer
can re-run the planning pipeline deterministically to isolate whether an
anomalous motion originated from policy instability, sensor noise, or an
unmodeled physical disturbance.

Finally, the record preserves a sequence of monotonic hardware
timestamps marking each stage of the plan lifecycle. These include the
planning request timestamp \(t_{\mathrm{req}}\), the generation
completion timestamp \(t_{\mathrm{gen}}\), the validation completion
timestamp \(t_{\mathrm{val}}\), the runtime admission timestamp
\(t_{\mathrm{adm}}\), the physical activation timestamp
\(t_{\mathrm{act}}\), and the formal expiration timestamp
\(t_{\mathrm{exp}}\). Preserving this timing chain enables the automated
detection of scheduling jitter, queue starvation, and communication
delays across the system bus. In fault-injection testing, where
synthetic delays are introduced to verify that safety monitors fire
before physical limits are exceeded (Chapter~\ref{sec-verification}),
this timestamp ledger provides the ground truth required to measure
exact detector latency. When an upstream computation stalls and
\(t_{\mathrm{adm}}\) slips past \(t_{\mathrm{commit}}\), comparing the
logged activation times against the planned replacement deadlines
demonstrates whether the nervous system transferred authority to the
terminal fallback path at the exact microsecond required by the body's
momentum.

\section{Systems Stress-Tests}\label{sec-planning-11-8}

A planned trajectory that is kinematically feasible and collision-free
at the moment of admission can become physically impossible while the
machine is in motion. Feasibility is not an intrinsic property of a
polynomial spline or a neural network rollout; it is a time-bound claim
that the state of the body and the geometry of the workcell will match
the assumptions recorded in the trajectory record. The following six
stress-tests evaluate whether an execution runtime and its underlying
safety monitors can detect when physical reality diverges from the
admitted plan, and whether the nervous system enforces safe termination
before the machine violates its hardware limits:

\subsection*{Stress-Test 1: Obstacle Displacement After Admission
(Dynamic Invalidation of
Clearance)}\label{stress-test-1-obstacle-displacement-after-admission-dynamic-invalidation-of-clearance}
\addcontentsline{toc}{subsection}{Stress-Test 1: Obstacle Displacement
After Admission (Dynamic Invalidation of Clearance)}

\emph{The scenario.} An articulated 6-axis arm sweeps an end-effector at
\(1.5\text{ m/s}\) along an admitted \(400.0\text{ ms}\) trajectory
certified with a \(50.0\text{ mm}\) clearance margin from an assembly
jig at \(t_{\mathrm{adm}} = 0.0\text{ ms}\). At execution time
\(t = 120.0\text{ ms}\), an external mechanism shifts the jig
\(70.0\text{ mm}\) directly into the tool flange path.

\emph{The systems dilemma.} The static clearance certificate is invalid.
The real-time perception pipeline must evaluate proximity invariants at
\(1000\text{ Hz}\), detecting spatial boundary penetration and
triggering hard deceleration before the remaining travel distance closes
to impact.

\emph{The systems verdict: Static Clearance Invalidation Law.} Clearance
guarantees expire the moment the physical environment changes; static
admission checks cannot protect a machine against dynamic workspace
alterations during active motion.

\subsection*{Stress-Test 2: Directional Oscillation Across the Seam
(Inter-Chunk Limit-Cycle
Chatter)}\label{stress-test-2-directional-oscillation-across-the-seam-inter-chunk-limit-cycle-chatter}
\addcontentsline{toc}{subsection}{Stress-Test 2: Directional Oscillation
Across the Seam (Inter-Chunk Limit-Cycle Chatter)}

\emph{The scenario.} A manipulator carrying a \(3.0\text{ kg}\) payload
executes consecutive \(200.0\text{ ms}\) trajectory chunks at
\(5\text{ Hz}\). Chunk \(k\) accelerates joint 1 to
\(\dot{q}_1 = +1.5\text{ rad/s}\) and stops at the seam. Chunk \(k+1\)
immediately accelerates joint 1 to \(\dot{q}_1 = -1.5\text{ rad/s}\).
Both chunks independently satisfy acceleration (\(12.0\text{ rad/s}^2\))
and jerk (\(60.0\text{ rad/s}^3\)) bounds.

\emph{The systems dilemma.} While each chunk is feasible in isolation,
stitching them together produces a \(2.5\text{ Hz}\) directional chatter
that excites gearbox structural resonances. Detecting limit cycles
requires the admission gate to maintain a sliding temporal window of
historical derivative metrics across seams.

\emph{The systems verdict: Inter-Chunk Continuity Invariant.}
Independent feasibility of individual trajectory chunks does not compose
into dynamic feasibility across their seam; multi-chunk history
verification is required to guarantee smooth physical execution.

\subsection*{Stress-Test 3: Tail Latency and the Fallback Commitment
Deadline (Hard Mechanical Time
Boundaries)}\label{stress-test-3-tail-latency-and-the-fallback-commitment-deadline-hard-mechanical-time-boundaries}
\addcontentsline{toc}{subsection}{Stress-Test 3: Tail Latency and the
Fallback Commitment Deadline (Hard Mechanical Time Boundaries)}

\emph{The scenario.} A diffusion planner generates \(300.0\text{ ms}\)
trajectory chunks with an analytical stopping duration
\(T_{\mathrm{stop}} = 120.0\text{ ms}\), establishing a fallback
commitment point at \(t_{\mathrm{commit}} = 180.0\text{ ms}\). Compute
latency degrades to \(210.0\text{ ms}\) under an interrupt storm. At
\(t = 180.0\text{ ms}\), the coprocessor signals inference is \(90\%\)
complete.

\emph{The systems dilemma.} Delaying fallback in the hope that compute
finishes guarantees an unmonitored overtravel if the coprocessor
deadlocks. The execution queue must strictly commit to the precomputed
terminal deceleration suffix at
\(t_{\mathrm{commit}} = 180.0\text{ ms}\), bringing the arm to rest
under maximum payload inertia.

\emph{The systems verdict: Fallback Commitment Law.} The fallback
commitment point is a hard mechanical boundary dictated by physical
stopping distance (\(T_{\mathrm{stop}} = v/a_{\max}\)), not a scheduling
heuristic that can yield to delayed computation.

\subsection*{Stress-Test 4: Intermediate Acceleration Spikes in
Trajectory Blends (The Midpoint Centripetal
Trap)}\label{stress-test-4-intermediate-acceleration-spikes-in-trajectory-blends-the-midpoint-centripetal-trap}
\addcontentsline{toc}{subsection}{Stress-Test 4: Intermediate
Acceleration Spikes in Trajectory Blends (The Midpoint Centripetal
Trap)}

\emph{The scenario.} A linear gantry moving an \(8.0\text{ kg}\) fixture
blends two perpendicular velocity segments
(\(\mathbf{v}_1 = [1.2, 0.0]\text{ m/s}\) to
\(\mathbf{v}_2 = [0.0, 1.2]\text{ m/s}\)) using a quadratic polynomial
over \(\Delta t = 50.0\text{ ms}\). Endpoint velocities and
accelerations meet the actuator limit (\(a_{\max} = 5.0\text{ m/s}^2\)).

\emph{The systems dilemma.} Evaluating the blend trajectory at the
midpoint reveals a peak centripetal acceleration of
\(a_{\mathrm{mid}} = 9.6\text{ m/s}^2\), demanding nearly double the
continuous force rating of the linear motors and tripping drive
overcurrent protection.

\emph{The systems verdict: Interior Path Feasibility Invariant.}
Boundary condition feasibility does not imply interior path feasibility.
Continuous dynamic envelope verification over the entire blend window is
required before admitting composite trajectories.

\subsection*{Stress-Test 5: Payload Drop During the Committed Horizon
(Feedforward Dynamic Model
Rupture)}\label{stress-test-5-payload-drop-during-the-committed-horizon-feedforward-dynamic-model-rupture}
\addcontentsline{toc}{subsection}{Stress-Test 5: Payload Drop During the
Committed Horizon (Feedforward Dynamic Model Rupture)}

\emph{The scenario.} An arm executes an admitted \(400.0\text{ ms}\)
trajectory lifting a \(6.0\text{ kg}\) casting at \(2.0\text{ m/s}^2\).
At \(t = 160.0\text{ ms}\), the casting slips from the gripper, reducing
payload mass to zero while motor drives continue applying full
feedforward torque.

\emph{The systems dilemma.} The sudden mass loss invalidates the
rigid-body dynamic model, causing an acceleration spike. Closed-loop
tracking error monitors detect the velocity divergence within
\(2.0\text{ ms}\), intercepting excess current commands before position
errors grow dangerous.

\emph{The systems verdict: Model Invalidation Response Law.} Sudden
contact and mass transitions instantly rupture feedforward dynamic
models; closed-loop feedback monitors must isolate disturbances before
unmodeled forces violate physical bounds.

\subsection*{Stress-Test 6: Infeasible Terminal Stopping Suffix (Missing
Deceleration
Guarantees)}\label{stress-test-6-infeasible-terminal-stopping-suffix-missing-deceleration-guarantees}
\addcontentsline{toc}{subsection}{Stress-Test 6: Infeasible Terminal
Stopping Suffix (Missing Deceleration Guarantees)}

\emph{The scenario.} A Cartesian gantry operates at
\(v = 2.0\text{ m/s}\) with a physical deceleration limit
\(a_{\mathrm{dec}} = 4.0\text{ m/s}^2\), requiring a minimum stopping
time \(T_{\mathrm{min}} = v / a_{\mathrm{dec}} = 500.0\text{ ms}\)
(\(s = 0.50\text{ m}\)). A planner emits a candidate plan with an active
horizon of \(350.0\text{ ms}\) and an undersized terminal suffix of
\(T_{\mathrm{stop}} = 200.0\text{ ms}\).

\emph{The systems dilemma.} Executing the \(200.0\text{ ms}\)
deceleration requires \(10.0\text{ m/s}^2\) braking---exceeding motor
capabilities by \(250\%\). The static admission gate must evaluate
\(T_{\mathrm{stop}} \ge v_{\mathrm{boundary}} / a_{\mathrm{dec}}\) and
reject the packet before execution.

\emph{The systems verdict: Terminal Suffix Completeness Law.} A
trajectory lacking a physically achievable terminal stopping suffix is
invalid at the moment of creation, regardless of whether a replacement
plan arrives on time.

\section{Summary}\label{sec-planning-11-9}

The physical machine does not execute a target coordinate or a geometric
curve. It consumes voltage, current, and torque commands indexed against
an uncompromising clock. An autonomous planner cannot emit a spatial
path and leave the timing to downstream interpretation. The concrete
artifact handed to the tracking hardware must be a time-parameterized
trajectory, defined as an explicit sequence of states, feedforward
inputs, and bounding envelopes specified at discrete time steps. A
geometric line in configuration space carries no information about
acceleration, jerk, or thermal dissipation in the motor windings. Only
when every waypoint is bound to an explicit timestamp can the
lower-level nervous system evaluate whether motor coils will overheat or
whether an instantaneous acceleration step demands more torque than the
power stage can deliver.

Sizing the temporal duration of this trajectory is an exercise in
latency budgeting rather than algorithmic preference. The active horizon
cannot be chosen arbitrarily or tuned by intuition. It is dictated by
the measured distribution of end-to-end replacement latency, the total
operational exposure time of the mission, and the permitted probability
of trajectory exhaustion. If the planning pipeline exhibits a
\(P_{99.9}\) replacement latency of \(180\text{ ms}\) under peak compute
load, an active horizon of \(150\text{ ms}\) guarantees periodic
starvation, forcing the machine to run out of instructions during
ordinary operation. Sizing the plan horizon requires calculating the
tail percentile of compute, bus transfer, and queue delays across the
entire mission duration, ensuring that the probability of the machine
outrunning its buffer remains below the target safety threshold.

A plan that satisfies every mechanical constraint in isolation can still
destroy the machine when spliced against its successor. Feasibility
inside a single trajectory requires that every interior state transition
respects torque limits, velocity ceilings, and kinematic reach. Yet the
primary failure mode in physical AI systems occurs at the seam between
two individually valid plans. When a newly generated trajectory arrives,
its initial state must match the physical machine's projected state at
the exact timestamp of handoff, preserving position, velocity, and
acceleration continuity. A discontinuity of \(0.05\text{ m/s}\) in joint
velocity across the seam forces an infinite instantaneous acceleration
demand, causing motor current loops to saturate and tripping overcurrent
protection hardware. Within-plan smoothness is insufficient without
explicit seam invariance.

Even when internally smooth and seam-continuous, trajectory feasibility
remains strictly conditional on the state measurements and physical
assumptions recorded at the moment of plan synthesis. A trajectory
generated under the premise of a rigid \(5.0\text{ kg}\) payload becomes
dynamic fiction if the gripper drops the workpiece or if surface
friction drops by half mid-motion. The planned feedforward torques then
act as unmodeled physical disturbances rather than tracking aids.
Feasibility is not a static property stamped onto a mathematical
function at creation time. It is a perishable contract that degrades as
real-world states drift away from the state estimate and environmental
model captured at the planning epoch.

Because compute stalls and communication delays are inevitable in real
computing hardware, the plan must define what happens when a replacement
fails to arrive on time. Terminal behavior is not an external exception
handler or an asynchronous recovery routine patched onto the runtime. It
is an integral, mandatory suffix of the trajectory itself. Every
admissible plan must terminate in a dynamically achievable stopping
trajectory or an invariant holding state, sized from first principles
using maximum deceleration limits and current velocities. If an
end-effector moves at \(v = 1.5\text{ m/s}\) and the axis deceleration
limit is \(a_{\mathrm{dec}} = 3.0\text{ m/s}^2\), the minimum physical
stopping duration is \(T_{\mathrm{min}} = 500\text{ ms}\), requiring
\(0.375\text{ m}\) of clear deceleration travel. A plan that omits this
safe terminal suffix is invalid before execution begins, because when a
deadline is missed, the nervous system has no time to deliberate.

The planning subsystem concludes its task by delivering a fully
specified artifact for independent verification rather than by asserting
its own safety. This handoff artifact contains five concrete elements,
comprising a dated, time-indexed trajectory of reference states and
feedforward inputs, declared physical bounds on velocity, acceleration,
and torque, boundary conditions enforcing continuity at the handoff
seam, explicit environmental validity assumptions, and a finite,
mechanically achievable terminal behavior.

\protect\phantomsection\label{callout-takeawaysux2a-11.7}
\begin{fbxSimple}{callout-takeaways}{Key Chapter Takeaways}{}
\phantomsection\label{callout-takeaways*-11.7}

\begin{itemize}
\tightlist
\item
  \textbf{Trajectories Are Time-Parameterized Contracts, Not Geometric
  Paths:} Physical actuators consume current, voltage, and torque
  indexed against an uncompromising clock. An admissible plan must
  explicitly bind every waypoint to timestamps, state vectors
  (\(\mathbf{q}, \dot{\mathbf{q}}, \ddot{\mathbf{q}}\)), and feedforward
  torques to evaluate winding heating and acceleration limits.
\item
  \textbf{Action Chunk Horizons Are Dimensioned by Tail Replacement
  Latency:} Planning horizons (\(H\)) cannot be chosen arbitrarily.
  Sizing requires bounding the \(P_{99.9}\) sense-to-proposal
  replacement latency across the total mission duration to prevent
  periodic trajectory buffer starvation
  (\(t_{\text{horizon}} \ge t_{\text{replace}}^{P_{99.9}}\)).
\item
  \textbf{Seam Invariance Trumps Internal Plan Smoothness:} Splicing two
  individually valid trajectories together without matching boundary
  derivatives (\(\Delta \dot{\mathbf{q}} = 0.05\text{ m/s}\)) demands
  infinite instantaneous acceleration
  (\(\ddot{\mathbf{q}} \to \infty\)), causing current loop saturation
  and inverter trips. Trajectories must enforce \(\mathcal{C}^2\)
  continuity at handoff seams.
\item
  \textbf{Terminal Stopping Suffixes Are Mandatory Architectural
  Components:} A trajectory lacking an explicit, physically achievable
  stopping suffix
  (\(T_{\text{stop}} \ge v_{\text{bound}}/a_{\text{dec}}\)) is invalid
  at creation. When an upstream replanning deadline is missed, the
  nervous system executes the terminal stopping deceleration without
  deliberating.
\item
  \textbf{Feedforward Torques Degrade into Disturbances Upon Model
  Rupture:} Feasibility is conditional on payload and contact
  assumptions. A sudden loss of payload mass
  (\(6.0\text{ kg} \to 0\text{ kg}\)) turns planned feedforward torques
  into destructive disturbance forces that must be isolated by real-time
  tracking error monitors within milliseconds.
\end{itemize}

\end{fbxSimple}

Planning proposes this contract, but it cannot be permitted to grade its
own work. Chapter~\ref{sec-enforcement} constructs the runtime enforcer,
the deterministic nervous system layer that validates these five
components against physical truth before a single microsecond of
execution is admitted to the actuators.

\chapter{\texorpdfstring{Deterministic Safety
Enforcement\chaptersubtitle{From Command to Force, and the Part That Can Say No}}{Deterministic Safety Enforcement}}\label{sec-enforcement}

\begin{figure}[H]

\centering{

\includegraphics[width=1\linewidth,height=\textheight,keepaspectratio]{index_files/mediabag//Users/VJ/GitHub/MLSysBook-integrate-physical/physical/book/chapters/12-enforcement/figures/fig_locator.pdf}

}

\caption{\label{fig-locator}\textbf{Systems Organ Focus: Nervous.} In
the physical AI systems architecture, the deterministic nervous system
sits between the unverified deliberative brain and the dynamical body,
serving as the sole authority permitted to translate advisory commands
into physical force or refuse them outright. Operating at microsecond
control cadences, the enforcer projects candidate actions onto
forward-invariant safe sets and manages a multi-tier fallback hierarchy
to guarantee physical containment.}

\end{figure}%

\emph{How does an unprivileged software proposal become physical torque,
and what guarantees that the gatekeeper can always say no in time?}

When a high-capacity learned policy encounters an out-of-distribution
sensory input, it can output an erratic, high-energy control command in
a single forward pass. Waiting for the model to diagnose its own failure
or relying on Linux user-space exceptions is physically impossible once
setpoints reach motor drivers. Real-time safety enforcement is the
deterministic hardware and mathematical barrier that stands between
unverified neural proposals and physical actuators. Operating at
microsecond cadences on dedicated microcontroller silicon, the enforcer
uses Control Barrier Functions and minimal-intervention quadratic
programs (CBF-QP) to project candidate actions onto forward-invariant
safe sets---guaranteeing that the machine never breaches physical
constraints while preserving learned task competence.

\newpage{}

\protect\phantomsection\label{callout-objectiveux2a-12.1}
\begin{fbxSimple}{callout-objective}{Learning Objectives}{}
\phantomsection\label{callout-objective*-12.1}

After studying this chapter, you will be able to:

\begin{enumerate}
\def\labelenumi{\arabic{enumi}.}
\tightlist
\item
  Explain why a correction rate is set by how fast the plant deviates,
  not by preference.
\item
  Predict what an actuator at saturation does to a loop that does not
  know it has saturated.
\item
  Explain why delay in a loop causes oscillation rather than lag alone.
\item
  State the tracking error bound a controller guarantees, and inset a
  safety margin by it.
\end{enumerate}

\emph{Core chapter takeaway.} The nervous system's job end to end: track
a permitted command into force, and refuse one that should not become
force. The enforcer's margin is only as good as the tracking error of
the loop beneath it. Safety enforcers depend on inner feedback loops
tracking commanded braking currents within bounded error; tracking and
refusal are a unified mathematical contract.

\end{fbxSimple}

\section{The Deterministic Gatekeeper}\label{sec-enforcement-12-1}

\lettrine{A}{learned policy produces numbers,}\lettrine{b}{ut} numbers
do not move mass. In classical computing, software executes ``behind
glass'': algorithmic errors raise exceptions, invalid matrix
multiplications return NaN tokens, and failed operations can be retried
without altering the physical universe. In Physical AI, computation
directly controls moving mass, dissipates electrical power, and commands
mechanical force. Once an actuator energizes a motor winding, momentum
integrates, kinetic energy transfers, and mechanical deformation cannot
be undone by a software catch block. Between the output tensor of an
upstream neural network and the magnetic field inside an electric motor
sits the nervous system of the machine. As mapped in the four-tier
architectural hierarchy of Figure~\ref{fig-locator}, this subsystem
occupies the critical boundary where unverified computational proposals
meet physical mechanics. It carries two distinct responsibilities that
define its physical role: it must translate a permitted command into
physical force under deterministic timing deadlines, and it must
maintain the independent authority to refuse execution entirely.
Upstream neural networks, vision-language-action (VLA) models, and
high-level trajectory optimizers are unverified advisory processes. They
can propose a target velocity, generate a motion waypoint, or request an
end-effector wrench, but they do not possess the hardware authority to
drive an actuator. The \textbf{permission path} is the sole component on
the machine permitted to energize a motor winding, and it is the only
component endowed with the capability to refuse an upstream request.

Tracking and safety enforcement must be co-located at this physical
machine boundary rather than delegated to upstream software layers.
High-level learned models operate over deep computational pipelines with
variable execution times, often exhibiting inference latencies between
\(100\text{ ms}\) and \(500\text{ ms}\). In contrast, the physical plant
is governed by mechanical momentum, where a delay of \(5\text{ ms}\) in
a high-speed motion can allow an unmonitored mass to exceed its
allowable travel envelope. If the authority to refuse were placed
upstream, network dropouts, bus contention spikes, or memory management
pauses in an inference runtime would silently disable safety arbitration
while the physical mechanism continues its motion. The nervous system
executes on local microcontroller or real-time processor hardware
directly wired to motor drives and joint encoders, evaluating commands
at update rates such as \(1000\text{ Hz}\). At this boundary, safety is
evaluated against the physical state that exists in the current
millisecond rather than a stale estimate captured in an earlier sensor
frame.

An enforcer cannot guarantee safety beyond the physical fidelity of the
control loop beneath it. Every mathematical safety barrier, whether
formulated as a geometric spatial boundary or a maximum torque limit,
assumes that the tracking controller follows a commanded trajectory
within a bounded error. Consider an analytical example where a robotic
manipulator must maintain a minimum clearance of \(20\text{ mm}\) from a
rigid fixture while moving at \(1.0\text{ m/s}\). If the underlying
tracking loop exhibits a worst-case position tracking error of
\(e_{\max} = 15\text{ mm}\) under load disturbances, the safety enforcer
must evaluate its barriers against a boundary expanded by that full
\(15\text{ mm}\). A safety certificate that assumes perfect tracking
will permit a commanded path that physically collides with the obstacle.
The certified clearance of the enforcer is strictly bounded by the
tracking performance of the underlying loop. When motor saturation,
joint compliance, or unmodeled friction causes the tracking error to
grow, the available safety margin contracts by the same amount.

Even when mathematically sound, deterministic enforcers can fail because
of physical and computational constraints in the systems architecture.
An enforcer requires a finite processing budget to evaluate barrier
functions, solve quadratic programs, or project proposals onto an
admissible action set. If the constraint solver experiences a \(P_{99}\)
latency tail that exceeds its allocated \(1\text{ ms}\) execution
budget, the system faces an immediate failure mode, choosing between
executing an unverified proposal or applying an uncoordinated emergency
brake that may cause structural damage. Enforcers also fail when
supplied with stale or phase-lagged state estimates. If a low-pass
sensor filter introduces a \(20\text{ ms}\) delay while an actuator
accelerates at \(5.0\text{ m/s}^2\), the computed state lags true
position by \(1.0\text{ mm}\), systematically corrupting barrier
evaluations. Furthermore, shared hardware modes, such as common power
supplies, shared communication buses, or single-point clock generators,
can cause the tracking loop, the enforcer, and the sensor pipeline to
fail simultaneously.

Constructing an enforcement architecture requires ordering the machine's
software layers from inner physical reflexes to outer policy interfaces.
The lowest layer governs the inner feedback mechanics that turn
high-frequency setpoints into motor currents. Above this tracking
foundation sits the barrier evaluation logic, which checks incoming
proposals against invariant state constraints. When a proposal violates
a constraint, a minimal intervention filter modifies the command just
enough to restore feasibility, preserving the policy's intended motion
to the extent physical limits allow. If no safe projection exists or if
state estimation uncertainty exceeds verified thresholds, the
architecture shifts authority to a deterministic fallback handler to
bring the machine to a controlled stop.

The construction of this entire chain begins at the interface between
the safety filter and the tracking controller. To determine whether an
action is safe to execute, the enforcer does not need to analyze inner
control loop gains, motor winding dynamics, or pulse-width modulation
switching frequencies. It consumes a single contract from the controller
beneath it, which is an explicit bound on the worst-case tracking error
between where the machine was commanded to be and where its physical
mass actually resides.

\section{What Control Hands Over}\label{sec-enforcement-12-2}

The tracking contract isolates the enforcer from the internal mechanics
of the plant. A feedback controller operating over continuous physical
dynamics makes many promises internally, but only one number passes
upward across the systems interface. For any commanded trajectory
\(r(t)\) satisfying admissible velocity and acceleration limits, the
physical state \(x(t)\) remains within a bounded distance
\(\|x(t) - r(t)\| \le \epsilon_{\text{track}}\). The enforcer does not
inspect transfer functions, root loci, or
proportional-integral-derivative gains. The methods used to synthesize
feedback gains and prove asymptotic stability belong to classical
control theory (\citeproc{ref-slotine1991applied}{Slotine and Li 1991};
\citeproc{ref-astrom2021feedback}{Åström and Murray 2021}). To the
software sitting above the loop, the controller is a black box that
turns reference setpoints into physical motion, guaranteed to stay
within the error bound \(\epsilon_{\text{track}}\) of the commanded
path.

This bound is a physical distance rather than a mathematical
abstraction. In a closed-loop mechanical system with active stiffness
\(K\) subject to an unmodeled external disturbance force \(d(t)\), the
steady-state position error balances the disturbance against the
restoring control effort. When disturbances are bounded such that
\(|d(t)| \le d_{\text{max}}\), the worst-case steady-state tracking
error bound \(\epsilon_{\text{track}}\) derives directly from the
stiffness of the closed loop.
\[\epsilon_{\text{track}} = \frac{d_{\text{max}}}{K}\] Consider an
analytical example of a linear positioning stage with closed-loop
stiffness \(K = 5000\text{ N/m}\) operating in an environment where
friction and payload variations introduce an unmodeled disturbance force
bounded by \(d_{\text{max}} = 25.0\text{ N}\). The maximum steady-state
error is
\(\epsilon_{\text{track}} = 25.0\text{ N} / 5000\text{ N/m} = 0.005\text{ m}\),
which corresponds to \(5.0\text{ mm}\). When the reference command
executes a step of \(100.0\text{ mm}\), the carriage accelerates,
settles into the bounded corridor, and comes to rest. At steady state,
the physical carriage does not sit precisely at \(100.0\text{ mm}\). It
rests anywhere within the interval
\([95.0\text{ mm}, 105.0\text{ mm}]\).

That guarantee holds only within a defined envelope of physical
operating conditions. The controller provides the bound
\(\epsilon_{\text{track}}\) under three physical conditions. First, the
reference trajectory cannot demand accelerations beyond what the
actuators can deliver. Second, external disturbances must remain
strictly below \(d_{\text{max}}\). Third, the actuator must operate
below its internal current and voltage limits. If an upstream policy
commands a trajectory that demands an acceleration of
\(15.0\text{ m/s}^2\) from an actuator whose current limit caps
acceleration at \(10.0\text{ m/s}^2\), the feedback loop saturates. The
integrator winds up, phase lag accumulates, and the physical mechanism
falls behind the reference by an amount that grows with every
millisecond of saturation. If an unexpected external load exceeds
\(25.0\text{ N}\), the disturbance overpowers the restoring force of
stiffness \(K\), pulling the carriage outside the \(\pm 5.0\text{ mm}\)
corridor. Outside its rated envelope, the controller provides no bound
at all.

Because physical reality can diverge from the reference command by up to
\(\epsilon_{\text{track}}\), every geometric and dynamic margin
evaluated by an enforcer must be inset by that exact distance. If a
robot arm must maintain a physical clearance of at least
\(20.0\text{ mm}\) from a fragile fixture, and the underlying tracking
loop guarantees an error bound
\(\epsilon_{\text{track}} = 5.0\text{ mm}\), the enforcer cannot
evaluate safety against the nominal geometry. It must inset the
admissible state space by adding a virtual buffer of \(5.0\text{ mm}\),
checking that the commanded reference maintains a clearance of at least
\(25.0\text{ mm}\) from the obstacle surface. If worn bearings, thermal
expansion, or altered motor parameters loosen the true tracking error
from \(5.0\text{ mm}\) to \(12.0\text{ mm}\) without the enforcer
updating its internal model, the true physical clearance to the fixture
shrinks from \(15.0\text{ mm}\) to \(8.0\text{ mm}\). The enforcer
believes it is operating with a comfortable safety margin while the
physical hardware moves dangerously close to collision. A degraded
tracking bound silently erodes safety margins across the entire software
stack.

A tracking bound that is never checked at run time is an assumption
rather than a contract. Because environmental conditions fluctuate and
mechanical components degrade, the nervous system must publish
\(\epsilon_{\text{track}}\) alongside the explicit operational
conditions that validate it, while monitoring the instantaneous tracking
error \(e(t) = \|x(t) - r(t)\|\) on every tick of the control cycle. If
sensor feedback reveals that \(e(t)\) approaches or exceeds
\(\epsilon_{\text{track}}\), the tracking contract has broken. When a
disturbance force spikes to \(40.0\text{ N}\) or an actuator enters
thermal throttling, the nervous system cannot treat the resulting
divergence as a normal transient. It must detect the contract violation
immediately, invalidating the spatial assumptions used by downstream
filters before unmodeled deflection causes physical damage.

\protect\phantomsection\label{callout-ruleux2a-12.2}
\begin{fbxSimple}{callout-rule}{Architect’s Rule of Thumb:}{The Inset Margin Invariant}
\phantomsection\label{callout-rule*-12.2}

Every spatial and dynamic safety boundary evaluated by an enforcer must
be inset by the underlying tracking loop's verified worst-case tracking
error plus the transport lag displacement:
\(\delta_{\text{margin}} \ge \epsilon_{\text{track}} + v \cdot \tau_{\text{delay}} + \epsilon_{\text{est}}\).
An enforcer that evaluates barrier certificates against nominal
geometric boundaries without insetting will permit trajectories that
physically collide with obstacles whenever tracking error reaches its
certified bound.

\end{fbxSimple}

The mathematical synthesis of tracking controllers, observer design, and
disturbance rejection proofs are treated in depth by the control
literature (\citeproc{ref-slotine1991applied}{Slotine and Li 1991}). For
the systems engineer building physical AI architectures, the
responsibility begins where control guarantees end. The tracking
controller provides the mechanism to translate permitted references into
physical currents, but it has no semantic awareness of obstacles, task
goals, or system-level hazards. It will execute a command into a solid
wall with the same closed-loop fidelity that it brings to free-space
motion. Before any proposed action is handed to this tracking machinery,
the system must determine whether the proposal is permitted to become
force at all.

\section{The Authority to Refuse}\label{sec-enforcement-12-3}

The permission path is not a transmission line that relays actions while
logging safety statistics in the background. It is a safety gate that
admits actuation only when specific physical criteria are satisfied,
treating every incoming proposal as unsafe until proven otherwise. In
conventional automation, control pipelines operate on trust, executing
commands under the assumption that the upstream software component
produces valid setpoints. In physical AI, that assumption fails by
design. A learned policy operating in high-dimensional perception space
has no mathematical guarantees of stability, no bounded output variance,
and no internal representation of actuator current limits or mechanical
stress. Because the learned model cannot be held accountable for
physical destruction, the nervous system must not treat its outputs as
authoritative directives. Upstream models produce proposals, but only
the nervous system possesses the authority to turn a proposal into
force.

This division creates an asymmetric authority model between the brain
and the nervous system, formalized in systems engineering as the
\textbf{Simplex Architecture Pattern} (\citeproc{ref-sha2001using}{Sha
2001}). The Simplex pattern resolves the tension between high
performance and safety by decomposing the control system into four
structural components:

\begin{enumerate}
\def\labelenumi{\arabic{enumi}.}
\tightlist
\item
  \textbf{The Advanced Controller}: An unverified, high-capacity learned
  model (such as a Vision-Language-Action network or diffusion policy
  executing on a Linux MPU) that optimizes for complex, open-ended task
  performance.
\item
  \textbf{The Baseline Safety Controller}: A certified, high-assurance
  deterministic controller (such as a Lyapunov-based posture regulator
  or precomputed quintic deceleration spline executing on bare-metal MCU
  firmware) that guarantees stability and physical containment within a
  known envelope.
\item
  \textbf{The Safety Manager}: Deterministic arbitration logic that
  continuously evaluates the state of the machine against
  forward-invariant barrier sets, tracking error bounds, and execution
  deadlines.
\item
  \textbf{The Simplex Switch}: An atomic hardware or real-time firmware
  multiplexer that routes actuation signals to the motor bridge,
  selecting the Advanced Controller's output when certified safe and
  instantaneously falling back to the Baseline Controller when safety
  envelopes are endangered.
\end{enumerate}

The \textbf{Simplex architecture pattern} is a safety-critical systems
structure featuring an Advanced (unverified) Controller, a certified
Baseline (Safety) Controller, a Safety Manager, and an arbitration
switch (Simplex Switch) that autonomously strips the Advanced Controller
of authority and transfers control to the Baseline Controller before the
physical state breaches recoverable bounds.

Under this Simplex contract, the upstream policy never commands the
plant directly. Instead, communication across the Dual-Brain boundary
operates through time-bounded \textbf{intent leases}. The Linux MPU
transmits candidate control vectors as structured packets over shared
SRAM mailboxes using the Remote Processor Messaging (RPMSG) protocol.
Each proposal payload contains a monotonic sequence identifier, a
microsecond timestamp, an explicit lease duration
(\(\tau_{\text{lease}}\)), candidate setpoints, and a cyclic redundancy
check (CRC32). The real-time microcontroller inspects the lease on every
control cycle. If the lease is valid, the state is within safe set
boundaries, and tracking error remains bounded, the Simplex switch
permits the proposal (or its minimally projected variant) to drive the
actuators. If the lease expires without renewal, tracking diverges, or
the proposal threatens physical boundaries, the Simplex switch revokes
permission and engages the Baseline Safety Controller
unilaterally.\sidenote{\footnotesize In canonical safety microcontrollers, hardware
  trip zones (\texttt{EPWM\_TZFRC}, \texttt{TIMx\_BKIN}) operate
  independently of the CPU. Hardware break circuits force complementary
  PWM outputs into high-impedance or low-side clamp states within
  \(<50\text{ ns}\), preventing shoot-through before software ISRs
  execute.} The upstream policy is not consulted, does not negotiate,
and possesses no hardware mechanism to override the switch.

When an incoming proposal conflicts with physical constraints, the
enforcer chooses between two distinct interventions: command filtering
and an outright veto. Command filtering applies when the proposal is
structurally sound but marginally exceeds an admissible kinematic or
dynamic boundary. If an analytical model of a six-axis manipulator shows
that a proposed end-effector velocity vector demands a joint speed
\(4\%\) above a rated motor limit of \(3.0\text{ rad/s}\), the nervous
system can project the reference onto the boundary of the safe set using
control barrier functions or quadratic programming. This projection
preserves the geometric direction of the motion while scaling down
magnitude to obey the \(3.0\text{ rad/s}\) ceiling. Projection is
permissible only when the deviation is local, continuous, and cleanly
separated from catastrophic hazards. When a proposal exhibits structural
invalidity---such as a discontinuous jump in commanded position across
consecutive frames, a request for acceleration that exceeds maximum
physical torque capabilities, or a trajectory vector aimed directly into
an immovable obstacle---geometric projection becomes dangerous. A
discontinuous or wildly infeasible proposal indicates that the upstream
perception or inference pipeline has entered an unmodeled failure mode.
Attempting to project an ungrounded hallucination onto the nearest safe
boundary executes a synthetic motion that has no relationship to task
intent. In such cases, the enforcer must veto the proposal entirely,
tripping the Simplex switch and dropping authority to a local fallback.

To maintain authority when the upstream system collapses, the refusal
mechanism must operate with complete local autonomy. The nervous system
cannot depend on the brain to report its own failures or request a
shutdown. Consider an analytical example where an upstream policy
publishes trajectory setpoints at \(20\text{ Hz}\), or one frame every
\(50\text{ ms}\), across an inter-process communication bus to a nervous
system running a joint tracking loop at \(1000\text{ Hz}\). The nervous
system couples incoming frames to a hardware-timed watchdog timer. If
the upstream inference engine experiences an operating system scheduling
stall, an out-of-memory exception, or a CUDA kernel panic, the flow of
setpoints halts. If the intent lease duration is set to
\(\tau_{\text{lease}} = 40\text{ ms}\) and the watchdog timeout is
configured for \(t_{\text{timeout}} = 10\text{ ms}\) of heartbeat
silence, the nervous system detects the fault within milliseconds. It
does not freeze the last commanded actuator currents, which would cause
an arm moving at \(1.2\text{ m/s}\) to drift under momentum and strike
environmental boundaries. Instead, the local enforcer immediately severs
the upstream command channel and transitions to an autonomous position
hold, synthesizing a smooth deceleration trajectory that brings joint
velocities to zero at a controlled rate of \(8.0\text{ m/s}^2\) while
remaining within current limits.

Severing upstream authority requires a well-defined physical state at
the lowest level of the hardware interface, which forms the
\textbf{zero-command floor}. In software engineering, halting an
execution pipeline often means returning a null pointer or clearing a
register to zero. In physical systems, writing a zero to a motor drive
register produces radically different mechanics depending on actuator
topology. Setting current commands to zero in a low-friction,
direct-drive motor removes all resistive torque, causing a robot arm
holding a \(5.0\text{ kg}\) payload to drop under gravity in free fall.
Conversely, commanding zero velocity in a stiff position controller with
high proportional gains commands maximum peak current to arrest motion
instantaneously, which can shear gear teeth or buckle structural
linkages if momentum is high. The system designer must select the
zero-command state to match the mechanical physics of the plant. In
geared transmissions with high backdrivability, dynamic braking shorts
the motor windings through low-resistance shunts or MOSFET bridges,
converting the back-electromotive force of the spinning rotor into
electrical current that generates a braking torque proportional to
velocity. This dynamic damping smoothly dissipates kinetic energy as
heat without requiring active software control, though it cannot resist
static gravitational torque once motion ceases. When static holding or
power-off safety is required, the zero-command state de-energizes
electromechanical spring-applied friction brakes, clamping the joints
mechanically after dynamic braking has reduced angular velocity below a
threshold of \(0.1\text{ rad/s}\) to prevent excessive friction lining
wear.

Refusal is therefore not an emergency exception handler invoked only
during rare failures. It is the steady-state operating discipline of the
entire machine. On every millisecond tick of the control clock, the
nervous system verifies that an incoming proposal is topologically
valid, that communication heartbeats remain unbroken, that the commanded
reference lies strictly inside invariant safety sets, and that tracking
error remains bounded. If any condition fails, the enforcer says no,
substituting an autonomous fallback or engaging the zero-command floor.
Because the enforcer holds the ultimate authority to deny actuation, its
implementation must not depend on the very components it is designed to
constrain.

\section{Where the Enforcer Sits}\label{sec-enforcement-12-4}

The enforcer must sit at a specific physical and logical boundary in the
signal pipeline, determined by the timescale of the dynamics it governs
and the latencies of the channels feeding it. Two natural boundaries
exist in every machine that translates learned proposals into motion.
The first boundary sits upstream of the feedback control loop, where the
system executes reference checking. A reference checker validates
planned trajectory waypoints, coordinate targets, or velocity vectors
emitted by a high-level planner before those targets enter the tracking
loop. Because it operates in kinematic space over a forward horizon, an
analytical reference checker operating at \(50\text{ Hz}\) can evaluate
a \(300\text{ ms}\) trajectory segment against static geometric
envelopes, workspace boundaries, and self-collision models. The strength
of reference checking is wide spatial foresight. Its limitation is
complete blindness to the actual state of the plant. A reference checker
verifies where the machine intends to go, not where it goes. If a
planned path maintains a clearance of \(5.0\text{ mm}\) from a steel
bracket, the reference checker marks the command valid. Yet if the
downstream tracking controller exhibits a transient position overshoot
of \(10.0\text{ mm}\) when accelerating a heavy load, the physical arm
strikes the bracket despite passing every upstream validation step.

The second boundary sits downstream of the feedback controller, directly
in front of the power electronics, where the system executes command
checking. A command checker inspects the raw physical quantities
produced by the tracking loop, evaluating motor phase currents,
hydraulic valve spool displacements, pulse-width modulation duty cycles,
and joint torque references immediately before writing them to hardware
registers. Operating at the inner control loop frequency, typically
\(1000\text{ Hz}\) with an execution budget under \(100\text{ \mu s}\),
the command checker enforces hard instantaneous limits. It prevents
electrical overcurrent, structural over-torque, and thermal runaway by
clipping or zeroing signals that exceed the physical ratings of the
plant. What the command checker gains in physical certainty, it
surrenders in temporal foresight. Operating on a single discrete time
step \(\Delta t = 1.0\text{ ms}\), it cannot evaluate future state
evolution. An instantaneous torque command of
\(18\text{ N}\cdot\text{m}\), which is well within an analytical motor
rating of \(45\text{ N}\cdot\text{m}\), appears entirely safe to a
post-control monitor. However, if that torque accelerates a
\(10\text{ kg}\) linkage toward a rigid stop at \(1.5\text{ m/s}\), the
kinetic energy of the mechanism will exceed the maximum deceleration
capacity of the actuator over the remaining travel distance. The command
checker detects no fault until the linkage is inches from impact, at
which point no allowable reverse torque can prevent a collision.

This split reveals the fundamental trade between latency and visibility.
Reference checkers possess multi-step geometric visibility but suffer
from high latency and blind spots regarding closed-loop tracking
dynamics. Operating at planning rates between \(10\text{ Hz}\) and
\(50\text{ Hz}\), a reference filter cannot respond to unmodeled plant
resonance, sensor noise amplification, or sudden control instability
occurring at frequencies above its update rate. If a feedback controller
develops a \(60\text{ Hz}\) limit cycle oscillation, the reference
setpoints remain stationary while the physical linkages shake.
Conversely, command checkers respond within microseconds on every
millisecond clock tick, eliminating latency, but their one-step horizon
creates dynamic traps where every instantaneous torque is legal until an
irrecoverable state is reached.

The placement of the enforcer also determines how it integrates into the
computing fabric. In a sequential in-line architecture, every signal
passes synchronously through the enforcement logic before reaching the
next stage. A reference checker filters the trajectory before the
controller reads it, or a command clamp modifies the torque register
before the motor drive samples it. Sequential execution guarantees zero
uninspected commands, but it inserts the worst-case execution time of
the enforcer directly into the critical latency path. For an analytical
\(1000\text{ Hz}\) loop with a total deadline of \(1000\text{ \mu s}\),
allocating a \(P_{99}\) latency of \(250\text{ \mu s}\) to an in-line
safety check shrinks the budget available for feedback computation. If
the enforcer overruns its execution budget, the entire cycle faults. In
a parallel asynchronous architecture, the enforcer executes on a
separate co-processor or an isolated processor core, monitoring the
state of the machine and the command stream concurrently without
intercepting every cycle. When the parallel monitor detects an invariant
violation, it asserts a hardware line to open a safety contactor or
switch the motor drive to dynamic braking. Parallel execution preserves
control loop timing, but the asynchronous detection and communication
introduce an enforcement lag bounded by the monitor cycle time and bus
jitter, analytical examples showing \(3.0\text{ ms}\) to
\(8.0\text{ ms}\) of delay. To compensate for this reaction delay, the
system designer must expand the physical safety margins, slowing the
machine or increasing its minimum clearance buffers.

Because neither single location provides both dynamic foresight and
instantaneous protection, physical AI architectures resolve the tension
through a \textbf{dual-stage enforcement pattern}. This pattern maps
cleanly onto a canonical heterogeneous dual-brain System-on-Chip (SoC)
architecture:

\begin{itemize}
\tightlist
\item
  \textbf{Upstream Stage (Pre-Control Corridor Filter on Application
  MPU):} Operating at \(20\text{--}50\text{ Hz}\) on an application
  microprocessor core (a high-capacity unprivileged application
  processor), this stage accepts multi-step trajectory chunks emitted by
  the learned policy. It verifies that the planned path maintains an
  analytical stopping corridor under worst-case braking parameters and
  validates geometric clearance against static workspace envelopes
  before packing the waypoints into an RPMSG intent lease.
\item
  \textbf{Downstream Stage (Instantaneous CBF-QP Safety Shield on
  Bare-Metal Real-Time MCU):} Operating inside a strict
  \(1000\text{ Hz}\) (\(1000\,\mu\text{s}\)) timer interrupt loop on a
  deterministic microcontroller core (an isolated hard real-time safety
  microcontroller), this stage receives the proposed setpoint from
  shared SRAM, verifies tracking error bounds, solves an active-set
  Control Barrier Function quadratic program in pre-allocated static
  SRAM, and clamps raw current and torque commands directly before
  writing to PWM timer registers.
\end{itemize}

The pre-control stage prevents the system from entering dynamic traps by
ensuring that a valid deceleration trajectory always exists, while the
post-control stage catches loop instability, mechanical overshoots, and
transient sensor dropouts before current reaches the motor windings.

Placing the enforcer correctly is useless if its inputs share failure
modes with the component it checks. An enforcer may only evaluate
quantities that it can establish independently of the proposal path. If
a pre-control corridor filter calculates collision distances using a
depth map generated by the same neural perception model that planned the
motion, an out-of-distribution visual artifact will deceive both the
planner and the enforcer simultaneously. Independent enforcement
requires independent sensing and isolated computation. On heterogeneous
dual-brain SoCs and industrial robotics controllers, the bare-metal
safety MCU reads secondary optical encoders and hardware limit switches
wired directly to its own GPIO and hardware timer capture registers over
dedicated buses, completely isolated from the Linux virtual memory space
and operating system scheduler. A collision enforcer must read dedicated
safety lidar or ultrasonic proximity sensors whose signal conditioning
relies on deterministic threshold logic rather than learned inference.
Where no physically independent sensor exists, such as when estimating
the friction coefficient of a wet roadway or the variable mass of an
unmodeled payload, the systems engineer cannot construct an independent
physical guarantee. In those operating regimes, the safety assertion is
conditional, its validity bounded by the statistical uncertainty of the
learned estimator rather than guaranteed by the nervous system. Every
invariant that the enforcer checks, whether evaluated upstream across a
geometric corridor or downstream at the motor terminals, ultimately
assumes that the physical plant retains the control authority to arrest
its own motion before crossing a forbidden boundary.

\protect\phantomsection\label{callout-ruleux2a-12.3}
\begin{fbxSimple}{callout-rule}{Architect’s Rule of Thumb:}{The Simplex Independence Rule}
\phantomsection\label{callout-rule*-12.3}

A safety monitor must never evaluate a state estimate produced by the
component it is designed to constrain, nor share an un-preemptible
communication bus, memory matrix, or clock tree with unverified
processes. Where independent sensing is physically impossible, the
safety claim is conditional on model uncertainty rather than guaranteed
by hardware.

\end{fbxSimple}

\section{Stopping Envelopes}\label{sec-enforcement-12-5}

To guarantee that a machine never enters a forbidden state, the nervous
system must know at every moment that an unassisted stop remains
physically achievable. This guarantee defines the \textbf{stopping
envelope}, the set of all state trajectories from which maximum
available braking can bring the system to rest without violating a
spatial or kinematic boundary. Maintaining this guarantee requires
reserving both distance and time for an emergency deceleration maneuver
that the machine may never execute during nominal operation. Knowing
continuously that a stop is still available costs achievable speed, and
that trade is the design. Every increase in operating velocity expands
the distance required to arrest momentum, consuming the spatial buffer
between the machine and its physical constraints.

The geometry of the stopping envelope derives directly from Newtonian
kinematics combined with the signal latency of the nervous system.
Consider a body moving along a single axis at instantaneous velocity
\(v\). When an enforcer detects a boundary violation or receives an
unsafe proposal, a finite transport delay \(\tau_{\text{delay}}\)
elapses before actuators can generate opposing force. This total delay
compounds perception exposure time, transport latency over the internal
communications bus, enforcer evaluation cycle time, and the physical
current rise time in the motor windings. During \(\tau_{\text{delay}}\),
the body continues to travel under its pre-existing velocity, traversing
a reaction distance of \(v \cdot \tau_{\text{delay}}\). Once braking
torque engages, assuming a constant maximum deceleration
\(a_{\text{max}}\), the time required to arrest the remaining velocity
is \(v / a_{\text{max}}\), during which the body traverses a braking
distance of \(v^2 / (2 a_{\text{max}})\).

A pure kinematic calculation assumes that the underlying feedback
controller tracks the reference trajectory with zero error. In physical
reality, external disturbances, unmodeled joint friction, and sensor
quantization cause the physical plant to diverge from the reference
trajectory by an instantaneous tracking error bounded by
\(\epsilon_{\text{track}}\), established by the low-level loop in
Section~\ref{sec-enforcement-12-2}. If an obstacle sits at a coordinate
\(x_{\text{obs}}\), placing the nominal stopping target exactly at
\(x_{\text{obs}}\) causes a collision whenever the tracking error
reaches its positive bound. The enforcer must therefore inset the
allowable reference boundary away from the obstacle by
\(\epsilon_{\text{track}}\). Combining reaction displacement, braking
displacement, and the tracking bound yields the total minimum stopping
distance,

\[d_{\text{stop}}(v) = v \cdot \tau_{\text{delay}} + \frac{v^2}{2 a_{\text{max}}} + \epsilon_{\text{track}}\]

Any trajectory proposal that places the machine at a distance
\(d < d_{\text{stop}}(v)\) from a physical constraint is unsafe, even if
the instantaneous velocity vector points away from the obstacle.

Consider an analytical example of an automated transport platform moving
through an indoor industrial facility. The platform runs a
\(50\text{ Hz}\) trajectory planner, communicates over an internal bus
with deterministic message delivery, and uses an enforcer co-processor
operating at \(1000\text{ Hz}\). Summing the perception pipeline latency
(\(30\text{ ms}\)), the worst-case message transport interval
(\(20\text{ ms}\)), and the motor drive torque response time
(\(10\text{ ms}\)) yields a total nervous system delay
\(\tau_{\text{delay}} = 60\text{ ms}\). The mechanical brakes and motor
drives provide a maximum deceleration
\(a_{\text{max}} = 2.0\text{ m/s}^2\), and encoder verification
establishes a worst-case tracking error bound
\(\epsilon_{\text{track}} = 0.050\text{ m}\). At a moderate operational
velocity \(v = 1.5\text{ m/s}\), the reaction distance during delay is
\(1.5\text{ m/s} \times 0.060\text{ s} = 0.090\text{ m}\), and the
braking distance is
\((1.5\text{ m/s})^2 / (2 \times 2.0\text{ m/s}^2) = 0.563\text{ m}\).
Adding the \(0.050\text{ m}\) tracking inset yields a total required
stopping clearance \(d_{\text{stop}}(1.5\text{ m/s}) = 0.703\text{ m}\).
If the system increases nominal speed by two-thirds to
\(v = 2.5\text{ m/s}\), the delay displacement rises to
\(0.150\text{ m}\) and the braking displacement scales quadratically to
\(1.563\text{ m}\), expanding the required envelope to
\(1.763\text{ m}\). A \(67\%\) increase in velocity demands a \(151\%\)
increase in clear clearance ahead of the chassis. In physical autonomous
mobile robots and automated guided vehicles, this dynamic clearance
envelope is continuously scanned and enforced by dedicated hardware,
such as the time-of-flight safety laser scanner shown in
Figure~\ref{fig-12-real-laser-scanner}.

\begin{figure}[t]

\centering{

\includegraphics[width=0.85\linewidth,height=\textheight,keepaspectratio]{/Users/VJ/GitHub/MLSysBook-integrate-physical/physical/book/chapters/12-enforcement/figures/fig12_real_laser_scanner.jpg}

}

\caption{\label{fig-12-real-laser-scanner}\textbf{Industrial Safety
Laser Scanner and Dynamic Protective Field Enforcer.} A physical
time-of-flight (ToF) safety laser scanner commonly mounted on the
leading bumper of automated guided vehicles (AGVs) and autonomous mobile
robots (AMRs). The unit utilizes an internal motor-driven mirror
rotating at \(20\text{--}50\text{ Hz}\) to sweep pulsed infrared laser
beams across a \(270^\circ\) planar scanning plane, resolving obstacle
distances via time-of-flight measurements. To enforce the dynamic
stopping distance
\(d_{\text{stop}}(v) = v \cdot \tau_{\text{delay}} + \frac{v^2}{2 a_{\max}} + \epsilon_{\text{track}}\),
the scanner continuously switches between velocity-dependent zone
profiles: outer \emph{warning fields} that prompt soft advisory
deceleration in software when breached, and inner \emph{protective
fields} wired through dual-channel safety relays to trigger immediate
Category 0 (Safe Torque Off) or Category 1 hardware stops when
penetrated.}

\end{figure}%

Because velocity and obstacle distances change continuously, the
enforcer cannot evaluate feasibility only when a new trajectory begins.
It must verify feasibility at every cycle of the low-level loop. In a
phase-plane representation where velocity is plotted against distance to
the nearest constraint, the stopping equation defines a parabolic
boundary that separates safe states from inescapable collision states.
The region between the physical obstacle boundary and the tracking-inset
curve represents the margin consumed by lower-level control uncertainty.
At every cycle, the enforcer projects the current state \((x, v)\)
forward along an emergency stop trajectory. If the projected final rest
position crosses the inset boundary, or if the current operating state
crosses into the forbidden zone of the phase plane, the enforcer vetoes
the planner's proposal and commands maximum deceleration.

The physical necessity of maintaining a clear stopping envelope levies
an inevitable tax on operational throughput. In any deployment where
sensing range is finite or environmental visibility is restricted by
geometry, the maximum detection horizon \(d_{\text{sense}}\) dictates a
hard ceiling on legal velocity. If an onboard optical sensor can confirm
an unobstructed corridor only up to \(2.0\text{ m}\) ahead around a
blind warehouse corner, the system cannot permit a speed exceeding
\(2.68\text{ m/s}\) under the parameters of the earlier example,
regardless of how quickly the underlying machine learning model can
compute actions. To run faster, the engineer must either increase
braking torque \(a_{\text{max}}\), compress signal latency
\(\tau_{\text{delay}}\), or tighten the tracking error
\(\epsilon_{\text{track}}\) through higher loop gains. When hardware
limits prevent further reduction in those terms, the only remaining
parameter is velocity.

The parameters governing the stopping envelope are rarely stationary
during physical operation. When a robotic manipulator grasps an
unmodeled \(15\text{ kg}\) steel billet, the combined inertia of the arm
increases, reducing the maximum angular deceleration \(a_{\text{max}}\)
that joint actuators can produce without exceeding their peak mechanical
torque ratings. Similarly, on a mobile robot, moisture or oil on a
concrete floor reduces the tire friction coefficient, lowering the
maximum tractive deceleration before wheel slip occurs and widening the
tracking error \(\epsilon_{\text{track}}\). Repeated hard decelerations
also drive up rotor and inverter temperatures, triggering thermal
derating logic that temporarily lowers available braking torque. An
enforcer that treats \(a_{\text{max}}\) as a static datasheet constant
will underestimate the stopping envelope precisely when the machine is
under heavy load. The enforcer must adapt \(d_{\text{stop}}(v)\)
dynamically, pulling updated friction estimates, payload mass estimates,
and thermal derating factors from the nervous system at runtime.

Dynamic envelope adaptation functions only as long as the enforcer
possesses accurate state estimates and the actuators retain the physical
capacity to supply the requested forces. If surface contamination drops
effective traction to near zero, or if a sustained duty cycle pushes
winding temperatures past their thermal ceilings, the deceleration
capability assumed in the safety calculation evaporates. Under those
conditions, the mathematical envelope expands toward infinity while the
physical machine continues along its ballistic path. The moment the
required deceleration torque exceeds what the physical hardware can
produce, the enforcer transitions from a deterministic guardian into an
impotent observer.

\section{When the Enforcer Runs Out of
Authority}\label{sec-enforcement-12-6}

In control engineering, \textbf{actuator saturation} is treated as a
bounded nonlinearity that introduces phase lag, degrades loop stability,
and leads to windup in feedback controllers. In the architecture of a
physical AI system, saturation is more immediate. It is the instant the
permission path loses the physical capacity required to keep its safety
warranty. The enforcer makes guarantees on the premise that it can
interpose a corrective control action \(u_{\text{corr}}\) whenever the
machine approaches a forbidden state boundary. When an actuator reaches
its mechanical, electrical, or thermal limit, that corrective force
cannot be delivered at any commanded magnitude. The enforcer cannot
protect a boundary using force that the physical body cannot generate.

The consequence is that a stopping envelope computed under nominal
assumptions is invalidated the moment authority is consumed. Consider an
analytical example of a mobile base with total mass
\(m = 100\text{ kg}\) moving at a forward velocity
\(v_0 = 2.0\text{ m/s}\). Under nominal conditions, the drive motors
supply a peak braking force \(F_{\text{nom}} = 400\text{ N}\), producing
a maximum deceleration
\(a_{\text{nom}} = F_{\text{nom}} / m = 4.0\text{ m/s}^2\). The nominal
stopping distance is
\[d_{\text{stop,nom}} = \frac{v_0^2}{2 a_{\text{nom}}} = \frac{(2.0\text{ m/s})^2}{2 \times 4.0\text{ m/s}^2} = 0.50\text{ m}.\]
If steady-state resistance, such as climbing an incline or overcoming
friction, consumes an authority fraction \(\alpha \in [0, 1)\) of peak
actuator capacity, the remaining deceleration falls to
\(a_{\text{eff}} = (1 - \alpha) a_{\text{nom}}\). The actual stopping
distance expands to
\[d_{\text{stop}}(\alpha) = \frac{v_0^2}{2 (1 - \alpha) a_{\text{nom}}} = \frac{d_{\text{stop,nom}}}{1 - \alpha}.\]
The loss of stopping margin is the difference
\[\Delta d_{\text{stop}} = d_{\text{stop}}(\alpha) - d_{\text{stop,nom}} = d_{\text{stop,nom}} \left( \frac{\alpha}{1 - \alpha} \right).\]
If the vehicle uses \(75\%\) of its motor torque to maintain steady
speed on a slope (\(\alpha = 0.75\)), only \(100\text{ N}\) of force
remains for deceleration, yielding
\(a_{\text{eff}} = 1.0\text{ m/s}^2\). The stopping distance expands
from \(0.50\text{ m}\) to \(2.00\text{ m}\). The \(1.50\text{ m}\) of
safety margin was spent during steady-state motion, long before the
enforcer issued a deceleration command.

When tracking controllers beneath the enforcer employ integral action,
saturation introduces an additional latency penalty. If an actuator
reaches its limit while tracking error
\(\epsilon(t) = x_{\text{des}}(t) - x(t)\) persists, the integrator
continues accumulating error, driving the internal control state far
beyond the physical ceiling. When the enforcer intervenes to command
full braking, the motor does not reverse force immediately. The physical
output remains pinned at the forward saturation limit while the
accumulated integrator error unwinds below the clamping threshold. In an
analytical example where tracking error accumulates over
\(200\text{ ms}\), the integrator unwinding phase can introduce a
recovery delay \(\tau_{\text{unwind}} = 150\text{ ms}\) before net
negative force develops. At \(v_0 = 2.0\text{ m/s}\), this delay adds an
unbraked coasting distance
\(\Delta d_{\text{delay}} = v_0 \tau_{\text{unwind}} = 0.30\text{ m}\)
directly to the stopping envelope. Control literature addresses this
behavior through anti-windup clamping, back-calculation, and conditional
integration (see Åström and Hägglund
(\citeproc{ref-astrom2006advanced}{2006}); Franklin et al.
(\citeproc{ref-franklin1998digital}{1998})). From the perspective of
systems enforcement, the requirement is not merely to stabilize the
control loop, but to measure this unwinding delay and budget its
distance penalty inside the forward stopping horizon.

Actuator saturation is one of the few hardware failures with a direct,
deterministic detector inside the nervous system. At each execution
cycle of the motor tracking loop, typically running at
\(1000\text{ Hz}\) (\(1\text{ ms}\) tick interval), the nervous system
compares the commanded effort \(u_{\text{cmd}}\) with the realizable
actuator limit
\(u_{\text{act}} = \text{clip}(u_{\text{cmd}}, u_{\text{min}}, u_{\text{max}})\).
A non-zero clipping difference,
\(|u_{\text{cmd}} - u_{\text{act}}| > 0\), or an inverter status
register indicating current limit saturation, provides an immediate
signal that the actuator has exhausted its authority. By tracking the
duration and magnitude of this clamping signal, the nervous system
monitors available control authority on every millisecond tick without
relying on noisy external state estimation.

Because authority consumption alters stopping distance immediately, the
enforcer cannot treat actuator limits as static parameters in an offline
calculation. Instead, the nervous system must fold real-time authority
consumption directly into the stopping envelope. When an actuator
allocates baseline torque \(u_{\text{bias}}\) to hold a payload against
gravity or overcome an opposing grade, the enforcer recalculates
effective deceleration as
\(a_{\text{eff}} = (u_{\text{max}} - |u_{\text{bias}}|) / m\) before
verifying the next trajectory segment. As authority is consumed, the
allowable speed ceiling decreases continuously. If authority depletion
causes the calculated stopping distance to exceed the clear sensing
horizon, the system reaches an authority deficit. The enforcer cannot
allow the machine to maintain its speed. It must command deceleration
while sufficient physical authority still remains.

Authority depletion is state information that must be exported upstream
to the planner. A learned policy operating at \(20\text{ Hz}\) that
continues commanding accelerations to a motor saturated at
\(1000\text{ Hz}\) is executing without physical grounding. If the
planner does not receive saturation feedback, it will continue
generating trajectories based on acceleration profiles the body cannot
execute, widening the divergence between model state and physical
reality. The nervous system communicates saturation events to the policy
layer as an active constraint violation. When an authority alarm
triggers, the planner's active trajectory is revoked, requiring the
model to re-plan from the machine's true, authority-constrained state.

Every dynamic stopping envelope and safety margin depends on the
physical capacity of the machine to generate corrective force within the
remaining distance. When that capacity is restricted by load, grade, or
thermal derating, the space of executable safe trajectories contracts.
An enforcer cannot verify safety simply by checking whether the current
state is collision-free. It must confirm that the current state permits
an executable sequence of future control actions that terminate safely
under the actuator authority that actually exists.

\protect\phantomsection\label{callout-ruleux2a-12.4}
\begin{fbxSimple}{callout-rule}{Architect’s Rule of Thumb:}{The Actuator Authority Invariant}
\phantomsection\label{callout-rule*-12.4}

A Control Barrier Function is mathematically valid only while physical
actuators retain sufficient current, voltage, and thermal headroom to
deliver the computed corrective force: \(u^* \in [u_{\min}, u_{\max}]\).
When steady-state gravity or friction consumes the actuator's operating
envelope, the enforcer must dynamically contract the safe set before
control authority at the margin vanishes.

\end{fbxSimple}

\section{Safe Sets as Conditional
Permission}\label{sec-enforcement-12-7}

A safe set gives the enforcer a concrete predicate that can be evaluated
on every tick of the control clock. In control theory, a \textbf{safe
set} \(\mathcal{C}\) is defined as the zero-superlevel set of a
continuously differentiable scalar certificate
\(h(x): \mathcal{X} \to \mathbb{R}\), such that
\[\mathcal{C} = \{x \in \mathcal{X} \mid h(x) \ge 0\}, \quad \partial \mathcal{C} = \{x \in \mathcal{X} \mid h(x) = 0\}.\]
The mathematical foundation of \textbf{forward invariance}, established
by Blanchini (\citeproc{ref-blanchini1999set}{Blanchini 1999}) and
extended through Control Barrier Functions by Ames et al.
(\citeproc{ref-ames2014control}{Ames et al. 2014},
\citeproc{ref-ames2019control}{2019}),\sidenote{\footnotesize Nagumo (1942)
  established that dynamic systems remain within closed sets
  \(\mathcal{C}\) if vector fields lie in Bouligand tangent cones at
  boundaries. Ames et al.~synthesized Nagumo's theorem with Lyapunov
  methods, creating Control Barrier Functions (CBFs) that enforce
  invariance via convex affine control constraints.} guarantees that if
the initial physical state begins inside \(\mathcal{C}\), there exists
an admissible control input sequence that prevents the trajectory from
ever leaving \(\mathcal{C}\).

A \textbf{forward-invariant safe set} is a designated sub-region of the
state space
\(\mathcal{C} = \{\mathbf{x} \in \mathcal{X} \mid h(\mathbf{x}) \ge 0\}\)
possessing the property that any state trajectory originating within
\(\mathcal{C}\) remains strictly within \(\mathcal{C}\) for all forward
time under admissible control inputs, despite bounded disturbances and
model discrepancies.

\protect\phantomsection\label{callout-mathux2a-12.5}
\begin{fbxSimple}{callout-math}{Fermi Math:}{Forward Invariance and Lie Derivatives}
\phantomsection\label{callout-math*-12.5}

Consider a continuous-time physical plant whose dynamics are affine in
the control input: \[\dot{x} = f(x) + g(x)u,\] where
\(x \in \mathcal{X} \subseteq \mathbb{R}^n\) denotes the mechanical
state vector (positions, velocities) and
\(u \in \mathcal{U} \subseteq \mathbb{R}^m\) represents the commanded
actuator efforts (phase currents, motor torques). The time derivative of
the safety certificate along the system trajectory is computed using Lie
derivatives:
\[\dot{h}(x, u) = \nabla h(x) \cdot \dot{x} = \frac{\partial h}{\partial x} f(x) + \frac{\partial h}{\partial x} g(x) u = L_f h(x) + L_g h(x) u.\]
By Nagumo's viability theorem and Ames's Control Barrier Function
framework, the safe set \(\mathcal{C}\) is rendered forward-invariant if
the control input \(u\) satisfies the differential inequality:
\[\dot{h}(x, u) \ge -\gamma(h(x)),\] where
\(\gamma: \mathbb{R} \to \mathbb{R}\) is an extended
class-\(\mathcal{K}_\infty\) function (typically selected as a linear
contraction rate \(\gamma(h) = \gamma_0 h\) with \(\gamma_0 > 0\)).
Substituting the Lie derivatives into this condition yields:
\[L_f h(x) + L_g h(x) u \ge -\gamma_0 h(x).\] Because the dynamics are
affine in \(u\), this condition transforms the nonlinear set-invariance
requirement into an instantaneous affine linear half-space inequality on
the control input: \[a(x)^T u \le b(x),\] where the constraint normal
vector \(a(x) \in \mathbb{R}^m\) and boundary threshold
\(b(x) \in \mathbb{R}\) are given by:
\[a(x) \triangleq -L_g h(x)^T, \quad b(x) \triangleq L_f h(x) + \gamma_0 h(x).\]

\end{fbxSimple}

\protect\phantomsection\label{callout-definitionux2a-12.6}
\begin{fbxSimple}{callout-definition}{Definition:}{Control Barrier Function (CBF-QP)}
\phantomsection\label{callout-definition*-12.6}

A \textbf{Control Barrier Function (CBF-QP)} is a certified,
forward-invariant safety filter formulated as a sub-millisecond convex
quadratic program
(\(\min_{\mathbf{u}} \frac{1}{2}\|\mathbf{u} - \mathbf{u}_{\text{nom}}\|^2 \text{ s.t. } L_f h(\mathbf{x}) + L_g h(\mathbf{x})\mathbf{u} \ge -\gamma(h(\mathbf{x})), \mathbf{u}_{\min} \le \mathbf{u} \le \mathbf{u}_{\max}\))
that continuously evaluates state-space constraints
\(h(\mathbf{x}) \ge 0\) along control-affine plant dynamics, admitting
unverified learned policy proposals \(\mathbf{u}_{\text{nom}}\) without
modification in the safe set interior and computing the least-squares
minimal orthogonal projection \(\mathbf{u}^*\) on the boundary to
guarantee physical containment.

\end{fbxSimple}

When a learned neural policy proposes an unverified action
\(u_{\text{nom}}\), the enforcer does not need to forecast the entire
future path of the machine. It evaluates this single half-space
inequality at the current state \(x\). If the proposed action satisfies
\(a(x)^T u_{\text{nom}} \le b(x)\), the enforcer passes the command to
the motor untouched. If the proposed action vectors the trajectory
across the boundary \(\partial \mathcal{C}\) into the unsafe domain, the
enforcer filters the command, deflecting the state trajectory
tangentially along the boundary.

The arithmetic of this per-cycle barrier check is direct. Consider an
analytical example of a cart with mass \(m = 10.0\text{ kg}\)
translating along a linear track toward a hard stop at position
\(p_{\text{max}} = 2.0\text{ m}\). The actuator can deliver a maximum
braking force of \(u_{\text{max}} = 40.0\text{ N}\), corresponding to a
maximum deceleration of
\(a_{\text{max}} = u_{\text{max}} / m = 4.0\text{ m/s}^2\). For a
second-order mechanical system (\(\dot{p} = v, \dot{v} = u/m\)), a pure
position constraint \(h(p) = p_{\text{max}} - p\) has relative degree
\(r = 2\) because control input \(u\) does not appear in
\(\dot{h}(p) = -v\).\sidenote{\footnotesize Zeroing CBFs require relative degree one
  (\(L_g h \neq 0\)). For second-order systems where constraints govern
  position, filters construct velocity-dependent kinetic stopping
  boundaries (\(h = p_{\max} - p - v^2/(2a_{\max})\)) or High-Order CBFs
  to ensure control inputs appear in first derivatives.} To recover
relative degree one (\(r = 1\)), the barrier function incorporates the
kinetic stopping distance:
\[h(p, v) = (p_{\text{max}} - p) - \frac{v^2}{2 a_{\text{max}}}.\]
Differentiating with respect to time along the dynamics yields:
\[\dot{h}(p, v, u) = -\dot{p} - \frac{v \dot{v}}{a_{\text{max}}} = -v - \frac{v (u / m)}{a_{\text{max}}} = -v \left(1 + \frac{u}{u_{\text{max}}}\right).\]
Suppose the cart is currently at position \(p = 1.50\text{ m}\)
traveling forward at velocity \(v = 1.50\text{ m/s}\). The required
stopping distance under maximum deceleration is
\((1.50)^2 / (2 \times 4.0) = 0.281\text{ m}\), leaving a barrier value
of \(h = (2.0 - 1.50) - 0.281 = 0.219\text{ m}\). Choosing a linear
contraction rate \(\gamma_0 = 2.0\text{ s}^{-1}\), the enforcer requires
\(\dot{h} \ge -2.0 \times 0.219 = -0.438\text{ m/s}\). Substituting the
state values yields the inequality \(-1.50 (1 + u / 40.0) \ge -0.438\),
which simplifies to \(1 + u / 40.0 \le 0.292\), or
\(u \le -28.3\text{ N}\). If the policy proposes an acceleration command
of \(u_{\text{nom}} = +10.0\text{ N}\), the enforcer overrides the
proposal and commands at least \(28.3\text{ N}\) of braking force.

A real machine does not execute continuous, ideal forces on an
uncorrupted plant. The nominal barrier assumes exact state feedback and
instantaneous torque realization, but the nervous system operates
subject to lower-level tracking errors and sensor uncertainty. If the
tracking loop exhibits a bounded velocity error
\(\epsilon_v = 0.10\text{ m/s}\) and the state estimator exhibits a
position uncertainty \(\epsilon_p = 0.02\text{ m}\), the enforcer cannot
evaluate the barrier at the geometric edge of the safe set. It must
evaluate an inset barrier
\(h_{\text{safe}}(x) = h(x) - \delta_{\text{margin}}\). In this
analytical example, the worst-case stopping distance under velocity
error expands by
\(\Delta d = (v + \epsilon_v)^2 / (2 a_{\text{max}}) - v^2 / (2 a_{\text{max}}) \approx v \epsilon_v / a_{\text{max}} = (1.50 \times 0.10) / 4.0 = 0.0375\text{ m}\).
Combining this stopping deficit with the position uncertainty yields a
required barrier inset of
\(\delta_{\text{margin}} = \epsilon_p + \Delta d = 0.02 + 0.0375 = 0.0575\text{ m}\).
Tightening the barrier by this margin reduces the effective barrier
value from \(0.219\text{ m}\) to \(0.1615\text{ m}\), shifting the
required intervention threshold earlier in the trajectory.

The safety warrant provided by the barrier calculation is strictly
conditional. Forward invariance guarantees that the state remains within
\(\mathcal{C}\) if and only if five underlying conditions hold
simultaneously:

\begin{enumerate}
\def\labelenumi{\arabic{enumi}.}
\tightlist
\item
  \textbf{Model Fidelity:} The nominal plant dynamics \(f(x)\) and
  \(g(x)\) must capture the true system dynamics within a certified
  parameter error bound.
\item
  \textbf{Disturbance Bounds:} External forces and unmodeled friction
  acting on the mass must remain strictly within the modeled disturbance
  envelope \(|d(t)| \le d_{\max}\).
\item
  \textbf{Actuation Authority:} Physical actuators must possess
  sufficient voltage, current, and thermal headroom to realize the
  commanded corrective force without saturation
  (\(u^* \in [u_{\min}, u_{\max}]\)).
\item
  \textbf{Solver Determinism:} The optimization solver or projection
  algorithm must terminate strictly within its allocated execution
  budget (\(T_{\text{solve}} \le T_{\text{budget}} < T_{\text{loop}}\)).
\item
  \textbf{Discretization Invariance:} The discrete sampling interval
  \(T_{\text{ctrl}}\) must be sufficiently small that inter-sample state
  drift does not breach the continuous-time invariance boundary.
\end{enumerate}

Each of these five conditions maps directly to the physical system
budgets established in Chapter~\ref{sec-body}, separating empirical
bench characterization from assumptions that the release argument must
explicitly carry. Model fidelity and disturbance bounds are measured on
the bench through system identification and load testing, yet the
release argument must assume that operational environments do not exceed
the tested envelope. Actuation authority is governed by the energy and
thermal budgets of the drive electronics, verified by dynamometer curves
under continuous load. Solver execution is governed by the latency
budget of the nervous system processor, verified through static
worst-case execution time analysis. Sampling interval validity is
governed by the loop frequency budget, ensuring that the discrete step
\(T_{\text{ctrl}} = 1.0\text{ ms}\) limits zero-order hold integration
error to a fraction of \(\delta_{\text{margin}}\).

Because a conditional guarantee whose conditions are unmonitored is an
ungrounded assertion, the nervous system must attach a dedicated
hardware or software detector to each condition. A disturbance observer
calculates the residual between predicted acceleration and measured
accelerometer data on every cycle, flagging unmodeled external forces
that exceed the disturbance budget. Inverter status registers and
current shunt monitors track actuator saturation, reporting when
commanded torque cannot be delivered. A hardware watchdog and
high-resolution cycle counter monitor solver execution latency against
the sub-millisecond deadline. Timestamp counters on the sensor bus
monitor communication jitter and missed sample deadlines.

When any runtime detector trips, the mathematical premise of the safe
set collapses. If an actuator overheats and drops its torque limit, or
if an unexpected payload disturbance exceeds the observer threshold, the
enforcer cannot continue claiming that tangential projection along
\(\partial \mathcal{C}\) guarantees safety. The correct system response
is not to recompute the barrier with optimistic parameters, but to
execute an immediate transition to a dedicated fallback mode via the
Simplex switch. The enforcer revokes the policy's authority, disengages
normal tracking, and commands an emergency stop sequence or mechanical
brake engagement. The safe set provides a valid permission path only
while its operating contract is intact, and the nervous system enforces
safety by knowing precisely when that contract has failed.

\section{Minimal Intervention}\label{sec-enforcement-12-8}

The principle of \textbf{minimal intervention} grants the learned policy
complete operational authority across the interior of the safe set,
modifying a command if and only if the proposed action violates a
barrier constraint. When an autonomous machine executes a task,
replacing the neural network's commands with a hand-crafted conservative
controller during nominal operation discards the behavioral competence
that justified deploying a learned model. Conversely, overriding a
proposed action by forcing an aggressive, full-magnitude braking command
at the first encounter with a boundary creates large acceleration
spikes, destabilizes dynamic balance, and interrupts task progress.
Minimal intervention treats safety enforcement as a \textbf{safety
shield} (projection operator). The nervous system accepts the policy's
proposal as the default intent and computes the smallest physical
correction that keeps the system within the invariant set.

A \textbf{minimal-intervention QP projection (safety shield)} is a
sub-millisecond convex optimization safety filter
(\(\min_{\mathbf{u}} \frac{1}{2}\|\mathbf{u} - \mathbf{u}_{\text{nom}}\|^2 \text{ s.t. } a(\mathbf{x})^T \mathbf{u} \le b(\mathbf{x}), \mathbf{u}_{\min} \le \mathbf{u} \le \mathbf{u}_{\max}\))
that passes unverified learned policy proposals unaltered in the
interior of the safe set and computes the least-squares orthogonal
modification on the boundary, preserving task intent without
uncoordinated scalar clipping.

Formally, this safety filtering mechanism is framed as a \textbf{Control
Barrier Function Quadratic Program (CBF-QP)}
(\citeproc{ref-ames2016control}{Ames et al. 2016}). At each discrete
control interval \(T_{\text{ctrl}} = 1.0\text{ ms}\) on the bare-metal
microcontroller, the enforcer receives a nominal torque command
\(u_{\text{nom}} \in \mathbb{R}^m\) from the policy mailbox and solves a
convex optimization problem.

\protect\phantomsection\label{callout-mathux2a-12.7}
\begin{fbxSimple}{callout-math}{Fermi Math:}{CBF-QP Formulation}
\phantomsection\label{callout-math*-12.7}

\textbf{Objective:} Compute the minimal intervention \(u^*\):
\[\begin{aligned}
u^* = \arg\min_{u \in \mathbb{R}^m} \quad & \frac{1}{2} \|u - u_{\text{nom}}\|^2 \\
\text{subject to} \quad & a_k(x)^T u \le b_k(x), \quad \forall k \in \{1, \dots, K\}, \\
& u_{\min} \le u \le u_{\max}.
\end{aligned}\] The objective minimizes the Euclidean norm of the
intervention, \(\frac{1}{2}\|u - u_{\text{nom}}\|^2\), subject to linear
inequality constraints derived from the \(K\) active barrier
certificates and physical actuator saturation limits. Because the
objective function is strictly convex and the constraints form a closed
convex polyhedron in control space, the optimization problem yields a
unique global minimum that can be solved within deterministic
sub-millisecond execution bounds.\sidenote{\footnotesize On bare-metal safety MCUs, QP
  safety filters use static memory arrays with hard iteration ceilings.
  While typical active-set solve times range from
  \(60\text{--}120\ \mu\text{s}\), worst-case execution time spikes to
  \(600\ \mu\text{s}\) under nearly degenerate constraints, requiring
  static WCET budgeting.} For a single active barrier certificate
\(a^T u \le b\), the quadratic program admits an exact closed-form
analytical solution that eliminates numerical solver iterations
entirely. When the nominal input satisfies \(a^T u_{\text{nom}} \le b\),
the constraint is inactive, the correction is zero, and
\(u^* = u_{\text{nom}}\). When the proposal violates the boundary such
that \(a^T u_{\text{nom}} > b\), the optimal command \(u^*\) lies
directly on the bounding hyperplane \(a^T u = b\). Constructing the
Lagrangian
\(\mathcal{L}(u, \lambda) = \frac{1}{2} \|u - u_{\text{nom}}\|^2 + \lambda (a^T u - b)\)
with non-negative multiplier \(\lambda \ge 0\), setting the gradient
\(\nabla_u \mathcal{L} = u - u_{\text{nom}} + \lambda a\) to zero yields
\(u = u_{\text{nom}} - \lambda a\). Substituting this expression into
the boundary equality \(a^T (u_{\text{nom}} - \lambda a) = b\) produces
\(a^T u_{\text{nom}} - \lambda \|a\|^2 = b\), which resolves to
\(\lambda = (a^T u_{\text{nom}} - b) / \|a\|^2\). Substituting
\(\lambda\) back into the stationarity equation establishes the
orthogonal projection:
\[u^* = u_{\text{nom}} - \frac{\max(0, a^T u_{\text{nom}} - b)}{\|a\|^2} a.\]

\end{fbxSimple}

Consider an analytical example for a two-axis planar mechanism where the
learned policy proposes a nominal joint torque command
\(u_{\text{nom}} = [4.0, 3.0]^T\text{ N}\cdot\text{m}\). The barrier
certificate enforces the half-space constraint \(a^T u \le b\) with
normal vector \(a = [1.0, 2.0]^T\) and boundary limit
\(b = 5.0\text{ N}\cdot\text{m}\). Evaluating the nominal proposal
yields
\(a^T u_{\text{nom}} = (1.0)(4.0) + (2.0)(3.0) = 10.0\text{ N}\cdot\text{m}\),
which exceeds the admissible threshold by \(5.0\text{ N}\cdot\text{m}\).
The squared norm of the constraint vector is
\(\|a\|^2 = 1.0^2 + 2.0^2 = 5.0\). The Lagrange multiplier evaluates to
\(\lambda = (10.0 - 5.0) / 5.0 = 1.0\), producing a correction vector
\(\lambda a = [1.0, 2.0]^T\text{ N}\cdot\text{m}\). Subtracting this
correction from the nominal input produces the filtered command
\(u^* = [4.0 - 1.0, 3.0 - 2.0]^T = [3.0, 1.0]^T\text{ N}\cdot\text{m}\).
The magnitude of the required physical intervention is
\(\|u^* - u_{\text{nom}}\| = \sqrt{(-1.0)^2 + (-2.0)^2} = \sqrt{5.0} \approx 2.24\text{ N}\cdot\text{m}\).
This orthogonal modification satisfies the boundary condition
\(a^T u^* = (1.0)(3.0) + (2.0)(1.0) = 5.0\text{ N}\cdot\text{m}\) while
preserving the maximum possible component of the policy's intended
effort.

As illustrated in Figure~\ref{fig-12-cbf-safety-filter}, this
minimal-intervention architecture operates simultaneously across state
space \(\mathcal{X}\) and control input space \(\mathcal{U}\). In state
space (panel a), unconstrained policy proposals vectoring across the
barrier into the unsafe region are deflected tangentially along the
forward-invariant safe set boundary \(\partial \mathcal{C}\) inset by
the tracking error margin \(\delta_{\text{margin}}\). In control space
(panel b), the enforcer orthogonally projects the unconstrained command
\(u_{\text{nom}}\) onto the linear half-space boundary \(a^T u \le b\),
preserving multi-axis dynamic coupling where heuristic scalar clipping
fails to guarantee safety.

\begin{figure}[t]

\centering{

\includegraphics[width=1\linewidth,height=\textheight,keepaspectratio]{index_files/mediabag//Users/VJ/GitHub/MLSysBook-integrate-physical/physical/book/chapters/12-enforcement/figures/fig12_cbf_safety_filter.pdf}

}

\caption{\label{fig-12-cbf-safety-filter}\textbf{Control Barrier
Function Safe Set Invariance and Minimal-Intervention Quadratic Program
Projection.} (a) State Space Invariance (\(\mathcal{X}\)): The safe set
\(\mathcal{C} = \{x \in \mathcal{X} \mid h(x) \ge 0\}\) is rendered
forward-invariant by evaluating the Lie derivative condition
\(\dot{h}(x, u) = L_f h(x) + L_g h(x) u \ge -\gamma h(x)\). To guarantee
physical safety under lower-level tracking error
\(\epsilon_{\text{track}}\) and sensing latency, the enforcer insets the
boundary by
\(\delta_{\text{margin}} = \epsilon_p + (v \epsilon_v)/a_{\max}\). When
an unconstrained learned policy proposal vectors toward the forbidden
set (red dashed vector \(f(x, u_{\text{nom}})\)), the enforcer filters
the command, deflecting the state trajectory tangentially along the
inset boundary (teal vector \(f(x, u^*)\)) in compliance with Nagumo's
viability theorem (\(f \in T_{\mathcal{C}}(x)\)). (b) Control Input
Space Minimal Intervention (\(\mathcal{U}\)): In the 2D torque space
\((u_1, u_2)\), the active barrier certificate establishes a linear
half-space boundary \(a^T u \le b\). An unsafe policy proposal
\(u_{\text{nom}} = [4.0, 3.0]^T\text{ N}\cdot\text{m}\) is orthogonally
projected via a sub-millisecond quadratic program onto the bounding
hyperplane at \(u^* = [3.0, 1.0]^T\text{ N}\cdot\text{m}\) with minimal
intervention correction \(\Delta u = -\lambda a\). In contrast,
heuristic per-axis scalar clipping
\(u_{\text{clip}} = [3.5, 2.0]^T\text{ N}\cdot\text{m}\) distorts the
net Cartesian wrench angle and remains strictly outside the admissible
control half-space (\(a^T u_{\text{clip}} = 7.5 > 5.0\)).}

\end{figure}%

The geometric coordination of this projection distinguishes principled
enforcement from heuristic clipping. A common implementation flaw in
embedded motor controllers is independent per-axis saturation, where
each joint command is clamped to an isolated scalar limit without
evaluating coupling across degrees of freedom. In a multi-link
articulated body, joint accelerations are coupled through the
off-diagonal terms of the mass matrix. Clamping joint torques
independently rotates the net Cartesian wrench vector away from its
intended line of action. When an uncoordinated saturation clips torque
on one axis while leaving adjacent axes unrestricted, the resulting
wrench generates unmodeled twisting moments across structural linkages,
introduces severe lateral vibrations, and increases mechanical fatigue
on gear teeth and output bearings. Constrained quadratic projection
preserves dynamic coupling by computing the modification across all
actuation axes simultaneously.

The selection of an enforcement filtering paradigm establishes the
boundary between microsecond real-time determinism, mathematical safety
guarantees, and the preservation of learned task competence. As
summarized in Table~\ref{tbl-12-safety-filter-architectures},
architectural choices vary dramatically in their handling of relative
degree, solver worst-case execution time (WCET), forward invariance
guarantees, and directional fidelity. While naive per-axis clamping
introduces uncoordinated net wrenches and predictive filters exceed
microsecond motor loop deadlines, Control Barrier Function quadratic
programs (ZCBF-QP and HOCBF-QP) provide the mathematical structure
required for minimal intervention on embedded real-time silicon.

\begin{longtable}[]{@{}
  >{\raggedright\arraybackslash}p{(\linewidth - 12\tabcolsep) * \real{0.1429}}
  >{\raggedright\arraybackslash}p{(\linewidth - 12\tabcolsep) * \real{0.1429}}
  >{\raggedright\arraybackslash}p{(\linewidth - 12\tabcolsep) * \real{0.1429}}
  >{\raggedright\arraybackslash}p{(\linewidth - 12\tabcolsep) * \real{0.1429}}
  >{\raggedright\arraybackslash}p{(\linewidth - 12\tabcolsep) * \real{0.1429}}
  >{\raggedright\arraybackslash}p{(\linewidth - 12\tabcolsep) * \real{0.1429}}
  >{\raggedright\arraybackslash}p{(\linewidth - 12\tabcolsep) * \real{0.1429}}@{}}
\caption{\textbf{Comparative Systems Architecture of Real-Time Safety
Enforcement and Filter Paradigms.} Comparison of mathematical
guarantees, relative degree handling, algorithmic optimization
structure, worst-case execution time (WCET) bounds on embedded
bare-metal silicon, intervention directional fidelity, and primary
failure modes across five candidate safety filtering architectures.
Zeroing CBFs (ZCBF) and High-Order CBFs (HOCBF) provide the optimal
trade-off between microsecond determinism and minimal intervention,
whereas uncoordinated clamping compromises multi-axis dynamics and
predictive filters exceed bare-metal control loop
deadlines.}\label{tbl-12-safety-filter-architectures}\tabularnewline
\toprule\noalign{}
\begin{minipage}[b]{\linewidth}\raggedright
Safety Filter Paradigm
\end{minipage} & \begin{minipage}[b]{\linewidth}\raggedright
Relative Degree (\(r\))
\end{minipage} & \begin{minipage}[b]{\linewidth}\raggedright
Mathematical Guarantee
\end{minipage} & \begin{minipage}[b]{\linewidth}\raggedright
Optimization Formulation
\end{minipage} & \begin{minipage}[b]{\linewidth}\raggedright
Embedded Latency (WCET)
\end{minipage} & \begin{minipage}[b]{\linewidth}\raggedright
Intervention Fidelity
\end{minipage} & \begin{minipage}[b]{\linewidth}\raggedright
Primary Failure Mode
\end{minipage} \\
\midrule\noalign{}
\endfirsthead
\toprule\noalign{}
\begin{minipage}[b]{\linewidth}\raggedright
Safety Filter Paradigm
\end{minipage} & \begin{minipage}[b]{\linewidth}\raggedright
Relative Degree (\(r\))
\end{minipage} & \begin{minipage}[b]{\linewidth}\raggedright
Mathematical Guarantee
\end{minipage} & \begin{minipage}[b]{\linewidth}\raggedright
Optimization Formulation
\end{minipage} & \begin{minipage}[b]{\linewidth}\raggedright
Embedded Latency (WCET)
\end{minipage} & \begin{minipage}[b]{\linewidth}\raggedright
Intervention Fidelity
\end{minipage} & \begin{minipage}[b]{\linewidth}\raggedright
Primary Failure Mode
\end{minipage} \\
\midrule\noalign{}
\endhead
\bottomrule\noalign{}
\endlastfoot
\textbf{Per-Axis Heuristic Clamping} & \(r = 0\) / \(1\) (Algebraic) &
None (Violates invariant set geometry) &
\(u_i^* = \text{clip}(u_{\text{nom}, i}, u_{\min, i}, u_{\max, i})\) &
\(< 1\,\mu\text{s}\) (Vectorized ALU) & Uncoordinated (Distorts net
wrench direction) & High-frequency chattering; unmodeled joint stress \\
\textbf{Zeroing Control Barrier (ZCBF-QP)} & \(r = 1\)
(\(L_g h \neq \mathbf{0}\)) & Forward invariance \(\mathcal{C}\)
(\(\dot{h} \ge -\gamma h\)) &
\(\min_{\mathbf{u}} \frac{1}{2}\|\mathbf{u} - \mathbf{u}_{\text{nom}}\|^2 \text{ s.t. } \mathbf{A}_{\text{cbf}}\mathbf{u} \le \mathbf{b}_{\text{cbf}}\)
& \(15\text{--}120\,\mu\text{s}\) (\(P_{99.9} \le 175\,\mu\text{s}\)) &
Optimal (Orthogonal half-space projection) & Infeasible when actuators
saturate (\(\mathcal{U}_{\text{safe}} = \emptyset\)) \\
\textbf{High-Order CBF (HOCBF / Kinetic Inset)} & \(r \ge 2\) (Position
limits) & Higher-order invariant set \(\mathcal{C}_r\) via Lie
derivatives &
\(\min_{\mathbf{u}} \frac{1}{2}\|\mathbf{u} - \mathbf{u}_{\text{nom}}\|^2 \text{ s.t. } L_f^r h + L_g L_f^{r-1} h\,\mathbf{u} \ge -\boldsymbol{\alpha}(\boldsymbol{\psi})\)
& \(45\text{--}250\,\mu\text{s}\) (\(P_{99.9} \le 350\,\mu\text{s}\)) &
Near-Optimal (Dissipation-aligned projection) & Highly sensitive to
plant inertia / mass matrix error \\
\textbf{Explicit Polyhedral Reachability} & \(r \in \mathbb{N}\) (Linear
systems) & Robust positively invariant (RPI) set containment &
Multi-parametric QP lookup:
\(\mathbf{H}_i \mathbf{x} + \mathbf{G}_i \mathbf{u} \le \mathbf{f}_i\) &
\(5\text{--}30\,\mu\text{s}\) (Polytopic region search) & Exact on
polyhedral boundaries & Polytopic complexity explosion (\(>10^4\)
regions for \(n>6\)) \\
\textbf{Predictive Safety Filter (MPSF)} & Arbitrary (\(N\)-step
horizon) & Recursive feasibility with terminal invariant set
\(\mathcal{X}_f\) &
\(\min_{\mathbf{U}} \sum_{k=0}^{N-1}\|\mathbf{u}_k - \mathbf{u}_{\text{nom}, k}\|^2 + V_f(\mathbf{x}_N)\)
& \(2.0\text{--}15.0\text{ ms}\) (Real-time SQP solver) & Global
foresight (Avoids dynamic traps) & Deadline overruns on microsecond
motor control loops \\
\end{longtable}

Operating directly on the boundary of the safe set introduces a distinct
physical failure mode known as boundary chattering. If a learned policy
continuously drives the system against an active barrier, discrete-time
enforcement creates a high-frequency limit cycle. At
\(T_{\text{ctrl}} = 1.0\text{ ms}\), the enforcer applies an orthogonal
correction on one time step to satisfy the barrier, moving the state
slightly into the safe set interior. On the subsequent time step, the
policy again commands an unsafe action, triggering another correction.
This alternating sequence switches the active constraint at frequencies
approaching the \(500\text{ Hz}\) Nyquist limit of the
\(1000\text{ Hz}\) control loop. Rapid torque reversals excite unmodeled
structural resonances in the mechanical transmission, cause rapid
thermal dissipation in the drive inverters, and degrade surface contact
in harmonic drives. To eliminate boundary chattering, the nervous system
incorporates boundary layer smoothing, scaling the projection through a
thin margin parameter \(\epsilon\) that continuously damps the command
as the state approaches the boundary rather than applying a
discontinuous step correction.

Beyond enforcing safety on each instantaneous cycle, the intervention
vector \(\Delta u = u^* - u_{\text{nom}}\) provides a continuous
diagnostic measurement of upstream policy health. During nominal
operation within the training distribution, a well-trained policy
operates inside the safe set, producing zero intervention norm across
most cycles. When sensory degradation, model drift, or novel
environmental dynamics cause the policy to produce unsafe proposals, the
frequency and magnitude of the enforcer's corrections increase. By
tracking the cumulative intervention metric over a rolling window of
\(W = 500\text{ ms}\), calculated as the discrete sum
\(\sum_{k=0}^{N-1} \|u_k^* - u_{\text{nom}, k}\| \Delta t\), the nervous
system detects upstream failure before the machine reaches an
irrecoverable state. When the \(P_{95}\) intervention metric over the
window exceeds a predetermined threshold, the system flags policy
divergence, allowing the host architecture to trigger retraining data
logging or operator notification.

The theoretical guarantee of this minimal intervention projection is an
imported result that holds under three specific engineering assumptions.
First, the plant dynamics must remain affine in the control input over
the control period. Second, the underlying current tracking loop must
achieve sufficient closed-loop bandwidth (\(10\text{ kHz}\) current
regulation beneath the \(1000\text{ Hz}\) torque command loop) so that
commanded torque matches delivered motor torque. Third, the intersection
of the barrier half-spaces and physical actuator saturation limits must
remain non-empty (\(\mathcal{U}_{\text{safe}} \neq \emptyset\)).

\protect\phantomsection\label{callout-mathux2a-12.8}
\begin{fbxSimple}{callout-math}{Fermi Math:}{Bare-Metal CPU Cycle Budget for Sub-Millisecond CBF-QP Solving}
\phantomsection\label{callout-math*-12.8}

\textbf{Question:} \emph{Can a bare-metal \(480\text{ MHz}\)
microcontroller execute an active-set Quadratic Program safety filter
inside a \(1000\text{ Hz}\) (\(1000\,\mu\text{s}\)) control loop
alongside Field-Oriented Control (FOC) current regulation without
missing real-time timer deadlines?}

\textbf{Hardware Substrate:} A canonical hard real-time safety
microcontroller operating at \(f_{\text{CPU}} = 480\text{ MHz}\) with
hardware double-precision Floating-Point Unit (FPU), \(64\text{ KB}\)
Data Tightly-Coupled Memory (DTCM) running at single-cycle access (\(0\)
wait states).

\textbf{1. Available Real-Time Clock Cycle Budget:} For a control loop
period \(T_{\text{loop}} = 1.0\text{ ms} = 1000\,\mu\text{s}\):
\[N_{\text{total\_cycles}} = 480\text{ MHz} \times 1.0\text{ ms} = 480{,}000\text{ clock cycles}.\]

\textbf{2. Sub-Task Execution Cycle Accounting (Worst-Case Analysis):}

\begin{itemize}
\item
  \textbf{State Acquisition \& Sensor Ingestion (\(T_1\)):} Read optical
  encoder SPI registers and 3-phase ADC current shunts via direct DTCM
  buffers:
  \[N_1 \approx 12{,}000\text{ cycles} \quad (25.0\,\mu\text{s}).\]
\item
  \textbf{Lease Monotonicity \& Tracking Contract Audit (\(T_2\)):}
  Verify CRC32, timestamp freshness
  (\(\tau_{\text{lease}} \le 40\text{ ms}\)), and position error
  \(\|x - r\| \le \epsilon_{\text{track}}\):
  \[N_2 \approx 4{,}800\text{ cycles} \quad (10.0\,\mu\text{s}).\]
\item
  \textbf{Lie Derivative \& Constraint Half-Space Synthesis (\(T_3\)):}
  Evaluate \(L_f h_k(x)\) and \(L_g h_k(x)\) for \(K = 6\) barrier
  certificates (joint position limits, velocity ceilings, and kinematic
  stopping boundaries):
  \[N_3 \approx 28{,}800\text{ cycles} \quad (60.0\,\mu\text{s}).\]
\item
  \textbf{Active-Set QP Projection Solve (\(T_4\)):} Primal-Dual
  active-set solver in pre-allocated static SRAM with \(m = 6\) control
  inputs, \(K = 6\) constraints, taking at most \(8\) active-set
  iterations (Cholesky factorizations \(\mathcal{O}(m^3)\) and
  back-solves \(\mathcal{O}(m^2)\)):
  \[N_4 \approx 57{,}600\text{ cycles} \quad (120.0\,\mu\text{s}).\]
\item
  \textbf{Simplex Boundary Arbitration \& PWM Compare Register Reload
  (\(T_5\)):} Atomic write of \(u^*\) to timer compare registers
  (\texttt{TIMx\_CCR1}--\texttt{CCR6}):
  \[N_5 \approx 4{,}800\text{ cycles} \quad (10.0\,\mu\text{s}).\]
\end{itemize}

\textbf{3. Total Worst-Case Execution Time (WCET):}
\[T_{\text{WCET}} = 25.0 + 10.0 + 60.0 + 120.0 + 10.0 = 225.0\,\mu\text{s} \quad (108{,}000\text{ cycles}).\]
\[\text{CPU Utilization}_{\text{enforcer}} = \frac{108{,}000}{480{,}000} = 22.5\%.\]

\textbf{Verdict:} The entire CBF-QP safety shield consumes only
\(225\,\mu\text{s}\) (\(22.5\%\)) of the \(1\text{ ms}\) tick interval,
leaving \(775\,\mu\text{s}\) (\(372{,}000\text{ cycles}\), or \(77.5\%\)
CPU headroom) for background \(20\text{ kHz}\) current loop regulation,
hardware watchdog heartbeat servicing, and telemetry serialization.

\end{fbxSimple}

On a canonical heterogeneous dual-brain System-on-Chip (SoC), this
safety shield executes inside the bare-metal \(1000\text{ Hz}\)
(\(1000\,\mu\text{s}\)) timer interrupt service routine on the real-time
microcontroller core. The execution follows a strict deterministic
sequence:

\begin{enumerate}
\def\labelenumi{\arabic{enumi}.}
\tightlist
\item
  \textbf{Interrupt Entry \& State Acquisition
  (\(0\text{--}15\,\mu\text{s}\)):} Hardware timer triggers ISR;
  firmware reads raw optical encoder registers and ADC current shunts.
\item
  \textbf{Proposal \& Lease Verification
  (\(15\text{--}25\,\mu\text{s}\)):} The MCU inspects the shared SRAM
  ring buffer populated by the Linux MPU over RPMSG or shared memory. It
  validates the CRC32 checksum, confirms monotonic sequence order, and
  verifies that the intent lease has not expired
  (\(\tau_{\text{lease}} \le 40\text{ ms}\)). It resets the hardware
  countdown watchdog (\(\tau_{\text{watchdog}} \le 10\text{ ms}\)).
\item
  \textbf{Tracking Contract Verification
  (\(25\text{--}35\,\mu\text{s}\)):} Firmware evaluates instantaneous
  tracking error \(e(t) = \|x(t) - r(t)\|\) against the certified bound
  \(\epsilon_{\text{track}}\).
\item
  \textbf{CBF-QP Safety Shield Solve (\(35\text{--}150\,\mu\text{s}\)):}
  An active-set QP solver running in pre-allocated static DTCM SRAM
  evaluates active barrier inequalities and computes
  \(u^* = \arg\min_u \frac{1}{2}\|u - u_{\text{nom}}\|^2\). Dynamic
  memory allocation (\texttt{malloc}) is strictly forbidden to prevent
  non-deterministic heap fragmentation.
\item
  \textbf{Simplex Switch Arbitration
  (\(150\text{--}165\,\mu\text{s}\)):} If the QP is feasible and
  constraints hold, \(u^*\) is written directly to the PWM timer compare
  registers. If the intent lease has expired, tracking error exceeds
  \(\epsilon_{\text{track}}\), or actuator saturation renders the
  feasible control set empty
  (\(\mathcal{U}_{\text{safe}} = \emptyset\)), the Simplex switch trips,
  immediately transferring authority from the unverified policy to the
  deterministic Fallback Ladder.
\end{enumerate}

The deterministic execution loop that unifies barrier evaluation,
active-set quadratic optimization, and Simplex fallback arbitration on
bare-metal microcontroller silicon is formalized below:

\subsection{Bare-Metal 1 kHz Control Barrier Function Safety Shield
Protocol}\label{bare-metal-1-khz-control-barrier-function-safety-shield-protocol}

\textbf{Cadence \& Invariants:}

\begin{itemize}
\tightlist
\item
  \emph{Execution Cadence:} Hard real-time periodic timer interrupt
  (\(f = 1000\text{ Hz}\), period \(T_{\text{loop}} = 1.0\text{ ms}\),
  maximum allowable hardware timing jitter \(< 1.0\,\mu\text{s}\)).
\item
  \emph{Memory Invariant:} Zero dynamic heap allocation
  (\texttt{malloc}/\texttt{free} strictly forbidden; all state matrices
  and solver workspaces statically allocated in DTCM SRAM).
\item
  \emph{Hardware Authority:} Exclusive direct write access to motor PWM
  timer compare registers (\texttt{TIMx\_CCR} / \texttt{ePWM\_CMPA}) and
  hardware brake trip GPIOs.
\end{itemize}

\textbf{Execution Protocol:}

\begin{enumerate}
\def\labelenumi{\arabic{enumi}.}
\item
  \textbf{State Acquisition \& Innovation Gating (Stage 1):} Sample raw
  encoder registers \(\mathbf{q}_{\text{enc}}\) and phase currents
  \(\mathbf{i}_{\text{phase}}\) via SPI/ADC DMA. Compute reconstructed
  physical state
  \(\hat{\mathbf{x}}_t = (\hat{\mathbf{q}}_t, \dot{\hat{\mathbf{q}}}_t)\).
  Evaluate measurement innovation
  \(\mathbf{r}_t \leftarrow \|\mathbf{y}_t - \mathbf{C}\hat{\mathbf{x}}_t\|\).
  If \(\mathbf{r}_t > \gamma_{\text{sensor}}\), assert
  \texttt{ERR\_SENSOR\_CORRUPTION} and transition to Stage 5 (Fallback
  Escalation).
\item
  \textbf{Intent Lease \& Tracking Contract Audit (Stage 2):} Verify
  CRC32 checksum of \(\mathcal{C}_{\text{prop}}\) in shared SRAM.
  Compute lease age
  \(\Delta t_{\text{lease}} \leftarrow t_{\text{now}} - \mathcal{C}_{\text{prop}}.t_{\text{stamp}}\).
  If \(\Delta t_{\text{lease}} > \tau_{\text{lease}}\) or CRC mismatch,
  assert \texttt{ERR\_LEASE\_EXPIRED} and jump to Stage 5. Compute
  instantaneous tracking error
  \(e(t) \leftarrow \|\hat{\mathbf{q}}_t - \mathcal{C}_{\text{prop}}.\mathbf{r}_t\|\);
  if \(e(t) > \epsilon_{\text{crit}}\), assert
  \texttt{ERR\_TRACKING\_BREACH} and jump to Stage 5.
\item
  \textbf{Lie Derivative \& Inset Barrier Evaluation (Stage 3):} For
  each active safety barrier \(k \in \{1, \dots, K\}\):

  \begin{itemize}
  \item
    Compute Lie derivatives along control-affine plant dynamics:
    \[L_f h_k(\hat{\mathbf{x}}_t) = \nabla h_k(\hat{\mathbf{x}}_t)^T f(\hat{\mathbf{x}}_t), \quad L_g h_k(\hat{\mathbf{x}}_t) = \nabla h_k(\hat{\mathbf{x}}_t)^T g(\hat{\mathbf{x}}_t).\]
  \item
    Construct half-space constraint row:
    \[\mathbf{a}_k^T \leftarrow -L_g h_k(\hat{\mathbf{x}}_t), \quad b_k \leftarrow L_f h_k(\hat{\mathbf{x}}_t) + \gamma_k (h_k(\hat{\mathbf{x}}_t) - \delta_k).\]
    If \(\mathbf{a}_k^T \mathbf{u}_{\text{nom}} \le b_k, \forall k\) and
    \(\mathbf{u}_{\min} \le \mathbf{u}_{\text{nom}} \le \mathbf{u}_{\max}\):
    set \(\mathbf{u}_t^* \leftarrow \mathbf{u}_{\text{nom}}\),
    \(\mathrm{Regime} \leftarrow \text{NOMINAL\_PASS}\), write
    \(\mathbf{u}_t^*\) to \texttt{TIMx\_CCR} registers, kick hardware
    watchdog timer, and return.
  \end{itemize}
\item
  \textbf{Active-Set Minimal-Intervention Quadratic Program Solve (Stage
  4):} Initialize active constraint set
  \(\mathcal{W} \leftarrow \{k \mid \mathbf{a}_k^T \mathbf{u}_{\text{nom}} > b_k\} \cup \{i \mid u_{\text{nom}, i} \notin [u_{\min, i}, u_{\max, i}]\}\).
  Solve convex quadratic program in static DTCM workspace:
  \[\mathbf{u}^* = \arg\min_{\mathbf{u} \in \mathbb{R}^m} \frac{1}{2}\|\mathbf{u} - \mathbf{u}_{\text{nom}}\|^2 \quad \text{s.t.} \quad \mathbf{a}_k^T \mathbf{u} \le b_k, \; \forall k \in \{1, \dots, K\}, \quad \mathbf{u}_{\min} \le \mathbf{u} \le \mathbf{u}_{\max}.\]
  If solver converges (\(N_{\text{iter}} \le 15\)): set
  \(\mathbf{u}_t^* \leftarrow \mathbf{u}^*\),
  \(\mathrm{Regime} \leftarrow \text{QP\_PROJECTED}\), write
  \(\mathbf{u}_t^*\) to \texttt{TIMx\_CCR} registers, kick watchdog, and
  return. Else (timeout or infeasibility
  \(\mathcal{U}_{\text{safe}} = \emptyset\)), assert
  \texttt{ERR\_QP\_INFEASIBLE} and escalate to Stage 5.
\item
  \textbf{Simplex Deterministic Fallback Escalation (Stage 5):} Trip
  Simplex Switch to revoke policy execution authority:

  \begin{itemize}
  \tightlist
  \item
    Moderate fault (\(e(t) \le \epsilon_{\text{crit}}\) and lease
    expired): set \(\mathrm{Regime} \leftarrow \text{FALLBACK\_HOLD}\),
    compute local position hold torque
    \(\mathbf{u}_{\text{hold}} = \mathbf{K}_p (\mathbf{q}_{\text{hold}} - \hat{\mathbf{q}}_t) - \mathbf{K}_d \dot{\hat{\mathbf{q}}}_t + \mathbf{g}(\hat{\mathbf{q}}_t)\),
    and write to \texttt{TIMx\_CCR}.
  \item
    Severe fault (\(e(t) > \epsilon_{\text{crit}}\) or \(h_k \le 0\)):
    set \(\mathrm{Regime} \leftarrow \text{FALLBACK\_DYNAMIC\_STOP}\),
    command dynamic braking deceleration
    \(a_{\text{brake}} = a_{\max}\).
  \item
    Emergency fault (watchdog timeout or over-current): set
    \(\mathrm{Regime} \leftarrow \text{ESTOP\_TRIPPED}\), assert
    hardware break pin (\texttt{TIMx\_BKIN}), engaging Safe Torque Off
    (STO) and spring friction brakes in \(< 5\,\mu\text{s}\).
  \end{itemize}
\end{enumerate}

\section{The Fallback Ladder}\label{sec-enforcement-12-9}

When nothing admissible remains within the barrier constraints, the
response is graded, and each rung costs more than the one before. A
physical machine cannot treat safety as a single binary switch. If an
enforcer cuts bus power whenever an upstream proposal brushes against a
barrier margin, the system suffers uncoordinated deceleration, thermal
cycling, and continuous mission aborts. If it instead relies on software
optimization when an inverter bridge experiences shoot-through current,
the gate drivers burn before an operating system can schedule an
interrupt handler. The nervous system resolves this tension through an
escalation ladder of deterministic defensive actions, moving from
microsecond mathematical projection down to galvanic power isolation.

The \textbf{four-tier fallback ladder} is a graded, deterministic
escalation hierarchy of defensive actions (Level 1: Minimal-Intervention
QP Projection; Level 2: Active Position Hold; Level 3: Controlled
Dynamic Stop; Level 4: Galvanic Safe Torque Off) that transitions across
independent execution substrates as control authority, solver
feasibility, or hardware invariants degrade.

The first rung is least-squares projection, executed on every control
cycle while the feasible control set remains non-empty
(\(\mathcal{U}_{\text{safe}} \neq \emptyset\)). When a nominal proposal
\(u_{\text{nom}}\) points outside the admissible half-spaces, the
enforcer computes the minimal-norm modification
\(u^* = \arg\min_u \|u - u_{\text{nom}}\|^2\) subject to the active
barrier inequalities. In an analytical \(1000\text{ Hz}\) control loop
with a worst-case execution budget of \(175\,\mu\text{s}\), this
projection runs entirely in pre-allocated static SRAM. It guarantees
forward invariance while maintaining active mission progress, charging
zero mechanical wear and zero mission downtime. The projection succeeds
only while the intersection of the barrier constraints and the actuator
torque envelope contains at least one solution. When physical
saturation, conflicting obstacle boundaries, or numerical divergence
render the admissible control set empty
(\(\mathcal{U}_{\text{safe}} = \emptyset\)), the mathematical premise of
the projection fails, tripping the Simplex switch to escalate to the
second rung.

The second rung is an active position hold, corresponding to an IEC
60204-1 Category 2 stop (Safe Stop 2).\sidenote{\footnotesize IEC 60204-1 defines
  functional stops: Category 0 (Safe Torque Off, STO; immediate
  unpowered mechanical stop), Category 1 (Safe Stop 1, SS1; controlled
  deceleration then power removal), and Category 2 (Safe Stop 2, SS2;
  powered standstill). ISO 13849-1 PLe mandates dual-channel hardware
  safety paths.} Rather than attempting to navigate along the planned
trajectory, the nervous system commands the actuators to arrest motion
and maintain the current coordinate \(x_{\text{hold}}\) under
closed-loop feedback. Inverter gate drivers remain actively modulated by
pulse-width modulation, driving stator currents to generate the reaction
torques needed to resist gravitational load and payload inertia. This
rung costs electrical energy and generates continuous Joule heating
(\(I^2 R\)) across the motor windings, but it incurs zero friction wear
on mechanical brake pads and leaves the kinematic state fully preserved.
If an upstream planner recovers within its lease timeout and provides a
valid trajectory, the enforcer resumes nominal execution with zero
homing overhead. If the intent lease expires without renewal, or if
tracking error exceeds the closed-loop holding tolerance, the system
abandons station-keeping and drops to the third rung.

The third rung is a controlled dynamic stop, classified under IEC
60204-1 as Category 1 (Safe Stop 1). The nervous system overrides all
task goals and commands maximum certified counter-electromotive braking
along a dedicated deceleration profile
(\(a_{\text{brake}} = a_{\max}\)). For an arm operating at a linear
velocity of \(2.0\text{ m/s}\) with a certified braking acceleration of
\(a_{\max} = 4.0\text{ m/s}^2\) and an actuation response delay of
\(12\text{ ms}\), stopping requires a deterministic clearance budget of
\(d_{\text{stop}} = v \tau + v^2 / (2 a_{\max}) = 0.524\text{ m}\). As
counter-torque slows the moving mass, mechanical kinetic energy
(\(E_k = \frac{1}{2} m v^2\)) regenerates through the inverter bridge
into the DC bus capacitors.\sidenote{\footnotesize During high-deceleration braking
  (Category 1 SS1), kinetic energy charges DC bus capacitors
  (\(\Delta E = \frac{1}{2} m v^2\)). Hardware comparators activate
  dynamic braking chopper MOSFETs to dissipate regenerative current
  through ceramic resistors, preventing bus overvoltage faults.} The
power stage must absorb this electrical surge through sized bus
capacitance and dynamic brake chopper resistors to prevent DC-bus
overvoltage faults. A Category 1 stop imposes peak torque loads across
gear teeth, aborts the active task, and requires a full state
re-initialization sequence before motion can resume, but it brings the
machine to a controlled standstill before isolating inverter power.

The fourth and final rung is a hard power cutoff, classified as an IEC
60204-1 Category 0 stop, or Safe Torque Off. This tier bypasses software
execution entirely. As illustrated in
Figure~\ref{fig-12-real-estop-interlock}, when a physical emergency stop
pushbutton is depressed, an over-current comparator trips, or a hardware
watchdog timer exceeds its \(50\text{ ms}\) timeout, hardware
optocouplers and safety relays de-energize the gate-driver bias supply
in under \(5\,\mu\text{s}\). The inverter MOSFET gates drop below their
conduction threshold, forcing the bridge into high-impedance cutoff and
eliminating electromagnetic torque regardless of processor register
state. Electrically released, spring-actuated friction brakes clamp the
motor shafts to arrest mechanical rotation. This rung provides galvanic
and hardware-level isolation against silicon failure, but it exacts the
highest physical penalty. Clamping friction brakes at high rotational
speeds causes rapid lining ablation, transmits undamped shock loads
directly through harmonic drive flexsplines, and can induce
gravitational payload drop during the millisecond mechanical dead-time
before brake pads seat.

\begin{figure}[t]

\centering{

\includegraphics[width=0.85\linewidth,height=\textheight,keepaspectratio]{/Users/VJ/GitHub/MLSysBook-integrate-physical/physical/book/chapters/12-enforcement/figures/fig12_real_estop_interlock.jpg}

}

\caption{\label{fig-12-real-estop-interlock}\textbf{Industrial Emergency
Stop Pushbutton and Hardware Cutoff Gate.} A positive-break,
mechanically latching mushroom emergency stop (E-stop) switch mounted on
an industrial machine tool console. The prominent red actuator head and
yellow high-visibility safety bezel govern dual-channel,
positive-opening normally closed (NC) contacts wired directly into
certified safety relays conforming to IEC 60204-1 (Category 0) and ISO
13849-1 (Category 4 / Performance Level e). Depressing the button
physically breaks the electrical coil circuit to main safety contactors
and interrupts the gate-driver power rail in under \(5\,\mu\text{s}\),
executing immediate Safe Torque Off (STO) and engaging mechanical
spring-applied friction brakes without depending on microcontrollers,
operating system schedulers, or software decision logic.}

\end{figure}%

The justification for structuring enforcement as a ladder lies in the
fundamental trade-off between engagement latency, structural stress,
component wear, and mission reversibility. Rung 1 resolves deviations in
microseconds with zero hardware degradation and complete reversibility.
Rung 2 arrests motion in tens of milliseconds, consuming electrical
power and thermal headroom while preserving the mission context. Rung 3
commits the machine to a complete stop over hundreds of milliseconds,
dissipating regenerative energy and discarding the trajectory plan to
guarantee collision avoidance. Rung 4 acts in microseconds at the
silicon layer, sacrificing brake linings, gear life, and operational
continuity to prevent electrical fires or unmitigated kinetic strikes.

The deterministic escalation hierarchy across these four defensive tiers
is structured in Figure~\ref{fig-12-fallback-ladder}. As execution
descends from software-level quadratic optimization to galvanic analog
contactors, reaction latency drops and physical isolation increases,
while mechanical wear costs and mission recovery penalties escalate
irreversibly.

\begin{figure}[t]

\centering{

\includegraphics[width=1\linewidth,height=\textheight,keepaspectratio]{index_files/mediabag//Users/VJ/GitHub/MLSysBook-integrate-physical/physical/book/chapters/12-enforcement/figures/fig12_fallback_ladder.pdf}

}

\caption{\label{fig-12-fallback-ladder}\textbf{The Deterministic
Four-Tier Fallback Escalation Ladder.} The physical AI nervous system
structures defensive actions into a graded hierarchy of four
deterministic rungs across distinct execution substrates. Rung 1
executes sub-millisecond quadratic programming (QP) projections in
static microcontroller SRAM, maintaining active mission progress with
zero mechanical wear. Rung 2 executes an active closed-loop position
hold (IEC 60204-1 Category 2 / SS2) on a real-time RTOS core,
dissipating electrical Joule heating while preserving trajectory
context. Rung 3 executes a controlled dynamic stop (IEC 60204-1 Category
1 / SS1) via dedicated safety MCU firmware, absorbing regenerative
DC-bus energy into brake choppers at the cost of task abort and required
re-homing. Rung 4 executes a hardware-level Safe Torque Off (IEC 60204-1
Category 0 / STO / ISO 13849-1 PLe) via optical isolators and
spring-actuated mechanical friction brakes in under \(5\,\mu\text{s}\),
providing galvanic isolation against silicon and power rail failures at
the expense of mechanical lining ablation and mandatory manual reset.}

\end{figure}%

To bridge abstract control tiers with statutory machinery compliance,
Table~\ref{tbl-12-iec60204-stopping-categories} maps the four rungs of
the fallback ladder to the standardized emergency stopping categories of
IEC 60204-1, ISO 13849-1, and ISO 26262. The table articulates the
physical trade-offs governing electrical power status, deceleration
mechanisms, engagement latencies, mechanical wear penalties, and
post-fault mission recovery workflows across each stopping tier.

\begin{longtable}[]{@{}
  >{\raggedright\arraybackslash}p{(\linewidth - 12\tabcolsep) * \real{0.1429}}
  >{\raggedright\arraybackslash}p{(\linewidth - 12\tabcolsep) * \real{0.1429}}
  >{\raggedright\arraybackslash}p{(\linewidth - 12\tabcolsep) * \real{0.1429}}
  >{\raggedright\arraybackslash}p{(\linewidth - 12\tabcolsep) * \real{0.1429}}
  >{\raggedright\arraybackslash}p{(\linewidth - 12\tabcolsep) * \real{0.1429}}
  >{\raggedright\arraybackslash}p{(\linewidth - 12\tabcolsep) * \real{0.1429}}
  >{\raggedright\arraybackslash}p{(\linewidth - 12\tabcolsep) * \real{0.1429}}@{}}
\caption{\textbf{IEC 60204-1 Emergency Stopping Categories and Hardware
Implementation Standards.} Functional mapping of international safety
stopping categories to electrical power states, braking mechanisms,
activation latencies, mechanical component degradation, and mission
recovery workflows. Category 2 (SS2) and Category 1 (SS1) utilize
controlled electronic deceleration to eliminate friction brake ablation
and structural shock, whereas Category 0 (STO) enforces uncoordinated
galvanic isolation via dual-channel hardware safety paths to guarantee
fail-safe de-energization during catastrophic hardware or timing
faults.}\label{tbl-12-iec60204-stopping-categories}\tabularnewline
\toprule\noalign{}
\begin{minipage}[b]{\linewidth}\raggedright
Stop Category \& Standard
\end{minipage} & \begin{minipage}[b]{\linewidth}\raggedright
Functional Definition \& Power Status
\end{minipage} & \begin{minipage}[b]{\linewidth}\raggedright
Deceleration \& Braking Mechanism
\end{minipage} & \begin{minipage}[b]{\linewidth}\raggedright
Triggering Fault Condition
\end{minipage} & \begin{minipage}[b]{\linewidth}\raggedright
Engagement Latency (\(\tau_{\text{engage}}\))
\end{minipage} & \begin{minipage}[b]{\linewidth}\raggedright
Mechanical Wear \& Dynamic Stress
\end{minipage} & \begin{minipage}[b]{\linewidth}\raggedright
Mission Context \& Recovery Workflow
\end{minipage} \\
\midrule\noalign{}
\endfirsthead
\toprule\noalign{}
\begin{minipage}[b]{\linewidth}\raggedright
Stop Category \& Standard
\end{minipage} & \begin{minipage}[b]{\linewidth}\raggedright
Functional Definition \& Power Status
\end{minipage} & \begin{minipage}[b]{\linewidth}\raggedright
Deceleration \& Braking Mechanism
\end{minipage} & \begin{minipage}[b]{\linewidth}\raggedright
Triggering Fault Condition
\end{minipage} & \begin{minipage}[b]{\linewidth}\raggedright
Engagement Latency (\(\tau_{\text{engage}}\))
\end{minipage} & \begin{minipage}[b]{\linewidth}\raggedright
Mechanical Wear \& Dynamic Stress
\end{minipage} & \begin{minipage}[b]{\linewidth}\raggedright
Mission Context \& Recovery Workflow
\end{minipage} \\
\midrule\noalign{}
\endhead
\bottomrule\noalign{}
\endlastfoot
\textbf{Safely-Limited Speed / Force (SLS/SLF)}\emph{IEC 61800-5-2 / ISO
10218} & Power continuously applied; active closed-loop PWM modulation &
Real-time orthogonal QP projection
(\(\mathbf{u}^* = \text{proj}(\mathbf{u}_{\text{nom}})\)) & Nominal
proposal encroaching on barrier margin
(\(a^T \mathbf{u}_{\text{nom}} > b\)) & \(< 175\,\mu\text{s}\) (Inline
single-cycle DTCM projection) & Zero friction wear; nominal inverter
thermal dissipation & Continuous; zero mission interruption; transparent
to upstream planner \\
\textbf{Category 2 (Safe Stop 2, SS2)}\emph{IEC 60204-1 §9.2.2 / ISO
13849-1} & Power continuously applied to motor drives; PWM switching
active & Controlled deceleration along safe profile to stationary
position hold & Intent lease expiration
(\(\tau_{\text{lease}} > 40\text{ ms}\)); minor tracking error
(\(\epsilon_{\text{track}} < e \le \epsilon_{\text{crit}}\)) &
\(1.0\text{--}3.0\text{ ms}\) (Deterministic firmware interrupt) & Zero
friction wear; continuous Joule heating (\(I^2 R\)) in windings & Fully
reversible; state space preserved; instant motion resumption on
re-lease \\
\textbf{Category 1 (Safe Stop 1, SS1)}\emph{IEC 60204-1 §9.2.2 / SIL 3 /
PLe} & Power maintained during braking, then galvanically isolated (STO
engaged) & Maximum counter-electromotive braking (\(a_{\max}\)) via
inverter or dynamic chopper resistors & Critical tracking error
(\(e > \epsilon_{\text{crit}} = 8.0\text{ mm}\)); barrier breach
(\(h \le 0\)); QP infeasibility & \(2.0\text{--}15.0\text{ ms}\) (MCU
interrupt; braking duration \(50\text{--}300\text{ ms}\)) & Peak
gear-tooth torque; DC-bus regenerative surge absorbed by shunts &
Mission aborted; trajectory discarded; requires supervisor state-machine
reset \\
\textbf{Category 0 (Safe Torque Off, STO)}\emph{IEC 60204-1 §9.2.2 /
ASIL-D / PLe} & Immediate, uncoordinated removal of electrical power
from motor windings & Uncontrolled coasting arrested by passive
spring-applied friction brakes & Hardware watchdog timeout
(\(t_{\text{timeout}} \ge 3.0\text{ ms}\)); physical E-stop; rail limit
switch & \(< 5\,\mu\text{s}\) silicon gate cutoff;
\(20\text{--}80\text{ ms}\) mechanical brake pad drop & Severe friction
pad ablation; structural shock on harmonic flexsplines & Irreversible
abort; possible lost absolute commutation; manual homing calibration \\
\end{longtable}

Every rung of the ladder assumes that the hardware executing the
fallback remains independent of the fault that triggered the escalation.
If an enforcer shares its power rail, crystal oscillator, or memory bus
with the host processor producing the unsafe commands, a single
electrical disturbance can compromise both the nominal planner and the
defensive fallback. A deterministic state machine in firmware cannot
provide safety if the physical substrate executing it lacks genuine
isolation.

\section{How Enforcement Goes Wrong}\label{sec-enforcement-12-10}

Placing a secondary processor alongside the primary compute module
creates an illusion of architectural independence that physical
packaging routinely undermines. When a machine hosts a trajectory
planner on a high-throughput system-on-chip and an enforcement monitor
on a real-time microcontroller, both chips almost always share a common
circuit board. That shared board introduces common-mode coupling through
power distribution networks, clock sources, and ground planes. If a
single \(5.0\text{ V}\) to \(3.3\text{ V}\) step-down regulator powers
both processors, an electrical short, transient latch-up, or thermal
shutdown on the host side drops the supply rail below the
microcontroller operating threshold of \(2.7\text{ V}\), resetting both
chips simultaneously. Similarly, routing clock lines from a single
crystal oscillator or distribution tree creates a shared failure point
where a crystal fracture or line ringing halts both execution pipelines
on the exact same cycle. High return currents from pulse-width-modulated
motor drivers elevate local potential across the digital ground plane,
inducing ground bounce that corrupts logic thresholds on shared
board-level communication buses like Inter-Integrated Circuit or Serial
Peripheral Interface lines. These hardware failure modes are completely
invisible to software running on the chips. A microcontroller brownout
detector can pull a reset line when rail voltage collapses, but if the
reference supply to the detector itself drops, the enforcer fails
silently without executing its defensive routine.

Even when physical isolation across silicon and power supplies is
complete, an enforcer fails if its computational delay falls out of
phase with the dynamics of the plant. A safety filter evaluates
admissibility by projecting state trajectories or evaluating control
barrier functions at each discrete control tick. If the calculation
budget is allocated \(1.0\text{ ms}\) within a \(1000\text{ Hz}\)
control loop, but complex boundary geometry or numerical conditioning
issues push solver execution to \(12.0\text{ ms}\), the resulting
command arrives eleven cycles late. For an articulated link moving at an
angular velocity of \(3.0\text{ rad/s}\), a \(12.0\text{ ms}\) delay
represents an unmodeled angular displacement of \(0.036\text{ rad}\), or
approximately \(2.1^{\circ}\). When the enforcer finally asserts an
emergency braking torque computed for the earlier state, that torque
acts against physical coordinates the machine has already passed. In an
underdamped mechanical assembly, an intervention executed out of phase
with plant velocity injects energy into the structural resonance instead
of damping it, accelerating the body toward the obstacle the filter was
constructed to avoid. A hardware watchdog or real-time timer interrupt
detects the deadline overrun and aborts the cycle, but it cannot undo
the physical momentum accrued during the overrun period.

An enforcer cannot observe physical mechanics directly; it observes
digital state estimates reconstructed from sensor data. When the sensing
pipeline introduces phase lag, calibration drift, or undetected
mechanical slip, the enforcer becomes blind to boundary violations
occurring in the physical world. Consider an optical encoder measuring
motor shaft angle across an elastomeric timing belt. If dynamic belt
stretch under high torque produces \(4.0\text{ mm}\) of linear
compliance, or if the drive pulley slips by several gear teeth during an
aggressive acceleration transient, the joint state reported to the
enforcer remains within allowable geometric limits while the physical
link has breached a virtual wall. Low-pass filters applied to eliminate
high-frequency sensor noise introduce group delay. A second-order filter
with a \(20\text{ Hz}\) cutoff frequency delays the velocity estimate by
approximately \(15\text{ ms}\). If an unexpected contact begins at
\(t = 0\text{ ms}\), the filtered velocity signal does not cross the
over-speed fault threshold until impact forces have already peaked and
loaded the mechanical transmission. Kinematic cross-validation against
redundant absolute encoders or joint torque sensors can detect large
structural discrepancies, but slow thermal bias drift in inertial
measurement units and sub-threshold mechanical slip remain undetected
until physical contact occurs.

A common engineering misconception equates deterministic software
execution with system safety. Determinism guarantees that a computation
finishes in an exact number of clock cycles and produces bit-identical
outputs for identical inputs, but it provides no guarantee regarding the
physical validity of the underlying equations. A real-time operating
system executing a control barrier filter on an isolated core with zero
microsecond scheduling jitter is fully deterministic. Yet if the plant
model programmed into that filter assumes a nominal friction coefficient
of \(\mu = 0.80\) when oil contamination on the rail has dropped the
actual friction to \(\mu = 0.15\), the enforcer calculates an analytical
stopping distance of \(0.40\text{ m}\) for a deceleration maneuver that
physically requires \(2.13\text{ m}\) of travel. The enforcer executes
its intervention with microsecond precision, asserts the calculated
braking current on schedule, and deterministically crashes the carriage
into the mechanical end-stop. Determinism guarantees repeatable
behavior, but when paired with an unmodeled payload mass, an inaccurate
inertia tensor, or an uncalibrated force constant, it simply guarantees
that the failure will be perfectly repeatable. Open-loop timing monitors
cannot detect model inaccuracy, which reveals itself only when
closed-loop tracking error diverges in downstream physical states.

Incident analysis demonstrates how software interlock elimination, bus
contention, and resource starvation defeat safety enforcement in
production hardware. In multi-processor embedded designs, the safety
enforcer frequently shares a physical communication bus with background
telemetry, logging, or debugging subsystems. When an asynchronous burst
of diagnostic logging saturates the shared Serial Peripheral Interface
bus, the enforcer core can become starved of direct memory access
bandwidth or trapped in an un-preemptible interrupt service routine
while awaiting a bus acknowledge signal.\sidenote{\footnotesize On SoC interconnects
  (ARM AHB matrix), round-robin DMA arbitration allows bulk telemetry
  dumps (8 KB flash writes) to saturate the bus matrix. This starves
  real-time SPI DMA channels driving motor DACs, inducing millisecond
  enforcement overruns despite idle processor cores.} During this
communication dead-time, the actuator receives no updated torque
commands, retaining its last valid setpoint while momentum carries the
machine toward a physical barrier.

\protect\phantomsection\label{callout-incidentux2a-12.9}
\begin{fbxSimple}{callout-incident}{Landmark Engineering Incident:}{Therac-25 Software Interlock Elimination and Table Positioning Over-Torque}
\phantomsection\label{callout-incident*-12.9}

\textbf{System:} AECL Therac-25 dual-mode (\(25\text{ MeV}\) electron /
X-ray) medical linear accelerator and automated radiotherapy patient
table, controlled by a single PDP-11 computer running uniprocessor
assembly firmware (\citeproc{ref-leveson1993investigating}{Leveson and
Turner 1993}).

\textbf{Incident:} Between 1985 and 1987, six patients received massive
radiation overdoses (up to \(20{,}000\text{ rad}\), over \(100\times\)
the prescribed therapeutic dose) and multiple patients suffered fatal
radiation-induced injuries accompanied by uncommanded motorized patient
positioning table over-torque collisions.

\textbf{Telemetry \& Software Analysis:}

\begin{itemize}
\tightlist
\item
  In preceding models (Therac-6 and Therac-20), physical safety was
  guaranteed by independent hardware interlocks: physical microswitches,
  mechanical turntable detents, and analog diode steering circuits that
  physically prevented the linear accelerator from firing in
  high-current mode unless the beam-flattening target and scattering
  foil were mechanically locked in place.
\item
  In the Therac-25, engineers removed all physical hardware interlocks
  to reduce manufacturing cost and simplify hardware design, delegating
  exclusive safety authority to software running on the PDP-11 computer.
\item
  \textbf{Root Cause 1 (8-Bit Counter Rollover):} A software safety
  interlock routine incremented an 8-bit shared variable
  (\texttt{Class3\_done}) on every evaluation loop. When an operator
  completed setup at the exact moment the counter incremented from 255
  to 256 (\(0\text{xFF} \to 0\text{x00}\)), the software evaluated
  \texttt{Class3\_done\ ==\ 0} as ``checks complete and safe,''
  bypassing upper-column interlock verification entirely.
\item
  \textbf{Root Cause 2 (Concurrent Keypress Race Condition):} When an
  experienced technician rapidly corrected an operating mode selection
  on the VT100 terminal in under 8 seconds (typing \texttt{X},
  cursor-up, \texttt{E}, and \texttt{Return}), the input parser updated
  the beam energy register to 25 MeV electron mode while a separate
  background task was still positioning the mechanical turntable
  magnets. The software fired the high-current beam into the patient
  without the scattering foil or tungsten target in position.
\item
  \textbf{Root Cause 3 (Shared Software Motor Table Over-Torque):} Table
  positioning motors and turntable drives shared processor execution
  cycles without independent closed-loop current limiting or end-stop
  limit switch cutoffs. When command queues overflowed, table drives
  slammed patients into machine housings under full motor stall torque.
\end{itemize}

\textbf{Physical AI Architectural Lesson:} Software-only safety filters
executing on complex, non-deterministic compute stacks cannot replace
independent bare-metal hardware trip zones. The Therac-25 disaster is
the canonical historical proof that removing physical interlocks in
favor of software decision logic transforms transient concurrency bugs
and arithmetic overflows directly into catastrophic kinetic and physical
energy releases. In modern Physical AI architectures, the nervous system
must isolate safety enforcement onto dedicated microcontrollers with
hardware break circuits (\texttt{TIMx\_BKIN}, Safe Torque Off relays)
that physically de-energize gate drivers within microseconds, completely
bypassing operating system schedulers and application
software.\sidenote{\footnotesize Leveson and Turner (1993) demonstrated the fatal
  hazard of replacing physical hardware interlocks with software logic.
  Therac-25 suffered arithmetic counter overflows, missing mechanical
  turntable feedback, and shared UI/actuator execution threads, causing
  catastrophic radiation overdoses.}

\end{fbxSimple}

\protect\phantomsection\label{callout-incidentux2a-12.10}
\begin{fbxSimple}{callout-incident}{Landmark Engineering Incident:}{Autonomous Rail Gantry Emergency Brake Lockout via SPI Bus Contention}
\phantomsection\label{callout-incident*-12.10}

\textbf{System:} Heavy industrial rail-mounted automated storage gantry
(\(850\text{ kg}\) carriage) driven by linear synchronous motors,
controlled by an embedded Linux trajectory planner over SPI to a
bare-metal safety co-processor running a \(1000\text{ Hz}\) barrier
enforcement loop.

\textbf{Telemetry Snapshot:} At \(t = 0.000\text{ s}\), host Linux
processor initiates an asynchronous \(16\text{ KB}\) diagnostic flash
memory dump over the shared SPI bus; at \(t = 0.004\text{ s}\), the
gantry approaches the dynamic braking boundary
(\(d_{\text{clear}} = 0.45\text{ m}\)); at \(t = 0.005\text{ s}\), the
safety enforcer attempts to transmit emergency dynamic braking registers
over SPI; from \(t = 0.005\text{ s}\) to \(0.016\text{ s}\), the SPI DMA
bus arbiter blocks the safety transaction due to non-preemptible flash
DMA transfer; at \(t = 0.018\text{ s}\), carriage breaches the
mechanical boundary and strikes the end-stop at \(3.2\text{ m/s}\).

\textbf{Root Cause:} DMA priority inversion and bus starvation on a
shared SPI peripheral bus locked out the safety enforcer from updating
motor drive registers during the critical \(11\text{ ms}\) stopping
window.

\textbf{Architectural Flaw:} Sharing a physical communication bus
between non-real-time diagnostic telemetry and safety-critical motor
command channels without dedicated hardware preemption or galvanically
independent emergency stop routing.

\textbf{Provenance:} European Machinery Safety Directive (2006/42/EC)
statutory industrial robotics incident investigation archive and
manufacturer failure analysis report.

\end{fbxSimple}

Every failure mode in the enforcement layer traces back to a breakdown
of independence or a divergence between the internal plant
representation and physical reality. The nervous system cannot treat its
safety claims as unconditional truths. A mathematical barrier
certificate, a dual-core silicon package, and a zero-jitter real-time
scheduler establish safety only within the bounded envelope where power
supplies remain isolated, state estimates match physical mechanics, and
execution deadlines are met without exception.

\section{Enforcer Telemetry}\label{sec-enforcement-12-11}

Because safety guarantees cannot hold unconditionally across arbitrary
physical states, the nervous system must formalize the exact boundary
where its protections apply. The \textbf{enforcement record} is the data
structure that exports this boundary to system integration, deployment,
and task-placement layers. Rather than treating safety enforcement as an
ambient property of the software stack, the record defines an explicit
specification contract. It establishes what the permission path
guarantees, the physical parameters it assumes, the execution budgets it
requires, and the failure states that cause it to refuse action.
Following the baseline specification schema and provenance conventions
established in earlier chapters (such as Chapter~\ref{sec-nervous}),
this record appends the temporal, spatial, and dynamic fields unique to
closed-loop force regulation and intervention.

The \textbf{enforcement record} is the immutable specification contract
and runtime audit ledger that exports the verified tracking error bounds
(\(\epsilon_{\text{track}}\)), stopping envelopes
(\(t_{\text{stop}}, d_{\text{stop}}\)), Control Barrier Function
parameters, discrete fallback escalation thresholds, and certified
environmental operating assumptions across the permission boundary.

The specification contract grounds its guarantees in four primary
quantified parameters: the verified tracking error bound
\(\epsilon_{\text{track}}\), the maximum stopping time
\(t_{\text{stop}}\), the barrier function coefficients, and the loop
period \(T_{\text{loop}}\). Consider an analytical example of a linear
gantry stage that regulates contact force before translating a payload.
The nervous system executes its inner regulation loop at a fixed period
of \(T_{\text{loop}} = 1.0\text{ ms}\) (\(1000\text{ Hz}\)). Across
nominal operations, closed-loop tracking remains within
\(\epsilon_{\text{track}} = 4.2\text{ mm}\) at \(P_{99.99}\) under
maximum rated load disturbance. If the safety filter intervenes, the
downstream braking mechanism brings the carriage to rest from its
maximum allowable velocity of \(1.5\text{ m/s}\) within
\(t_{\text{stop}} = 85\text{ ms}\), traversing an analytical stopping
distance of \(0.064\text{ m}\) under peak deceleration. When the
intervention relies on a control barrier function \(h(x) \ge 0\) with
class-\(\mathcal{K}\) linear parameter \(\alpha = 12.0\text{ s}^{-1}\),
the record binds these exact values directly to the physical plant
model. Upstream trajectory generators cannot treat the enforcer as an
unconstrained safety net, because they must budget a minimum spatial
clearance of
\(\Delta x_{\text{margin}} \ge \epsilon_{\text{track}} + v_{\text{max}} t_{\text{stop}}\)
to prevent an intervention from contacting a physical limit.

Every numerical entry in the record carries an explicit provenance tag
documenting whether the bound originates from analytical derivation,
physics simulation, or hardware worst-case stress testing. These three
sources offer vastly different levels of real-world confidence. An
analytical derivation based on rigid-body equations provides a
mathematical baseline, such as an ideal stopping time of
\(t_{\text{stop}} = 62\text{ ms}\) that assumes instantaneous current
reversal and constant traction. A multibody physics simulation sweeping
parameter variations across fifty thousand runs might incorporate
thermal winding degradation and calculate a \(P_{99.9}\) tail of
\(t_{\text{stop}} = 78\text{ ms}\). Direct hardware injection testing on
the physical rail under cold-temperature lubricant conditions reveals
contactor bounce and electromagnetic latching latency, establishing the
true worst-case bound of \(t_{\text{stop}} = 85\text{ ms}\). Documenting
provenance prevents an integrator from mistaking an uncalibrated
simulation boundary for an empirically validated hardware invariant.

The record formalizes the deterministic fallback ladder by mapping
discrete numerical thresholds to specific protective actions. Escalation
does not rely on coarse error strings or unquantified health flags. In
our analytical positioning stage, Level 1 nominal operation permits
quadratic programming filter adjustments as long as tracking error
satisfies \(e(t) \le \epsilon_{\text{track}} = 4.2\text{ mm}\) and
optimization solver latency remains below
\(T_{\text{solve}} \le 650\ \mu\text{s}\). If tracking error diverges
into the intermediate window between \(4.2\text{ mm}\) and
\(\epsilon_{\text{crit}} = 8.0\text{ mm}\), or if solver execution
exceeds \(800\ \mu\text{s}\), the system escalates to Level 2,
overriding policy proposals with a deterministic safe setpoint that
decelerates along a precomputed trajectory. If tracking error exceeds
\(\epsilon_{\text{crit}} = 8.0\text{ mm}\) or the barrier margin
violates \(h(x) \le 0\), the controller escalates to Level 3, shorting
the motor windings through dynamic braking resistors to dissipate
kinetic energy. Finally, if hardware watchdog heartbeats cease for more
than \(t_{\text{timeout}} = 3.0\text{ ms}\) or mechanical overtravel
switches trip, the circuit escalates to Level 4, de-energizing the
primary safety contactor to cut motor bus power completely.

To formalize this specification contract across the dual-brain boundary,
Table~\ref{tbl-12-safety-schema} defines the abstract data schema for
the safety invariant envelope. Every field is bound to an explicit
physical SI unit, a real-time validation invariant, an independent
hardware verification substrate, and a deterministic fallback escalation
tier.

\begin{longtable}[]{@{}
  >{\raggedright\arraybackslash}p{(\linewidth - 10\tabcolsep) * \real{0.1667}}
  >{\raggedright\arraybackslash}p{(\linewidth - 10\tabcolsep) * \real{0.1667}}
  >{\raggedright\arraybackslash}p{(\linewidth - 10\tabcolsep) * \real{0.1667}}
  >{\raggedright\arraybackslash}p{(\linewidth - 10\tabcolsep) * \real{0.1667}}
  >{\raggedright\arraybackslash}p{(\linewidth - 10\tabcolsep) * \real{0.1667}}
  >{\raggedright\arraybackslash}p{(\linewidth - 10\tabcolsep) * \real{0.1667}}@{}}
\caption{\textbf{Abstract Data Schema for the Real-Time Safety Invariant
Envelope (\protect\texttt{SafetyInvariantEnvelope}).} Specification
contract layout (\(68\text{ bytes}\) total), binary data types, physical
SI units, mathematical verification invariants, hardware audit
substrates, and deterministic escalation tiers governing real-time
safety enforcement across the proposal--permission
boundary.}\label{tbl-12-safety-schema}\tabularnewline
\toprule\noalign{}
\begin{minipage}[b]{\linewidth}\raggedright
Schema Field Identifier
\end{minipage} & \begin{minipage}[b]{\linewidth}\raggedright
Binary Type \& Layout
\end{minipage} & \begin{minipage}[b]{\linewidth}\raggedright
Physical SI Units
\end{minipage} & \begin{minipage}[b]{\linewidth}\raggedright
Verification Invariant
\end{minipage} & \begin{minipage}[b]{\linewidth}\raggedright
Hardware Audit Substrate
\end{minipage} & \begin{minipage}[b]{\linewidth}\raggedright
Fallback Escalation Tier
\end{minipage} \\
\midrule\noalign{}
\endfirsthead
\toprule\noalign{}
\begin{minipage}[b]{\linewidth}\raggedright
Schema Field Identifier
\end{minipage} & \begin{minipage}[b]{\linewidth}\raggedright
Binary Type \& Layout
\end{minipage} & \begin{minipage}[b]{\linewidth}\raggedright
Physical SI Units
\end{minipage} & \begin{minipage}[b]{\linewidth}\raggedright
Verification Invariant
\end{minipage} & \begin{minipage}[b]{\linewidth}\raggedright
Hardware Audit Substrate
\end{minipage} & \begin{minipage}[b]{\linewidth}\raggedright
Fallback Escalation Tier
\end{minipage} \\
\midrule\noalign{}
\endhead
\bottomrule\noalign{}
\endlastfoot
\texttt{seq\_id} & \texttt{uint64\_t} (\(8\text{ B}\)) & Monotonic
counter & \(s_k = s_{k-1} + 1\) & MCU DMA Packet Parser & Level 2 (Hold
on gap) \\
\texttt{timestamp\_us} & \texttt{uint64\_t} (\(8\text{ B}\)) &
Microseconds (\(\mu\text{s}\)) &
\(t_{\text{now}} - t_{\text{stamp}} \le \tau_{\text{lease}}\) & Hardware
Timer 64-bit & Level 2 (\(\tau_{\text{lease}} > 40\text{ ms}\)) \\
\texttt{pos\_tracking\_bound} & \texttt{float32\_t} (\(4\text{ B}\)) &
Meters (\(\text{m}\)) &
\(\|x_{\text{meas}} - r_{\text{cmd}}\| \le \epsilon_{\text{track}}\) &
Dual Optical Encoders & Level 2 / Level 3 \\
\texttt{vel\_ceiling\_limit} & \texttt{float32\_t} (\(4\text{ B}\)) &
Radians/sec (\(\text{rad/s}\)) & \(\|\dot{q}\| \le \dot{q}_{\max}\) &
Timer Input Capture & Level 1 (QP Clamp) \\
\texttt{torque\_slew\_ceiling} & \texttt{float32\_t} (\(4\text{ B}\)) &
\(\text{N}\cdot\text{m/s}\) & \(\|\dot{\tau}\| \le \dot{\tau}_{\max}\) &
Inverter Current Loop & Level 1 (Slew Limiter) \\
\texttt{cbf\_barrier\_margin} & \texttt{float32\_t{[}K{]}}
(\(24\text{ B}\)) & Normalized distance &
\(h_k(\mathbf{x}) \ge \delta_{\text{margin}, k}\) & Active-Set QP Engine
& Level 1 / Level 3 \\
\texttt{actuator\_curr\_limit} & \texttt{float32\_t} (\(4\text{ B}\)) &
Amperes (\(\text{A}\)) & \(|I_{\text{phase}}| \le I_{\text{peak}}\) &
ADC Current Shunts & Level 3 (Dynamic Brake) \\
\texttt{dc\_bus\_overvolt\_thresh} & \texttt{float32\_t}
(\(4\text{ B}\)) & Volts (\(\text{V}\)) &
\(V_{\text{bus}} \le V_{\text{max\_safe}}\) & Analog Comparator
(\texttt{COMP1}) & Level 3 (Brake Chopper) \\
\texttt{watchdog\_timeout\_us} & \texttt{uint32\_t} (\(4\text{ B}\)) &
Microseconds (\(\mu\text{s}\)) &
\(t_{\text{silent}} \le t_{\text{timeout}}\) & Silicon Windowed Watchdog
& Level 4 (Safe Torque Off) \\
\texttt{crc32\_checksum} & \texttt{uint32\_t} (\(4\text{ B}\)) &
CRC32-IEEE 802.3 & \(\text{CRC}(\text{payload}) == \text{checksum}\) &
Hardware CRC Accelerator & Level 2 (Reject packet) \\
\end{longtable}

The complete specification contract for this multi-tiered
defense-in-depth strategy is formalized in
Table~\ref{tbl-12-fallback-hierarchy}. By binding discrete quantitative
thresholds directly to execution substrates, power-stage actuation
states, stopping distance budgets, and mission recovery protocols, the
enforcement record provides the deterministic boundary required for safe
system integration.

\begin{longtable}[]{@{}
  >{\raggedright\arraybackslash}p{(\linewidth - 10\tabcolsep) * \real{0.1667}}
  >{\raggedright\arraybackslash}p{(\linewidth - 10\tabcolsep) * \real{0.1667}}
  >{\raggedright\arraybackslash}p{(\linewidth - 10\tabcolsep) * \real{0.1667}}
  >{\raggedright\arraybackslash}p{(\linewidth - 10\tabcolsep) * \real{0.1667}}
  >{\raggedright\arraybackslash}p{(\linewidth - 10\tabcolsep) * \real{0.1667}}
  >{\raggedright\arraybackslash}p{(\linewidth - 10\tabcolsep) * \real{0.1667}}@{}}
\caption{\textbf{Deterministic Enforcer Fallback Hierarchy and
Specification Contract Matrix.} Formal specification of the multi-tiered
defense-in-depth fallback hierarchy recorded in the chapter enforcement
record. The table details quantitative numerical triggers (\(e(t)\),
\(T_{\text{solve}}\), \(t_{\text{timeout}}\), \(h(\mathbf{x})\),
\(I_{\text{cmd}}\)), physical silicon execution substrates, inverter
power-stage actions, certified worst-case stopping distance budgets from
\(v_{\max} = 1.5\text{ m/s}\), and recovery protocols across nominal,
degraded, emergency, and catastrophic operating
regimes.}\label{tbl-12-fallback-hierarchy}\tabularnewline
\toprule\noalign{}
\begin{minipage}[b]{\linewidth}\raggedright
Escalation Level \& Action
\end{minipage} & \begin{minipage}[b]{\linewidth}\raggedright
Quantitative Trigger Bounds
\end{minipage} & \begin{minipage}[b]{\linewidth}\raggedright
Execution Substrate \& Silicon Isolation
\end{minipage} & \begin{minipage}[b]{\linewidth}\raggedright
Actuation State \& Power Stage Action
\end{minipage} & \begin{minipage}[b]{\linewidth}\raggedright
Physical Arrest Budget (\(v_{\max} = 1.5\text{ m/s}\))
\end{minipage} & \begin{minipage}[b]{\linewidth}\raggedright
Mission Reversibility \& Protocol
\end{minipage} \\
\midrule\noalign{}
\endfirsthead
\toprule\noalign{}
\begin{minipage}[b]{\linewidth}\raggedright
Escalation Level \& Action
\end{minipage} & \begin{minipage}[b]{\linewidth}\raggedright
Quantitative Trigger Bounds
\end{minipage} & \begin{minipage}[b]{\linewidth}\raggedright
Execution Substrate \& Silicon Isolation
\end{minipage} & \begin{minipage}[b]{\linewidth}\raggedright
Actuation State \& Power Stage Action
\end{minipage} & \begin{minipage}[b]{\linewidth}\raggedright
Physical Arrest Budget (\(v_{\max} = 1.5\text{ m/s}\))
\end{minipage} & \begin{minipage}[b]{\linewidth}\raggedright
Mission Reversibility \& Protocol
\end{minipage} \\
\midrule\noalign{}
\endhead
\bottomrule\noalign{}
\endlastfoot
\textbf{Level 1: Minimal Intervention} & \(e(t) \le 4.2\text{ mm}\),
\(T_{\text{solve}} \le 650\,\mu\text{s}\), \(h(\mathbf{x}) > 0\),
\(\vert I_{\text{cmd}}\vert \le 12.0\text{ A}\) & Primary Real-Time Core
(bare-metal MCU, static SRAM) & Orthogonal QP projection
\(\mathbf{u}^*\); active closed-loop FOC PWM modulation at
\(20\text{ kHz}\) & Continuous tracking
(\(d_{\text{error}} \le \epsilon_{\text{track}} = 4.2\text{ mm}\)) &
Fully transparent; automatic resumption; logs \(\Delta \mathbf{u}\)
diagnostic metric \\
\textbf{Level 2: Safe Trajectory Decel} &
\(4.2\text{ mm} < e(t) \le 8.0\text{ mm}\),
\(650\,\mu\text{s} < T_{\text{solve}} \le 800\,\mu\text{s}\),
\(\tau_{\text{lease}} > 40\text{ ms}\) & Local Safety MCU (Autonomous
bare-metal execution; static ROM splines) & Override policy setpoints;
execute smooth quintic deceleration (\(a = -8.0\text{ m/s}^2\)) to
position hold & \(t_{\text{stop}} \le 190.5\text{ ms}\),
\(d_{\text{stop}} \le 144.6\text{ mm}\) (including \(3\text{ ms}\)
detection) & Reversible; auto-resumes if policy renews lease within
\(200\text{ ms}\), else escalates to Level 3 \\
\textbf{Level 3: Dynamic Shunt Braking} & \(e(t) > 8.0\text{ mm}\),
\(h(\mathbf{x}) \le 0\), current saturation \(> 5\text{ ms}\),
\(T_{\text{solve}} > 800\,\mu\text{s}\) & Independent Safety
Co-Processor / CPLD (Isolated power rail, direct SPI/PWM override) &
Inhibit upstream PWM; switch low-side MOSFET bridge or engage dynamic
braking chopper resistors & \(t_{\text{stop}} \le 85.0\text{ ms}\),
\(d_{\text{stop}} \le 64.0\text{ mm}\) (under peak electromagnetic
counter-torque) & Non-reversible in active cycle; mission aborted;
requires safety latch reset and re-zeroing \\
\textbf{Level 4: Galvanic Power Cutoff} & Watchdog trip
(\(t_{\text{timeout}} \ge 3.0\text{ ms}\)), physical E-stop switch,
overtravel limit switch & Galvanically Isolated Discrete Safety Relay
Logic (Hardware interlock; zero software) & De-energize gate-driver
optocouplers (\(< 5\,\mu\text{s}\)); de-energize contactor coils; drop
spring brakes & \(t_{\text{stop}} \le 290.0\text{ ms}\),
\(d_{\text{stop}} \le 256.0\text{ mm}\) (including \(30\text{ ms}\)
brake pad latching dead-time) & Latched hardware fault; requires
physical inspection, manual reset button, and homing calibration \\
\end{longtable}

No enforcement guarantee holds when the machine operates outside its
designed physical environment, so the record bounds the plant dynamics
within an explicit operating envelope. For the gantry stage, the record
specifies a Coulomb friction coefficient bounded by
\(\mu \in [0.65, 0.90]\), a total payload mass of
\(m_{\text{load}} \in [0.0\text{ kg}, 5.0\text{ kg}]\) attached to a
carriage mass of \(12.0\text{ kg}\), and an ambient operating
temperature range between \(5\ ^\circ\text{C}\) and
\(45\ ^\circ\text{C}\). If water contamination drops rail friction to
\(\mu = 0.20\), or if an unmodeled payload mass of \(9.0\text{ kg}\) is
mounted to the end-effector, the available actuator torque becomes
mathematically incapable of satisfying the barrier contraction condition
\(\dot{h}(x) \ge -\alpha h(x)\). The record exposes these constraints
directly to supervisory monitors, which must revoke policy dispatch and
trigger a controlled shutdown whenever sensors detect that physical
parameters have crossed outside the certified envelope.

An enforcement record complete with verified mathematical parameters,
strict escalation thresholds, and documented operating boundaries
provides the formal contract required to integrate learned control into
physical hardware. Yet a specification contract written to disk
guarantees safety only if the underlying execution platform can satisfy
its timing and computational constraints under physical load. If an
enforcer is mathematically verified but stalls on a saturated memory bus
during a high-speed transit maneuver, the failure occurs in time rather
than in logic.

\section{Systems Stress-Tests}\label{sec-enforcement-12-12}

An enforcer that is mathematically correct in logic but too slow to
execute in physical time provides zero protection to an operating
machine. When a barrier solver overruns its execution budget during a
high-speed transit maneuver, the moving mass does not pause to wait for
the computation to finish. Momentum continues to integrate velocity
along the unbraked trajectory, carrying the body across the safe
boundary before the microcontroller can assert counter-torque. Verifying
the enforcement organ requires evaluating its behavior under deliberate
computational and physical faults. The following four analytical
stress-tests evaluate the safety filter under synthetic latency jitter,
mechanical shock loads, process termination, and corrupted state
feedback:

\subsection*{Stress-Test 1: Real-Time Deadline Violation (Barrier QP
Solver
Jitter)}\label{stress-test-1-real-time-deadline-violation-barrier-qp-solver-jitter}
\addcontentsline{toc}{subsection}{Stress-Test 1: Real-Time Deadline
Violation (Barrier QP Solver Jitter)}

\emph{The scenario.} A linear gantry carrying a \(4.5\text{ kg}\)
payload operates at \(v = 2.4\text{ m/s}\) toward an end-of-travel
buffer. The real-time safety enforcer executes an active-set Quadratic
Program at \(1000\text{ Hz}\) (\(T_{\text{period}} = 1.0\text{ ms}\),
nominal solve time \(420.0\ \mu\text{s}\)). Synthetic memory contention
and cache line invalidations cause solver latency to stretch to
\(910.0\ \mu\text{s}\), exceeding the hard \(800.0\ \mu\text{s}\) Level
2 timeout.

\emph{The systems dilemma.} Sizing safety around mean solve times leads
to overtravel. Permitting the solver to run past the deadline misses the
timer interrupt update reload, driving the linear motor under a stale
torque command while the carriage moves \(2.4\text{ mm/ms}\). The
real-time executive must abort the solver precisely at
\(800.0\ \mu\text{s}\), injecting a precomputed Level 2 deceleration
setpoint from static local SRAM that halts the carriage before the
end-stop buffer.

\emph{The systems verdict: Bounded Execution Deadline Law.} A safety
proof is valid only within the clock cycle in which it completes. An
enforcer that computes a safe action \(50.0\ \mu\text{s}\) after the
timer interrupt has fired does not prevent a mechanical collision;
real-time safety filters must fail closed via deterministic hardware
timeouts.

\subsection*{Stress-Test 2: Actuator Saturation under Shock Load
(Dynamic Viability Kernel
Contraction)}\label{stress-test-2-actuator-saturation-under-shock-load-dynamic-viability-kernel-contraction}
\addcontentsline{toc}{subsection}{Stress-Test 2: Actuator Saturation
under Shock Load (Dynamic Viability Kernel Contraction)}

\emph{The scenario.} A Cartesian gantry (\(m = 16.5\text{ kg}\))
traverses at \(v_{\text{max}} = 2.8\text{ m/s}\) toward a boundary with
remaining clearance \(d_{\text{clear}} = 65.0\text{ mm}\). The drive
motor has a peak saturation limit of \(I_{\text{peak}} = 28.0\text{ A}\)
(\(F_{\text{peak}} = 350.0\text{ N}\)). An unmodeled external step load
of \(F_{\text{ext}} = 180.0\text{ N}\) opposes motion, driving the
tracking controller into full current saturation.

\emph{The systems dilemma.} Demanded current hits the \(28.0\text{ A}\)
limit. If tracking integral error accumulates while saturated,
discharging the wound-up state when resistance clears will drive
\(350.0\text{ N}\) forward thrust toward the boundary during maximum
braking. The enforcer must detect current saturation within
\(1.0\text{ ms}\), clamp integral accumulators, and contract the barrier
viability kernel to reflect physical current ceilings.

\emph{The systems verdict: Control Authority Saturation Invariant.} A
Control Barrier Function assumes commanded control outputs can be
converted into physical force. When actuators saturate, control
authority at the margin vanishes; safety filters must dynamically
contract admissible sets to match physical voltage and current limits.

\subsection*{Stress-Test 3: Sudden Upstream Policy Death (Autonomous
Low-Level Stopping
Authority)}\label{stress-test-3-sudden-upstream-policy-death-autonomous-low-level-stopping-authority}
\addcontentsline{toc}{subsection}{Stress-Test 3: Sudden Upstream Policy
Death (Autonomous Low-Level Stopping Authority)}

\emph{The scenario.} A neural policy on an application processor
generates trajectory chunks at \(50\text{ Hz}\) for a linear stage
moving at \(v = 2.0\text{ m/s}\) (\(E_k = 33.0\text{ J}\)) governed by
an intent lease of \(\tau_{\text{lease}} = 40.0\text{ ms}\) and a
\(100\text{ Hz}\) heartbeat. At \(t = 1.200\text{ s}\), an unhandled
kernel panic kills the host policy process.

\emph{The systems dilemma.} The real-time MCU detects heartbeat
cessation within \(3.0\text{ ms}\). If the MCU cuts contactors via Safe
Torque Off, the carriage coasts \(290.0\text{ mm}\) under friction and
crashes into hardstops. Instead, the MCU autonomously executes Level 2
controlled deceleration (\(a_{\text{brake}} = -8.0\text{ m/s}^2\))
stored in local SRAM, arresting the mass in
\(d_{\text{arrest}} = 256.0\text{ mm}\) and preserving
\(18.4\text{ mm}\) of clearance before the boundary.

\emph{The systems verdict: Autonomous Safety Authority Invariant.}
High-level learned planners are unverified processes that can crash or
deadlock without notice. Physical protection survives only when the
deterministic low-level enforcer possesses independent stopping
trajectories and the exclusive authority to execute them.

\subsection*{Stress-Test 4: Sensor Estimate Corruption (Residual
Innovation
Bounds)}\label{stress-test-4-sensor-estimate-corruption-residual-innovation-bounds}
\addcontentsline{toc}{subsection}{Stress-Test 4: Sensor Estimate
Corruption (Residual Innovation Bounds)}

\emph{The scenario.} A precision motion stage uses dual feedback
sensing: an optical encoder (\(0.5\ \mu\text{m}\) resolution) and an
accelerometer filtered at \(1000\text{ Hz}\). During a
\(2.2\text{ m/s}\) sweep, an optical read-head contamination fault
introduces a progressive position bias of
\(\Delta x_{\text{drift}} = 0.8\text{ mm/ms}\) with
\(\pm 3.5\text{ mm}\) noise spikes.

\emph{The systems dilemma.} Evaluating barrier inequalities over the
corrupted estimate leads the filter to believe the carriage is further
from the boundary than its true position, delaying braking until impact.
The enforcer evaluates residual innovations
(\(\|r_k\| = \|\mathbf{y}_k - \hat{\mathbf{y}}_k\| \le \gamma\)) at
\(1000\text{ Hz}\), detecting sensor divergence within \(2.0\text{ ms}\)
and switching to redundant sensing or dynamic braking before drift
exhausts the barrier buffer.

\emph{The systems verdict: Innovation Gating Law.} An enforcer evaluates
barrier certificates against state estimates rather than physical ground
truth. If perception fails, the enforcer guarantees safety for a phantom
system while real mass drifts into danger; physical margins must swallow
worst-case unobservable estimation error.

\section{Summary}\label{sec-enforcement-12-13}

A physical AI system resolves the tension between learned intent and
mechanical reality inside the nervous system. The brain proposes actions
through high-dimensional representations, but the nervous system alone
holds the physical authority to transform a proposal into phase current
or to reject it outright. Tracking and refusal form a single
mathematical contract. Tracking guarantees that when an action is
permitted, the low-level control loop converts that command into
mechanical force within a bounded tracking error. Refusal guarantees
that whenever a proposed command threatens the physical boundaries of
the machine, the permission path intercepts the signal and asserts a
safe alternative before momentum carries the mass into an obstacle. The
integrity of an enforcer depends on the tracking performance of the
underlying loop. A barrier certificate derived on paper becomes invalid
if the drive electronics cannot track the commanded braking current
within the time and error bounds assumed during verification.

The physical clearance protecting a machine is never a static geometric
measurement. It is a composite safety margin that compounds feedback
tracking bounds, transport latency, actuator saturation limits, and
state estimation error. As an analytical example, consider a linear
carriage translating at \(v = 1.5\text{ m/s}\) toward an end-stop. If
the current loop guarantees a maximum tracking deviation
\(\epsilon_{\text{track}} = 4.0\text{ mm}\), the transport delay across
the communication bus introduces a \(t_{\text{delay}} = 2.0\text{ ms}\)
lag during which the carriage travels
\(d_{\text{delay}} = v \cdot t_{\text{delay}} = 3.0\text{ mm}\) before
the inverter responds, the motor current limit caps braking acceleration
at \(a_{\text{max}} = 5.0\text{ m/s}^2\) requiring a stopping distance
\(d_{\text{brake}} = v^2 / (2 a_{\text{max}}) = 225.0\text{ mm}\), and
the state estimator carries an unobservable residual bound
\(\epsilon_{\text{est}} = 3.5\text{ mm}\), the minimum certified
clearance is the arithmetic sum of these terms,
\(d_{\text{clear}} \ge \epsilon_{\text{track}} + d_{\text{delay}} + d_{\text{brake}} + \epsilon_{\text{est}} = 235.5\text{ mm}\).
If an engineer budgets only the nominal stopping distance
\(d_{\text{brake}}\), the unmodeled transport lag, tracking jitter, and
sensor uncertainty consume the remaining \(10.5\text{ mm}\) of margin,
and the carriage strikes the end-stop at velocity.

This deterministic containment is achieved through forward invariance
rather than multi-step trajectory prediction. Upstream planners attempt
to ensure safety by projecting candidate trajectories forward over
finite horizons, but trajectory optimization remains vulnerable to
distribution shifts, model inaccuracies, and compute exhaustion. When an
unmodeled load appears or numerical optimization fails to converge
before the sample deadline, a predictive planner cannot certify safety.
In contrast, set-theoretic barrier enforcement evaluates an
instantaneous inequality directly on the proposed control input at the
current state. By constraining the instantaneous time derivative of a
barrier function
\(\dot{h}(\mathbf{x}, \mathbf{u}) \ge -\alpha(h(\mathbf{x}))\), the
enforcer maintains the forward invariance of the safe set at every
control cycle. The nervous system does not need to forecast the
trajectory ten steps into the future to guarantee that the machine
remains within admissible physical bounds right now.

To enforce invariance under physical fault conditions, the permission
path must operate with complete architectural autonomy at the hardware
floor. If the enforcer relies on the high-level policy computer for
execution or heartbeats, any operating system stall, memory exhaustion,
or thermal throttling event compromises mechanical safety. The nervous
system isolates enforcement onto dedicated real-time microcontrollers or
programmable logic running synchronously with the motor inverters at
\(1000\text{ Hz}\) or higher. This bare-metal layer treats the upstream
learned model as an untrusted client. When transport monitors detect
dropped packets, watchdog timers expire, or state estimates diverge
beyond residual thresholds, the local permission path revokes the
policy's authority and asserts deterministic braking profiles
autonomously. Safety is maintained not by hoping that the upstream
software recovers, but by designing the nervous system to govern
physical forces when upstream software crashes.

\protect\phantomsection\label{callout-takeawaysux2a-12.11}
\begin{fbxSimple}{callout-takeaways}{Key Chapter Takeaways}{}
\phantomsection\label{callout-takeaways*-12.11}

\begin{itemize}
\tightlist
\item
  \textbf{Tracking and Refusal Form a Single Mathematical Contract:} The
  brain proposes actions (\(p_t\)), but the deterministic nervous system
  holds exclusive hardware authority over motor gate drivers. Safety
  enforcement requires an inner tracking loop that can reliably convert
  permitted commands into physical counter-torque within bounded
  tracking error (\(\epsilon_{\text{track}}\)).
\item
  \textbf{Composite Clearance Bounds Must Compound All Inelastic
  Delays:} Physical stopping clearance is never just nominal stopping
  distance (\(v^2/2a_{\max}\)). The certified margin must strictly sum
  tracking error bounds (\(\epsilon_{\text{track}}\)), transport bus
  latency (\(v \cdot t_{\text{delay}}\)), kinematic deceleration
  (\(v^2/2a_{\max}\)), and unobservable state estimation residual
  (\(\epsilon_{\text{est}}\)).
\item
  \textbf{Forward Invariance Outperforms Multi-Step Trajectory
  Prediction:} Receding-horizon trajectory optimizers are vulnerable to
  solver timeouts and model divergence. Control Barrier Functions
  evaluate instantaneous set-theoretic constraints
  (\(\dot{h}(\mathbf{x}, \mathbf{u}) \ge -\alpha h(\mathbf{x})\)) at
  \(1000\text{ Hz}\), ensuring safety at the current microsecond step
  regardless of downstream replanning stalls.
\item
  \textbf{The Deterministic 4-Tier Fallback Ladder:} When perturbations
  strike, the system escalates deterministically: Level 1 (QP orthogonal
  projection), Level 2 (autonomous MCU deceleration spline), Level 3
  (dynamic shunt braking via low-side MOSFET shorting), and Level 4
  (galvanic power cutoff via discrete safety relays).
\item
  \textbf{Hardware Watchdogs Fail Closed on Upstream Policy Crashes:}
  Application processors executing learned policies are
  non-deterministic and crash-prone. When Linux kernel panics or mailbox
  heartbeats cease (\(t_{\text{timeout}} \ge 3.0\text{ ms}\)), the
  isolated bare-metal MCU autonomously arrests the plant without waiting
  for upstream recovery.
\end{itemize}

\end{fbxSimple}

The architecture of a physical AI system must reconcile two distinct
computational entities that share a single machine. The high-level brain
explores, plans, and optimizes across open-ended operational domains
without safety guarantees, while the low-level nervous system tracks
verified commands, evaluates barrier certificates, and retains the
absolute authority to say no. These two workloads impose fundamentally
opposing demands on hardware. A learned policy requires parallel tensor
compute with massive memory bandwidth but operates under loose deadlines
where occasional latency jitter is acceptable. Conversely, barrier
evaluation and phase-current regulation demand minimal arithmetic
throughput but impose hard, sub-millisecond real-time deadlines with
zero tolerance for tail latency. Attempting to co-locate these divergent
workloads on a single monolithic processor compromises both determinism
and throughput. Chapter~\ref{sec-placement} turns to this physical
partitioning, examining how the computational and bandwidth demands of
perception, policy evaluation, and barrier enforcement determine the
placement of silicon across the edge, the nervous system, and the
central brain.

\protect\phantomsection\label{callout-labux2a-12.12}
\begin{fbxSimple}{callout-lab}{Hands-On Lab:}{The 1 kHz MCU Safety Enforcer}
\phantomsection\label{callout-lab*-12.12}

Test the boundary between unverified learned proposals and deterministic
physical authority on the dual-brain kit:

\begin{enumerate}
\def\labelenumi{\arabic{enumi}.}
\tightlist
\item
  \textbf{Bare-Metal CBF-QP Implementation:} Deploy an active-set
  Control Barrier Function safety filter on the bare-metal MCU running
  at \(1000\text{ Hz}\) in zero-allocation static SRAM.
\item
  \textbf{Proposal Filtering on the Bench:} Stream unconstrained,
  out-of-envelope trajectory proposals from the Linux MPU over shared
  SRAM via RPMSG mailboxes and verify that the MCU orthogonally projects
  commands onto the safe set boundary \(\partial \mathcal{C}\) without
  distorting net Cartesian direction.
\item
  \textbf{Seeded Crash Injection:} Trigger a synthetic Linux kernel
  panic during a high-speed motion sweep. Measure the time required for
  the MCU hardware watchdog to detect heartbeat cessation
  (\(t_{\text{timeout}} \le 3.0\text{ ms}\)) and verify autonomous Level
  2 deceleration to position hold before carriage travel reaches
  physical end-stops.
\end{enumerate}

\emph{Hardware bench guide and starter firmware:} Lab 08: MCU Enforcer.

\end{fbxSimple}

\chapter{\texorpdfstring{Heterogeneous Silicon
Placement\chaptersubtitle{What Happens When Everything Shares One Chip}}{Heterogeneous Silicon Placement}}\label{sec-placement}

\begin{figure}[H]

\centering{

\includegraphics[width=1\linewidth,height=\textheight,keepaspectratio]{index_files/mediabag//Users/VJ/GitHub/MLSysBook-integrate-physical/physical/book/chapters/13-placement/figures/fig_locator.pdf}

}

\caption{\label{fig-locator}\textbf{Systems Organ Focus: Nervous.} The
physical co-location boundary where learned policy deliberation and
deterministic safety enforcement converge on shared silicon. On a
heterogeneous system-on-chip, high-throughput cognitive bursts and hard
real-time reflex loops compete for memory controllers, on-chip
interconnect crossbars, power distribution rails, and thermal mass.
Attribution: Book architecture diagram, CC BY 4.0.}

\end{figure}%

\emph{The boundaries were drawn on a diagram. Do they exist on the die?}

In abstract system diagrams, neural perception, trajectory planning, and
safety enforcement appear as clean, independent blocks connected by
ideal mathematical arrows. In a physical machine, all of these workloads
are forced onto shared system-on-chip silicon to meet strict size,
weight, power, and cost constraints. When a dense neural accelerator
bursts across a unified memory bus, it can saturate DRAM controllers and
delay a safety-critical encoder interrupt by several milliseconds.
Heterogeneous compute placement is the engineering discipline of
allocating cognitive and reflex workloads across CPUs, NPUs, GPUs, and
real-time microcontrollers while using hardware-enforced AXI QoS
arbitration, power rail decoupling, and thermal management to prove that
real-time safety guarantees survive execution on shared silicon.

\newpage{}

\protect\phantomsection\label{callout-objectiveux2a-13.1}
\begin{fbxSimple}{callout-objective}{Learning Objectives}{}
\phantomsection\label{callout-objective*-13.1}

After studying this chapter, you will be able to:

\begin{enumerate}
\def\labelenumi{\arabic{enumi}.}
\tightlist
\item
  Name every claimant on each shared resource in a given heterogeneous
  design.
\item
  Measure the tail latency distribution of the enforcement path under
  synthetic contention on the proposal path.
\item
  State what experimental result would falsify an independence claim,
  then construct the test harness to run it.
\item
  Predict how a power step or thermal excursion in the accelerator
  domain alters the execution timing and clock stability of the safety
  core.
\item
  Evaluate the four remediation strategies---static reservation, spatial
  partitioning, workload reduction, and separate silicon---against the
  physical dynamics of the plant.
\end{enumerate}

\emph{Core chapter takeaway.} Independence claimed on a block diagram is
not independence measured under load. Two paths on one die share a bus,
a rail, and a thermal mass, and physical separation must be demonstrated
through worst-case empirical measurement rather than asserted through
logical privilege.

\end{fbxSimple}

\section{The Silicon Substrate}\label{sec-placement-13-1}

\lettrine{C}{onsider an industrial bypass}\lettrine{v}{alve} regulating
high-pressure fuel flow into a gas turbine combustion chamber. Upstream
sensors monitor manifold pressure, mass flow rate, and fuel temperature.
A learned policy running on an integrated neural accelerator evaluates
these telemetry streams every \(20\text{ ms}\), proposing continuous
micro-adjustments to the valve actuator to optimize combustion
efficiency. Parallel to this proposal pipeline, a deterministic safety
enforcer executes on a separate real-time processor core, evaluating
flow rate limits and chamber pressure thresholds every
\(1.0\text{ ms}\). If chamber pressure exceeds a critical limit, the
enforcer must immediately override the policy, command the valve servo
to clamp shut within \(2.0\text{ ms}\), and prevent an overpressure
blowout. In the four-tier architectural locator shown in
Figure~\ref{fig-locator}, the nervous system serves as the definitive
temporal barrier between the stochastic, high-capacity cognition of the
brain and the irreversible dynamics of the physical body. On an
architectural block diagram, the proposal path and the enforcement path
appear as two completely isolated functional blocks connected only by a
logical interface. On the physical system-on-chip installed inside the
valve actuator housing---such as a canonical heterogeneous dual-brain
System-on-Chip (SoC)---both compute engines reside on the same square
centimeter of silicon, sharing memory buses, power rails, clock trees,
and peripheral controllers.

This shared-silicon reality governs every domain of physical embodiment
across all Three Physical Machine Archetypes:

\begin{itemize}
\tightlist
\item
  In \emph{Mobile systems} (Class 1), an autonomous chassis or mobile
  delivery robot co-locates multi-camera vision-language-action (VLA)
  navigation models on an embedded GPU with \(1000\text{ Hz}\)
  steer-by-wire and emergency braking reflexes on an adjacent real-time
  core. A memory bus stall during feature extraction can delay an
  obstacle-avoidance brake command, allowing moving kinetic mass
  (\(E_k = \frac{1}{2}mv^2\)) to breach safety stopping margins
  (\(d_{\text{stop}}\)).
\item
  In \emph{Manipulators} (Class 2), an articulated robot arm executes
  high-dimensional diffusion grasping policies
  (\(20\text{--}50\text{ Hz}\)) alongside high-gain Cartesian impedance
  and joint torque loops (\(1000\text{ Hz}\)). If an accelerator tensor
  burst evicts safety state tables from the shared last-level cache, the
  resulting microsecond jitter degrades feedback phase margins
  (\(\Delta \phi = -\frac{\omega T_s}{2}\)), inducing acoustic chatter,
  motor torque ripple, or fixture collisions.
\item
  In \emph{Process and energy systems} (Class 3), industrial chemical
  bypass valves and high-voltage grid battery inverters co-locate
  continuous thermodynamic optimization with hard real-time overpressure
  and power-stage shoot-through interlocks. A delayed safety interrupt
  allows manifold pressure to spike beyond structural burst limits or
  power transistors to enter catastrophic thermal runaway.
\end{itemize}

In traditional software systems operating ``behind glass,'' computing is
idempotent and isolation is easily abstracted through virtual memory,
kernel time-slicing, and software containers. A delayed database query
or an unhandled exception causes temporary software stalls without
altering physical reality. In Physical AI, however, computation exerts
physical force, moves physical mass, and dissipates electrical energy.
Safety arguments constructed in Chapter~\ref{sec-nervous} and
Chapter~\ref{sec-enforcement} depend on maintaining strict physical
separation between the untrusted cognitive authority that proposes
action and the deterministic reflex authority that permits it. On a
shared heterogeneous die, this separation requires \textbf{operational
independence}, which can be defined through two concrete physical
requirements. First, no computational activity, memory transfer, or
burst traffic originating from the proposal path may cause the execution
latency of the enforcement path to exceed its bounded deadline. Second,
no failure mode, software crash, or privilege escalation within the
proposal path may inhibit the enforcer from evaluating safety invariants
or asserting a safe default state. If an unconstrained tensor
calculation on the Linux Microprocessor Unit (MPU) delays a safety
evaluation on the Real-Time Microcontroller Unit (MCU) long enough for
pressure in the valve manifold to spike past its structural limit, the
logical isolation depicted on the architecture diagram has failed in
hardware.

To understand how logical isolation breaks down on silicon, consider an
analytical example of this flow-regulating system implemented on a
single system-on-chip. The deterministic enforcer runs at
\(1000\text{ Hz}\), initiating a state evaluation cycle every
\(1.0\text{ ms}\). In each cycle, the enforcer reads \(64\text{ KB}\) of
calibrated sensor history from main memory, verifies boundary
conditions, and issues an actuator command, requiring
\(120\text{ }\mu\text{s}\) of baseline execution time when memory access
is uncontended. The chip uses a shared 64-bit LPDDR4 memory subsystem
running at a data rate of \(1600\text{ MT/s}\), delivering a peak
theoretical bandwidth of \(12.8\text{ GB/s}\). In isolation, streaming
the enforcer's \(64\text{ KB}\) sensor buffer over this bus requires
\(64\text{ KB} / (12.8\text{ GB/s}) = 5.0\text{ }\mu\text{s}\).

Contention arises when the learned policy executes on the adjacent
neural core. Running an image-based flow estimation model at
\(50\text{ Hz}\), the neural accelerator streams weights and activation
maps totaling \(180\text{ MB}\) per inference pass across a sustained
\(15\text{ ms}\) compute window. This inference burst demands an average
bandwidth of \((180\text{ MB}) / (15\text{ ms}) = 12.0\text{ GB/s}\),
consuming \(93.75\%\) of the total available DRAM throughput. When the
memory controller arbitrates between the two masters, the enforcer's
memory requests wait in controller queues behind long bursts of matrix
multiplication reads. If the effective bandwidth allocated to the
enforcer falls to the residual \(0.8\text{ GB/s}\), the sensor read
transfer alone expands from \(5.0\text{ }\mu\text{s}\) to
\(80.0\text{ }\mu\text{s}\). When combined with last-level cache line
evictions, crossbar arbitration stalls, and interrupt handling jitter,
the enforcer's observed tail latency at \(P_{99.9}\) reaches
\(1055\text{ }\mu\text{s}\). The enforcer misses its \(1.0\text{ ms}\)
physical deadline, not because its own algorithmic workload changed, but
because the proposal path occupied the shared memory controller.

Main memory bandwidth is merely one member of an extensive list of
shared on-chip resources. A full accounting across the die reveals
physical couplings at every layer of the hardware platform. The proposal
engine and the safety enforcer share the last-level cache, where large
activation matrices evict the enforcer's pinned lookup tables and sensor
history buffers. They share the on-chip crossbar interconnect, where
arbitration policies often favor burst throughput over bounded latency.
They share the power distribution network, where the sudden activation
of billions of switching transistors inside the neural engine produces
an inductive and resistive voltage drop across the internal power rails
(\(\Delta V = L \cdot dI/dt + IR\)), risking clock instability and
timing closure failures in adjacent real-time cores. They share the
phase-locked loops that generate clock signals, the hardware reset
lines, and the common thermal package. If the neural accelerator
dissipates \(15\text{ W}\) of continuous power, the junction temperature
across the silicon substrate can rise from \(45^\circ\text{C}\) to
\(95^\circ\text{C}\), prompting thermal management logic to cut
processor clock frequencies in half and double execution times across
every core on the die.

Because these interactions are governed by physical coupling rather than
software abstractions, architectural claims of independence must be
treated as hypotheses to be falsified through rigorous measurement.
Verifying a system requires four distinct steps. First, an engineer
builds a claimant inventory, identifying every physical resource shared
between proposal and enforcement domains, from memory channels and
interconnect queues down to voltage regulators and thermal interfaces.
Second, the system undergoes loaded measurement, subjecting the proposal
path to synthetic worst-case compute, memory, and bus traffic while
logging the full tail latency distribution of the enforcement path at
\(P_{99}\) and \(P_{99.9}\). Third, the system undergoes warm
measurement, operating under peak ambient temperatures and sustained
thermal dissipation to measure clock throttling and voltage droop under
load. Fourth, the engineer records a formal placement decision,
verifying that the worst-case response time of the safety loop remains
strictly within physical requirements under all possible operating
states.

When empirical testing reveals that shared hardware invalidates safety
margins, modifying software task priorities or adjusting thread
scheduling cannot resolve the underlying physical contention. The
coupling exists in the silicon wiring, the queue depths of the crossbar
arbiters, and the thermal resistance of the integrated circuit package.
Determining which compute engines can safely reside on the same die, how
shared interconnects must be partitioned, and where physical isolation
boundaries must be enforced represents the fundamental engineering
problem of physical placement.

\section{Two Paths on One Die}\label{sec-placement-13-2}

Placement cannot be deferred to integration, because the physical
assignment of software functions to hardware execution engines
determines whether separation exists at all. Consider a valve regulating
high-pressure fluid flow in a chemical processing line, where a learned
model adjusts the valve trim position to track a variable target flow
rate while an enforcement filter prevents downstream pressure transients
that would rupture the pipe. The machine contains two paths with
conflicting requirements. The proposal path ingests downstream pressure
histories, acoustic vibration signals, and temperature logs over a
moving window of one hundred samples, running a learned neural policy at
\(20\text{ Hz}\) to calculate a recommended valve displacement. It
operates under a soft deadline of \(45\text{ ms}\), holds no direct
hardware authority over the motor driver, and runs in an unprivileged
user space. The enforcement path samples raw differential pressure
transducers at \(1000\text{ Hz}\), evaluates physical boundary
constraints on pressure slew rates and cavitation limits, and generates
electrical drive commands for the valve servo. It executes under a hard
physical deadline of \(1.0\text{ ms}\), requires an execution budget
bounded at \(200\text{ }\mu\text{s}\) to leave sufficient margin for
mechanical valve transit, and runs in a privileged real-time execution
domain. Only the enforcement path holds memory-mapped input-output
access to the pulse-width modulation registers that command the valve
motor.

\protect\phantomsection\label{callout-definitionux2a-13.2}
\begin{fbxSimple}{callout-definition}{Definition:}{Heterogeneous Silicon Partitioning}
\phantomsection\label{callout-definition*-13.2}

The \textbf{heterogeneous silicon partitioning} pattern is the
architectural allocation of stochastic, high-capacity cognitive
workloads (perception, learned policy evaluation) and deterministic
reflex loops (safety invariant checks, actuator commutation) across
specialized execution domains---typically an unprivileged Application
Processor (MPU) and a privileged Real-Time Microcontroller
(MCU)---isolated via hardware memory protection units (MPUs),
tightly-coupled static RAM (TCM), and dedicated peripheral controllers.

\end{fbxSimple}

To assign concrete performance numbers, these functions must be mapped
onto physical silicon domains. In a heterogeneous dual-brain
system-on-chip---such as the canonical dual-brain architecture
introduced in Chapter~\ref{sec-nervous}---proposal inference maps onto
an unprivileged Linux Application Processor (MPU) or dedicated neural
accelerator. State exchange between the paths maps onto a shared static
memory mailbox or ring buffer in dynamic random-access memory. Safety
enforcement and sensor validation map onto a dedicated Real-Time
Controller Subsystem (MCU) running bare-metal firmware or a
deterministic real-time operating system. Device input-output maps onto
dedicated serial peripheral interface controllers, analog-to-digital
converters, and timer channels directly wired to the MCU. Telemetry
logging, network communications, and fault recording map onto the Linux
MPU. As detailed in Table~\ref{tbl-13-compute-allocation}, each compute
domain on a modern physical AI platform exhibits distinct trade-offs
across nominal clock rate, tail latency determinism, operating system
isolation, and memory bus coupling.

\begin{longtable}[]{@{}
  >{\raggedright\arraybackslash}p{(\linewidth - 10\tabcolsep) * \real{0.1667}}
  >{\raggedright\arraybackslash}p{(\linewidth - 10\tabcolsep) * \real{0.1667}}
  >{\raggedright\arraybackslash}p{(\linewidth - 10\tabcolsep) * \real{0.1667}}
  >{\raggedright\arraybackslash}p{(\linewidth - 10\tabcolsep) * \real{0.1667}}
  >{\raggedright\arraybackslash}p{(\linewidth - 10\tabcolsep) * \real{0.1667}}
  >{\raggedright\arraybackslash}p{(\linewidth - 10\tabcolsep) * \real{0.1667}}@{}}
\caption{\textbf{Heterogeneous Compute Allocation and Execution Domain
Characteristics.} Comparative analysis of compute domains on modern
physical AI platforms. Allocation decisions dictate latency determinism,
failure isolation, and memory contention, establishing why
safety-critical reflex loops cannot execute on general-purpose
application cores or GPUs without hardware-enforced
isolation.}\label{tbl-13-compute-allocation}\tabularnewline
\toprule\noalign{}
\begin{minipage}[b]{\linewidth}\raggedright
Compute Engine Domain
\end{minipage} & \begin{minipage}[b]{\linewidth}\raggedright
Nominal Clock \& Power
\end{minipage} & \begin{minipage}[b]{\linewidth}\raggedright
Latency Determinism (\(P_{99.99}\))
\end{minipage} & \begin{minipage}[b]{\linewidth}\raggedright
OS \& Privilege Isolation
\end{minipage} & \begin{minipage}[b]{\linewidth}\raggedright
Memory \& Interconnect Coupling
\end{minipage} & \begin{minipage}[b]{\linewidth}\raggedright
Target Workload Assignment
\end{minipage} \\
\midrule\noalign{}
\endfirsthead
\toprule\noalign{}
\begin{minipage}[b]{\linewidth}\raggedright
Compute Engine Domain
\end{minipage} & \begin{minipage}[b]{\linewidth}\raggedright
Nominal Clock \& Power
\end{minipage} & \begin{minipage}[b]{\linewidth}\raggedright
Latency Determinism (\(P_{99.99}\))
\end{minipage} & \begin{minipage}[b]{\linewidth}\raggedright
OS \& Privilege Isolation
\end{minipage} & \begin{minipage}[b]{\linewidth}\raggedright
Memory \& Interconnect Coupling
\end{minipage} & \begin{minipage}[b]{\linewidth}\raggedright
Target Workload Assignment
\end{minipage} \\
\midrule\noalign{}
\endhead
\bottomrule\noalign{}
\endlastfoot
\textbf{Application Processor Subsystem (MPU)} &
\(1.0\text{--}2.2\text{ GHz}\), \(2\text{--}10\text{ W}\) &
Non-deterministic (Jitter \(> 20\text{ ms}\); MMU paging, kernel
preemption) & Rich OS (Linux / POSIX), Ring 3 user space vs.~Ring 0
kernel & Shared coherent L3 cache, virtual memory via MMU, unified DRAM
& Telemetry logging, edge VLMs/VLAs, PyTorch, ROS2 policy evaluation
(\(10\text{--}50\text{ Hz}\)) \\
\textbf{GPU Tensor Core Engine} & \(0.8\text{--}1.3\text{ GHz}\),
\(10\text{--}40\text{ W}\) & Semi-deterministic (Jitter
\(\approx 200\text{--}800\text{ }\mu\text{s}\); warp scheduling) & Host
driver queue, user-space runtime, unified virtual memory &
High-throughput LPDDR5/HBM, wide AXI bursts (\(256\text{ B}\)), high
DRAM load & High-throughput perception, multi-camera feature extraction,
dense VLM inference \\
\textbf{Dedicated Neural Processing Unit} (NPU / Systolic Array) &
\(0.6\text{--}1.0\text{ GHz}\), \(2\text{--}8\text{ W}\) & High
determinism (Jitter \(< 50\text{ }\mu\text{s}\); static graph execution)
& Hardware mailbox, fixed-depth systolic pipeline, pinned memory buffers
& Dedicated on-chip SRAM (\(8\text{--}32\text{ MB}\)), periodic DMA
bursts to DRAM & Learned policy evaluation, low-latency trajectory chunk
generation (\(20\text{--}100\text{ Hz}\)) \\
\textbf{Real-Time Controller Subsystem (MCU)} &
\(200\text{--}800\text{ MHz}\), \(0.1\text{--}1.5\text{ W}\) & Strictly
bounded (Jitter \(< 5.0\text{ }\mu\text{s}\); deterministic in-order
core) & RTOS (FreeRTOS / Zephyr) or Bare-Metal, MPU memory regions &
Tightly-Coupled Memory (TCM / SRAM), private peripheral bus, zero paging
& Deterministic safety invariants, \(1\text{ kHz}\) CBF safety filters,
direct PWM timers \\
\textbf{Discrete Safety Microcontroller} &
\(200\text{--}400\text{ MHz}\), \(0.2\text{--}1\text{ W}\) & Fully
invariant (Jitter \(< 50\text{ ns}\); cycle-accurate lockstep core) &
Bare-metal / AUTOSAR Classic, hardware Fault Collection Unit (FCCU) &
Isolated SRAM/Flash, independent external power rails, dedicated crystal
& Emergency mechanical clamping, hardware interlocks, supervisory
watchdog trips \\
\end{longtable}

In an analytical evaluation of the unloaded system, the computational
budgets appear generous. In the enforcement path on the real-time MCU,
reading four pressure sensors over a high-speed peripheral bus requires
\(40\text{ }\mu\text{s}\), evaluating the constraint envelope and
computing the actuator signal takes \(80\text{ }\mu\text{s}\), and
writing the resulting register values to the motor controller timer
takes \(15\text{ }\mu\text{s}\). The total unloaded execution demand is
\(135\text{ }\mu\text{s}\). Against the \(400\text{ }\mu\text{s}\)
computational deadline mandated by the \(1.0\text{ ms}\) physical cycle
(which reserves \(600\text{ }\mu\text{s}\) for servo motor settling),
the enforcement path exhibits an unloaded slack of
\(265\text{ }\mu\text{s}\). In the proposal path on the Linux MPU, a
policy requiring \(3.2\text{ GFLOPs}\) per step executes on an
accelerator capable of \(2.0\text{ TFLOPS}\) sustained throughput,
yielding an inference time of \(1.6\text{ ms}\). Adding
\(0.8\text{ ms}\) for tensor formatting and \(0.1\text{ ms}\) to publish
the proposal to the shared ring buffer, the proposal path requires
\(2.5\text{ ms}\) of computation per \(50\text{ ms}\) cycle, leaving
\(47.5\text{ ms}\) of unloaded slack. Evaluated in isolation, both loops
operate with comfortable margins.

Assigning these tasks to separate processor cores does not establish
independent failure domains if the underlying silicon infrastructure
remains shared.\sidenote{\footnotesize FAA AC 20-193 / CAST-32A mandates multicore
  interference analysis, requiring bounded WCET proofs under contention
  across shared caches, buses, and memory controllers.} While the
real-time core executes instructions independently from the application
processor, both cores rely on common hardware services:

\begin{itemize}
\tightlist
\item
  \emph{Shared interrupt routing:} The Generic Interrupt Controller
  (GIC) routes peripheral and inter-processor notifications across a
  shared fabric, where a flood of Ethernet packet interrupts on the
  Linux MPU can delay the timer interrupt that triggers the
  \(1000\text{ Hz}\) MCU safety loop.
\item
  \emph{Shared memory controllers:} Both paths access the same dynamic
  memory controller and interconnect crossbar, where large matrix tile
  loads generated by the neural accelerator saturate DRAM command
  queues, turning an expected \(20\text{ ns}\) level-two cache refill on
  the real-time core into a \(350\text{ ns}\) round-trip to main memory.
\item
  \emph{Shared clock trees:} The cores share phase-locked loops (PLLs)
  for clock generation, where dynamic voltage and frequency scaling
  (DVFS) transients on the application processor introduce clock jitter
  across adjacent real-time domains.
\item
  \emph{Shared power rails:} A single Power Distribution Network (PDN)
  delivers current to both domains, where switching surges from the
  neural accelerator drop the internal voltage rail, risking logic state
  corruption or brownout resets.
\item
  \emph{Shared reset domains:} A kernel panic on the Linux MPU that
  triggers a hardware watchdog reset can pull down the shared peripheral
  bus or system interconnect, disabling pulse-width modulation outputs
  and leaving the valve unpowered.
\end{itemize}

Evaluating candidate placements reveals how physical coupling governs
system feasibility. In a unified placement where proposal, enforcement,
and logging execute as threads under a single symmetric multiprocessing
operating system on shared application cores, communication between
paths takes less than \(1\text{ }\mu\text{s}\) via shared pointer
passing. However, cache eviction and kernel preemption cause the tail
latency of the enforcement path to reach \(2.8\text{ ms}\) at
\(P_{99.9}\), repeatedly missing the \(1.0\text{ ms}\) hard deadline and
risking pipe overpressure. In a heterogeneous single-die placement
(e.g., dual-brain SoC with Linux MPU + Real-Time MCU on one chip),
proposal runs on the MPU/NPU and enforcement runs on an isolated MCU
core; communication across on-chip shared static random-access memory
requires \(12\text{ }\mu\text{s}\). Under synthetic worst-case memory
contention, enforcement latency rises to \(380\text{ }\mu\text{s}\),
which still fits within the \(400\text{ }\mu\text{s}\) compute budget.
However, when the accelerator dissipates sustained peak power, local
thermal elevation triggers automatic clock frequency halving, stretching
enforcement execution to \(540\text{ }\mu\text{s}\) and violating the
physical deadline. In a dual-die placement where enforcement resides on
a physically isolated microcontroller with its own power regulation,
crystal oscillator, and memory, communication across an external
serialized peripheral interface costs \(150\text{ }\mu\text{s}\).
Because the failure domains are physically decoupled, the enforcement
execution time remains invariant at \(135\text{ }\mu\text{s}\),
guaranteeing that if the proposal path stalls, crashes, or overheats,
the enforcement core detects the missing heartbeat and transitions the
flow valve to a safe mechanical relief position within a single
\(1.0\text{ ms}\) cycle.

\section{Contention for Shared Resources}\label{sec-placement-13-3}

Every shared resource has more than one claimant, and a number that does
not say what it competes with is not yet a systems number. On a
system-on-chip regulating a fluid pipeline, physical wires do not
respect functional boundaries drawn on a software architecture diagram.
The proposal path, the enforcement path, and the system infrastructure
share the same silicon, competing for the same memory controllers,
crossbar switches, and power distribution rails. An inventory of digital
claimants on the main memory bus includes sensor direct memory access
streams that deposit raw pressure transducer readings and acoustic flow
telemetry into dynamic random-access memory, the neural accelerator
fetching model weight matrices and spilling intermediate activation
tensors, the real-time core reading pipeline boundary invariants and
historic sensor ring buffers, actuator write buffers pushing pulse-width
modulation commands to the flow valve driver, cache controllers
executing dirty line writebacks, and diagnostic logging processes
writing state records to nonvolatile storage buffers.

In physical AI hardware design, \emph{crossbar bus and DRAM bank
contention} is the physical queuing delay and priority inversion induced
when high-throughput memory or direct memory access (DMA) bursts from
unconstrained compute masters saturate shared on-chip interconnect
crossbars and DRAM controllers, inflating the worst-case read and write
latencies of high-priority real-time safety transactions due to
non-preemptive packet burst arbitration and DRAM row-buffer bank
conflicts.

Each claimant imposes a demand that can be derived from its transfer
size per event and its event rate. In an analytical example of a
flow-regulation controller sharing a single \(64\text{-bit}\) LPDDR4
memory channel rated at \(12.8\text{ GB/s}\) peak bandwidth, sensor
direct memory access transfers a \(64\text{ kB}\) buffer at
\(1000\text{ Hz}\), generating an average demand of
\(64.0\text{ MB/s}\). The learned proposal model executes an inference
pass every \(20\text{ ms}\) (\(50\text{ Hz}\)), streaming
\(60\text{ MB}\) of parameters and activation data for an average demand
of \(3.0\text{ GB/s}\). The enforcement path executes every
\(1.0\text{ ms}\) (\(1000\text{ Hz}\)), reading \(4\text{ kB}\) of
pipeline state and historical sensor logs, consuming \(4.0\text{ MB/s}\)
on average. Actuator writes require \(1\text{ kB}\) per millisecond
(\(1.0\text{ MB/s}\)), cache maintenance evictions contribute
\(128.0\text{ MB/s}\) in writebacks, and diagnostic telemetry logging
transfers \(32\text{ kB}\) records at \(100\text{ Hz}\)
(\(3.2\text{ MB/s}\)). Summing these rates gives an aggregate demand of
\(3.2002\text{ GB/s}\), which represents \(25.0\%\) of the channel
capacity. An engineer who inspects only average utilization would
conclude that three-quarters of the bus capacity remains available. But
memory arbiters do not schedule average rates, they arbitrate discrete
bursts, and two workloads with identical average throughput create
distinct interference profiles if one streams continuously in small
packets while the other arrives in long, saturating blocks.

Arithmetic intensity, defined as the ratio of arithmetic operations to
memory bytes transferred, indicates whether a computational routine is
bounded by execution units or by memory bandwidth during steady-state
execution. For a convolution or dense matrix multiplication on the
proposal path, an arithmetic intensity of \(120\text{ FLOP/byte}\)
indicates that arithmetic units will be fully occupied once data resides
in local cache or static memory. However, during the initial phase where
weights are transferred from dynamic memory, the operational arithmetic
intensity drops to zero, converting the accelerator into a high-priority
direct memory access engine that saturates the memory bus. Arithmetic
intensity reveals whether a workload can hide memory transfer latency
behind arithmetic work in isolation, but it provides no information
about queue contention, transfer priority, burst duration, or which
request waits when two paths access the bus simultaneously.

The latency experienced by the enforcement path depends directly on the
longest uninterrupted burst permitted on the shared interconnect and the
internal banking dynamics of the DRAM subsystem
(Figure~\ref{fig-13-soc-contention-droop} a). In this analytical system,
the neural accelerator transfers weights in \(256\text{ kB}\) tiles.
Across a \(12.8\text{ GB/s}\) interface, an uninterrupted
\(256\text{ kB}\) burst occupies the physical bus for
\(20.0\text{ }\mu\text{s}\). When the on-chip interconnect uses
packet-atomic non-preemptive arbitration (even with AMBA AXI
Quality-of-Service \texttt{AxQOS} priority tags), an enforcement request
that arrives immediately after a tile transfer begins must wait the full
\(20.0\text{ }\mu\text{s}\) transfer duration. Contention is further
compounded inside the dynamic memory controller. Dynamic random-access
memory is physically organized into independent banks, where each bank
contains a two-dimensional array of storage cells managed by a single
row buffer. When an access targets a row already loaded into the buffer
(a \emph{row-buffer hit}), data is transferred in column access time
(\(t_{\text{CAS}} \approx 14\text{ ns}\)). If an access targets a
different row in the same bank (a \emph{bank conflict} or
\emph{row-buffer miss}), the controller must first restore and close the
open row by precharging the bank
(\(t_{\text{RP}} \approx 14\text{ ns}\)), activate the newly requested
row (\(t_{\text{RCD}} \approx 14\text{ ns}\)), and execute the column
read (\(t_{\text{CAS}} \approx 14\text{ ns}\)), tripling baseline access
latency.

Crucially, modern DRAM memory controllers employ First-Ready
First-Come-First-Served (FR-FCFS) scheduling algorithms designed to
maximize aggregate bus utilization by prioritizing row-buffer hits over
raw transaction arrival order. When the neural engine streams a
continuous matrix tile, it holds open active row buffers across multiple
DRAM banks. If a high-priority real-time safety read targeting a closed
row arrives from the MCU, the DRAM controller repeatedly reorders queued
transactions to prioritize remaining hits from the neural stream. If the
controller queue contains a pending \(64\text{ kB}\) sensor direct
memory access burst (\(5.0\text{ }\mu\text{s}\)) and a \(32\text{ kB}\)
logging burst (\(2.5\text{ }\mu\text{s}\)), plus multiple bank conflict
precharge cycles, total queuing delay reaches
\(27.6\text{ }\mu\text{s}\). If the enforcement routine has a total
compute allocation of \(100\text{ }\mu\text{s}\) within a
\(1.0\text{ ms}\) physical deadline and was provisioned with
\(20.0\text{ }\mu\text{s}\) of timing slack, this
\(27.6\text{ }\mu\text{s}\) memory stall consumes the entire margin and
delays the valve closing command by \(7.6\text{ }\mu\text{s}\), causing
an actuation deadline violation.

\protect\phantomsection\label{callout-ruleux2a-13.3}
\begin{fbxSimple}{callout-rule}{Architect’s Rule of Thumb:}{The Memory Contention Invariant}
\phantomsection\label{callout-rule*-13.3}

Average memory bus bandwidth is a vanity metric; real-time safety is
dictated entirely by worst-case burst length, DRAM bank conflict
precharge latencies, and non-preemptive crossbar queue depths.

\end{fbxSimple}

Contention extends to the power distribution network, where claimants
compete for charge rather than bandwidth
(Figure~\ref{fig-13-soc-contention-droop} b). When the neural
accelerator starts a matrix multiply across thousands of parallel
multiply-accumulate units, its current draw steps by
\(\Delta I = 6.0\text{ A}\) with a transition time of
\(\Delta t = 2.0\text{ ns}\), producing a current slew rate of
\(dI/dt = 3.0\text{ A/ns}\). Across a shared power distribution network
with an effective loop inductance of \(L = 30\text{ pH}\) and an
equivalent series resistance of \(R = 8.0\text{ m}\Omega\), the
transient voltage disturbance across the rail is derived from the
fundamental circuit relationship:
\[\Delta V = L \cdot \frac{dI}{dt} + I \cdot R\] The inductive drop
contributes
\(L(dI/dt) = (30 \times 10^{-12}\text{ H}) \times (3.0 \times 10^9\text{ A/s}) = 90\text{ mV}\),
while the resistive drop contributes
\(I \cdot R = (6.0\text{ A}) \times (8.0 \times 10^{-3}\ \Omega) = 48\text{ mV}\),
producing a total rail droop of \(\Delta V = 138\text{ mV}\).

\emph{PDN voltage droop coupling} is the instantaneous electrical rail
collapse (\(\Delta V = L \cdot \frac{dI}{dt} + I \cdot R\)) induced
across shared on-die power distribution networks when fast, large
current steps from parallel neural execution units deplete local
decoupling capacitance, stretching logic gate propagation delays
(\(t_{\text{pd}}\)) and triggering clock frequency degradation, timing
closure violations, or brownout resets across co-located real-time
cores.

On a nominal \(0.85\text{ V}\) core supply, this \(138\text{ mV}\) drop
represents a \(16.2\%\) voltage collapse, falling well below the
\(0.78\text{ V}\) minimum operating voltage (\(V_{\min}\)) of the
co-located real-time core. The off-chip Voltage Regulator Module (VRM)
operates with a control loop bandwidth below \(1\text{ MHz}\),
responding within \(2.0\text{ }\mu\text{s}\)---three orders of magnitude
too slow to counteract a nanosecond transient. First-stage charge
delivery relies entirely on on-die deep-trench decoupling capacitors
(DTCs). If local capacitance is depleted before the current transient
settles, supply voltage collapses. In digital CMOS circuits, logic gate
propagation delay (\(t_{\text{pd}}\)) follows the Alpha-Power Law:
\[t_{\text{pd}} \propto \frac{C_L V_{dd}}{(V_{dd} - V_{\text{th}})^\alpha}\]
where \(C_L\) is load capacitance,
\(V_{\text{th}} \approx 0.35\text{ V}\) is transistor threshold voltage,
and \(\alpha \approx 1.3\) is the velocity saturation index. As
\(V_{dd}\) collapses from \(0.85\text{ V}\) to \(0.712\text{ V}\), the
overdrive term \((V_{dd} - V_{\text{th}})\) drops by \(27.6\%\),
stretching gate propagation delays by \(+38\%\). Logic paths fail
flip-flop setup-time margins (\(t_{\text{setup}}\)), causing pipeline
timing faults, dynamic frequency throttling, or hardware brownout resets
that immediately drop valve actuation drive signals.

\protect\phantomsection\label{callout-mathux2a-13.4}
\begin{fbxSimple}{callout-math}{Fermi Math:}{Inductive Power Rail Voltage Droop and Timing Headroom Collapse}
\phantomsection\label{callout-math*-13.4}

\emph{Practical Engineering Question:} Why does activating a neural
accelerator's systolic array cause an adjacent real-time safety
microcontroller core on the same die to fail setup-time timing closure,
even when the power supply is rated for \(100\text{ W}\)?

\emph{1. Current Slew Rate (\(dI/dt\)) and Parasitic Rail Inductance
(\(L\)):} Suppose an unclocked \(8\text{--}\text{core}\) neural
processing unit suddenly enables \(4096\) multiply-accumulate (MAC)
units to execute a transformer linear projection layer. Current draw
steps from \(I_0 = 1.0\text{ A}\) (quiescent) to
\(I_{\text{peak}} = 7.0\text{ A}\) (\(\Delta I = 6.0\text{ A}\)) within
a transition rise time of \(\Delta t = 2.0\text{ ns}\).
\[\frac{dI}{dt} = \frac{6.0\text{ A}}{2.0 \times 10^{-9}\text{ s}} = 3.0 \times 10^9\text{ A/s} = 3.0\text{ A/ns}\]
The on-chip power distribution grid and ball-grid array (BGA) package
exhibit an effective parasitic loop inductance
\(L = 30\text{ pH} = 30 \times 10^{-12}\text{ H}\) and equivalent series
resistance \(R = 8.0\text{ m}\Omega = 0.008\ \Omega\).

\emph{2. Rail Voltage Droop Calculation:}
\[\Delta V = L \cdot \frac{dI}{dt} + I_{\text{peak}} \cdot R = (30 \times 10^{-12}\text{ H})(3.0 \times 10^9\text{ A/s}) + (6.0\text{ A})(0.008\ \Omega) = 90\text{ mV} + 48\text{ mV} = 138\text{ mV}\]
On a nominal core supply \(V_{dd} = 0.850\text{ V}\), the rail sags to
\(V_{\text{core}} = 0.850\text{ V} - 0.138\text{ V} = 0.712\text{ V}\),
breaching the core's minimum operating voltage ceiling
(\(V_{\min} = 0.780\text{ V}\)).

\emph{3. Alpha-Power Law Gate Delay Stretch:} With threshold voltage
\(V_{\text{th}} = 0.350\text{ V}\) and velocity saturation parameter
\(\alpha = 1.30\):

\begin{itemize}
\tightlist
\item
  Nominal overdrive:
  \((V_{dd} - V_{\text{th}})^\alpha = (0.850 - 0.350)^{1.3} = (0.500)^{1.3} \approx 0.406\)
\item
  Drooped overdrive:
  \((V_{\text{droop}} - V_{\text{th}})^\alpha = (0.712 - 0.350)^{1.3} = (0.362)^{1.3} \approx 0.266\)
\item
  Propagation delay scaling ratio:
  \[\frac{t_{\text{pd,droop}}}{t_{\text{pd,nom}}} = \frac{V_{\text{droop}} / (V_{\text{droop}} - V_{\text{th}})^\alpha}{V_{dd} / (V_{dd} - V_{\text{th}})^\alpha} = \frac{0.712 / 0.266}{0.850 / 0.406} = \frac{2.677}{2.094} \approx 1.278 \implies +27.8\% \text{ gate delay stretch}\]
  When combined with high-temperature metallization resistance
  (\(+10\%\)), critical logic path delay expands by \(+38\%\).
\end{itemize}

\emph{Systems Verdict:} A logic pipeline designed with a \(15\%\)
setup-time margin (\(t_{\text{setup}}\)) suffers catastrophic timing
violations on the real-time core. To survive, the system-on-chip's power
management unit must either enforce dedicated deep-trench capacitors
(DTC) with localized power islands or force emergency dynamic frequency
halving (\(1.2\text{ GHz} \to 600\text{ MHz}\)), doubling enforcement
execution time and overrunning control deadlines.

\end{fbxSimple}

\begin{figure}[t]

\centering{

\includegraphics[width=0.95\linewidth,height=\textheight,keepaspectratio]{index_files/mediabag//Users/VJ/GitHub/MLSysBook-integrate-physical/physical/book/chapters/13-placement/figures/fig13_soc_contention_droop.pdf}

}

\caption{\label{fig-13-soc-contention-droop}\textbf{Heterogeneous SoC
Resource Contention and Power-Rail Interference.} (a) On-chip memory bus
and cache contention. A high-throughput proposal engine (NPU/MPU)
streams dense parameter tensors (\(180\text{ MB}\), \(256\text{ kB}\)
bursts at \(12.0\text{ GB/s}\)), evicting safety state tables from
shared L3 cache and saturating the non-preemptive DRAM controller queue
(\(t_{\text{wait}} = 27.6\text{ }\mu\text{s}\)). In the timing waterfall
below, queuing stalls inflate the real-time enforcer's sensor read
latency from \(5.0\text{ }\mu\text{s}\) to \(32.6\text{ }\mu\text{s}\),
consuming available timing slack (\(+20\text{ }\mu\text{s}\)) and
overrunning the \(120\text{ }\mu\text{s}\) compute budget by
\(7.6\text{ }\mu\text{s}\) within the \(1.0\text{ ms}\) control loop.
(b) Power Distribution Network (PDN) transient voltage droop and clock
degradation. Fast current steps (\(\Delta I = 6.0\text{ A}\),
\(dI/dt = 3.0\text{ A/ns}\)) excite package loop parasitics
(\(L = 30\text{ pH}\), \(R = 8.0\text{ m}\Omega\)), inducing an
instantaneous rail droop \(\Delta V = L(dI/dt) + IR = 138\text{ mV}\)
(\(16.2\%\) collapse on \(0.85\text{ V}\) rail). The resulting voltage
sag breaches \(V_{\min} = 0.78\text{ V}\), stretching gate propagation
delays by \(+38\%\) and triggering DVFS emergency clock halving
(\(1.2\text{ GHz} \to 600\text{ MHz}\)) that expands enforcement
execution beyond physical deadlines. Attribution: Original diagram
created for this textbook (CC BY 4.0).}

\end{figure}%

Detecting these interactions requires dedicated hardware
instrumentation, because software running on a core that has been
stalled cannot record the duration of its own stall. Memory-controller
performance counters log queue occupancy, bank conflicts, and per-master
read and write transaction counts. Performance monitoring units on the
real-time core record cache miss counts and stall cycles against an
invariant free-running counter. On the electrical side, on-die
analog-to-digital converters and fast comparator circuits record voltage
excursions on internal rails, thermal diodes log junction temperature
gradients across the floorplan, and power-management logs record dynamic
frequency throttling states. When an uncommanded reset occurs,
reset-cause registers distinguish a software watchdog expiration from an
under-voltage brownout or a thermal emergency trip.

When contention is budgeted through measured burst times, queue
structures, and power supply impedances rather than assumed from
isolated benchmarks, the boundary between proposal and enforcement
ceases to be a conceptual line on a block diagram. As summarized in
Table~\ref{tbl-13-contention-channels}, multi-master contention on
shared silicon operates across digital memory queues, cache hierarchies,
crossbar switches, and analog power/thermal networks, requiring active
hardware interceptors to guarantee real-time safety.

\begin{longtable}[]{@{}
  >{\raggedright\arraybackslash}p{(\linewidth - 10\tabcolsep) * \real{0.1667}}
  >{\raggedright\arraybackslash}p{(\linewidth - 10\tabcolsep) * \real{0.1667}}
  >{\raggedright\arraybackslash}p{(\linewidth - 10\tabcolsep) * \real{0.1667}}
  >{\raggedright\arraybackslash}p{(\linewidth - 10\tabcolsep) * \real{0.1667}}
  >{\raggedright\arraybackslash}p{(\linewidth - 10\tabcolsep) * \real{0.1667}}
  >{\raggedright\arraybackslash}p{(\linewidth - 10\tabcolsep) * \real{0.1667}}@{}}
\caption{\textbf{Shared SoC Resource Contention Channels and Physical
Coupling Mechanisms.} Breakdown of five physical interference vectors on
unified silicon. Contention operates at microsecond and nanosecond
timescales beneath operating system visibility, requiring architectural
partitioning or dedicated hardware isolation to prevent cognitive bursts
from violating reflex timing
bounds.}\label{tbl-13-contention-channels}\tabularnewline
\toprule\noalign{}
\begin{minipage}[b]{\linewidth}\raggedright
Contention Channel
\end{minipage} & \begin{minipage}[b]{\linewidth}\raggedright
Physical Coupling Mechanism
\end{minipage} & \begin{minipage}[b]{\linewidth}\raggedright
Contention Trigger \& Workload Dynamics
\end{minipage} & \begin{minipage}[b]{\linewidth}\raggedright
Nominal vs.~Contended Delay
\end{minipage} & \begin{minipage}[b]{\linewidth}\raggedright
Systems Failure Manifestation
\end{minipage} & \begin{minipage}[b]{\linewidth}\raggedright
Hardware Mitigation \& Interceptor
\end{minipage} \\
\midrule\noalign{}
\endfirsthead
\toprule\noalign{}
\begin{minipage}[b]{\linewidth}\raggedright
Contention Channel
\end{minipage} & \begin{minipage}[b]{\linewidth}\raggedright
Physical Coupling Mechanism
\end{minipage} & \begin{minipage}[b]{\linewidth}\raggedright
Contention Trigger \& Workload Dynamics
\end{minipage} & \begin{minipage}[b]{\linewidth}\raggedright
Nominal vs.~Contended Delay
\end{minipage} & \begin{minipage}[b]{\linewidth}\raggedright
Systems Failure Manifestation
\end{minipage} & \begin{minipage}[b]{\linewidth}\raggedright
Hardware Mitigation \& Interceptor
\end{minipage} \\
\midrule\noalign{}
\endhead
\bottomrule\noalign{}
\endlastfoot
\textbf{DRAM Memory Bus \& Controller} & Non-preemptive burst
arbitration, bank conflicts (\(t_{\text{RCD}}\), \(t_{\text{RP}}\)) &
Accelerator weight tile prefetch (\(256\text{ kB}\) bursts) overlapping
sensor history reads &
\(5.0\text{ }\mu\text{s} \rightarrow 80.0\text{ }\mu\text{s}\)
(\(16\times\) inflation per \(64\text{ kB}\)) & Sensor buffer
starvation, \(1.0\text{ ms}\) control loop overrun, delayed clamp & AXI
QoS priority tags (\texttt{AxQOS}), dedicated DRAM channels, DMA token
buckets \\
\textbf{Shared Last-Level Cache (L3)} & Cache line capacity thrashing
via pseudo-LRU replacement policy & Large parameter tensor streaming
(\(>10\text{ MB}\)) evicting safety state tables & \(20\text{ ns}\) (L2
hit) \(\rightarrow 350\text{ ns}\) (DRAM round-trip penalty) & Pipeline
stalls in invariant checkers, tail latency jitter at \(P_{99.9}\) &
Cache way-partitioning (ARM MPAM, Intel CAT), cache-line locking,
private TCM \\
\textbf{On-Chip Interconnect Crossbar} & Arbiter queue head-of-line
blocking on multi-layer AXI/NoC fabrics & Concurrent camera DMA streams
and diagnostic circular ring-buffer flushes &
\(35\text{ ns} \rightarrow 840\text{ ns}\) transaction stall
(\(24\times\) inflation) & Actuator packet dispatch delay, PWM register
jitter, missed heartbeat checks & Segregated virtual channels,
strict-priority crossbar routing, decoupled slave ports \\
\textbf{Power Distribution Network (PDN)} & Inductive (\(L\cdot dI/dt\))
and resistive (\(I\cdot R\)) rail voltage collapse & Simultaneous
unclocked MAC array activation (\(\Delta I = 6\text{--}22\text{ A}\),
\(t_r < 5\text{ ns}\)) & \(\Delta V = 138\text{--}165\text{ mV}\) droop
(\(>16\%\) drop on \(0.85\text{ V}\) rail) & Gate propagation stretch,
timing closure violations, under-voltage reset & On-package deep-trench
capacitors (DTC), current-slew limiters, isolated PMIC rails \\
\textbf{Shared Thermal Package} & Substrate heat diffusion through
package spreader (\(\theta_{\text{JA}} = 2.0\text{ K/W}\)) & Sustained
peak inference (\(15\text{--}25\text{ W}\)) driving junction past
\(T_{\text{trip}} = 95^\circ\text{C}\) & Clock halved
(\(1.2\text{ GHz} \rightarrow 600\text{ MHz}\)), execution \(2.0\times\)
& Enforcement execution expands
(\(0.70\text{ ms} \rightarrow 1.40\text{ ms}\)), deadline missed &
Independent PLL domains, NPU thermal throttling caps, physical dual-die
split \\
\end{longtable}

The interaction between two paths on one die cannot be left to
uncoordinated arbitration. What crosses the physical boundary must be
defined by explicit guarantees on production time, validity, and failure
response.

\section{The Boundary Contract}\label{sec-placement-13-4}

What crosses between the proposal path and the enforcement path is a
contract carrying a production timestamp, a sequence identifier, a
validity interval, and a defined response for every failure mode, not an
unqualified pointer into shared memory. Passing a raw memory address
across the trust boundary couples the reader directly to the writer's
internal memory layout, execution timing, and failure domain. If a
vision-language model or deep neural policy running on an accelerator
provides only a pointer to its private output buffer, the real-time core
that regulates a gas valve has no mechanism to determine whether the
writer is still modifying the buffer, has stalled mid-inference, or has
encountered an unhandled fault. The boundary payload must therefore be
an immutable, self-contained record. For a fluid flow regulator, this
record contains the target mass flow rate command
\(\dot{m}_{\text{cmd}}\), a monotonically increasing sequence counter
\(k\), a generation timestamp \(t_{\text{prod}}\) sampled from an
invariant hardware timer, an expiration lease duration
\(\Delta t_{\text{valid}}\), and an integrity checksum calculated over
the payload bytes. In a canonical heterogeneous dual-brain architecture
(as formalized in the intent contract in Chapter~\ref{sec-nervous}),
this interface manifests as a packed struct exchanged across zero-copy
shared SRAM mailboxes, binding spatial target vectors,
velocity/acceleration limits, epistemic model confidence
\(\gamma_{\text{conf}}\), and an IEEE 802.3 32-bit CRC checksum. The
enforcement path reads the record as an isolated snapshot, evaluating
its temporal freshness and physical validity before any value reaches
the valve actuator.

Preserving temporal isolation across a shared die requires eliminating
all blocking synchronization primitives at the exchange interface. If
the real-time enforcement loop acquires a mutual exclusion lock held by
the proposal generator, any execution stall in the neural pipeline
propagates directly into the physical plant. In an analytical example of
a flow regulation loop operating at \(1000\text{ Hz}\) with a
\(1.0\text{ ms}\) tick budget, a \(25\text{ ms}\) inference tail stall
on the host processor would cause the real-time core to miss 25
consecutive actuation updates while gas pressure in the manifold
continues to rise. Dynamic memory allocation introduces the same
vulnerability, because allocating or freeing buffer structures from a
shared heap contends for allocator spinlocks and triggers unpredictable
memory reclamation latency. Backpressure is equally prohibited. The
proposal path cannot throttle the rate of the enforcement loop when the
plant requires continuous regulation, and the enforcement path cannot
stall the policy pipeline when fresh sensor frames arrive. The
interaction must be one-way, non-blocking, and asynchronous, ensuring
that the enforcement core runs to completion on every tick regardless of
the operational state of the publisher.

When the proposal path emits trajectory chunks rather than single
setpoints, the exchange boundary requires a bounded ring buffer sized to
decouple the two execution rates. Consider an analytical example where
the learned policy produces flow setpoint updates at a rate of
\(f_{\text{prop}} = 50\text{ Hz}\), giving a production period of
\(T_{\text{prop}} = 20\text{ ms}\), while the enforcement loop services
the valve at \(f_{\text{enf}} = 1000\text{ Hz}\) with a period of
\(T_{\text{enf}} = 1.0\text{ ms}\). If the real-time core experiences a
maximum servicing delay of \(\Delta t_{\text{stall}} = 60\text{ ms}\)
during emergency sensor diagnostics or flash logging, the producer emits
proposals that the consumer cannot immediately drain. The minimum buffer
capacity \(N\) required to preserve all unread proposals during this
worst-case consumer stall is derived from the producer rate and the
stall duration,
\(N = \lceil f_{\text{prop}} \cdot \Delta t_{\text{stall}} \rceil + 1 = \lceil 50\text{ Hz} \cdot 0.060\text{ s} \rceil + 1 = 4\text{ slots}\).
If the buffer fills beyond \(N\) slots, the exchange policy must define
whether to discard the oldest unread proposal, reject the incoming
write, or initiate a fallback. In a flow regulation system where
physical valve dynamics evolve forward in time, retaining stale
trajectory targets while rejecting new state corrections degrades
tracking accuracy. The exchange contract therefore overwrites the oldest
slot on overflow while setting a queue-overrun flag, forcing the
consumer to recognize that intermediate proposals were dropped.

Preventing the enforcement reader from observing a partially written or
mixed-version record requires a lock-free publication protocol governed
by memory barriers (Figure~\ref{fig-13-lockfree-boundary-contract}
a).\sidenote{\footnotesize On weakly ordered SoCs, cross-core memory visibility
  requires hardware barriers. Writers issue Data Memory Barriers
  (\texttt{DMB\ ISH}) after updating sequence counters and flushing
  cache lines. Readers issue \texttt{DMB\ ISH} before and after reading
  payloads to prevent out-of-order execution past sequence checks.} The
writer and reader coordinate through a single shared sequence counter
\(S\). Before modifying the payload buffer, the writer increments \(S\)
from an even value to an odd value, executes a store-release memory
barrier to ensure that the updated sequence is visible to interconnect
arbiters, writes the payload fields, executes a second store-release
barrier, and increments \(S\) to the next even integer. The enforcement
reader reads \(S\) using a load-acquire operation, storing the initial
value \(S_1\). If \(S_1\) is odd (\(S_1 \pmod 2 = 1\)), the writer is
actively modifying the buffer, and the reader rejects the slot without
touching the payload. If \(S_1\) is even, the reader executes a
load-acquire barrier, reads the payload fields, executes another
load-acquire barrier, and samples the counter a second time as \(S_2\).
The reader accepts the data if and only if \(S_1 = S_2\). If
\(S_1 \neq S_2\), the writer updated or started modifying the record
during the read operation, indicating that the data is torn. Tracing the
writer across all preemption points demonstrates why a mixed-version
record cannot be exposed. If the writer halts before the first
increment, the reader samples the previous even sequence and obtains the
prior intact record. If the writer halts after setting \(S\) odd or
during payload writes, the reader observes an odd sequence on every
subsequent cycle and rejects the buffer. If the writer halts after the
final even increment, the reader receives the new, fully committed
record. At no point can a stalled or preempted writer force the reader
to consume inconsistent bytes.

\protect\phantomsection\label{callout-ruleux2a-13.5}
\begin{fbxSimple}{callout-rule}{Architect’s Rule of Thumb:}{Lock-Free Asymmetric Decoupling}
\phantomsection\label{callout-rule*-13.5}

Never share a blocking synchronization primitive across cognitive
deliberation and hard real-time safety domains; asynchronous lock-free
sequence mailboxes guarantee that reader execution latency is bounded
(\(< 5\,\mu\text{s}\)) regardless of writer preemption, crashes, or
unhandled faults.

\end{fbxSimple}

\protect\phantomsection\label{callout-algorithmux2a-13.6}
\begin{fbxSimple}{callout-algorithm}{Algorithm:}{Algorithm 13.1: Hardware Crossbar QoS Arbitration and Thermal Throttling Supervisor}
\phantomsection\label{callout-algorithm*-13.6}

\textbf{Inputs:} Memory transaction request
\(\text{req} = \{\text{master\_id}, \text{addr}, \text{burst\_len}, \text{is\_write}, \text{qos\_tag}\}\),
silicon junction temperature \(T_j\) from on-die thermal diodes.

\textbf{State:} Crossbar arbiter queues
\(Q_{\text{realtime}}, Q_{\text{besteffort}}\); PMU clock frequencies
\(f_{\text{mcu}}, f_{\text{npu}}\).

\textbf{Structural Invariants:}

\begin{enumerate}
\def\labelenumi{\arabic{enumi}.}
\tightlist
\item
  \emph{QoS Priority Invariant:} Real-time safety transactions
  (\(\text{qos\_tag} == \text{QOS\_CRITICAL}\)) strictly preempt pending
  best-effort bursts.
\item
  \emph{Real-Time Clock Invariant:} Real-time core clock
  \(f_{\text{mcu}} \ge f_{\text{nom}}\) is invariant; thermal DVFS
  throttling applies exclusively to \(f_{\text{npu}}\).
\end{enumerate}

\textbf{Procedure 1: Crossbar Transaction Arbitration
(\texttt{ARBITRATE\_CROSSBAR\_TRANSACTION}):}

\begin{itemize}
\item
  If
  \(\text{req}.\text{qos\_tag} == \text{QOS\_CRITICAL} \implies \text{EnqueueHead}(Q_{\text{realtime}}, \text{req})\)
\item
  Else
  \(\implies \text{EnqueueTail}(Q_{\text{besteffort}}, \text{req})\)
  (rate-limited token bucket queue).
\item
  If \(\text{IsCrossbarIdle}()\) and not
  \(\text{IsEmpty}(Q_{\text{realtime}})\):

  \begin{itemize}
  \tightlist
  \item
    \(\text{txn} \leftarrow \text{Dequeue}(Q_{\text{realtime}})\);
    \(\text{DispatchAxiBurst}(\text{txn}, \text{channel}=\text{CH\_DIRECT})\)
    (dedicated DRAM bank).
  \end{itemize}
\item
  Else if \(\text{IsCrossbarIdle}()\) and not
  \(\text{IsEmpty}(Q_{\text{besteffort}})\):

  \begin{itemize}
  \item
    \(\text{txn} \leftarrow \text{Dequeue}(Q_{\text{besteffort}})\)
  \item
    If \(\text{TokenBucketAvailable}(\text{txn}.\text{burst\_len})\):

    \begin{itemize}
    \tightlist
    \item
      \(\text{ConsumeTokens}(\text{txn}.\text{burst\_len})\);
      \(\text{DispatchAxiBurst}(\text{txn}, \text{channel}=\text{CH\_SHARED})\).
    \end{itemize}
  \item
    Else:

    \begin{itemize}
    \tightlist
    \item
      \(\text{StallMaster}(\text{txn}.\text{master\_id}, \text{duration}=50\text{ ns})\)
      (throttle unconstrained DMA bursts).
    \end{itemize}
  \end{itemize}
\end{itemize}

\textbf{Procedure 2: Thermal Budget Supervision
(\texttt{SUPERVISE\_THERMAL\_BUDGET}):}

\begin{itemize}
\item
  If \(T_j \ge T_{\text{trip}}\) (e.g., \(95.0^\circ\text{C}\)):

  \begin{itemize}
  \tightlist
  \item
    \(\text{ThrottleClock}(\text{DOMAIN\_NPU}, f_{\text{target}} = f_{\text{npu}} / 2)\);
    \(\text{AssertThermalFlag}(\text{STATUS\_DEGRADED})\).
  \item
    \(\text{LockClockDomain}(\text{DOMAIN\_MCU}, f_{\text{target}} = f_{\text{nom}})\)
    (protect safety loop clock).
  \end{itemize}
\item
  Else if \(T_j \le T_{\text{clear}}\) (e.g., \(85.0^\circ\text{C}\)):

  \begin{itemize}
  \tightlist
  \item
    \(\text{RestoreClock}(\text{DOMAIN\_NPU}, f_{\text{target}} = f_{\text{npu\_nom}})\);
    \(\text{ClearThermalFlag}(\text{STATUS\_DEGRADED})\).
  \end{itemize}
\end{itemize}

\end{fbxSimple}

Every read operation on the enforcement side must evaluate the published
record against four distinct failure classes and execute a pre-computed
fallback within the single-tick control window
(Figure~\ref{fig-13-lockfree-boundary-contract} b). A record is
classified as malformed if the payload checksum fails or if the
commanded flow \(\dot{m}_{\text{cmd}}\) violates physical bounds, such
as a negative mass flow or a valve displacement exceeding mechanical
stops. A record is classified as stale when the local monotonic hardware
clock \(t_{\text{now}}\) exceeds the expiration limit,
\(t_{\text{now}} > t_{\text{prod}} + \Delta t_{\text{valid}}\). A record
is classified as missing when the sequence counter skips an expected
index (\(k_{\text{new}} > k_{\text{prev}} + 1\)), and classified as
overwritten when the sequence check fails (\(S_1 \neq S_2\)). In every
failure case, the enforcement loop cannot pause to request
retransmission or wait for the host processor to recover. Instead, the
loop branches to a deterministic fallback action parameterized by the
duration of the outage. For transient dropouts lasting less than a hold
threshold \(\tau_{\text{hold}} = 10\text{ ms}\), the controller clamps
the valve at its last verified safe orifice position. If the fault
persists beyond \(\tau_{\text{hold}}\), the controller executes a
predefined deceleration ramp to close the valve at a bounded rate of
\(d\dot{m}/dt = -5.0\text{ (kg/s)/s}\), preventing pressure transients
in the upstream piping before de-energizing the actuator drive.

\begin{figure}[t]

\centering{

\includegraphics[width=0.95\linewidth,height=\textheight,keepaspectratio]{index_files/mediabag//Users/VJ/GitHub/MLSysBook-integrate-physical/physical/book/chapters/13-placement/figures/fig13_lockfree_boundary_contract.pdf}

}

\caption{\label{fig-13-lockfree-boundary-contract}\textbf{Lock-Free
Seqlock Boundary Contract and Preemption Recovery Hierarchy.} (a)
Inter-core lock-free state exchange protocol between an untrusted
proposal writer (NPU/MPU) and a deterministic real-time safety reader
(MCU). Atomic \texttt{store-release} and \texttt{load-acquire} memory
barriers coordinate updates through a shared monotonic sequence counter
\(S\) without acquiring mutual exclusion locks or incurring heap
allocation. A self-contained \(256\text{ byte}\) payload binds commanded
mass flow \(\dot{m}_{\text{cmd}}\), sequence index \(k\), hardware
timestamp \(t_{\text{prod}}\), lease validity
\(\Delta t_{\text{valid}}\), and CRC-32 checksum. (b) Writer preemption
matrix and deterministic fallbacks. The safety reader samples sequence
snapshots (\(S_1, S_2\)); odd values (\(S \pmod 2 = 1\)) or mismatched
sequences (\(S_1 \neq S_2\)) instantly trap in-flight or torn writes
within \(<5\text{ }\mu\text{s}\) with zero blocking. Four boundary
failure classes (malformed, stale, missing, torn) map directly to
bounded real-time fallback actions, guaranteeing continuous
\(1000\text{ Hz}\) actuator regulation during host crashes. Attribution:
Original diagram created for this textbook (CC BY 4.0).}

\end{figure}%

Enforcing these guarantees in hardware memory and software state ensures
that software crashes, memory corruptions, and pipeline stalls on the
proposal side cannot breach the timing bounds of the physical
controller. The contract decouples the logical validity of each command
from the operational state of the generating core, allowing the flow
regulator to maintain safe containment even when the neural network
enters an uncommanded reset. However, establishing clean boundaries in
memory and interconnect arbitration solves only the logical half of
co-location. When the accelerator and the real-time processor share the
same silicon die, the high power dissipated during dense matrix
operations conducts directly through the substrate, elevating local
junction temperatures and altering transistor switching speeds across
adjacent cores.

\section{Thermal Coupling and Timing}\label{sec-placement-13-5}

Every watt dissipated on a silicon die flows into a single thermal
domain. In a co-located architecture, heat generation is not confined to
the compute cores executing the proposal policy. Power dissipation
divides across four coupled mechanisms, comprising dense matrix units
performing neural inference, general-purpose cores executing the
real-time enforcement loop, memory interface transceivers driving
high-bandwidth external memory traffic, and parasitic losses inside
integrated voltage regulators. While the enforcement loop consumes a
predictable, low-power baseline, the proposal path exhibits large
dynamic swings depending on batch dimensions and model activation
patterns. Because the package heat spreader, thermal interface material,
and protective enclosure present a finite thermal conductivity to
ambient air, heat produced by peak accelerator activity raises the
common junction temperature across all adjacent circuits on the die.

In physical AI systems engineering, \emph{thermal clock throttling} is
the automatic, hardware-enforced reduction of processor clock frequency
triggered when sustained high-power cognitive inference elevates silicon
junction temperature (\(T_j\)) past thermal trip thresholds
(\(T_{\text{trip}}\)), introducing state-dependent hysteretic execution
time inflation that forces co-located real-time safety loops to overrun
hard physical deadlines.

Device temperature cannot be approximated by ambient room conditions. A
machine operating in an industrial flow-control enclosure at an ambient
temperature of \(45.0\,^{\circ}\text{C}\) experiences an internal
temperature rise governed by total power dissipation and the thermal
resistance of the assembly. On a modern heterogeneous robotics platform,
such as the high-density heterogeneous SoC architecture shown in
Figure~\ref{fig-13-jetson-orin}, high-density neural accelerators,
general-purpose application cores, and real-time processing clusters
reside within millimeters of each other under a common heatsink and fan
assembly. Consider an analytical example where a shared system-on-chip
is mounted on a conduction-cooled enclosure with an effective
junction-to-ambient thermal resistance of
\(\theta_{\text{JA}} = 2.0\text{ K/W}\) and a thermal time constant of
\(\tau_{\text{th}} = 8.0\text{ s}\). When the enforcement loop runs in
isolation, drawing \(P_{\text{enf}} = 1.5\text{ W}\) alongside baseline
idle and memory power of \(P_{\text{base}} = 2.5\text{ W}\), total
dissipation is \(4.0\text{ W}\), producing a steady-state temperature
rise of \(8.0\text{ K}\) and maintaining the junction at
\(53.0\,^{\circ}\text{C}\). When a high-rate vision-language-action
policy begins continuous proposal inference, the accelerator draws an
additional \(P_{\text{prop}} = 16.0\text{ W}\), memory traffic adds
\(P_{\text{mem}} = 3.5\text{ W}\), and internal power regulator
inefficiencies contribute \(P_{\text{reg}} = 2.5\text{ W}\). Total power
steps to \(P_{\text{total}} = 26.0\text{ W}\). The steady-state junction
temperature rises by
\(\Delta T = P_{\text{total}} \times \theta_{\text{JA}} = 26.0\text{ W} \times 2.0\text{ K/W} = 52.0\text{ K}\),
elevating the silicon substrate to \(T_j = 97.0\,^{\circ}\text{C}\).

\begin{figure}[t]

\centering{

\includegraphics[width=0.85\linewidth,height=\textheight,keepaspectratio]{/Users/VJ/GitHub/MLSysBook-integrate-physical/physical/book/chapters/13-placement/figures/fig13_real_jetson_orin.jpg}

}

\caption{\label{fig-13-jetson-orin}\textbf{Heterogeneous Edge Robotics
Compute Architecture and Thermal Management.} High-density edge robotics
System-on-Module (SoM) and carrier board architecture. An integrated SoC
combines a high-throughput GPU, dedicated deep learning accelerators
(e.g., systolic arrays), and real-time application and safety processor
clusters sharing a unified LPDDR5 memory subsystem. Under peak cognitive
inference workloads (\(15\text{--}60\text{ W}\) TDP), localized
substrate heat dissipation necessitates a massive finned aluminum
heatsink with PWM fan cooling to prevent junction temperatures (\(T_j\))
from crossing thermal trip thresholds (\(T_{\text{trip}}\)) that
forcibly throttle real-time safety loop clock frequencies. Peripheral
interfaces (PCIe Gen4, 10GbE, USB 3.2, and 40-pin nervous system
headers) illustrate the physical coupling between high-bandwidth sensor
streams and deterministic actuator lines. Attribution: Hardware photo
courtesy of NVIDIA Corporation (CC BY-NC 4.0).}

\end{figure}%

Silicon thermal management controllers protect the die from permanent
structural degradation through automated clock throttling. When the
internal temperature sensor registers a junction temperature crossing a
hardware trip threshold \(T_{\text{trip}} = 95.0\,^{\circ}\text{C}\),
the clock distribution network forcibly reduces operating frequencies.
In this analytical example, the real-time core cluster shifts from a
nominal frequency of \(f_{\text{nom}} = 1.20\text{ GHz}\) down to a
throttled frequency of \(f_{\text{throt}} = 600\text{ MHz}\). For an
enforcement cycle requiring
\(N_{\text{cycles}} = 8.4 \times 10^5\text{ cycles}\) to parse
telemetry, verify sequence records, compute pressure derivatives, and
evaluate safety polygons, nominal execution time is
\(t_{\text{exec}} = 0.70\text{ ms}\). Within a \(1000\text{ Hz}\)
control period (\(1.00\text{ ms}\) tick deadline), nominal execution
leaves a timing margin of \(0.30\text{ ms}\). Under the throttled state,
execution time scales inversely with frequency, expanding to
\(t_{\text{exec,throt}} = 8.4 \times 10^5 / (600 \times 10^6\text{ s}^{-1}) = 1.40\text{ ms}\).
The enforcement task overruns its period by \(0.40\text{ ms}\), dropping
control cycles and triggering uncommanded safety trips.

Thermal throttling introduces state-dependent hysteresis that extends
deadline violations long after the triggering load subsides. Hardware
thermal managers do not restore nominal clock frequencies immediately
when temperature drops below \(T_{\text{trip}}\). To prevent
high-frequency thermal oscillation around the trip point, the controller
enforces a clearance threshold, such as
\(T_{\text{clear}} = 85.0\,^{\circ}\text{C}\), requiring the junction to
cool by \(10.0\text{ K}\) before restoring the \(1.20\text{ GHz}\)
clock. Because the thermal mass of the package and heat sink cools at a
rate dictated by its exponential time constant
\(\tau_{\text{th}} = 8.0\text{ s}\), reducing processor power after a
missed deadline cannot produce an instantaneous return to safe timing.
Even if software halts the proposal path completely, the die requires
several seconds to dissipate stored heat and cross \(T_{\text{clear}}\).
During this recovery window, the enforcement loop remains trapped at
\(600\text{ MHz}\), systematically violating consecutive real-time
deadlines and forcing the valve into hardware-latched emergency
containment.

\protect\phantomsection\label{callout-ruleux2a-13.7}
\begin{fbxSimple}{callout-rule}{Architect’s Rule of Thumb:}{Thermal Mass as Latency Memory}
\phantomsection\label{callout-rule*-13.7}

Thermal mass acts as a low-pass filter on power, but a long-duration
memory on execution timing: once a heavy cognitive workload trips
thermal clock throttling, the real-time core remains trapped at degraded
clock frequencies for multiple seconds after the cognitive workload has
halted.

\end{fbxSimple}

Diagnosing execution-time expansion requires separating thermal
degradation from electrical power-rail coupling. Both failure modes
originate in heavy accelerator activity, but they operate on distinct
physical timescales. Power-rail coupling manifests as resistive voltage
drops and inductive transients across shared power distribution networks
within microseconds of a matrix multiplication burst. If supply voltage
sags below minimum operating margins, logic gates experience propagation
delays within a single clock cycle, causing pipeline timing violations
or immediate brownout resets. Thermal coupling develops over hundreds of
milliseconds to several seconds as heat diffuses through silicon and
packaging. Correlating execution-time drift against high-frequency
oscilloscope captures of core supply voltage \(V_{\text{core}}(t)\) and
low-frequency telemetry from on-die thermal diodes \(T_j(t)\) isolates
the root cause. A timing fault that tracks the envelope of temperature
rather than instantaneous current transients confirms thermal throttling
rather than rail collapse.

Relocating the proposal model to a physically separate processor or
discrete companion accelerator eliminates on-die thermal interference,
but it alters the timing and resource balance of the entire machine.
Moving the neural workload off-chip replaces intra-die shared memory
transfers with board-level serial buses such as PCIe or Ethernet.
Serialization delays, packet framing, and bus arbitration add tens of
milliseconds of latency to proposal delivery, increasing the required
validity window \(\Delta t_{\text{valid}}\) defined in the boundary
contract. Off-chip transceivers also increase board-level power
consumption, shift heat dissipation toward peripheral power converters,
and introduce new contention points on external memory buses.
Partitioning computation to resolve a thermal deadline violation
redistributes electrical, communication, and timing constraints,
requiring the complete placement budget to be recalculated from physical
measurements.

\section{Testing the Separation}\label{sec-placement-13-6}

An independence claim on shared silicon is not a structural property
established by a block diagram. It is an empirical hypothesis, and the
verification procedure is an experiment designed to refute it. In a
machine that regulates a fluid or gas flow, the nervous system
guarantees that a throttling valve actuates within a deterministic
deadline regardless of what the learned policy executes on the shared
processor. Stating this guarantee requires expressing independence as a
measurable bound on enforcement latency and refusal availability across
all permitted proposal workloads, chip temperatures, and operating
modes. If memory contention or thermal throttling from the proposal
engine delays a safety refusal past the deadline, the two paths are
coupled, and the logical isolation drawn in the system architecture is
false.

The pass threshold for this experiment derives directly from physical
dynamics. Consider a high-pressure gas regulation manifold where
downstream pipe compliance and fluid compressibility dictate a hard
physical deadline \(t_{\text{deadline}} = 15.0\text{ ms}\) to initiate
valve closure before line pressure breaches safe containment limits.
Fixed sensing overhead consumes \(t_{\text{sense}} = 3.5\text{ ms}\) for
pressure transducer sampling, analog-to-digital conversion, and direct
memory access ingest into local buffers. Fixed output overhead consumes
\(t_{\text{act}} = 2.1\text{ ms}\) for serial bus transmission to the
valve drive, command cyclic redundancy checks, and actuator coil
energization. Setting aside an instrumentation uncertainty margin
\(t_{\text{uncert}} = 0.4\text{ ms}\) to account for measurement jitter
and timer resolution, the maximum duration permitted for enforcement
computation is derived by subtraction,
\[t_{\text{enf,max}} = t_{\text{deadline}} - t_{\text{sense}} - t_{\text{act}} - t_{\text{uncert}} = 15.0\text{ ms} - 3.5\text{ ms} - 2.1\text{ ms} - 0.4\text{ ms} = 9.0\text{ ms}\]
This \(9.0\text{ ms}\) boundary is the \textbf{falsification threshold}.
If any enforcement cycle under test exceeds this value, the independence
hypothesis is refuted.

Refuting an independence claim requires synthetic workloads engineered
to induce worst-case resource contention rather than standard benchmarks
that report average throughput.\sidenote{\footnotesize In out-of-order multicore SoCs,
  timing composability collapses due to \emph{timing anomalies}.
  Speculative execution and non-blocking caches mean local optimizations
  (e.g., cache hits) can alter instruction issue order and paradoxically
  increase total execution time under concurrent execution.} High
average utilization schedules memory accesses evenly across clock
cycles, whereas actual interference peaks during synchronized bursts.
The test framework must deliberately flood shared memory channels with
unaligned direct memory access transfers, thrash shared cache capacity
with dirty line writebacks from transformer weight tensors, trigger
accelerator matrix multiplications that induce rapid supply rail voltage
droops, and stream peak-rate sensor telemetry while diagnostic logging
writes to non-volatile flash storage. These stress profiles reproduce
the transient contention that occurs when a proposal model initiates a
large attention pass at the exact microsecond the flow enforcement loop
reads pressure telemetry.

Measuring end-to-end latency across the complete control loop conceals
the physical site of interference. When total response time drifts, an
external pin measurement cannot distinguish between memory bus stalls
during invariant evaluation, pipeline stalls in the central processor,
or queue buildup on the peripheral bus. Instrumentation must record
hardware timestamps at three distinct boundaries, specifically
enforcement-path entry when sensor frames arrive in local memory
(\(t_0\)), decision completion when the safety contract resolves
(\(t_1\)), and actuator publication when the command byte commits to the
transmit queue (\(t_2\)). The computation interval \(t_1 - t_0\)
isolates core execution, cache misses, and local memory stalls, while
the dispatch interval \(t_2 - t_1\) isolates bus arbitration and output
queue contention caused by concurrent background traffic.

Coupling mechanisms compound across physical, electrical, and logical
domains, requiring a full factorial experimental matrix. The evaluation
sweeps across cold junction conditions (\(T_j = 35^\circ\text{C}\)) and
warm junction conditions (\(T_j = 85^\circ\text{C}\)), an idle proposal
path versus a saturated proposal engine executing maximum-batch
inference, diagnostic circular logging disabled versus streaming at full
non-volatile memory bandwidth, and nominal tracking commands versus
boundary-violating commands that trigger refusal routines. Testing
refusal paths under full stress is mandatory because executing a safe
clamping trajectory or entering emergency shutdown involves cold branch
paths, non-cached error-handling code, and secondary bus transactions
that never execute during steady-state flow regulation.

Consider an analytical example of a flow regulation processor evaluated
across \(10^7\) control cycles. In the baseline configuration with an
idle proposal path and a cold die at \(T_j = 35^\circ\text{C}\), the
enforcement path exhibits tight timing bounds, with a median latency
\(P_{50} = 1.8\text{ ms}\), a ninety-ninth percentile
\(P_{99} = 2.4\text{ ms}\), and a maximum observed excursion of
\(2.9\text{ ms}\), well below the \(9.0\text{ ms}\) falsification
threshold. When the proposal path runs sustained vision-language
inference with full diagnostic logging active on a warm die at
\(T_j = 82^\circ\text{C}\), shared cache evictions and memory bus
arbitration expand the distribution. In this loaded state, median
latency increases to \(P_{50} = 4.6\text{ ms}\), while the tail latency
reaches \(P_{99} = 8.1\text{ ms}\) and \(P_{99.9} = 9.7\text{ ms}\),
with a worst-case peak excursion of \(11.4\text{ ms}\). Although the
median remains within the acceptable window, the \(P_{99.9}\) metric and
the peak excursion cross the \(9.0\text{ ms}\) threshold, refuting the
independence hypothesis. Reporting average execution time or passing
percentiles obscures the reality that one out of every thousand cycles
violates the physical deadline and risks fluid overpressure.

Completing millions of test cycles without a timing violation does not
prove absolute independence. It proves only that interference remained
below the falsification threshold across the specific trajectory of
states exercised during that test window. A rigorous verification report
presents the full empirical latency distribution, the maximum observed
excursion, the exact sample count, the total test duration, and the
explicit list of unexercised operational conditions. Operating states
excluded from the bench setup, including electromagnetic interference
from nearby pump motors, degraded power supply rails, and silicon aging
effects, remain unverified risks rather than proven boundaries.

\protect\phantomsection\label{callout-incidentux2a-13.8}
\begin{fbxSimple}{callout-incident}{Landmark Engineering Incident:}{The 1996 Ariane 5 Flight 501 and Shared Avionics Bus Starvation}
\phantomsection\label{callout-incident*-13.8}

On June 4, 1996, the maiden flight of the European Ariane 5 launch
vehicle (Flight 501) was catastrophically terminated 37 seconds after
liftoff, destroying the rocket and its four scientific
satellites.\sidenote{\footnotesize The official Inquiry Board report
  (\citeproc{ref-lions1996ariane5}{{Lions et al.} 1996}) established
  that the SRI software exception, combined with diagnostic data on the
  shared 1553 bus and lack of OBC validation, caused vehicle aerodynamic
  destruction.}

The primary root cause originated in the Inertial Reference System (SRI)
software, which had been reused directly from the Ariane 4 platform. The
software contained a fatal numeric exception: a 64-bit floating-point
variable representing horizontal bias velocity (\(BH\)) was converted
into a 16-bit signed integer without software overflow protection.
Because Ariane 5's initial flight trajectory exhibited substantially
higher horizontal velocities than Ariane 4, the conversion overflowed at
\(t = +36.7\text{ s}\).

When the arithmetic exception occurred, the primary SRI computer halted
its normal navigation loop and transmitted an active diagnostic bit
pattern and error debugging dump across the shared MIL-STD-1553 avionics
data bus. The backup SRI unit, executing identical software on identical
hardware, suffered the exact same arithmetic overflow within
milliseconds and also ceased normal navigation operations.

Critically, the shared MIL-STD-1553 avionics bus lacked hardware
Quality-of-Service (QoS) bandwidth reservation and strict message
typing. The rocket's main On-Board Computer (OBC) misinterpreted the
diagnostic bit pattern and error telemetry stream dumped onto the shared
bus as valid flight attitude coordinates. Interpreting the corrupted
data as an extreme angular deviation, the OBC commanded full gimbal
deflection of the solid rocket booster nozzles and the Vulcain main
engine (\(20^\circ\) full mechanical travel). The sudden extreme angle
of attack at Mach 0.7 induced massive aerodynamic shear forces that
sheared the solid rocket boosters from the core fuselage, triggering the
automatic flight termination system and self-destruction of the
launcher.

\emph{Systems significance for Physical AI:} Shared communication buses
and unthrottled diagnostic logging across shared computing substrates
convert software exceptions directly into catastrophic physical
actuation. Proving Freedom From Interference on shared silicon and bus
fabrics demands strict hardware isolation, dedicated bus channels, and
invariant bandwidth contracts where diagnostic failure logging can never
starve, corrupt, or deceive physical safety reflexes.

\end{fbxSimple}

The failure analysis of this incident illustrates how unexercised
coupling paths breach physical boundaries. On the bench, the proposal
path executed steady-state flow optimization while the enforcement loop
processed nominal valve adjustments. In the field, an abrupt upstream
pressure spike caused the enforcement monitor to trigger an emergency
refusal and command rapid valve closure. Simultaneously, the policy
engine detected an anomaly and initiated a high-priority diagnostic
memory dump over the shared direct memory access bus. The diagnostic
transfer saturated the memory arbiter for \(4.8\text{ ms}\), expanding
the refusal calculation time from its nominal \(3.2\text{ ms}\) to
\(12.1\text{ ms}\). Because the bench test suite had evaluated refusal
logic and diagnostic memory logging in isolation rather than in
combination, the arbiter lockup went undetected until the delayed valve
command allowed line pressure to rupture the downstream burst disc.

Falsifying an independence claim localizes the precise mechanism of
interference, identifying whether isolation collapsed in the shared
cache banks, across the memory bus arbiter, through thermal clock
throttling, or within peripheral serialization queues. Once measurement
proves that the logical boundary drawn on the block diagram does not
exist on the die, the engineering task shifts from characterization to
remediation, determining which shared resources must be partitioned,
reserved, or physically separated across distinct packages.

\section{Isolation and Re-placement Decisions}\label{sec-placement-13-7}

When an empirical test falsifies an independence claim, it demonstrates
that a boundary drawn on an architectural diagram does not exist in the
physical silicon. Resolving that failure requires an explicit
engineering choice among four distinct interventions, comprising static
reservation, spatial partitioning, workload reduction, and physical
relocation to separate silicon. As structured in the decision framework
in Table~\ref{tbl-13-replacement-decisions}, each intervention operates
at a distinct layer of the computing stack, trading execution
throughput, memory capacity, transport latency, board area, or
bill-of-materials cost to restore predictable bounds to the nervous
system.

\begin{longtable}[]{@{}
  >{\raggedright\arraybackslash}p{(\linewidth - 10\tabcolsep) * \real{0.1667}}
  >{\raggedright\arraybackslash}p{(\linewidth - 10\tabcolsep) * \real{0.1667}}
  >{\raggedright\arraybackslash}p{(\linewidth - 10\tabcolsep) * \real{0.1667}}
  >{\raggedright\arraybackslash}p{(\linewidth - 10\tabcolsep) * \real{0.1667}}
  >{\raggedright\arraybackslash}p{(\linewidth - 10\tabcolsep) * \real{0.1667}}
  >{\raggedright\arraybackslash}p{(\linewidth - 10\tabcolsep) * \real{0.1667}}@{}}
\caption{\textbf{Workload Re-Placement and Isolation Decision Matrix.}
Systematic evaluation of remediation pathways when empirical testing
falsifies independence on shared silicon. Each tier trades distinct
engineering currencies (throughput, capacity, latency, board footprint)
to eliminate contention channels before escalating to physically
segregated silicon.}\label{tbl-13-replacement-decisions}\tabularnewline
\toprule\noalign{}
\begin{minipage}[b]{\linewidth}\raggedright
Remediation Strategy
\end{minipage} & \begin{minipage}[b]{\linewidth}\raggedright
Target Coupling Layer
\end{minipage} & \begin{minipage}[b]{\linewidth}\raggedright
Implementation Mechanism
\end{minipage} & \begin{minipage}[b]{\linewidth}\raggedright
Systems Engineering Trade-Off
\end{minipage} & \begin{minipage}[b]{\linewidth}\raggedright
Residual Vulnerability
\end{minipage} & \begin{minipage}[b]{\linewidth}\raggedright
Physical Silicon Migration Trigger
\end{minipage} \\
\midrule\noalign{}
\endfirsthead
\toprule\noalign{}
\begin{minipage}[b]{\linewidth}\raggedright
Remediation Strategy
\end{minipage} & \begin{minipage}[b]{\linewidth}\raggedright
Target Coupling Layer
\end{minipage} & \begin{minipage}[b]{\linewidth}\raggedright
Implementation Mechanism
\end{minipage} & \begin{minipage}[b]{\linewidth}\raggedright
Systems Engineering Trade-Off
\end{minipage} & \begin{minipage}[b]{\linewidth}\raggedright
Residual Vulnerability
\end{minipage} & \begin{minipage}[b]{\linewidth}\raggedright
Physical Silicon Migration Trigger
\end{minipage} \\
\midrule\noalign{}
\endhead
\bottomrule\noalign{}
\endlastfoot
\textbf{Static Reservation} & CPU execution time, OS scheduling
preemption & Fixed-priority preemptive scheduling
(\texttt{SCHED\_FIFO}), static time-slot slicing (ARINC 653) & Idle
capacity (\(20\text{--}40\%\) throughput loss); unallocated cycles
wasted during quiescent states & Micro-architectural contention beneath
OS visibility (shared L3, DRAM bank conflicts) & Memory bus or cache
stalls violate real-time deadlines despite \(100\%\) CPU core
allocation \\
\textbf{Spatial Partitioning} & Cache capacity thrashing, DRAM bus
contention & Cache way-partitioning (ARM MPAM, Intel CAT), MPU/IOMMU
domains, AXI QoS rate limiters & Structural fragmentation; reduced L3
working set for perception (\(+15\text{--}30\text{ ns}\) DRAM access
penalty) & Unpartitionable DRAM controller queues, row-buffer locality
reordering, PDN voltage droop & Matrix engine current steps
(\(\Delta I > 10\text{ A}\)) trigger rail droop, or DRAM queue stalls
exceed slack \\
\textbf{Workload Reduction} & Memory throughput saturation, peak thermal
power & Model quantization (INT8/FP8), context downsampling, reduced
policy rate (\(50\text{ Hz} \rightarrow 20\text{ Hz}\)) & Degraded
policy tracking fidelity; increased control phase lag near dynamic
stability boundaries & Common clock distribution trees, shared Power
Management Units (PMU), unified reset domains & Model accuracy drops
below operational limits, or shared clock/reset faults propagate across
domains \\
\textbf{Physical Relocation} & Power rail collapse, clock loss, thermal
propagation, common reset & Segregate safety reflex to discrete lockstep
MCU (e.g., dual-core lockstep processor) with isolated PMIC and crystal
& Increased board area, BOM cost, and inter-chip transport latency
(\(0.2\text{ }\mu\text{s} \rightarrow 1.5\text{ ms}\) over SPI/CAN) &
External harness EMI, PCB transceiver failures, inter-chip serialization
delay overhead & Mandatory for ISO 26262 ASIL D / IEC 61508 SIL 3 where
substrate-level common-cause faults are prohibited \\
\end{longtable}

Static reservation preserves the physical placement of workloads on the
shared chip while allocating a guaranteed fraction of a shared resource
to the critical path. In a flow regulation machine where nominal valve
tracking runs alongside an adaptive neural model optimizing pump
efficiency, reservation assigns dedicated execution slots or fixed
memory bus bandwidth to the tracking loop. The cost of reservation is
idle capacity. If the tracking loop requires \(2.0\text{ ms}\) of memory
arbiter access during every \(10\text{ ms}\) control window to guarantee
deadline compliance during fluid pressure transients, that \(20\%\)
slice of memory bandwidth remains unavailable to the background model
during steady-state operation, even when the valve position is
stationary and no fluid correction is needed.

Spatial partitioning divides physical hardware structures so that
workloads no longer compete for the same physical cells.\sidenote{\footnotesize ISO
  26262-9 Clause 5 defines Freedom From Interference (FFI) across ASIL
  boundaries. Proving FFI requires verified spatial partitioning
  (MPU/IOMMU memory protection), temporal partitioning (hardware
  bandwidth limiters and watchdog timers), and authenticated
  communication contracts.} Cache way-partitioning, core pinning, and
dedicated memory channels establish hardware barriers within a single
die. The cost of partitioning is lost flexibility and structural
fragmentation. Splitting an eight-megabyte shared cache into a
six-megabyte partition for the policy model and a two-megabyte partition
for flow safety logic prevents cache line evictions across the boundary.
However, it permanently restricts the working set of the learned model,
forcing higher DRAM access rates and increasing mean memory latency by
tens of nanoseconds on every inference step.

\protect\phantomsection\label{callout-ruleux2a-13.9}
\begin{fbxSimple}{callout-rule}{Architect’s Rule of Thumb:}{Substrate Failure Domain Invariant}
\phantomsection\label{callout-rule*-13.9}

When functional safety standards (ISO 26262 ASIL D, IEC 61508 SIL 3)
prohibit common-cause power, clock, and reset failures, software
partitioning on a single SoC is legally and physically insufficient;
safety reflexes must execute on a physically segregated microcontroller
package.

\end{fbxSimple}

Workload reduction resolves contention by shrinking the computational or
memory footprint of the non-critical workload until its peak demand
falls within the uncontended margins of the shared system. An engineer
might downsample telemetry logging from \(1000\text{ Hz}\) to
\(100\text{ Hz}\), quantize policy weights from sixteen-bit floating
point to eight-bit integers to cut memory bus traffic in half, or lower
the execution frequency of the flow optimization model. The cost of
reduction is degraded functional authority. Lowering the update rate of
the policy increases the phase lag of its control proposals, while
weight quantization can degrade the accuracy of flow predictions near
non-linear valve operating limits.

When software reservations, spatial partitioning, and workload reduction
fail to guarantee isolation, the remaining remedy is physical relocation
to separate silicon. Software reservations and cache partitioning
succeed when interference is confined to processor cores and cache
hierarchies where the operating system or hypervisor can enforce strict
scheduling policies. They fail when contention occurs in unpartitionable
hardware structures beneath the visibility of software. DRAM bank
arbiters, command queues, memory bus crossbars, and peripheral direct
memory access controllers process transactions based on physical arrival
times and low-level bus protocols. If a learned policy model initiates a
wide burst transfer that occupies the memory bus while a valve refusal
monitor attempts an emergency register write, the hardware arbiter will
complete the in-flight transfer before servicing the write. If the
resulting queuing delay expands the monitor response time beyond the
physical tolerance of the hydraulic line, no software priority
assignment can prevent the line from overpressuring.

Physical coupling beyond the digital logic also forces a move to
separate silicon. When high-performance neural compute cores draw
transient currents of tens of amperes during complex inference bursts,
localized voltage droop on shared power rails can force the
system-on-chip power management unit to throttle clock frequencies
across the entire die. If that thermal or electrical throttling slows
the nervous system's flow monitoring loop from \(1000\text{ Hz}\) to
\(500\text{ Hz}\), the safety mechanism violates its control deadline
through physical rail coupling rather than bus contention. Furthermore,
sharing a single substrate couples the failure domains. A hardware
fault, memory parity panic, or watchdog reset triggered by the learned
policy restarts the entire system-on-chip, extinguishing the enforcement
monitor at the exact moment when valve control authority is lost.
Placing the nervous system on a physically isolated microcontroller or
provisioning dual redundant compute nodes on a single domain controller
board (Figure~\ref{fig-13-fsd-board}) restores independent power,
clocking, and reset circuitry.

\begin{figure}[t]

\centering{

\includegraphics[width=0.85\linewidth,height=\textheight,keepaspectratio]{/Users/VJ/GitHub/MLSysBook-integrate-physical/physical/book/chapters/13-placement/figures/fig13_real_fsd_board.jpg}

}

\caption{\label{fig-13-fsd-board}\textbf{Dual-Silicon Automotive Domain
Controller Motherboard and Spatial Partitioning.} Dual-SoC
fail-operational compute motherboard from a canonical autonomous vehicle
domain controller. Two fully redundant System-on-Chip packages (Node A
and Node B) are placed side-by-side on a single automotive-grade printed
circuit board. Each compute node is provisioned with its own discrete
LPDDR4 DRAM chips (four packages per SoC), isolated power management
integrated circuits (PMICs), dedicated power rails, and localized flash
storage. Threaded mounting standoffs provide a perimeter interface for a
shared liquid-cooling cold plate, dissipating thermal energy from
simultaneous high-rate vision inference. Peripheral interfaces along the
board perimeter integrate multi-channel coaxial FPD-Link/GMSL camera
deserializers, CAN bus transceivers, and redundant vehicle power feeds,
physically enforcing independent failure domains across the two
autonomous compute engines. Attribution: Hardware photo courtesy of
Tesla, Inc.~/ Ingineerix teardown analysis (fair use for educational
critique).}

\end{figure}%

Relocating a workload to separate silicon is rarely a localized
modification. It triggers a cascade of placement and communication
adjustments across the entire machine. Moving the flow enforcement
monitor off the primary system-on-chip replaces an internal, zero-copy
shared-memory pointer exchange that completed in
\(0.2\text{ }\mu\text{s}\) with an external serial or packetized bus
interface. A point-to-point serial peripheral interface or controller
area network bus operating across board traces introduces serialization
delays, packet framing overhead, and transceiver latency, expanding the
transport time of each sensor reading and valve command to
\(1.5\text{ ms}\) or more.

This added transport latency consumes a substantial portion of the total
reaction budget. In a system with a strict \(10.0\text{ ms}\) deadline
to arrest a fluid pressure surge, spending \(3.0\text{ ms}\) on
round-trip cross-chip communication leaves only \(7.0\text{ ms}\) for
sensor conversion, state estimation, and physical valve actuation. The
engineering team must then re-evaluate the placement of upstream sensor
interfaces, deciding whether pressure transducers should wire directly
to the safety microcontroller or remain on the primary board. Moving the
sensor interface offloads the main chip and eliminates communication
latency for the safety loop, but it simultaneously forces the learned
policy model to receive its observation stream over the slower external
link. Every re-placement decision propagates through the physical
interconnects, shifting bottlenecks from on-die memory arbiters to
board-level pin counts, wiring runs, and serial buses.

\section{Hardware Allocation}\label{sec-placement-13-8}

Every placement choice settles a trade-off between communication latency
and resource contention, but an engineering team cannot rely on
unwritten assumptions to track where those boundaries fall. The
\textbf{placement record} formalizes this physical topology in the
systems engineering notebook. Building upon the baseline schema,
versioning rules, and provenance requirements established in
Chapter~\ref{sec-nervous}, this chapter introduces the specific
descriptors required to capture on-die co-location and silicon
contention.

The \textbf{placement record} is the formal systems specification
contract and empirical verification ledger that documents the physical
execution domain, privilege level, activation period, loaded execution
bounds, shared resource inventory, hardware arbitration mechanisms, and
pre-declared falsification criteria across co-located silicon domains.
As formalized in the memory-mapped specification contract in
Table~\ref{tbl-13-placement-record-schema}
(\texttt{struct\ HeterogeneousPlacementRecord\_t}), the placement record
binds ten core hardware and timing descriptors across every function in
the machine, spanning both the high-level policy regulating flow demand
and the hard-real-time monitor managing valve safety.

\begin{longtable}[]{@{}
  >{\raggedright\arraybackslash}p{(\linewidth - 8\tabcolsep) * \real{0.2000}}
  >{\raggedright\arraybackslash}p{(\linewidth - 8\tabcolsep) * \real{0.2000}}
  >{\raggedright\arraybackslash}p{(\linewidth - 8\tabcolsep) * \real{0.2000}}
  >{\raggedright\arraybackslash}p{(\linewidth - 8\tabcolsep) * \real{0.2000}}
  >{\raggedright\arraybackslash}p{(\linewidth - 8\tabcolsep) * \real{0.2000}}@{}}
\caption{\textbf{Abstract Data Schema for the Heterogeneous SoC
Placement Record (\protect\texttt{HeterogeneousPlacementRecord\_t}).}
Formal systems specification contract and empirical verification ledger
layout (\(128\text{ bytes}\) total), field types, physical SI units,
structural invariants, and hardware audit substrates. This record
establishes Freedom From Interference (FFI) across co-located silicon
domains and is passed directly to the verification and release
pipelines.}\label{tbl-13-placement-record-schema}\tabularnewline
\toprule\noalign{}
\begin{minipage}[b]{\linewidth}\raggedright
Offset \& Field Name
\end{minipage} & \begin{minipage}[b]{\linewidth}\raggedright
Type \& Size
\end{minipage} & \begin{minipage}[b]{\linewidth}\raggedright
Units \& Topology
\end{minipage} & \begin{minipage}[b]{\linewidth}\raggedright
Structural Invariant \& Bound
\end{minipage} & \begin{minipage}[b]{\linewidth}\raggedright
Hardware Audit Substrate
\end{minipage} \\
\midrule\noalign{}
\endfirsthead
\toprule\noalign{}
\begin{minipage}[b]{\linewidth}\raggedright
Offset \& Field Name
\end{minipage} & \begin{minipage}[b]{\linewidth}\raggedright
Type \& Size
\end{minipage} & \begin{minipage}[b]{\linewidth}\raggedright
Units \& Topology
\end{minipage} & \begin{minipage}[b]{\linewidth}\raggedright
Structural Invariant \& Bound
\end{minipage} & \begin{minipage}[b]{\linewidth}\raggedright
Hardware Audit Substrate
\end{minipage} \\
\midrule\noalign{}
\endhead
\bottomrule\noalign{}
\endlastfoot
\texttt{0x00}: \texttt{execution\_domain\_id} & \texttt{uint64\_t} (8B)
& Core Enum & \texttt{domain\_id}
\(\in \{\text{MPU}, \text{NPU}, \text{GPU}, \text{MCU}, \text{LOCKSTEP}\}\)
& Memory Protection Unit (MPU) \\
\texttt{0x08}: \texttt{privilege\_and\_ring} & \texttt{uint64\_t} (8B) &
Privilege Level & Proposal \(\le \text{EL0}\), Enforcer
\(\ge \text{EL1}/\text{Secure}\) & ARM TrustZone / HSM \\
\texttt{0x10}: \texttt{nominal\_period\_ns} & \texttt{uint64\_t} (8B) &
\(\text{ns}\) & \(T_s \in [10^6, 5 \times 10^7]\text{ ns}\)
(\(1.0\text{ ms}\) enforcer) & Free-running monotonic timer \\
\texttt{0x18}: \texttt{hard\_deadline\_ns} & \texttt{uint64\_t} (8B) &
\(\text{ns}\) &
\(t_{\text{deadline}} = T_{\text{period}} - t_{\text{actuator\_transit}}\)
(\(1.0\text{ ms}\)) & Hardware watchdog timer (WDT) \\
\texttt{0x20}: \texttt{wcet\_unloaded\_ns} & \texttt{uint64\_t} (8B) &
\(\text{ns}\) & \(t_{\text{exec,nom}} \le 135\,000\text{ ns}\)
(\(135\,\mu\text{s}\) baseline) & PMU cycle counter
(\texttt{PMCCNTR}) \\
\texttt{0x28}: \texttt{wcet\_loaded\_synthetic\_ns} & \texttt{uint64\_t}
(8B) & \(\text{ns}\) & \(t_{\text{exec,max}} \le 380\,000\text{ ns}\)
(under full AXI load) & Factorial contention test harness \\
\texttt{0x30}: \texttt{shared\_axi\_channel\_qos} & \texttt{uint64\_t}
(8B) & AMBA \texttt{AxQOS} & \(\text{QoS}_{\text{MCU}} = \text{0xF}\),
\(\text{QoS}_{\text{NPU}} \le \text{0x4}\) & Crossbar QoS priority
registers \\
\texttt{0x38}: \texttt{max\_pdn\_droop\_thermal\_trip} &
\texttt{uint32\_t{[}2{]}} (8B) & \(\text{mV}\) \& \(^\circ\text{C}\) &
\(\Delta V_{\text{droop}} \le 70\text{ mV}\),
\(T_{\text{trip}} \le 95.0^\circ\text{C}\) & On-die comparator \&
thermal diode \\
\texttt{0x40}: \texttt{hardware\_mitigation\_flags} &
\texttt{uint8\_t{[}32{]}} (32B) & Bitfield & Active: TCM SRAM, MPAM,
Isolated Rail & Register lock (\texttt{SYSCON\_LOCK}) \\
\texttt{0x60}: \texttt{model\_firmware\_sha256} &
\texttt{uint8\_t{[}32{]}} (32B) & SHA-256 & Cryptographic digest
(weights, kernel, RTOS) & Secure Boot ROM / RoT signature \\
\end{longtable}

These parameters define the physical execution domain (such as a neural
processing unit cluster, an application processor core, or a discrete
lockstep microcontroller), the hardware privilege level, the nominal
activation period, the hard physical deadline, and the measured loaded
execution bound. When an analytical example allocates an upstream
pipeline pressure estimator to a core running at \(1000\text{ Hz}\) with
a \(1.0\text{ ms}\) deadline, the record documents both the unloaded
execution time of \(0.18\text{ ms}\) and the worst-case bound of
\(0.62\text{ ms}\) observed when all co-located execution units operate
under maximum synthetic contention.

Beyond per-function properties, the record constructs an inventory of
shared physical resources across the die and board. This inventory names
every memory interconnect, shared last-level cache, external memory
channel, direct memory access engine, and power rail that connects two
or more execution domains. For each shared element, the entry lists the
competing claimants, the hardware arbitration or partitioning mechanism,
and the maximum interference observed during validation. In a shared-die
system where a \(50\text{ Hz}\) flow regulation policy and a
\(1000\text{ Hz}\) valve actuation loop access the same physical dynamic
random-access memory bus, memory bus contention can inflate memory read
latencies. If quality-of-service channel priorities or cache-line
locking mechanisms are engaged, the record logs their hardware register
configurations alongside the resulting tail latency. In an analytical
test scenario, logging an unmitigated read stall inflation from
\(45\text{ ns}\) to \(310\text{ ns}\) establishes an empirical baseline
that explains why the safety loop required dedicated static
random-access memory buffers or separate memory controllers.

An interference measurement is incomplete without the exact operational
state of the silicon when the sample was taken. The placement record
therefore preserves the complete test environment, binding the data to
immutable hardware and software identifiers. This metadata includes the
exact neural network weight checksums, firmware versions, kernel
configuration hashes, and clock management policies, specifying whether
dynamic voltage and frequency scaling was active or locked to fixed
clock frequencies. It documents the physical test fixture, including the
enclosure geometry, the thermal interface material, the ambient room
temperature, and the physical instrumentation attached to the board,
such as logic analyzer channels or current sense resistors on power
rails. The workload description must record the synthetic traffic
generators used to saturate memory channels, the sensor sampling rates
driving input buffers, and the statistical sample count, requiring a
minimum of \(10^7\) operational cycles to capture rare tail events at
\(P_{99.99}\) confidence.

Every placement entry requires a predeclared falsification criterion
written before the test executes. This criterion specifies the physical
threshold that constitutes architectural failure, such as any instance
where a valve current clamp command is delayed beyond \(0.80\text{ ms}\)
during simultaneous neural accelerator bursts and direct memory access
transfers. The record logs the measured empirical distribution against
this criterion, noting whether the hypothesis of non-interference was
sustained or falsified. When interference violates the target budget and
forces an engineering change, the record details the exact mitigation
applied, such as enabling memory bus bandwidth throttling on the neural
accelerator or moving the pressure sensor direct memory access buffer to
an isolated internal memory bank. The entry must also document the
boundary conditions that were not exercised during the test run, such as
operating states above \(45^\circ\text{C}\) ambient temperature,
multi-bit memory error injection, or simultaneous sensor transceiver
fault states.

Documenting these boundaries during physical integration provides the
structural substrate for the safety analysis in subsequent chapters. The
placement record cannot be generated during the logical architecture
design in Chapter~\ref{sec-enforcement}, because logical block diagrams
do not contain physical pin assignments, clock trees, or memory arbiter
contention profiles. Instead, this record is handed forward directly to
the verification and release workflows. Chapter~\ref{sec-verification}
consumes the recorded arbitration points and shared power rails as the
target map for hardware and software fault injection, deliberately
corrupting shared structures to prove that fault propagation matches the
predicted failure domains. Chapter~\ref{sec-release} incorporates the
recorded loaded execution bounds and unexercised test limits into the
formal safety case, transforming empirical measurements into evidentiary
arguments for operational deployment.

A placement record produced under standard test conditions demonstrates
that physical separation mechanisms function when the hardware operates
within nominal parameters. However, an interference envelope measured on
a cold bench or across a short sample sequence does not guarantee
separation when physical thermal mass saturates and junction
temperatures climb.

\section{Systems Stress-Tests}\label{sec-placement-13-9}

A physical separation that holds at room temperature can fail when the
die warms to operational equilibrium. On a cold test bench at
\(25.0^{\circ}\text{C}\), power distribution networks operate with
minimal resistance, transistor switching margins remain wide, and memory
controllers spend negligible bandwidth on refresh commands. Once an
autonomous machine runs under continuous load, junction temperatures
reach \(85.0^{\circ}\text{C}\) to \(105.0^{\circ}\text{C}\). At these
temperatures, metal line resistance increases, propagation delays
stretch, and resource contention sharpens. A placement claim cannot rest
on nominal traces. The following four analytical stress-tests evaluate
where logical boundaries collapse into physical interference under
thermal and electrical stress:

\subsection*{Stress-Test 1: High-Temperature Current Transient
(Shared-Rail Power
Droop)}\label{stress-test-1-high-temperature-current-transient-shared-rail-power-droop}
\addcontentsline{toc}{subsection}{Stress-Test 1: High-Temperature
Current Transient (Shared-Rail Power Droop)}

\emph{The scenario.} A natural gas flow unit integrates an 8-core neural
accelerator and an isolated lockstep MCU core on one SoC. The enforcer
must read line pressure (\(12.0\text{ MPa}\)) and verify valve current
within \(1.20\text{ ms}\) at \(1000\text{ Hz}\). At a junction
temperature of \(95.0^{\circ}\text{C}\), the accelerator draws a
\(22.0\text{ A}\) current step (\(35.0\text{ ns}\) rise time) on the
shared \(0.85\text{ V}\) rail.

\emph{The systems dilemma.} The transient load induces a
\(165.0\text{ mV}\) droop across package parasitics, triggering on-chip
DVFS circuits to scale enforcer clock frequency from
\(400.0\text{ MHz}\) to \(200.0\text{ MHz}\) to prevent timing faults.
At \(200.0\text{ MHz}\), enforcement execution time stretches beyond the
\(1.20\text{ ms}\) pressure rise deadline, causing physical
over-pressurization.

\emph{The systems verdict: Power Grid Isolation Law.} Electrical rails
do not honor logical privilege levels. An accelerator burst inside a
software container can collapse the power grid of the safety enforcer,
turning a compute spike into actuator control failure unless protected
by on-package decoupling and hardware current-slew limiters.

\subsection*{Stress-Test 2: Clock Jitter and Shared Reset Domains
(Common-Mode Fault
Propagation)}\label{stress-test-2-clock-jitter-and-shared-reset-domains-common-mode-fault-propagation}
\addcontentsline{toc}{subsection}{Stress-Test 2: Clock Jitter and Shared
Reset Domains (Common-Mode Fault Propagation)}

\emph{The scenario.} A cryogenic propellant regulator uses an on-chip
PLL to synthesize clocks for vision pipelines, policy runtimes, and the
real-time enforcement bus. While regulating liquid oxygen at
\(6.50\text{ kg/s}\), a thermal transient unlocks the primary PLL, or a
watchdog asserts a warm reset of the stalled policy cluster across a
shared reset controller.

\emph{The systems dilemma.} Reset assertions in the proposal domain
induce transient electrical glitches and metastable states in the shared
enforcement register file, disabling the enforcer. The enforcer requires
a dedicated crystal oscillator and independent power-on-reset rail,
transitioning valve drivers to a passive spring-return bypass state
within \(0.60\text{ ms}\) of clock instability.

\emph{The systems verdict: Clock and Reset Domain Invariant.} A safety
enforcer cannot outlive the clock distribution network that drives its
flip-flops. When clock trees or reset lines are shared across failure
domains, a fault in an unverified policy engine propagates directly into
the nervous system.

\subsection*{Stress-Test 3: Phase-Aligned Bus Collision and DRAM Refresh
Contention (Interconnect
Starvation)}\label{stress-test-3-phase-aligned-bus-collision-and-dram-refresh-contention-interconnect-starvation}
\addcontentsline{toc}{subsection}{Stress-Test 3: Phase-Aligned Bus
Collision and DRAM Refresh Contention (Interconnect Starvation)}

\emph{The scenario.} A chemical metering machine streams camera DMA
traffic at \(950.0\text{ MB/s}\) into shared DRAM while logging flushes
\(128.0\text{ KB}\) buffers. The enforcer reads calibration tables and
pressure across the main crossbar. Although average bandwidth
utilization is \(28\%\), DMA bursts align with enforcer reads at
\(T_j = 90.0^{\circ}\text{C}\), where DRAM refresh overhead consumes
\(12\%\) of total channel availability.

\emph{The systems dilemma.} The phase-aligned collision expands enforcer
read latency from \(35.0\ \mu\text{s}\) to a \(P_{99.99}\) tail of
\(840.0\ \mu\text{s}\), starving the \(1000\text{ Hz}\) control loop.
Relocating safety-critical state and calibration tables to
tightly-coupled SRAM (TCM) bounds read latency below
\(100.0\ \text{ns}\) under full bus saturation.

\emph{The systems verdict: Peak Contention Law.} Average memory
bandwidth is meaningless for real-time safety. Interference on shared
interconnects is instantaneous and phase-dependent; safety state must
reside in tightly-coupled static memory isolated from shared DRAM
traffic.

\subsection*{Stress-Test 4: Abrupt Policy Crash During Record
Serialization (Asymmetric Mailbox
Boundaries)}\label{stress-test-4-abrupt-policy-crash-during-record-serialization-asymmetric-mailbox-boundaries}
\addcontentsline{toc}{subsection}{Stress-Test 4: Abrupt Policy Crash
During Record Serialization (Asymmetric Mailbox Boundaries)}

\emph{The scenario.} The policy governing a gas blending manifold
experiences an unhandled numeric exception while writing a
\(256\text{ byte}\) flow schedule into a shared lock-free ring buffer,
terminating with sequence counters half-updated while gas flows at
\(25.0\text{ m/s}\).

\emph{The systems dilemma.} If the interface relies on software mutexes,
the crashed producer leaves locks permanently acquired, deadlocking the
enforcer. The boundary detector must identify malformed or partial
records in under \(5.0\ \mu\text{s}\), dropping proposals without
blocking and executing deterministic tracking fallback.

\emph{The systems verdict: Adversarial Producer Law.} An untrusted
proposal engine must be treated as an adversarial producer. The nervous
system must reject corrupt or partial publications with zero wait
states, ensuring that a policy crash never consumes enforcement
execution budget.

\section{Summary}\label{sec-placement-13-10}

A placement result is an empirical claim whose validity extends only as
far as the measured envelope. It does not establish that two software
components are independent in an absolute sense. Instead, the placement
record demonstrates that under an analytical stress workload executing
at a specific ambient temperature, with fixed clock governance policies
and documented fault domains, the shared physical substrates of the die
did not cause the untrusted path to violate the timing constraints of
the safety path. On a gas blending manifold regulating high-velocity
flows at \(25.0\text{ m/s}\), drawing the policy and the safety enforcer
in separate boxes on an architectural diagram does not isolate them from
a shared voltage regulator, a common memory crossbar, or the heat sink
clamped across both processor clusters. If the matrix multiplication
core executing the learned flow proposal draws peak current while all
four memory channels service direct memory access transfers, the
resulting voltage drop and thermal dissipation affect the enforcer core
immediately. The boundaries drawn during system design exist on the
silicon only when verified against these shared electrical, thermal, and
architectural limits.

The placement record converts qualitative architectural intent into five
budgetable physical quantities. First, it defines the loaded enforcement
latency, measured not in isolation but during concurrent
maximum-bandwidth direct memory access transfers and peak inference
activity, establishing an analytical \(P_{99.99}\) execution bound of
\(85\text{ }\mu\text{s}\) against a hard \(100\text{ }\mu\text{s}\)
control deadline. Second, it documents shared-resource headroom,
confirming that average memory bus saturation remains below \(68\%\)
under peak burst traffic, which preserves the \(12\%\) channel
availability required for high-temperature DRAM refresh cycles. Third,
it identifies common-mode failure mechanisms, including frequency
throttling triggered by thermal sensors that slow down deterministic
safety cores alongside policy engines, or phase-aligned bus collisions
that lock arbitration logic. Fourth, it measures detector coverage,
proving that corrupt packets, deadlocks, and stalled ring buffers are
trapped within a deterministic window of \(5.0\text{ }\mu\text{s}\)
without requiring the enforcer to wait on an unresponsive thread. Fifth,
it catalogs the untested conditions, explicitly specifying the operating
temperatures above \(85^\circ\text{C}\), the unprofiled sensor burst
patterns, and the degraded power supply voltages where separation
guarantees have not been established.

With the completion of this placement analysis, the five architectural
handoffs of Part III form a single measured machine. The flow of data
from sensor ingress into state estimation, from estimated state into
learned policy proposal, from proposal into reflex enforcement, and from
enforcement into physical valve actuation can no longer be treated as
independent software modules communicating across idealized channels.
The logical boundaries established across earlier chapters provide the
structure needed to reason about authority and failure propagation, but
physical reality governs their execution. A memory controller does not
recognize software privilege levels when arbitrating row-buffer
conflicts, and an interrupt line asserts whether the processor is
executing user-space code or an enforcer loop. When a valve controller
commands a proportional solenoid to alter gas mixture ratios, the action
is the culmination of a pipeline whose timing, power, and memory
footprints have been measured together on one die.

\protect\phantomsection\label{callout-takeawaysux2a-13.10}
\begin{fbxSimple}{callout-takeaways}{Key Chapter Takeaways}{}
\phantomsection\label{callout-takeaways*-13.10}

\begin{itemize}
\tightlist
\item
  \emph{Logical privilege rings do not isolate shared electrical rails:}
  An accelerator power surge (\(22.0\text{ A}\) current transient,
  \(35\text{ ns}\) rise time) drops shared silicon rail voltage
  (\(L \, dI/dt\)), triggering DVFS clock throttling that doubles safety
  enforcer execution latency and causes deadline overruns unless
  protected by isolated power domains.
\item
  \emph{Clock distribution networks and reset lines are common-mode
  vectors:} Sharing on-chip Phase-Locked Loops (PLLs) or reset
  controllers across proposal and permission domains allows an unhandled
  exception in an unverified policy engine to reset or glitch the
  real-time enforcer's register files.
\item
  \emph{Phase-aligned bus collisions corrupt real-time determinism:}
  Average memory bus utilization (\(28\%\)) is meaningless for real-time
  safety. Simultaneous DMA bursts and DRAM refresh cycles inflate
  enforcer read latency from \(35\ \mu\text{s}\) to \(840\ \mu\text{s}\)
  (\(P_{99.99}\)), requiring safety-critical state to be locked in
  Tightly-Coupled SRAM (TCM).
\item
  \emph{Untrusted policy engines must be treated as adversarial
  producers:} Upstream neural policies crash, hang, or leave sequence
  counters half-updated. The enforcer must consume shared ring buffers
  via lock-free seqlocks and validate header checksums in
  \(<5.0\ \mu\text{s}\) without blocking on locks.
\item
  \emph{The placement deliverable quintet:} Before deploying an
  architecture, five physical invariants must be measured under
  concurrent maximum load: loaded enforcement latency (\(P_{99.99}\)),
  shared-resource headroom, common-mode failure paths, detector fault
  coverage, and cataloged untested boundaries.
\end{itemize}

\end{fbxSimple}

Part IV takes this bounded machine and confronts the problem of
deployment. Having established what the hardware and software layers
cost under load, the remaining questions concern justification,
verification, and operational governance. The analysis shifts from
whether the enforcer can execute within \(100\text{ }\mu\text{s}\) to
how human authority is preserved when an operator overrides a learned
trajectory, what empirical verification evidence is sufficient to
justify operating a non-deterministic model over combustible physical
flows, and what criteria determine a formal release verdict. Because the
architectural limits and shared-resource contention have been measured
rather than assumed, the safety case in Part IV can proceed from
physical facts rather than architectural optimism.

\protect\phantomsection\label{callout-labux2a-13.11}
\begin{fbxSimple}{callout-lab}{Hands-On Lab:}{Placement Ripple and Shared-Bus Contention}
\phantomsection\label{callout-lab*-13.11}

Measure the timing of the real-time enforcement path under heavy
cognitive and direct-memory-access load on the dual-brain kit:

\begin{enumerate}
\def\labelenumi{\arabic{enumi}.}
\tightlist
\item
  \emph{Unloaded baseline:} Run the deterministic \(1\text{ kHz}\)
  safety enforcer on the ARM Cortex-M controller and establish baseline
  execution latency and timer jitter
  (\(t_{\text{exec}} \le 120\,\mu\text{s}\), jitter
  \(< 5\,\mu\text{s}\)) with an idle Linux MPU.
\item
  \emph{Synthetic contention injection:} Execute a high-throughput
  vision-language policy or dense matrix-multiplication benchmark on the
  Linux MPU while streaming maximum-bandwidth MIPI-CSI camera DMA
  transfers into shared DRAM. Log the empirical tail latency
  distribution (\(P_{95}, P_{99}, P_{99.9}\)) of the MCU safety loop.
\item
  \emph{Lock-free preemption testing:} Trigger simulated MPU kernel
  panics and out-of-memory crashes mid-publication in the shared SRAM
  mailbox (\texttt{PhysicalAgent\_IntentProposal\_t}). Verify that the
  MCU detects stale and torn records within \(< 5\,\mu\text{s}\) and
  transitions to deterministic Level 2 safe stop without timing
  overruns.
\end{enumerate}

\emph{Hardware bench guide and starter firmware:} Lab 09: Placement
Ripple.

\end{fbxSimple}

\part{Governing the Machine}

\lettrine{D}{emonstrating} a physical AI model on a curated video reel
or in a controlled laboratory demo is an exciting milestone, but it is
not the destination.

The true test of a physical AI machine begins when it leaves the
workbench and enters the open, unstructured physical world. In that
world, operating conditions are uncooperative: optical lenses become
obscured by dust, lighting transitions violently, wireless networks drop
packets, mechanical gearboxes wear, and human coworkers share the
operating envelope.

Dependability in physical AI does not come from blind optimism in neural
weights or unverified claims of statistical generalization. It comes
from \textbf{The Adult in the Room: Accountability in the Wild}:

\begin{itemize}
\tightlist
\item
  \textbf{Shared Authority with Humans (Chapter 14: Intervention):}
  Treating human operators not as an afterthought or an emergency master
  switch, but as an integral, latency-bounded authority channel. We
  explore how to manage the Four-Phase Handshake Protocol, execute
  \(\mathcal{C}^1\) bumpless control transfers, eliminate operator mode
  confusion, and log tamper-evident forensic black-box telemetry under
  delayed network links.
\item
  \textbf{Breaking the Machine on Purpose (Chapter 15: Verification):}
  Ascending the Four-Stage Qualification Ladder from Simulation to
  Processor-in-the-Loop (PIL), Hardware-in-the-Loop (HIL), and In-Situ
  physical testbeds. We systematically inject hardware and software
  faults---dropping camera frames, starving memory buses, corrupting
  sensor calibration, and inducing power rail droop---to measure exact
  physical survival boundaries before encountering them in production.
\item
  \textbf{The Decision to Ship and Operate (Chapter 16: Release):}
  Grounding production release in empirical, auditable engineering
  evidence. We construct formal Claim-Argument-Evidence (CAE) safety
  cases using Goal Structuring Notation (GSN) under ANSI/UL 4600 and ISO
  26262/21448 standards. We formalize Operational Design Domain (ODD)
  invariants, instrument fleet telemetry, deploy secure A/B partition
  rollbacks, and issue defensible, cryptographic deployment verdicts.
\item
  \textbf{The Open Frontier (Chapter 17: Frontier):} Confronting the
  foundational epistemic limits of embodied intelligence. We analyze
  observational indistinguishability, detectorless failure modes, and
  shortcut learning, providing a capstone synthesis of the entire
  Physical AI curriculum.
\end{itemize}

Part IV provides the engineering discipline needed to take physical AI
from an experimental research prototype to a dependable,
production-certified machine that operates accountably among human
beings.

\chapter{\texorpdfstring{Supervisory
Intervention\chaptersubtitle{Taking Authority Away From the Machine}}{Supervisory Intervention}}\label{sec-intervention}

\begin{figure}[H]

\centering{

\includegraphics[width=1\linewidth,height=\textheight,keepaspectratio]{index_files/mediabag//Users/VJ/GitHub/MLSysBook-integrate-physical/physical/book/chapters/14-intervention/figures/fig_locator.pdf}

}

\caption{\label{fig-locator}\textbf{Systems Organ Focus: Nervous.}
Revoking authority from an autonomous policy requires an independent
real-time arbitration layer to manage discrete state handshakes and
enforce bumpless actuator transitions. Without a deterministic
nervous-system arbiter, concurrent human and machine commands produce
split-brain execution and catastrophic mechanical shock.}

\end{figure}%

\emph{When a human grabs the controls of a running machine, how do you
transfer kinetic authority without causing the very collision you meant
to prevent?}

No autonomous physical machine operates in total isolation from human
authority. Whether through an onboard operator taking a steering wheel,
a technician hitting an emergency stop, or a remote operator engaging
teleoperation, authority over physical actuators must transition
dynamically between human and machine. If this takeover is handled as an
informal software blend, the resulting hybrid command represents neither
intent and risks instant collision. Runtime intervention requires
deterministic state-machine arbitration, synchronized four-phase
handshakes, and higher-order \(\mathcal{C}^2\) bumpless transfer that
respects physical inertia, while cryptographically isolating
intervention traces so that human corrections do not contaminate
autonomous training data.

\newpage{}

\protect\phantomsection\label{callout-objectiveux2a-14.1}
\begin{fbxSimple}{callout-objective}{Learning Objectives}{}
\phantomsection\label{callout-objective*-14.1}

After studying this chapter, you will be able to:

\begin{enumerate}
\def\labelenumi{\arabic{enumi}.}
\tightlist
\item
  Specify which operations a person may command, under what conditions,
  within what deadline.
\item
  Hand over a moving machine without a step change in command, and
  measure the result.
\item
  Mark the moment authority changed hands, and the earlier moment the
  policy diverged, so Chapter~\ref{sec-data} can tell a rescue from a
  nominal episode.
\item
  State what an operational log must contain for a claim to rest on it
  months later.
\end{enumerate}

\emph{Core chapter takeaway.} Authority attaches to named operations
with deadlines, not to a person drawn in a loop. Taking control is a
different mechanism from refusing a command, and a handover that leaves
a moving machine with no coherent command is a hazard rather than a
rescue.

\end{fbxSimple}

\section{Shared Authority}\label{sec-intervention-14-1}

\lettrine{A}{s mapped}\lettrine{a}{cross} the nervous-system tier in
Figure~\ref{fig-locator}, a physical AI system operates with two
distinct internal holders of authority over its actuators: the learned
policy that plans motion over extended horizons and the deterministic
nervous system that enforces immediate safety barriers. When an operator
attempts to halt a moving platform or correct an articulated arm, a
third holder of authority enters the architecture. This human operator
is neither an infallible supervisor nor a synchronous component of the
real-time loop. Human perception, cognitive arbitration, and physical
actuation introduce substantial latency, operating under an analytical
baseline of \(400\text{ ms}\) to \(1200\text{ ms}\) between visual
stimulus and physical input. Across that latency window, an autonomous
vehicle traveling at \(2.0\text{ m/s}\) covers up to \(2.4\text{ m}\) of
unmonitored distance before human intent converts into an electrical
signal at the receiver. Furthermore, the external wireless link carrying
the operator's command remains subject to packet drop, jitter, and
signal attenuation. Because the human operates outside the synchronous
bus and communicates across an untrusted channel, authority cannot be
granted to an abstract human role. Authority attaches to concrete
operations executed under strict, verifiable deadlines.

Treating the human operator as a generic fail-safe layer creates a
dangerous abstraction. Architectural diagrams frequently position a
human in the loop above the autonomy stack, implying that manual
oversight provides a continuous safety net whenever the learned model
encounters unfamiliar operating conditions. In physical hardware, an
ungrounded supervisory role cannot dissipate kinetic energy or
counteract mechanical momentum. If a \(500\text{ kg}\) industrial
payload begins to oscillate while an automated crane translates at
\(1.0\text{ m/s}\), invoking human authority is meaningless unless that
authority maps directly to an explicit operational primitive. An
intervention primitive must define a concrete actuator response, such as
commanding a controlled deceleration profile of \(2.5\text{ m/s}^2\)
with a maximum nervous-system dispatch deadline of \(10\text{ ms}\).
Without a named operation and an enforceable time bound, the machine
drifts under inertia while software services attempt to resolve which
process holds the lock. Authority exists only when the low-level
controller can evaluate an incoming instruction, verify its operational
scope, and execute it before the physical plant breaches a safety
envelope.

Taking control of a machine already in motion forces the nervous system
to resolve two tightly coupled challenges, continuous physical
transition without mechanical shock and discrete state arbitration
without ambiguity. In continuous domain dynamics, abruptly zeroing motor
current or injecting a step change in steering angle damages planetary
gearboxes, saturates tire traction, or destabilizes an articulated
payload. The transition from autonomous trajectory tracking to manual
guidance requires what control engineering terms a \textbf{bumpless
transfer}, where the nervous system blends active reference
trajectories, ramps joint torques within calibrated jerk limits, and
matches initial actuator velocities before handing full command to the
operator. Simultaneously, in the discrete software domain, the
architecture must prevent split-brain execution where the learned model
and the human operator issue concurrent, conflicting setpoints. The
state arbiter must atomically revoke the model's execution token,
acknowledge the operator's override, and switch tracking multiplexers
within a single control period of \(1\text{ ms}\) on a
\(1000\text{ Hz}\) loop. A handover that drops setpoints mid-stride
creates an immediate physical hazard rather than a recovery.

Intervention events introduce a secondary failure mode inside the data
pipeline that trains subsequent model iterations. When an operator
intervenes to steer a vehicle away from a barrier or steady a slipping
gripper, the telemetry logger captures an abrupt discontinuity in the
state trajectory. If the training pipeline ingests these rescued
episodes as standard behavioral demonstrations without recording the
exact moment of truncation, imitation learning policies and
reinforcement algorithms learn to reproduce the failure states that
necessitated the rescue. The training distribution becomes biased toward
pre-failure trajectories, optimizing the policy to navigate toward the
boundary where the human took over because the subsequent human
demonstration contains high rewards. Preventing this policy corruption
requires strict data hygiene. The logging pipeline must timestamp the
precise millisecond of authority transfer, mark the pre-intervention
trajectory as a terminated failure, and isolate manual recovery
trajectories from standard policy training datasets.

An engineered intervention architecture spans the complete operational
lifecycle of authority, from the physical input device down to the motor
drivers and backward into the data pipelines. That lifecycle requires
unambiguous operational scoping so that operators command specific
mechanical primitives without disrupting unmanaged degrees of freedom,
deterministic state arbitration that resolves conflicting commands in
fixed time, continuous trajectory blending that protects mechanical
transmissions from shock, and authenticated communication paths that
reject corrupted override packets. It requires non-volatile evidence
logging that records which entity held authority during every
millisecond of a deployment, paired with data segregation pipelines that
prevent rescued runs from poisoning model retraining. A machine cannot
be safely operated under the authority of a learned model until the
systems engineer can prove how that authority is stripped away when
physical reality diverges from the plan.

\section{Human Authority Over Actuation}\label{sec-intervention-14-2}

System block diagrams frequently insert an icon of a human supervisor
with an arrow pointing toward the plant, treating human oversight as an
unexamined safety backstop. In the physical world, drawing a person in
the loop is not an authority model. Authority cannot attach to a generic
human role or an abstract supervision block. It must attach to specific
named operations with explicit timing budgets, physical degrees of
freedom, and deterministic preemption rules. An operational machine does
not simply grant an operator broad administrative authority. Instead,
the nervous system maintains an explicit capability matrix that maps
individual operations (such as an emergency power cutoff, a velocity
trim adjustment, a trajectory path veto, or full manual steering) to
authenticated roles.\sidenote{\footnotesize Sheridan and Verplank (1978) formulated
  the ten levels of supervisory control, showing human authority
  degrades into ambiguous hazard without explicit decision boundaries
  and timing contracts. In autonomous architectures (SAE J3016),
  supervisory authority decomposes into operational design domains with
  bounded fallback windows.} Each entry in this matrix defines the exact
actuator channels the operation governs, the preemption priority it
carries over the learned policy, and the maximum propagation latency
permitted from physical input to actuator response.

Naming the operations separates intervention from safety enforcement,
two mechanisms that system architectures routinely conflate. Enforcement
withholds permission from a candidate action proposed by a learned
model, modifying, clamping, or rejecting the setpoint while leaving the
machine under autonomous control. Intervention transfers authority away
from the model to an external entity, terminating autonomous operation
and placing actuation under direct supervisory command. An enforcement
filter sits in the fast control path to prevent the policy from
demanding an over-current maneuver or violating a barrier certificate
during nominal tracking. An intervention mechanism strips the policy of
its execution token entirely when an operational boundary is breached or
an operator demands control. A system equipped with enforcement but
lacking intervention cannot be redirected when human intent changes or
external conditions fall outside the policy design envelope. Conversely,
a system equipped with intervention but lacking enforcement permits a
lagging operator to drive the plant into structural failure during
nominal manual flight. A system with one mechanism and not the other is
incomplete in a way no amount of the other fixes.

When an operator attempts an intervention, human biology imposes a hard
real-time latency budget before any physical corrective action can
occur. Human response time is not a single point event, but a cascaded
pipeline of physiological and mechanical stages. In an analytical
breakdown of an operator monitoring a vehicle telemetry stream, the
sensory recognition stage requires \(150\text{ ms}\) to
\(300\text{ ms}\) for light to strike the retina, travel through the
visual pathway, and register as an optical anomaly. The situational
assessment stage, in which the operator recognizes that the learned
policy is drifting toward a hazard and selects an appropriate
counter-action, demands \(500\text{ ms}\) to \(1500\text{ ms}\).
Finally, the motor actuation stage, covering the transmission of motor
signals down the spinal cord, muscle recruitment in the hand, and the
physical displacement of an override switch, consumes \(200\text{ ms}\)
to \(400\text{ ms}\). Summing these sequential budgets establishes the
physiological \textbf{takeover reaction floor}, defined as
\(T_{\text{reaction}} = T_{\text{sensory}} + T_{\text{cognitive}} + T_{\text{motor}}\)
(\(850\text{--}2200\text{ ms}\) under active visual-manual monitoring,
and expanding beyond \(3500\text{ ms}\) when an operator is
out-of-the-loop), during which a moving machine traverses an unmonitored
blind drift distance \(d_{\text{drift}} = v \cdot T_{\text{reaction}}\)
under unchecked policy authority before physical braking or evasive
actuation can commence.

Throughout this reaction budget, the machine continues moving under the
unchecked authority of the learned policy. During the reaction interval
\(T_{\text{reaction}}\), a vehicle traveling at forward velocity \(v\)
covers a blind distance
\(d_{\text{human}} = v \cdot T_{\text{reaction}}\) before the
intervention signal reaches the hardware arbiter. In an analytical
scenario where a mobile machine travels at \(v = 5.0\text{ m/s}\) with a
moderate takeover budget of \(T_{\text{reaction}} = 1.5\text{ s}\), the
vehicle covers \(d_{\text{human}} = 7.5\text{ m}\) while executing the
unmonitored policy trajectory. Once the operator engages physical
braking, the vehicle requires an additional stopping distance governed
by its maximum deceleration \(a_{\text{max}}\). Integrating the constant
velocity phase during human reaction with the uniform deceleration phase
produces the minimum safe intervention distance,
\[d_{\text{intervene}} = v \cdot T_{\text{reaction}} + \frac{v^2}{2 a_{\text{max}}}.\]
For the same machine operating at \(5.0\text{ m/s}\) with a maximum
plant braking deceleration of \(a_{\text{max}} = 2.5\text{ m/s}^2\), the
physical braking phase requires an additional \(5.0\text{ m}\). The
total physical distance required for a human operator to detect an
obstacle, actuate an override, and bring the plant to rest is
\(12.5\text{ m}\), of which sixty percent is consumed entirely by human
neurological and muscular delay before the brakes engage. Because the
reaction distance scales linearly with velocity while the deceleration
distance scales quadratically, human intervention becomes physically
incapable of averting close-range collisions as operational speed
increases.

\protect\phantomsection\label{callout-mathux2a-14.2}
\begin{fbxSimple}{callout-math}{Fermi Math:}{Operator Takeover Budget and Spatial Clearance Margin}
\phantomsection\label{callout-math*-14.2}

\textbf{Question:} An autonomous transport vehicle
(\(m = 2200\text{ kg}\)) travels at \(v_0 = 22.0\text{ m/s}\)
(\(79.2\text{ km/h}\)) on a dry asphalt roadway (\(\mu = 0.80\), maximum
braking deceleration \(a_{\text{max}} = 6.0\text{ m/s}^2\)). An
attentive human safety driver monitors the forward path. If an unmapped
obstacle is perceived at a forward sensing horizon of
\(d_{\text{sense}} = 65.0\text{ m}\), can a human takeover avert
collision if the driver has an active reaction budget
\(T_{\text{react}} = 1.0\text{ s}\), the real-time safety enforcer
requires a four-phase handshake \(T_{\text{handshake}} = 20\text{ ms}\),
and the drive-by-wire steering/braking system uses a \(\mathcal{C}^2\)
bumpless blend window \(\tau_{\text{blend}} = 75\text{ ms}\)?

\textbf{1. Reaction Drift Distance
(\(T_{\text{react}} = 1.00\text{ s}\)):}
\[d_{\text{react}} = v_0 \cdot T_{\text{react}} = 22.0\text{ m/s} \times 1.00\text{ s} = 22.00\text{ m}.\]

\textbf{2. Handshake Protocol Drift Distance
(\(T_{\text{handshake}} = 0.020\text{ s}\)):}
\[d_{\text{handshake}} = v_0 \cdot T_{\text{handshake}} = 22.0\text{ m/s} \times 0.020\text{ s} = 0.44\text{ m}.\]

\textbf{3. Bumpless Actuator Blending Displacement
(\(\tau_{\text{blend}} = 0.075\text{ s}\)):} Over the \(\mathcal{C}^2\)
quintic spline blend, braking deceleration ramps from \(0\) to
\(-a_{\text{max}}\). Integrating
\(\ddot{x}(t) = -\alpha(t) a_{\text{max}}\) yields:
\[\Delta v_{\text{blend}} = \int_0^{\tau_{\text{blend}}} \alpha(t) a_{\text{max}} dt = \frac{1}{2} a_{\text{max}} \tau_{\text{blend}} = \frac{1}{2}(6.0)(0.075) = 0.225\text{ m/s},\]
\[v_{\text{post-blend}} = 22.0 - 0.225 = 21.775\text{ m/s},\]
\[\begin{aligned}
d_{\text{blend}} &= v_0 \tau_{\text{blend}} - 0.15 a_{\text{max}} \tau_{\text{blend}}^2 \\
&= (22.0)(0.075) - (0.15)(6.0)(0.075)^2 = 1.650 - 0.005 = 1.645\text{ m}.
\end{aligned}\]

\textbf{4. Full Deceleration Braking Distance:}
\[d_{\text{brake}} = \frac{v_{\text{post-blend}}^2}{2 a_{\text{max}}} = \frac{(21.775\text{ m/s})^2}{2 \times 6.0\text{ m/s}^2} = \frac{474.15}{12.0} = 39.51\text{ m}.\]

\textbf{5. Total Stopping Distance vs.~Available Clearance:}
\[d_{\text{total}} = d_{\text{react}} + d_{\text{handshake}} + d_{\text{blend}} + d_{\text{brake}} = 22.00 + 0.44 + 1.65 + 39.51 = 63.60\text{ m}.\]
\[\Delta d_{\text{clearance}} = d_{\text{sense}} - d_{\text{total}} = 65.00\text{ m} - 63.60\text{ m} = +1.40\text{ m}.\]

\textbf{Verdict:} The vehicle stops with a razor-thin safety clearance
margin of \(1.40\text{ m}\). If human situational awareness decays by
even \(100\text{ ms}\) (increasing \(T_{\text{react}}\) to
\(1.10\text{ s}\)), the vehicle travels an additional \(2.20\text{ m}\),
striking the obstacle at \(v_{\text{impact}} = 5.1\text{ m/s}\)
(\(18.4\text{ km/h}\)).

\end{fbxSimple}

\protect\phantomsection\label{callout-ruleux2a-14.3}
\begin{fbxSimple}{callout-rule}{Architect’s Rule of Thumb:}{The Spatial Takeover Budget Law}
\phantomsection\label{callout-rule*-14.3}

Handover is an allocation of physical distance, not digital clock
cycles. If the spatial travel during human reaction time
\(v \cdot T_{\text{reaction}}\) plus the kinematic braking distance
\(\frac{v^2}{2 a_{\text{max}}}\) exceeds the sensing clearance margin,
no software handshake or arbitration algorithm can avert collision. The
handover deadline \(T_{\text{timeout}}\) must be dynamically bounded by
physical stopping distance, not fixed in configuration files.

\end{fbxSimple}

As quantified across intervention modalities in
Table~\ref{tbl-14-teleop-stability}, this latency penalty scales
directly with communication topology and operator cognitive engagement.
In high-bandwidth bilateral haptic teleoperation, proprioceptive
reflexes and deterministic local buses bound closed-loop delay
\(T_{\Sigma}\) to under \(185\text{ ms}\), preserving dynamic passivity.
In contrast, remote supervisory dispatch and out-of-the-loop takeover
inflate total reaction delay \(T_{\Sigma}\) to upwards of
\(3700\text{ ms}\), expanding the stopping drift envelope
\(d_{\text{drift}}\) to over \(17.5\text{ m}\) at nominal speed. When
communication or cognitive delays exceed the plant's dynamic stability
margin, human manual commands induce severe phase lag and pilot-induced
oscillations, requiring local autonomous safety interceptors to clamp
velocity and enforce barrier certificates.

\begin{longtable}[]{@{}
  >{\raggedright\arraybackslash}p{(\linewidth - 12\tabcolsep) * \real{0.1429}}
  >{\raggedright\arraybackslash}p{(\linewidth - 12\tabcolsep) * \real{0.1429}}
  >{\raggedright\arraybackslash}p{(\linewidth - 12\tabcolsep) * \real{0.1429}}
  >{\raggedright\arraybackslash}p{(\linewidth - 12\tabcolsep) * \real{0.1429}}
  >{\raggedright\arraybackslash}p{(\linewidth - 12\tabcolsep) * \real{0.1429}}
  >{\raggedright\arraybackslash}p{(\linewidth - 12\tabcolsep) * \real{0.1429}}
  >{\raggedright\arraybackslash}p{(\linewidth - 12\tabcolsep) * \real{0.1429}}@{}}
\caption{\textbf{Teleoperation Latency, Human Neuromuscular Delay, and
Dynamic Control Stability Bounds.} Quantitative timing and stability
budget across supervisory, teleoperated, and manual intervention
modalities. Total closed-loop delay
\(T_{\Sigma} = T_{\text{RTT}} + T_{\text{human}} + T_{\text{act}}\)
combines network transport, physiological response (sensory, cognitive,
motor), and actuator mechanical rise time
(\(T_{\text{act}} \approx 10\text{--}60\text{ ms}\)). Stopping drift
\(d_{\text{drift}} = v_0 \cdot T_{\Sigma} + v_0^2 / (2 a_{\text{max}})\)
is evaluated at nominal operational velocity \(v_0 = 5.0\text{ m/s}\)
(\(18.0\text{ km/h}\)) under maximum deceleration
\(a_{\text{max}} = 2.5\text{ m/s}^2\). As network and cognitive
latencies grow, manual guidance loses dynamic passivity and induces
pilot-induced oscillations, requiring local autonomous safety
interceptors to enforce kinematic barrier certificates independently of
remote operator input.}\label{tbl-14-teleop-stability}\tabularnewline
\toprule\noalign{}
\begin{minipage}[b]{\linewidth}\raggedright
Control Modality \& Operational Tier
\end{minipage} & \begin{minipage}[b]{\linewidth}\raggedright
Network \& Transport Latency (\(T_{\text{RTT}}\))
\end{minipage} & \begin{minipage}[b]{\linewidth}\raggedright
Neuromuscular \& Cognitive Budget (\(T_{\text{human}}\))
\end{minipage} & \begin{minipage}[b]{\linewidth}\raggedright
Total Closed-Loop Delay (\(T_{\Sigma}\))
\end{minipage} & \begin{minipage}[b]{\linewidth}\raggedright
Stopping Drift (\(d_{\text{drift}}\) @ \(5\text{ m/s}\))
\end{minipage} & \begin{minipage}[b]{\linewidth}\raggedright
Dynamic Instability \& Failure Mode
\end{minipage} & \begin{minipage}[b]{\linewidth}\raggedright
Hardware \& Safety Interceptor
\end{minipage} \\
\midrule\noalign{}
\endfirsthead
\toprule\noalign{}
\begin{minipage}[b]{\linewidth}\raggedright
Control Modality \& Operational Tier
\end{minipage} & \begin{minipage}[b]{\linewidth}\raggedright
Network \& Transport Latency (\(T_{\text{RTT}}\))
\end{minipage} & \begin{minipage}[b]{\linewidth}\raggedright
Neuromuscular \& Cognitive Budget (\(T_{\text{human}}\))
\end{minipage} & \begin{minipage}[b]{\linewidth}\raggedright
Total Closed-Loop Delay (\(T_{\Sigma}\))
\end{minipage} & \begin{minipage}[b]{\linewidth}\raggedright
Stopping Drift (\(d_{\text{drift}}\) @ \(5\text{ m/s}\))
\end{minipage} & \begin{minipage}[b]{\linewidth}\raggedright
Dynamic Instability \& Failure Mode
\end{minipage} & \begin{minipage}[b]{\linewidth}\raggedright
Hardware \& Safety Interceptor
\end{minipage} \\
\midrule\noalign{}
\endhead
\bottomrule\noalign{}
\endlastfoot
\textbf{Direct Bilateral Teleoperation} (Haptic / Force-Reflective) &
\(\le 5\text{--}20\text{ ms}\) (Local deterministic bus / 5G URLLC) &
\(100\text{--}150\text{ ms}\) (Proprioceptive reflex \& tactile
feedback) & \(110\text{--}185\text{ ms}\)
(\(f_{\text{ctrl}} \ge 500\text{ Hz}\)) & \(5.75\text{ m}\)
(\(0.75\text{ m}\) drift \(+\) \(5.0\text{ m}\) brake) & Loss of
passivity, pilot-induced oscillations (PIO), high-frequency haptic
shudder & Time-domain passivity observer (TDPO), scattering wave
variables, motor current clamps \\
\textbf{Continuous Direct Velocity Guidance} (Rate-Assisted Steering /
Joystick) & \(30\text{--}80\text{ ms}\) (Private wireless / 5G SA edge
slice) & \(250\text{--}450\text{ ms}\) (Active visual-manual tracking) &
\(300\text{--}570\text{ ms}\)
(\(f_{\text{ctrl}} \approx 50\text{--}100\text{ Hz}\)) &
\(7.25\text{ m}\) (\(2.25\text{ m}\) drift \(+\) \(5.0\text{ m}\) brake)
& Phase lag instability, steering hunting limit cycles
(\(2\text{--}4\text{ Hz}\)), lateral overshoot & Onboard velocity
envelope limiter, control barrier certificates, deadman lease timeout
(\(100\text{ ms}\)) \\
\textbf{Supervisory Waypoint \& Corridor Dispatch} (Shared Autonomy /
Path Veto) & \(100\text{--}250\text{ ms}\) (Public cellular / IP broker
link) & \(600\text{--}1200\text{ ms}\) (Situational assessment \&
re-planning) & \(750\text{--}1550\text{ ms}\)
(\(f_{\text{ctrl}} \approx 5\text{--}10\text{ Hz}\)) &
\(10.75\text{ m}\) (\(5.75\text{ m}\) drift \(+\) \(5.0\text{ m}\)
brake) & Stale corridor execution, spatial collision during uplink
delay, command starvation freeze & Local receding-horizon obstacle veto,
time-to-collision (TTC) governor, autonomous fallback \\
\textbf{Out-of-the-Loop Takeover} (Passive Monitor / Autonomous Rescue)
& \(50\text{--}150\text{ ms}\) (Cabin bus / near-edge telemetry) &
\(1500\text{--}3500\text{ ms}\) (Gaze re-engagement \& cognitive
recovery) & \(1580\text{--}3710\text{ ms}\)
(\(f_{\text{ctrl}} \le 1\text{ Hz}\)) & \(17.50\text{ m}\)
(\(12.5\text{ m}\) drift \(+\) \(5.0\text{ m}\) brake) & Unmonitored
blind trajectory drift, panic over-correction, mechanical column lockup
& Monotonic Category 1 stop upon \(T_{\text{timeout}}\) expiry, DMS
gaze-decay alarm, torque rate limiter \\
\textbf{Direct Hardware Emergency Stop} (Physical Switch / Safety
Interlock) & \(\le 1\text{ ms}\) (Hardwired dual-channel safety line) &
\(150\text{--}300\text{ ms}\) (Reflexive palm slap / deadman release) &
\(160\text{--}325\text{ ms}\) (\(f_{\text{ctrl}} \ge 1000\text{ Hz}\)) &
\(6.25\text{ m}\) (\(1.25\text{ m}\) drift \(+\) \(5.0\text{ m}\) brake)
& Infinite jerk shock, tire adhesion breakaway, mechanical driveline
overload & Safe Torque Off (STO / SS1 ramp), dynamic friction braking,
spring-applied holding brakes \\
\end{longtable}

Because human response is slow and prone to interruption, manual
authority cannot be granted as an open-ended state.\sidenote{\footnotesize NTSB Report
  HAR-19/03 (\citeproc{ref-ntsb2019uber}{National Transportation Safety
  Board 2019}) revealed that disabling the factory autonomous emergency
  braking system eliminated the deterministic physical safety enforcer,
  relying on unmonitored human supervision.} If an operator assumes
control and subsequently experiences wireless packet loss, physical
incapacitation, or cognitive distraction, the machine cannot remain
adrift under stale manual setpoints. The architecture protects against
operator dropouts by bounding manual authority with strict temporal
leases and hardware deadman switches
(Figure~\ref{fig-14-teach-pendant}). In a lease-based communication
model, the vehicle nervous system issues a bounded delegation window of
\(50\text{ ms}\) to \(100\text{ ms}\). The operator console must
continuously stream authenticated keep-alive packets to renew the lease
before the deadline expires. If a single lease window lapses without
renewal, or if an operator releases a spring-loaded physical deadman
switch on the controller, manual authority terminates immediately, and
the low-level nervous system transitions the actuators into a
deterministic safe stop.\sidenote{\footnotesize ISO 13849-1 (PLe) and ISO 10218-1
  mandate three-position enabling devices for manual robot control.
  Releasing (Position 1) or over-squeezing in panic (Position 3)
  activates dual-channel safety contacts, cutting power via Safe Torque
  Off (STO) to prevent involuntary muscle spasms holding motion
  channels.} Authority is not a persistent condition. It is a perishable
resource that requires active, periodic renewal to survive.

\begin{figure}[t]

\centering{

\includegraphics[width=0.85\linewidth,height=\textheight,keepaspectratio]{/Users/VJ/GitHub/MLSysBook-integrate-physical/physical/book/chapters/14-intervention/figures/fig14_real_teach_pendant.jpg}

}

\caption{\label{fig-14-teach-pendant}\textbf{Industrial Robot Teach
Pendant and Physical Intervention Hardware.} Handheld teach pendant (ABB
IRB series) incorporating multi-layered physical authority mechanisms: a
key-lockable mode selector enforcing hardware mutual exclusion between
Automatic and Manual (T1/T2) states, a hardwired mushroom-head Emergency
Stop (E-Stop) button wired directly into dual-channel safety relays to
de-energize actuator contactors, a 3-axis proportional joystick for
manual coordinate jogging, and an ergonomic rear grip housing a
three-position enabling device (deadman switch). Under ISO 10218-1
safety standards, manual authority requires continuous holding in the
center position; releasing the grip or panicking with an over-squeeze
immediately drops the safety lease and forces an actuator Category 1
safe stop.}

\end{figure}%

Even while an operator holds an active lease, human authority remains
subordinate to the kinematic and dynamic limits enforced at the hardware
boundary. A common engineering mistake assumes that an emergency manual
override must bypass all downstream safety filters to guarantee direct,
unmediated actuator control. Bypassing safety enforcement at the motor
level creates severe mechanical failure modes. If an operator panics and
commands a steering rate that induces vehicle rollover, or requests a
joint torque that strips the gear teeth of a robotic actuator, the
low-level safety enforcer sitting at the hardware floor rejects the
command. The enforcer evaluates human setpoints against the same control
barrier certificates, current limits, and kinematic floor boundaries
applied to the learned model. Human intervention transfers operational
discretion away from the learned policy, but it cannot suspend the laws
of physics or dismantle the hardware-enforced envelopes that preserve
the physical integrity of the machine.

\section{Authority States and Transitions}\label{sec-intervention-14-3}

Control authority at any instant must be a discrete state, never an
inference or a mathematical blend. When an engineering team attempts to
soften transitions by continuously averaging the control vectors of an
autonomous policy and a human operator, the resulting command represents
neither intent. If a human steers left to avoid a ditch while a neural
network policy steers right to follow a lane marker, a linear
interpolation between the two commands directs the steering angle
straight forward, driving the vehicle into the obstacle at full speed.
This authority conflict was tragically demonstrated in the Boeing 737
MAX MCAS flight control failures (\citeproc{ref-jatr2019boeing}{Joint
Authorities Technical Review 2019}), where automated stabilizer trim
repeatedly overrode pilot control column inputs based on a single faulty
sensor without explicit, unambiguous handover states.\sidenote{\footnotesize Airbus
  fly-by-wire sums stick inputs algebraically unless a takeover button
  latches priority, whereas Boeing back-drives columns with active
  servos. In multi-brain physical AI, inter-core authority multiplexing
  relies on hardware semaphores and memory barriers (\texttt{DMB\ ISH})
  to ensure atomic ownership transfers across control ticks.} To prevent
this ambiguity, control authority at every discrete time tick \(t\) must
resolve to a single, mutually exclusive state
\(S_t \in \{\text{Autonomous}, \text{Handover-Pending}, \text{Human-Manual}, \text{Degraded-Fallback}\}\).
The nervous system evaluates arbitration logic, determines the active
state, and commits that identity to the bus before dispatching setpoints
to low-level motor controllers. At no point does the system compute a
weighted average of conflicting masters or leave ownership open to
statistical interpretation.

The complete state transition logic, guard preconditions, execution time
bounds (\(T_{\text{trans}}\)), and bumpless actuator blending profiles
governing this arbiter are formalized in
Table~\ref{tbl-14-authority-state-matrix}. Every edge in the state
machine defines an unambiguous trigger condition, mutual exclusion
enforcement across actuator channels, and a deterministic fallback path
should an active communication lease lapse.

\begin{longtable}[]{@{}
  >{\raggedright\arraybackslash}p{(\linewidth - 12\tabcolsep) * \real{0.1429}}
  >{\raggedright\arraybackslash}p{(\linewidth - 12\tabcolsep) * \real{0.1429}}
  >{\raggedright\arraybackslash}p{(\linewidth - 12\tabcolsep) * \real{0.1429}}
  >{\raggedright\arraybackslash}p{(\linewidth - 12\tabcolsep) * \real{0.1429}}
  >{\raggedright\arraybackslash}p{(\linewidth - 12\tabcolsep) * \real{0.1429}}
  >{\raggedright\arraybackslash}p{(\linewidth - 12\tabcolsep) * \real{0.1429}}
  >{\raggedright\arraybackslash}p{(\linewidth - 12\tabcolsep) * \real{0.1429}}@{}}
\caption{\textbf{Authority State Machine Transition Matrix.}
Deterministic transition matrix for the real-time nervous system arbiter
across discrete authority states
\(S_t \in \{\text{Autonomous}, \text{Handover-Pending}, \text{Human-Manual}, \text{Degraded-Fallback}, \text{Emergency-Stop}\}\).
Each entry defines the entry guard condition, Worst-Case Execution Time
(WCET) or protocol latency bound \(T_{\text{trans}}\), channel
arbitration rule, bumpless actuator blending profile, and fault recovery
path. All downward transitions toward higher safety containment operate
under a strict monotonic constraint, preventing upward state
re-engagement while the machine remains in
motion.}\label{tbl-14-authority-state-matrix}\tabularnewline
\toprule\noalign{}
\begin{minipage}[b]{\linewidth}\raggedright
Source State (\(S_{\text{src}}\))
\end{minipage} & \begin{minipage}[b]{\linewidth}\raggedright
Target State (\(S_{\text{tgt}}\))
\end{minipage} & \begin{minipage}[b]{\linewidth}\raggedright
Transition Trigger \& Guard Condition
\end{minipage} & \begin{minipage}[b]{\linewidth}\raggedright
Latency Bound (\(T_{\text{trans}}\) / WCET)
\end{minipage} & \begin{minipage}[b]{\linewidth}\raggedright
Handshake \& Arbitration Policy
\end{minipage} & \begin{minipage}[b]{\linewidth}\raggedright
Bumpless Transfer \& Dynamic Profile
\end{minipage} & \begin{minipage}[b]{\linewidth}\raggedright
Fallback \& Fault Mitigation
\end{minipage} \\
\midrule\noalign{}
\endfirsthead
\toprule\noalign{}
\begin{minipage}[b]{\linewidth}\raggedright
Source State (\(S_{\text{src}}\))
\end{minipage} & \begin{minipage}[b]{\linewidth}\raggedright
Target State (\(S_{\text{tgt}}\))
\end{minipage} & \begin{minipage}[b]{\linewidth}\raggedright
Transition Trigger \& Guard Condition
\end{minipage} & \begin{minipage}[b]{\linewidth}\raggedright
Latency Bound (\(T_{\text{trans}}\) / WCET)
\end{minipage} & \begin{minipage}[b]{\linewidth}\raggedright
Handshake \& Arbitration Policy
\end{minipage} & \begin{minipage}[b]{\linewidth}\raggedright
Bumpless Transfer \& Dynamic Profile
\end{minipage} & \begin{minipage}[b]{\linewidth}\raggedright
Fallback \& Fault Mitigation
\end{minipage} \\
\midrule\noalign{}
\endhead
\bottomrule\noalign{}
\endlastfoot
\textbf{Autonomous} & \textbf{Human-Manual} & Manual column torque
\(|\tau_{\text{manual}}| > 4.0\text{ Nm}\) or authenticated override
button assert & \(\le 1.0\text{ ms}\) (1 tick @ \(1000\text{ Hz}\)) &
Policy execution token atomically revoked; human granted exclusive write
lock & \(\mathcal{C}^2\) quintic spline blend
(\(\tau_{\text{blend}} = 50\text{ ms}\)); manual integrators pre-biased
& Drop to \textbf{Degraded-Fallback} if manual keep-alive lease lapses
(\(> 100\text{ ms}\)) \\
\textbf{Autonomous} & \textbf{Handover-Pending} & Perception confidence
\(\gamma < \gamma_{\text{min}}\), ODD geofence boundary, or non-critical
sensor fault & \(\le 10.0\text{ ms}\) (Takeover broadcast frame) &
Policy maintains nominal trajectory holding while awaiting operator ACK
& Nominal tracking hold; velocity rate-clamping enabled; acoustic/visual
alerts active & If unacknowledged within
\(T_{\text{timeout}} = 500\text{ ms}\), transition to
\textbf{Degraded-Fallback} \\
\textbf{Handover-Pending} & \textbf{Human-Manual} & Operator capacitive
grip confirmed AND valid cryptographic token received &
\(\le 20.0\text{ ms}\) (4-phase protocol round-trip) & Human granted
exclusive channel authority; autonomous policy locked out of bus &
\(\mathcal{C}^2\) quintic spline blend
(\(\tau_{\text{blend}} = 75\text{ ms}\)) over active actuator reference
& Abort transfer on token signature failure; raise security alert to
supervisor \\
\textbf{Handover-Pending} & \textbf{Degraded-Fallback} & Handover timer
expired (\(t > T_{\text{timeout}}\)) with zero operator acknowledgment &
\(\le 1.0\text{ ms}\) (Timer interrupt dispatch) & Autonomous policy and
operator both locked out; nervous system asserts control & Deterministic
Category 1 stop (\(a_{\text{decel}} = -2.5\text{ m/s}^2\)); stability
assist engaged & Monotonic safety lockout; upward transition prohibited
until full standstill \\
\textbf{Human-Manual} & \textbf{Degraded-Fallback} & Deadman release
(\(> 100\text{ ms}\)), DMS gaze loss (\(> 2.0\text{ s}\)), or link loss
(\(> 50\text{ ms}\)) & \(\le 5.0\text{ ms}\) (Watchdog lease expiration)
& Human authority revoked; hardware multiplexer switches to fallback
profile & Controlled deceleration ramp to zero velocity; hazard flashers
energized & Vehicle brought to full stop; requires manual physical
key-cycle to re-arm \\
\textbf{Any Active State} & \textbf{Emergency-Stop} & Hardwired E-stop
line open, critical CBF breach (\(h(\mathbf{x}) < 0\)), or power rail
fault & \(\le 0.1\text{ ms}\) (Hardware interrupt / gate cutoff) &
Unconditional hardware preemption; bypasses software bus and
microcontrollers & Safe Torque Off (STO) / mechanical spring holding
brake actuation (\(j \to \infty\)) & Actuator drive de-energized;
non-volatile freeze of pre-trigger evidence ring buffer \\
\end{longtable}

Shifting authority between these discrete states requires a
deterministic \textbf{four-phase handshake protocol} (\texttt{Request}
\(\to\) \texttt{Acknowledge} \(\to\) \texttt{Commit} \(\to\)
\texttt{Confirm}). Executed by the real-time safety enforcer, this
four-stage transaction guarantees mutual exclusion, validates
cryptographic freshness, and atomically shifts actuation authority
between discrete holders across a synchronous tick boundary without
split-brain command contention. In the first phase, designated as the
\emph{Request} phase, one entity signals an intent to transfer control.
This occurs either when an operator asserts an override through a
physical input like a button press or steering torque exceeding a
calibrated threshold, or when the autonomous policy detects that sensor
confidence has degraded and requests human assumption of control. In the
second phase, the \emph{Acknowledge} phase, the safety enforcer inspects
the target state prerequisites, verifying that sensor data is valid,
communications are active, and the recipient is physically capable of
receiving command authority. In the third phase, the \emph{Commit}
phase, the enforcer executes the state transition at a synchronous
control tick boundary, mapping the target input source to the actuator
queue and locking the relinquishing source out of the bus. In the final
phase, the \emph{Confirm} phase, the nervous system broadcasts a state
confirmation packet across the network and triggers physical feedback,
such as haptic vibration or optical indicators, to verify to the
operator that the transfer of authority is complete.

The protocol primitives, timing deadlines, cryptographic verification
artifacts, and fallback actions for each stage of this four-phase
sequence are specified in Table~\ref{tbl-14-handshake-protocol}. By
structuring authority transfer as an atomic, multi-stage transaction
verified by the real-time safety enforcer, the architecture prevents
split-brain command contention, rejects spoofed override injections, and
guarantees that interrupted handovers fail safely into deterministic
holding states.

\begin{longtable}[]{@{}
  >{\raggedright\arraybackslash}p{(\linewidth - 12\tabcolsep) * \real{0.1429}}
  >{\raggedright\arraybackslash}p{(\linewidth - 12\tabcolsep) * \real{0.1429}}
  >{\raggedright\arraybackslash}p{(\linewidth - 12\tabcolsep) * \real{0.1429}}
  >{\raggedright\arraybackslash}p{(\linewidth - 12\tabcolsep) * \real{0.1429}}
  >{\raggedright\arraybackslash}p{(\linewidth - 12\tabcolsep) * \real{0.1429}}
  >{\raggedright\arraybackslash}p{(\linewidth - 12\tabcolsep) * \real{0.1429}}
  >{\raggedright\arraybackslash}p{(\linewidth - 12\tabcolsep) * \real{0.1429}}@{}}
\caption{\textbf{Four-Phase Authority Handshake Protocol Timing and
Verification Contract.} Real-time communication and verification
specification governing the synchronous transfer of actuation authority.
The protocol decomposes into four deterministic phases executed across
the isolated safety enforcer microcontroller and distributed actuator
motor drives. Cumulative protocol execution budget
\(T_{\text{handshake}} = t_1 + t_2 + t_3 + t_4 \le 31.0\text{ ms}\)
guarantees deterministic preemption within two \(50\text{ Hz}\) motion
control intervals, while cryptographic freshness validation and atomic
mailbox multiplexing eliminate split-brain command
execution.}\label{tbl-14-handshake-protocol}\tabularnewline
\toprule\noalign{}
\begin{minipage}[b]{\linewidth}\raggedright
Handshake Phase
\end{minipage} & \begin{minipage}[b]{\linewidth}\raggedright
Protocol Primitive \& Direction
\end{minipage} & \begin{minipage}[b]{\linewidth}\raggedright
Executing Entity
\end{minipage} & \begin{minipage}[b]{\linewidth}\raggedright
Timing Bound (\(t_{\text{phase}}\))
\end{minipage} & \begin{minipage}[b]{\linewidth}\raggedright
Safety Invariant \& Precondition
\end{minipage} & \begin{minipage}[b]{\linewidth}\raggedright
Cryptographic \& Telemetry Artifact
\end{minipage} & \begin{minipage}[b]{\linewidth}\raggedright
Fault \& Anomaly Mitigation
\end{minipage} \\
\midrule\noalign{}
\endfirsthead
\toprule\noalign{}
\begin{minipage}[b]{\linewidth}\raggedright
Handshake Phase
\end{minipage} & \begin{minipage}[b]{\linewidth}\raggedright
Protocol Primitive \& Direction
\end{minipage} & \begin{minipage}[b]{\linewidth}\raggedright
Executing Entity
\end{minipage} & \begin{minipage}[b]{\linewidth}\raggedright
Timing Bound (\(t_{\text{phase}}\))
\end{minipage} & \begin{minipage}[b]{\linewidth}\raggedright
Safety Invariant \& Precondition
\end{minipage} & \begin{minipage}[b]{\linewidth}\raggedright
Cryptographic \& Telemetry Artifact
\end{minipage} & \begin{minipage}[b]{\linewidth}\raggedright
Fault \& Anomaly Mitigation
\end{minipage} \\
\midrule\noalign{}
\endhead
\bottomrule\noalign{}
\endlastfoot
\textbf{Phase 1: Request (Propose)} & \texttt{REQ\_TAKEOVER} (Operator
\(\to\) Enforcer or Policy \(\to\) Operator) & Operator Console or
Policy Planner & \(t_1 \le 10.0\text{ ms}\) (CAN-FD queue + tx) & PTP
timestamp freshness \(|\Delta t_{\text{PTP}}| \le \tau_{\text{fresh}}\)
(\(50\text{ ms}\)); entity authenticated in capability matrix & 64-bit
IEEE 1588 timestamp \(+\) 64-bit monotonic nonce \(+\) AES-128-GMAC tag
& Reject stale/forged packet; latch bus error; increment intrusion
counter \\
\textbf{Phase 2: Acknowledge (Verify)} & \texttt{ACK\_TAKEOVER}
(Enforcer \(\to\) Operator \& Arbiter) & Real-Time Safety Enforcer MCU
(Cortex-R5) & \(t_2 \le 5.0\text{ ms}\) (Crypto check + check
invariants) & Kinematic state
\(\mathbf{x}(t) \in \mathcal{S}_{\text{safe}}\); manual controller
integrators pre-biased
(\(\mathbf{u}_{\text{human}} \approx \mathbf{u}_{\text{plant}}\)) &
Ephemeral 32-bit transfer token \(+\) HMAC-SHA256 enforcer ticket & If
plant outside safe envelope, reject request and command degraded holding
mode \\
\textbf{Phase 3: Commit (Switch)} & \texttt{COMMIT\_TRANSFER} (Enforcer
\(\to\) Actuator Inverters) & Hardware Arbiter / Inverter Multiplexer &
\(t_3 \le 1.0\text{ ms}\) (Synchronous tick boundary) & Atomic register
swap; mutually exclusive write lock; active authority weights
\(\sum \alpha_i = 1.0\) & Serialized evidence tuple
\(\langle t_{\text{PTP}}, S_{\text{auth}}, \mathbf{x}, \mathbf{u}_{\text{prop}}, \mathbf{u}_{\text{enf}}, \mathbf{u}_{\text{act}}\rangle\)
to SRAM buffer & If PTP synchronization error \(> 100\,\mu\text{s}\),
abort switch and trigger Category 1 stop \\
\textbf{Phase 4: Confirm (Feedback)} & \texttt{CONFIRM\_ACTIVE}
(Enforcer \(\to\) HMI \& Fleet Log) & Nervous System Bus \& HMI
Controller & \(t_4 \le 15.0\text{ ms}\) (HMI indicator + haptic pulse) &
Closed-loop telemetry echoes active authority state
\(S_{\text{auth}} = S_{\text{tgt}}\); manual keep-alive lease
initialized (\(100\text{ ms}\)) & Chained SHA-256 log block root \(+\)
tactile haptic pattern dispatched to operator controls & If confirm
packet lost, sound acoustic alert; operator lease watchdog active \\
\end{longtable}

Because human attention is volatile, a handover request initiated by the
autonomous policy cannot remain open indefinitely. If a machine
encounters an operational design domain boundary at speed and requests
intervention, an operator who is distracted, asleep, or incapacitated
will fail to complete the handshake. The state machine enforces an
unacknowledged takeover deadline governed by a hard timeout
\(T_{\text{timeout}}\). As an analytical example, consider a machine
operating on a \(100\text{ Hz}\) control loop (\(10\text{ ms}\) tick
period) with a configured handover timeout
\(T_{\text{timeout}} = 500\text{ ms}\). If fifty consecutive control
cycles elapse in the \(\text{Handover-Pending}\) state without a valid
acknowledge signal from the operator, the nervous system terminates the
handshake. The state machine transitions immediately to
\(\text{Degraded-Fallback}\), executing an automated Category 1 stop
that brings the plant to rest under controlled deceleration while
maintaining actuator power and stability augmentation.

Once authority resides in \(\text{Human-Manual}\), the machine must
continuously verify that an operator remains actively engaged. Passive
occupancy is insufficient. The nervous system monitors engagement
through physical sensor streams, including capacitive touch sensors in
the steering wheel or control grips, optical gaze tracking directed at
the operator, and continuous torque sensing on the steering column. If
these sensors detect a loss of operator presence, such as hands leaving
the controls for longer than a delegation lease of \(100\text{ ms}\),
the system initiates an automated fallback ladder. Fallback transitions
operate under a strict monotonic constraint toward increasing safety
containment. Authority falls from \(\text{Human-Manual}\) to
\(\text{Degraded-Fallback}\), and the state machine locks out upward
transitions back into \(\text{Autonomous}\) or \(\text{Human-Manual}\)
while the machine remains in motion. Resuming normal operation requires
bringing the plant to a full stop, clearing the fault log, and executing
a complete handshake re-initialization sequence.

When multiple control sources generate commands simultaneously, the
nervous system resolves the contention through a fixed arbitration
hierarchy evaluated every cycle. At the top of the hierarchy sits the
low-level safety enforcer, which holds absolute authority to clamp,
modify, or reject any command that violates kinematic boundaries,
thermal thresholds, or control barrier certificates. Below the safety
enforcer sits the human operator, whose manual setpoints unconditionally
supersede learned policy proposals whenever the system is in
\(\text{Human-Manual}\) and an active lease is held. At the bottom of
the hierarchy sits the learned policy, whose outputs execute only when
the state is explicitly \(\text{Autonomous}\) and the safety enforcer
grants admission. This deterministic hierarchy guarantees that authority
conflicts resolve in \(O(1)\) time on the real-time controller,
eliminating arbitration deadlock without relying on consensus heuristics
or online negotiation between software layers.

As synthesized in Figure~\ref{fig-14-authority-handshake-fsm}, the
complete runtime intervention interface integrates the four-phase
synchronous handshake protocol pipeline, the five discrete authority
states with deterministic timeout and monotonic fallback guards, and the
three-tier preemption priority hierarchy backed by an immutable 64-byte
evidence audit tuple.

\begin{figure}[t]

\centering{

\includegraphics[width=0.95\linewidth,height=\textheight,keepaspectratio]{index_files/mediabag//Users/VJ/GitHub/MLSysBook-integrate-physical/physical/book/chapters/14-intervention/figures/fig14_authority_handshake_fsm.pdf}

}

\caption{\label{fig-14-authority-handshake-fsm}\textbf{Four-Phase
Authority Handshake Protocol, State Transition Arbiter, and Multi-Tier
Priority Stack.} Top: Synchronous four-phase handshake pipeline
decomposing authority transfer into Request (operator override torque
\(|\tau| > 4.0\text{ N}\cdot\text{m}\) or policy ODD departure),
Intercept \& Auth (\(1\text{ kHz}\) microcontroller hardware token and
freshness validation \(\Delta t \le 100\text{ ms}\)), Commit \& Blend
(atomic register swap at \(1\text{ ms}\) tick with \(\mathcal{C}^2\)
quintic blending \(\alpha(t)\) and integrator pre-biasing), and Confirm
\& Active (\(\alpha = 1.0\) active human command under continuous
deadman lease \(T_{\text{lease}} \le 100\text{ ms}\)). Bottom left:
Discrete state machine across
\(S_t \in \{\text{Autonomous}, \text{Handover-Pending}, \text{Bumpless-Blend}, \text{Human-Manual}, \text{Degraded-Fallback}\}\)
showing unacknowledged timeout (\(T_{\text{timeout}} = 500\text{ ms}\))
and monotonic safe stop edges. Bottom right: Real-time three-tier
arbitration hierarchy (Tier 1 safety barrier enforcer
\(h(\mathbf{x}) \ge 0\) overriding Tier 2 human manual, overriding Tier
3 learned policy) paired with a \(64\text{-byte}\) IEEE 1588 timestamped
SHA-256 HMAC evidence audit tuple. Architectural diagram modeled after
ASIL D / SIL 3 multi-core arbitration architectures.}

\end{figure}%

While a discrete state machine guarantees that authority is unambiguous
at every control tick, instantaneous switching between independent
controllers introduces a distinct physical challenge. If an autonomous
policy is commanding \(15.0\text{ Nm}\) of drive torque to accelerate up
an incline at the exact millisecond the commit phase routes manual input
to the motor controllers, and the operator's physical throttle is at
rest, commanding \(0.0\text{ Nm}\), the actuator experiences a step
discontinuity in commanded torque. That sudden step produces mechanical
jerk, driveline shock, and potential loss of surface traction. Defining
discrete authority states is therefore only the first half of the
intervention mechanism. The second half requires executing the switch
across control boundaries without injecting discontinuous forces into
the moving body.

\section{Transfer Without Discontinuity}\label{sec-intervention-14-4}

When authority transfers between controllers at time \(t_0\), the
learned policy outputs a command vector
\(\mathbf{u}_{\text{policy}}(t_0)\) while the incoming human input
device provides \(\mathbf{u}_{\text{human}}(t_0)\). Switching
instantaneously between these independent command streams introduces a
step discontinuity
\(\Delta \mathbf{u} = \mathbf{u}_{\text{human}}(t_0) - \mathbf{u}_{\text{policy}}(t_0)\).
In a physical plant, applying a step change in commanded force or torque
across zero elapsed time demands an infinite time derivative of force.
On a wheeled chassis, this instantaneous torque gradient exceeds the
shear limit of the tire contact patch, breaking adhesion and causing
drive slip. In an articulated manipulator or robotic powertrain, that
same step shock excites unmodeled structural resonances in the gearboxes
and risks tooth damage on reduction drives. In a process fluid or steam
power plant, an instantaneous valve closure induces high-pressure
acoustic shock waves (hydraulic water hammer) that rupture pipeline
fittings. A discrete state transition in software that executes within a
single clock cycle cannot be mapped directly to actuators that operate
under the laws of continuum mechanics.

Eliminating this step discontinuity requires dynamic blending across a
finite time window \(\tau_{\text{blend}}\), a technique known as
\textbf{bumpless transfer}. The blended command routed to the actuator
driver is expressed as
\(\mathbf{u}(t) = (1 - \alpha(t))\mathbf{u}_{\text{policy}}(t) + \alpha(t)\mathbf{u}_{\text{human}}(t)\),
where \(\alpha(t)\) is an interpolation factor defined over the interval
\(t \in [t_0, t_0 + \tau_{\text{blend}}]\) satisfying the boundary
conditions \(\alpha(t_0) = 0\) and
\(\alpha(t_0 + \tau_{\text{blend}}) = 1\). A linear interpolation ramp
produces discontinuous derivatives at both endpoints, injecting
secondary jerk spikes at the onset and conclusion of the handover. To
preserve mechanical smoothness, the schedule must form a
\textbf{\(\mathcal{C}^2\) bumpless transfer ramp} possessing continuous
first and second derivatives with vanishing boundaries,
\(\dot{\alpha}(t_0) = 0\), \(\ddot{\alpha}(t_0) = 0\),
\(\dot{\alpha}(t_0 + \tau_{\text{blend}}) = 0\), and
\(\ddot{\alpha}(t_0 + \tau_{\text{blend}}) = 0\). A standard polynomial
that satisfies these boundary conditions is the quintic spline
\(\alpha(t) = 6\left(\frac{t - t_0}{\tau_{\text{blend}}}\right)^5 - 15\left(\frac{t - t_0}{\tau_{\text{blend}}}\right)^4 + 10\left(\frac{t - t_0}{\tau_{\text{blend}}}\right)^3\),
which yields continuous acceleration and bounded jerk profiles.

\protect\phantomsection\label{callout-definitionux2a-14.4}
\begin{fbxSimple}{callout-definition}{Definition:}{The $\mathcal{C}^2$ Bumpless Authority Transfer Ramp}
\phantomsection\label{callout-definition*-14.4}

An active control blending architecture that transitions actuator
authority between two conflicting or distinct controllers
\(\mathbf{u}_1(t)\) and \(\mathbf{u}_2(t)\) via a strictly monotonic,
thrice-differentiable weighting profile \(\alpha(t) \in [0, 1]\),
guaranteeing \(\mathcal{C}^2\) continuity across the transition window
\(T_{\mathrm{trans}}\) such that actuator acceleration and jerk remain
within physical inverter and gearhead saturation envelopes.

\end{fbxSimple}

The blending duration \(\tau_{\text{blend}}\) is bounded below by the
physical jerk limit of the plant. Consider a vehicle chassis where the
command \(u(t)\) corresponds to linear acceleration \(\ddot{x}(t)\),
such that the physical jerk along the path is
\(j(t) = \frac{d^3 x}{dt^3} = \dot{u}(t)\). Assuming for this analytical
derivation that the target manual command and policy command remain
approximately stationary over the brief transition window,
differentiating the blended command gives
\(\dot{u}(t) = \dot{\alpha}(t)(\mathbf{u}_{\text{human}} - \mathbf{u}_{\text{policy}}) = \dot{\alpha}(t)\Delta u\).
For the quintic interpolation schedule, the peak derivative occurs at
the midpoint \(t = t_0 + \frac{\tau_{\text{blend}}}{2}\), where
evaluating \(\dot{\alpha}(t)\) gives
\(\dot{\alpha}_{\text{max}} = \frac{6}{\tau_{\text{blend}}}\left(\frac{1}{2} - \frac{1}{4}\right) = \frac{1.5}{\tau_{\text{blend}}}\).
To ensure that peak actuator and chassis jerk does not exceed a maximum
allowable structural or traction threshold \(j_{\text{allowable}}\), the
condition
\(j_{\text{max}} = \frac{1.5 |\Delta u|}{\tau_{\text{blend}}} \le j_{\text{allowable}}\)
must hold. Rearranging this inequality yields the minimum allowable
blending duration
\[\tau_{\text{blend}} \ge \frac{1.5 |\mathbf{u}_{\text{human}} - \mathbf{u}_{\text{policy}}|}{j_{\text{allowable}}}\]
This first-principles bound establishes that no moving machine can
transfer control authority faster than the ratio of the command
discrepancy to its maximum allowable jerk.

The derivation exposes a latency-jerk tradeoff that governs every
physical takeover. Increasing \(\tau_{\text{blend}}\) reduces peak jerk
and keeps driveline stress within safe margins, but it delays full human
authority and lengthens the stopping trajectory. Consider an analytical
example of an autonomous vehicle moving at \(v_0 = 10.0\text{ m/s}\)
along a flat surface. The learned policy is commanding cruising
acceleration \(u_{\text{policy}} = 0.0\text{ m/s}^2\) when an operator
abruptly intervenes with full emergency braking, commanding
\(u_{\text{human}} = -6.0\text{ m/s}^2\), creating a command gap
\(|\Delta u| = 6.0\text{ m/s}^2\). If the vehicle chassis has a traction
jerk limit of \(j_{\text{allowable}} = 12.0\text{ m/s}^3\) to prevent
tire breakaway, the minimum blending window is
\(\tau_{\text{blend}} = \frac{1.5 \times 6.0\text{ m/s}^2}{12.0\text{ m/s}^3} = 0.750\text{ s}\),
or \(750\text{ ms}\). Over this transition window, integrating the
acceleration trajectory under quintic blending yields a total velocity
reduction of
\(\Delta v = \int_{0}^{\tau_{\text{blend}}} \alpha(t) u_{\text{human}} dt = \frac{1}{2} u_{\text{human}} \tau_{\text{blend}} = \frac{1}{2}(-6.0\text{ m/s}^2)(0.750\text{ s}) = -2.25\text{ m/s}\).
The distance traversed during the handover window is
\(d_{\text{blend}} = v_0 \tau_{\text{blend}} + \int_{0}^{\tau_{\text{blend}}} \left(\int_{0}^{t} \alpha(\sigma) u_{\text{human}} d\sigma\right) dt = (10.0\text{ m/s})(0.750\text{ s}) + 0.15 u_{\text{human}} \tau_{\text{blend}}^2 = 7.50\text{ m} + (0.15)(-6.0\text{ m/s}^2)(0.5625\text{ s}^2) = 7.50\text{ m} - 0.51\text{ m} = 6.99\text{ m}\).
At the end of the blend window, the vehicle reaches
\(v(\tau_{\text{blend}}) = 7.75\text{ m/s}\), after which full manual
deceleration of \(-6.0\text{ m/s}^2\) brings it to rest over an
additional distance of
\(d_{\text{post}} = \frac{(7.75\text{ m/s})^2}{2(6.0\text{ m/s}^2)} = 5.01\text{ m}\),
giving a total stopping distance of \(12.00\text{ m}\). In contrast, a
theoretical step transition that achieved \(-6.0\text{ m/s}^2\) with
zero transition delay would require
\(d_{\text{step}} = \frac{(10.0\text{ m/s})^2}{2(6.0\text{ m/s}^2)} = 8.33\text{ m}\).
The \(750\text{ ms}\) jerk-limited blend therefore adds
\(3.67\text{ m}\) to the stopping distance. However, attempting to apply
the \(-6.0\text{ m/s}^2\) step over a single \(1\text{ ms}\) control
cycle would demand a jerk of \(6000\text{ m/s}^3\), which exceeds tire
adhesion, locks the wheels, and induces sliding friction that lengthens
the stopping distance far beyond the blended path.

\protect\phantomsection\label{callout-ruleux2a-14.5}
\begin{fbxSimple}{callout-rule}{Architect’s Rule of Thumb:}{The Smoothstep Jerk-Latency Invariant}
\phantomsection\label{callout-rule*-14.5}

Actuator continuity enforces a hard lower bound on takeover rate:
\(\tau_{\text{blend}} \ge \frac{1.5 |\Delta \mathbf{u}|}{j_{\text{allowable}}}\).
Handover is an allocation of physical distance, not time; taking over
without bumpless transfer creates worse driveline instability than the
fault itself. Shortening the blending window below this physical floor
does not grant the operator faster physical authority; it excites
mechanical driveline resonances, triggers tire slip breakaway, or trips
motor current limits, lengthening the actual physical stopping envelope.

\end{fbxSimple}

Dynamic command blending prevents discontinuities at the output of the
arbitrator, but it leaves an internal discontinuity inside the human
tracking loop. If the human input channel uses a feedback controller or
an impedance filter that runs in the background while the policy holds
authority, its integral error terms accumulate against a stationary
setpoint or a disconnected plant state. When manual authority engages,
this accumulated integrator state discharges as an immediate command
spike that overrides the smooth output blend. Preventing this transient
requires \textbf{initial condition matching},\sidenote{\footnotesize In digital motor
  drives, initial condition matching pre-biases current loops
  (\(I_q, I_d\)) to suppress inductive switching spikes.} wherein the
internal states and integrators of the manual control loop are
continuously pre-biased to track the true plant state \(\mathbf{x}(t)\)
during autonomous operation. At the instant of takeover, the manual loop
initializes with zero tracking error, ensuring that the feedforward and
feedback components evaluate to the current plant equilibrium such that
\(\mathbf{u}_{\text{human}}(t_0) \approx \mathbf{u}_{\text{policy}}(t_0)\).
Pre-biasing minimizes the command delta \(|\Delta \mathbf{u}|\) before
blending begins, reducing both the required blending duration and the
resulting transition delay.

Verifying that a handover satisfies physical constraints requires
measuring the acceleration derivative \(j(t) = \frac{d^3 x}{dt^3}\) on
the physical chassis using high-bandwidth inertial measurement units,
matched against motor current telemetry and wheel encoders. In an
unblended step handover, the telemetry displays immediate high-frequency
ringing in motor torque, a sharp spike in chassis jerk that exceeds
hundreds of \(\text{m/s}^3\), and a transient divergence between wheel
rotational speed and ground speed that marks tire micro-slip. In a
\(\mathcal{C}^2\) blended transition, motor torque follows a smooth
sigmoid curve, chassis jerk stays strictly bounded below
\(j_{\text{allowable}}\), and wheel slip ratio remains within the linear
traction envelope. Handover is therefore budgeted as a physical maneuver
with measurable costs in time, jerk, and displacement, rather than
treated as a software event that completes in a single CPU cycle.

The physical waveforms comparing an unblended step handover against
\(\mathcal{C}^2\) quintic bumpless blending are illustrated in
Figure~\ref{fig-14-bumpless-transfer-dynamics}.

\begin{figure}[t]

\centering{

\includegraphics[width=0.95\linewidth,height=\textheight,keepaspectratio]{index_files/mediabag//Users/VJ/GitHub/MLSysBook-integrate-physical/physical/book/chapters/14-intervention/figures/fig14_bumpless_transfer_dynamics.pdf}

}

\caption{\label{fig-14-bumpless-transfer-dynamics}\textbf{Mechanical
Dynamics of Control Handover: Step Discontinuity versus
\(\mathcal{C}^2\) Cubic Bumpless Blending.} Synchronized oscilloscope
waveforms across an emergency takeover window (\(t_0 = 200\text{ ms}\),
\(\Delta u = 25\text{ N}\cdot\text{m}\)). Left (Panel A): Naïve
unblended step handover (\(\tau_{\text{blend}} = 0\text{ ms}\)) creates
an infinite theoretical torque gradient, exciting \(25\text{ Hz}\)
drivetrain ringing, a massive chassis jerk spike
(\(j_{\text{max}} > 6000\text{ m/s}^3 \gg j_{\text{allowable}} = 12\text{ m/s}^3\)),
and tire slip breakaway (\(s(t) > s_{\text{crit}} = 0.15\)), lengthening
the physical stopping distance by \(+3.67\text{ m}\) under sliding
friction. Right (Panel B): \(\mathcal{C}^2\) quintic bumpless blending
over \(\tau_{\text{blend}} = 750\text{ ms}\) with pre-biased tracking
integrators ensures smooth sigmoidal torque transition, strictly bounds
chassis jerk below \(j_{\text{allowable}}\), and maintains wheel slip in
the linear elastic adhesion zone (\(s(t) \le 0.04\)), eliminating
mechanical drivetrain shock and ensuring deterministic trajectory
stability. Waveforms captured on high-bandwidth dynamometer test bench.}

\end{figure}%

The complete algorithmic governor executing discrete four-phase state
handshakes, integrator pre-biasing, and continuous \(\mathcal{C}^2\)
quintic spline blending on bare-metal microcontroller silicon is
formalized below:

\subsection{Four-Phase Authority Handshake and Bumpless Transfer
Protocol}\label{four-phase-authority-handshake-and-bumpless-transfer-protocol}

\textbf{Cadence \& Invariants:}

\begin{itemize}
\tightlist
\item
  \emph{Execution Cadence:} Hard real-time periodic timer interrupt
  (\(f = 1000\text{ Hz}\), period \(T_{\text{loop}} = 1.0\text{ ms}\),
  hardware timing jitter \(< 1.0\,\mu\text{s}\)).
\item
  \emph{Memory Invariant:} Zero dynamic heap allocation
  (\texttt{malloc}/\texttt{free} strictly forbidden; all state
  variables, mailboxes, and circular log buffers statically
  pre-allocated in DTCM SRAM).
\item
  \emph{Hardware Authority:} Exclusive direct register write lock to
  motor inverter PWM compare registers (\texttt{TIMx\_CCR} /
  \texttt{ePWM\ CMPA}), hydraulic proportional valve DACs, and Safe
  Torque Off (STO) hardware cutoffs.
\end{itemize}

\textbf{Execution Protocol:}

\begin{enumerate}
\def\labelenumi{\arabic{enumi}.}
\item
  \textbf{Phase 1: Emergency \& Preemption Intercept (Stage 1):}

  \begin{itemize}
  \tightlist
  \item
    If hardwired E-stop line opens or critical barrier certificate is
    breached (\(h(\mathbf{x}) < 0\)): set
    \(S_{\text{auth}} \leftarrow \text{ESTOP}\), de-energize gate
    drivers via Safe Torque Off (STO), and freeze the evidence ring
    buffer.
  \item
    Else if \(S_{\text{auth}} == \text{AUTONOMOUS}\):

    \begin{itemize}
    \tightlist
    \item
      If manual column torque \(|\tau_{\text{meas}}| > 4.0\text{ Nm}\)
      or authenticated override button assert is received: pre-bias
      manual tracking loop integrators
      (\(\hat{\mathbf{u}}_{\text{manual\_bias}} \leftarrow \mathbf{u}_{\text{policy}}(t)\)),
      latch transition onset (\(t_0 \leftarrow t_{\text{now}}\)),
      compute discrepancy
      \(\Delta \mathbf{u} = \|\mathbf{u}_{\text{human}} - \mathbf{u}_{\text{policy}}\|\),
      size the physical blending window
      \(\tau_{\text{blend}} \leftarrow \max(50\text{ ms}, \frac{1.5 \|\Delta \mathbf{u}\|}{j_{\text{allowable}}})\),
      and transition to
      \(S_{\text{auth}} \leftarrow \text{BUMPLESS\_BLEND}\).
    \item
      Else if policy confidence \(\gamma < \gamma_{\text{min}}\) or ODD
      geofence departure occurs: set
      \(t_{\text{req}} \leftarrow t_{\text{now}}\), emit alert to
      operator, and transition to
      \(S_{\text{auth}} \leftarrow \text{HANDOVER\_PENDING}\).
    \end{itemize}
  \end{itemize}
\item
  \textbf{Phase 2: Handover Acknowledgment \& Timeout Governor (Stage
  2):} If \(S_{\text{auth}} == \text{HANDOVER\_PENDING}\):

  \begin{itemize}
  \tightlist
  \item
    If capacitive grip confirmed and valid AES-128-GMAC operator packet
    received: pre-bias manual integrators, initialize
    \(t_0 \leftarrow t_{\text{now}}\), and transition to
    \(S_{\text{auth}} \leftarrow \text{BUMPLESS\_BLEND}\).
  \item
    Else if \((t_{\text{now}} - t_{\text{req}}) > T_{\text{timeout}}\)
    (unacknowledged takeover): lock out upward transitions and
    transition to
    \(S_{\text{auth}} \leftarrow \text{DEGRADED\_FALLBACK}\).
  \end{itemize}
\item
  \textbf{Phase 3: Synchronous Commit \& \(\mathcal{C}^2\) Smoothstep
  Blending (Stage 3):} If \(S_{\text{auth}} == \text{BUMPLESS\_BLEND}\):

  \begin{itemize}
  \tightlist
  \item
    Normalized blend progress:
    \(\hat{t} \leftarrow \min(1.0, \frac{t_{\text{now}} - t_0}{\tau_{\text{blend}}})\).
  \item
    Cubic smoothstep weight:
    \(\alpha(\hat{t}) \leftarrow 3\hat{t}^2 - 2\hat{t}^3\).
  \item
    Blended actuator command:
    \(\mathbf{u}_{\text{act}}(t) \leftarrow (1 - \alpha(\hat{t}))\mathbf{u}_{\text{policy}}(t) + \alpha(\hat{t})\mathbf{u}_{\text{human}}(t)\).
  \item
    If \(\hat{t} \ge 1.0\): revoke policy execution token atomically and
    transition to \(S_{\text{auth}} \leftarrow \text{HUMAN\_MANUAL}\).
  \end{itemize}
\item
  \textbf{Phase 4: Confirm \& Active Lease Watchdog (Stage 4):} If
  \(S_{\text{auth}} == \text{HUMAN\_MANUAL}\):

  \begin{itemize}
  \tightlist
  \item
    Evaluate keep-alive lease age:
    \(\Delta t_{\text{lease}} \leftarrow t_{\text{now}} - t_{\text{lease\_human}}\).
  \item
    If \(\Delta t_{\text{lease}} > \tau_{\text{lease}}\) or enabling
    device (deadman switch) released: transition to
    \(S_{\text{auth}} \leftarrow \text{DEGRADED\_FALLBACK}\).
  \item
    Else: direct manual actuation
    \(\mathbf{u}_{\text{act}}(t) \leftarrow \mathbf{u}_{\text{human}}(t)\).
  \end{itemize}
\item
  \textbf{Stage 5: Fallback Execution \& Immutable Telemetry Logging
  (Stage 5):}

  \begin{itemize}
  \tightlist
  \item
    If \(S_{\text{auth}} == \text{DEGRADED\_FALLBACK}\): dispatch
    controlled deceleration profile
    (\(a_{\text{decel}} = -2.5\text{ m/s}^2\)) and maintain dynamic
    stability assist.
  \item
    Enforce hardware current clamps and barrier certificate projection:
    \(\mathbf{u}_{\text{act}} \leftarrow \text{CBF\_Project}(\mathbf{u}_{\text{act}}, \mathbf{x}_{\text{plant}})\).
  \item
    Write \(\mathbf{u}_{\text{act}}(t)\) directly to timer compare
    registers (\texttt{TIMx\_CCR}), serialize \(\mathcal{E}_t\) into
    pre-allocated DTCM SRAM ring buffers, and update chained SHA-256
    HMAC digests.
  \end{itemize}
\end{enumerate}

Calculating blending filters and pre-biasing controller states assumes
that the incoming signal originates from a verified, legitimate operator
entitled to assume physical control. If the nervous system accepts a
transfer request without validating identity, an unauthenticated command
injection or a malfunctioning subsystem can seize authority and override
the machine while it is in motion. The mechanisms that make an
intervention smooth must therefore be paired with mechanisms that
establish who is allowed to intervene.

\section{Authenticating an Intervention}\label{sec-intervention-14-5}

Authenticating an intervention requires preventing an unauthorized
entity from seizing physical authority while ensuring that an authorized
operator can halt the machine before collision. These two requirements
pull in opposite directions, and resolving the conflict requires
classifying intervention operations by the mechanical work they can
induce. Removing power or asserting a mechanical brake is physically
distinct from injecting active motion commands, and the two actions do
not warrant identical verification mechanisms. Removing kinetic energy
forces the machine toward a static equilibrium state governed by passive
mechanics, where energy dissipates through friction and contactors open
by spring return. In contrast, injecting active steering angles, wheel
torques, or joint velocity setpoints introduces fresh mechanical power
into the system. An unauthorized or corrupted motion command can
accelerate a moving mass into an obstacle, whereas an unauthorized
emergency stop halts the machine where it stands.

Because motion commands introduce active kinetic hazards, the
verification mechanism must satisfy both a latency bound and a freshness
bound derived directly from the kinematics of the physical maneuver.
Consider an analytical mobile platform of mass \(m = 250\text{ kg}\)
traveling at a forward velocity of \(v = 2.0\text{ m/s}\) toward an
obstruction detected at an initial clearance distance of
\(d = 1.20\text{ m}\). If the mechanical brakes apply a constant
deceleration of \(a = 4.0\text{ m/s}^2\), the physical stopping distance
during active braking is
\(d_{\text{brake}} = v^2 / (2a) = (2.0\text{ m/s})^2 / (2 \times 4.0\text{ m/s}^2) = 0.50\text{ m}\).
The displacement available for detection, message transit, signature
verification, and brake activation is
\(d_{\text{margin}} = 1.20\text{ m} - 0.50\text{ m} = 0.70\text{ m}\).
At \(v = 2.0\text{ m/s}\), the total elapsed time before braking must
commence is
\(T_{\text{max}} = 0.70\text{ m} / 2.0\text{ m/s} = 350\text{ ms}\). If
mechanical actuator lag and hydraulic pressure buildup consume
\(160\text{ ms}\), and internal bus serialization consumes
\(40\text{ ms}\), the remaining deadline for cryptographic signature
verification and replay checking is bounded by
\(T_{\text{auth}} \le 150\text{ ms}\) at \(P_{99}\).

\marginnote{\begin{footnotesize}

\textbf{CAN Bus Priority:} In CAN controllers lacking transmit-abort
registers, queued low-priority telemetry blocks emergency override
frames. Modern IP (Bosch M\_CAN) resolves this with transmit-buffer
cancellation interrupts (\texttt{TXBCF}) and hardware priority queuing
(\texttt{TXBAR}), ensuring high-priority safety frames preempt bulk
telemetry within one bit-time.

\end{footnotesize}}

Freshness requirements arise from the same kinematics. If an
authentication window permits a timestamp drift or message age threshold
of \(\tau_{\text{fresh}} = 100\text{ ms}\), an authenticated packet
replayed from the boundary of that window commands an action based on a
chassis position displaced by
\(\Delta x = v \cdot \tau_{\text{fresh}} = (2.0\text{ m/s})(0.10\text{ s}) = 0.20\text{ m}\).
In a precision navigation envelope, that spatial displacement turns a
safe clearance margin into a physical collision.

Standard information security protocols fail closed when authentication
fails, dropping invalid packets and maintaining the existing operating
mode. In a physical machine, failing closed during an emergency
intervention leaves an uncontrolled mass moving under the command of a
failing or compromised policy. The emergency path must therefore
\textbf{fail safe} rather than fail closed. If an emergency stop signal
arrives with a corrupted checksum, an expired sequence nonce, or a
signature that fails cryptographic verification, the nervous system
cannot discard the message and await a corrected retransmission. It must
immediately de-energize the motor drives and assert mechanical holding
brakes. The cost of this design is complete vulnerability to
denial-of-service disruptions. An adversary broadcasting malformed
packets across a wireless link, or electrical noise induced on an
unshielded CAN bus wire, halts the machine and terminates operation.
This trade is made explicitly, accepting a loss of operational
availability to ensure that no legitimate stop signal is delayed by
cryptographic negotiation while a machine is in motion.

Designing intervention channels around unverified, zero-latency
assumptions leads to systematic failures during policy training. In
interactive learning regimes such as dataset aggregation, human
supervisors monitor autonomous rollouts and seize control whenever the
learned model approaches an unsafe state, appending the corrective
trajectory to the training set. Across ten to fifteen analytical
training rounds, this protocol teaches the learned policy that
aggressive drift toward physical boundaries carries no operational
penalty because an external rescue occurs instantaneously before
contact. Because the training environment executes supervisory handovers
without cryptographic signing or freshness validation, the policy learns
to operate right at the margin of the physical safety envelope. When
that policy is deployed on physical hardware where intervention packets
must pass through cryptographic verification, timestamp checks, and
replay filters, the resulting latency adds \(100\text{ ms}\) to
\(200\text{ ms}\) of unmodeled delay. During an actual field rescue, the
supervisor initiates a takeover, but the authenticated command arrives
after the machine has already traversed the clearance margin. This
failure mode cannot be detected by monitoring offline training loss,
which continues to decrease as the policy memorizes rescue trajectories.
It is detected only by injecting realistic cryptographic latency into
hardware-in-the-loop simulation, or it goes entirely undetected until
the machine collides with an obstacle during a real-world intervention.

Enforcing cryptographic authentication at the primary telemetry gateway
is ineffective if unauthenticated communication channels bypass the
nervous system. Every physical machine carries \textbf{shadow paths}
that can alter actuator states without traversing the safety monitor.
Diagnostic interfaces, such as unauthenticated CAN or Ethernet ports
intended for maintenance calibration, frequently expose memory-mapped
write commands directly to motor controllers. Microcontroller debug
ports, including JTAG and serial-wire debug interfaces left active on
production circuit boards, allow an external probe to write directly to
pulse-width modulation registers, asserting full motor torque within
microseconds. Bench testing fixtures spliced into actuator wiring looms
can pull analog enable lines to high logic levels, overriding software
interlocks. In addition, microcontroller bootloaders that accept
unauthenticated firmware updates over internal buses during power-on
initialization can replace motor control logic entirely. Securing the
machine requires identifying every physical path to actuation and either
burning one-time programmable fuses to permanently disable hardware
debug ports on production boards, or routing all internal bus traffic
through a dedicated hardware root of trust.

\textbf{Cryptographic Primitives:} The mathematical formulation of
cryptographic primitives, including public key signatures such as
Ed25519, symmetric message authentication codes such as AES-GMAC, and
key distribution hierarchies, is treated comprehensively in foundational
security literature. Systems engineering for physical AI is responsible
for mapping those primitives onto physical deadlines, managing the trade
between latency and authority, and identifying the residual risks that
cryptographic verification cannot address. The architecture presented
here defends against remote packet injection over wireless links, replay
attacks on internal buses, and unauthorized commands originating from
compromised high-level application processors. Defending against
physical hardware tampering within the chassis requires potting
electronics in epoxy, integrating cryptographic co-processors into every
individual motor stator, and encasing wiring looms in armored conduit.

When an intervention occurs, whether triggered by a legitimate operator,
an unauthenticated command injection, or electrical noise on a safety
line, the nervous system must record the transition alongside the
physical state of the plant. Diagnosing why an intervention occurred,
establishing whether an emergency stop was caused by a genuine boundary
violation or an invalid cryptographic token, and proving that the
machine behaved correctly under handover deadlines requires an immutable
record preserved across physical power cycles.

\section{Trustworthy Operational Logs}\label{sec-intervention-14-6}

A record that has to support an argument months later has requirements a
debug log does not. In standard software development, logging is
best-effort, buffered in volatile user space, dropped under
backpressure, and formatted for human inspection during live debugging.
If an asynchronous disk write drops three lines of trace data during an
operating system stall, an engineer loses minor diagnostic context while
the application continues running. In a physical machine, a lossy log
destroys the causal chain between command and physical consequence. When
a moving manipulator strikes an assembly jig, establishing whether the
learned policy commanded an illegal torque, whether the safety enforcer
failed to clamp it, or whether an operator intervened too late requires
a deterministic, legally defensible, and tamper-evident
record.\sidenote{\footnotesize Statutory standards (49 CFR Part 563, IEEE 1616)
  mandate pre-crash logs (\(5\text{ s}\) at \(\ge 10\text{ Hz}\)), but
  physical AI requires sub-microsecond state tuples. Ethernet MAC
  registers (\texttt{PTP\_TSL}/\texttt{TSH}) latch physical bit-arrival
  timestamps at the MII interface to satisfy ISO 26262 ASIL D audit
  trails.} A \textbf{forensic operational log} therefore does not exist
to assist a developer with line-by-line code tracing. Its purpose is to
preserve mathematical proof of what the machine perceived, what its
subsystems decided, and what physical forces reached the actuators
across every control cycle.

Reconstructing that causal chain requires every discrete control tick to
produce a structured, immutable record of authority and action. The
minimal \textbf{evidence tuple} takes the form
\(\langle t_{\text{PTP}}, S_{\text{auth}}, \mathbf{x}_{\text{state}}, \mathbf{u}_{\text{prop}}, \mathbf{u}_{\text{enf}}, \mathbf{u}_{\text{act}}, \text{ID}_{\text{operator}}\rangle\).
In this structure, \(t_{\text{PTP}}\) provides a hardware-synchronized
timestamp, \(S_{\text{auth}}\) records the active authority state
indicating whether autonomous policy, manual override, or fault
mitigation holds control, \(\mathbf{x}_{\text{state}}\) captures the
measured kinematic and dynamic state of the plant from physical sensors,
\(\mathbf{u}_{\text{prop}}\) contains the unconstrained action vector
proposed by the learned policy, \(\mathbf{u}_{\text{enf}}\) records the
action vector permitted or modified by the safety enforcer,
\(\mathbf{u}_{\text{act}}\) represents the actual command delivered over
the physical bus to motor drives, and \(\text{ID}_{\text{operator}}\)
identifies the authenticated source of any active intervention. Omitting
any element of this tuple breaks post-incident analysis. If the log
records \(\mathbf{u}_{\text{act}}\) without
\(\mathbf{u}_{\text{prop}}\), investigators cannot determine whether an
improper motion originated in a policy inference failure or in
downstream actuator driver corruption. If \(\mathbf{u}_{\text{enf}}\) is
omitted, the log cannot prove whether the safety enforcer actively
clamped a hazardous command or failed to evaluate its safety invariant
within the control deadline.

A \textbf{forensic authority record} is a statically verified,
cryptographically signed, and tamper-evident ledger entry---synchronized
via hardware timestamping (IEEE 1588 PTP)---that logs the discrete
authority state \(S_{\text{auth}}\), raw proposed action
\(\mathbf{u}_{\text{prop}}\), enforcer-modified action
\(\mathbf{u}_{\text{enf}}\), delivered actuator command
\(\mathbf{u}_{\text{act}}\), and operator identity across every control
cycle to establish deterministic post-incident liability and isolate
uncurated rescue trajectories from offline training sets.

Ordering these tuples requires physical timestamp integrity rather than
software clock queries. Software timestamps derived from operating
system calls inherit scheduling jitter, context-switch delays, and
variable interrupt latency that can exceed \(5\text{ ms}\) on loaded
application processors. Over a \(5\text{ ms}\) interval, a mobile base
traveling at \(2.0\text{ m/s}\) covers \(10\text{ mm}\), and an internal
control loop running at \(1000\text{ Hz}\) executes five complete
control cycles. If an operating system logs the arrival of a manual
intervention command with a \(5\text{ ms}\) delay, the record places the
operator's command after the actuator response, inverting physical cause
and effect. To guarantee causal ordering between operator input, bus
transit, and motor current, the nervous system relies on
hardware-timestamped Precision Time Protocol (IEEE 1588 PTP). By
stamping Ethernet frames and CAN FD packets at the physical layer
transceiver as bits cross the physical interface, the system bounds
timestamp uncertainty to under \(100\text{ ns}\). This sub-microsecond
resolution ensures that if an operator asserts an emergency stop line
while an actuator is in motion, the log proves definitively whether the
bus delivered the stop command before or after the mechanism made
contact with an obstacle.

Persisting full sensor streams and high-rate telemetry continuously to
non-volatile storage exceeds both the bandwidth and the endurance limits
of embedded hardware. Consider an analytical mobile platform equipped
with four depth sensors generating \(1280 \times 720\) 16-bit depth
frames at \(30\text{ Hz}\). Each camera produces approximately
\(55.3\text{ MB/s}\), totaling \(221\text{ MB/s}\) across four sensors.
Adding the control telemetry of the 128-byte evidence tuple recorded at
\(1000\text{ Hz}\) contributes another \(128\text{ kB/s}\). In
uncompressed form, writing this continuous stream directly to
non-volatile flash storage would saturate storage controllers and
exhaust solid-state write endurance within days. To retain complete
forensic fidelity without storage exhaustion, the nervous system employs
a dual-rate architecture centered on a high-frequency circular ring
buffer allocated in fast volatile memory. The machine writes full plant
state, policy proposals, enforcer evaluations, and raw sensor frames
into this ring buffer continuously at \(100\text{ Hz}\) to
\(1000\text{ Hz}\), overwriting the oldest entries. When a triggering
condition occurs (such as an operator takeover, an enforcer
intervention, a cryptographic authentication failure, or a physical
limit trip), the nervous system freezes the buffer. It commits a bounded
window, such as \(10\text{ s}\) of pre-trigger history and
\(10\text{ s}\) of post-trigger trajectory, to non-volatile flash while
background telemetry continues logging summary metrics at
\(10\text{ Hz}\).

The pre-trigger window must capture the moment of policy divergence
rather than merely the moment of physical takeover. When a learned
policy commands an erroneous trajectory, human operators and external
monitoring processes exhibit reaction latencies ranging from hundreds of
milliseconds to multiple seconds before recognizing the deviation and
asserting authority. If the ring buffer only recorded data from the
instant an operator pulled an override switch, the telemetry would
capture the rescue while omitting the anomalous commands that made the
rescue necessary. Capturing the preceding \(10\text{ s}\) of execution
preserves the precise instant where \(\mathbf{u}_{\text{prop}}\)
diverged from the safe manifold or nominal distribution. As established
in Section~\ref{sec-data-5-5}, without the complete pre-divergence state
history, evaluation pipelines cannot separate a necessary rescue from an
unnecessary or premature human interruption. Losing that divergence
window corrupts the statistical metrics used to validate policy safety
cases and distorts the imitation data collected for offline model
retraining.

Preserving this data against post-incident modification, power loss, or
operating system compromise requires cryptographic append-only
structures. Within the nervous system, logging routines batch evidence
tuples into fixed blocks corresponding to \(100\text{ ms}\) of
execution. Each block is hashed using SHA-256, and its hash is
concatenated with the hash of the preceding block to form an append-only
cryptographic chain. Modifying or deleting any single historical tuple
invalidates the hash of its containing block and breaks the chain for
all subsequent entries. To prevent an attacker or a compromised
application processor from rewriting the log after an unauthorized
command or collision, the nervous system periodically writes root hashes
to write-once media or anchors them inside a secure enclave equipped
with hardware monotonic counters. When a physical collision causes
sudden power termination, the append-only structure ensures that all
blocks committed prior to the power failure remain cryptographically
verifiable, while partial writes occurring during voltage collapse are
unambiguously detected and isolated.

A verifiable operational log guarantees that months after an incident,
an engineering team can prove what authority state governed the machine,
what values the policy proposed, and what forces the actuators applied
at every millisecond of an operation. Yet possessing a mathematically
sound record of authority does not prevent physical instability during
the transition itself. When authority transfers while a mechanical body
is carrying momentum and actuators are delivering active torque, the
machine faces the immediate physical problem of coordinating two
independent control sources. The central danger of an intervention is
not that the log will fail to record the event, but that the handover
itself will leave the machine in a state where neither the policy nor
the operator has coherent control.

\section{How Handover Goes Wrong}\label{sec-intervention-14-7}

The most common handover failure is that nobody was sure who had
control. When an autonomous policy and a human operator share authority
over an embodied platform, system failures rarely stem from broken
mechanical linkages or unhandled software exceptions. They stem from
mode confusion and authority ambiguity, where the machine assumes the
human is managing an operational axis while the human assumes the
machine retains jurisdiction. In multi-axis physical systems, authority
is frequently decomposed across degrees of freedom, separating lateral
steering from longitudinal throttle and braking. If a learned policy
encounters an unmodeled surface reflection and quietly yields lateral
tracking without disengaging the longitudinal speed regulator, an
operator who notices the lateral drift may grasp the wheel to steer
while assuming the automated system will continue to brake for
downstream traffic. Conversely, if the operator taps the brake pedal to
initiate longitudinal deceleration, firmware that disconnects only the
longitudinal controller leaves the lateral steering policy actively
hunting for lane boundaries against the driver's intended exit
trajectory. In telemetry logs, this failure presents as disjointed
control flags where authority bits toggle asynchronously across
execution threads. At the human-machine interface, however, this state
is undetected until kinematic error exceeds the physical clearance
margin, because neither claimant receives an immediate, unambiguous
signal of the other's partial withdrawal.

Even when authority transfer protocols are defined, the human nervous
system cannot match the response deadlines of moving mass. An active
human operator in direct physical control reacts to sudden tactile or
visual hazards with a motor reflex latency of roughly \(0.5\text{ s}\).
When the same operator is placed in a passive supervisory role while a
neural policy guides the platform, situational awareness decays. In an
analytical scenario of an autonomous vehicle traveling at
\(30\text{ m/s}\) (\(108\text{ km/h}\)), nominal cruising for several
minutes causes operator visual gaze to fixate away from the forward
collision envelope, while cognitive readiness drops. When the onboard
system encounters an unresolvable perceptual distribution shift and
issues an emergency handover request, human reaction time inflates from
\(0.5\text{ s}\) to between \(3.0\text{ s}\) and \(5.0\text{ s}\).
During this transition gap, a platform moving at \(30\text{ m/s}\)
covers \(90\text{ m}\) to \(150\text{ m}\) without coherent guidance. By
the time the operator perceives the obstacle, re-establishes spatial
orientation, and applies physical input, the available stopping distance
has been entirely consumed. Cabin-facing infrared cameras can track
pupil dilation and gaze direction at \(60\text{ Hz}\), detecting eye
closure or head rotation. They cannot measure cognitive comprehension of
dynamic scene geometry, meaning this latency expansion is detected only
after the response deadline has elapsed and impact occurs.

When an intervention occurs while mechanical actuators are delivering
active force, software filters and human limbs can enter violent
opposition. If an operator applies an abrupt
\(15\text{ N}\cdot\text{m}\) manual steering torque to swerve around a
road hazard, a downstream safety enforcer executing a Control Barrier
Function or acceleration rate limiter may classify the resulting lateral
jerk as a dynamic rollover hazard. The enforcer calculates a
minimal-intervention correction and commands the electric power steering
motor to apply counter-torque. Because the human neuromuscular system
acts as a mechanical spring-damper while the motor drive closes its
current loop at \(20\text{ kHz}\), the two opposing controllers excite
the compliance of the steering column. The mechanical assembly enters a
high-frequency limit-cycle oscillation between \(4\text{ Hz}\) and
\(8\text{ Hz}\). Rather than executing the evasive maneuver, the motor
inverter rapidly hits its thermal and current saturation limits, the
effective net steering angle drops to near zero, and the machine
continues along its original collision trajectory. In the nervous
system, this condition is detected immediately by comparing the sign of
the measured column torque against motor current feedback
(\(\tau_{\text{human}} \cdot \tau_{\text{motor}} < 0\)). Yet without an
explicit preemption contract in firmware that instantly zeroes motor
torque or yields compliance upon manual threshold crossing, the control
loop fights the human operator until electrical overload trips the drive
or mechanical stops shear.

The damage from an uncoordinated handover extends beyond the physical
interaction into the learned models trained on fleet logs. When an
intervention occurs, the telemetry stream contains several seconds of
anomalous states where the automated policy drifted, followed by
high-amplitude, high-jerk recovery actions executed by the human
operator. If an automated data ingestion pipeline incorporates these raw
telemetry logs into offline behavioral cloning or diffusion training
datasets without explicit intervention masking, the training
distribution is corrupted. The learned policy optimizes its weights
against a dataset where nominal lane features are paired with
aggressive, non-equilibrium steering inputs. In subsequent training
epochs, the model associates ordinary highway geometry with emergency
recovery dynamics, causing the next-generation policy to hallucinate
abrupt steering steps in nominal operating conditions. This failure is
undetected during training, as standard validation losses (such as mean
squared error on trajectory chunks) show nominal numerical convergence.
The dataset corruption is discovered only during deployment, when
production machines execute sudden, uncommanded swerves on open paths.

\begin{figure}[t]

\centering{

\includegraphics[width=0.9\linewidth,height=\textheight,keepaspectratio]{/Users/VJ/GitHub/MLSysBook-integrate-physical/physical/book/chapters/14-intervention/figures/fig14_real_ntsb_telemetry.png}

}

\caption{\label{fig-14-ntsb-telemetry}\textbf{NTSB Telemetry Playback of
Autonomous Disengagement and Delayed Handover.} Official National
Transportation Safety Board (NTSB) data playback from the fatal March
2018 Uber automated test vehicle collision in Tempe, Arizona
(\citeproc{ref-ntsb2019uber}{National Transportation Safety Board
2019}). The telemetry trace illustrates the failure cascade of informal
supervisory handover: the onboard perception system detected the
obstacle 6.0 seconds prior to impact but repeatedly toggled its
classification between vehicle, bicycle, and unknown, resetting tracking
filters; autonomous emergency braking had been disabled in software to
prevent false-positive braking jitter; and decayed operator situational
awareness delayed physical intervention until 1.3 seconds before
collision (when steering wheel torque was applied) and 0.7 seconds (when
the brake pedal was depressed), far exceeding the physical stopping
distance budget. Telemetry trace archived in statutory NTSB docket.}

\end{figure}%

An autopsy of physical handover collisions reveals the compound failure
when these mechanisms align in a single event
(Figure~\ref{fig-14-ntsb-telemetry}). In an analytical reconstruction of
a high-speed mobile platform approaching an unmapped stationary barrier,
the timeline decomposes into four distinct phases. In Phase 1, the
learned vision policy fails to segment the barrier under low-contrast
illumination, maintaining nominal speed while the safety margin shrinks.
In Phase 2, the onboard monitoring loop detects kinematic divergence
\(1.8\text{ s}\) before projected impact and issues an acoustic takeover
alert, but passive supervisory gaze decay delays operator perception by
\(1.2\text{ s}\). In Phase 3, with only \(0.6\text{ s}\) remaining, the
operator applies emergency manual steering torque, but the active
low-level stability enforcer clamps the rate of change and applies
counter-torque to prevent chassis slip. The resulting \(6\text{ Hz}\)
torque fighting locks the steering assembly, preventing lateral evasion.
In Phase 4, the vehicle strikes the obstacle at \(22\text{ m/s}\),
shearing the front suspension and fracturing the steering gearbox drive
pins. The failure occurred not because the operator was absent or
because the safety enforcer was inactive, but because the architecture
possessed no unified contract governing how authority transfers between
them under deadlines.

\protect\phantomsection\label{callout-incidentux2a-14.6}
\begin{fbxSimple}{callout-incident}{Landmark Engineering Incident:}{Boeing 737 MAX MCAS and Airbus A320 Warsaw Mode Confusion \& Authority Contention}
\phantomsection\label{callout-incident*-14.6}

\emph{Systems.}

\begin{enumerate}
\def\labelenumi{\arabic{enumi}.}
\tightlist
\item
  Boeing 737 MAX 8 Flight Control Computer with Maneuvering
  Characteristics Augmentation System (MCAS)
  (\citeproc{ref-jatr2019boeing}{Joint Authorities Technical Review
  2019}).
\item
  Airbus A320-211 Flight Augmentation and Braking Computer on Lufthansa
  Flight 2904 (\citeproc{ref-bfu1994warsaw}{Main Commission Aircraft
  Accident Investigation Warsaw and German Federal Bureau of Aircraft
  Accident Investigation 1994}).
\end{enumerate}

\emph{Incidents.}

\begin{itemize}
\tightlist
\item
  \textbf{Boeing 737 MAX (Lion Air JT610 \& Ethiopian ET302,
  2018--2019):} Repetitive uncommanded automated nose-down stabilizer
  trim inputs resulting in loss of aircraft pitch control, structural
  overspeed, and fatal high-velocity dive impacts.
\item
  \textbf{Airbus A320 (Warsaw Okęcie Runway Overrun, 1993):} Complete
  lockout of manual wheel braking and engine thrust reversers for 9
  seconds after runway touchdown in heavy rain, resulting in runway
  overrun and embankment collision at \(133\text{ km/h}\).
\end{itemize}

\emph{Telemetry \& systems dynamics analysis.}

\begin{itemize}
\tightlist
\item
  \textbf{MCAS Single-Sensor Authority Hijack:} MCAS was designed to
  automatically deflect the horizontal stabilizer at
  \(0.27^\circ/\text{s}\) (up to \(2.5^\circ\) maximum per cycle) to
  linearize stick-force gradients at elevated angles of attack (AoA).
  Relying on a single unvalidated AoA sensor vane that failed on
  takeoff, MCAS repeatedly engaged. When pilots pulled back on the
  control column with \(>50\text{ lb}\) manual force, electric trim
  cutout was not asserted; MCAS re-triggered 5 seconds after each manual
  thumb-switch trim input, accumulating full nose-down stabilizer trim.
  Aerodynamic loading on the horizontal stabilizer jackscrew created
  mechanical column lockup, making manual manual pitch recovery
  physically impossible.\sidenote{\footnotesize The JATR investigation
    (\citeproc{ref-jatr2019boeing}{Joint Authorities Technical Review
    2019}) found MCAS lacked cross-channel AoA redundancy, omitted
    cockpit annunciation, and repeatedly reset nose-down trim authority
    upon thumb-switch release. This allowed automated flight software to
    overpower pilot control forces without explicit authorization.}
\item
  \textbf{Warsaw Dual-Condition Inhibit Lockout:} The A320 braking
  computer prohibited ground spoiler deployment and wheel braking unless
  two hardware conditions were met: main gear shock absorbers compressed
  with \(\ge 6.3\text{ tons}\) force per strut (squat switches) OR main
  wheel tachometers measured angular velocity \(\ge 72\text{ knots}\).
  Due to a \(20\text{ knot}\) tailwind and hydroplaning on waterlogged
  asphalt, the left gear touched down lightly while the right gear
  floated for 9 seconds. The flight computers concluded the aircraft was
  airborne, rejecting pilot pedal braking and locking the mechanical
  thrust reverser levers. By the time the right gear compressed and
  brakes unlocked, \(1525\text{ m}\) of the \(2800\text{ m}\) runway had
  been consumed without deceleration.\sidenote{\footnotesize Accident investigation
    into Lufthansa Flight 2904 (\citeproc{ref-bfu1994warsaw}{Main
    Commission Aircraft Accident Investigation Warsaw and German Federal
    Bureau of Aircraft Accident Investigation 1994}) revealed that
    flight software required simultaneous wheel load and spin-up to
    confirm ground contact. Aquaplaning prevented both conditions,
    locking out pilot manual braking and thrust reversers for 9 seconds
    during runway overrun.}
\end{itemize}

\emph{Physical AI architectural lesson.} Mode confusion and automated
lockout of human authority transform transient sensor errors into fatal
kinetic catastrophes. If an automated policy can repeatedly override
human manual inputs without clear state annunciation, or if hardware
interlocks refuse human emergency stopping commands based on faulty
environmental estimation, the machine enters a fatal authority deadlock.
Actuator arbitration layers must enforce deterministic, mutually
exclusive state transitions, provide active tactile feedback on control
surfaces, and guarantee that physical emergency stopping paths remain
fail-safe and human-accessible regardless of software sensor
classification.

\end{fbxSimple}

\protect\phantomsection\label{callout-ruleux2a-14.7}
\begin{fbxSimple}{callout-rule}{Architect’s Rule of Thumb:}{The Invariant Priority Invariant}
\phantomsection\label{callout-rule*-14.7}

Manual takeover grants authority to select among safe physical
trajectories, never authority to command the plant outside its physical
survival envelope. Human inputs remain unprivileged proposals subject to
deterministic barrier gating and kinematic floor limits.

\end{fbxSimple}

Preventing these concurrent failures requires an explicit, unambiguous
mechanism that adjudicates command authority before physical signals
reach the actuator stage. The nervous system cannot rely on implicit
handoffs, human attentiveness, or ad hoc torque overrides. It requires
an immutable authority record that specifies exactly who holds
permission over each physical axis, under what operational state
conditions that permission transfers, and what strict temporal deadlines
govern the release of mechanical control.

\section{The Authority Log}\label{sec-intervention-14-8}

The machine notebook established in Section~\ref{sec-nervous-4-9}
defines the base schema, hardware provenance records, cryptographic
hashing, and monotonic versioning for all machine configurations.
Chapter~\ref{sec-intervention} adds a single table to this manifest, the
\textbf{forensic authority record}. This entry is not a high-level
policy document or a runtime advisory flag. It is the deterministic,
statically verified configuration block loaded into the nervous system
before any actuator driver energizes. It binds every physical actuator
channel to an explicit matrix of authorized actors, operational
preconditions, preemption priorities, and transition deadlines.

The structure of the authority record specifies every permitted
interaction as a hard contract between the physical platform and its
potential operators. Each entry identifies a distinct physical authority
domain through a hardware actuator mask, an authorized entity class
(such as a local physical operator, a remote teleoperation console, an
autonomous trajectory policy, or an emergency stop loop), and an integer
preemption tier. An entity presenting a command vector must prove its
authority through cryptographic role credentials verified by the nervous
system before its input can enter the actuation pipeline. Alongside the
credential, the entry specifies the mutual exclusion policy for the
channel, the maximum allowed command gradient, the active blending
duration \(\tau_{\text{blend}}\), and the terminal fallback state that
the hardware enters if an active session is severed.

To formalize this verification and forensic contract across the hardware
boundary, Table~\ref{tbl-14-authority-log-schema} defines the abstract
data schema for the forensic authority and takeover record. Every field
is bound to an explicit physical SI unit, a real-time validation
invariant, an independent hardware verification substrate, and a
forensic verification purpose.

\begin{longtable}[]{@{}
  >{\raggedright\arraybackslash}p{(\linewidth - 10\tabcolsep) * \real{0.1667}}
  >{\raggedright\arraybackslash}p{(\linewidth - 10\tabcolsep) * \real{0.1667}}
  >{\raggedright\arraybackslash}p{(\linewidth - 10\tabcolsep) * \real{0.1667}}
  >{\raggedright\arraybackslash}p{(\linewidth - 10\tabcolsep) * \real{0.1667}}
  >{\raggedright\arraybackslash}p{(\linewidth - 10\tabcolsep) * \real{0.1667}}
  >{\raggedright\arraybackslash}p{(\linewidth - 10\tabcolsep) * \real{0.1667}}@{}}
\caption{\textbf{Abstract Data Schema for the Forensic Authority and
Takeover Record (\protect\texttt{ForensicAuthorityRecord\_t}).} Memory
layout (\(118\text{ bytes}\) packed, \(128\text{ bytes}\) aligned),
binary data types, physical SI units, validation invariants, hardware
audit substrates, and forensic verification contracts. Statically
allocated in DTCM SRAM and committed to write-once media upon preemption
trigger.}\label{tbl-14-authority-log-schema}\tabularnewline
\toprule\noalign{}
\begin{minipage}[b]{\linewidth}\raggedright
Schema Field Identifier
\end{minipage} & \begin{minipage}[b]{\linewidth}\raggedright
Binary Type \& Layout
\end{minipage} & \begin{minipage}[b]{\linewidth}\raggedright
Physical SI Units
\end{minipage} & \begin{minipage}[b]{\linewidth}\raggedright
Verification \& Range Invariant
\end{minipage} & \begin{minipage}[b]{\linewidth}\raggedright
Hardware Audit Substrate
\end{minipage} & \begin{minipage}[b]{\linewidth}\raggedright
Forensic \& Systems Purpose
\end{minipage} \\
\midrule\noalign{}
\endfirsthead
\toprule\noalign{}
\begin{minipage}[b]{\linewidth}\raggedright
Schema Field Identifier
\end{minipage} & \begin{minipage}[b]{\linewidth}\raggedright
Binary Type \& Layout
\end{minipage} & \begin{minipage}[b]{\linewidth}\raggedright
Physical SI Units
\end{minipage} & \begin{minipage}[b]{\linewidth}\raggedright
Verification \& Range Invariant
\end{minipage} & \begin{minipage}[b]{\linewidth}\raggedright
Hardware Audit Substrate
\end{minipage} & \begin{minipage}[b]{\linewidth}\raggedright
Forensic \& Systems Purpose
\end{minipage} \\
\midrule\noalign{}
\endhead
\bottomrule\noalign{}
\endlastfoot
\texttt{timestamp\_ptp\_ns} & \texttt{uint64\_t} (\(8\text{ B}\)) &
Nanoseconds (\(\text{ns}\)) & Monotonic;
\(|\Delta t_{\text{PTP}}| \le 100\text{ ns}\) & Ethernet PHY / PTP
hardware timer & Causal ordering against external reference clock \\
\texttt{actuator\_channel\_mask} & \texttt{uint16\_t} (\(2\text{ B}\)) &
Bitmask (\(16\text{ channels}\)) & \(\sum \text{channel}_i \ge 1\);
non-zero active mask & Hardware Memory Protection Unit (MPU) & Binds
command to physical motor inverters \\
\texttt{authority\_state} & \texttt{uint8\_t} (\(1\text{ B}\)) & Enum
(\(0\text{--}5\)) &
\(S_{\text{auth}} \in \{\text{AUTO}, \text{PEND}, \text{BLEND}, \text{MAN}, \text{FALLBACK}, \text{ESTOP}\}\)
& Enforcer State Arbiter Register & Discrete mutually exclusive
authority state \\
\texttt{operator\_role\_id} & \texttt{uint8\_t} (\(1\text{ B}\)) & Enum
/ Role ID & Authenticated role \(\in \text{CapabilityMatrix}\) & Crypto
Hardware Co-Processor & Proof of legitimate intervention role \\
\texttt{preempt\_trigger\_code} & \texttt{uint16\_t} (\(2\text{ B}\)) &
Bitfield & Exact trigger (Torque, Timeout, Deadman, E-stop) & Nested
Vectored Interrupt Controller (\texttt{NVIC}) & Deterministic cause of
state transition \\
\texttt{blend\_factor\_alpha} & \texttt{float32\_t} (\(4\text{ B}\)) &
Unitless \([0.0, 1.0]\) &
\(\alpha(t) = 3\hat{t}^2 - 2\hat{t}^3 \in [0.0, 1.0]\) & Microcontroller
FPU Blending Governor & Continuous blending factor progress \\
\texttt{policy\_proposed\_cmd} & \texttt{float32\_t{[}4{]}}
(\(16\text{ B}\)) &
\(\text{N}\cdot\text{m}, \text{rad/s}, \text{m/s}^2\) & Unconstrained
policy output \(\mathbf{u}_{\text{prop}}\) & Lock-Free Inter-Core SRAM
Mailbox & Autonomous model intent verification \\
\texttt{enforcer\_clamped\_cmd} & \texttt{float32\_t{[}4{]}}
(\(16\text{ B}\)) &
\(\text{N}\cdot\text{m}, \text{rad/s}, \text{m/s}^2\) &
\(h(\mathbf{x}, \mathbf{u}_{\text{enf}}) \ge 0\); barrier-projected &
Cortex-R5 Real-Time Safety Filter & Proof of barrier filter
intervention \\
\texttt{actuator\_delivered\_cmd} & \texttt{float32\_t{[}4{]}}
(\(16\text{ B}\)) & \(\text{N}\cdot\text{m}, \text{A}, \text{PWM \%}\) &
Inverter commutation setpoint \(\mathbf{u}_{\text{act}}\) & PWM Timer
Compare Registers (\texttt{TIMx\_CCR}) & Actual command dispatched to
motor phases \\
\texttt{measured\_plant\_state} & \texttt{float32\_t{[}4{]}}
(\(16\text{ B}\)) &
\(\text{m/s}, \text{rad}, \text{A}, \text{N}\cdot\text{m}\) & Ground
truth telemetry \(\mathbf{x}_{\text{plant}}(t)\) & Dual Optical Encoders
\& Shunt ADCs & Physical dynamic response and slip tracking \\
\texttt{prev\_block\_hmac\_chain} & \texttt{uint8\_t{[}32{]}}
(\(32\text{ B}\)) & SHA-256 Digest &
\(H_k = \text{HMAC-SHA256}(H_{k-1} \parallel \text{Payload}_k)\) &
Secure Enclave Hardware Monotonic Root & Cryptographic tamper-evidence
across power cycles \\
\end{longtable}

Mapping these entries to the physical dynamics of the machine requires
bounding four timing parameters against platform limits across the three
archetypes: the operator reaction budget \(T_{\text{react}}\), the
handshake protocol latency \(T_{\text{handshake}}\), the timeout
threshold \(T_{\text{timeout}}\), and the minimum blending window
\(\tau_{\text{blend}}\). The formal kinematic and timing specifications
for Mobile Systems, Manipulators, and Process/Energy Systems are
detailed in Table~\ref{tbl-14-authority-timing-contract}.

\begin{longtable}[]{@{}
  >{\raggedright\arraybackslash}p{(\linewidth - 12\tabcolsep) * \real{0.1429}}
  >{\raggedright\arraybackslash}p{(\linewidth - 12\tabcolsep) * \real{0.1429}}
  >{\raggedright\arraybackslash}p{(\linewidth - 12\tabcolsep) * \real{0.1429}}
  >{\raggedright\arraybackslash}p{(\linewidth - 12\tabcolsep) * \real{0.1429}}
  >{\raggedright\arraybackslash}p{(\linewidth - 12\tabcolsep) * \real{0.1429}}
  >{\raggedright\arraybackslash}p{(\linewidth - 12\tabcolsep) * \real{0.1429}}
  >{\raggedright\arraybackslash}p{(\linewidth - 12\tabcolsep) * \real{0.1429}}@{}}
\caption{\textbf{Authority Timing and Kinematic Verification Contracts
Across the Three Archetypes.} Deterministic physical timing bounds,
actuator jerk constraints, handshake execution budgets, and fallback
timeout contracts mapped across Mobile Systems, Manipulators, and
Process/Energy
Systems.}\label{tbl-14-authority-timing-contract}\tabularnewline
\toprule\noalign{}
\begin{minipage}[b]{\linewidth}\raggedright
Physical Machine Archetype
\end{minipage} & \begin{minipage}[b]{\linewidth}\raggedright
Actuator Channel \& Inertia
\end{minipage} & \begin{minipage}[b]{\linewidth}\raggedright
Max Allowable Jerk (\(j_{\text{allowable}}\))
\end{minipage} & \begin{minipage}[b]{\linewidth}\raggedright
Operator Reaction Floor (\(T_{\text{react}}\))
\end{minipage} & \begin{minipage}[b]{\linewidth}\raggedright
Handshake Latency (\(T_{\text{handshake}}\))
\end{minipage} & \begin{minipage}[b]{\linewidth}\raggedright
Max Blending Window (\(\tau_{\text{blend}}\))
\end{minipage} & \begin{minipage}[b]{\linewidth}\raggedright
Timeout / Fallback Lease (\(T_{\text{timeout}}\))
\end{minipage} \\
\midrule\noalign{}
\endfirsthead
\toprule\noalign{}
\begin{minipage}[b]{\linewidth}\raggedright
Physical Machine Archetype
\end{minipage} & \begin{minipage}[b]{\linewidth}\raggedright
Actuator Channel \& Inertia
\end{minipage} & \begin{minipage}[b]{\linewidth}\raggedright
Max Allowable Jerk (\(j_{\text{allowable}}\))
\end{minipage} & \begin{minipage}[b]{\linewidth}\raggedright
Operator Reaction Floor (\(T_{\text{react}}\))
\end{minipage} & \begin{minipage}[b]{\linewidth}\raggedright
Handshake Latency (\(T_{\text{handshake}}\))
\end{minipage} & \begin{minipage}[b]{\linewidth}\raggedright
Max Blending Window (\(\tau_{\text{blend}}\))
\end{minipage} & \begin{minipage}[b]{\linewidth}\raggedright
Timeout / Fallback Lease (\(T_{\text{timeout}}\))
\end{minipage} \\
\midrule\noalign{}
\endhead
\bottomrule\noalign{}
\endlastfoot
\textbf{Mobile System} (Autonomous Vehicle / Rover) & Electric Power
Steering \& Traction Motors (\(m = 2400\text{ kg}\)) &
\(12.0\text{ m/s}^3\) (Chassis) / \(500\text{ N}\cdot\text{m/s}\)
(Column) & \(1.0\text{ s}\) (Attentive) / \(3.0\text{ s}\) (Gaze Decay)
& \(\le 20.0\text{ ms}\) (CAN-FD PTP) & \(75\text{ ms}\)
(\(\Delta \tau = 25\text{ Nm}\)) & \(100\text{ ms}\) (Deadman lease) /
\(500\text{ ms}\) (ODD Timeout) \\
\textbf{Manipulator} (6-DOF Cobot / Industrial Arm) & Harmonic Drive
Joint Actuators (\(J = 0.45\text{ kg}\cdot\text{m}^2\)) &
\(85.0\text{ rad/s}^3\) (Joint Jerk) / \(250\text{ N/s}\) (End-Effector)
& \(0.25\text{ s}\) (Direct Teach Pendant) & \(\le 2.0\text{ ms}\)
(EtherCAT Distributed Clock) & \(30\text{ ms}\)
(\(\Delta \tau = 15\text{ Nm}\)) & \(50\text{ ms}\) (Teach pendant
3-position deadman release) \\
\textbf{Process / Energy System} (Hydro-Turbine / Steam Valve) &
High-Pressure Hydraulic Proportional Valve (\(P = 210\text{ bar}\)) &
\(15.0\text{ MPa/s}\) (Acoustic Water Hammer Limit) & \(2.5\text{ s}\)
(Supervisory SCADA Desk) & \(\le 10.0\text{ ms}\) (Profinet IRT / TSN) &
\(250\text{ ms}\) (\(\Delta P = 2.5\text{ MPa}\)) & \(1000\text{ ms}\)
(Remote link loss before Category 1 safe venting) \\
\end{longtable}

Consider an analytical ground vehicle operating at \(22\text{ m/s}\).
The operator reaction budget \(T_{\text{react}}\) sets the maximum
duration the nervous system permits between issuing a takeover demand
and detecting manual steering torque, budgeted at
\(T_{\text{react}} = 1.0\text{ s}\) for an attentive seated driver
(derived from the supervisory timing bounds in
Table~\ref{tbl-14-teleop-stability}). The handshake latency
\(T_{\text{handshake}}\) defines the round-trip deadline for mutual
authentication and state acknowledgment over the bus, budgeted at
\(T_{\text{handshake}} = 20\text{ ms}\) per the four-phase protocol in
Table~\ref{tbl-14-handshake-protocol}. The timeout threshold
\(T_{\text{timeout}}\) dictates the maximum communication silence
tolerated from an overriding entity before the nervous system declares a
link failure and reverts to the fallback trajectory, set to
\(T_{\text{timeout}} = 50\text{ ms}\) per the state transition rules in
Table~\ref{tbl-14-authority-state-matrix}. Finally, the blending window
\(\tau_{\text{blend}}\) governs the rate of authority transfer to
prevent mechanical shock. For an electric power steering column with a
maximum allowable jerk limit of \(500\text{ N}\cdot\text{m/s}\),
transitioning an automated \(25\text{ N}\cdot\text{m}\) torque bias to
zero requires an analytical blending duration of
\(\tau_{\text{blend}} \ge \frac{1.5 \times 25\text{ N}\cdot\text{m}}{500\text{ N}\cdot\text{m/s}} = 75\text{ ms}\).

These four timing parameters directly determine the spatial margin
available to avoid hazards. At a velocity of \(22\text{ m/s}\), the
combined transition interval
\(T_{\text{react}} + T_{\text{handshake}} + \tau_{\text{blend}} = 1000\text{ ms} + 20\text{ ms} + 75\text{ ms} = 1095\text{ ms}\)
consumes \(23.54\text{ m}\) of travel before the human operator commands
full physical steering authority. If the forward perception stack
detects a stationary obstacle at a sensing horizon of \(65\text{ m}\)
and maximum emergency braking requires \(35\text{ m}\), the total budget
of \(23.54\text{ m} + 35\text{ m} = 58.54\text{ m}\) leaves a clearance
margin of \(6.46\text{ m}\). Any unbudgeted jitter or dropped handshake
packet consumes this remaining distance, turning an evasion into a
collision.

This explicit schema forms the foundation for automated verification.
Chapter~\ref{sec-verification} uses the authority record to generate
automated hardware-in-the-loop fault-injection suites that
systematically probe every failure mode of the transition interface. The
verification suite tests four deterministic invariants. First, no two
entities may hold non-zero authority weights summing to more than
\(1.0\) on any shared axis. Second, no command vector from an entity on
a lower preemption tier may reach an actuator while a higher-tier entity
is authenticated and within its active session. Third, an unacknowledged
handshake must trigger the deterministic fallback state precisely at
\(t = T_{\text{timeout}}\). Fourth, any cryptographic signature failure
or bitfield corruption must immediately lock out the unauthenticated
channel. By injecting synthetic fault conditions (such as dropping
twenty consecutive CAN frames, simulating a frozen operator holding
constant torque during a planned evasion, or replaying forged override
packets), the test framework verifies that the nervous system maintains
its physical safety boundaries without software deadlocks or torque
contention.

When assembling the safety case for operational certification,
Chapter~\ref{sec-release} evaluates the authority record as primary
evidence for release. Regulators and safety assessors do not audit
neural network weights or heuristic arbitration code; they evaluate the
mathematical proof that the nervous system enforces the bounds specified
in the record under all operating regimes. The certification package
pairs the static authority ledger with empirical execution traces
gathered from the fault-injection campaign, demonstrating that the
maximum measured preemption latency \(P_{99.99}\) never exceeds
\(T_{\text{handshake}} + \tau_{\text{blend}}\), and that every
intervention leaves an unbroken, verifiable audit trail.

A verified authority record guarantees that the machine resolves command
contention deterministically in silicon and code. But verifying that the
nervous system switches authority according to specification does not
guarantee that the resulting physical vehicle dynamics remain stable
under real-world tarmac friction, worn actuators, and split-second
driver reflexes. A handover that executes flawlessly at the bus protocol
level can still induce dynamic rollover or road departure if the
physical machine and the human operator are pushed beyond their combined
operational limits.

\section{Systems Stress-Tests}\label{sec-intervention-14-9}

A handover that succeeds mechanically can still fail operationally. When
the nervous system switches authority registers without bus errors,
dropped packets, or torque deadlocks, the physical machine and its human
operator may still enter unrecoverable failure modes if the timing,
physical dynamics, adversarial communication loads, or downstream data
pipelines violate the assumptions of the control architecture. The
following four stress-tests evaluate where clean algorithmic handovers
break down in physical execution:

\subsection*{Stress-Test 1: The Incapacitated
Operator}\label{stress-test-1-the-incapacitated-operator}
\addcontentsline{toc}{subsection}{Stress-Test 1: The Incapacitated
Operator}

\emph{The scenario.} An autonomous transport platform
(\(m = 2400.0\text{ kg}\)) travels at \(v = 15.0\text{ m/s}\). Occlusion
drops perception confidence, triggering a handover request at
\(t = 0.0\text{ s}\)
(\(T_{\text{react}} = 1.50\text{ s}, T_{\text{timeout}} = 2.00\text{ s}\)).
An incapacitated human applies zero torque. Clearance to a stationary
barrier ahead is \(42.0\text{ m}\).

\emph{The systems dilemma.} Waiting \(2.00\text{ s}\) for acknowledgment
traverses
\(d_{\text{wait}} = 15.0\text{ m/s} \times 2.00\text{ s} = 30.0\text{ m}\).
At \(t = 2.00\text{ s}\), remaining clearance is \(12.0\text{ m}\).
Maximum braking on dry concrete
(\(\mu = 0.70, a_{\text{brake}} = 6.87\text{ m/s}^2\)) requires stopping
distance
\(d_{\text{stop}} = v^2 / (2 a_{\text{brake}}) = 16.38\text{ m}\),
causing a \(v_{\text{impact}} = 7.75\text{ m/s}\) barrier crash before
the vehicle can halt.

\emph{The systems verdict: Spatial Handover Budget Law.} Handover is an
allocation of physical distance, not digital time. If the time allocated
for human response exceeds the spatial stopping budget of the remaining
clearance, the handshake protocol guarantees an impact;
\(T_{\text{timeout}}\) must be dynamically bounded by physical stopping
margins.

\subsection*{Stress-Test 2: Destructive Panic
Override}\label{stress-test-2-destructive-panic-override}
\addcontentsline{toc}{subsection}{Stress-Test 2: Destructive Panic
Override}

\emph{The scenario.} A delivery robot (\(m = 350.0\text{ kg}\)) travels
at \(v = 4.0\text{ m/s}\). A pedestrian steps out \(6.0\text{ m}\)
ahead. A panicking operator applies
\(\tau_{\text{driver}} = 18.0\text{ N}\cdot\text{m}\), taking full
manual steering (\(\alpha = 1.0\)) and steering \(\delta = 35.0^\circ\)
directly toward a plate-glass partition \(2.5\text{ m}\) away.

\emph{The systems dilemma.} The operator command avoids the pedestrian
but guarantees structural impact in \(<650.0\text{ ms}\). The real-time
enforcer evaluates Control Barrier Functions at \(1000\text{ Hz}\).
Rather than granting unconstrained register writes, the enforcer
computes a minimum-norm projection of the manual steering vector onto
the half-space of admissible inputs
(\(\nabla h(\mathbf{x}) \cdot \dot{\mathbf{x}} \ge -\gamma h(\mathbf{x})\)),
steering along the safe tangent without inducing column fighting.

\emph{The systems verdict: Invariant Priority Invariant.} Manual
takeover grants authority to select among safe trajectories, never
authority to command the plant outside its physical survival envelope;
human inputs remain unprivileged proposals subject to deterministic
barrier gating.

\subsection*{Stress-Test 3: Babbling Field
Bus}\label{stress-test-3-babbling-field-bus}
\addcontentsline{toc}{subsection}{Stress-Test 3: Babbling Field Bus}

\emph{The scenario.} A heavy manipulator (\(m = 4500.0\text{ kg}\))
executes safety invariants at \(1000\text{ Hz}\) on an isolated
real-time core over a CAN-FD bus. Takeover frames require cryptographic
AES-128-GMAC verification (\(18.5\ \mu\text{s}\) hardware cycle). A
babbling node floods the bus with spoofed frames at
\(8000\text{ frames/s}\).

\emph{The systems dilemma.} Servicing interrupts (\(45.0\ \mu\text{s}\))
and verifying tags (\(18.5\ \mu\text{s}\)) across
\(8000\text{ frames/s}\) consumes \(50.8\%\) of CPU time
(\(508.0\text{ ms/s}\)), inflating safety loop execution from
\(620.0\ \mu\text{s}\) to \(1380.0\ \mu\text{s}\) (\(P_{99.9}\)) and
tripping hardware watchdogs. Hardware mailbox filtering and silicon
token-bucket rate limiters must drop unauthenticated packets before CPU
interrupt assertion.

\emph{The systems verdict: Hardware Policing Invariant.} Cryptographic
authentication consumes silicon cycles and bus bandwidth. Security
mechanisms lacking hardware-level traffic policing defend against
unauthorized commands by creating a real-time availability failure.

\subsection*{Stress-Test 4: Cascading Retraining
Degradation}\label{stress-test-4-cascading-retraining-degradation}
\addcontentsline{toc}{subsection}{Stress-Test 4: Cascading Retraining
Degradation}

\emph{The scenario.} A continuous learning pipeline ingests \(4200\)
human intervention logs from depot tractors (\(500{,}000\text{ h}\)
operation), training imitation models on the \(30\text{-second}\)
segments preceding each takeover across 10 generations without causal
filtering.

\emph{The systems dilemma.} Only \(38\%\) of interventions represent
true policy failures; \(44\%\) are cautious early taps and \(18\%\) are
driver mistakes. Ingesting the \(3.0\text{--}5.0\text{ s}\) pre-takeover
reaction window teaches the policy to reproduce human hesitation,
micro-steering twitches, and erratic braking, increasing intervention
rates by \(340\%\) by generation 8. The pipeline must truncate
pre-takeover reaction windows and train strictly on post-takeover
recovery trajectories where \(\alpha = 1.0\).

\emph{The systems verdict: Intervention Data Curation Law.}
Interventions are not expert demonstrations. Training on raw takeover
logs teaches the model to imitate the human hesitation and delayed
reflexes that precipitated the intervention.

\section{Summary}\label{sec-intervention-14-10}

Human authority over a moving machine cannot function as an informal
safety net or an unmodeled override switch. Authority is an engineered
real-time interface defined by explicit operational scope, cryptographic
identity, and deterministic timing bounds. While enforcement in
Chapter~\ref{sec-enforcement} provides the reflex mechanisms that refuse
an unsafe model action, intervention addresses the complementary problem
of acquiring authority and revoking the autonomy of an executing policy.
When an operator touches a steering yoke, depresses an emergency pedal,
or issues an override command across an authenticated network link, the
nervous system must arbitrate that request within strict timing
deadlines before the physical state diverges. Consider an analytical
example of a vehicle traveling at \(15.0\text{ m/s}\)
(\(54\text{ km/h}\)). If the nervous system requires \(200\text{ ms}\)
of round-trip bus negotiation to authenticate an operator token and
resolve authority conflicts, the vehicle travels \(3.0\text{ m}\) along
an unmanaged trajectory before the takeover begins. Treating
intervention as an external exception rather than an explicit state
transition in the nervous system converts the takeover mechanism itself
into an unbudgeted failure mode.

Transferring control away from an active policy is never an
instantaneous step function. If an actuator abruptly swaps a learned
policy's torque command for a human operator's raw input within a single
\(1\text{ ms}\) control cycle, the sudden force discontinuity induces
driveline shock, wheel slip, and structural vibration. Bumpless transfer
requires continuous actuator blending, where an authority parameter
\(\alpha(t)\) transitions smoothly from zero to one along a trajectory
governed by the physical plant's maximum jerk limits. In an analytical
example of a \(2000\text{ kg}\) vehicle chassis with an actuator rate
limit of \(35\text{ Nm/s}\) and a lateral jerk bound of
\(2.5\text{ m/s}^3\), blending the learned policy and the human override
across a \(150\text{ ms}\) window prevents the transient torque spike
that would otherwise break tire traction. The nervous system enforces
these plant dynamics throughout the handover, guaranteeing that neither
the policy's residual momentum nor the human operator's sudden
counter-steering forces the physical plant beyond its envelope of
dynamic stability.

Every intervention event generates operational telemetry that directly
governs both post-incident liability analysis and model retraining.
Because an intervention marks the exact boundary where autonomous
control was revoked, the resulting logs require cryptographic
tamper-evidence. Hardware security modules must sign the Authority
Record at the bus level, writing timestamped hash chains to non-volatile
storage to create an immutable log of arbitration decisions. In a
post-incident forensic investigation, the absence of cryptographic
integrity reduces engineering claims to unverified assertions. Signed
telemetry provides the only empirical proof of whether a collision
resulted from an unhandled model failure, a late enforcer trigger, or an
erratic human override that forced an actuator into saturation.

For the machine learning training pipeline, raw intervention logs
represent a severe hazard if ingested without causal filtering. Human
operators do not intervene with the precision of optimal feedback
controllers. During the \(3.0\text{ s}\) to \(5.0\text{ s}\) reaction
window immediately preceding an intervention, human situational
awareness decays, producing hesitation, inconsistent pedal taps, and
erratic micro-steering inputs. Feeding these uncurated pre-takeover
windows into an imitation learning pipeline corrupts the state-action
distribution, training future model generations to replicate the very
human instabilities that precipitated the failure. An automated data
curation pipeline must apply causal truncation to discard the
pre-intervention reaction window, isolating only the verified recovery
trajectory where human authority reached full control and stabilized the
physical state.

\protect\phantomsection\label{callout-takeawaysux2a-14.8}
\begin{fbxSimple}{callout-takeaways}{Key Chapter Takeaways}{}
\phantomsection\label{callout-takeaways*-14.8}

\begin{itemize}
\tightlist
\item
  \textbf{Handover is an Allocation of Physical Distance, Not Digital
  Time:} Waiting \(2.0\text{ s}\) for human operator acknowledgment at
  \(15\text{ m/s}\) consumes \(30\text{ m}\) of clearance. If remaining
  clearance to an obstacle is less than
  \(d_{\text{wait}} + d_{\text{brake}}\), handover guarantees a
  collision; fallback braking must trigger dynamically based on spatial
  stopping budgets.
\item
  \textbf{Manual Takeover Remains Unprivileged (Invariant Priority):}
  Human operators panicking during emergencies can command inputs
  directly into obstacles. Operator inputs must be treated as
  unprivileged proposals (\(u_{\text{human}}\)) projected onto the
  admissible Control Barrier Function half-space to prevent structural
  collisions.
\item
  \textbf{Bumpless Transfer Requires Continuous Jerk-Bounded Blending:}
  Abruptly swapping control authority in a single \(1.0\text{ ms}\) step
  causes torque step discontinuities that excite driveline resonance and
  break tire traction. The authority parameter \(\alpha(t) \in [0, 1]\)
  must blend commands along a \(\mathcal{C}^2\) polynomial bounded by
  plant jerk limits.
\item
  \textbf{Hardware Policing Prevents Security DoS on Fieldbuses:}
  Cryptographic authentication consumes CPU cycles
  (\(18.5\ \mu\text{s}\) per AES-GMAC frame). Flooding buses with
  spoofed takeover requests starves the \(1000\text{ Hz}\) safety loop
  unless hardware CAN controllers drop unauthenticated frames before CPU
  interrupt assertion.
\item
  \textbf{Intervention Retraining Demands Causal Pre-Takeover
  Truncation:} Ingesting raw telemetry from the \(3\text{--}5\text{ s}\)
  window preceding human intervention teaches neural policies to imitate
  driver hesitation, micro-steering twitches, and erratic braking.
  Pipelines must truncate pre-takeover windows and train strictly on
  verified recovery trajectories (\(\alpha = 1.0\)).
\end{itemize}

\end{fbxSimple}

A physical AI system cannot be verified or certified until its authority
boundaries are formalized. By defining human authority as a
deterministic, cryptographically authenticated, and physically
rate-limited interface, the Authority Record provides the concrete
specifications required for verification and safety assurance. In
Chapter~\ref{sec-verification}, this record defines the explicit
fault-injection targets needed to test whether the machine maintains
physical stability under dropped override packets, delayed tokens, and
actuator saturation. In Chapter~\ref{sec-release}, the tamper-evident
records and bounded transfer guarantees supply the empirical evidence
required to construct a defensible safety case. Without an engineered
intervention interface, a physical AI system remains an unconstrained
experimental prototype. With it, the machine possesses a verifiable
boundary of control.

\chapter{\texorpdfstring{Adversarial
Verification\chaptersubtitle{Breaking the Assembled Machine to See If Its Claims Hold}}{Adversarial Verification}}\label{sec-verification}

\begin{figure}[H]

\centering{

\includegraphics[width=1\linewidth,height=\textheight,keepaspectratio]{index_files/mediabag//Users/VJ/GitHub/MLSysBook-integrate-physical/physical/book/chapters/15-verification/figures/fig_locator.pdf}

}

\caption{\label{fig-locator}\textbf{Systems Organ Focus: Nervous.} The
fault list derived from the records, as a traceable mapping from design
assertions to physical fault injection vectors across the qualification
ladder. Whole-system verification deterministically attacks each
cyber-physical boundary to falsify safety claims before deployment.}

\end{figure}%

\emph{Why is a robot that has operated for weeks without a single crash
not evidence that the machine is safe?}

A physical machine cannot be certified safe by running unperturbed
operational hours under nominal conditions. Fleet testing exercises only
the common interior of the operational design domain, leaving
catastrophic corner cases and hidden failure modes unprobed.
Whole-system verification is the adversarial engineering discipline of
subjecting the integrated machine to systematic, deterministic fault
injection across every stage of the qualification ladder (SIL \(\to\)
PIL \(\to\) HIL \(\to\) In-Situ). By deliberately corrupting memory
buffers, dropping sensor packets, inducing brownouts, and locking
actuator shafts on dynamometers, engineers measure the exact physical
boundaries of the survival envelope and falsify safety claims before
hardware is deployed into unconstrained environments.

\newpage{}

\protect\phantomsection\label{callout-objectiveux2a-15.1}
\begin{fbxSimple}{callout-objective}{Learning Objectives}{}
\phantomsection\label{callout-objective*-15.1}

After studying this chapter, you will be able to:

\begin{enumerate}
\def\labelenumi{\arabic{enumi}.}
\tightlist
\item
  State an operating envelope precisely enough that a condition outside
  it can be deliberately created.
\item
  Derive a fault list by walking the records from Chapter~\ref{sec-data}
  through Chapter~\ref{sec-intervention} and violating each stated
  property.
\item
  Inject a stale observation, a late plan, and an expired intent on real
  hardware, and record what happened.
\item
  Say what a given rung of the qualification ladder can establish and
  what it cannot.
\end{enumerate}

\emph{Core chapter takeaway.} The fault list is derived from what the
design claims, not chosen from a menu of techniques. Every property
recorded in Parts II and III is a fault waiting to be injected, and the
intervention mechanism is under test alongside everything else.

\end{fbxSimple}

\section{Fault Injection}\label{sec-verification-15-1}

\lettrine{A}{ flow controller}\lettrine{t}{hat} has operated inside an
industrial facility for three weeks without an operator intervention has
demonstrated exactly one property. It has proven that there exists at
least one sequence of benign physical states, undisturbed sensor
readings, and nominal actuator commands under which the policy
stabilizes the liquid level. This is an existence proof, and an
existence proof cannot establish a reliability metric, an operating
envelope, or a guarantee of fault tolerance. A machine that worked in
front of an observer is not evidence that it works, because an
uninstrumented demonstration records only a sample from an
uncharacterized, narrow slice of reality. If the ambient temperature
remained within \(20.0^\circ\text{C}\) and \(24.0^\circ\text{C}\), the
supply pressure stayed at \(3.0\text{ bar}\), and the upstream pump
never cavitated, the test stand did not evaluate the controller. It
evaluated the benign operating environment.

As mapped in the four-tier architectural locator
(Figure~\ref{fig-locator}), whole-system verification sits in the
nervous system tier as the adversarial auditor of the entire stack.
Every design property established across the Body (mechanical inertia,
thermal dissipation, actuator slew limits), the Brain (training
distribution bounds, inference latencies, intent validity windows), and
the Nervous System (bus arbitration, clock synchronization, safety
enforcers, authority transfer) is not a permanent truth, but an explicit
claim that must be attacked. The fault list is not selected arbitrarily
from a generic menu of techniques; it is derived directly by inverting
the architectural records and physical constraints of the assembled
machine.

The mathematical reason a successful trajectory cannot support a claim
over unobserved conditions rests on distribution support. Consider an
analytical example where a machine executes \(n = 10^6\) control cycles
at \(100\text{ Hz}\), corresponding to \(2.8\text{ hours}\) of
continuous operation, with zero observed violations of its state
constraints. Under a stationary operational distribution \(D_0\), this
sequence establishes an empirical failure rate of zero, allowing an
engineer to place a statistical upper bound on the per-cycle failure
probability on \(D_0\). But physical environments are neither stationary
nor confined to the center of their probability distributions. The true
operating environment \(D\) includes rare supply surges, sensor
degradation, and thermal drift whose joint probability mass lies in the
tail region \(D \setminus D_0\). Because the empirical evaluation
collected zero samples in this region, the observed data provides zero
information about the policy gradient, the actuator saturation limits,
or the numerical stability of the state estimator when forced outside
\(D_0\). A neural network policy that maps state observations to valve
positions is an unconstrained function outside the convex hull of its
training and testing trajectories. Its output on an unvisited state is
not merely uncertain; it is arbitrary.

Consider a flow regulator designed to maintain liquid in a
\(100.0\text{ L}\) buffer tank below an overflow weir at
\(95.0\text{ L}\). Under nominal conditions, liquid enters at
\(40.0\text{ L/min}\), an ultrasound sensor reports the liquid height
with a \(10\text{ ms}\) bus latency, and an electric actuator adjusts
the discharge valve at \(50\text{ Hz}\). In routine operation, the
liquid height hovers at \(60.0\text{ L}\), well clear of the boundary.
The machine appears successful because the disturbances never exceed
\(\pm 5.0\text{ L/min}\), and the discharge valve operates within its
linear flow regime between \(30\%\) and \(50\%\) open. Now introduce an
unobserved condition where an upstream valve switch causes a supply
pressure spike, driving the inlet flow to \(80.0\text{ L/min}\).
Simultaneously, a memory bus conflict on the controller board delays the
ultrasound reading by \(140\text{ ms}\), bringing the total sensing
latency to \(150\text{ ms}\). With a net fill rate of
\(40.0\text{ L/min}\), which is \(0.67\text{ L/s}\), the tank
accumulates volume rapidly. Because the policy receives delayed level
data, it underestimates the rate of rise near \(90.0\text{ L}\) and
commands the discharge valve to open late. When the stale reading
finally updates, the policy outputs a maximum-rate slew command that
slams the valve open, inducing a hydraulic pressure transient of
\(12.0\text{ bar}\) that cracks the downstream pipe coupling. The liquid
breaches the \(95.0\text{ L}\) weir, triggering an emergency float
switch that cuts pump power, but the structural damage to the pipe
fitting is detected by no runtime sensor at all until fluid pools on the
floor.

Converting an execution trace into verification evidence requires
structured methods designed to challenge specific system properties.
Five ingredients turn a physical run into a test result that can support
a safety case. First is a bounded claim, which states the exact
invariant the machine must maintain, such as keeping the liquid volume
below \(95.0\text{ L}\) and the transient line pressure below
\(6.0\text{ bar}\). Second are defined test conditions, specifying the
initial state space, fluid viscosity, ambient temperature, and supply
pressure range. Third is a precommitted oracle, an automated evaluator
that inspects the recorded telemetry against the claim without human
intervention or post-hoc adjustments. Fourth is a stated sampling or
injection procedure that deliberately drives the machine toward its
performance limits and injects faults into its sensors, buses, and
actuators. Fifth is provenance, the complete metadata record of model
weights, calibration tables, sensor serial numbers, and raw sensor logs
that permits an external engineer to reconstruct the execution and
reproduce the result.

Without these five elements, a thousand hours of spotless test stand
operation provide no evidentiary value. A machine that worked in front
of you is not evidence that it works; it is only evidence that it has
not yet encountered the combination of state, disturbance, and latency
that breaks it. To construct evidence that satisfies a safety argument,
the engineering task shifts from observing nominal behavior to
systematically constructing tests that attempt to violate every design
claim. That construction begins with the formal boundary that separates
the conditions where the machine is guaranteed to perform from the
conditions where it is permitted to fail.

\section{The Operating Envelope}\label{sec-verification-15-2}

The \textbf{operating envelope} defines the complete set of simultaneous
physical states, environmental disturbances, sensing conditions,
execution latencies, compute loads, electrical supply variations, and
operator response times under which the machine is claimed to satisfy
its performance specifications. Stating this envelope is not an exercise
in cataloging typical operating points. Instead, the boundary must be
defined with enough numerical precision that an automated test framework
can drive the system past each edge on purpose. For a fluid regulation
system, this means specifying not merely a nominal liquid level, but the
exact allowable ranges for fluid temperature, kinematic viscosity,
supply line pressure fluctuations, ambient temperature, acoustic
transducer signal-to-noise ratios, Controller Area Network (CAN) bus
arbitration jitter, central processing unit (CPU) thermal throttling
states, supply rail voltage droop, and the maximum permissible time
before an operator acknowledges an alert.

An operating envelope is distinct from the larger \textbf{survival
envelope} that encompasses it. Inside the operating envelope, the
learned policy and primary tracking controllers have full authority to
regulate the physical process, maintaining closed-loop tracking within
tight error tolerances. When external disturbances, sensor failures, or
compute stalls push the system across the boundary of the operating
envelope, nominal regulation is conceded. At that threshold, the machine
enters the survival envelope, where the primary objective shifts
entirely to bringing the physical body to a bounded, safe state before
structural limits are breached. In the flow regulator, the survival
envelope defines the broader region of volume, pressure, and temperature
within which emergency shutoff valves and mechanical relief mechanisms
can successfully vent pressure and arrest inflow before a pipe ruptures
or fluid overflows secondary containment. Beyond the boundary of the
survival envelope lies catastrophic physical failure, where stored
energy or incoming mass exceeds the mechanical dissipation capacity of
the hardware. This fundamental boundary distinction is visually and
physically manifested in destructive vehicle crash qualification
(Figure~\ref{fig-15-crash-test-instrumentation}): while nominal
autonomous policies operate strictly within the dynamic control
envelope, catastrophic perception or planning failures push the machine
across the boundary into the survival envelope, where structural crumple
zones and deterministic passive or pyrotechnic enforcers must physically
dissipate kinetic energy before passenger compartment survival limits
are breached.

\begin{figure}[t]

\centering{

\includegraphics[width=0.85\linewidth,height=\textheight,keepaspectratio]{/Users/VJ/GitHub/MLSysBook-integrate-physical/physical/book/chapters/15-verification/figures/fig15_real_crash_test_instrumentation.jpg}

}

\caption{\label{fig-15-crash-test-instrumentation}\textbf{Physical
Survival Envelope Verification and Impact Instrumentation.} Full-scale
vehicle crashworthiness and dynamic boundary qualification under
controlled laboratory impact conditions (55.28 km/h 90-degree Moving
Deformable Barrier side-impact test, NHTSA NCAP Test No.~10132).
High-contrast photogrammetric optical targets (yellow-and-black circular
fiducials) affixed along the vehicle B-pillar, door frames, and rocker
panels enable microsecond-resolution spatial deformation tracking via
high-speed photometric cameras. Heavy yellow sensor umbilical harnesses
emerging through the forward pillar route raw high-rate telemetry from
internal Anthropomorphic Test Device (ATD) accelerometers, chest
deflection potentiometers, and structural load cells to isolated
digitizers. When an unmodeled disturbance or policy failure pushes the
machine past its operating envelope, physical safety is no longer
governed by software optimization, but by the mechanical dissipation
capacity and deterministic passive/pyrotechnic enforcers within the
survival envelope. Attribution: U.S. National Highway Traffic Safety
Administration (NHTSA), Public Domain.}

\end{figure}%

Defining the boundary of the operating envelope requires deriving
physical limits from first principles rather than asserting convenient
software timeouts. Consider a pumped reservoir with an overflow weir at
\(V_{\text{max}} = 95.0\text{ L}\) and a nominal operating target of
\(80.0\text{ L}\). When the system operates at an instantaneous volume
of \(V(t) = 90.0\text{ L}\), the remaining volume margin is
\(\Delta V = V_{\text{max}} - V(t) = 5.0\text{ L}\). If the supply pump
delivers a maximum inflow of \(Q_{\text{in}} = 1.25\text{ L/s}\) while a
downstream valve failure drops discharge to
\(Q_{\text{out}} = 0.0\text{ L/s}\), conservation of volume dictates a
net accumulation rate of
\(dV/dt = Q_{\text{in}} - Q_{\text{out}} = 1.25\text{ L/s}\). The
remaining reaction time before liquid breaches the weir is
\(t_{\text{react}} = \Delta V / (dV/dt) = 5.0\text{ L} / (1.25\text{ L/s}) = 4.00\text{ s}\).
Every source of delay in the machine consumes a fraction of this
four-second budget. If an emergency isolation valve requires
\(1.00\text{ s}\) of mechanical transit time to seat against line
pressure, and fluid acoustic settling consumes \(0.20\text{ s}\), the
combined budget for sensing delay, inference execution, bus transport,
and safety enforcer evaluation cannot exceed \(2.80\text{ s}\) at
\(P_{99.9}\). If the instantaneous volume rises to
\(V(t) = 93.0\text{ L}\), the remaining volume margin shrinks to
\(2.0\text{ L}\), and the total reaction time drops to
\(1.60\text{ s}\). Subtracting the same \(1.20\text{ s}\) of mechanical
and hydraulic response leaves a total computation and communication
budget of only \(0.40\text{ s}\) before an unmitigated overflow becomes
physically inevitable.

Tabulating separate minimum and maximum values for individual variables
fails to capture the true operating boundary because physical and
computational dimensions are coupled. In isolation, an inlet pressure of
\(6.0\text{ bar}\), a fluid viscosity of
\(45.0\text{ mPa}\cdot\text{s}\), a valve slew rate of
\(30.0\text{ deg/s}\), and an observation age of \(120\text{ ms}\) might
each fall within their respective single-variable ranges. When these
conditions coincide, the higher viscosity increases fluid shear across
the valve spool, slowing mechanical travel below the nominal slew rate.
Simultaneously, the elevated inlet pressure accelerates volume
accumulation, while the \(120\text{ ms}\) observation age hides the
initial pressure surge from the policy. The physical state crosses into
the survival envelope long before any single parameter exceeds its
isolated maximum. The true operating boundary is an interconnected
hypersurface in state, parameter, and latency space. A valid
verification suite must probe this combined surface rather than checking
isolated scalar limits.

To prevent envelope boundaries from deteriorating into optimistic
assumptions, every defined parameter must specify its physical units,
its measurement provenance, its calibrated uncertainty, and a designated
subsystem engineer who remains accountable for its validity. A boundary
stated as a maximum pressure of \(8.0\text{ bar}\) is incomplete without
recording whether that limit derives from hydrostatic burst testing of
the flanged joints, finite-element pipe stress models, or pump stall
characteristics. The boundary specification must include the measurement
uncertainty of the sensors used during characterization, such as
\(\pm 0.15\text{ bar}\) at \(P_{99}\), and must name the individual
engineer authorized to accept or narrow the claim. When system
components change, such as swapping an elastomeric seal or upgrading a
CAN transceiver, the designated owner must formally re-evaluate whether
the physical boundary has shifted.

Verification of the operating envelope requires a disciplined sampling
strategy that deliberately attacks the boundary surface. Testing nominal
interior points, where the tank sits at \(50.0\text{ L}\), ambient
temperature is \(21.0^\circ\text{C}\), and sensing latency is
\(12\text{ ms}\), demonstrates only that the machine functions under
ideal conditions. To generate rigorous verification evidence, the test
framework must place test points directly on, immediately inside, and
immediately outside each consequential boundary. Evaluating points just
inside the boundary verifies that the policy maintains closed-loop
stability under maximum permissible stress. Evaluating points on the
boundary confirms that safety margins do not collapse under nominal
sensor noise. Evaluating points just outside the boundary confirms that
safety enforcers intervene predictably, taking authority from the
learned policy and bringing the system to a safe halt within the
survival envelope before containment is lost.

\protect\phantomsection\label{callout-ruleux2a-15.2}
\begin{fbxSimple}{callout-rule}{Architect’s Rule of Thumb:}{The Hardware Falsification Principle}
\phantomsection\label{callout-rule*-15.2}

If a safety property cannot be deterministically falsified by injecting
a physical or electrical fault on target hardware, the property is an
unverified assumption rather than an engineering invariant.

\end{fbxSimple}

Formally documenting the operating and survival boundaries converts the
design claims of the preceding chapters into concrete targets for fault
injection. When every physical constraint, latency budget, and
environmental limit is stated as a verifiable number, the safety
engineer knows precisely what conditions the machine must withstand,
what thresholds must trigger protective intervention, and what dynamic
claims must be subjected to deliberate failure.

\section{Deriving the Fault List}\label{sec-verification-15-3}

Deriving the initial fault list begins by inverting the design records
compiled across the architecture. Chapter~\ref{sec-data} through
Chapter~\ref{sec-evaluation} supply the properties of the learned
policy, including training distribution bounds, state-action
assumptions, and inference latency ceilings.
Chapter~\ref{sec-perception} through Chapter~\ref{sec-placement} provide
interface and execution properties, such as bus bandwidth allocations,
serialization overheads, observation freshness guarantees, and state
estimation bounds. Chapter~\ref{sec-intervention} supplies the
properties of the authority-transfer mechanism, specifically the maximum
latency to disconnect a failing policy and the independent override
authority of the safety enforcer. Inverting these records systematically
generates an initial candidate set for verification. If an analytical
model of a flow control loop records that the learned policy produces a
valid valve command every \(20\text{ ms}\) under all flow rates between
\(0.10\text{ m}^3\text{/s}\) and \(2.50\text{ m}^3\text{/s}\), the
verification suite inverts that claim by injecting stalled command
streams, delayed observations, and out-of-range sensor readings to test
whether the downstream nervous system detects the violation within its
budget.

Every recorded design property becomes a testable fault only when
structured as a five-element chain comprising the negated condition, the
physical injection point, the expected detector, the bounded response,
and an observable pass condition. In a high-pressure fluid delivery
system, consider a flow meter that publishes volumetric discharge rates
across a Controller Area Network (CAN) fieldbus with an analytical
latency bound of \(5.0\text{ ms}\). The negated property injects an
artificial transport delay of \(25\text{ ms}\) at the network interface
buffer. The expected detector is the stale-frame monitor in the safety
enforcer, the required response is a transition to a recirculating
bypass within \(10\text{ ms}\), and the observable pass condition is
that fluid pressure at the manifold remains below \(6.0\text{ bar}\) at
\(P_{99}\). When any element of this chain cannot be defined or
observed, the system contains a \textbf{verification gap}. A
verification gap indicates that the claim cannot be falsified by test.
Such gaps cannot be ignored; they require dedicated hardware
instrumentation, a narrower operational claim, or a redesign of the
execution path.

This inversion process must challenge data provenance, simulation
validity, and offline evaluation assumptions alongside physical runtime
packets. Distribution shift and misleading optimization metrics are
system faults even when every communication frame arrives without
corruption. If a learned policy was trained on hydrodynamic simulations
that assumed a constant fluid viscosity of
\(1.0\times 10^{-3}\text{ Pa}\cdot\text{s}\), cooling the fluid to
increase viscosity to \(3.5\times 10^{-3}\text{ Pa}\cdot\text{s}\)
alters the pump head curve and introduces cavitation noise. Inverting
the evaluation claim requires injecting these unmodeled fluid dynamics,
simulated sensor bias, and synthetic data perturbations where a high
nominal reward hides high-frequency actuator jitter. A physical AI
system does not fail solely through dropped packets; it fails when
valid, uncorrupted commands drive the physical body into an unstable
mechanical regime because the training assumptions were violated.

At the structural boundaries established across
Chapter~\ref{sec-perception} through Chapter~\ref{sec-placement}, the
fault generator inverts each inter-module contract at its physical point
of handoff. These tests negate observation freshness, state validity,
intent lifetime, planning deadlines, enforcement independence, and
resource allocations. If a trajectory planner issues an actuation
profile with an intent lifetime of \(100\text{ ms}\), the test
environment injects an otherwise valid actuation frame at
\(125\text{ ms}\) to confirm that the actuator interface rejects it as
expired. If the safety enforcer claims execution independence from the
primary policy computer, the test saturates the shared memory bus with
maximum direct memory access (DMA) traffic to confirm that the enforcer
still meets its \(1.0\text{ ms}\) control deadline. Conversely, fault
injection must remain disciplined, because bus starvation, clock skew,
and supply voltage fluctuations should be injected only where the design
explicitly claims tolerance. Testing arbitrary failure modes from a
generic menu invents requirements that the machine was never specified
to meet, diverting engineering effort away from validating actual
architectural contracts.

The structural limit of record inversion is absolute. Inverting recorded
properties produces only the faults that the design team already
modeled, which is precisely the set of behaviors that will not surprise
the creators of the machine. The completeness of an inverted fault list
is bounded by the completeness of the original design claims. If the
engineering records never documented how an electric motorized valve
behaves when ambient temperatures freeze the mechanical seals, no record
exists to invert, and the failure mode will not appear in the
verification suite. Inversion guarantees that every documented claim is
tested, but it remains blind to unstated assumptions, overlooked
physical hazards, and unmodeled failure dependencies.

Overcoming this limitation requires a second, independent source that
works forward from the physical machine rather than backward from design
documents. This forward analysis enumerates what the physical body, the
power distribution network, the fluid dynamics, and the operating
environment can do regardless of what the software believes. Physical
pipes experience pressure spikes from rapid valve closure known as water
hammer, mechanical valve seats accumulate particulate deposits that
prevent full seating, elastomeric seals wear under repeated pressure
cycling, and pump vibration induces harmonic resonance in downstream
sensor brackets. These physical phenomena originate in mechanics,
materials science, and fluid dynamics, operating entirely outside the
digital abstractions of the nervous system.\sidenote{\footnotesize Leveson (2011)
  demonstrated via STAMP/STPA that in complex cyber-physical
  architectures, catastrophic hazards often arise not from individual
  component failures, but from unsafe, unhandled interactions among
  correctly operating subsystems whose local control assumptions
  conflict.}

Working forward from the physical hardware exposes entire fault classes
that no software property record implies. These include total clock
loss, supply rail brownouts where logic levels float indeterminately
between \(0.8\text{ V}\) and \(2.0\text{ V}\), electromagnetic
interference from high-current pump contactors coupling into analog
pressure lines, floating output pins during power-on reset transitions,
corrupted non-volatile configuration memory, and shared reset lines that
bridge independent domains. If an emergency pressure-relief circuit and
a learned policy accelerator share a single power management integrated
circuit, an undervoltage transient on the accelerator can reset the
safety circuit, disabling the relief mechanism at the exact moment the
policy loses control. No software record contains this failure mode
because each software module assumes its underlying power rail is an
ideal, independent step function.

Physical forward analysis also challenges the common-cause independence
claims established in Chapter~\ref{sec-nervous} and
Chapter~\ref{sec-placement}. Two control paths that appear isolated in
an architectural diagram often share a single physical vulnerability. A
primary flow sensor and a redundant pressure switch may route their
wiring through a single conduit exposed to mechanical abrasion, or share
a common analog reference voltage within a multi-channel converter, or
rely on identical timer silicon vulnerable to the same manufacturer
errata. If an electromagnetic transient induces simultaneous latch-up
across both channels, treating them as statistically independent
failures invalidates the entire safety argument. Common-cause analysis
traces physical routing, power paths, and shared clock trees to ensure
that no single physical event can disable both the operational policy
and the intervention mechanism simultaneously.

The physical source must also encompass operator error modes and
foreseeable misuse from Chapter~\ref{sec-intervention}, treating human
interaction as an operational hazard of the machine. Field operators
force bypass switches during maintenance, misconfigure manual throttling
valves upstream of automated manifolds, disconnect telemetry cables to
silence acoustic alarms, and request setpoint changes that exceed the
thermal dissipation capacity of pump motors. When an operator commands
an instantaneous step increase in flow from \(0.20\text{ m}^3\text{/s}\)
to \(2.00\text{ m}^3\text{/s}\), the verification framework must verify
that the nervous system clamps the command to a safe acceleration rate,
preventing cavitation and hydraulic shock while maintaining closed-loop
stability.

Finally, the fault inventory must directly challenge the
authority-transfer claims established in Chapter~\ref{sec-intervention}.
That chapter claims that authority can be taken from a learned policy
within a strict physical deadline while the machine is in motion, and
that the policy cannot resist or delay its own eviction. A policy
driving a motorized valve at full speed creates mechanical momentum and
electrical back-EMF that resist rapid deceleration. In an analytical
test scenario where a valve actuator rotates at \(90^\circ\text{/s}\),
the fault injection framework triggers an unrecoverable policy fault and
measures whether the safety enforcer de-energizes the motor driver
within \(4.0\text{ ms}\), engages the mechanical brake, and halts flow
before line pressure exceeds the pipe burst threshold of
\(10.0\text{ bar}\). The intervention mechanism is not an unassailable
arbiter; it is a physical subsystem whose claims of speed, independence,
and override capability must be deliberately attacked, measured, and
falsified before the assembled machine can be judged ready for
structured qualification.

Working forward from the physical machine establishes an exhaustive
cross-layer taxonomy that connects each physical failure mechanism
directly to its target manifestation, detection latency budget, enforced
safety response, and observable pass criterion, as detailed in
Table~\ref{tbl-15-fault-classification}.

\begin{longtable}[]{@{}
  >{\raggedright\arraybackslash}p{(\linewidth - 10\tabcolsep) * \real{0.1667}}
  >{\raggedright\arraybackslash}p{(\linewidth - 10\tabcolsep) * \real{0.1667}}
  >{\raggedright\arraybackslash}p{(\linewidth - 10\tabcolsep) * \real{0.1667}}
  >{\raggedright\arraybackslash}p{(\linewidth - 10\tabcolsep) * \real{0.1667}}
  >{\raggedright\arraybackslash}p{(\linewidth - 10\tabcolsep) * \real{0.1667}}
  >{\raggedright\arraybackslash}p{(\linewidth - 10\tabcolsep) * \real{0.1667}}@{}}
\caption{\textbf{Exhaustive Cross-Layer Hardware and Cyber-Physical
Fault Classification Matrix.} Structural mapping of physical and logical
failure modes across injection mechanisms, target system manifestations,
hardware and software detection primitives, enforced safety recovery
responses, and unambiguous empirical pass
criteria.}\label{tbl-15-fault-classification}\tabularnewline
\toprule\noalign{}
\begin{minipage}[b]{\linewidth}\raggedright
Failure Domain \& Fault Mode
\end{minipage} & \begin{minipage}[b]{\linewidth}\raggedright
Physical Mechanism \& Injection Point
\end{minipage} & \begin{minipage}[b]{\linewidth}\raggedright
Target Manifestation \& Causal Impact
\end{minipage} & \begin{minipage}[b]{\linewidth}\raggedright
Detection Primitive \& Latency Budget
\end{minipage} & \begin{minipage}[b]{\linewidth}\raggedright
Enforced Safe Response \& Bounded Recovery
\end{minipage} & \begin{minipage}[b]{\linewidth}\raggedright
Observable Pass Criterion
\end{minipage} \\
\midrule\noalign{}
\endfirsthead
\toprule\noalign{}
\begin{minipage}[b]{\linewidth}\raggedright
Failure Domain \& Fault Mode
\end{minipage} & \begin{minipage}[b]{\linewidth}\raggedright
Physical Mechanism \& Injection Point
\end{minipage} & \begin{minipage}[b]{\linewidth}\raggedright
Target Manifestation \& Causal Impact
\end{minipage} & \begin{minipage}[b]{\linewidth}\raggedright
Detection Primitive \& Latency Budget
\end{minipage} & \begin{minipage}[b]{\linewidth}\raggedright
Enforced Safe Response \& Bounded Recovery
\end{minipage} & \begin{minipage}[b]{\linewidth}\raggedright
Observable Pass Criterion
\end{minipage} \\
\midrule\noalign{}
\endhead
\bottomrule\noalign{}
\endlastfoot
\textbf{Sensor Corruption \& Signal Bias} & • Piezoelectric transducer
drift• Stuck-at-zero ADC register• Injected via DAC inline interceptor &
Stale or unphysical state estimate \(\hat{\mathbf{x}}(t)\); policy
commands erroneous slew rate & \(\chi^2\) innovation / plausibility
monitor (\(\Delta t_{\text{det}} \le 20\text{ ms}\)) & Isolate faulty
sensor channel; switch to dead-reckoning or secondary estimator &
Process state remains inside survival envelope:
\(P(t) \le P_{\text{max}} + \Delta P_{\text{margin}}\) \\
\textbf{Bus Contention \& Frame Loss} & • Babbling-idiot CAN
transceiver• Real-time Ethernet queue overflow• Injected via CAN
disturbance node & Missing observation frames; heartbeat starvation;
outdated plan execution & Cyclic redundancy check (CRC) \& frame arrival
watchdog (\(\Delta t_{\text{det}} \le 10\text{ ms}\)) & Reject expired
proposal; freeze actuator to zero-torque safe hold or clamp valve &
End-to-end reaction time \(t_{\text{react}} \le t_{\text{budget}}\);
zero unhandled bus exceptions \\
\textbf{Power Brownout \& Voltage Droop} & • Motor contactor inrush
transient• LDO rail sag (\(V_{\text{core}} < 0.85\text{ V}\))• Injected
via programmable DC load & Indeterminate logic states; corrupted SRAM
registers; asynchronous MCU reset & Hardware Undervoltage-Lockout (UVLO)
monitor (\(\Delta t_{\text{det}} \le 15\ \mu\text{s}\)) & Hardware
crowbar circuit fires; de-energize gate drivers; assert fail-safe spring
return & Logic pins never float in undefined state; secondary relief
valve seats within \(50\text{ ms}\) \\
\textbf{Actuator Jam \& Mechanical Stall} & • Particulate contamination
in valve• Bearing seizure under shock load• Injected via hydraulic brake
clamp & Locked rotor; severe back-EMF collapse; rapid stator winding
heating (\(I^2 R\)) & Motor phase overcurrent \& Hall-effect velocity
stall monitor (\(\Delta t_{\text{det}} \le 5\text{ ms}\)) & De-energize
H-bridge power stage; engage electromechanical brake; route bypass &
Stator temperature \(T_{\text{coil}} < T_{\text{melt}}\); line pressure
relieved before burst disc rupture \\
\textbf{Watchdog Panic \& Priority Inversion} & • DMA bus-master lockup
on GPU• Deadlock in high-priority RTOS task• Injected via bus stall
emulator & Interrupted control loop cadence; un-serviced safety
heartbeat; hung NPU & Independent windowed hardware watchdog timer
(\(\Delta t_{\text{det}} \le 2.0\text{ ms}\)) & Assert Non-Maskable
Interrupt (NMI); hardwired authority transfer to safety MCU & Takeover
completed in \(\le 4.0\text{ ms}\); zero command bleed-through from hung
processor \\
\end{longtable}

\section{The Qualification Ladder}\label{sec-verification-15-4}

Subjecting an assembled machine to an inventory of faults requires
choosing an execution environment for each test. In safety engineering,
these environments form a structured progression known as the
\textbf{qualification ladder}, spanning Software-in-the-Loop (SIL),
Processor-in-the-Loop (PIL), Hardware-in-the-Loop (HIL), and In-Situ
Physical Fault Injection (Figure~\ref{fig-15-qualification-ladder}).
Each rung of this ladder exercises a different subset of the physical
and computational causal loop that connects sensor inputs, inference
pipelines, communication buses, actuator drives, plant dynamics, and
environmental feedback. The central engineering reality of qualification
is that these environments do not form a strict hierarchy where higher
rungs subsume lower ones. Each mode closes, breaks, or substitutes
specific causal links, establishing distinct categories of evidence
while remaining blind to others.

\protect\phantomsection\label{callout-definitionux2a-15.3}
\begin{fbxSimple}{callout-definition}{Definition:}{The Four-Stage Qualification Ladder}
\phantomsection\label{callout-definition*-15.3}

The structured progression of cyber-physical evaluation
environments---spanning Software-in-the-Loop (SIL),
Processor-in-the-Loop (PIL), Hardware-in-the-Loop (HIL), and In-Situ
Physical Fault Rigs---that systematically closes, breaks, or substitutes
specific arcs in the closed causal loop to balance combinatorial test
throughput against physical and temporal execution fidelity.

\end{fbxSimple}

Offline trace replay feeds recorded sensor data into the computational
pipeline, evaluating how the software responds to historical operational
sequences. Because the input signals come from physical sensors
operating in real environments, replay preserves the true noise
profiles, quantization anomalies, and electrical transients of
operational hardware, such as the high-frequency ripple from a turbine
flowmeter operating at \(1000\text{ Hz}\). Replay verifies whether the
software produces deterministic outputs, meets internal latency
deadlines, and handles historical edge cases without crashing. Yet
replay leaves the feedback loop broken. The recorded inputs are frozen
in time, completely decoupled from the commands the policy emits under
test. If an updated policy alters a valve opening command at an
analytical test offset of \(t = 50\text{ ms}\), the replay log continues
to supply the pressure measurements that resulted from the original
historical command. Replay establishes algorithmic reproducibility and
computation time under real input distributions, but it cannot establish
closed-loop dynamic stability or evaluate how the physical plant
responds when the policy diverges from the recorded path.

Software-in-the-loop (SIL) simulation restores the closed feedback loop
by substituting mathematical models for the physical plant, actuators,
sensors, and environmental dynamics on high-performance workstation or
cloud clusters. Within this synthetic environment, the simulation clock
is decoupled from wall-clock time, allowing an engineering team to
execute massive parallel parameter sweeps exceeding
\(10^8\text{ cycles/day}\). SIL enables dense adversarial testing across
millions of operating conditions, systematically sweeping line pressures
from \(1.0\text{ bar}\) to \(25.0\text{ bar}\), injecting tensor
bit-flips in neural network weight matrices, or commanding instantaneous
fluid flow steps without risking physical hardware destruction. SIL
establishes that the integrated control logic, state estimation filters,
and safety enforcers achieve mathematical stability across an expansive
parameter space. However, the validity of this evidence is bounded
strictly by the fidelity of the underlying numerical Ordinary
Differential Equation (ODE) or Partial Differential Equation (PDE)
models. If the simulator models a hydraulic valve as a pure first-order
lag with a fixed time constant of \(30\text{ ms}\), it omits physical
phenomena such as fluid-induced hydrodynamic torque, seal friction
hysteresis, and cavitation-induced vibration. Furthermore, running
software on a workstation host abstracts away the physical clock drift,
direct memory access contention, cache thrashing, and thermal throttling
of embedded target silicon. Simulation proves the mathematical
correctness of the control algorithm against an idealized model, but it
cannot prove that the model matches physical reality.

Processor-in-the-loop (PIL) evaluation closes the architectural gap
between host compilation and embedded execution by cross-compiling the
neural policy and safety control algorithms for the target Instruction
Set Architecture (ISA)---such as ARM Cortex-A/M cores or RISC-V
microcontrollers---and executing the resulting binary on actual target
silicon or cycle-accurate core emulators. The target processor is
interfaced to a host-simulated plant via high-speed co-simulation links
(such as JTAG, background debug interfaces, or cycle-synchronized shared
memory). PIL verification directly validates compiler optimization
levels (\texttt{-O3} versus \texttt{-Os}), target vector acceleration
libraries (such as ARM Helium or CMSIS-NN), fixed-point and quantized
tensor scaling artifacts (e.g., int8 versus float32 truncation error),
stack memory headroom, register spills, and worst-case execution time
(WCET) bounds. PIL exposes subtle software faults that never appear in
host SIL, such as compiler-induced register aliasing, unaligned memory
access faults, and arithmetic overflow in fixed-point state estimators.
Nevertheless, PIL operates across a synthetic peripheral boundary: it
simulates external buses and timers, remaining blind to physical
transceiver electrical contention, bus voltage droop, electromagnetic
interference, and dynamic actuator back-EMF.

Hardware-in-the-loop (HIL) bench testing mounts the physical target
computing hardware (the Linux MPU and real-time MCU Dual-Brain
architecture), communication fieldbuses, power supplies, and actuator
drive stages onto an instrumented laboratory stand, interfacing them in
hard real-time to high-bandwidth FPGA plant emulators and dynamic
dynamometer testbeds. This mode subjects the real target electronics to
physical electrical loads and authentic temporal jitter. HIL testing
exposes real-world bus arbitration delays across physical CAN and
Ethernet transceivers, microcontroller interrupt preemption latencies,
direct memory access (DMA) bus lockups, and the thermal rise of motor
driver power stages delivering \(10\text{ A}\) under continuous load.

\marginnote{\begin{footnotesize}

\textbf{Dynamometer Specifications}

Physical HIL benches use high-bandwidth four-quadrant regenerative
dynamometers (\(f > 1\text{ kHz}\)) to inject active counter-torques,
simulating aerodynamics and friction reversals. This physical loading
induces realistic motor back-EMF, inverter thermal dissipation, and
gearbox compliance.

\end{footnotesize}}

By coupling the actuator drive to a four-quadrant regenerative
dynamometer, HIL benches dynamically synthesize mechanical
counter-torques, injecting sudden joint stalls and high-rate load steps
to measure physical back-EMF collapse and gate-driver switching spikes.
HIL establishes that the nervous system executes reliably on actual
silicon, that clock drift between distributed microcontrollers remains
within budgeted bounds, and that physical safety switches disconnect
actuator power within the required physical deadline, such as
\(4.0\text{ ms}\) following an overpressure trigger. The limitation of
the HIL bench lies in its synthesized plant boundary: an FPGA model or
mechanical dynamometer cannot reproduce microscopic material fatigue,
fluid viscosity degradation under chemical shear, seal wear, or
destructive structural burst limits.

In-situ physical fault rigs and shadow deployments evaluate the fully
assembled physical machine---complete with production mechanics,
hydraulic conduits, wiring harnesses, and primary sensors---operating
under authentic operational stresses. Physical fault rigs operate within
reinforced laboratory containment cells equipped with dedicated
destructive actuators, including pyrotechnic valve jammers, pneumatic
line clamps, and magnetic brake clamps that physically seize rotating
shafts or obstruct high-pressure fluid lines. In-situ testing captures
the full unmodeled complexity of physical reality: non-linear Coulomb
friction reversals, turbulent fluid cavitation noise, structural
resonant vibrations, and thermal runaway. Alongside destructive rigs,
shadow operation deploys the fully assembled software and computing
hardware into the field alongside an incumbent controller, processing
live sensor telemetry from an operating industrial plant without
exerting physical control over the actuators. Shadow deployment
establishes that the perception and policy stack operates reliably under
true production workloads without software exceptions or false-positive
fault declarations. Yet shadow testing cannot establish the physical
consequences of the commands it generates. Because the shadow policy
commands are intercepted and discarded at the nervous system boundary
while the legacy controller drives the plant, the plant dynamics reflect
the legacy controller actions rather than the candidate policy. If the
candidate policy issues an unstable valve trajectory that would induce
severe water hammer, the shadow log records only the mathematical
command, while the physical pipe remains protected by the legacy
controller.

Because every qualification mode alters or severs different arcs in the
causal loop, treating these rungs as a cumulative staircase where higher
levels replace lower ones introduces dangerous verification gaps. A live
shadow deployment cannot intentionally command a valve to slam shut
against a \(15.0\text{ bar}\) head to measure the resulting pressure
surge, because doing so would destroy the production facility.
Conversely, a software simulator cannot reveal a race condition in
physical CAN transceiver silicon, and a hardware bench cannot execute
the tens of millions of randomized parameter variations required to
establish statistical confidence in edge-case handling. High-assurance
qualification relies on treating these four modes as complementary
evidence scopes, directing each fault injection test to the specific
environment capable of exercising the relevant causal mechanisms
(Figure~\ref{fig-15-qualification-ladder}). As structured in
Table~\ref{tbl-15-qualification-ladder}, the four qualification rungs
differ fundamentally in their compute substrates, causal-loop realism,
fault injection primitives, throughput scaling, and evidentiary limits.

\begin{figure}[t]

\centering{

\includegraphics[width=0.95\linewidth,height=\textheight,keepaspectratio]{index_files/mediabag//Users/VJ/GitHub/MLSysBook-integrate-physical/physical/book/chapters/15-verification/figures/fig15_qualification_ladder.pdf}

}

\caption{\label{fig-15-qualification-ladder}\textbf{The Four-Stage
Qualification Ladder and Causal-Loop Verification Matrix.} Systematic
progression across the throughput--fidelity continuum:
Software-in-the-Loop (SIL, \(\sim 10^6\text{ runs/day}\), \(0\%\)
hardware fidelity) executes massive parallel parameter sweeps over
idealized ODE dynamics; Processor-in-the-Loop (PIL,
\(\sim 10^4\text{ runs/day}\), \(40\%\) fidelity) validates target SoC
registers, interrupt latency tails, and zero-malloc RTOS execution
against FPGA plant models; Hardware-in-the-Loop Dynamo Bench (HIL,
\(\sim 10^2\text{ runs/day}\), \(85\%\) fidelity) exposes real inverter
switching, power rail droop, and motor back-EMF under 4-quadrant dynamic
loading; and In-Situ Physical Fault Injection
(\(\sim 10^0\text{ runs/day}\), \(100\%\) physical reality) subjects the
fully assembled machine inside a containment cell to destructive
hydraulic water hammer, mechanical jamming, and structural burst limits.
Each stage establishes a distinct evidentiary scope while remaining
blind to mechanisms outside its causal loop. Attribution: Original
textbook graphic, CC BY 4.0.}

\end{figure}%

\begin{longtable}[]{@{}
  >{\raggedright\arraybackslash}p{(\linewidth - 10\tabcolsep) * \real{0.1667}}
  >{\raggedright\arraybackslash}p{(\linewidth - 10\tabcolsep) * \real{0.1667}}
  >{\raggedright\arraybackslash}p{(\linewidth - 10\tabcolsep) * \real{0.1667}}
  >{\raggedright\arraybackslash}p{(\linewidth - 10\tabcolsep) * \real{0.1667}}
  >{\raggedright\arraybackslash}p{(\linewidth - 10\tabcolsep) * \real{0.1667}}
  >{\raggedright\arraybackslash}p{(\linewidth - 10\tabcolsep) * \real{0.1667}}@{}}
\caption{\textbf{The Four-Stage Qualification Ladder across
Cyber-Physical Boundaries.} Systematic decomposition of qualification
environments across computing substrates, exercised causal-loop
elements, fault injection primitives, execution throughput, physical and
temporal fidelity, and definitive test oracle
mechanisms.}\label{tbl-15-qualification-ladder}\tabularnewline
\toprule\noalign{}
\begin{minipage}[b]{\linewidth}\raggedright
Qualification Rung \& Substrate
\end{minipage} & \begin{minipage}[b]{\linewidth}\raggedright
Causal Loop Realism (Compute, Bus, Actuator, Plant)
\end{minipage} & \begin{minipage}[b]{\linewidth}\raggedright
Fault Injection Mechanism \& Scope
\end{minipage} & \begin{minipage}[b]{\linewidth}\raggedright
Throughput \& Sample Scaling
\end{minipage} & \begin{minipage}[b]{\linewidth}\raggedright
Physical \& Temporal Fidelity
\end{minipage} & \begin{minipage}[b]{\linewidth}\raggedright
Primary Test Oracle \& Evidentiary Limit
\end{minipage} \\
\midrule\noalign{}
\endfirsthead
\toprule\noalign{}
\begin{minipage}[b]{\linewidth}\raggedright
Qualification Rung \& Substrate
\end{minipage} & \begin{minipage}[b]{\linewidth}\raggedright
Causal Loop Realism (Compute, Bus, Actuator, Plant)
\end{minipage} & \begin{minipage}[b]{\linewidth}\raggedright
Fault Injection Mechanism \& Scope
\end{minipage} & \begin{minipage}[b]{\linewidth}\raggedright
Throughput \& Sample Scaling
\end{minipage} & \begin{minipage}[b]{\linewidth}\raggedright
Physical \& Temporal Fidelity
\end{minipage} & \begin{minipage}[b]{\linewidth}\raggedright
Primary Test Oracle \& Evidentiary Limit
\end{minipage} \\
\midrule\noalign{}
\endhead
\bottomrule\noalign{}
\endlastfoot
\textbf{Software-in-the-Loop (SIL)}\emph{(Host Workstation / Cloud
Cluster)} & • \textbf{Compute:} Synthetic / Host x86• \textbf{Bus:}
Idealized zero-copy buffer• \textbf{Actuator:} Linear / saturation
model• \textbf{Plant:} Numerical ODE/PDE solver & Tensor bit-flips,
synthetic sensor dropout, hydrodynamic parameter sweeps
(\(\mu, \rho, T\)) & \(10^3\text{--}10^5\times\) real-time parallel
sweeps (\(>10^8\text{ cycles/day}\)) & High mathematical flexibility;
zero target silicon, bus timing, or electrical realism &
\textbf{Oracle:} Mathematical assertion corridors and state
invariants.\textbf{Limit:} Cannot validate target execution timing or
unmodeled physical dynamics. \\
\textbf{Processor-in-the-Loop (PIL)}\emph{(Target SoC / MCU Core
Emulation)} & • \textbf{Compute:} Real target silicon (ARM/RISC-V)•
\textbf{Bus:} Simulated bus interconnect• \textbf{Actuator:} Emulated
command register• \textbf{Plant:} Real-time host model & Instruction
cache evictions, register corruption via JTAG/BIST, interrupt jitter
injection & \(0.1\text{--}1.0\times\) real-time; bound by cross-platform
synchronization & Instruction-set and compiler accurate; abstracts
analog bus lines and actuator back-EMF & \textbf{Oracle:} Target
register state and WCET execution bounds.\textbf{Limit:} Ignores
physical transceiver bus contention, voltage droop, and thermal rise. \\
\textbf{Hardware-in-the-Loop (HIL)}\emph{(Target ECUs + Real-Time Plant
Rig)} & • \textbf{Compute:} Production ECU / NPU• \textbf{Bus:} Real
physical CAN/Ethernet• \textbf{Actuator:} Real motor drives / dyno•
\textbf{Plant:} Hard real-time FPGA plant rig & Pin break/short, CAN
frame stuffing / babbling idiot, supply rail droop, load steps &
\(1.0\times\) strict real-time (\(24\text{ h/day}\) wall-clock limit) &
High electrical, timing, and protocol fidelity; bounded by FPGA plant
model accuracy & \textbf{Oracle:} Protocol CRC, bus deadline monitors,
and interlock timings.\textbf{Limit:} Cannot execute destructive
physical overload or simulate micro-cavitation wear. \\
\textbf{Physical Fault Rig \& Shadow}\emph{(Dynamometer Cell / Field
Plant)} & • \textbf{Compute:} Production ECU / NPU• \textbf{Bus:}
Production harness \& sensors• \textbf{Actuator:} Real physical
mechanisms• \textbf{Plant:} Real fluid / mechanical body & Pyrotechnic
valve jammers, magnetic brake stall, thermal heating, transducer bias
drift & \(1.0\times\) wall-clock; constrained by mechanical wear and
thermal reset (\(<10^3\text{ cycles/day}\)) & Full authentic physics
(Coulomb friction, acoustic wave shock, thermal dissipation, wear) &
\textbf{Oracle:} Independent physical containment instrumentation and
yield transducers.\textbf{Limit:} High financial cost per trial;
destructive tests cannot scale to statistical tail bounds. \\
\end{longtable}

To prevent false confidence, every reported verification result must
explicitly state which elements of the causal loop were physically real,
which were simulated, which were replayed from static logs, which were
bypassed, and which remained unobserved. An analytical engineering claim
that an intervention mechanism halted an overpressure event within
\(12\text{ ms}\) means nothing without specifying whether the pressure
signal was generated by a differential equation, transmitted over a
physical analog current loop, or replayed from a digital buffer.
Specifying this causal mapping exposes the assumptions underlying each
test, ensuring that unmodeled physics or artificial boundary conditions
do not disguise latent hazards.

The \textbf{causal-loop fidelity boundary} defines the structural
demarcation between the physical mechanisms and digital abstractions
exercised within a given test environment, defining the exact threshold
where synthetic models, emulated buses, or replayed sensor traces sever
the continuous bi-directional energy and information exchange between
compute and mechanics.

Once the test environment is defined and its causal boundaries are
accounted for, the qualification process requires concrete mechanisms to
evaluate whether the machine behavior during a fault injection
constitutes a pass or a failure.

\section{Fault Oracles and Acceptance
Criteria}\label{sec-verification-15-5}

A fault test without a precommitted detector and an acceptable response
is an anecdote about a machine that did something. When a fault is
injected into a running machine, the system will inevitably react,
producing telemetry transients, actuator motions, and acoustic
emissions. Without an unambiguous, pre-established oracle to evaluate
those reactions, engineers evaluate the result by looking at time-series
plots after the run, deciding whether the machine survived based on
whether any component burst or threw an unhandled exception. An
\textbf{oracle} is the mathematical predicate that evaluates a test
execution and returns a binary verdict of pass or fail, defined in full
before the execution begins.

In physical AI systems, evaluating whether a machine satisfied its
safety claims requires moving beyond static scalar bounds to continuous
temporal specifications. Physical systems do not fail at single isolated
instants; they fail through dynamic transients over time. To express and
verify these dynamic properties rigorously, verification engineering
uses \textbf{Signal Temporal Logic (STL)}. Given a continuous-time
cyber-physical trajectory \(\mathbf{x}: [0, T] \to \mathbb{R}^n\)
capturing state variables (e.g., fluid pressure \(P(t)\), flow rate
\(Q(t)\), actuator temperature \(T_{\text{coil}}(t)\)), an STL formula
\(\phi\) is constructed from atomic predicates
\(\mu \equiv (g(\mathbf{x}(t)) \ge 0)\), standard Boolean connectives,
and metric real-time temporal operators:
\[\phi := \mu \mid \neg \phi \mid \phi_1 \wedge \phi_2 \mid \phi_1 \vee \phi_2 \mid \square_{[a, b]} \phi \mid \Diamond_{[a, b]} \phi \mid \phi_1 \mathcal{U}_{[a, b]} \phi_2\]
where \(\square_{[a, b]}\) denotes the \emph{Always} operator (the
formula must hold at every time \(t' \in [t+a, t+b]\)),
\(\Diamond_{[a, b]}\) denotes the \emph{Eventually} operator (the
formula must hold at some time \(t' \in [t+a, t+b]\)), and
\(\mathcal{U}_{[a, b]}\) denotes the timed \emph{Until} operator.

Crucially, STL equips verification with \textbf{quantitative robustness
semantics}, defining a real-valued metric
\(\rho(\phi, \mathbf{x}, t) \in \mathbb{R}\) that measures not merely
\emph{whether} the specification held, but \emph{how far} the trajectory
remained from violating it: \[\begin{aligned}
\rho(\mu, \mathbf{x}, t) &= g(\mathbf{x}(t)) \\
\rho(\neg \phi, \mathbf{x}, t) &= -\rho(\phi, \mathbf{x}, t) \\
\rho(\phi_1 \wedge \phi_2, \mathbf{x}, t) &= \min\left(\rho(\phi_1, \mathbf{x}, t), \rho(\phi_2, \mathbf{x}, t)\right) \\
\rho(\phi_1 \vee \phi_2, \mathbf{x}, t) &= \max\left(\rho(\phi_1, \mathbf{x}, t), \rho(\phi_2, \mathbf{x}, t)\right) \\
\rho(\square_{[a, b]} \phi, \mathbf{x}, t) &= \min_{t' \in [t+a, t+b]} \rho(\phi, \mathbf{x}, t') \\
\rho(\Diamond_{[a, b]} \phi, \mathbf{x}, t) &= \max_{t' \in [t+a, t+b]} \rho(\phi, \mathbf{x}, t')
\end{aligned}\]

The quantitative robustness metric satisfies the foundational soundness
property: \[\begin{cases}
\rho(\phi, \mathbf{x}, 0) > 0 \implies \mathbf{x} \models \phi & (\text{Specification satisfied with physical margin } \rho) \\
\rho(\phi, \mathbf{x}, 0) < 0 \implies \mathbf{x} \not\models \phi & (\text{Specification violated by physical depth } |\rho|) \\
\rho(\phi, \mathbf{x}, 0) = 0 & (\text{Trajectory on critical failure boundary})
\end{cases}\]

Consider a high-pressure manifold regulating coolant flow at
\(10.0\text{ bar}\) with a nominal throughput of \(5.0\text{ L/s}\).
When a downstream valve jams under a \(15.0\text{ bar}\) head, the
physical safety requirement is formalized in STL as:
\[\phi_{\text{safe}} = \square_{[0, T]} \left( P(t) \le 16.5\text{ bar} \right) \wedge \square_{[0, T - 200\text{ ms}]} \left( \Diamond_{[0, 45\text{ ms}]} \left( Q_{\text{bypass}}(t) \ge 3.5\text{ L/s} \right) \right)\]
This formula states that at all times during the trial, line pressure
must remain below the \(16.5\text{ bar}\) containment limit, and
whenever a disturbance occurs, the bypass flow must exceed
\(3.5\text{ L/s}\) within \(45\text{ ms}\). If an injected fault causes
the peak pressure to reach \(17.2\text{ bar}\), the quantitative
robustness evaluates to
\(\rho(\phi_{\text{safe}}, \mathbf{x}, 0) = 16.5 - 17.2 = -0.7\text{ bar}\),
immediately quantifying the failure severity without manual inspection.

Quantitative semantics transforms whole-system verification from
unguided random testing into \textbf{temporal logic falsification as
adversarial optimization}. Because neural network policies operate
across high-dimensional continuous state spaces, uniform random sampling
(Monte Carlo fuzzing) requires an intractable number of trials to
encounter rare failure modes. By formulating the verification goal as
minimizing the STL robustness metric over a parameterized disturbance
space \(\mathcal{W}\), an automated falsification engine actively
searches for the worst-case disturbance vector \(\mathbf{w}^*\):
\[\mathbf{w}^* = \arg\min_{\mathbf{w} \in \mathcal{W}} \rho\left(\phi, \mathbf{x}_{\mathbf{w}}, 0\right)\]
where \(\mathbf{w}\) parameterizes disturbance parameters such as sensor
noise variance, network packet arrival jitter, supply voltage droop
depth, initial fluid temperature, and valve spool friction coefficients.

Falsification engines employ global black-box optimization
algorithms---including Covariance Matrix Adaptation Evolution Strategy
(CMA-ES), Bayesian Optimization with Gaussian Process surrogates, and
the Cross-Entropy Method---to navigate non-convex, discontinuous
cyber-physical landscapes. In differentiable simulation environments,
analytical gradients \(\nabla_{\mathbf{w}} \rho\) can be computed via
automatic differentiation through the neural policy and differentiable
physics simulator, enabling gradient-directed falsification. When the
optimization algorithm finds
\(\rho(\phi, \mathbf{x}_{\mathbf{w}^*}, 0) < 0\), it produces a
concrete, deterministic counterexample trajectory that falsifies the
safety claim, exposing the exact combination of physical disturbance and
latency that breaks the machine.

Alongside the pass and fail criteria, automated verification requires an
\textbf{invalid-test rule} that identifies when an experimental run must
be discarded rather than scored. A test run is invalid whenever the
execution environment fails to deliver the specified fault conditions or
violates the required operational preconditions. If an upstream supply
pump cavitates before the planned fault injection timestamp, lowering
the baseline manifold pressure to \(6.0\text{ bar}\) instead of the
required \(10.0\text{ bar} \pm 0.2\text{ bar}\), the test did not
evaluate the safety envelope under high-head conditions. Similarly, if
an injection framework injecting corrupted CAN frames suffers an
internal queue overflow and drops the fault packet before it reaches the
transceiver, the system under test was never subjected to the
disturbance. Scoring an unexercised run as a pass manufactures false
confidence, while scoring it as a failure penalizes the machine for an
error in the test fixture. The invalid-test rule codifies the boundary
conditions, rejecting the run and resetting the test sequence whenever
preconditions fail.

Precommitment is the engineering discipline that separates verifiable
evidence from post-hoc storytelling. When an assembled machine undergoes
physical or electrical fault injection, the resulting behavior is almost
always noisy, complex, and surprising. If acceptance criteria are
negotiated after observing the machine react, human engineers
systematically rationalize anomalies. An unexpected pressure oscillation
of \(17.5\text{ bar}\) is explained away as an acceptable structural
surge, or a \(60\text{ ms}\) reaction delay is excused because the
system eventually brought the flow back to steady state without tripping
the mechanical burst disc. Post-hoc evaluation converts safety
verification into a negotiation, where the definition of success drifts
to match whatever the machine happened to do.

Precommitment demands that the STL specification \(\phi\), the
robustness evaluation function \(\rho\), the response-time bounds, and
the invalidation triggers are compiled into immutable test code before
the first injection run is scheduled. When the test executes, the
automated evaluation framework samples the instrumentation, computes
\(\rho(\phi, \mathbf{x}, 0)\), and outputs a deterministic verdict
without human intervention. If the machine breaches the precommitted
state boundary by \(0.1\text{ bar}\) or exceeds the response deadline by
\(2\text{ ms}\), the run is recorded as a failure, pointing directly to
an inadequate control margin or a slow diagnostic filter in the nervous
system.

Establishing precommitted oracles and quantitative temporal logic
boundaries transforms fault injection from a qualitative demonstration
into a quantitative measurement of system margins. Yet the validity of
that measurement remains entirely dependent on whether the injected
disturbances and the monitored reactions interact through real physical
dynamics or through mathematical approximations. Defining the exact
boundaries of a valid test and the criteria for success prepares the
engineering team to confront the physical machine itself, where
synthetic models give way to the unyielding mechanics of real hardware.

\section{Fault Injection on Hardware}\label{sec-verification-15-6}

A fault injected in simulation tests the simulator. When a disturbance
is synthesized inside a mathematical model, the software reacts against
the numerical assumptions that the modeler wrote down. Digital
simulations assume ideal clock synchronization, perfectly isolated
memory addresses, deterministic context switches, and linear
communication delays across shared interconnects. In the physical
machine, faults do not execute inside isolated functions. They propagate
across physical silicon, where memory access contention on a shared
system bus, clock drift between asynchronous microcontrollers, direct
memory access stalls, and voltage drops across inductive actuator
drivers dictate whether an intervention executes in time. To establish
evidence that survives scrutiny, fault injection must be performed
directly on target hardware through the production handoff, driving
disturbances through physical transceivers, target timers, shared buses,
and the hardware enforcement path, while an independent containment rig
protects the test facility from physical destruction.

\textbf{In-situ fault injection} is the deliberate, deterministic
introduction of physical, electrical, and computational disturbances
(including sensor register latches, clock drift, bus contention, supply
brownouts, and actuator load stalls) directly into target production
silicon and mechanics while operating under authentic electrical,
thermal, and dynamic loads.

Full-scale hardware-in-the-loop (HIL) test facilities
(Figure~\ref{fig-15-hil-avionics-testbed}) bridge this reality gap by
mounting production Electronic Control Units (ECUs) and actuator drivers
in their authentic physical flight geometry, connecting real
line-replaceable units through physical wiring harnesses, power
distribution channels, and breakout interfaces to real-time
fault-injection racks.

\begin{figure}[t]

\centering{

\includegraphics[width=0.85\linewidth,height=\textheight,keepaspectratio]{/Users/VJ/GitHub/MLSysBook-integrate-physical/physical/book/chapters/15-verification/figures/fig15_real_hil_avionics_testbed.jpg}

}

\caption{\label{fig-15-hil-avionics-testbed}\textbf{Full-Scale
Hardware-in-the-Loop (HIL) Avionics Test Ring and Fault Injection
Testbed.} Authentic aerospace and autonomous vehicle qualification
facility at NASA Marshall Space Flight Center (NASA ID: 1001823).
Production Line-Replaceable Units (LRUs)---including Actuator Control
Units (ACUs), Booster Control and Power Distribution Units (BCPDUs),
Recovery Control Units (RCUs), Ignition and Staging Controllers (ISCs),
and Data Acquisition and Recording Units (DARUs)---are mounted around a
full-scale circular structural ring matching the physical flight
geometry. Heavy-gauge DC power distribution lines, MIL-STD communication
fieldbuses, and shielded analog/digital signal harnesses interconnect
the embedded target processors. Dedicated breakout interfaces and
real-time simulator racks inject deterministic signal corruptions, clock
drift, power bus brownouts, and actuator stall transients into physical
silicon, measuring microsecond-level detection and takeover latencies
under authentic electromagnetic interference and thermal loads.
Attribution: NASA Marshall Space Flight Center, Public Domain.}

\end{figure}%

\protect\phantomsection\label{callout-decisionux2a-15.4}
\begin{fbxSimple}{callout-decision}{Design Decision:}{Testbed Instrumentation Architecture}
\phantomsection\label{callout-decision*-15.4}

Physical testbed verification requires dedicated bench instrumentation
engineered to intercept, corrupt, and observe signals without altering
circuit impedance or timing characteristics. A fully instrumented
physical AI verification testbed integrates five specialized hardware
subsystems:

\begin{enumerate}
\def\labelenumi{\arabic{enumi}.}
\tightlist
\item
  \textbf{High-Bandwidth Digital Storage Oscilloscopes (DSOs) and Logic
  Analyzers:} Synchronized across common \(10\text{ MHz}\) reference
  clocks to capture microsecond-level transitions across SPI, CAN, and
  gate-driver PWM lines.
\item
  \textbf{Inline Breakout Injection Boxes:} Equipped with solid-state
  Field-Effect Transistor (FET) switches capable of sub-microsecond
  pin-level open-circuit disconnections, short-to-ground, and
  short-to-power faults.
\item
  \textbf{Programmable Fieldbus Disturbance Nodes:} Dedicated
  microcontrollers capable of injecting bit-level CAN frame stuffing,
  cyclic redundancy check (CRC) corruptions, and unthrottled
  ``babbling-idiot'' message floods.
\item
  \textbf{Precision Programmable DC Power Supplies and Electronic
  Loads:} Capable of synthesizing fast supply rail droops
  (\(V_{\text{core}} < 0.85\text{ V}\)), ground bounce, and inrush
  current transients under dynamic motor loading.
\item
  \textbf{External Truth Transducers:} Laser Doppler vibrometers,
  optical shaft encoders, and high-speed photogrammetric cameras wired
  to isolated digitizers, recording ground-truth physical motion
  independent of the embedded controllers under test.
\end{enumerate}

\end{fbxSimple}

Consider a high-pressure fluid loop where an electric control valve
regulates a mass flow rate of \(12.0\text{ kg/s}\) across a differential
pressure of \(30.0\text{ bar}\). In this analytical example, the
downstream accumulator has a volume of \(0.04\text{ m}^3\) and an
effective fluid bulk modulus of \(1.6\times 10^9\text{ Pa}\). If the
valve stem shifts unexpectedly by \(2.0\text{ mm}\), the resulting mass
flow imbalance causes downstream pressure to rise at an analytical rate
of \(80.0\text{ bar/s}\), which consumes the entire \(0.8\text{ bar}\)
allowable transient margin in exactly \(10\text{ ms}\). In a digital
simulator, verifying that the nervous system catches this excursion
requires only checking an assertion statement. On production silicon,
verifying the response requires injecting the disturbance into the
physical sensor interface. When the test fixture injects a stale
observation by holding the analog-to-digital converter (ADC) output
constant while maintaining the original acquisition timestamp, the test
evaluates whether the freshness detector on the real-time safety
microcontroller (MCU) identifies the stall within its budgeted
\(P_{99}\) latency of \(4.0\text{ ms}\). If the detector relies on a
software timer that shares an interrupt line with a busy network
interface, the arrival of background telemetry can delay the interrupt
handler by \(8\text{ ms}\), allowing the stale state estimate to persist
until the pressure breaches the physical containment boundary.

Testing the heterogeneous Dual-Brain architecture---where an
unprivileged Linux host (MPU) generates learned neural proposals while a
deterministic real-time microcontroller (MCU) enforces safety
bounds---demands physical fault injection directly at their
inter-processor boundary:

\begin{itemize}
\tightlist
\item
  \textbf{Stale Observation Injection:} An inline FPGA interceptor
  latches an ADC register or suppresses CAN frame updates while
  preserving the original hardware timestamp. This evaluates whether the
  MCU's temporal freshness monitor detects the frozen observation within
  its \(\Delta t_{\text{det}} \le 4.0\text{ ms}\) budget before stale
  state estimates degrade actuator commands.
\item
  \textbf{Late Plan Injection:} The injection testbed delays neural
  trajectory proposals across the MPU-to-MCU shared memory bridge. If a
  proposal scheduled for execution at \(t = 100\text{ ms}\) arrives at
  \(t = 106\text{ ms}\), the test verifies that the MCU's deterministic
  scheduler instantly refuses the late submission, clamps the actuator
  to a precommitted safe hold position, and transitions the loop to
  deterministic fallback before physical vessel margins expire.
\item
  \textbf{Expired Intent Lease Revocation:} The test framework
  deliberately suppresses the cryptographic lease renewal packet
  transmitted by the Linux MPU. The verification oracle confirms that
  the MCU hardware gate revokes proposal authority on the exact clock
  cycle of lease expiry (\(t = 20.0\text{ ms}\) in
  Figure~\ref{fig-15-hardware-fault-injection}), immediately dropping
  un-leased commands even if their setpoints appear numerically
  plausible and smooth.
\end{itemize}

Recording the resulting state trajectory, the detector identity, the
exact refusal timestamp, and the time required to restore the bounded
state requires instrumentation that is electrically and logically
independent of the machine under test. If the test engineer relies on
the onboard logging daemon to measure reaction times, the measurement is
corrupted by the very failure it attempts to capture, because a faulted
CPU or a saturated bus delays log serialization.\sidenote{\footnotesize On safety
  microcontrollers (ARM Cortex-R5F, Infineon AURIX), dedicated fault
  injection registers and BIST controllers flip SRAM bits to verify that
  ECC SECDED triggers an immediate Non-Maskable Interrupt (NMI) before
  corrupted coordinates reach actuator PWM registers.} Furthermore,
independent measurement prevents the acceptance of false passes where
the machine survived for reasons absent from the safety claim. A
documented example occurred during ground qualification testing of the
Apollo guidance computer, where simulated testing failed to uncover an
asynchronous memory-bus stealing flaw caused by the rendezvous radar
interface; in flight, the hardware stole \(15\%\) of all memory cycles,
triggering the famous 1202 and 1201 alarm conditions.\sidenote{\footnotesize During
  Apollo 11's lunar descent, asynchronous radar pulses flooded the AGC
  with unbudgeted interrupts, stealing \(15\%\) of CPU cycles and
  triggering 1202 alarms. The lander survived because priority-scheduled
  core memory dropped low-priority display tasks to maintain thruster
  cycles.} An even simpler failure mode on modern test benches occurs
when an injected overpressure fault triggers an unmonitored thermal
breaker or an auxiliary ground loop that drops power to the valve drive
coil. The valve springs shut mechanically, keeping pressure inside the
containment bound, while the software enforcer has actually deadlocked
in a priority inversion. An automated oracle that only inspects
downstream pressure registers a pass, manufacturing false confidence
from an unverified hardware artifact. The test framework must verify
that the specific enforcer path described in the safety claim executed
the corrective action.

\protect\phantomsection\label{callout-incidentux2a-15.5}
\begin{fbxSimple}{callout-incident}{Landmark Engineering Incident:}{Toyota Unintended Acceleration Bit-Flip \& Watchdog Death}
\phantomsection\label{callout-incident*-15.5}

\emph{System:} Toyota Electronic Throttle Control System (ETCS-i)
operating on an integrated Denso/Hitachi 32-bit microcontroller without
a hardware Memory Protection Unit (MPU) or task-isolated memory spaces
(\citeproc{ref-barr2013toyota}{Barr 2013};
\citeproc{ref-nasa2011toyota}{National Aeronautics and Space
Administration Engineering and Safety Center 2011}).

\emph{Incident:} Widespread unintended acceleration events (notably
\emph{Bookout v. Toyota}, 2007 Camry) resulting in catastrophic
high-speed collisions. Vehicles accelerated to wide-open throttle (WOT)
uncommanded; pumping or standing on the mechanical brake pedal was
insufficient to overcome full engine torque, and the throttle actuator
refused to close.

\emph{Root cause \& physics of failure:} Cosmic-ray Single-Event Upsets
(SEUs) or unhandled stack memory overwrites corrupted the Task Control
Block (TCB) and operating system stack pointers for the critical
real-time throttle control task (\texttt{Task\_Throttle}). The
corruption caused \texttt{Task\_Throttle} to terminate or hang in an
infinite loop, freezing its last commanded throttle angle at \(100\%\)
open. Crucially, the hardware watchdog timer was designed with a
catastrophic architectural flaw: it was kicked not by verifying the
deterministic execution and deadline compliance of
\texttt{Task\_Throttle}, but unconditionally by an asynchronous
background OS monitor task (\texttt{Task\_Monitor}). Because
\texttt{Task\_Monitor} continued running normally, the watchdog timer
never timed out and never asserted a hardware system reset.\sidenote{\footnotesize Forensic
  analysis in \emph{Bookout v. Toyota} (2013) revealed unbounded stack
  growth and a fatal architectural flaw: the hardware watchdog timer was
  refreshed by an asynchronous OS monitor task that continued executing
  after the primary throttle control task had died.} Furthermore,
because the failsafe task shared global memory variables without memory
protection, it too suffered memory corruption and failed to engage the
throttle fail-safe return spring.

\emph{Systems resolution:} Enforce strict spatial memory isolation via
hardware MPUs/MMUs; decouple safety reflexes onto an independent
real-time microcontroller (the Dual-Brain architecture); implement
windowed hardware watchdogs that require deadman tokens from
\emph{every} active hard real-time task; and hardwire independent
brake-override logic directly into power stage gate drivers.

\emph{Systems takeaway:} A software watchdog that monitors overall CPU
execution rather than individual task invariants is a fatal
architectural illusion: a dead control task can leave actuators locked
at maximum destructive power while the watchdog continues reporting
nominal health.

\end{fbxSimple}

\protect\phantomsection\label{callout-ruleux2a-15.6}
\begin{fbxSimple}{callout-rule}{Architect’s Rule of Thumb:}{The Watchdog Liveness Invariant}
\phantomsection\label{callout-rule*-15.6}

A software watchdog timer that monitors only overall CPU execution
verifies merely that silicon has electrical power; safety requires
independent windowed watchdogs that verify the deterministic execution
cadence and invariant assertions of every individual critical task.

\end{fbxSimple}

The intervention mechanism must be subjected to fault injection under
the exact operational conditions that make intervention necessary. A
safety enforcer that takes over cleanly when the system sits idle on a
test bench provides no evidence of safety in production. During an
operational crisis, the learned policy is typically running at maximum
computational throughput, memory buses are saturated with
high-resolution sensory streams, direct memory access channels are
transferring frame buffers, and the actuator is experiencing maximum
hydrodynamic load. The test framework must inject faults while the
processor is under peak thermal and electrical stress, verifying that
the hardware multiplexer can preempt the proposal path, flush the
command buffers, and assert control over the valve driver within the
\(P_{99}\) deadline of \(2.0\text{ ms}\)
(Figure~\ref{fig-15-hardware-fault-injection}). If the intervention path
shares an internal bus arbiter with the neural accelerator, the hardware
takeover can stall behind an un-preemptible burst transfer. A machine
whose safety case assumes instantaneous takeover will fail in the field
if that takeover was only ever verified on a quiet bus.

The complete deterministic protocol for conducting automated
hardware-in-the-loop fault injection, verifying real-time detection
invariants, and evaluating quantitative Signal Temporal Logic robustness
on target silicon is formalized below:

\subsection{Automated Hardware-in-the-Loop Fault Injection
Protocol}\label{automated-hardware-in-the-loop-fault-injection-protocol}

\textbf{Substrate \& Invariants:}

\begin{itemize}
\tightlist
\item
  \emph{Execution Substrate:} Hard real-time FPGA plant rig and testbed
  orchestrator (\(f = 1000\text{ Hz}\), hardware timestamping jitter
  \(< 500\text{ ns}\)).
\item
  \emph{Memory Invariant:} Pre-allocated static DMA buffers in host test
  orchestrator; zero runtime dynamic heap allocation on safety MCU.
\item
  \emph{Safety Interlock Authority:} Independent optical/laser truth
  sensors wired to hardware crowbar relays capable of cutting target
  actuator power within \(2.0\text{ ms}\) upon physical envelope breach.
\end{itemize}

\emph{Execution protocol:}

\begin{enumerate}
\def\labelenumi{\arabic{enumi}.}
\item
  \textbf{Pre-Execution Invariant \& Envelope Precondition Audit (Stage
  1):} Sample baseline environmental sensors via
  \(\mathcal{I}_{\text{truth}}\) (ambient temperature
  \(T_{\text{amb}}\), supply rail voltage \(V_{\text{bus}}\), baseline
  fluid pressure \(P_0\)). Verify environmental vector matches declared
  operating point:
  \(\|\mathbf{w}_{\text{meas}} - \mathbf{w}_{\text{env}}\| \le \boldsymbol{\epsilon}_{\text{env}}\).
  If preconditions are violated, log
  \texttt{ERR\_PRECONDITION\_OUT\_OF\_SPEC}, set
  \(\text{Verdict} \leftarrow \text{INVALID\_PRECONDITION}\), and abort.
\item
  \textbf{Baseline Closed-Loop Synchronization \& Steady-State Arming
  (Stage 2):} Arm FPGA plant simulator and power up target Dual-Brain
  SoC; initiate periodic control loop (\(1000\text{ Hz}\)). Establish
  synchronized reference clock across logic analyzer, FPGA plant, and
  truth digitizers via \(10\text{ MHz}\) PXI clock. Allow physical plant
  state \(\mathbf{x}(t)\) to stabilize in nominal operating envelope for
  \(t_{\text{settle}} = 2.000\text{ s}\). If baseline tracking error
  \(\|e_{\text{track}}\| > \epsilon_{\text{nom}}\), assert
  \texttt{ERR\_BENCH\_INSTABILITY}, set
  \(\text{Verdict} \leftarrow \text{INVALID\_PRECONDITION}\), and abort.
\item
  \textbf{Deterministic In-Situ Fault Injection Trigger (Stage 3):} At
  monotonic test epoch \(t_{\text{inject}}\), hardware test orchestrator
  asserts trigger line to inline injection rig:

  \begin{itemize}
  \tightlist
  \item
    \emph{Sensor Latency/Bias:} FPGA interceptor latches ADC register or
    holds stale CAN payload while preserving valid hardware timestamp.
  \item
    \emph{Bus Contention / Frame Stuffing:} Fieldbus disturbance node
    floods CAN bus with dominant identifier bits (\(98\%\) bus load).
  \item
    \emph{Power Rail Droop:} Programmable electronic load pulls
    \(V_{\text{core}}\) down from \(1.10\text{ V}\) to \(0.82\text{ V}\)
    for \(\Delta t_{\text{duration}} = 5.0\text{ ms}\).
  \item
    \emph{Actuator Mechanical Stall:} 4-quadrant dynamometer applies
    step counter-torque
    \(\tau_{\text{stall}} = 3.5 \times \tau_{\text{rated}}\). Record
    microsecond injection timestamp
    \(t_{\text{fault\_start}} \leftarrow t_{\text{now}}\).
  \end{itemize}
\item
  \textbf{Real-Time Safety Enforcer Response \& Takeover Audit (Stage
  4):} Monitor target MCU safety enforcer pins, CAN bus telemetry, and
  gate-driver PWM lines via isolated logic analyzer:

  \begin{itemize}
  \tightlist
  \item
    Detect assertion of Non-Maskable Interrupt (NMI) / diagnostic flag:
    \(t_{\text{det}} \leftarrow t_{\text{now}} \implies \Delta t_{\text{det}} = t_{\text{det}} - t_{\text{fault\_start}}\).
  \item
    Detect hardware multiplexer preemption of motor driver:
    \(t_{\text{takeover}} \leftarrow t_{\text{now}} \implies \Delta t_{\text{takeover}} = t_{\text{takeover}} - t_{\text{fault\_start}}\).
  \item
    Detect physical plant transition to safe bounded regime:
    \(t_{\text{bounded}} \leftarrow t_{\text{now}} \implies \Delta t_{\text{bounded}} = t_{\text{bounded}} - t_{\text{fault\_start}}\).
    If external containment rig trips due to structural envelope breach:
    hardware crowbar cuts actuator power; assert
    \texttt{ERR\_PHYSICAL\_BREACH}.
  \end{itemize}
\item
  \textbf{Quantitative STL Robustness Evaluation \& Ledger Serialization
  (Stage 5):} Ingest continuous truth trajectory
  \(\mathbf{x}_{\text{obs}}(t)\) from independent instrumentation
  \(\mathcal{I}_{\text{truth}}\) over \(t \in [0, T_{\text{trial}}]\).
  Check telemetry integrity; if truth frames were dropped or DMA buffers
  overflowed, set
  \(\text{Verdict} \leftarrow \text{UNOBSERVABLE\_TELEMETRY}\),
  serialize partial trace, and return. Compute quantitative Signal
  Temporal Logic robustness:
  \[\rho^* = \rho(\phi_{\text{safe}}, \mathbf{x}_{\text{obs}}, 0).\] If
  \(\rho^* > 0\), \(\Delta t_{\text{det}} \le \tau_{\text{det\_max}}\),
  and \(\Delta t_{\text{takeover}} \le \tau_{\text{takeover\_max}}\),
  set \(\text{Verdict} \leftarrow \text{PASS}\); else set
  \(\text{Verdict} \leftarrow \text{FAIL}\). Assemble immutable
  \texttt{FaultExecutionRecord\_t} record, sign with Ed25519 hardware
  key, append to machine notebook, and return.
\end{enumerate}

\begin{figure}[t]

\centering{

\includegraphics[width=0.95\linewidth,height=\textheight,keepaspectratio]{index_files/mediabag//Users/VJ/GitHub/MLSysBook-integrate-physical/physical/book/chapters/15-verification/figures/fig15_hardware_fault_injection.pdf}

}

\caption{\label{fig-15-hardware-fault-injection}\textbf{Real-Time
Hardware Fault Injection and Deterministic Intervention Timeline.}
Multi-channel logic analyzer and pressure sensor trace capturing the
microsecond fault containment cascade on target hardware. At
\(t_0 = 0.0\text{ ms}\), the test fixture injects a sensor stall by
latching the analog-to-digital converter register at a static reading
(\(4.2\text{ kg/s}\)). At \(t_1 = 20.0\text{ ms}\), the cryptographic
intent lease expires without renewal, revoking host MPU proposal
authority. At \(t_2 = 24.0\text{ ms}\), the MCU safety enforcer's
diagnostic monitor detects the stale observation
(\(\Delta t_{\text{det}} = 4.0\text{ ms}\)) and asserts a Non-Maskable
Interrupt (NMI). At \(t_3 = 26.0\text{ ms}\), the hardware multiplexer
preempts the gate drivers
(\(\Delta t_{\text{takeover}} = 2.0\text{ ms}\)), forcing a
\(-12.0\text{ A}\) regenerative dynamic braking profile. The downstream
manifold pressure \(P(t)\) surges at \(80.0\text{ bar/s}\) but is safely
arrested at \(P_{\text{peak}} = 16.2\text{ bar}\), remaining strictly
within the \(16.5\text{ bar}\) containment boundary and far below the
\(22.0\text{ bar}\) pipe burst limit, reaching a bounded safe state at
\(t_4 = 42.8\text{ ms}\), well inside the \(45.0\text{ ms}\) physical
vessel deadline. Attribution: Original textbook graphic, CC BY 4.0.}

\end{figure}%

\section{The Fault Manifest}\label{sec-verification-15-7}

The machine notebook tracks components, software artifacts, and
execution traces under the provenance and versioning rules established
in Section~\ref{sec-nervous-4-10}. To document verification evidence,
the notebook requires an explicit data structure that binds every
injected fault directly to the claim it challenges. The \textbf{fault
record} records nine mandatory fields for each experimental trial,
comprising the safety claim under attack, the precise operating envelope
point, the injected fault magnitude and duration, the target hardware
and software build versions, the raw sensor and bus telemetry trace, the
primary detector that caught the anomaly, the measured intervention
response time, the final physical outcome, and any secondary hardware
anomalies. In an active flow regulator, for instance, a complete entry
records that safety claim C-04 (containment under upstream supply surge)
was attacked at an operating envelope point of \(350\text{ kPa}\) inlet
pressure and \(45^\circ\text{C}\) fluid temperature. The fault was an
analytical step command forcing a \(24\text{ V}\) override on the valve
driver for \(120\text{ ms}\), applied to controller firmware revision
3.4.1 running on board revision B. The primary detector was an
independent pressure-gradient monitor, the response took
\(1.8\text{ ms}\), the physical outcome was containment maintained below
\(380\text{ kPa}\), and no thermal breaker trips or bus retries
occurred.

A verification log that only records passes and failures creates an
incentive to discard corrupted experiments, which systematically biases
the resulting evidence. If an injected communication fault on the
controller bus corrupts the logging daemon or causes an external
digitizer to drop frames, the experiment is unobservable. If an
auxiliary ground-fault breaker cuts bench power before the safety
enforcer can evaluate the flow boundary, the experiment is invalid.
Treating an unobservable or invalid run as a failed test misidentifies
bench artifacts as machine defects, while treating it as an absent trial
silently discards the conditions under which telemetry failed. The
notebook therefore segregates invalid trials, unobservable trials, and
valid executions into distinct categories. Every trial that fails to
yield complete telemetry is logged as unobservable alongside the raw
partial buffer, ensuring that missing telemetry cannot artificially
inflate apparent test coverage. The binary memory layout and invariant
contracts governing this data structure are formalized in
Table~\ref{tbl-15-verification-manifest-schema}.

\begin{longtable}[]{@{}
  >{\raggedright\arraybackslash}p{(\linewidth - 10\tabcolsep) * \real{0.1667}}
  >{\raggedright\arraybackslash}p{(\linewidth - 10\tabcolsep) * \real{0.1667}}
  >{\raggedright\arraybackslash}p{(\linewidth - 10\tabcolsep) * \real{0.1667}}
  >{\raggedright\arraybackslash}p{(\linewidth - 10\tabcolsep) * \real{0.1667}}
  >{\raggedright\arraybackslash}p{(\linewidth - 10\tabcolsep) * \real{0.1667}}
  >{\raggedright\arraybackslash}p{(\linewidth - 10\tabcolsep) * \real{0.1667}}@{}}
\caption{\textbf{Abstract Data Schema for the Verification Test Manifest
and Fault Execution Record (\protect\texttt{FaultExecutionRecord\_t}).}
Specification contract layout (\(147\text{ bytes}\) total), binary data
types, physical SI units, verification invariants, hardware audit
substrates, and deterministic invalidation triggers governing automated
HIL fault injection
logging.}\label{tbl-15-verification-manifest-schema}\tabularnewline
\toprule\noalign{}
\begin{minipage}[b]{\linewidth}\raggedright
Schema Field Identifier
\end{minipage} & \begin{minipage}[b]{\linewidth}\raggedright
Binary Type \& Memory Layout
\end{minipage} & \begin{minipage}[b]{\linewidth}\raggedright
Physical SI Units
\end{minipage} & \begin{minipage}[b]{\linewidth}\raggedright
Structural Invariant \& Verification Scope
\end{minipage} & \begin{minipage}[b]{\linewidth}\raggedright
Hardware Audit Substrate \& Injection Site
\end{minipage} & \begin{minipage}[b]{\linewidth}\raggedright
Deterministic Invalidation \& Trip Trigger
\end{minipage} \\
\midrule\noalign{}
\endfirsthead
\toprule\noalign{}
\begin{minipage}[b]{\linewidth}\raggedright
Schema Field Identifier
\end{minipage} & \begin{minipage}[b]{\linewidth}\raggedright
Binary Type \& Memory Layout
\end{minipage} & \begin{minipage}[b]{\linewidth}\raggedright
Physical SI Units
\end{minipage} & \begin{minipage}[b]{\linewidth}\raggedright
Structural Invariant \& Verification Scope
\end{minipage} & \begin{minipage}[b]{\linewidth}\raggedright
Hardware Audit Substrate \& Injection Site
\end{minipage} & \begin{minipage}[b]{\linewidth}\raggedright
Deterministic Invalidation \& Trip Trigger
\end{minipage} \\
\midrule\noalign{}
\endhead
\bottomrule\noalign{}
\endlastfoot
\texttt{manifest\_id} & \texttt{uint64\_t} (8 bytes) & dimensionless &
Unique monotonic 64-bit trial identifier & Test orchestrator static RAM
& Zero value or duplicate identifier in ledger \\
\texttt{claim\_ref\_id} & \texttt{uint32\_t} (4 bytes) & dimensionless &
Traceable safety claim index (e.g., C-04) & Engineering notebook record
pointer & Unmapped claim index in requirements graph \\
\texttt{envelope\_point} &
\texttt{struct\ \{\ float32\_t\ p,\ T,\ mu,\ v;\ \}} (16 bytes) &
\(\text{Pa}, \text{K}, \text{Pa}\cdot\text{s}, \text{m/s}\) & Declared
physical operating point & Independent environmental sensors &
\(\|\mathbf{w}_{\text{meas}} - \mathbf{w}_{\text{env}}\| > \boldsymbol{\epsilon}_{\text{env}}\)
before injection \\
\texttt{fault\_injection\_mode} & \texttt{enum\ uint16\_t} (2 bytes) &
dimensionless & Fault class (SENSOR\_BIAS, BUS\_FLOOD, DROOP, STALL) &
Inline breakout switch / CAN disturbance node & Unrecognized fault mode
code \\
\texttt{fault\_magnitude} & \texttt{float32\_t} (4 bytes) &
\(\text{V}, \text{bar}, \text{ms}, \text{N}\cdot\text{m}\) & Calibrated
disturbance amplitude & Programmable DC load / dynamometer & Injection
actuator output saturation \\
\texttt{injection\_target\_node} & \texttt{enum\ uint16\_t} (2 bytes) &
dimensionless & Target hardware boundary (MPU\_SHM, MCU\_SPI, PWM) & Pin
breakout connector / crossbar tap & High-impedance disconnected harness
pin \\
\texttt{stl\_spec\_hash} & \texttt{uint8\_t{[}32{]}} (32 bytes) &
dimensionless & SHA-256 digest of immutable STL formula \(\phi\) &
Precommitted test oracle binary & Hash mismatch with test manifest
registry \\
\texttt{measured\_t\_det} & \texttt{uint32\_t} (4 bytes) &
\(\mu\text{s}\) (\(10^{-6}\text{ s}\)) & Latency from fault trigger to
enforcer NMI & Isolated 10 MHz logic analyzer channel &
\(\Delta t_{\text{det}} > \tau_{\text{det\_max}}\) (triggers fail
verdict) \\
\texttt{measured\_t\_takeover} & \texttt{uint32\_t} (4 bytes) &
\(\mu\text{s}\) (\(10^{-6}\text{ s}\)) & Latency from NMI to gate-driver
override & Current probe / gate voltage digitizer &
\(\Delta t_{\text{takeover}} > \tau_{\text{takeover\_max}}\) (triggers
fail) \\
\texttt{stl\_robustness\_rho} & \texttt{float32\_t} (4 bytes) & physical
margin (\(\text{bar}, \text{m/s}, \text{N}\)) & Continuous metric
\(\rho(\phi_{\text{safe}}, \mathbf{x}, 0)\) & Independent optical/laser
truth sensors & \(\rho^* < 0\) (falsifies safety claim) \\
\texttt{verdict\_status} & \texttt{enum\ uint8\_t} (1 byte) &
dimensionless & Trial outcome (PASS, FAIL, INVALID, UNOBSERVABLE) &
Automated post-trial evaluator & Missing truth packets \(\to\)
UNOBSERVABLE \\
\texttt{provenance\_sig} & \texttt{uint8\_t{[}64{]}} (64 bytes) &
dimensionless & Ed25519 signature of complete trial binary &
Cryptographic hardware security module (HSM) & Signature verification
failure on ingestion \\
\end{longtable}

Reporting verification progress as a single aggregate percentage
obscures critical engineering vulnerabilities. A safety case cannot rely
on a summary metric stating that a machine survived ninety-eight percent
of five thousand fault trials if all five thousand injections were
conducted at room temperature on a quiescent memory bus. Coverage must
instead be reported across separate operational dimensions, including
functional safety requirements, physical envelope boundaries, physical
and logical injection sites, timing tails, hardware build revisions, and
multi-fault combinations. For a flow regulator, boundary coverage tracks
whether faults were injected at maximum fluid viscosity, minimum
back-pressure, and maximum pump ripple. Timing tail coverage isolates
injections targeting the \(P_{99}\) scheduling latency of the nervous
system and peak direct memory access utilization. A machine achieves
acceptable verification not when an average coverage index exceeds an
arbitrary threshold, but when the minimum coverage across every
individual dimension satisfies the requirements of the safety
case.\sidenote{\footnotesize ANSI/UL 4600 mandates goal-based safety cases linking
  claims to traceable evidence rather than raw operational hours.
  Similarly, ISO 26262-5 mandates hardware fault injection to verify
  ASIL D diagnostic coverage (Single-Point Fault Metric \(\ge 99\%\),
  Latent Fault Metric \(\ge 90\%\)).}

Interpreting repeated trials requires distinguishing between
deterministic execution and stochastic sampling. Repeating the exact
same deterministic fault script five hundred times on a digital test
bench yields no more information than running it once. If initial
conditions, processor register states, clock dividers, and input
sequences are fixed, identical deterministic software executes the
identical branch sequence every time. Repetition adds evidentiary value
only when sampling physical and computational distributions that exhibit
stochastic variance. In a physical flow loop, repeated trials sample
turbulent velocity fluctuations, asynchronous interrupt arrivals, bus
arbitration contention, and phase jitter between the high-frequency
sensor clock and the policy inference cycle. In these stochastic
regimes, repetition measures the spread and tail of the response-time
distribution rather than re-confirming a deterministic path.

When a stochastic trial runs \(n\) consecutive times without a single
safety violation, the count of zero-failure trials bounds the maximum
underlying failure probability \(p\) for that specific test
distribution.\sidenote{\footnotesize For zero-failure testing, standard Wald
  confidence intervals degenerate to zero width. The exact
  Clopper-Pearson interval yields
  \(p_{\text{upper}} = 1 - \alpha^{1/n}\), preventing asymptotic
  Gaussian approximations from understating catastrophic risk in
  rare-event tail distributions.} If each trial has an independent
failure probability \(p\), the probability of observing zero failures
across \(n\) independent trials is \((1 - p)^n\). To establish that the
failure probability is bounded above by \(p\) at a statistical
confidence level of \(1 - \alpha\), the probability of observing an
all-pass outcome by chance from an unsafe system must not exceed
\(\alpha\), yielding \((1 - p)^n \le \alpha\). Taking the natural
logarithm of both sides gives \(n \ln(1 - p) \le \ln(\alpha)\). For rare
failure rates where \(p \ll 1\), the first-order Taylor expansion
\(\ln(1 - p) \approx -p\) provides the standard engineering
approximation \(n \ge -\ln(\alpha) / p\). For an exact calculation to
bound the failure probability below one percent (\(p \le 0.01\)) at
ninety-five percent confidence (\(\alpha = 0.05\)), the inequality is
\((1 - 0.01)^n \le 0.05\). Taking logarithms produces
\(n \ge \ln(0.05) / \ln(0.99) \approx -2.9957 / -0.01005 \approx 298.07\),
which requires exactly \(299\) independent, zero-failure trials.

The statistical validity of this one-sided binomial bound depends
entirely on two physical assumptions, namely trial independence and
operational representativeness. Independence requires that the outcome
of trial \(k\) has no effect on trial \(k+1\). If repeated fault
injections cause resistive heating in a valve solenoid coil, increasing
coil resistance and slowing valve actuation across successive tests, the
trials are coupled by a hidden thermal state, and the binomial formula
is invalid. Representativeness requires that the test conditions
accurately reflect the true operational environment. Running 299
successful fault injections under steady, laminar flow at
\(20^\circ\text{C}\) proves that the failure probability is below one
percent at ninety-five percent confidence only for laminar flow at
\(20^\circ\text{C}\). It provides zero empirical or mathematical basis
for claiming safety during turbulent flow, severe cavitation, or
cold-start operations at \(-10^\circ\text{C}\).

\protect\phantomsection\label{callout-mathux2a-15.7}
\begin{fbxSimple}{callout-math}{Fermi Math:}{Hardware Fault Injection Coverage and Mean Time Between Unmitigated Faults (MTBUF)}
\phantomsection\label{callout-math*-15.7}

\emph{Scenario:} An autonomous physical AI system operates with a
cumulative raw hardware component failure rate \(\lambda_{\text{raw}}\)
across its perception sensors, communication transceivers,
microcontrollers, and inverter gate drivers. We calculate the diagnostic
coverage \(C_{\text{diag}}\) that Hardware-in-the-Loop (HIL) fault
injection must empirically prove to achieve a target Mean Time Between
Unmitigated Faults (\(\text{MTBUF}\)) of \(10^7\text{ hours}\)
(equivalent to an unmitigated failure rate
\(\lambda_{\text{unmitigated}} \le 10^{-7}\text{ h}^{-1} = 100\text{ FIT}\)),
and derive the required number of zero-failure HIL test trials
\(N_{\text{trials}}\).

\textbf{System Failure Rate \& Unmitigated Failure Relation:}
\[\lambda_{\text{unmitigated}} = \lambda_{\text{raw}} \cdot (1 - C_{\text{diag}}) \implies \text{MTBUF} = \frac{1}{\lambda_{\text{unmitigated}}} = \frac{1}{\lambda_{\text{raw}} (1 - C_{\text{diag}})}\]
Solving for required diagnostic coverage:
\[C_{\text{diag}} \ge 1 - \frac{\lambda_{\text{unmitigated}}}{\lambda_{\text{raw}}} = 1 - \frac{1}{\text{MTBUF} \cdot \lambda_{\text{raw}}}\]

\textbf{1. Cumulative Raw Hardware Failure Rate:} Suppose an embedded
computing and actuation stack incorporates:

\begin{itemize}
\tightlist
\item
  4 CMOS cameras + 2 LiDAR modules:
  \(\lambda_{\text{sensors}} = 6 \times 250\text{ FIT} = 1,500\text{ FIT}\)
\item
  2 Dual-Brain SoCs (Linux MPUs + Safety MCUs):
  \(\lambda_{\text{compute}} = 2 \times 500\text{ FIT} = 1,000\text{ FIT}\)
\item
  CAN/Ethernet Transceivers + Power Distribution PMICs:
  \(\lambda_{\text{bus/power}} = 2,500\text{ FIT}\)
\item
  6 Inverter Power Stages and Gate Drivers:
  \(\lambda_{\text{inverters}} = 6 \times 800\text{ FIT} = 4,800\text{ FIT}\)
\item
  Total raw hardware failure rate: \[\begin{aligned}
  \lambda_{\text{raw}} &= 1,500 + 1,000 + 2,500 + 4,800 = 9,800\text{ FIT} \\
  &= 9.8 \times 10^{-6}\text{ faults/h} \quad (\approx 1\text{ raw fault per } 102,000\text{ hours})
  \end{aligned}\]
\end{itemize}

\textbf{2. Required Diagnostic Coverage Calculation:} To guarantee an
unmitigated catastrophic failure rate
\(\lambda_{\text{unmitigated}} \le 10^{-7}\text{ h}^{-1}\)
(\(\text{MTBUF} \ge 10^7\text{ hours}\)):
\[C_{\text{diag}} \ge 1 - \frac{10^{-7}\text{ h}^{-1}}{9.8 \times 10^{-6}\text{ h}^{-1}} = 1 - 0.0102 = 0.9898 \implies C_{\text{diag}} \ge 99.0\%\]

\textbf{3. Required Statistical Zero-Failure HIL Test Count:} Using the
exact Clopper-Pearson binomial tail bound at \(95\%\) confidence
(\(\alpha = 0.05\)): \[\begin{aligned}
(1 - p_{\text{unmitigated}})^N &\le \alpha \implies N_{\text{trials}} \ge \frac{\ln(\alpha)}{\ln(C_{\text{diag}})} \\
&= \frac{\ln(0.05)}{\ln(0.990)} \approx \frac{-2.9957}{-0.01005} \approx 298.1 \implies 299\text{ trials}
\end{aligned}\] If safety integrity escalates to ASIL D
(\(C_{\text{diag}} \ge 99.9\%\), \(p_{\text{unmitigated}} \le 0.001\)):
\[N_{\text{trials}} \ge \frac{-\ln(0.05)}{1 - 0.999} = \frac{2.9957}{0.001} \approx 2,996\text{ consecutive zero-failure HIL injection trials}\]

\emph{Archetype coverage comparison:}

\begin{itemize}
\tightlist
\item
  \emph{Mobile System (Autonomous Ground Transport):} Total raw rate
  \(\lambda_{\text{raw}} = 1.2 \times 10^{-5}\text{ h}^{-1}\); targeting
  ASIL D unmitigated rate (\(10^{-8}\text{ h}^{-1}\)) requires
  \(C_{\text{diag}} \ge 99.92\%\), demanding \(3,745\) zero-failure HIL
  brake/steer brownout injections.
\item
  \emph{Manipulator (6-DoF Articulated Arm):} Raw joint inverter rate
  \(\lambda_{\text{raw}} = 6.0 \times 10^{-6}\text{ h}^{-1}\); targeting
  SIL 3 (\(10^{-7}\text{ h}^{-1}\)) requires
  \(C_{\text{diag}} \ge 98.33\%\), requiring \(179\) zero-failure
  joint-stall brake clamp trials.
\item
  \emph{Process System (High-Pressure Flow Regulator):} Raw
  transducer/valve rate
  \(\lambda_{\text{raw}} = 8.5 \times 10^{-6}\text{ h}^{-1}\); targeting
  \(10^{-7}\text{ h}^{-1}\) requires \(C_{\text{diag}} \ge 98.82\%\),
  demanding \(253\) zero-failure cavitation/stale-pressure injection
  trials.
\end{itemize}

\end{fbxSimple}

When target reliability requirements scale from advisory levels to
life-critical integrity levels, the sample size required by binomial and
exponential tail bounds becomes physically impossible to achieve through
black-box operational exposure alone. This fundamental
limitation---formalized by Butler and Finelli
(\citeproc{ref-butler1993infeasibility}{Butler and Finelli 1993}) as the
\textbf{Butler \& Finelli Infeasibility Bound}---dictates that
certifying ultra-high cyber-physical reliability (e.g., \(10^{-9}\)
catastrophic failures per hour) through unperturbed operational exposure
alone is physically impossible due to astronomical time requirements
(e.g., \(>3.0 \times 10^9\text{ hours}\) of continuous zero-failure
testing), mathematically necessitating deterministic fault injection,
layered runtime invariant enforcement, and formal architectural
partitioning.

As quantified in Table~\ref{tbl-15-statistical-testing-bounds},
certifying an avionics DAL A or ISO 26262 ASIL D catastrophic failure
rate of \(10^{-9}\text{ h}^{-1}\) at \(95\%\) confidence requires over
\(3.0 \times 10^9\text{ hours}\) (more than 342,000 years) of
continuous, zero-failure physical operation. Even a distributed fleet of
100 dedicated test vehicles operating continuously 24 hours a day would
require over 3,400 years to establish this statistical claim.
Consequently, physical AI systems cannot be validated solely by
accumulating millions of unperturbed operational fleet miles; safety
assurance must instead be established through deterministic fault
injection, layered runtime invariant enforcement, and formal
architectural partitioning.

\begin{longtable}[]{@{}
  >{\raggedright\arraybackslash}p{(\linewidth - 12\tabcolsep) * \real{0.1429}}
  >{\raggedright\arraybackslash}p{(\linewidth - 12\tabcolsep) * \real{0.1429}}
  >{\raggedright\arraybackslash}p{(\linewidth - 12\tabcolsep) * \real{0.1429}}
  >{\raggedright\arraybackslash}p{(\linewidth - 12\tabcolsep) * \real{0.1429}}
  >{\raggedright\arraybackslash}p{(\linewidth - 12\tabcolsep) * \real{0.1429}}
  >{\raggedright\arraybackslash}p{(\linewidth - 12\tabcolsep) * \real{0.1429}}
  >{\raggedright\arraybackslash}p{(\linewidth - 12\tabcolsep) * \real{0.1429}}@{}}
\caption{\textbf{Statistical Confidence vs.~Physical Testing Exposure
Bounds across Safety Integrity Levels.} Quantification of required
zero-failure test exposure (\(T = -\ln(1-C)/\lambda\) at \(C=0.95\)
confidence) and physical operational fleet hours per Butler \& Finelli
(\citeproc{ref-butler1993infeasibility}{Butler and Finelli 1993}),
demonstrating the mathematical infeasibility of certifying life-critical
cyber-physical systems through black-box operational burn-in
alone.}\label{tbl-15-statistical-testing-bounds}\tabularnewline
\toprule\noalign{}
\begin{minipage}[b]{\linewidth}\raggedright
Target Safety Integrity Level \& Standard
\end{minipage} & \begin{minipage}[b]{\linewidth}\raggedright
Target Failure Rate (\(\lambda\), \(\text{h}^{-1}\))
\end{minipage} & \begin{minipage}[b]{\linewidth}\raggedright
Permissible Per-Cycle Failure Prob. (\(p\) at \(100\text{ Hz}\))
\end{minipage} & \begin{minipage}[b]{\linewidth}\raggedright
Zero-Failure Test Hours (\(T\), \(C=0.95\))
\end{minipage} & \begin{minipage}[b]{\linewidth}\raggedright
Single-Unit Physical Exposure Time
\end{minipage} & \begin{minipage}[b]{\linewidth}\raggedright
100-Unit Fleet Exposure Duration
\end{minipage} & \begin{minipage}[b]{\linewidth}\raggedright
Required Verification Strategy (Butler \& Finelli Implication)
\end{minipage} \\
\midrule\noalign{}
\endfirsthead
\toprule\noalign{}
\begin{minipage}[b]{\linewidth}\raggedright
Target Safety Integrity Level \& Standard
\end{minipage} & \begin{minipage}[b]{\linewidth}\raggedright
Target Failure Rate (\(\lambda\), \(\text{h}^{-1}\))
\end{minipage} & \begin{minipage}[b]{\linewidth}\raggedright
Permissible Per-Cycle Failure Prob. (\(p\) at \(100\text{ Hz}\))
\end{minipage} & \begin{minipage}[b]{\linewidth}\raggedright
Zero-Failure Test Hours (\(T\), \(C=0.95\))
\end{minipage} & \begin{minipage}[b]{\linewidth}\raggedright
Single-Unit Physical Exposure Time
\end{minipage} & \begin{minipage}[b]{\linewidth}\raggedright
100-Unit Fleet Exposure Duration
\end{minipage} & \begin{minipage}[b]{\linewidth}\raggedright
Required Verification Strategy (Butler \& Finelli Implication)
\end{minipage} \\
\midrule\noalign{}
\endhead
\bottomrule\noalign{}
\endlastfoot
\textbf{Consumer Robotics / Advisory AI}\emph{(Unregulated / Soft Safe
State)} & \(10^{-3}\text{ h}^{-1}\)\emph{(1 failure /
\(10^3\text{ h}\))} & \(p \approx 2.78 \times 10^{-9}\) &
\(3.0 \times 10^3\text{ h}\) & \(3,000\text{ h}\)\emph{(125 days
continuous)} & \(30\text{ h}\)\emph{(\textless{} 1 week)} &
\textbf{Black-Box Testing Feasible:} Direct operational sampling and
nominal trace replay sufficient to validate empirical claims. \\
\textbf{Industrial Automation (SIL 2 / PL d)}\emph{(IEC 61508 / ISO
13849)} & \(10^{-6}\text{ h}^{-1}\)\emph{(1 failure /
\(10^6\text{ h}\))} & \(p \approx 2.78 \times 10^{-12}\) &
\(3.0 \times 10^6\text{ h}\) & \(3.0 \times 10^6\text{ h}\)\emph{(342
years continuous)} & \(3.4\text{ years}\)\emph{(at 24/7 continuous
operation)} & \textbf{Hybrid Verification:} Statistical sampling paired
with accelerated fault injection and HIL parameter stress sweeps. \\
\textbf{Automotive Safety (ASIL D / SIL 3)}\emph{(ISO 26262 / IEC
61508)} & \(10^{-8}\text{ h}^{-1}\)\emph{(10 FIT / \(10^8\text{ h}\))} &
\(p \approx 2.78 \times 10^{-14}\) & \(3.0 \times 10^8\text{ h}\) &
\(3.0 \times 10^8\text{ h}\)\emph{(34,200 years continuous)} &
\(342\text{ years}\)\emph{(100 vehicles 24/7)} & \textbf{Architectural
Containment Mandatory:} Formal invariant synthesis, hardware interlocks,
and systematic fault injection; fleet testing cannot prove safety. \\
\textbf{Life-Critical Avionics (DAL A / SIL 4)}\emph{(DO-178C / FAA AC
20-115D)} & \(10^{-9}\text{ h}^{-1}\)\emph{(1 failure /
\(10^9\text{ h}\))} & \(p \approx 2.78 \times 10^{-15}\) &
\(3.0 \times 10^9\text{ h}\) &
\(3.0 \times 10^9\text{ h}\)\emph{(342,465 years continuous)} &
\(3,425\text{ years}\)\emph{(100 units 24/7 continuous)} &
\textbf{Mathematical Impossibility of Black-Box Test:} Safety
\emph{cannot} be established by fleet exposure. Provenance, formal
methods, and isolated intervention are mathematically required. \\
\end{longtable}

A defensible safety case requires that untested operational envelopes,
unobservable test results, and unsupported probability claims be
presented with the exact same visual and structural prominence as
successful verifications. Hiding an untested envelope boundary or an
unverified hardware permutation in an appendix manufactures a false
impression of completeness. The fault record must explicitly list every
parameter combination where fault injection was omitted due to test
bench limitations, sensor bandwidth constraints, or destructive hardware
limits. Documenting what was never tested prevents unverified
operational zones from silently masquerading as verified safety margins
when the record is transferred to the safety case in
Chapter~\ref{sec-release}. Furthermore, every isolated entry in the
record documents the response of the machine to a single, localized
disturbance. In physical operation, however, disturbances rarely occur
in strict isolation, and verifying how a machine responds to an
individual fault leaves open the question of how it behaves when
multiple nominal disturbances strike the system at the same time.

\section{Systems Stress-Tests}\label{sec-verification-15-8}

A machine that passes every isolated fault injection test can fail when
two nominal disturbances occur at the same time. If a system possesses
\(N\) cataloged failure modes, testing every arbitrary pair requires
\(\mathcal{O}(N^2)\) injection trials, and higher-order combinations
expand exponentially. Exhaustive combinatorial injection is physically
intractable, and unguided random sampling wastes testing resources on
independent components that share no physical or logical connection.
Systems stress-testing avoids combinatorial explosion by deriving
compound fault vectors directly from shared physical resources, common
electrical paths, cascading detector dependencies, and latency budgets
that accumulate across handoffs. The following four analytical
stress-tests evaluate machines where individually compliant subsystems
combine into catastrophic physical failures:

\subsection*{Stress-Test 1: Compound Latency Accumulation across
Distributed
Handoffs}\label{stress-test-1-compound-latency-accumulation-across-distributed-handoffs}
\addcontentsline{toc}{subsection}{Stress-Test 1: Compound Latency
Accumulation across Distributed Handoffs}

\emph{The scenario.} An industrial liquid-propellant flow regulator
(\(P_{\text{inlet}} = 12.0\text{ MPa}\), nominal flow
\(15.0\text{ kg/s}\)) requires flow transients to be arrested within an
end-to-end deadline of \(250.0\text{ ms}\) to prevent chamber rupture at
\(14.5\text{ MPa}\). Execution is distributed across three nodes on a
shared \(10\text{ Mbps}\) CAN bus: an optical densitometer (ingest
budget \(80.0\text{ ms}\)), a GPU inference node (execution budget
\(120.0\text{ ms}\)), and a valve actuator node (budget
\(40.0\text{ ms}\)). In standalone tests, sensor delay
(\(75.0\text{ ms}\)), inference delay (\(115.0\text{ ms}\)), and valve
delay (\(38.0\text{ ms}\)) all pass local watchdog limits.

\emph{The systems dilemma.} When an upstream cavitation surge occurs,
sensor, planner, and actuator simultaneously experience their worst-case
individually compliant delays, while unthrottled diagnostic logging adds
\(45.0\text{ ms}\) of bus queuing. Total cumulative latency reaches
\(75.0\text{ ms} + 115.0\text{ ms} + 38.0\text{ ms} + 45.0\text{ ms} = 273.0\text{ ms}\),
exceeding the \(250.0\text{ ms}\) physical vessel limit by
\(23.0\text{ ms}\) and bursting the vessel while every local watchdog
reports nominal health.

\emph{The systems verdict: Additive Timing Composition Law.} Local
watchdog timers are mathematically incapable of guaranteeing end-to-end
physical safety when handoff delays accumulate. Verification must
enforce distributed end-to-end deadline tracking rather than isolated
component timeouts.

\subsection*{Stress-Test 2: Shared Bus Saturation and Control Loop
Starvation}\label{stress-test-2-shared-bus-saturation-and-control-loop-starvation}
\addcontentsline{toc}{subsection}{Stress-Test 2: Shared Bus Saturation
and Control Loop Starvation}

\emph{The scenario.} An autonomous platform coordinates sensing, neural
inference, safety reflexes, and evidence logging across a single shared
field bus. An unthrottled diagnostic flight-recorder memory dump pushes
bus utilization to \(98\%\).

\emph{The systems dilemma.} Bus contention does not produce an isolated
delay in a single subsystem; it simultaneously delays raw sensor packets
to the planner, setpoints to valve actuators, heartbeats to the safety
enforcer, and crash logs to the evidence ledger. The simultaneous loss
of actuator commands and enforcement telemetry starves the safety loop,
preventing the monitor from detecting an actuator freeze.

\emph{The systems verdict: Shared Interconnect Contention Invariant.}
Common communication media couple previously independent subsystems
during transients. Critical safety heartbeats and actuator telemetry
must reside on isolated physical channels or hardware-prioritized
time-triggered slots.

\subsection*{Stress-Test 3: Asymmetric Detector Dependency under Supply
Transients}\label{stress-test-3-asymmetric-detector-dependency-under-supply-transients}
\addcontentsline{toc}{subsection}{Stress-Test 3: Asymmetric Detector
Dependency under Supply Transients}

\emph{The scenario.} A power supply ripple causes an optical flow meter
to experience a digital register freeze, latching output telemetry at a
steady nominal \(4.2\text{ kg/s}\). Concurrently, a mechanical valve
blockage restricts physical flow to \(0.8\text{ kg/s}\).

\emph{The systems dilemma.} Because the latched value
(\(4.2\text{ kg/s}\)) falls within nominal bounds and generates no
framing errors, the anomaly detector never flags a fault. The neural
planner reads nominal flow and issues no corrective command, preventing
the safety enforcer from opening the auxiliary relief line. Testing the
blockage alone or the sensor freeze alone passes; only the compound test
exposes the disabled detector.

\emph{The systems verdict: Asymmetric Detector Dependency Law.} A
primary electrical fault can disable the exact sensory detector required
to contain a secondary mechanical failure. Verification must audit the
causal graph to ensure safety monitors rely on physically independent,
cross-modal transducers.

\subsection*{Stress-Test 4: Unobservable Trial Outcomes under Diagnostic
Reset}\label{stress-test-4-unobservable-trial-outcomes-under-diagnostic-reset}
\addcontentsline{toc}{subsection}{Stress-Test 4: Unobservable Trial
Outcomes under Diagnostic Reset}

\emph{The scenario.} During an aggressive multi-fault injection test on
high-speed motion hardware, bus contention causes the testbench data
logger to drop packets, and an emergency watchdog reset wipes volatile
diagnostic trace RAM.

\emph{The systems dilemma.} Because the recording apparatus failed, the
physical trajectory, peak reaction forces, and timing margins throughout
the transient cannot be reconstructed. The test team attempts to score
the trial as a ``pass'' because the physical machine did not strike
mechanical stops.

\emph{The systems verdict: Unobservable Trial Invalidation Law.} An
unobservable outcome is an invalid test and an instrumentation defect,
never evidence that the machine survived. If the testbed cannot
deterministically measure state throughout a fault transient, the run
yields zero evidentiary value.

\section{Summary}\label{sec-verification-15-9}

A demonstration shows that a machine can function under favorable
circumstances, while verification generates structured evidence that a
design claim holds when those circumstances fail. In an analytical
example of an industrial flow regulator, allowing an observer to watch
an automated control loop maintain line pressure at \(4.0\text{ MPa}\)
across three uninterrupted operating shifts proves only that the nominal
trajectory did not encounter an unhandled disturbance. It does not
establish what happens when a proportional relief valve jams at \(15\%\)
travel during a sudden pump surge, or whether the pressure transducer's
degraded bandwidth violates the intervention deadline. A demonstration
produces confidence, which is psychological and decays with the first
unexpected physical transient. Verification produces evidence, which is
an auditable record of bounded claims tested against deliberate physical
and computational stress. Downstream systems and safety cases receive
this concrete evidence rather than subjective confidence.

The operating envelope defines the physical boundary within which safety
claims are made, evaluated, and reproduced. Outside this envelope,
physical dynamics deviate from the models used during design, turning
test results into uninterpretable noise. Inside it, every parameter
boundary corresponds to a specific physical constraint, such as an inlet
fluid velocity ceiling of \(6.0\text{ m/s}\) or an ambient temperature
limit of \(65.0^\circ\text{C}\). When a stress-test injects an upstream
pressure surge of \(1.5\text{ MPa/s}\) into the regulator, the test is
meaningful only because the design explicitly claimed stability across
pressure derivatives up to \(2.0\text{ MPa/s}\). If a test exposes a
pressure overshoot exceeding the \(5.0\text{ MPa}\) structural vessel
limit while operating within this boundary, the failure is reproducible
and falsifies the design claim. If an identical pressure overshoot
occurs only after the fluid velocity exceeds \(9.0\text{ m/s}\), the
failure indicates that the test apparatus breached the operating
envelope, measuring physics that the controller was never budgeted to
reject.

Verification derives its structure from the properties established
across the preceding design stages. Every contract stated in the
pipeline becomes an explicit fault waiting to be injected into the
assembled machine. For the flow regulator, the design claims from
earlier chapters dictate a strict chain of traceability. A sensor
calibration claim from the perception pipeline becomes an injected drift
fault of \(0.2\text{ MPa}\), testing whether the diagnostic estimator
detects the discrepancy before flow estimation errors propagate to the
policy. A policy latency claim of \(115\text{ ms}\) at \(P_{99}\)
becomes an injected inference jitter fault, verifying that the nervous
system detects the schedule overrun and executes an immediate fallback
trajectory. An actuator authority claim becomes an injected valve slew
saturation fault, confirming that the mechanical intervention mechanism
isolates the line within the \(250\text{ ms}\) physical vessel budget.
Each verification vector links five distinct elements in an unbroken
chain, connecting the original design property, the injected
disturbance, the runtime detector, the nervous system response, and the
final measured physical state.

Every qualification mode operates under explicit evidentiary limits.
Hardware-in-the-loop simulation permits dense sampling of rare boundary
crossings, running thousands of synthetic cavitation transients without
wearing out physical seals, but it remains blind to unmodeled
hydrodynamic turbulence and bus transceiver heating. Physical
dynamometer testing captures true thermal dissipation, electromagnetic
interference, and mechanical backlash, but it cannot safely execute
destructive overpressurization tests to failure across hundreds of
consecutive runs. Treating a simulated pass as proof of physical
compliance confuses the fidelity of a model with physical reality.
Conversely, treating a physical test under a single ambient condition as
a universal guarantee ignores environmental coupling. Any operating
condition, input distribution, or hardware degradation state that has
not been subjected to deliberate fault injection remains an open claim
rather than an established property.

\protect\phantomsection\label{callout-takeawaysux2a-15.8}
\begin{fbxSimple}{callout-takeaways}{Key Chapter Takeaways}{}
\phantomsection\label{callout-takeaways*-15.8}

\begin{itemize}
\tightlist
\item
  \textbf{Demonstrations Build Psychological Confidence; Verification
  Generates Auditable Evidence:} Watching a machine run without failure
  proves only that nominal conditions held. Verification requires
  deliberate physical fault injection across boundary envelopes to prove
  that safety invariants hold when subsystems fail.
\item
  \textbf{Local Watchdogs Cannot Prevent Additive Timing Failure:} If
  sensor (\(75\text{ ms}\)), GPU inference (\(115\text{ ms}\)), valve
  actuator (\(38\text{ ms}\)), and diagnostic logging (\(45\text{ ms}\))
  all operate within local limits, their sum (\(273\text{ ms}\))
  breaches the physical vessel limit (\(250\text{ ms}\)). Safety
  verification must track end-to-end distributed deadline accumulation.
\item
  \textbf{Compound Fault Vectors Must Trace Shared Physical Resources:}
  Testing single faults in isolation misses catastrophic interactions.
  Electrical voltage ripples can freeze digital register outputs at
  nominal values, disabling the exact software anomaly detector needed
  to catch a concurrent mechanical valve jam.
\item
  \textbf{Unobservable Test Outcomes Yield Zero Evidentiary Value:} If
  testbed data loggers drop packets or watchdog resets wipe diagnostic
  trace RAM during a fault transient, the run is an invalid test, never
  evidence that the machine survived the fault.
\item
  \textbf{The Versioned Fault Record is the Sole Prerequisite for
  Release:} The output of verification is an immutable fault record
  cataloging verified passes, hard failures, invalid runs, and explicit
  coverage gaps handed directly to the release assurance safety case.
\end{itemize}

\end{fbxSimple}

The output of this verification process is a versioned fault record
delivered directly to the safety case assembled in
Chapter~\ref{sec-release}. This record documents the outcome of every
injected vector without aggregation or statistical smoothing. It details
the verified passes where measured bounds held with quantifiable
margins, the hard failures where detection or intervention deadlines
were missed, and the invalid tests where testbed artifacts or envelope
violations corrupted the data. It highlights the remaining coverage gaps
where test vectors could not be safely or physically executed, and it
explicitly flags unsupported probability claims where rare-event tail
distributions were estimated without sufficient sample depth.
Chapter~\ref{sec-release} builds the formal assurance argument that
justifies operating the machine, converting the fault record into
bounded claims for release. When those physical evidence bounds confront
the open-world distribution shift and rare, unmodeled hazards of
deployment, Chapter~\ref{sec-frontier} synthesizes the overarching
epistemic boundary of physical AI, establishing how runtime systems
enforce safety even when verification reaches its fundamental limits.

\chapter{\texorpdfstring{Deployment Release \& Safety
Cases\chaptersubtitle{Deciding Whether It Should Operate}}{Deployment Release \& Safety Cases}}\label{sec-release}

\begin{figure}[H]

\centering{

\includegraphics[width=1\linewidth,height=\textheight,keepaspectratio]{index_files/mediabag//Users/VJ/GitHub/MLSysBook-integrate-physical/physical/book/chapters/16-release/figures/fig_locator.pdf}

}

\caption{\label{fig-locator}\textbf{Systems Organ Focus: Nervous.}
Claims, arguments, and evidence linking empirical measurements and
runtime enforcers to operational authority across the safety lifecycle.
Deciding whether an assembled machine should operate requires auditable
proof rather than statistical optimism. Attribution: Textbook team (CC
BY 4.0).}

\end{figure}%

\emph{Under what auditable chain of physical evidence is an autonomous
machine permitted to move among humans?}

Deploying a physical AI system into human environments is an accountable
engineering decision requiring auditable evidence rather than
statistical optimism. Whether releasing an autonomous mobile
transporter, a collaborative robot arm, or a continuous chemical
reactor, a safety case cannot rely on benchmark accuracy. It demands a
structured Claim-Argument-Evidence tree connecting every operational
claim to reproducible hardware measurements and deterministic bounds. A
defensible verdict defines an explicit operational design domain,
deploys real-time monitors for distribution shift, and specifies strict
triggers that revoke operating authority the moment physical assumptions
fail.

\newpage{}

\protect\phantomsection\label{callout-objectiveux2a-16.1}
\begin{fbxSimple}{callout-objective}{Learning Objectives}{}
\phantomsection\label{callout-objective*-16.1}

After studying this section, you will be able to:

\begin{enumerate}
\def\labelenumi{\arabic{enumi}.}
\tightlist
\item
  Assemble a claim, an argument, and the evidence supporting each link
  across physical machine archetypes.
\item
  Trace every condition in a verdict back to a calibrated sensor
  measurement, physical force limit, or timing bound.
\item
  Render operate, operate-with-conditions, or do not operate verdicts,
  and state what each obliges during subsequent deployment.
\item
  Attack an argument you wrote using adversarial red-team probes and
  identify its weakest link.
\end{enumerate}

\emph{Core physical AI insight.} A verdict is an argument from evidence
to a claim, and it is only as strong as its weakest link. Physical AI
systems are never `generally safe'; operating authority is a decaying
cryptographic lease strictly conditioned on bounded ODD evidence.
Release is an accountable systems decision rather than another benchmark
test, and what it decides is one machine, inside a stated envelope, for
a stated period.

\end{fbxSimple}

\section{The Decision to Deploy}\label{sec-release-16-1}

\lettrine{E}{very design choice}\lettrine{,}{} latency budget, fallback
reflex, and fault injection test recorded across the preceding systems
architecture layers was built to withstand physical reality. But
physical reality does not inspect logs or grant operating authority. At
some point, the development phase concludes, the engineering team
prepares to hand the machine over to its operators, and an engineer must
sign a statement asserting that the machine is fit to operate. As
highlighted in the nervous system organ locator
(Figure~\ref{fig-locator}), everything the engineering team has
measured, bounded, and verified across the body, the nervous system, the
learned policy, and the runtime architecture must now be compiled into
an explicit, traceable argument that an adversarial auditor could
attempt to dismantle. The strength of that argument dictates whether the
machine is allowed to draw power, move its joints, and interact with the
physical world outside the laboratory.

Release introduces no new physical mechanisms, adds no fallback
controllers, and retrains no neural weights. It is an act of systems
adjudication over existing evidence. Until this point, engineering
effort was directed at creating capabilities and bounding their hazards.
Engineers budgeted sensor jitter on the communication bus, measured tail
latencies under distributions such as \(P_{99.9}\), implemented watchdog
timers to preempt frozen inference pipelines, and characterized
behavioral stability when the policy encounters unfamiliar physical
contact. The release process does not alter any of these measurements.
Instead, it aggregates the disparate verification artifacts (timing
logs, hardware stress tests, empirical fault rates, and formal
containment proofs) into a connected argument graph. The graph links
low-level evidence to a high-level operational claim. If the evidence is
sound, and if the assumptions regarding the operating environment remain
valid, the machine satisfies its safety envelope. If any single link in
that chain fails under inspection, the entire claim collapses.

A release verdict is not a prediction of future perfection, nor is it a
statistical guarantee that a machine will never encounter an unhandled
condition. A learned model deployed in an unconstrained physical
environment cannot be proven correct in the sense of a mathematical
theorem. Instead, the release verdict functions as an explicit, bounded
contract between an operational claim, empirical evidence, and runtime
enforcement. The operational claim defines precisely what the machine is
authorized to do, the specific physical environment in which it is
permitted to act, and the exact time interval over which the
authorization remains valid. The empirical evidence demonstrates that
the learned brain acts within acceptable behavioral envelopes during
nominal operations. When the learned model inevitably makes an error,
runtime enforcement guarantees that the deterministic nervous system
detects the transgression and forces the body into a safe state before
kinetic energy, thermal dissipation, or mechanical stress exceeds
tolerable limits.

Adjudication evaluates the compiled evidence against the operational
claim and produces one of three distinct verdicts, designated as
\textbf{Operate}, \textbf{Operate with Conditions}, and \textbf{Do Not
Operate}. A verdict of operate grants full authorization within the
requested envelope, certifying that every link in the argument graph is
supported by verified data. A verdict of operate with conditions grants
restricted authorization, narrowing the operating envelope to match the
limits of the evidence. For instance, the adjudicator might reduce
maximum allowable end-effector velocity from \(1.5\text{ m/s}\) to
\(0.5\text{ m/s}\), mandate a physical barrier between the workspace and
human personnel, or shorten the mandatory sensor recalibration interval
from thirty days to seven. A verdict of do not operate denies
authorization entirely. In traditional software engineering, refusing to
deploy code is often viewed as a scheduling failure or a lost release
cycle. In physical AI, where an uncontained failure exerts physical work
against the world, refusal is a successful engineering outcome. The
safety case functioned correctly by identifying an unsupported claim
before the machine could cause physical harm.

The rigor of this adjudication becomes clear when applied to a concrete
physical system. Throughout this text, the release process is grounded
in an autonomous tactile manipulation cell designed to sort, finish, and
transfer heavy metal castings on an active manufacturing line. The
machine comprises an articulated arm capable of exerting
\(450\text{ N}\) of grip force, carrying a \(12\text{ kg}\) payload at
velocities up to \(1.8\text{ m/s}\), and interacting directly with human
assembly workers during part handovers. The primary operational claim
asserts that the cell will perform its sorting and finishing tasks
without exceeding \(50\text{ N}\) of transient contact force against a
human worker, even during unpredicted collisions, model stalls, or
sudden power interruptions.

Establishing that claim requires an unbroken evidentiary chain. The
adjudicator must verify that the low-level tactile reflex executes at
\(1000\text{ Hz}\) with a deterministic worst-case latency of
\(P_{99.99} \le 1.8\text{ ms}\), that the motor drives physically cut
drive current within \(15\text{ ms}\) of a threshold violation, and that
the vision-language-action policy cannot command a joint trajectory
outside the mechanical work envelope even when visual sensing is
degraded by dust or fluctuating illumination. If the vision pipeline
exhibits a \(240\text{ ms}\) tail stall that delays object detection, or
if thermal buildup in the actuator windings alters torque delivery over
a four-hour operating shift, the argument fails. Adjudication strips
away optimism, replacing informal confidence with an auditable record of
who measured each quantity, what instrumentation was used, and who
accepts formal accountability for the verdict.

\section{The Release Question}\label{sec-release-16-2}

Deciding whether a physical machine should operate is fundamentally
different from evaluating the statistical performance of a model or
witnessing a successful laboratory demonstration. High validation
accuracy on an offline dataset demonstrates only that a learned policy
can reproduce training correlations under stationary data distributions.
Similarly, a flawless live demonstration shows only that one specific
sequence of states executed without failure under the particular
environmental conditions present during that test run. Neither result
answers the release question. The engineering release question asks
whether the accumulated evidence justifies delegating physical authority
over mass, force, and velocity to an autonomous system under a defined
set of operating conditions. A physical AI system does not merely
produce informational outputs that a human supervisor can review and
discard. It converts computational outputs into mechanical work against
physical objects and human bodies. When authority is granted to a
learned policy whose internal representations are unaccountable, release
is the formal engineering decision to permit that policy to actuate the
hardware because deterministic safeguards guarantee that the machine
remains within safe physical bounds under all single-point failures.

The boundaries of this delegation are formalized through the
\textbf{operational design domain}, defined as an explicit set of
measurable physical invariants. An \textbf{operational design domain
(ODD) invariant} represents a measurable physical, environmental, or
kinematic parameter boundary (such as payload mass, workspace volume,
friction coefficients, or thermal limits) within which a physical AI
system's learned policies and deterministic reflexes are certified to
maintain safe operation. In the tactile manipulation cell, these
invariants specify the geometric workspace volume as a hemisphere of
radius \(1.2\text{ m}\) centered at the arm base, payload mass between
\(0.5\text{ kg}\) and \(12.0\text{ kg}\), contact friction coefficients
satisfying \(\mu \ge 0.35\), ambient operating temperatures from
\(15^\circ\text{C}\) to \(40^\circ\text{C}\), ambient lighting between
\(300\text{ lx}\) and \(2000\text{ lx}\), and human co-presence states
restricted to a single trained operator within the collaborative
handover envelope.

Inside these boundaries lies the nominal operating zone where the
learned policy retains full authority. When environmental factors drift
toward the perimeter, such as illuminance dropping toward
\(200\text{ lx}\) or surface moisture reducing friction below
\(\mu = 0.35\), the system enters degraded conditional zones where
operating speeds are automatically curtailed or secondary sensor
modalities are engaged. Outside the boundary lies the hard refusal
envelope. If a sensor measurement violates any physical invariant, such
as detecting a \(15\text{ kg}\) payload or ambient heat exceeding
\(40^\circ\text{C}\), the nervous system revokes policy authority,
commands mechanical brakes, and prevents further operation. In heavy
autonomous fleets, such as the autonomous mining haul truck shown in
Figure~\ref{fig-16-real-autonomous-haul-truck}, this boundary
enforcement is multi-tiered: high-precision GNSS antennas and radio
telemetry masts enforce spatial geofencing within the pit corridor,
bumper radar and LiDAR pods detect foreign obstacles and dust
degradation, and dedicated safety controllers trip electro-hydraulic
service brakes the instant perimeter containment is lost.

\begin{figure}[t]

\centering{


\includegraphics[width=0.85\linewidth,height=\textheight,keepaspectratio]{/Users/VJ/GitHub/MLSysBook-integrate-physical/physical/book/chapters/16-release/figures/fig16_real_autonomous_haul_truck.jpg}

}

\caption{\label{fig-16-real-autonomous-haul-truck}\textbf{Operational
Design Domain (ODD) Enforcement in Heavy Autonomous Systems.} A
240-metric-ton Caterpillar 793D autonomous mining haul truck operating
inside a strictly geofenced open-pit mine corridor. Heavy physical AI
systems operating with tens of megajoules of kinetic energy enforce
explicit, multi-layered ODD boundaries: centimeter-accurate dual GNSS
antenna masts and wireless telemetry on the cab roof enforce spatial
corridor geofencing; bumper-mounted radar and LiDAR sensor pods detect
foreign obstacles and dust/fog atmospheric degradation; and hardwired
emergency-stop tripwires de-energize electro-hydraulic steering and
braking actuators if perimeter containment is breached. Release
authority is contingent on physical boundary guarantees rather than
statistical navigation policies. Attribution: Caterpillar Inc.~/
MineStar System Field Trial Archive (Public Domain / Fair Use under 17
U.S.C. § 107).}

\end{figure}%

As formalized in Table~\ref{tbl-16-odd-invariants}, every physical
parameter governing the operational design domain must be partitioned
into three explicit subsets: nominal operating states
(\(\mathcal{S}_{\mathrm{nom}}\)), degraded conditional states
(\(\mathcal{S}_{\mathrm{cond}}\)), and hard refusal boundaries
(\(\mathcal{S}_{\mathrm{refuse}}\)). Rather than treating environmental
boundaries as descriptive guidelines, the release contract couples each
physical invariant to a calibrated runtime sensor modality, a guaranteed
sampling period, and a deterministic hardware enforcer action. If any
environmental parameter drifts into \(\mathcal{S}_{\mathrm{cond}}\), the
nervous system automatically clamps kinematic authority; if it
penetrates \(\mathcal{S}_{\mathrm{refuse}}\), operating authority is
immediately revoked before uncontained kinetic energy can accumulate.

\begin{longtable}[]{@{}
  >{\raggedright\arraybackslash}p{(\linewidth - 10\tabcolsep) * \real{0.1667}}
  >{\raggedright\arraybackslash}p{(\linewidth - 10\tabcolsep) * \real{0.1667}}
  >{\raggedright\arraybackslash}p{(\linewidth - 10\tabcolsep) * \real{0.1667}}
  >{\raggedright\arraybackslash}p{(\linewidth - 10\tabcolsep) * \real{0.1667}}
  >{\raggedright\arraybackslash}p{(\linewidth - 10\tabcolsep) * \real{0.1667}}
  >{\raggedright\arraybackslash}p{(\linewidth - 10\tabcolsep) * \real{0.1667}}@{}}
\caption{\textbf{Operational Design Domain (ODD) Invariant Boundary
Schema and Runtime Monitor Triggers.} Quantification of physical
boundaries for the autonomous manipulation cell across nominal
(\(\mathcal{S}_{\mathrm{nom}}\)), degraded conditional
(\(\mathcal{S}_{\mathrm{cond}}\)), and hard refusal
(\(\mathcal{S}_{\mathrm{refuse}}\)) operational partitions. Each
invariant is tied to real-time physical instrumentation, deterministic
sampling rates, and hard real-time nervous system enforcer trip actions
to guarantee containment under environmental
drift.}\label{tbl-16-odd-invariants}\tabularnewline
\toprule\noalign{}
\begin{minipage}[b]{\linewidth}\raggedright
ODD Dimension / Invariant
\end{minipage} & \begin{minipage}[b]{\linewidth}\raggedright
Nominal Envelope (\(\mathcal{S}_{\mathrm{nom}}\))
\end{minipage} & \begin{minipage}[b]{\linewidth}\raggedright
Degraded Boundary (\(\mathcal{S}_{\mathrm{cond}}\))
\end{minipage} & \begin{minipage}[b]{\linewidth}\raggedright
Hard Refusal (\(\mathcal{S}_{\mathrm{refuse}}\))
\end{minipage} & \begin{minipage}[b]{\linewidth}\raggedright
Runtime Monitor \& Rate
\end{minipage} & \begin{minipage}[b]{\linewidth}\raggedright
Enforcer Action \& Latency
\end{minipage} \\
\midrule\noalign{}
\endfirsthead
\toprule\noalign{}
\begin{minipage}[b]{\linewidth}\raggedright
ODD Dimension / Invariant
\end{minipage} & \begin{minipage}[b]{\linewidth}\raggedright
Nominal Envelope (\(\mathcal{S}_{\mathrm{nom}}\))
\end{minipage} & \begin{minipage}[b]{\linewidth}\raggedright
Degraded Boundary (\(\mathcal{S}_{\mathrm{cond}}\))
\end{minipage} & \begin{minipage}[b]{\linewidth}\raggedright
Hard Refusal (\(\mathcal{S}_{\mathrm{refuse}}\))
\end{minipage} & \begin{minipage}[b]{\linewidth}\raggedright
Runtime Monitor \& Rate
\end{minipage} & \begin{minipage}[b]{\linewidth}\raggedright
Enforcer Action \& Latency
\end{minipage} \\
\midrule\noalign{}
\endhead
\bottomrule\noalign{}
\endlastfoot
\textbf{Workspace Reach \& Height} & \(r \le 1.00\text{ m}\),
\(z \in [0.10, 1.20]\text{ m}\) &
\(1.00\text{ m} < r \le 1.20\text{ m}\) (buffer) & \(r > 1.20\text{ m}\)
or \(z < 0.05\text{ m}\) & Absolute optical encoders @
\(1000\text{ Hz}\) & Clamp \(v \le 0.20\text{ m/s}\) in
\(\mathcal{S}_{\mathrm{cond}}\); apply holding brakes
(\(t_{\mathrm{brake}} \le 18\text{ ms}\)) \\
\textbf{Payload Mass \& Inertia} & \(m \in [0.5, 8.0]\text{ kg}\),
\(I_{zz} \le 0.08\text{ kg}\cdot\text{m}^2\) &
\(8.0\text{ kg} < m \le 12.0\text{ kg}\) & \(m > 12.0\text{ kg}\) or
negative residual & Wrist 6-axis load cell @ \(1000\text{ Hz}\) & Derate
Cartesian speed to \(0.40\text{ m/s}\) in
\(\mathcal{S}_{\mathrm{cond}}\); freeze trajectory on refusal \\
\textbf{Contact Surface Friction} & \(\mu \ge 0.50\) (clean elastomeric
grip) & \(0.35 \le \mu < 0.50\) (minor sheen) & \(\mu < 0.35\)
(lubricant film / slip) & Optical-tactile shear tracking @
\(500\text{ Hz}\) & Boost normal grip force (to \(450\text{ N}\)); abort
grasp to catch tray if slip persists \(> 10\text{ ms}\) \\
\textbf{Ambient Illumination} & \(E \in [500, 2000]\text{ lx}\) (factory
floor) & \(300\text{ lx} \le E < 500\text{ lx}\) or
\(2000 < E \le 3500\text{ lx}\) & \(E < 300\text{ lx}\) or
\(E > 3500\text{ lx}\) & Calibrated photodiode array @ \(60\text{ Hz}\)
& Switch to tactile-proprioceptive reflex (\(v \le 0.20\text{ m/s}\));
park arm if saturated \\
\textbf{Thermal Actuator State} &
\(T_{\mathrm{amb}} \in [15, 35]^\circ\text{C}\),
\(T_{\mathrm{coil}} \le 353\text{ K}\) &
\(35^\circ\text{C} < T_{\mathrm{amb}} \le 40^\circ\text{C}\),
\(T_{\mathrm{coil}} \in (353, 383]\text{ K}\) &
\(T_{\mathrm{amb}} > 40^\circ\text{C}\) or
\(T_{\mathrm{coil}} > 393\text{ K}\) & Stator winding thermistors @
\(10\text{ Hz}\) & Dynamic torque derating (\(50\%\) limit); hardware
gate shutdown (\(t_{\mathrm{gate}} \le 1\text{ ms}\)) on refusal \\
\textbf{Human Co-Presence} & Separation
\(d_{\mathrm{sep}} \ge 1.50\text{ m}\) &
\(0.30\text{ m} \le d_{\mathrm{sep}} < 1.50\text{ m}\) (handover) &
\(d_{\mathrm{sep}} < 0.30\text{ m}\) (unsanctioned breach) & Dual safety
pulsed ToF lidar @ \(100\text{ Hz}\) & Speed and Separation Monitoring
(\(v \propto d_{\mathrm{sep}}\)); Category 0 emergency stop
(\(t_{\mathrm{stop}} \le 40\text{ ms}\)) \\
\end{longtable}

A defensible release claim must be stated as a falsifiable physical
proposition bounded by mechanics and human injury criteria.\sidenote{\footnotesize Biomechanical
  limits for cobots are codified in ISO/TS 15066
  (\citeproc{ref-isots15066}{\emph{ISO/TS 15066} 2016}) across 29 body
  regions (\(140\text{ N}\) chest, \(65\text{ N}\) hands). UL 4600
  (\citeproc{ref-ul4600}{\emph{UL 4600} 2023}) formalizes goal-based
  safety argumentation for autonomous systems.} For the manipulation
cell handling metal castings alongside workers, the claim asserts that
under all operating conditions and all single-point sensor or compute
failures within the operational design domain, the peak contact force
delivered to a human body will not exceed \(F_{\max} \le 50\text{ N}\),
and the system will arrest all motion within a stopping time of
\(t_{\text{stop}} \le 40\text{ ms}\). This force limit derives directly
from contact mechanics and biomechanical tissue thresholds. During an
impact against compliant human tissue, the contact force is governed by
the rate of change of momentum,
\(F = \Delta p / \Delta t = m_{\text{eff}} \Delta v / \Delta t_{\text{contact}}\),
where \(m_{\text{eff}}\) is the effective moving mass of the arm and
workpiece, and \(\Delta t_{\text{contact}}\) is the duration of the
contact deformation. In an analytical example where the combined arm and
casting present an effective mass of \(m_{\text{eff}} = 18\text{ kg}\)
and compress soft tissue over a contact duration of
\(\Delta t_{\text{contact}} = 30\text{ ms}\), keeping the impact force
below the tissue pain threshold of \(50\text{ N}\) requires that the
relative contact velocity not exceed
\(v_{\text{contact}} \le F_{\max} \Delta t_{\text{contact}} / m_{\text{eff}} = (50\text{ N})(0.030\text{ s}) / (18\text{ kg}) = 0.083\text{ m/s}\).

When the machine moves at its peak transit velocity of
\(1.8\text{ m/s}\), the kinetic energy stored in the moving mass is
\(E_k = \frac{1}{2} m_{\text{eff}} v^2 = \frac{1}{2} (18\text{ kg})(1.8\text{ m/s})^2 = 29.16\text{ J}\).
An unbraked collision at this velocity would deliver a peak force far
above the allowable threshold. To maintain the release claim, the safety
system must detect unexpected contact through tactile sensing within
\(t_{\text{detect}} \le 2\text{ ms}\) and engage mechanical brakes with
a sustained deceleration of \(a_{\text{brake}} = 45\text{ m/s}^2\). The
resulting braking time is
\(t_{\text{brake}} = v / a_{\text{brake}} = (1.8\text{ m/s}) / (45\text{ m/s}^2) = 0.040\text{ s}\)
(\(40\text{ ms}\)), and the stopping distance is
\(d_{\text{stop}} = v^2 / (2 a_{\text{brake}}) = (1.8\text{ m/s})^2 / (2 \times 45\text{ m/s}^2) = 0.036\text{ m}\)
(\(36\text{ mm}\)). If the spatial clearance between the moving casting
and a worker is less than \(36\text{ mm}\), or if the braking latency
exceeds \(40\text{ ms}\), the physical energy cannot be dissipated in
time, and the release claim is falsified.

Accountability for the release decision must reside with an identifiable
lead systems adjudicator rather than being diffused across development
teams. In complex physical AI architectures, diffusion of responsibility
is a primary organizational failure mode. A vision team signs off on
mean average precision scores, an actuator supplier signs off on motor
torque curves, and a reinforcement learning team signs off on policy
convergence. Each group assumes that another layer provides the margin
of safety. The lead systems adjudicator alone assumes personal technical
authority for the release verdict, attesting that the integrated system
satisfies the safety case. This attestation binds the verdict to a set
of immutable technical artifacts, including the cryptographic hash of
the verified policy weights, the calibration logs of the tactile
sensors, the execution trace demonstrating
\(P_{99.99} \le 1.8\text{ ms}\) latency in the low-level reflex loop,
and the physical dynamometer records establishing the
\(45\text{ m/s}^2\) braking deceleration. Signing the release verdict
certifies that these artifacts form a complete evidentiary basis for
operating within the specified domain.

Every release verdict carries an expiration boundary defined by physical
degradation and environmental change. A machine that was provably safe
at commissioning can become unsafe through operational drift,
out-of-distribution physical interactions, and unmodeled mechanical
wear. Over months of operation in a manufacturing plant, particulate
contamination on optical tactile covers degrades force sensing
resolution, ambient temperature swings alter motor coil resistance, and
gear backlash in the arm joints increases tracking errors while
degrading braking efficiency. Out-of-distribution interactions, such as
an operator presenting an unmodeled workpiece with unfamiliar surface
compliance, invalidate the momentum transfer assumptions of the contact
proof. When sensor diagnostics detect parameter drift beyond tolerance,
or when total operating cycles reach the service limit of the mechanical
friction brakes, the release verdict is invalidated automatically. The
release question must then be asked again, demanding a renewed chain of
evidence before physical authority can be restored.

\section{Claims, Argument, and Evidence}\label{sec-release-16-3}

A safety case is a structured argument demonstrating that a physical
machine satisfies its safety claims within a specified operational
design domain. In systems engineering, this argument takes the form of a
\textbf{claim-argument-evidence (CAE) safety case}, adapting the
principles of Goal Structuring Notation (GSN)
(\citeproc{ref-kelly1998arguing}{Kelly 1998};
\citeproc{ref-kelly2004goal}{Kelly and Weaver 2004}) into a rigorous,
auditable analytical tool. A claim-argument-evidence safety case
decomposes a top-level physical safety proposition (such as
non-injurious kinetic force dissipation or chemical fluid containment)
into falsifiable sub-claims, connecting them via architectural and
inductive warrants to empirical test measurements, hardware stress logs,
and formal containment proofs.

The root of the tree is the top-level claim, which asserts a measurable
physical property required to prevent harm, such as guaranteeing that
the kinetic energy of an articulated manipulator dissipates before the
tool center point penetrates a human safety zone, or that an autonomous
haul truck brings its 240-ton moving mass to rest within corridor
containment. This top-level claim cannot be validated directly by
inspection. Instead, it is broken down through strategies or
intermediate arguments that partition the verification problem across
operational sub-regimes, physical hazard classes, or architectural
subsystems. Each sub-claim terminates either in supporting sub-claims or
in concrete evidentiary artifacts, bounded by explicit context and
assumptions that define the physical conditions under which the links
hold.

The validity of the entire structure depends on the distinction between
genuine evidence and demonstration. A demonstration is a record of
successful task execution over an uncharacterized test distribution,
such as recording ten thousand error-free bin-picking cycles under
nominal factory conditions. Demonstrations measure the mean performance
of a learned policy, but they provide no bound on tail behavior and
conceal failure modes that lie outside the sampled trajectory envelope.
Genuine evidence consists of falsifiable, calibrated artifacts with
quantified bounds. In a physical AI system, evidence comprises
worst-case timing distributions from real-time hardware execution
traces, calibrated sensor noise profiles under degraded environmental
states, dynamic load cell measurements recorded across deceleration
regimes, and targeted fault injection trials that deliberately force the
policy to the boundary of its training distribution. If an evidentiary
artifact lacks a quantified measurement uncertainty or fails to test the
physical failure boundary, it remains a demonstration and cannot support
an inferential link.

Connecting an item of evidence to a claim requires a warrant, which is
the inferential rule stating why the data proves the assertion. In
software systems governed by learned policies, this link is inductive,
asserting that because the system avoided failure across finite test
episodes, it will continue to avoid failure during future deployment. In
physical systems operating in continuous state spaces, this inductive
link is fundamentally incomplete without a deterministic enforcer.

To understand why finite black-box testing cannot validate
high-consequence physical claims, consider a target hazard rate
\(\lambda\) expressed in failures per operating hour. If a test rig runs
for \(N\) failure-free operating hours, and failures follow a Poisson
arrival process with rate \(\lambda\), the probability of observing zero
failures during testing is \(P(\text{zero failures}) = e^{-\lambda N}\).
To establish that the true hazard rate does not exceed \(\lambda\) at a
statistical confidence level of \(1 - \alpha\), this survival
probability must be bounded by the significance level \(\alpha\),
requiring \(e^{-\lambda N} \le \alpha\). Solving for the minimum
required test duration yields the bound
\[N = -\frac{\ln(\alpha)}{\lambda}\] For an industrial manipulator
operating beside human assembly workers, an acceptable catastrophic
hazard rate might be set at \(\lambda = 10^{-6}\text{ failures/hour}\),
representing one catastrophic failure per million operating hours. To
demonstrate this bound empirically with a standard \(95\%\) confidence
level (\(\alpha = 0.05\)), the test campaign must complete
\[N = -\frac{\ln(0.05)}{10^{-6}\text{ hours}^{-1}} \approx 2{,}995{,}732\text{ hours}\]
This analytical example requires approximately \(342\text{ years}\) of
continuous, failure-free testing on a single machine. Even a fleet of
one hundred test cells running continuously in parallel would require
more than three years of uninterrupted operation without encountering a
single unmitigated fault. Because empirical testing cannot cover the
continuous state space within practical project schedules, an inductive
link based solely on empirical sampling cannot establish safety. The
argument must bridge the inductive gap through an architectural warrant,
grounding the claim in a deterministic physical enforcer
(Chapter~\ref{sec-enforcement}) that monitors reflex latencies and
kinematics in hardware, interrupting motor power whenever learned
commands violate physical safety envelopes.

\protect\phantomsection\label{callout-mathux2a-16.2}
\begin{fbxSimple}{callout-math}{Fermi Math:}{ODD Boundary Coverage, Finite Failure-Free Hours, and Bayesian Safety Claim Confidence}
\phantomsection\label{callout-math*-16.2}

\textbf{Engineering Question:} How many failure-free test hours \(N\)
must an autonomous machine operate to substantiate a catastrophic hazard
rate of \(\lambda = 10^{-6}\text{ /h}\) at \(95\%\) statistical
confidence through empirical black-box testing alone, and how does
adding a deterministic hardware enforcer with coverage \(c = 0.9999\)
alter the required testing duration?

\textbf{1. Classical Poisson Zero-Failure Testing Model:} For a
continuous-time deployment governed by a Poisson failure arrival process
with constant hazard rate \(\lambda\), the probability of observing
\(k = 0\) failures in continuous testing time \(N\) is
\(P(k = 0 \mid \lambda) = e^{-\lambda N}\). To establish that the upper
\((1-\alpha)\) confidence bound on the true hazard rate satisfies
\(\lambda \le \lambda_{\text{target}}\), we set
\(e^{-\lambda_{\text{target}} N} \le \alpha\). Taking the natural
logarithm:
\[N = -\frac{\ln(\alpha)}{\lambda_{\text{target}}} = \frac{\ln(1/\alpha)}{\lambda_{\text{target}}}\]
For \(\alpha = 0.05\) (\(95\%\) confidence level) and
\(\lambda_{\text{target}} = 10^{-6}\text{ failures/hour}\) (one failure
per million operating hours):
\[N = \frac{-\ln(0.05)}{10^{-6}\text{ h}^{-1}} = \frac{2.99573}{10^{-6}\text{ h}^{-1}} \approx 2{,}995{,}732\text{ hours} \approx 342.0\text{ machine-years}\]
Even a testing fleet of \(100\) parallel robotic cells running
\(24\text{ hours/day}\) without interruption would require:
\[T_{\text{fleet}} = \frac{2{,}995{,}732\text{ h}}{100 \times 8760\text{ h/year}} \approx 3.42\text{ continuous calendar years without a single failure.}\]
If even a single failure occurs (\(k = 1\)), the required test hours
under the Poisson cumulative distribution
\(\sum_{i=0}^{k} \frac{(\lambda N)^i}{i!} e^{-\lambda N} \le \alpha\)
jumps to
\(N = \frac{4.74386}{10^{-6}} \approx 4{,}743{,}860\text{ hours}\)
(\(541.5\text{ years}\)).

\textbf{2. Non-Asymptotic Clopper-Pearson Exact Binomial Bounds:} In
discrete operational cycles (e.g., bin-picking grasp cycles, trajectory
handovers), let \(n\) be the number of observed cycles with \(k = 0\)
catastrophic failures. The exact Clopper-Pearson upper bound on the
per-cycle failure probability \(p\) at confidence level \(1 - \alpha\)
is: \[p_{\text{upper}} = 1 - \alpha^{1/n}\] To certify \(p \le 10^{-6}\)
per cycle at \(95\%\) confidence:
\[n = \frac{\ln(\alpha)}{\ln(1 - p_{\text{target}})} \approx \frac{\ln(0.05)}{-p_{\text{target}}} = \frac{2.99573}{10^{-6}} \approx 2{,}995{,}732\text{ discrete cycles.}\]

\textbf{3. Bayesian Posterior Risk with Conjugate Beta Prior:} Under a
Bayesian framing with prior distribution
\(p \sim \mathrm{Beta}(\alpha_0, \beta_0)\) over per-hour hazard rate
\(p\), observing \(n\) test hours with \(k = 0\) failures yields the
posterior distribution:
\[p \mid (k=0, n) \sim \mathrm{Beta}(\alpha_0, \beta_0 + n)\] Using an
uninformative Jeffreys prior (\(\alpha_0 = 0.5, \beta_0 = 0.5\)), the
posterior mean is \(\mathbb{E}[p \mid k=0, n] = \frac{0.5}{n + 1}\), and
the \(95\%\) credible upper bound \(p_{0.95}\) satisfies
\(\int_{0}^{p_{0.95}} \mathrm{Beta}(0.5, n+0.5) \, dp = 0.95\), yielding
\(p_{0.95} \approx \frac{1.92}{n}\). To reach \(p_{0.95} \le 10^{-6}\),
the test duration is \(n \ge 1{,}920{,}000\text{ hours}\).

\textbf{4. The Architectural Enforcer Multiplier:} Now consider a
dual-brain architecture where the learned policy possesses an
uncharacterized hazard rate
\(\lambda_{\text{policy}} = 10^{-2}\text{ /h}\) (one anomalous
trajectory command every 100 hours), but is enclosed by a deterministic,
hardware-isolated safety reflex enforcer with guaranteed diagnostic
coverage \(c = 0.9999\) and independent physical failure rate
\(\lambda_{\text{enforcer}} = 10^{-6}\text{ /h}\). The composite system
catastrophic hazard rate becomes:
\[\lambda_{\text{system}} = \lambda_{\text{policy}} \cdot (1 - c) + \lambda_{\text{enforcer}} = (10^{-2}) \times (10^{-4}) + 10^{-6} = 10^{-6} + 10^{-6} = 2 \times 10^{-6}\text{ /h}\]
Verifying the learned policy's nominal failure rate of
\(\lambda_{\text{policy}} \le 10^{-2}\text{ /h}\) requires only:
\[N_{\text{policy}} = \frac{-\ln(0.05)}{10^{-2}\text{ h}^{-1}} \approx 300\text{ test hours (12.5 days on one machine).}\]
Meanwhile, the enforcer's deterministic coverage \(c = 0.9999\) is
verified not through black-box endurance testing, but through formal
architectural bounds, static timing analysis on bare-metal silicon
(\(P_{100} \le 1.8\text{ ms}\)), and targeted pin-level fault injection
(\(N_{\text{inject}} = 10{,}000\) synthetic trials in minutes).

\textbf{Key Systems Insight:} Bounding safety through pure empirical
testing of learned policies in high-consequence domains is
mathematically impossible within human project schedules. An
architectural warrant backed by a deterministic hardware enforcer
shrinks required empirical test time from \(342\text{ years}\) to under
two weeks.

\end{fbxSimple}

\protect\phantomsection\label{callout-ruleux2a-16.3}
\begin{fbxSimple}{callout-rule}{Architect’s Rule of Thumb:}{The Inductive Safety Fallacy}
\phantomsection\label{callout-rule*-16.3}

Finite empirical testing proves only what occurred in the past under
sampled conditions; delegating physical authority over mass and momentum
requires deterministic architectural warrants that bound physical
consequences when inductive models inevitably encounter unmapped states.

\end{fbxSimple}

Every argument link relies on contextual assumptions about the operating
environment, and unstated assumptions represent the most common
structural vulnerability in a safety case. A warrant that derives safe
stopping distance from measured optical camera frame rates implicitly
assumes that ambient illumination remains above the sensor noise
threshold, that lens covers remain free from airborne grease, and that
internal bus arbitration introduces negligible jitter. If ambient light
drops below \(100\text{ lux}\), camera exposure times automatically
lengthen from \(5\text{ ms}\) to \(30\text{ ms}\), adding
\(25\text{ ms}\) of unmodeled latency to the perception pipeline. At a
tool speed of \(1.8\text{ m/s}\), this unstated latency shift adds
\(0.045\text{ m}\) (\(45\text{ mm}\)) to the displacement of the arm
before the braking command can be issued. The link between the vision
benchmark and the stopping distance claim breaks silently because the
physical environment drifted across an unstated boundary. A sound
argument makes every physical dependency an explicit, bounded context
element attached to the warrant, coupled to runtime diagnostic monitors
that detect when an assumption is violated.

The rule of the weakest link governs every claim-argument-evidence tree.
A safety case is not an aggregate metric where statistical
over-performance in one branch compensates for evidentiary deficits in
another. High accuracy on object recognition benchmarks does not offset
an uncharacterized \(P_{99.9}\) latency spike in the CAN bus
transceiver, nor does extensive policy simulation offset the absence of
physical dynamometer brake wear data. If a top-level claim decomposes
into sub-claims for kinetic energy dissipation, thermal limit
protection, and trajectory bounding, a failure in the evidence or
warrant of any single branch invalidates the root claim. When an
adjudicator examines the tree, the goal is not to count the number of
supporting artifacts, but to identify the single weakest link whose
failure would permit physical damage. Evaluating whether every branch
maintains an unbroken chain of evidence from physical measurements to
verified limits determines the verdict that governs machine release.

The structured decomposition of physical claims, inferential warrants,
empirical evidence artifacts, and audit acceptance thresholds is
formalized in the Claim-Argument-Evidence scorecard of
Table~\ref{tbl-16-cae-scorecard}. By auditing each branch against
rigorous physical bounds and distinguishing between architectural
warrants (enforced deterministically in hardware) and inductive warrants
(derived from finite empirical sampling), the scorecard exposes
vulnerabilities that aggregate accuracy metrics conceal. As demonstrated
in Claim C3, even when a learned model demonstrates high benchmark
accuracy under clean laboratory conditions, the emergence of an
unmonitored physical condition (such as oil contamination reducing
interface friction) breaks the inductive warrant, requiring the
adjudicator to reject unconditional release and mandate secondary
hardware tripwires.

\begin{longtable}[]{@{}
  >{\raggedright\arraybackslash}p{(\linewidth - 10\tabcolsep) * \real{0.1667}}
  >{\raggedright\arraybackslash}p{(\linewidth - 10\tabcolsep) * \real{0.1667}}
  >{\raggedright\arraybackslash}p{(\linewidth - 10\tabcolsep) * \real{0.1667}}
  >{\raggedright\arraybackslash}p{(\linewidth - 10\tabcolsep) * \real{0.1667}}
  >{\raggedright\arraybackslash}p{(\linewidth - 10\tabcolsep) * \real{0.1667}}
  >{\raggedright\arraybackslash}p{(\linewidth - 10\tabcolsep) * \real{0.1667}}@{}}
\caption{\textbf{Claim-Argument-Evidence (CAE) Safety Case Scorecard
Matrix.} Decomposition of top-level physical AI safety claims across
architectural and inductive warrants, contextual physical assumptions,
empirical evidentiary artifacts, audit acceptance thresholds, and final
adjudication verdicts. High inductive accuracy on nominal benchmarks
cannot compensate for a broken warrant, such as tactile slip
classification failing under fluid contamination (Claim C3), requiring
secondary deterministic fallback
interceptors.}\label{tbl-16-cae-scorecard}\tabularnewline
\toprule\noalign{}
\begin{minipage}[b]{\linewidth}\raggedright
Safety Sub-Claim (ID)
\end{minipage} & \begin{minipage}[b]{\linewidth}\raggedright
Argument Strategy \& Warrant Type
\end{minipage} & \begin{minipage}[b]{\linewidth}\raggedright
Underlying Assumptions \& Context
\end{minipage} & \begin{minipage}[b]{\linewidth}\raggedright
Verified Evidence Artifacts
\end{minipage} & \begin{minipage}[b]{\linewidth}\raggedright
Audit Threshold \& Confidence
\end{minipage} & \begin{minipage}[b]{\linewidth}\raggedright
Adjudication Verdict \& Interceptor
\end{minipage} \\
\midrule\noalign{}
\endfirsthead
\toprule\noalign{}
\begin{minipage}[b]{\linewidth}\raggedright
Safety Sub-Claim (ID)
\end{minipage} & \begin{minipage}[b]{\linewidth}\raggedright
Argument Strategy \& Warrant Type
\end{minipage} & \begin{minipage}[b]{\linewidth}\raggedright
Underlying Assumptions \& Context
\end{minipage} & \begin{minipage}[b]{\linewidth}\raggedright
Verified Evidence Artifacts
\end{minipage} & \begin{minipage}[b]{\linewidth}\raggedright
Audit Threshold \& Confidence
\end{minipage} & \begin{minipage}[b]{\linewidth}\raggedright
Adjudication Verdict \& Interceptor
\end{minipage} \\
\midrule\noalign{}
\endhead
\bottomrule\noalign{}
\endlastfoot
\textbf{C1: Kinetic Energy Dissipation} & Dissipate kinetic energy via
mechanical friction braking before tissue pain limit
(\(F \le 50\text{ N}\)) is reached \emph{(Architectural)} & Moving mass
\(m_{\mathrm{eff}} \le 18\text{ kg}\); tissue stiffness
\(k \le 15\text{ N/mm}\); dry workpiece interface & Dynamometer
deceleration profiles (\(N = 500\) runs); 6-axis load cell impact logs &
Peak force \(F_{\max} \le 42.5\text{ N}\) (\(15\%\) margin);
\(t_{\mathrm{stop}} \le 40\text{ ms}\) at \(P_{99.99}\) & \textbf{Pass}
/ Spring-applied electromechanical brakes \\
\textbf{C2: Real-Time Reflex Latency} & Nervous-system enforcer preempts
policy command and cuts inverter drive current within deadline
\emph{(Architectural)} & Dedicated real-time CPU core; zero DMA
starvation; isolated peripheral bus & Logic analyzer traces
(\(1\times 10^7\) cycles) under synthetic \(95\%\) PCIe bus load & Loop
jitter \(< 45\,\mu\text{s}\); total response
\(t_{\mathrm{sense+gate}} \le 2.0\text{ ms}\) at \(P_{100}\) &
\textbf{Pass} / FPGA-mapped latching fault interlock \\
\textbf{C3: Optical-Tactile Slip Detection} & Visuo-tactile network
classifies shear displacement before workpiece loss or shear force
exceeds \(20\text{ N}\) \emph{(Inductive)} & Clean elastomeric contact
interface (\(\mu \ge 0.50\)); unmarred silicone skin & High-speed camera
marker tracking and tactile benchmark (\(N = 100{,}000\) cycles) & Slip
detection latency \(t_{\mathrm{detect}} \le 12\text{ ms}\);
\(\mathrm{FNR} \le 10^{-4}\) under clean \(\mu\) & \textbf{Fail (Oily
Parts)} / Secondary current tripwire
(\(F_{\mathrm{trip}} \le 18\text{ N}\)) \\
\textbf{C4: Perception Throughput \& Jitter} & Policy inference
maintains throughput to support dynamic collision avoidance
\emph{(Inductive)} & Illuminance \(E \in [300, 2000]\text{ lx}\);
optical haze \(< 5\%\); nominal supply voltage & Hardware-in-the-loop
profiling (\(250\text{ h}\)) across thermal envelope
(\(15\text{--}40^\circ\text{C}\)) & \(P_{99.9}\) inference latency
\(\le 12.0\text{ ms}\); tail stall \(P_{\max} \le 16.0\text{ ms}\) &
\textbf{Pass} / Synchronous inference watchdog timer \\
\textbf{C5: Permission Path Independence} & High-priority policy threads
cannot corrupt, starve, or bypass safety assertions or E-stops
\emph{(Architectural)} & Hardware memory domain separation; isolated
power rails; dedicated E-stop lines & Fault-injection suite (kernel
panics, DMA bus flooding, voltage droop testing) & Zero unserviced
E-stop assertions across \(1\times 10^4\) trials; switchover
\(t \le 1.0\text{ ms}\) & \textbf{Pass} / Galvanic safety relay
de-energizing DC bus \\
\end{longtable}

\section{The Three Verdicts}\label{sec-release-16-4}

The adjudication of a safety case resolves the claim-argument-evidence
tree into three distinct verdicts, comprising unconditional operation,
conditional operation, and refusal to operate.\sidenote{\footnotesize ISO 26262 Part
  10 (\citeproc{ref-iso26262}{\emph{ISO 26262} 2018}) mandates safety
  cases under hardware faults. Because neural policies also fail from
  distribution shift, ISO 21448 (SOTIF
  (\citeproc{ref-iso21448}{\emph{ISO 21448} 2022})) requires
  verification to shrink unknown unsafe states.} Every verdict
establishes an explicit operating contract and obliges concrete
engineering actions afterwards. The first verdict, unconditional
operation, applies when every hazard identified during verification
possesses an unbroken chain of evidence grounded in both deterministic
nervous-system enforcement and immutable physical limits. Under this
verdict, no credible failure mode relies on the statistical benevolence
or unverified generalization of a learned policy. Kinetic energy
dissipation remains within thermal and mechanical brake ratings under
worst-case payload velocity, and maximum tracking deviations are bounded
by hardware interlocks. Unconditional release does not imply that the
machine cannot fail, but rather that every known hazard is governed by a
verified physical or architectural mitigation that preserves safety
across the entire operational design domain.

The second verdict, conditional operation, applies when the evidence
demonstrates that the machine is safe within a restricted operational
subset, but lacks the warrants required for unconstrained deployment.
Rather than treating release as an all-or-nothing threshold, engineering
release can enforce physical, kinematic, or environmental boundaries to
eliminate unverified hazard states. These operational boundaries include
de-rating peak actuator torque or joint velocity, restricting ambient
operating conditions such as temperature ranges or illumination levels,
imposing physical exclusion barriers that keep human personnel outside
the maximum reach radius of the manipulator, or activating runtime
telemetry guardrails within the nervous system. In heavy industrial
robotics, when kinetic energy and force dissipation cannot be bounded
unconditionally for unseparated human co-presence, conditional release
requires hard physical fencing and certified safety interlocks
(Figure~\ref{fig-16-real-robot-safety-cell}). Dual-channel perimeter
interlock switches and optical safety curtains ensure that the hazardous
workspace is de-energized before human workers can enter, converting an
unmodeled contact risk into a deterministic physical exclusion boundary.

\begin{figure}[t]

\centering{

\includegraphics[width=0.85\linewidth,height=\textheight,keepaspectratio]{/Users/VJ/GitHub/MLSysBook-integrate-physical/physical/book/chapters/16-release/figures/fig16_real_robot_safety_cell.jpg}

}

\caption{\label{fig-16-real-robot-safety-cell}\textbf{Physical Exclusion
Barriers and Certified Safety Interlocks for Conditional Release.} Heavy
industrial robotic manipulators operating inside a high-integrity safety
enclosure equipped with interlocked perimeter fencing, warning placards,
access interlocks, and automated shutdown circuits at the Pueblo
Chemical Agent-Destruction Pilot Plant. When a physical AI system's
kinetic energy exceeds human biomechanical injury thresholds
(\(F_{\max} > 50\text{ N}\) per ISO/TS 15066) and learned policies
cannot provide deterministic non-collision guarantees under all sensory
conditions, an \emph{Operate with Conditions} verdict mandates hard
physical exclusion boundaries. Safety-rated interlock switches on
perimeter gates de-energize actuator drive stages before personnel can
penetrate the hazardous kinematic envelope. Attribution: U.S. Army
Chemical Materials Activity / PEO ACWA (Public Domain).}

\end{figure}%

Every condition attached to a release verdict is a binding runtime
contract rather than an operational recommendation. A conditional
release is valid only if every restriction maps directly to a
measurable, real-time observable in hardware or the nervous system,
governed by a deterministic trigger threshold that transitions the
machine to a safe stop state if violated. Consider an analytical example
of an articulated manipulator with a total reach of \(0.80\text{ m}\)
designed for automated packaging. In the nominal specification, the
end-effector translates at \(1.5\text{ m/s}\) with an effective moving
mass of \(40\text{ kg}\), yielding a peak translational kinetic energy
of \(45\text{ J}\). If contact verification data establishes that
transient impact energy must remain below \(15\text{ J}\) to prevent
injury during unseparated human collaboration, and the system lacks
certified dynamic separation sensors, the release authority issues a
conditional verdict that de-rates the maximum velocity to
\(0.86\text{ m/s}\). Because kinetic energy scales quadratically with
velocity, this reduction constrains peak kinetic energy to
\((0.86 / 1.5)^2 \times 45\text{ J} = 14.8\text{ J}\). The
nervous-system enforcer monitors optical encoder velocity at
\(1000\text{ Hz}\) with a hard trip limit set at \(0.90\text{ m/s}\). If
an anomalous policy command or external disturbance drives the tool
velocity above \(0.90\text{ m/s}\) for more than \(2\text{ ms}\) (two
consecutive control cycles), the enforcer interrupts motor drive current
and engages mechanical friction brakes within \(18\text{ ms}\), bounding
the total stopping displacement to less than \(0.020\text{ m}\)
(\(20\text{ mm}\)). The condition is not left to operator compliance or
policy training, but is enforced by real-time hardware execution.

The third verdict, refusal to operate, is the mandatory outcome whenever
the safety case contains missing warrants, uncharacterized tail
latencies, incomplete fault-injection coverage, or unmonitored failure
modes. If a system exhibits a \(P_{99.99}\) control loop latency of
\(85\text{ ms}\) against a physical stability deadline of
\(20\text{ ms}\), or if memory corruption tests under bus contention
were omitted, the hazard remains unquantified and cannot support
release. Refusal is likewise required when empirical observations
contradict analytical physical models, such as measured motor winding
temperatures rising \(15\text{ K}\) above lumped-parameter thermal
predictions during repetitive braking cycles. Under these conditions,
operating the machine is an unmodeled gamble rather than an engineered
deployment.

A refusal verdict is an engineering deliverable rather than an
administrative failure. In industrial environments, release evaluations
can degrade into political friction between delivery schedules and
unquantified safety apprehensions. A formal refusal verdict eliminates
this friction by isolating the precise broken link in the argument tree,
converting organizational impasse into an actionable, falsifiable
testing mandate. If release is refused because transient bus arbitration
causes uncharacterized packet loss during simultaneous camera
transmission and motor feedback updates, the engineering team does not
debate subjective risk tolerances. Instead, they construct a
deterministic bus schedule or partition direct memory access channels,
and measure communication performance until they demonstrate a
\(P_{100}\) transmission latency below \(4.0\text{ ms}\) under maximum
synthetic bus load. The refusal verdict defines the exact evidentiary
test that must pass before re-adjudication can occur.

None of the three verdicts represents a permanent, static grant of
authority. Operating under unconditional or conditional release creates
ongoing obligations that must be maintained throughout the active
deployment lifecycle. Mechanical linkages develop backlash, friction
surfaces wear, optical sensor lenses accumulate particulate film, and
ambient distributions change. A valid release verdict specifies
mandatory recalibration intervals, such as recalibrating joint torque
sensors every \(500\text{ operating hours}\) or verifying brake pad
clearance every \(250000\text{ actuation cycles}\). It also establishes
explicit drift budgets for runtime telemetry. If the \(P_{99}\)
trajectory tracking error of the nervous system drifts from
\(0.004\text{ m}\) at initial release to \(0.012\text{ m}\) under
accumulated mechanical wear, the drift budget is exhausted, and the
release authorization automatically lapses until maintenance restores
the baseline. High-frequency ring buffers in the nervous system must
continuously record pre-trigger and post-trigger states during every
enforcer trip, preserving the audit trail required to confirm that the
empirical assumptions underlying the verdict remain true over time.

The operational criteria, standards mappings, enforcement bindings, and
post-release obligations governing each of the three release outcomes
are systematized in Table~\ref{tbl-16-verdicts-taxonomy}. Rather than
treating safety adjudication as a subjective management sign-off, the
taxonomy establishes precise mathematical and empirical entry conditions
anchored in international safety standards, including ISO 26262 for
functional safety under hardware faults, ISO 21448 (SOTIF) for
functional insufficiencies in learned models, ISO/TS 15066 for
collaborative robot biomechanical limits, and UL 4600 for goal-based
autonomous safety cases. Each verdict binds the physical machine to a
distinct level of operational authority and mandates continuous
telemetry verification to preserve that authority throughout its
lifecycle.

\begin{longtable}[]{@{}
  >{\raggedright\arraybackslash}p{(\linewidth - 10\tabcolsep) * \real{0.1667}}
  >{\raggedright\arraybackslash}p{(\linewidth - 10\tabcolsep) * \real{0.1667}}
  >{\raggedright\arraybackslash}p{(\linewidth - 10\tabcolsep) * \real{0.1667}}
  >{\raggedright\arraybackslash}p{(\linewidth - 10\tabcolsep) * \real{0.1667}}
  >{\raggedright\arraybackslash}p{(\linewidth - 10\tabcolsep) * \real{0.1667}}
  >{\raggedright\arraybackslash}p{(\linewidth - 10\tabcolsep) * \real{0.1667}}@{}}
\caption{\textbf{The Three Release Verdicts and Statutory Adjudication
Criteria.} Comprehensive taxonomy mapping the three release verdicts to
their required evidentiary criteria, statutory safety standards,
real-time hardware enforcement mechanisms, operational envelope
constraints, and mandatory post-release maintenance obligations. Release
is a bounded, enforceable contract whose validity expires if standing
conditions or telemetry drift budgets are
violated.}\label{tbl-16-verdicts-taxonomy}\tabularnewline
\toprule\noalign{}
\begin{minipage}[b]{\linewidth}\raggedright
Release Verdict Class
\end{minipage} & \begin{minipage}[b]{\linewidth}\raggedright
Evidentiary \& Physical Entry Criteria
\end{minipage} & \begin{minipage}[b]{\linewidth}\raggedright
Statutory \& Standards Mapping
\end{minipage} & \begin{minipage}[b]{\linewidth}\raggedright
Runtime Enforcement Mechanism
\end{minipage} & \begin{minipage}[b]{\linewidth}\raggedright
Operational Restrictions \& Envelope
\end{minipage} & \begin{minipage}[b]{\linewidth}\raggedright
Post-Release Obligations \& Expiry
\end{minipage} \\
\midrule\noalign{}
\endfirsthead
\toprule\noalign{}
\begin{minipage}[b]{\linewidth}\raggedright
Release Verdict Class
\end{minipage} & \begin{minipage}[b]{\linewidth}\raggedright
Evidentiary \& Physical Entry Criteria
\end{minipage} & \begin{minipage}[b]{\linewidth}\raggedright
Statutory \& Standards Mapping
\end{minipage} & \begin{minipage}[b]{\linewidth}\raggedright
Runtime Enforcement Mechanism
\end{minipage} & \begin{minipage}[b]{\linewidth}\raggedright
Operational Restrictions \& Envelope
\end{minipage} & \begin{minipage}[b]{\linewidth}\raggedright
Post-Release Obligations \& Expiry
\end{minipage} \\
\midrule\noalign{}
\endhead
\bottomrule\noalign{}
\endlastfoot
\textbf{Operate (Unconditional Release)} & Unbroken CAE tree across all
hazards; all physical bounds (\(F_{\max}, E_k, t_{\mathrm{stop}}\))
verified on hardware; \(P_{100}\) reflex latency bounded under bus load;
zero unmonitored premises & ISO 26262 ASIL-D; ISO 21448 / SOTIF Clause
8; UL 4600 Clause 5; ISO/TS 15066 & Hardware-isolated nervous-system
safety enforcer @ \(1000\text{ Hz}\); windowed watchdog timers; inverter
interlocks & Full ODD operational envelope
(\(v_{\max} = 1.8\text{ m/s}\), \(m \le 12.0\text{ kg}\), collaborative
workspace) & Telemetry ring-buffer logging; torque recalibration every
\(500\text{ h}\); drift audit every \(250\text{k}\) cycles; hard expiry
@ \(2000\text{ h}\) \\
\textbf{Operate with Conditions (Restricted Release)} & Safety argument
valid within bounded sub-envelope, but exhibits unverified warrants or
narrow margins at physical extremes (e.g., wet surfaces, tail perception
jitter) & ISO 21448 Clause 9; ISO 10218-2 / ISO/TS 15066 Clause 5.5;
OSHA 1910.212 & Multi-tier dynamic enforcer clamping torque/speed
registers; optical dryness station interlock; light curtains & De-rated
Cartesian speed (\(v_{\max} \le 0.40\text{ m/s}\)); payload mass
\(m \le 1.5\text{ kg}\); mandatory surface-dryness check;
\(F_{\mathrm{trip}} \le 18\text{ N}\) & Daily sensor zero-point drift
check; weekly brake lining inspection; high-frequency telemetry stream;
accelerated expiry @ \(500\text{ h}\) \\
\textbf{Do Not Operate (Refusal / Rejection)} & One or more critical CAE
links broken; uncharacterized tail latency
(\(P_{99.99} > t_{\mathrm{margin}}\)); missing fault injection coverage;
empirical data contradicts analytical models & UL 4600 Clause 14; ISO
26262 Part 4 Clause 7; Machinery Directive 2006/42/EC Annex I & Hardware
OTP fuse latching; cryptographic bootloader inhibit; DC bus de-energized
via safety contactor & Zero operational authority; machine prohibited
from drawing actuator power or executing joint motion & Formal refusal
defect brief issued; remedial test mandate defined (e.g., bus
partitioning, feeder redesign); full re-adjudication \\
\end{longtable}

\section{The Verdict's Standing Conditions}\label{sec-release-16-5}

A release verdict is not an intrinsic property of the machine. It is a
proposition conditional on explicit physical and operational premises
that must hold across every execution cycle. When an engineering team
grants release, they assert that empirical and analytical evidence
justifies operating the machine if and only if those underlying premises
remain valid. Decomposing the safety argument isolates three distinct
classes of premises, comprising those that are continuously monitored at
runtime, those that are periodically re-established through physical
inspection or recalibration, and those that are assumed for the life of
the release. In an analytical example of a mobile manipulator handling
\(10.0\text{ kg}\) steel billets, the assumption that the concrete floor
friction coefficient \(\mu\) remains above \(0.60\) is treated as an
assumed premise for the warehouse facility, requiring no onboard sensor.
The premise that joint torque sensor drift remains within
\(\pm 0.5\text{ N}\cdot\text{m}\) is periodically re-established through
a calibration fixture every \(500\text{ operating hours}\). The premise
that the payload mass does not exceed \(12.0\text{ kg}\) is a
continuously monitored premise, evaluated by the nervous system from
load cell telemetry during every grasp attempt before arm acceleration
begins.

Every monitored premise requires both a deterministic runtime detector
and a pre-allocated response. If a safety case claims that image
exposure latency remains under \(16.0\text{ ms}\) or that CAN bus frame
drops remain below \(1\times 10^{-5}\), but the embedded software lacks
a hardware timestamp comparator or a frame error counter, that premise
has not been monitored. It has silently degraded into an unverified
assumption. An unmonitored premise blinds the machine to its own
boundary crossings, converting an explicit safety argument into an
unmodeled gamble. The runtime detector must evaluate the premise against
physical bounds at a frequency set by the physical dynamics of the
failure mode. In an analytical example of a high-torque brushless
actuator, the thermal model predicts stator winding insulation breakdown
at temperatures exceeding \(393\text{ K}\). If the motor thermistor
samples temperature at \(10\text{ Hz}\), the detector measures the
thermal rise within \(100\text{ ms}\), identifying excessive current
draw before the copper coil reaches destructive plastic deformation and
commanding an immediate current derating.

Evaluating monitored premises at runtime exposes a fundamental
trichotomy in envelope membership. Software systems outside physical
domains often treat safety conditions as binary booleans, presuming that
in the absence of an explicit fault flag, execution proceeds normally.
In physical AI, envelope membership occupies three distinct states,
namely known-true, known-false, and unknown. When sensor telemetry is
verified, timestamps fall within deadline bounds, and state estimators
converge with bounded covariance, envelope membership is known-true.
When a physical limit is breached, such as trajectory tracking error
exceeding \(0.050\text{ m}\), the state is known-false. When an optical
encoder drops communication, a camera frame arrives with a cyclic
redundancy check mismatch, or an extended Kalman filter diverges, the
machine cannot determine whether it is inside or outside the verified
boundary. In the presence of mass and momentum, an unknown state must be
treated as outside the envelope rather than inside. Presuming safety
during an epistemic gap allows unmeasured kinetic energy to accumulate
in unmodeled configurations.

\protect\phantomsection\label{callout-definitionux2a-16.4}
\begin{fbxSimple}{callout-definition}{Definition:}{The Envelope Trichotomy}
\phantomsection\label{callout-definition*-16.4}

A three-state classification of runtime system health relative to its
certified safety envelope---\emph{known-true} (verified nominal state
with bounded telemetry covariance), \emph{known-false} (verified safety
boundary violation triggering immediate hardware refusal), and
\emph{unknown} (epistemic state uncertainty due to sensor dropouts, bus
jitter, CRC corruption, or state estimator divergence, which must be
strictly adjudicated as outside the operational envelope).

\end{fbxSimple}

\protect\phantomsection\label{callout-ruleux2a-16.5}
\begin{fbxSimple}{callout-rule}{Architect’s Rule of Thumb:}{The Epistemic Invariant}
\phantomsection\label{callout-rule*-16.5}

In physical AI, ignorance is not neutral: any state whose envelope
membership is unknown carries unmeasured momentum and must be treated as
known-false by immediately triggering deterministic safe-state
containment.

\end{fbxSimple}

The response to an unevaluable premise must be derived directly from the
physical consequence of continuing motion without telemetry, rather than
from subjective severity classifications. If a mobile base travels at
\(1.5\text{ m/s}\) with a total mass of \(80\text{ kg}\), it carries
\(90\text{ J}\) of kinetic energy. If its primary lidar sensor stops
transmitting updates, each millisecond of blind travel carries the
machine \(1.5\text{ mm}\) closer to an unmapped obstacle, exhausting the
stopping distance available to the mechanical brakes. The system must
select among three deterministic fallback responses, established during
the safety case rather than at the moment of confusion. A degraded mode
continues operation when redundant sensors preserve sufficient state
observability to bound positional uncertainty, such as falling back to
wheel odometry and halving the maximum commanded velocity to
\(0.75\text{ m/s}\). Restricting authority clamps actuator torque limits
or constrains policy outputs to a precomputed passive corridor when
secondary telemetry fails. An emergency stop cuts drive power and
applies mechanical friction brakes when state observability drops below
the minimum required to guarantee collision avoidance within the
available braking clearance.

Even when every monitored premise remains known-true and periodic
recalibrations pass, the verdict itself carries an explicit expiration
timestamp and cycle budget. Operating hours and mechanical cycles
introduce irreversible physical transformations that runtime telemetry
cannot observe. Gear teeth develop surface pitting, elastomer seals
stiffen, structural fasteners lose clamp preload, and subtle
environmental shifts alter the distribution of visual inputs presented
to learned models. A software modification or firmware update
immediately revokes the active release verdict, because changing a
single neural network checkpoint or communication driver invalidates the
timing, tracking, and latency bounds established during prior empirical
evaluation. Establishing an analytical expiry of
\(2000\text{ operating hours}\) or
\(1\times 10^{6}\text{ manipulation cycles}\) forces a formal
re-adjudication of the safety case, ensuring that the evidence base is
systematically re-examined before physical degradation outpaces the
safety margins of the machine.

Enforcing these standing conditions cannot be delegated to an
asynchronous user-space daemon or a post-facto logging service. A
telemetry entry recording a timing violation or an over-temperature
state provides zero protection if the robotic link has already deflected
into an unmodeled workpiece. Enforcement must be routed directly through
the nervous system architecture of Chapter~\ref{sec-enforcement}, where
the safety enforcer intercepts control signals before they reach the
actuator power stages.\sidenote{\footnotesize Runtime standing condition enforcement
  is anchored in hardware to bypass OS kernel panics. Windowed watchdog
  timers and memory-mapped registers gate motor inverters via hardware
  interlocks within microseconds.} The nervous system executes the
deterministic detector logic synchronously within the hard real-time
loop, checking monitored premises and envelope membership at the primary
motor update rate of \(1000\text{ Hz}\). If a standing condition
evaluates to unknown or known-false, the enforcer does not request the
learned policy to generate an alternate trajectory. It clamps actuator
commands, engages the mechanical holding brakes within
\(1.0\text{ ms}\), and asserts a safe stop. By placing standing
conditions inside the only subsystem capable of refusing motion, the
release verdict remains an active physical contract enforced across
every microsecond of machine operation.

\section{A Case Worked in Full}\label{sec-release-16-6}

Consider an autonomous robotic workstation tasked with inserting
machined steel transmission splines into an engine housing fixture
adjacent to human assembly workers. The physical body comprises a
six-axis articulated arm equipped with soft elastomeric optical-tactile
fingertips, drawing its mechanical compliance parameters from
Chapter~\ref{sec-body} and dual tactile sensor specifications from
Chapter~\ref{sec-nervous}. Trajectory generation is governed by a
learned visuo-tactile insertion policy (Chapter~\ref{sec-planning})
running on an embedded accelerator with a worst-case inference latency
of \(12\text{ ms}\) (Chapter~\ref{sec-perception}). The nervous system
enforces reflexes through a deterministic safety engine operating at
\(1000\text{ Hz}\) (Chapter~\ref{sec-enforcement}), backed by empirical
stress boundaries derived from the fault-injection matrix of
Chapter~\ref{sec-verification}. As structured in the Goal Structuring
Notation (GSN) (\citeproc{ref-kelly1998arguing}{Kelly 1998}) safety case
tree in Figure~\ref{fig-16-cae-tree}, the top-level release claim
submitted for adjudication asserts that the manipulator will never exert
a normal contact force exceeding \(50\text{ N}\) or a shear force
exceeding \(20\text{ N}\) on an unexpected physical obstacle or an
operator's hand across any operational state, including learned model
misclassifications and visual sensor dropouts.

\begin{figure}[t]

\centering{

\includegraphics[width=1\linewidth,height=\textheight,keepaspectratio]{index_files/mediabag//Users/VJ/GitHub/MLSysBook-integrate-physical/physical/book/chapters/16-release/figures/fig16_cae_tree.pdf}

}

\caption{\label{fig-16-cae-tree}\textbf{Claim-Argument-Evidence (CAE)
Goal Structuring Tree for the Autonomous Tactile Workstation.} The
top-level operational safety claim (\(C_1\)) decomposes into three
sub-claims governing soft tactile reflex containment (\(C_{1.1}\)),
motor torque tripwire protection (\(C_{1.2}\)), and deterministic
hardware execution (\(C_{1.3}\)). An adversarial red-team audit exposes
a critical defeater (\(D_1\)): residual cutting oil lowers contact
friction (\(\mu = 0.85 \to 0.12\)), allowing the workpiece to slip
across the silicone pad without internal marker deformation and silently
disabling the tactile reflex. Because the unmitigated secondary current
tripwire allows peak contact forces of \(88\text{ N} > 50\text{ N}\),
unconditional release is refused. Adjudication establishes an
\emph{Operate with Conditions} verdict by enforcing three binding
runtime contracts: restricting payload mass to \(m \le 1.5\text{ kg}\),
inserting an automated optical surface-dryness check at the parts
feeder, and tightening the real-time torque observer tripwire to
\(F_{\text{trip}} = 18\text{ N}\), restoring peak impact forces to
\(F_{\text{peak}} = 38\text{ N} \le 50\text{ N}\) and contact impulse to
\(1.1\text{ N}\cdot\text{s}\). Attribution: Authors / GSN Safety
Architecture Case Study (CC BY-NC-ND 4.0).}

\end{figure}%

Under nominal operating assumptions, the evidence chain linking this
claim to physical reality appears complete. When the learned policy
moves the arm at its maximum allowable velocity of \(0.40\text{ m/s}\),
an unexpected collision with a stationary object compresses the
fingertip elastomer, which has an analytical stiffness of
\(k = 4.0\text{ N/mm}\). The optical-tactile perception pipeline detects
internal elastomer shear displacement within a worst-case perception
window of \(12\text{ ms}\) at \(P_{99.9}\). Upon receiving this signal,
the safety enforcer in the nervous system trips a reflex in
\(8\text{ ms}\), clamping joint torque commands and commanding the
mechanical brakes. Because braking deceleration halts the arm in an
additional \(5\text{ ms}\), total displacement during the
\(25\text{ ms}\) response window is \(10\text{ mm}\), of which the soft
elastomer absorbs \(2.5\text{ mm}\) before the joint structure arrests.
The peak normal force reaches \(10\text{ N}\), well below the
\(50\text{ N}\) limit. Fault-injection trials across \(1\times 10^{5}\)
synthetic sensor dropouts confirm that whenever tactile perception
functions, the physical reflex halts the manipulator before normal force
reaches \(14\text{ N}\).

The evidence chain breaks when subjected to boundary analysis of the
operating envelope. The optical-tactile slip detector operates by
tracking the motion of printed surface markers embedded inside the
transparent silicone fingertip skin. That mechanism relies on surface
friction to transfer shear strain from the workpiece to the elastomer.
In an analytical machining cell environment, machined splines carry
residual cutting fluid that reduces the contact friction coefficient
from \(\mu = 0.85\) to \(\mu = 0.12\). When an oily spline contacts an
obstacle, the metal workpiece slides freely across the silicone pad
without inducing internal marker deformation. The learned perception
model observes no shear strain and fails to declare a slip event.
Because the tactile reflex relies entirely on the perception model's
slip classification, the primary safety loop fails silently. The
manipulator continues driving forward with its full actuator authority,
converting what was modeled as a soft, monitored contact into an
unmonitored rigid impact between steel and the surrounding environment.

With the primary tactile reflex disabled, the safety of the workstation
falls to a secondary safety mechanism, specifically a motor-current
threshold tripwire running in the joint servo drives at
\(1000\text{ Hz}\). When mechanical resistance impedes motion, motor
current rises until the measured joint torque exceeds expected
feedforward torque by \(\Delta \tau = 3.5\text{ N}\cdot\text{m}\), which
translates to an end-effector trip force threshold of
\(F_{\text{trip}} = 35\text{ N}\) across an effective moment arm of
\(0.50\text{ m}\). To evaluate whether this secondary mechanism
satisfies the safety claim, the physical stopping profile must be
derived from first principles. Consider an analytical impact against a
human hand backed by a rigid jig, modeled with a contact stiffness of
\(k_{\text{contact}} = 15\text{ N/mm}\), with the arm moving at
\(v_0 = 0.40\text{ m/s}\) and carrying an effective inertia of
\(m_{\text{eff}} = 4.5\text{ kg}\) (\(2.5\text{ kg}\) of reflected arm
mass plus a \(2.0\text{ kg}\) workpiece).

From the instant of initial contact at \(t = 0\), the contact force
rises at a rate of \(dF/dt = k_{\text{contact}} v_0 = 6.0\text{ N/ms}\).
The contact force reaches the tripwire threshold of
\(F_{\text{trip}} = 35\text{ N}\) at \(t_1 = 5.83\text{ ms}\). Once
tripped, current-filtering latency adds
\(t_{\text{sense}} = 6\text{ ms}\), the inverter gate driver cuts power
in \(t_{\text{gate}} = 1\text{ ms}\), and the electromechanical holding
brakes take \(t_{\text{brake}} = 18\text{ ms}\) to reach full braking
torque, during which the arm continues forward. Peak normal force occurs
as kinetic energy is absorbed by the combined braking deceleration of
\(a_{\text{brake}} = 8.0\text{ m/s}^2\) and contact deflection, reaching
a peak force of \(88\text{ N}\) at \(t = 28\text{ ms}\). Integrating the
force over the total contact duration of
\(t_{\text{end}} = 62\text{ ms}\) yields a contact impulse integral of
\(\int_{0}^{t_{\text{end}}} F(t) \, dt = 2.7\text{ N}\cdot\text{s}\).
Because the peak force of \(88\text{ N}\) exceeds the \(50\text{ N}\)
ceiling and the unmonitored sliding contact produces a shear force of
\(24\text{ N}\), the secondary protection mechanism fails to enforce the
top-level claim.

When presented with this evidence tree, an adjudicator cannot issue an
unconditional release verdict. The primary evidence link supporting the
tactile slip reflex is broken by fluid contamination, and the secondary
motor-current tripwire cannot prevent excessive contact impulse under
the nominal payload mass of \(2.0\text{ kg}\). The adjudication
therefore issues a rejection of the unconditional proposal and
constructs an amended contract for an \textbf{Operate with Conditions}
verdict. To restore mathematical bounds on contact force, the
operational envelope is restricted along three dimensions. First, the
maximum allowable payload mass is capped at \(1.5\text{ kg}\), lowering
the total moving mass to \(m_{\text{eff}} = 4.0\text{ kg}\) and reducing
system kinetic energy. Second, an optical station is added to the
feeding fixture to perform a surface-dryness check before part
acquisition, rejecting any component displaying fluid sheen. Third, the
motor-current tripwire threshold is tightened from \(35\text{ N}\) to
\(18\text{ N}\) by running joint torque estimation directly on the
real-time safety bus, cutting inverter power within \(2\text{ ms}\) if
tracking error deviates from nominal.

Recalculating the physical stopping profile under these three standing
conditions restores the safety argument. With a \(1.5\text{ kg}\)
payload, an \(18\text{ N}\) torque-observer tripwire, and a guaranteed
dry contact surface, the tactile slip detector functions reliably within
its \(12\text{ ms}\) latency budget. Even in the event of an unmodeled
tactile sensor hardware failure, the faster \(18\text{ N}\) secondary
tripwire halts motion with a peak normal force of \(38\text{ N}\) and an
impulse integral of \(1.1\text{ N}\cdot\text{s}\), remaining safely
below the \(50\text{ N}\) limit. The release verdict is granted, but its
validity is tied strictly to the runtime enforcement of the optical part
inspection, the payload mass ceiling, and the hard real-time execution
of the torque observer. The argument holds because every link from
physical mechanics to sensor latency has been quantified and bounded.
Whether that argument can survive contact with an unsympathetic reviewer
depends entirely on where that reviewer chooses to press against its
assumptions.

\section{Adversarial Review}\label{sec-release-16-7}

The useful question during adjudication is not whether the engineering
team believes its own safety argument, but where an independent, hostile
reviewer can break it. An \textbf{adversarial review} is a structured
audit that treats the safety case not as a defense brief compiled to
justify deployment, but as an attack surface designed to be
systematically dismantled. When engineers assemble a safety case,
confirmation bias steers data collection toward nominal operating
envelopes where learned models succeed and monitoring tripwires trigger
on schedule. The adversarial reviewer inverts this posture, searching
for unsupported warrants, uncharacterized physical transitions, and edge
conditions that invalidate the mathematical bounds of the release
contract. If an argument cannot survive an adversary whose explicit
objective is to force a physical violation on paper, the machine will
not survive an operating environment governed by mechanics and wear.

The review evaluates every evidentiary node in the argument tree using
four targeted probe vectors. The first probe challenges empirical
sufficiency by testing whether verification datasets covered the
physical extremes of the operational envelope. For an analytical
example, consider a benchmark demonstrating ten thousand successful
manipulation cycles at an ambient temperature of \(20^\circ\text{C}\) on
clean workpieces, supporting a warrant that grip slip probability
remains below \(P_{99.9} = 0.001\). The reviewer probes the boundary by
introducing airborne oil mist and an ambient temperature of
\(40^\circ\text{C}\), conditions that reduce the friction coefficient of
the elastomer gripper pads from \(\mu = 0.80\) to \(\mu = 0.35\). If the
test suite contains zero data under degraded friction, the warrant is
unsupported. This failure mode is detected through distribution coverage
audits of physical parameters during review, but on the factory floor it
goes undetected until a dropped payload strikes the workstation. The
second probe examines domain transfer, identifying claims where
simulation data substitutes for hardware measurement. An optical
classifier may demonstrate a \(99.9\%\) detection rate for surface
contamination under synthetic rendering, but specular reflections from
micro-machining marks on physical steel surfaces can drop the hardware
accuracy to \(86\%\). This defect is detected by paired
hardware-in-the-loop cross-validation on physical workpieces, and
without those physical benchmarks, it remains unmonitored until
contaminated stock enters the cell.

The third probe attacks structural independence, probing whether the
learned model and the real-time safety enforcer share common hardware
resources that introduce single points of failure. If the
vision-language policy and the joint torque estimator execute on a
shared system-on-chip or communicate over a single PCIe bus, memory
traffic from a sudden visual re-planning cycle can introduce a
\(14\text{ ms}\) direct memory access stall into the torque tripwire
loop. An argument that models the safety enforcer as an independent
\(2\text{ ms}\) observer collapses under this shared bus contention.
Timing degradation of this nature is detected through stress testing
under synthetic worst-case bus saturation, but standard unit tests
running on an idle memory bus fail to expose it. The fourth probe
identifies unmonitored state drift, where slow physical degradation
alters plant dynamics without crossing diagnostic thresholds. Over six
months of operation, mechanical backlash in a cycloidal drive gearbox
can expand from \(0.01^\circ\) to \(0.12^\circ\). This play adds an
unmodeled \(15\text{ ms}\) phase lag during emergency deceleration,
increasing the peak impact force during a workpiece collision from
\(38\text{ N}\) to \(64\text{ N}\). Standard motor-current tripwires do
not detect backlash during free-space motion, meaning the plant drifts
outside its certified safety envelope without triggering an alarm.

Reviewers operationalize these vulnerabilities by constructing explicit
\textbf{defeaters}, which are physically realizable scenarios that
invalidate a specific argument warrant. Standard diagnostic error
handlers identify binary faults such as disconnected sensors or memory
segmentation panics, but they cannot detect coupled, sub-threshold
degradations. An adversary constructs combinations of minor anomalies
that each remain within allowable operational bands yet collectively
defeat a safety claim. Consider an optical lens with a \(10\%\)
particulate haze paired with a \(4\text{ mm}\) mechanical misalignment
from fixture bolt relaxation. In isolation, the optical sensor reports
normal visibility, the fixture monitor registers within tolerance, and
the execution thread meets its timing budget. Concurrently, the degraded
image and positional offset force the policy into an uncharacterized
latent state, extending inference latency by \(20\text{ ms}\) and
positioning the gripper off-center. The slip reflex arrives too late to
halt the falling mass. Because no individual sensor crossed an alarm
threshold, software diagnostics report normal operation while the
physical invariant fails.

The adversarial review also audits the execution integrity of the
permission path, verifying that neither model inference tasks nor
auxiliary background services can bypass, corrupt, or starve safety
enforcer assertions (Chapter~\ref{sec-enforcement}) or human authority
overrides (Chapter~\ref{sec-intervention}). As established in the system
architecture, safety assertions and human overrides must maintain
non-preemptible execution priority over learned models. An adversary
evaluates whether a runaway policy thread, an unhandled exception in an
accelerator runtime, or a high-volume telemetry stream can saturate
memory channels, lock a kernel spinlock, or exhaust power delivery
rails. If an onboard graphics accelerator draws peak transient current
that causes a voltage droop, resetting an adjacent micro-controller on
the same power plane, the nervous system loses its ability to command
actuator brakes within the required \(5\text{ ms}\) deadline. Validating
permission path independence requires physical fault injection,
including bus saturation, power rail glitching, and forced kernel
panics. If the permission path is implemented in user-space software
rather than isolated hard real-time threads and dedicated hardware
interrupt lines, the reviewer marks the link as broken.

When an adversarial attack succeeds, the safety case cannot be salvaged
by editing claim descriptions or inserting arbitrary safety margins. The
protocol requires decomposing the broken warrant into its underlying
systems failure, whether that represents an uncharacterized physical
boundary, an inadequate verification dataset, or an unmodeled hardware
latency. Remediation routes the defect directly back to the owning
subsystem. If an adversary proves that the motor-current observer fails
under mechanical backlash, the mechanical specification must introduce
hard wear limits and runtime backlash estimation. If shared bus
contention compromises enforcer latency, the compute architecture must
enforce physical memory partitioning and real-time scheduling
guarantees. If perception transfer fails across lighting variations, the
verification pipeline must collect hardware-in-the-loop datasets under
physical extremes. The safety case returns to the engineering team with
broken warrants explicitly identified, requiring updated hardware
boundaries, re-derived stopping distances, and renewed empirical
verification before the machine can be submitted for a new release
verdict.

\section{How Release Arguments Go Wrong}\label{sec-release-16-8}

Arguments that were accepted and should not have been share a small
number of patterns. When a flawed safety case passes adjudication and
permits a hazardous machine to operate, the defect is almost never an
unclassifiable theoretical anomaly. It traces back to one of four
recurring structural pathologies in the claim-argument-evidence tree:
the simulation-to-reality leap, the proxy metric fallacy, the
unmonitored condition, and the organizational normalization of deviance.
Each pathology introduces a specific flaw into the argument graph,
creating an illusion of empirical verification while leaving the
physical machine unprotected against mechanics, timing jitter, and wear.

The \emph{simulation-to-reality leap} occurs when high-volume synthetic
or simulated pass rates are substituted for physical verification
without bounding unmodeled contact dynamics or sensor noise floors. An
engineering team may run \(1\times 10^{6}\) collision-free simulated
episodes with a \(99.9\%\) task completion metric in a rigid-body
physics engine. Yet virtual physics engines calculate contact forces
using idealized Coulomb friction, rigid geometries, instantaneous
surface normals, and noise-free coordinate frames. On physical hardware,
mechanical interaction involves elastomeric deformation, lubricant film
breakdown, micro-slip, optical sensor noise, and thermal expansion in
actuator housings. If an argument treats simulation rollouts as direct
empirical evidence for a physical force or trajectory bound without
measuring the unmodeled dynamics on hardware test benches, the warrant
is invalid. When the deployed machine encounters physical contact,
unmodeled compliance and sensor latencies cause state trajectories to
diverge exponentially, converting a simulated gentle touch into a
damaging mechanical impact.

The \emph{proxy metric fallacy} occurs when statistical loss functions,
mean average precision, or cumulative reward are presented as evidence
for physical safety, ignoring the fat-tailed distribution of hazardous
physical states. In statistical machine learning, an object detector
that achieves a mean average precision of \(98.5\%\) on a validation
benchmark is considered high-performing. But physical damage does not
depend on average performance. An autonomous platform moving mass under
momentum is governed by worst-case tail states (\(P_{99.99}\)). In a
software application, an isolated misclassification produces an
incorrect bounding box for a single frame. In an embodied vehicle
traveling at \(19.2\text{ m/s}\), a \(1.2\text{ s}\) classification
oscillation delays brake actuation across \(23.0\text{ m}\) of
displacement. By evaluating the system against statistical averages
rather than physical worst-case limits, the argument conceals severe
physical hazards behind high average accuracy metrics.

The \emph{unmonitored condition} arises when an adjudication issues a
conditional release verdict without deploying continuous,
hardware-enforced runtime monitors to detect when operational boundaries
are exceeded. A safety case may stipulate that a robotic workstation
operates safely provided that ambient temperature remains below
\(35^\circ\text{C}\), workpiece mass does not exceed \(2.0\text{ kg}\),
or a human operator remains attentive during collaborative tasks. If the
physical machine lacks a hardware thermistor, a calibrated load cell
check, or an automated gaze monitor with direct authority to interrupt
motor power, the condition is an unmonitored assumption. When
environmental conditions drift or an operator looks away, the
unmonitored premise fails silently. Without a deterministic
nervous-system reflex to detect the boundary crossing and force a safe
stop, the machine continues actuating under the authority of an
invalidated verdict.

The \emph{normalization of deviance} occurs when safety margins are
degraded over time because non-catastrophic edge-case anomalies are
observed without triggering argument invalidation or formal re-review
(\citeproc{ref-leveson2011engineering}{Leveson 2011};
\citeproc{ref-vaughan1996challenger}{Vaughan 1996}). In deployed fleets,
machines routinely encounter minor boundary excursions, such as
transient \(40\text{ ms}\) communication jitter spikes on an internal
bus, momentary false-positive sensor drops, or sub-threshold tracking
errors during high-speed maneuvers. Because these anomalies did not
immediately cause a collision or mechanical fracture, engineering teams
gradually treat the anomalies as acceptable operational noise rather
than evidence of broken warrants. Software thresholds are widened,
safety margins are trimmed to improve task throughput, and warning
alarms are suppressed. Over months of field operation, the gap between
the verified safety envelope and actual machine behavior widens until a
routine operational disturbance exhausts the remaining physical margin.

\protect\phantomsection\label{callout-incidentux2a-16.6}
\begin{fbxSimple}{callout-incident}{Landmark Engineering Incident:}{Space Shuttle Challenger O-Ring Cold Launch ODD Envelope Violation}
\phantomsection\label{callout-incident*-16.6}

\emph{System:} NASA Space Shuttle Solid Rocket Booster (SRB) field
joints utilizing dual Viton-A fluoroelastomer O-rings with zinc chromate
putty thermal insulation
(\citeproc{ref-rogers1986challenger}{Presidential Commission on the
Space Shuttle Challenger Accident 1986};
\citeproc{ref-vaughan1996challenger}{Vaughan 1996};
\citeproc{ref-leveson2011engineering}{Leveson 2011}).

\emph{Incident:} Catastrophic in-flight structural breakup of STS-51-L
at \(T+73.1\text{ s}\) (altitude \(14.6\text{ km}\), Mach 1.92, dynamic
pressure \(34.5\text{ kPa}\)), resulting in the total loss of the
orbiter vehicle and all seven crew members.

\emph{Root Cause \& Physics of Failure:} Ambient launch temperature at
Pad 39B was \(31^\circ\text{F}\) (\(-0.5^\circ\text{C}\)), plunging the
SRB field joint temperature far below the certified minimum
qualification baseline of \(53^\circ\text{F}\) (\(12^\circ\text{C}\))
and near the glass transition temperature of Viton
(\(T_g \approx -1^\circ\text{C}\) to \(+4^\circ\text{C}\)). Cold soak
increased elastomer resilience recovery time from milliseconds to over
\(2.0\text{ s}\). Upon ignition, rapid motor pressurization to
\(6.2\text{ MPa}\) (\(900\text{ psi}\)) caused instantaneous joint
casing expansion (``hole opening'' gap deflection of
\(0.76\text{ mm}\)). Because the frozen primary and secondary O-rings
failed to expand and seat in time, \(3100^\circ\text{C}\) combustion gas
escaped through the joint seal gap (``blowby''). Hot gas eroded
\(100\%\) of the primary O-ring and scorched the secondary O-ring,
creating a continuous blowtorch flame. The plume impinged directly on
the lower SRB attachment strut and the external liquid hydrogen tank
(\(\mathrm{LH}_2\)), causing structural strut detachment, uncontrollable
thrust vector yaw, aerodynamic vehicle breakup under \(1.8\text{ GJ}\)
of uncontrolled kinetic energy, and external tank structural
rupture.\sidenote{\footnotesize The Rogers Commission
  (\citeproc{ref-rogers1986challenger}{Presidential Commission on the
  Space Shuttle Challenger Accident 1986}) documented that management
  treated O-ring blowby as acceptable risk, illustrating Vaughan's
  \emph{normalization of deviance}
  (\citeproc{ref-vaughan1996challenger}{Vaughan 1996}).}

\emph{Systems \& Adjudication Failure:} Shifting the burden of proof
from \emph{proving safety within evidence} to \emph{proving impending
disaster without prior failure data} under management launch pressure
(organizational normalization of deviance). Prior flights exhibiting
sub-critical O-ring erosion and blowby were treated as ``acceptable
operational variance'' rather than active falsifications of the joint
containment warrant. Release authority operated outside the qualified
ODD temperature envelope without empirical hardware qualification data
at \(31^\circ\text{F}\).

\emph{Systems Resolution:} Enforce strict mathematical and physical ODD
boundary tripwires; treat unmonitored or unverified operational domains
as hard refusal states; mandate independent technical authority for
launch adjudication; and redesign mechanical field joints with capture
features (heaters and interference fit) that guarantee gas containment
regardless of environmental temperature.

\emph{Systems Takeaway:} In physical systems, operating outside an
empirically certified operational design domain without hardware
evidence is not an act of calculated risk---it is operating in an
unmonitored unknown state where physical laws, material phase changes,
and kinetic energies will ruthlessly enforce reality.

\end{fbxSimple}

The physical consequences of compound perception and supervisory
pathologies are further demonstrated in the autopsy of the 2018 Uber
Advanced Technologies Group (ATG) autonomous vehicle collision in Tempe,
Arizona, investigated by the National Transportation Safety Board
(\citeproc{ref-ntsb2019uber}{National Transportation Safety Board
2019}). The automated driving system detected an unmodeled pedestrian
wheeling a bicycle across the roadway \(5.6\text{ s}\) before impact at
a range of approximately \(82\text{ m}\), while traveling at
\(19.2\text{ m/s}\) (\(43\text{ mph}\)). Because the perception model
had not been trained on pedestrians outside crosswalks with unmounted
bicycles, the tracker toggled classifications between vehicle, bicycle,
and unknown object. Each classification switch reset the tracking
history, re-initializing trajectory predictions and preventing the
system from identifying the crossing path. To prevent false-positive
braking maneuvers that degraded ride comfort, the software architecture
had inserted a \(1.2\text{ s}\) suppression delay between fault
detection and emergency braking, while engineers had disabled the
vehicle manufacturer's factory autonomous emergency braking system to
avoid control contention. The release argument relied on the standing
condition that a human safety driver would provide continuous visual
monitoring. Yet no hardware gaze-tracking monitor was deployed to verify
driver attentiveness, and safety driver supervision procedures had been
relaxed over time. The safety driver was looking away, autonomous
braking was suppressed, and manual braking was initiated only
\(0.2\text{ s}\) before impact, resulting in a fatal collision at
\(17.4\text{ m/s}\) (\(39\text{ mph}\)).\sidenote{\footnotesize NTSB Report HAR-19/03
  (\citeproc{ref-ntsb2019uber}{National Transportation Safety Board
  2019}) revealed that disabling the factory autonomous emergency
  braking system eliminated the deterministic physical safety enforcer,
  relying on unmonitored human supervision.}

\protect\phantomsection\label{callout-incidentux2a-16.7}
\begin{fbxSimple}{callout-incident}{Landmark Engineering Incident:}{Autonomous Vehicle Detection Suppression and Handover Collision}
\phantomsection\label{callout-incident*-16.7}

\emph{System:} Autonomous test vehicle with rooftop lidar, radar,
forward cameras, learned perception and tracking stack, and retrofitted
drive-by-wire actuation.

\emph{Incident:} Fatal collision with a pedestrian pushing a bicycle at
\(17.4\text{ m/s}\) following perception classification switching,
suppressed autonomous braking, and delayed human takeover.

\emph{Telemetry Snapshot:} At \(t = -5.6\text{ s}\), sensors detect
obstacle at range \(82\text{ m}\); from \(t = -5.6\text{ s}\) to
\(-1.2\text{ s}\), classification switches between vehicle, bicycle, and
unknown, resetting track history; at \(t = -1.2\text{ s}\), system
determines collision is imminent but suppresses braking for
\(1.2\text{ s}\); at \(t = -0.2\text{ s}\), operator applies manual
brakes; at \(t = 0.0\text{ s}\), impact occurs at \(17.4\text{ m/s}\).

\emph{Root Cause:} Object classification toggling repeatedly reset
tracking state vectors, while a \(1.2\text{ s}\) suppression delay
prevented autonomous deceleration and human supervisory attention was
not continuously verified.

\emph{Architectural Flaw:} Disabling the factory autonomous emergency
braking system eliminated the independent physical enforcer, while
release authority was delegated to an unmonitored human supervisory
condition.

\emph{Provenance:} National Transportation Safety Board (NTSB) statutory
highway accident report NTSB/HAR-19/03
(\citeproc{ref-ntsb2019uber}{National Transportation Safety Board
2019}).

\end{fbxSimple}

An adversarial review auditing this safety case would have dismantled
the argument before deployment. Probing the permission path would have
revealed that disabling the factory emergency braking system eliminated
the independent enforcer, leaving moving mass without a deterministic
physical safeguard. Probing the standing conditions would have exposed
that relying on human intervention without continuous, hardware-enforced
gaze tracking transformed an operational requirement into an unmonitored
assumption. Probing the proxy metrics would have shown that suppressing
braking to optimize ride smoothness sacrificed stopping clearance.
Preventing these structural failures requires that every release verdict
be compiled into an immutable engineering artifact that links
operational claims directly to verified physical bounds and runtime
monitoring contracts.

\section{The Release Manifest}\label{sec-release-16-9}

A release verdict is only as durable as the immutable artifact that
records what was decided, under what physical conditions, and when that
authority expires. The machine notebook tracks components, software
artifacts, and execution traces under the provenance and versioning
rules established in Section~\ref{sec-nervous-4-10}. This section adds
the final schema layer to that ledger, incorporating twelve explicit
structural fields (Table~\ref{tbl-16-release-manifest-schema}),
comprising the unique claim identifier, the argument tree topology hash,
the verified evidence pointer array, the cryptographic hardware and
software manifests, the signed verdict classification, the active
standing conditions, the telemetry monitoring contracts, and the
expiration schedule. These additions transform the notebook from a
retrospective log of development tests into an active operational
contract governing physical execution.

The compiled result is the \textbf{cryptographic release record} (or
\textbf{release manifest}), the formal, immutable engineering artifact
that encapsulates the operational claim, the complete
claim-argument-evidence tree, the verified evidence traces, the exact
software and model hashes, the hardware component revisions, and the
signed verdict of the lead systems adjudicator.\sidenote{\footnotesize Deployments
  enforce release integrity via hardware Root of Trust (TPMs, OTP
  eFuses). Dual-bank flash boots new policies; if standing condition
  telemetry fails, hardware rolls back to the certified partition.}

As illustrated in Figure~\ref{fig-16-release-record}, stored as a
cryptographically signed manifest within non-volatile storage, the
release record forms a tamper-evident chain of custody between
verification data and physical execution. At system boot, the nervous
system bootloader verifies the cryptographic signature of the release
record against a public key burned into hardware fuses. If a policy
neural network weight file differs by a single bit from the hash
recorded during adjudication, or if an FPGA bitstream revision does not
match the certified nervous-system safety enforcer, the boot sequence
halts, the motor inverters remain unpowered, and the machine refuses to
enter an operational state.

\begin{figure}[t]

\centering{

\includegraphics[width=1\linewidth,height=\textheight,keepaspectratio]{index_files/mediabag//Users/VJ/GitHub/MLSysBook-integrate-physical/physical/book/chapters/16-release/figures/fig16_release_record.pdf}

}

\caption{\label{fig-16-release-record}\textbf{Architecture of the
Cryptographic Release Record and Runtime Invalidation Pipeline.} The
release record compiles the complete systems verification basis into an
immutable engineering manifest (left), recording unique claim
identifiers, GSN topology hashes, verified evidence pointers, and
cryptographic hashes for policy weights, FPGA bitstreams, and RTOS
kernels, signed by the lead systems adjudicator. At boot, the hardware
Root of Trust and silicon OTP eFuses (center) verify manifest integrity
before energizing motor inverters. During execution, the hard real-time
nervous system monitors standing conditions across the envelope
trichotomy (\emph{known-true}, \emph{known-false}, and \emph{unknown},
where epistemic gaps are treated strictly as outside the envelope). Any
violation of software hashes, mechanical wear limits, or environmental
ODD boundaries trips automatic invalidation (right), latching hardware
interlocks and revoking operating authority until formal
re-adjudication. Attribution: Textbook production team (CC BY-NC-ND
4.0).}

\end{figure}%

The formal procedure for attesting hardware silicon state, verifying
cryptographic manifest signatures, validating firmware image hashes, and
executing pre-flight standing condition audits prior to energizing motor
power stages is formalized below:

\subsection{Cryptographic Release Manifest Verification
Protocol}\label{cryptographic-release-manifest-verification-protocol}

\textbf{Substrate \& Invariants:}

\begin{itemize}
\tightlist
\item
  \emph{Execution Substrate:} Secure Hardware Root of Trust (RoT)
  cryptographic coprocessor / isolated safety core
  (\(f = 400\text{ MHz}\), write-protected OTP eFuse bank, tamper-detect
  enclosure).
\item
  \emph{Memory Invariant:} Read-only execution memory mapping
  (\(\text{RX}\) permissions); manifest verification completed entirely
  in internal SRAM before motor drive DC bus pre-charge relays are
  commanded.
\item
  \emph{Safety Interlock Authority:} Hardware-latching crowbar and
  inverter gate pull-down lines; de-asserting hardware permit requires
  valid cryptographic signature, uncorrupted hashes, active standing
  condition contracts, and verified telemetry drift margins.
\item
  \emph{Signature \& Hash Invariants:}
  \(\text{VerifyEd25519}(\mathbf{K}_{\mathrm{pub\_adj}}, \text{SHA384}(\mathcal{M}_{\mathrm{rel}}), \boldsymbol{\sigma}_{\mathrm{adj}}) = \text{VALID}\),
  and
  \(\text{SHA384}(\mathbf{W}_{\mathrm{neural}}) == \mathbf{H}_{\mathrm{policy}} \land \text{SHA384}(\mathbf{B}_{\mathrm{FPGA}}) == \mathbf{H}_{\mathrm{FPGA}} \land \text{SHA384}(\mathbf{K}_{\mathrm{kernel}}) == \mathbf{H}_{\mathrm{RTOS}}\).
\end{itemize}

\textbf{Execution Protocol:}

\begin{enumerate}
\def\labelenumi{\arabic{enumi}.}
\item
  \textbf{Silicon Attestation \& Key Recovery:} Read immutable
  adjudicator public key \(\mathbf{K}_{\mathrm{pub\_adj}}\) and
  revocation mask \(\text{OTP}_{\mathrm{rev}}\) from hardware eFuses.
\item
  \textbf{Cryptographic Manifest Verification:} Compute cryptographic
  digest
  \(\mathbf{d}_{\mathcal{M}} \leftarrow \text{SHA-384}(\mathcal{M}_{\mathrm{rel}})\).
  Verify digital signature
  \(s \leftarrow \text{Ed25519\_Verify}(\mathbf{K}_{\mathrm{pub\_adj}}, \mathbf{d}_{\mathcal{M}}, \boldsymbol{\sigma}_{\mathrm{adj}})\).
  If signature fails, log signature forgery attempt, assert
  \(\text{Permit}_{\mathrm{gate}} \leftarrow 0\), and return
  \texttt{REFUSED\_LATCHED}.
\item
  \textbf{Binary Image \& Weight Hash Verification:} Compute active
  image digests
  \(\mathbf{h}_{\mathbf{W}} \leftarrow \text{SHA-384}(\mathbf{W}_{\mathrm{neural}})\),
  \(\mathbf{h}_{\mathbf{B}} \leftarrow \text{SHA-384}(\mathbf{B}_{\mathrm{FPGA}})\),
  and
  \(\mathbf{h}_{\mathbf{K}} \leftarrow \text{SHA-384}(\mathbf{K}_{\mathrm{kernel}})\).
  If any digest does not match the manifest hashes, assert hardware
  rollback to Bank B golden image, latch
  \(\text{Permit}_{\mathrm{gate}} \leftarrow 0\), and return
  \texttt{REFUSED\_LATCHED}.
\item
  \textbf{Temporal \& Mechanical Lifecycle Audit:} Query tamper-proof
  real-time clock \(t_{\mathrm{RTC}}\) and cycle registers
  \(\tau_{\mathrm{hw}}\). If
  \((t_{\mathrm{RTC}} > t_{\mathrm{expiry}}) \lor (T_{\mathrm{operating\_hours}} > \mathcal{M}_{\mathrm{rel}}.T_{\max}) \lor (N_{\mathrm{cycles}} > \mathcal{M}_{\mathrm{rel}}.N_{\max})\),
  raise manifest expiration fault, latch
  \(\text{Permit}_{\mathrm{gate}} \leftarrow 0\), and return
  \texttt{REFUSED\_LATCHED}.
\item
  \textbf{Sensor Calibration \& Hardware Trace Audit:} Validate
  cryptographic timestamps on load cell dynamometer traces and
  optical-tactile calibration certificates
  \(\mathcal{L}_{\mathrm{calib}}\). If calibration age exceeds
  \(T_{\mathrm{calib\_interval}}\), restrict operating authority by
  downgrading to \texttt{CONDITIONAL\_ENFORCED} and clamping Cartesian
  speed \(v_{\max} \leftarrow 0.20\text{ m/s}\).
\item
  \textbf{Standing Condition \& Envelope Pre-Flight Check:} Sample
  real-time environmental sensors (ambient temperature, supply voltage,
  bus jitter, enclosure pressure). For each standing condition premise
  \(p_i \in \mathcal{P}_{\mathrm{stand}}\), evaluate state
  \(S_i \leftarrow \text{EvaluatePremise}(p_i)\). If
  \(S_i == \text{KNOWN\_FALSE}\) or \(S_i == \text{UNKNOWN}\), log
  unverified premise, latch
  \(\text{Permit}_{\mathrm{gate}} \leftarrow 0\), and return
  \texttt{REFUSED\_LATCHED}.
\item
  \textbf{Operating Verdict Authorization \& Hardware Interlock
  Release:}

  \begin{itemize}
  \tightlist
  \item
    If manifest is marked \texttt{OPERATE\_UNCONDITIONAL} and all checks
    pass: configure safety enforcer registers to full envelope
    (\(v_{\max} \leftarrow 1.8\text{ m/s}, F_{\mathrm{trip}} \leftarrow 50\text{ N}, \Delta t_{\mathrm{watchdog}} \leftarrow 2.0\text{ ms}\)),
    assert \(\text{Permit}_{\mathrm{gate}} \leftarrow 1\) (energize DC
    power stages), and return \texttt{AUTHORIZED}.
  \item
    Else if manifest is marked \texttt{OPERATE\_WITH\_CONDITIONS}: apply
    hardware derating
    (\(v_{\max} \leftarrow 0.40\text{ m/s}, F_{\mathrm{trip}} \leftarrow 18\text{ N}, \Delta t_{\mathrm{watchdog}} \leftarrow 1.0\text{ ms}\)),
    engage auxiliary interlock monitors, assert
    \(\text{Permit}_{\mathrm{gate}} \leftarrow 1\), and return
    \texttt{CONDITIONAL\_ENFORCED}.
  \item
    Else: assert \(\text{Permit}_{\mathrm{gate}} \leftarrow 0\) and
    return \texttt{REFUSED\_LATCHED}.
  \end{itemize}
\end{enumerate}

\begin{longtable}[]{@{}
  >{\raggedright\arraybackslash}p{(\linewidth - 8\tabcolsep) * \real{0.2000}}
  >{\raggedright\arraybackslash}p{(\linewidth - 8\tabcolsep) * \real{0.2000}}
  >{\raggedright\arraybackslash}p{(\linewidth - 8\tabcolsep) * \real{0.2000}}
  >{\raggedright\arraybackslash}p{(\linewidth - 8\tabcolsep) * \real{0.2000}}
  >{\raggedright\arraybackslash}p{(\linewidth - 8\tabcolsep) * \real{0.2000}}@{}}
\caption{\textbf{Abstract Cryptographic Release Record Schema and
Hardware Storage Substrates.} The twelve mandatory structural fields
comprising the cryptographic release manifest. Each record binds
physical AI software binaries, nervous-system FPGA bitstreams, empirical
sensor calibration traces, and runtime standing condition contracts into
an immutable manifest verified by hardware Root of Trust prior to
actuator power stage
energization.}\label{tbl-16-release-manifest-schema}\tabularnewline
\toprule\noalign{}
\begin{minipage}[b]{\linewidth}\raggedright
Field Identifier
\end{minipage} & \begin{minipage}[b]{\linewidth}\raggedright
Data / Representation Type
\end{minipage} & \begin{minipage}[b]{\linewidth}\raggedright
Physical SI Unit / Invariant
\end{minipage} & \begin{minipage}[b]{\linewidth}\raggedright
Hardware Storage Substrate
\end{minipage} & \begin{minipage}[b]{\linewidth}\raggedright
Architectural Verification \& Trip Invariant
\end{minipage} \\
\midrule\noalign{}
\endfirsthead
\toprule\noalign{}
\begin{minipage}[b]{\linewidth}\raggedright
Field Identifier
\end{minipage} & \begin{minipage}[b]{\linewidth}\raggedright
Data / Representation Type
\end{minipage} & \begin{minipage}[b]{\linewidth}\raggedright
Physical SI Unit / Invariant
\end{minipage} & \begin{minipage}[b]{\linewidth}\raggedright
Hardware Storage Substrate
\end{minipage} & \begin{minipage}[b]{\linewidth}\raggedright
Architectural Verification \& Trip Invariant
\end{minipage} \\
\midrule\noalign{}
\endhead
\bottomrule\noalign{}
\endlastfoot
\texttt{claim\_manifest\_id} & UUID-128 / Hex Digest & Top-level safety
claim binding & Secure Flash (Bank A/B) & Cryptographically linked to
GSN root node hash \\
\texttt{gsn\_topology\_digest} & SHA-384 Hash (48 bytes) & Complete CAE
argument graph & Secure Flash / TPM NVRAM & Verified against lead
adjudicator digital signature \\
\texttt{policy\_weights\_sha384} & SHA-384 Hash (48 bytes) & Neural
network model weights & Silicon OTP Fuse / TPM PCR & Mismatch halts
inference engine boot sequence \\
\texttt{enforcer\_bitstream\_hash} & SHA-384 Hash (48 bytes) & FPGA
real-time safety logic & Silicon OTP Fuse / Secure Boot & Mismatch
prevents power inverter relay engagement \\
\texttt{rtos\_kernel\_hash} & SHA-384 Hash (48 bytes) & Microkernel
image \& drivers & Secure Flash (A/B) & Verified by hardware Root of
Trust at power-on \\
\texttt{calib\_trace\_pointers} & URI Array + HMAC-SHA256 & Dynamometer
/ Load cell certs & Machine Ledger / Secure NVRAM & Calibration age
\(> 500\text{ h}\) forces conditional mode \\
\texttt{verdict\_classification} & Enum (\texttt{UNCOND}, \texttt{COND},
\texttt{REFUSE}) & Granted operational authority & TPM PCR Register 14 &
Read-only register gates DC bus contactors \\
\texttt{standing\_conditions} & Struct Array
\(\langle \text{Premise}, \text{Rate}, \text{Trip} \rangle\) & Physical
sensor trip thresholds & FPGA Block RAM / Real-Time MCU & Evaluated at
\(1000\text{ Hz}\); unknown state asserts safe stop \\
\texttt{telemetry\_drift\_budget} & Covariance Bound
\(\boldsymbol{\Sigma}_{\max} \in \mathbb{R}^{d\times d}\) & Tracking \&
timing error drift & Circular NVRAM Buffer & \(P_{99}\) drift
\(> 3\sigma\) triggers maintenance shutdown \\
\texttt{verdict\_expiry\_timestamp} & UTC Timestamp (ISO 8601) & Seconds
(\(s \le t_{\mathrm{expiry}}\)) & Tamper-Proof RTC / Crypto-Timer & RTC
\(> t_{\mathrm{expiry}}\) voids operating permit automatically \\
\texttt{max\_cycle\_budget} & \texttt{uint64\_t} scalar & Actuation
cycles (\(N \le N_{\max}\)) & Monotonic Hardware Counter & Counter
overflow revokes motion authority \\
\texttt{adjudicator\_signature} & Ed25519 Signature (64 bytes) &
Personal legal \& technical attestation & OTP / TPM Attestation Block &
Verified against root CA key burned in silicon eFuses \\
\end{longtable}

\protect\phantomsection\label{callout-ruleux2a-16.8}
\begin{fbxSimple}{callout-rule}{Architect’s Rule of Thumb:}{The Silicon Root-of-Trust Invariant}
\phantomsection\label{callout-rule*-16.8}

A safety case stored in mutable user-space flash is merely an
operational suggestion; operating authority must be cryptographically
locked into write-protected silicon OTP fuses and verified by an
isolated hardware Root of Trust before motor power stages are energized.

\end{fbxSimple}

Every leaf node in the argument tree must carry verifiable metadata
linking back to calibrated test bench logs, simulation run identifiers,
and hardware serial numbers. An assertion that a low-level reflex loop
executes with a worst-case latency of \(P_{99.99} \le 1.8\text{ ms}\)
cannot reference an informal summary table. It must record the unique
execution trace identifier from a real-time logic analyzer, the serial
numbers of the timing probes, the calibration certificate of the
oscilloscope, and the test script git commit that generated the
synthetic bus saturation. Similarly, claims bounding peak contact force
below \(50\text{ N}\) must link to raw load cell time-series data,
dynamometer calibration records, and the exact software environment
seeds used in hardware-in-the-loop stress testing. This provenance
guarantees end-to-end reproducibility, allowing an adversarial auditor
to replay any verification trial or trace an unexpected field behavior
back to the exact test bench configuration that certified the subsystem.

Because a release verdict is valid only for the exact physical and
computational plant evaluated during adjudication, the release record
enumerates explicit invalidation triggers that instantly void operating
authority. Software triggers invalidate the release upon any retraining
or fine-tuning of neural policy weights, any modification to real-time
thread priorities, or any patch applied to lower-level sensor drivers.
Hardware triggers fire when any physical component is replaced with an
unverified part revision, when a sensor vendor alters an internal
optical filter, or when mechanical wear exceeds verified tolerances,
such as cycloidal drive backlash widening past \(0.10^\circ\) or
friction brake lining thickness wearing below \(2.0\text{ mm}\).
Operational triggers execute whenever environmental parameters exceed
the certified operational design domain, such as ambient temperature
climbing above \(40^\circ\text{C}\), line supply voltage drooping below
\(21.6\text{ V}\) (a \(10\%\) drop on a nominal \(24\text{ V}\) bus), or
workpiece mass exceeding \(1.5\text{ kg}\). The occurrence of any
invalidation trigger immediately revokes the operating verdict, forcing
the nervous system to inhibit autonomous motion until a formal
re-adjudication is completed.

Operating authority is maintained through a continuous re-decision
protocol that evaluates streaming field telemetry against the baseline
distributions captured in the release record. The nervous system
continuously logs high-frequency observables, including joint torque
estimation residuals, control cycle execution jitter, and enforcer trip
counts, into non-volatile circular buffers. During scheduled maintenance
intervals, such as every \(500\text{ operating hours}\) or
\(250{,}000\text{ manipulation cycles}\), automated analysis routines
execute an empirical safety audit, comparing field telemetry
distributions against the original verification baseline. If long-term
drift analysis reveals that optical tactile noise variance has increased
by \(15\%\), or if the \(P_{99}\) braking response time drifts from
\(18\text{ ms}\) to \(23\text{ ms}\) due to contact surface polishing,
the protocol triggers an automatic transition. The machine either
downgrades its authority to a restricted conditional mode (for instance,
capping maximum tool velocity at \(0.20\text{ m/s}\)) or executes a
controlled operational shutdown, preventing physical degradation from
quietly consuming the design's safety margins.

The release record closes the systems engineering framework developed
across Chapter~\ref{sec-boundary} through Chapter~\ref{sec-release} and
establishes the operational baseline for long-term fleet monitoring,
periodic drift audits, and lifecycle maintenance
(Chapter~\ref{sec-frontier}). In Chapter~\ref{sec-boundary}, the causal
boundary established the divide between software commands that can be
revoked and physical actions governed irreversibly by mass, force, and
momentum. Across the subsequent layers, the body was constrained by
mechanical dynamics and thermal ratings (Chapter~\ref{sec-body}),
sensing and estimation were grounded in hardware noise floors and
physical provenance (Chapter~\ref{sec-nervous}, Chapter~\ref{sec-data}),
the learned brain was bounded by inference budgets and fault
distributions (Chapter~\ref{sec-perception}, Chapter~\ref{sec-memory},
Chapter~\ref{sec-intent}, Chapter~\ref{sec-planning}), and the nervous
system was engineered to enforce real-time reflexes with deterministic
priority over learned policies (Chapter~\ref{sec-enforcement},
Chapter~\ref{sec-placement}, Chapter~\ref{sec-intervention},
Chapter~\ref{sec-verification}). The release record integrates these
physical limits, timing budgets, and architectural safeguards into an
auditable, enforceable contract. By binding learned behavior to
verifiable physical bounds and making operating authority conditional on
continuous empirical proof, the machine operates not on ungrounded
optimism, but on an accountable, reproducible foundation that withstands
the physical world.

\section{Summary}\label{sec-release-16-10}

This text established the engineering and governance discipline required
to transition a physical AI system from a development artifact into a
production-certified machine:

\begin{itemize}
\tightlist
\item
  \textbf{The Release Decision:} Operating authority is not earned
  through ungrounded statistical optimism or benchmark scores, but
  through a defensible, evidence-backed safety case that bounds physical
  risk within a defined Operational Design Domain (ODD).
\item
  \textbf{Claim-Argument-Evidence (CAE) Formalism:} Decomposing
  high-level safety claims into verifiable sub-claims, inferential
  warrants, and calibrated empirical test traces using Goal Structuring
  Notation (GSN).
\item
  \textbf{The Verdict Trichotomy:} Classifying operational authority
  into \emph{Authorized}, \emph{Conditional}, or \emph{Refused},
  anchored by explicit standing conditions and hardware-enforced
  invalidation triggers.
\item
  \textbf{The Cryptographic Release Record:} Binding policy weights,
  FPGA bitstreams, real-time safety enforcers, and telemetry monitoring
  contracts into an immutable manifest verified by hardware Root of
  Trust.
\end{itemize}

\protect\phantomsection\label{callout-takeawaysux2a-16.9}
\begin{fbxSimple}{callout-takeaways}{Key Chapter Takeaways}{}
\phantomsection\label{callout-takeaways*-16.9}

\begin{itemize}
\tightlist
\item
  \textbf{Safety is Never an Inherent Global Attribute (ODD Bounds):} A
  physical AI system is never ``generally safe.'' Operational authority
  is strictly bounded within an explicit Operational Design Domain (ODD)
  defining physical limits on mass, friction, velocity, and temperature.
\item
  \textbf{Claims Without Calibrated Hardware Traces Are Opinions:}
  High-level safety arguments in Goal Structuring Notation (GSN) require
  auditable leaf evidence linking directly to calibrated dynamometer
  traces, logic analyzer captures, and test scripts under
  hardware-in-the-loop stress.
\item
  \textbf{Explicit Invalidation Triggers Instantly Void Authority:} Any
  unauthorized modification (fine-tuning policy weights, updating
  firmware drivers, or mechanical wear exceeding backlash thresholds
  \(\Delta \theta > 0.10^\circ\)) triggers immediate cryptographic
  invalidation, forcing the nervous system to inhibit motor
  energization.
\item
  \textbf{The Verdict Trichotomy (Authorized, Conditional, Refused):}
  Release authority is not binary. Operating permissions must support
  conditional and refused states, backed by automated standing
  conditions evaluated by hardware supervisors at \(1000\text{ Hz}\).
\item
  \textbf{Operating Authority is a Decaying Cryptographic Lease:}
  Authority is an active lease verified by hardware Root of Trust (TPM /
  OTP eFuses) and maintained through continuous field telemetry audits,
  automatically expiring if drift envelopes (\(>3\sigma\)) or cycle
  counts (\(N > N_{\max}\)) are exceeded.
\end{itemize}

\end{fbxSimple}

\bookmarksetup{startatroot}

\chapter{\texorpdfstring{The Epistemic
Frontier\chaptersubtitle{Open Frontiers in Physical AI}}{The Epistemic Frontier}}\label{sec-frontier}

\begin{figure}[H]

\centering{

\includegraphics[width=1\linewidth,height=\textheight,keepaspectratio]{index_files/mediabag//Users/VJ/GitHub/MLSysBook-integrate-physical/physical/book/chapters/17-frontier/figures/fig_locator.pdf}

}

\caption{\label{fig-locator}\textbf{Systems Stage Focus: The Epistemic
Frontier.} Mapping cataloged physical failure modes against runtime
detectors across the machine lifecycle, isolating the detectorless gaps
where empirical evidence ends and architectural containment must
govern.}

\end{figure}%

\emph{What physical failure modes remain invisible to modern runtime
detectors, and how must an engineer govern what evidence cannot settle?}

A systems engineering framework that teaches practitioners to demand
rigorous evidence must be equally transparent about the boundaries of
its own methodology. No finite amount of empirical testing can certify
open-world safety against catastrophic rare events, and several critical
failure modes in physical systems lack practical real-time detectors.
The frontier of physical AI is not a speculative list of future
algorithms, but the formal register where unproven premises and
detectorless gaps are openly documented. Recognizing what current
evidence cannot settle---and enforcing the physical authority limits and
operational restrictions needed to render the unknown harmless---is the
definitive mark of professional engineering discipline.

\newpage{}

\protect\phantomsection\label{callout-objectiveux2a-17.1}
\begin{fbxSimple}{callout-objective}{Learning Objectives}{}
\phantomsection\label{callout-objective*-17.1}

After studying this chapter, you will be able to:

\begin{enumerate}
\def\labelenumi{\arabic{enumi}.}
\tightlist
\item
  Identify, for a machine of your own, which claims your evidence cannot
  currently support.
\item
  Name a failure mode from this book for which no practical detector
  exists.
\item
  State what advance would close a specific gap, and how you would know
  it had.
\item
  Apply the four properties to a machine this book never described.
\end{enumerate}

\emph{Core chapter takeaway.} A complete record still leaves claims
unsupported. Several failure modes named in this book have no practical
detector, and some claims cannot be supported by any evidence obtainable
today. Naming those is the register of the discipline, not a weakness in
it.

\end{fbxSimple}

\section{Epistemic Limits}\label{sec-frontier-17-1}

\lettrine{A}{ safety case}\lettrine{a}{ssembled} through rigorous
systems engineering concludes with an explicit verdict. When an
engineering review accepts that verdict, it does not certify that the
physical machine is unconditionally safe. The verdict authorizes a
specific hardware revision running a designated software build,
restricted to a defined operating envelope, validated for an explicit
operational duration such as \(2000\text{ h}\), and conditioned on the
continuous functioning of specified runtime monitors. As illustrated in
the systems stage locator (Figure~\ref{fig-locator}), a rigorous
engineering framework maps every cataloged physical failure mode against
available runtime detectors across the machine's architecture. Where
verified detectors and deterministic reflexes exist, the nervous system
can enforce physical boundaries, clip unsafe commands, and arrest motion
within budgeted latency windows. But where detectors are absent, the
matrix reveals empty cells---unclosable epistemic gaps where empirical
evidence ends and operational restrictions must govern. It establishes
that across the validated parameter space, every cataloged hazard has an
assigned mitigation path with verified latency budgets and bounded
actuation forces. The authority of that authorization depends entirely
on the validity of its underlying physical premises. A book that teaches
engineers to demand evidence has to be honest about the evidence its own
method cannot produce.

Consider three machines, each possessing an apparently complete safety
record whose line items close without a single arithmetic deficit. The
first is a mobile transport platform moving at \(1.5\text{ m/s}\), whose
stopping distance claim of \(0.49\text{ m}\) assumes an available
tire-to-floor friction coefficient of \(\mu \ge 0.60\). If unmonitored
liquid on the concrete floor reduces that coefficient to \(\mu = 0.15\),
the kinematic braking distance derived from friction limits quadruples
from \(0.19\text{ m}\) to \(0.76\text{ m}\). Combined with the
\(0.30\text{ m}\) traversed during a \(200\text{ ms}\) perception and
actuation latency, the total stopping distance expands to
\(1.06\text{ m}\), breaching the \(0.80\text{ m}\) clearance buffer
before the vehicle halts. The second is an articulated manipulator
commanded to limit contact force to \(40\text{ N}\) during workpiece
assembly, where the force-limiting guarantee assumes structural
stiffness and material integrity across the end-effector mount. If
fatigue micro-cracking introduces unmodeled compliance into the mount,
mechanical deflection stores strain energy under load and releases it
abruptly during surface transitions, delivering an unmeasured
\(120\text{ N}\) transient peak to the workpiece while the joint torque
sensors register nominal values. The third is a chemical flow regulator
maintaining downstream line pressure below \(200\text{ kPa}\), whose
control loop assumes an undrifted diaphragm pressure transducer. If
chemical deposition on the diaphragm induces a steady calibration drift
of \(0.5\text{ kPa/h}\), the nervous system throttles the valve against
a phantom pressure spike or opens it into an actual overpressure state
of \(240\text{ kPa}\), unaware that its reference sensor has drifted
from physical ground truth.

These failures expose the difference between completeness within an
engineering method and completeness of physical knowledge. In all three
systems, every required cell in the verification matrix was populated
with valid data. The timing budgets closed, the safety reflexes passed
fault injection, the thermal dissipation of the actuators remained below
the \(80\,^{\circ}\text{C}\) winding limit, and the runtime enforcer
executed deterministically at \(1000\text{ Hz}\). Yet a necessary
physical premise in each system remained unsupported by direct,
continuous evidence. The record documented compliance with the
specification, but the specification treated environmental friction,
structural mechanics, and sensor calibration as static invariants rather
than dynamic physical states. When an unmonitored invariant shifts, the
physical machine departs from the verified envelope while telemetry
continues to report normal operation.

Treating these blind spots with systems engineering rigor requires
defining an \textbf{evidentiary gap} as an operational quantity rather
than a general caveat: an operational state in a safety case where a
safety-critical claim relies on a physical premise that cannot be
continuously verified during execution or established through prior
testing. To be actionable, every gap must satisfy a four-part
specification: (1) state the affected claim, such as maintaining
deceleration within a designated spatial buffer; (2) identify the
missing observable or unattainable evidence, such as the instantaneous
friction coefficient between the wheel contact patch and an oily
surface; (3) state the present operational containment response, such as
derating the maximum transit speed to \(0.6\text{ m/s}\) or restricting
the machine to dry indoor flooring; and (4) define the predeclared
closure test---the physical measurement, sensor modality, or analytical
result that would close the gap and permit the operational restriction
to be lifted.

Framing the frontier in these terms clarifies how to evaluate new
machine learning models, training regimes, and empirical benchmarks. In
software machine learning, progress is often tracked by aggregate scores
on standardized test datasets. In physical AI, a new model architecture,
an expanded parameter count, or an improved benchmark score is
irrelevant to a deployment authority unless it alters one of the four
entries in an evidentiary gap. If a learned policy achieves higher
average accuracy on an offline manipulation dataset but still requires
an external mechanical guard to prevent over-torque, the safety case
remains identical. Conversely, if an algorithmic advance provides
verifiable bounds on out-of-distribution visual inputs, or if an
improved model reduces inference latency from \(80\text{ ms}\) to
\(10\text{ ms}\) at the \(P_{99}\) tail and thereby recovers
\(0.10\text{ m}\) of stopping margin, that improvement directly alters
the release verdict.

The frontier of the discipline is not a survey of emerging software
techniques, which date quickly as tooling evolves. It is the boundary
where measurable evidence ends and operational containment begins. An
engineer who builds machines that act on the physical world must know
exactly which claims their safety record proves, which claims it merely
assumes, and what instrumentation or physical insight would be required
to transform an assumption into evidence.

\section{What the Method Cannot Establish}\label{sec-frontier-17-2}

A safety release claim for a physical AI system is not a monolithic
assertion of competence. It is a conjunction of discrete, necessary
links that must all hold simultaneously for a physical guarantee to
remain valid. A complete argument requires establishing domain
membership for incoming sensor observations, measured physical bounds on
actuator authority and stopping distances, formal component contracts
across software boundaries, strict execution independence between the
learned policy and the safety enforcement path, quantified detector
coverage across all identified fault modes, and continuous runtime
monitoring of system invariants. Because these premises form a logical
conjunction, the strength of the overall claim is governed by its
weakest constituent link. Generating extensive, high-confidence
empirical evidence for five of these links cannot compensate for a sixth
link that remains unsupported, because a failure in any single premise
causes the deductive safety argument to fail entirely.\sidenote{\footnotesize Karl
  Popper formalized the asymmetry of empirical induction: no finite
  sequence of safe observations (\(\forall t \le T\)) can prove
  universal safety (\(\forall t\)), whereas one counterexample refutes
  it. Operating hours establish survival rates under stationary
  distributions, not out-of-distribution guarantees.}

Every link in a safety case requires an explicit epistemic
classification, identifying each premise as measured directly through
physical instrumentation, derived analytically from measured premises,
assumed on physical or operational grounds, decided by accountable
judgment, left unsupported by existing data, or explicitly declared
outside the threat model. In the three reference machines examined
throughout this text, these classifications reveal the exact boundary
where evidence ends. For the wheeled mobile robot, the stopping-distance
guarantee relies on an assumed friction coefficient between tire and
floor that cannot be measured before braking begins. For the articulated
manipulator, the collision-containment argument depends on an assumed
material-damage threshold for surrounding equipment, which is an
engineering judgment about structural yield rather than a measured
property of the controller. For the aerial vehicle, the state-estimation
envelope assumes statistical independence between optical flow and
inertial sensors, an assumption that collapses under sudden illumination
changes or frame resonance. Naming these premises explicitly does not
weaken the release record. It establishes the physical envelope within
which the release verdict is justified, and outside of which the verdict
ceases to apply.

Separating measurement from judgment is fundamental to systems
engineering discipline. Empirical testing and statistical
characterization can estimate physical exposure, quantify sensor noise,
and bound parametric uncertainty, but no measurement can choose an
acceptable residual risk or decide which individuals in an operational
environment must bear that risk. Those choices are accountable human
decisions that must be recorded as explicit judgments rather than
disguised as empirical findings. In the same way, the threat model
establishes an absolute boundary on the validity of the record.
Empirical testing across nominal fleet operations and accidental
component faults provides no evidence regarding intentional adversarial
attacks. A safety record demonstrating that a vision pipeline handles
sensor noise, lens blur, and partial occlusions provides zero
statistical support for claims about resilience against adversarial
optical projection, model parameter tampering, or intentional denial of
service across the internal communication bus.

When an engineering claim relies on statistical testing to bound the
probability of a hazardous event, the claim requires rigorous formal
definition before conducting trials. The engineering specification must
define the target failure probability \(p\), the physical unit of
exposure, the required statistical confidence level \(1-\alpha\), and
the necessary independence assumptions before counting a single hour or
trial. If a dangerous failure occurs with constant probability \(p\) per
independent, identically distributed exposure unit, the probability of
observing zero failures across \(n\) units of exposure is \((1-p)^n\).
To assert with confidence \(1-\alpha\) that the true failure rate does
not exceed \(p\), the probability of observing zero failures purely by
chance must not exceed \(\alpha\), yielding the inequality
\((1-p)^n \le \alpha\). Taking the natural logarithm of both sides gives
\(n \ln(1-p) \le \ln(\alpha)\). Because \(p \in (0, 1)\), the quantity
\(\ln(1-p)\) is negative, and dividing by it reverses the inequality,
yielding the required exposure bound:
\[n \ge \frac{\ln(\alpha)}{\ln(1-p)}\] For rare failure events where
\(p \ll 1\), the first-order Taylor approximation
\(\ln(1-p) \approx -p\) simplifies this relation to
\(n \ge \ln(1/\alpha)/p\).

As an analytical example, consider a safety requirement specifying that
a high-consequence physical failure must occur with probability
\(p \le 10^{-9}\) per operating hour at a 95 percent confidence level
(\(1-\alpha = 0.95\), so \(\alpha = 0.05\)). Substituting these values
into the exact relation demonstrates the scale of empirical testing
required to support the claim. With \(\ln(0.05) \approx -2.9957\) and
\(\ln(1 - 10^{-9}) \approx -1.0000 \times 10^{-9}\), the required number
of independent, failure-free operational hours is:
\[n \ge \frac{-2.9957}{-1.0000 \times 10^{-9}} \approx 3.0 \times 10^{9}\text{ operating hours}\]
Accumulating \(3.0 \times 10^{9}\text{ operating hours}\) without a
single observed failure requires approximately 342,000 machine-years of
continuous execution (Figure~\ref{fig-17-epistemic-limits}, Panel 2).
This threshold represents the \textbf{astronomical exposure wall}: the
mathematical scaling limit (\(n \ge \frac{\ln(1/\alpha)}{p}\)) where
certifying ultra-rare tail failure probabilities
(\(p \le 10^{-9}/\text{h}\)) at statistical confidence (\(1-\alpha\)) in
open-world physical domains requires billions of independent,
failure-free operational exposure hours (\(>10^9\text{ h}\) or hundreds
of thousands of machine-years), rendering pure empirical black-box fleet
testing logically and economically insufficient.\sidenote{\footnotesize Standards like
  IEC 61508 (SIL 4) target failure rates of
  \(10^{-9}\text{ failures/h}\). Butler and Finelli (1993) proved
  verifying such rates via black-box testing requires billions of hours
  (\(>10^9\text{ h}\)), necessitating formal architectural containment
  and deterministic runtime enforcers.}

\protect\phantomsection\label{callout-mathux2a-17.2}
\begin{fbxSimple}{callout-math}{Fermi Math:}{Observational Deadlines, Tail Exposure Scaling, and Energy Dissipation}
\phantomsection\label{callout-math*-17.2}

Three fundamental calculations demonstrate why empirical fleet testing
cannot substitute for deterministic physical containment:

\begin{enumerate}
\def\labelenumi{\arabic{enumi}.}
\item
  \textbf{Observational Indistinguishability Deadline Deficit:} Consider
  a high-speed robotic gripper closing on a brittle ceramic workpiece
  where brittle fracture occurs at normal load
  \(F_{\text{crit}} = 15\text{ N}\). Given a force ramp rate
  \(\dot{F} = 400\text{ N/s}\), the physical time-to-harm from initial
  contact is:
  \[t_{\text{harm}} = \frac{F_{\text{crit}}}{\dot{F}} = \frac{15\text{ N}}{400\text{ N/s}} = 37.5\text{ ms}\]
  The real-time nervous system and plant introduce actuator current
  decay latency \(\tau_{\text{act}} = 25\text{ ms}\) and fieldbus
  serialization latency \(\tau_{\text{bus}} = 2\text{ ms}\). For a
  runtime detection filter to prevent fracture, detection latency must
  satisfy:
  \[\tau_{\text{det}} \le t_{\text{harm}} - \tau_{\text{act}} - \tau_{\text{bus}} = 37.5\text{ ms} - 25\text{ ms} - 2\text{ ms} = 10.5\text{ ms}\]
  If an exteroceptive camera pipeline operates at \(30\text{ Hz}\)
  (\(\tau_{\text{frame}} = 33.3\text{ ms}\)) with neural inference
  latency \(\tau_{\text{infer}} = 20\text{ ms}\), total perception
  latency is \(\tau_{\text{det}} = 53.3\text{ ms}\). The timing budget
  deficit is:
  \[\Delta \tau = \tau_{\text{det}} - \tau_{\text{det,max}} = 53.3\text{ ms} - 10.5\text{ ms} = +42.8\text{ ms}\]
  Because \(\Delta \tau > 0\), algorithmic detection is physically
  impossible before fracture occurs, mandating passive spring compliance
  (\(k = 2.5\text{ N/mm}\)) rather than vision software.
\item
  \textbf{Astronomical Exposure Scaling vs.~Fleet Testing:} To certify a
  catastrophic autonomy failure rate \(p \le 10^{-9}/\text{h}\) at
  \(95\%\) confidence (\(\alpha = 0.05\)):
  \[n \ge \frac{\ln(1/0.05)}{10^{-9}} = \frac{2.9957}{10^{-9}} \approx 3.0 \times 10^9\text{ hours} \approx 342{,}465\text{ machine-years}\]
  For an industrial fleet of \(1{,}000\) robots deployed
  \(8\text{ h/day}\) (\(2.92 \times 10^6\text{ fleet-h/year}\)),
  accumulating \(3.0 \times 10^9\text{ h}\) of independent, failure-free
  data would require:
  \[T_{\text{test}} = \frac{3.0 \times 10^9\text{ h}}{2.92 \times 10^6\text{ h/year}} \approx 1{,}027\text{ calendar years}\]
  Any software checkpoint update resets this counter to zero,
  demonstrating why safety cases must rest on deterministic barrier
  shields rather than fleet sampling.
\item
  \textbf{Kinetic Energy Clamping and Friction Collapse:} A
  \(450\text{ kg}\) transport robot traveling at nominal velocity
  \(v_{\text{nom}} = 2.2\text{ m/s}\) possesses kinetic energy:
  \[E_{k,\text{nom}} = \frac{1}{2} m v_{\text{nom}}^2 = \frac{1}{2}(450\text{ kg})(2.2\text{ m/s})^2 = 1089\text{ J}\]
  When environmental surface friction collapses from clean concrete
  (\(\mu = 0.70\)) to a thin lubricant spill (\(\mu = 0.12\)), stopping
  distance with sensorimotor reaction delay \(\tau = 50\text{ ms}\)
  expands to:
  \[d_{\text{stop,nom}} = v_{\text{nom}} \tau + \frac{v_{\text{nom}}^2}{2 \mu g} = (2.2)(0.050) + \frac{2.2^2}{2(0.12)(9.81)} = 0.110 + 2.056 = 2.166\text{ m}\]
  By clamping velocity in firmware to
  \(v_{\text{clamp}} = 0.90\text{ m/s}\), kinetic energy falls to:
  \[E_{k,\text{clamp}} = \frac{1}{2}(450\text{ kg})(0.90\text{ m/s})^2 = 182.25\text{ J} \quad (-83.3\%)\]
  The clamped stopping distance becomes:
  \[d_{\text{stop,clamp}} = (0.90)(0.050) + \frac{0.90^2}{2(0.12)(9.81)} = 0.045 + 0.344 = 0.389\text{ m} \le 0.50\text{ m}\]
  This structural velocity clamp preserves the \(0.50\text{ m}\)
  clearance buffer without requiring real-time friction sensing.
\end{enumerate}

\end{fbxSimple}

Operational fleet data rarely satisfies the statistical assumptions
underlying this derivation. Aggregating runtime across a fleet of one
thousand machines operating in parallel does not yield one thousand
independent hours per wall-clock hour, because shared environmental
conditions, correlated fleet-wide software updates, and common spatial
layouts introduce systematic dependencies. Furthermore, learned systems
undergo continuous modification; deploying an updated neural network
checkpoint, retuned sensor filter, or modified calibration table resets
the statistical trial counter, because the new software constitutes a
distinct statistical distribution. Human safety interventions further
distort empirical records through censoring. When an attentive human
operator intercepts an erratic motion before contact occurs, the
intervention truncates a potential accident, converting an uncontained
failure into an unrecorded counterfactual. Similarly, unvalidated
physics simulations cannot substitute for physical exposure, because
simulation trials test only the disturbance distributions and contact
dynamics that the modeler explicitly programmed, leaving real-world
transfer discrepancies unmeasured.

An engineer must therefore distinguish between evidence that is
unaffordable and evidence that is logically insufficient. Demonstrating
a failure rate bound of \(p \le 10^{-4}\) per operating hour on a
mechanical brake actuator may be unaffordable within an eight-week test
campaign, but it remains a tractable physical hypothesis that additional
test benches and capital can establish. Conversely, attempting to prove
open-world safety bounds for an end-to-end vision-language-action model
through empirical sampling is logically insufficient, because no finite
dataset can span the unbounded combinatorial possibilities of
unconstrained physical environments. When a required premise cannot be
settled by evidence, accumulating further empirical test hours is
futile. The defensible engineering responses are structural. An engineer
can narrow the operational design domain to eliminate untestable
environmental states, bound the mechanical authority of the actuators so
that maximum kinetic energy remains below structural damage limits,
install an independent deterministic safety enforcer in the nervous
system with veto authority over the learned policy, shorten the temporal
duration of the release verdict to mandate frequent physical
re-inspection, assign the residual risk explicitly to a designated human
operator through formal interface contracts, or refuse to issue a
release verdict.\sidenote{\footnotesize In safety architectures, physical authority is
  enforced at the register level. Hardware fault registers (DSP
  \texttt{TZFLG}, MCU \texttt{TIMx\_BDTR}) force PWM outputs into
  shutdown in \(<200\text{ ns}\) upon comparator trip, bypassing
  software pipelines and isolating real-time DMA traffic.}

This epistemic boundary is vividly illustrated in extreme-environment
exploration platforms, such as planetary rovers drilling and navigating
unstructured extraterrestrial terrain
(Figure~\ref{fig-17-mars-perseverance-drill}). Operating millions of
kilometers from Earth under multi-minute communication latencies that
preclude human teleoperation, the machine must interact with
uncharacterized geological formations whose internal fracture dynamics
and surface friction cannot be certified through prior fleet testing.
Mission survivability is achieved not by asserting statistical
competence over unknown terrain distributions, but by enforcing
deterministic physical authority limits: hardware torque clutches,
passive suspension compliance, and hard-coded nervous-system reflexes
that halt actuator drives the instant contact impedance departs from
safe envelopes.

\begin{figure}[t]

\centering{


\includegraphics[width=0.85\linewidth,height=\textheight,keepaspectratio]{/Users/VJ/GitHub/MLSysBook-integrate-physical/physical/book/chapters/17-frontier/figures/fig17_real_perseverance_mars_drill.jpg}

}

\caption{\label{fig-17-mars-perseverance-drill}\textbf{Autonomous
Robotic Manipulation and Core Sampling in Unstructured Extraterrestrial
Terrain.} NASA's Perseverance Mars Rover operating at the ``Rochette''
bedrock outcrop in Jezero Crater (Sol 198), displaying its 2.1-meter
five-degree-of-freedom robotic arm and 45-kilogram turret equipped with
a rotary-percussive coring drill, Mastcam-Z stereoscopic navigation
cameras, and proximity sensors. Two visible coring boreholes
(\emph{Montdenier} and \emph{Montagnac}) illustrate direct mechanical
contact interaction with uncharacterized geological strata. Operating
under a 4-to-24-minute one-way speed-of-light communication latency to
Earth, the rover cannot rely on human teleoperation to catch contact
bifurcations or slip hazards. Because empirical testing across open
Martian environments is logically incapable of exhausting all
unmonitored physical invariants (such as subsurface rock fractures,
grain cohesion, or unexpected drill-bit binding), safety is guaranteed
structurally: passive mechanical compliance, deterministic motor current
torque clamps, and autonomous nervous-system reflexes bound energy
transfer and arrest motion before irreversible structural yield can
occur. Image credit: NASA/JPL-Caltech/MSSS (public domain).}

\end{figure}%

\section{A Residual-Claims Register}\label{sec-frontier-17-3}

An unproven premise does not vanish when an engineer omits it from a
safety report. When empirical testing cannot establish a safety bound,
the unproven premise must be entered into a formal
\textbf{residual-claims register}: an auditable systems engineering
ledger that maps every unproven physical premise or unmonitored
invariant in a safety case to its missing physical observable, active
mechanical or architectural containment constraint, named accountable
human owner, and predeclared closure test. An engineering team cannot
treat an unverified physical assumption as an unwritten risk or a
temporary modeling convenience. Every entry in the register translates
an epistemic shortfall into four mandatory components, comprising the
unsupported claim and its missing physical observable, the structural
operational response that contains the hazard, the named engineer who
owns the operational risk, and the closure test that defines what
evidence would retire the entry. Without these four components, an
unresolved claim is simply an unmonitored defect.

Every entry begins by linking the unsupported claim to the specific
physical observable that cannot be measured. In an embodied system with
physical authority over mass, a learned policy fails to settle a safety
property because sensors cannot transduce the relevant environmental
state, the computing substrate cannot process the measurement within the
control deadline, or the physical interaction cannot be observed before
contact occurs. An engineer must define the missing observable with
physical precision rather than statistical abstractions. Characterizing
an unknown failure as an out-of-distribution latent state provides no
actionable signal for an enforcement loop. Stating an unmeasured surface
friction coefficient, an undetected optical flare across a camera lens
array, or an unobservable micro-fracture in an actuator housing
identifies the exact physical deficit. Naming the missing observable
establishes the boundary where learned inference stops and physical
uncertainty begins.

In an analytical example, consider a \(450\text{ kg}\) autonomous mobile
transport robot moving palletized goods across an industrial logistics
depot. The safety case asserts that the learned navigation policy will
never command a turning maneuver that causes wheel slip and lateral
skidding on smooth concrete floors contaminated with spilled lubricant.
The missing observable is the local surface friction coefficient
\(\mu\), which varies from \(\mu = 0.70\) on clean dry concrete down to
\(\mu = 0.12\) under a lubricant slick. Standard RGB cameras and
structured-light depth sensors cannot detect a transparent,
two-millimeter layer of synthetic machine oil before the drive wheels
make contact. Because this observable cannot be transduced prior to
traversal, empirical fleet hours cannot establish that the learned
policy will maintain traction during high-speed turns.

An open entry in the register requires an immediate operational response
enforced in the nervous system or the physical plant. The response
cannot consist of a warning label, an advisory note in an operator
manual, or a software retry loop. It must be a structural constraint
that prevents the unmeasured observable from causing damage. In the
transport robot example, the nervous system microcontroller clamps the
maximum translational velocity from the nominal \(2.2\text{ m/s}\) down
to \(v_{\text{clamp}} = 0.9\text{ m/s}\), and derates maximum allowable
lateral acceleration to \(a_{\text{lat}} \le 1.0\text{ m/s}^2\) within
its static configuration table. This velocity restriction drops the
vehicle's kinetic energy from
\(\frac{1}{2}(450\text{ kg})(2.2\text{ m/s})^2 = 1089\text{ J}\) down to
\(\frac{1}{2}(450\text{ kg})(0.9\text{ m/s})^2 = 182\text{ J}\), an
\(83\text{ percent}\) reduction in destructive potential. At
\(0.9\text{ m/s}\), the worst-case stopping distance on an oil slick
with \(\mu = 0.12\) (producing a deceleration of
\(a = \mu g \approx 1.18\text{ m/s}^2\)) under a \(50\text{ ms}\)
sensorimotor reaction delay is
\(d_{\text{stop}} = (0.9 \times 0.050) + \frac{0.9^2}{2 \times 1.18} = 0.045 + 0.343 = 0.388\text{ m}\).
This stopping distance keeps the machine within the \(0.50\text{ m}\)
safety margin maintained by its proximity sensors.

Every operational response requires a named owner with the authority to
maintain the constraint against organizational pressure. If ownership is
assigned to an entire development team or an anonymous safety committee,
the constraint will erode whenever deployment schedules tighten or
customer throughput targets slip. The owner is an individual engineer or
an operational manager whose signature appears on the release verdict.
In the transport robot system, the drivetrain engineering lead and the
depot operations manager serve as co-owners. The drivetrain lead owns
the firmware parameter lock on the velocity clamp, ensuring that no
software update can override the \(0.9\text{ m/s}\) ceiling. The
operations manager owns the physical operating environment, guaranteeing
that pedestrian walkways remain separated by physical barriers wherever
transport paths intersect. Both owners legally accept the throughput
penalty as the direct operational price of operating with an unproven
claim.

Every entry must define its closure test at the moment the gap is
identified. If an engineering team defines the closure test after
gathering new data, the criteria will inevitably shift to match whatever
measurements were inexpensive to collect. A closure test specifies the
verifiable physical measurement, sensor modality, or mathematical
invariant that eliminates the dependency on the operational constraint.
For the mobile transport robot, the closure test requires integrating an
onboard polarimetric surface reflectivity sensor validated across
\(N = 1000\) empirical trials under controlled surface contaminants,
combined with a firmware-level wheel-encoder differential slip detector
that executes at \(1000\text{ Hz}\) and triggers emergency braking
within \(10\text{ ms}\) of slip onset. Once that instrumentation is
installed, calibrated, and proven to satisfy the \(P_{99.9}\) detection
bound, the entry is retired and the velocity clamp is removed.

Specifying a closure test forces an engineering organization to
determine whether a gap can be resolved with improved instrumentation
and testing, or whether it represents a permanent unobservable state.
When a closure test requires a physical transducer that cannot be
manufactured or a computational budget that the edge platform cannot
support, the register entry remains open for the operating lifetime of
the machine. The operational response becomes a permanent architectural
boundary. The danger arises when an organization treats an unclosable
entry as a temporary backlog item, assuming that every physical failure
mode can eventually be trapped by an onboard runtime detector.

The structured anatomy of these register entries is formalized across
representative physical AI archetypes in
Table~\ref{tbl-17-claims-register}, illustrating how unmonitored hazards
are systematically bounded by operational constraints until falsifiable
closure tests are satisfied.

\begin{longtable}[]{@{}
  >{\raggedright\arraybackslash}p{(\linewidth - 8\tabcolsep) * \real{0.2000}}
  >{\raggedright\arraybackslash}p{(\linewidth - 8\tabcolsep) * \real{0.2000}}
  >{\raggedright\arraybackslash}p{(\linewidth - 8\tabcolsep) * \real{0.2000}}
  >{\raggedright\arraybackslash}p{(\linewidth - 8\tabcolsep) * \real{0.2000}}
  >{\raggedright\arraybackslash}p{(\linewidth - 8\tabcolsep) * \real{0.2000}}@{}}
\caption{\textbf{The Residual-Claims Register Matrix.} Auditable ledger
of unproven physical premises across representative physical AI
archetypes. For each open claim, the register defines the missing
physical observable, the physical hazard mechanism, the active
operational containment enforced by the nervous system or mechanical
hardware, the accountable human owners, and the falsifiable closure test
required to retire the operational
restriction.}\label{tbl-17-claims-register}\tabularnewline
\toprule\noalign{}
\begin{minipage}[b]{\linewidth}\raggedright
System Archetype \& Release Claim
\end{minipage} & \begin{minipage}[b]{\linewidth}\raggedright
Missing Physical Observable
\end{minipage} & \begin{minipage}[b]{\linewidth}\raggedright
Physical Hazard Mechanism
\end{minipage} & \begin{minipage}[b]{\linewidth}\raggedright
Operational Containment \& Accountable Owner
\end{minipage} & \begin{minipage}[b]{\linewidth}\raggedright
Predeclared Closure Test
\end{minipage} \\
\midrule\noalign{}
\endfirsthead
\toprule\noalign{}
\begin{minipage}[b]{\linewidth}\raggedright
System Archetype \& Release Claim
\end{minipage} & \begin{minipage}[b]{\linewidth}\raggedright
Missing Physical Observable
\end{minipage} & \begin{minipage}[b]{\linewidth}\raggedright
Physical Hazard Mechanism
\end{minipage} & \begin{minipage}[b]{\linewidth}\raggedright
Operational Containment \& Accountable Owner
\end{minipage} & \begin{minipage}[b]{\linewidth}\raggedright
Predeclared Closure Test
\end{minipage} \\
\midrule\noalign{}
\endhead
\bottomrule\noalign{}
\endlastfoot
\textbf{Autonomous Mobile Robot (450 kg)}\emph{Claim}: Stopping distance
\(d \le 0.50\text{ m}\) on depot concrete
(\(v_{\text{nom}} = 2.2\text{ m/s}\)). & Forward surface friction
coefficient \(\mu \in [0.12, 0.70]\) under thin fluid/lubricant films
(\(h \le 2\text{ mm}\)). & Reduced friction quadruples kinematic braking
distance (\(d = \frac{v^2}{2\mu g} + v\tau\)), expanding stopping
envelope from \(0.39\text{ m}\) to \(1.64\text{ m}\) and breaching
clearance buffer. & Drivetrain firmware clamps max velocity to
\(v_{\text{clamp}} \le 0.90\text{ m/s}\) (\(E_{\text{kin}}\) drops from
\(1089\text{ J}\) to \(182\text{ J}\), \(-83.3\%\)) and lateral
acceleration \(a_{\text{lat}} \le 1.0\text{ m/s}^2\); physical depot
pedestrian barriers.\emph{Owners}: Drivetrain Lead \& Depot Operations
Manager. & Integrate forward polarimetric scatterometer (\(P_{99.9}\)
oil detection at \(3.0\text{ m}\)) + \(1000\text{ Hz}\) wheel-encoder
differential slip reflex (\(t_{\text{arrest}} \le 10\text{ ms}\))
verified across \(N \ge 1000\) contaminated track trials. \\
\textbf{Articulated Manipulator (6-DoF)}\emph{Claim}: Bounded contact
force \(F \le 40\text{ N}\) during brittle ceramic workpiece assembly. &
Subsurface voids or micro-cracks (\(V_{\text{void}} > 0.5\text{ mm}^3\))
within uninspected ceramic workpieces. & Normal force ramps at
\(400\text{ N/s}\), crossing the \(15\text{ N}\) brittle fracture
threshold in \(38\text{ ms}\); actuator current-decay delay
(\(\tau_{\text{decay}} = 25\text{ ms}\)) prevents closed-loop force
arrest. & Passive mechanical compliance interface with calibrated spring
detent (\(k = 2.5\text{ N/mm}\)) and pneumatic breakaway relief limiting
peak force to \(F_{\text{peak}} \le 18\text{ N}\); derate jaw speed to
\(5\text{ mm/s}\).\emph{Owner}: Tooling \& End-Effector Lead. &
Integrate jaw-embedded ultrasonic pulse-echo transducer
(\(f = 5\text{ MHz}\), \(\tau_{\text{meas}} \le 12\text{ ms}\))
detecting structural voids prior to reaching \(5\text{ N}\) clamping
force with false-alarm rate \(\alpha_{\text{FA}} \le 10^{-4}\). \\
\textbf{Chemical Flow Regulator}\emph{Claim}: Downstream line pressure
bound \(P \le 200\text{ kPa}\) over \(2000\text{ h}\) continuous
campaign. & Sensor zero-point baseline drift
(\(\kappa = 0.5\text{ kPa/h}\)) from chemical deposition on silicon
diaphragm. & Downward drift causes controller to open valve against
phantom under-pressure, driving physical manifold to \(240\text{ kPa}\)
(rupture risk) while telemetry reads nominal \(195\text{ kPa}\). &
Hard-wired analog rupture burst-disc
(\(P_{\text{burst}} = 215\text{ kPa}\)) and mandatory scheduled physical
zero-point proof-tests every
\(\Delta t_{\text{proof}} = 72\text{ h}\).\emph{Owners}: Process Safety
Lead \& Plant Reliability Engineer. & Install dual-dissimilar sensing:
piezoresistive transducer paired with non-intrusive ultrasonic Coriolis
flow meter, triggering cross-channel disparity trip when
\(\lvert \Delta P_{\text{disparity}} \rvert \ge 5.0\text{ kPa}\). \\
\textbf{Multi-Rotor Aerial Drone}\emph{Claim}: Positional drift
\(e_{\text{pos}} \le 0.10\text{ m}\) during rapid transitions across
indoor corridors. & Visual feature track persistence under
high-dynamic-range lighting shifts (\(>10^4\text{ lux/s}\)). & Sudden
feature loss triggers visual odometry divergence; IMU double-integration
error drifts quadratically
(\(\sigma_x \propto \frac{1}{2} a_{\text{bias}} t^2 \approx 0.8\text{ m}\)
in \(400\text{ ms}\)), causing perimeter collision. & Autonomous
geofence safety bubble (\(d_{\text{standoff}} \ge 1.5\text{ m}\))
enforced by independent time-of-flight (ToF) pulsed-lidar safety shield
triggering auto-hover retro-thrust within \(20\text{ ms}\).\emph{Owner}:
Flight Autonomy Lead. & Integrate multi-spectral event-based
neuromorphic vision sensor (\(>120\text{ dB}\) dynamic range, latency
\(\tau \le 1\text{ ms}\)) maintaining state error
\(e \le 0.05\text{ m}\) over \(N = 500\) abrupt blackout transitions. \\
\end{longtable}

\subsection*{Frontier Synthesis and Residual-Claims Manifest
Schema}\label{frontier-synthesis-and-residual-claims-manifest-schema}
\addcontentsline{toc}{subsection}{Frontier Synthesis and Residual-Claims
Manifest Schema}

To ensure auditable tracking of unmonitored invariants across software
deployments and physical revisions, safety-critical systems represent
open residual claims through a formal, structured data contract
(Table~\ref{tbl-17-residual-claims-schema}). The abstract schema below
defines the mandatory fields, scalar types, physical SI units, and
structural invariants required for every unresolved physical assumption:

\begin{longtable}[]{@{}
  >{\raggedright\arraybackslash}p{(\linewidth - 6\tabcolsep) * \real{0.2500}}
  >{\raggedright\arraybackslash}p{(\linewidth - 6\tabcolsep) * \real{0.2500}}
  >{\raggedright\arraybackslash}p{(\linewidth - 6\tabcolsep) * \real{0.2500}}
  >{\raggedright\arraybackslash}p{(\linewidth - 6\tabcolsep) * \real{0.2500}}@{}}
\caption{Concrete Data Schema: Epistemic Residual-Claims
Register.}\label{tbl-17-residual-claims-schema}\tabularnewline
\toprule\noalign{}
\begin{minipage}[b]{\linewidth}\raggedright
Manifest Category
\end{minipage} & \begin{minipage}[b]{\linewidth}\raggedright
Parameter \& Type
\end{minipage} & \begin{minipage}[b]{\linewidth}\raggedright
Physical Unit
\end{minipage} & \begin{minipage}[b]{\linewidth}\raggedright
Invariant \& Operational Contract
\end{minipage} \\
\midrule\noalign{}
\endfirsthead
\toprule\noalign{}
\begin{minipage}[b]{\linewidth}\raggedright
Manifest Category
\end{minipage} & \begin{minipage}[b]{\linewidth}\raggedright
Parameter \& Type
\end{minipage} & \begin{minipage}[b]{\linewidth}\raggedright
Physical Unit
\end{minipage} & \begin{minipage}[b]{\linewidth}\raggedright
Invariant \& Operational Contract
\end{minipage} \\
\midrule\noalign{}
\endhead
\bottomrule\noalign{}
\endlastfoot
\textbf{Manifest Header} & \texttt{system\_identifier} (\texttt{str}) &
--- & Unique embodied platform ID (e.g., \texttt{AMR-450-DEPOT}) \\
\textbf{Manifest Header} & \texttt{hardware\_build\_rev} (\texttt{str})
& --- & Target silicon and chassis revision (\texttt{HW-REV-3.2}) \\
\textbf{Manifest Header} & \texttt{software\_digest} (\texttt{str}) &
--- & Cryptographic binary hash (SHA-256) \\
\textbf{Manifest Header} & \texttt{validity\_limit\_hours}
(\texttt{float}) & hours & Maximum statutory operating window before
review \\
\textbf{Residual Claims} & \texttt{entry\_id} (\texttt{str}) & --- &
Unique claim tag (e.g., \texttt{CLAIM-FRIC-001}) \\
\textbf{Residual Claims} & \texttt{target\_safety\_claim} (\texttt{str})
& --- & Explicit safety invariant
(\(d_{\text{stop}} \le 0.50\text{ m}\)) \\
\textbf{Residual Claims} & \texttt{missing\_observable} (\texttt{str}) &
SI units & Unmeasured physical state (e.g., surface friction \(\mu\)) \\
\textbf{Physical Hazard} & \texttt{governing\_law} (\texttt{str}) & ---
& Newton-Euler, Navier-Stokes, Joule heating \\
\textbf{Physical Hazard} & \texttt{time\_to\_harm} (\texttt{float}) &
seconds & Time from unmodeled transition to irreversible yield \\
\textbf{Operational Containment} & \texttt{containment\_type}
(\texttt{enum}) & --- & \texttt{FIRMWARE\_PARAM\_CLAMP},
\texttt{PASSIVE\_MECHANICAL}, \texttt{PROOF\_TEST} \\
\textbf{Operational Containment} & \texttt{authority\_limit}
(\texttt{float}) & \(\text{m/s}, \text{N}\) & Hard numeric setpoint
ceiling (\(v_{\max}, F_{\max}\)) \\
\textbf{Operational Containment} & \texttt{kinetic\_energy\_ceiling}
(\texttt{float}) & \(\text{J}\) & Maximum allowable kinetic energy under
fault (\(E_k \le E_{\max}\)) \\
\textbf{Predeclared Closure} & \texttt{closure\_modality} (\texttt{str})
& --- & Transducer technology or formal method required \\
\textbf{Predeclared Closure} & \texttt{sample\_trials} (\texttt{int}) &
count & Minimum empirical sample exposure (\(N \ge 1000\)) \\
\textbf{Predeclared Closure} & \texttt{confidence\_level}
(\texttt{float}) & --- & Statistical confidence threshold
(\(1 - \alpha \ge 0.999\)) \\
\end{longtable}

\section{Failures Without Detectors}\label{sec-frontier-17-4}

A runtime detector is not an abstract classifier. A \textbf{practical
detector} is a physical subsystem governed by three coupled parameters,
namely coverage across the failure distribution, a bounded false-alarm
rate, and a detection latency strictly shorter than the physical
system's time-to-harm. If a mobile transport base traveling at
\(1.5\text{ m/s}\) requires \(180\text{ ms}\) of braking actuation to
avoid colliding with an obstacle, a perception pipeline requiring
\(250\text{ ms}\) of sensor aggregation and compute latency is not a
detector for that hazard. It is a post-incident logging mechanism.
Producing an accurate classification label after the physical body has
crossed an irreversible threshold generates diagnostic evidence for an
investigation, but it cannot participate in the safety case as a runtime
safeguard.

The fundamental limit of runtime detection arises from observational
indistinguishability. Consider two physical states of the world, \(S_A\)
and \(S_B\). In state \(S_A\), a commanded actuator trajectory \(u(t)\)
is safe. In state \(S_B\), the identical trajectory \(u(t)\) causes
structural damage or human injury, requiring the nervous system to
intervene and execute a stopping reflex. The dynamics of the body impose
a hard deadline \(t_{\text{deadline}}\), determined by actuator response
time and system momentum, after which physical contact or mechanical
yield cannot be prevented. If the telemetry history \(y(t)\) provided by
onboard instrumentation is identical for \(S_A\) and \(S_B\) across the
entire interval \(t \le t_{\text{deadline}}\), no algorithm operating on
those observations can reliably distinguish the benign state from the
hazardous one. When two physical realities produce identical sensor
traces up to the action deadline, the probability of safe intervention
cannot exceed the unconditioned prior, regardless of model scale or
compute capacity.

\protect\phantomsection\label{callout-definitionux2a-17.3}
\begin{fbxSimple}{callout-definition}{Definition:}{Observational Indistinguishability}
\phantomsection\label{callout-definition*-17.3}

The physical condition where two distinct world states---one nominal
(\(S_A\)) and one catastrophic (\(S_B\))---generate identical pre-harm
sensor telemetry (\(y_A(t) \equiv y_B(t)\)) up to the system's physical
action deadline
(\(t_{\text{deadline}} = t_{\text{harm}} - t_{\text{act}}\)), rendering
timely algorithmic differentiation physically and mathematically
impossible regardless of model capacity, inference compute, or training
dataset volume.

\end{fbxSimple}

In learned perception-action policies, this detectability limit
manifests as silent distribution shift and embodied shortcut learning.
High-capacity neural networks, such as Vision-Language-Action (VLA)
models or deep imitation policies, minimize empirical loss by
discovering opportunistic statistical correlations across
high-dimensional demonstration datasets. In embodied tasks, unmeasured
physical properties---such as contact friction, material yield strength,
or payload inertia---are unobserved latent variables. Deep models
routinely learn to exploit non-causal visual features (such as ambient
illumination angles, floor specular highlights, or camera frame-rate
synchronization) as computational proxies for these physical properties.
When operating under open-world distribution shift, these non-causal
correlations collapse. However, because high-dimensional non-linear
encoders compress the out-of-distribution observation \(\mathbf{x}\)
directly into a nominal latent manifold
\(\mathbf{z} = \phi(\mathbf{x}) \in \mathcal{Z}_{\text{nom}}\), the
policy outputs action distributions \(\pi(\mathbf{a} \mid \mathbf{z})\)
with near-zero softmax entropy (\(H(\pi) \approx 0\)) and minimal
ensemble disagreement. The model commands aggressive, damaging physical
trajectories while its internal representations indicate nominal
performance. Crucially, any epistemic uncertainty estimator implemented
in software (such as latent Mahalanobis distance, Monte Carlo dropout,
or autoencoder reconstruction residuals) evaluates uncertainty using the
very same corrupted representations \(\mathbf{z}\). Because the
confidence estimator and the policy share the broken perceptual
representations, internal software metrics cannot serve as an
independent detector of shift. This failure mode is termed
\textbf{embodied shortcut learning and latent epistemic shift}: the
process where a high-capacity neural policy exploits non-causal
statistical correlations in training demonstrations as computational
proxies for unmeasured physical states, causing the model to command
destructive actions with maximal confidence while rendering internal
uncertainty metrics completely blind to the failure. The learned Brain
is fundamentally blind to its own shortcut failure, necessitating an
independent deterministic Nervous System to enforce physical
containment.

Physical sensing hardware exhibits an identical lack of identifiability
over longer operating durations. In a fluid regulation system, an inline
thermal mass flow sensor subjected to surface fouling or thermal aging
drifts downward at an analytical example rate of \(0.02\text{ L/min}\)
per \(100\text{ h}\) of operation. Over the same timescale, an incipient
pipe constriction or supply pressure drop produces an identical, genuine
reduction in mass flow. A single measurement channel measuring heat
dissipation reports a declining scalar signal \(q(t)\). The nervous
system receives the telemetry stream, but no algorithm processing that
single signal can identify whether the physical process has changed or
the transducer calibration has degraded. Both physical conditions
generate the exact same observation sequence. Without an independent
reference transducer or a direct differential measurement, the drift is
undetectable at runtime.\sidenote{\footnotesize Piezoresistive diaphragms and strain
  gauges undergo mechanical creep and lattice fatigue over thermal
  cycles, shifting baseline offset voltages
  (\(V_{\text{offset}} = V_0 + \kappa t\)). Single-channel software
  filters cannot distinguish transducer aging from genuine physical
  changes without out-of-band calibration.}

Mechanical manipulation reveals failures where the threshold for safe
action cannot be known before physical contact occurs. Consider a
robotic manipulator grasping ceramic workpieces. A sound workpiece
safely tolerates a normal gripping force of \(50\text{ N}\), whereas a
workpiece containing an internal structural void fractures when normal
force exceeds \(15\text{ N}\). Exterior optical cameras and tactile skin
arrays inspect the workpiece before grasping, but internal voids produce
zero surface signatures. Prior to contact, the force sensors report
\(0.0\text{ N}\). Once contact occurs, the gripper jaws close, ramping
normal force at an analytical example rate of \(400\text{ N/s}\). The
boundary between safe load application and brittle fracture is crossed
in \(38\text{ ms}\), which is shorter than the combined latency of the
force sensor filter, the serial bus transfer, and the motor drive's
current-decay response. The safe force cannot be determined before
contact, and fracture occurs before the nervous system can arrest jaw
motion. This scenario defines a \textbf{detectorless failure mode}: a
physical degradation mechanism, sensor aging effect, or contact
bifurcation whose observation latency, sensing blind spot, or
transduction physics prevents detection within the body's time-to-harm
(\(\tau_{\text{det}} + \tau_{\text{act}} > t_{\text{harm}}\)),
necessitating structural physical authority limits, mechanical
compliance, or scheduled proof-tests rather than software detection
loops.\sidenote{\footnotesize Rigid-body dynamics with Coulomb friction exhibit
  non-smooth velocity jumps (\(\dot{q}^+ \neq \dot{q}^-\)) and Painlevé
  paradoxes. Neural policies lacking uniform Lipschitz bounds cannot
  guarantee closed-loop reachability across impact manifolds without
  mechanical compliance or hardware limiters.}

\protect\phantomsection\label{callout-incidentux2a-17.4}
\begin{fbxSimple}{callout-incident}{Landmark Engineering Incident:}{Deepwater Horizon Blind Shear Ram Structural Degradation}
\phantomsection\label{callout-incident*-17.4}

On April 20, 2010, the Macondo well in the Gulf of Mexico experienced a
catastrophic blowout, leading to the loss of 11 lives, the sinking of
the \emph{Deepwater Horizon} drilling rig, and the discharge of 4.9
million barrels of crude oil. The final physical safeguard intended to
prevent catastrophic containment loss was the subsea Blowout Preventer
(BOP)---a \(450\text{-ton}\) stack of hydraulic valves and mechanical
cutting rams positioned on the seabed \(1{,}500\text{ m}\) below the
surface (\citeproc{ref-csb2014macondo}{U.S. Chemical Safety and Hazard
Investigation Board 2014}).

When gas kicks breached well casing barriers and ignited the drilling
platform, automated emergency systems (the Emergency Disconnect System
and acoustic automatic-mode function) commanded the Blind Shear Rams
(BSR) to close. Surface telemetry and subsea accumulator pressure
transducers reported nominal hydraulic pressure delivery
(\(3{,}000\text{ psi}\) / \(20.7\text{ MPa}\)) and indicated full piston
stroke travel. Control logic recorded the actuation as successful.

In reality, physical containment had completely failed due to an
unmonitored, detectorless physical state transition:

\begin{enumerate}
\def\labelenumi{\arabic{enumi}.}
\tightlist
\item
  \textbf{Flow-Induced Axial Compressive Load:} The uncontained surge of
  high-velocity multiphase oil and gas flowing up the drill pipe created
  an extreme pressure differential (\(\Delta P > 35\text{ MPa}\)) and
  hydrodynamic shear friction. This upward force exerted massive
  compressive loading along the elastic drill string.
\item
  \textbf{Elastic-Plastic Euler Buckling:} Under severe axial
  compression, the drill pipe exceeded its critical Euler buckling load
  (\(P_{\text{cr}} = \pi^2 E I / L_{\text{eff}}^2\)). The pipe buckled
  laterally, deflecting out of the wellbore centerline and pinning
  itself firmly against the inner wall of the blowout preventer casing.
\item
  \textbf{Shear Blade Geometric Incompatibility:} The Cameron blind
  shear rams were designed, tested, and certified assuming a
  concentrically aligned drill pipe positioned squarely in the center of
  the bore. When the hardened V-shaped shear blades advanced, the
  off-center, buckled pipe was crushed and folded against the casing
  wall rather than cleanly severed. The flattened pipe jammed between
  the ram blocks, preventing the elastomeric seals from seating together
  and leaving a massive annular cross-section open for hydrocarbons to
  continue surging upward.
\end{enumerate}

\emph{Systems significance.} The failure of the Deepwater Horizon blind
shear ram is a definitive real-world manifestation of observational
indistinguishability (\(y_A(t) \equiv y_B(t)\)). The subsea BOP
instrumentation monitored only hydraulic supply pressure and piston
stroke position; it possessed zero internal proximity sensors, magnetic
resonance coils, or ultrasound transducers to observe the lateral
spatial displacement or cross-sectional deformation of the drill pipe
inside the casing. Control software operated under the unmonitored
physical premise that the drill pipe remained concentric. Because the
physical time-to-harm from pipe buckling to containment breach was
instantaneous under blowout flow, and because the failure was completely
detectorless to the control loop, software commands could not prevent
disaster. Containment in such extreme regimes cannot rely on blind
hydraulic actuation; it requires structural containment envelopes,
multi-plane redundant shear geometries capable of severing off-center
cross-sections, and hard physical proof-tests.

\end{fbxSimple}

An identical detectorless challenge governs dynamic legged locomotion
across unstructured disaster debris
(Figure~\ref{fig-17-humanoid-disaster-rubble}). While exteroceptive
vision and LiDAR can construct dense point clouds of irregular rubble,
they cannot determine whether a fractured cinderblock will shear, tip,
or crumble under footfall loading prior to touchdown. Because
ground-reaction collapse occurs within tens of milliseconds---far faster
than an end-to-end vision network can re-infer footstep
trajectories---safety cannot rely on vision software. Traversability
depends on low-level joint torque reflexes, passive foot compliance, and
deterministic center-of-mass momentum bounds that arrest or adapt gait
dynamics without waiting for high-level neural replanning.

\begin{figure}[t]

\centering{

\includegraphics[width=0.85\linewidth,height=\textheight,keepaspectratio]{/Users/VJ/GitHub/MLSysBook-integrate-physical/physical/book/chapters/17-frontier/figures/fig17_real_humanoid_disaster_rubble.jpg}

}

\caption{\label{fig-17-humanoid-disaster-rubble}\textbf{Physical AI at
the Disaster Frontier: Humanoid Locomotion Across Unstructured Debris.}
The WARNER humanoid robot (Boston Dynamics Atlas platform with
autonomous balance and footstep planning developed by WPI and Carnegie
Mellon University) negotiating an irregular cinderblock and rubble
debris field during the DARPA Robotics Challenge Finals. In unstructured
disaster environments, exteroceptive sensing (head-mounted multi-sense
LiDAR and stereo cameras) can map surface elevation but cannot measure
the internal structural integrity, frictional stability, or settling
compliance of individual loose cinderblocks before foot touchdown.
Because the contact-to-instability time-to-harm
(\(t_{\text{harm}} \le 30\text{ ms}\)) is substantially shorter than
high-level vision replanning cycles
(\(\tau_{\text{vision}} \ge 100\text{ ms}\)), visual detection of
foothold collapse is physically impossible before balance is lost.
Robust traversal requires low-level joint torque reflexes, passive foot
compliance, and deterministic center-of-mass momentum bounds that arrest
or adapt gait dynamics without waiting for vision software permission.
Image credit: DARPA / WPI-CMU (public domain).}

\end{figure}%

When a failure mode lacks a practical detector, software engineering
teams often attempt to construct heuristic proxies, such as autoencoder
reconstruction errors or sliding-window telemetry variance. These
proxies offer the appearance of safety while leaving the physical hazard
unmitigated. When a failure cannot be detected within the time-to-harm,
only five architectural responses are physically sound. The machine can
reduce actuator authority, capping velocities, forces, or pressures so
that undetected transitions remain non-destructive. It can narrow the
operational design domain, physically excluding conditions where
unobservable states occur. It can install a complementary sensor based
on an independent physical principle, such as pairing optical cameras
with acoustic emission sensors. It can establish scheduled proof tests
and physical inspections to bound the exposure duration of silent
degradation. If none of these measures is feasible, the only valid
engineering decision is to refuse operation.\sidenote{\footnotesize GMSL2
  deserializers introduce fixed serialization latencies
  (\(16\text{--}33\text{ ms}\)). POSIX IPC stacks suffer tail latency
  (\(P_{99.9} > 50\text{ ms}\)), creating dangerous sensory freshness
  deficits.}

The safety case must document these gaps explicitly rather than
concealing them beneath unverified detection claims. As illustrated in
Figure~\ref{fig-17-epistemic-limits}, when safe state \(S_A\) and
hazardous state \(S_B\) produce identical sensor observations prior to
the physical deadline \(t_{\text{deadline}}\), software cannot intervene
before irreversible damage occurs, and statistical sampling cannot
certify ultra-rare failure rates across open-world distributions. The
residual-claims register records the detector's blind region, defining
the set of physical states that produce indistinguishable observations,
along with the maximum exposure duration permitted between physical
proof tests. Treating an undetectable failure as a solved runtime
problem undermines the validity of the safety case. Preserving the exact
boundary where sensing stops forces the engineering organization to rely
on physical constraints rather than software optimism.

\begin{figure}[t]

\centering{


\includegraphics[width=1\linewidth,height=\textheight,keepaspectratio]{index_files/mediabag//Users/VJ/GitHub/MLSysBook-integrate-physical/physical/book/chapters/17-frontier/figures/fig17_epistemic_limits.pdf}

}

\caption{\label{fig-17-epistemic-limits}\textbf{The Dual Epistemic
Boundaries: Observability Deadlines, Astronomical Exposure Scaling, and
the Residual-Claims Register.} (1) \emph{Observational
Indistinguishability:} When safe state \(S_A\) (nominal friction / sound
core) and hazardous state \(S_B\) (lubricant slick / internal void)
yield identical pre-harm telemetry \(y_A(t) \equiv y_B(t) = y_0\),
detection latency \(\tau_{\text{detect}} = 64\text{ ms}\) exceeds the
action deadline
\(t_{\text{deadline}} = t_{\text{harm}} - t_{\text{brake}}\), rendering
safe intervention impossible without mechanical force/velocity limits.
(2) \emph{Astronomical Exposure Wall:} The zero-failure exposure bound
\(n \ge \frac{\ln(1/\alpha)}{p}\) at \(95\%\) confidence requires
\(3.0 \times 10^9\text{ h}\) (\(342{,}000\text{ machine-years}\)) to
support \(p \le 10^{-9}/\text{h}\), proving that open-world
combinatorial spaces cannot be certified by empirical fleet testing
alone. (3) \emph{The Residual-Claims Register:} Every unsettled premise
is mapped to an operational containment boundary and an accountable
owner until falsifiable independent evidence satisfies a predeclared
closure test.}

\end{figure}%

\subsection*{Active Epistemic Probing and Invariant
Monitoring}\label{active-epistemic-probing-and-invariant-monitoring}
\addcontentsline{toc}{subsection}{Active Epistemic Probing and Invariant
Monitoring}

To mitigate detectorless failure modes before physical limits are
breached, safety-critical nervous systems can execute active epistemic
probing. In this pattern, the controller injects a high-frequency,
sub-perceptual mechanical or electrical dither \(\delta u(t)\) (such as
a \(125\text{ Hz}\) torque ripple
\(\delta \tau \le 0.03\,\tau_{\text{max}}\)) and measures the dynamic
plant impedance response. If the measured transfer function departs from
nominal stiffness or flux bounds, the nervous system trips a
deterministic fallback before irreversible yield occurs, as formalized
in Algorithm 17.1.

\protect\phantomsection\label{callout-algorithmux2a-17.5}
\begin{fbxSimple}{callout-algorithm}{Algorithm:}{Algorithm 17.1: Active Epistemic Probe and Invariant Degradation Monitor}
\phantomsection\label{callout-algorithm*-17.5}

\textbf{Inputs:} Hardware telemetry stream \(\mathbf{y}_t\), learned
action proposal \(\mathbf{u}_{\text{policy}}\), nominal plant parameters
\(\mathcal{M}_{\text{nominal}} = \{k_0, K_{t0}, \mu_0\}\), threshold
limits
\(\{\Delta_{\text{disparity}}, \tau_{\text{deadline}}, E_{\text{kin},\max}\}\).

\textbf{Outputs:} Permitted actuator command
\(\mathbf{u}_{\text{cmd}}\), hardware trip flag
\(\text{register\_trip\_flag}\).

\textbf{Structural Invariants:}

\begin{enumerate}
\def\labelenumi{\arabic{enumi}.}
\tightlist
\item
  \emph{Synchronous Reaction Bound:} Execution loop strictly periodic at
  \(1000\text{ Hz}\) (\(\Delta t = 1.0\text{ ms} \pm 2\,\mu\text{s}\)).
\item
  \emph{Sub-Threshold Probing:} Injected perturbation energy
  \(\|\delta \mathbf{u}\| \le 0.03 \, \|\mathbf{u}_{\text{nominal}}\|\)
  (zero operational disruption).
\item
  \emph{Fail-Safe Primacy:} If estimated plant disparity exceeds
  threshold, trip hardware crowbar within \(\le 500\,\mu\text{s}\).
\end{enumerate}

\textbf{Execution Pipeline:}

\begin{enumerate}
\def\labelenumi{\arabic{enumi}.}
\tightlist
\item
  \textbf{Synchronous Sample \& Timestamp Acquisition:}

  \begin{itemize}
  \tightlist
  \item
    \(\mathbf{y}_t \leftarrow \text{ReadAdcsAndEncodersDma}()\);
    \(t_{\text{now}} \leftarrow \text{ReadMonotonicHardwareTimer}()\).
  \end{itemize}
\item
  \textbf{Epistemic Dither Injection:}

  \begin{itemize}
  \tightlist
  \item
    \(\delta \mathbf{u}_t \leftarrow \text{GenerateOrthogonalDither}(\text{freq}=125\text{ Hz}, \text{amp}=0.03\,u_{\max}, \phi_t)\).
  \item
    \(\mathbf{u}_{\text{probe}} \leftarrow \mathbf{u}_{\text{policy}} + \delta \mathbf{u}_t\).
  \end{itemize}
\item
  \textbf{Empirical Transfer Impedance Extraction:}

  \begin{itemize}
  \tightlist
  \item
    \(\hat{Z}(t) \leftarrow \text{ComputePhaseLockedImpedance}(\delta \mathbf{u}_t, \mathbf{y}_t.\text{high\_rate\_response})\).
  \item
    \(\Delta_{\text{disp}} \leftarrow \|\hat{Z}(t) - \mathcal{M}_{\text{nominal}}.\text{impedance}\|\).
  \end{itemize}
\item
  \textbf{Observational Boundary \& Invariant Evaluation:}

  \begin{itemize}
  \item
    If \(\Delta_{\text{disp}} > \Delta_{\text{disparity}}\):

    \begin{itemize}
    \tightlist
    \item
      \(\mathbf{u}_{\text{cmd}} \leftarrow \text{EnforceDeterministicSafeStop}(\text{profile}=\text{ESTOP\_SPRING\_BRAKE})\).
    \item
      \(\text{register\_trip\_flag} \leftarrow \text{TRUE}\);
      \(\text{SetHardwareFaultRegister}(\text{ACTIVE\_HIGH})\).
    \item
      \(\text{LogClaimsRegisterEvent}(\text{"DEGRADATION\_TRIP"}, \Delta_{\text{disp}}, t_{\text{now}})\).
    \item
      \textbf{return}
      \((\mathbf{u}_{\text{cmd}}, \text{register\_trip\_flag})\).
    \end{itemize}
  \end{itemize}
\item
  \textbf{Dynamic Kinematic Authority Clamping:}

  \begin{itemize}
  \item
    \(E_{\text{kin}} \leftarrow \frac{1}{2} m \, \|\mathbf{v}_t\|^2\).
  \item
    If
    \((E_{\text{kin}} > E_{\text{kin},\max}) \lor (\text{clearance} < d_{\text{stop}}(\mathbf{v}_t, \mu_{\min}))\):

    \begin{itemize}
    \tightlist
    \item
      \(\mathbf{u}_{\text{cmd}} \leftarrow \text{ProjectControlBarrierFunction}(\mathbf{u}_{\text{probe}}, \mathbf{y}_t, \mathcal{C}_{\text{safe}})\).
    \end{itemize}
  \item
    Else:

    \begin{itemize}
    \tightlist
    \item
      \(\mathbf{u}_{\text{cmd}} \leftarrow \mathbf{u}_{\text{probe}}\).
    \end{itemize}
  \item
    \(\text{register\_trip\_flag} \leftarrow \text{FALSE}\).
  \end{itemize}
\item
  \textbf{Register-Level Actuator Drive Write:}

  \begin{itemize}
  \tightlist
  \item
    \(\text{WriteHardwarePWMRegisters}(\mathbf{u}_{\text{cmd}})\).
  \item
    \textbf{return}
    \((\mathbf{u}_{\text{cmd}}, \text{register\_trip\_flag})\).
  \end{itemize}
\end{enumerate}

\end{fbxSimple}

A taxonomy of these unobservable physical failure modes across
transduction, contact mechanics, electromagnetic actuation, and learned
perception is compiled in Table~\ref{tbl-17-detectorless-modes},
highlighting the timing budget deficits and why heuristic software
proxies fail.\sidenote{\footnotesize Stator resistance rises with temperature
  (\(R(T) = R_0 [1 + \alpha_{\text{Cu}}\Delta T]\)), escalating
  \(I^2 R\) heating and slowing emergency braking current decay.}

\begin{longtable}[]{@{}
  >{\raggedright\arraybackslash}p{(\linewidth - 10\tabcolsep) * \real{0.1667}}
  >{\raggedright\arraybackslash}p{(\linewidth - 10\tabcolsep) * \real{0.1667}}
  >{\raggedright\arraybackslash}p{(\linewidth - 10\tabcolsep) * \real{0.1667}}
  >{\raggedright\arraybackslash}p{(\linewidth - 10\tabcolsep) * \real{0.1667}}
  >{\raggedright\arraybackslash}p{(\linewidth - 10\tabcolsep) * \real{0.1667}}
  >{\raggedright\arraybackslash}p{(\linewidth - 10\tabcolsep) * \real{0.1667}}@{}}
\caption{\textbf{Taxonomy of Physical Failure Modes Lacking Real-Time
Runtime Detectors.} Comparison of critical failure modes across
transduction, contact mechanics, electromagnetic drives, and learned
vision-action policies where physical observables are indistinguishable
before the time-to-harm. The table details the physical degradation
mechanism, the root cause of observational blindness, the timing budget
deficit (\(\tau_{\text{det}} + \tau_{\text{act}} > t_{\text{harm}}\)),
the failure of software proxies, and the required physical hardware or
architectural
interceptors.}\label{tbl-17-detectorless-modes}\tabularnewline
\toprule\noalign{}
\begin{minipage}[b]{\linewidth}\raggedright
Failure Mode \& Physical Domain
\end{minipage} & \begin{minipage}[b]{\linewidth}\raggedright
Underlying Physical Degradation Mechanism
\end{minipage} & \begin{minipage}[b]{\linewidth}\raggedright
Root Cause of Observational Indistinguishability
\end{minipage} & \begin{minipage}[b]{\linewidth}\raggedright
Timing Budget Mismatch (\(t_{\text{harm}}\)
vs.~\(\tau_{\text{det}} + \tau_{\text{act}}\))
\end{minipage} & \begin{minipage}[b]{\linewidth}\raggedright
Ineffective Software Proxy
\end{minipage} & \begin{minipage}[b]{\linewidth}\raggedright
Defensible Hardware / Architectural Interceptor
\end{minipage} \\
\midrule\noalign{}
\endfirsthead
\toprule\noalign{}
\begin{minipage}[b]{\linewidth}\raggedright
Failure Mode \& Physical Domain
\end{minipage} & \begin{minipage}[b]{\linewidth}\raggedright
Underlying Physical Degradation Mechanism
\end{minipage} & \begin{minipage}[b]{\linewidth}\raggedright
Root Cause of Observational Indistinguishability
\end{minipage} & \begin{minipage}[b]{\linewidth}\raggedright
Timing Budget Mismatch (\(t_{\text{harm}}\)
vs.~\(\tau_{\text{det}} + \tau_{\text{act}}\))
\end{minipage} & \begin{minipage}[b]{\linewidth}\raggedright
Ineffective Software Proxy
\end{minipage} & \begin{minipage}[b]{\linewidth}\raggedright
Defensible Hardware / Architectural Interceptor
\end{minipage} \\
\midrule\noalign{}
\endhead
\bottomrule\noalign{}
\endlastfoot
\textbf{Piezoresistive Creep}\emph{(Pressure / Strain Transduction)} &
Mechanical crystal fatigue and surface passivation on silicon strain
diaphragm, inducing offset drift \(V_{\text{off}}(t) = V_0 + \kappa t\).
& Analog-to-digital converter digitizes slow monotonic baseline drift
into valid, low-noise codes indistinguishable from genuine process
pressure changes. & Drift accumulates over
\(10^2\text{--}10^3\text{ h}\); process overpressure rupture occurs in
\(t_{\text{harm}} \approx 100\text{ ms}\), while detection latency is
unbounded (\(\tau_{\text{det}} = \infty\)). & Moving-average smoothing
and autoencoder reconstruction residuals classify slow drift as normal
process dynamics. & Dual-dissimilar redundant sensors (piezoresistive +
resonant quartz); mechanical burst discs; scheduled zero-point
proof-tests. \\
\textbf{Non-Smooth Contact Bifurcation}\emph{(High-Speed Manipulation /
Locomotion)} & Rigid-body impact transitions and Coulomb friction
stick-slip boundaries exhibiting discontinuous velocity jumps
(\(\dot{q}^+ \neq \dot{q}^-\)) and Painlevé acceleration non-uniqueness.
& Pre-contact telemetry (\(F_{\text{ext}} = 0\), \(d > 0\)) provides
zero gradient of the non-smooth impact manifold; microscopic surface
asperities are unobservable prior to collision. & Contact-to-fracture
time \(t_{\text{harm}} \le 15\text{--}38\text{ ms}\); sensor filter
(\(\tau = 15\text{ ms}\)) + bus transfer (\(\tau = 2\text{ ms}\)) +
motor inductive decay (\(\tau = 25\text{ ms}\)) yields
\(42\text{ ms} > t_{\text{harm}}\). & High-gain impedance controllers
and neural contact policies lack uniform Lipschitz continuity across
collision manifolds, causing limit cycles or saturation. &
Series-elastic actuators (SEA) providing passive physical compliance;
mechanical spring detents; calibrated breakaway tool mounts. \\
\textbf{Silent Rotor Demagnetization}\emph{(Permanent Magnet Synchronous
Motors)} & NdFeB rotor magnets exceed maximum thermal ceiling
(\(T_{\text{rotor}} > 140^\circ\text{C}\) under stall torque), causing
irreversible flux loss and reduction in torque constant \(K_t\). &
Thermocouples monitor stator windings only; motor drive reports nominal
current \(I_q\), but actual delivered electromagnetic torque
\(\tau = K_t I_q\) drops silently by \(20\text{--}40\%\). &
Demagnetization develops during multi-second thermal soak; catastrophic
load drop occurs in \(t_{\text{harm}} \approx 50\text{ ms}\) upon
gravity loading with zero runtime warning. & Stator-based
lumped-parameter thermal models underestimate rotor eddy-current heating
under high-frequency PWM switching harmonics. & High-Curie-temperature
magnets (\(T_C > 310^\circ\text{C}\)); spring-applied fail-safe
electromagnetic friction brakes (de-energize to lock); motor
over-sizing. \\
\textbf{Latent Epistemic Shift}\emph{(Vision-Language-Action Models)} &
Out-of-distribution optical artifacts (caustic glares, specular
reflections) project onto nominal intermediate activation manifolds in
deep latent layers. & Learned network outputs action distributions with
high confidence (entropy \(H \approx 0\)) and low ensemble variance,
driving actuators along an unsafe trajectory. & Trajectory departure
breaches safe clearance in \(t_{\text{harm}} \approx 120\text{ ms}\);
high-dimensional density estimation suffers the curse of dimensionality
and cannot run within the \(10\text{ ms}\) loop. & Softmax temperature
scaling, Mahalanobis feature distance, and VAE reconstruction loss
exhibit severe false negatives on structured shifts. &
Proposal--Permission architecture; independent deterministic kinematic
safety shield with hardware rangefinders (ToF, sonar) holding veto
authority. \\
\end{longtable}

\section{What Would Change These Answers}\label{sec-frontier-17-5}

Every entry in the residual-claims register marks an unsupported link in
the safety case where argument gives way to assumption. Transforming
these recorded caveats into an engineering agenda requires sorting them
by what is physically or theoretically missing. A recorded gap falls
into one of four classes, defined by the missing element. It stems from
an unmeasured observable when the physical state cannot be distinguished
from sensor data before the deadline for action; from an evaluation
deficit when the tail probability of failure cannot be bounded because
rare events require intractable physical exposure; from a theoretical
deficit when the learned policy lacks a bounded relationship between
training conditions and execution stability; or from a scope mismatch
when the target claim exceeds what the system architecture can
physically support. Naming the missing capability defines the exact
physical or mathematical requirement that would discharge the entry.

Closing an observability gap requires a new physical measurement that
separates previously indistinguishable states within the nervous
system's deadline. For the ceramic grasping failure where an internal
void is undetectable prior to contact, an advance cannot be a software
filter or a learned estimator operating on the same surface sensors. The
advance must introduce a transducer operating on an independent physical
channel, such as an ultrasonic pulse-echo transducer or an active
acoustic resonance detector. To close the gap, the measurement must
satisfy a budgetable engineering threshold. It must resolve the defect
within a latency that leaves sufficient margin for the actuator to halt,
while maintaining a bounded false-alarm rate and demonstrating failure
independence from the primary optical stream. If the transducer takes an
analytical example duration of \(12\text{ ms}\) to measure acoustic
damping, and the mechanical brake halts the jaw in \(10\text{ ms}\), the
total detection-and-arrest budget of \(22\text{ ms}\) fits within the
\(38\text{ ms}\) time-to-fracture. The gap closes because the physical
deadline is met, not because the estimator improved in expectation.

Closing a rare-event evaluation gap requires replacing raw sample counts
with structured estimation whose bias and variance can be analytically
bounded. A safety claim asserting a tail failure rate below
\(10^{-6}\text{ failures/hr}\) cannot be verified by logging
\(10^{6}\text{ hours}\) of unperturbed physical operation, nor by
running \(10^{8}\) simulated episodes whose fidelity in the tail is
unknown. An advance that closes this gap must construct an importance
sampling or adaptive stress-testing scheme that demonstrates
conservative probability estimation on known, analytical benchmark
distributions. The evaluation must report effective independent exposure
rather than raw trial counts. If ten thousand simulated trajectories
explore the identical local basin of a collision surface, the effective
independent sample size is one, not ten thousand. An advance closes the
evaluation gap only when it provides an estimator whose variance
decreases at a verifiable rate and whose assumptions about the tail
geometry are independently checkable.

Closing a theory gap requires a bounded analytical claim paired with a
runtime mechanism that checks its premises. When a policy is trained by
imitation or reinforcement learning, the gap is the absence of an
inductive bias that guarantees bounded outputs under distribution shift.
A theoretical advance cannot close this gap by proving asymptotic
convergence or by deriving bounds that depend on uncheckable Lipschitz
constants across high-dimensional activation spaces. The theory must
state the precise invariance it guarantees, define the support over
which that invariance holds, and provide an operational detector that
flags when physical state trajectories leave that support. If a learned
tracking controller is proven stable assuming the payload mass remains
within \([1.0\text{ kg}, 2.5\text{ kg}]\) and actuator bandwidth exceeds
\(40\text{ Hz}\), the theoretical advance closes the safety link only if
an online estimator monitors payload inertia at \(P_{99}\) intervals and
transitions to a mechanical brake if the mass assumption is violated.

To prevent research optimism from corrupting the safety case, every
candidate advance must be evaluated against a predeclared closure test.
The test defines the missing argument link, the numerical threshold
required for sufficiency, and the independent evidence necessary to
verify it. An advance fails the test if it merely relocates an
assumption to a less visible layer of the system. Replacing a neural
network controller with a control-barrier function whose certificate
assumes perfectly known linear dynamics does not close a theory gap. It
relocates the unmodeled dynamics from the network weights to the barrier
derivative. A candidate advance is stopped before it can alter the
safety verdict whenever its validation relies on the exact unverified
properties it was designed to establish.\sidenote{\footnotesize Richard Hamming
  observed that true reliability arises from architecting systems where
  unmitigated errors are bounded by structural invariants before
  propagating. In Physical AI, this establishes Proposal--Permission
  architecture: learned models propose actions, but nervous systems
  enforce invariant limits.}

Research directions in physical machine learning must be judged solely
by whether they convert an open residual claim into verified evidence
capable of altering the operational verdict. General improvements on
benchmark leaderboards, such as a three percent reduction in average
tracking error or higher reward on a standard task suite, provide zero
evidence toward closing an unsupported safety link. A research project
becomes relevant to physical deployment only when its deliverable
targets a specific entry in the claims register, whether by providing
the missing transducer, bounding the tail estimator, or furnishing the
runtime-checked stability proof. If an advance improves mean performance
but leaves the detector blind spot or the tail uncertainty unchanged,
the authority limits of the machine remain exactly where they were.

Certain gaps cannot be closed by any known sensor, sampling regime, or
mathematical framework. When physical laws, computational limits, or
economic constraints preclude the necessary measurement or proof, the
safety case cannot remain suspended in anticipation of future discovery.
Narrowing the actuator's physical authority, restricting the operational
design domain, or refusing to operate the machine is not an admission of
research defeat. It is the definitive engineering resolution. A machine
whose authority is physically bounded to non-destructive force levels
does not require an ultrasonic sensor to survive an undetected void. The
residual-claims register records the boundary of verifiable knowledge,
and the nervous system enforces the physical constraints that render the
remaining unknown harmless.

\section{Generalizing to Unfamiliar Systems}\label{sec-frontier-17-6}

Applying the framework developed across this book to a machine not
previously examined demonstrates how the four properties, consisting of
physical bounds, information-age budgets, the authority split, and the
evidence record, isolate hazards before hardware is energized. Consider
an automated packaging cell designed for prepared meals. The machine
moves rigid plastic trays along a linear track, presses a heated sealing
head onto tray perimeters to bond plastic barrier film, and injects
pressurized steam into product cavities to blanch ingredients
immediately prior to sealing. This cell exercises all three modalities
of physical consequence within a single enclosure, spanning kinetic mass
transport, mechanical contact force, and thermodynamic energy transfer.
Evaluating this system does not require specialized food-processing
heuristics. It requires evaluating the physical envelope, computing the
timing deadlines imposed by the body, enforcing strict boundaries on
what learned software can touch, and determining what the available
evidence can prove.

Every evaluation begins by cataloging the irreversible state changes the
cell can produce. A loaded food tray traversing a linear transfer rail
carries momentum. When a tray strikes an obstruction, the kinetic energy
is dissipated through plastic deformation of the container and
structural deflection of the frame, a state transition that cannot be
reversed by commanding the drive motor in reverse. The sealing mechanism
introduces a second irreversible boundary. An actuator drives an
aluminum sealing head at \(180^\circ\text{C}\) against the tray flange
to melt polymer bonding layers. If the actuator exerts force beyond the
compressive yield strength of the container wall, the flange buckles,
destroying the hermetic geometry. The steam subsystem introduces a third
irreversible boundary through enthalpy exchange. Saturated steam
injected at \(120^\circ\text{C}\) delivers heat directly into biological
tissue. Once proteins denature and cell walls rupture under excessive
thermal exposure, cooling the chamber cannot restore product texture.
Software retries do not exist in thermodynamic or structural domains.
Once an actuator delivers energy beyond the yield threshold of a
material, the physical state is altered permanently.

Because state changes cannot be undone, the nervous system must
intervene before physical clearances are consumed. In an analytical
baseline, let the linear carriage transport a loaded tray at a nominal
velocity of \(v = 1.5\text{ m/s}\) toward an inspection station whose
mechanical gate provides a clearance margin of
\(d_{\text{clear}} = 60\text{ mm}\) before contacting a rigid stop.
Determining whether an onboard vision system can halt the carriage prior
to an impact requires computing the end-to-end information age across
the full observation-to-enforcement path. Optical image capture and
sensor readout at \(60\text{ Hz}\) consume \(16\text{ ms}\), direct
memory access serialization and image normalization consume
\(8\text{ ms}\), neural inference on an edge accelerator takes
\(25\text{ ms}\) at \(P_{99}\), nervous-system trajectory arbitration
consumes \(3\text{ ms}\), and communication across the controller area
network bus to the drive consumes \(2\text{ ms}\). The motor drive
introduces an additional \(10\text{ ms}\) of current rise time before
the mechanical brake engages, yielding a total information age of
\(\tau = 64\text{ ms}\). Over this interval, the carriage travels
unchecked for a distance of
\(d_{\text{lag}} = v \cdot \tau = (1.5\text{ m/s})(0.064\text{ s}) = 96\text{ mm}\).
If the emergency brake decelerates the carriage at
\(a = 15.0\text{ m/s}^2\), the subsequent braking distance is
\(d_{\text{brake}} = v^2 / (2a) = (1.5\text{ m/s})^2 / (30.0\text{ m/s}^2) = 75\text{ mm}\).
The total stopping envelope is
\(d_{\text{total}} = d_{\text{lag}} + d_{\text{brake}} = 171\text{ mm}\).
Because \(171\text{ mm}\) exceeds the \(60\text{ mm}\) clearance by
\(111\text{ mm}\), the machine cannot be protected against an
obstruction at this velocity by relying on perception software. The
physical deadline is missed before the brake begins to generate torque.

Resolving this hazard requires separating proposal from permission
through an explicit authority boundary. In this cell, learned models
propose actions across all three domains. A vision model infers tray
contents to propose carriage transit speeds, a multimodal network
estimates packaging film thickness to propose sealing forces and dwell
durations, and a thermal estimator proposes steam valve pulse widths
based on incoming product temperature. The nervous system prevents these
proposals from reaching actuators directly. An independent kinematic
governor operating at \(1000\text{ Hz}\) monitors optical encoders along
the track, clipping carriage velocity setpoints so that the stopping
envelope never exceeds the instantaneous clearance measured by hardware
rangefinders. For the sealing head, an analog pressure relief valve
mechanically vents pneumatic pressure above \(220\text{ N}\), preventing
structural flange collapse regardless of what force command the policy
issues. For the steam delivery system, a redundant thermocouple and
hard-wired bimetallic switch interrupt power to the solenoid valve if
chamber temperature exceeds \(105^\circ\text{C}\) or if dwell time
exceeds \(400\text{ ms}\). The learned policy retains authority to
optimize throughput and seal quality within these bounds, but the
nervous system retains sole authority to permit physical execution.

With the physical bounds, information-age budget, and authority split
established, the safety case is compiled into a formal release claim.
The evidence record supports specific operational guarantees through
deterministic test logs and physical proofs. The kinetic hazard is
mitigated by constraining carriage velocity to
\(v \le 0.50\text{ m/s}\), where \(d_{\text{lag}} = 32\text{ mm}\) and
\(d_{\text{brake}} = 8.3\text{ mm}\), yielding a total stopping distance
of \(40.3\text{ mm}\) that fits within the \(60\text{ mm}\) clearance.
Mechanical flange crushing is bounded by the relief valve spring
constant, and gross thermal destruction is bounded by the electrical
cutoff. However, assembling the evidence record also uncovers claims
that remain unsupported. To verify package sterility and seal
hermeticity, the engineering team evaluates a vision model that inspects
optical scattering through the steam chamber to detect micro-tears in
the sealing film. Within the chamber, localized steam condensation
creates optical refraction patterns that mimic film defects while
masking genuine punctures. Because no independent physical sensor exists
in the cell to measure film seal integrity in real time under dense
vapor, this claim has no practical detector. Marking this link as
detectorless and unsupported in the claims register is not an omission
in the documentation, but an accurate statement of what the evidence can
justify.

The verdict permitted by this record divides the machine into approved
and restricted operating modes. The physical cell is released for
automated transport, mechanical sealing, and thermal processing because
all kinetic, crushing, and thermal hazards are bounded by independent
nervous-system mechanisms that do not rely on the learned policy. The
cell is denied release for autonomous packaging sign-off. The machine
cannot certify its own output because the evidence record contains an
unclosed, detectorless link regarding seal integrity under vapor.
Closing this gap cannot be achieved by gathering more training images or
increasing the capacity of the perception model. The verdict can change
only through one of four engineering actions. The team can add an
independent physical transducer, such as a post-cooling ultrasonic leak
detector or downstream vacuum-decay sensor, that verifies seal integrity
after steam clears. They can construct a thermodynamic proof
demonstrating that the mechanical dwell and temperature envelope
guaranteed by the hardware switches ensures a complete seal across all
specified film thicknesses. They can alter the packaging chemistry to
use a self-indicating colorimetric barrier film. Or they can narrow the
operational scope, routing every packaged unit to an external secondary
inspection line before distribution.

This evaluation illustrates the discipline that governs every physical
machine operating under learned control. The engineering process does
not begin with model architecture, training losses, or benchmark scores.
It begins at the physical body, identifying every point where mass,
force, or energy creates an irreversible state change. From those
physical mechanics, the engineer derives the freshness deadlines that
the nervous system must meet, verifies that the end-to-end information
age fits within physical clearances, and enforces a strict authority
split that denies learned software unmediated access to actuators.
Finally, the engineer assembles the evidence record, acknowledging where
mathematical proofs and sensor capabilities end, marking open links
without equivocation, and restricting operational authority to match the
boundaries of what has been proven. Building physical systems that learn
and act requires accepting that gravity, inertia, and thermodynamics do
not yield to statistical expectations.

\section{Architectural Synthesis Across the Physical AI
Stack}\label{sec-frontier-17-7}

Across the seventeen chapters and four parts of this book, this
discipline has unfolded across every layer of the physical AI stack,
synthesized in Figure~\ref{fig-17-causal-loop-synthesis}. In Part I
(Chapter~\ref{sec-boundary} through Chapter~\ref{sec-nervous}), we
established the foundational systems premise: \emph{the Causal
Boundary}. Crossing this boundary marks the transition from decoupled
software living safely behind glass---where errors are caught by runtime
exception handlers and transactions are rolled back without
consequence---to continuous physical reality, where bits modulate motor
drives, accelerate inertial mass, and dissipate irreversible Joule heat.
We decomposed the embodied machine into its three constitutive organs:
the \emph{Body} governed by mechanics and thermodynamics, the
\emph{Brain} built from high-capacity learned policies, and the
deterministic \emph{Nervous System} that enforces real-time timing
deadlines and holds sole authority over safety refusal.

In Part II (Chapter~\ref{sec-data} through
Chapter~\ref{sec-evaluation}), we demonstrated why closing the
sensorimotor loop breaks the foundational i.i.d. assumptions of
classical machine learning. Because every physical action alters the
pose of the hardware and shapes future sensor inputs, observation
streams are endogenous and non-stationary, requiring rigorous
demonstration budgeting, distribution-shift characterization, and
finite-sample exposure arithmetic.

In Part III (Chapter~\ref{sec-perception} through
Chapter~\ref{sec-placement}), we constructed the online real-time
execution engine. We traced perception latencies across direct memory
access transfers, maintained temporal state freshness in structured
memory, negotiated multi-cadence intent horizons, planned kinematically
bounded trajectories, enforced deterministic invariant shields at the
hardware interface, and placed compute across heterogeneous edge silicon
fabrics to satisfy hard physical deadlines.

In Part IV (Chapter~\ref{sec-intervention} through
Chapter~\ref{sec-frontier}), we governed the physical machine in
operation: managing shared human-machine interventions, verifying fault
tolerance through empirical stress testing, compiling
Claim-Argument-Evidence trees into auditable release verdicts, and now,
at the Frontier, confronting what finite empirical evidence can never
settle.

\begin{figure}[t]

\centering{

\includegraphics[width=1\linewidth,height=\textheight,keepaspectratio]{index_files/mediabag//Users/VJ/GitHub/MLSysBook-integrate-physical/physical/book/chapters/17-frontier/figures/fig17_closed_loop_synthesis.pdf}

}

\caption{\label{fig-17-causal-loop-synthesis}\textbf{The Closed Causal
Loop and Epistemic Boundary Synthesis.} Unifying the seventeen chapters
across the four parts of the physical AI curriculum. Cognitive
deliberation in the high-capacity Brain (heterogeneous application MPU /
NPU) emits candidate action trajectory chunks \(\hat{u}_{t:t+H}\) across
the Proposal--Permission causal interface. The deterministic Nervous
System (real-time bare-metal MCU) projects proposals onto Control
Barrier Functions \(h(x) \ge 0\) at \(1000\text{ Hz}\), enforcing
physical stopping distances \(d_{\text{stop}}(v) \le d_{\text{gap}}\)
before commanding actuator currents to the physical plant
(\(W_t \to W_{t+1}\)). Continuous dynamics, Coulomb friction limits
(\(F \le \mu N\)), non-smooth contacts, and thermal limits dissipate
energy without software rollback, closing the loop through sensory
transduction DMA ring buffers back to memory. Chapter~\ref{sec-frontier}
isolates the epistemic frontier: where detectorless blind spots and
astronomical exposure walls
(\citeproc{ref-butler1993infeasibility}{Butler and Finelli 1993};
\citeproc{ref-koopman2016challenges}{Koopman and Wagner 2016}) prevent
empirical certification, the nervous system enforces structural velocity
and force clamps to render the unknown harmless.}

\end{figure}%

The complete architectural synthesis across all seventeen chapters and
four parts is codified in Table~\ref{tbl-17-master-synthesis}, mapping
each layer of the stack to its target organ, foundational systems
problem, bedrock physical law, governing operational invariant, and
non-negotiable hardware floor.

\begin{longtable}[]{@{}
  >{\raggedright\arraybackslash}p{(\linewidth - 10\tabcolsep) * \real{0.1667}}
  >{\raggedright\arraybackslash}p{(\linewidth - 10\tabcolsep) * \real{0.1667}}
  >{\raggedright\arraybackslash}p{(\linewidth - 10\tabcolsep) * \real{0.1667}}
  >{\raggedright\arraybackslash}p{(\linewidth - 10\tabcolsep) * \real{0.1667}}
  >{\raggedright\arraybackslash}p{(\linewidth - 10\tabcolsep) * \real{0.1667}}
  >{\raggedright\arraybackslash}p{(\linewidth - 10\tabcolsep) * \real{0.1667}}@{}}
\caption{\textbf{Master Synthesis of Physical AI Across Seventeen
Chapters and Four Parts.} Comprehensive architectural blueprint mapping
every chapter in the text to its target organ, foundational systems
problem, bedrock physical law, governing operational invariant, and
non-negotiable hardware floor. The taxonomy illustrates how physical
constraints dictate software boundaries from foundational thermodynamic
laws to deployment release
verdicts.}\label{tbl-17-master-synthesis}\tabularnewline
\toprule\noalign{}
\begin{minipage}[b]{\linewidth}\raggedright
Chapter \& Subject
\end{minipage} & \begin{minipage}[b]{\linewidth}\raggedright
Architectural Organ \& Part
\end{minipage} & \begin{minipage}[b]{\linewidth}\raggedright
Core Systems Problem
\end{minipage} & \begin{minipage}[b]{\linewidth}\raggedright
Bedrock Physical Law / Governing Principle
\end{minipage} & \begin{minipage}[b]{\linewidth}\raggedright
Primary Operational Invariant
\end{minipage} & \begin{minipage}[b]{\linewidth}\raggedright
Hardware Floor / Physical Limit
\end{minipage} \\
\midrule\noalign{}
\endfirsthead
\toprule\noalign{}
\begin{minipage}[b]{\linewidth}\raggedright
Chapter \& Subject
\end{minipage} & \begin{minipage}[b]{\linewidth}\raggedright
Architectural Organ \& Part
\end{minipage} & \begin{minipage}[b]{\linewidth}\raggedright
Core Systems Problem
\end{minipage} & \begin{minipage}[b]{\linewidth}\raggedright
Bedrock Physical Law / Governing Principle
\end{minipage} & \begin{minipage}[b]{\linewidth}\raggedright
Primary Operational Invariant
\end{minipage} & \begin{minipage}[b]{\linewidth}\raggedright
Hardware Floor / Physical Limit
\end{minipage} \\
\midrule\noalign{}
\endhead
\bottomrule\noalign{}
\endlastfoot
\textbf{01: Boundary} & Whole Machine \emph{(Part I)} & Non-idempotent
action in continuous physical reality & First \& Second Laws of
Thermodynamics; Conservation of Momentum & Proposal--Permission
separation; No unmediated actuator writes & Actuator kinetic energy
dissipation \& thermal limits \\
\textbf{02: Body} & Body (Plant) \emph{(Part I)} & Actuator bandwidth,
compliance, and torque saturation & Newton-Euler rigid-body dynamics;
Joule heating (\(P = I^2 R\)) & Mechanical impedance \& torque-speed
envelope (\(T \le T_{\text{stall}}\)) & Motor winding thermal limit
(\(T_{\text{max}} \le 130^\circ\text{C}\)); Backlash \\
\textbf{03: Brain} & Brain (Policy) \emph{(Part I)} & High-capacity
representation \& inference latency scaling & Universal approximation;
Roofline compute vs.~memory bandwidth & Action distribution support
(\(\pi(a \mid s) \in \mathcal{U}_{\text{valid}}\)); Latency variance &
Edge GPU/NPU SRAM size \& DRAM memory bandwidth \\
\textbf{04: Nervous} & Nervous System \emph{(Part I)} & Deterministic
real-time timing \& fail-safe arbitration & Hard real-time
schedulability (RM/EDF); Synchronous reactivity & Safety reflex
determinism (\(t_{\text{reflex}} \le t_{\text{deadline}}\)); Shielding &
Deterministic MCU/FPGA execution; Sub-microsecond clock jitter \\
\textbf{05: Data} & Closed Loop \emph{(Part II)} & Non-stationary
endogenous observation distributions & Markov decision process state
visitation distribution (\(d^{\pi}\)) & Covariate shift bounds;
Policy-induced distribution containment & Telemetry bus throughput
(PCIe/MIPI); Storage write endurance \\
\textbf{06: Training} & Policy Optimization \emph{(Part II)} &
Compounding imitation error and distribution drift & Bellman optimality;
Lipschitz bounds under shift (\(O(T^2 \epsilon)\)) & Bounded Bellman
residual; Expert-policy divergence budget & Accelerator inter-node
interconnect bandwidth (NVLink/InfiniBand) \\
\textbf{07: Evaluation} & Empirical Testing \emph{(Part II)} &
Rare-event tail estimation in open physical domains & Law of Large
Numbers; Zero-failure exposure bound
(\(n \ge \frac{\ln \alpha}{\ln(1-p)}\)) & Statistical confidence
certificate (\(1-\alpha\) on tail rate \(p\)) & Finite fleet testing
hours; Accelerated life-test physical bounds \\
\textbf{08: Perception} & Sensing Subsystem \emph{(Part III)} & Sensor
readout latency, jitter, and multi-modal alignment & Shannon-Hartley
theorem; Speed-of-light / acoustic propagation & Temporal freshness
deadline (\(\tau_{\text{age}} \le \tau_{\text{max}}\)); Timestamp skew &
Sensor exposure time \& MIPI CSI-2 DMA serialization pipeline \\
\textbf{09: Memory} & State Representation \emph{(Part III)} & Epistemic
state staleness under continuous ego-motion & Information decay; Kalman
filter covariance propagation (\(\mathbf{P}_{k \mid k-1}\)) & Occupancy
grid spatial freshness; Bounded estimation covariance & DRAM memory
footprint \& cache hierarchy hit latency for voxels \\
\textbf{10: Intent} & Goal Arbitration \emph{(Part III)} & Multi-cadence
planning horizons and intent preemption & Hierarchical control theory;
Reachability analysis & Intent lease validity contract; Bounded
preemption latency & Asynchronous IPC message queue depth \& RTOS
context-switch latency \\
\textbf{11: Planning} & Trajectory Synthesis \emph{(Part III)} &
Kinodynamic feasibility under continuous obstacle constraints &
Pontryagin's Maximum Principle; Kinematic jerk/acceleration bounds &
Collision-free corridor (\(\mathcal{X}_{\text{free}}\)); Curvature bound
(\(\kappa \le \kappa_{\text{max}}\)) & Peak actuator jerk capacity
(\(\dot{a}_{\text{max}}\)); Re-planning rate
(\(10\text{--}50\text{ Hz}\)) \\
\textbf{12: Enforcement} & Runtime Safety Shield \emph{(Part III)} &
Deterministic clipping and safe fallback intervention & Set invariance;
Control Barrier Functions (\(\dot{B} + \alpha(B) \ge 0\)) & Forward
invariance of safe set \(\mathcal{C}\); Hardware veto independence &
Hardware interlock latency (\(\tau \le 1\text{ ms}\)); Analog crowbar /
relay trip \\
\textbf{13: Placement} & Compute Topology \emph{(Part III)} &
Heterogeneous pipeline partitioning across edge fabrics & Amdahl's Law;
Power-Performance-Area (PPA) Pareto frontier & End-to-end
sensor-to-actuator budget (\(\sum \tau_i \le \tau_{\text{harm}}\)) &
Edge thermal design power (TDP \(\le 35\text{ W}\)); Bus communication
delays \\
\textbf{14: Intervention} & Human-Machine Team \emph{(Part IV)} & Mode
confusion, authority handoff, and cognitive takeover lag & Human
sensorimotor latency
(\(t_{\text{takeover}} \approx 0.5\text{--}2.0\text{ s}\)); Shared
control & Unambiguous authority lease; Safe minimal risk maneuver (MRM)
& Physical haptic actuator bandwidth; Human cognitive reaction floor \\
\textbf{15: Verification} & Rigorous Testing \emph{(Part IV)} & Coverage
completeness under non-linear physical dynamics & Formal methods; Hybrid
system reachability; Boundary stress testing & Fault-injection
containment matrix; Pass/fail oracle determinism & Dynamometer \&
Hardware-in-the-Loop test-rig fidelity and bandwidth \\
\textbf{16: Release} & Deployment Governance \emph{(Part IV)} &
Claim-Argument-Evidence auditability \& statutory compliance & Deductive
logic of necessary premises; SIL/ASIL risk calculus & Structured Safety
Case closure; Cryptographic build manifest & Statutory legal liability
bounds; ODD physical envelope limits \\
\textbf{17: Frontier} & Epistemic Governance \emph{(Part IV)} &
Detectorless physical failure modes \& unclosable claims & Popperian
empirical falsification; Observational indistinguishability &
Residual-claims register closure test; Structural containment & Physical
transducer sensitivity limits; Fundamental time-to-harm \\
\end{longtable}

\protect\phantomsection\label{callout-ruleux2a-17.6}
\begin{fbxSimple}{callout-rule}{Architect’s Rule of Thumb:}{Epistemic Containment and Physical AI Governance}
\phantomsection\label{callout-rule*-17.6}

\begin{enumerate}
\def\labelenumi{\arabic{enumi}.}
\tightlist
\item
  \textbf{The Observability Invariant:} If the sensor-to-actuator
  detection latency exceeds the body's time-to-harm
  (\(\tau_{\text{det}} + \tau_{\text{act}} > t_{\text{harm}}\)),
  delegate safety to passive mechanical compliance, firmware parameter
  clamps, and scheduled proof-tests; never rely on learned inference or
  software exception handlers.
\item
  \textbf{The Exposure Invariant:} Never use operational fleet hours to
  certify tail safety below the inverse fleet lifetime
  (\(p \le 1/T_{\text{fleet}}\)); open-world combinatorial safety cannot
  be proven by empirical sampling, but must be guaranteed by
  deterministic nervous-system shields.
\item
  \textbf{The Proposal--Permission Invariant:} High-capacity learned
  Brain models may propose complex actions across multi-cadence intent
  horizons, but the deterministic Nervous System must retain sole and
  non-overridable authority to permit, clip, or veto commands at the
  hardware register level.
\end{enumerate}

\end{fbxSimple}

\section{Epistemic Governance and Operational
Limits}\label{sec-frontier-17-8}

The ultimate maturity of physical AI engineering lies in how an
organization governs the claims its empirical records cannot settle. In
digital software, unhandled edge cases are often treated as minor bugs
to be patched in subsequent agile sprints. In physical systems governing
moving mass, high voltages, or chemical flows, unhandled edge cases
cause permanent material fracture, thermal runaway, and loss of human
life. Statutory safety standards---such as ISO 26262 ASIL D, IEC 61508
SIL 3, and modern autonomous governance standards such as UL 4600 and
ISO 21448 (SOTIF)---explicitly require that unmonitored assumptions be
cataloged, bounded, and audited throughout the entire operational
lifecycle.\sidenote{\footnotesize While ISO 26262/IEC 61508 address hardware and
  software bugs, ISO 21448 (SOTIF) and UL 4600 mandate safety cases
  evaluating performance insufficiencies, sensor blind zones, and
  unknown edge cases, prohibiting deployment based on statistical
  extrapolation.}

When empirical fleet exposure cannot breach the Astronomical Exposure
Wall (\(n \ge \frac{\ln(1/\alpha)}{p}\)), the engineering team must
enforce four non-negotiable operational disciplines:

\begin{enumerate}
\def\labelenumi{\arabic{enumi}.}
\tightlist
\item
  \textbf{Permanent Residual Registers:} Epistemic blind spots are not
  temporary backlog items. Every entry in
  Table~\ref{tbl-17-claims-register} represents a permanent
  architectural constraint until an independent physical measurement or
  proven analytical invariant discharges it.
\item
  \textbf{Deterministic Nervous-System Primacy:} High-capacity neural
  policies running on application processors must never possess
  unmediated actuator write privileges. Actuator drive registers
  (\texttt{TIMx\_CCR}, \texttt{ePWM\ CMPA}) must be governed by
  bare-metal real-time microcontrollers executing deterministic safety
  filters with sub-millisecond reaction times.
\item
  \textbf{Passive Hardware Bounding:} Where observation latencies exceed
  the body's time-to-harm
  (\(\tau_{\text{det}} + \tau_{\text{act}} > t_{\text{harm}}\)), safety
  cannot rely on software. Physical containment must be delegated to
  passive hardware: series-elastic springs, pneumatic relief valves,
  torque-limiting slip clutches, and spring-applied fail-safe brakes.
\item
  \textbf{Accountable Human Signatures:} Operational risk cannot be
  dispersed across an anonymous team. Every release claim with an open
  residual gap must carry the explicit, legally accountable signature of
  named engineering leads who verify that the physical containment
  envelope renders the unknown harmless.
\end{enumerate}

\section{Systems Stress-Tests}\label{sec-frontier-17-9}

With the core epistemic limits, residual-claims register schemas,
detectorless failure mechanisms, and master architectural taxonomy
established, an engineer must be prepared to audit complex physical
systems at the boundary of empirical verification. The following four
systems stress-tests serve as capstone analytical exercises, challenging
you to diagnose observational indistinguishability, detectorless
physical aging, astronomical exposure traps, and latent manifold
shortcuts across embodied platforms:

\subsection*{Stress-Test 1: The Oily Surface Gradient (Observational
Indistinguishability)}\label{stress-test-1-the-oily-surface-gradient-observational-indistinguishability}
\addcontentsline{toc}{subsection}{Stress-Test 1: The Oily Surface
Gradient (Observational Indistinguishability)}

\emph{The scenario.} An autonomous \(600\text{ kg}\) pallet-transport
robot navigates a smooth concrete warehouse concourse at
\(2.2\text{ m/s}\). An unmapped spill of transparent synthetic hydraulic
fluid (\(\mu = 0.10\)) covers a \(3\text{ m}\) patch of floor directly
ahead, presenting zero visual contrast to onboard RGB-D and LiDAR
sensors.

\emph{The systems dilemma.} Before the drive wheels enter the slick,
sensor telemetry is identical to clean dry concrete
(\(y_A(t) \equiv y_B(t)\)). By the time wheel encoders detect
differential angular slip, the kinematic stopping distance on
\(\mu = 0.10\) has expanded from \(0.45\text{ m}\) to \(2.51\text{ m}\),
breaching the \(0.80\text{ m}\) obstacle clearance buffer before the
vehicle can halt.

\emph{The systems verdict: Law of Physical Causality and
Irreversibility.} Safe stopping cannot be certified through vision-based
perception alone when the physical observable (friction coefficient
\(\mu\)) is unmeasurable prior to contact. The engineering authority
must derate maximum velocity to \(v \le 0.8\text{ m/s}\)
(\(E_{\text{kin}} \le 192\text{ J}\)) in mixed-traffic zones or install
forward polarimetric scatterometry, recording the unobservable surface
condition in the residual-claims register.

\subsection*{Stress-Test 2: Silent Thermal Demagnetization (Detectorless
Actuator
Aging)}\label{stress-test-2-silent-thermal-demagnetization-detectorless-actuator-aging}
\addcontentsline{toc}{subsection}{Stress-Test 2: Silent Thermal
Demagnetization (Detectorless Actuator Aging)}

\emph{The scenario.} A 6-DoF robotic manipulator performs high-cadence
pick-and-place operations with heavy steel payloads. Over repeated duty
cycles, rotor eddy currents heat the NdFeB permanent magnets in Joint
3's motor past \(145^\circ\text{C}\), exceeding the magnet's thermal
ceiling and causing an irreversible \(30\%\) reduction in remanent
magnetic flux (\(B_r\)) and torque constant \(K_t\).

\emph{The systems dilemma.} Stator resistance temperature detectors
(RTDs) report nominal stator temperatures (\(75^\circ\text{C}\)), and
motor inverter current shunts report full commanded quadrature current
(\(I_q = 25\text{ A}\)). However, actual output electromagnetic torque
(\(\tau = K_t I_q\)) silently drops from \(50\text{ N}\cdot\text{m}\) to
\(35\text{ N}\cdot\text{m}\). When the arm transitions to cantilevered
gravity holding, the joint sags uncontrollably, striking a human
operator station in \(t_{\text{harm}} = 45\text{ ms}\).

\emph{The systems verdict: Law of Governed Change and Defensible
Release.} Stator-only thermal telemetry cannot detect rotor magnet flux
degradation at runtime. The machine must incorporate
high-Curie-temperature magnets (\(T_C > 310^\circ\text{C}\)),
spring-applied fail-safe friction brakes (de-energize to lock via
hardware interlock lines), and mandatory scheduled offline
torque-constant proof-tests on a dynamometer every 500 operating hours.

\subsection*{Stress-Test 3: The Astronomical Exposure Trap (Certifying
Open-World Edge
Cases)}\label{stress-test-3-the-astronomical-exposure-trap-certifying-open-world-edge-cases}
\addcontentsline{toc}{subsection}{Stress-Test 3: The Astronomical
Exposure Trap (Certifying Open-World Edge Cases)}

\emph{The scenario.} An autonomous delivery vehicle developer logs 10
million accident-free operational miles (\(500{,}000\text{ hours}\))
across multiple urban deployments and submits a safety case claiming a
catastrophic collision probability \(p \le 10^{-8}/\text{hour}\) at
\(95\%\) statistical confidence for its end-to-end multimodal foundation
model.

\emph{The systems dilemma.} Under the zero-failure exposure bound
\(n \ge \frac{\ln(1/\alpha)}{p}\), certifying
\(p \le 10^{-8}/\text{hour}\) at \(95\%\) confidence (\(\alpha = 0.05\))
mathematically requires \(n \ge 3.0 \times 10^8\text{ hours}\) (over
34,200 machine-years) of independent, identically distributed exposure.
Furthermore, periodic over-the-air neural weight updates and correlated
weather patterns across the fleet continuously reset and invalidate
sample independence.

\emph{The systems verdict: Law of Information Freshness \& Zero-Failure
Exposure Bound.} Empirical black-box fleet testing is mathematically and
logically insufficient to certify ultra-rare tail safety in open
combinatorial environments. The release claim must be rejected unless
safety is guaranteed by an independent deterministic nervous-system
shield (CBF-QP enforcer running on an isolated MCU) holding exclusive
hardware veto authority over motor inverters.

\subsection*{Stress-Test 4: The Latent Manifold Shortcut (Specular Glare
in Brittle
Grasping)}\label{stress-test-4-the-latent-manifold-shortcut-specular-glare-in-brittle-grasping}
\addcontentsline{toc}{subsection}{Stress-Test 4: The Latent Manifold
Shortcut (Specular Glare in Brittle Grasping)}

\emph{The scenario.} A high-precision dual-arm manipulator employs an
end-to-end Vision-Language-Action (VLA) model to insert thin glass
pharmaceutical ampoules into packaging trays. Direct late-afternoon
sunlight reflects off a stainless-steel workbench, casting caustic
specular glares across the camera sensor array.

\emph{The systems dilemma.} The high-dimensional visual encoder projects
the out-of-distribution optical flare onto the nominal latent manifold
(\(\mathbf{z} \in \mathcal{Z}_{\text{nom}}\)). The policy commands a
rapid downward plunging trajectory with near-zero softmax entropy
(\(H(\pi) \approx 0\)). Latent-space Mahalanobis distance uncertainty
estimators fail to flag the anomaly because the model is confident in
its corrupted representation. The gripper drives the glass ampoule into
the metal tray at \(300\text{ mm/s}\), shattering the vial in
\(14\text{ ms}\).

\emph{The systems verdict: Law of Proposal and Permission Privilege.}
High-capacity learned models cannot serve as their own safety
supervisors. Actuator commands must pass through a Proposal--Permission
boundary where a deterministic bare-metal nervous-system enforcer
monitors hardware load cells and laser proximity rangefinders at
\(1000\text{ Hz}\), backed by mechanical spring detents that decouple
force before glass fracture limits are reached.

The final frontier of physical AI is not algorithmic perfection, but
epistemic governance: establishing rigorous institutional and
architectural controls over what cannot be known. When an engineering
team acknowledges that open-world complexity, catastrophic tail risks,
and detectorless failure modes preclude absolute empirical certainty,
the safety argument does not collapse into paralysis. Instead, it
transitions from inductive confidence to deductive structural
containment. The governing instrument of this discipline is the
\textbf{Residual-Claims Register}---an immutable ledger that explicitly
records every unclosable claim, unmodeled physical parameter, and
unobservable state transition, pairing each with a permanent,
hardware-enforced containment boundary.

Where software cannot observe, hardware must constrain. If friction
coefficients on unmapped terrain cannot be guaranteed, velocity must be
physically clamped so that stopping distance remains within the shortest
possible line-of-sight clearance. If internal structural voids in
handled materials cannot be detected before contact, robotic
end-effectors must incorporate passive mechanical compliance---such as
spring-loaded flexures or torque-limiting slip clutches---that limit
peak impact forces below destructive thresholds regardless of software
commands. If neural network latent spaces cannot guarantee
out-of-distribution detection, the machine must operate behind physical
exclusion barriers or within certified geofenced corridors. Structural
containment renders epistemic uncertainty harmless by ensuring that when
the unknown occurs, the physical energy released remains strictly below
the threshold of catastrophic damage.

Ultimately, epistemic governance is a human accountability contract. An
autonomous system, no matter how sophisticated its neural architecture
or how extensive its training corpus, is a mathematical artifact lacking
legal personhood, moral agency, and the capacity to bear liability. When
a physical AI system damages property, disrupts critical infrastructure,
or harms a human being, responsibility cannot be diffused into the
opacity of high-dimensional latent spaces. The Cryptographic Release
Manifest established in Chapter~\ref{sec-release} binds the operational
authority of the machine directly to the cryptographic signatures of
designated systems engineers and adjudicators. Epistemic humility is
therefore not merely a philosophical virtue, but the core ethical and
professional discipline of physical AI systems engineering: recognizing
the limits of empirical proof, bounding physical authority through
deterministic hardware shields, and accepting personal accountability
for the machines we release into the physical world.

The synthesis of physical AI systems engineering is illustrated across
the three pillars in Figure~\ref{fig-17-closed-loop-synthesis}: the
\textbf{Brain} (which explores, plans, and optimizes across open-world
representations), the \textbf{Nervous System} (which arbitrates
authority, verifies intent leases, and enforces real-time barrier
certificates), and the \textbf{Body} (which interacts with continuum
mechanics under strict conservation laws).

\begin{figure}[t]

\centering{

\includegraphics[width=0.95\linewidth,height=\textheight,keepaspectratio]{index_files/mediabag//Users/VJ/GitHub/MLSysBook-integrate-physical/physical/book/chapters/17-frontier/figures/fig17_closed_loop_synthesis.pdf}

}

\caption{\label{fig-17-closed-loop-synthesis}\textbf{The Tripartite
Architecture of Physical AI Systems Engineering.} High-level synthesis
illustrating the fundamental division of responsibility across the three
physical AI pillars. (Left) The Brain operates in the unprivileged
cognitive domain, synthesizing open-vocabulary goals, volumetric memory,
and continuous trajectory proposals under soft real-time deadlines.
(Center) The Nervous System operates in the privileged, deterministic
domain, enforcing hard real-time CBF-QP safety shields, evaluating
intent leases, and executing hardware trip-zone vetoes at
\(1000\text{ Hz}\). (Right) The Body executes forces and torques under
immutable physical conservation laws (time, energy, mass, momentum,
thermal dissipation). Safety is achieved not by hoping the Brain never
fails, but by guaranteeing that the Nervous System strictly bounds the
kinetic authority of the Body when the Brain inevitably encounters the
unknown. Attribution: Original textbook diagram (CC BY 4.0).}

\end{figure}%

\subsection*{Capstone Synthesis: The Three
Archetypes}\label{capstone-synthesis-the-three-archetypes}
\addcontentsline{toc}{subsection}{Capstone Synthesis: The Three
Archetypes}

To close the curriculum, we revisit the Three Archetypes introduced in
earlier chapters, demonstrating how the complete physical AI framework
governs real hardware across diverse physical embodiments:

\begin{enumerate}
\def\labelenumi{\arabic{enumi}.}
\tightlist
\item
  \textbf{The Autonomous Mobile Manipulator (Mobile Archetype):}

  \begin{itemize}
  \tightlist
  \item
    \emph{Physical Substrate:} Differential-drive AGV with a 7-DOF
    articulated arm (\(m_{\text{total}} = 185\text{ kg}\),
    \(v_{\max} = 1.8\text{ m/s}\), payload \(10\text{ kg}\)).
  \item
    \emph{The Failure Mode:} Navigating a warehouse corridor under
    changing lighting, a vision-language policy misidentifies a thin,
    oil-slicked metal ramp as dry concrete, commanding a high-speed
    turning trajectory.
  \item
    \emph{The Systems Defense:} The exteroceptive Brain is blind to the
    low-friction hazard (observational indistinguishability). However,
    the deterministic Nervous System evaluates wheel slip ratios and IMU
    yaw rates at \(1000\text{ Hz}\). Detecting a micro-slip anomaly
    (\(s > 0.08\)) within \(4\text{ ms}\), the real-time MCU safety
    shield instantly overrides the policy's torque vectoring, clamps
    maximum velocity to \(0.4\text{ m/s}\) (firmware velocity derating),
    and engages \(\mathcal{C}^1\) bumpless deceleration along a
    straight-line stopping corridor.
  \item
    \emph{The Systems Verdict: Law of Structural Containment.} Safety
    did not depend on the VLA model recognizing oil; safety was
    guaranteed by the nervous system's deterministic traction barrier
    invariant.
  \end{itemize}
\item
  \textbf{The High-Speed Collaborative Delta Robot (Manipulator
  Archetype):}

  \begin{itemize}
  \tightlist
  \item
    \emph{Physical Substrate:} 4-axis parallel kinematic delta robot
    (\(v_{\text{tool}} \le 4.0\text{ m/s}\), cycle rate
    \(120\text{ picks/min}\)) sorting food products alongside human
    packing personnel.
  \item
    \emph{The Failure Mode:} A worker reaches into the shared pick zone
    to clear a jammed carton. The robot's learned object-tracking head
    experiences a transient \(45\text{ ms}\) GPU kernel scheduling
    stall, causing an intent lease timeout.
  \item
    \emph{The Systems Defense:} The nervous system's windowed hardware
    watchdog detects heartbeat absence within \(3.0\text{ ms}\). The
    Simplex switch trips, immediately transferring motor authority from
    the stalled Linux MPU to a precomputed quintic deceleration spline
    resident in microcontroller SRAM TCM. Actuators smoothly arrest
    motion within \(18\text{ ms}\) (\(d_{\text{stop}} = 36\text{ mm}\)),
    strictly respecting ISO/TS 15066 transient impact force limits
    (\(F \le 50\text{ N}\)).
  \item
    \emph{The Systems Verdict: Law of Asymmetric Authority.} The
    high-level neural policy proposes; the bare-metal safety enforcer
    permits. When deliberation stalls, the nervous system enforces
    zero-command containment without waiting for software recovery.
  \end{itemize}
\item
  \textbf{The Continuous Chemical Reaction Furnace (Process Archetype):}

  \begin{itemize}
  \tightlist
  \item
    \emph{Physical Substrate:} Exothermic batch reactor with
    variable-speed reactant injection pumps and proportional cooling
    jacket valves (\(P_{\max} = 4.5\text{ MPa}\),
    \(T_{\max} = 450^\circ\text{C}\)).
  \item
    \emph{The Failure Mode:} An adaptive neural optimization policy
    adjusts reactant feed rates to maximize yield. A sudden supply batch
    variance introduces an unmodeled catalytic promoter, initiating
    thermal runaway with temperature rising at
    \(\dot{T} = 12.5^\circ\text{C/s}\).
  \item
    \emph{The Systems Defense:} The neural optimizer misinterprets the
    rapid temperature rise as a productive yield surge. However, an
    independent Control Barrier Function running on an isolated lockstep
    safety controller continuously evaluates the thermal-hydraulic
    barrier \(h(T, P) = T_{\max} - T - \tau_{\text{cool}}\dot{T}\). The
    moment \(h(T, P) \le 0\), the enforcer asserts hardwired Safe Torque
    Off (STO) to the feed pumps and commands 100\% duty cycle to the
    emergency quench coolant valve via direct memory-mapped PWM
    registers, bypassing the Linux optimizer entirely.
  \item
    \emph{The Systems Verdict: Law of Proposal and Permission
    Privilege.} High-capacity learned models cannot serve as their own
    safety supervisors. Actuator commands must pass through a
    Proposal--Permission boundary where a deterministic bare-metal
    nervous-system enforcer monitors hardware load cells and laser
    proximity rangefinders at \(1000\text{ Hz}\), backed by mechanical
    spring detents that decouple force before glass fracture limits are
    reached.
  \end{itemize}
\end{enumerate}

\section{Summary}\label{sec-frontier-17-10}

The central question posed at the opening of this book---\emph{why does
an error that costs nothing on a screen destroy hardware in the physical
world?}---finds its ultimate resolution at this frontier. Software
operating behind glass enjoys the luxury of idempotency. An embodied
machine operating in the physical world is bound by conservation laws:
delay is paid in physical displacement (\(d = v \tau\)), kinetic energy
cannot be recalled once released, and unmonitored physical invariants
drift regardless of model confidence. The frontier of physical AI is
therefore not a catalog of speculative algorithms, but the permanent
boundary where measurable evidence ends and structural containment
begins. A machine is safe to operate not when its learned models appear
competent, but when its nervous system ensures that failure remains
physically bounded.

\protect\phantomsection\label{callout-takeawaysux2a-17.7}
\begin{fbxSimple}{callout-takeaways}{Key Chapter Takeaways}{}
\phantomsection\label{callout-takeaways*-17.7}

\begin{itemize}
\tightlist
\item
  \textbf{The Epistemic Boundary is Permanent (Residual-Claims
  Register):} Empirical verification can never test the full
  combinatorial open world. Every physical AI release requires an
  explicit Residual-Claims Register that identifies unclosable
  hypotheses and binds each to a deterministic physical or mechanical
  containment mechanism.
\item
  \textbf{Observational Indistinguishability Demands Structural
  Containment:} When safe and catastrophic physical states emit
  identical pre-harm sensor signals (\(y_A(t) \equiv y_B(t)\)) until the
  reaction window is exhausted
  (\(\tau_{\text{det}} + \tau_{\text{act}} > t_{\text{harm}}\)),
  software cannot prevent damage. Safety must be enforced through
  passive mechanical compliance, slip clutches, or velocity derating.
\item
  \textbf{The Astronomical Exposure Wall Precludes Black-Box Statistical
  Proofs:} Validating an ultra-reliable safety claim
  (\(p \le 10^{-8}/\text{hour}\)) requires
  \(>3.0 \times 10^8\text{ hours}\) (34,200 machine-years) of
  failure-free exposure---a mathematical impossibility. High-assurance
  physical safety rests exclusively on architectural isolation and
  deterministic runtime enforcers.
\item
  \textbf{The Proposal--Permission Architectural Invariant:} Across all
  embodiments (mobile robots, manipulators, continuous chemical plants),
  high-capacity neural foundation models propose actions (\(p_t\)),
  while bare-metal real-time microcontrollers hold exclusive hardware
  veto authority over power inverters (\(1000\text{ Hz}\)).
\item
  \textbf{The Sovereign Master Law of Physical AI:} Software behind
  glass enjoys idempotency; in the physical universe, computational
  delay is paid in millimeters (\(\Delta x = \int v \, dt\)), kinetic
  energy cannot be recalled, and safety is achieved only when failure
  remains physically bounded.
\end{itemize}

\end{fbxSimple}

\cleardoublepage
\phantomsection
\addcontentsline{toc}{part}{Appendices}
\appendix

\chapter*{Mathematical and Physical Foundations for Systems
Engineers}\label{sec-appendix-math}
\addcontentsline{toc}{chapter}{Mathematical and Physical Foundations for
Systems Engineers}

\markboth{Mathematical and Physical Foundations for Systems
Engineers}{Mathematical and Physical Foundations for Systems Engineers}

This appendix provides a self-contained, systems-oriented reference for
the mathematical and control-theoretic formulations used across this
book.

For systems, embedded, and software engineers, mathematical formulas in
Physical AI are not abstract analytical exercises; they are the
\textbf{specifications for real-time silicon execution}. Every equation
maps directly into memory access patterns, execution cycle budgets
(\(\text{WCET}\)), numerical precision requirements, and physical safety
bounds.

\begin{center}\rule{0.5\linewidth}{0.5pt}\end{center}

\section*{1. Kinematics and Stopping
Envelopes}\label{kinematics-and-stopping-envelopes}
\addcontentsline{toc}{section}{1. Kinematics and Stopping Envelopes}

\markright{1. Kinematics and Stopping Envelopes}

\subsection*{The Physical Derivation of Stopping
Distance}\label{the-physical-derivation-of-stopping-distance}
\addcontentsline{toc}{subsection}{The Physical Derivation of Stopping
Distance}

\lettrine{W}{hen} an embodied robot maneuvers at linear velocity
\(v(t)\), the distance required to bring its physical mass to a
complete, controlled stop consists of two decoupled components:

\begin{enumerate}
\def\labelenumi{\arabic{enumi}.}
\tightlist
\item
  \textbf{Reaction Drift (\(d_{\text{drift}}\)):} The distance traversed
  while sensory signals propagate through transduction, bus
  serialization, neural inference, and communication queues before the
  actuators can apply braking torque.
\item
  \textbf{Kinematic Braking (\(d_{\text{brake}}\)):} The distance
  traversed while physical friction or counter-electromotive braking
  forces decelerate the moving mass at maximum safe deceleration
  \(a_{\text{max}}\).
\end{enumerate}

\begin{longtable}[]{@{}
  >{\raggedright\arraybackslash}p{(\linewidth - 6\tabcolsep) * \real{0.2500}}
  >{\raggedright\arraybackslash}p{(\linewidth - 6\tabcolsep) * \real{0.2500}}
  >{\raggedright\arraybackslash}p{(\linewidth - 6\tabcolsep) * \real{0.2500}}
  >{\raggedright\arraybackslash}p{(\linewidth - 6\tabcolsep) * \real{0.2500}}@{}}
\caption{Physical Decomposition of Total Stopping Distance
(\(d_{\text{stop}} = d_{\text{drift}} + d_{\text{brake}}\)).}\label{tbl-stopping-distance-decomp}\tabularnewline
\toprule\noalign{}
\begin{minipage}[b]{\linewidth}\raggedright
Phase
\end{minipage} & \begin{minipage}[b]{\linewidth}\raggedright
Physical Mechanism
\end{minipage} & \begin{minipage}[b]{\linewidth}\raggedright
Distance Formulation
\end{minipage} & \begin{minipage}[b]{\linewidth}\raggedright
Time Horizon
\end{minipage} \\
\midrule\noalign{}
\endfirsthead
\toprule\noalign{}
\begin{minipage}[b]{\linewidth}\raggedright
Phase
\end{minipage} & \begin{minipage}[b]{\linewidth}\raggedright
Physical Mechanism
\end{minipage} & \begin{minipage}[b]{\linewidth}\raggedright
Distance Formulation
\end{minipage} & \begin{minipage}[b]{\linewidth}\raggedright
Time Horizon
\end{minipage} \\
\midrule\noalign{}
\endhead
\bottomrule\noalign{}
\endlastfoot
\textbf{Phase 1: Reaction Drift} & Sensor exposure, neural inference,
network transport, inverter rise &
\(d_{\text{drift}} = v \cdot \tau_{\text{latency}}\) &
\(t \in [0, \tau_{\text{latency}}]\) \\
\textbf{Phase 2: Kinematic Braking} & Counter-EMF braking, friction,
ground adhesion & \(d_{\text{brake}} = \dfrac{v^2}{2 a_{\text{max}}}\) &
\(t \in [\tau_{\text{latency}}, t_{\text{stop}}]\) \\
\end{longtable}

Integrating Newtonian dynamics under constant deceleration
\(a(t) = -a_{\text{max}}\):

\[v(t) = v_0 - a_{\text{max}} t \implies t_{\text{brake}} = \frac{v_0}{a_{\text{max}}}\]

\[d_{\text{brake}} = \int_0^{t_{\text{brake}}} (v_0 - a_{\text{max}} t) \, dt = v_0 t_{\text{brake}} - \frac{1}{2} a_{\text{max}} t_{\text{brake}}^2 = \frac{v_0^2}{2 a_{\text{max}}}\]

Accounting for total sensorimotor transport latency
\(\tau_{\text{latency}}\), the \textbf{Total Dynamic Stopping Envelope}
is:

\[d_{\text{stop}}(v) = \underbrace{v \cdot \tau_{\text{latency}}}_{\text{Reaction Drift}} + \underbrace{\frac{v^2}{2 a_{\text{max}}}}_{\text{Kinematic Braking}}\]

where: * \(v\): Current operational velocity (\(\text{m/s}\)). *
\(\tau_{\text{latency}}\): Total end-to-end transport delay
(\(\text{s}\)), encompassing camera exposure, bus serialization
(CAN-FD/PCIe), operating system scheduling tail latency, and inverter
current rise time. * \(a_{\text{max}}\): Maximum permissible physical
braking deceleration (\(\text{m/s}^2\)), constrained by tire friction
coefficient (\(\mu g\)), payload tipping limits, or motor peak torque.

\subsection*{Parabolic Safe Velocity
Clamping}\label{parabolic-safe-velocity-clamping}
\addcontentsline{toc}{subsection}{Parabolic Safe Velocity Clamping}

Given a measured physical clearance \(d_{\text{clearance}}\) to an
obstacle and a safety margin \(d_{\text{margin}}\), the maximum safe
operational velocity \(v_{\text{safe}}\) is derived by solving the
quadratic inequality
\(d_{\text{stop}}(v) \le d_{\text{clearance}} - d_{\text{margin}}\):

\[\frac{1}{2 a_{\text{max}}} v^2 + \tau_{\text{latency}} v - (d_{\text{clearance}} - d_{\text{margin}}) \le 0\]

Applying the quadratic formula and selecting the positive root:

\[v_{\text{safe}} = \sqrt{a_{\text{max}}^2 \tau_{\text{latency}}^2 + 2 a_{\text{max}} (d_{\text{clearance}} - d_{\text{margin}})} - a_{\text{max}} \tau_{\text{latency}}\]

\subsection*{Microcontroller Execution Complexity and Timing
Bounds}\label{microcontroller-execution-complexity-and-timing-bounds}
\addcontentsline{toc}{subsection}{Microcontroller Execution Complexity
and Timing Bounds}

In real-time firmware, evaluating \(v_{\text{safe}}\) and
\(d_{\text{stop}}(v)\) executes at \(1000\text{ Hz}\) in static SRAM
with zero dynamic memory allocation (\texttt{malloc\ =\ 0}). Because the
formulation requires only a single square-root operation and basic
floating-point arithmetic, its Worst-Case Execution Time (WCET) on a
\(480\text{ MHz}\) ARM Cortex-M7 with a hardware Floating Point Unit
(FPU) is bounded to \(\text{WCET} \le 1.2\,\mu\text{s}\), consuming less
than \(0.15\%\) of the \(1.0\text{ ms}\) control interval.

\begin{center}\rule{0.5\linewidth}{0.5pt}\end{center}

\section*{2. Spatial Geometry and Lie
Algebra}\label{spatial-geometry-and-lie-algebra}
\addcontentsline{toc}{section}{2. Spatial Geometry and Lie Algebra}

\markright{2. Spatial Geometry and Lie Algebra}

Physical AI agents operate in three-dimensional Euclidean space
\(\mathbb{R}^3\). To transform sensor observations, object affordances,
and robot kinematics across coordinate frames, systems rely on the
Special Euclidean Group \(SE(3)\).

\subsection*{Homogeneous Transformation
Matrices}\label{homogeneous-transformation-matrices}
\addcontentsline{toc}{subsection}{Homogeneous Transformation Matrices}

A rigid-body transformation from coordinate frame \(A\) to coordinate
frame \(B\) consists of a 3D rotation \(\mathbf{R}_A^B \in SO(3)\) and a
3D translation \(\mathbf{p}_A^B \in \mathbb{R}^3\), unified into a
\(4 \times 4\) homogeneous matrix \(\mathbf{T}_A^B \in SE(3)\):

\[\mathbf{T}_A^B = \begin{bmatrix} \mathbf{R}_A^B & \mathbf{p}_A^B \\ \mathbf{0}_{1 \times 3} & 1 \end{bmatrix} \in \mathbb{R}^{4 \times 4}\]

Points in 3D space are transformed using homogeneous coordinates
\(\tilde{\mathbf{x}} = [\mathbf{x}^T, 1]^T\):

\[\tilde{\mathbf{x}}_B = \mathbf{T}_A^B \tilde{\mathbf{x}}_A \implies \begin{bmatrix} \mathbf{x}_B \\ 1 \end{bmatrix} = \begin{bmatrix} \mathbf{R}_A^B \mathbf{x}_A + \mathbf{p}_A^B \\ 1 \end{bmatrix}\]

Matrix inversion in \(SE(3)\) has a closed-form structure that avoids
general \(4 \times 4\) Gaussian elimination:

\[(\mathbf{T}_A^B)^{-1} = \mathbf{T}_B^A = \begin{bmatrix} (\mathbf{R}_A^B)^T & -(\mathbf{R}_A^B)^T \mathbf{p}_A^B \\ \mathbf{0}_{1 \times 3} & 1 \end{bmatrix}\]

\subsection*{Quaternions for Singularity-Free
Rotations}\label{quaternions-for-singularity-free-rotations}
\addcontentsline{toc}{subsection}{Quaternions for Singularity-Free
Rotations}

While Euler angles (Roll, Pitch, Yaw) suffer from kinematic
singularities (\textbf{Gimbal Lock} when pitch approaches
\(\pm 90^\circ\), causing division by zero in coordinate
transformations), unit quaternions provide a continuous, computationally
robust 4-parameter representation:

\[\mathbf{q} = [w, x, y, z]^T = [\cos(\theta / 2), \mathbf{u}^T \sin(\theta / 2)]^T, \quad \|\mathbf{q}\|_2^2 = w^2 + x^2 + y^2 + z^2 = 1\]

where \(\mathbf{u} \in \mathbb{R}^3\) is the unit rotation axis and
\(\theta\) is the rotation angle.

The \textbf{Hamilton Quaternion Product}
(\(\mathbf{q}_c = \mathbf{q}_a \otimes \mathbf{q}_b\)) compounds
successive 3D rotations:

\[\mathbf{q}_a \otimes \mathbf{q}_b = \begin{bmatrix}
w_a w_b - x_a x_b - y_a y_b - z_a z_b \\
w_a x_b + x_a w_b + y_a z_b - z_a y_b \\
w_a y_b - x_a z_b + y_a w_b + z_a x_b \\
w_a z_b + x_a y_b - y_a x_b + z_a w_b
\end{bmatrix}\]

\subsection*{Lie Algebra and Tangent Space
Integration}\label{lie-algebra-and-tangent-space-integration}
\addcontentsline{toc}{subsection}{Lie Algebra and Tangent Space
Integration}

In high-frequency control loops (\(1000\text{ Hz}\)), velocities are
expressed as continuous \textbf{spatial twists}
\(\boldsymbol{\xi} \in \mathbb{R}^6\) combining linear velocity
\(\mathbf{v} \in \mathbb{R}^3\) and angular velocity
\(\boldsymbol{\omega} \in \mathbb{R}^3\):

\[\boldsymbol{\xi} = \begin{bmatrix} \mathbf{v} \\ \boldsymbol{\omega} \end{bmatrix} \in \mathbb{R}^6\]

The wedge operator \((\cdot)^\wedge\) maps
\(\boldsymbol{\xi} \in \mathbb{R}^6\) to the matrix Lie algebra
\(\mathfrak{se}(3)\):

\[\boldsymbol{\xi}^\wedge = \begin{bmatrix} [\boldsymbol{\omega}]_\times & \mathbf{v} \\ \mathbf{0}_{1 \times 3} & 0 \end{bmatrix} \in \mathbb{R}^{4 \times 4}, \quad \text{where } [\boldsymbol{\omega}]_\times = \begin{bmatrix} 0 & -\omega_z & \omega_y \\ \omega_z & 0 & -\omega_x \\ -\omega_y & \omega_x & 0 \end{bmatrix}\]

Integrating a constant twist \(\boldsymbol{\xi}\) over a discrete time
step \(\Delta t\) updates the robot pose without numerical distortion
using the \textbf{Exponential Map}:

\[\mathbf{T}(t + \Delta t) = \mathbf{T}(t) \cdot \exp\left( \boldsymbol{\xi}^\wedge \Delta t \right)\]

\[\exp\left( \boldsymbol{\xi}^\wedge \Delta t \right) = \begin{bmatrix} \exp([\boldsymbol{\omega}]_\times \Delta t) & \mathbf{V}\mathbf{v}\Delta t \\ \mathbf{0}_{1 \times 3} & 1 \end{bmatrix}\]

where the left-Jacobian matrix \(\mathbf{V}\) and rotation matrix
\(\exp([\boldsymbol{\omega}]_\times \Delta t)\) are evaluated in closed
form via Rodrigues' formula with
\(\theta = \|\boldsymbol{\omega}\| \Delta t\):

\[\exp([\boldsymbol{\omega}]_\times \Delta t) = \mathbf{I}_{3 \times 3} + \frac{\sin(\theta)}{\theta}[\boldsymbol{\omega}]_\times + \frac{1 - \cos(\theta)}{\theta^2}[\boldsymbol{\omega}]_\times^2\]

\[\mathbf{V} = \mathbf{I}_{3 \times 3} + \frac{1 - \cos(\theta)}{\theta^2}[\boldsymbol{\omega}]_\times + \frac{\theta - \sin(\theta)}{\theta^3}[\boldsymbol{\omega}]_\times^2\]

\subsection*{Numerical Normalization
Invariant}\label{numerical-normalization-invariant}
\addcontentsline{toc}{subsection}{Numerical Normalization Invariant}

Under discrete integration, floating-point round-off causes unit
quaternions to drift (\(\|\mathbf{q}\| \neq 1\)), corrupting the
\(SO(3)\) rotation manifold. In real-time \(1000\text{ Hz}\) state
estimators, the quaternion is renormalized every cycle
(\(\mathbf{q} \gets \mathbf{q} / \|\mathbf{q}\|\)), falling back to
identity \([1, 0, 0, 0]^T\) if \(\|\mathbf{q}\|^2 < 10^{-12}\),
executing in \(< 0.1\,\mu\text{s}\) with zero dynamic memory allocation.

\begin{center}\rule{0.5\linewidth}{0.5pt}\end{center}

\section*{3. Control Barrier Functions and Quadratic
Programs}\label{control-barrier-functions-and-quadratic-programs}
\addcontentsline{toc}{section}{3. Control Barrier Functions and
Quadratic Programs}

\markright{3. Control Barrier Functions and Quadratic Programs}

\subsection*{Forward Invariance and Nagumo's
Theorem}\label{forward-invariance-and-nagumos-theorem}
\addcontentsline{toc}{subsection}{Forward Invariance and Nagumo's
Theorem}

Consider a control-affine physical dynamical system:

\[\dot{\mathbf{x}} = f(\mathbf{x}) + g(\mathbf{x})\mathbf{u}\]

where \(\mathbf{x} \in \mathbb{R}^n\) is the physical state (positions,
velocities), \(\mathbf{u} \in \mathbb{R}^m\) is the control input (motor
torques, phase voltages), and \(f: \mathbb{R}^n \to \mathbb{R}^n\) and
\(g: \mathbb{R}^n \to \mathbb{R}^{n \times m}\) describe the autonomous
drift dynamics and control input mapping. Both \(f\) and \(g\) are
assumed to be locally Lipschitz continuous.

A physical safe operating set \(\mathcal{C} \subset \mathbb{R}^n\) is
defined as the 0-superlevel set of a continuously differentiable scalar
barrier certificate function
\(h(\mathbf{x}): \mathbb{R}^n \to \mathbb{R}\):

\[\mathcal{C} = \{ \mathbf{x} \in \mathbb{R}^n \mid h(\mathbf{x}) \ge 0 \}\]

\[\partial \mathcal{C} = \{ \mathbf{x} \in \mathbb{R}^n \mid h(\mathbf{x}) = 0 \} \quad (\text{Safety Boundary})\]

\[\text{Int}(\mathcal{C}) = \{ \mathbf{x} \in \mathbb{R}^n \mid h(\mathbf{x}) > 0 \} \quad (\text{Safe Interior})\]

The safe set \(\mathcal{C}\) is \textbf{forward invariant} if for every
initial physical state starting within the set
(\(\mathbf{x}(0) \in \mathcal{C}\)), the trajectory remains safe for all
future time: \(\mathbf{x}(t) \in \mathcal{C}, \forall t \ge 0\).

\subsubsection*{Formal Proof of Forward Invariance via Nagumo's
Condition}\label{formal-proof-of-forward-invariance-via-nagumos-condition}
\addcontentsline{toc}{subsubsection}{Formal Proof of Forward Invariance
via Nagumo's Condition}

Nagumo's theorem for set invariance states that a closed set
\(\mathcal{C}\) is forward invariant under the vector field
\(\dot{\mathbf{x}} = F(\mathbf{x})\) if and only if the velocity vector
\(F(\mathbf{x})\) points into the tangent cone (subtangent vectors) of
\(\mathcal{C}\) at every boundary point
\(\mathbf{x} \in \partial \mathcal{C}\):

\[\forall \mathbf{x} \in \partial \mathcal{C}, \quad \nabla h(\mathbf{x}) \cdot F(\mathbf{x}) \ge 0\]

To guarantee asymptotic stability of the safe set and avoid edge
chatter, we enforce the generalized Control Barrier Function condition
with an extended class-\(\mathcal{K}\) function \(\alpha(\cdot)\) (where
\(\alpha(0) = 0\) and \(\alpha\) is strictly increasing):

\[\dot{h}(\mathbf{x}, \mathbf{u}) = \nabla h(\mathbf{x}) [f(\mathbf{x}) + g(\mathbf{x})\mathbf{u}] \ge -\alpha(h(\mathbf{x}))\]

Choosing the canonical linear class-\(\mathcal{K}\) form
\(\alpha(h) = \gamma h\) with rate parameter \(\gamma > 0\), the
differential inequality becomes:

\[\frac{d h(\mathbf{x}(t))}{dt} \ge -\gamma h(\mathbf{x}(t))\]

Applying the Comparison Lemma (Grönwall's Inequality) to this scalar
differential inequality yields:

\[h(\mathbf{x}(t)) \ge h(\mathbf{x}(0)) e^{-\gamma t}, \quad \forall t \ge 0\]

Because \(h(\mathbf{x}(0)) \ge 0\) for any state starting in
\(\mathcal{C}\) and the exponential decay term \(e^{-\gamma t} > 0\) for
all finite \(t \in [0, \infty)\), it follows that:

\[h(\mathbf{x}(t)) \ge 0, \quad \forall t \ge 0\]

Thus, the state trajectory \(\mathbf{x}(t)\) can never cross the zero
boundary \(\partial \mathcal{C}\), proving rigorous mathematical forward
set invariance.

\subsection*{Lie Derivatives and Affine Control
Invariants}\label{lie-derivatives-and-affine-control-invariants}
\addcontentsline{toc}{subsection}{Lie Derivatives and Affine Control
Invariants}

Expanding \(\dot{h}(\mathbf{x}, \mathbf{u})\) via the chain rule:

\[\dot{h}(\mathbf{x}, \mathbf{u}) = \nabla h(\mathbf{x}) \dot{\mathbf{x}} = \nabla h(\mathbf{x}) [f(\mathbf{x}) + g(\mathbf{x})\mathbf{u}] = \underbrace{\nabla h(\mathbf{x}) f(\mathbf{x})}_{L_f h(\mathbf{x})} + \underbrace{\nabla h(\mathbf{x}) g(\mathbf{x})}_{L_g h(\mathbf{x})} \mathbf{u}\]

\begin{itemize}
\tightlist
\item
  \textbf{\(L_f h(\mathbf{x}) \in \mathbb{R}\):} The \emph{Lie
  derivative along \(f\)}, representing the natural physical rate of
  change of safety clearance under zero actuator power (unforced
  momentum drift).
\item
  \textbf{\(L_g h(\mathbf{x}) \in \mathbb{R}^{1 \times m}\):} The
  \emph{Lie derivative along \(g\)}, representing the control authority
  over safety clearance per unit actuator input.
\end{itemize}

The Control Barrier Function condition becomes an \textbf{affine linear
inequality} in the control input vector \(\mathbf{u}\):

\[L_f h(\mathbf{x}) + L_g h(\mathbf{x})\mathbf{u} + \gamma h(\mathbf{x}) \ge 0\]

\begin{itemize}
\tightlist
\item
  \textbf{The Speed Governor Invariant:} When far inside the safe set
  (\(h(\mathbf{x}) \gg 0\)), \(-\gamma h(\mathbf{x})\) evaluates to a
  large negative number, leaving the candidate control input
  \(\mathbf{u}\) unconstrained. As physical clearance collapses toward
  the boundary (\(h(\mathbf{x}) \to 0\)),
  \(-\gamma h(\mathbf{x}) \to 0\), forcing the approach speed
  \(\dot{h} \ge 0\). At the exact boundary (\(h=0\)), any action that
  produces negative clearance rate is strictly vetoed, rendering the
  boundary mathematically impenetrable.
\end{itemize}

\subsection*{The Minimal-Intervention Quadratic Program
(CBF-QP)}\label{the-minimal-intervention-quadratic-program-cbf-qp}
\addcontentsline{toc}{subsection}{The Minimal-Intervention Quadratic
Program (CBF-QP)}

In the Proposal--Permission architecture, upstream neural policies
(System 2 / System 1.5) emit unverified candidate action proposals
\(\mathbf{u}_{\text{cand}} \in \mathbb{R}^m\). The bare-metal
microcontroller (System 1) solves the following convex Quadratic Program
at \(1000\text{ Hz}\):

\[\mathbf{u}^* = \arg\min_{\mathbf{u} \in \mathbb{R}^m} \frac{1}{2} (\mathbf{u} - \mathbf{u}_{\text{cand}})^T \mathbf{W} (\mathbf{u} - \mathbf{u}_{\text{cand}})\]

\[\begin{aligned}
\text{subject to:} \quad &L_f h_i(\mathbf{x}) + L_g h_i(\mathbf{x})\mathbf{u} + \gamma_i h_i(\mathbf{x}) \ge 0, \quad i = 1, \dots, M \\
&\mathbf{u}_{\text{min}} \le \mathbf{u} \le \mathbf{u}_{\text{max}}
\end{aligned}\]

where \(\mathbf{W} \succ 0\) is a positive-definite weighting metric.

\begin{itemize}
\tightlist
\item
  \textbf{Minimal Intervention Principle:} When
  \(\mathbf{u}_{\text{cand}}\) satisfies all barrier inequalities, the
  objective is minimized at \(\mathbf{u}^* = \mathbf{u}_{\text{cand}}\)
  with zero active constraints, passing the AI policy's action
  unmodified. When \(\mathbf{u}_{\text{cand}}\) violates a barrier, the
  QP executes an orthogonal projection onto the safe half-space boundary
  \(\partial \mathcal{U}_{\text{safe}}\), modifying the command by the
  minimum Euclidean perturbation required to preserve invariance.
\end{itemize}

\subsection*{Multi-Variable KKT Matrix System and Active-Set
Derivation}\label{multi-variable-kkt-matrix-system-and-active-set-derivation}
\addcontentsline{toc}{subsection}{Multi-Variable KKT Matrix System and
Active-Set Derivation}

Writing the CBF-QP in standard matrix inequality form:

\[\min_{\mathbf{u}} \frac{1}{2} \mathbf{u}^T \mathbf{W} \mathbf{u} - \mathbf{u}_{\text{cand}}^T \mathbf{W} \mathbf{u} \quad \text{subject to} \quad \mathbf{A} \mathbf{u} \ge \mathbf{b}\]

where the constraint matrix \(\mathbf{A} \in \mathbb{R}^{K \times m}\)
and vector \(\mathbf{b} \in \mathbb{R}^K\) stack both the \(M\) barrier
constraints and the \(2m\) box limits (\(K = M + 2m\)):

\[\mathbf{A} = \begin{bmatrix} L_g h_1(\mathbf{x}) \\ \vdots \\ L_g h_M(\mathbf{x}) \\ \mathbf{I}_{m \times m} \\ -\mathbf{I}_{m \times m} \end{bmatrix}, \qquad \mathbf{b} = \begin{bmatrix} -L_f h_1(\mathbf{x}) - \gamma_1 h_1(\mathbf{x}) \\ \vdots \\ -L_f h_M(\mathbf{x}) - \gamma_M h_M(\mathbf{x}) \\ \mathbf{u}_{\text{min}} \\ -\mathbf{u}_{\text{max}} \end{bmatrix}\]

The Lagrangian function for this convex optimization problem is:

\[\mathcal{L}(\mathbf{u}, \boldsymbol{\lambda}) = \frac{1}{2} (\mathbf{u} - \mathbf{u}_{\text{cand}})^T \mathbf{W} (\mathbf{u} - \mathbf{u}_{\text{cand}}) - \boldsymbol{\lambda}^T (\mathbf{A} \mathbf{u} - \mathbf{b})\]

where \(\boldsymbol{\lambda} \in \mathbb{R}^K\) is the vector of
non-negative Lagrange multipliers. The Karush-Kuhn-Tucker (KKT)
optimality conditions are:

\begin{enumerate}
\def\labelenumi{\arabic{enumi}.}
\tightlist
\item
  \textbf{Stationarity:}
  \(\nabla_{\mathbf{u}} \mathcal{L} = \mathbf{W} (\mathbf{u}^* - \mathbf{u}_{\text{cand}}) - \mathbf{A}^T \boldsymbol{\lambda}^* = \mathbf{0} \implies \mathbf{u}^* = \mathbf{u}_{\text{cand}} + \mathbf{W}^{-1} \mathbf{A}^T \boldsymbol{\lambda}^*\)
\item
  \textbf{Primal Feasibility:}
  \(\mathbf{A} \mathbf{u}^* - \mathbf{b} \ge \mathbf{0}\)
\item
  \textbf{Dual Feasibility:} \(\boldsymbol{\lambda}^* \ge \mathbf{0}\)
\item
  \textbf{Complementary Slackness:}
  \(\lambda_k^* (\mathbf{a}_k^T \mathbf{u}^* - b_k) = 0, \quad \forall k \in \{1, \dots, K\}\)
\end{enumerate}

Let \(\mathcal{A} \subseteq \{1, \dots, K\}\) denote the active
constraint set at the optimal solution (where
\(\mathbf{a}_k^T \mathbf{u}^* = b_k\) and \(\lambda_k^* > 0\)).
Restricting \(\mathbf{A}\) and \(\mathbf{b}\) to the active rows
\(\mathbf{A}_{\mathcal{A}}\):

\[\mathbf{A}_{\mathcal{A}} \mathbf{u}^* = \mathbf{A}_{\mathcal{A}} \left( \mathbf{u}_{\text{cand}} + \mathbf{W}^{-1} \mathbf{A}_{\mathcal{A}}^T \boldsymbol{\lambda}_{\mathcal{A}}^* \right) = \mathbf{b}_{\mathcal{A}}\]

Rearranging yields the symmetric positive-definite active-set linear
system for the optimal dual variables:

\[\left( \mathbf{A}_{\mathcal{A}} \mathbf{W}^{-1} \mathbf{A}_{\mathcal{A}}^T \right) \boldsymbol{\lambda}_{\mathcal{A}}^* = \mathbf{b}_{\mathcal{A}} - \mathbf{A}_{\mathcal{A}} \mathbf{u}_{\text{cand}}\]

Solving for \(\boldsymbol{\lambda}_{\mathcal{A}}^*\) via Cholesky
decomposition
\(\mathbf{L} \mathbf{L}^T = \mathbf{A}_{\mathcal{A}} \mathbf{W}^{-1} \mathbf{A}_{\mathcal{A}}^T\),
the optimal projected control command is:

\[\mathbf{u}^* = \mathbf{u}_{\text{cand}} + \mathbf{W}^{-1} \mathbf{A}_{\mathcal{A}}^T \boldsymbol{\lambda}_{\mathcal{A}}^*\]

\subsection*{Microcontroller Real-Time Solvers and Complexity
Bounds}\label{microcontroller-real-time-solvers-and-complexity-bounds}
\addcontentsline{toc}{subsection}{Microcontroller Real-Time Solvers and
Complexity Bounds}

On safety-critical microcontrollers (such as a dual-core lockstep ARM
Cortex-M7 at \(480\text{ MHz}\)), solving the active-set CBF-QP is
executed using static memory allocations in Data Tightly-Coupled Memory
(DTCM) with zero heap fragmentation (\texttt{malloc\ =\ 0}).

For a 6-DoF manipulator with \(M=8\) active barrier certificates, the
active-set algorithm resolves the optimal projection in \(3\text{--}6\)
iterations. Each iteration performs a Cholesky factorization of the
small active-set Gram matrix
\((\mathbf{A}_{\mathcal{A}} \mathbf{W}^{-1} \mathbf{A}_{\mathcal{A}}^T) \in \mathbb{R}^{|\mathcal{A}| \times |\mathcal{A}|}\)
(where \(|\mathcal{A}| \le 6\)). On target silicon, nominal solve time
is under \(45\,\mu\text{s}\), with a hard worst-case execution time
bounded to \(\text{WCET} \le 175\,\mu\text{s}\) at maximum iteration
count (\(k_{\text{max}} = 10\)). If the active-set fails to converge
within \(10\) iterations, the microcontroller deterministically aborts
to the Category 1 dynamic stopping trajectory (\(a = -a_{\text{max}}\))
within the remaining \(825\,\mu\text{s}\) of the \(1.0\text{ ms}\)
control cycle.

\begin{center}\rule{0.5\linewidth}{0.5pt}\end{center}

\section*{4. Continuous-Discrete Extended Kalman
Filtering}\label{continuous-discrete-extended-kalman-filtering}
\addcontentsline{toc}{section}{4. Continuous-Discrete Extended Kalman
Filtering}

\markright{4. Continuous-Discrete Extended Kalman Filtering}

Physical AI pipelines fuse high-rate (\(1000\text{ Hz}\)) noisy
proprioceptive inertial measurements with delayed, multi-modal
exteroceptive observations (\(10\text{--}30\text{ Hz}\)) using the
\textbf{Continuous-Discrete Extended Kalman Filter (EKF)}.

\subsection*{Continuous Physics Error Dynamics and Riccati ODE
Derivation}\label{continuous-physics-error-dynamics-and-riccati-ode-derivation}
\addcontentsline{toc}{subsection}{Continuous Physics Error Dynamics and
Riccati ODE Derivation}

Between discrete measurement arrivals, the continuous physical state
evolves under non-linear dynamics with continuous process noise:

\[\dot{\mathbf{x}}(t) = f(\mathbf{x}(t), \mathbf{u}(t)) + \mathbf{w}(t), \qquad \mathbf{w}(t) \sim \mathcal{N}(\mathbf{0}, \mathbf{Q}_c(t))\]

The estimated state mean propagates via the nominal dynamics:

\[\dot{\hat{\mathbf{x}}}(t) = f(\hat{\mathbf{x}}(t), \mathbf{u}(t))\]

Defining estimation error
\(\mathbf{e}(t) = \mathbf{x}(t) - \hat{\mathbf{x}}(t)\) and linearizing
\(f(\mathbf{x}, \mathbf{u})\) via Taylor expansion about
\(\hat{\mathbf{x}}(t)\):

\[f(\mathbf{x}(t), \mathbf{u}(t)) = f(\hat{\mathbf{x}}(t), \mathbf{u}(t)) + \mathbf{F}(t)\mathbf{e}(t) + \mathcal{O}(\|\mathbf{e}\|^2), \qquad \mathbf{F}(t) = \left. \frac{\partial f}{\partial \mathbf{x}} \right|_{\hat{\mathbf{x}}(t), \mathbf{u}(t)}\]

The continuous error dynamics satisfy:

\[\dot{\mathbf{e}}(t) = \dot{\mathbf{x}}(t) - \dot{\hat{\mathbf{x}}}(t) = \mathbf{F}(t)\mathbf{e}(t) + \mathbf{w}(t)\]

The spatial error covariance matrix is defined as the expectation
\(\mathbf{P}(t) = \mathbb{E}[\mathbf{e}(t)\mathbf{e}(t)^T]\). Applying
the product rule for differentiation:

\[\dot{\mathbf{P}}(t) = \frac{d}{dt}\mathbb{E}[\mathbf{e}(t)\mathbf{e}(t)^T] = \mathbb{E}[\dot{\mathbf{e}}(t)\mathbf{e}(t)^T] + \mathbb{E}[\mathbf{e}(t)\dot{\mathbf{e}}(t)^T]\]

Substituting the error differential equation:

\[\mathbb{E}[\dot{\mathbf{e}}(t)\mathbf{e}(t)^T] = \mathbb{E}[(\mathbf{F}(t)\mathbf{e}(t) + \mathbf{w}(t))\mathbf{e}(t)^T] = \mathbf{F}(t)\mathbf{P}(t) + \mathbb{E}[\mathbf{w}(t)\mathbf{e}(t)^T]\]

Because white Gaussian process noise \(\mathbf{w}(t)\) is uncorrelated
with past estimation errors,
\(\mathbb{E}[\mathbf{w}(t)\mathbf{e}(t)^T] = \frac{1}{2}\mathbf{Q}_c(t)\).
Adding the transpose term yields the \textbf{Continuous Matrix Riccati
Differential Equation}:

\[\dot{\mathbf{P}}(t) = \mathbf{F}(t)\mathbf{P}(t) + \mathbf{P}(t)\mathbf{F}(t)^T + \mathbf{Q}_c(t)\]

Discretizing across integration time step \(\Delta t = 1.0\text{ ms}\):

\[\hat{\mathbf{x}}_{k|k-1} = \hat{\mathbf{x}}_{k-1|k-1} + \int_{t_{k-1}}^{t_k} f(\hat{\mathbf{x}}(\tau), \mathbf{u}_{k-1})\,d\tau \approx \hat{\mathbf{x}}_{k-1|k-1} + f(\hat{\mathbf{x}}_{k-1|k-1}, \mathbf{u}_{k-1})\Delta t\]

\[\mathbf{P}_{k|k-1} = \boldsymbol{\Phi}_k \mathbf{P}_{k-1|k-1} \boldsymbol{\Phi}_k^T + \mathbf{Q}_d\]

where
\(\boldsymbol{\Phi}_k \approx \mathbf{I} + \mathbf{F}(t_{k-1})\Delta t\)
is the state transition matrix and
\(\mathbf{Q}_d \approx \mathbf{Q}_c \Delta t\) is discrete process noise
covariance.

\subsection*{Optimal Kalman Gain Derivation via Minimum Trace
Estimation}\label{optimal-kalman-gain-derivation-via-minimum-trace-estimation}
\addcontentsline{toc}{subsection}{Optimal Kalman Gain Derivation via
Minimum Trace Estimation}

When a discrete measurement \(\mathbf{z}_k \in \mathbb{R}^m\) arrives at
time \(t_k\):

\[\mathbf{z}_k = h(\mathbf{x}_k) + \mathbf{v}_k, \qquad \mathbf{v}_k \sim \mathcal{N}(\mathbf{0}, \mathbf{R}_k)\]

The linear estimator updates the state estimate according to:

\[\hat{\mathbf{x}}_{k|k} = \hat{\mathbf{x}}_{k|k-1} + \mathbf{K}_k (\mathbf{z}_k - h(\hat{\mathbf{x}}_{k|k-1}))\]

Linearizing the observation function
\(h(\mathbf{x}_k) \approx h(\hat{\mathbf{x}}_{k|k-1}) + \mathbf{H}_k \mathbf{e}_{k|k-1}\)
with
\(\mathbf{H}_k = \left. \frac{\partial h}{\partial \mathbf{x}} \right|_{\hat{\mathbf{x}}_{k|k-1}}\),
the a posteriori error is:

\[\mathbf{e}_{k|k} = \mathbf{x}_k - \hat{\mathbf{x}}_{k|k} = (\mathbf{I} - \mathbf{K}_k \mathbf{H}_k)\mathbf{e}_{k|k-1} - \mathbf{K}_k \mathbf{v}_k\]

Evaluating the a posteriori covariance
\(\mathbf{P}_{k|k} = \mathbb{E}[\mathbf{e}_{k|k}\mathbf{e}_{k|k}^T]\)
under mutually independent state error and sensor noise:

\[\mathbf{P}_{k|k} = (\mathbf{I} - \mathbf{K}_k \mathbf{H}_k)\mathbf{P}_{k|k-1}(\mathbf{I} - \mathbf{K}_k \mathbf{H}_k)^T + \mathbf{K}_k \mathbf{R}_k \mathbf{K}_k^T\]

To find the gain \(\mathbf{K}_k\) minimizing the mean squared estimation
error, we minimize the scalar cost function
\(J = \text{Tr}(\mathbf{P}_{k|k})\). Expanding the trace:

\[J = \text{Tr}(\mathbf{P}_{k|k-1}) - 2\text{Tr}(\mathbf{K}_k \mathbf{H}_k \mathbf{P}_{k|k-1}) + \text{Tr}\left(\mathbf{K}_k (\mathbf{H}_k \mathbf{P}_{k|k-1} \mathbf{H}_k^T + \mathbf{R}_k)\mathbf{K}_k^T\right)\]

Applying matrix differentiation identities
\(\frac{\partial \text{Tr}(\mathbf{A}\mathbf{B})}{\partial \mathbf{A}} = \mathbf{B}^T\)
and
\(\frac{\partial \text{Tr}(\mathbf{A}\mathbf{C}\mathbf{A}^T)}{\partial \mathbf{A}} = 2\mathbf{A}\mathbf{C}\)
(for symmetric \(\mathbf{C}\)):

\[\frac{\partial J}{\partial \mathbf{K}_k} = -2 (\mathbf{H}_k \mathbf{P}_{k|k-1})^T + 2 \mathbf{K}_k (\mathbf{H}_k \mathbf{P}_{k|k-1} \mathbf{H}_k^T + \mathbf{R}_k) = \mathbf{0}\]

Defining innovation covariance
\(\mathbf{S}_k = \mathbf{H}_k \mathbf{P}_{k|k-1} \mathbf{H}_k^T + \mathbf{R}_k\):

\[\mathbf{K}_k \mathbf{S}_k = \mathbf{P}_{k|k-1} \mathbf{H}_k^T \implies \mathbf{K}_k = \mathbf{P}_{k|k-1} \mathbf{H}_k^T \mathbf{S}_k^{-1}\]

\subsection*{Structural Proof of the Joseph Stabilized Covariance
Form}\label{structural-proof-of-the-joseph-stabilized-covariance-form}
\addcontentsline{toc}{subsection}{Structural Proof of the Joseph
Stabilized Covariance Form}

On 32-bit embedded architectures (FP32), the simplified covariance
update \(\mathbf{P} = (\mathbf{I} - \mathbf{K}\mathbf{H})\mathbf{P}\)
causes catastrophic divergence due to numerical asymmetry and lost
positive definiteness.

The \textbf{Joseph stabilized form} guarantees positive
semi-definiteness:

\[\mathbf{P}_{k|k} = (\mathbf{I} - \mathbf{K}_k \mathbf{H}_k) \mathbf{P}_{k|k-1} (\mathbf{I} - \mathbf{K}_k \mathbf{H}_k)^T + \mathbf{K}_k \mathbf{R}_k \mathbf{K}_k^T\]

\textbf{Mathematical Proof of Global Stability:} Let
\(\mathbf{y} \in \mathbb{R}^n\) be an arbitrary non-zero test vector.
The quadratic form evaluates to:

\[\mathbf{y}^T \mathbf{P}_{k|k} \mathbf{y} = \mathbf{y}^T (\mathbf{I} - \mathbf{K}_k \mathbf{H}_k) \mathbf{P}_{k|k-1} (\mathbf{I} - \mathbf{K}_k \mathbf{H}_k)^T \mathbf{y} + \mathbf{y}^T \mathbf{K}_k \mathbf{R}_k \mathbf{K}_k^T \mathbf{y}\]

Defining transformed vectors
\(\mathbf{u}_1 = (\mathbf{I} - \mathbf{K}_k \mathbf{H}_k)^T \mathbf{y}\)
and \(\mathbf{u}_2 = \mathbf{K}_k^T \mathbf{y}\):

\[\mathbf{y}^T \mathbf{P}_{k|k} \mathbf{y} = \mathbf{u}_1^T \mathbf{P}_{k|k-1} \mathbf{u}_1 + \mathbf{u}_2^T \mathbf{R}_k \mathbf{u}_2\]

Because \(\mathbf{P}_{k|k-1} \succeq 0\) and \(\mathbf{R}_k \succ 0\):

\[\mathbf{u}_1^T \mathbf{P}_{k|k-1} \mathbf{u}_1 \ge 0, \qquad \mathbf{u}_2^T \mathbf{R}_k \mathbf{u}_2 \ge 0 \implies \mathbf{y}^T \mathbf{P}_{k|k} \mathbf{y} \ge 0 \quad \forall \mathbf{y} \in \mathbb{R}^n\]

Thus, \(\mathbf{P}_{k|k}\) is strictly positive semi-definite
(\(\mathbf{P}_{k|k} \succeq 0\)) regardless of floating-point truncation
errors or suboptimal gain matrices.

\subsection*{Mahalanobis Innovation
Gating}\label{mahalanobis-innovation-gating}
\addcontentsline{toc}{subsection}{Mahalanobis Innovation Gating}

The measurement innovation
\(\tilde{\mathbf{y}}_k = \mathbf{z}_k - h(\hat{\mathbf{x}}_{k|k-1})\)
represents the new information introduced by the sensor. Under Gaussian
assumptions,
\(\tilde{\mathbf{y}}_k \sim \mathcal{N}(\mathbf{0}, \mathbf{S}_k)\).

The squared Mahalanobis distance scales the innovation residual by its
expected uncertainty:

\[d_M^2 = \tilde{\mathbf{y}}_k^T \mathbf{S}_k^{-1} \tilde{\mathbf{y}}_k \sim \chi_m^2\]

where \(m = \dim(\mathbf{z}_k)\). The gating threshold is evaluated
against the cumulative Chi-Square distribution:

\[\text{Decision} = \begin{cases} \text{Accept Measurement and Apply Update}, & d_M^2 \le \chi_{m, 1-\alpha}^2 \\\\ \text{Reject Outlier and Propagate Open-Loop}, & d_M^2 > \chi_{m, 1-\alpha}^2 \end{cases}\]

\subsection*{Real-Time Kalman Filter Determinism and
Complexity}\label{real-time-kalman-filter-determinism-and-complexity}
\addcontentsline{toc}{subsection}{Real-Time Kalman Filter Determinism
and Complexity}

In embedded state estimation pipelines, the Continuous-Discrete EKF
executes on dedicated microcontroller hardware or real-time DSPs in
static DTCM SRAM with zero dynamic memory allocation
(\texttt{malloc\ =\ 0}).

For a 6-dimensional kinematic state (\(n=6\)) with 3-dimensional
position observations (\(m=3\)), the propagation step requires
\(O(n^3)\) operations (\(216\) multiply-accumulates), and the
Joseph-form measurement update requires matrix inversion of the
\(3 \times 3\) innovation covariance \(\mathbf{S}_k\), followed by outer
product updates
\(\mathbf{P}_k = (\mathbf{I} - \mathbf{K}_k\mathbf{H}_k)\mathbf{P}_k^-(\mathbf{I} - \mathbf{K}_k\mathbf{H}_k)^T + \mathbf{K}_k\mathbf{R}_k\mathbf{K}_k^T\).

On an ARM Cortex-M7 with hardware SIMD and single-precision FPU, the
complete propagation and Joseph update cycle executes in under
\(18.5\,\mu\text{s}\) (\(\text{WCET} \le 25.0\,\mu\text{s}\)), ensuring
that state estimates and uncertainty bounds are updated with
deterministic microsecond precision before feeding downstream Control
Barrier Functions.

\begin{center}\rule{0.5\linewidth}{0.5pt}\end{center}

\section*{5. Field-Oriented Motor
Control}\label{field-oriented-motor-control}
\addcontentsline{toc}{section}{5. Field-Oriented Motor Control}

\markright{5. Field-Oriented Motor Control}

To translate high-level torque commands into physical electromagnetic
forces, brushless permanent-magnet synchronous motors (PMSM) rely on
\textbf{Field-Oriented Control (FOC)} executing in dedicated hardware
timer peripherals at \(20\text{ kHz}\) (\(\Delta t = 50\,\mu\text{s}\)).

\subsection*{Clarke Transformation}\label{clarke-transformation}
\addcontentsline{toc}{subsection}{Clarke Transformation}

Transforms balanced three-phase stator currents \((i_a, i_b, i_c)\) into
a stationary two-phase orthogonal coordinate system:

\[\begin{bmatrix} i_\alpha \\ i_\beta \end{bmatrix} = \frac{2}{3} \begin{bmatrix} 1 & -\frac{1}{2} & -\frac{1}{2} \\ 0 & \frac{\sqrt{3}}{2} & -\frac{\sqrt{3}}{2} \end{bmatrix} \begin{bmatrix} i_a \\ i_b \\ i_c \end{bmatrix}\]

Using the zero-sum property \(i_a + i_b + i_c = 0\), this simplifies to:

\[i_\alpha = i_a, \quad i_\beta = \frac{1}{\sqrt{3}} (i_a + 2 i_b)\]

\subsection*{Park Transformation}\label{park-transformation}
\addcontentsline{toc}{subsection}{Park Transformation}

Rotates the stationary \((\alpha, \beta)\) currents into the rotor's
rotating reference frame using the electrical rotor angle
\(\theta_e = P \cdot \theta_{\text{mech}}\) (where \(P\) is the motor
pole-pair count):

\[\begin{bmatrix} i_d \\ i_q \end{bmatrix} = \begin{bmatrix} \cos\theta_e & \sin\theta_e \\ -\sin\theta_e & \cos\theta_e \end{bmatrix} \begin{bmatrix} i_\alpha \\ i_\beta \end{bmatrix}\]

\begin{itemize}
\tightlist
\item
  \textbf{Direct-Axis Current (\(i_d\)):} Aligns parallel with the rotor
  permanent magnet flux. Setting \(i_d^* = 0\) eliminates wasteful
  radial forces and maximizes electrical efficiency.
\item
  \textbf{Quadrature-Axis Current (\(i_q\)):} Aligns \(90^\circ\)
  electrical ahead of rotor flux, producing pure orthogonal
  electromagnetic torque:
\end{itemize}

\[\tau_{\text{motor}} = \frac{3}{2} P \lambda_{\text{pm}} i_q = K_t \cdot i_q\]

where \(\lambda_{\text{pm}}\) is the permanent magnet flux linkage
(\(\text{Wb}\)) and \(K_t\) is the motor torque constant
(\(\text{N}\cdot\text{m/A}\)).

\subsection*{Thermal Energy
Accumulator}\label{thermal-energy-accumulator}
\addcontentsline{toc}{subsection}{Thermal Energy Accumulator}

To prevent motor stator insulation breakdown
(\(T > 155^\circ\text{C}\)), the motor driver integrates Joulean thermal
dissipation in real time:

\[E_{\text{thermal}}(t) = \int_0^t I_{\text{phase}}^2(\tau) R_{\text{stator}} \, d\tau\]

When \(E_{\text{thermal}}\) exceeds the continuous thermal rating
\(E_{\text{continuous}}\), the controller linearly derates peak torque
current
\(i_q^{\text{limit}} \leftarrow i_q^{\text{nominal}} \cdot \left(1 - \frac{T_j - 85^\circ\text{C}}{155^\circ\text{C} - 85^\circ\text{C}}\right)\).

\begin{center}\rule{0.5\linewidth}{0.5pt}\end{center}

\section*{6. Frequency-Domain Stability and Phase
Margin}\label{frequency-domain-stability-and-phase-margin}
\addcontentsline{toc}{section}{6. Frequency-Domain Stability and Phase
Margin}

\markright{6. Frequency-Domain Stability and Phase Margin}

In classical digital control, converting continuous commands to discrete
hardware outputs introduces uncompensated phase lag that erodes the
system's stability margin.

\subsection*{Zero-Order Hold (ZOH) Phase
Lag}\label{zero-order-hold-zoh-phase-lag}
\addcontentsline{toc}{subsection}{Zero-Order Hold (ZOH) Phase Lag}

A digital-to-analog converter holding a commanded setpoint constant over
sample period \(T\) is modeled by the Zero-Order Hold transfer function:

\[G_{\text{ZOH}}(s) = \frac{1 - e^{-sT}}{s}\]

Evaluating its frequency response \(s = j\omega\):

\[G_{\text{ZOH}}(j\omega) = T \cdot \text{sinc}\left(\frac{\omega T}{2}\right) e^{-j \frac{\omega T}{2}}\]

The magnitude attenuation is
\(\left| \text{sinc}\left(\frac{\omega T}{2}\right) \right|\), while the
\textbf{exact phase lag} is:

\[\Delta \phi_{\text{ZOH}} = -\frac{\omega T}{2} \quad (\text{radians}) = -\omega \cdot \frac{T}{2} \cdot \frac{180^\circ}{\pi} \quad (\text{degrees})\]

\subsection*{The Phase Margin Stability
Criterion}\label{the-phase-margin-stability-criterion}
\addcontentsline{toc}{subsection}{The Phase Margin Stability Criterion}

Let \(\omega_c\) be the open-loop gain crossover frequency
(\(\|G(j\omega_c)\| = 1\)) of the physical mechanism, and
\(\text{PM}_{\text{nominal}}\) be the nominal phase margin. Total
sensorimotor transport delay
\(\tau_{\text{delay}} = \tau_{\text{compute}} + \frac{1}{2} T_{\text{sample}}\)
introduces phase lag:

\[\Delta \phi = -\omega_c \cdot \tau_{\text{delay}} \cdot \frac{180^\circ}{\pi}\]

The effective closed-loop phase margin is:

\[\text{PM}_{\text{effective}} = \text{PM}_{\text{nominal}} - \omega_c \tau_{\text{delay}} \cdot \frac{180^\circ}{\pi}\]

\begin{itemize}
\tightlist
\item
  \textbf{Stability Limit:} If
  \(\tau_{\text{delay}} \ge \frac{\text{PM}_{\text{nominal}}}{\omega_c} \cdot \frac{\pi}{180^\circ}\),
  then \(\text{PM}_{\text{effective}} \le 0^\circ\). Stabilizing
  negative feedback flips into positive feedback, triggering explosive
  limit-cycle oscillations. As quantified in
  Table~\ref{tbl-math-phase-lag}, execution frequency directly dictates
  the preserved phase margin across operational cadences.
\end{itemize}

\begin{longtable}[]{@{}
  >{\raggedright\arraybackslash}p{(\linewidth - 8\tabcolsep) * \real{0.2000}}
  >{\raggedright\arraybackslash}p{(\linewidth - 8\tabcolsep) * \real{0.2000}}
  >{\raggedright\arraybackslash}p{(\linewidth - 8\tabcolsep) * \real{0.2000}}
  >{\raggedright\arraybackslash}p{(\linewidth - 8\tabcolsep) * \real{0.2000}}
  >{\raggedright\arraybackslash}p{(\linewidth - 8\tabcolsep) * \real{0.2000}}@{}}
\caption{\textbf{Frequency-Domain Phase Lag Impact across Physical AI
Cadences.} Demonstrating why high-frequency execution is mathematically
mandatory for physical
stability.}\label{tbl-math-phase-lag}\tabularnewline
\toprule\noalign{}
\begin{minipage}[b]{\linewidth}\raggedright
Sample Frequency (\(f_s\))
\end{minipage} & \begin{minipage}[b]{\linewidth}\raggedright
Sample Period (\(T\))
\end{minipage} & \begin{minipage}[b]{\linewidth}\raggedright
ZOH Delay (\(\tau_d = T/2\))
\end{minipage} & \begin{minipage}[b]{\linewidth}\raggedright
Phase Lag at \(\omega_c = 25\text{ rad/s}\) (\(4\text{ Hz}\))
\end{minipage} & \begin{minipage}[b]{\linewidth}\raggedright
Impact on \(35^\circ\) Phase Margin
\end{minipage} \\
\midrule\noalign{}
\endfirsthead
\toprule\noalign{}
\begin{minipage}[b]{\linewidth}\raggedright
Sample Frequency (\(f_s\))
\end{minipage} & \begin{minipage}[b]{\linewidth}\raggedright
Sample Period (\(T\))
\end{minipage} & \begin{minipage}[b]{\linewidth}\raggedright
ZOH Delay (\(\tau_d = T/2\))
\end{minipage} & \begin{minipage}[b]{\linewidth}\raggedright
Phase Lag at \(\omega_c = 25\text{ rad/s}\) (\(4\text{ Hz}\))
\end{minipage} & \begin{minipage}[b]{\linewidth}\raggedright
Impact on \(35^\circ\) Phase Margin
\end{minipage} \\
\midrule\noalign{}
\endhead
\bottomrule\noalign{}
\endlastfoot
\textbf{\(10\text{ Hz}\)} (Slow VLM) & \(100\text{ ms}\) &
\(50\text{ ms}\) & \textbf{\(-71.6^\circ\)} & \textbf{\(-36.6^\circ\)
(Destructive Instability / Crash)} \\
\textbf{\(50\text{ Hz}\)} (Action Chunk) & \(20\text{ ms}\) &
\(10\text{ ms}\) & \textbf{\(-14.3^\circ\)} & \textbf{\(+20.7^\circ\)
(Reduced Damping, Mild Overshoot)} \\
\textbf{\(1000\text{ Hz}\)} (MCU Reflex) & \(1.0\text{ ms}\) &
\(0.5\text{ ms}\) & \textbf{\(-0.72^\circ\)} & \textbf{\(+34.3^\circ\)
(Robust Closed-Loop Stability)} \\
\end{longtable}

\begin{center}\rule{0.5\linewidth}{0.5pt}\end{center}

\section*{7. Systems Translation
Reference}\label{systems-translation-reference}
\addcontentsline{toc}{section}{7. Systems Translation Reference}

\markright{7. Systems Translation Reference}

{\def\LTcaptype{none} % do not increment counter
\begin{longtable}[]{@{}
  >{\raggedright\arraybackslash}p{(\linewidth - 8\tabcolsep) * \real{0.2000}}
  >{\raggedright\arraybackslash}p{(\linewidth - 8\tabcolsep) * \real{0.2000}}
  >{\raggedright\arraybackslash}p{(\linewidth - 8\tabcolsep) * \real{0.2000}}
  >{\raggedright\arraybackslash}p{(\linewidth - 8\tabcolsep) * \real{0.2000}}
  >{\raggedright\arraybackslash}p{(\linewidth - 8\tabcolsep) * \real{0.2000}}@{}}
\toprule\noalign{}
\begin{minipage}[b]{\linewidth}\raggedright
Mathematical Concept
\end{minipage} & \begin{minipage}[b]{\linewidth}\raggedright
Core Equation / Invariant
\end{minipage} & \begin{minipage}[b]{\linewidth}\raggedright
Systems Meaning
\end{minipage} & \begin{minipage}[b]{\linewidth}\raggedright
Implementation Standard
\end{minipage} & \begin{minipage}[b]{\linewidth}\raggedright
Target Silicon
\end{minipage} \\
\midrule\noalign{}
\endhead
\bottomrule\noalign{}
\endlastfoot
\textbf{Dynamic Stopping} &
\(d_{\text{stop}} = v \tau + \frac{v^2}{2a_{\text{max}}}\) & Minimum
clearance to prevent impact & Closed-form velocity clamp
(\(v_{\text{safe}}\)) & Real-Time MCU (\(1000\text{ Hz}\)) \\
\textbf{Spatial Pose} &
\(\mathbf{T} \in SE(3), \, \mathbf{q} \in \mathbb{S}^3\) & 6-DoF
position and rotation & Quaternion Hamilton product + Normalization &
All Silicon Tiers \\
\textbf{Twist Integration} &
\(\mathbf{T}_{k+1} = \mathbf{T}_k \exp(\boldsymbol{\xi}^\wedge \Delta t)\)
& Singularity-free motion update & Rodrigues formula in C
(\texttt{float{[}4{]}{[}4{]}}) & MCU / Host MPU \\
\textbf{Control Barrier} & \(\dot{h} + \gamma h \ge 0\) & Forward
invariant safe operating set & Affine constraint in microsecond QP &
Real-Time MCU (\texttt{malloc\ =\ 0}) \\
\textbf{CBF-QP Solver} &
\(\min \frac{1}{2} \|\mathbf{u} - \mathbf{u}_{\text{cand}}\|^2\) &
Safety projection of neural proposals & Closed-form Active Set
(\texttt{qpOASES}/\texttt{OSQP}) & Static DTCM SRAM
(\(<175\,\mu\text{s}\)) \\
\textbf{State Estimation} &
\(\mathbf{P} = (\mathbf{I}-\mathbf{K}\mathbf{H})\mathbf{P}(\mathbf{I}-\mathbf{K}\mathbf{H})^T + \mathbf{K}\mathbf{R}\mathbf{K}^T\)
& Fusing noisy IMUs and delayed cameras & Joseph-form
continuous-discrete EKF & Real-Time MCU / Edge NPU \\
\textbf{Motor Commutation} & \(i_d = 0, \, \tau = K_t i_q\) & Decoupled
electromagnetic torque & Clarke/Park + Space Vector PWM & Inverter Timer
(\(20\text{ kHz}\)) \\
\textbf{Phase Margin} & \(\Delta \phi = -\omega_c \tau_{\text{delay}}\)
& Transport delay eroding control stability & \(1000\text{ Hz}\) control
loops to bound \(\Delta \phi \le 1.5^\circ\) & Bare-Metal MCU \\
\textbf{Camera Unprojection} &
\(\mathbf{P}_B = \mathbf{R}_C^B(Z \mathbf{K}^{-1}\tilde{\mathbf{x}}) + \mathbf{p}_C^B\)
& 2D pixel to metric 3D base coordinate & Closed-form unprojection
(\texttt{float{[}3{]}}) & MCU / Edge NPU \\
\textbf{Epipolar Invariant} &
\(\mathbf{x}_2^T \mathbf{F} \mathbf{x}_1 = 0, \, \epsilon_{\text{Sampson}}\)
& Stereo consistency \& extrinsic drift tracking & Online Sampson
residual filtering & Real-Time MPU / NPU \\
\textbf{Manipulator Jacobian} &
\(\boldsymbol{\xi} = \mathbf{J}(\mathbf{q})\dot{\mathbf{q}}, \, \mathbf{J}^\dagger = \mathbf{J}^T(\mathbf{J}\mathbf{J}^T)^{-1}\)
& Differential mapping from joint rates to Cartesian twist & Closed-form
column cross-products & Real-Time MCU / Host MPU \\
\textbf{Manipulability \& Singularity} &
\(w = \sqrt{\det(\mathbf{J}\mathbf{J}^T)} \ge 0.05, \, \kappa(\mathbf{J}) \le 50.0\)
& Dexterity volume and velocity amplification bound & SVD singular value
screening in DTCM SRAM & Real-Time MCU (\texttt{malloc\ =\ 0}) \\
\end{longtable}
}

\begin{center}\rule{0.5\linewidth}{0.5pt}\end{center}

\section*{8. Camera Projection and Epipolar
Geometry}\label{camera-projection-and-epipolar-geometry}
\addcontentsline{toc}{section}{8. Camera Projection and Epipolar
Geometry}

\markright{8. Camera Projection and Epipolar Geometry}

In physical perception, an optical camera does not produce 3D
coordinates directly. It captures a 2D projection of 3D Euclidean space
onto a semiconductor photodiode grid. To transform pixel observations
into metric affordance tokens in \(SE(3)\) that an actuator can reach,
systems rely on calibrated projective geometry, depth unprojection, and
epipolar multi-view constraints.

\subsection*{Pinhole Projection and Intrinsic
Matrix}\label{pinhole-projection-and-intrinsic-matrix}
\addcontentsline{toc}{subsection}{Pinhole Projection and Intrinsic
Matrix}

Consider a 3D point
\(\mathbf{P}_C = [X_C, Y_C, Z_C]^T \in \mathbb{R}^3\) in the camera
optical reference frame, where the \(+Z_C\) axis points forward along
the optical axis, \(+X_C\) points to the right, and \(+Y_C\) points
downward. Under ideal perspective projection, light rays pass through
the optical center (pinhole aperture) and project onto the image sensor
plane at focal distance \(f\):

\[\tilde{x} = \frac{X_C}{Z_C}, \quad \tilde{y} = \frac{Y_C}{Z_C}\]

where \((\tilde{x}, \tilde{y})\) are dimensionless normalized image
plane coordinates.

To convert normalized coordinates into discrete sensor pixel indices
\((u, v)\), the optical projection is parameterized by the
\textbf{Camera Intrinsic Matrix}
\(\mathbf{K} \in \mathbb{R}^{3 \times 3}\):

\[\mathbf{K} = \begin{bmatrix} f_x & s & c_x \\ 0 & f_y & c_y \\ 0 & 0 & 1 \end{bmatrix}\]

where: * \(f_x = \frac{f}{p_w}, \; f_y = \frac{f}{p_h}\): Effective
focal lengths expressed in pixel units, where \(p_w, p_h\) are the
physical horizontal and vertical pixel pitches (\(\mu\text{m/pixel}\))
on the CMOS silicon die. * \((c_x, c_y)\): Pixel coordinates of the
\textbf{principal point} (the intersection of the optical axis with the
sensor array, ideally \((W/2, H/2)\)). * \(s\): Pixel grid skew
coefficient (identically zero, \(s = 0\), for modern lithographic
rectangular silicon sensors).

Under zero skew (\(s=0\)), the projective mapping from a 3D camera point
\(\mathbf{P}_C\) to homogeneous pixel coordinates
\(\tilde{\mathbf{x}} = [u, v, 1]^T\) is:

\[Z_C \begin{bmatrix} u \\ v \\ 1 \end{bmatrix} = \mathbf{K} \mathbf{P}_C = \begin{bmatrix} f_x & 0 & c_x \\ 0 & f_y & c_y \\ 0 & 0 & 1 \end{bmatrix} \begin{bmatrix} X_C \\ Y_C \\ Z_C \end{bmatrix} = \begin{bmatrix} f_x X_C + c_x Z_C \\ f_y Y_C + c_y Z_C \\ Z_C \end{bmatrix}\]

Expanding into inhomogeneous Cartesian pixel coordinates:

\[u = f_x \frac{X_C}{Z_C} + c_x, \qquad v = f_y \frac{Y_C}{Z_C} + c_y\]

\subsection*{Optical Ray Unprojection and Metric Cartesian
Coordinates}\label{optical-ray-unprojection-and-metric-cartesian-coordinates}
\addcontentsline{toc}{subsection}{Optical Ray Unprojection and Metric
Cartesian Coordinates}

Because prospective projection collapses 3D rays onto a 2D plane, a
single pixel coordinate \((u, v)\) defines an infinite angular ray
rather than a discrete physical location. Recovering the 3D metric
coordinate requires unprojecting along the ray using registered
orthogonal depth \(Z = d(u, v) \in \mathbb{R}^+\) (measured via active
stereo, structured light, or time-of-flight).

The upper-triangular intrinsic matrix \(\mathbf{K}\) has a direct
closed-form inverse:

\[\mathbf{K}^{-1} = \begin{bmatrix} \frac{1}{f_x} & 0 & -\frac{c_x}{f_x} \\ 0 & \frac{1}{f_y} & -\frac{c_y}{f_y} \\ 0 & 0 & 1 \end{bmatrix}\]

Multiplying the homogeneous pixel vector by \(\mathbf{K}^{-1}\)
reconstructs the normalized unit-depth ray vector
\(\mathbf{r}_{\text{cam}}\) in the camera optical frame:

\[\mathbf{r}_{\text{cam}} = \mathbf{K}^{-1} \begin{bmatrix} u \\ v \\ 1 \end{bmatrix} = \begin{bmatrix} \frac{u - c_x}{f_x} \\ \frac{v - c_y}{f_y} \\ 1 \end{bmatrix}\]

Scaling this ray vector by the registered depth \(Z = d(u, v)\) yields
the 3D metric position \(\mathbf{P}_C\) in the camera optical frame:

\[\mathbf{P}_C = Z \cdot \mathbf{r}_{\text{cam}} = \begin{bmatrix} X_C \\ Y_C \\ Z_C \end{bmatrix} = \begin{bmatrix} \frac{u - c_x}{f_x} Z \\ \frac{v - c_y}{f_y} Z \\ Z \end{bmatrix}\]

\subsection*{Rigid Extrinsic Transform to Robot Base
Frame}\label{rigid-extrinsic-transform-to-robot-base-frame}
\addcontentsline{toc}{subsection}{Rigid Extrinsic Transform to Robot
Base Frame}

To enable physical manipulation, navigation, or obstacle avoidance, the
camera-frame point \(\mathbf{P}_C\) must be transformed into the robot's
stationary physical base frame \(B\). Applying the extrinsic homogeneous
transformation matrix
\(\mathbf{T}_C^B = (\mathbf{R}_C^B, \mathbf{p}_C^B) \in SE(3)\):

\[\mathbf{P}_B = \mathbf{R}_C^B \mathbf{P}_C + \mathbf{p}_C^B = \mathbf{R}_C^B \left( Z \cdot \mathbf{K}^{-1} \begin{bmatrix} u \\ v \\ 1 \end{bmatrix} \right) + \mathbf{p}_C^B\]

where \(\mathbf{R}_C^B \in SO(3)\) is the \(3 \times 3\) rotation matrix
aligning the camera optical axes to the robot base coordinate axes, and
\(\mathbf{p}_C^B \in \mathbb{R}^3\) is the physical lever-arm
displacement vector from the base origin to the optical center.

\subsection*{Surface Normal Estimation from Discrete Depth
Gradients}\label{surface-normal-estimation-from-discrete-depth-gradients}
\addcontentsline{toc}{subsection}{Surface Normal Estimation from
Discrete Depth Gradients}

For robotic grasping and contact mechanics, the agent requires the
outward unit surface normal vector \(\hat{\mathbf{n}} \in \mathbb{S}^2\)
at the contact target. Considering unprojected points
\(\mathbf{P}_B(u, v)\) as a 2D parameterized surface, the tangent
vectors spanning the local contact plane are given by spatial partial
derivatives:

\[\mathbf{t}_u = \frac{\partial \mathbf{P}_B}{\partial u} \approx \frac{\mathbf{P}_B(u + \Delta u, v) - \mathbf{P}_B(u - \Delta u, v)}{2 \Delta u}\]

\[\mathbf{t}_v = \frac{\partial \mathbf{P}_B}{\partial v} \approx \frac{\mathbf{P}_B(u, v + \Delta v) - \mathbf{P}_B(u, v - \Delta v)}{2 \Delta v}\]

The outward unit normal vector \(\hat{\mathbf{n}}\) facing the robot
sensor is computed via the normalized vector cross-product:

\[\hat{\mathbf{n}} = \frac{\mathbf{t}_v \times \mathbf{t}_u}{\|\mathbf{t}_v \times \mathbf{t}_u\|_2} \in \mathbb{S}^2\]

\subsection*{Multi-View Epipolar Geometry and Sampson Residual
Metric}\label{multi-view-epipolar-geometry-and-sampson-residual-metric}
\addcontentsline{toc}{subsection}{Multi-View Epipolar Geometry and
Sampson Residual Metric}

When fusing multi-camera observations or tracking stereo visual
odometry, corresponding pixel points across two calibrated viewpoints
must satisfy geometric coplanarity.

Let \(\mathbf{P} \in \mathbb{R}^3\) be a physical 3D scene point
observed by Camera 1 as \(\mathbf{P}_1\) and by Camera 2 as
\(\mathbf{P}_2\). The two cameras are related by relative rigid-body
rotation \(\mathbf{R} \in SO(3)\) and baseline translation
\(\mathbf{t} \in \mathbb{R}^3\), such that
\(\mathbf{P}_2 = \mathbf{R} \mathbf{P}_1 + \mathbf{t}\).

The vectors \(\mathbf{P}_2\), \(\mathbf{t}\), and
\(\mathbf{R}\mathbf{P}_1\) are coplanar. Expressing this geometric
condition as a scalar triple product:

\[\mathbf{P}_2^T \left( \mathbf{t} \times (\mathbf{R} \mathbf{P}_1) \right) = 0 \implies \mathbf{P}_2^T [\mathbf{t}]_\times \mathbf{R} \mathbf{P}_1 = 0\]

Defining the \textbf{Essential Matrix}
\(\mathbf{E} \in \mathbb{R}^{3 \times 3}\):

\[\mathbf{E} = [\mathbf{t}]_\times \mathbf{R}\]

Substituting the normalized image coordinates
\(\mathbf{x}_1 = \mathbf{K}_1^{-1} \tilde{\mathbf{x}}_1\) and
\(\mathbf{x}_2 = \mathbf{K}_2^{-1} \tilde{\mathbf{x}}_2\) defines the
\textbf{Fundamental Matrix} \(\mathbf{F} \in \mathbb{R}^{3 \times 3}\):

\[\tilde{\mathbf{x}}_2^T \mathbf{F} \tilde{\mathbf{x}}_1 = 0, \qquad \mathbf{F} = \mathbf{K}_2^{-T} \mathbf{E} \mathbf{K}_1^{-1} = \mathbf{K}_2^{-T} [\mathbf{t}]_\times \mathbf{R} \mathbf{K}_1^{-1}\]

\begin{itemize}
\tightlist
\item
  \textbf{Epipolar Line Invariant:} A pixel point
  \(\tilde{\mathbf{x}}_1\) in Camera 1 defines an epipolar line
  \(\mathbf{l}_2 = \mathbf{F} \tilde{\mathbf{x}}_1\) in Camera 2 upon
  which the corresponding point \(\tilde{\mathbf{x}}_2\) must lie:
  \(\tilde{\mathbf{x}}_2^T \mathbf{l}_2 = 0\).
\end{itemize}

During active robot operation, mechanical shocks and thermal expansion
degrade camera extrinsic parameters. To quantify geometric consistency
without re-triangulating 3D points, the online calibration filter
computes the \textbf{Sampson Epipolar Error}
\(\epsilon_{\text{Sampson}}\) (the first-order geometric distance
approximation from the points to their respective epipolar lines):

\[\epsilon_{\text{Sampson}}(\tilde{\mathbf{x}}_1, \tilde{\mathbf{x}}_2) = \frac{(\tilde{\mathbf{x}}_2^T \mathbf{F} \tilde{\mathbf{x}}_1)^2}{(\mathbf{F}\tilde{\mathbf{x}}_1)_1^2 + (\mathbf{F}\tilde{\mathbf{x}}_1)_2^2 + (\mathbf{F}^T\tilde{\mathbf{x}}_2)_1^2 + (\mathbf{F}^T\tilde{\mathbf{x}}_2)_2^2}\]

where \((\mathbf{v})_i\) denotes the \(i\)-th element of vector
\(\mathbf{v}\). If
\(\epsilon_{\text{Sampson}} > \epsilon_{\text{threshold}}\) across
tracked feature pairs, the estimator flags extrinsic baseline drift and
triggers real-time bundle adjustment.

\subsection*{Thermal Baseline Expansion and Triangulation Error
Sensitivity}\label{thermal-baseline-expansion-and-triangulation-error-sensitivity}
\addcontentsline{toc}{subsection}{Thermal Baseline Expansion and
Triangulation Error Sensitivity}

In stereo perception, metric depth \(Z\) is triangulated from stereo
disparity \(d = u_1 - u_2\) across physical baseline
\(b = \|\mathbf{t}\|\):

\[Z = \frac{f \cdot b}{d}\]

Differentiating with respect to disparity \(d\) and baseline length
\(b\):

\[\delta Z = \frac{\partial Z}{\partial d} \delta d + \frac{\partial Z}{\partial b} \delta b = -\frac{f \cdot b}{d^2} \delta d + \frac{f}{d} \delta b = -\frac{Z^2}{f \cdot b} \delta d + \frac{Z}{b} \delta b\]

\begin{itemize}
\tightlist
\item
  \textbf{Quadratic Distance Sensitivity:} Triangulation depth
  uncertainty grows quadratically (\(\propto Z^2\)) with distance. At
  short ranges (\(Z=0.5\text{ m}\)), a \(0.5\text{ px}\) disparity noise
  induces millimeter error; at far ranges (\(Z=5.0\text{ m}\)), the
  exact same noise induces tens of centimeters of spatial uncertainty.
\end{itemize}

Chassis aluminum expands under operational temperature rise \(\Delta T\)
with thermal coefficient
\(\alpha_{\text{Al}} \approx 23 \times 10^{-6}\text{ K}^{-1}\):

\[\delta b = b_0 \cdot \alpha_{\text{Al}} \cdot \Delta T\]

Additionally, uneven thermal expansion or motor torque vibration induces
an angular baseline tilt \(\delta \theta\), shifting feature disparity
by
\(\delta d_{\text{tilt}} = f \cdot \tan(\delta \theta) \approx f \cdot \delta \theta\).
The resulting depth error is:

\[\delta Z_{\text{tilt}} \approx -\frac{Z^2}{b} \delta \theta\]

For a stereo camera with baseline \(b = 0.20\text{ m}\), an angular
deflection of just \(\delta \theta = 0.5^\circ = 0.00873\text{ rad}\)
produces: * At working distance \(Z = 1.0\text{ m}\):
\(\delta Z = \frac{1.0^2}{0.20} \times 0.00873 = 0.0436\text{ m} \approx 4.4\text{ cm}\).
* At obstacle distance \(Z = 5.0\text{ m}\):
\(\delta Z = \frac{5.0^2}{0.20} \times 0.00873 = 1.091\text{ m} \approx 109.1\text{ cm}\).

This reveals why static calibration matrices fail in physical deployment
and mandates continuous online epipolar tracking.

\subsection*{Real-Time Metric Unprojection
Complexity}\label{real-time-metric-unprojection-complexity}
\addcontentsline{toc}{subsection}{Real-Time Metric Unprojection
Complexity}

In embedded perception and spatial affordance engines, metric
unprojection from 2D depth arrays into 3D robot base frames
(\(\mathbf{p}_B = \mathbf{R}_C^B \mathbf{K}^{-1} [u, v, 1]^T Z + \mathbf{p}_C^B\))
is implemented using pre-computed inverse focal length constants
(\(f_x^{-1}, f_y^{-1}\)) in static memory (\texttt{malloc\ =\ 0}).

On an ARM Cortex-M7 at \(480\text{ MHz}\), evaluating a single point
unprojection requires \(9\) multiply-accumulate operations, completing
in \(< 0.4\,\mu\text{s}\). For local patch surface normal estimation via
cross products across a \(3 \times 3\) depth kernel, computation is
bounded to \(< 1.8\,\mu\text{s}\), enabling real-time contact normal
extraction directly within high-frequency tactile and grasping reflex
loops.

\begin{center}\rule{0.5\linewidth}{0.5pt}\end{center}

\section*{9. Differential Kinematics and
Jacobians}\label{differential-kinematics-and-jacobians}
\addcontentsline{toc}{section}{9. Differential Kinematics and Jacobians}

\markright{9. Differential Kinematics and Jacobians}

In articulated manipulation, high-level intent proposals specify desired
operational-space trajectories for the end-effector
(\(\mathbf{x}_e \in SE(3)\)), while physical motors exert joint torques
and velocities (\(\mathbf{q}, \dot{\mathbf{q}} \in \mathbb{R}^n\)). The
fundamental differential mapping between joint velocity space and
Cartesian spatial twist is governed by the \textbf{Geometric Jacobian}.

\subsection*{Forward Differential Kinematics and Geometric Jacobian
Formulation}\label{forward-differential-kinematics-and-geometric-jacobian-formulation}
\addcontentsline{toc}{subsection}{Forward Differential Kinematics and
Geometric Jacobian Formulation}

Consider an \(n\)-DoF open kinematic chain composed of revolute joints
with generalized joint coordinates
\(\mathbf{q} = [q_1, q_2, \dots, q_n]^T \in \mathbb{R}^n\). The position
and orientation of each link frame \(i \in \{1, \dots, n\}\) relative to
the stationary base frame \(0\) is computed via cumulative homogeneous
transformation matrices:

\[\mathbf{T}_i^0(\mathbf{q}) = \mathbf{T}_1^0(q_1) \mathbf{T}_2^1(q_2) \cdots \mathbf{T}_i^{i-1}(q_i) = \begin{bmatrix} \mathbf{R}_i^0(\mathbf{q}) & \mathbf{p}_i^0(\mathbf{q}) \\ \mathbf{0}_{1 \times 3} & 1 \end{bmatrix}\]

where \(\mathbf{R}_i^0 \in SO(3)\) represents the rotation matrix and
\(\mathbf{p}_i^0 \in \mathbb{R}^3\) denotes the origin position of frame
\(i\).

The instantaneous spatial twist of the end-effector
\(\boldsymbol{\xi}_e = [\mathbf{v}_e^T, \boldsymbol{\omega}_e^T]^T \in \mathbb{R}^6\)
(combining linear velocity
\(\mathbf{v}_e = \dot{\mathbf{p}}_e \in \mathbb{R}^3\) and angular
velocity \(\boldsymbol{\omega}_e \in \mathbb{R}^3\)) is linearly related
to the joint velocity vector \(\dot{\mathbf{q}}\):

\[\boldsymbol{\xi}_e = \begin{bmatrix} \mathbf{v}_e \\ \boldsymbol{\omega}_e \end{bmatrix} = \mathbf{J}(\mathbf{q}) \dot{\mathbf{q}}\]

where \(\mathbf{J}(\mathbf{q}) \in \mathbb{R}^{6 \times n}\) is the
\textbf{Geometric Jacobian Matrix}.

For a kinematic chain consisting entirely of revolute joints, each joint
\(i\) contributes to the end-effector velocity by rotating about its
instantaneous unit axis
\(\mathbf{z}_{i-1} = \mathbf{R}_{i-1}^0 \hat{\mathbf{z}}\). The \(i\)-th
column \(\mathbf{J}_i(\mathbf{q}) \in \mathbb{R}^6\) is computed using
the vector cross product:

\[\mathbf{J}_i(\mathbf{q}) = \begin{bmatrix} \mathbf{J}_{v, i} \\ \mathbf{J}_{\omega, i} \end{bmatrix} = \begin{bmatrix} \mathbf{z}_{i-1} \times (\mathbf{p}_e - \mathbf{p}_{i-1}) \\ \mathbf{z}_{i-1} \end{bmatrix}\]

where \(\mathbf{p}_e = \mathbf{p}_n^0\) is the current position of the
end-effector flange, and \(\mathbf{p}_{i-1}\) is the position of the
origin of joint frame \(i-1\) expressed in the base frame.

\subsection*{Inverse Differential Kinematics and the Singularity
Barrier}\label{inverse-differential-kinematics-and-the-singularity-barrier}
\addcontentsline{toc}{subsection}{Inverse Differential Kinematics and
the Singularity Barrier}

To compute the joint velocities required to achieve a commanded
end-effector twist \(\boldsymbol{\xi}_{\text{cmd}}\), the system inverts
the differential kinematic relationship:

\[\dot{\mathbf{q}} = \mathbf{J}^\dagger(\mathbf{q}) \boldsymbol{\xi}_{\text{cmd}}\]

For redundant manipulators with more joint degrees of freedom than
task-space dimensions (\(n > 6\)), the standard minimum-norm solution
uses the Moore-Penrose right pseudo-inverse:

\[\mathbf{J}^\dagger(\mathbf{q}) = \mathbf{J}^T (\mathbf{J} \mathbf{J}^T)^{-1} \in \mathbb{R}^{n \times 6}\]

When the manipulator approaches a \textbf{Kinematic Singularity} (such
as full arm reach extension, shoulder-wrist coaxial alignment, or wrist
gimbal lock), the Jacobian matrix loses rank
(\(\text{rank}(\mathbf{J}) < 6\)). As a consequence:

\[\det\left( \mathbf{J}(\mathbf{q}) \mathbf{J}^T(\mathbf{q}) \right) \longrightarrow 0\]

The matrix \((\mathbf{J}\mathbf{J}^T)\) becomes singular, and its
smallest singular value collapses
(\(\sigma_{\text{min}}(\mathbf{J}) \to 0\)). Attempting to invert an
ill-conditioned Jacobian forces the pseudo-inverse gain toward infinity
(\(\|\mathbf{J}^\dagger\| \to \infty\)). Even an infinitesimal Cartesian
tracking error or commanded velocity results in unbounded joint velocity
commands (\(\|\dot{\mathbf{q}}\| \to \infty\)), causing immediate
actuator phase current saturation, violent mechanical whiplash, and
transmission gear shear.

\subsection*{Damped Least-Squares (Levenberg-Marquardt)
Regularization}\label{damped-least-squares-levenberg-marquardt-regularization}
\addcontentsline{toc}{subsection}{Damped Least-Squares
(Levenberg-Marquardt) Regularization}

To preserve numerical stability near singular configurations, real-time
trajectory runtimes replace the unconstrained pseudo-inverse with
\textbf{Damped Least-Squares (DLS)} inversion:

\[\mathbf{J}^*(\mathbf{q}) = \mathbf{J}^T \left( \mathbf{J} \mathbf{J}^T + \lambda^2 \mathbf{I}_{6 \times 6} \right)^{-1}\]

The damping parameter \(\lambda \in \mathbb{R}^+\) regularizes the
inversion by reformulating the problem as an unconstrained optimization
trading off tracking precision against joint velocity magnitude:

\[\min_{\dot{\mathbf{q}}} \left( \|\boldsymbol{\xi}_{\text{cmd}} - \mathbf{J}\dot{\mathbf{q}}\|_2^2 + \lambda^2 \|\dot{\mathbf{q}}\|_2^2 \right)\]

To maintain exact kinematic tracking in well-conditioned regions while
providing robust damping near singularities, the runtime modulates
\(\lambda\) dynamically based on the instantaneous Yoshikawa
manipulability index \(w(\mathbf{q})\):

\[\lambda^2 = \begin{cases} 0 & \text{if } w(\mathbf{q}) \ge w_0 \\ \lambda_0^2 \left( 1 - \frac{w(\mathbf{q})}{w_0} \right)^2 & \text{if } w(\mathbf{q}) < w_0 \end{cases}\]

where \(w_0\) is the manipulability threshold below which damping
engages, and \(\lambda_0\) is the maximum damping factor.

\subsection*{Singular Value Decomposition, Manipulability Ellipsoids,
and Condition
Numbers}\label{singular-value-decomposition-manipulability-ellipsoids-and-condition-numbers}
\addcontentsline{toc}{subsection}{Singular Value Decomposition,
Manipulability Ellipsoids, and Condition Numbers}

Applying Singular Value Decomposition (SVD) to the Jacobian matrix
provides complete geometric insight into the manipulator's instantaneous
dexterity:

\[\mathbf{J}(\mathbf{q}) = \mathbf{U} \boldsymbol{\Sigma} \mathbf{V}^T = \sum_{i=1}^6 \sigma_i \mathbf{u}_i \mathbf{v}_i^T\]

where
\(\mathbf{U} = [\mathbf{u}_1, \dots, \mathbf{u}_6] \in \mathbb{R}^{6 \times 6}\)
and
\(\mathbf{V} = [\mathbf{v}_1, \dots, \mathbf{v}_n] \in \mathbb{R}^{n \times n}\)
are orthonormal matrices spanning Cartesian twist directions and joint
velocity directions, respectively, and
\(\boldsymbol{\Sigma} = [\text{diag}(\sigma_1, \dots, \sigma_6) \mid \mathbf{0}_{6 \times (n-6)}] \in \mathbb{R}^{6 \times n}\)
contains the ordered singular values
\(\sigma_1 \ge \sigma_2 \ge \dots \ge \sigma_6 \ge 0\).

Mapping a unit hypersphere in joint velocity space
(\(\|\dot{\mathbf{q}}\|_2 \le 1\)) into Cartesian space defines the
\textbf{Manipulability Ellipsoid}:

\[\boldsymbol{\xi}^T (\mathbf{J}\mathbf{J}^T)^{-1} \boldsymbol{\xi} \le 1 \implies \sum_{i=1}^6 \frac{(\mathbf{u}_i^T \boldsymbol{\xi})^2}{\sigma_i^2} \le 1\]

The principal axes of the ellipsoid align with the eigenvectors
\(\mathbf{u}_i\) of \(\mathbf{J}\mathbf{J}^T\), with semi-axis lengths
exactly equal to the singular values \(\sigma_i\).

We define two fundamental metrics for real-time safety gating:

\begin{enumerate}
\def\labelenumi{\arabic{enumi}.}
\tightlist
\item
  \textbf{Yoshikawa Manipulability Index (\(w(\mathbf{q})\)):} Measures
  the total volume of the manipulability ellipsoid:
\end{enumerate}

\[w(\mathbf{q}) = \sqrt{\det\left( \mathbf{J}(\mathbf{q}) \mathbf{J}^T(\mathbf{q}) \right)} = \prod_{i=1}^6 \sigma_i(\mathbf{J})\]

When \(w(\mathbf{q}) \to 0\), the ellipsoid collapses along at least one
axis, indicating a loss of instantaneous mobility along that direction.

\begin{enumerate}
\def\labelenumi{\arabic{enumi}.}
\setcounter{enumi}{1}
\tightlist
\item
  \textbf{Jacobian Condition Number (\(\kappa(\mathbf{J})\)):}
  Quantifies the ratio between maximum and minimum dexterity axes:
\end{enumerate}

\[\kappa(\mathbf{J}) = \frac{\sigma_{\text{max}}(\mathbf{J})}{\sigma_{\text{min}}(\mathbf{J})} = \frac{\sigma_1}{\sigma_6} \in [1, \infty)\]

\begin{itemize}
\tightlist
\item
  When \(\kappa(\mathbf{J}) = 1.0\), the configuration is
  \textbf{isotropic} (the manipulability ellipsoid is a perfect sphere,
  providing uniform velocity generation in all directions).
\item
  When \(\kappa(\mathbf{J}) \gg 50.0\), the configuration is severely
  \textbf{anisotropic and ill-conditioned}, meaning small Cartesian
  velocity commands along the minor axis demand massive joint
  velocities.
\end{itemize}

\subsection*{Real-Time Kinematic Evaluation and Deterministic
Bounds}\label{real-time-kinematic-evaluation-and-deterministic-bounds}
\addcontentsline{toc}{subsection}{Real-Time Kinematic Evaluation and
Deterministic Bounds}

In embedded motion controllers, calculating the \(6 \times n\) Geometric
Jacobian \(\mathbf{J}(\mathbf{q})\), Yoshikawa manipulability
\(w(\mathbf{q}) = \sqrt{\det(\mathbf{J}\mathbf{J}^T)}\), and condition
metric \(\kappa(\mathbf{J})\) is performed using static array
allocations in SRAM with zero dynamic memory allocation
(\texttt{malloc\ =\ 0}).

On an ARM Cortex-M7 at \(480\text{ MHz}\), evaluating forward
kinematics, cross-product Jacobian columns
(\(\mathbf{z}_{i-1} \times (\mathbf{p}_e - \mathbf{p}_{i-1})\)), and
\(6 \times 6\) Gram matrix determinant requires under
\(4.5\,\mu\text{s}\) for a 7-DoF manipulator. When
\(w(\mathbf{q}) < 0.05\) or \(\kappa(\mathbf{J}) > 50.0\), the
microcontroller immediately flags singularity proximity and attenuates
Cartesian velocity proposals within the same \(1000\text{ Hz}\) control
cycle.

\begin{center}\rule{0.5\linewidth}{0.5pt}\end{center}

\begin{center}\rule{0.5\linewidth}{0.5pt}\end{center}

\section*{10. Trajectory Splines and Variational
Planning}\label{trajectory-splines-and-variational-planning}
\addcontentsline{toc}{section}{10. Trajectory Splines and Variational
Planning}

\markright{10. Trajectory Splines and Variational Planning}

In System 1.5 trajectory synthesis, mathematical derivations guarantee
physical feasibility and mechanical transmission protection. This
section catalogs the closed-form boundary-value matrix algebra for
quintic splines, the variational score-matching formulation for
diffusion action chunking, and the real-time kinodynamic direct
collocation program.

\subsection*{Closed-Form Quintic Boundary-Value Matrix
Inversion}\label{closed-form-quintic-boundary-value-matrix-inversion}
\addcontentsline{toc}{subsection}{Closed-Form Quintic Boundary-Value
Matrix Inversion}

A fifth-order polynomial segment \(s(\tau)\) parameterized over local
time \(\tau \in [0, T]\) is defined by:

\[s(\tau) = \sum_{i=0}^5 a_i \tau^i = a_0 + a_1 \tau + a_2 \tau^2 + a_3 \tau^3 + a_4 \tau^4 + a_5 \tau^5\]

The first three time derivatives correspond to velocity
\(\dot{s}(\tau)\), acceleration \(\ddot{s}(\tau)\), and jerk
\(\dddot{s}(\tau)\):

\[\dot{s}(\tau) = a_1 + 2 a_2 \tau + 3 a_3 \tau^2 + 4 a_4 \tau^3 + 5 a_5 \tau^4\]

\[\ddot{s}(\tau) = 2 a_2 + 6 a_3 \tau + 12 a_4 \tau^2 + 20 a_5 \tau^3\]

\[\dddot{s}(\tau) = 6 a_3 + 24 a_4 \tau + 60 a_5 \tau^2\]

Matching the six boundary conditions at \(\tau = 0\) and \(\tau = T\):

\[\mathbf{b} = \begin{bmatrix} s(0) \\ \dot{s}(0) \\ \ddot{s}(0) \\ s(T) \\ \dot{s}(T) \\ \ddot{s}(T) \end{bmatrix} = \begin{bmatrix} s_0 \\ \dot{s}_0 \\ \ddot{s}_0 \\ s_1 \\ \dot{s}_1 \\ \ddot{s}_1 \end{bmatrix} = \mathbf{M}(T) \mathbf{a}\]

where the Vandermonde-like mapping matrix
\(\mathbf{M}(T) \in \mathbb{R}^{6 \times 6}\) is:

\[\mathbf{M}(T) = \begin{bmatrix}
1 & 0 & 0 & 0 & 0 & 0 \\
0 & 1 & 0 & 0 & 0 & 0 \\
0 & 0 & 2 & 0 & 0 & 0 \\
1 & T & T^2 & T^3 & T^4 & T^5 \\
0 & 1 & 2T & 3T^2 & 4T^3 & 5T^4 \\
0 & 0 & 2 & 6T & 12T^2 & 20T^3
\end{bmatrix}\]

Because \(\mathbf{M}(T)\) is triangular in its top block, the lower
\(3 \times 3\) sub-block can be inverted symbolically without numerical
Gaussian elimination. The closed-form algebraic inverse
\(\mathbf{a} = \mathbf{M}^{-1}(T) \mathbf{b}\) yields:

\[\begin{aligned}
a_0 &= s_0 \\
a_1 &= \dot{s}_0 \\
a_2 &= \frac{1}{2} \ddot{s}_0 \\
a_3 &= \frac{20(s_1 - s_0) - (8\dot{s}_1 + 12\dot{s}_0)T - (3\ddot{s}_0 - \ddot{s}_1)T^2}{2T^3} \\
a_4 &= \frac{30(s_0 - s_1) + (14\dot{s}_1 + 16\dot{s}_0)T + (3\ddot{s}_0 - 2\ddot{s}_1)T^2}{2T^4} \\
a_5 &= \frac{12(s_1 - s_0) - 6(\dot{s}_1 + \dot{s}_0)T - (\ddot{s}_1 - \ddot{s}_0)T^2}{2T^5}
\end{aligned}\]

Enforcing boundary acceleration matching
(\(\ddot{s}(0) = \ddot{s}_{\text{prev}}(T)\)) guarantees global
\(\mathcal{C}^2\) continuity across all concatenated trajectory
segments, bounding jerk
(\(\|\dddot{\mathbf{q}}(t)\| \le j_{\text{max}}\)) and eliminating Dirac
acceleration impulses.

\subsection*{Stochastic Denoising Diffusion Trajectory Proofs and Score
Matching}\label{stochastic-denoising-diffusion-trajectory-proofs-and-score-matching}
\addcontentsline{toc}{subsection}{Stochastic Denoising Diffusion
Trajectory Proofs and Score Matching}

In Conditional Diffusion Policies, action chunk synthesis is structured
as learning the reverse conditional distribution
\(p_\theta(\mathbf{A}_{0:H} \mid \mathbf{O}_t)\) over trajectory
horizons \(\mathbf{A} \in \mathbb{R}^{H \times d_a}\).

\subsubsection*{Forward Noising Process}\label{forward-noising-process}
\addcontentsline{toc}{subsubsection}{Forward Noising Process}

The forward process corrupts ground-truth demonstration trajectories
\(\mathbf{A}^0\) with Gaussian noise over discrete diffusion steps
\(k \in \{1, \dots, K\}\) according to variance schedule
\(\beta_1, \dots, \beta_K\):

\[q(\mathbf{A}^k \mid \mathbf{A}^{k-1}) = \mathcal{N}\left( \mathbf{A}^k; \sqrt{1 - \beta_k} \mathbf{A}^{k-1}, \beta_k \mathbf{I} \right)\]

Using the substitution \(\alpha_k = 1 - \beta_k\) and cumulative product
\(\bar{\alpha}_k = \prod_{s=1}^k \alpha_s\), marginalizing the
intermediate distributions yields the closed-form marginal:

\[q(\mathbf{A}^k \mid \mathbf{A}^0) = \mathcal{N}\left( \mathbf{A}^k; \sqrt{\bar{\alpha}_k} \mathbf{A}^0, (1 - \bar{\alpha}_k) \mathbf{I} \right)\]

\[\mathbf{A}^k = \sqrt{\bar{\alpha}_k} \mathbf{A}^0 + \sqrt{1 - \bar{\alpha}_k} \boldsymbol{\epsilon}, \quad \boldsymbol{\epsilon} \sim \mathcal{N}(\mathbf{0}, \mathbf{I})\]

\subsubsection*{Variational Lower Bound (ELBO) and Score
Matching}\label{variational-lower-bound-elbo-and-score-matching}
\addcontentsline{toc}{subsubsection}{Variational Lower Bound (ELBO) and
Score Matching}

The neural network parameterized by weights \(\theta\) is trained by
optimizing the variational lower bound on data log-likelihood
\(\mathbb{E}[\log p_\theta(\mathbf{A}^0 \mid \mathbf{O})]\). Following
Ho et al. (\citeproc{ref-ho2020denoising}{2020}), the objective reduces
to mean-squared error score matching:

\[\mathcal{L}_{\text{diffusion}}(\theta) = \mathbb{E}_{k, \mathbf{A}^0, \boldsymbol{\epsilon}}\left[ \left\| \boldsymbol{\epsilon} - \boldsymbol{\epsilon}_\theta\left( \sqrt{\bar{\alpha}_k} \mathbf{A}^0 + \sqrt{1 - \bar{\alpha}_k} \boldsymbol{\epsilon}, k, \mathbf{O}_t \right) \right\|_2^2 \right]\]

By Tweedie's Formula, the estimated noise vector
\(\boldsymbol{\epsilon}_\theta\) directly estimates the Stein score
function (gradient of the data log-density):

\[\nabla_{\mathbf{A}^k} \log p(\mathbf{A}^k \mid \mathbf{O}_t) = -\frac{\boldsymbol{\epsilon}_\theta(\mathbf{A}^k, k, \mathbf{O}_t)}{\sqrt{1 - \bar{\alpha}_k}}\]

\subsubsection*{DDIM Accelerated Deterministic
Sampling}\label{ddim-accelerated-deterministic-sampling}
\addcontentsline{toc}{subsubsection}{DDIM Accelerated Deterministic
Sampling}

To execute real-time inference within \(20\text{--}50\text{ ms}\),
systems deploy Denoising Diffusion Implicit Models (DDIM), skipping
steps along a sub-sequence \(\tau_1, \dots, \tau_S\) with \(S \ll K\):

\[\mathbf{A}^{\tau_{i-1}} = \sqrt{\bar{\alpha}_{\tau_{i-1}}} \left( \frac{\mathbf{A}^{\tau_i} - \sqrt{1 - \bar{\alpha}_{\tau_i}} \boldsymbol{\epsilon}_\theta(\mathbf{A}^{\tau_i}, \tau_i, \mathbf{O}_t)}{\sqrt{\bar{\alpha}_{\tau_i}}} \right) + \sqrt{1 - \bar{\alpha}_{\tau_{i-1}} - \sigma_{\tau_i}^2} \boldsymbol{\epsilon}_\theta(\mathbf{A}^{\tau_i}, \tau_i, \mathbf{O}_t) + \sigma_{\tau_i} \mathbf{z}\]

Setting \(\sigma_{\tau_i} = 0\) yields a deterministic implicit ODE
solver that reconstructs candidate trajectories in \(S = 4\text{--}8\)
steps.

\subsection*{Multi-Variable Kinodynamic Collocation Quadratic
Programs}\label{multi-variable-kinodynamic-collocation-quadratic-programs}
\addcontentsline{toc}{subsection}{Multi-Variable Kinodynamic Collocation
Quadratic Programs}

When refining candidate action chunks under non-linear rigid body
dynamics and workspace obstacle envelopes, the system solves a
discrete-time kinodynamic direct collocation program:

\[\min_{\mathbf{x}_{0:H}, \mathbf{u}_{0:H-1}} \sum_{k=0}^{H-1} \left( \|\mathbf{x}_k - \mathbf{x}_k^{\text{ref}}\|_{\mathbf{Q}}^2 + \|\mathbf{u}_k\|_{\mathbf{R}}^2 + \|\mathbf{u}_k - \mathbf{u}_{k-1}\|_{\mathbf{S}}^2 \right) + \|\mathbf{x}_H - \mathbf{x}_H^{\text{ref}}\|_{\mathbf{Q}_f}^2\]

\[\begin{aligned}
\text{subject to:} \quad &\mathbf{x}_{k+1} = \mathbf{x}_k + \frac{\Delta t}{2} \left[ f(\mathbf{x}_k, \mathbf{u}_k) + f(\mathbf{x}_{k+1}, \mathbf{u}_{k+1}) \right] \quad &&(\text{Hermite-Simpson Collocation}) \\
&\mathbf{x}_0 = \hat{\mathbf{x}}(t_{\text{now}}) \quad &&(\text{Initial State Consensus}) \\
&\mathbf{q}_{\text{min}} \le \mathbf{q}_k \le \mathbf{q}_{\text{max}} \quad &&(\text{Joint Limit Convex Box}) \\
&\dot{\mathbf{q}}_{\text{min}} \le \dot{\mathbf{q}}_k \le \dot{\mathbf{q}}_{\text{max}} \quad &&(\text{Kinematic Velocity Bound}) \\
&|\boldsymbol{\tau}_k| \le \boldsymbol{\tau}_{\text{max}} \quad &&(\text{Motor Inverter Torque Saturation}) \\
&\|\dot{\boldsymbol{\tau}}_k\| \le \dot{\tau}_{\text{max}} \quad &&(\text{Gearbox Transmission Protection}) \\
&\phi_{\text{SDF}}(\mathbf{p}_{\text{link}}(\mathbf{q}_k)) \ge d_{\text{margin}} \quad &&(\text{Signed Distance Clearance})
\end{aligned}\] Warm-starting this program with the action chunk decoded
by ACT or Diffusion Policy allows Sequential Quadratic Programming (SQP)
solvers to converge within 2 to 3 iterations, guaranteeing physical
safety envelopes before commands enter the shared SRAM execution queue.

\begin{center}\rule{0.5\linewidth}{0.5pt}\end{center}

\section*{11. Silicon Interconnects and Thermal
Dynamics}\label{silicon-interconnects-and-thermal-dynamics}
\addcontentsline{toc}{section}{11. Silicon Interconnects and Thermal
Dynamics}

\markright{11. Silicon Interconnects and Thermal Dynamics}

In Physical AI systems, cognitive stages are partitioned across
heterogeneous silicon fabrics based on execution cadence, memory access
patterns, and thermal dissipation envelopes. This section establishes
the mathematical formulations governing shared interconnect bus
arbitration, arithmetic intensity rooflines, and dynamic junction
temperature evolution.

\subsection*{AXI Bus Quality-of-Service Priority
Arbitration}\label{axi-bus-quality-of-service-priority-arbitration}
\addcontentsline{toc}{subsection}{AXI Bus Quality-of-Service Priority
Arbitration}

In Unified Memory Architecture (UMA) Systems-on-Chip, heterogeneous
masters (CPU cores, Neural Processing Units, and direct-memory-access
camera peripherals) issue concurrent read and write bursts across a
shared AXI interconnect crossbar.

Let \(\mathcal{M} = \{1, \dots, M\}\) denote the set of bus masters
contending for memory crossbar bandwidth \(B_{\text{total}}\)
(\(\text{GB/s}\)). Each transaction from master \(m \in \mathcal{M}\)
carries a 4-bit Quality-of-Service priority value
\(q_m \in \{0, 1, \dots, 15\}\), where \(q = 15\) designates the highest
real-time priority.

The effective service latency \(L_m\) for a memory read request of size
\(S_m\) bytes from master \(m\) under weighted priority arbitration is
modeled by:

\[L_m = \tau_{\text{arb}} + \frac{S_m}{B_{\text{crossbar}}} + \sum_{k \in \mathcal{M}, q_k > q_m} \frac{S_k}{B_{\text{crossbar}}} + \mathbb{I}(q_m < 15) \cdot \tau_{\text{burst\_in\_flight}}\]

where: * \(\tau_{\text{arb}}\): Crossbar arbitration switching latency
(typically \(4\text{--}12\text{ ns}\)). *
\(\tau_{\text{burst\_in\_flight}}\): Latency penalty incurred when
waiting for an ongoing atomic non-preemptible burst from a
lower-priority master to conclude
(\(S_{\text{burst}} / B_{\text{crossbar}}\)).

To guarantee that real-time microcontroller reads (\(1000\text{ Hz}\)
safety loops) never miss control deadlines, the silicon arbiter enforces
hard preemption for \(q_{\text{MCU}} = 15\):

\[L_{\text{MCU}} \le \tau_{\text{arb}} + \frac{S_{\text{payload}}}{B_{\text{crossbar}}} + \frac{S_{\text{burst\_max}}}{B_{\text{crossbar}}} \le 50.0\text{ ns}\]

For background camera sensor DMA streams (\(q_{\text{cam}} = 4\)),
transactions are admitted only during bus idle cycles when
higher-priority queues (\(q \ge 10\)) are empty.

\subsection*{Arithmetic Intensity and the Memory Power
Roofline}\label{arithmetic-intensity-and-the-memory-power-roofline}
\addcontentsline{toc}{subsection}{Arithmetic Intensity and the Memory
Power Roofline}

The computational efficiency of an embedded neural model is governed by
its arithmetic intensity \(I_{\text{arith}}\), defined as the ratio of
floating-point arithmetic operations to off-chip DRAM memory transfers:

\[I_{\text{arith}} = \frac{N_{\text{FLOPs}}}{N_{\text{Bytes\_DRAM}}} \quad \left[ \frac{\text{FLOPs}}{\text{Byte}} \right]\]

Total electrical power dissipated by the compute subsystem
\(P_{\text{compute}}\) is the sum of dynamic gate switching power,
static transistor leakage, and DRAM bus input/output pad driving power:

\[P_{\text{compute}} = P_{\text{dynamic}} + P_{\text{leakage}} + P_{\text{DRAM\_IO}}\]

\[P_{\text{dynamic}} = \alpha \cdot C_{\text{eff}} \cdot V_{\text{core}}^2 \cdot f_{\text{clock}}\]

\[P_{\text{DRAM\_IO}} = E_{\text{DRAM\_byte}} \cdot \text{BW}_{\text{DRAM}} = E_{\text{DRAM\_byte}} \cdot \frac{\text{Throughput}_{\text{FLOPs/s}}}{I_{\text{arith}}}\]

where \(E_{\text{DRAM\_byte}} \approx 150\text{ pJ/Byte}\) for off-chip
LPDDR5 DRAM, compared to
\(E_{\text{SRAM\_byte}} \approx 2\text{ pJ/Byte}\) for on-chip static
SRAM.

\begin{itemize}
\tightlist
\item
  \textbf{SRAM Model Tiling Invariant:} When trajectory decoders are
  quantized to INT8 and fit entirely within on-chip SRAM
  (\(N_{\text{Bytes\_DRAM}} \to 0\)), \(I_{\text{arith}} \to \infty\).
  DRAM power collapses (\(P_{\text{DRAM\_IO}} \to 0\)), reducing total
  inference energy consumption by up to \(98.5\%\).
\end{itemize}

\subsection*{Lumped-Parameter Thermodynamic Junction Differential
Equations}\label{lumped-parameter-thermodynamic-junction-differential-equations}
\addcontentsline{toc}{subsection}{Lumped-Parameter Thermodynamic
Junction Differential Equations}

The thermal dynamics of a silicon die mounted inside an IP67-sealed,
passively cooled enclosure are modeled by a coupled two-node thermal
resistance-capacitance (RC) network:

\[\begin{aligned}
C_{\text{die}} \frac{d T_j(t)}{dt} &= P_{\text{compute}}(t) - \frac{T_j(t) - T_{\text{case}}(t)}{R_{\theta,\text{jc}}} \\
C_{\text{case}} \frac{d T_{\text{case}}(t)}{dt} &= \frac{T_j(t) - T_{\text{case}}(t)}{R_{\theta,\text{jc}}} - \frac{T_{\text{case}}(t) - T_{\text{ambient}}}{R_{\theta,\text{ca}}}
\end{aligned}\]

where: * \(T_j\): Silicon junction temperature (\(^\circ\text{C}\)). *
\(T_{\text{case}}\): Chassis metal case temperature
(\(^\circ\text{C}\)). * \(C_{\text{die}}, C_{\text{case}}\): Thermal
heat capacities of the silicon package and aluminum chassis
(\(\text{J/K}\)). * \(R_{\theta,\text{jc}}\): Thermal resistance from
junction to case (\(\text{K/W}\)). * \(R_{\theta,\text{ca}}\): Thermal
resistance from case to ambient air (\(\text{K/W}\)).

Solving the transient step response under constant high-load neural
inference power \(P_{\text{step}}\):

\[T_j(t) = T_{\text{ambient}} + P_{\text{step}} (R_{\theta,\text{jc}} + R_{\theta,\text{ca}}) \left[ 1 - \frac{\tau_1 e^{-t/\tau_1} - \tau_2 e^{-t/\tau_2}}{\tau_1 - \tau_2} \right]\]

where \(\tau_1, \tau_2\) are the dominant thermal time constants of the
silicon package (\(1\text{--}5\text{ s}\)) and the physical chassis
(\(60\text{--}300\text{ s}\)).

\begin{itemize}
\tightlist
\item
  \textbf{Proactive Thermal Derating Boundary:} If
  \(T_j(t) \ge T_{\text{warn}} = 75.0^\circ\text{C}\), the firmware
  triggers proactive cadence scaling (\(50\text{ Hz} \to 20\text{ Hz}\),
  \(H=16 \to 32\)), dropping \(P_{\text{compute}}\) by \(55\%\) before
  junction temperature reaches the destructive hardware DVFS throttling
  knee (\(T_{\text{trip}} = 85.0^\circ\text{C}\)).
\end{itemize}

\begin{center}\rule{0.5\linewidth}{0.5pt}\end{center}

\section*{12. Authority Dynamics and Bumpless
Transfer}\label{authority-dynamics-and-bumpless-transfer}
\addcontentsline{toc}{section}{12. Authority Dynamics and Bumpless
Transfer}

\markright{12. Authority Dynamics and Bumpless Transfer}

In collaborative and supervisory robotics, physical control authority
transfers dynamically between autonomous neural policies and human
safety operators. This section formalizes the authority state space,
derives \(\mathcal{C}^2\) quintic smoothstep blending polynomials, and
establishes statistical bounds on policy drift under human intervention.

\subsection*{The 4-State Authority Machine and Preemption
Invariant}\label{the-4-state-authority-machine-and-preemption-invariant}
\addcontentsline{toc}{subsection}{The 4-State Authority Machine and
Preemption Invariant}

The state transition mapping
\(\delta: \mathcal{S}_{\text{auth}} \times \{0, 1\}^4 \to \mathcal{S}_{\text{auth}}\)
enforces strict priority:

\[\delta(s, \mathbf{e}) = \begin{cases}
\text{ESTOP}, & \text{if } e_{\text{estop}} = 1 \\
\text{MANUAL}, & \text{if } e_{\text{human}} = 1, \neg e_{\text{estop}} \\
\text{SHARED}, & \text{if } e_{\text{shared}} = 1, \neg (e_{\text{estop}} \lor e_{\text{human}}) \\
\text{AUTONOMOUS}, & \text{if } e_{\text{lease}} \land e_{\text{cbf}}, \neg (e_{\text{estop}} \lor e_{\text{human}}) \\
\text{ESTOP}, & \text{otherwise}
\end{cases}\]

Hardware Safe Torque Off lines enforce the preemption invariant: if
\(e_{\text{estop}} = 1\), gate driver power is physically unlatched
within \(< 5.0\,\mu\text{s}\), regardless of processor state.

\subsection*{\texorpdfstring{Derivation of the \(\mathcal{C}^2\) Quintic
Smoothstep
Polynomial}{Derivation of the \textbackslash mathcal\{C\}\^{}2 Quintic Smoothstep Polynomial}}\label{derivation-of-the-mathcalc2-quintic-smoothstep-polynomial}
\addcontentsline{toc}{subsection}{Derivation of the \(\mathcal{C}^2\)
Quintic Smoothstep Polynomial}

When transitioning between autonomous policy command
\(\mathbf{u}_{\text{policy}}(t)\) and human operator input
\(\mathbf{u}_{\text{human}}(t)\) over transition window
\(T_{\text{blend}} \in [100, 200]\text{ ms}\), the blended command is:

\[\mathbf{u}_{\text{cmd}}(t) = (1 - S(\alpha(t))) \mathbf{u}_{\text{policy}}(t) + S(\alpha(t)) \mathbf{u}_{\text{human}}(t), \quad \text{with } \alpha(t) = \frac{t - t_{\text{takeover}}}{T_{\text{blend}}} \in [0, 1]\]

To eliminate torque derivative shocks (\(\dot{\mathbf{u}} \to \infty\))
and mechanical jerk impulses (\(\dddot{\mathbf{q}} \to \infty\)), the
blending polynomial \(S(\alpha) = \sum_{k=0}^5 c_k \alpha^k\) must
satisfy six physical boundary conditions:

\[\begin{aligned}
S(0) &= 0, \quad &S(1) &= 1 \quad &&(\text{Zero and Unity Amplitudes}) \\
\dot{S}(0) &= 0, \quad &\dot{S}(1) &= 0 \quad &&(\text{Zero Initial and Final Velocities}) \\
\ddot{S}(0) &= 0, \quad &\ddot{S}(1) &= 0 \quad &&(\text{Zero Initial and Final Accelerations})
\end{aligned}\]

Setting \(S(0) = \dot{S}(0) = \ddot{S}(0) = 0\) directly forces
\(c_0 = c_1 = c_2 = 0\). The polynomial simplifies to:

\[S(\alpha) = c_3 \alpha^3 + c_4 \alpha^4 + c_5 \alpha^5\]

Applying the remaining three boundary conditions at \(\alpha = 1\):

\[\begin{aligned}
S(1) &= c_3 + c_4 + c_5 = 1 \\
\dot{S}(1) &= 3 c_3 + 4 c_4 + 5 c_5 = 0 \\
\ddot{S}(1) &= 6 c_3 + 12 c_4 + 20 c_5 = 0
\end{aligned}\]

Solving this \(3 \times 3\) linear system:

\[\begin{bmatrix} 1 & 1 & 1 \\ 3 & 4 & 5 \\ 6 & 12 & 20 \end{bmatrix} \begin{bmatrix} c_3 \\ c_4 \\ c_5 \end{bmatrix} = \begin{bmatrix} 1 \\ 0 \\ 0 \end{bmatrix} \implies \begin{bmatrix} c_3 \\ c_4 \\ c_5 \end{bmatrix} = \begin{bmatrix} 10 \\ -15 \\ 6 \end{bmatrix}\]

The unique \(\mathcal{C}^2\) quintic smoothstep polynomial is:

\[S(\alpha) = 6\alpha^5 - 15\alpha^4 + 10\alpha^3\]

The peak derivative occurs at the midpoint \(\alpha = 0.5\):

\[\dot{S}_{\text{max}} = \left. (30\alpha^4 - 60\alpha^3 + 30\alpha^2) \right|_{\alpha=0.5} = 30(0.0625) - 60(0.125) + 30(0.25) = 1.875\]

The maximum commanded torque rate of change during human takeover is
strictly bounded:

\[\|\dot{\mathbf{u}}_{\text{cmd}}\| \le \frac{1.875}{T_{\text{blend}}} \|\mathbf{u}_{\text{human}} - \mathbf{u}_{\text{policy}}\| + \|\dot{\mathbf{u}}_{\text{policy}}\|\]

For \(T_{\text{blend}} = 150.0\text{ ms}\) and a maximum torque
discrepancy of \(40.0\text{ N}\cdot\text{m}\), the peak torque rate is
bounded to \(\le 500.0\text{ N}\cdot\text{m/s}\), well within the
continuous capability of brushless motor gearboxes.

\subsection*{Compounding Covariate Shift Bounds Under Behavioral
Cloning}\label{compounding-covariate-shift-bounds-under-behavioral-cloning}
\addcontentsline{toc}{subsection}{Compounding Covariate Shift Bounds
Under Behavioral Cloning}

When training a policy \(\pi_\theta\) via supervised behavioral cloning
on logged human teleoperation data \(\mathcal{D} = \{(s_t, a_t^*)\}\),
let \(\epsilon\) denote the expected single-step error rate under the
training distribution:

\[\mathbb{E}_{s \sim d_{\pi^*}} \left[ \ell(\pi_\theta(s), \pi^*(s)) \right] \le \epsilon\]

Under autonomous closed-loop execution, small errors shift the state
distribution away from the demonstration support. Following Ross,
Gordon, and Bagnell (\citeproc{ref-ross2011reduction}{2011}), the
compounding total variation distance between the expert state
distribution \(d_{\pi^*}\) and the learned policy distribution
\(d_{\pi_\theta}\) over horizon \(T\) is bounded by:

\[D_{\text{TV}}(d_{\pi^*}, d_{\pi_\theta}) \le 2 (1 - (1 - \epsilon)^T) \le 2 \epsilon T\]

The cumulative task loss over horizon \(T\) compounds quadratically:

\[J(\pi_\theta) \le J(\pi^*) + \epsilon \cdot T^2\]

\begin{itemize}
\tightlist
\item
  \textbf{DAgger Policy Gradient Formulation:} To prevent quadratic
  error compounding, dataset aggregation (DAgger) samples states from
  the running policy's own induced distribution \(d_{\pi_\theta}\) and
  queries an expert oracle for corrections \(a_t^* = \pi^*(s_t)\).
  Truncating the pre-takeover drift phase prevents the policy from
  learning spurious associations between dangerous drift and recovery
  maneuvers.
\end{itemize}

\begin{center}\rule{0.5\linewidth}{0.5pt}\end{center}

\section*{13. STPA Hazard Analysis and Reliability
Proofs}\label{stpa-hazard-analysis-and-reliability-proofs}
\addcontentsline{toc}{section}{13. STPA Hazard Analysis and Reliability
Proofs}

\markright{13. STPA Hazard Analysis and Reliability Proofs}

In whole-system verification and safety assurance, formal safety cases
require connecting top-level system claims to rigorous control-theoretic
hazard analyses and statistical extreme-value reliability proofs. This
section establishes the mathematical foundations of STPA hazard networks
and derives tail probability bounds using the Generalized Pareto
Distribution.

\subsection*{Formal Mathematical Formulation of STPA Hierarchical
Control
Structures}\label{formal-mathematical-formulation-of-stpa-hierarchical-control-structures}
\addcontentsline{toc}{subsection}{Formal Mathematical Formulation of
STPA Hierarchical Control Structures}

System-Theoretic Process Analysis (STPA) models safety as a hierarchical
dynamic control problem (\citeproc{ref-leveson2011engineering}{Leveson
2011}).

Let \(\mathcal{L} = \{L_1, L_2, \dots, L_p\}\) denote the set of
unacceptable \textbf{System Losses} (such as human injury, structural
equipment destruction, or environmental damage).

A \textbf{System Hazard} \(H_k \in \mathcal{H}\) is defined as a system
state or set of conditions that, together with a worst-case set of
environmental conditions, leads to a system loss:

\[H_k = \langle \mathcal{S}_k, \mathcal{C}_k \rangle \implies \bigcup_{i} L_i\]

where \(\mathcal{S}_k \subset \mathcal{X}\) is the unsafe physical state
space and \(\mathcal{C}_k\) is the active environmental context.

The multi-rate, heterogeneous Physical AI architecture is modeled as a
directed hierarchical control graph
\(\mathcal{G} = (\mathcal{V}, \mathcal{E})\): * Vertices
\(\mathcal{V} = \{\text{Cloud}, \text{Host MPU}, \text{Edge NPU}, \text{Bare-Metal MCU}, \text{Actuators}, \text{Physical Plant}\}\)
represent control entities and physical substrates. * Directed edges
\(\mathcal{E} = \mathcal{U}_{\text{ctrl}} \cup \mathcal{Y}_{\text{fb}}\)
represent control actions \(\mathbf{u} \in \mathcal{U}_{\text{ctrl}}\)
descending the hierarchy and feedback measurements
\(\mathbf{y} \in \mathcal{Y}_{\text{fb}}\) ascending the hierarchy.

Every control action \(\mathbf{u}_j\) issued by controller \(V_i\) is
evaluated across the four canonical classes of \textbf{Unsafe Control
Actions (UCAs)}:

\[\begin{aligned}
\text{UCA}_{j,1}(\mathbf{x}, t) &= (\mathbf{u}_j(t) = \emptyset) \land \mathcal{P}_{\text{req}}(\mathbf{x}(t)) &&\text{[Not Provided]} \\
\text{UCA}_{j,2}(\mathbf{x}, \mathbf{u}, t) &= (\mathbf{u}_j(t) = \mathbf{u}_{\text{unsafe}}) \land \mathcal{P}_{\text{hazard}}(\mathbf{x}(t), \mathbf{u}_j) &&\text{[Unsafe Action]} \\
\text{UCA}_{j,3}(\mathbf{x}, \mathbf{u}, t) &= (\mathbf{u}_j(t) \neq \emptyset) \land (t < t_{\text{early}} \lor t > t_{\text{late}}) &&\text{[Wrong Timing]} \\
\text{UCA}_{j,4}(\mathbf{x}, \mathbf{u}, \Delta t) &= (\Delta t < \Delta t_{\text{min}} \lor \Delta t > \Delta t_{\text{max}}) &&\text{[Wrong Duration]}
\end{aligned}\]

where \(\mathcal{P}(\mathbf{x})\) denotes a Boolean predicate over
continuous physical state space \(\mathcal{X}\).

The STPA safety verification contract dictates that for every identified
\(\text{UCA}_{j,k}\), the system must incorporate a deterministic,
hardware-isolated invariant enforcer:

\[\forall \text{UCA}_{j,k}, \quad \exists \mathcal{C}_{\text{barrier}}: \quad \text{Condition}(\text{UCA}_{j,k}) \implies \text{Trigger}(\text{Fallback Safety Interlock})\]

\subsection*{Extreme Value Theory (EVT) and Peak-Over-Threshold
Reliability
Proofs}\label{extreme-value-theory-evt-and-peak-over-threshold-reliability-proofs}
\addcontentsline{toc}{subsection}{Extreme Value Theory (EVT) and
Peak-Over-Threshold Reliability Proofs}

In safety-critical physical systems, nominal statistical distributions
(such as Gaussian or Student-\(t\) models) systematically underestimate
tail risk. Catastrophic collisions occur at the extreme tail of the
latency and tracking error distributions (\(P_{99.999}\)).

To rigorously prove that the probability of catastrophic hazard
occurrence satisfies the top-level safety claim:

\[P(\text{Hazard}) < 10^{-6} \text{ per operating hour}\]

Physical AI applies \textbf{Extreme Value Theory (EVT)} using the
\textbf{Peak-Over-Threshold (POT)} methodology
(\citeproc{ref-coles2001introduction}{Coles 2001}).

Let \(X_1, X_2, \dots, X_N\) be a sequence of independent and
identically distributed physical error measurements (such as
sense-to-actuation latency \(\tau_{\text{delay}}\) or spatial clearance
violation \(\Delta d\)) with underlying continuous distribution function
\(F(x)\).

Given a high threshold \(u \in \mathbb{R}\), define the threshold
exceedance variable \(Y = X - u\) for \(X > u\). The conditional excess
distribution function \(F_u(y)\) is:

\[F_u(y) = P(X - u \le y \mid X > u) = \frac{F(u + y) - F(u)}{1 - F(u)}\]

\subsubsection*{The Pickands-Balkema-de Haan
Theorem}\label{the-pickands-balkema-de-haan-theorem}
\addcontentsline{toc}{subsubsection}{The Pickands-Balkema-de Haan
Theorem}

The fundamental theorem of Extreme Value Theory states that for a broad
class of underlying distributions \(F(x)\), as threshold \(u\)
approaches the upper endpoint \(x_F\), the conditional excess
distribution \(F_u(y)\) converges uniformly to the \textbf{Generalized
Pareto Distribution (GPD)}:

\[\lim_{u \to x_F} \sup_{0 \le y < x_F - u} \left| F_u(y) - G(y; \sigma, \xi) \right| = 0\]

where the Generalized Pareto cumulative distribution function
\(G(y; \sigma, \xi)\) is defined by:

\[G(y; \sigma, \xi) = \begin{cases}
1 - \left( 1 + \frac{\xi y}{\sigma} \right)^{-1/\xi} & \text{for } \xi \neq 0 \\
1 - \exp\left( -\frac{y}{\sigma} \right) & \text{for } \xi = 0
\end{cases}\]

parameterized by scale parameter \(\sigma > 0\) and shape parameter
(tail index) \(\xi \in \mathbb{R}\), defined on support \(y \ge 0\) for
\(\xi \ge 0\), and \(0 \le y \le -\sigma/\xi\) for \(\xi < 0\).

\begin{itemize}
\tightlist
\item
  \textbf{\(\xi > 0\) (Fréchet / Heavy-Tailed):} Models
  non-deterministic systems with operating system scheduling stalls and
  memory bus contention where tail latencies decay polynomially.
\item
  \textbf{\(\xi = 0\) (Gumbel / Exponential):} Models memoryless random
  processes.
\item
  \textbf{\(\xi < 0\) (Weibull / Bounded):} Models hard real-time
  systems with deterministic physical upper limits.
\end{itemize}

\subsubsection*{Derivation of the Tail Return Level for Safety
Assurance}\label{derivation-of-the-tail-return-level-for-safety-assurance}
\addcontentsline{toc}{subsubsection}{Derivation of the Tail Return Level
for Safety Assurance}

Let \(N_u = \sum_{i=1}^N \mathbb{I}(X_i > u)\) denote the number of
recorded exceedances above threshold \(u\) in a sample of size \(N\).
The empirical probability of exceeding the threshold is
\(\hat{\zeta}_u = N_u / N\).

Combining the empirical threshold probability with the Generalized
Pareto model yields the cumulative tail distribution for \(x > u\):

\[P(X > x) = P(X > u) \cdot P(X - u > x - u \mid X > u) = \frac{N_u}{N} \left( 1 + \frac{\xi (x - u)}{\sigma} \right)^{-1/\xi}\]

Inverting this equation yields the \textbf{High-Quantile Safety Return
Level} \(x_p\) corresponding to extreme target failure probability
\(p\):

\[x_p = u + \frac{\sigma}{\xi} \left[ \left( \frac{N}{N_u} p \right)^{-\xi} - 1 \right]\]

For an embodied robot operating at control cadence
\(f_s = 1000\text{ Hz}\), one operating hour contains
\(N_{\text{hour}} = 3.6 \times 10^6\) execution cycles. A target hazard
rate of \(10^{-6}\text{ / hour}\) corresponds to a per-cycle failure
probability of:

\[p_{\text{cycle}} = \frac{10^{-6}}{3.6 \times 10^6} = 2.78 \times 10^{-13}\]

Evaluating \(x_{p_{\text{cycle}}}\) on empirical Hardware-in-the-Loop
fault injection datasets provides the exact mathematical upper bound on
worst-case latency or stopping distance required to substantiate the
Claims-Arguments-Evidence safety case.

\subsection*{Runtime STPA Monitor Architecture and Tail Estimation
Bounds}\label{runtime-stpa-monitor-architecture-and-tail-estimation-bounds}
\addcontentsline{toc}{subsection}{Runtime STPA Monitor Architecture and
Tail Estimation Bounds}

In whole-system assurance architectures, runtime STPA monitors execute
directly on the safety microcontroller in static DTCM SRAM
(\texttt{malloc\ =\ 0}). The monitor evaluates the four Unsafe Control
Action (UCA) conditions every \(1000\text{ Hz}\) tick: 1.
\textbf{Dynamic Clearance Violation (UCA-01):} Verified via
\(d_{\text{clearance}} \ge d_{\text{stop}}(v, \hat{\mu}, \tau_{\text{lat}})\).
2. \textbf{Heartbeat Loss (UCA-02):} Verified via
\(\Delta t_{\text{heartbeat}} \le 35.0\text{ ms}\). 3. \textbf{Tail
Latency Derate (UCA-03):} Verified via online Peak-Over-Threshold (POT)
Generalized Pareto return-level tracking. 4. \textbf{Trajectory Jitter
(UCA-04):} Verified via continuous torque derivative bounds
\(\|\dot{\boldsymbol{\tau}}\| \le \dot{\tau}_{\text{max}}\).

Evaluating all four invariants executes in \(< 1.2\,\mu\text{s}\),
guaranteeing immediate Category 1 stopping engagement upon any safety
invariant violation.

\begin{center}\rule{0.5\linewidth}{0.5pt}\end{center}

\section*{14. Decentralized Barrier
Certificates}\label{decentralized-barrier-certificates}
\addcontentsline{toc}{section}{14. Decentralized Barrier Certificates}

\markright{14. Decentralized Barrier Certificates}

This section provides formal proofs and embedded implementations for
distributed multi-agent safety certificates and geometric symmetry
constraints in physical foundation models.

\subsection*{Multi-Agent Decentralized Control Barrier Functions
(D-CBFs)}\label{multi-agent-decentralized-control-barrier-functions-d-cbfs}
\addcontentsline{toc}{subsection}{Multi-Agent Decentralized Control
Barrier Functions (D-CBFs)}

Consider a fleet of \(N\) autonomous physical agents operating in a
shared physical workspace. Each agent \(i \in \{1, \dots, N\}\)
possesses continuous control-affine dynamics:

\[\dot{\mathbf{x}}_i = f_i(\mathbf{x}_i) + g_i(\mathbf{x}_i)\mathbf{u}_i\]

where
\(\mathbf{x}_i = [\mathbf{p}_i^T, \mathbf{v}_i^T]^T \in \mathbb{R}^{n_i}\)
is the agent state (position \(\mathbf{p}_i \in \mathbb{R}^3\) and
velocity \(\mathbf{v}_i \in \mathbb{R}^3\)) and
\(\mathbf{u}_i \in \mathcal{U}_i \subset \mathbb{R}^{m_i}\) is the
control input (acceleration or motor force).

For any pair of agents \((i, j)\) with \(i \ne j\), let \(r_i\) and
\(r_j\) denote their respective physical bounding radii, and let
\(d_{\text{margin}}\) be a safety buffer. The minimum required physical
separation is \(R_{ij} = r_i + r_j + d_{\text{margin}}\). The pairwise
spatial barrier certificate function
\(h_{ij}(\mathbf{x}_i, \mathbf{x}_j): \mathbb{R}^{n_i} \times \mathbb{R}^{n_j} \to \mathbb{R}\)
is defined as:

\[h_{ij}(\mathbf{x}_i, \mathbf{x}_j) = \|\mathbf{p}_i - \mathbf{p}_j\|_2^2 - R_{ij}^2\]

The pairwise collision-free safe set is
\(\mathcal{C}_{ij} = \{ (\mathbf{x}_i, \mathbf{x}_j) \mid h_{ij}(\mathbf{x}_i, \mathbf{x}_j) \ge 0 \}\).
The collective fleet safe operating set \(\mathcal{C}_{\text{fleet}}\)
is the intersection across all pairs:

\[\mathcal{C}_{\text{fleet}} = \bigcap_{i=1}^N \bigcap_{j > i}^N \mathcal{C}_{ij}\]

\subsection*{Reciprocating Collision Avoidance and Decentralized Safety
Guarantees}\label{reciprocating-collision-avoidance-and-decentralized-safety-guarantees}
\addcontentsline{toc}{subsection}{Reciprocating Collision Avoidance and
Decentralized Safety Guarantees}

Differentiating \(h_{ij}(\mathbf{x}_i, \mathbf{x}_j)\) along the joint
trajectories yields:

\[\dot{h}_{ij}(\mathbf{x}_i, \mathbf{x}_j, \mathbf{u}_i, \mathbf{u}_j) = 2(\mathbf{p}_i - \mathbf{p}_j)^T (\mathbf{v}_i - \mathbf{v}_j) = \nabla_{\mathbf{x}_i} h_{ij} \left( f_i(\mathbf{x}_i) + g_i(\mathbf{x}_i)\mathbf{u}_i \right) + \nabla_{\mathbf{x}_j} h_{ij} \left( f_j(\mathbf{x}_j) + g_j(\mathbf{x}_j)\mathbf{u}_j \right)\]

Using Lie derivative notation:

\[\dot{h}_{ij} = L_{f_i} h_{ij} + L_{g_i} h_{ij} \mathbf{u}_i + L_{f_j} h_{ij} + L_{g_j} h_{ij} \mathbf{u}_j\]

To maintain forward invariance of \(\mathcal{C}_{ij}\), the joint
control input \((\mathbf{u}_i, \mathbf{u}_j)\) must satisfy:

\[L_{f_i} h_{ij} + L_{g_i} h_{ij} \mathbf{u}_i + L_{f_j} h_{ij} + L_{g_j} h_{ij} \mathbf{u}_j + \alpha(h_{ij}) \ge 0\]

where \(\alpha(\cdot)\) is an extended class-\(\mathcal{K}\) function
(chosen as linear \(\alpha(h) = \gamma h\) with \(\gamma > 0\)).

\subsubsection*{Reciprocating Responsibility
Decomposition}\label{reciprocating-responsibility-decomposition}
\addcontentsline{toc}{subsubsection}{Reciprocating Responsibility
Decomposition}

In a decentralized fleet lacking a central coordinator, agent \(i\)
cannot command agent \(j\)'s control input \(\mathbf{u}_j\). Under the
\textbf{Reciprocating Collision Avoidance Principle}, each agent assumes
equal responsibility for preserving the pairwise safety barrier:

\[\alpha_{ij}(h_{ij}) = \frac{1}{2} \alpha(h_{ij})\]

Agent \(i\) enforces its local decentralized barrier constraint:

\[L_{f_i} h_{ij}(\mathbf{x}_i, \hat{\mathbf{x}}_j) + L_{g_i} h_{ij}(\mathbf{x}_i, \hat{\mathbf{x}}_j) \mathbf{u}_i + \frac{1}{2} \gamma h_{ij}(\mathbf{x}_i, \hat{\mathbf{x}}_j) \ge 0\]

Simultaneously, agent \(j\) enforces its reciprocal constraint:

\[L_{f_j} h_{ij}(\hat{\mathbf{x}}_i, \mathbf{x}_j) + L_{g_j} h_{ij}(\hat{\mathbf{x}}_i, \mathbf{x}_j) \mathbf{u}_j + \frac{1}{2} \gamma h_{ij}(\hat{\mathbf{x}}_i, \mathbf{x}_j) \ge 0\]

\subsubsection*{Theorem: Collective Forward Set Invariance under
Reciprocating
D-CBFs}\label{theorem-collective-forward-set-invariance-under-reciprocating-d-cbfs}
\addcontentsline{toc}{subsubsection}{Theorem: Collective Forward Set
Invariance under Reciprocating D-CBFs}

If every agent \(i \in \{1, \dots, N\}\) selects its local control
action \(\mathbf{u}_i\) from its admissible set:

\[\mathcal{U}_i(\mathbf{x}_i) = \left\{ \mathbf{u}_i \in \mathbb{R}^{m_i} \;\middle|\; L_{f_i} h_{ij} + L_{g_i} h_{ij} \mathbf{u}_i + \frac{1}{2}\gamma h_{ij} \ge 0, \quad \forall j \ne i \right\}\]

then the joint system trajectory satisfies:

\[\begin{aligned}
\dot{h}_{ij} + \gamma h_{ij} &= \left( L_{f_i} h_{ij} + L_{g_i} h_{ij} \mathbf{u}_i + \tfrac{1}{2}\gamma h_{ij} \right) + \left( L_{f_j} h_{ij} + L_{g_j} h_{ij} \mathbf{u}_j + \tfrac{1}{2}\gamma h_{ij} \right) \\
&\ge 0 + 0 = 0
\end{aligned}\]

Applying the Comparison Lemma,
\(h_{ij}(t) \ge h_{ij}(0) e^{-\gamma t} \ge 0\) for all \(t \ge 0\). The
collective safe operating set \(\mathcal{C}_{\text{fleet}}\) is strictly
forward invariant, guaranteeing zero inter-agent collisions across the
entire fleet without centralized arbitration.

\subsection*{Asynchronous State Estimation and Network Latency Margin
Bounds}\label{asynchronous-state-estimation-and-network-latency-margin-bounds}
\addcontentsline{toc}{subsection}{Asynchronous State Estimation and
Network Latency Margin Bounds}

In realistic radio-frequency networks, agent \(i\) receives position
broadcasts from neighbor \(j\) with transport delay
\(\tau_{ij} \in [0, \tau_{\text{max}}]\) and bounded packet drop
probability.

Let
\(\hat{\mathbf{p}}_j(t) = \mathbf{p}_j(t - \tau_{ij}) + \mathbf{v}_j(t - \tau_{ij})\tau_{ij}\)
be the dead-reckoned position estimate. Given a known maximum physical
acceleration \(a_{\text{max}}\), the state estimation error is bounded:

\[\|\mathbf{p}_j(t) - \hat{\mathbf{p}}_j(t)\|_2 \le \frac{1}{2} a_{\text{max}} \tau_{ij}^2 \le \frac{1}{2} a_{\text{max}} \tau_{\text{max}}^2 = \epsilon_{\text{pos}}(\tau_{\text{max}})\]

To guarantee forward invariance under network latency, agent \(i\)
inflates the effective pairwise collision radius:

\[R_{ij}^{\text{net}}(\tau_{\text{max}}) = r_i + r_j + d_{\text{margin}} + 2 \epsilon_{\text{pos}}(\tau_{\text{max}}) = r_i + r_j + d_{\text{margin}} + a_{\text{max}} \tau_{\text{max}}^2\]

Evaluating the D-CBF against the inflated radius \(R_{ij}^{\text{net}}\)
mathematically guarantees collision avoidance despite asynchronous
communications.

\subsection*{\texorpdfstring{\(SE(3)\) Lie Group Equivariance in Neural
Physical
Policies}{SE(3) Lie Group Equivariance in Neural Physical Policies}}\label{se3-lie-group-equivariance-in-neural-physical-policies}
\addcontentsline{toc}{subsection}{\(SE(3)\) Lie Group Equivariance in
Neural Physical Policies}

In 3D physical manipulation and locomotion, the underlying laws of
Newtonian mechanics are invariant under rigid coordinate
transformations. A policy that learns to grasp an object at pose
\(\mathbf{T}_1 \in SE(3)\) should generalize directly when that same
object is rotated and translated to pose
\(\mathbf{T}_2 = g \cdot \mathbf{T}_1\), where \(g \in SE(3)\).

\subsubsection*{Mathematical Group Action on State and Action
Spaces}\label{mathematical-group-action-on-state-and-action-spaces}
\addcontentsline{toc}{subsubsection}{Mathematical Group Action on State
and Action Spaces}

Let \(G = SE(3) = SO(3) \ltimes \mathbb{R}^3\) denote the Special
Euclidean group. An element \(g = (\mathbf{R}, \mathbf{t}) \in G\)
consists of a 3D rotation matrix \(\mathbf{R} \in SO(3)\) and a
translation vector \(\mathbf{t} \in \mathbb{R}^3\).

Let \(\mathcal{X}\) denote the sensory observation space (point clouds,
3D feature embeddings, or depth tokens) and let \(\mathcal{A}\) denote
the physical action space (Cartesian end-effector waypoints or spatial
wrench vectors). We define the group representations
\(\rho_{\mathcal{X}}(g): \mathcal{X} \to \mathcal{X}\) and
\(\rho_{\mathcal{A}}(g): \mathcal{A} \to \mathcal{A}\):

\begin{enumerate}
\def\labelenumi{\arabic{enumi}.}
\tightlist
\item
  \textbf{On Observation Points
  \(\mathbf{P} \in \mathbb{R}^{N \times 3}\):}
  \(\rho_{\mathcal{X}}(g) \mathbf{P} = \mathbf{P} \mathbf{R}^T + \mathbf{1}_{N \times 1} \mathbf{t}^T\).
\item
  \textbf{On Commanded Trajectory Waypoints
  \(\mathbf{p}_k \in \mathbb{R}^3\) and Orientations
  \(\mathbf{R}_k \in SO(3)\):}
  \[\rho_{\mathcal{A}}(g)(\mathbf{p}_k, \mathbf{R}_k) = (\mathbf{R}\mathbf{p}_k + \mathbf{t}, \mathbf{R}\mathbf{R}_k)\]
\end{enumerate}

\subsubsection*{\texorpdfstring{Formal Definition of
\(SE(3)\)-Equivariance}{Formal Definition of SE(3)-Equivariance}}\label{formal-definition-of-se3-equivariance}
\addcontentsline{toc}{subsubsection}{Formal Definition of
\(SE(3)\)-Equivariance}

A neural policy \(\pi_\theta: \mathcal{X} \to \mathcal{A}\) is strictly
\textbf{\(SE(3)\)-equivariant} if the mapping commutes with the group
action:

\[\pi_\theta(\rho_{\mathcal{X}}(g) \cdot \mathbf{O}) = \rho_{\mathcal{A}}(g) \cdot \pi_\theta(\mathbf{O}), \quad \forall g \in SE(3), \quad \forall \mathbf{O} \in \mathcal{X}\]

\begin{verbatim}
                rho_X(g)
       O ──────────────────────> g · O
       │                           │
       │                           │
pi_theta(O)                    pi_theta(g · O)
       │                           │
       v                           v
       A ──────────────────────> g · A
                rho_A(g)
\end{verbatim}

When an unconstrained transformer or MLP processes physical
observations, it must learn the effect of rotations empirically by
consuming thousands of data augmentations. An \(SE(3)\)-equivariant
architecture enforces this commutation structurally in its layer
weights, guaranteeing that rotating the input scene rotates the
predicted action trajectory identically by construction.

\subsection*{Steerable Tensor Representations and Clebsch-Gordan
Decompositions}\label{steerable-tensor-representations-and-clebsch-gordan-decompositions}
\addcontentsline{toc}{subsection}{Steerable Tensor Representations and
Clebsch-Gordan Decompositions}

To build neural layers that respect \(SO(3)\) rotational symmetries,
feature vectors at each 3D point are decomposed into direct sums of
irreducible representations (irreps) of \(SO(3)\):

\[V = \bigoplus_{l=0}^{l_{\text{max}}} \mathbb{R}^{2l + 1}\]

where \(l\) is the rotation order: * \textbf{\(l = 0\) (Scalars,
dimension 1):} Invariant under rotation (temperatures, masses,
distances, grasp probabilities). \(\rho_0(\mathbf{R}) = 1\). *
\textbf{\(l = 1\) (Vectors, dimension 3):} Standard 3D vectors
(positions, velocities, forces). \(\rho_1(\mathbf{R}) = \mathbf{R}\). *
\textbf{\(l = 2\) (Tensors, dimension 5):} Symmetric traceless
\(3 \times 3\) matrices (stress tensors, quadric surfaces, inertia
moments). \(\rho_2(\mathbf{R}) = \mathbf{D}^2(\mathbf{R})\), where
\(\mathbf{D}^l(\mathbf{R})\) is the \((2l+1) \times (2l+1)\) Wigner
D-matrix.

When combining geometric features \(\mathbf{v}_1 \in V_{l_1}\) and
\(\mathbf{v}_2 \in V_{l_2}\), the tensor product is projected back into
irreducible channels using \textbf{Clebsch-Gordan coefficients}
\(C_{(l_1, m_1)(l_2, m_2)}^{(l, m)}\):

\[(\mathbf{v}_1 \otimes_{\text{CG}} \mathbf{v}_2)_m^l = \sum_{m_1=-l_1}^{l_1} \sum_{m_2=-l_2}^{l_2} C_{(l_1, m_1)(l_2, m_2)}^{(l, m)} v_{1, m_1}^{l_1} v_{2, m_2}^{l_2}\]

for all target orders \(|l_1 - l_2| \le l \le l_1 + l_2\).

By constructing neural network layers strictly from Clebsch-Gordan
tensor products and radial scalar non-linearities, the entire network
satisfies exact \(SE(3)\)-equivariance at every hidden layer.

\subsection*{Sample Complexity Scaling Proof in Physical
Space}\label{sample-complexity-scaling-proof-in-physical-space}
\addcontentsline{toc}{subsection}{Sample Complexity Scaling Proof in
Physical Space}

Let \(\mathcal{F}\) be the class of \(L\)-Lipschitz continuous robotic
manipulation policies mapping a bounded workspace
\(\Omega \subset \mathbb{R}^3\) to bounded action trajectories in
\([-1, 1]^d\).

\subsubsection*{Non-Equivariant Hypothesis Space
Complexity}\label{non-equivariant-hypothesis-space-complexity}
\addcontentsline{toc}{subsubsection}{Non-Equivariant Hypothesis Space
Complexity}

For a standard unconstrained neural network (such as a Vision
Transformer or MLP), the \(\epsilon\)-covering number of the policy
hypothesis space \(\mathcal{H}_{\text{standard}}\) under the infinity
norm satisfies:

\[\log \mathcal{N}\left(\epsilon, \mathcal{H}_{\text{standard}}, \|\cdot\|_\infty\right) = \mathcal{O}\left( \left(\frac{L}{\epsilon}\right)^{d_{\text{in}}} \right)\]

where \(d_{\text{in}} = 3K\) represents the dimensionality of \(K\)
spatial input points. The number of physical demonstrations
\(M_{\text{standard}}\) required to achieve generalization error
\(\le \epsilon\) scales as:

\[M_{\text{standard}} = \Omega\left( \frac{1}{\epsilon^2} \left(\frac{L}{\epsilon}\right)^{3K} \right)\]

\subsubsection*{Equivariant Hypothesis Space
Complexity}\label{equivariant-hypothesis-space-complexity}
\addcontentsline{toc}{subsubsection}{Equivariant Hypothesis Space
Complexity}

When the policy space is constrained to the \(SE(3)\)-equivariant
subspace
\(\mathcal{H}_{\text{equiv}} \subset \mathcal{H}_{\text{standard}}\),
any function \(f \in \mathcal{H}_{\text{equiv}}\) is uniquely determined
by its values on the orbit quotient space \(\Omega / SE(3)\).

Because the group \(SE(3)\) has dimension \(\dim(SE(3)) = 6\) (3
translational and 3 rotational degrees of freedom), the effective input
manifold dimensionality is reduced from \(3K\) to \(3K - 6\) for
\(K \ge 3\). Furthermore, the volume of the search space is divided by
the Haar volume of the compact rotation group \(|SO(3)| = 8\pi^2\):

\[\log \mathcal{N}\left(\epsilon, \mathcal{H}_{\text{equiv}}, \|\cdot\|_\infty\right) = \mathcal{O}\left( \frac{1}{|SO(3)|} \left(\frac{L}{\epsilon}\right)^{3K - 6} \right)\]

For 3D robotic grasping where the target object can appear in arbitrary
\(SE(3)\) poses, the sample complexity ratio evaluates to:

\[\frac{M_{\text{equiv}}}{M_{\text{standard}}} \approx \frac{1}{|SO(3)| \cdot (\text{discretization steps})^{3}} \approx \frac{1}{100} \text{ to } \frac{1}{1000}\]

This mathematical bound explains why \(SE(3)\)-equivariant foundation
policies (such as Equivariant Diffusion Policy and Perceiver-Actor)
achieve \(95\%\) pick-and-place success rates with only
\(20\text{--}50\) physical demonstrations, whereas non-equivariant
models require \(2{,}000\text{--}5{,}000\) trials to learn spatial
rotational invariance.

\subsection*{Distributed Safety Filter Determinism and Complexity
Bounds}\label{distributed-safety-filter-determinism-and-complexity-bounds}
\addcontentsline{toc}{subsection}{Distributed Safety Filter Determinism
and Complexity Bounds}

In multi-agent fleets, each agent's bare-metal microcontroller evaluates
reciprocating Decentralized Control Barrier Functions (D-CBFs) across
all neighboring peers within its RF sensing radius (e.g.~up to
\(N_{\text{nbr}} = 8\) neighbors) at \(1000\text{ Hz}\).

Because pairwise half-space projection calculations are decoupled
(\(L_f h_{ij} + L_g h_{ij}\mathbf{u}_i + \frac{1}{2}\gamma h_{ij} \ge 0\)),
filtering candidate 3D acceleration vectors
\(\mathbf{u}_{\text{cand}} \in \mathbb{R}^3\) requires only
\(O(N_{\text{nbr}})\) dot products in static memory
(\texttt{malloc\ =\ 0}). On an ARM Cortex-M7 at \(480\text{ MHz}\),
evaluating 8 neighbor barrier certificates and projecting onto the safe
velocity half-space executes in under \(2.8\,\mu\text{s}\)
(\(\text{WCET} \le 5.0\,\mu\text{s}\)), providing distributed forward
set invariance with minimal compute overhead.

\begin{center}\rule{0.5\linewidth}{0.5pt}\end{center}

\chapter*{Canonical Systems Glossary}\label{sec-appendix-glossary}
\addcontentsline{toc}{chapter}{Canonical Systems Glossary}

\markboth{Canonical Systems Glossary}{Canonical Systems Glossary}

This glossary provides the authoritative operational definitions for the
foundational concepts of Physical AI. Each term represents a precise
architectural invariant, physical boundary, or mathematical guarantee.

\begin{center}\rule{0.5\linewidth}{0.5pt}\end{center}

\section*{Foundational Physical AI Architecture \&
Invariants}\label{foundational-physical-ai-architecture-invariants}
\addcontentsline{toc}{section}{Foundational Physical AI Architecture \&
Invariants}

\markright{Foundational Physical AI Architecture \& Invariants}

{\def\LTcaptype{none} % do not increment counter
\begin{longtable}[]{@{}
  >{\raggedright\arraybackslash}p{(\linewidth - 4\tabcolsep) * \real{0.3333}}
  >{\raggedright\arraybackslash}p{(\linewidth - 4\tabcolsep) * \real{0.3333}}
  >{\raggedright\arraybackslash}p{(\linewidth - 4\tabcolsep) * \real{0.3333}}@{}}
\toprule\noalign{}
\begin{minipage}[b]{\linewidth}\raggedright
Term
\end{minipage} & \begin{minipage}[b]{\linewidth}\raggedright
Operational Systems Definition
\end{minipage} & \begin{minipage}[b]{\linewidth}\raggedright
Primary Reference
\end{minipage} \\
\midrule\noalign{}
\endhead
\bottomrule\noalign{}
\endlastfoot
\textbf{Physical AI System} & An autonomous computational system that
couples unspecifiable learned neural policies to real-time physical
actuators, exercising delegated authority over continuous matter,
momentum, and thermodynamic energy under non-zero mass and finite
latency. & Chapter~\ref{sec-boundary} \\
\textbf{Causal Boundary} & The physical and temporal threshold at the
memory-mapped actuator register write (DAC, PWM capture/compare latch,
gate-driver control word) where digital bits are irreversibly
transformed into physical energy, momentum, or thermal heat. Beyond this
boundary, computational errors cannot be undone by software rollback. &
Chapter~\ref{sec-boundary} \\
\textbf{Endogenous Sensory Distribution Shift} & The fundamental
closed-loop feedback phenomenon where a machine's own physical actions
modify the external environment, thereby reshaping its own future sensor
distribution (\(s_{t+1} \sim P(\cdot \mid s_t, a_t)\)). A single policy
deviation alters all subsequent observations, causing compounding
quadratic error (\(O(T^2 \epsilon)\)) under behavioral cloning. &
Chapter~\ref{sec-boundary}, Chapter~\ref{sec-data},
Chapter~\ref{sec-training} \\
\textbf{Proposal--Permission Privilege Boundary} & The architectural
separation of privilege establishing that learned models (System 2 /
MPU) generate unverified candidate proposals (\(u_{\text{prop}}\)),
while an independent, deterministic supervisor (System 1 / MCU) running
at \(1000\text{ Hz}\) retains exclusive authority to permit, project, or
veto execution. & Chapter~\ref{sec-boundary}, Chapter~\ref{sec-brain},
Chapter~\ref{sec-enforcement} \\
\textbf{Two-Speed Embodied Compute Architecture} & An asymmetric compute
topology partitioning a machine into a high-TDP, stochastic deliberative
tier (1--50 Hz on Linux/accelerators) for semantic reasoning and
perception, and a hard real-time, deterministic tier (1000 Hz on
bare-metal MCU/FPGA) for closed-loop barrier enforcement. &
Chapter~\ref{sec-brain}, Chapter~\ref{sec-placement} \\
\textbf{Expiring Intent Lease} & A strictly typed, time-bounded contract
(\(\langle \mathcal{G}, \tau_{\text{expire}}, \mathcal{E}_{\text{tol}} \rangle\))
issued by the cognitive tier defining a spatial target envelope. An
intent lease carries its own expiration deadline; if not renewed by the
Brain before expiry, the real-time nervous system treats silence as an
automatic fail-safe trigger and initiates controlled stopping. &
Chapter~\ref{sec-intent}, Chapter~\ref{sec-planning} \\
\textbf{Action Chunk Seam Discontinuity} & The \(C^0\) velocity or
acceleration mismatch occurring at the boundary between successive
multi-step receding-horizon trajectory chunks (\(\hat{u}_{t:t+H}\)).
Splicing unblended chunks transmits destructive impulsive torque spikes
(\(\tau = J_{\text{eff}} \frac{\Delta\dot{q}}{T_c}\)) through gearboxes,
necessitating \(C^2\) quintic polynomial blending. &
Chapter~\ref{sec-planning} \\
\textbf{Minimal-Intervention Safety Shield} & A real-time optimization
filter that evaluates candidate neural actions at \(1000\text{ Hz}\) via
Control Barrier Functions (CBF). When a proposal threatens forward
invariance (\(h(x) < 0\)), the shield solves a microsecond Quadratic
Program (QP) that projects the command onto the boundary of the safe set
(\(\partial \mathcal{C}\)) with minimal deviation from the user's
intended heading. & Chapter~\ref{sec-enforcement} \\
\textbf{Multidimensional Sim-to-Real Reality Gap Vector
(\(\boldsymbol{\Delta}_{\text{gap}}\))} & The coupled vector of physical
discrepancies between simulation and reality across dynamics
(\(\Delta_{\text{dyn}}\)), sensor transduction
(\(\Delta_{\text{sens}}\)), visual appearance (\(\Delta_{\text{vis}}\)),
and compute/transport latency (\(\Delta_{\text{lat}}\)). It cannot be
collapsed into an abstract scalar. & Chapter~\ref{sec-training} \\
\textbf{Observational Indistinguishability Limit} & The physical
epistemic horizon where two distinct world states (one benign, one
fatal) generate mathematically identical sensor telemetry prior to the
braking deadline (\(y_A(t) \equiv y_B(t)\) for
\(t \le t_{\text{deadline}}\)). Once crossed, no downstream compute or
neural architecture can differentiate the states in time to prevent
physical harm. & Chapter~\ref{sec-frontier} \\
\textbf{Astronomical Exposure Wall (Butler \& Finelli Bound)} & The
mathematical theorem proving that certifying life-critical autonomous
systems (\(10^{-9}\text{ failures/h}\)) purely through empirical fleet
testing requires unfeasible testing durations
(\(>3.0 \times 10^9\text{ hours} \approx 342{,}000\text{ machine-years}\)).
High assurance must be achieved through deterministic architectural
containment rather than accumulated road hours. &
Chapter~\ref{sec-evaluation}, Chapter~\ref{sec-verification},
Chapter~\ref{sec-frontier} \\
\end{longtable}
}

\begin{center}\rule{0.5\linewidth}{0.5pt}\end{center}

\section*{Real-Time Silicon \& Nervous System
Mechanics}\label{real-time-silicon-nervous-system-mechanics}
\addcontentsline{toc}{section}{Real-Time Silicon \& Nervous System
Mechanics}

\markright{Real-Time Silicon \& Nervous System Mechanics}

{\def\LTcaptype{none} % do not increment counter
\begin{longtable}[]{@{}
  >{\raggedright\arraybackslash}p{(\linewidth - 4\tabcolsep) * \real{0.3333}}
  >{\raggedright\arraybackslash}p{(\linewidth - 4\tabcolsep) * \real{0.3333}}
  >{\raggedright\arraybackslash}p{(\linewidth - 4\tabcolsep) * \real{0.3333}}@{}}
\toprule\noalign{}
\begin{minipage}[b]{\linewidth}\raggedright
Term
\end{minipage} & \begin{minipage}[b]{\linewidth}\raggedright
Operational Systems Definition
\end{minipage} & \begin{minipage}[b]{\linewidth}\raggedright
Primary Reference
\end{minipage} \\
\midrule\noalign{}
\endhead
\bottomrule\noalign{}
\endlastfoot
\textbf{Lock-Free Seqlock Buffer} & A Single-Producer Single-Consumer
(SPSC) shared-memory synchronization protocol utilizing an atomic
monotonic sequence counter. Writers increment the counter before and
after copying data; readers check parity before and after reading to
detect torn writes without ever blocking the real-time execution path. &
Chapter~\ref{sec-nervous}, Chapter~\ref{sec-placement} \\
\textbf{Worst-Case Execution Time (WCET)} & The provable upper bound on
the time required for a software task to execute on target hardware
(\(R = C \cdot \frac{f_0}{f_1} + B + I_{\text{bus}} \le D\)), accounting
for core downclocking, interrupt blocking, and memory crossbar bus
contention. & Chapter~\ref{sec-nervous}, Chapter~\ref{sec-placement} \\
\textbf{Zero-Allocation SRAM Invariant} & The architectural rule
strictly prohibiting dynamic heap memory allocation (\texttt{malloc},
\texttt{free}, garbage collection) on the hard real-time safety path.
All data buffers, active-set QP matrices, and ring structures must be
statically allocated in on-chip SRAM/TCM to eliminate non-deterministic
timing spikes. & Chapter~\ref{sec-nervous},
Chapter~\ref{sec-enforcement}, Chapter~\ref{sec-placement} \\
\textbf{Spatial Information Age (\(\tau_{\text{age}}\))} & The total
physical duration elapsed between the midpoint of sensor optical
exposure (\(t_{\text{exp}}/2\)) and the instant the resulting actuator
command produces physical torque
(\(\tau_{\text{age}} = \frac{t_{\text{exp}}}{2} + t_{\text{readout}} + t_{\text{ISP}} + t_{\text{DMA}} + t_{\text{infer}} + t_{\text{enforce}} + t_{\text{rise}}\)).
& Chapter~\ref{sec-body}, Chapter~\ref{sec-perception} \\
\textbf{Belief State Invalidation Horizon (\(t_{\text{exp}}\))} & The
finite duration over which an unobserved state estimate remains within
acceptable task tolerance
(\(t_{\text{exp}} = \inf \{ t \ge 0 : E(t) \ge E_{\max} \}\)). Governed
by physical drift rates and external disturbance acceleration. &
Chapter~\ref{sec-memory} \\
\textbf{Volumetric Free-Space Raycasting} & The process of casting
negative-evidence sightlines through 3D spatial occupancy structures
(TSDF, Sparse VDB) to clear previously occupied voxels. Prevents phantom
obstacle persistence while tracking dynamic obstacles. &
Chapter~\ref{sec-memory} \\
\textbf{Frozen Transform Trap} & A failure mode where an intermediate
software transport bridge republishes cached spatial coordinate matrices
(\(T_A^B\)) with updated wall-clock timestamps, misleading downstream
nodes into treating stale spatial geometry as fresh telemetry. &
Chapter~\ref{sec-memory} \\
\end{longtable}
}

\begin{center}\rule{0.5\linewidth}{0.5pt}\end{center}

\section*{Robotics, Safety Assurance \& Statutory
Standards}\label{robotics-safety-assurance-statutory-standards}
\addcontentsline{toc}{section}{Robotics, Safety Assurance \& Statutory
Standards}

\markright{Robotics, Safety Assurance \& Statutory Standards}

{\def\LTcaptype{none} % do not increment counter
\begin{longtable}[]{@{}
  >{\raggedright\arraybackslash}p{(\linewidth - 4\tabcolsep) * \real{0.3333}}
  >{\raggedright\arraybackslash}p{(\linewidth - 4\tabcolsep) * \real{0.3333}}
  >{\raggedright\arraybackslash}p{(\linewidth - 4\tabcolsep) * \real{0.3333}}@{}}
\toprule\noalign{}
\begin{minipage}[b]{\linewidth}\raggedright
Term
\end{minipage} & \begin{minipage}[b]{\linewidth}\raggedright
Operational Systems Definition
\end{minipage} & \begin{minipage}[b]{\linewidth}\raggedright
Primary Reference
\end{minipage} \\
\midrule\noalign{}
\endhead
\bottomrule\noalign{}
\endlastfoot
\textbf{Dynamic Stopping Distance (\(d_{\text{stop}}\))} & The total
physical clearance required to bring a moving mass to a complete halt
(\(d_{\text{stop}}(v) = v \cdot \tau_{\text{latency}} + \frac{v^2}{2 a_{\max}} + \delta_{\text{margin}}\)),
combining linear reaction latency displacement with quadratic braking
travel. & Chapter~\ref{sec-body}, Chapter~\ref{sec-enforcement} \\
\textbf{Inevitable Collision State (ICS)} & A physical state \((x, v)\)
from which no admissible control trajectory exists that can prevent a
collision with obstacles or environment boundaries, given the robot's
finite acceleration limits \(a_{\max}\). & Chapter~\ref{sec-body},
Chapter~\ref{sec-planning}, Chapter~\ref{sec-enforcement} \\
\textbf{Dynamic Fallback Ladder (MRM)} & A 4-tier structured escalation
protocol executed by the real-time supervisor when anomalies occur:
Level 1 (Minimal Intervention QP), Level 2 (Category 2 Position Hold),
Level 3 (Category 1 Controlled Dynamic Stop), Level 4 (Category 0 Safe
Torque Off with spring-applied mechanical brakes). &
Chapter~\ref{sec-enforcement}, Chapter~\ref{sec-intervention} \\
\textbf{Category 0 Stop (IEC 60204-1)} & An uncontrolled emergency stop
achieved by immediately disconnecting electrical power to machine
actuators, allowing gravity or spring-applied friction brakes to
dissipate energy. & Chapter~\ref{sec-body},
Chapter~\ref{sec-enforcement} \\
\textbf{Category 1 Stop (IEC 60204-1)} & A controlled stop where power
is maintained to actuators to execute a maximum-deceleration dynamic
brake trajectory until standstill, followed by complete power cutoff. &
Chapter~\ref{sec-body}, Chapter~\ref{sec-enforcement} \\
\textbf{Category 2 Stop (IEC 60204-1)} & A controlled stop where power
remains continuously supplied to machine actuators to maintain an active
position hold against external gravity or process loads. &
Chapter~\ref{sec-body}, Chapter~\ref{sec-enforcement} \\
\textbf{Envelope Trichotomy} & The formal safety case principle
partitioning runtime state queries into \emph{known-true} (verified
safe, permit action), \emph{known-false} (boundary breach, immediate
trip), and \emph{unknown} (epistemic uncertainty, sensor dropout).
Crucially, \textbf{unknown must be treated as strictly outside the
envelope}, triggering immediate degraded control. &
Chapter~\ref{sec-release} \\
\textbf{Claim-Argument-Evidence (CAE)} & A formal assurance framework
structuring safety cases: top-level claims are broken down into logical
sub-arguments supported by empirical test evidence from
hardware-in-the-loop and in-situ fault injection records. &
Chapter~\ref{sec-release} \\
\textbf{Residual-Claims Register} & An auditable engineering ledger
documenting every operational claim that cannot be empirically verified
by testing or proven by software contracts, binding each residual risk
to active mechanical containment and accountable human ownership. &
Chapter~\ref{sec-frontier} \\
\textbf{Detectorless Failure Mode} & A catastrophic physical failure
mechanism that produces no detectable change in electronic telemetry
before physical damage occurs (e.g., thermal rotor demagnetization,
subsurface gear pitting, piezoresistive sensor diaphragm fatigue).
Requires physical proof-testing rather than software monitoring. &
Chapter~\ref{sec-frontier} \\
\textbf{Covariate Shift Flywheel} & The vicious cycle in imitation
learning where compounding execution errors push the robot into
unfamiliar states outside demonstration support, triggering larger
errors that steer the machine directly toward hazards
(\(O(T^2 \epsilon)\) error accumulation). & Chapter~\ref{sec-data},
Chapter~\ref{sec-training} \\
\textbf{Demonstration Provenance Tuple} & The 4-stream action logging
protocol
\((a_{\text{req}}, a_{\text{cmd}}, a_{\text{enf}}, a_{\text{meas}})\)
that distinguishes intentional human teleop commands
(\(a_{\text{req}}\)) from kinematic transformations
(\(a_{\text{cmd}}\)), safety governor clamps (\(a_{\text{enf}}\)), and
structural compliance/backlash (\(a_{\text{meas}}\)). &
Chapter~\ref{sec-data} \\
\end{longtable}
}

\chapter*{Cross-Disciplinary Translation
Lexicon}\label{sec-appendix-lexicon}
\addcontentsline{toc}{chapter}{Cross-Disciplinary Translation Lexicon}

\markboth{Cross-Disciplinary Translation Lexicon}{Cross-Disciplinary
Translation Lexicon}

Physical AI sits at the intersection of four distinct engineering
traditions: \textbf{Modern Deep Learning (AI)}, \textbf{Classical
Robotics \& Non-Linear Control}, \textbf{Real-Time Embedded Systems \&
Silicon Architecture}, and \textbf{Safety-Critical Systems Assurance}.

Because each discipline developed its own vocabulary over decades,
engineers from different backgrounds frequently use the same words to
mean completely different things---or use different words for the exact
same underlying physical concept. This appendix serves as a
cross-disciplinary Rosetta Stone.

\begin{center}\rule{0.5\linewidth}{0.5pt}\end{center}

\section*{1. The Four-Tribe Rosetta
Stone}\label{the-four-tribe-rosetta-stone}
\addcontentsline{toc}{section}{1. The Four-Tribe Rosetta Stone}

\markright{1. The Four-Tribe Rosetta Stone}

{\def\LTcaptype{none} % do not increment counter
\begin{longtable}[]{@{}
  >{\raggedright\arraybackslash}p{(\linewidth - 8\tabcolsep) * \real{0.2000}}
  >{\raggedright\arraybackslash}p{(\linewidth - 8\tabcolsep) * \real{0.2000}}
  >{\raggedright\arraybackslash}p{(\linewidth - 8\tabcolsep) * \real{0.2000}}
  >{\raggedright\arraybackslash}p{(\linewidth - 8\tabcolsep) * \real{0.2000}}
  >{\raggedright\arraybackslash}p{(\linewidth - 8\tabcolsep) * \real{0.2000}}@{}}
\toprule\noalign{}
\begin{minipage}[b]{\linewidth}\raggedright
Physical Concept
\end{minipage} & \begin{minipage}[b]{\linewidth}\raggedright
Modern Deep Learning (AI)
\end{minipage} & \begin{minipage}[b]{\linewidth}\raggedright
Robotics \& Control Theory
\end{minipage} & \begin{minipage}[b]{\linewidth}\raggedright
Embedded Systems \& Silicon
\end{minipage} & \begin{minipage}[b]{\linewidth}\raggedright
Safety Assurance \& Standards
\end{minipage} \\
\midrule\noalign{}
\endhead
\bottomrule\noalign{}
\endlastfoot
\textbf{The Action Command} & Policy Output / Action Token / Diffusion
Chunk \(\hat{\mathbf{a}}_{t:t+H}\) & Control Input \(\mathbf{u}(t)\) /
Generalized Joint Torque \(\boldsymbol{\tau}\) & DAC Register / Timer
Capture-Compare Value / PWM Duty Cycle & Commanded Actuation Value /
Actuator Demand \\
\textbf{The Observation} & Input Tokens / Visual Embedding
\(\mathbf{z}_t\) / Context History & Generalized Coordinates
\(\mathbf{q}, \dot{\mathbf{q}}\) / State Vector \(\mathbf{x}(t)\) & ADC
Raw Buffer / SPI Sensor Frame / DMA Ring Packet & Measured Process
Variable / Transduced Feedback \\
\textbf{System Governing Rules} & Loss Function \(\mathcal{L}(\theta)\)
/ Objective Function / Reward \(R(s, a)\) & Plant Dynamics
\(\dot{\mathbf{x}} = f(\mathbf{x}, \mathbf{u})\) / Mass Matrix
\(\mathbf{M}(\mathbf{q})\) & Hardware Execution Model / Memory Map /
Interrupt Vector & System Specification / Operating Constraints \\
\textbf{Execution Horizon} & Autoregressive Sequence / Action Chunk
Length \(H\) & Receding Horizon \(N\) / Planning Horizon \(T_h\) &
Execution Period \(T_{\text{tick}}\) / Task Schedule Quantum &
Operational Cycle / Task Deadline \\
\textbf{Boundary Violation} & Out-of-Distribution (OOD) Shift / Mode
Collapse & Inevitable Collision State (ICS) / Loss of Stability &
Watchdog Trip / Hard Fault / Bus Error Exception & Hazard Event / Safety
Goal Violation (ISO 26262) \\
\textbf{Runtime Enforcement} & Action Space Clipping / Rejection
Sampling & Control Barrier Function (CBF-QP) / Reference Governor &
Memory Protection Unit (MPU) / Hardware Watchdog (WDT) & Safety Shield /
Independent Monitor Channel (E-Gas) \\
\textbf{Execution Delay} & Forward Pass Duration / Model Latency &
Transport Lag / Phase Delay \(e^{-s \tau}\) & Worst-Case Execution Time
(WCET) / Bus Contention & Fault Tolerant Time Interval (FTTI) / Latency
Budget \\
\textbf{Uncertainty Representation} & Softmax Entropy / Epistemic
Ensemble Variance \(\sigma^2\) & Estimation Covariance \(\mathbf{P}(t)\)
/ Disturbance Bound \(\mathcal{W}\) & ADC Quantization Noise / Sensor
Drift / Jitter & Confidence Bound (Clopper-Pearson) / SOTIF Triggering
Event \\
\textbf{Coordinate Transforms} & Cross-Attention Spatial Projection /
Camera Positional Embeddings & \(SE(3)\) Rigid Body Transform Tree
(\texttt{/tf}) / Lie Algebra \(\mathfrak{se}(3)\) & Extrinsic Matrix
Register / Sensor Offset Calibration ROM & Sensor Spatial Alignment /
Coordinate Reference Datum \\
\textbf{Optimization Method} & Stochastic Gradient Descent (SGD) / AdamW
Backpropagation & Quadratic Programming (QP) / Sequential Quadratic
Programming (SQP) & Fixed-Point Lookup Table / Iterative CORDIC Hardware
Core & Deterministic Algorithmic Solver / Validated Function Block \\
\textbf{Memory Architecture} & GPU High-Bandwidth Memory (HBM3 / LPDDR5)
/ KV Cache & Continuous State Covariance Matrix
\(\mathbf{P} \in \mathbb{R}^{n \times n}\) & Zero-Allocation DTCM SRAM /
Ring Buffer FIFO & Non-Volatile Fault Log / Redundant Shadow Memory \\
\textbf{Feedback Correction} & Policy Gradient / Temporal Difference
Error \(\delta_t\) & Error Dynamics
\(\mathbf{e}(t) = \mathbf{r}(t) - \mathbf{y}(t)\) / State Feedback
\(\mathbf{K}\mathbf{x}\) & Closed-Loop PID Interrupt / Current Shunt
Sense Loop & Closed-Loop Regulated Output / Mitigation Feedback \\
\end{longtable}
}

\begin{center}\rule{0.5\linewidth}{0.5pt}\end{center}

\section*{2. False Friends: Words That Cause Catastrophic
Misunderstandings}\label{false-friends-words-that-cause-catastrophic-misunderstandings}
\addcontentsline{toc}{section}{2. False Friends: Words That Cause
Catastrophic Misunderstandings}

\markright{2. False Friends: Words That Cause Catastrophic
Misunderstandings}

\subsection*{``Deterministic''}\label{deterministic}
\addcontentsline{toc}{subsection}{``Deterministic''}

\begin{itemize}
\tightlist
\item
  \textbf{In Deep Learning:} Setting \texttt{torch.manual\_seed(42)} and
  \texttt{cudnn.deterministic\ =\ True} so the matrix multiplication
  forward pass produces bit-identical tensor values across offline runs.
\item
  \textbf{In Embedded \& Real-Time Systems:} Guaranteeing that every
  periodic task completes in hardware time \(t \le T_{\text{deadline}}\)
  with bounded jitter (\(<1\,\mu\text{s}\)), regardless of CPU cache
  misses, interrupt storms, or bus contention. \emph{An algorithm that
  always outputs the same tensor 10 milliseconds too late is
  non-deterministic in physical systems.}
\end{itemize}

\subsection*{``Real-Time''}\label{real-time}
\addcontentsline{toc}{subsection}{``Real-Time''}

\begin{itemize}
\tightlist
\item
  \textbf{In Web \& Cloud Software:} A web service or UI that updates
  within \(100\text{--}200\text{ ms}\) (``real-time chat'', ``real-time
  streaming''), where occasional delays cause dropped frames or
  buffering spinners without danger.
\item
  \textbf{In Physical AI Systems:} Hard deadline determinism where
  missing a single \(1.0\text{ ms}\) execution deadline
  (\(1000\text{ Hz}\)) produces an inverter current loop runaway,
  mechanical resonance, or irreversible collision.
\end{itemize}

\subsection*{``State''}\label{state}
\addcontentsline{toc}{subsection}{``State''}

\begin{itemize}
\tightlist
\item
  \textbf{In Deep Learning:} The internal hidden activations of an RNN,
  the latent state of an autoencoder, or the Key-Value (KV) cache of an
  attention layer (\(\mathbf{h}_t \in \mathbb{R}^d\)). It has no
  explicit spatial units.
\item
  \textbf{In Robotics \& Physics:} The complete, physically grounded
  coordinate representation of the mechanical system
  (\(\mathbf{x} = [\mathbf{q}^T, \dot{\mathbf{q}}^T]^T \in \mathbb{R}^{2n}\))
  representing positions in meters, angles in radians, velocities in
  \(\text{m/s}\), and temperatures in Kelvin.
\end{itemize}

\subsection*{``Latency''}\label{latency}
\addcontentsline{toc}{subsection}{``Latency''}

\begin{itemize}
\tightlist
\item
  \textbf{In Machine Learning Benchmarks:} The average execution time of
  \texttt{model.forward()} on an isolated GPU under synthetic batching
  (\(P_{50} \approx 15\text{ ms}\)), ignoring camera frame exposure, OS
  scheduler queueing, and DMA memory bus transfer.
\item
  \textbf{In Physical AI:} The \textbf{Spatial Information Age}
  (\(\tau_{\text{age}}\)): the total wall-clock duration from the photon
  striking the camera photodiode at exposure midpoint
  \(t_{\text{exp}}/2\) to the instant mechanical torque is exerted on
  the physical environment. In real hardware, \(\tau_{\text{age}}\) is
  often \(3\times\) to \(5\times\) larger than the GPU forward pass.
\end{itemize}

\subsection*{``Safe''}\label{safe}
\addcontentsline{toc}{subsection}{``Safe''}

\begin{itemize}
\tightlist
\item
  \textbf{In AI Alignment:} The model avoids generating toxic language,
  hate speech, or copyrighted text.
\item
  \textbf{In Physical Systems:} The system maintains forward invariance
  with respect to physical kinetic boundaries (\(h(\mathbf{x}) \ge 0\)),
  guaranteeing that actuator torque, velocity, and deceleration remain
  mathematically incapable of inflicting mechanical trauma or structural
  destruction.
\end{itemize}

\subsection*{``Robustness''}\label{robustness}
\addcontentsline{toc}{subsection}{``Robustness''}

\begin{itemize}
\tightlist
\item
  \textbf{In Deep Learning:} Maintaining classification accuracy under
  adversarial pixel noise (\(\ell_\infty\) epsilon-balls) or image
  corruption (fog, blur).
\item
  \textbf{In Control \& Embedded Systems:} Maintaining closed-loop
  stability across plant parameter variations (gain/phase margins,
  \(\mathcal{H}_\infty\) norm bounds) and hardware physical resilience
  against electrical noise, bus EMI, power supply voltage droop, and
  thermal throttling.
\end{itemize}

\subsection*{``Policy''}\label{policy}
\addcontentsline{toc}{subsection}{``Policy''}

\begin{itemize}
\tightlist
\item
  \textbf{In Reinforcement Learning:} The parameterized neural network
  \(\pi_\theta(a \mid s)\) mapping states to action distributions.
\item
  \textbf{In Systems Engineering:} The top-level administrative or
  operational policy governing access rights, operator qualifications,
  and deployment sign-offs.
\end{itemize}

\begin{center}\rule{0.5\linewidth}{0.5pt}\end{center}

\section*{3. Practical Translation Guides for
Students}\label{practical-translation-guides-for-students}
\addcontentsline{toc}{section}{3. Practical Translation Guides for
Students}

\markright{3. Practical Translation Guides for Students}

\subsection*{If you are coming from Computer Science / Deep
Learning:}\label{if-you-are-coming-from-computer-science-deep-learning}
\addcontentsline{toc}{subsection}{If you are coming from Computer
Science / Deep Learning:}

\begin{enumerate}
\def\labelenumi{\arabic{enumi}.}
\tightlist
\item
  \textbf{Forget \texttt{malloc()} on the fast loop:} You cannot
  allocate dynamic memory or resize lists inside a \(1000\text{ Hz}\)
  control loop. Heap fragmentation causes random multi-millisecond
  pauses that crash physical robots.
\item
  \textbf{Your model is an unprivileged proposer:} A neural network must
  never have direct write access to motor registers. It proposes intent;
  a deterministic mathematical filter running on bare metal permits,
  clamps, or rejects the command.
\item
  \textbf{Loss \(= 0\) in simulation does not mean it works:} Physics
  engines approximate contact dynamics as linear complementarity
  problems (LCP). Real contact friction is non-smooth,
  stiction-dominated, and thermally variable.
\item
  \textbf{Bandwidth matters more than FLOPS:} Streaming multi-gigabyte
  neural weights from LPDDR5 DRAM on edge SoCs creates a memory wall
  (\(t_{\text{stream}} = M_{\text{weights}} / B_{\text{mem}}\)) that
  limits inference cadence, regardless of how many TFLOPS the NPU matrix
  core claims.
\end{enumerate}

\subsection*{If you are coming from Mechanical Engineering / Control
Theory:}\label{if-you-are-coming-from-mechanical-engineering-control-theory}
\addcontentsline{toc}{subsection}{If you are coming from Mechanical
Engineering / Control Theory:}

\begin{enumerate}
\def\labelenumi{\arabic{enumi}.}
\tightlist
\item
  \textbf{Uncertainty is multimodal:} Classical Kalman filters assume
  unimodal Gaussian noise. Perception models in the real world face
  discrete semantic hypotheses (e.g., an object is either in front of or
  behind a partition), requiring multimodal policy representations
  (Diffusion / Mixture Density Networks).
\item
  \textbf{You cannot write down the world dynamics \(f(\mathbf{x})\) for
  an unconstrained kitchen or highway:} Analytical ODEs work for steel
  linkages in a factory cell, but fail when interacting with deformable
  objects, varying lighting, and open-world human behavior. Deep models
  provide the semantic perception front-end that classical control
  lacks.
\item
  \textbf{Control Barrier Functions are the bridge:} You do not have to
  discard control theory when using deep learning. Instead, treat the
  learned policy as a nominal reference generator
  \(\mathbf{u}_{\text{ref}}\) and use a Control Barrier Function
  (CBF-QP) as an active safety filter that minimally adjusts
  \(\mathbf{u}\) only when boundary invariants are threatened.
\end{enumerate}

\subsection*{If you are coming from Embedded Systems / Electrical
Engineering:}\label{if-you-are-coming-from-embedded-systems-electrical-engineering}
\addcontentsline{toc}{subsection}{If you are coming from Embedded
Systems / Electrical Engineering:}

\begin{enumerate}
\def\labelenumi{\arabic{enumi}.}
\tightlist
\item
  \textbf{AI compute is dynamic and stochastic:} You cannot treat neural
  inference like a fixed-cycle DSP FIR filter. Execution time varies
  based on image content, caching, and thermal throttling. You must
  design asynchronous two-speed architectures with expiring leases to
  absorb timing jitter safely.
\item
  \textbf{Data provenance is as critical as voltage integrity:} Logging
  teleoperation data without tracking operator response delays
  (\(340\text{ ms}\)) or safety governor clamps poisons training
  pipelines via covariate shift.
\item
  \textbf{Memory crossbars are a shared resource:} When an NPU bursts
  \(25.6\text{ GB/s}\) of DRAM traffic during a vision transformer
  forward pass, it can starve real-time CPU DMA channels unless hardware
  AXI Quality of Service (QoS) bus arbitration is configured.
\end{enumerate}

\begin{center}\rule{0.5\linewidth}{0.5pt}\end{center}

\section*{4. The Three Canonical Archetypes
Quick-Reference}\label{the-three-canonical-archetypes-quick-reference}
\addcontentsline{toc}{section}{4. The Three Canonical Archetypes
Quick-Reference}

\markright{4. The Three Canonical Archetypes Quick-Reference}

Every systems concept in this textbook is grounded across the same three
physical machines:

\begin{Shaded}
\begin{Highlighting}[]
\NormalTok{flowchart LR}
\NormalTok{    subgraph A1["1. Mobile Archetype"]}
\NormalTok{        M1["Differential AGV / Mobile Robot"]}
\NormalTok{        M2["Dominant Limit: Momentum \& Stopping Distance"]}
\NormalTok{        M3["Key Hazard: Rolling{-}shutter skew \& traction loss"]}
\NormalTok{    end}
\NormalTok{    subgraph A2["2. Manipulator Archetype"]}
\NormalTok{        R1["6/7{-}DOF Articulated Arm"]}
\NormalTok{        R2["Dominant Limit: Reflected Inertia \& Contact Yield"]}
\NormalTok{        R3["Key Hazard: Seam torque impulse \& human impact"]}
\NormalTok{    end}
\NormalTok{    subgraph A3["3. Process / Desk Twin"]}
\NormalTok{        P1["Arduino UNO Q Dual{-}Brain Kit"]}
\NormalTok{        P2["Dominant Limit: Thermal Accumulation \& Bus Contention"]}
\NormalTok{        P3["Key Hazard: Memory starve \& watchdog timeout"]}
\NormalTok{    end}
\end{Highlighting}
\end{Shaded}

{\def\LTcaptype{none} % do not increment counter
\begin{longtable}[]{@{}
  >{\raggedright\arraybackslash}p{(\linewidth - 6\tabcolsep) * \real{0.2500}}
  >{\raggedright\arraybackslash}p{(\linewidth - 6\tabcolsep) * \real{0.2500}}
  >{\raggedright\arraybackslash}p{(\linewidth - 6\tabcolsep) * \real{0.2500}}
  >{\raggedright\arraybackslash}p{(\linewidth - 6\tabcolsep) * \real{0.2500}}@{}}
\toprule\noalign{}
\begin{minipage}[b]{\linewidth}\raggedright
Physical Metric
\end{minipage} & \begin{minipage}[b]{\linewidth}\raggedright
Mobile Archetype (AGV / Quadruped)
\end{minipage} & \begin{minipage}[b]{\linewidth}\raggedright
Manipulator Archetype (7-DOF Arm)
\end{minipage} & \begin{minipage}[b]{\linewidth}\raggedright
Process / Bench Twin (Arduino UNO Q)
\end{minipage} \\
\midrule\noalign{}
\endhead
\bottomrule\noalign{}
\endlastfoot
\textbf{Typical Total Mass (\(m\))} & \(15\text{--}250\text{ kg}\) &
\(15\text{--}45\text{ kg}\) (moving mass) & \(<1\text{ kg}\) (benchtop
kit) \\
\textbf{Operating Velocity (\(v\))} & \(1.0\text{--}3.5\text{ m/s}\) &
\(0.5\text{--}2.0\text{ m/s}\) (end-effector) & N/A (stationary
servo/stepper) \\
\textbf{Critical Bottleneck} & Braking distance
\(d_{\text{stop}} = \frac{v^2}{2 \mu g}\) under unknown tire friction
\(\mu\). & Reflected rotor inertia \(J_{\text{eff}} = J_L + N^2 J_m\)
under gearbox compliance. & Thermal dissipation
\(\Delta T = P \cdot \theta_{\text{JA}}\) and AXI QoS bus contention. \\
\textbf{Inner Enforcer Rate} & \(100\text{--}500\text{ Hz}\) (Traction
\& Yaw stability) & \(1000\text{ Hz}\) (Joint torque \& CBF safety
shield) & \(1000\text{ Hz}\) (Cortex-M4 PWM \& trip-zone registers) \\
\textbf{Outer Brain Rate} & \(5\text{--}20\text{ Hz}\) (VLA / Diffusion
path navigation) & \(10\text{--}50\text{ Hz}\) (ACT / Diffusion
trajectory chunking) & \(2\text{--}10\text{ Hz}\) (Edge VLM vision
tokenization) \\
\end{longtable}
}

\chapter*{Hardware, Fieldbus \& Actuator
Reference}\label{sec-appendix-hardware}
\addcontentsline{toc}{chapter}{Hardware, Fieldbus \& Actuator Reference}

\markboth{Hardware, Fieldbus \& Actuator Reference}{Hardware, Fieldbus
\& Actuator Reference}

This appendix provides a consolidated physical reference for the
electronic, communication, and electromechanical hardware substrates
encountered in Physical AI deployments and educational laboratory kits.

\begin{center}\rule{0.5\linewidth}{0.5pt}\end{center}

\section*{1. Dual-Brain System
Configurations}\label{dual-brain-system-configurations}
\addcontentsline{toc}{section}{1. Dual-Brain System Configurations}

\markright{1. Dual-Brain System Configurations}

\lettrine{P}{hysical} AI architectures decouple stochastic cognitive
reasoning from deterministic safety enforcement across two distinct
physical compute domains:

{\def\LTcaptype{none} % do not increment counter
\begin{longtable}[]{@{}
  >{\raggedright\arraybackslash}p{(\linewidth - 6\tabcolsep) * \real{0.2500}}
  >{\raggedright\arraybackslash}p{(\linewidth - 6\tabcolsep) * \real{0.2500}}
  >{\raggedright\arraybackslash}p{(\linewidth - 6\tabcolsep) * \real{0.2500}}
  >{\raggedright\arraybackslash}p{(\linewidth - 6\tabcolsep) * \real{0.2500}}@{}}
\toprule\noalign{}
\begin{minipage}[b]{\linewidth}\raggedright
Architecture Tier
\end{minipage} & \begin{minipage}[b]{\linewidth}\raggedright
Educational Dual-Brain Kit (Arduino UNO Q)
\end{minipage} & \begin{minipage}[b]{\linewidth}\raggedright
Industrial Robotics Platform (e.g., KUKA / Franka)
\end{minipage} & \begin{minipage}[b]{\linewidth}\raggedright
Autonomous Vehicle Compute Unit (e.g., Drive Orin)
\end{minipage} \\
\midrule\noalign{}
\endhead
\bottomrule\noalign{}
\endlastfoot
\textbf{Cognitive Brain (System 2)} & Qualcomm DragonBoard / Multi-core
Cortex-A (Debian Linux) & Industrial PC / Intel Xeon + NVIDIA RTX GPU
(Ubuntu RT-PREEMPT) & NVIDIA Tegra Orin SoC (12x Cortex-A78AE + Ampere
GPU + DLA) \\
\textbf{Safety Enforcer (System 1)} & ARM Cortex-M4 / STM32 (FreeRTOS /
Bare-Metal @ 1 kHz) & ARM Cortex-R5 / TI TMS570 (SafeRTOS / ASIL D MCU @
1 kHz) & Infineon AURIX TC397 / TC4x (TriCore Safe Lockstep MCU @ 1
kHz) \\
\textbf{Inter-Brain Physical Bus} & Shared SRAM Mailbox via RPMSG /
OpenAMP & Isolated PCIe Gen 4 with DMA or Dedicated Ethernet &
High-speed PCIe Gen 4 x4 + Dedicated Dual-SPI Heartbeat \\
\textbf{Safety Watchdog Mechanism} & Hardware GPIO pulse toggle
(\(t_{\text{timeout}} \le 3.0\text{ ms}\)) & Independent Dual-Channel
Safety Relay (ISO 13849 PLe) & Hardware System Safety Manager (SSM) with
Failsafe Watchdog \\
\textbf{Power Consumption (TDP)} & \(5\text{--}15\text{ W}\) &
\(150\text{--}450\text{ W}\) & \(60\text{--}250\text{ W}\) \\
\end{longtable}
}

\begin{center}\rule{0.5\linewidth}{0.5pt}\end{center}

\section*{2. Robotics Fieldbus \& Interconnect Comparison
Matrix}\label{robotics-fieldbus-interconnect-comparison-matrix}
\addcontentsline{toc}{section}{2. Robotics Fieldbus \& Interconnect
Comparison Matrix}

\markright{2. Robotics Fieldbus \& Interconnect Comparison Matrix}

{\def\LTcaptype{none} % do not increment counter
\begin{longtable}[]{@{}
  >{\raggedright\arraybackslash}p{(\linewidth - 12\tabcolsep) * \real{0.1429}}
  >{\raggedright\arraybackslash}p{(\linewidth - 12\tabcolsep) * \real{0.1429}}
  >{\raggedright\arraybackslash}p{(\linewidth - 12\tabcolsep) * \real{0.1429}}
  >{\raggedright\arraybackslash}p{(\linewidth - 12\tabcolsep) * \real{0.1429}}
  >{\raggedright\arraybackslash}p{(\linewidth - 12\tabcolsep) * \real{0.1429}}
  >{\raggedright\arraybackslash}p{(\linewidth - 12\tabcolsep) * \real{0.1429}}
  >{\raggedright\arraybackslash}p{(\linewidth - 12\tabcolsep) * \real{0.1429}}@{}}
\toprule\noalign{}
\begin{minipage}[b]{\linewidth}\raggedright
Fieldbus Protocol
\end{minipage} & \begin{minipage}[b]{\linewidth}\raggedright
Physical Layer
\end{minipage} & \begin{minipage}[b]{\linewidth}\raggedright
Max Payload
\end{minipage} & \begin{minipage}[b]{\linewidth}\raggedright
Bitrate / Bandwidth
\end{minipage} & \begin{minipage}[b]{\linewidth}\raggedright
\(P_{99}\) Latency
\end{minipage} & \begin{minipage}[b]{\linewidth}\raggedright
Hardware Determinism \& Jitter
\end{minipage} & \begin{minipage}[b]{\linewidth}\raggedright
Typical Physical AI Function
\end{minipage} \\
\midrule\noalign{}
\endhead
\bottomrule\noalign{}
\endlastfoot
\textbf{SPI (Serial Peripheral Interface)} & 4-wire CMOS (SCLK, MOSI,
MISO, CS) & Arbitrary byte stream & Up to \(50\text{ Mbps}\) &
\(<10\,\mu\text{s}\) & Hard Real-Time (\(<0.1\,\mu\text{s}\) jitter) &
On-board IMUs, joint encoders, ADC chips \\
\textbf{CAN 2.0B} & Differential pair (ISO 11898-2) & 8 bytes &
\(1.0\text{ Mbps}\) & \(0.5\text{--}2.0\text{ ms}\) &
Priority-arbitrated (\(<50\,\mu\text{s}\) jitter) & Legacy automotive
sensors, low-speed motor drives \\
\textbf{CAN-FD} & Differential pair & 64 bytes &
\(5.0\text{--}8.0\text{ Mbps}\) & \(100\text{--}350\,\mu\text{s}\) &
Priority-arbitrated (\(<20\,\mu\text{s}\) jitter) & Robot joint drives,
battery management, chassis control \\
\textbf{EtherCAT} & 100BASE-TX Ethernet & Up to 1486 bytes (on-the-fly)
& \(100\text{ Mbps}\) & \(50\text{--}125\,\mu\text{s}\) & Hard Real-Time
(\(<1\,\mu\text{s}\) distributed clock) & High-speed multi-axis
industrial motor synchronization \\
\textbf{MIPI CSI-2} & Differential D-PHY / C-PHY & Uncompressed frame &
Up to \(10\text{ Gbps}\) & Frame rate (\(16\text{--}33\text{ ms}\)) &
Deterministic DMA push & Edge vision and depth sensor raw ingestion \\
\textbf{PCIe Gen 4} & Point-to-point differential lanes & Packets up to
4 KB & \(16\text{ GT/s}\) per lane (\(2.0\text{ GB/s/lane}\)) &
\(1.5\text{--}5.0\,\mu\text{s}\) & High-throughput DMA (Subject to UMA
contention) & Host AP to GPU / NPU accelerator tensor handoffs \\
\end{longtable}
}

\begin{center}\rule{0.5\linewidth}{0.5pt}\end{center}

\section*{3. Actuator Transmission Mechanics \& Reflected
Inertia}\label{actuator-transmission-mechanics-reflected-inertia}
\addcontentsline{toc}{section}{3. Actuator Transmission Mechanics \&
Reflected Inertia}

\markright{3. Actuator Transmission Mechanics \& Reflected Inertia}

In physical actuation, a motor's rotor inertia \(J_{\text{rotor}}\) is
amplified by the square of the gear reduction ratio \(N\):

\[J_{\text{eff}} = J_{\text{load}} + N^2 J_{\text{rotor}}\]

{\def\LTcaptype{none} % do not increment counter
\begin{longtable}[]{@{}
  >{\raggedright\arraybackslash}p{(\linewidth - 10\tabcolsep) * \real{0.1667}}
  >{\raggedright\arraybackslash}p{(\linewidth - 10\tabcolsep) * \real{0.1667}}
  >{\raggedright\arraybackslash}p{(\linewidth - 10\tabcolsep) * \real{0.1667}}
  >{\raggedright\arraybackslash}p{(\linewidth - 10\tabcolsep) * \real{0.1667}}
  >{\raggedright\arraybackslash}p{(\linewidth - 10\tabcolsep) * \real{0.1667}}
  >{\raggedright\arraybackslash}p{(\linewidth - 10\tabcolsep) * \real{0.1667}}@{}}
\toprule\noalign{}
\begin{minipage}[b]{\linewidth}\raggedright
Transmission Type
\end{minipage} & \begin{minipage}[b]{\linewidth}\raggedright
Gear Ratio Range (\(N\))
\end{minipage} & \begin{minipage}[b]{\linewidth}\raggedright
Backdrivability \& Compliance
\end{minipage} & \begin{minipage}[b]{\linewidth}\raggedright
Reflected Inertia Impact (\(N^2\))
\end{minipage} & \begin{minipage}[b]{\linewidth}\raggedright
Torque Density (\(\text{N}\cdot\text{m/kg}\))
\end{minipage} & \begin{minipage}[b]{\linewidth}\raggedright
Typical Physical AI Application
\end{minipage} \\
\midrule\noalign{}
\endhead
\bottomrule\noalign{}
\endlastfoot
\textbf{Direct Drive} & \(1:1\) & Perfect (No friction) & Minimal
(\(1\times\)) & Low (\(1\text{--}3\)) & High-bandwidth haptic
teleoperation, direct-drive turntables \\
\textbf{Quasi-Direct Drive (QDD)} & \(4:1\text{--}10:1\) & High (Low
friction) & Low (\(16\times\text{--}100\times\)) & Moderate
(\(15\text{--}30\)) & Dynamic legged locomotion (quadrupeds), agile
humanoids \\
\textbf{Planetary Reducer} & \(10:1\text{--}50:1\) & Moderate & High
(\(100\times\text{--}2500\times\)) & High (\(30\text{--}60\)) & Mobile
robot wheels, industrial robotic base joints \\
\textbf{Strain Wave (Harmonic)} & \(50:1\text{--}160:1\) & Low /
Non-backdrivable & Severe (\(2500\times\text{--}25600\times\)) & Very
High (\(80\text{--}150\)) & High-precision articulated robot arms,
surgical manipulators \\
\textbf{Cycloidal Drive} & \(30:1\text{--}100:1\) & Moderate /
Shock-resistant & High (\(900\text{--}10000\times\)) & Exceptional
(\(100\text{--}200\)) & Heavy-payload manipulator wrists, industrial
AGVs \\
\end{longtable}
}

\begin{center}\rule{0.5\linewidth}{0.5pt}\end{center}

\section*{4. Power Delivery Network (PDN) \& Voltage
Droop}\label{power-delivery-network-pdn-voltage-droop}
\addcontentsline{toc}{section}{4. Power Delivery Network (PDN) \&
Voltage Droop}

\markright{4. Power Delivery Network (PDN) \& Voltage Droop}

When a deep neural network accelerator switches from an idle state to
full-load matrix tensor computation, the step current demand
\(\Delta I\) creates a transient voltage droop across the power
distribution network:

\[\Delta V_{\text{droop}} = L_{\text{trace}} \frac{dI}{dt} + R_{\text{trace}} \Delta I\]

\begin{itemize}
\tightlist
\item
  \textbf{The Brownout Risk:} If an edge GPU draws a
  \(\Delta I = 80\text{ A}\) current step in \(10\text{ ns}\) across an
  unconditioned parasitic inductance
  \(L_{\text{trace}} = 0.5\text{ nH}\), the transient droop
  \(\Delta V = 40.0\text{ V}\) can collapse a shared \(48\text{ V}\) DC
  bus below the minimum operating threshold of the microcontroller's
  low-dropout (LDO) regulator (\(<3.3\text{ V}\)), triggering a local
  processor reset.
\item
  \textbf{Architectural Mitigation:} Physical AI systems must employ
  galvanic DC-DC power isolation and dedicated bulk decoupling
  capacitors
  (\(C_{\text{decouple}} \ge \frac{\Delta I \cdot \Delta t}{\Delta V_{\text{allowable}}}\))
  between cognitive accelerators and real-time microcontroller buses.
\end{itemize}

\begin{center}\rule{0.5\linewidth}{0.5pt}\end{center}

\section*{5. Actuator Thermal Rise \& Resistance
Drift}\label{actuator-thermal-rise-resistance-drift}
\addcontentsline{toc}{section}{5. Actuator Thermal Rise \& Resistance
Drift}

\markright{5. Actuator Thermal Rise \& Resistance Drift}

The electrical resistance of copper motor windings increases linearly
with operating temperature:

\[R_{\text{winding}}(T) = R_0 \left( 1 + \alpha_{\text{Cu}} (T - T_0) \right)\]

where \(\alpha_{\text{Cu}} \approx 0.00393\text{ K}^{-1}\) at
\(T_0 = 20^\circ\text{C}\).

\begin{itemize}
\tightlist
\item
  \textbf{Consequence for Control:} As a motor heats from
  \(20^\circ\text{C}\) to its maximum thermal limit
  (\(120^\circ\text{C}\)), coil resistance increases by \(+39.3\%\).
  Under voltage-mode control, the achievable phase current and output
  torque drop proportionally by \(28.2\%\), violating nominal kinematic
  reachability assumptions unless closed-loop Field-Oriented Current
  Control (FOC) is enforced at \(>10\text{ kHz}\).
\end{itemize}

\chapter*{Incident Autopsies \& Statutory Standards
Reference}\label{sec-appendix-standards}
\addcontentsline{toc}{chapter}{Incident Autopsies \& Statutory Standards
Reference}

\markboth{Incident Autopsies \& Statutory Standards Reference}{Incident
Autopsies \& Statutory Standards Reference}

This appendix provides a structured reference summarizing the forensic
incident autopsies and statutory functional safety standards analyzed
throughout this text.

\begin{center}\rule{0.5\linewidth}{0.5pt}\end{center}

\section*{1. Master Compendium of Incident
Autopsies}\label{master-compendium-of-incident-autopsies}
\addcontentsline{toc}{section}{1. Master Compendium of Incident
Autopsies}

\markright{1. Master Compendium of Incident Autopsies}

\lettrine{E}{very} engineering principle in Physical AI is illuminated
by real-world physical failures where software abstractions collided
with physical dynamics.

{\def\LTcaptype{none} % do not increment counter
\begin{longtable}[]{@{}
  >{\raggedright\arraybackslash}p{(\linewidth - 10\tabcolsep) * \real{0.1667}}
  >{\raggedright\arraybackslash}p{(\linewidth - 10\tabcolsep) * \real{0.1667}}
  >{\raggedright\arraybackslash}p{(\linewidth - 10\tabcolsep) * \real{0.1667}}
  >{\raggedright\arraybackslash}p{(\linewidth - 10\tabcolsep) * \real{0.1667}}
  >{\raggedright\arraybackslash}p{(\linewidth - 10\tabcolsep) * \real{0.1667}}
  >{\raggedright\arraybackslash}p{(\linewidth - 10\tabcolsep) * \real{0.1667}}@{}}
\toprule\noalign{}
\begin{minipage}[b]{\linewidth}\raggedright
Case Name \& Date
\end{minipage} & \begin{minipage}[b]{\linewidth}\raggedright
System \& Application
\end{minipage} & \begin{minipage}[b]{\linewidth}\raggedright
Primary Root Cause
\end{minipage} & \begin{minipage}[b]{\linewidth}\raggedright
Critical Architectural Flaw
\end{minipage} & \begin{minipage}[b]{\linewidth}\raggedright
Provable Physical Lesson
\end{minipage} & \begin{minipage}[b]{\linewidth}\raggedright
Primary Reference
\end{minipage} \\
\midrule\noalign{}
\endhead
\bottomrule\noalign{}
\endlastfoot
\textbf{Uber ATG Tempe Crash} (Mar 2018) & Autonomous Test Vehicle
(Volvo XC90) & Semantic classification toggled between vehicle, bicycle,
and unknown, resetting track history on each cycle & Automatic Emergency
Braking (AEB) suppressed for \(1.2\text{ s}\) to prevent false-alarm
jerk; no deterministic low-level override & High-level perception must
never gate or suppress the low-level kinematic stopping envelope
(\(d_{\text{stop}}\)). & Chapter~\ref{sec-boundary},
Chapter~\ref{sec-brain}, Chapter~\ref{sec-intervention},
Chapter~\ref{sec-release} \\
\textbf{Tesla Autopilot Williston Crash} (May 2016) & Semi-Autonomous
Highway Autopilot & Monocular vision saturated against bright midday
sky; radar reflection rejected as overhead sign & System lacked
cross-modal sensor voting and failed to enforce human gaze attention
invariants & Sensor fusion must explicitly bound epistemic uncertainty;
vision cannot be assumed universally valid under high-dynamic-range
illumination. & Chapter~\ref{sec-perception},
Chapter~\ref{sec-intervention} \\
\textbf{Autonomous Gantry Milling Collision} (Industrial Autopsy) &
5-Axis Heavy CNC Gantry Milling Center & Flash trace logging triggered
high-priority DMA bursts across shared AHB crossbar, starving real-time
SPI bus & Emergency stop task suffered \(600\,\mu\text{s}\) bus
contention, missing \(3.0\text{ ms}\) deceleration deadline &
Safety-critical enforcer code and I/O registers must be isolated in
zero-contention Data Tightly-Coupled Memory (DTCM). &
Chapter~\ref{sec-nervous}, Chapter~\ref{sec-enforcement},
Chapter~\ref{sec-placement} \\
\textbf{Gas Turbine Hydraulic Manifold Rupture} (Energy Plant) &
Closed-Loop Fuel Manifold Pressure Regulator & Low-pass anti-aliasing
filter added \(40\text{ ms}\) phase lag to pressure telemetry, causing
PI controller windup & Control loop lacked dynamic saturation clamping
and rate-of-rise barrier certificates & Low-pass filtering trades noise
reduction for phase lag; transport latency consumes physical stability
margin. & Chapter~\ref{sec-body}, Chapter~\ref{sec-nervous} \\
\textbf{Mars Climate Orbiter Disintegration} (Sep 1999) & Interplanetary
Spacecraft Trajectory Controller & Ground software emitted angular
momentum thruster impulses in pound-force seconds
(\(\text{lbf}\cdot\text{s}\)), while onboard software expected
Newton-seconds (\(\text{N}\cdot\text{s}\)) & Lack of statically typed
dimensional units and coordinate frame contracts in IPC schemas & Every
data structure crossing an interface must explicitly bind physical SI
units and reference frames into immutable schemas. &
Chapter~\ref{sec-data}, Chapter~\ref{sec-perception},
Chapter~\ref{sec-intent} \\
\textbf{Therac-25 Radiation Overdoses} (1985--1987) & Medical Linear
Accelerator & Microsecond operator race conditions and arithmetic
counter overflows bypassed software position checks & Software-only
interlocks replaced physical mechanical interlock switches on the
turntable & Software decision logic cannot replace independent hardware
break circuits and physical trip zones. &
Chapter~\ref{sec-enforcement} \\
\textbf{Ariane 5 Flight 501 Destruction} (Jun 1996) & Heavy Space Launch
Vehicle & 64-bit floating-point horizontal bias converted to 16-bit
signed integer without overflow protection & Reused flight software
emitted diagnostic dump onto shared 1553 bus, interpreted as valid
flight guidance & Reused software components must have explicit
precondition assertions; shared communication buses require strict
spatial and temporal isolation. & Chapter~\ref{sec-placement} \\
\textbf{Boeing 737 MAX MCAS Fatal Crashes} (2018--2019) & Commercial Jet
Transport Flight Control & Single unvalidated Angle-of-Attack (AoA)
sensor vane failure repeatedly triggered automated nose-down trim & Lack
of cross-channel sensor validation; lack of cockpit annunciation; MCAS
overrode pilot control column inputs & Automated flight software must
not possess unmediated authority to overpower physical human control
forces without explicit authorization and disengagement. &
Chapter~\ref{sec-intervention} \\
\textbf{Airbus A320 Warsaw Runway Overrun} (Sep 1993) & Commercial Jet
Aircraft Braking System & Flight computers required simultaneous squat
switch compression and wheel spin-up to confirm ground contact &
Aquaplaning prevented both sensor conditions; flight computer locked out
pilot manual braking and thrust reversers & Rigid boolean software
interlocks must not lock out human emergency intervention during dynamic
edge-case scenarios. & Chapter~\ref{sec-intervention} \\
\textbf{Toyota Unintended Acceleration} (2007--2010) & Electronic
Throttle Control System (ETCS-i) & Single-Event Upset or stack overflow
killed the real-time throttle control task (\texttt{Task\_Throttle}) &
Hardware watchdog was unconditionally kicked by background OS task; dead
throttle task left engine locked at wide-open throttle & Hardware
watchdogs must monitor task-specific progress tokens across all
real-time tasks, not generic CPU liveness. &
Chapter~\ref{sec-verification} \\
\textbf{Space Shuttle Challenger Disaster} (Jan 1986) & Solid Rocket
Booster (SRB) Field Joint & Launch temperature of \(31^\circ\text{F}\)
exceeded Viton O-ring glass transition temperature, delaying resilience
recovery & Engineering management normalized secondary blowby as
acceptable risk under schedule pressure & Statistical survival of past
out-of-envelope operations does not expand certified physical design
envelopes (\emph{normalization of deviance}). &
Chapter~\ref{sec-release} \\
\textbf{Deepwater Horizon BOP Containment Failure} (Apr 2010) & Subsea
Blowout Preventer (BOP) Blind Shear Rams & High-velocity blowout flow
caused drill pipe Euler buckling, deflecting pipe laterally against
casing wall & Blind shear rams were designed only for centered pipe;
pressure telemetry reported nominal actuation while pipe jammed
unsevered & Observational indistinguishability (\(y_A \equiv y_B\))
renders detectorless failure modes invisible to software, requiring
structural containment and physical proof-tests. &
Chapter~\ref{sec-frontier} \\
\textbf{Apollo 11 Radar Cycle Stealing} (Jul 1969) & Apollo Guidance
Computer (AGC) Lunar Module & Rendezvous radar in manual slew flooded
AGC with 800-Hz asynchronous reference pulses (\(15\%\) CPU stolen) &
Priority-scheduled core memory dropped low-priority display tasks to
preserve throttle/thruster control cycles & Deterministic fixed-priority
scheduling and graceful task shedding preserve real-time control
deadlines under bus overload. & Chapter~\ref{sec-verification} \\
\end{longtable}
}

\begin{center}\rule{0.5\linewidth}{0.5pt}\end{center}

\section*{2. Statutory Functional Safety Standards
Alignment}\label{statutory-functional-safety-standards-alignment}
\addcontentsline{toc}{section}{2. Statutory Functional Safety Standards
Alignment}

\markright{2. Statutory Functional Safety Standards Alignment}

{\def\LTcaptype{none} % do not increment counter
\begin{longtable}[]{@{}
  >{\raggedright\arraybackslash}p{(\linewidth - 6\tabcolsep) * \real{0.2500}}
  >{\raggedright\arraybackslash}p{(\linewidth - 6\tabcolsep) * \real{0.2500}}
  >{\raggedright\arraybackslash}p{(\linewidth - 6\tabcolsep) * \real{0.2500}}
  >{\raggedright\arraybackslash}p{(\linewidth - 6\tabcolsep) * \real{0.2500}}@{}}
\toprule\noalign{}
\begin{minipage}[b]{\linewidth}\raggedright
Safety Standard
\end{minipage} & \begin{minipage}[b]{\linewidth}\raggedright
Target Industry \& Domain
\end{minipage} & \begin{minipage}[b]{\linewidth}\raggedright
Core Architectural Requirements \& Integrity Metrics
\end{minipage} & \begin{minipage}[b]{\linewidth}\raggedright
Physical AI Mapping \& Implementation
\end{minipage} \\
\midrule\noalign{}
\endhead
\bottomrule\noalign{}
\endlastfoot
\textbf{ISO 26262} & Road Vehicles (Functional Safety) & • Automotive
Safety Integrity Levels (ASIL A to ASIL D)• Single-Point Fault Metric
(SPFM \(\ge 99\%\) for ASIL D)• Fault Tolerant Time Interval (FTTI
\(\le 50\text{ ms}\))• E-Gas 3-Level Monitoring Architecture & Maps
directly to the \textbf{Proposal--Permission Split}: Level 1 (Learned
Brain proposal), Level 2 (Trajectory cross-check), Level 3 (Hardware MCU
watchdog and galvanic E-Stop). \\
\textbf{IEC 61508 / IEC 62061} & Industrial Automation \&
Safety-Critical E/E/PE & • Safety Integrity Levels (SIL 1 to SIL 4)•
Target failure rate \(<10^{-9}\text{ failures/hour}\) (SIL 4)• Safe
Failure Fraction (\(\text{SFF} \ge 99\%\)) & Governs the hardware
implementation of the deterministic safety enforcer and physical safe
torque off (STO) stages. \\
\textbf{ISO 21448 (SOTIF)} & Safety of the Intended Functionality & •
Mitigating hazards caused by performance limitations, sensor blindness,
and complex neural perception algorithms without hardware faults•
Shrinking Area 2 (Known Unsafe) and Area 3 (Unknown Unsafe) & Governs
the \textbf{Envelope Trichotomy} and Claim-Argument-Evidence safety
cases in Chapter~\ref{sec-release}. \\
\textbf{ISO 10218-1/2 \& ISO/TS 15066} & Collaborative Robots (Cobots) &
• Power and Force Limiting (PFL: Transient contact \(<50\text{ N}\))•
Speed and Separation Monitoring (SSM: Dynamic distance
\(S \ge v_r T_r + v_h T_h + B\))• 3-position enabling devices (teach
pendants) & Supplies the quantitative force and velocity barrier bounds
for Control Barrier Function (CBF-QP) filters in
Chapter~\ref{sec-enforcement}. \\
\textbf{UL 4600} & Autonomous Products \& Systems & • Goal Structuring
Notation (GSN) safety cases• Prohibits purely statistical / road-hour
safety claims• Mandates explicit residual-claims ledgers and proof of
non-probabilistic safety envelopes & Provides the philosophical
foundation for the \textbf{Astronomical Exposure Wall} and deductive
Claim-Argument-Evidence safety cases. \\
\textbf{IEC 60204-1} & Safety of Machinery (Electrical Equipment) & •
\textbf{Stop Category 0:} Immediate power cutoff (Safe Torque Off /
mechanical brake drop)• \textbf{Stop Category 1:} Controlled
deceleration to stop with power on, then cutoff• \textbf{Stop Category
2:} Controlled position hold with power maintained & Defines the formal
3-tier physical fallback actions executed by the enforcer during anomaly
escalation in Chapter~\ref{sec-enforcement}. \\
\end{longtable}
}

\chapter*{Engineering Notebook
Templates}\label{sec-appendix-notebook-templates}
\addcontentsline{toc}{chapter}{Engineering Notebook Templates}

\markboth{Engineering Notebook Templates}{Engineering Notebook
Templates}

This appendix provides clean, standardized templates for the core
engineering schemas and contracts that compose the \textbf{Physical AI
engineering notebook}. Every chapter in the curriculum defines or
freezes one of these schemas.

\section*{Loop Charter --- Chapter 1}\label{loop-charter-chapter-1}
\addcontentsline{toc}{section}{Loop Charter --- Chapter 1}

\markright{Loop Charter --- Chapter 1}

\begin{Shaded}
\begin{Highlighting}[]
\FunctionTok{loop\_charter}\KeywordTok{:}
\AttributeTok{  }\FunctionTok{version}\KeywordTok{:}\AttributeTok{ }\StringTok{"1.0"}
\AttributeTok{  }\FunctionTok{system\_name}\KeywordTok{:}\AttributeTok{ }\StringTok{"Physical Agent Handling Station"}
\AttributeTok{  }\FunctionTok{physical\_boundary}\KeywordTok{:}
\AttributeTok{    }\FunctionTok{operating\_space}\KeywordTok{:}\AttributeTok{ }\StringTok{"Tabletop 50cm x 50cm workspace"}
\AttributeTok{    }\FunctionTok{actuation\_degrees\_of\_freedom}\KeywordTok{:}\AttributeTok{ }\DecValTok{2}
\AttributeTok{    }\FunctionTok{maximum\_kinetic\_energy\_joules}\KeywordTok{:}\AttributeTok{ }\FloatTok{1.2}
\AttributeTok{  }\FunctionTok{learned\_components}\KeywordTok{:}
\AttributeTok{    }\KeywordTok{{-}}\AttributeTok{ }\FunctionTok{name}\KeywordTok{:}\AttributeTok{ }\StringTok{"VisionAffordanceEncoder"}
\AttributeTok{      }\FunctionTok{role}\KeywordTok{:}\AttributeTok{ }\StringTok{"Propose 3D grasp coordinates"}
\AttributeTok{      }\FunctionTok{privilege\_level}\KeywordTok{:}\AttributeTok{ }\StringTok{"UNTRUSTED\_PROPOSAL"}
\AttributeTok{  }\FunctionTok{authority\_boundary}\KeywordTok{:}
\AttributeTok{    }\FunctionTok{autonomous\_action\_allowed}\KeywordTok{:}\AttributeTok{ }\CharTok{true}
\AttributeTok{    }\FunctionTok{human\_in\_proximity}\KeywordTok{:}\AttributeTok{ }\CharTok{true}
\AttributeTok{    }\FunctionTok{e\_stop\_mechanism}\KeywordTok{:}\AttributeTok{ }\StringTok{"Hardware{-}wired cut to motor driver"}
\AttributeTok{  }\FunctionTok{invariants}\KeywordTok{:}
\AttributeTok{    }\KeywordTok{{-}}\AttributeTok{ }\StringTok{"No motor PWM permitted without active MCU watchdog tick"}
\AttributeTok{    }\KeywordTok{{-}}\AttributeTok{ }\StringTok{"Motion prohibited outside 50cm bounding cube"}
\end{Highlighting}
\end{Shaded}

\section*{Requirements and Physical Budgets Ledger --- Chapter
2}\label{requirements-and-physical-budgets-ledger-chapter-2}
\addcontentsline{toc}{section}{Requirements and Physical Budgets Ledger
--- Chapter 2}

\markright{Requirements and Physical Budgets Ledger --- Chapter 2}

\begin{Shaded}
\begin{Highlighting}[]
\FunctionTok{physical\_budgets\_ledger}\KeywordTok{:}
\AttributeTok{  }\FunctionTok{version}\KeywordTok{:}\AttributeTok{ }\StringTok{"1.0"}
\AttributeTok{  }\FunctionTok{freshness\_deadline\_delta\_t\_fresh\_ms}\KeywordTok{:}\AttributeTok{ }\FloatTok{40.0}
\AttributeTok{  }\FunctionTok{maximum\_velocity\_v\_max\_m\_s}\KeywordTok{:}\AttributeTok{ }\FloatTok{1.80}
\AttributeTok{  }\FunctionTok{emergency\_brake\_decel\_a\_brake\_m\_s2}\KeywordTok{:}\AttributeTok{ }\FloatTok{3.00}
\AttributeTok{  }\FunctionTok{reflected\_inertia\_multiplier\_N2}\KeywordTok{:}\AttributeTok{ }\FloatTok{2500.0}
\AttributeTok{  }\FunctionTok{stator\_continuous\_temp\_ceiling\_C}\KeywordTok{:}\AttributeTok{ }\FloatTok{140.0}
\AttributeTok{  }\FunctionTok{maximum\_bus\_droop\_delta\_V\_V}\KeywordTok{:}\AttributeTok{ }\FloatTok{10.0}
\AttributeTok{  }\FunctionTok{assumptions}\KeywordTok{:}
\AttributeTok{    }\KeywordTok{{-}}\AttributeTok{ }\StringTok{"Ambient internal chassis temperature \textless{}= 45°C"}
\AttributeTok{    }\KeywordTok{{-}}\AttributeTok{ }\StringTok{"Surface friction coefficient mu \textgreater{}= 0.40"}
\end{Highlighting}
\end{Shaded}

\section*{Nervous System Inter-Processor Mailbox Schema --- Chapter
4}\label{nervous-system-inter-processor-mailbox-schema-chapter-4}
\addcontentsline{toc}{section}{Nervous System Inter-Processor Mailbox
Schema --- Chapter 4}

\markright{Nervous System Inter-Processor Mailbox Schema --- Chapter 4}

\begin{Shaded}
\begin{Highlighting}[]
\FunctionTok{nervous\_ipc\_mailbox}\KeywordTok{:}
\AttributeTok{  }\FunctionTok{version}\KeywordTok{:}\AttributeTok{ }\StringTok{"2.1.0"}
\AttributeTok{  }\FunctionTok{header}\KeywordTok{:}
\AttributeTok{    }\FunctionTok{sequence\_version}\KeywordTok{:}\AttributeTok{ }\DecValTok{10429}
\AttributeTok{    }\FunctionTok{timestamp\_epoch\_ns}\KeywordTok{:}\AttributeTok{ }\DecValTok{1724832000123456789}
\AttributeTok{    }\FunctionTok{sender\_core\_id}\KeywordTok{:}\AttributeTok{ }\StringTok{"MPU\_LINUX\_CORE\_0"}
\AttributeTok{    }\FunctionTok{lease\_validity\_us}\KeywordTok{:}\AttributeTok{ }\DecValTok{25000}
\AttributeTok{  }\FunctionTok{payload}\KeywordTok{:}
\AttributeTok{    }\FunctionTok{target\_joint\_positions\_rad}\KeywordTok{:}\AttributeTok{ }\KeywordTok{[}\FloatTok{0.12}\KeywordTok{,}\AttributeTok{ }\FloatTok{{-}0.45}\KeywordTok{,}\AttributeTok{ }\FloatTok{1.20}\KeywordTok{,}\AttributeTok{ }\FloatTok{0.05}\KeywordTok{,}\AttributeTok{ }\FloatTok{{-}0.80}\KeywordTok{,}\AttributeTok{ }\FloatTok{0.32}\KeywordTok{]}
\AttributeTok{    }\FunctionTok{feedforward\_torques\_nm}\KeywordTok{:}\AttributeTok{ }\KeywordTok{[}\FloatTok{1.2}\KeywordTok{,}\AttributeTok{ }\FloatTok{4.5}\KeywordTok{,}\AttributeTok{ }\FloatTok{2.1}\KeywordTok{,}\AttributeTok{ }\FloatTok{0.4}\KeywordTok{,}\AttributeTok{ }\FloatTok{0.8}\KeywordTok{,}\AttributeTok{ }\FloatTok{0.2}\KeywordTok{]}
\AttributeTok{    }\FunctionTok{cbf\_safety\_margin\_m}\KeywordTok{:}\AttributeTok{ }\FloatTok{0.085}
\AttributeTok{  }\FunctionTok{integrity}\KeywordTok{:}
\AttributeTok{    }\FunctionTok{crc32\_checksum}\KeywordTok{:}\AttributeTok{ }\StringTok{"0x4B3A8F91"}
\end{Highlighting}
\end{Shaded}

\section*{Demonstration Dataset Provenance Schema --- Chapter
5}\label{demonstration-dataset-provenance-schema-chapter-5}
\addcontentsline{toc}{section}{Demonstration Dataset Provenance Schema
--- Chapter 5}

\markright{Demonstration Dataset Provenance Schema --- Chapter 5}

\begin{Shaded}
\begin{Highlighting}[]
\FunctionTok{dataset\_provenance}\KeywordTok{:}
\AttributeTok{  }\FunctionTok{version}\KeywordTok{:}\AttributeTok{ }\StringTok{"1.0"}
\AttributeTok{  }\FunctionTok{dataset\_id}\KeywordTok{:}\AttributeTok{ }\StringTok{"DS{-}ALOHA{-}INSERTION{-}V2"}
\AttributeTok{  }\FunctionTok{collection\_modality}\KeywordTok{:}\AttributeTok{ }\StringTok{"BILATERAL\_TELEOPERATION"}
\AttributeTok{  }\FunctionTok{teleoperation\_loop\_latency\_ms}\KeywordTok{:}\AttributeTok{ }\FloatTok{42.0}
\AttributeTok{  }\FunctionTok{state\_action\_stream\_taps}\KeywordTok{:}
\AttributeTok{    }\FunctionTok{logged\_streams}\KeywordTok{:}\AttributeTok{ }\KeywordTok{[}\StringTok{"a\_req"}\KeywordTok{,}\AttributeTok{ }\StringTok{"a\_cmd"}\KeywordTok{,}\AttributeTok{ }\StringTok{"a\_enf"}\KeywordTok{,}\AttributeTok{ }\StringTok{"a\_meas"}\KeywordTok{]}
\AttributeTok{    }\FunctionTok{synchronization\_clock}\KeywordTok{:}\AttributeTok{ }\StringTok{"IEEE\_1588\_PTP"}
\AttributeTok{  }\FunctionTok{intervention\_policy}\KeywordTok{:}
\AttributeTok{    }\FunctionTok{causal\_divergence\_cutoff\_t\_div\_ms}\KeywordTok{:}\AttributeTok{ }\FloatTok{300.0}
\AttributeTok{  }\FunctionTok{environment\_envelope}\KeywordTok{:}
\AttributeTok{    }\FunctionTok{ambient\_temp\_range\_C}\KeywordTok{:}\AttributeTok{ }\KeywordTok{[}\FloatTok{18.0}\KeywordTok{,}\AttributeTok{ }\FloatTok{26.0}\KeywordTok{]}
\AttributeTok{    }\FunctionTok{payload\_mass\_kg}\KeywordTok{:}\AttributeTok{ }\FloatTok{0.25}
\end{Highlighting}
\end{Shaded}

\section*{Physical Policy Training Lineage Manifest --- Chapter
6}\label{physical-policy-training-lineage-manifest-chapter-6}
\addcontentsline{toc}{section}{Physical Policy Training Lineage Manifest
--- Chapter 6}

\markright{Physical Policy Training Lineage Manifest --- Chapter 6}

\begin{Shaded}
\begin{Highlighting}[]
\FunctionTok{policy\_training\_manifest}\KeywordTok{:}
\AttributeTok{  }\FunctionTok{version}\KeywordTok{:}\AttributeTok{ }\StringTok{"1.0"}
\AttributeTok{  }\FunctionTok{policy\_id}\KeywordTok{:}\AttributeTok{ }\StringTok{"POL{-}DIFFUSION{-}ACT{-}V4"}
\AttributeTok{  }\FunctionTok{synthesis\_regime}\KeywordTok{:}\AttributeTok{ }\StringTok{"DIFFUSION\_POLICY\_CHUNKED"}
\AttributeTok{  }\FunctionTok{action\_chunk\_horizon\_H}\KeywordTok{:}\AttributeTok{ }\DecValTok{32}
\AttributeTok{  }\FunctionTok{temporal\_ensembling\_k}\KeywordTok{:}\AttributeTok{ }\FloatTok{0.01}
\AttributeTok{  }\FunctionTok{training\_dataset\_id}\KeywordTok{:}\AttributeTok{ }\StringTok{"DS{-}ALOHA{-}INSERTION{-}V2"}
\AttributeTok{  }\FunctionTok{sim\_to\_real\_reality\_gap\_bounds}\KeywordTok{:}
\AttributeTok{    }\FunctionTok{latency\_delta\_ms}\KeywordTok{:}\AttributeTok{ }\FloatTok{12.0}
\AttributeTok{    }\FunctionTok{stiction\_friction\_margin\_pct}\KeywordTok{:}\AttributeTok{ }\FloatTok{15.0}
\AttributeTok{  }\FunctionTok{target\_hardware\_deployment}\KeywordTok{:}
\AttributeTok{    }\FunctionTok{host\_mpu}\KeywordTok{:}\AttributeTok{ }\StringTok{"Jetson\_AGX\_Orin\_64GB"}
\AttributeTok{    }\FunctionTok{inference\_cadence\_hz}\KeywordTok{:}\AttributeTok{ }\DecValTok{20}
\end{Highlighting}
\end{Shaded}

\section*{Policy Pre-Deployment Evaluation Record --- Chapter
7}\label{policy-pre-deployment-evaluation-record-chapter-7}
\addcontentsline{toc}{section}{Policy Pre-Deployment Evaluation Record
--- Chapter 7}

\markright{Policy Pre-Deployment Evaluation Record --- Chapter 7}

\begin{Shaded}
\begin{Highlighting}[]
\FunctionTok{evaluation\_record}\KeywordTok{:}
\AttributeTok{  }\FunctionTok{version}\KeywordTok{:}\AttributeTok{ }\StringTok{"1.0"}
\AttributeTok{  }\FunctionTok{policy\_id}\KeywordTok{:}\AttributeTok{ }\StringTok{"POL{-}DIFFUSION{-}ACT{-}V4"}
\AttributeTok{  }\FunctionTok{trial\_suite}\KeywordTok{:}
\AttributeTok{    }\FunctionTok{independent\_trials\_count}\KeywordTok{:}\AttributeTok{ }\DecValTok{500}
\AttributeTok{    }\FunctionTok{zero\_failure\_confidence\_level}\KeywordTok{:}\AttributeTok{ }\FloatTok{0.95}
\AttributeTok{    }\FunctionTok{exact\_clopper\_pearson\_lower\_bound}\KeywordTok{:}\AttributeTok{ }\FloatTok{0.994}
\AttributeTok{  }\FunctionTok{worst\_case\_tail\_latency\_p99\_9\_ms}\KeywordTok{:}\AttributeTok{ }\FloatTok{38.2}
\AttributeTok{  }\FunctionTok{runtime\_reject\_option}\KeywordTok{:}
\AttributeTok{    }\FunctionTok{entropy\_threshold}\KeywordTok{:}\AttributeTok{ }\FloatTok{0.85}
\AttributeTok{    }\FunctionTok{fallback\_action}\KeywordTok{:}\AttributeTok{ }\StringTok{"TRIGGER\_HOLDING\_REFLEX"}
\end{Highlighting}
\end{Shaded}

\section*{Observation Contract Schema --- Chapter
8}\label{observation-contract-schema-chapter-8}
\addcontentsline{toc}{section}{Observation Contract Schema --- Chapter
8}

\markright{Observation Contract Schema --- Chapter 8}

\begin{Shaded}
\begin{Highlighting}[]
\FunctionTok{observation\_contract}\KeywordTok{:}
\AttributeTok{  }\FunctionTok{version}\KeywordTok{:}\AttributeTok{ }\StringTok{"1.0"}
\AttributeTok{  }\FunctionTok{sensors}\KeywordTok{:}
\AttributeTok{    }\KeywordTok{{-}}\AttributeTok{ }\FunctionTok{sensor\_id}\KeywordTok{:}\AttributeTok{ }\StringTok{"CAM\_WRIST\_RGB"}
\AttributeTok{      }\FunctionTok{modality}\KeywordTok{:}\AttributeTok{ }\StringTok{"RGB"}
\AttributeTok{      }\FunctionTok{resolution}\KeywordTok{:}\AttributeTok{ }\KeywordTok{[}\DecValTok{640}\KeywordTok{,}\AttributeTok{ }\DecValTok{480}\KeywordTok{]}
\AttributeTok{      }\FunctionTok{frame\_rate\_hz}\KeywordTok{:}\AttributeTok{ }\DecValTok{60}
\AttributeTok{      }\FunctionTok{transport}\KeywordTok{:}\AttributeTok{ }\StringTok{"MIPI\_CSI2\_DMA"}
\AttributeTok{      }\FunctionTok{timestamp\_sync}\KeywordTok{:}\AttributeTok{ }\StringTok{"PTP\_IEEE1588"}
\AttributeTok{      }\FunctionTok{latency\_budget\_ms}\KeywordTok{:}\AttributeTok{ }\FloatTok{16.6}
\AttributeTok{    }\KeywordTok{{-}}\AttributeTok{ }\FunctionTok{sensor\_id}\KeywordTok{:}\AttributeTok{ }\StringTok{"JOINT\_ENCODER\_BUS"}
\AttributeTok{      }\FunctionTok{modality}\KeywordTok{:}\AttributeTok{ }\StringTok{"Absolute Optical"}
\AttributeTok{      }\FunctionTok{sample\_rate\_hz}\KeywordTok{:}\AttributeTok{ }\DecValTok{1000}
\AttributeTok{      }\FunctionTok{transport}\KeywordTok{:}\AttributeTok{ }\StringTok{"EtherCAT"}
\AttributeTok{      }\FunctionTok{latency\_budget\_ms}\KeywordTok{:}\AttributeTok{ }\FloatTok{1.0}
\AttributeTok{  }\FunctionTok{fallbacks}\KeywordTok{:}
\AttributeTok{    }\FunctionTok{max\_dropped\_frames\_before\_halt}\KeywordTok{:}\AttributeTok{ }\DecValTok{2}
\AttributeTok{    }\FunctionTok{sensor\_dropout\_action}\KeywordTok{:}\AttributeTok{ }\StringTok{"TRIGGER\_STATE\_DRIFT\_ALARM"}
\end{Highlighting}
\end{Shaded}

\section*{Spatial State Belief Record --- Chapter
9}\label{spatial-state-belief-record-chapter-9}
\addcontentsline{toc}{section}{Spatial State Belief Record --- Chapter
9}

\markright{Spatial State Belief Record --- Chapter 9}

\begin{Shaded}
\begin{Highlighting}[]
\FunctionTok{spatial\_state\_record}\KeywordTok{:}
\AttributeTok{  }\FunctionTok{version}\KeywordTok{:}\AttributeTok{ }\StringTok{"1.0"}
\AttributeTok{  }\FunctionTok{state\_id}\KeywordTok{:}\AttributeTok{ }\StringTok{"EST\_END\_EFFECTOR\_SE3"}
\AttributeTok{  }\FunctionTok{evidence\_epoch\_t0\_ns}\KeywordTok{:}\AttributeTok{ }\DecValTok{1724832000100000000}
\AttributeTok{  }\FunctionTok{reference\_frame}\KeywordTok{:}\AttributeTok{ }\StringTok{"REP105\_BASE\_LINK"}
\AttributeTok{  }\FunctionTok{state\_estimate}\KeywordTok{:}\AttributeTok{ }\KeywordTok{[}\FloatTok{0.450}\KeywordTok{,}\AttributeTok{ }\FloatTok{0.120}\KeywordTok{,}\AttributeTok{ }\FloatTok{0.310}\KeywordTok{,}\AttributeTok{ }\FloatTok{0.0}\KeywordTok{,}\AttributeTok{ }\FloatTok{0.707}\KeywordTok{,}\AttributeTok{ }\FloatTok{0.0}\KeywordTok{,}\AttributeTok{ }\FloatTok{0.707}\KeywordTok{]}
\AttributeTok{  }\FunctionTok{covariance\_diagonal\_P}\KeywordTok{:}\AttributeTok{ }\KeywordTok{[}\FloatTok{0.0004}\KeywordTok{,}\AttributeTok{ }\FloatTok{0.0004}\KeywordTok{,}\AttributeTok{ }\FloatTok{0.0009}\KeywordTok{,}\AttributeTok{ }\FloatTok{0.001}\KeywordTok{,}\AttributeTok{ }\FloatTok{0.001}\KeywordTok{,}\AttributeTok{ }\FloatTok{0.001}\KeywordTok{]}
\AttributeTok{  }\FunctionTok{lyapunov\_growth\_rate\_q\_c}\KeywordTok{:}\AttributeTok{ }\FloatTok{0.002}
\AttributeTok{  }\FunctionTok{invalidation\_timeout\_ms}\KeywordTok{:}\AttributeTok{ }\FloatTok{95.0}
\AttributeTok{  }\FunctionTok{lifecycle\_status}\KeywordTok{:}\AttributeTok{ }\StringTok{"VALID"}
\end{Highlighting}
\end{Shaded}

\section*{Intent Lease Contract Schema --- Chapter
10}\label{intent-lease-contract-schema-chapter-10}
\addcontentsline{toc}{section}{Intent Lease Contract Schema --- Chapter
10}

\markright{Intent Lease Contract Schema --- Chapter 10}

\begin{Shaded}
\begin{Highlighting}[]
\FunctionTok{intent\_lease}\KeywordTok{:}
\AttributeTok{  }\FunctionTok{version}\KeywordTok{:}\AttributeTok{ }\StringTok{"1.0"}
\AttributeTok{  }\FunctionTok{lease\_id}\KeywordTok{:}\AttributeTok{ }\StringTok{"LEASE{-}2026{-}08{-}28{-}0042"}
\AttributeTok{  }\FunctionTok{issuing\_model}\KeywordTok{:}\AttributeTok{ }\StringTok{"VLA\_Spatial\_Planner"}
\AttributeTok{  }\FunctionTok{target\_pose\_se3}\KeywordTok{:}\AttributeTok{ }\KeywordTok{[}\FloatTok{0.520}\KeywordTok{,}\AttributeTok{ }\FloatTok{0.000}\KeywordTok{,}\AttributeTok{ }\FloatTok{0.150}\KeywordTok{,}\AttributeTok{ }\FloatTok{0.0}\KeywordTok{,}\AttributeTok{ }\FloatTok{1.0}\KeywordTok{,}\AttributeTok{ }\FloatTok{0.0}\KeywordTok{,}\AttributeTok{ }\FloatTok{0.0}\KeywordTok{]}
\AttributeTok{  }\FunctionTok{spatial\_tolerance\_radius\_m}\KeywordTok{:}\AttributeTok{ }\FloatTok{0.020}
\AttributeTok{  }\FunctionTok{wrench\_limit\_normal\_force\_N}\KeywordTok{:}\AttributeTok{ }\FloatTok{45.0}
\AttributeTok{  }\FunctionTok{time\_to\_live\_ttl\_ms}\KeywordTok{:}\AttributeTok{ }\FloatTok{250.0}
\AttributeTok{  }\FunctionTok{silence\_action}\KeywordTok{:}\AttributeTok{ }\StringTok{"FAILSAFE\_DECELERATION\_TO\_ZERO"}
\end{Highlighting}
\end{Shaded}

\section*{Trajectory Proposal and Execution Contract --- Chapter
11}\label{trajectory-proposal-and-execution-contract-chapter-11}
\addcontentsline{toc}{section}{Trajectory Proposal and Execution
Contract --- Chapter 11}

\markright{Trajectory Proposal and Execution Contract --- Chapter 11}

\begin{Shaded}
\begin{Highlighting}[]
\FunctionTok{trajectory\_proposal\_chunk}\KeywordTok{:}
\AttributeTok{  }\FunctionTok{version}\KeywordTok{:}\AttributeTok{ }\StringTok{"1.0"}
\AttributeTok{  }\FunctionTok{chunk\_horizon\_steps}\KeywordTok{:}\AttributeTok{ }\DecValTok{32}
\AttributeTok{  }\FunctionTok{step\_period\_dt\_ms}\KeywordTok{:}\AttributeTok{ }\FloatTok{10.0}
\AttributeTok{  }\FunctionTok{spline\_continuity\_order}\KeywordTok{:}\AttributeTok{ }\StringTok{"C2\_CONTINUOUS"}
\AttributeTok{  }\FunctionTok{jerk\_limit\_rad\_s3}\KeywordTok{:}\AttributeTok{ }\FloatTok{150.0}
\AttributeTok{  }\FunctionTok{precomputed\_stopping\_suffix}\KeywordTok{:}
\AttributeTok{    }\FunctionTok{braking\_duration\_ms}\KeywordTok{:}\AttributeTok{ }\FloatTok{85.0}
\AttributeTok{    }\FunctionTok{terminal\_deceleration\_rad\_s2}\KeywordTok{:}\AttributeTok{ }\FloatTok{{-}35.0}
\AttributeTok{    }\FunctionTok{safe\_holding\_torque\_nm}\KeywordTok{:}\AttributeTok{ }\FloatTok{0.0}
\end{Highlighting}
\end{Shaded}

\section*{Enforcement Design and Simplex Filter --- Chapter
12}\label{enforcement-design-and-simplex-filter-chapter-12}
\addcontentsline{toc}{section}{Enforcement Design and Simplex Filter ---
Chapter 12}

\markright{Enforcement Design and Simplex Filter --- Chapter 12}

\begin{Shaded}
\begin{Highlighting}[]
\FunctionTok{enforcement\_design}\KeywordTok{:}
\AttributeTok{  }\FunctionTok{version}\KeywordTok{:}\AttributeTok{ }\StringTok{"1.0"}
\AttributeTok{  }\FunctionTok{mcu\_filter\_frequency\_hz}\KeywordTok{:}\AttributeTok{ }\DecValTok{1000}
\AttributeTok{  }\FunctionTok{barrier\_functions}\KeywordTok{:}
\AttributeTok{    }\KeywordTok{{-}}\AttributeTok{ }\FunctionTok{name}\KeywordTok{:}\AttributeTok{ }\StringTok{"WorkspaceGeofence"}
\AttributeTok{      }\FunctionTok{type}\KeywordTok{:}\AttributeTok{ }\StringTok{"CONTROL\_BARRIER\_FUNCTION"}
\AttributeTok{      }\FunctionTok{inviolable\_bound}\KeywordTok{:}\AttributeTok{ }\StringTok{"x\^{}2 + y\^{}2 \textless{}= (0.50)\^{}2"}
\AttributeTok{    }\KeywordTok{{-}}\AttributeTok{ }\FunctionTok{name}\KeywordTok{:}\AttributeTok{ }\StringTok{"DynamicStoppingClearance"}
\AttributeTok{      }\FunctionTok{inviolable\_bound}\KeywordTok{:}\AttributeTok{ }\StringTok{"d\_clearance \textgreater{}= v * t\_delay + v\^{}2 / (2 * a\_max) + delta\_margin"}
\AttributeTok{  }\FunctionTok{tracking\_error\_inset\_margin\_m}\KeywordTok{:}\AttributeTok{ }\FloatTok{0.015}
\AttributeTok{  }\FunctionTok{fallbacks}\KeywordTok{:}
\AttributeTok{    }\FunctionTok{tier\_1}\KeywordTok{:}\AttributeTok{ }\StringTok{"CBF\_QP\_ACTIVE\_INTERVENTION"}
\AttributeTok{    }\FunctionTok{tier\_2}\KeywordTok{:}\AttributeTok{ }\StringTok{"BASELINE\_DECELERATION\_SPLINE"}
\AttributeTok{    }\FunctionTok{tier\_3}\KeywordTok{:}\AttributeTok{ }\StringTok{"DYNAMIC\_BACK\_EMF\_BRAKING"}
\AttributeTok{    }\FunctionTok{tier\_4}\KeywordTok{:}\AttributeTok{ }\StringTok{"STO\_HARDWARE\_POWER\_TRIP"}
\end{Highlighting}
\end{Shaded}

\section*{Heterogeneous Silicon Placement Matrix --- Chapter
13}\label{heterogeneous-silicon-placement-matrix-chapter-13}
\addcontentsline{toc}{section}{Heterogeneous Silicon Placement Matrix
--- Chapter 13}

\markright{Heterogeneous Silicon Placement Matrix --- Chapter 13}

\begin{Shaded}
\begin{Highlighting}[]
\FunctionTok{silicon\_placement\_matrix}\KeywordTok{:}
\AttributeTok{  }\FunctionTok{version}\KeywordTok{:}\AttributeTok{ }\StringTok{"1.0"}
\AttributeTok{  }\FunctionTok{target\_soc}\KeywordTok{:}\AttributeTok{ }\StringTok{"HETERO\_DUAL\_BRAIN\_V3"}
\AttributeTok{  }\FunctionTok{core\_allocations}\KeywordTok{:}
\AttributeTok{    }\FunctionTok{unprivileged\_mpu}\KeywordTok{:}\AttributeTok{ }\KeywordTok{[}\StringTok{"Linux\_Cortex\_A78AE\_0\_3"}\KeywordTok{,}\AttributeTok{ }\StringTok{"Inference\_NPU"}\KeywordTok{]}
\AttributeTok{    }\FunctionTok{real\_time\_mcu}\KeywordTok{:}\AttributeTok{ }\KeywordTok{[}\StringTok{"Bare\_Metal\_Cortex\_R52\_DualLockstep"}\KeywordTok{]}
\AttributeTok{  }\FunctionTok{memory\_partitioning}\KeywordTok{:}
\AttributeTok{    }\FunctionTok{arm\_mpam\_npu\_bandwidth\_limit\_pct}\KeywordTok{:}\AttributeTok{ }\FloatTok{70.0}
\AttributeTok{    }\FunctionTok{dedicated\_mcu\_sram\_tcm\_kb}\KeywordTok{:}\AttributeTok{ }\DecValTok{512}
\AttributeTok{  }\FunctionTok{power\_rail\_integrity}\KeywordTok{:}
\AttributeTok{    }\FunctionTok{maximum\_transient\_droop\_mv}\KeywordTok{:}\AttributeTok{ }\FloatTok{50.0}
\AttributeTok{    }\FunctionTok{maximum\_di\_dt\_A\_us}\KeywordTok{:}\AttributeTok{ }\FloatTok{15.0}
\end{Highlighting}
\end{Shaded}

\section*{Forensic Authority Record Schema --- Chapter
14}\label{forensic-authority-record-schema-chapter-14}
\addcontentsline{toc}{section}{Forensic Authority Record Schema ---
Chapter 14}

\markright{Forensic Authority Record Schema --- Chapter 14}

\begin{Shaded}
\begin{Highlighting}[]
\FunctionTok{forensic\_authority\_record}\KeywordTok{:}
\AttributeTok{  }\FunctionTok{version}\KeywordTok{:}\AttributeTok{ }\StringTok{"1.0"}
\AttributeTok{  }\FunctionTok{takeover\_event\_id}\KeywordTok{:}\AttributeTok{ }\StringTok{"TO{-}2026{-}08{-}28{-}0914"}
\AttributeTok{  }\FunctionTok{arbitration\_state}\KeywordTok{:}\AttributeTok{ }\StringTok{"BLEND\_TO\_MANUAL"}
\AttributeTok{  }\FunctionTok{physiological\_reaction\_lag\_ms}\KeywordTok{:}\AttributeTok{ }\FloatTok{920.0}
\AttributeTok{  }\FunctionTok{bumpless\_transfer\_ramp\_duration\_ms}\KeywordTok{:}\AttributeTok{ }\FloatTok{180.0}
\AttributeTok{  }\FunctionTok{pre\_takeover\_divergence\_window\_ms}\KeywordTok{:}\AttributeTok{ }\FloatTok{500.0}
\AttributeTok{  }\FunctionTok{cryptographic\_auth}\KeywordTok{:}
\AttributeTok{    }\FunctionTok{hmac\_sha256\_digest}\KeywordTok{:}\AttributeTok{ }\StringTok{"0xE9A4...3B7C"}
\AttributeTok{    }\FunctionTok{authenticated\_operator\_id}\KeywordTok{:}\AttributeTok{ }\StringTok{"OP{-}CERT{-}904"}
\end{Highlighting}
\end{Shaded}

\section*{Hardware Qualification \& Fault Execution Record --- Chapter
15}\label{hardware-qualification-fault-execution-record-chapter-15}
\addcontentsline{toc}{section}{Hardware Qualification \& Fault Execution
Record --- Chapter 15}

\markright{Hardware Qualification \& Fault Execution Record --- Chapter
15}

\begin{Shaded}
\begin{Highlighting}[]
\FunctionTok{fault\_execution\_record}\KeywordTok{:}
\AttributeTok{  }\FunctionTok{version}\KeywordTok{:}\AttributeTok{ }\StringTok{"1.0"}
\AttributeTok{  }\FunctionTok{test\_campaign\_id}\KeywordTok{:}\AttributeTok{ }\StringTok{"HIL{-}QUAL{-}CAMPAIGN{-}08"}
\AttributeTok{  }\FunctionTok{qualification\_stage}\KeywordTok{:}\AttributeTok{ }\StringTok{"HIL\_HARDWARE\_IN\_LOOP"}
\AttributeTok{  }\FunctionTok{injected\_fault\_type}\KeywordTok{:}\AttributeTok{ }\StringTok{"SHARED\_AXI\_BUS\_SATURATION\_AND\_FRAME\_DROP"}
\AttributeTok{  }\FunctionTok{injection\_point}\KeywordTok{:}\AttributeTok{ }\StringTok{"MIPI\_CSI2\_DMA\_CONTROLLER"}
\AttributeTok{  }\FunctionTok{detector\_channel}\KeywordTok{:}\AttributeTok{ }\StringTok{"MCU\_HEARTBEAT\_WATCHDOG\_INTERRUPT"}
\AttributeTok{  }\FunctionTok{detection\_latency\_us}\KeywordTok{:}\AttributeTok{ }\FloatTok{840.0}
\AttributeTok{  }\FunctionTok{containment\_action}\KeywordTok{:}\AttributeTok{ }\StringTok{"TIER\_2\_SAFE\_BRAKING\_SPLINE"}
\AttributeTok{  }\FunctionTok{containment\_verdict}\KeywordTok{:}\AttributeTok{ }\StringTok{"PASS\_ZERO\_BOUNDARY\_BREACH"}
\end{Highlighting}
\end{Shaded}

\section*{Claim-Argument-Evidence Release Dossier --- Chapter
16}\label{claim-argument-evidence-release-dossier-chapter-16}
\addcontentsline{toc}{section}{Claim-Argument-Evidence Release Dossier
--- Chapter 16}

\markright{Claim-Argument-Evidence Release Dossier --- Chapter 16}

\begin{Shaded}
\begin{Highlighting}[]
\FunctionTok{release\_dossier}\KeywordTok{:}
\AttributeTok{  }\FunctionTok{version}\KeywordTok{:}\AttributeTok{ }\StringTok{"1.0"}
\AttributeTok{  }\FunctionTok{system\_name}\KeywordTok{:}\AttributeTok{ }\StringTok{"Collaborative Tactile Manipulator Cell"}
\AttributeTok{  }\FunctionTok{release\_verdict}\KeywordTok{:}\AttributeTok{ }\StringTok{"OPERATE\_WITH\_CONDITIONS"}
\AttributeTok{  }\FunctionTok{odd\_specification}\KeywordTok{:}
\AttributeTok{    }\FunctionTok{domain\_id}\KeywordTok{:}\AttributeTok{ }\StringTok{"ODD{-}ISO15066{-}COLLAB{-}01"}
\AttributeTok{    }\FunctionTok{contact\_force\_limit\_N}\KeywordTok{:}\AttributeTok{ }\FloatTok{50.0}
\AttributeTok{    }\FunctionTok{cartesian\_speed\_cap\_m\_s}\KeywordTok{:}\AttributeTok{ }\FloatTok{0.35}
\AttributeTok{  }\FunctionTok{evidence\_summary}\KeywordTok{:}
\AttributeTok{    }\FunctionTok{hil\_fault\_injection\_pass\_rate}\KeywordTok{:}\AttributeTok{ }\FloatTok{1.0}\CommentTok{ \# 500/500 faults safely contained}
\AttributeTok{    }\FunctionTok{zero\_allocation\_audit}\KeywordTok{:}\AttributeTok{ }\StringTok{"CONFIRMED\_PASS"}
\AttributeTok{    }\FunctionTok{tail\_latency\_p99\_99\_ms}\KeywordTok{:}\AttributeTok{ }\FloatTok{1.80}
\AttributeTok{  }\FunctionTok{adjudicator\_signoff}\KeywordTok{:}
\AttributeTok{    }\FunctionTok{systems\_lead}\KeywordTok{:}\AttributeTok{ }\StringTok{"APPROVED\_SHA256\_0x8F3D7B02"}
\AttributeTok{    }\FunctionTok{safety\_lead}\KeywordTok{:}\AttributeTok{ }\StringTok{"APPROVED\_SHA256\_0xA1C09E44"}
\end{Highlighting}
\end{Shaded}

\section*{Residual-Claims Register Schema --- Chapter
17}\label{residual-claims-register-schema-chapter-17}
\addcontentsline{toc}{section}{Residual-Claims Register Schema ---
Chapter 17}

\markright{Residual-Claims Register Schema --- Chapter 17}

\begin{Shaded}
\begin{Highlighting}[]
\FunctionTok{residual\_claims\_register}\KeywordTok{:}
\AttributeTok{  }\FunctionTok{version}\KeywordTok{:}\AttributeTok{ }\StringTok{"1.0"}
\AttributeTok{  }\FunctionTok{register\_entries}\KeywordTok{:}
\AttributeTok{    }\KeywordTok{{-}}\AttributeTok{ }\FunctionTok{entry\_id}\KeywordTok{:}\AttributeTok{ }\StringTok{"RCR{-}01{-}SURFACE{-}FRICTION"}
\AttributeTok{      }\FunctionTok{missing\_physical\_observable}\KeywordTok{:}\AttributeTok{ }\StringTok{"Sub{-}millisecond dynamic friction coefficient on contaminated epoxy floor"}
\AttributeTok{      }\FunctionTok{active\_containment\_constraint}\KeywordTok{:}\AttributeTok{ }\StringTok{"Firmware speed cap v \textless{}= 0.80 m/s ensuring stopping distance \textless{}= 0.40 m even under mu = 0.15"}
\AttributeTok{      }\FunctionTok{named\_accountable\_owner}\KeywordTok{:}\AttributeTok{ }\StringTok{"Lead Controls Architect"}
\AttributeTok{      }\FunctionTok{predeclared\_closure\_test}\KeywordTok{:}\AttributeTok{ }\StringTok{"Tactile acoustic slip sensing HIL characterization campaign"}
\end{Highlighting}
\end{Shaded}

\chapter*{Heterogeneous Dual-Brain SoC Reference
Architecture}\label{sec-appendix-heterogeneous-soc}
\addcontentsline{toc}{chapter}{Heterogeneous Dual-Brain SoC Reference
Architecture}

\markboth{Heterogeneous Dual-Brain SoC Reference
Architecture}{Heterogeneous Dual-Brain SoC Reference Architecture}

This appendix defines the canonical reference architecture for
heterogeneous dual-brain edge systems in Physical AI. It specifies the
hardware topology, core partitioning, inter-processor communication
(IPC) protocols, and hardware safety watchdog contracts that separate
unprivileged stochastic policy deliberation from deterministic real-time
safety enforcement.

\section*{Dual-Brain System Topology}\label{dual-brain-system-topology}
\addcontentsline{toc}{section}{Dual-Brain System Topology}

\markright{Dual-Brain System Topology}

\lettrine{I}{n} an embodied machine, cognitive deliberation and
mechanical actuation operate in fundamentally incompatible timing and
determinism domains. The dual-brain architecture bridges this divide by
mapping functional roles onto asymmetric processing domains on a shared
System-on-Chip (SoC) or coupled compute module:

{\def\LTcaptype{none} % do not increment counter
\begin{longtable}[]{@{}
  >{\raggedright\arraybackslash}p{(\linewidth - 4\tabcolsep) * \real{0.3333}}
  >{\raggedright\arraybackslash}p{(\linewidth - 4\tabcolsep) * \real{0.3333}}
  >{\raggedright\arraybackslash}p{(\linewidth - 4\tabcolsep) * \real{0.3333}}@{}}
\toprule\noalign{}
\begin{minipage}[b]{\linewidth}\raggedright
Processor Subsystem
\end{minipage} & \begin{minipage}[b]{\linewidth}\raggedright
Architectural Core and OS
\end{minipage} & \begin{minipage}[b]{\linewidth}\raggedright
Primary Systems Role and Execution Contract
\end{minipage} \\
\midrule\noalign{}
\endhead
\bottomrule\noalign{}
\endlastfoot
\textbf{Application Processor (MPU)} & Multi-core 64-bit Application
CPURich OS (Embedded Linux / RT-Linux) & \textbf{Untrusted Proposal
Service:} Runs multimodal perception backbones, foundation models, ROS 2
stacks, and generative trajectory planners; emits timestamped, expiring
candidate proposals (\(p_t\)). \\
\textbf{Real-Time Safety Controller (MCU)} & Deterministic
Microcontroller (ARM Cortex-M / RISC-V)Bare-Metal / Hard RTOS (Zephyr /
FreeRTOS) & \textbf{Trusted Permission Authority:} Runs deterministic
\(1000\text{ Hz}\) control loops, Control Barrier Function safety
filters (\(h(x) \ge 0\)), dynamic stopping enforcers
(\(d_{\text{stop}}\)), and hardware PWM/timer registers. \\
\textbf{Inter-Processor Link (IPC)} & Shared SRAM Mailbox \& Hardware
SemaphoreRPMSG / OpenAMP Framework & Zero-copy lock-free circular ring
buffers with atomic sequence counters, cache-coherent headers, and
hardware mailbox interrupts. \\
\end{longtable}
}

\section*{Hardware Resource
Partitioning}\label{hardware-resource-partitioning}
\addcontentsline{toc}{section}{Hardware Resource Partitioning}

\markright{Hardware Resource Partitioning}

{\def\LTcaptype{none} % do not increment counter
\begin{longtable}[]{@{}
  >{\raggedright\arraybackslash}p{(\linewidth - 4\tabcolsep) * \real{0.3333}}
  >{\raggedright\arraybackslash}p{(\linewidth - 4\tabcolsep) * \real{0.3333}}
  >{\raggedright\arraybackslash}p{(\linewidth - 4\tabcolsep) * \real{0.3333}}@{}}
\toprule\noalign{}
\begin{minipage}[b]{\linewidth}\raggedright
Subsystem
\end{minipage} & \begin{minipage}[b]{\linewidth}\raggedright
Application Processor (MPU)
\end{minipage} & \begin{minipage}[b]{\linewidth}\raggedright
Real-Time Safety Controller (MCU)
\end{minipage} \\
\midrule\noalign{}
\endhead
\bottomrule\noalign{}
\endlastfoot
\textbf{Core Architecture} & Multi-core Out-of-Order Application CPU +
NPU/GPU & Deterministic In-Order Microcontroller (Zero-Wait-State) \\
\textbf{Operating System} & Embedded Linux (Preempt-RT / Debian / Yocto)
& Bare-Metal Firmware / Hard Real-Time OS \\
\textbf{Memory Subsystem} & 2 GB to 32 GB LPDDR4/LPDDR5 DRAM (Shared) &
256 KB to 1 MB On-Chip SRAM (TCM, Zero-Wait-State) \\
\textbf{Memory Management} & Virtual memory, TLB, demand paging &
Physical addressing only, zero dynamic allocation (\texttt{malloc}
prohibited) \\
\textbf{Execution Determinism} & Best-effort with tail jitter
(\(P_{99} \approx 20\text{--}80\text{ ms}\)) & Hard real-time cycle
predictability (WCET bounded, jitter \(< 1\,\mu\text{s}\)) \\
\textbf{Physical I/O Authority} & High-speed buses: MIPI CSI-2, USB 3.0,
PCIe, GbE & Dedicated safety I/O: Timer PWM, Quadrature Encoders,
Hardware E-Stop GPIO \\
\end{longtable}
}

\section*{Shared Memory RPMSG Message
Contract}\label{shared-memory-rpmsg-message-contract}
\addcontentsline{toc}{section}{Shared Memory RPMSG Message Contract}

\markright{Shared Memory RPMSG Message Contract}

The application processor transmits timestamped intent proposals to the
real-time safety controller over non-blocking shared SRAM mailboxes
using the Remote Processor Messaging (RPMSG) protocol. The shared
message schema (\texttt{PhysicalAgent\_IntentProposal\_t}) enforces a
packed binary layout across cores:

\begin{longtable}[]{@{}
  >{\raggedright\arraybackslash}p{(\linewidth - 6\tabcolsep) * \real{0.2500}}
  >{\raggedright\arraybackslash}p{(\linewidth - 6\tabcolsep) * \real{0.2500}}
  >{\raggedright\arraybackslash}p{(\linewidth - 6\tabcolsep) * \real{0.2500}}
  >{\raggedright\arraybackslash}p{(\linewidth - 6\tabcolsep) * \real{0.2500}}@{}}
\caption{Physical AI Dual-Brain IPC Intent Proposal Schema. Binary
layout of the packed shared-memory RPMSG
payload.}\label{tbl-ipc-intent-contract}\tabularnewline
\toprule\noalign{}
\begin{minipage}[b]{\linewidth}\raggedright
Field Name
\end{minipage} & \begin{minipage}[b]{\linewidth}\raggedright
Type
\end{minipage} & \begin{minipage}[b]{\linewidth}\raggedright
Size
\end{minipage} & \begin{minipage}[b]{\linewidth}\raggedright
Physical Units \& Description
\end{minipage} \\
\midrule\noalign{}
\endfirsthead
\toprule\noalign{}
\begin{minipage}[b]{\linewidth}\raggedright
Field Name
\end{minipage} & \begin{minipage}[b]{\linewidth}\raggedright
Type
\end{minipage} & \begin{minipage}[b]{\linewidth}\raggedright
Size
\end{minipage} & \begin{minipage}[b]{\linewidth}\raggedright
Physical Units \& Description
\end{minipage} \\
\midrule\noalign{}
\endhead
\bottomrule\noalign{}
\endlastfoot
\texttt{message\_id} & \texttt{uint32\_t} & 4 bytes & Monotonic proposal
sequence counter. \\
\texttt{timestamp\_us} & \texttt{uint64\_t} & 8 bytes & 64-bit IEEE 1588
PTP synchronized hardware timestamp (\(\mu\text{s}\)). \\
\texttt{lease\_duration\_ms} & \texttt{uint32\_t} & 4 bytes & Maximum
validity lease duration before proposal expires (\(\text{ms}\)). \\
\texttt{target\_pos\_{[}x,y,z{]}} & \texttt{float{[}3{]}} & 12 bytes &
3D Cartesian spatial target coordinates in robot base frame
(\(\text{m}\)). \\
\texttt{max\_velocity\_m\_s} & \texttt{float} & 4 bytes & Maximum
permissible Cartesian velocity envelope (\(\text{m/s}\)). \\
\texttt{max\_acceleration\_m\_s2} & \texttt{float} & 4 bytes & Maximum
permissible Cartesian acceleration bound (\(\text{m/s}^2\)). \\
\texttt{model\_confidence} & \texttt{float} & 4 bytes & Epistemic model
confidence score \(\gamma_{\text{conf}} \in [0.0, 1.0]\). \\
\texttt{crc32} & \texttt{uint32\_t} & 4 bytes & IEEE 802.3 32-bit CRC
payload integrity checksum. \\
\end{longtable}

\section*{Hardware Safety Watchdog
Protocol}\label{hardware-safety-watchdog-protocol}
\addcontentsline{toc}{section}{Hardware Safety Watchdog Protocol}

\markright{Hardware Safety Watchdog Protocol}

The real-time controller monitors the application processor via a
dedicated hardware countdown watchdog and heartbeat protocol:

\begin{enumerate}
\def\labelenumi{\arabic{enumi}.}
\tightlist
\item
  \textbf{Heartbeat Packet:} The application processor must transmit an
  authenticated heartbeat packet every
  \(\tau_{\text{ping}} \le 10\text{ ms}\).
\item
  \textbf{Watchdog Timeout:} If the real-time controller receives no
  valid heartbeat for \(\tau_{\text{timeout}} = 50\text{ ms}\), an
  independent hardware timer triggers a Non-Maskable Interrupt (NMI).
\item
  \textbf{Hardware Action:} The real-time controller immediately
  de-asserts motor gate driver enable lines and transitions inverter
  power stages into a deterministic safe stop (Category 1 / Safe Torque
  Off) within \(250\,\mu\text{s}\).
\end{enumerate}

\chapter*{References}\label{references}
\addcontentsline{toc}{chapter}{References}

\markboth{References}{References}

\protect\phantomsection\label{refs}
\begin{CSLReferences}{1}{1}
\bibitem[\citeproctext]{ref-ahn2022saycan}
{Ahn, Michael, Anthony Brohan, Yevgen Chebotar, et al.} 2023. {``Do as
{I} Can, Not as {I} Say: Grounding Language in Robotic Affordances.''}
\emph{Proceedings of the 6th Conference on Robot Learning (CoRL)},
Proceedings of machine learning research, vol. 205: 287--318.

\bibitem[\citeproctext]{ref-amanatides1987fast}
Amanatides, John, and Andrew Woo. 1987. {``A Fast Voxel Traversal
Algorithm for Ray Tracing.''} \emph{Eurographics '87}, 3--10.

\bibitem[\citeproctext]{ref-ames2019control}
Ames, Aaron D, Samuel Coogan, Magnus Egerstedt, Gennaro Notomista,
Koushil Sreenath, and Paulo Tabuada. 2019. {``Control Barrier Functions:
Theory and Applications.''} \emph{European Control Conference (ECC)},
3420--31.

\bibitem[\citeproctext]{ref-ames2014control}
Ames, Aaron D, Jessy W Grizzle, and Paulo Tabuada. 2014. {``Control
Barrier Function Based Quadratic Programs with Application to Adaptive
Cruise Control.''} \emph{53rd IEEE Conference on Decision and Control},
6271--78.

\bibitem[\citeproctext]{ref-ames2016control}
Ames, Aaron D, Xiangru Xu, Jessy W Grizzle, and Paulo Tabuada. 2016.
{``Control Barrier Function Based Quadratic Programs for Safety Critical
Systems.''} \emph{IEEE Transactions on Automatic Control} 62 (8):
3861--76.

\bibitem[\citeproctext]{ref-andrychowicz2020learning}
{Andrychowicz, Marcin, Bowen Baker, Maciek Chociej, Rafał J'ozefowicz,
et al.} 2020. {``Learning Dexterous in-Hand Manipulation.''} \emph{The
International Journal of Robotics Research} 39 (1): 3--20.

\bibitem[\citeproctext]{ref-astrom2006advanced}
Åström, Karl Johan, and Tore Hägglund. 2006. \emph{Advanced PID
Control}. ISA-The Instrumentation, Systems,; Automation Society.

\bibitem[\citeproctext]{ref-astrom2021feedback}
Åström, Karl Johan, and Richard M Murray. 2021. \emph{Feedback Systems:
An Introduction for Scientists and Engineers}. Princeton university
press.

\bibitem[\citeproctext]{ref-astrom1997computer}
Åström, Karl Johan, and Björn Wittenmark. 1997.
\emph{Computer-Controlled Systems: Theory and Design}. 3rd ed. Prentice
Hall.

\bibitem[\citeproctext]{ref-bahl2024okrobot}
Bahl, Shikhar, Mustafa Mukadam, Abhinav Gupta, and Deepak Pathak. 2024.
{``{OK-Robot}: What Really Matters in Integrating Open-Knowledge Models
for Robotics.''} \emph{arXiv Preprint arXiv:2401.12202}.

\bibitem[\citeproctext]{ref-barr2013toyota}
Barr, Michael. 2013. {``Unintended Acceleration and Other Embedded
Software Defects.''} \emph{Barr Group Technical Report and Expert
Testimony in Bookout v. Toyota}.

\bibitem[\citeproctext]{ref-biagiotti2008trajectory}
Biagiotti, Luigi, and Claudio Melchiorri. 2008. \emph{Trajectory
Planning for Automatic Machines and Robots}. Springer Science \&
Business Media.

\bibitem[\citeproctext]{ref-blanchini1999set}
Blanchini, Franco. 1999. {``Set Invariance in Control.''}
\emph{Automatica} 35 (11): 1747--67.

\bibitem[\citeproctext]{ref-brohan2023rt1}
{Brohan, Anthony, Noah Brown, Justice Carbajal, Yevgen Chebotar, Joseph
Dabis, et al.} 2023. {``{RT-1}: Robotics Transformer for Large-Scale
Real-World Control.''} \emph{Robotics: Science and Systems (RSS)}.

\bibitem[\citeproctext]{ref-brohan2023rt2}
{Brohan, Anthony, Noah Brown, Justice Carbajal, Yevgen Chebotar, Xi
Chen, et al.} 2023. {``Rt-2: Vision-Language-Action Models Transfer Web
Knowledge to Robotic Control.''} \emph{arXiv Preprint arXiv:2307.15818}.

\bibitem[\citeproctext]{ref-brooks1991intelligence}
Brooks, Rodney A. 1991. {``Intelligence Without Representation.''}
\emph{Artificial Intelligence} 47 (1--3): 139--59.
\url{https://doi.org/10.1016/0004-3702(91)90053-M}.

\bibitem[\citeproctext]{ref-bryant2015computer}
Bryant, Randal E., and David R. O'Hallaron. 2015. \emph{Computer
Systems: A Programmer's Perspective}. 3rd ed. Pearson.

\bibitem[\citeproctext]{ref-butler1993infeasibility}
Butler, Ricky W, and George B Finelli. 1993. {``The Infeasibility of
Quantifying the Reliability of Life-Critical Software.''} \emph{IEEE
Transactions on Software Engineering} 19 (1): 3--12.

\bibitem[\citeproctext]{ref-chi2023diffusionpolicy}
Chi, Cheng, Siyuan Feng, Yilun Pan, et al. 2025. {``Diffusion Policy:
Visuomotor Policy Learning via Action Diffusion.''} \emph{The
International Journal of Robotics Research} 44 (10): 1684--704.

\bibitem[\citeproctext]{ref-clopper1934use}
Clopper, C. J., and Egon S. Pearson. 1934. {``The Use of Confidence or
Fiducial Limits Illustrated in the Case of the Binomial.''}
\emph{Biometrika} 26 (4): 404--13.

\bibitem[\citeproctext]{ref-coles2001introduction}
Coles, Stuart. 2001. \emph{An Introduction to Statistical Modeling of
Extreme Values}. Springer.

\bibitem[\citeproctext]{ref-craig2005introduction}
Craig, John J. 2005. \emph{Introduction to Robotics: Mechanics and
Control}. 3rd ed. Pearson Prentice Hall.

\bibitem[\citeproctext]{ref-curless1996volumetric}
Curless, Brian, and Marc Levoy. 1996. {``A Volumetric Method for
Building Complex Models from Range Images.''} \emph{Proceedings of the
23rd Annual Conference on Computer Graphics and Interactive Techniques
(SIGGRAPH '96)}, 303--12.

\bibitem[\citeproctext]{ref-dosovitskiy2020image}
{Dosovitskiy, Alexey, Lucas Beyer, Alexander Kolesnikov, et al.} 2020.
{``An Image Is Worth 16x16 Words: Transformers for Image Recognition at
Scale.''} \emph{arXiv Preprint arXiv:2010.11929}.

\bibitem[\citeproctext]{ref-driess2023palme}
{Driess, Danny, Fei Xia, Mehdi SM Sajjadi, et al.} 2023. {``PaLM-e: An
Embodied Multimodal Language Model.''} \emph{International Conference on
Machine Learning (ICML)}, 8469--88.

\bibitem[\citeproctext]{ref-forster2017onmanifold}
Forster, Christian, Luca Carlone, Frank Dellaert, and Davide Scaramuzza.
2017. {``On-Manifold Preintegration for Real-Time Visual--Inertial
Odometry.''} \emph{IEEE Transactions on Robotics} 33 (1): 1--21.

\bibitem[\citeproctext]{ref-franklin1998digital}
Franklin, Gene F, J David Powell, and Michael L Workman. 1998.
\emph{Digital Control of Dynamic Systems}. 3rd ed. Addison-Wesley.

\bibitem[\citeproctext]{ref-furgale2013unified}
Furgale, Paul, Joern Rehder, and Roland Siegwart. 2013. {``Unified
Temporal and Spatial Calibration for Multi-Sensor Systems.''} \emph{2013
IEEE/RSJ International Conference on Intelligent Robots and Systems
(IROS)}, 1280--86.

\bibitem[\citeproctext]{ref-gu2022open}
Gu, Xiuye, Tsung-Yi Lin, Weicheng Kuo, and Yin Cui. 2022.
{``Open-Vocabulary Object Detection via Vision and Language Knowledge
Distillation.''} \emph{International Conference on Learning
Representations}.

\bibitem[\citeproctext]{ref-hartley2003multiple}
Hartley, Richard, and Andrew Zisserman. 2003. \emph{Multiple View
Geometry in Computer Vision}. Cambridge university press.

\bibitem[\citeproctext]{ref-hennessy2019computer}
Hennessy, John L., and David A. Patterson. 2019. \emph{Computer
Architecture: A Quantitative Approach}. 6th ed. Morgan Kaufmann.

\bibitem[\citeproctext]{ref-hermann1977nonlinear}
Hermann, Robert, and Arthur J. Krener. 1977. {``Nonlinear
Controllability and Observability.''} \emph{IEEE Transactions on
Automatic Control} 22 (5): 728--40.

\bibitem[\citeproctext]{ref-ho2020denoising}
Ho, Jonathan, Ajay Jain, and Pieter Abbeel. 2020. {``Denoising Diffusion
Probabilistic Models.''} \emph{Advances in Neural Information Processing
Systems} 33: 6840--51.

\bibitem[\citeproctext]{ref-hornung2013octomap}
Hornung, Armin, Kai M Wurm, Maren Bennewitz, Cyrill Stachniss, and
Wolfram Burgard. 2013. {``Octo{M}ap: An Efficient Probabilistic {3D}
Mapping Framework Based on Octrees.''} \emph{Autonomous Robots} 34 (3):
189--206.

\bibitem[\citeproctext]{ref-huang2023voxposer}
Huang, Wenlong, Chen Wang, Ruohan Zhang, et al. 2023. {``{VoxPoser}:
Composable {3D} Value Maps for Everyday Robot Manipulation with Language
Models.''} \emph{arXiv Preprint arXiv:2307.05973}.

\bibitem[\citeproctext]{ref-hughes2022hydra}
Hughes, Nathan, Yun Chang, and Luca Carlone. 2022. {``Hydra: A Real-Time
Spatial Perception Engine for {3D} Scene Graph Construction and
Optimization.''} \emph{Proceedings of Robotics: Science and Systems
(RSS)}.

\bibitem[\citeproctext]{ref-iso21448}
\emph{ISO 21448: Road Vehicles, Safety of the Intended Functionality}.
2022. International Organization for Standardization.

\bibitem[\citeproctext]{ref-iso26262}
\emph{ISO 26262: Road Vehicles, Functional Safety}. 2018. International
Organization for Standardization.

\bibitem[\citeproctext]{ref-isots15066}
\emph{ISO/TS 15066: Robots and Robotic Devices, Collaborative Robots}.
2016. International Organization for Standardization.

\bibitem[\citeproctext]{ref-jatavallabhula2023conceptfusion}
{Jatavallabhula, Krishna Murthy, Alihusein Kuwajerwala, Qiao Gu, et al.}
2023. {``{ConceptFusion}: Open-Set Multimodal {3D} Mapping.''}
\emph{Proceedings of Robotics: Science and Systems (RSS)}.

\bibitem[\citeproctext]{ref-jatr2019boeing}
Joint Authorities Technical Review. 2019. \emph{Boeing 737 MAX Flight
Control System: Observations, Findings, and Recommendations}. Federal
Aviation Administration.

\bibitem[\citeproctext]{ref-kalman1960new}
Kalman, Rudolph Emil. 1960. {``A New Approach to Linear Filtering and
Prediction Problems.''} \emph{Journal of Basic Engineering} 82 (1):
35--45.

\bibitem[\citeproctext]{ref-kamath2021mdetr}
Kamath, Aishwarya, Mannat Singh, Yann LeCun, Gabriel Synnaeve, Ishan
Misra, and Nicolas Carion. 2021. {``Mdetr-Modulated Detection for
End-to-End Multi-Modal Understanding.''} \emph{Proceedings of the
IEEE/CVF International Conference on Computer Vision}, 1780--90.

\bibitem[\citeproctext]{ref-kelly1998arguing}
Kelly, Timothy Patrick. 1998. {``Arguing Safety: A Systematic Approach
to Managing Safety Cases.''} PhD thesis, Department of Computer Science,
University of York.

\bibitem[\citeproctext]{ref-kelly2004goal}
Kelly, Tim, and Rob Weaver. 2004. {``The Goal Structuring Notation: A
Safety Argument Notation.''} \emph{Proceedings of the Dependable Systems
and Networks Workshop on Assurance Cases}.

\bibitem[\citeproctext]{ref-kerbl20233d}
Kerbl, Bernhard, Georgios Kopanas, Thomas Leimk"uhler, and George
Drettakis. 2023. {``{3D} {G}aussian Splatting for Real-Time Radiance
Field Rendering.''} \emph{ACM Transactions on Graphics} 42 (4):
139:1--14.

\bibitem[\citeproctext]{ref-kerr2023lerf}
Kerr, Justin, Chung Min Kim, Ken Goldberg, Angjoo Kanazawa, and Matthew
Tancik. 2023. {``{LERF}: Language Embedded Radiance Fields.''}
\emph{Proceedings of the IEEE/CVF International Conference on Computer
Vision (ICCV)}, 19729--39.

\bibitem[\citeproctext]{ref-khazatsky2024droid}
{Khazatsky, Alexander, Karl Pertsch, Suraj Nair, Ashwin Balakrishna,
Sudeep Dasari, et al.} 2024. {``{DROID}: A Large-Scale in-the-Wild Robot
Manipulation Dataset.''} \emph{Robotics: Science and Systems (RSS)}.

\bibitem[\citeproctext]{ref-kim2024openvla}
{Kim, Moo Jin, Karl Pertsch, Siddharth Karamcheti, et al.} 2024.
{``OpenVLA: An Open-Source Vision-Language-Action Model.''} \emph{arXiv
Preprint arXiv:2406.09246}.

\bibitem[\citeproctext]{ref-kirillov2023segment}
{Kirillov, Alexander, Eric Mintun, Nikhila Ravi, Hanzi Mao, et al.}
2023. {``Segment Anything.''} \emph{Proceedings of the IEEE/CVF
International Conference on Computer Vision (ICCV)}, 4015--26.

\bibitem[\citeproctext]{ref-koopman2016challenges}
Koopman, Philip, and Michael Wagner. 2016. {``Challenges in Autonomous
Vehicle Testing and Validation.''} \emph{SAE International Journal of
Transportation Safety} 4 (1): 15--24.

\bibitem[\citeproctext]{ref-kuwajerwala2024conceptgraphs}
{Kuwajerwala, Alihusein, Qiao Gu, Sacha Morin, et al.} 2024.
{``{ConceptGraphs}: Open-Vocabulary {3D} Scene Graphs for Perception and
Planning.''} \emph{2024 IEEE International Conference on Robotics and
Automation (ICRA)}, 14170--77.

\bibitem[\citeproctext]{ref-leveson2011engineering}
Leveson, Nancy. 2011. \emph{Engineering a Safer World: Systems Thinking
Applied to Safety}. MIT Press.

\bibitem[\citeproctext]{ref-leveson1993investigating}
Leveson, Nancy G., and Clark S. Turner. 1993. {``An Investigation of the
Therac-25 Accidents.''} \emph{Computer} 26 (7): 18--41.

\bibitem[\citeproctext]{ref-levine2018learning}
Levine, Sergey, Peter Pastor, Alex Krizhevsky, Julian Ibarz, and Deirdre
Quillen. 2018. {``Learning Hand-Eye Coordination for Robotic Grasping
with Deep Learning and Large-Scale Data Collection.''} \emph{The
International Journal of Robotics Research} 37 (4-5): 421--36.

\bibitem[\citeproctext]{ref-li2014online}
Li, Mingyang, and Anastasios I. Mourikis. 2014. {``Online Temporal
Calibration for Camera--{IMU} Systems: Theory and Algorithms.''}
\emph{The International Journal of Robotics Research} 33 (7): 947--64.

\bibitem[\citeproctext]{ref-liang2023code}
{Liang, Jacky, Wenlong Huang, Fei Xia, et al.} 2023. {``Code as
Policies: Language Model Programs for Embodied Control.''} \emph{2023
IEEE International Conference on Robotics and Automation (ICRA)},
9493--500.

\bibitem[\citeproctext]{ref-lin2018architectural}
{Lin, Shih-Chieh, Yunqi Zhang, Chang-Hong Hsu, Matt Skach, Md E. Haque,
et al.} 2018. {``The Architectural Implications of Autonomous Driving:
Constraints and Acceleration.''} \emph{Proceedings of the Twenty-Third
International Conference on Architectural Support for Programming
Languages and Operating Systems (ASPLOS)}, 751--66.

\bibitem[\citeproctext]{ref-lions1996ariane5}
{Lions, Jacques-Louis et al.} 1996. \emph{Ariane 5 Flight 501 Failure
Report by the Inquiry Board}. European Space Agency (ESA).

\bibitem[\citeproctext]{ref-liu2023grounding}
{Liu, Shilong, Zhaoyang Zeng, Tianhe Ren, et al.} 2023. {``Grounding
DINO: Marrying DINO with Grounded Pre-Training for Open-Set Object
Detection.''} \emph{arXiv Preprint arXiv:2303.05499}.

\bibitem[\citeproctext]{ref-lynch2017modern}
Lynch, Kevin M, and Frank C Park. 2017. \emph{Modern Robotics:
Mechanics, Planning, and Control}. Cambridge University Press.

\bibitem[\citeproctext]{ref-bfu1994warsaw}
Main Commission Aircraft Accident Investigation Warsaw and German
Federal Bureau of Aircraft Accident Investigation. 1994. \emph{Report on
the Accident to Airbus A320-211 Aircraft in Warsaw on 14 September
1993}. Polish Ministry of Transport; Maritime Economy / BFU.

\bibitem[\citeproctext]{ref-matsuki2024gaussian}
Matsuki, Hidenobu, Ryo Murai, Paul H. J. Kelly, and Andrew J. Davison.
2024. {``Gaussian Splatting {SLAM}.''} \emph{Proceedings of the IEEE/CVF
Conference on Computer Vision and Pattern Recognition (CVPR)},
18039--48.

\bibitem[\citeproctext]{ref-maxwell1868governors}
Maxwell, J Clerk. 1868. {``On Governors.''} \emph{Proceedings of the
Royal Society of London} 16: 270--83.

\bibitem[\citeproctext]{ref-moravec1985high}
Moravec, Hans, and Alberto Elfes. 1985. {``High Resolution Maps from
Wide Angle Sonar.''} \emph{IEEE International Conference on Robotics and
Automation} 2: 116--21.

\bibitem[\citeproctext]{ref-nasa1999mco}
NASA Mishap Investigation Board. 1999. \emph{Mars Climate Orbiter Mishap
Investigation Board Phase i Report}. National Aeronautics; Space
Administration.

\bibitem[\citeproctext]{ref-nasa2011toyota}
National Aeronautics and Space Administration Engineering and Safety
Center. 2011. \emph{Technical Assessment of Toyota Electronic Throttle
Control (ETCS-i) Systems}. NASA NESC.

\bibitem[\citeproctext]{ref-nhtsa2023cruise}
National Highway Traffic Safety Administration (NHTSA). 2023.
\emph{Special Crash Investigation: Automated Driving System Post-Crash
Maneuver Collision and Pedestrian Injury (San Francisco, CA)}. Nos.
PE23-018. U.S. Department of Transportation.

\bibitem[\citeproctext]{ref-ntsb2017tesla}
National Transportation Safety Board. 2017. \emph{Collision Between a
Car Operating with Automated Control Systems and a Tractor-Semitrailer
Truck Near Williston, Florida, May 7, 2016}. HAR-17/02. NTSB.

\bibitem[\citeproctext]{ref-ntsb2019uber}
National Transportation Safety Board. 2019. \emph{Collision Between a
Self-Driving Car and a Pedestrian, Tempe, Arizona, March 18, 2018}.
HAR-19/03. NTSB.

\bibitem[\citeproctext]{ref-newcombe2011kinectfusion}
Newcombe, Richard A, Shahram Izadi, Otmar Hilliges, et al. 2011.
{``KinectFusion: Real-Time Dense Surface Mapping and Tracking.''}
\emph{IEEE International Symposium on Mixed and Augmented Reality
(ISMAR)}, 127--36.

\bibitem[\citeproctext]{ref-oleynikova2017voxblox}
Oleynikova, Helen, Zachary Taylor, Marius Fehr, Roland Siegwart, and
Juan Nieto. 2017. {``Voxblox: Incremental {3D} {E}uclidean Signed
Distance Fields for on-Board {MAV} Planning.''} \emph{IEEE/RSJ
International Conference on Intelligent Robots and Systems (IROS)},
1366--73.

\bibitem[\citeproctext]{ref-open_x_embodiment_collaboration2024open}
{Open X-Embodiment Collaboration, Abhishek Padalkar, Acorn Pooley, et
al.} 2024. {``Open {X-Embodiment}: Robotic Learning Datasets and {RT-X}
Models.''} \emph{2024 IEEE International Conference on Robotics and
Automation (ICRA)}, 6892--903.

\bibitem[\citeproctext]{ref-oquab2023dinov2}
{Oquab, Maxime, Timothée Darcet, Théo Moutakanni, et al.} 2023.
{``Dinov2: Learning Robust Visual Features Without Supervision.''}
\emph{arXiv Preprint arXiv:2304.07193}.

\bibitem[\citeproctext]{ref-pomerleau1989alvinn}
Pomerleau, Dean A. 1989. {``{ALVINN}: An Autonomous Land Vehicle in a
Neural Network.''} \emph{Advances in Neural Information Processing
Systems (NeurIPS)} 1: 305--13.

\bibitem[\citeproctext]{ref-rogers1986challenger}
Presidential Commission on the Space Shuttle Challenger Accident. 1986.
\emph{Report of the Presidential Commission on the Space Shuttle
Challenger Accident}. Cmd. 9811. NASA.

\bibitem[\citeproctext]{ref-radford2021learning}
{Radford, Alec, Jong Wook Kim, Chris Hallacy, et al.} 2021. {``Learning
Transferable Visual Models from Natural Language Supervision.''}
\emph{International Conference on Machine Learning}, 8748--63.

\bibitem[\citeproctext]{ref-ross2011reduction}
Ross, Stéphane, Geoffrey Gordon, and Drew Bagnell. 2011. {``A Reduction
of Imitation Learning and Structured Prediction to No-Regret Online
Learning.''} \emph{Proceedings of the Fourteenth International
Conference on Artificial Intelligence and Statistics}, 627--35.

\bibitem[\citeproctext]{ref-sha2001using}
Sha, Lui. 2001. {``Using Simplicity to Control Complexity.''} \emph{IEEE
Software} 18 (4): 20--28.

\bibitem[\citeproctext]{ref-siciliano2009robotics}
Siciliano, Bruno, Lorenzo Sciavicco, Luigi Villani, and Giuseppe Oriolo.
2009. \emph{Robotics: Modelling, Planning and Control}. Springer Science
\& Business Media.

\bibitem[\citeproctext]{ref-slotine1991applied}
Slotine, Jean-Jacques E, and Weiping Li. 1991. \emph{Applied Nonlinear
Control}. Prentice Hall.

\bibitem[\citeproctext]{ref-smith1986estimating}
Smith, Randall C., and Peter Cheeseman. 1986. {``Estimating Uncertain
Spatial Relationships in Robotics.''} \emph{Autonomous Robot Vehicles},
435--61.

\bibitem[\citeproctext]{ref-spencer2021feedback}
Spencer, Jonathan, Sanjiban Choudhury, Arun Venkatraman, Brian Ziebart,
and J. Andrew Bagnell. 2021. {``Feedback in Imitation Learning: The
Three Regimes of Covariate Shift.''} \emph{arXiv Preprint
arXiv:2102.02872}.

\bibitem[\citeproctext]{ref-sundaralingam2023curobo}
Sundaralingam, Balakumar, Siva Kumar Sastry Hari, Adam Fishman, et al.
2023. {``Cu{R}obo: Parallelized Collision-Free Robot Motion
Generation.''} \emph{IEEE International Conference on Robotics and
Automation (ICRA)}, 11124--31.

\bibitem[\citeproctext]{ref-sutton2019bitter}
Sutton, Richard. 2019. {``The Bitter Lesson.''} \emph{Incomplete Ideas
(Blog)} 13 (1): 38.

\bibitem[\citeproctext]{ref-szeliski2022computer}
Szeliski, Richard. 2022. \emph{Computer Vision: Algorithms and
Applications}. 2nd ed. Springer.

\bibitem[\citeproctext]{ref-tedrake2023manipulation}
Tedrake, Russ. 2023. \emph{Robotic Manipulation: Perception, Planning,
and Control}. Course Notes, MIT.

\bibitem[\citeproctext]{ref-ul4600}
\emph{UL 4600: Standard for Safety for the Evaluation of Autonomous
Products}. 2023. Underwriters Laboratories.

\bibitem[\citeproctext]{ref-csb2014macondo}
U.S. Chemical Safety and Hazard Investigation Board. 2014.
\emph{Investigation Report: Drilling Rig Explosion and Fire at the
Macondo Well (Deepwater Horizon)}. Report No. 2010-10-I-OS. U.S.
Chemical Safety Board.

\bibitem[\citeproctext]{ref-vaughan1996challenger}
Vaughan, Diane. 1996. \emph{The Challenger Launch Decision: Risky
Technology, Culture, and Deviance at NASA}. University of Chicago Press.

\bibitem[\citeproctext]{ref-wiener1948cybernetics}
Wiener, Norbert. 1948. \emph{Cybernetics: Or Control and Communication
in the Animal and the Machine}. John Wiley \& Sons.

\bibitem[\citeproctext]{ref-williams2009roofline}
Williams, Samuel, Andrew Waterman, and David Patterson. 2009.
{``Roofline: An Insightful Visual Performance Model for Multicore
Architectures.''} \emph{Communications of the ACM} 52 (4): 65--76.

\bibitem[\citeproctext]{ref-wu2023gello}
Wu, Philipp, Yiyang Shentu, Zhonghe Yi, Xingyu Lin, Ken Goldberg, and
Pieter Abbeel. 2024. {``GELLO: A General, Low-Cost, and Intuitive
Teleoperation Framework for Robot Manipulation.''} \emph{IEEE/RSJ
International Conference on Intelligent Robots and Systems (IROS)}.

\bibitem[\citeproctext]{ref-zhao2023learning}
Zhao, Tony Z, Vikash Kumar, Sergey Levine, and Chelsea Finn. 2023.
{``Learning Fine-Grained Bimanual Manipulation with Low-Cost
Hardware.''} \emph{Robotics: Science and Systems (RSS)}.

\end{CSLReferences}


\backmatter


\end{document}
\end{document}